% 本文件由 LaTeX 排版 Agent 根据有效 translation.json 自动生成。
% author=Augustine of Hippo; work_id=1801; title=圣咏释义

\chapter{\WorkTitle{圣咏释义}}

% structure: p0001 source=h1 target=omitted reason=page-title-already-rendered-by-parent

圣咏 1\newline{}圣咏 2\newline{}圣咏 3\newline{}圣咏 4\newline{}圣咏 5\newline{}圣咏 6\newline{}圣咏 7\newline{}圣咏 8\newline{}圣咏 9\newline{}圣咏 10\newline{}圣咏 11\newline{}圣咏 12\newline{}圣咏 13\newline{}圣咏 14\newline{}圣咏 15\newline{}圣咏 16\newline{}圣咏 17\newline{}圣咏 18\newline{}圣咏 19\newline{}圣咏 20\newline{}圣咏 21\newline{}圣咏 22\newline{}圣咏 23\newline{}圣咏 24\newline{}圣咏 25\newline{}圣咏 26\newline{}圣咏 27\newline{}圣咏 28\newline{}圣咏 29\newline{}圣咏 30\newline{}圣咏 31\newline{}圣咏 32\newline{}圣咏 33\newline{}圣咏 34\newline{}圣咏 35\newline{}圣咏 36\newline{}圣咏 37\newline{}圣咏 38\newline{}圣咏 39\newline{}圣咏 40\newline{}圣咏 41\newline{}圣咏 42\newline{}圣咏 43\newline{}圣咏 44\newline{}圣咏 45\newline{}圣咏 46\newline{}圣咏 47\newline{}圣咏 48\newline{}圣咏 49\newline{}圣咏 50\newline{}圣咏 51\newline{}圣咏 52\newline{}圣咏 53\newline{}圣咏 54\newline{}圣咏 55\newline{}圣咏 56\newline{}圣咏 57\newline{}圣咏 58\newline{}圣咏 59\newline{}圣咏 60\newline{}圣咏 61\newline{}圣咏 62\newline{}圣咏 63\newline{}圣咏 64\newline{}圣咏 65\newline{}圣咏 66\newline{}圣咏 67\newline{}圣咏 68\newline{}圣咏 69\newline{}圣咏 70\newline{}圣咏 71\newline{}圣咏 72\newline{}圣咏 73\newline{}圣咏 74\newline{}圣咏 75\newline{}圣咏 76\newline{}圣咏 77\newline{}圣咏 78\newline{}圣咏 79\newline{}圣咏 80\newline{}圣咏 81\newline{}圣咏 82\newline{}圣咏 83\newline{}圣咏 84\newline{}圣咏 85\newline{}圣咏 86\newline{}圣咏 87\newline{}圣咏 88\newline{}圣咏 89\newline{}圣咏 90\newline{}圣咏 91\newline{}圣咏 92\newline{}圣咏 93\newline{}圣咏 94\newline{}圣咏 95\newline{}圣咏 96\newline{}圣咏 97\newline{}圣咏 98\newline{}圣咏 99\newline{}圣咏 100\newline{}圣咏 101\newline{}圣咏 102\newline{}圣咏 103\newline{}圣咏 104\newline{}圣咏 105\newline{}圣咏 106\newline{}圣咏 107\newline{}圣咏 108\newline{}圣咏 109\newline{}圣咏 110\newline{}圣咏 111\newline{}圣咏 112\newline{}圣咏 113\newline{}圣咏 114\newline{}圣咏 115\newline{}圣咏 116\newline{}圣咏 117\newline{}圣咏 118\newline{}圣咏 119\newline{}圣咏 120\newline{}圣咏 121\newline{}圣咏 122\newline{}圣咏 123\newline{}圣咏 124\newline{}圣咏 125\newline{}圣咏 126\newline{}圣咏 127\newline{}圣咏 128\newline{}圣咏 129\newline{}圣咏 130\newline{}圣咏 131\newline{}圣咏 132\newline{}圣咏 133\newline{}圣咏 134\newline{}圣咏 135\newline{}圣咏 136\newline{}圣咏 137\newline{}圣咏 138\newline{}圣咏 139\newline{}圣咏 140\newline{}圣咏 141\newline{}圣咏 142\newline{}圣咏 143\newline{}圣咏 144\newline{}圣咏 145\newline{}圣咏 146\newline{}圣咏 147\newline{}圣咏 148\newline{}圣咏 149\newline{}圣咏 150\par

\SourceRecord{Translated by J.E. Tweed.；Nicene and Post-Nicene Fathers, First Series；Vol. 8.；Edited by Philip Schaff.；Buffalo, NY: Christian Literature Publishing Co.,；1888.；<http://www.newadvent.org/fathers/1801.htm>.}

% structure: p0001 source=h1 target=section reason=page-title

\section{圣咏第一篇诠释}

1. \Quote{不随从恶人的计谋，这人是有福的}\BibleRef{咏1:1}。这是指我们的主耶稣基督，那位主之人。\Quote{不随从恶人的计谋，这人是有福的}，正如\Quote{那属土的人}所做的\BibleRef{格前15:47}，他同意了被蛇欺骗的妻子，违背了天主的诫命。\Quote{也不站在罪人的道路上}。因为祂确实藉着像罪人一样出生而来到罪人的道路上，但祂没有\Quote{站}在其中，因为世上的诱惑没有抓住祂。\Quote{也不坐在瘟疫的座位上}。祂没有意愿于地上的国度和骄傲，这正可视为\Quote{瘟疫的座位}；因为几乎无人能免于统治的欲望，渴求人的荣耀。因为\Quote{瘟疫}是广泛传播的疾病，波及所有或几乎所有人。而\Quote{瘟疫的座位}更适宜理解为有害的教义；\Quote{他们的言语如同毒疮}蔓延\BibleRef{弟后2:17}。也必须注意这些词语的顺序：\Quote{随从}、\Quote{站}、\Quote{坐}。因为当人背离天主时，他就\Quote{随从}了；当人在罪中取乐时，他就\Quote{站}了；当人固执于骄傲中，除非被那既没有\Quote{随从}恶人的计谋、也没有\Quote{站}在罪人的道路上、也没有\Quote{坐}在瘟疫的座位上者所释放，否则不能回头时，他就\Quote{坐}了。\par

2. \Quote{他喜爱上主的法律，昼夜默思祂的诫律}\BibleRef{咏1:2}。宗徒说：\Quote{法律不是为义人立的}\BibleRef{弟前1:9}。但在法律之内与在法律之下是不同的。在法律之内的，按法律行事；在法律之下的，按法律受管束：因此前者是自由的，后者是奴隶。再者，法律有两种：一种是写下来并加诸于仆人的；另一种是由那不需要其\Quote{字面}者从心中领悟的。\Quote{昼夜默思}可以理解为持续不断，或\Quote{白天}在喜乐中，\Quote{黑夜}在苦难中。因为经上记载：\Quote{亚巴郎看见了我的日子，就欢喜}\BibleRef{若8:5}；论到苦难时说：\Quote{我的五内也教训了我，直到黑夜}。\par

3. \Quote{他像一棵树，栽在溪水旁}\BibleRef{咏1:3}。这可以指那为了我们的救恩而屈尊取了人性的智慧本身\EditorialNote{箴言第八章}，使祂作为人成为\Quote{栽在溪水旁的树}；因为在此意义上，另一首圣咏所说的话也可如此理解：\Quote{天主的河流充满了水}。或者指的是圣神，关于祂曾说：\Quote{祂要用圣神给你们授洗}\BibleRef{玛3:11}；又说：\Quote{谁若渴，到我这里来喝罢}\BibleRef{若7:37}；又说：\Quote{若是你知道天主的恩赐和那对你说\InnerQuote{给我水喝}的是谁，你或许会求祂，而祂必会给你活水；凡喝这水的，将永远不渴，并且这水要在祂内成为涌到永生的水泉}。或者，\Quote{溪水旁}可以指人民的罪恶，因为在默示录中，水被称为\Quote{人民}\BibleRef{默17:15}；再者，\Quote{溪流}合理地理解为\Quote{堕落}，这与罪恶相关。那么那\Quote{树}，即我们的主，从溪水旁，即从罪恶的人民被沿途吸引到祂训诲的根中，将\Quote{结果实}，即建立教会；\Quote{按季节}，即祂被光荣地复活升天之后。因为那时，藉着派遣圣神到宗徒们那里，并坚定他们对祂的信德，以及派遣他们到世界去，祂使教会\Quote{结果实}。\Quote{它的叶子也不凋落}，即祂的圣言绝不落空。因为\Quote{凡有血肉的，尽都如草，他的一切荣美，如同草上的花；草必枯萎，花必凋谢，但上主的话永远长存}\BibleRef{依40:6}。\Quote{凡祂所做的，尽都顺利}，即那树所结的一切；这必须理解为果实和叶子，即行为和言语。\par

4. \BibleQuote{不虔敬者并非如此}，他们并非如此，\BibleQuote{却像糠秕，被风吹散}。\BibleRef{咏 1:4} 此处\OriginalTerm{大地}{earth}应理解为在天主内的稳固，正如经上所说：\BibleQuote{上主是我的产业，我确有美好的产业。}为此经上说：\BibleQuote{你要等候上主，遵守祂的道路，祂必高举你，使你承受大地。}为此经上说：\BibleQuote{温良的人是有福的，因为他们要承受大地。}\BibleRef{玛 5:5} 由此也引申出一个比较：就如这可见的大地支持并包含外在的人，那不可见的大地也支持\OriginalTerm{内在的人}{inner man}。从\Quote{那大地的面前}，风将\OriginalTerm{不虔敬者}{ungodly}吹散，这风就是\OriginalTerm{骄傲}{pride}，因它使人自高。那因主家的丰裕而沉醉、因祂欢乐的急流而酣畅的人，曾警醒地说：\BibleQuote{不要让骄傲的脚踩到我。}骄傲从这大地驱散了那说\BibleQuote{我要将我的宝座安置在北边，我要相似至高者}的人。\BibleRef{依 14:13} 从大地的面前，它也驱散了那在同意并尝了禁树之果、妄想如同神之后，躲避天主面的人。\BibleRef{创 3:8} 这大地指向内在的人，而人因骄傲被从那里驱散，这在经上特别显明：\BibleQuote{为什么尘土和灰烬骄傲？因为他在有生之年，就已吐出了自己的脏腑。}\BibleRef{德 10:9} 因为，从他被驱散之处，说他驱散了自己，并非不合理。\par

5. \BibleQuote{因此，不虔敬者在审判中不得兴起}\BibleRef{咏 1:5}：\Quote{因此}，即因为\Quote{他们像尘土，被从大地的面前吹散。}他说这话很好，因为这剥夺了他们所渴求的，即他们可以审判；因此同一思想在下一句中更清晰地表达出来：\BibleQuote{罪人也不得进入义人的议决。}因为通常前面的话会这样更清楚地重复。所以，通过\Quote{罪人}应理解\Quote{不虔敬者}；前面\Quote{在审判中}，这里应是\Quote{在义人的议决中}。或者，如果不虔敬者与罪人是不同的事物，尽管每个不虔敬者都是罪人，但并非每个罪人都是不虔敬者；\BibleQuote{不虔敬者在审判中不得兴起}，意思是他们确实要兴起，但不是为了受审判，因为他们已被判定最确定的惩罚。但\Quote{罪人}不得兴起\Quote{在义人的议决中}，即他们可以审判，但也许他们要被审判；正如对他们所说的：\BibleQuote{火要试验各人的工程是怎样的。如果人的工程存留，他必得到报酬；如果人的工程被烧毁，他必受亏损，但他自己却要得救，只是像从火里经过一样。}\par

6. \BibleQuote{因为上主认识义人的道路}\BibleRef{咏 1:6}。正如人们所说，医学认识健康，但不认识疾病，然而疾病被医学技艺所识别。同样可以说，\Quote{上主认识义人的道路}，但不虔敬者的道路祂不认识。并非上主一无所知，然而祂对罪人说：\BibleQuote{我从来不认识你们}\BibleRef{玛 7:23}。\Quote{但不虔敬者的道路必要灭亡；}这等于说，不虔敬者的道路上主不认识。但更清楚地表达，这即是不被上主认识，就是\Quote{灭亡}；而这即是被上主认识，就是\Quote{存留}；如此，存在属于天主的认识，而不存在不属于祂的认识。因为上主说：\BibleQuote{我是自有者}，以及\BibleQuote{那自有者派遣了我。}\par

\SourceRecord{Translated by an anonymous scholar.；Nicene and Post-Nicene Fathers, First Series；Vol. 8.；Edited by Philip Schaff.；Buffalo, NY: Christian Literature Publishing Co.,；1888.；<http://www.newadvent.org/fathers/1801001.htm>.}

% structure: p0001 source=h1 target=section reason=page-title

\section{释义圣咏第二篇}

1. \BibleQuote{为何外邦人喧嚷，万民谋算虚妄的事？}\BibleRef{咏 2:1}。\BibleQuote{世上的君王起来，首领一同商议，反对主，反对祂的基督}\BibleRef{咏 2:2}。经文说：\Quote{为何？}意思是说，是枉然的。因为他们所希望的，即基督的毁灭，并未实现；因为这段话是针对我们主的迫害者说的，在\WorkTitle{宗徒大事录}中也曾提及。\BibleRef{宗 4:26}\par

2. \BibleQuote{让我们挣断他们的捆绑，脱去他们的轭}\BibleRef{咏 2:3}。虽然这段经文可以另作解释，但更适宜理解为那些被称为\Quote{谋算虚妄的事}的人的话。因此\Quote{让我们挣断他们的捆绑，脱去他们的轭}的意思是，让我们尽力使基督信仰不束缚我们，也不加在我们身上。\par

3. \BibleQuote{那坐在天上的必嘲笑他们，主必嗤笑他们}\BibleRef{咏 2:4}。这句话是重复的；因为\Quote{那坐在天上的}之后又说\Quote{主}；而\Quote{必嘲笑他们}之后又说\Quote{必嗤笑他们}。但这一切不可按肉体来理解，好像天主用面颊发笑，或用鼻孔嗤笑；而是应理解为祂赐给圣人的能力，使他们预见将来的事，即基督的名号和统治将传于后代并拥有万民，从而明白那些人\Quote{谋算虚妄的事}。因为这预知的能力就是天主的\Quote{嘲笑}和\Quote{嗤笑}。\Quote{那坐在天上的必嘲笑他们。}如果我们以\Quote{天上}理解为圣善的灵魂，那么天主藉着他们，预知将来，必\Quote{嘲笑他们，并嗤笑他们}。\par

4. \BibleQuote{那时祂要在烈怒中对他们说话，在怒气中使他们惊惶}\BibleRef{咏 2:5}。为了更清楚地说明祂如何\Quote{对他们说话}，他接着说，祂要\Quote{使他们惊惶}；因此\Quote{在祂的烈怒中}就是\Quote{在祂的怒气中}。但上主天主的\Quote{烈怒和怒气}不可理解为任何内心的扰动，而是祂最公正地施行报复的能力，藉着使一切受造物服从祂的服务。因为应当注意并记念在\WorkTitle{智慧篇}中所记载的：\BibleQuote{但祢，掌权的主，以平静审判，以极大的恩惠治理我们。}\BibleRef{智 12:18}。天主的\Quote{怒气}因此是那认识天主法律的灵魂，当它看到同一法律被罪人违犯时，在灵魂内产生的一种情感。因为藉着义人灵魂的这种情感，许多事得以报复。虽然天主的\Quote{怒气}也可以理解为那违犯天主法律的人所遭遇的理智的昏暗。\par

5. \BibleQuote{然而我已由祂立为君王，在熙雍，祂的圣山上，宣讲祂的谕旨}\BibleRef{咏 2:6}。这段话显然是出自我们救主基督本人的口。但如果熙雍如某些人所解释的是\Quote{观看}之意，我们不应理解为别的，而应理解为教会，在那里每日都激起渴望观看天主荣耀的光辉，正如宗徒所说的：\BibleQuote{我们众人以揭开的帕子，观看主的荣耀。}\BibleRef{格后 3:18}。因此这句话的意思是：然而我已由祂立为君王，统御祂的圣教会；因其卓越和稳固，祂称之为山。\Quote{然而我已由祂立为君王。}我，就是那些人所谋算要\Quote{挣断}其\Quote{捆绑}、要\Quote{脱去}其\Quote{轭}的那一位。\Quote{宣讲祂的谕旨。}既然这是每日实行的，谁不明白它的意思呢？\par

6. \BibleQuote{上主曾对我说：祢是我的儿子，我今日生了祢}\BibleRef{咏 2:7}。虽然那一日似乎也可以预言式地指耶稣基督按肉身诞生的那一日；在永恒中，没有过去好像已不存在，也没有未来好像尚未存在，只有现在，因为凡是永恒的，永远是现在；然而，正如\Quote{今日}暗示了在场的性质，这句话\Quote{我今日生了祢}被赋予了一种属神的解释，借此纯洁无瑕的公教信仰宣告了天主的德能和智慧——祂是独生子——的永恒生育。\par

7. \BibleQuote{祢向我请求，我必将万民赐给祢作为产业}\BibleRef{咏 2:8}。这段话立即有一种时间性的含义，指向祂所取的人性，祂将自己作为祭献献上，代替了一切祭献，祂也为我们转求；因此\Quote{祢向我请求}这句话可指为人类而设立的全部时间性安排，即万民应联结于基督的名下，这样从死亡中被救赎，并被天主所拥有。\Quote{我必将万民赐给祢作为产业}，祂拥有他们是为了他们的得救，并为祢结出属灵的果实。\Quote{并将地极赐给祢作为基业。}这是重复，\Quote{地极}代替了\Quote{万民}；但更清楚的是，好使我们明白所有的民族。而\Quote{祢的基业}代替了\Quote{祢的产业}。\par

8. \Quote{祢必用铁杖治理他们，}以毫不妥协的正义，\Quote{祢必将他们如同陶匠的器皿打碎} \BibleRef{咏 2:9}；这就是说，\Quote{祢要打碎}他们身上属地的情欲和旧人的污秽行为，以及一切从有罪的泥土中衍生和习得的东西。\Quote{现今，你们君王应当觉悟} \BibleRef{咏 2:10}。\Quote{现今；}即现今已被更新，你们所覆盖的泥土外衣已被除去，也就是属于你们过去生命的肉体的错误容器，\Quote{现今觉悟吧，}你们如今是\Quote{君王；}即如今能够治理你们内一切奴役和兽性的部分，如今也能够战斗，不像\Quote{那些打空气的，而是刻苦己身，叫身服我。} \BibleRef{格前 9:26} \Quote{你们审判世人的，当受管教。}这又是重复；\Quote{当受管教}代替了\Quote{觉悟；}而\Quote{你们审判世人的}代替了\Quote{你们君王。}因为祂以\Quote{审判世人的人}来表示属灵的人。因为我们所审判的一切，都在我们之下；凡在属灵的人之下的，有充分的理由被称为\Quote{地；}因为它被属地的败坏所玷污。\par

9. \Quote{当以敬畏事奉上主；}免得那话，\Quote{你们世上的君王和审判官，}变成骄傲：\Quote{且要战兢喜乐} \BibleRef{咏 2:11}。极佳地加上了\Quote{喜乐}，免得\Quote{当以敬畏事奉上主}看来趋向痛苦。但同样，免得这同一喜乐流于无节制的轻率，就加上了\Quote{战兢}，以作为警告并谨慎守护圣德。也可以这样理解：\Quote{现今，你们君王应当觉悟；}即现今我已被立为王，世上的君王们，不要忧愁，仿佛你们的尊荣被夺去了，反而要\Quote{觉悟并受管教。}因为你们在祂之下，对你们有益，祂赐给你们领悟和训诲。这对你们有益，使你们不轻率地主宰，而是\Quote{以敬畏事奉}万有的\Quote{上主，}并\Quote{喜乐}于最确定、最纯洁的真福，以一切谨慎和小心，免得你们因此堕入骄傲。\par

10. \Quote{你们应接受训诫，免得上主发怒，你们便由正义的道路上丧亡} \BibleRef{咏 2:12}。这与\Quote{觉悟}和\Quote{受管教}相同。因为觉悟和受管教，就是接受训诫。然而，说\Quote{接受}，已足够清楚地表明有一种保护和防御，抵御一切可能造成伤害的事物，除非以极大的谨慎去接受它。\Quote{免得上主发怒，}以一种不确定的方式表达，不是就先知所见而言，因为对他来说是确定的，而是就那些受警告的人而言；因为那些未得明确启示的人，通常以疑惑之心思虑天主的愤怒。因此他们应当对自己说：让我们\Quote{接受训诫，免得上主发怒，我们便由正义的道路上丧亡。}至于\Quote{上主发怒}应如何理解，以上已经说过。\Quote{你们便由正义的道路上丧亡。}这是一个严重的惩罚，为那些对正义的甘甜稍有觉察的人所惧怕；因为谁若由正义的道路上丧亡，必在极大的痛苦中徘徊于不义的道路上。\par

11. \Quote{当祂的怒火转眼燃烧时，凡投靠祂的人真是有福；}即当为不敬虔者和罪人所预备的惩罚临到时，它不仅不会落在那些\Quote{投靠}上主的人身上，反而为他们奠定并提升一个国度。因为祂不是说：\Quote{当祂的怒火转眼燃烧时，}投靠祂的人都安全，好像他们仅得免除惩罚；而是说：\Quote{有福；}其中包含一切美善的集大成。现在我认为\Quote{转眼}的意思是这样的：它将是一件突然的事，而罪人却以为它遥远且迟迟不来。\par

\SourceRecord{Translated by an anonymous scholar.；Nicene and Post-Nicene Fathers, First Series；Vol. 8.；Edited by Philip Schaff.；Buffalo, NY: Christian Literature Publishing Co.,；1888.；<http://www.newadvent.org/fathers/1801002.htm>.}

% structure: p0001 source=h1 target=section reason=page-title

\section{圣咏第三篇释义}

达味的圣咏，当他逃避他的儿子阿贝沙隆的面时。\par

1. \Quote{我睡了，安歇了；我醒来，因为主要扶持我}这句话，引导我们相信这篇圣咏应在基督的位格中来理解；因为这些词句更适用于我们主的苦难与复活，而不适用于描述达味逃避他叛逆之子的面的那段历史。况且，由于经上论及基督的门徒说：\Quote{新郎的朋友岂能当新郎还在他们身边时禁食呢？}\BibleRef{玛 9:15}，所以，若这里的不肖子意指那位不肖的门徒——出卖祂的人，便不足为奇了。虽然从历史角度可以理解说，当祂离开时，祂同其余的宗徒退到山上，从而可以说祂逃避了那人的面；然而在属灵的意义上，当天主子，即天主的德能与智慧，放弃了犹达斯的心思，当魔鬼完全占据了他时，正如经上所记：\Quote{魔鬼进入了他的心}\BibleRef{若 13:27}，那么可以很好地理解说基督逃避了他的面；并非因为基督给魔鬼让位，而是因为基督离去，魔鬼便占据了。这个离去，我想，在此圣咏中被称为\Quote{逃}，乃因其迅速；这也由我们主的话所表明：\Quote{你所要做的，快做吧！}\BibleRef{若 13:27}。同样，在日常谈话中，当我们说某件事想不起来时，会说\Quote{它从我这里逃走了}；对一位博学之人，我们说\Quote{没有什么能从他那里逃走}。因此，真理从犹达斯的心思中逃走了，当它不再光照他时。但阿贝沙隆，就如有些人所解释的，在拉丁文中意为\OriginalTerm{父亲之和平}{Patris pax}，即父亲的和平。那么，无论是在列王纪的历史中，当阿贝沙隆与父亲作战时；还是在新约的历史中，当犹达斯出卖我们主时，似乎奇怪，\Quote{父亲的和平}如何能得到理解。但在前者，仔细阅读的人会看到，达味在那场战争中对儿子怀着和平，甚至极度悲痛地哀悼其死亡，说：\Quote{我儿阿贝沙隆！巴不得我替你死了！}\BibleRef{撒下 18:33}。而在新约的历史中，由于我们主那极伟大、极奇妙的宽容：祂长久容忍他，仿佛他是个好人，尽管祂并非不知他的心思；祂接纳他参加晚餐，在晚餐中祂交付并赐予门徒祂体血的表象；最后，在祂被出卖的那一刻，祂接受了他和平的亲吻；由此很容易理解基督如何向出卖祂的人显示了和平，尽管那人因如此可憎的阴谋的内战而荒废了。因此阿贝沙隆被称为\Quote{父亲的和平}，因为他的父亲拥有那和平，而他却没有。\par

2. \Quote{主啊，迫害我的人何其众多！}\BibleRef{咏 3:1}。他们确实众多，甚至从祂门徒的数量中也有一人没有缺失，被加到了迫害者的行列。\Quote{有许多人起来攻击我；有许多人对我的灵魂说：在他天主那里没有他的救恩}\BibleRef{咏 3:2}。显然，如果他们有任何祂将要复活的念头，他们肯定不会杀害祂。正是为此，有那些话：\Quote{如果祂是天主子，就从十字架上下来吧！}以及\Quote{祂救了别人，却不能救自己}\BibleRef{玛 27:42}。因此，犹达斯也不会出卖祂，除非他属于那些藐视基督的人，说：\Quote{他在他的天主那里没有救恩。}\par

3. \Quote{然而你，主啊，你是我的扶持者。}这是在人的本性中向天主说的，因为人的扶持是圣言成了血肉。\Quote{我的荣耀。}连祂也称天主为祂的荣耀，因为天主的圣言如此取了祂，以致天主与祂成为一体。愿骄傲的人学习，他们不愿听人对自己说：\Quote{你有什么不是领受的呢？既然领受了，为什么还夸口，仿佛没有领受呢？}\BibleRef{格前 4:7}\Quote{和举扬我头的主}\BibleRef{咏 3:3}。我认为此处应指人的心智，它并非不合理地被称作灵魂的头；它如此固有于，并在某种程度上与取人的圣言的卓绝优越性结合，以至于并未因受难如此重大的贬抑而被摒弃。\par

4. \Quote{我用我的声音向上主呼求}\BibleRef{咏 3:4}；这不是肉体的声音，即由空气反振发出的声音；而是内心的声音，它不对人说话，但在天主面前却如同呐喊。苏撒纳便是以这声音被垂听；主自己也命令要以这声音在密室中祈祷，\BibleRef{玛 6:6}即在内心里无声地。也许有人会说，若肉体不发出词语的声音，就不是用这声音祈祷；因为即使在静默中我们内心祈祷，若思念干扰了祈祷者的心意，仍不能说是\Quote{我用我的声音向上主呼求}。唯有当灵魂单独地，不取任何肉性之物，也不取任何肉性的意图，在祈祷中对天主说话，而天主独自垂听时，这话才说得正确。但即使如此，也因意向的力度而被称为呐喊。\Quote{祂从祂的圣山上应允了我}。在先知书中，主本身被称为山，正如经上所记：\Quote{有一块非人手凿成的石头，长大成为一座山}\BibleRef{达 2:34}。但这不能理解为祂的位格，除非也许祂这样说话：凭我自己，就如我从我的圣山上应允了我，当祂住在我内时，也就是在这座山上。但如果我们理解为天主因祂的公义而应允，那就更清楚、更无困扰了。因为公义要求祂使那被杀的无辜者从死者中复活，祂行善却受恶报；祂也应向那以恶报善的迫害者给予适当的报酬。因为我们读到：\Quote{你的正义有如天主的山岳}。\par

5. \BibleQuote{我睡下，且安歇}\BibleRef{咏 3:5}。不妨合宜地指出，经上特意说\BibleQuote{我}，为表明祂是自愿受死，正如那话所说：\BibleQuote{为此我父爱我，因为我舍弃我的性命，为能再取回它。没有人夺去我的性命，而是我自愿舍弃它；我有权舍弃它，也有权再取回它。}\BibleRef{若 10:17}。因此，祂说：你们并非违我之意捉拿我、杀害我，而是\BibleQuote{我睡下，且安歇；我醒来，因为上主必扶持我。}圣经中有无数处用睡眠代替死亡的例子，如宗徒说：\BibleQuote{弟兄们，关于安息的人，我不愿意你们不知道。}\BibleRef{得前 4:13}。既然已说了\Quote{我睡下}，为何又加上\Quote{且安歇}？我们无需质疑。这类重复在圣经中常见，正如我们在第二篇圣咏中指出的多处。但有些抄本作\Quote{我睡下，且陷入沉睡}。根据不同抄本，按照希腊词\OriginalTerm{我}{\GreekText{ἐ} \GreekText{γὼ} \GreekText{δ}}的可能译法，也有不同的表达方式。\par

\SourceRecord{Translated by an anonymous scholar.；Nicene and Post-Nicene Fathers, First Series；Vol. 8.；Edited by Philip Schaff.；Buffalo, NY: Christian Literature Publishing Co.,；1888.；<http://www.newadvent.org/fathers/1801003.htm>.}

% structure: p0001 source=h1 target=section reason=page-title

\section{\\WorkTitle{圣咏第四篇释义}}

致终末：大卫的圣咏之歌。\par

1. \\BibleQuote{基督就是法律的终向，使一切信者获得正义。} \\BibleRef{罗 10:4} 因为此处的\\Quote{终向}意指圆满，而非消灭。现在或许有人会问：是否每一首\\Quote{咏}都是一篇\\Quote{圣咏}，或者反过来，每一篇\\Quote{圣咏}都是一首\\Quote{咏}？是否有些\\Quote{咏}不能称为\\Quote{圣咏}，而有些\\Quote{圣咏}不能称为\\Quote{咏}？但必须留意圣经，也许\\Quote{咏}并不是指欢愉的主题。而那些称为\\Quote{圣咏}的，是伴以索尔特里琴而唱的；历史告诉我们，这是先知达味所使用的，具有高深的奥秘。但此处不宜讨论此事，因为它需要长时间的探究和大量的辩论。目前，我们应关注主的人性在复活后的言语，或在教会中相信并寄望于祂的人的言语。\par

2. \\BibleQuote{当我呼号时，我正义的天主应允了我。} \\BibleRef{咏 4:1} \\Quote{我在困苦中，你曾使我宽广。} 你引领我从忧苦的窄路进入喜乐的宽广之路。因为，\\BibleQuote{困苦和窘迫临到每一个作恶的人。} \\BibleRef{罗 2:9} 但那说\\Quote{我们在患难中欢欣，因为知道患难生忍耐，}直到他说\\Quote{因为天主的爱借着所赐给我们的圣神，已倾注在我们心中。} 的人，心中没有窘迫，窘迫只由迫害他们的人从外在加给他们。人称的变化，即他从第三人称说\\Quote{祂应允了我，}突然转到第二人称说\\Quote{你使我宽广了；}如果这不是为了变化和文雅，那么圣咏作者先想向人宣告他已被垂听，后来又对垂听他的主说话，这很奇怪。除非，当他宣告了他如何被垂听时，在这心灵的宽广中，他更愿与天主交谈；由此显示什么是心灵的宽广，即让天主已经倾注在心中，以至能与祂内心交谈。这正应理解为一个相信基督、受光照的人所说的话；但如果理解为被天主的智慧所取的主的人性所说，我看不出这是合适的。因为主的人性从未被天主的智慧抛弃。但正如祂在患难中的祈祷正是我们软弱的标记，同样，这突然的心灵宽广，主也可以为祂的忠信者说话，祂也曾代替他们说：\\BibleQuote{我饿了，你们没有给我吃的；我渴了，你们没有给我喝的，} \\BibleRef{玛 25:42} 等等。因此，此处祂也可以说：\\Quote{你使我宽广了，}这是为祂最小的弟兄之一，与天主交谈，其\\Quote{爱}已\\BibleQuote{借所赐给我们的圣神倾注在他心中。} \\BibleRef{罗 5:5} \\Quote{求你怜悯我，应允我的祈祷。} 为何他已宣告被垂听且又被宽广后又再次祈求？这是为我们的缘故，关于我们，经上说：\\BibleQuote{但我们若盼望那未见之物，就得耐心等待；} \\BibleRef{罗 8:25} 还是说，信者中那已开始的，也许要得以完成？\par

3. \\BibleQuote{人子啊，你们心重几时？} \\BibleRef{咏 4:2} 他说，你们的错误至少应持续到天主之子的来临；为何你们还要心重？你们何时才停止诡计，如果真理已经来临，你们却不接受？\\Quote{你们为何喜爱虚妄，寻求谎言？} 你们为何想从最低下的事物中获得幸福？唯独真理——万物由之而真——使人幸福。因为，\\BibleQuote{虚妄出于欺骗者，一切皆是虚妄。} \\BibleRef{训 1:2} \\Quote{人在日光之下劳碌，有什么益处呢？} 那么，你们为何被对暂时之物的爱所束缚？为何追求那末后者，如同追求首生者，而这本是虚妄和谎言？因为你们想要它们与你们长存，但它们全都像影子一样消逝。\par

4. \\BibleQuote{你们要知道，主已显扬了祂的圣者。} \\BibleRef{咏 4:3} 除了祂——那从下地升起、置于天上天主右手的——还有谁呢？因此，祂斥责人类，要他们终能从此世的爱情转向祂。但如果连词的添加（因祂说：\\Quote{你们要知道}）对某人造成困难，他容易在圣经中注意到，这种表达方式在先知们所说的语言中是常见的。因为你常看到这样的开头：\\Quote{\\InnerQuote{且主对他说}}，\\Quote{\\InnerQuote{且主的话临到他}}。这种连词连接，在没有先行句子来衔接下文时，可能美妙地告诉我们：用言语表达真理，是与内心进行的默观相联的。不过在此处，可以说前一句\\Quote{你们为何喜爱虚妄，寻求谎言？} 仿佛是写成了：不要喜爱虚妄，不要寻求谎言。若这样读，最直接的衔接就是\\Quote{你们要知道，主已显扬了祂的圣者。} 但\\OriginalTerm{咏调停顿}{Diapsalma}的插入，阻止我们将此句与前句相连。因为无论这是希伯来语词，如有些人认为，意为\\Quote{但愿如此}；还是希腊语词，表示歌咏中的停顿（即\\Quote{咏歌}是指咏唱的部分，\\Quote{停顿}则指咏唱中的静默间隔；正如歌唱中声音的合拢称为\\OriginalTerm{合咏}{Sympsalma}，分开称为\\OriginalTerm{分咏}{Diapsalma}，标记了某种连续中断的停顿）：我说，无论是前者、后者，还是其他，至少这一点是可能的：在有此插曲之处，文意不能正确地延续和连接。\par

5. \BibleQuote{主必听我的呼求，当我向祂呼号时。} 我相信这是在此警告我们，应以极大的心灵热忱，即以内在而无形体的呼号，向天主恳求帮助。因为正如我们必须为今生所受的光照感谢天主，同样我们也必须为今生之后的安息祈祷。因此，无论从忠信的福音传道者，还是从我们的主自己的位格，都可以理解为好像是写着：\Quote{主必听你们的呼求，当你们向祂呼号时。}\par

6. \BibleQuote{你们应发怒，却不可犯罪} \BibleRef{咏 4:4}。因为有人会想，谁堪当被垂听？或是罪人如何向主呼号而不至于徒劳？因此，他说：\Quote{你们应发怒，}\Quote{却不可犯罪。}这可以有两种解释：要么，即使你们发怒，也不要犯罪；也就是说，即使灵魂内产生一种情绪（这情绪如今因罪的惩罚而不在我们控制之内），至少不要让理性与心智（这是灵魂内按天主重生的部分，使我们以心智服事天主的法律，虽然以肉身我们仍服事罪的法律\BibleRef{罗 7:25}）同意它；要么，你们要悔改，即对你们过去的罪发怒，并从此停止犯罪。\Quote{你们心里说什么：}这里理解为\Quote{你们要说：}，所以完整的句子是\Quote{你们心里说什么，你们就要说什么；}也就是说，不要成为那批人，关于他们曾说：\Quote{他们用嘴唇尊敬我，他们的心却远离我。\BibleRef{依 29:13}} \Quote{在你们的卧室里，要被刺透。}这就是前面已经说过的\Quote{在心里。}因为这就是那卧室，我们的主曾警告我们，要在里面关上门祈祷。\BibleRef{玛 6:6} 但是，\Quote{要被刺透}要么指悔改的痛苦，使灵魂在惩罚中自我刺透，以免在天主的审判中被定罪和折磨；要么指唤醒，使我们醒来以看见基督的光，就像使用刺那样。但有些人说，不是\Quote{要被刺透}，而是\Quote{要被打开}是更好的读法；因为在希腊文圣咏集中是\OriginalTerm{被刺透}{\GreekText{κατανύγητε}}，这指的是心灵的扩张，以便接受由圣神所倾注的爱。\par

7. \BibleQuote{你们要献上公义的祭，并仰望主} \BibleRef{咏 4:5}。他在另一篇圣咏中也说了同样的话：\BibleQuote{天主所要的祭，就是痛悔的精神。}因此，这里理解为这就是藉着悔改所献上的公义之祭，并非不合理。因为有什么比各人向自己的罪发怒，而非向别人的罪发怒，并在自我惩罚中将自己祭献给天主更为公义呢？或者，悔改之后的义行就是公义之祭吗？因为插入的\OriginalTerm{休止符}{Diapsalma}或许不无道理地暗示了从旧生命到新生命的过渡：即当旧人因悔改而被摧毁或削弱时，按照新人的重生，公义之祭可以献给天主；此时心灵已洁净，将自己放在信德的祭台上，被天上的火，即圣神所环绕。因此，这可能是其含义：\Quote{你们要献上公义的祭，并仰望主；}也就是说，要正直地生活，并盼望圣神的恩赐，好使你所信的真理由此照耀在你身上。\par

8. 但是，\Quote{望主}这句话尚未解释清楚。现在所望的是什么呢？岂不就是美善的事吗？但既然各人都要从天主那里获得他所爱的美善，而爱内在美善的人却不容易找到，这内在美善即属于内心之人，唯独这美善应当被爱，其余的事只应为需要而使用，不应为享乐而贪求；所以当他已说了\BibleQuote{望主}\BibleRef{咏 4:6}之后，极妙地补充道：\BibleQuote{许多人说：谁向我们显示幸福？}这就是所有愚妄不义之人的话，也是他们每日的询问：无论是那些向往世俗生活的平安与宁静，却因人类的乖戾而不得安宁的人；他们甚至在盲目中胆敢责怪事态的秩序，当他们因自己的罪过而陷入困境时，竟以为现在的时光比过去更糟；或是那些怀疑并绝望于那应许给我们的未来生命的人，他们常说：\Quote{谁知道这是真的？}或者\Quote{谁曾从下面上来告诉我们这事？}因此，他极其巧妙而简短地（即向那些拥有内在眼光的人）指出了应追求的美善是什么；回答了那些说\BibleQuote{谁向我们显示幸福？}的人。他说：\BibleQuote{你慈颜的光辉已印在我们身上，上主}。这光就是人全部且真正的美善，不是用眼看，而是用理智看的。但他说的\BibleQuote{印在我们身上}，如同钱币印有君王的肖像。因为人是按天主的肖像和模样受造的\BibleRef{创 1:26}，却因罪而玷污了它；因此，他真正永恒的美善在于藉新生而重新印上这肖像。我相信这就是某些人巧妙领会的意思：即主看见\Person{凯撒}的税银时所说的：\BibleQuote{凯撒的归凯撒，天主的归天主}\BibleRef{玛 22:21}。好像祂说：同样，\Person{凯撒}从你们身上索取他肖像的印记，天主也如此；就如税银归还给\Person{凯撒}，灵魂也当归还给天主，被祂慈颜的光辉照亮和印上标记。\BibleRef{咏 4:7}\BibleQuote{你使我心中充满喜乐}。因此，那些仍然心地沉重、\BibleQuote{喜爱虚幻，寻求谎话}的人，不应在外寻求喜乐，而应在内，即在天主慈颜之光被印刻的地方。因为基督住在内心之人\BibleRef{弗 3:16}，正如宗徒所说；因为看见真理属于祂，因为祂说过：\BibleQuote{我是真理}\BibleRef{若 14:6}。再者，当祂在宗徒内说话时说：\BibleQuote{你们愿意接受那位在我内说话的基督的证明吗？}\BibleRef{格后 13:3}。当然，祂不是从外面向他说话，而是在他的内心深处，即我们应当祈祷的内室。\par

9. 但是那些追求世俗之事的人（无疑很多），除了说\BibleQuote{谁向我们显示幸福？}之外，不知道该说什么，因为他们看不到自己内心那真实而确定的善。因此，对于他们，下面的话说得极其恰当：\BibleQuote{从祂的谷物、酒和油的时候起，他们便增多。}因为\Quote{祂的}这个添加并非多余。谷物属于天主：因为祂是\BibleQuote{从天上降下的生命之粮}\BibleRef{若 6:51}。酒也属于天主：因为他说：\BibleQuote{他们必因你殿中的丰盛而沉醉。}油也属于天主：关于这点经上说：\BibleQuote{你用油润泽了我的头。}但是那些说\BibleQuote{谁向我们显示幸福？}并且看不到天国在他们内的人，这些人\BibleQuote{从祂的谷物、酒和油的时候起，便增多。}因为增多并不总是意味着丰盈，而往往意味着匮乏：当灵魂沉溺于世俗享乐，永远燃烧着欲望，无法满足；且因纷繁忧虑而分心，不得看见单一的美善。经上论及这样的灵魂说：\BibleQuote{因为可朽坏的身体压着灵魂，这尘世的居所沉重地压着那思虑万端的心智}\BibleRef{智 9:15}。这样的灵魂，由于世俗之物的消逝与更替，即\BibleQuote{从祂的谷物、酒和油的时候起}，充满了无数空虚幻想，是如此之多，以致无法遵行那命令：\BibleQuote{思念主，在美善中，以纯朴的心寻求祂}\BibleRef{智 1:1}。因为这繁多与那纯朴严重对立。因此，作者舍弃这些因渴望世俗之物而增多的人，即那些说\BibleQuote{谁向我们显示幸福？}的人——这幸福并非用肉眼在外寻求，而是以内心的纯朴在内寻求——信者欢乐地说：\BibleQuote{在平安中，我将同眠安息}\BibleRef{咏 4:8}。因为这样的人合法地盼望远离必死之物的一切思想，并忘记现世的痛苦；这美妙而先知地以睡眠和安息为象征，在那状态中最完美的平安不会被任何骚动打断。但这在今生尚未获得，应在来生盼望。就连这些本身是未来时态的词语也向我们显示。因为经上不是说\Quote{我已睡眠安息}或\Quote{我正睡眠安息}，而是说\BibleQuote{我将睡眠安息}。那时，\BibleQuote{这必朽坏的要穿上不朽坏，这必死的要穿上不死；那时死亡将被胜利吞噬}\BibleRef{格前 15:54}。因此说：\BibleQuote{但如果我们盼望那看不见的，就要忍耐等候}\BibleRef{罗 8:25}。\par

10. 因此，与此相符，他添加了最后的话，说：\Quote{主啊，祢曾使我在单纯中安居于望德。}此处他说的不是\Quote{将使}，而是\Quote{已经使}。那么，现在望德在谁那里，谁就必定将得到所期望的。他说的\Quote{在单纯中}很恰当。因为这可能指那些众多的人，他们从祢的谷物、酒和油的时候起就增多，说：\Quote{谁给我们显示好事？}因为这众多的事物会消逝，而在圣人们中则持守着单纯：关于他们，在宗徒大事录中记载：\BibleQuote{众信徒都一心一意。}\BibleRef{宗 4:32}因此，我们应当在\OriginalTerm{单纯}{singleness}与质朴中，即远离那些生灭之物的众多与喧嚣，成为永恒与合一的爱好者，如果我们渴求依附于唯一的天主和我们的主。\par

\SourceRecord{Translated by an anonymous scholar.；Nicene and Post-Nicene Fathers, First Series；Vol. 8.；Edited by Philip Schaff.；Buffalo, NY: Christian Literature Publishing Co.,；1888.；<http://www.newadvent.org/fathers/1801004.htm>.}

% structure: p0001 source=h1 target=section reason=page-title

\section{圣咏第五篇释义}

1. 这篇圣咏的标题是：\Quote{为那承受产业的她}。教会即指此而言，她藉我们的主耶稣基督承受永生为产业；好使她拥有天主本身，在依恋祂时得以有福，如经上所载：\BibleQuote{温良的人是有福的，因为他们要承受土地}\BibleRef{玛 5:5}。是什么土地？岂不是那被称为\Quote{祢是我的希望，我在活人之地所得的份}的土地？又更清楚地：\Quote{上主是我产业的份，是我杯爵的份。}反过来说，教会也被称为天主的产业，如经上所载：\Quote{祢向我请求，我必将万民赐你为产业。}因此，天主被称为我们的产业，因为祂养育并扶持我们；而我们被称为天主的产业，因为祂治理并管理我们。因此，这篇圣咏中教会的声音被召叫承受产业，好让她自己也成为上主的产业。\par

2. \BibleQuote{上主，请听我的言语}\BibleRef{咏 5:1}。她既被召叫，便呼求上主；好使同一上主作她的助佑，使她得以穿越这世界的邪恶，而达于祂面前。\Quote{请理会我的哀号。}圣咏作者清楚地表明这哀号是什么：如何从内在、从心灵的内室，无需肉体的发声，而直达天主；因为肉体的声音被听见，而属神的声音被领会。虽然天主的听闻也可能是这样，不是凭着肉体的耳朵，而是凭祂至尊的无所不在。\par

3. \Quote{请垂听我祈祷的声音}；即那声音，他恳求天主理会的声音；关于这声音的性质，他已在说\Quote{请理会我的哀号}时暗示过了。\BibleQuote{请垂听我祈祷的声音，我的君王，我的天主}\BibleRef{咏 5:2}。虽然圣子是天主，圣父也是天主，圣父与圣子同是一位天主；若问及圣神，我们只能回答祂是天主；且当提及圣父、圣子、圣神时，我们只能理解为一位天主；然而圣经习惯将\Quote{君王}的称谓给予圣子。因此，按经上所说\BibleQuote{我是道路，藉着我，人来到父面前}\BibleRef{若 14:6}，正确地先是\Quote{我的君王}，然后\Quote{我的天主}。然而圣咏作者并未说\Quote{请祢垂听}，而是说\Quote{请垂听}。因为公教信仰宣认的不是两位或三位神，而是天主圣三，一位天主。并非如同\OriginalTerm{撒伯里乌}{Sabellius}所信的，同一圣三可时而为圣父，时而为圣子，时而为圣神；而是圣父只能是圣父，圣子只能是圣子，圣神只能是圣神，而这圣三仅是一位天主。因此，当宗徒说\BibleQuote{万有出于祂，万有藉着祂，万有归于祂}\BibleRef{罗 11:36}时，人们相信他暗示了天主圣三；然而他没有说\Quote{愿荣耀归于他们}，而是说\Quote{愿荣耀归于祂}。\par

4. \BibleQuote{因为我必要向祢祈祷。上主，在早晨祢必听我的声音。}\BibleRef{咏 5:3}。他上面所说的\Quote{请听}是什么意思？仿佛他渴望立刻被垂听。但现在他说：\Quote{在早晨祢必听}；不是\Quote{请听}；又说：\Quote{我必要向祢祈祷}；不是\Quote{我向祢祈祷}；随后又说：\Quote{在早晨我必要侍立在祢面前，并要观看}；不是\Quote{我侍立在祢面前并观看}。除非他先前的祈祷指的是呼求本身；但处在今世风暴的黑暗中，他意识到他尚未看到他渴望的，却仍不停止希望，\BibleQuote{因为所见的希望，不是希望}\BibleRef{罗 8:24}。然而，他明白为什么他看不见，因为黑夜尚未过去，即我们的罪所招致的黑暗。因此他说：\Quote{因为我必要向祢祈祷，上主}；即因为祢是如此大能，我必要向祢申诉，\Quote{在早晨祢必听我的声音。}他说，祢不是那能被那些尚未除去罪恶黑夜的人看见的；当我错误的黑夜过去，黑暗消散，即我因我的罪而招致的黑暗，那时\Quote{祢必听我的声音}。那么为何他上面不说\Quote{祢必听}而说\Quote{请听}？是否因为教会呼喊\Quote{请听}却未被垂听，她便意识到必须经历什么才能被垂听？或者，她已被垂听于高天，但尚未明白她已被垂听，因为她尚未看见垂听她的那一位；而她现在所说的\Quote{在早晨祢必听}，应当这样理解：在早晨我将明白我已蒙垂听？正如那句话：\Quote{上主，求祢兴起}，使我兴起。无论如何，经文\BibleQuote{上主你们的天主试验你们，要知道你们是否爱祂}\BibleRef{申 13:3}，只能这样理解：使你们藉着祂得以知道，并向你们自己显明，你们在祂的爱中有了何等进步。\par

5. \BibleQuote{早晨我將站在你面前，我將瞻仰} \BibleRef{詠 5:3}。所謂\Quote{我將站}，不就是\Quote{我不躺下}嗎？躺下是什麼，不就是安息於地上，即在世俗的快樂中尋求幸福嗎？他說：\Quote{我將站在你面前}，\Quote{我將瞻仰}。因此，我們不可依附於世俗的事物，如果我們想看天主，祂是由潔淨的心所看見的。\BibleQuote{因為你不是喜愛罪惡的天主。惡人不能與你同住，不義的人也不能在你眼前存留。你惱恨所有作惡的人，你要消滅所有說謊的人。上主憎惡流人血的和詭詐的人。} \BibleRef{詠 5:4} 罪惡、惡意、說謊、殺人、詭詐，以及一切類似的事，就是我們所說的夜；當這夜過去，清晨來臨，好使天主被看見。因此，他揭示了為什麼他要在早晨站在你面前並瞻仰：\BibleQuote{因為}，他說，\BibleQuote{你不是喜愛罪惡的天主。} 因為如果祂是一位喜愛罪惡的天主，祂甚至可以被不義的人看見，如此祂就不會在早晨被看見，即當罪惡之夜過去之時。\par

6. \BibleQuote{惡人不能與你同住：} 即他不能如此看見，以至於依附於你。接著是：\BibleQuote{不義的人也不能在你眼前存留。} 因為他們的眼睛，即他們的理智，被真理之光擊退，因他們罪惡的黑暗；由於罪惡的習慣，他們不能承受正確理解的明亮。因此，即使他們有時看見，即理解真理，但他們仍是不義的，因他們喜愛那使人背離真理的事物，而不能存留其中。因為他們攜帶自己的夜，即不僅是犯罪習慣，甚至是對犯罪的愛。但如果這夜過去了，即如果他們停止犯罪，且這愛及習慣被驅散，清晨來臨，使他們不僅理解，而且依附於真理。\par

7. \BibleQuote{你惱恨所有作惡的人。} 天主的惱恨可從一種表達方式理解，即每個罪人惱恨真理。因為似乎真理也惱恨那些不允許存留於她內的人。他們不存留，因為他們不能承受真理。\BibleQuote{你要消滅所有說謊的人。} 因為這是與真理相對的。但為避免有人認為任何實體或本性與真理相對，他應理解\Quote{謊言}關乎那不存在者，而非存在者。因為如果說存在者，就是說真理；但如果說不存在者，就是謊言。因此他說：\BibleQuote{你要消滅所有說謊的人；} 因為他們背離那存在者，轉向那不存在者。的確，許多謊言似乎是為了某人的安全或利益而說的，不是出於惡意，而是出於善意：就如\BibleBook{出谷}{紀}中那些收生婆所說的，\BibleRef{出 1:19} 她們向法郎報告假消息，為使以色列的嬰孩不被殺害。但即使是這些，她們受讚美也不是因為事實，而是因為所顯示的心態；因為那些僅以此方式說謊的人，終將達到脫離一切謊言的地步。因為在成全的人身上，甚至這些謊言也找不到。因為對這些人說：\BibleQuote{你們的話該當是：是就說是，非就說非；其他多餘的，便是出於邪惡。} \BibleRef{瑪 5:37} 而且在另一處經文也合宜地寫著：\BibleQuote{說話虛偽的，是殺害靈魂。} \BibleRef{智 1:11} 以免有人以為成全的、屬神的人，為了這暫時的生命應當說謊，在這生命中沒有靈魂被殺害，無論是他自己的還是他人的。但既然說謊是一回事，隱瞞真理是另一回事（如果說假話是一回事，不說真話是另一回事），如果有人不願出賣一個人甚至到可見的死亡，他應準備隱瞞真理，而不說假話；這樣他既不交出那人，也不說謊，免得為他人的身體殺害自己的靈魂。但如果他還不能做到，至少他只應承認出於如此必須的謊言，以使他能脫離甚至這些謊言，如果只剩下這些，並接受聖神的德能，藉此他能輕視一切為真理必須忍受的事。總之，有兩種謊言，其中沒有大罪，但也不是無罪：要麼當我們開玩笑時，要麼當我們為了行善而說謊時。第一種，開玩笑，不那麼有害，因為沒有欺騙。因為被說的人知道那是為了開玩笑而說的。第二種，則較無害，因為它帶有某種善意。老實說，那沒有雙重性的事，甚至不能被稱為謊言。例如，如果有人把劍託付給某人，他答應當託付者要求時就歸還；如果碰巧他發狂時要求他的劍，顯然那時不應歸還，免得他殺了自己或他人，直到他恢復清醒。這裡沒有雙重性，因為接受劍的人，當他答應在對方要求時歸還時，沒有預想到他會在發狂時要求。甚至主也隱瞞了真理，當他對尚未堅強的門徒說：\BibleQuote{我還有許多事要告訴你們，但你們現在不能承擔。} \BibleRef{若 16:12} 以及宗徒保祿說：\BibleQuote{我從前不能對你們說話，如同對屬神的人，只能如同對屬血肉的人。} \BibleRef{格前 3:1} 由此可見，有時不說真話是無可指責的。但說假話從未被允許給成全的人。\par

8. \BibleQuote{嗜血的人，和诡诈的人，主要憎恶。} 祂以上所说的，\BibleQuote{祢恼恨一切作恶的人，祢要消灭一切说谎的人，} 似乎在此重复了：因此可以将\Quote{嗜血的人}归于\Quote{作恶的人}，将\Quote{诡诈的人}归于\Quote{谎言}。因为诡诈就是做一件事，却假装另一件事。祂也恰当地用了一个词说：\Quote{要憎恶。}因为被剥夺继承权的人通常被称为被憎恶的。而这圣咏是\BibleQuote{为承受产业者}；她满怀希望，欣喜地加上说：\BibleQuote{至于我，因祢丰厚的慈爱，我要进入祢的殿宇} \BibleRef{咏 5:7}。\Quote{因祢丰厚的慈爱：} 祂也许是指那圆满而有福的众人，那将构成那座城邑的众人，如今教会正在产痛中，逐个地孕育他们。那许多重生并成全的人，正当地被称为天主丰厚的慈爱，谁能否认呢？因为最真实地说：\BibleQuote{世人算什么，祢竟眷顾他？人子算什么，祢竟垂顾他？} \BibleQuote{我要进入祢的殿宇：} 我想，意思就像一块石头进入建筑物中。因为天主的殿宇若不是天主的圣殿，又是什么呢？关于这圣殿，经上说：\BibleQuote{因为天主的圣殿是圣的，这圣殿就是你们} \BibleRef{格前 3:17}。在这建筑物中，祂是角石 \BibleRef{弗 2:20}，那与圣父同永的天主的能力和智慧取了祂。\par

9. \BibleQuote{我要在祢的圣殿前，在祢的敬畏中朝拜祢。} \Quote{在圣殿前}，我们理解为\Quote{靠近}圣殿。因为祂并没有说，我要在祢的圣殿\Quote{内}朝拜；而是\Quote{我要在祢的圣殿前朝拜祢。}这也要理解为不是指成全，而是指趋向成全的进步：因此\BibleQuote{我要进入祢的殿宇}这句话应表示成全。但为了使这顺利达成，祂先说\Quote{我要在祢的圣殿前朝拜。}也许因此祂加上了\Quote{在祢的敬畏中}，这对那些迈向救恩的人是极大的保护。但当人到达那里时，在他身上就发生了经上所写的：\BibleQuote{完美的爱驱除恐惧} \BibleRef{若一 4:18}。因为他们不再畏惧那如今已成为他们朋友的那一位，对祂说：\BibleQuote{从此我不再称你们为仆人，而称你们为朋友} \BibleRef{若 15:15}，当他们被引领到达所应许的境地时。\par

10. \BibleQuote{主，求祢因我的仇敌，以祢的公义引导我} \BibleRef{咏 5:8}。祂在这里足够清楚地表明，祂是在前进的道路上，即是在趋向成全的进步中，尚未在成全本身中，当祂热切渴望被引导时。祂说\Quote{因祢的公义}，不是那在人看来似乎是公义的。因为以恶报恶似乎就是公义，但那不是祂的公义——关于祂，经上说：祂使太阳升起，照耀善人，也照耀恶人；因为即使当天主惩罚罪人时，祂也不是把自己的恶加在他们身上，而是任凭他们陷入自己的恶。圣咏作者说：\BibleQuote{看，他怀了不义，怀了劳苦，生了恶：他掘了坑，挖了它，却掉在自己所挖的陷阱里；他的劳苦必归到自己头上，他的不义必降到自己头顶。} 因此，当天主惩罚时，祂是作为审判者惩罚那些违犯法律的人，不是从自己身上带恶给他们，而是驱使他们走向自己所选择的，以填满他们痛苦的总和。但人，当他以恶报恶时，是带着恶意的；因此他自己先是恶的，当他惩罚恶时。\par

11. \Quote{求祢在我眼前修直我的道路。} 再清楚不过的是，祂在此陈述了祂在旅途中的那段时期。因为这条道路不是行在地球的任何区域，而是在心灵的爱情中。祂说：\Quote{求祢在我眼前修直我的道路：} 即在没有人看见的地方，人的称赞或责备是不可信的。因为他们绝不能判断别人的良心，而正是在良心中，通往天主的道路被行走。因此接着又说：\BibleQuote{因为他们的口中没有真诚} \BibleRef{咏 5:9}。当然，对于这样的判断，我们不能信任，因此我们必须逃往良心和天主的注视。\BibleQuote{他们的心是虚妄的。} 那么，他们的口中怎能会有真理呢？他们的心被罪和罪的惩罚所欺骗。由此，人们被那声音唤回：\BibleQuote{你们为什么喜爱虚妄，寻求谎言呢？}\par

12. \BibleQuote{他们的喉咙是敞开的坟墓。} 这可以解释为指贪饕，为了它，人常常以谄媚说谎。祂巧妙地说：\Quote{敞开的坟墓}，因为这贪饕总是张着大口，不像坟墓，在收受尸体后就封闭了。也可以这样理解：人以谎言和盲目的谄媚将那些他们引诱犯罪的人吸引到自己身边，并如同吞噬他们一样，当他们把他们转向自己的生活方式时。而当事发生在他们身上，因为他们因罪而死亡，那些引领他们的人，正当地被称为敞开的坟墓：因为这些人本身也是毫无生气的，缺乏真理的生命；他们接纳那些被他们以谎言和虚妄之心杀死并转向他们自己的死人。\BibleQuote{他们用舌头行诡诈：} 即用邪恶的舌头。因为当祂说\Quote{他们的}时，似乎表示了这一点。因为恶人有邪恶的舌头，即当他们行诡诈时，他们说的是恶言。主对他们说：\BibleQuote{你们既然是恶的，怎能说出善来？} \BibleRef{玛 12:34}\par

13. \Quote{天主，求祢审判他们，使他们因自己的计谋而跌倒} \BibleRef{咏 5:10}。这是预言，而非咒骂。因为他并不希望此事发生，而是预见将要发生的事。这事降临于他们，并非因为他似乎曾盼望如此，而是因为他们理应承受这样的结局。因为随后他所说的\Quote{愿所有寄望于祢的人都欢喜}，也是以预言的方式说出，因他预见他们会欢喜。同样，\Quote{奋起祢的威能，降临吧}也是预言性的，因为他看见主会降临。虽然\Quote{使他们因自己的计谋而跌倒}这句话也可如此理解：更可能的是作者怀着善意，希望他们从恶念中跌倒，即不再图谋恶事。但接下来的\Quote{驱逐他们}却否定了这种解释。因为\Quote{被天主驱逐}无论如何不能视为好意。因此，这应理解为预言，而非出于恶意；当说出这些对于那些执意坚持前述之罪的人必然发生的事时，便是如此。所以，\Quote{使他们因自己的计谋而跌倒}，就是让他们因自我谴责的意念而跌倒，\Quote{他们自己的良心也作证}，如宗徒所说：\Quote{他们的思念互相控告或辩护，在天主公义审判的启示中。} \BibleRef{罗 2:15}\par

14. \Quote{按他们罪孽的众多驱逐他们}：即是远远地驱逐他们。因为\Quote{按他们罪孽的众多}，他们应被远远驱逐。于是，恶人从这产业中被驱逐，这产业是通过认识并看见天主而拥有的；如同病眼被驱逐，离开光明的照耀，当别人视为喜乐时，它们却感到痛苦。因此，这些人不能在清晨站立并看见。而这样的说法是极大的惩罚，正如\Quote{至于我，亲近天主对我是好的}是极大的赏报。与这惩罚相对的是\Quote{进入你主人的欢乐中} \BibleRef{玛 25:21}；因为与这驱逐相似的是\Quote{把他丢在外面的黑暗中} \BibleRef{玛 25:30}。\par

15. \Quote{他们使祢，上主，感到苦涩}：祂说：\Quote{我是从天上降下来的食粮} \BibleRef{若 6:51}；又说：\Quote{你们不要为那必坏的食物劳碌} \BibleRef{若 6:27}；又说：\Quote{你们要尝尝并观看，上主是何等甘饴。} 但对于罪人，真理的食粮是苦涩的。因此他们憎恨那说真理的口。这些人因犯罪而陷入如此病态，以致健康灵魂所喜悦的真理食粮，他们却如苦胆一样不能忍受。\par

16. \Quote{愿所有寄望于祢的人都欢喜}；当然是指那些尝到主是甘饴的人。\Quote{他们要永远欢欣，祢要住在其间} \BibleRef{咏 5:11}。这便是永恒的欢欣：当义人成为天主的圣殿，而祂，他们的内住者，成为他们的喜乐。\Quote{凡爱慕祢圣名的人，必因祢而夸耀}：如同当他们所爱的对象临在，供他们享受时一样。\Quote{因祢}说得很好，宛如拥有这产业的继承权，正如诗篇标题所说：当他们也是祂的产业时，这由\Quote{祢要住在其间}暗示。从这美善中被阻隔的，正是那些天主按他们罪孽的众多所驱逐的人。\par

17. \Quote{因为祢必祝福义人} \BibleRef{咏 5:12}。这祝福就是因天主而夸耀，并被天主所居住。这样的圣化赐予义人。但为使他们成义，先有召叫；这召叫并非出于功德，而是出于天主的恩宠。\Quote{因为所有人都犯了罪，缺乏天主的荣耀。} \BibleRef{罗 3:23} \Quote{祂所召叫的人，祂也使他们成义；祂所成义的人，祂也使他们光荣。} \BibleRef{罗 8:30} 既然召叫不是出于我们的功德，而是出于天主的良善和仁慈，祂接着说：\Quote{上主，祢以祢恩惠的盾牌为我们加冕。} 因为天主的恩惠先于我们的善愿，召叫罪人悔改。这就是战胜仇敌的武器，如经上所说：\Quote{谁能控告天主所拣选的人呢？} 又说：\Quote{若天主偕同我们，谁能反对我们呢？祂没有怜惜自己的独生子，反而为我们众人把祂交出来。} \Quote{因为当我们还是仇敌时，基督为我们死了；何况我们既和好，更要藉着祂得救，免于天主的义怒。} \BibleRef{罗 5:10} 这便是那无敌的盾牌，当仇敌因诸多磨难与诱惑而使我们怀疑救恩时，它便击退仇敌。\par

18. 整篇圣咏的内容，从\BibleQuote{主，请听我的言语}直到\BibleQuote{我的君王，我的天主}，都是祈求她被垂听的祈祷。接着是那些阻碍看见天主之事的观点，即她被垂听的知识，从\BibleQuote{因为我早晨向祢祈祷，祢必听我的声音}直到\BibleQuote{流人血和诡诈的人，主必憎恶}。第三，她希望那将成为天主殿宇的，甚至在现在那驱逐恐惧的成全之前，就开始在敬畏中亲近祂，从\BibleQuote{但我赖祢丰厚的慈爱}直到\BibleQuote{我要因祢的敬畏向祢的圣殿朝拜}。第四，当她正在她感到阻碍的那些事中进步和前进时，她祈祷能在内里得到帮助，无人看见，以免被恶言所误导，因为从\BibleQuote{主，因我的仇敌，请以祢的公义引导我}直到\BibleQuote{他们用舌头弄诡诈}的话。第五，是预言将临到不敬虔者的惩罚，那时义人仅能得救；以及义人将得到的赏报，他们蒙召就来，并勇毅地承受一切，直到被带到终点，从\BibleQuote{天主，请审判他们}直到圣咏的结束。\par

\SourceRecord{Translated by an anonymous scholar.；Nicene and Post-Nicene Fathers, First Series；Vol. 8.；Edited by Philip Schaff.；Buffalo, NY: Christian Literature Publishing Co.,；1888.；<http://www.newadvent.org/fathers/1801005.htm>.}

% structure: p0001 source=h1 target=section reason=page-title

\section{圣咏第六篇释义}

\Emph{直到终结，在第八篇的诗歌中}，\Emph{一篇达味的圣咏}。\par

1. \Quote{论第八篇}在此似乎晦暗不明。因为此标题的其余部分更为清晰。如今，有人以为它暗示了审判之日，即我们的主来临之时，祂将审判生者与死者。这来临，据信，将在从亚当算起七千年后，即七千年如七日过去，然后那时间如同第八日来临。但既然主曾说过：\BibleQuote{圣父凭自己的权柄所定的时候，不是你们可以知道的。}\BibleRef{宗 1:7} 又说：\BibleQuote{但关于那日子和那时刻，没有人知道，连天上的天使和权能也不知道，圣子也不知道，惟独圣父知道。}\BibleRef{谷 13:32} 并且经上又写着：\BibleQuote{主的日子像贼一样来到。}\BibleRef{得前 5:2} 这就清楚地表明，谁也不能凭任何年数计算而自诩知道那时间。因为如果那日要在七千年后到来，那么每个人都可以通过计算年份知道它的来临。那么，圣子甚至不知道这一点，又作何解释呢？当然，这话的意思是说，人并非由圣子得知此事，而非祂自己不知道；正如那种说法：\Quote{上主你们的天主试探你们，为要知道；}\BibleRef{申 13:3} 意思是说，为使你们知道；以及\Quote{上主，求祢兴起}，意思是说，使我们兴起。因此，当圣子如此被说成不知道那日子，不是因为祂不知道，而是因为祂使那些不宜知道的人不知道，即祂不向他们显示；那么，那种奇怪的狂妄，竟凭计算年份，预料主的日子在七千年后必定到来，有何意义呢？\par

2. 那么，我们甘愿不知道主不愿我们知道的事吧；让我们查考这标题\Quote{论第八篇}的含义。审判之日确实可以，甚至无需轻率计算年份，由第八篇来理解，因为今世结束后，获得永生，义人的灵魂将不再受时间的支配；而且，既然所有时间都以七日重复的循环运转，那或许被称为第八日，将不会有这种变化。还有另一个理由，可以合理接受，即为何审判应被称为第八，因为它将在两代之后发生，一代关乎身体，另一代关乎灵魂。因为从亚当到梅瑟，人类是照着身体生活的，即按照血肉：这被称为外在的人和旧人，旧约即赐予此人，以便通过虽是宗教的、却是属肉体的行为，预表将来的属神之事。在整个这段时期，当人照着身体生活时，\Quote{死亡为王，}如宗徒所说，\Quote{甚至是在那些没有犯过罪的人之上。}如今它为王\Quote{是照着亚当违命的样式，}\BibleRef{罗 5:14} 如同宗徒所说；因为这必须理解为直到梅瑟的时期，直到那时，法律的行为，即那些肉体遵守的圣事，因某种奥秘的缘故，也束缚着那些敬畏唯一天主的人。但从主的来临开始，从他那里经历了从肉体的割损到心灵的割损的转变，召叫发出，要人按照灵魂生活，即按照内在的人，这人也被称为\Quote{新人}\BibleRef{哥 3:10} 因着重生和属神交谈的更新。如今，显而易见，数字四与身体有关，因为身体由四种众所周知的元素构成，以及四种性质：干、湿、暖、冷。因此身体也由四季掌管：春、夏、秋、冬。这一切都是众所周知的。因为关于数字四与身体的关系，我们已在别处有些巧妙但隐晦地论述过：这在本书中应避免，因为我们希望适应未受过教育的人。但数字三与心灵有关，可以从我们被命令以三重方式爱天主，即全心、全灵、全意来理解：其中每一项我们都必须分别论述，不是在圣咏中，而是在福音中：就目前而言，为了证明数字三与心灵的关系，我认为已经说够了。那么，那些与旧人和旧约相关的身体数字，既已过去；那些与新人及新约相关的灵魂数字，也既已过去，可以说一个七数已经过去；因为一切都发生在时间中，四已分配给身体，三已分配给心灵；第八将到来，即审判之日：它将按功过分配当得之物，立时将圣人转入，不是属世的工作，而是永生；但将判罚不虔诚者进入永罚。\par

3. 教会因畏惧这种定罪，在此圣咏中祈祷说：\Quote{上主，求祢不要在震怒中责斥我}\BibleRef{咏 6:1}。宗徒也提到审判的忿怒：\Quote{你为自己积蓄忿怒，}他说，\Quote{为在忿怒的日子，即上主公义审判彰显的日子}\BibleRef{罗 2:5}。凡渴望在此生获得医治的人，都不愿在那审判中被责斥。\Quote{也不要在祢的烈怒中惩罚我}。\Quote{惩罚}一词似乎较为温和，因为它有助于改正。因为被责斥，即被控告的人，恐怕其结局就是定罪。但由于\Quote{烈怒}似乎比\Quote{愤怒}更强烈，那么将较温和的\Quote{惩罚}与较严厉的\Quote{烈怒}并列，可能令人困惑。但我认为这两个词指的是同一件事。因为在希腊文中，第一节的\OriginalTerm{愤怒}{\GreekText{θυμὸς}}与第二节的\OriginalTerm{怒气}{\GreekText{ὀργὴ}}含义相同。但拉丁人自己希望使用两个不同的词时，他们寻找了与\Quote{愤怒}相近的词，于是用了\Quote{烈怒}。因此抄本存在差异：有的先出现\Quote{愤怒}，后出现\Quote{烈怒}；有的则用\Quote{愤慨}或\Quote{怒气}代替\Quote{烈怒}。但无论何种读法，它都是灵魂的一种情绪，促使施加惩罚。然而，这种情绪不可归于天主，仿佛天主具有灵魂一般，关于天主曾说：\Quote{但是，上主，祢以安宁执行审判}\BibleRef{智 12:18}。安宁意味着不被扰乱。因此，扰乱并不属于作为审判者的天主；但祂的仆人按祂的法律所行的事，被称为祂的愤怒。在这愤怒中，那现在祈祷的灵魂不仅不愿被责斥，也不愿被惩罚，即被改正或训诫。因为在希腊文中是\OriginalTerm{训诫}{\GreekText{παιδεύσῃς}}，即教导。在审判之日，凡不以基督为根基的人，都要被\Quote{责斥}。但那些\Quote{在这根基上建造木、草、禾秸}的人，则要被改正，即被炼净，\Quote{他们将要受亏损，自己却要得救，像从火中经过一样}。那么，不愿在主愤怒中被责斥或改正的人祈祷什么？无非是获得医治。因为健康时，既不怕死亡，也不怕医生的烙铁或手术刀。\par

4. 接着他说：\Quote{上主，求祢怜悯我，因为我软弱；上主，求祢医治我，因为我骨骸战栗}\BibleRef{咏 6:2}，\Quote{骨骸}在此指灵魂的支撑或力量。因此，灵魂在提到骨骸时，是说她的力量在战栗。因为不能认为灵魂有我们身体中所见的骨骸。故此，下面的话有助于解释：\Quote{我的灵魂也万分惊惶}\BibleRef{咏 6:3}，以免因提到骨骸而被理解为身体上的。\Quote{上主，祢要到何时？}谁看不出这里描绘了一个与自身疾病挣扎的灵魂，却被医生长久拖延，好让她明白自己因罪陷入何等邪恶的境地？因为容易治愈的，人们不太回避；但治愈的艰难，会使康复后的健康更加谨慎。那么，天主，那向祂说\Quote{上主，祢要到何时？}的灵魂，不可被视为残忍，而是仁慈地使她明白自己招致的祸患。因为这灵魂尚未完美地祈祷，以致能对她说：\Quote{你还在说话时，我必应允}\BibleRef{依 65:24}。同时，她也得以知道，如果悔改者遭遇如此大的困难，那么不转向天主的恶人将受到何等大的惩罚，正如另一处所写：\Quote{义人尚且难以得救，而罪人和恶人又将如何？}\BibleRef{伯前 4:18}\par

5. \Quote{上主，求祢回心转意，拯救我的灵魂}\BibleRef{咏 6:4}。她自身回转时，祈求天主也转向她；正如经上说：\Quote{你们归向我，我也必归向你们——万军的上主说}\BibleRef{匝 1:3}。或者应该按照这种说法理解：\Quote{上主，求祢回心转意}，即使我回转，因为在这回转中，灵魂感到困难和劳苦？因为当我们完全回转时，必发现天主是准备好的，正如先知所说：\Quote{我们必发现祂如晨曦般随时准备着}。并非那位无所不在者离开了，而是我们的背离使我们失去了祂；\Quote{祂在世界，世界原是藉祂造成的，但世界却不认识祂}\BibleRef{若 1:10}。如果祂在世界，而世界不认识祂，那么我们的不洁无法承受看见祂。当我们回转时，即通过改变旧生活而塑造我们的精神，我们感到从尘世情欲的黑暗中被拽回神圣光明的宁静、安详与平安是艰难而劳苦的。在这困境中，我们说：\Quote{上主，求祢回心转意}，即帮助我们，使那回转在我们中得以完成，这回转发现祢是准备好的，并把自己献给爱祢的人享用。因此，他说了\Quote{上主，求祢回心转意}之后，又加上\Quote{拯救我的灵魂}，仿佛灵魂仍纠缠于此世的网罗，并在回转时遭受撕裂纠缠欲望的荆棘之苦。\Quote{求祢因祢的慈爱拯救我}。他知道自己得医治并非凭自己的功德：因为他犯了罪，违背了所赐的命令，应受公义的定罪。因此他说：求祢医治我，并非因我的功德，而是因祢的慈爱。\par

6. \BibleQuote{因为死亡之中，没有人记念祢} \BibleRef{咏 6:5}。他也知道，现在正是归向天主的时机：因为当此生过去后，只剩下我们应得的报应。\BibleQuote{但在地狱中，谁能承认祢？} \EditorialNote{路加福音第十六章} 主所提及的那个富户，他看见\Person{拉匝禄}安息，自己却在痛苦中哀号，就在地狱里承认了，甚至希望他的弟兄们被警告，使他们远离罪恶，因为那地狱中不为人相信的惩罚。因此，虽然毫无用处，但他承认那些折磨是他应得的，因为他甚至希望他的弟兄们受到教导，免得他们陷入同样的境地。那么，\BibleQuote{但在地狱中，谁能承认祢？} 是什么意思呢？地狱是否应理解为审判后不敬神的人将被投入的地方，那时因那更深的黑暗，他们将再也看不见天主的光明，无从向祂承认什么？因为那个富户举目观看，虽然有一巨大的深渊隔开，仍能看见\Person{拉匝禄}安息；与后者相比，他被迫承认自己的罪有应得。还可以这样理解：圣咏作者将违反天主法律的罪称为死亡，以至于我们应将死亡之刺称为死亡，因为它导致死亡。\BibleQuote{因为死亡之刺是罪。} \BibleRef{格前 15:56} 在这种死亡中，忘记天主、蔑视祂的法律和诫命，就是如此：因此圣咏作者用地狱意指那笼罩并包裹罪人（即垂死者）的灵魂的盲目。\BibleQuote{他们既然不认为好}，宗徒说，\BibleQuote{把天主保留在} \BibleQuote{他们的知识中，天主就任凭他们陷于邪恶的心思。} \BibleRef{罗 1:28} 灵魂渴望从这种死亡和这种地狱中得到保全，当她努力归向天主，并感受到自身的困难时，她恳切祈求。\par

7. 因此他接着说：\BibleQuote{我因呻吟而疲惫。} 好像这还不够，他接着说：\BibleQuote{我每夜要洗我的床铺} \BibleRef{咏 6:6}。这里所称的床铺，是指病弱灵魂安息之处，即肉体的享乐和一切世俗的快乐。凡努力摆脱这快乐的人，便以眼泪洗床铺。因他看出自己已谴责肉体的情欲，然而他的软弱仍被快乐所束缚，并乐意躺在其中，惟有痊愈的灵魂才能从中崛起。至于他说\Quote{每夜}，也许要这样理解：那精神上准备好、并察觉到一些真理之光的人，由于肉体的软弱，有时会停留在此世的快乐中，就不得不经历像昼夜交替的感觉：当他说\BibleQuote{以我的心思奉事天主的法律}时，他感到像是白天；当他说\BibleQuote{但以我的肉身为罪恶的法律服役} \BibleRef{罗 7:25} 时，他降入黑夜；直到所有的黑夜过去，那一日到来，经上说：\BibleQuote{早晨我必侍立在祢面前，并要观看}。那时他将站立不动，但现在他躺卧，在他的床铺上；他每夜要洗这床铺，要用如此丰富的眼泪从天主的仁慈中获得最可靠的医治。\BibleQuote{用眼泪浸透我的床榻}。这是重复。因为当他用\Quote{眼泪}一词，表明他上述\BibleQuote{我洗}的意思。我们这里将\Quote{床榻}与上述\Quote{床铺}视为相同。虽然\Quote{浸透}比\Quote{洗}更进一步：因为洗只是表面洁净，而浸透深入更内部；这表示流泪直达心底。他使用不同的时态：过去时\BibleQuote{我因呻吟而疲惫}；将来时\BibleQuote{我每夜要洗我的床铺}；将来时\BibleQuote{用眼泪浸透我的床榻}；这显示出当人徒然呻吟疲惫时，应当对自己说什么。好像他说：我这样做没有益处，因此我要做另一件事。\par

8. \BibleQuote{我的眼睛因忿怒而昏花} \BibleRef{咏 6:7}：是因他的忿怒，还是天主的忿怒？后者是他恳求不要被责斥或惩罚的。但若那里的忿怒暗示审判的日子，现在怎能理解？是不是它的开始：人在此世受苦、遭折磨，尤其是丧失对真理的理解；正如我已引用的：\BibleQuote{天主就任凭他们陷于邪恶的心思} \BibleRef{罗 1:28}。因为这就是心灵的盲目。凡被交付于此的人，就被隔绝于天主内在的光明之外；但此生在世时尚未完全隔绝。因为有\BibleQuote{外面的黑暗} \BibleRef{玛 25:30}，这被理解为属于审判的日子；凡趁还有时间拒绝改正的人，将完全失去天主。完全无天主，岂不是极端的盲目？如果天主\BibleQuote{住在不可接近的光中} \BibleRef{弟前 6:16}，进入那光的，就是那些被说\BibleQuote{进入你主人的福乐}的人。这就是此忿怒的开始，每个罪人在此世都经历。因此，在恐惧审判的日子时，他经受考验和悲伤；免得他陷入那灾难性的开始，他现在正经历其萌芽。所以他没有说\Quote{我的眼睛灭了}，而是\BibleQuote{我的眼睛因忿怒而昏花}。但如果他意指他的眼睛因自己的忿怒而昏花，这也并不奇怪。因为也许为此才说：\BibleQuote{不要让太阳在你们的忿怒中落下} \BibleRef{弗 4:26}；因为心灵因自己的混乱而不能看见天主，便以为内在的太阳，即天主的智慧，在她内像太阳落下一样。\par

9. \Quote{我因我的一切仇敌而衰老了。} 他先前只提到了愤怒（若他谈的真是他自己的愤怒），但想到其他恶习，他发现他被这一切所包围。这些恶习，既属于旧生活和旧人——我们必须脱去旧人，以便穿上新人（\BibleRef{哥 3:9}）——因此说得好：\Quote{我衰老了。} 但\Quote{在我的一切仇敌中}，他意指要么处在这些恶习之中，要么处在那些不愿归向天主的人群之中。因为这些人，即使他们不认识这些恶习，即使他们容忍它们，即使他们共用同样的餐桌、房屋和城市，彼此之间没有冲突，且经常交谈，看似和睦，但因他们目标的悖逆，他们成为那些归向天主者的仇敌。因为既然一些人爱慕并贪求这个世界，另一些人则希望摆脱这个世界，谁看不出前者是后者的仇敌呢？因为如果他们能够，他们便将后者也拖入惩罚之中。能每日与他们的言语打交道却不偏离天主诫命的道路，实乃极大的恩宠。因为那奋力迈向天主的灵魂，在路上受到粗暴对待时，常常惊惶；且通常不能实现其善愿，以免冒犯那些与之同居共处、喜爱并追求其他可朽坏且短暂之物的人。凡健全之人，与这样的人分离，不在空间上，而在灵魂上。因为身体处于空间之中，而灵魂的空间乃是其爱慕。\par

10. 因此，在劳苦、叹息和频繁的眼泪之后——因为如此恳切地向祂（祂是一切仁慈的泉源）所求的，不会没有效果——并且最真实地说：\Quote{主亲近那些心碎的人。} 在如此大的困难之后，敬虔的灵魂（我们也可以借此理解教会）暗示她已被垂听，请看她又加上了什么：\Quote{你们这些作恶的人，离开我吧，因为上主已听见我哭泣的声音。} \BibleRef{咏 6:8} 这或是预言性的，因为他们将要离开，即恶人将在审判之日与义人分离；或是针对当下的时刻。因为尽管两者同样出现在同一集会中，但在打谷场上，麦子已与糠秕分离，尽管它藏在糠秕之中。因此他们可以在一起，却不能一同被风吹去。\par

11. \Quote{主已听见我哭泣的声音；主已听见我的恳求；主已收纳了我的祈祷。} \BibleRef{咏 6:9} 相同情感的多次重复，并非表示叙述者的缺乏，而是他喜乐的炽热之情。因为喜乐的人常常如此说话，以至于一次说出喜乐的对象还不够。这就是在叹息中的劳苦和那泪水的果实——泪水洗了床榻，浸透了床铺。因为：\Quote{流泪撒种的，必将欢乐收割。} 以及\Quote{哀恸的人是有福的，因为他们要受安慰。}\par

12. \Quote{愿我所有的仇敌都羞愧又惊惶。} \BibleRef{咏 6:10} 他上面说：\Quote{你们所有这些人都离开我。} 这可以如所解释的，甚至在此生发生；但至于他所说的\Quote{愿他们羞愧又惊惶}，我看不出这如何能发生，除非在那一日，即义人的赏报和罪人的惩罚显明的时候。因为现今，恶人非但不羞愧，反而不断侮辱我们。而且他们的嘲弄大多如此有效，以至于使软弱的人因基督的名字而羞愧。因此经上说：\Quote{凡在人前以我为羞耻的，我在我父面前也要以他为羞耻。} 但如今，凡愿意履行那些崇高命令——散财、赒济穷人，使自己的义德永远长存，并且卖掉所有世俗财物、分施给穷人——而跟随基督的人，说：\Quote{我们没有带什么到世上来，也不能带什么去；只要有衣有食，就当知足。} \BibleRef{弟前 6:7} 他遭受那些人的亵渎嘲笑，并被那些不愿得痊愈的人称为疯子；并且常常为了避免被绝望的人这样称呼，他害怕去做，并推迟那最忠实、最有力的医生所吩咐的事。因此，现今这些不能羞愧，我们反倒要祈求不被他们弄得羞愧，以免被从既定的旅程中叫回、阻碍或延迟。但时候将到，他们要羞愧，正如经上记载：\Quote{这些人是他们曾一度讥笑和污蔑的对象。我们愚人曾以为他们的生活是疯狂，他们的结局是羞辱；如今他们怎么被列在天主的儿女中，与圣者同享分？如此，我们迷失了真理的道路，正义的光没有照耀我们，太阳也没有为我们升起。我们充满了邪恶和毁灭的道路，走过崎岖的荒野，却没有认识主的道路。骄傲给了我们什么好处？夸耀财富又给我们带了什么？这一切都像阴影一样消逝了。} \BibleRef{智 5:3}\par

13. 但是关于祂所说的话：\BibleQuote{愿他们转向并感到羞愧}，谁不认为这是一个极其公义的惩罚：他们该当得到一个转向混乱，因为他们不愿有一个转向救恩？此后他加上了\Quote{极其迅速}。因为当审判的日子不再被期待的时候，当他们会说：\BibleQuote{和平，然后突然的毁灭就会临到他们身上}（\BibleRef{得前 5:3}）。现在，无论它何时到来，它都来得非常迅速——我们对它的到来完全不再期待；而除了活下去的希望之外，没有任何事物能使此生之漫长被感受到。因为没有比其中已经过去的一切更显得迅速的了。因此，当审判的日子到来时，那时罪人们将会感到所有逝去的生命并非长久。那在他们未经渴望、确切说是未经相信的情况下就已来临的，将丝毫不会显得来迟了。虽然此处也可这样理解：既然天主已听到了——可以说——她的呻吟和长久频繁的眼泪，她可以被理解为从她的罪中得释放，并驯服了肉性情欲的每一个失序冲动；正如她所说：\BibleQuote{离开我吧，所有行不义的人，因为主已听到了我哭泣的声音}。当她有了这幸福的结局，她若已如此成全以至于为仇敌祈祷，就不足为奇了。那么，\BibleQuote{愿我所有的仇敌都蒙羞并苦恼}这话可能有此含义：即他们应悔改自己的罪，没有混乱和苦恼便无法成就此事。那么，没有什么阻止我们也将下文作此理解：\BibleQuote{愿他们转向并感到羞愧}，即愿他们转向天主，并为他们曾在以往罪的黑暗中夸耀而感到羞愧；正如宗徒所说：\BibleQuote{那时你们在那些如今羞愧之事上有什么光荣呢？}（\BibleRef{罗 6:21}）。但关于他加上的\Quote{极其迅速}，它必须或者归于她愿望的热切情怀，或者归于基督的权能；祂以如此迅速的时间，将那些曾为着偶像的缘故而迫害教会的万民，归化于福音的信德。\par

\SourceRecord{Translated by an anonymous scholar.；Nicene and Post-Nicene Fathers, First Series；Vol. 8.；Edited by Philip Schaff.；Buffalo, NY: Christian Literature Publishing Co.,；1888.；<http://www.newadvent.org/fathers/1801006.htm>.}

% structure: p0001 source=h1 target=section reason=page-title

\section{第七篇圣咏释义}

\BibleQuote{达味自己的圣咏，他因革米尼人\Person{胡史}的话而向上主歌唱。}\par

1. 现在，赋予此篇预言的机会，在\BibleBook{列王纪}{下}中很容易辨认。\BibleRef{撒下 15:34} 因为在那里，达味王的友人\Person{胡史}投奔了其子\Person{阿贝沙隆}一方，他正与父亲作战，目的是为了发现并报告他正针对父亲所采取的计谋，为\Person{阿希托费耳}所煽动，此人背叛了\Person{达味}的友谊，并以他的谋略尽最大努力教唆儿子反对父亲。但在此篇圣咏中，所应当考虑的并非故事本身（先知由此取用了奥秘的遮蔽），如果我们已转向基督，就让这遮蔽除去。\BibleRef{格后 3:16} 首先，让我们探究这些名字本身的含义。因为不乏有诠释者，他们不是按文字、肉体地，而是属神地探究这些相同的词语，宣告给我们：\Person{胡史}应被解释为寂静，\Person{革米尼}为右手，\Person{阿希托费耳}为兄弟的毁灭。在这些解释中，那叛徒\Person{犹达斯}又出现在我们面前，而\Person{阿贝沙隆}应承继他的形象，根据对其作为父之和平的解释：因为他的父亲对他满怀和平的意念，尽管他在诡诈中内心怀有战争，这已在\BibleBook{圣咏}{第三篇}中论及。正如我们在福音中看到，主\Person{耶稣基督}的门徒被称为子女，\BibleRef{玛 9:15} 同样在福音中我们发现他们也被称为兄弟。因为主在复活时说：\BibleQuote{去告诉我的弟兄们。}\BibleRef{若 20:17} 而宗徒也称祂为\BibleQuote{在许多弟兄中作首生者。}因此，那位出卖祂的门徒的毁灭，被恰当理解为兄弟的毁灭，即我们所指的\Person{阿希托费耳}的解释。至于\Person{胡史}，从寂静的解释，可以恰当理解为主在寂静中对抗那诡计，即在最深奥的奥秘中，藉此\BibleQuote{以色列的一部分变得眼瞎，}\BibleRef{罗 11:25} 当他们迫害主时，为的是让外邦人的丰盈进入，\BibleQuote{如此全以色列就得救。}当宗徒触及这深奥的奥秘和深邃的寂静时，他惊叹道，仿佛被其深度所震慑：\BibleQuote{啊，天主的智慧和知识的财富是多么深奥！祂的审判是多么不可测度，祂的道路是多么不可追寻！谁曾知道主的心意，或是谁作过祂的谋士？}\BibleRef{罗 11:33} 因此，那伟大的寂静，他并非通过解释而发现，而是通过赞叹彰显了其伟大。在这寂静中，主隐藏了祂可敬受难的圣事，将兄弟自愿的毁灭，即那出卖者不虔敬的邪恶，转变成祂仁慈和眷顾的秩序：使那怀着悖逆心意旨在毁灭一人之物，经由眷顾的掌管而被安排为所有人的救恩。因此，那完美的灵魂，已堪当认识天主的奥秘，向上主歌唱圣咏，她\BibleQuote{因\Person{胡史}的话}而歌唱，因为她已获知那寂静的话语：因为在不信者和迫害者中，有那寂静和奥秘。但在祂自己的人中间，他们对祂说：\BibleQuote{现在我不再称你们为仆人，因为仆人不知道主人所做的事；但我已称你们为朋友，因为我从我的父那里所听到的一切，我都已使你们知道。}\BibleRef{若 15:15} 在祂的朋友中间，我再说，没有寂静，但有寂静的话语，即那寂静的意义被陈述并显明。这寂静，即\Person{胡史}，被称为\Person{革米尼}的儿子，即右手之子。因为为圣徒所做的事，不应向他们隐藏。然而祂说：\BibleQuote{不要让左手知道右手所做的事。}\BibleRef{玛 6:3} 因此，那完美的灵魂，那奥秘已使她知道，她在预言中\BibleQuote{因\Person{胡史}的话}而歌唱，即因那同一奥秘的知识。那奥秘，天主在她右边，即对她恩慈和眷顾，已作成。因此这寂静被称为右手之子，即\Quote{\Person{胡史}，\Person{革米尼}之子。}\par

2. \BibleQuote{上主，我的天主，我投靠了祢；求祢救我脱离所有迫害我的人，并拯救我}\BibleRef{咏 7:1} 如同一个已经成全的人，克服了所有恶习的战争和敌意后，再也没有敌人只剩下嫉妒的魔鬼，他说：\BibleQuote{求祢救我脱离所有迫害我的人，并拯救我}\BibleRef{咏 7:2} \BibleQuote{免得他像狮子一样撕裂我的灵魂。}宗徒说：\BibleQuote{你们的仇敌魔鬼，如同咆哮的狮子巡游，寻找可吞食的人。}\BibleRef{伯前 5:8} 因此，当圣咏作者用复数说\BibleQuote{求祢救我脱离所有迫害我的人}时，随后引入单数，说\BibleQuote{免得他像狮子一样撕裂我的灵魂。}因为他没有说：免得他们撕裂；他知道那完美的灵魂所剩下的敌人和凶恶的对手是谁。\BibleQuote{若无人救赎，也无人拯救}，即免得他撕裂我，而祢不救赎，也不拯救。因为，如果天主不救赎，也不拯救，他便撕裂。\par

3. 为了表明那已经成全的灵魂——她只需提防魔鬼最阴险的陷阱——所说的这话，请看下文。\BibleQuote{主、我的天主，如果我做了这事} \BibleRef{咏 7:3}。他称什么为\Quote{这事}？既然他没有指名那罪，我们是否要理解为一般的罪？如果这个意思令我们不悦，我们可以理解为以下所指：就好像我们问过：你所说的\Quote{这事}是什么？他回答说：\BibleQuote{如果我手中有不义}。那么显而易见，这话是针对一切罪恶说的：\BibleQuote{如果我报复了那以恶报我的人} \BibleRef{咏 7:4}。除了成全的人，没有人能真实地说这话。因为主这样说：\BibleQuote{你们应当是成全的，如同你们在天上的圣父那样；祂使太阳上升，光照善人也光照恶人，降雨给义人也给不义的人。} \BibleRef{玛 5:43}。那么，那不以恶报恶的人，便是成全的。因此，当那成全的灵魂祈求\Quote{关于\Person{耶米尼}的儿子\Person{库史}的话}时，也就是说，祈求认识那个秘密和沉默——主慈悯我们、恩待我们，为了我们的救恩而成就了它，以便容忍并忍耐地承受这个背叛者的诡计——如同主向这个成全的灵魂解释这个秘密的用意时所说：因为你是不敬虔的罪人，为要使你的不义被我的流血洗去，我在极大的沉默和极大的忍耐中容忍了我的背叛者；你难道不效法我，好叫你也以恶报恶吗？于是，圣咏作者思量并理解主为他所做的，并因主的榜样而走向成全，他说：\BibleQuote{如果我报复了那以恶报我的人}，也就是说，如果我没有做你以榜样教训我的事：\BibleQuote{就让我因此空虚地倒在我的仇敌面前}。他说得好，不是\Quote{如果我报复了那对我作恶的人}；而是\Quote{报答的人}。因为那报答的人，已经先领受了某些东西。现在，一个更大的忍耐的例子是，甚至不报复那在领受恩惠之后还以恶报善的人，胜过若他先前没有领受任何恩惠却有意伤害别人。因此，如果他说：\BibleQuote{我报复了那以恶报我的人}，也就是说，如果我没有在那沉默中效法你，就是在你所为我成就的忍耐中效法你，\BibleQuote{就让我空虚地倒在我的仇敌面前}。因为他是空虚的自夸者，他自己是人，却想要向人报复；当他公开试图胜过一个人时，暗中却被魔鬼胜过，被虚妄而骄傲的喜乐变为空虚，因为他仿佛不能被战胜。圣咏作者知道在何处可以获得更大的胜利，也知道\BibleQuote{那在暗中看见的圣父将报答}之处。\BibleRef{玛 6:6}。免得他报复那以恶报他的人，他宁愿胜过自己的愤怒而不是别人，也被那些圣书所教导，其中写道：\BibleQuote{那克服自己愤怒的，胜过取城的人。} \BibleRef{箴 16:32}。\BibleQuote{如果我报复了那以恶报我的人，就让我因此空虚地倒在我的仇敌面前。}他似乎是借诅咒的方式起誓，这是最重的一种誓愿，如同有人说：如果我做了某事，就让我受某种苦。但在发誓者口中起誓是一回事，在先知的意思中是另一回事。因为在此他提及实际会临到那报复以恶报他们的人身上的事，而不是像起誓那样，为自己或别人招来诅咒。\par

4. 因此，\BibleQuote{愿仇敌迫害我的灵魂并捕捉它} \BibleRef{咏 7:5}。他再次以单数提及仇敌，更清楚地指出他上文称为狮子的那一位。因为他迫害灵魂，如果欺骗了它，就会捕捉它。因为人愤怒的极限是身体的毁灭；但灵魂在可见的死亡之后，他们不能将它留在自己的权下；然而魔鬼通过它的迫害所捕捉到的任何灵魂，它都会留住。\BibleQuote{愿他在地上践踏我的生命}：也就是说，通过践踏使我的生命成为尘土，也就是它的食物。因为它不仅被称为狮子，也被称为蛇，对它曾有话说：\BibleQuote{你要吃土} \BibleRef{创 3:14}。对罪人也有话说：\BibleQuote{你是土，你也要归于土} \BibleRef{创 3:19}。\BibleQuote{愿他将我的荣耀带到尘土中}。这就是\Quote{风从地面吹起}的尘土，即骄傲者虚妄而愚蠢的自夸，膨胀而没有实质重量，如同一团被风卷起的灰尘。所以，他在这里恰当地提到了荣耀，他不愿意将它带到尘土中。因为他希望荣耀在良心面前坚定地建立在天主面前，那里没有自夸。宗徒说：\BibleQuote{那夸耀的，要在主内夸耀} \BibleRef{格前 1:31}。如果一个人因骄傲而轻视良心的隐秘处——在那里只有天主考验人——并希望在世人面前夸耀，这坚定性就被带到尘土中。因此，圣咏作者在别处说：\Quote{天主将打碎取悦于人的骨头}。现在，那已很好学习或体验到克服恶习步骤的人就知道，这虚荣的恶习要么是唯独，要么是首要的，成全的人必须避免它。因为灵魂最初跌倒之处，她最后才得胜。\Quote{因为一切罪恶的开端是骄傲}；又说：\Quote{人骄傲的开端是离弃天主}。\par

5. \BibleQuote{主，求祢在愤怒中兴起}\BibleRef{咏 7:6}。为什么我们称为成全的人，还要激怒天主呢？我们岂不应当看看，那被石头砸死时说\BibleQuote{主，不要将这罪归于他们}？\BibleRef{宗 7:60} 或者，圣咏作者这样祈祷，不是反对人，而是反对魔鬼及其使者；罪人和不敬神的人乃是魔鬼的产业吗？那么，凡是祈祷说那产业从魔鬼手中被夺去的人，不是在愤怒中向他祈祷，而是在怜悯中，因那\BibleQuote{使不敬神的人成义}\BibleRef{罗 4:5}的主。因为当不敬神的人成义时，他从不敬神变为正义，从魔鬼的产业变为天主的圣殿。既然一个产业（某人渴望在其中统治）被夺走是一种惩罚：圣咏作者称这惩罚——即魔鬼不再占有他现在所占有的人——为天主对魔鬼的愤怒。\Quote{主，求祢在愤怒中兴起。}\Quote{兴起}（他用作\Quote{显现}），这是用了人的、隐晦的说法；好像天主睡着了，当祂在祂隐秘的作为中不被认识、不被看见时。\Quote{求祢在我仇敌的边界上被高举。}他所说的边界就是指那产业本身，他希望天主在那产业中被高举，即被尊崇、受光荣，而不是魔鬼；当不敬神的人成义并赞美天主时。\Quote{主，我的天主，求祢在祢所定的诫命中兴起：}这就是说，既然祢命令了谦卑，求祢在谦卑中显现；并首先履行祢所命令的；好使人因祢的榜样克服骄傲，而不被那违背祢诫命劝人骄傲的魔鬼所占有，魔鬼说：\BibleQuote{吃吧，你们的眼就会开了，你们就会像天主一样。}\BibleRef{创 3:5}\par

6. \Quote{百姓的集会要环绕祢。}这可以有两种理解。因为百姓的集会可以指信从的人，也可以指迫害的人，两者都在我们主的同一谦卑中发生了：迫害的人群因鄙视祂而环绕祂；关于这，经上说：\Quote{外邦人为什么喧嚷，万民为什么谋算虚妄的事？}但藉着祂的谦卑而信从的人如此众多地环绕祂，以至于可以最真实地说：\BibleQuote{部分硬心临到了以色列，直到外邦人的总数全数进来；}\BibleRef{罗 11:25} 又说：\Quote{你向我请求，我必将外邦人赐你作产业，将地极赐你作基业。}\Quote{祢要为他们的缘故回到至高之处：}也就是说，为了这集会的缘故，祢回到至高之处：这被理解成祂复活升天时所行的。因为这样受光荣后，祂赐下了圣神；在祂被高举之前，圣神是不能赐下的，正如福音所记载的：\BibleQuote{那时圣神还没有赐下，因为耶稣还没有受光荣。}\BibleRef{若 7:39} 然后，祂为了百姓的集会而回到至高之处，就派遣了圣神；传福音的人被圣神充满，从而用教会充满了全世界。\par

7. 这也可从另一意义来理解：\BibleQuote{主，求祢在祢的义怒中奋起，在我仇敌的领土上被高举。}也就是说，求祢在祢的义怒中奋起，不让我的仇敌领悟祢；这样，所谓\BibleQuote{被高举}，就是成为崇高，使祢不被理解；这与前面所说的静默有关。因为另一篇圣咏论及这高举时说：\BibleQuote{祂乘驾革鲁宾，飞翔。}以及\BibleQuote{祂以黑暗为祂的隐秘处。}在这高举或掩藏中，当他们因自己应得的罪罚而不领悟祢——那些将要钉死祢的人——信者的会众\BibleQuote{必环绕祢。}因为祂正是在祂的卑抑中受到高举，亦即不被理解。所以，\BibleQuote{主，我的天主，求祢在祢所赐的诫命中奋起。}可能指这一点：当祢彰显自己时，成为高处或深处，使我的仇敌不能领悟祢。罪人是义人的仇敌，不虔敬的人是虔诚人的仇敌。\BibleQuote{万民的会众必环绕祢。}这就是说，正因那些钉死祢的人不了解祢，外邦人将信从祢，如此\BibleQuote{万民的会众必环绕祢。}但接下来的话，若这真是其含义，则已开始感受到的痛苦多于因理解而得的喜乐。因为接下来是：\BibleQuote{求祢因他们而返回高处。}也就是说，因这人类会众的缘故——教会正为此充塞——求祢返回高处，即再次不被理解。那么，\BibleQuote{因他们}是什么意思呢？无非是这同样的会众也将得罪祢，以致祢最能真实地预告并说：\BibleQuote{你想，当人子来临时，祂在世上能找到信德吗？}\BibleRef{路 18:8}再者，关于被视为异端者的假先知，祂说：\BibleQuote{由于罪恶增多，许多人的爱德渐渐冷淡。}\BibleRef{玛 24:12}既然即使在教会内——即在基督之名已极其广泛传播的万民万邦的会众中——也将有如此众多的罪人，这在很大程度上已可察觉；岂不正是此处所预言、另一先知也曾警告过的圣言的饥荒\BibleRef{亚 8:11}吗？难道不也是因这会众的缘故——他们因自己的罪而使自己远离真理的光明——天主返回高处，即信德纯洁、涤除一切邪僻见解的腐化后，要么完全不被持守和接受，要么仅被极少数人持守，关于这些人曾说过：\BibleQuote{惟有坚持到底的，才能得救。}\BibleRef{谷 13:13}所以并非无因地说：\BibleQuote{因这}会众\BibleQuote{求祢返回高处。}也就是说，再次隐退到祢隐秘的深处，甚至是因这万民会众的缘故，他们拥有祢的名，却不实行祢的作为。\par

8. 但无论是此处的第一种解释，还是最后一种更为合适，且不影响任何更好、相同或同等好的解释，下文都非常一致地说：\BibleQuote{主审判万民。}因为无论祂是在复活后升天回到高处，下文都很合适地接着说：\BibleQuote{主审判万民。}因为祂将从那里来审判活人死人。或者，祂返回高处，是指当对真理的领悟离开了有罪的基督徒时，因为论及祂的来临曾说过：\BibleQuote{你想，当人子来临时，祂在世上能找到信德吗？}\BibleRef{路 18:8}然后\BibleQuote{主}就\BibleQuote{审判万民。}这是什么主？岂不是耶稣基督？\BibleQuote{因为父不审判任何人，但把一切审判权交给了子。}\BibleRef{若 5:22}因此，这完美祈祷的灵魂，看她是多么不惧怕审判之日，并怀着真正稳妥的渴望在祈祷中说：\BibleQuote{愿祢的国来临：审判我，}她说，\BibleQuote{主，请按我的义德审判我。}在前一篇圣咏中，一个软弱的人恳求，与其说是提及自己的功劳，不如说是祈求天主的慈悲：因为天主子来，是\BibleQuote{为召叫罪人悔改。}\BibleRef{玛 9:13}因此他在那里曾说过：\BibleQuote{主，求祢因祢的慈悲拯救我。}也就是说，不是因我的功德。但现在，既然被召叫后，他持守并遵守了所领受的诫命，就敢说：\BibleQuote{主，请按我的义德，并按我的无害来审判我，这义德与无害在我身上。}这是真正的无害，连仇敌也不伤害。因此，那能真实地说\BibleQuote{我若以恶报恶}的人，要求按他的无害受审判是合理的。至于他所加的\BibleQuote{在我身上}，不仅可以指无害，也可以理解为指义德；意思是：主，请按我的义德，并按我的无害来审判我，这义德与无害在我身上。通过这个附加语，他表明：灵魂之所以是义和无害的，并非出于自身，而是出于那赐予光明和亮光的天主。因为关于这一点，他在另一篇圣咏中说：\BibleQuote{主，祢必照亮我的烛光。}关于若望，经上说：\BibleQuote{他不是那光，但为光作证。}\BibleRef{若 1:8}\BibleQuote{他是点着和照耀的灯。}\BibleRef{若 5:35}那光，就是灵魂如灯被点燃的源头，所发出的不是借来的光明，而是本源的光明，这光就是真理本身。所以这话是这样说的：\BibleQuote{按我的义德，并按我的无害，这在我身上。}好像点着和照耀的灯在说：请按在我身上的火焰审判我，这不是我自己的火焰，而是你点燃我使我发光的火焰。\par

9. \BibleQuote{但願惡人的邪惡終結}（見：\BibleRef{詠 7:9}）。他說，\Quote{終結}，即完成，正如\WorkTitle{默示錄}所言：\BibleQuote{讓義人更義，讓污穢者仍舊污穢}（見：\BibleRef{默 22:11}）。因為那些釘死天主之子的人，他們的邪惡似乎終極了；但那些不願正直生活、憎恨真理誡命的人，他們的邪惡更大——因為天主之子也為他們被釘死。\Quote{但願惡人的邪惡}，他接著說，\Quote{終結}，即達到邪惡的頂點，以便正義的審判能隨即到來。但既然不僅說\BibleQuote{讓污穢者仍舊污穢}（見：\BibleRef{默 22:11}），也說\BibleQuote{讓義人更義}（見：\BibleRef{默 22:11}），他接著說：\BibleQuote{而祢，察看人心腎腑的天主，將引導義人。} 那麼，義人如何被引導，豈不是暗中進行？當基督徒時代之初，聖人們尚受世俗之人的迫害時，在人看來似乎奇妙的事，如今基督徒名號已開始獲得如此崇高的尊嚴，假冒為善（即偽裝）增多；我所指的是那些藉著基督徒身份寧願討人喜歡而不是天主喜歡的人。那麼，在如此混亂的偽裝中，義人如何被引導？唯有當天主察看人心腎腑，看見所有人的思念（以\Quote{心}一字表示），和他們的喜樂（以\Quote{腎}一字理解）。因為對暫時和屬地事物的喜樂正當地歸於腎；因為它是人較低的部分，也是肉體生殖快感的所在，藉此人的本性藉著種族延續轉入這充滿憂慮和欺騙性喜樂的生命中。所以，天主察看我們的心，看出我們的財寶在那裡，即在天上；也察看腎，看出我們不隨從血肉，而在主內喜樂。祂就在內在良心中引導義人——在祂面前，沒有人看見，唯有祂獨自察知各人思念什麼、喜樂什麼。因為喜樂是憂慮的終點；因為每個人藉著憂慮和思念努力達到那終點，即他的喜樂。所以，那察看心的，看見我們的憂慮。那細察腎的，也看見憂慮的終點（即喜樂）；當祂發現我們的憂慮既不傾向\BibleQuote{肉體的貪慾，也不傾向眼目的貪慾，也不傾向生活的驕傲}（見：\BibleRef{若一 2:16}）——這些都如影消逝——而是提升到永恆事物的喜樂（不受任何變化損害）時，祂，察看人心腎腑的天主，將引導義人。因為我們的行為（在言行中所做的）或許可以被人知道；但以怎樣的心態去做，以及藉此我們企圖達到怎樣的終點，只有祂知道，那察看人心腎腑的天主。\par

10. \BibleQuote{我的正義幫助來自上主，祂使心裡正直的人完全}（見：\BibleRef{詠 7:10}）。醫治的職能是雙重的：一是醫治軟弱，二是保持健康。根據前者，在前一篇聖詠中說：\BibleQuote{上主，求祢憐憫我，因為我軟弱}；根據後者，在本聖詠中說：\BibleQuote{若我手中有不義，若我報復了那以惡待我的人，就讓我因敵人而空手敗亡}。因為那裡軟弱的人祈求被解救，這裡已經完全的人祈求不再變壞。根據前者，那裡說：\BibleQuote{求祢因祢的仁慈使我痊癒}；根據後者，這裡說：\BibleQuote{上主，請按我的正義審判我}。因為那裡他請求藥方以脫離疾病；這裡則請求保護以免陷入疾病。根據前者說：\BibleQuote{上主，求祢按祢的仁慈使我痊癒}；根據後者說：\BibleQuote{我的正義幫助來自上主，祂使心裡正直的人完全}。兩者都使人完全；但前者使人脫離疾病進入健康，後者則保守人在健康中。因此那裡的幫助是仁慈的，因為罪人沒有功勞，他尚渴望成義，\BibleQuote{信賴那使不虔敬者成義的天主}（見：\BibleRef{羅 4:5}）；但這裡的幫助是正義的，因為賜給了已義的人。因此，那說過\BibleQuote{我軟弱}的罪人，首先說：\BibleQuote{上主，求祢因祢的仁慈使我痊癒}；而在這裡，那說過\BibleQuote{若我報復了那以惡待我的人}的義人說：\BibleQuote{我的正義幫助來自上主，祂使心裡正直的人完全}。因為如果他提供了藥方，使我們在軟弱時得醫治，何況那使我們保持健康的藥方呢？因為\BibleQuote{當我們還是罪人時，基督為我們死了，何況現在既已成義，豈不更要藉著祂免去天主的義怒嗎？}（見：\BibleRef{羅 5:8}）\par

11. \Quote{我的义援来自主，祂使心里正直的人痊愈。} 天主，祂洞察人心和肾脏，引导义人；但祂以义援使心里正直的人痊愈。祂并不像祂洞察人心和肾脏那样使心里正直的人和肾脏痊愈；因为在一个败坏的心中，思想是恶的，在一个正直的心中，思想是善的；但不好的快乐属于肾脏，因为它们更为低级和属世；而那些好的快乐不属于肾脏，而是属于心本身。因此，人不能被称为肾脏正直，如同他们被称为心里正直，因为哪里有了思想，哪里立刻也有快乐；除非当思想神性和永恒的事物时，这才能成立。他说：\Quote{祢赐给我心中的喜乐}，当他说了：\Quote{主，祢面容的光辉已印在我们身上}。因为虽然对暂时事物的幻影，当灵魂被虚妄和必死的希望抛来抛去时，常常给虚妄的想象带来狂乱和疯狂的快乐；但这快乐必须归于肾脏，而非心；因为这些想象都来自低级的事物，即属地和属肉性的事物。因此，那洞察人心和肾脏的天主，在内心察觉正直的思想，在肾脏中察觉没有快乐，就将义援赐给心里正直的人，在那里属天的快乐与清洁的思想结合在一起。因此，当他在另一篇圣咏中说：\Quote{此外，今夜我的肾脏也责备了我}；他就接着说到援助：\Quote{我常将主摆在我眼前，因祂在我右边，我便不致摇动}。这里他表明自己只受到来自肾脏的暗示，而非快乐；因为如果他受到了这些快乐，那么他当然会被摇动。但他说：\Quote{主在我右边，我便不致摇动}；然后他补充道：\Quote{因此我的心喜乐}；这样肾脏能够责备他，却不能使他快乐。因此快乐不是产生在肾脏中，而是在那里，即心里，在那里，面对肾脏的责备，天主被预见到在右边。\par

12. \BibleQuote{天主是正义的审判者，坚强（忍耐）而宽仁} \BibleRef{咏 7:11}。什么审判者是主，审判万民的那位？祂是正义的；祂\BibleQuote{要按各人的行为报应各人} \BibleRef{玛 16:27}。祂是坚强（忍耐）的；祂虽全能，却为了我们的救恩容忍了不敬神的迫害者。祂是宽仁的；祂复活后并没有立即加速惩罚那些迫害祂的人，而是容忍他们，好让他们最终从不敬神中转回而得救：至今祂仍然容忍他们，将最后的惩罚保留给最后审判，直到此刻仍召叫罪人悔改。\Quote{不天天发怒}。也许\Quote{发怒}是比\Quote{愤怒}更显著的表达（就如我们在希腊文抄本中发现的）；这样怒气（祂用以惩罚的）不在祂内，而在那些服从真理诫命的执行者心中，他们被命令甚至下级的仆役（称为愤怒的天使）去惩罚罪：这些人现在也以惩罚罪人为乐，不是为了公义（他们不以此为乐），而是为了恶意。因此，天主并不\Quote{天天发怒}，即祂不是每天召集祂的仆役来报复。因为现在天主的忍耐召叫人悔改；但在最后时刻，当人们\BibleQuote{因着他们的硬心和不肯悔改的心，为自己积蓄忿怒，直到忿怒的日子，和天主公义审判显示的日子} \BibleRef{罗 2:5}，那时祂将挥动祂的剑。\par

13. \BibleQuote{你们若不悔改}，祂说，\BibleQuote{祂必挥动祂的剑} \BibleRef{咏 7:12}。主自己——那人——可以被视为天主双刃的剑，即祂的枪，在祂第一次来临时祂不挥动它，而是好像藏在谦卑的鞘中；但当祂第二次来临审判活人死人时，在祂荣耀的显赫光辉中，祂将挥动它，向义人发出光明，向不敬神的人发出恐怖。因为在其他抄本中，代替\Quote{祂必挥动祂的剑}，写作了\Quote{祂必使祂的枪明亮}：我认为这个词语最恰当地表明了主荣耀的最后来临：因为另一篇圣咏所说的：\Quote{主，求祢从恶人手中拯救我的灵魂，祢的枪从敌人的手中拯救我}，也被理解为指祂的位格。\Quote{祂已拉紧祂的弓，并预备好了}。词语的时态不能完全忽略，他如何用将来时提到\Quote{剑}（\Quote{祂必挥动}），用过去时提到\Quote{弓}（\Quote{祂已拉紧}）：这些过去时的词语随后出现了。\par

14. \BibleQuote{祂在其中预备了死亡的器具，祂为燃烧制成了祂的箭} \BibleRef{咏 7:13}。我乐意将那弓理解为圣经，在其中，藉着新约的力量（如同弓弦），旧约的坚硬被弯曲并制服。从那里，宗徒们被遣发如同箭矢，或神圣的宣讲被射出。这些箭矢，\BibleQuote{祂为燃烧制成了}，即这些箭矢，使人被击中后为天上的爱所燃烧。因为那说\BibleQuote{领我进入酒房，将我置于香料中，使我拥挤在蜜中，因我被爱所伤}的妇人，是被别的什么箭击中的吗？那渴望回归天主、从迷途返回、祈求帮助对抗诡诈的舌头的人，以及那听到\BibleQuote{要给你什么，或加给你什么以对抗诡诈的舌头？强者的利箭，伴有毁灭性的火炭}的人，是被别的什么箭点燃的吗？即火炭，当你被击中并点燃时，会如此热爱天国，以致蔑视所有抵抗你、试图将你从目标中召回之人的舌头，并嘲笑他们的迫害，说：\BibleQuote{谁能使我们与基督的爱隔绝呢？是患难吗？是困苦吗？是迫害吗？是饥饿吗？是赤身露体吗？是危险吗？是刀剑吗？因为我深信}，他说，\BibleQuote{无论是死，是生，是天使，是掌权者，是现在的事，是将来的事，是有能者，是高处的，是深处的，或是任何别的受造之物，都不能使我们与天主的爱隔绝，这爱是在我们的主耶稣基督内的。}因此，祂为燃烧制成了祂的箭。因为在希腊文抄本中是这样写的：\BibleQuote{祂为燃烧制成了祂的箭。}但大多数拉丁文抄本写作\BibleQuote{燃烧的箭。}然而，无论是箭本身燃烧，还是使他人燃烧（除非自身燃烧，否则无法使他人燃烧），这两种意义都是完整的。\par

15. 但既然他说主不仅准备了箭，还在弓中准备了\BibleQuote{死亡的器具}，那么可以问：什么是\BibleQuote{死亡的器具}？或许是异端者？因为同样从那弓，即从那同圣经，他们射中灵魂，不是为以爱点燃，而是为以毒药毁灭：这并非无故发生，而是按他们的功过；因此，这安排也应归于天主眷顾，并非祂使人成为罪人，而是按他们犯罪后的状态来安排。因为罪使他们怀着恶意理解圣经，这本身成为罪的惩罚；然而，藉着他们的死亡，天主教会的儿子们，如同被某些荆棘（这样说吧）从沉睡中唤醒，并在理解圣经上取得进步。\BibleQuote{因为必须也有异端，使那些被认可的人}，他说，\BibleQuote{在你们中间显明出来} \BibleRef{格前 11:19}，即在人中间显明，因为他们在天主面前原是显明的。或者，祂是否特意将同一箭矢安排为既是不信者之毁灭的死亡器具，又制作为燃烧的或为燃烧的箭，以锻炼信友？因为宗徒所说的并非虚假：\BibleQuote{对于后者，我们是出于死的香气使他死；对于前者，我们是出于活的香气使他活；对这些事，谁可胜任呢？} \BibleRef{格后 2:16}。所以，如果同一宗徒们对于那些迫害他们的人是死亡器具，对于信友的心是燃烧的箭，这并不奇怪。\par

16. 在此安排之后，公义的审判将到来：圣咏作者如此谈论，使我们明白，每个人的惩罚来自他自己的罪，他的不义转为报复；使我们不认为天主那宁静和不可言说的光从自身产生惩罚罪的手段，而是祂如此安排罪，使人在犯罪时所喜悦的，成为主报复的工具。\BibleQuote{看啊}，他说，\BibleQuote{他怀了不义。}他怀了什么，以致怀了不义？\BibleQuote{他怀了}，他说，\BibleQuote{劳苦。}由此而来：\BibleQuote{你必须流汗才能吃你的饼。} \BibleRef{创 3:17}。由此也有：\BibleQuote{凡劳苦和负重担的，你们都到我这里来；因为我的轭是柔和的，我的担子是轻的。}因为除非一个人爱那不能违背其意愿被夺走的东西，否则劳苦永不停止。因为当人所爱的是那些我们可能违背意愿失去的东西时，他们必须最悲惨地劳苦；为获得它们，在人世忧虑的困厄中，每个人都想为自己抢夺，想抢先他人，或从他人手中夺取，必然图谋不义。所以，他怀了不义是恰如其分而又很有次序的，因为他怀了劳苦。他生下什么？不就是他所怀的吗？虽然他所怀的不是他怀的。因为那未怀的不会出生；但怀的是种子，从种子形成并出生。那么劳苦是不义的种子，而罪是劳苦的孕育，即那最初的罪——\BibleQuote{离开天主}。\BibleRef{德 10:12}因此，那怀了劳苦的人，便怀了不义。\BibleQuote{他生下了不义。}\Quote{不义}与\Quote{不公}相同；所以他生下了他所怀的。接下来是什么？\par

17. \BibleQuote{他挖了坑，也掘深了它} \BibleRef{咏 7:15}。挖坑，就是在世俗事务中，也就是在地上，预备诡计，使不义者所愿欺骗的人掉在其中。这坑在同意世俗情欲的邪恶建议时被挖开；但在同意之后我们着手实际行骗时，它才被掘深。然而，不义岂非更多地伤害它所针对的义人，而非产生它的不义之心吗？例如，偷钱的人，当他想要对他人造成痛苦伤害时，他自己却被贪婪之伤所残害。谁，即使心智正常，看不出这两人之间的巨大差别：一人损失金钱，另一人损失清白？\BibleQuote{他必掉进} 然后 \BibleQuote{自己所挖的坑中}。正如另一篇圣咏所说：\BibleQuote{主在施行审判时被人认识；罪人被自己手所作的缠住。}\par

18. \BibleQuote{他的劳苦必转到他头上，他的不义必落在他头顶上} \BibleRef{咏 7:16}。因为他无意逃脱罪恶，而是被带在罪恶之下为奴，可以这样说，正如主说：\BibleQuote{凡犯罪的，就是奴隶} \BibleRef{若 8:34}。他的不义就必在他身上，当他受制于其不义时；因为他不能向主说，如同无辜和正直的人所说：\Quote{我的荣耀，使我抬起头来的那位}。那时他将这样在下面，以致他的不义在上面，并落到他身上；因为这不义压住他，拖累他，不容他飞回圣人们的安息。当在一个未受规训的人身上，理性是奴隶，情欲掌权时，就会发生这事。\par

19. \BibleQuote{我要依照祂的正义向主称谢} \BibleRef{咏 7:17}。这不是罪人的告解：因为说这话的人，是上面极真诚地说过\Quote{若我手中有不义}的那位；但这是对天主正义的宣认，我们如此说：的确，主啊，祢是正义的，因为祢既这样保护义人，使祢亲自光照他们；又如此安排罪人，叫他们受惩罚，不是出于祢的恶意，而是出于他们自己的恶意。这种宣认如此赞美主，致使不敬虔者的亵渎毫无效力；他们愿意为自己的恶行开脱，不愿将犯罪归咎于自己的过错，也就是不愿将自己的错归咎于自己的错。因此他们要么指控命运或运数，要么指控魔鬼——创造我们的主愿意我们有权拒绝同意魔鬼；要么他们引入另一种本性，不是出于天主：可怜动摇的人，迷途的人，而不是向天主承认，好使祂宽恕他们。因为除了说\Quote{我犯了罪}的人，没有人配得宽恕。因此，凡看见灵魂的功过被天主如此安排，以致每个人都被赐予他自己的份，宇宙的美善在任何部分都没有被破坏的人，就在一切事上赞美天主：这不是罪人的告解，而是义人的告解。因为当主说\BibleQuote{父啊，天地的主宰，我称谢祢，因为祢将这些事瞒住了智慧人，而启示给了小孩子} \BibleRef{玛 11:25}时，这并不是罪人的告解。同样在德训篇中记载：\BibleQuote{你们要在主的一切工程上称谢主；在称谢时你们要说：主的一切工程都是极好的。}在这篇圣咏中可以看出，如果有人怀着虔诚的心，藉着主的帮助，分辨义人的赏报和罪人的惩罚，就会发现，在这二者中，天主所造并治理的整个受造界，被一种奇妙而少为人知的美所装饰。因此他说：\BibleQuote{我要依照祂的正义向主称谢}，如同看见了黑暗不是天主造的，却仍被安排的人。因为天主说：\BibleQuote{有光，就有了光} \BibleRef{创 1:3}。祂没有说：有黑暗，就有了黑暗；但祂却安排了黑暗。因此经上说：\BibleQuote{天主分开光明和黑暗，天主称光明为昼，称黑暗为夜} \BibleRef{创 1:4}。这就是区别：祂造了光明并安排了它；但黑暗祂却没有造，却也安排了它。如今罪恶由黑暗所象征，正如先知所说：\BibleQuote{你的黑暗将如正午} \BibleRef{依 58:10}；以及宗徒所说：\BibleQuote{恼恨自己弟兄的，就是在黑暗中} \BibleRef{若一 2:11}；尤其是那段经文：\BibleQuote{让我们脱去黑暗的行为，穿上光明的武器} \BibleRef{罗 13:12}。这并不是说黑暗有任何本性。因为一切本性，就其是本性而言，必然是存在。存在属于光明，不存在属于黑暗。因此，那离开造他的主，而转向他受造的来源，即虚无的人，就在这罪中陷入黑暗；然而他并未完全毁灭，而是被安置在最卑微的事物中。因此，圣咏作者说了\BibleQuote{我要向主称谢}之后，免得我们理解为罪人的告解，他最后加上\BibleQuote{我要歌颂至高主的名}。因为歌颂与喜乐有关，而忏悔则与悲伤有关。\par

20. 这篇圣咏也可以取主基督的人性角度：只要其中在卑微中所说的话，被归因于祂所承担我们的软弱。\par

\SourceRecord{Translated by an anonymous scholar.；Nicene and Post-Nicene Fathers, First Series；Vol. 8.；Edited by Philip Schaff.；Buffalo, NY: Christian Literature Publishing Co.,；1888.；<http://www.newadvent.org/fathers/1801007.htm>.}

% structure: p0001 source=h1 target=section reason=page-title

\section{论圣咏第八篇}

\Emph{致终末，为榨酒池，大卫自己的圣咏。}\par

对于这首标题为\Quote{为\OriginalTerm{榨酒池}{wine-presses}}的圣咏正文，他似乎并未提及榨酒池。由此可见，圣经中常以众多不同的比喻来指代同一事物。因此我们可以将榨酒池理解为教会，正如我们也藉\OriginalTerm{打谷场}{threshing-floor}来理解教会。因为在打谷场或榨酒池中，所做的不过是清除产物的外壳；这对于农作物的初生、生长以及达到收割或葡萄收成的成熟期都是必要的。于是，这些外壳或支撑物——即打谷场上的\OriginalTerm{糠秕}{chaff}与谷粒，榨酒池中的\OriginalTerm{皮壳}{husks}与葡萄酒——被剥去。同样，在教会中，从与善人一同聚集的众多世俗之人中（因那众多的人对于他们的诞生和适应天主圣言是必要的），正发生这样的事：藉着天主仆役的运作，他们因属神的爱而被分离。因为目前的情形是，善人暂时与恶人分离，不是空间上的，而是情感上的；尽管在身体临在方面，他们在教会中彼此往来。但另一个时候将要来到，谷物将被分别存放在谷仓里，葡萄酒在酒窖里。他说：\BibleQuote{麦子要收入仓里，但糠秕要用不灭的火焚烧。}\BibleRef{路 3:17}同样的事可以用另一个比喻来理解：祂将葡萄酒收藏在酒窖里，却把皮壳丢给牲畜；这样，藉着牲畜的肚腹，我们得以用比喻理解地狱的苦痛。\par

关于榨酒池还有另一种解释，但依然保持教会的意义。因为连天主圣言也可由葡萄来理解：主甚至被称为一挂葡萄；以色列人民先前所派的人从预许之地带来挂在棍子上的，仿佛被钉在十字架上。\BibleRef{户 13:23}因此，当圣言为了彰显自己而必须使用声音，藉以传达给听者时，知识像酒一样被密封在声音的同一外壳中；于是这葡萄进入耳朵，如同进入榨酒者的压榨机。在那里分离发生，声音直达耳朵；而知识被收在听者的记忆中，如同某种酒池；再从那里进入生活的纪律和心灵的习性，如同从酒池进入酒窖：若不被疏忽而变酸，就会因年月而变得醇厚。因为它在犹太人中变酸了，他们把这酸醋给主喝。\BibleRef{若 19:29}因为那酒，就是主从新约葡萄园的出产中，要在祂父的国里与祂的圣徒一同饮用的，必定是最甘甜、最醇厚的。\BibleRef{玛 26:29}\par

榨酒池也常被视为殉道，似乎当那些宣认基督之名的人被迫害的打击所践踏时，他们必死的遗骸像皮壳般留在地上，而他们的灵魂则流向天堂居所的安息。但这种解释也未使我们脱离教会的丰饶。于是，为榨酒池——即为教会的建立——而歌唱；那时我们的主复活后升了天。因为那时祂派遣了圣神；门徒们被圣神充满，满怀信心地宣讲天主圣言，使教会得以聚集。\par

因此说：\BibleQuote{上主，我们的主，祢的名在全地何其可敬！}\BibleRef{咏 8:1}我问，祂的名如何在全地受到赞美？回答是：\BibleQuote{因为祢的光荣已被高举于诸天之上。}因此意思是：上主，我们的主，凡居住在全地的人何等敬佩祢！因为祢的光荣已从地上的卑微被提升至诸天之上。因为由此显明了祢是谁降下，而当有人看见、其余人相信祢升往何处时，这便显明了。\par

\BibleQuote{从婴孩和吃奶者的口中，祢得以完善了赞美，为因祢的仇敌}\BibleRef{咏 8:2}我不能将婴孩和吃奶者理解为别的，只理解为宗徒所说的：\BibleQuote{如对基督内的婴孩，我将奶给你们喝，不是食物。}\BibleRef{格前 3:1}这些人是那些在主前行在祂前面赞美祂的人，主曾亲自用这见证回答那些叫祂责备他们的犹太人说：\BibleQuote{你们没有读过：从婴孩和吃奶者的口中，祢得以完善了赞美？}\BibleRef{玛 21:16}祂说得合情合理，祂不是说\Quote{祢得以了赞美}，而是\Quote{祢得以完善了赞美}。因为教会中也有那些如今不再喝奶，而是吃干粮的人；同一位宗徒指出他们说：\BibleQuote{我们在完人中讲论智慧}\BibleRef{格前 2:6}；但教会并非仅靠这些人得以完善；因为如果只有这些人，对人类族类的顾及就会很少。然而，顾及发生了，当那些尚不能领会属灵和永恒之事的人，也被关于时间性历史的信德所滋养——这历史是为我们的得救，在众圣祖和众先知之后，由天主最强的能力和智慧所施行，甚至在所摄取的人性的圣事中（在其中每个相信的人都得救）；为的是每个受其权威感动的人都服从其规诫，由此被净化并\BibleQuote{在爱德中扎根奠基}，使他能不再是一个喝奶的婴孩，而是一个吃干粮的少年，与圣徒一同奔跑，\BibleQuote{领悟广度、长度、高度和深度，并认识基督那超越知识的爱}\BibleRef{弗 3:17}\par

6. \BibleQuote{由孩童和吃奶者的口中，祢建立了完美的赞颂，为了祢的仇敌。} 所谓敌人，是指那些反对通过耶稣基督和祂被钉十字架而实现的这项安排的人，我们通常应当理解，那些禁止相信未知之事（\BibleRef{格前 2:6}）并许诺确切知识的人，正如所有异端者以及那些在外邦人的迷信中被称为哲学家的人。并非许诺知识本身应受责备，而是因为他们认为最有益且必要的信德步骤应被忽略，我们必须藉此上升到某确定之事，而此事除永恒者外别无可能。由此看来，他们甚至并不拥有这种他们因鄙视信德而许诺的知识，因为他们不知道信德这个如此有用且必要的步骤。\Quote{由口中，} 因此 \BibleQuote{由孩童和吃奶者的口中，祢建立了完美的赞颂，} 祢，我们的主，首先藉宗徒宣告：\BibleQuote{除非你们相信，你们不会明白；} 又亲口说：\BibleQuote{那些没有看见而相信的，才是有福的。} \BibleRef{若 20:29} \BibleQuote{因为仇敌：} 这同样是对他们所说的：\BibleQuote{父啊！天地的主宰！我称谢祢，因为祢将这些事瞒住了智慧及明达的人，而启示给了小孩子。} \BibleRef{玛 11:25} \BibleQuote{从智慧的人，} 祂说，并非真智慧者，而是那些自以为智慧的人。\BibleQuote{使祢能摧毁敌人和辩护者。} 这不是指异端者吗？因为他既是敌人又是辩护者，当他攻击基督信仰时，似乎是在辩护。虽然这世界的哲学家也完全可以被视为敌人和辩护者；因为天主之子是天主的能力和智慧，凡由真理而成为智慧的每个人都被祂光照；这些哲学家自称是智慧的爱好者，并因此得名哲学家；因此他们似乎是在辩护智慧，而实际上是它的敌人，因为他们不断推荐有害的迷信，主张应当崇拜和敬畏这世界的元素。\par

7. \BibleQuote{我要观看祢的诸天，祢手指的作为。} \BibleRef{咏 8:3}。我们读到法律是用天主的指头写的，并通过祂的圣仆梅瑟赐下；许多人都将这\Quote{天主的指头}理解为圣神。因此，如果我们正确地将\Quote{天主的指头}理解为这些充满圣神的服务者，因这同一位圣神在他们内工作，因为全部圣经都是藉他们为我们写成；那么，我们在此处将两约的书卷称为\Quote{诸天}就与之一致。关于梅瑟本人，法郎王的术士在被他战胜时也说：\BibleQuote{这是天主的手指。} \BibleRef{出 8:19} 并且经上写道：\BibleQuote{诸天将要被卷起如同书卷。} 虽然这是指这以太天说的，但自然地，按同一形象，书卷之天也以寓意的方式被命名。\BibleQuote{我要观看，} 他说，\BibleQuote{诸天，祢手指的作为：} 也就是说，我将辨别并理解圣经，这些圣经是祢藉圣神的运作，通过祢的服务者写成的。\par

8. 因此，上面提到的\Quote{诸天}也可以解释为这些相同的书卷，当他说：\BibleQuote{因为祢的荣耀已被高举于诸天之上：} 这样，完整的意思应当是：\BibleQuote{因为祢的荣耀已被高举于诸天之上；} 因为祢的荣耀超越了所有圣经的宣告：\BibleQuote{由孩童和吃奶者的口中，祢建立了完美的赞颂，} 使那些将要达到祢荣耀之知识的人应从相信圣经开始；这荣耀被高举于圣经之上，因为它超越并超出了所有言语和语言的宣告。因此，天主甚至将圣经降低到婴孩和吃奶者的能力，正如在另一篇圣咏中所唱的：\BibleQuote{祂使天低垂，亲自降临，} 祂这样做是为了敌人，这些敌人因多言的骄傲，身为基督十字架的敌人，即使有时说一些真理，仍不能有益于婴孩和吃奶者。于是敌人和辩护者被消灭，无论他似乎是辩护智慧，还是甚至辩护基督的名，他仍然从信德的这一步攻击他所轻易许诺的真理。由此他也被证明并不拥有这真理；因为他攻击信德这一步，不知道如何上升到那里。因此，这轻率且盲目的真理许诺者，即敌人和辩护者，被消灭了，当诸天——天主手指的作为——被看见时，也就是说，当圣经甚至被降低到婴孩的迟缓而被理解时；并且藉着历史性信德的卑微（这信德是按时发生的），他们将那些在永恒事物理解的高峰上、在那些他们所建立的事物上被良好养育和坚固的人提升。因为这些诸天，即这些书卷，是天主手指的作为；因为藉着圣神在圣人们内的运作，它们得以完成。因为那些看重自己的荣耀而非人的得救的人，是离开了圣神而发言的，在圣神内有天主仁慈的怀抱。\par

9. \BibleQuote{我要观看祢的诸天，祢手指的作为，月亮和星辰，是祢所布置的。} 月亮和星辰被安置在诸天之中；因为普世教会（月亮常用来象征教会），以及各地方的个别教会（我想是以星辰之名暗示），都建立在这些同样的圣经之中，我们相信圣经是以\Quote{诸天}一词表示的。但为何月亮恰当地象征教会，将在另一篇圣咏中更为适当地加以思考，经上说：\BibleQuote{罪人已弯弓搭箭，要在昏暗的月亮中射击心地正直的人。}\par

10. \BibleQuote{什麼是世人，祢竟眷顧他？或人子，祢竟看顧他？}\BibleRef{詠 8:4}。或許有人會問，世人與人子之間有何區別。倘若沒有區別，就不會這樣分開表達為\Quote{世人，或人子}。因為若寫作\Quote{什麼是世人，祢竟眷顧他，以及人子，祢竟看顧他？}，那似乎是\Quote{世人}一詞的重複。但如今表達為\Quote{世人或人子}時，就更加清楚地暗示了區別。這確實值得記住：每個人子都是世人；但並非每個世人都可被視為人子。例如，亞當是世人，卻不是人子。因此，我們可以從此處思考並區分世人與人子之間的差異：即那些帶有屬地之人的肖像、而非人子的人，應以\Quote{世人}之名指代；而那些帶有屬天之人肖像（\BibleRef{格前 15:49}）的，則應更稱為\Quote{人子}；因為前者又稱為舊人，後者稱為新人；但新人是從舊人而生，因為屬神的重生始於對屬地與世俗生活的改變；因此後者被稱為人子。那麼，此處的\Quote{世人}是屬地的，而\Quote{人子}是屬天的；前者遠離天主，後者與天主同在；因此祂眷顧前者，如同他們在遠處；但祂看顧後者，因他們同在，並以祂的容顏光照他們。因為\Quote{救恩遠離罪人}；並且\Quote{主，祢容顏的光輝已印在我們身上}。因此，在另一篇聖詠中，他提到世人與獸類一同得救，並非藉由內在的直接光照，而是藉由天主慈悲的倍增，祂的良善甚至達於最卑微之物；因為屬血肉之人的健全是屬血肉的，如同獸類一般；但將人子與那些被祂與獸類並列的世人分別出來，他宣告他們藉由真理本身的光照，以及某種生命泉源的湧流，以一種更為崇高的方式獲享真福。因為他如此說：\Quote{主，祢要救護世人與獸類，天主，因祢的慈悲已倍增。但人子將信賴祢翅膀的蔭庇。他們必因祢殿中的豐盈而醉，並使他們飲祢快樂的湍流。因為在祢那裡有生命的泉源，在祢的光明中我們必得見光明。求祢向認識祢的人廣施慈悲。}因著慈悲的倍增，祂眷顧世人，如同眷顧獸類；因為那倍增的慈悲甚至達於遠處的人；但祂看顧人子，使他們置於祂翅膀的蔭庇之下，向他們廣施慈悲，在祂的光明中賜予光明，使他們飲祂的快樂，並以祂殿中的豐盈使他們沉醉，以致忘記先前生活的憂苦與流離。這個人子，即新人，是由舊人的悔改以痛苦和眼淚所生。他雖是新，卻仍被稱為屬血肉的，當他以奶為食時；宗徒說：\Quote{我無法對你們說話如同對屬神的人，只如同對屬血肉的人。}為表明他們已重生，他說：\Quote{如同對在基督內的嬰孩，我以奶餵你們，而非乾糧。}而當他如常發生般地退回到舊生活時，他在責備中聽到他被稱為世人；他說：\Quote{你們豈不是世人嗎？且像世人一樣行事嗎？}\par

11. 因此，人子首先在主的真人身上被看顧，祂由童貞瑪利亞所生。關於祂，因著肉身的軟弱——天主的智慧甘願承擔——以及苦難的屈辱，便正當地說：\BibleQuote{祢使他稍微遜於天使}\BibleRef{詠 8:5}。但隨之而來的是榮耀，祂復活並升天；他說：\BibleQuote{以光榮和尊貴，祢給他加冕；並立他管理祢雙手的化工}\BibleRef{詠 8:6}。既然連天使也是天主雙手的化工，我們理解獨生子甚至被立於天使之上；我們聽見並相信，因著肉身的降生與苦難的屈辱，祂曾稍微遜於天使。\par

12. \BibleQuote{祢将，}他说，\BibleQuote{万有都置于祂脚下。}当他说\BibleQuote{万有，}时，他什么都不排除。为使他不许另有理解，宗徒命令要如此相信，当他说：\BibleQuote{那将万有置于祂下的，不在内}。\BibleRef{格前 15:27} 他也在\WorkTitle{致希伯来人书}中引用了这篇圣咏的这同一见证，为使人明白万有都以如此方式置于我们的主耶稣基督之下，以致无一被排除。\BibleRef{希 2:8} 然而，当他引用\BibleQuote{所有羊群和牛群，是的，此外还有田野的走兽，天空的飞鸟，海里的鱼，以及那些在海中路径上行走的}时，似乎并未加添什么大事，\BibleRef{咏 8:7} 因为，撇开天上的尊威和权能以及所有天使军队，甚至撇开人自己，他似乎只是将走兽置于祂之下；除非我们从羊和牛理解神圣的灵魂，他们或结出纯真的果实，或甚至劳作使大地结果实，即，使属世的人重生获得属灵的富饶。因此，如果我们想由此推知万有都置于我们的主耶稣基督之下，那么这些神圣的灵魂不仅应当理解为人的灵魂，也应当理解为所有天使的灵魂。因为不会有任何受造物不被置于祂之下，祂是那（请允许我如此说）最高级的神体所服从的。但我们从何证明羊甚至可以被解释为不是人，而是高天天使受造物中的真福神体？可否从主的话中得知：祂把九十九只羊留在山中，即留在高处，而下来寻找那一只？因为如果我们把那迷失的一只羊理解为在亚当内的人的灵魂（因为厄娃甚至是从他的肋旁创造的，\BibleRef{创 2:21}），而所有这一切的属灵处理与思考此时无暇顾及，那么留山上的九十九只羊就应指非人的、而是天使的神体。至于牛，这句话很容易处理；因为人本身之所以被称为牛，没有别的原因，乃是因为他们通过宣讲天主圣言的福音而效法天使，正如经上说的：\BibleQuote{打谷的牛不可笼住嘴}。那么，我们岂不更容易将那些真理的使者——天使本身——视为牛，因为福音作者因参与他们的称号而被称为牛？\BibleQuote{祢将置于之下}因此，他说，\BibleQuote{所有羊群和牛群，}即所有圣洁属灵的受造物；其中包括在教会中的圣洁的人，即是在那些酒榨中，这酒榨以月亮和星辰的另一比喻而暗示。\par

13. \Quote{不但如此，}他说，\Quote{田野的走兽。}\Quote{不但}这个词的添加绝非偶然。首先，因为平原的走兽可以理解为绵羊和公牛：因此，如果山羊是岩石和山区的走兽，那么绵羊可以很好地被当作田野的走兽。因此，即使经文这样写：\Quote{所有绵羊、公牛和田野的走兽，}也可以合理地询问平原的走兽指的是什么，因为连绵羊和公牛都可以被理解为这样的走兽。但加上\Quote{不但}一词，迫使我们毫无疑问地认识到某种区别。但是在这个词\Quote{不但}之下，不仅\Quote{田野的走兽，}而且\Quote{天空的飞鸟、海里的鱼，以及行走在海中路径的}\BibleRef{咏 8:8}都被包括在内。那么这种区别是什么？回想一下\Quote{酒榨，}装有皮和酒；打谷场，含有糠秕和麦粒；网，其中网罗了好鱼和坏鱼；诺厄的方舟，其中有洁与不洁的动物：你就会看到教会暂时，在这现今的时期，直到最后的审判时期，不仅包含绵羊和公牛，即圣洁的平信徒和圣洁的圣职人员，而且还有\Quote{田野的走兽、天空的飞鸟，以及行走在海中路径的海鸟。}因为田野的走兽被非常恰当地理解为沉湎于肉体享乐的人，他们不追求任何崇高、任何劳苦的事。因为田野也是\Quote{那引向毁灭的宽路}\BibleRef{玛 7:13}；并且亚伯尔是在田野中被杀的\BibleRef{创 4:8}。因此有理由害怕，免得有人从天主的正义之山上下来（\Quote{因为你的正义，}他说，\Quote{如同天主的山}），选择肉体享乐那宽大而容易的道路，而被魔鬼杀害。现在也看看\Quote{天空的飞鸟，}即骄傲之人，关于他们曾说：\Quote{他们把自己的口设在天上。}看他们如何被风带到高处，\Quote{他们曾说：\InnerQuote{我们必夸大我们的舌头，我们的嘴唇是我们自己的，谁是我们的主？}}也看看海里的鱼，即好奇之人；他们行走在海中的路径，即探究这世上暂时的事物：这些事物如同海中的路径，在给船只或任何行走或游动的东西让开通道之后，很快就随水合拢而消失并灭亡。因为他没有仅仅说\Quote{行走在海的路径，}而是说\Quote{穿越，}显示那些追求虚空和短暂事物之人极为坚定的热忱。现在这三种恶习，即肉体的享乐、骄傲和好奇，包括所有的罪。在我看来，它们似乎被宗徒若望列举出来，当他说：\Quote{不要爱世界；因为凡在世界上的，就是肉体的贪欲、眼目的贪欲，以及生活的骄傲。}\BibleRef{若一 2:15}因为通过眼睛，好奇尤其盛行。其余属于什么确实清楚。而\OriginalTerm{主那人}{Lord Man}的诱惑是三重的：通过食物，即肉体的贪欲，在此被建议：\Quote{命令这些石头变成饼}\BibleRef{玛 4:3}；通过虚妄的夸耀，在此站在山上，世上万国都指给祂看，并许诺如果祂朝拜就赐给祂\BibleRef{玛 4:8}；通过好奇，在此从圣殿顶上被建议跳下去，以试验祂是否会被天使托起\BibleRef{玛 4:6}。因此，既然仇敌不能通过这些诱惑中的任何一个胜了祂，关于他这样说：\Quote{当魔鬼结束了一切诱惑的时候}\BibleRef{路 4:13}。因此，参照酒榨的意义，不仅酒，而且皮也被放在祂的脚下；即不仅绵羊和公牛，即信徒的圣洁灵魂，无论是在平信徒中还是在圣职中；而且还有享乐的走兽、骄傲的飞鸟和好奇的鱼。所有这些类型的罪人，我们现在在教会中看到与善良圣洁的人混杂在一起。愿祂因而在祂的教会中工作，把酒从皮中分离出来：让我们留意，使我们成为酒、绵羊或公牛；而不是皮、田野的走兽、天空的飞鸟或海里的鱼，即那些行走在海中路径的。并不是说这些名称只能以这种方式理解和解释，但对它们的解释必须根据它们所在的地方。因为在别处它们有别的含义。而且这条规则必须在每个寓意中遵守，即用比喻表达的内容应当根据特定地方的意思来考虑：因为这是主和宗徒们的教导方式。那么让我们重复最后一节，它也被放在圣咏的开头，让我们赞美天主说：\Quote{上主，我们的主！祢的名在全地何其美妙！}因为适当地，在论述内容之后，回到标题，所有论述都必须归向那里。\par

\SourceRecord{Translated by an anonymous scholar.；Nicene and Post-Nicene Fathers, First Series；Vol. 8.；Edited by Philip Schaff.；Buffalo, NY: Christian Literature Publishing Co.,；1888.；<http://www.newadvent.org/fathers/1801008.htm>.}

% structure: p0001 source=h1 target=section reason=page-title

\section{圣咏第九篇释义}

1. 这篇圣咏的题词是：\BibleQuote{交与终末，为圣子的隐秘之事，达味的诗。}关于圣子的隐秘之事，可能存有疑问；但由于他没有说明是\Quote{谁的}圣子，便应理解为天主独生子。因为当一篇圣咏题词为\Quote{达味的儿子}，并说\BibleQuote{当他躲避他儿子阿贝沙隆的面时}，尽管连名字都提到了，因此不可能不清楚是指谁说的，但也仅仅说\Quote{从儿子阿贝沙隆的面}，而是加了\Quote{他的}。但此处，由于既没有加\Quote{他的}，又大量提到外邦人，便不能恰当地理解为阿贝沙隆。\EditorialNote{撒慕尔纪下第十五章}因为那个叛离者与他父亲进行的战争，与外邦人毫无关系，因为那里只是以色列人民内部的分裂。所以这篇圣咏是为天主独生子的隐秘之事而咏唱的。因为主自己，在单独使用\Quote{子}这个词时，也愿意人理解为他——独生子；正如祂所说的：\BibleQuote{如果子使你们自由，你们就真自由了。}\BibleRef{若 8:36}因为祂没有说\Quote{天主子}；而只是说\Quote{子}，就使我们明白是\Quote{谁的子}。这种表达方式唯独适用于祂的尊贵，即我们这样谈论祂，虽然不提祂的名字，却可以理解。因为我们说\Quote{下雨了}、\Quote{天晴了}、\Quote{打雷了}，等等，并不添加是谁做的；因为做事者的尊贵自发地出现在所有人的心目中，不需要言语。那么，圣子的隐秘之事是什么？从这个表达我们首先必须理解，圣子有一些明显的事，与那些称为隐秘的事相区别。因此，既然我们相信主的两次来临——一次已过去，犹太人没有理解；另一次将来，我们都盼望——而且那一次犹太人不理解的，使外邦人受益，那么\Quote{为圣子的隐秘之事}并非不合适地理解为指这次来临，在这次来临中\BibleQuote{以色列有几分硬心，直到外邦人的数目添满了}\BibleRef{罗 11:25}。\par

因为有关两种审判的提示贯穿于圣经传给我们，如果有人留意的话，一种隐秘，一种显明。隐秘的审判现在正在发生，对此伯多禄宗徒说：\BibleQuote{时候到了，审判要从天主的家开始}\BibleRef{伯前 4:17}。因此，这隐秘的审判就是痛苦，现在借着痛苦，每人都或被锻炼以致净化，或被警告而悔改，或者若藐视天主的召叫和管教，便为沉沦而被蒙蔽。但显明的审判是那一次，主在来临时将要审判活人死人，所有人都承认正是祂，祂要将赏报分给善人，将惩罚分给恶人。但那时，那承认将无助于解除灾祸，反而增加定罪的积累。关于这两种审判，一种隐秘，一种显明，主似乎已说过，祂说：\BibleQuote{凡信我的，已从死亡进入生命，不会进入审判}\BibleRef{若 5:24}，即不会进入那显明的审判。因为借着某种磨难而从死亡进入生命的审判，就是祂\BibleQuote{鞭打所收纳的每一个儿子}\BibleRef{希 12:6}，这是隐秘的审判。\BibleQuote{凡不信的，已经被定罪了，}\BibleRef{若 3:18}祂说，也就是说，借着这隐秘的审判已经为那显明的审判预备好了。这两种审判我们也读于智慧篇，正如所记载的：\BibleQuote{因此，你对他们如同没有理性的孩子一样，给了一个嘲弄他们的审判；但未被这审判惩戒的人，感到了配得上天主的审判。}\BibleRef{智 12:25}凡未被天主这隐秘审判惩戒的，将极堪当地受那显明审判的惩罚……\par

2. \BibleQuote{我要全心称谢主}\BibleRef{咏 9:1}。全心称谢天主的人，并非在某一细节上怀疑天主的眷顾；而是那已经看到天主的智慧隐秘之事的人，祂的不可见的赏报是何等伟大，祂说\BibleQuote{我们在磨难中也欢跃}\BibleRef{罗 5:3}；以及所有施加于身体的折磨，要么是为锻炼已悔改归向天主的人，要么是为警告他们悔改，要么是为固执者最终受罚的公正准备：因此，现在一切都归于天主上智的安排，愚昧人认为这些事情好像是偶然巧合、随机发生，毫无天主的安排。\BibleQuote{我要述说你的一切奇事。}述说天主一切奇事的人，是看到它们不仅公开地施行在身体上，而且也在灵魂上无形地施行，但方式远为崇高卓越。因为属地的、完全凭肉眼生活的人，更惊奇拉匝禄的身体复活，而非保禄这位迫害者的灵魂复活。但既然可见的奇迹召唤灵魂进入光明，而不可见的奇迹光照那当被召唤时前来的灵魂，那么，通过相信可见的而达到对不可见的理解的人，就是述说天主一切奇事的人。\par

3. \Quote{我將喜樂，在祢內歡躍}\BibleRef{詠 9:2}。不再是在這個世界上，不在肉體的歡愉之中，不在口舌的滋味之中，不在香氣的甜味之中，不在轉瞬即逝的聲音的喜樂之中，不在形形色色的形象之中，不在人讚譽的虛榮之中，不在婚姻與朽壞的後代之中，不在暫時財富的豐裕之中，不在今世的獲得之中，無論它延伸於地域與空間，或延續於時間的相繼；而是：\Quote{我將喜樂，在祢內歡躍}，即是在\Person{聖子}的隱密之處，那裡 \Quote{祢面容的光明已印在我們身上，上主}；因為 \Quote{祢將他們隱藏}，他說，\Quote{在祢面容的隱藏處}。因此，凡講述祢一切奇事的人，將在祢內喜樂歡躍。祂將講述祢的一切奇事（因為如今是先知性地談論祂），\Quote{祂來不是為執行自己的旨意，而是為執行派遣祂者的旨意}。\BibleRef{若 6:38}\par

4. 因為如今主的位格開始出現在這篇聖詠中說話。因為接下來是：\Quote{我將歌頌祢的聖名，至高者，在使我的仇敵轉向後方時}。那麼，祂的仇敵是在何處被轉向後方的？是否是在對牠說 \Quote{退到我後面去，\Person{撒殫}} 的時候？\BibleRef{瑪 16:23} 因為那時，牠藉試探想要置身於前的，卻因未能騙住受試探者，且對牠一無所成，而被轉向後方。因為屬地的人是後方的；但屬天的人被置於前，雖然祂是後來才來的。因為 \Quote{第一個人是出於地，屬於土；第二個人是出於天，屬於天}。\BibleRef{格前 15:47} 但從這一族系來了那一位，祂曾說：\Quote{在我以後來的那一位，在我以前已存在。}\BibleRef{若 1:15} 而宗徒忘記 \Quote{背後的事，奮力向著前面的事}。\BibleRef{斐 3:13} 因此，仇敵被轉向後方，在牠不能騙住受試探的屬天之人以後；牠便轉向屬地的人，在那裡牠能統治……。事實上，魔鬼甚至在迫害義人時也被轉向後方，而牠作為迫害者，對他們更為有利，遠勝於牠作為領袖和君王的在前引導。因此，我們必須在使仇敵轉向後方時，歌頌至高者的聖名：因為我們寧可選擇逃避牠如同迫害者，也不願跟隨牠如同領袖。因為我們有可逃往之處，並隱藏在\Person{聖子}的隱密之處，因為 \Quote{主已成為我們的避難所}。\par

5. \Quote{他們將被削弱，從祢面前滅亡}\BibleRef{詠 9:3}。誰將被削弱並滅亡呢？豈不是不義和不敬虔的人嗎？\Quote{他們將被削弱}，因為他們將一無所能；\Quote{他們將滅亡}，因為不敬虔的人將不存在；\Quote{從祢面前}，即是從天主的認識面前，正如那說 \Quote{如今我活著，不是我，而是基督在我內活著} 的人所消亡的那樣。\BibleRef{迦 2:20} 但為何不敬虔的人 \Quote{將從祢面前被削弱並滅亡}？\Quote{因為}，他說，\Quote{祢已為我伸了冤，辯明了我的案情}：即是我似乎被審判的審判，祢使它成了我的；以及人們定我這義人無辜者有罪的案情，祢使它成了我的。因為這樣的事為我們的得救服務；正如水手們也稱他們的風為他們的風，他們藉助它順風航行。\par

6. \Quote{祢坐在寶座上，按正義審判}\BibleRef{詠 9:4}。無論是\Person{聖子}對\Person{聖父}說這話，祂也曾說 \Quote{若不是從上頭賜給你，你對我就沒有任何權柄}，\BibleRef{若 19:11} 將人類的審判者為人類的利益受審這件事本身，歸於\Person{聖父}的公義和祂自己的隱密之處；或是人對天主說 \Quote{祢坐在寶座上，按正義審判}，以天主的寶座稱呼他的靈魂，這樣他的身體或許是稱為天主 \Quote{腳凳:} 的塵土；\BibleRef{依 66:1} 因為 \Quote{天主在\Person{基督}內，使世界與自己和好:}\BibleRef{格後 5:19} 或是\Person{教會}的靈魂，現在完美無瑕、沒有皺紋，\BibleRef{弗 5:27} 配得上\Person{聖子}的隱密之處，因為 \Quote{\Person{君王}帶她進入了內室}，\BibleRef{歌 1:4} 對她的淨配說 \Quote{祢坐在寶座上，按正義審判}，因為祢從死者中復活，升了天，坐在\Person{聖父}的右邊：無論這節經文可以歸於這些意見中的哪一個，它都沒有逾越信德的規則。\par

7. \BibleQuote{你斥责了外邦人，不虔敬的人已灭亡} \BibleRef{咏 9:5}。我们认为这句话更适合向主耶稣基督说，而非由他说。因为除了祂，还有谁斥责了外邦人，使不虔敬的人灭亡？祂升天后，派遣了圣神，使宗徒们被圣神充满，大胆宣讲天主的圣言，自由地谴责人的罪恶。在这斥责下，不虔敬的人灭亡了；因为不虔敬的人被成义，成了虔敬的人。\BibleQuote{你已涂抹了他们的名，直到永远，且直到世世代代}。不虔敬者的名已被涂抹。因为凡信仰真天主的人不再被称为不虔敬者。如今他们的名被涂抹，\BibleQuote{直到永远}，即直到此世时间历程的终结。\BibleQuote{且直到世世代代}。什么是\BibleQuote{世世代代}，岂非此世界所拥有的形象和阴影？因为季节更替，月亮亏缺复圆，太阳每年回到原位，春、夏、秋、冬消逝又复返，在某种程度上是对永恒的模仿。但这世世代代是居于不变永恒之中的。如同心中的诗句与口中的诗句：前者被理解，后者被听见；前者塑造后者；因此前者在艺术中运作并持久，后者在空中发声并消逝。同样，这个可变世界的样式由那个不变的世界——即所谓的世世代代——所界定。因此，一个居于艺术中，即天主的智慧与能力中；另一个则在受造界的治理中消逝。如果这并非重复，那么既然已说了\BibleQuote{直到永远}，为避免被理解为指这消逝的世界，就加上了\BibleQuote{且直到世世代代}。因为在希腊文抄本中是这样写的：\OriginalTerm{直到永远}{\GreekText{εἰς} \GreekText{τὸν} \GreekText{αἰῶνα}}，\OriginalTerm{且直到世世代代}{\GreekText{καὶ} \GreekText{εἰς} \GreekText{τὸν} \GreekText{αἰῶνα} \GreekText{τοῦ} \GreekText{αἰῶνος}}。拉丁文译者们大多译为，不是\Quote{直到永远，且直到世世代代}，而是\Quote{直到永远，且直到世世代代}，以便用\Quote{直到世世代代}来解释\Quote{直到永远}。因此，\BibleQuote{不虔敬者的名，你已涂抹直到永远}，因为从此以后不虔敬者将不复存在。如果他们的名未延至这个世界，更不用说延至世世代代了。\par

8. \BibleQuote{敌人的剑已在尽头失效} \BibleRef{咏 9:6}。不是复数敌人，而是单数的这个敌人。如今什么敌人的剑失效了，除非是魔鬼的剑？这些剑被理解为各种错误的意见，魔鬼借此如同用剑毁灭灵魂。在战胜这些剑并使之失效的过程中，使用了那剑，在第七篇圣咏中论及此剑说：\BibleQuote{你若不肯回心转意，祂必磨快自己的剑}。也许这就是尽头，敌人的剑在此尽头失效；因为直到尽头，它们尚有几分效力。如今它秘密运作，但在最后审判中它将被公开挥舞。藉此剑，城市被毁灭。因为接下来如此说：\BibleQuote{敌人的剑已在尽头失效：你已毁灭了城市}。这些城市正是魔鬼在其中统治的地方，那里狡诈欺诈的谋略占据了如同宫廷的位置，所有肢体的服务作为官员和臣仆侍奉这至高的统治：眼睛为好奇，耳朵为放荡，或为任何乐于听闻的恶事，手为抢劫或任何暴力与污秽，其他所有肢体以类似方式侍奉暴虐的至高统治，即乖谬的谋略。在这城中，平民百姓仿佛都是心灵中所有软弱的感情和扰乱的情绪，每日在人身内掀起骚动。因此，凡有君王、有宫廷、有官员、有百姓之处，便有一个城市。同样，这些事物若不在个别的人——他们仿佛城市的元素与种子——之中，又怎会在坏城市中存在呢？当那首领——经上论及他说：\BibleQuote{此世界的首领}已被\BibleQuote{逐出} \BibleRef{若 12:31}——被逐出后，这些城市便被毁灭，这些王国被真理之言荒废，恶谋被平息，卑鄙的情感被驯服，肢体和感官的服役被掳获，被转移到正义与善行的服侍中：正如宗徒所说：\BibleQuote{罪恶不能再在我们的必死身体中为王} \BibleRef{罗 6:12} 等等。那时灵魂得享平安，人便预备好接受安息与福乐。\BibleQuote{他们的纪念与喧嚣一同消逝}：与不虔敬者的喧嚣一同消逝。但说\Quote{与喧嚣一同}，或许因为当不虔敬被推翻时，有喧嚣发出：因为无人能抵达那最深寂静的最高之处，除非他先以极大的喧嚣与自己的一切恶习争战；或者\Quote{与喧嚣一同}是说不虔敬者的纪念在喧嚣本身——即不虔敬在其中横行——的消逝中一同消逝。\par

9. \BibleQuote{主永远存留}\BibleRef{咏 9:7}。那麼，\BibleQuote{为什么外邦人妄图背叛，百姓妄想空虚的事，反对主，反对祂的受傅者？}因为\BibleQuote{主永远存留。祂预备了审判的宝座，祂将按公义审判世界。}祂在受审时预备了祂的宝座。因为借着那样的忍耐，人购得了天堂，而天主在人性中使信徒获益。这就是圣子隐藏的审判。但既然祂也将公开地、在众人眼前审判生者死者，祂已在隐藏的审判中预备了祂的宝座；祂也将公开地\BibleQuote{按公义审判世界}，即祂将按功绩分配恩赐，把绵羊放在右边，山羊放在左边。\BibleRef{玛 25:33}\BibleQuote{祂将公正地审判万民}\BibleRef{咏 9:8}。这与上面所说的\BibleQuote{祂将按公义审判世界}是同一回事。不像那些看不见内心的世人审判，常常是更坏的人被宣告无罪而非被定罪；但主将\BibleQuote{按公义}并\BibleQuote{公正地}审判，\BibleQuote{良心作证，思想互相控告或宽免。}\BibleRef{罗 2:15}\par

10. \BibleQuote{主成了穷人的避难所}\BibleRef{咏 9:9}。无论那已被转向背后的敌人如何迫害，对他们这些以主为避难所的人，他能造成什么伤害呢？但这将会是如此的：如果在此世界中——那敌人在其中拥有权势——他们选择成为穷人，不爱任何在此世离人而去或人离世时所留下的东西。因为对于这样的穷人，主已成了避难所，\BibleQuote{在困厄时，作及时的援助。}看，祂使人贫穷，因为\BibleQuote{祂鞭打祂所收纳的每一个儿子。}\BibleRef{希 12:6}因为祂藉着加上\BibleQuote{在困厄时}，解释了什么是\BibleQuote{及时的援助}。因为灵魂除非转离此世，否则不会转向天主；而除非劳苦与痛苦与那轻浮、有害且毁灭性的快乐相混合，它不会更适时地转离此世。\par

11. \BibleQuote{凡认识祢名号的人，请仰望祢}\BibleRef{咏 9:10}，当他们停止仰望财富和此世的其他诱惑时。因为灵魂确实在寻找何处寄托她的希望，当她被从此世撕离时，天主名号的知识适时地接纳了她。因为天主的单纯名号如今已经传遍各地：但名号的认识是指知道这名字属于谁。因为名字并非为了自身而是为了它所意指的。如今经上说：\BibleQuote{主是祂的名号。}\BibleRef{耶 33:2}因此，凡自愿顺服于天主作祂仆人的人，已经知道了这名号。\BibleQuote{凡认识祢名号的人，请仰望祢}\BibleRef{咏 9:10}。再者，主对梅瑟说：\BibleQuote{我是那我是；你要对以色列子民说：\Emph{我是}派遣了我。}\BibleRef{出 3:14}那麼，\BibleQuote{凡认识祢名号的人，请仰望祢}，好使他们不至于仰望那些在时间快速流转中流逝的东西，它们除了\BibleQuote{将要}和\BibleQuote{已经}之外一无所有。因为其中未来之物，一旦来到，立刻成为过去；它被热切地期待，却痛苦地失去。但在天主的本性中，没有什么是尚未存在的，也没有什么是已不再存在的：只有那现在是的，这就是永恒。让他们因此停止仰望和热爱暂时之物，转而致力于仰望永恒吧，他们认识那一位的名号，祂曾说：\BibleQuote{我是那我是}，并且关于祂经上记着：\BibleQuote{\Emph{我是}派遣了我。}\BibleRef{出 3:14}\BibleQuote{因为祢没有抛弃寻求祢的人，上主。}凡寻求祂的人，不再寻求暂时和可朽坏之物；\BibleQuote{因为没有人能事奉两个主人。}\BibleRef{玛 6:24}\par

12. \BibleQuote{请向上主歌唱，祂居住在熙雍}\BibleRef{咏 9:11}，这是对那些主不抛弃寻求祂的人说的。祂居住在熙雍，熙雍被解释为警醒，并且带有现今教会的形像；正如耶路撒冷带有将来教会的形像，即已经享受天使般生命的圣徒之城；因为耶路撒冷被解释为和平的异象。如今警醒先于异象，正如这教会先于那所应许的、不死而永生的城。但在时间上它居先，而非在尊贵上：因为我们要努力达到的目标比我们为达到目标而实践的事更为尊贵；如今我们实践警醒，为能到达异象。但这同一个现今的教会，除非主居住在她内，最热切的警醒可能陷入任何错误。对这教会经上说：\BibleQuote{因为天主的宫殿是圣的，这宫殿就是你们。}\BibleRef{格前 3:17}又说：\BibleQuote{使基督因信德住在你们心里，住在内心的人中。}\BibleRef{弗 3:17}于是我们被命令，向那位居住在熙雍的主歌唱，好使我们同心合意赞美主、教会的居住者。\BibleQuote{在列邦中显扬祂的奇事。}这事已经做了，也将不停地做下去。\par

13. \BibleQuote{因索讨他们的血，祂已记念} \BibleRef{咏 9:12}。如同那些被派去传扬福音的人，对上述的命令作出回应：\Quote{在异民中显扬祂的奇事}，并应说：\BibleQuote{上主，谁相信了我们的报告？} \BibleRef{依 53:1}，又如同说：\BibleQuote{为了祢，我们终日被杀害}；圣咏作者适宜地接着说，基督徒若在迫害中死去，并非没有永生的巨大赏报，\BibleQuote{因索讨他们的血，祂已记念}。但他为何选择说\Quote{他们的血}？仿佛一个知识不全、信德薄弱的人会问：他们如何\Quote{显扬祂的奇事}，既然异民的不信将向他们发怒；而他应得到回答：\BibleQuote{因索讨他们的血，祂已记念}，即最后的审判将要来临，在其中被杀者的荣耀和杀人者的惩罚都将彰显？但不要让任何人以为\BibleQuote{祂已记念}这样的用法是认为天主会遗忘；但由于审判将在漫长的间隔之后，这是按软弱之人的感觉而说的，他们以为天主忘记了，因为祂没有按他们所愿的那样迅速行事。对于这样的人，接下来所说的话也是：\BibleQuote{祂没有忘记穷人的呼号}：即祂并没有像你所认为的那样忘记。仿佛他们听到\BibleQuote{祂已记念}后会说：那么祂曾是忘记了；不，圣咏作者说：\BibleQuote{祂没有忘记}穷人的呼号。\par

14. 但我问，天主所不忘记的穷人的呼号是什么？是这样的呼号吗：\BibleQuote{可怜我，上主，请看我在敌人手下所受的屈辱} \BibleRef{咏 9:13}。那么他为什么不说：可怜\Quote{我们}，上主，请看我们在\Quote{我们的}敌人手下所受的屈辱，仿佛有许多穷人在呼号；而是像一个人说：可怜\Quote{我，} 上主？是否因为有一位为圣人们代祷，\BibleQuote{祂首先为了我们的缘故成了贫穷，祂原是富足的} \BibleRef{格后 8:9} 而正是祂说：\BibleQuote{使我从死亡之门被高举，好使我能熙雍女子之门的各门口宣扬祢的一切赞美} \BibleRef{咏 9:14}？因为人在祂内被高举，不仅是祂所承担的那个人性，那是教会的头；而且我们中间任何一位若属于其他肢体，也要从所有邪恶的欲望中被高举；这些邪欲就是死亡之门，因为经过它们就是通往死亡的道路。但享乐中的快乐本身就是死亡，当一个人获得他因放纵的意志所贪恋的东西时：因为\BibleQuote{贪财是万恶的根源} \BibleRef{弟前 6:10}，因此它是死亡之门，因为\BibleQuote{那好宴乐的寡妇，虽生犹死} \BibleRef{弟前 5:6}。我们通过欲望如同通过死亡之门到达这些快乐。但所有最高的目的都是熙雍女子之门，通过它们我们得以在圣教会中达到和平的默观……或者，是否死亡之门是肉体的感官和眼睛，在人尝了禁果时就打开了，\BibleRef{创 3:7}……而熙雍女子之门是圣事和信德的开端，为叩门者敞开，使他们得以到达圣子的隐密之事？……\par

15. 接着是：\Quote{我要因祢的救恩而欢欣}：即我将因祢的救恩而蒙福地抓住，这救恩就是我们的主耶稣基督，天主的能力和智慧。因此，现在处于苦难中、因望德而得救的教会，只要圣子的审判还是隐藏的，就在希望中说：\Quote{我要因祢的救恩而欢欣}：因为她现在要么被周围的暴力咆哮所压垮，要么被异民的错误所困扰。\BibleQuote{异民陷在自己所做的腐化中} \BibleRef{咏 9:15}。请思考罪人的惩罚是如何从他自己的行为中保留的；以及那些想要迫害教会的人，如何陷在他们原想施加的腐化中。因为他们渴望杀害身体，而他们自己却在灵魂上死亡。\BibleQuote{他们被自己所藏的网罗缠住了脚}。隐藏的网罗是诡诈的计谋。灵魂的脚被很好地理解为它的爱：当爱堕落时，被称为贪欲或情欲；但当正直时，称为爱德……宗徒说：\BibleQuote{为使你们在爱德上根深蒂固，得以领会} \BibleRef{弗 3:17}。因此罪人的脚，即他们的爱，被他们隐藏的网罗所缠住：因为当快乐跟随诡诈的行为，当天主将他们交给自己心里的情欲时，那快乐立刻捆绑他们，使他们不敢将爱从那里移开而应用于有益的事物；因为他们若试图这样做，心里就会痛苦，如同想要将脚从镣铐中解脱出来：在这种痛苦下屈服，他们拒绝退出有害的快乐。因此，\Quote{在网罗}，即\Quote{他们所隐藏的}，也就是说，在诡诈的计谋中，\Quote{他们的脚被缠住了}，即他们的爱，借着诡诈达到了那种购买痛苦的虚假喜乐。\par

16. \BibleQuote{上主彰显审判的执行} \BibleRef{咏 9:16}。这些是天主的审判。并非从祂真福的宁静中，也不是从智慧的秘密处所——那里接受有福的灵魂——带出刀剑、火、野兽或任何此类可使罪人受折磨的东西：但他们是怎样被折磨的，上主如何施行审判？他说：\BibleQuote{罪人被他自己的手所做的缠住了}。\par

17. 这里插入了\Quote{\OriginalTerm{细拉}{diapsalma}之歌}：\BibleRef{咏 9:16}；就如隐藏的喜乐，按我们所能想象的，如今在罪人与义人之间所作的分离，不是关于地点，而是关于心中的爱恋，犹如谷粒与糠秕，尚在场地上。随后又说：\BibleQuote{愿罪人转归地狱}：\BibleRef{咏 9:17}即当他们被宽恕时，让他们被交在自己手中，沉溺于致命的欢乐。\BibleQuote{所有忘记天主的万民。}因为\BibleQuote{他们既然不肯认天主为天主，天主就任凭他们陷于卑鄙的心。}\BibleRef{罗 1:28}\par

18. \BibleQuote{因为穷人终不会被遗忘}：\BibleRef{咏 9:18}穷人如今似乎被遗忘，当罪人被认为在今世的幸福中昌盛，而义人却受苦；但主说：\BibleQuote{穷人的忍耐不会永远消失。}\par

19. \BibleQuote{上主，起来，让人不得强横}：\BibleRef{咏 9:19}这是祈求未来的审判；但在它来到之前，他说：\BibleQuote{愿外邦人在祢面前受审判：}即暗中受审，这被称为在天主眼前，只有少数圣者和义人知道。\BibleQuote{上主，求祢给他们立一个立法者。}\BibleRef{咏 9:20}他似乎向我指明\OriginalTerm{假基督}{Antichrist}；关于他，宗徒说：\BibleQuote{当那罪人显露出来的时候。}\BibleRef{得后 2:3}\BibleQuote{愿外邦人知道他们不过是人。}即那些将因天主之子而得自由、属于人子、成为人之子（即新人们）的人，将服侍人，即那旧人罪人，\BibleQuote{因为他们不过是人。}\par

20. 并且因为相信他将达到如此巨大的虚妄荣耀，并被允许行如此大事，既反对众人又反对天主的圣徒，以致那时有些软弱者会以为天主不关心人间事务，圣咏作者插入一个\OriginalTerm{细拉}{diapsalma}，仿佛加上呻吟的人声，问道为何审判被延迟。\par

\SourceRecord{Translated by an anonymous scholar.；Nicene and Post-Nicene Fathers, First Series；Vol. 8.；Edited by Philip Schaff.；Buffalo, NY: Christian Literature Publishing Co.,；1888.；<http://www.newadvent.org/fathers/1801009.htm>.}

% structure: p0001 source=h1 target=section reason=page-title

\section{圣咏第十篇释义}

\Quote{主啊，祢为什么}他说，\Quote{远远地退避？}\BibleRef{咏 9:1}。然后这样询问的人，仿佛突然明白了，或者仿佛他明知故问以便教导，接着说：\Quote{祢在适当的季节、在患难中轻视：}也就是说，祢适时地轻视，并引起患难来点燃人心对祢来临的渴望。因为生命的泉源对极其干渴的人更为甘甜。因此他暗示了迟延的原因，说：\Quote{当恶人自夸时，穷人被点燃}\BibleRef{咏 9:2}。奇妙而真实的是，小子们因与罪人比较，对正直生活的渴慕被点燃，怀着何等热切的善望。在这奥秘中，甚至异端也被允许存在；并非异端者自己希望这样，而是因为天主眷顾从他们的罪中产生这一结果，祂既制造并安排了光明，却只安排黑暗，好叫相比之下光明更令人愉快，正如与异端者相比，真理的发现更为甘甜。因为通过这种比较，那些被认可的、为天主所认识的人，就在人前显明出来。\par

1. \Quote{他们被自己的思想所擒获，即他们所想的：}也就是说，他们的邪恶思想成了他们的锁链。但它们如何成了锁链？\Quote{罪人受赞美，}他说，\Quote{在他灵魂的欲望中}\BibleRef{咏 9:3}。奉承者的舌头将灵魂束缚在罪中。因为在做那些事时有快感，不仅不怕责备者，反而听到赞许者。\Quote{行不义的人受祝福。}因此\Quote{他们被自己的思想所擒获，即他们所想的。}\par

2. \Quote{罪人激怒了主}\BibleRef{咏 9:4}。没有人应当祝贺那道路亨通的人，在他的罪上无报应者临近，却有赞许者在旁。这是主更大的愤怒。因为罪人激怒了主，以致他该受这些事，即不受纠正的鞭打。\Quote{罪人激怒了主：祂必不按祂愤怒的众多而追究。}祂的愤怒是大的，当祂不追究，仿佛忘记、不记罪，而人借着欺诈和邪恶获得财富和尊荣时，尤其是在那假基督身上，他将被人视为如此有福，以致甚至被认为是神。但天主的愤怒有多大，我们从下文得知。\par

3. \Quote{天主不在他眼前，他的道路时常污秽}\BibleRef{咏 9:5}。那知道什么在灵魂中带来喜乐和欢愉的人，就知道被真理之光遗弃是多么大的灾祸：因为人们认为肉眼的失明是极大的灾祸，因为这光被撤去。那么，那借着罪恶的顺利进展而陷入如此境地的，即天主不在他眼前，他的道路时常污秽（即他的思想和计谋不洁），他所受的惩罚何其大！\Quote{祢的审判从他面前被夺去。}因为自觉有罪的良心，当它似乎没有遭受惩罚时，便相信天主不审判，于是天主的审判从它面前被夺去；而这本身就是极大的定罪。\Quote{他将统治他所有的仇敌。}因为如此传讲，他将胜过所有君王，并独自获得王国；因为按宗徒论到他所宣讲的：\Quote{他坐在天主的殿中，高举自己超过一切受人崇拜的和被称为神的。}\BibleRef{得后 2:4}\par

4. 并且，既然他被交给自己内心的情欲，并预定受极大的刑罚，他将借着邪恶的伎俩达到那虚空而徒有高度的统治；因此下文说：\Quote{因为他在心里说：我必不离开世代直到世代而没有邪恶}\BibleRef{咏 9:6}。也就是说，我的名声和我的名字不会从这一代传到后代，除非我借着邪恶的伎俩获得如此崇高的统治，以致后代不能对它缄默。因为被遗弃、缺乏善技、远离正义之光的头脑，借着恶技为自己谋得如此持久的声名，甚至在后代中被传颂。而那些不能因善而被认识的人，希望人们即使因恶也谈论他们，只要他们的名字传扬远近。我认为这是这里的意思：\Quote{我必不离开世代直到世代而没有邪恶。}还有另一种解释：若一个虚荣而充满错误的头脑认为，除非借着恶技，它不能从有死的一代到达永恒的一代：这确实也曾论及\Person{西满}，那时他以为借着邪恶的伎俩可以获得天堂，并借着魔术从人类的一代到达神性的一代。\BibleRef{宗 8:9}那么，如果那罪恶之人——他将填满所有假先知已开始的一切邪恶和不虔敬，并如此行\Quote{大的神迹，以致倘若可能，连被选者也要被迷惑}\BibleRef{玛 24:24}——也在心里说：\Quote{我必不离开世代直到世代而没有邪恶}，这有什么奇怪呢？\par

5. \BibleQuote{他的口充满咒骂、苦涩和诡诈} \BibleRef{咏 9:7}。因为用如此可憎的手段去寻求天堂，并聚敛这等财富以取得永恒之座，实在是大咒骂。但他的口充满了这种咒骂。因为这欲望不会生效，只会在他的口内消灭他，他竟敢以苦涩和诡诈，即以愤怒和阴险，来许诺自己这些事，借此拉拢群众归向他。\Quote{在他舌下是劳苦和忧愁。} 没有比不义和不虔更劳苦的了：劳苦之后便是忧愁；因为那劳苦不仅没有果实，甚至趋于毁灭。这劳苦和忧愁是指他在心里所说的：\Quote{我从这一代到那一代，没有邪恶就不会动摇。} 因此，\Quote{在他舌头底下}，而不是在他舌头上，因为他会默默策划这些事，对人却说别的话，好显得自己善良、正义，是天主之子。\par

6. \BibleQuote{他与富人一同埋伏} \BibleRef{咏 9:8}。什么是富人呢？不就是那些他将以此世恩赐重重压住的人吗？因此说他与他们一同埋伏，因为他要显示他们虚假的幸福来欺骗人；那些人因着扭曲的意志而渴望成为像他们一样的人，不寻求永恒的好事，就会落入他的网罗。\Quote{好能在暗中杀害无辜者。} 我认为，\Quote{在暗中} 是指难以理解什么是该寻求的、什么是该避免的地方。杀死无辜者，就是使无辜者成为有罪者。\par

7. \BibleQuote{他的眼目窥视穷人} 因为他主要迫害义人，经上论及他们说：\BibleQuote{神贫的人是有福的，因为天国是他们的} \BibleRef{玛 5:3} \BibleRef{咏 9:9}。\Quote{他像狮子在洞穴中，在隐密处埋伏。} 狮子在洞穴中，是指一个既有暴力又有诡诈的人。因为教会最初的迫害是暴力的，那时通过放逐、酷刑、屠杀，强迫基督徒献祭；另一种迫害是诡诈的，现在由各种异端者和假弟兄进行；还有第三种，将由假基督带来，比任何迫害更危险；因为它将同时具有暴力和诡诈。他会在帝国中施行暴力，在奇迹中施行诡诈。\Quote{狮子} 一词指暴力，\Quote{在他的洞穴中} 指诡诈。这些又以顺序变化重复。他说：\Quote{他埋伏，} \Quote{为了捕捉穷人；} 这指诡诈；但接下来的 \Quote{当他吸引穷人时，} 则归于暴力。因为 \Quote{吸引} 意思是，他通过暴力，用尽各种酷刑把他带向自己。\par

8. 同样，后面的两个也是如此。\BibleQuote{他在自己的罗网中使他谦卑} 是诡诈 \BibleRef{咏 9:10}。\Quote{他必衰败和跌倒，当他统治穷人时} 是暴力。因为\Quote{罗网}自然指向\Quote{埋伏}；但统治最明显地传达了恐怖的概念。他说得好：\Quote{他会在他的罗网中使他谦卑。} 因为当他开始行那些奇事时，它们在人们眼中显得越奇妙，当时那些圣人就越受轻视，甚至被看作无有；他以正义和纯洁抵抗他的人，却似乎因他所行的奇事而胜出。但\Quote{他将会衰败和跌倒，当他统治穷人时；} 也就是说，当他向抵抗他的天主的仆人施加任何他愿意的惩罚时。\par

9. 但他将如何衰败和跌倒呢？\BibleQuote{因为他心里说：天主忘记了；祂转面不看，直到末了} \BibleRef{咏 9:11}。这就是衰败和最悲惨的跌倒，当一个人的心智好似在其不义中亨通，以为自己被豁免；其实是被蒙蔽，被保留等待最终的及时惩罚。关于此，圣咏作者现在说：\BibleQuote{上主天主，求祢起来，高举祢的手} \BibleRef{咏 9:12}；也就是说，使祢的大能彰显。他之前曾说：\Quote{上主，求祢起来，不要让人得胜，愿外邦人在祢面前受审判}：即暗中，只有天主看见。当恶人达到了在人看来极大的幸福时，这事就发生；在他们之上立了一个立法者，正是他们应得的，经上论他说：\Quote{上主，求祢在他们之上立一个立法者，愿外邦人知道他们不过是人。} 但在那隐藏的惩罚和复仇之后，才说：\Quote{上主天主，求祢起来，高举祢的手；} 当然不是在暗中，而是在最显明的光荣中。\Quote{使祢不要永远忘记穷人；} 也就是说，正如恶人所想的，他们说：\BibleQuote{天主忘记了，祂转面不看，以致看不见直到末了} \BibleRef{咏 9:11}。他们否认天主看顾到底，他们说祂不关心人类和地上的事，因为大地可说是事物的终结；它是最后的元素，人在其中按最有序的方式劳苦，但他们看不见自己劳苦的秩序，这特别属于圣子的隐密之事。因此，在那样的时期劳苦的教会，如同在巨浪和风暴中的船，唤醒那似在沉睡的主，使祂命令风，恢复平静。\BibleRef{玛 8:24} 因此他说：\Quote{上主天主，求祢起来，高举祢的手，使祢不要永远忘记穷人。}\par

10. 因此，如今他们明白那显明的审判，并为此欢欣鼓舞，说道：\Quote{为何不虔敬的人惹怒了天主？}\BibleRef{咏 9:13}；就是说，他行了如此大的恶事，对他有什么益处呢？\Quote{因为他心里说：祂必不追究。}接着是：\Quote{因为祢细察劳苦和忿怒，要将他们交付在祢手中。}\BibleRef{咏 9:14} 这一句需要清楚解释，否则如有错误，便会变得隐晦。因为不虔敬的人心里说：天主必不追究，好像天主顾念劳苦和忿怒，要将他们交付在祂手中；也就是说，好像祂害怕劳苦和忿怒，因此会宽恕他们，免得惩罚他们成为祂的负担，或者免得祂被怒气的风暴所扰；这正如人通常的做法，以逃避劳苦或怒气为由，为自己不施行报复辩解。\par

11. \Quote{贫穷的人已托付于祢。}因此他贫穷，即轻视这世界一切暂时之福，好使祢独作他的盼望。\Quote{祢将作孤儿的扶助，}即对于那其父——这世界——已死去的人（他按肉身为此世界所生），并且他已然能够说：\Quote{世界已钉在十字架上于我，我也于世界。}\BibleRef{迦 6:14} 因为对于这样的孤儿，天主成为他们的父亲。主实实在在教导我们，祂的门徒会成为孤儿，祂对他们说：\Quote{不要称地上的人为父。}\BibleRef{玛 23:9} 祂自己首先在此事上以身作则，说：\Quote{谁是我的母亲？谁是我的兄弟？}\BibleRef{玛 12:48} 为此，某些极其恶毒的异端者竟敢断言祂没有母亲；他们却没有看到，如果留意这些话，便会得出其门徒也没有父亲的结论。因为正如祂说\Quote{谁是我的母亲？}，同样，祂教训他们说：\Quote{不要称地上的人为父。}\par

12. \Quote{求祢打断罪人和恶人的臂膀}\BibleRef{咏 9:15}；即那上面所论及的\Quote{他必统辖他的一切仇敌。}他把他的能力称为臂膀；与它相对的是基督的能力，关于这能力经上说：\Quote{上主天主，求祢起来，愿祢的手高举。他的过犯必要被追究，他必因之而不复存在；}即他要因自己的罪受审判，并因自己的罪而丧亡。此后，若接着出现\Quote{上主为王，直到永远，世世无穷；异民要从祂的地上灭亡}\BibleRef{咏 9:16}，便不足为奇。他在这里用\OriginalTerm{异民}{heathen}指罪人和不虔敬的人。\par

13. \Quote{上主俯听了穷人的渴望}\BibleRef{咏 9:17}：那是他们热切渴望的，当他们在今世的困苦和磨难中，切望上主的日子。\Quote{祢的耳朵听到了他们心中的准备}。这就是心中的准备，另一篇圣咏曾歌颂道：\Quote{我的心已准备妥当，天主，我的心已准备妥当：}宗徒也说：\Quote{如果我们希望那未看见的，我们就坚忍等待。}\BibleRef{罗 8:25} 按一般的解释规则，我们应将天主的耳朵理解为祂聆听的能力，而非身体器官；同样（以免重复），凡经上提及祂的肢体——那些在我们身上是可见、有形的——都应理解为运作的能力。因为上主天主听的不是声音的响动，而是心中的准备，我们切不可将这一点理解为任何有形之物。\par

14. \Quote{为孤儿和卑微人施行审判}\BibleRef{咏 9:18}：即不是为那同化于这世界的人，也不是为骄傲的人。因为审判孤儿是一回事，为孤儿审判是另一回事。即使定罪他的，也是在审判孤儿；但为孤儿审判的，是作出有利于他的判决。\Quote{免得人在地上更加自高自大。}因为他们就是那些论及\Quote{上主，求祢在他们之上立一个立法者：愿异民知道他们不过是人。}的人。但那在此同一段落中被理解为置于他们之上的人，也将是个人，现在经上论他说：\Quote{免得人在地上更加自高自大：}即当人子来为孤儿施行审判时——那孤儿已脱去旧人，并且因此好像埋葬了他的父亲。\par

15. 论及圣子的隐秘事——本篇圣咏中已多言及——此后将显明圣子的彰显事，于同篇圣咏结尾已略述一二。但标题取自前者，因前者在本篇中占较大篇幅。确实，主来临的那一天，可合理地列于圣子的隐秘事中，尽管主本身的临在将是显明的。因为论及那日，经上说：连人子也不知道，连天使、掌权者也不知道，人子也不知道。\BibleRef{谷 13:32} 那么，还有什么事比这更隐秘呢？这竟连审判者本人也隐秘——并非就知识而言，而是就启示而言？但论及圣子的隐秘事，即便有人不愿将其理解为天主子，而愿理解为达味本人——全部\WorkTitle{圣咏集}都归其名下，因为吾人知道圣咏皆称为达味圣咏——也请倾听主所说的话：\Quote{达味之子，可怜我们吧！} \BibleRef{玛 20:30} 如此，即便以此方式，他也能理解同一主基督，本篇圣咏的标题正是关于祂的隐秘事。因为天使也同样说：\Quote{天主要把祂祖先达味的御座赐给祂。} \BibleRef{路 1:32} 对此理解，以下句子并不抵触：同一主问犹太人：\Quote{如果基督是达味之子，祂怎么在圣神里称祂为主，说：\InnerQuote{主对我主说：你坐在我右边，等我使你的仇敌作你的脚凳。}} \BibleRef{玛 22:43} 这是对未受过教导的人说的，他们虽然期待基督来临，却只期待祂作为人，而非作为天主的德能与智慧。于是，在那处，祂教导最真实、最纯洁的信德：祂既是达味王的主，因为祂是起初的圣言，与天主同在的天主，\BibleRef{若 1:1} 万物由祂造成；又是子，因为按肉身，祂是由达味的后裔所生。因为祂并没有说基督不是达味之子；但既然你们已经认定祂是达味之子，就应学习祂如何是达味之主：不要只从基督的人子方面来看——按此祂是达味之子；\BibleRef{罗 1:3} 而忽略了祂是天主子——按此祂是达味之主。\par

\SourceRecord{Translated by an anonymous scholar.；Nicene and Post-Nicene Fathers, First Series；Vol. 8.；Edited by Philip Schaff.；Buffalo, NY: Christian Literature Publishing Co.,；1888.；<http://www.newadvent.org/fathers/1801010.htm>.}

% structure: p0001 source=h1 target=section reason=page-title

\section{对圣咏第十一篇的讲疏}

\Emph{致候词，达味自己的圣咏。}\par

1. 这个标题不需要重新考虑；因为\Quote{致候词}的意义已经充分处理过了。那么就让我们来看这篇圣咏的本文，在我看来，它是反对异端者而唱的；他们列举并夸大教会中许多人的罪过，仿佛他们自己中的全部或大多数是义人，竭力要使我们离开唯一真正的慈母教会的怀抱，转向并抢夺我们：断言基督与他们同在，并警告我们，似乎带着虔诚和恳切，说通过投向他们，我们就能投靠基督，而他们却虚假地宣称拥有基督。我们知道，在预言中，基督在众多以比喻传达关于祂的称呼的名字中，也被称为一座山。所以我们必须回答这些人，并说：\BibleQuote{我信赖主：你们怎么对我的灵魂说：要像麻雀一样，迁到山里去？} \BibleRef{咏 10:1}。我持守一座我所信赖的山，你们怎么说我应该投靠你们，好像有许多基督一样？或者，若你们因骄傲而说你们是山，我确实需要一只翅膀装备了天主的权能和诫命的麻雀：但正是这些事情阻碍我飞向这些山，并将我的信赖寄托在骄傲的人身上。我有一座房屋可以安歇，因为我信赖主。因为甚至\Quote{麻雀找到了一个家}，并且\Quote{主成了穷人的避难所。}那么让我们满怀信心地说，免得我们在异端者中寻找基督而失去祂：\Quote{我信赖主：你们怎么对我的灵魂说：要像麻雀一样，迁到山里去？}\par

2. \BibleQuote{因为，看，罪人已经张弓，他们在箭囊中预备了箭，要在昏暗的月亮下射穿心地正直的人。} \BibleRef{咏 10:2}这些是那些以罪人来威胁我们的人的恐惧，好使我们投向他们，作为义人。\Quote{看，}他们说，\Quote{罪人已经张弓：} 我想是圣经，他们按肉性的解释，从中射出恶毒的语句。\Quote{他们在箭囊中预备了箭：} 即同样的言语，他们将依圣经的权威射出，已预备在心里的隐秘处。\Quote{要在昏暗的月亮下射穿心地正直的人：} 即当他们看到，由于众多无知和属肉性的人，教会的光明被遮蔽，他们不能被定罪时，\BibleQuote{他们以不良的交往败坏良好的品行。}\BibleRef{格前 15:33}但面对所有这些恐惧，我们必须说：\Quote{我信赖主。}\par

3. 现在我记得，我承诺要在这篇圣咏中考虑月亮以怎样的适当性象征教会。关于月亮有两种可能的意见：但其中哪一个是真实的，我想人不可能或很难断定。因为当我们问月亮的光从何而来时，有些人说它是自己的，但它的球体一半光明，一半黑暗；当它在自己的轨道上旋转时，其中明亮的部分逐渐转向地球，以便我们能看到它；因此起初它的外观好像有角……按照这种意见，月亮在比喻中象征教会，因为教会在其属灵的部分是光明的，但在其属肉性的部分是黑暗的：有时属灵的部分因善行而被看见，但有时它隐藏在良心中，只有天主知道，因为只有在身体上它才被人看见……但按照另一种意见，月亮也被理解为象征教会，因为她没有自己的光，而是被天主的独生子照亮，在圣经的许多地方，祂被比喻地称为太阳。某些异端者不认识祂，也不能看见祂，就竭力把单纯人的思想转向这个有形可见的太阳，它是人及苍蝇肉身的共同光明，他们确实败坏了一些人，这些人既然不能用心灵看见真理的内在光明，就不会满足于纯朴的公教信仰；这信仰是婴孩的唯一安全，并且仅藉着奶，他们才能在确定的力气中达到更坚实食物的稳固支撑。那么，这两种意见中无论哪一种是真的，月亮在比喻中都被合适地理解为教会。或者，在这样一些令人烦恼而非有教育性的困难中，如果没有满意或没有闲暇来锻炼心智，或者心智本身不能胜任，则用普通的眼睛看月亮就足够了，不必寻求晦涩的原因，而是和所有人一起察觉她的增减和圆缺；她之所以亏缺，是为了重新圆满，即使对粗野的大众，她也展示教会的形象，在其中相信死人的复活。\par

4. 其次，我们必须探究本圣咏中\Quote{昏暗的月亮}意指什么，罪人正是在这月亮中预备射那心里正直的人？因为月亮被称为昏暗，不止一种方式：当她的月周期结束，当她的光亮被云遮断，当她满月时被遮蔽，月亮都可称为昏暗。那么，首先可以理解为殉道者的迫害者，因为他们想在那昏暗的月亮中射那心里正直的人；这或许是在教会青春时期，因为她尚未在地上大放光明，克服外邦迷信的黑暗；或借亵渎者及诽谤基督徒之名的舌头，当大地仿佛被阴云笼罩时，月亮，即教会，不能清晰可见；或当殉道者本人遭受屠杀、如此大量流血，如同月蚀和昏暗——此时月亮似乎显出血色的面容——软弱者被吓退，不敢信奉基督徒之名；在这恐惧中，罪人射出狡猾而亵渎的话语，甚至要迷惑心里正直的人。其次，可以理解为教会所包含的这些罪人，因为那时他们趁这月亮昏暗的机会犯下许多罪行，如今这些罪行被异端者讥诮地归咎于我们，然而据说他们的创始人曾犯下这些罪行。但不论那在昏暗的月亮中所行的是什么，既然如今天主教之名已传遍并称颂于整个世界，我何必为不知之事而烦扰？因为\BibleQuote{我倚靠主}；我也不听那些对我的灵魂说\BibleQuote{你当像麻雀一样逃往山中。因为，看，罪人已弯弓，要在昏暗的月亮中射那心里正直的人}的人。或者，若月亮如今在他们看来仍是昏暗的，因为他们要使天主教会变得不确定，并试图凭她所包含的许多属血肉之人的罪来定她的罪；那对诚然说\BibleQuote{我倚靠主}的人有何相干？借此话语，每个人都表明他自己是麦子，并以忍耐忍受糠秕，直到扬场之时。\par

5. \Quote{在主内}，因此，\Quote{我倚靠}。让人害怕吧，那些倚靠人的人，他们无法否认自己属于某一派，凭其白发人起誓；当在交谈中被问及属于共融体时，除非他们说自己是那一派的，否则不被承认。……或者你或许会说，经上记载：\BibleQuote{你们可凭他们的果实辨别他们}？\BibleRef{玛 7:16} 我确实看到奇妙的事迹：那些\OriginalTerm{巡行派}{Circumcelliones}的日常暴力，以主教和司铎为首领，四处飞窜，称他们可怕的棍棒为\Quote{以色列}；这些事现今活着的人日日看见并感受。但至于马加略时代，他们对此发出不平之鸣，大多数人并未见过，如今也没有人看见；任何看见过的天主教徒，若愿作天主的仆人，都可以说：\BibleQuote{我倚靠主。}……\par

6. 那么，让天主教灵魂说：\BibleQuote{我倚靠主；你们怎敢对我的灵魂说，你当像麻雀一样逃往山中？因为，看，罪人已弯弓，他们在箭囊中预备了箭，要在昏暗的月亮中射那心里正直的人。}并从他们转向主，说：\BibleQuote{因为他们已毁坏了你所完善的。}\BibleRef{咏 10:3} 她这样说，不仅针对这些人，也针对所有异端者。因为他们都尽其所能毁坏了天主从婴孩和吃奶者口中建立的赞美，当他们以虚妄而吹毛求疵的问题扰乱幼小者，不容他们以信德的乳哺养。那么，仿佛对这灵魂说，他们为什么对你说\BibleQuote{你当像麻雀一样逃往山中}？为什么他们以那些\BibleQuote{弯弓要在昏暗的月亮中射心里正直的人}的罪人来恐吓你？她回答：他们恐吓我，正是\BibleQuote{因为他们已毁坏了你所完善的。}在哪里毁坏的？只在他们的\OriginalTerm{集会所}{conventicles}里，在那里他们不是以乳哺养，而是以毒药杀害那些内在光明尚幼小无知的人。\Quote{但那位义者做了什么？}如果马加略、塞西利安得罪了你，基督做了什么得罪你？他曾说：\BibleQuote{我将平安赐给你们，我将我的平安留给你们}；\BibleRef{若 14:27} 这平安却被你们可憎的分裂所破坏。基督做了什么得罪你？祂以极大的忍耐容忍了出卖祂的人，甚至将亲手祝圣、亲口祝福的第一份圣体也给了他，如同给其他宗徒一样。基督做了什么得罪你？祂派遣同一个出卖者——祂称他为魔鬼，\BibleRef{若 6:70} 这人在出卖主之前甚至对主的钱囊也不忠诚\BibleRef{若 12:6}——与其他门徒一起宣讲天国；\BibleRef{玛 10:5} 为要表明天主的恩赐临到那些以信德领受的人，尽管他们通过之领受的人像犹大那样。\par

7. \BibleQuote{主在祂的圣殿中}\BibleRef{咏 10:4}，诚如宗徒所说：\BibleQuote{天主的圣殿是圣的，这圣殿就是你们。}\BibleRef{格前 3:17}\BibleQuote{若有人毁坏天主的圣殿，天主必要毁坏那人。}他毁坏天主的圣殿，就是破坏合一；因为他\BibleQuote{不持守元首，本于祂，全身都靠祂联络得合式，百节各按各职，照着各体的功用彼此相助，便叫身体渐渐长大，在爱中建立自己。}主就在祂这圣殿中，这圣殿由祂的许多肢体组成，各尽其职，藉着爱建造成一座建筑物。那为自己的优越而脱离天主教会团体的人，就是毁坏这圣殿。\BibleQuote{主在祂的圣殿中；主，祂的宝座在天上。}你若把义人理解为天，像把罪人理解为地，经上对罪人说：\BibleQuote{你是土，仍要归于土；}\BibleRef{创 3:19}那么，当说到\BibleQuote{主，祂的宝座在天上}时，你就会明白这是重复\BibleQuote{主在祂的圣殿中}这句话。\par

8. \BibleQuote{祂的眼目垂顾穷人。}属于那位为穷人留下了自己的主，祂成了穷人的避难所。因此，在这网中（\BibleRef{玛 13:47}）直到被拉上岸的一切骚动和纷扰——异端者对此责备我们，自取灭亡，却使我们受到纠正——都是由那些不愿作基督穷人的人所引起的。但他们会转离天主垂顾那些愿意作穷人的人的眼目吗？\BibleQuote{因为祂的眼目垂顾穷人。}难道我们担心，在富人的拥挤中，祂可能看不见那少数穷人——祂在天主教会怀抱中安全地养育他们？\BibleQuote{祂的眼睑审问人子。}在此，我愿按这规则将\BibleQuote{人子}理解为那些藉信德从旧人重生的人。因为圣经中某些晦涩的段落，如同天主闭合的眼睛，使他们操练而寻求；而另一些清晰的段落，如同天主睁开的眼睛，使他们得光照而喜乐。这种在圣经中频繁的闭合与张开，就像是天主的眼睑；它们审问，即考验\BibleQuote{人子}：他们既不因事物的晦暗而厌倦，反而得到操练；也不因知识而自高，反而得到坚固。\par

9. \BibleQuote{主审问义人和恶人}\BibleRef{咏 10:5}。那么，我们为何害怕恶人倘若以不真诚的心与我们共享圣事会伤害我们呢？因为祂\BibleQuote{审问义人和恶人}。\BibleQuote{但喜爱不义的人恨恶自己的灵魂：}那喜爱不义的人伤害的并不是那信靠天主、不寄望于人的人，而只是他自己的灵魂。\par

10. \BibleQuote{祂要向罪人降下网罗}\BibleRef{咏 10:6}。若将云彩理解为一般的先知，无论好坏，他们也被称为假先知：主天主如此安排假先知，藉着他们向罪人降下网罗。\BibleRef{玛 24:24}因为除非是罪人，没有人会跟从他们，无论是为了预备最后的惩罚（若他选择坚持犯罪），还是为了劝阻骄傲（若他及时以更真诚的意图寻求天主）。但若将云彩理解为只是良善真确的先知，那么显然天主也藉着这些先知向罪人降下网罗，虽然藉着他们，祂也灌溉虔诚的人，使之结出果实。宗徒说：\BibleQuote{对于某些人，我们是出于死的香气叫人死；对于另一些人，我们是出于活的香气叫人活。}\BibleRef{格后 2:16}因为不仅先知，凡以天主之言浇灌灵魂的，都可称为云彩。当人们错误理解他们时，天主向罪人降下网罗；但当人们正确理解时，祂使虔诚和信者的心结出果实。例如，经文\BibleQuote{二人成为一体}\BibleRef{弗 5:31}，若有人以情欲的眼光解释它，祂便向罪人降下网罗；但你若如那说\BibleQuote{但我是指着基督和教会说的}\BibleRef{弗 5:32}的人那样理解，祂便向沃土降下甘霖。这两者都是由同一片云彩，即圣经，产生的。主又说：\BibleQuote{不是入口的污秽人，而是出口的。}\BibleRef{玛 15:11}罪人听见这话，便为贪食预备口舌；义人听见，便对食物上迷信的分别保持警惕。因此，同样从圣经的云彩中，照各人的功过，对罪人降下网罗之雨，对义人降下丰盛之雨。\par

11. \Quote{火与硫磺和暴风是他们杯中的份。}这就是那些亵渎天主圣名之人的惩罚和结局：首先，他们被自身情欲之火消耗净尽；然后，因恶行的恶臭，被逐出真福者的团体；最后，被卷走并淹没，遭受不可言喻的刑罚。因为他们杯中的份正是如此，正如义人的份是：\Quote{你的杯爵是多么令人沉醉！因为他们必因你殿中的丰盛而沉醉。}我以为，此处提及杯爵，是为了让我们不认为天主的眷顾在惩罚罪人时也有任何超出适度与分寸之处。因此，仿佛要说明为何如此，他又加上：\Quote{因为主是公义的，他爱慕义德。}\BibleRef{咏 10:7}。复数形式并非毫无意义，而是因为此处谈论的是人，所以应理解为\Quote{义德}指义人。因为在许多义人身上，似乎可以说存在着多重的义德，而天主的义德却只有一个，众人都分沾这同一义德。正如一张脸面对许多镜子，镜中映出唯一的脸，却因众多镜子而显得多样。因此，他回到单数，说：\Quote{他的面容看见了正直。}或许，\Quote{他的面容看见了正直}，意思是说，在他的面容中看见了正直，即在认识他之中看见了正直。因为天主的容面就是那使相称的人得以认识他的德能。或者至少，\Quote{他的面容看见了正直}，因为他不容恶人认识他，只容善人认识他；这就是正直。\par

12. 但若有人愿意将这月亮理解为会堂，便可将这篇圣咏联系到主的苦难，并论及犹太人说：\Quote{因为他们毁坏了你所成全的；}而论及主自己说：\Quote{但那位义者作了什么？}他们控告他毁坏了法律。而他们自己却因败坏的生活、藐视法律并建立自己的规条，毁坏了法律的诫命。因此，主自己可以按人的方式讲话，正如他惯常所说的：\Quote{我依靠主；你们怎敢对我的灵魂说：你要像麻雀一样飞到山上去？}这是因为那些想要逮捕并钉死他的人所带来的恐惧。将罪人愿意\Quote{射那心里正直的人}（即相信基督的人）解释为\Quote{在这昏暗的月亮中}，并非不合理。于此相应的还有：\Quote{主在他的圣殿里；主的宝座在天上}，即圣言在人内，或说人子自己就在天上。\BibleRef{若 3:13}他的眼目垂顾穷人；或指他作为天主所选取的人性，或指他作为人所为之受苦的人性。\Quote{他的眼睑察问人子。}眼睛的闭合与张开，很可能就是用\Quote{眼睑}这个词所表达的意思，我们可以将此理解为他的死亡与复活，藉此他考验了人子们——他的门徒们，他们在他的苦难中惊惶，又因复活而喜乐。主管察义人和恶人，即使如今他也从天上治理教会。\Quote{但喜爱不义的人，恨恶自己的灵魂。}为何如此，下文教导了我们。因为\Quote{他要在罪人身上降下网罗：}这应按照上述的解释来理解，圣咏其余部分直至结尾也是如此。\par

\SourceRecord{Translated by an anonymous scholar.；Nicene and Post-Nicene Fathers, First Series；Vol. 8.；Edited by Philip Schaff.；Buffalo, NY: Christian Literature Publishing Co.,；1888.；<http://www.newadvent.org/fathers/1801011.htm>.}

% structure: p0001 source=h1 target=section reason=page-title

\section{\WorkTitle{圣咏第十二篇释义}}

交与乐官，为第八日，达味圣咏。\par

1. 在第六篇中已说过，\Quote{第八}可视为审判之日。\Quote{为第八}亦可以指\Quote{为永恒的时代}；因为在当前的时间（即七天循环）之后，它将被赐予圣人们。\par

2. \BibleQuote{上主，求你拯救我，因为虔诚人已断绝}；即，找不到：正如我们说话时所说的，粮食断绝了，或钱财断绝了。\BibleQuote{真理在人子中已稀少} \BibleRef{咏 11:1}。真理是唯一的，藉以光照圣洁的灵魂；但因灵魂众多，可以说在其中有许多真理：如同镜子中从同一面孔反射出许多影像。\par

3. \BibleQuote{世人彼此谈论虚妄} \BibleRef{咏 11:2}。所谓\Quote{邻人}必须理解为每一个人：因为没有人可以与之行恶；\BibleQuote{爱近人却不作恶} \BibleRef{罗 13:10}。\BibleQuote{诡诈的嘴唇，用一颗心与一颗心说出恶言。}重复\Quote{用一颗心与一颗心}表示双心。\par

4. \BibleQuote{愿上主剪除一切诡诈的嘴唇} \BibleRef{咏 11:3}。他说\Quote{一切}，以使无人以为自己例外：正如宗徒所说：\BibleQuote{对一切行恶的人，先是犹太人，后是希腊人} \BibleRef{罗 2:9}。\BibleQuote{说大话的舌头}：骄傲的舌头。\par

5. \BibleQuote{他们曾说：我们要张大我们的舌头，我们的嘴唇属于我们，谁为我们主呢？} \BibleRef{咏 11:4}。这里指的是骄傲的伪君子，他们自信靠言论欺骗人，并不服从天主。\par

6. \BibleQuote{因贫穷者的困苦和穷乏者的叹息，现在我必兴起，上主说} \BibleRef{咏 11:5}。因为主自己在福音中怜悯祂的子民，因为他们没有领袖，但他们本可善从。因此福音中也说：\BibleQuote{庄稼多，工人少} \BibleRef{玛 9:37}。但此处必须看作是天主圣父以第一人称说话，祂因困苦者和穷乏者（即那些在属灵美善上困苦和穷乏的人）的缘故，屈尊派遣了祂的独生子。由此开始玛窦福音中祂在山上的讲道，祂说：\BibleQuote{神贫的人是有福的，因为天国是他们的} \BibleRef{玛 5:3}。\Quote{我要安置在救恩中}。祂没有说安置什么；但\Quote{在救恩中}必须理解为在基督内，正如：\BibleQuote{因为我亲眼看见了你的救恩} \BibleRef{路 2:30}。由此可见，祂将解除贫穷者的困苦和安慰穷乏者的叹息所需的一切都安置在祂内。\Quote{我要在祂内放心行事}：正如福音中：\BibleQuote{因为他教训他们，像有权威的人，不像他们的经师} \BibleRef{玛 7:29}。\par

7. \BibleQuote{上主的话}是\BibleQuote{纯净的话} \BibleRef{咏 11:6}。这是先知本人的口气：\Quote{上主的话}是\Quote{纯净的话}。他说\Quote{纯净}，没有虚伪的掺杂。因为许多人宣讲真理并不纯洁（参见\BibleRef{斐 1:16}），他们为换取今生利益的贿赂而出卖真理。宗徒如此说：他们不纯洁地宣讲基督。\BibleQuote{经过火炼的银子，为大地而炼}：这些主的话藉着磨难被证明对罪人有益。\BibleQuote{七次炼净}：即因敬畏天主、虔诚、知识、能力、谋略、聪慧、智慧（\BibleRef{依 11:2}）。也有七种真福，主在玛窦记述的同一篇山上讲道中逐一列出：\BibleQuote{神贫的人是有福的，温良的人是有福的，哀恸的人是有福的，饥渴慕义的人是有福的，怜悯人的人是有福的，心里洁净的人是有福的，缔造和平的人是有福的} \BibleRef{玛 5:3}。这七句话中，可以观察那整篇长篇讲道是如何讲述的。因为第八句：\BibleQuote{为义而受迫害的人是有福的} \BibleRef{玛 5:10}，指的是那炼净银子的火本身，这银子被七次炼净。在这篇讲道结束时说：\BibleQuote{因为他教训他们，像有权威的人，不像他们的经师} \BibleRef{玛 7:29}。这正与本圣咏中所说的\Quote{我要在祂内放心行事}相呼应。\par

8. \BibleQuote{上主，你必护佑我们，保守我们脱离这世代，直到永远} \BibleRef{咏 11:7}：在此世如困苦穷乏者，在那世如丰裕富足者。\par

9. \BibleQuote{恶人却在周围徘徊} \BibleRef{咏 11:8}：即在可朽坏之物的欲望中徘徊，如同车轮在七日的周而复始中循环；因此他们不能到达第八日，即永恒，而这正是本圣咏标题所指的。所罗门也说：\Quote{智慧的君王簸扬恶人，并转动恶人的轮子。——你的高位使人的子孙倍增}。因为在可朽坏之物中也有一种倍增，使人远离天主的合一。因此\BibleQuote{可朽坏的身体使灵魂沉重，这泥土的帐幕使多思虑的心压抑} \BibleRef{智 9:15}。但义人的倍增是\Quote{因天主的高位}，当他们\Quote{因力量而进至力量}时。\par

\SourceRecord{Translated by an anonymous scholar.；Nicene and Post-Nicene Fathers, First Series；Vol. 8.；Edited by Philip Schaff.；Buffalo, NY: Christian Literature Publishing Co.,；1888.；<http://www.newadvent.org/fathers/1801012.htm>.}

% structure: p0001 source=h1 target=section reason=page-title

\section{《圣咏第十三篇释义》}

\Quote{交与乐官，达味圣咏。}\par

1. \Quote{因为基督是法律的终向，使所有信者获得成义。}\BibleRef{罗 10:4}\Quote{上主，祢忘记我要到何时？} \BibleRef{咏 12:1} 也就是说，祢将我从属灵上理解基督——祂是天主的智慧，是灵魂一切追求的真正终向——推迟到何时？\Quote{祢转面不顾我，要到何时？} 正如天主不忘记，祂也不转面：但圣经是按我们的方式说话。当天主不将祂自己的知识赐给那尚未拥有纯洁心灵之眼的灵魂时，就说祂转面不顾。\par

2. \Quote{我心中筹算，要到何时？}\BibleRef{咏 12:2} 只有在逆境中才需要筹算。因此，\Quote{我心中筹算，要到何时？} 就像是在说：我要处在逆境中到何时？或者至少这是一个回答，意思是：上主，祢忘记我直到终局，并转面不顾我，要等到我在自己灵魂中筹算为止：除非人在自己灵魂中筹算，以完全施行仁慈，否则天主不会引导他到终局，也不会赐予他那圆满的自我认识，即\Quote{面对面}的认知。\Quote{我心中终日悲伤？} 要理解\Quote{我要有……到何时}。而\Quote{终日}表示持续，即\Quote{日}被当作时间：每个人都渴望从这时间中解脱，因此他心中悲伤，恳求升至永恒的事物，不再忍受人的时日。\par

3. \Quote{我的仇敌高居我上，要到何时？} 或是魔鬼，或是肉欲的习性。\par

4. \Quote{上主，我的天主，求祢俯视我，应允我。}\BibleRef{咏 12:3} \Quote{俯视我}指的是所说的\Quote{祢转面不顾我，要到何时？}\Quote{应允我}指的是所说的\Quote{祢忘记我要到何时？}\Quote{求祢光照我的眼睛，免得我沉睡于死亡。} 必须理解心灵的眼睛，免得它们被罪孽的愉悦蒙蔽而闭合。\par

5. \Quote{免得我的仇敌说：我战胜了他。}\BibleRef{咏 12:4} 魔鬼的嘲笑是可怕的。\Quote{如果我被摇动，那迫害我的人必会欢欣。} 魔鬼和他的天使；他们虽迫害那义人约伯，却没有在他身上欢欣，因为他没有被摇动，即没有从信德的坚定中退缩。\BibleRef{约 2:3}\par

6. \Quote{但我依靠祢的仁慈。}\BibleRef{咏 12:5} 因为人没有被摇动、坚定地持守于主内这件事本身，他不应归功于自己：免得当他夸耀自己没有被摇动时，反因这骄傲而被摇动。\Quote{我的心因祢的救恩而欢欣。} 在基督内，在天主的智慧内。\Quote{我要歌颂上主，因祂厚赐于我。} 属灵的恩赐，不属于人的时日。\Quote{我要歌颂至高上主的名。}\BibleRef{咏 12:6} 也就是说，我以喜乐感恩，并极其适当地运用我的身体，这是属灵灵魂的歌咏。但如果在这里要区分什么的话，\Quote{我要歌颂}是用心，\Quote{我要歌咏}是用我的行动；\Quote{歌颂上主}是指祂独自看见的，而\Quote{歌颂上主的名}是指人中所知的，这并非为祂有益，而是为我们有益。\par

\SourceRecord{Translated by an anonymous scholar.；Nicene and Post-Nicene Fathers, First Series；Vol. 8.；Edited by Philip Schaff.；Buffalo, NY: Christian Literature Publishing Co.,；1888.；<http://www.newadvent.org/fathers/1801013.htm>.}

% structure: p0001 source=h1 target=section reason=page-title

\section{\WorkTitle{圣咏第十四篇释义}}

致候终末，达味自己的一篇圣咏。\par

1. 所谓\Quote{致候终末}是什么意思，不必过于重复地解释。\Quote{因为基督是法律的终末，使凡信祂的人都获得成义}，\BibleRef{罗 10:4}如宗徒所说。我们信从祂，是在开始踏上善道之时；我们将看见祂，是在抵达终末之时。因此，祂就是终末。\par

2. \Quote{愚妄人心里说：没有天主} \BibleRef{咏 13:1}。因为某些亵渎神明、可憎恶的哲学家，他们虽对天主存有乖谬虚伪的观念，却也不敢说\Quote{没有天主}。因此，是心里说\Quote{没有天主}；因为没有人敢说出口，即使他胆敢如此想。\Quote{他们都败坏了，因他们的情感而变得可憎恶：}也就是说，他们爱世界而不爱天主；这些是腐蚀灵魂的情感，使之昏昧，以至于愚妄人甚至心里说\Quote{没有天主}。\Quote{因为他们既不愿意认识天主，天主就任凭他们陷于可鄙的心境。} \BibleRef{罗 1:28} \Quote{没有行善的，连一个也没有。} \Quote{连一个}可以理解为\Quote{连同那一个}，即指没有一个人；或者\Quote{除了那一位}，即主基督除外。正如我们说：这块田地直到海边；我们当然不把海连同田地算在内。这后一种解释更好，即没有人行善直到基督为止；因为除非祂赐予，没有人能行善。这是真实的，因为人若不认识独一天主，就不能行善。\par

3. \Quote{主从天上垂顾世人，要看看有谁领悟，有谁寻求天主} \BibleRef{咏 13:2}。这可以解释为论及犹太人；因为他们崇拜独一天主，与外邦人相比，他们被给予更尊贵的\Quote{世人}之称；前面的话\Quote{愚妄人心里说：没有天主}等大概是指外邦人。主垂顾，是藉着祂神圣的灵魂来看：这就是所谓\Quote{从天上}。因为祂自身而言，没有什么是向祂隐藏的。\par

4. \Quote{他们都离弃正道，一同变为无用：}即犹太人与前面所说的外邦人一样。\Quote{没有行善的，连一个也没有} \BibleRef{咏 13:3}应按上文解释。\Quote{他们的咽喉是敞开的坟墓。}这或是象征永不满足的食欲；或是寓指那些杀害者，以及如同吞吃被杀害者的人，他们将自身行为的败坏灌输给他人。与此相反的意义是曾对伯多禄说的：\Quote{宰杀，吃吧！} \BibleRef{宗 10:13}即他应使外邦人归信于自己的信仰和良好的行为。\Quote{他们用舌头行诡诈。}阿谀奉承是贪婪者和一切恶人的同伴。\Quote{他们的嘴唇下有蛇的毒液。}以\Quote{毒液}指诡诈；\Quote{蛇}则是因为他们不愿听从法律的训诫，如同蛇\Quote{不愿听魔术者的声音}；这在另一篇圣咏中说得更清楚。\Quote{他们的口充满咒骂和苦毒：}这就是\Quote{蛇的毒液}。\Quote{他们的脚飞快地流人血。}他这里表明行恶的习惯。\Quote{毁灭与不幸}在他们的\Quote{道路上}。因为恶人的一切道路都充满劳苦和痛苦。因此主呼喊道：\Quote{凡劳苦和负重担的，你们都到我这里来，我要使你们安息。你们背起我的轭，向我学习，因为我是良善心谦的。因为我的轭是柔和的，我的担子是轻省的。} \BibleRef{玛 11:28} \Quote{和平的道路他们不认识：}即那道路，主所说的，在柔和的轭和轻省的担子中。\Quote{天主的敬畏不在他们眼前。}这些人不说\Quote{没有天主}，但他们并不敬畏天主。\par

5. \Quote{难道所有作恶的人都不明白吗？} \BibleRef{咏 13:4}他警告审判。\Quote{他们吞食我的百姓如同吃饭：}即每日如此。因为饭是每日的食物。那些吞食百姓的人，是从百姓身上谋求私利，不将他们的服务归于天主的荣耀，以及所辖之人的得救。\par

6. \Quote{他们不呼求主。}因为凡渴求祂所不喜悦之事的人，并非真心呼求祂。\Quote{他们在那里恐惧战栗，其实没有可恐惧的} \BibleRef{咏 13:5}：即害怕失去暂时之物。因为他们说：\Quote{如果让他这样下去，众人都会信从他；罗马人必来，夺去我们的地方和民族。} \BibleRef{若 11:48}他们害怕失去地上的国，其实无需害怕；而他们失去了天国，这本是他们应当害怕的。对一切暂时的福乐也应如此理解；人若害怕失去它们，就无法得到永恒的事物。\par

7. \Quote{因为天主在正义的世代中。}这关联到前面的话，意思是：\Quote{难道所有作恶的人不明白，主在正义的世代中吗？}即祂不在那些爱世界的人里面。因为不义在于离弃万有的创造者，\Quote{事奉受造物超过造物主。} \BibleRef{罗 1:25} \Quote{你们羞辱了穷人的计谋，因为主是他的盼望} \BibleRef{咏 13:6}：即你们轻视了圣子的谦卑来临，因为你们没有在祂身上看到世界的浮华；祂所召叫的人应当只寄望于天主，而非那些消逝之物。\par

8. \BibleQuote{谁将从熙雍赐予以色列救恩？} \BibleRef{咏 13:7}。除了那位你们所鄙视的卑微者，还有谁呢？这是不言而喻的。因为祂要在光荣中来审判活人死人，并建立义人的国度。为的是，正如在那卑微的来临时，\BibleQuote{部分以色列人硬了心，直到外邦人的数目满盈} \BibleRef{罗 11:25}，然后接着应验的是，\BibleQuote{这样全以色列都要得救。}因为宗徒也引用了依撒意亚的预言，那里说：\BibleQuote{拯救者要从熙雍出来，消除雅各伯的不敬虔。} \BibleRef{依 59:20}对于犹太人，正如这里所说：\BibleQuote{谁将从熙雍赐予以色列救恩？}\BibleQuote{当上主将祂子民被掳者领回时，雅各伯要欢腾，以色列要喜乐。}这是一种重复，如同常有的情形：因为我认为，\Quote{以色列要喜乐}等同于\Quote{雅各伯要欢腾。}\par

\SourceRecord{Translated by an anonymous scholar.；Nicene and Post-Nicene Fathers, First Series；Vol. 8.；Edited by Philip Schaff.；Buffalo, NY: Christian Literature Publishing Co.,；1888.；<http://www.newadvent.org/fathers/1801014.htm>.}

% structure: p0001 source=h1 target=section reason=page-title

\section{对圣咏第十五篇的释义}

达味自己的圣咏。\par

1. 关于这个标题，没有疑问。\BibleQuote{主，谁能住在祢的帐幕里？}\BibleRef{咏 14:1}。虽然\Quote{帐幕}有时也用来指永恒的居所，但当\Quote{帐幕}按本义理解时，它是战争之物。因此，士兵被称为\Quote{同帐者}，因为他们一起支搭帐幕。这种含义由\BibleQuote{谁将寄居？}这句话得到加强。因为我们暂时与魔鬼作战，那时我们需要一个帐幕，在其中我们可以得到休憩。这特别指出了暂时性救恩计划的信仰，这计划是通过主的降生成人为我们在时间内完成的。\BibleQuote{谁能住在祢的圣山上？}这里他或许立即指永恒的居所本身，\BibleRef{格后 5:1}，我们应当理解\Quote{山}意指在永生中基督之爱的至高无上。\par

2. \BibleQuote{行为无瑕，实行正义的人}\BibleRef{咏 14:2}。这里他提出了命题；接下来他详细阐述。\par

3. \BibleQuote{他心中说实话的人}。因为有些人口中有真理，心中却没有。好比有人欺骗性地指路，明知那里有强盗，却说：如果你走这条路，就不会遇到强盗；结果发现那里实际上没有强盗：他讲了真话，但不是出于真心。因为他以为情况相反，在无知中说了真话。因此，光说真话是不够的，除非心中也是如此。\BibleQuote{他舌头上没有行欺诈的人}\BibleRef{咏 14:3}。当人口里说一套，心里藏另一套时，就是用舌头行欺诈。\BibleQuote{没有加害邻人}。众所周知，\Quote{邻人}应理解为每一个人。\BibleQuote{也没有接受对邻人的毁谤}，即不轻易或轻率地相信控诉者。\par

4. \BibleQuote{恶人在他眼中已被消灭}\BibleRef{咏 14:4}。这就是成全，即恶人对人毫无力量；并且这\Quote{在他眼中}，意思是他确切知道，恶之所以存在，只是因为心灵背离了自己造物主永恒不变的形式，转向从无中被造的受造物的形式。\BibleQuote{但他荣耀那些敬畏主的人：}就是主自己。现在，\BibleQuote{敬畏主是智慧的开端。}因此，正如上述之事属于成全者，那么他接下来要说的属于初学者。\par

5. \BibleQuote{他向邻人起誓，并不欺骗他。}\BibleQuote{他没有放债取利，也没有受贿反对无辜者}\BibleRef{咏 14:5}。这些不是大事：但连这些都不能做到的人，更不用说在心中说实话、舌头上没有行欺诈，而是如何心里有真理，口里也承认并持有，\BibleQuote{是就说是，非就说非；}\BibleRef{玛 5:37}，并且不加害邻人，即不加害任何人；也不接受对邻人的毁谤：所有这些都属于成全者的德性，在他们眼中恶人已被消灭。然而，他连这些较小的事也这样总结：\BibleQuote{行这些事的人，将永远不会动摇。}也就是说，他将达到那些更大的事，其中有着伟大而不可动摇的稳定。因为甚至时态的变化也不是没有原因的，以致于在上述结论中用了过去时，而这里用了将来时。因为那里曾说：\Quote{恶人在他眼中已被消灭}，而这里说：\Quote{将永远不会动摇。}\par

\SourceRecord{Translated by an anonymous scholar.；Nicene and Post-Nicene Fathers, First Series；Vol. 8.；Edited by Philip Schaff.；Buffalo, NY: Christian Literature Publishing Co.,；1888.；<http://www.newadvent.org/fathers/1801015.htm>.}

% structure: p0001 source=h1 target=section reason=page-title

\section{圣咏第十六篇讲疏}

\Emph{题记，达味所作}\par

1. 我们的君王在这篇圣咏中以祂所摄取的人性发言，祂的君王称号在祂受难之时被显明地展示出来。\par

现在祂这样说：\BibleQuote{上主，求祢保全我，因为我全心依靠祢}\BibleRef{咏 16:1}；\BibleQuote{我曾对上主说：祢是我的天主，因为祢不需要我的财物}\BibleRef{咏 16:2}；因为祢并不因我的财物而蒙福。\par

3. \BibleQuote{在祂地上的圣人}\BibleRef{咏 16:3}：即那些寄望于永生之地的圣人，天上耶路撒冷的居民，他们的属灵交往，借着希望的锚，定锚在那片国土上，这国土被正确地称为天主的土地；虽然目前他们也在这地上生活在肉身内。\Quote{祂奇妙地在我身上成就了我的一切愿望。}那么，对这些圣人，祂奇妙地在我身上成就了我的一切愿望，因他们的长进，他们得以明白，我的神性的人性如何为了使我死亡而使他们获益，以及人性的神性如何为了使我复活而使他们获益。\par

4. \BibleQuote{他们的软弱倍增了}\BibleRef{咏 16:4}：他们的软弱倍增了，不是为了毁灭他们，而是为了使他们渴望医生。\Quote{后来他们急忙。}因此，在软弱倍增之后，他们急忙，好得医治。\Quote{我必不因血而聚集他们的集会}。因为他们的集会不再是属血肉的，我也不会像因牲畜的血得安抚而聚集他们。\BibleRef{依 1:11} \Quote{我的嘴唇也必不再提他们的名字。}但借着属灵的改变，他们过去是什么将被遗忘；我也不再称他们为罪人、敌人或人，而是因我的和平称他们为义人、我的弟兄和天主的儿子。\par

5. \BibleQuote{上主是我产业的份和我的杯爵}\BibleRef{咏 16:5}。因为他们必与我一同承受产业，就是上主自己。让别人为自己选择地上的、暂时的产业去享用吧；圣人的产业是永恒的上主。让别人饮用致命的快乐吧，我的杯爵的份是上主。我说\Quote{我的}时，也包含了教会：因为头在哪里，身体也在哪里。因为我将聚集他们的集会进入产业；借着杯爵的沉醉，我将忘记他们的旧名。\BibleQuote{祢必将把我的产业归还给我：}好使这些我所释放的人也能认识\BibleQuote{我在世界受造之前与祢所有的光荣}\BibleRef{若 17:5}。因为祢不会将我从未失去的还给我，而是将这些已经失去这光荣知识的人归还给他们：因我在这光荣中，祢将归还给我。\par

6. \BibleQuote{准绳落在佳美之处}\BibleRef{咏 16:6}。我产业的界线落在祢的荣耀中，如同抽签，正如天主是司铎和肋未人的产业\BibleRef{户 18:20}。\Quote{因为我的产业对我乃是佳美的。}\Quote{我的产业是佳美的}，不是对所有人，而是对看见的人；因我在这产业中，\Quote{对我乃是佳美的。}\par

7. \BibleQuote{我要颂赞上主，祂给了我明悟}\BibleRef{咏 16:7}：使这产业能被看见和拥有。\Quote{而且，甚至到了夜间，我的肾脏也管束了我。}是的，除了明悟以外，甚至到死，我较低的部分，即所摄取的肉身，也教训了我，使我经历死亡的黑暗，这黑暗是那明悟所没有的。\par

8. \BibleQuote{我常将上主置于我眼前}\BibleRef{咏 16:8}。但进入逝去的事物时，我不从祂面前挪开我的眼目，祂永远常存；预见这事，即我经过暂时的事物后必要回到祂那里。\BibleQuote{因为祂在我右边，使我不致动摇。}因为祂恩待我，使我坚定地住在祂内。\par

9. \BibleQuote{因此，我的心欢喜，我的舌快乐}\BibleRef{咏 16:9}。因此，我的思念中有喜乐，我的言语中有欢乐。\BibleQuote{而且，我的肉身也必安息于希望中。}而且，我的肉身不会因毁灭而朽坏，而是安眠于复活的希望中。\par

10. \BibleQuote{因为祢不会将我的灵魂遗弃在地狱}\BibleRef{咏 16:10}。因为祢不会将我的灵魂交予下面那些地方作为产业。\BibleQuote{祢也不会让祢的圣者见到腐朽。}祢也不会允许那圣化的身体（其他的人也将借此身体而被圣化）见到腐朽。\BibleQuote{祢已将生命的道路指示给我了}\BibleRef{咏 16:11}。祢已藉着我指出了谦卑的道路，好使人从他们因骄傲而堕落之处回归生命；因我在这道路中，祢已指示给我了。\Quote{祢要以祢的容仪充满我喜乐。}祢要充满他们喜乐，使他们别无所求，当他们\Quote{面对面}看见祢时；因我在这喜乐中，祢要充满我。\Quote{快乐在祢的右边，直到永远。}快乐在祢的恩宠和怜悯中，在此生的旅途中，引领他们直到祢容仪荣耀的终点。\par

\SourceRecord{Translated by an anonymous scholar.；Nicene and Post-Nicene Fathers, First Series；Vol. 8.；Edited by Philip Schaff.；Buffalo, NY: Christian Literature Publishing Co.,；1888.；<http://www.newadvent.org/fathers/1801016.htm>.}

% structure: p0001 source=h1 target=section reason=page-title

\section{《圣咏第十七篇讲疏》}

达味自己的祈祷。\par

1. 此祈祷当归于主位格，并加上作为祂身体的教会。\par

2. \BibleQuote{天主，请听我的正义，垂顾我的恳求}\BibleRef{咏 16:1}。\BibleQuote{请倾听我的祈祷，不出于欺诈的口唇：}不出于欺诈的口唇而走向祢。\BibleQuote{愿我的判断从祢的面容发出}\BibleRef{咏 16:2}。因祢光照的认识，让我判断真理。或至少，愿我的判断发出，不从欺诈的口唇，而是从祢的面容，即我在判断时不说话语以外我所领悟于祢的。\BibleQuote{愿我的眼睛看见正直：}当然，是心灵的眼睛。\par

3. \BibleQuote{祢在夜间考验并眷顾了我的心}\BibleRef{咏 16:3}。因为我的心曾因考验的苦难而被证实。\BibleQuote{祢以火锻炼我，在我身上找不到不义。}现在不仅是黑夜（因它常扰乱人），而且火（因它燃烧）也称为苦难；藉此我受考验时被找到是正义的。\par

4. \BibleQuote{使我的口不谈论人的行为}\BibleRef{咏 16:4}。使我的口不出任何言语，除非关乎祢的荣耀和赞美；而不谈论人的行为——那些违背祢旨意的事。\BibleQuote{因祢嘴唇的言语。}因祢和平或祢先知的言语。\BibleQuote{我持守了艰难的道路。}我持守了人类必死和受苦的劳苦道路。\par

5. \BibleQuote{使我的脚步在祢的路上得以完全}\BibleRef{咏 16:5}。为使教会的爱在狭窄的道路上得以完全，她由此抵达祢的安息。\BibleQuote{使我的脚步不致动摇。}使我的道路的标志——如同脚步，已印在圣事和宗徒著作上——不致动摇，让那些愿跟随我的人能标记它们。或至少，在我完成艰难道路、在祢路径的狭窄中走完我的脚步后，我能永恒不动地安住。\par

6. \BibleQuote{我呼求了，因为祢应允了我，天主}\BibleRef{咏 16:6}。我以自由而有力的努力向祢倾注我的祈祷：因我能有此能力，祢在我更软弱祈祷时已应允了我。\BibleQuote{求祢侧耳听我，垂听我的言语。}愿祢的听闻不抛弃我的卑微。\par

7. \BibleQuote{使祢的仁慈奇妙}\BibleRef{咏 16:7}。愿祢的仁慈不被轻看，免得它们被爱得太少。\par

8. \BibleQuote{祢拯救那投靠祢的人脱离反抗祢右手的人：}脱离反抗祢恩待我的恩宠之人。\BibleQuote{求祢保护我如眼中的瞳仁}\BibleRef{咏 16:8}：它看似非常微小，但藉它，眼睛的视线得以定向，从而区分光明与黑暗；正如藉基督的人性，审判区分义人与罪人的神性得以辨别。\BibleQuote{在祢翅膀的荫庇下保护我。}在祢慈爱和怜悯的庇护下保护我。\BibleQuote{脱离那些困扰我的恶人的面}\BibleRef{咏 16:9}。\par

9. \BibleQuote{我的仇敌围绕了我的灵魂；}\BibleQuote{他们封闭了自己的脂肪}\BibleRef{咏 16:10}。\BibleQuote{他们的口说了骄傲的话。}因此他们的口说了骄傲的话，说：\BibleQuote{犹太人的君王，万岁，}\BibleRef{玛 27:29}以及其他类似的话。\par

10. \BibleQuote{他们抛弃我，如今围绕了我}\BibleRef{咏 16:11}。他们把我抛弃在城外，如今在十字架上围绕了我。\BibleQuote{他们决意把眼睛转向地上。}他们决意把内心的倾向转向这些地上之事：认为被杀者遭受了极大的不幸，而他们这些杀他的人却安然无恙。\par

11. \BibleQuote{他们像预备捕食的狮子擒拿了我}\BibleRef{咏 16:12}。他们擒拿了我，正如那仇敌\BibleQuote{四处巡游，寻找可吞噬的人。}\BibleRef{伯前 5:8}\BibleQuote{又如一只住在隐密处的幼狮。}又如它的幼崽，就是那曾被告知\BibleQuote{你们是出于你们的父亲魔鬼}\BibleRef{若 8:44}的百姓，他们谋划着陷阱，用以算计并毁灭那义者。\par

12. \BibleQuote{上主，求祢起来，迎击他们，推翻他们}\BibleRef{咏 16:13}。上主，求祢起来——祢被他们以为在沉睡，不关心人的不义——愿他们先被自己的恶意弄瞎，以致惩罚先于他们的行为；然后推翻他们。\par

13. \BibleQuote{求祢从恶人手中拯救我的灵魂。}求祢拯救我的灵魂，在我死后恢复我——那恶人加给我的死亡。祢的武器：脱离祢手的仇敌\BibleRef{咏 16:14}。因为我的灵魂是祢的武器，祢的手，即祢永恒的大能，曾拿起它用以征服不义的国度，并将义人与恶人分开。那么，求祢将这武器从祢手的仇敌——即祢大能的仇敌，也就是我的仇敌——中解救出来。\BibleQuote{上主，求祢将他们从地上消灭，在他们今生中分散他们。}上主，求祢将他们从他们所居住的地上消灭，在他们仅视为生命的今生中分散他们于普世——那些绝望于永生的人。\BibleQuote{他们因祢隐藏的事物而肚腹饱足。}如今，不仅这可见的惩罚会临到他们，而且他们的记忆也被罪恶充满——这些罪恶如同黑暗，隐藏于祢真理之光之外，以致他们忘记天主。\BibleQuote{他们饱尝了猪肉。}他们被不洁充满，践踏了天主圣言的珍珠。\BibleQuote{并把剩余留给他们的孩子：}他们喊道：\BibleQuote{这罪归于我们和我们的孩子。}\BibleRef{玛 27:25}\par

14. \BibleQuote{但我将因祢的正义在祢面前显现}\BibleRef{咏 16:15}。但我——并未向那些以污秽暗淡之心不能见智慧之光的人显现——\BibleQuote{我将因祢的正义在祢面前显现。}\par

\Quote{当祢的光荣显现时，我将获得饱饫。}当他们因自己的不洁而饱饫，以致不能认识我时，当祢的光荣显现时，我将在认识我的人中获得饱饫。确实，在那句经文说\Quote{充满了猪的肉}的地方，有些抄本作\Quote{充满了孩子}：因为希腊文的歧义导致了两种解释。现在，藉着\Quote{孩子}我们理解为行为；正如好的孩子代表好的行为，同样恶的孩子代表恶的行为。\par

\SourceRecord{Translated by an anonymous scholar.；Nicene and Post-Nicene Fathers, First Series；Vol. 8.；Edited by Philip Schaff.；Buffalo, NY: Christian Literature Publishing Co.,；1888.；<http://www.newadvent.org/fathers/1801017.htm>.}

% structure: p0001 source=h1 target=section reason=page-title

\section{\WorkTitle{圣咏第十八篇释义}}

致终局，为上主的仆人，\Person{达味}自己。\par

1. 就是说，为那只强而有力的手，即在其人性内的基督。\Quote{这首歌的言语，就是他在上主救他脱离一切仇敌和撒乌耳之手的日子，向上主所说的；他说：在上主救他脱离一切仇敌和撒乌耳之手的日子：}即犹太人的君王，他们为自己所要求的。\BibleRef{撒上 8:5}因为正如\Person{达味}按解释被称为\Quote{强而有力的手}，同样\Person{撒乌耳}被称为\Quote{要求}。众所周知，那人民为自己要求了一位君王，并接受他作他们的君王，不是按照天主的意愿，而是按照他们自己的意愿。\par

2. 因此，基督和教会，即整个基督，头和身体，在这里说：\BibleQuote{上主，我的力量，我要爱慕你。}\BibleRef{咏 17:1}上主，我因你而坚强，我要爱慕你。\par

3. \BibleQuote{上主，我的磐石，我的城堡，我的救主}\BibleRef{咏 17:2}。上主，你扶持了我，因为我投靠了你；我投靠了你，因为你救了我。\BibleQuote{我的天主是我的助佑，我要寄望于他。}我的天主，你先以你的召叫赐我帮助，使我能寄望于你。\BibleQuote{我的护卫，我救恩的角，我的救赎者。}我的护卫，因为我未曾倚靠自己，在你面前抬起骄傲的角；但我找到了你作为真正的角，即确实的救恩之巅；而且为能找到它，你救赎了我。\par

4. \BibleQuote{我要以赞美称呼上主，我必安全脱离我的仇敌}\BibleRef{咏 17:3}。我不寻求自己的荣耀，只寻求上主的荣耀；我要呼求他，这样，不虔敬的错误就无法伤害我。\par

5. \BibleQuote{死亡的痛苦}即刻于肉身的痛苦，\BibleQuote{环绕了我。不虔敬的汹涌将我困扰}\BibleRef{咏 17:4}。一时兴起的困扰，如同即将消退的暴雨，来搅扰我。\par

6. \BibleQuote{地狱的痛苦环绕了我}\BibleRef{咏 17:5}。在那些环绕我、要毁灭我的人中，有嫉妒的痛苦，它产生死亡，并引向罪的地狱。\BibleQuote{死亡的罗网迎面而来}\BibleRef{咏 17:5}。他们抢先，想先伤害我，这以后将报应在他们身上。他们以正义的夸耀恶毒地劝服了某些人，就抓住这些人使之丧亡；这些人所夸耀的正义只在名义上，而非在实际上，他们以此攻击外邦人。\par

7. \BibleQuote{我在困苦中呼求上主，向我的天主哀呼。他从他的圣殿中垂听了我的声音}\BibleRef{咏 17:6}。他从我的心——他居住之处——垂听了我的声音。\BibleQuote{我的哀号在他面前传入他的耳中；}\BibleRef{咏 17:6}我的哀号，我不是在人耳前发出，而是内在地在你自己面前，\Quote{传入他的耳中。}\par

8. \BibleQuote{大地震动战栗}\BibleRef{咏 17:7}。当人子这样受荣耀时，罪人们震动战栗。\BibleQuote{山岳的基础动摇。}\BibleRef{咏 17:7}骄傲者在此世所存的希望动摇了。\BibleQuote{它们摇动，因为天主向他们发怒。}\BibleRef{咏 17:7}就是说，使现世财物的希望不再在人心中有任何坚固的根基。\par

9. \BibleQuote{在他震怒中冒起浓烟}\BibleRef{咏 17:8}。当悔罪者得知天主对不虔敬者的威胁时，他们泪流满面的恳求升上去了。\BibleQuote{火焰从他面前燃烧}\BibleRef{咏 17:8}。在悔改之后，因认识他而燃起爱德的热火。\BibleQuote{从他中点燃了火炭}\BibleRef{咏 17:8}。那些已经死了，被善愿之火和正义之光所遗弃，停留在寒冷和黑暗中的人，重新点燃并得光照，又复活了。\par

10. \BibleQuote{他弯下诸天，亲自降临}\BibleRef{咏 17:9}。他贬抑了那义者，以便降临到人的软弱中。\BibleQuote{黑暗在他脚下}\BibleRef{咏 17:9}。不虔敬的人，他们品味世俗之事，在他们自身恶毒的黑暗里不认识他；因为他脚下的地好似他的脚凳。\par

11. \BibleQuote{他乘坐革鲁宾，飞行}\BibleRef{咏 17:10}。他被高举超越知识的圆满，使无人能到他那里去，除非借着爱；因为\BibleQuote{爱是法律的满全}\BibleRef{罗 13:10}。他很快向爱他的人显示他是不可捉摸的，免得他们以为他是借着肉体的想象而被领悟的。\BibleQuote{他驾御风的翅膀飞行}\BibleRef{咏 17:10}。但这显示他不可捉摸的迅速，是超越灵魂的能力的；灵魂凭借这些能力如同翅膀，从属地的恐惧中提升到自由的空间。\par

12. \BibleQuote{他以黑暗作他的掩藏处}\BibleRef{咏 17:11}。他设立了圣事的隐晦，以及在信徒心中隐藏的希望，他在那里隐藏而不舍弃他们。在这黑暗里，\BibleQuote{我们如今仍借着信德行走，而非借着目睹}\BibleRef{格后 5:7}，只要我们\BibleQuote{希望那未见的，并以坚忍等待}\BibleRef{罗 8:25}。\BibleQuote{他的帐幕在他周围}\BibleRef{咏 17:11}。然而，信他的人转向他，围绕他；因为他处在他们中间，同样地是所有朋友的朋友，他如今如同在帐幕中居住在他们内。\BibleQuote{暗中的水云间}\BibleRef{咏 17:11}。为此，若有人理解圣经，不要以为他已经在那光明中——这光明将是我们走出信德进入目睹时所有的；因为在先知们和所有天主圣言的宣讲者中，都有着隐晦的教导。\par

13. \Quote{在祂光辉的面前}\BibleRef{咏 17:12}：这与祂显现之光辉相比。\Quote{祂的云彩已過。}祂圣言的宣讲者如今不再局限于犹太境地，而是已過到了外邦人那里。\Quote{冰雹与火炭。}责打被描绘为冰雹，藉此坚硬的心被击碎；但若一片肥沃温润的土壤，即虔诚的心灵，接纳它们，冰雹的坚硬便融化为水，即那满含闪电般可畏、宛如冻结的责打，融化为令人满足的教义；被爱火点燃的心灵便复苏。这一切都在祂的云彩中過到了外邦人那里。\par

14. \Quote{主从天上发出雷声}\BibleRef{咏 17:13}。因着福音的笃信，主从义人心中发出响声。\Quote{至高者发出祂的声音；}好使我们接纳它，并在人类事务的深处，听见天上的事。\par

15. \Quote{祂射出箭来，使他们四散}\BibleRef{咏 17:14}。祂派遣了福音传播者，他们沿着笔直的道路，藉着力量之翼疾行，并非凭自己的权能，而是凭那派遣他们的祂。而\Quote{祂使他们四散，}即那些被派遣所至的人，好叫他们中有些人成为\Quote{活的香气叫人活，}另一些人成为\Quote{死的香气叫人死。}\BibleRef{格后 2:16}\Quote{祂多发电光，扰乱他们。}祂多行奇迹，扰乱他们。\par

16. \Quote{水泉就出现了。}那涌流到永生的水泉\BibleRef{若 4:14}，在宣讲者身上造成，就出现了。\Quote{大地的根基也显露出}\BibleRef{咏 17:15}。那些不被理解的先知，以及信主之人的世界要建造在其上的先知，就显露出来了。\Quote{主啊，因祢的斥责：}呼喊说：\Quote{天主的国临近你们了。}\BibleRef{路 10:9}\Quote{因祢鼻孔所出的怒气；}说：\Quote{你们若不悔改，都要如此灭亡。}\BibleRef{路 13:5}\par

17. 祂从高天伸手，把我拉上来\BibleRef{咏 17:16}：藉着从外邦人中召叫作为产业，\Quote{荣耀的教会，毫无玷污、皱纹等类的事。}\BibleRef{弗 5:27}\Quote{祂把我从大水之中拉上来。}祂把我从众多民族中拉上来。\par

18. \Quote{祂救我脱离我的強敵}\BibleRef{咏 17:17}。祂救我脱离我的仇敌，他们曾得胜，压迫并倾覆我这暂时生命。\Quote{并脱离恨我的人，因为他们比我强盛；}当我还伏在他们之下，不认识天主时。\par

19. \Quote{他们在我的苦难之日抢先攻击我}\BibleRef{咏 17:18}。他们先加害于我，在我承受着必死而劳苦的身体之时。\Quote{但主成了我的支柱。}由于那属世快乐的支柱因苦难的苦涩而被摇撼撕裂，主就成了我的支柱。\par

20. \Quote{祂领我到宽阔之处}\BibleRef{咏 17:19}。因我正忍受肉身的狭窄，祂领我进入了信德属灵的宽广。\Quote{祂救拔了我，因祂喜爱我。}在我喜爱祂之前，祂已救我脱离我最强盛的仇敌（他们在我一度喜爱祂时嫉妒我），并脱离恨我的人，因为我确实喜爱祂。\par

21. \Quote{主必按我的公义报答我}\BibleRef{咏 17:20}。主必按我善意的公义报答我，祂先施怜悯，在我有善意之前。\Quote{并按我手中的清洁偿还我。}祂必按我行为的清洁偿还我，祂藉着领我到信德的宽广之处，赐我能力行善。\par

22. \Quote{因为我遵守了主的道路}\BibleRef{咏 17:21}。好使那因信德而行的善工之宽广，以及坚忍的恒久，得以跟随。\par

23. \Quote{我也没有不敬虔地偏离我的天主。}\Quote{因为祂的一切典章在我面前}\BibleRef{咏 17:22}。\Quote{因为}我以恒常的默想权衡\Quote{祂的一切典章}，即义人的赏报，不敬虔者的惩罚，应受管教者的鞭笞，以及应受考验者的试炼。\Quote{我也没有将祂的公义从我身上丢弃：}不像那些人在重担下昏晕，转回自己的呕吐。\par

24. \Quote{我在祂面前必纯洁无疵，并保守自己远离我的罪孽}\BibleRef{咏 17:23}。\par

25. 主必按我的公义报答我\BibleRef{咏 17:24}。因此，不仅因那藉着爱而工作的信德之宽广，也因那坚忍之长久，主必按我的公义报答我。\Quote{并按照我眼中的清洁。}并非如人所见，而是\Quote{在祂眼中。}因为\Quote{所见的是暂时的，所不见的是永远的：}\BibleRef{格后 4:18}这属于望德的高超。\par

26. \Quote{对圣洁的人，祢显示圣洁}\BibleRef{咏 17:25}。其中有一种隐含的深意，即祢在圣洁的人中被认出是圣洁的，因为祢使他们成圣。\Quote{对纯全的人，祢显示纯全。}因为祢不伤害任何人，但各人被自己的罪恶绳索捆绑。\BibleRef{箴 5:22}\par

27. \Quote{对纯洁的人，祢显示纯洁。}\BibleRef{咏 17:26}祢被祢所拣选的人拣选。\Quote{对乖僻的人，祢显示乖僻。}对乖僻的人，祢显得乖僻；因为他们说：\Quote{主的道路不公平：}\BibleRef{则 18:25}而他们的道路却不公平。\par

28. \BibleQuote{因为祢必使谦卑的人民痊愈} \BibleRef{咏 17:27}。现在，对那些乖戾之人，祢使那些承认自己罪过的人得痊愈，这似乎显得乖戾。\BibleQuote{但祢必贬抑骄傲者的眼目。}然而，那些\BibleQuote{不知天主的成义，而力图建立自己的成义} \BibleRef{罗 10:3}的人，祢必贬抑他们。\par

29. \BibleQuote{因为祢，上主，要照亮我的烛台} \BibleRef{咏 17:28}。因为我们的光并非来自我们自己，而是\BibleQuote{祢，上主，要照亮我的烛台。我的天主，祢要照明我的黑暗。}因为我们因自己的罪而成为黑暗，但\BibleQuote{祢，我的天主，必要照明我的黑暗。}\par

30. \BibleQuote{因为藉着祢，我必从诱惑中被拯救} \BibleRef{咏 17:29}。因为不是凭我自己，而是藉着祢，我必从诱惑中被拯救。\BibleQuote{并要藉着我的天主跃过墙壁。}不是凭我自己，而是藉着我的天主，我要跃过那罪在人与天上的耶路撒冷之间筑起的墙壁。\par

31. \BibleQuote{我的天主，祂的道路是纯洁的} \BibleRef{咏 17:30}。我的天主不来到人中间，除非他们净化了信德的道路，好使祂能来到他们那里，因为\BibleQuote{祂的道路是纯洁的。}\BibleQuote{上主的话，经火炼净。}上主的话，经受患难之火炼净。\BibleQuote{祂是那些仰望祂之人的保护者。}凡不仰望自己，而仰望祂的人，不被同样的患难所吞灭。因为望德紧随信德。\par

32. \BibleQuote{除了上主，谁是神？} \BibleRef{咏 17:31}我们所事奉的。\BibleQuote{除了我们的天主，谁是神？}除了上主，谁是神？在善事奉之后，我们这些儿子必得着祂作为所盼望的产业。\par

33. \BibleQuote{天主，祢以力量束我的腰} \BibleRef{咏 17:32}。天主，祢束我的腰，使我坚强，免得松散飘荡的私欲阻碍我的行为和步伐。\BibleQuote{并使我道路纯洁。}并使我藉以到达祢那里的爱德之路纯洁，正如信德之路纯洁，祂藉此来到我这里。\par

34. \BibleQuote{祢使我的脚快如母鹿的蹄} \BibleRef{咏 17:33}。祢使我的爱德完美，得以超越这世界荆棘密布和缠绕的黑暗。\BibleQuote{并将我安置在高处。}并使我定睛于天上的居所，好使\BibleQuote{我充满天主的一切圆满} \BibleRef{弗 3:19}。\par

35. \BibleQuote{祢教导我的手作战} \BibleRef{咏 17:34}。祢教导我努力推翻那些力图关闭天国以对抗我们的仇敌。\BibleQuote{祢使我的手臂如铜弓。}祢使我热切追求善工的志向不知疲倦。\par

36. \BibleQuote{祢赐给我救恩的盾牌，祢的右手扶持了我} \BibleRef{咏 17:35}。祢恩宠的眷顾扶持了我。\BibleQuote{祢的管教引导我直到终点。}祢的纠正，不容许我偏离正路，引导我使我所做的一切都归向那终点，使我得以与祢结合。\BibleQuote{这祢的管教，必要教导我。}祢那同样的纠正必教导我达到它所引导我到达之处。\par

37. \BibleQuote{祢在我脚下扩展了我的步伐} \BibleRef{咏 17:36}。肉身的狭窄不能阻碍我，因为祢扩展了我的爱德，使我甚至在这些必死的肢体和属于我的部分中喜乐地工作。\BibleQuote{我的脚踪未曾动摇。}无论是我的行径，还是我为效法者所留下的印记，都没有动摇。\par

38. \BibleQuote{我要追赶我的仇敌，并擒获他们} \BibleRef{咏 17:37}。我要追赶我肉体的情欲，不被它们擒获，反而擒获它们，使它们被消灭。\BibleQuote{我不转回，直到他们灭绝。}我不停止这目的，直到那些在我周围喧嚷的仇敌失败。\par

39. \BibleQuote{我要击打他们，他们不能站立} \BibleRef{咏 17:38}：他们不能抵挡我。\BibleQuote{他们必倒在我脚下。}当他们被击倒时，我将把那些爱摆在我面前，使我永远行走。\par

40. \BibleQuote{祢以力量束我的腰，使我作战} \BibleRef{咏 17:39}。祢将我的肉身松弛的欲望用力量束缚起来，使我在这样的争战中不拖累。\BibleQuote{祢使那些起来攻击我的人屈伏在我之下。}祢使那些追赶我的人受迷惑，使他们被置于我之下，他们曾想凌驾于我之上。\par

41. \BibleQuote{祢使我的仇敌在我面前转背} \BibleRef{咏 17:40}。祢转过来我的仇敌，使他们成为我的背，即跟随我。\BibleQuote{祢消灭了那些恨我的人。}但他们中其余的，若坚持仇恨，祢便消灭了。\par

42. \BibleQuote{他们呼号，却无人拯救} \BibleRef{咏 17:41}。因为谁能拯救那些祢不愿拯救的人呢？\BibleQuote{他们向主呼号，主却不俯听。}他们并非向随便什么人呼号，而是向上主呼号，但上主不认为他们堪当被俯听，因为他们不离开自己的邪恶。\par

43. \BibleQuote{我要将他们打碎，如同风前的尘土} \BibleRef{咏 17:42}。我要将他们打得粉碎，因为他们枯干，得不到天主仁慈的雨水；他们被高举，骄傲自负，因而从坚定不动摇的望德，如同从大地的坚实稳固中被刮走。\BibleQuote{我要毁灭他们，如同街上的泥土。}在他们放荡不羁、沿着许多人所行的灭亡宽路上奔跑时，我要毁灭他们。\par

44. \BibleQuote{祢必救我脱离民众的反对} \BibleRef{咏 17:43}。祢必救我脱离那些说过\Quote{如果我们放祂走，全世界都要跟随祂}之人的反对。\par

45. \BibleQuote{祢要使我成为外邦人的元首。我素不认识的百姓必事奉我。}外邦的百姓，我虽未在肉身临在去拜访他们，却事奉了我。\BibleQuote{他们一听见我的名声就顺从我}\BibleRef{咏 17:44}。他们并未以眼见我，但接待了我的宣讲者，一听见就顺从了我。\par

46. \Quote{外邦的子民向我撒谎。}那些不配称为我的子民，而是外邦子民——对之正可说是\BibleQuote{你们是出于你们的父亲魔鬼}\BibleRef{若 8:44}——向我撒谎。\BibleQuote{外邦的子民已经衰老}\BibleRef{咏 17:45}。\Quote{并且在自己的路上跛行。}如同那些一足软弱的人，因固守旧约而拒绝新约，在他们自己的旧律中也成了跛子，随从他们自己的传统，而非天主的。因为他们以未洗手等琐事提出指控\BibleRef{玛 15:2}，因为那是他们自己开辟并因长久使用而磨损的路，偏离了天主诫命的道路。\par

47. \BibleQuote{主永生，我的天主应受赞美。}\BibleQuote{但随从肉性的思念是死亡}\BibleRef{罗 8:6}；因为\BibleQuote{主永生，我的天主应受赞美。愿我救恩的天主受尊崇}\BibleRef{咏 17:46}。愿我不以属世的方式思想我救恩的天主，也不向祂企求这世上的救恩，而是那高天的救恩。\par

48. \BibleQuote{天主，祢赐我复仇，使万民屈服于我}\BibleRef{咏 17:47}。天主，祢藉使万民屈服于我而为我复仇。\Quote{祢救我脱离发怒的仇敌：}那些犹太人喊着说：\BibleQuote{钉死祂，钉死祂}\BibleRef{若 19:6}。\par

49. \BibleQuote{祢要从起来攻击我的人中高举我}\BibleRef{咏 17:48}。祢要从那些在我受难时起来攻击我的犹太人中，在我复活时高举我。\Quote{祢要救我脱离不义的人。}祢要从他们的不义统治中拯救我。\par

50. \BibleQuote{为此，上主，我将在外邦人中称谢祢}\BibleRef{咏 17:49}。为此，外邦人要藉我在祢面前称谢祢，上主。\Quote{我要歌颂祢的圣名。}祢将因我的善行而更广为人知。\par

51. \BibleQuote{将祂君王的救恩显大}\BibleRef{咏 17:50}。天主彰显并使之奇妙的，是祂的圣子赐予信徒的救恩。\Quote{向祂的受傅者广施慈爱：}天主向祂的基督广施慈爱；\Quote{于达味及其后裔，直到永远：}就是向那手强壮的拯救者自己，祂已战胜了这世界；也向那些因信福音而由祂永远所生的人。凡在这圣咏中不能适用于主自己，即教会的元首的话，都当应用于教会。因为说话的是整个基督，其中包含了祂所有的肢体。\par

\SourceRecord{Translated by an anonymous scholar.；Nicene and Post-Nicene Fathers, First Series；Vol. 8.；Edited by Philip Schaff.；Buffalo, NY: Christian Literature Publishing Co.,；1888.；<http://www.newadvent.org/fathers/1801018.htm>.}

% structure: p0001 source=h1 target=section reason=page-title

\section{咏第十九篇注释}

致候经，达味本人的诗篇。\par

1. 这是一个众所周知的标题；以下所言并非主耶稣基督亲自所说，而是论及祂的。\par

2. \BibleQuote{诸天述说天主的荣耀}\BibleRef{咏 18:1}。义人般的福音作者，天主居住在他们内如同在诸天中，他们彰显了我们的主耶稣基督的荣耀，或是圣子在地上光荣圣父的荣耀。\BibleQuote{穹苍宣扬祂手中的作为。}穹苍宣扬主大能的作为，祂如今藉圣神的保证使天成为天，而先前它因恐惧只是地。\par

3. \BibleQuote{日复一日发出言语}\BibleRef{咏 18:2}。圣神向属灵的人赐下不变的天主智慧之丰满，即那太初与天主同在的圣言。\BibleRef{若 1:1}\BibleQuote{夜复夜传报知识。}向属血肉的人，如同向远处的人，肉身的死亡性藉着传递信德，传报将来的知识。\par

4. \BibleQuote{没有话语，也没有语言，听不见它们的声音}\BibleRef{咏 18:3}。在这些之中，福音作者的声音没有被听见，因为福音是用每一种语言传讲的。\par

5. \BibleQuote{它们的声音传遍天下，它们的言语达到地极}\BibleRef{咏 18:4}。\par

6. \BibleQuote{祂在太阳中设置了帐幕。}如今，为了与暂时错误的权势作战，主将要在地上投下不是和平而是刀剑，\BibleRef{玛 10:34}在时间或显现中，祂设立了，可说是，祂的军事寓所，即祂降生成人的安排。\BibleQuote{祂如同新郎出洞房}\BibleRef{咏 18:5}。祂从童贞女的子宫中出来，在那里天主与人性结合，如同新郎与新娘。\BibleQuote{如同巨人欢然奔路。}祂欢然如同一位极其强壮者，以无可比拟的能力超越所有其他人，不是为了居住，而是为了奔跑祂的路。因为，\BibleQuote{祂没有站在罪人的道路上。}\par

7. \BibleQuote{祂的出发从最高的天}\BibleRef{咏 18:6}。祂的出发来自圣父，不是在时间中，而是从永远，即祂由圣父所生。\BibleQuote{祂的相遇直达高天。}在完全的天主性中，祂相遇达到与圣父平等。\BibleQuote{没有能躲避祂的热气的。}然而，既然\BibleQuote{圣言成了血肉，居住在我们中间，}\BibleRef{若 1:14}祂取了我们必死的本性，祂不容许任何人以死亡的阴影为自己开脱；因为圣言的热气甚至渗透了它。\par

8. \BibleQuote{上主的法律是完善的，能畅快人灵}\BibleRef{咏 18:7}。因此，上主的法律就是祂自己，祂来为成全法律，而不是废除；\BibleRef{玛 5:17}一条完善的法律，\BibleQuote{祂没有犯过罪，口中也没有诡诈，}\BibleRef{伯前 2:22}不用奴役的轭压迫人灵，而是使人归向祂，在自由中效法祂。\BibleQuote{上主的证言是可信的，使愚蒙获得智慧。}\BibleQuote{上主的证言是可信的；}因为\BibleQuote{除了子，没有人认识父，除了子愿意启示的人，也不认识祂，}\BibleRef{玛 11:27}这些事向智慧人隐藏，却启示给婴孩；\BibleRef{路 10:21}因为\BibleQuote{天主拒绝骄傲的人，却赐恩宠给谦卑的人。}\BibleRef{雅 4:6}\par

9. \BibleQuote{上主的规条是正直的，能悦乐人心}\BibleRef{咏 18:8}。上主的一切规条在祂内都是正直的，因为祂所教导的，祂都实行了；这样，那些效法祂的人可以心中喜乐，自由地出于爱行事，而不是奴役般地出于恐惧。\BibleQuote{上主的诫命是光明的，能照亮眼睛。}\BibleQuote{上主的诫命是光明的，}没有肉体礼仪的幔子，能照亮内心之人的眼目。\par

10. \BibleQuote{对上主的敬畏是纯洁的，永远常存}\BibleRef{咏 18:9}。\BibleQuote{对上主的敬畏；}不是那法律下痛苦的恐惧，因极度害怕暂时之物的丧失，而灵魂因爱慕这暂时之物行淫；而是纯洁的敬畏，教会越热烈地爱她的净配，就越小心地避免冒犯祂，因此，\BibleQuote{完美的爱}不\BibleQuote{除去}这\BibleQuote{敬畏，}\BibleRef{若一 4:18}而是它永远常存。\par

11. \BibleQuote{上主的典章是真实的，全然公义。}那\BibleQuote{不审判任何人，而将一切审判交给了子}\BibleRef{若 5:22}的祂的典章，是在不变的真实中公义的。因为无论是在威胁还是应许中，天主都不欺骗任何人，也没有人能从不敬虔者身上夺去祂的惩罚，或从敬虔者身上夺去祂的赏报。\BibleQuote{比黄金和许多宝石更可羡慕}\BibleRef{咏 18:10}。无论是\BibleQuote{黄金和宝石本身珍贵，}还是\BibleQuote{极其珍贵，}或是\BibleQuote{极其可羡慕；}总之，天主的典章比这世界的浮华更可羡慕；由于对这浮华的贪求，人便不羡慕天主的典章，反而惧怕、鄙视或不相信它们。但如果有人本身是黄金和宝石，以致不被火焚烧，而被收进天主的宝库，那么他必羡慕天主的典章胜过自己，因为他的旨意优先于自己的旨意。\BibleQuote{比蜂房的蜜更甘甜。}无论有人现在已经是蜜，已从今生的锁链中解脱，等候日子来到天主的筵席；还是他如同蜂房，被今生包裹着如同被蜡包裹，尚未与它混合而成为一体，而是充满它，需要天主手的压榨，不是压迫而是挤出，使之从暂时的人生中被过滤到永生中：对于这样的人，天主的典章比他本人更甘甜，因为它们是\BibleQuote{比蜂房的蜜更甘甜。}\par

12. \BibleQuote{因为祢的仆人遵守了它们}\BibleRef{咏 18:11}。因为那不遵守它们的人，上主的日子是苦涩的。\BibleQuote{在遵守它们上有大赏报。}这赏报不在任何外在的利益，而在于这事本身，即天主的判断被遵守，就有大赏报；之所以大，是因为人以此欢欣。\par

13. \BibleQuote{谁能察觉自己的过犯？}\BibleRef{咏 18:13} 但在罪恶中能有什么甘甜呢？那里没有领悟。因为谁能了解罪恶呢？罪恶遮蔽了眼睛，真理对之是愉悦的，天主的判断对之是可欲且甘甜的。是的，如同黑暗遮蔽眼睛，罪恶也遮蔽心灵，不使它看见光明或自身。\par

14. \BibleQuote{求祢洗净我隐密的过失。} 从隐藏在我内的私欲中，求祢洗净我，上主。\BibleQuote{并从他人的}\OriginalTerm{过失}{faults}\BibleQuote{中保守祢的仆人}\BibleRef{咏 18:14}。求祢不要让我被他人引入歧途。因为那已被洗净自己过失的人，不会成为他人过失的猎物。因此，保守祢的仆人，不要受他人私欲的侵扰——不是骄傲的人，也不是那愿作自己主人的人，而是祢的仆人。\BibleQuote{若它们不得在我身上作主，我便无瑕。} 若我自己的隐密罪恶和他人的罪恶都不得在我身上作主，我便无瑕。因为罪恶没有第三个来源，只有自己的隐密罪（魔鬼由此堕落）和他人的罪（人由此被诱惑，以致因同意而使之成为自己的罪）。\BibleQuote{我将从大过犯中被洗净。} 这大过犯不是骄傲又是什么？因为没有比离弃天主更大的罪，这正是：\BibleQuote{人骄傲的开端}\BibleRef{德 10:12}。确实，那也摆脱了此过犯的人，便是无瑕的；因为这对于归向天主的人是最后的，对于离开祂的人却是最初的。\par

15. \BibleQuote{愿我口中的言辞和心中的默祷常在祢面前蒙悦纳}\BibleRef{咏 18:15}。我心中的默祷不是为了取悦人的虚荣，因为如今已无骄傲，而是常在祢面前，祢关注纯洁的良心。\BibleQuote{上主，我的援助者，我的救赎者}\BibleRef{咏 18:16}。上主，我的援助者，在我趋向祢时；因为祢是我的救赎者，使我得以出发走向祢：免得有人将他的归向祢归因于自己的智慧，或将他的到达祢归因于自己的力量，反被祢驱逐，因为祢抵制骄傲的人；因为这人不曾被洗净那个大过犯，也不在祢面前蒙悦纳，是祢救赎了我们使我们得以归向，祢援助了我们使我们得以到达祢。\par

\SourceRecord{Translated by an anonymous scholar.；Nicene and Post-Nicene Fathers, First Series；Vol. 8.；Edited by Philip Schaff.；Buffalo, NY: Christian Literature Publishing Co.,；1888.；<http://www.newadvent.org/fathers/1801019.htm>.}

% structure: p0001 source=h1 target=section reason=page-title

\section{论圣咏第二十篇}

交与乐官，达味圣咏。\par

这是一个耳熟能详的标题；说话的不是基督，而是先知以祝愿的形式向基督发言，预言未来的事。\par

\BibleQuote{愿上主在困苦之日垂听祢} \BibleRef{咏 19:1} 。上主在祢说\Quote{父啊，光荣祢的子}的那日垂听祢。\BibleQuote{愿雅各伯的天主的名保护祢}。因为那较年轻的民族属于祢。既然\Quote{年长的要服侍年幼的}。\par

\BibleQuote{从圣所派遣援助给祢，从熙雍扶持祢} \BibleRef{咏 19:2} 。为祢造成一个圣化的身体，即教会，从安全的守望中，等待祢从婚筵归来。\par

\BibleQuote{愿祢纪念祢的一切祭献} \BibleRef{咏 19:3} 。使我们纪念祢为我们所承受的一切伤害和凌辱。\Quote{愿祢的全燔祭成为丰腴}。并将祢完全奉献给天主的十字架，转化为复活的喜乐。\par

\Emph{\OriginalTerm{细拉}{Diapsalma}}。\BibleQuote{愿上主照祢的心意报答祢} \BibleRef{咏 19:4} 。愿上主报答祢，不是按那些人的心意，他们以为借着迫害能消灭祢；而是按祢的心意，祢知道祢的苦难将有何益处。\BibleRef{若 12:32} \BibleQuote{并且成就祢的一切计谋}。\BibleQuote{并且成就祢的一切计谋}，不仅祢为朋友舍命的那计谋，使腐败的麦粒复活而结出更多果实；\BibleRef{若 15:13} \BibleRef{若 12:24} 而且也照那计谋，即\Quote{以色列一部分硬了心，直到外邦人全数进入，这样全体以色列就必得救}。\BibleRef{罗 11:25}\par

\BibleQuote{我们要因祢的救恩欢欣} \BibleRef{咏 19:5} 。我们要因死亡绝不能伤害祢而欢欣；因为祢也要表明它也不能伤害我们。\Quote{并因我们天主的名而被尊大}。承认祢的名不仅不会毁灭我们，反而会尊大我们。\par

\BibleQuote{愿上主满全祢的一切请求}。愿上主不仅满全祢在地上所作的请求，也包括祢在天上为我们代求的那些。\BibleQuote{现在我知道上主已拯救了祂的受傅者} \BibleRef{咏 19:6} 。现在在预言中已向我显示，上主将再次复活祂的基督。\BibleQuote{祂必从祂的圣天上垂听祂}。祂不仅从地上垂听祂（在那里祂祈求受光荣），\BibleRef{若 17:1} 也从天上垂听祂，在那里祂坐在父的右边为我们代祷，\BibleRef{希 7:25} 并从那里将圣神倾注于信祂的人。\BibleQuote{祂右手的拯救大有能力}。我们的力量在于祂恩宠的拯救，即使在患难中祂也赐予帮助，使我们\Quote{软弱时，反而刚强}。\BibleRef{格后 12:10} 因为\Quote{人的拯救}是虚假的，它并非来自祂的右手，而是来自左手；因为那些在罪中获得暂时安全的人，因此被高举到极大的骄傲中。\par

\BibleQuote{有些人依靠战车，有些人依靠战马} \BibleRef{咏 19:7} 。有些人被时间性财物的不断流转所牵引；有些人被提升至骄傲的荣誉，并以此夸耀。\BibleQuote{但我们依靠我们天主的名而夸耀}。但我们，将希望寄托于永恒之事，不寻求自己的荣耀，要依靠我们天主的名而夸耀。\par

\BibleQuote{他们已被捆绑，跌倒了} \BibleRef{咏 19:8} 。因此他们被时间性事物的欲望所束缚，害怕宽恕上主，以免因\Quote{罗马人}失去地位；\BibleRef{若 11:48} 并且猛烈撞击那绊脚石和使人失足的磐石，从天上的希望坠落了；以色列一部分硬了心，他们对天主的正义无知，想要建立自己的正义。\BibleQuote{但我们起来，站得正直}。但我们，为使外邦人得以进入，从石头中给亚巴郎兴起子孙，\BibleRef{玛 3:9} 他们不追求正义，却获得了正义，并且起来了；\BibleRef{罗 9:30} 不是靠自己的力量，而是因信德成义，我们站得正直。\par

\BibleQuote{上主，请拯救君王}：使祂，在受难中给我们展示了战斗的榜样，也当呈上我们的祭品，即从死者中复活并升天的司祭。\BibleQuote{并在我们呼求祢的日子垂听我们} \BibleRef{咏 19:9} 。正如祂现在为我们献上祈祷，\Quote{在我们呼求祢的日子垂听我们}。\par

\SourceRecord{Translated by an anonymous scholar.；Nicene and Post-Nicene Fathers, First Series；Vol. 8.；Edited by Philip Schaff.；Buffalo, NY: Christian Literature Publishing Co.,；1888.；<http://www.newadvent.org/fathers/1801020.htm>.}

% structure: p0001 source=h1 target=section reason=page-title

\section{圣咏第二十一篇释义}

致终末，达味自己的圣咏。\par

这个标题是熟悉的；这篇圣咏关乎基督。\par

\BibleQuote{上主，君王要在祢的威能中喜乐} \BibleRef{咏 20:1}。主啊，在祢的威能中，圣言成了血肉，人而基督耶稣要喜乐。\Quote{且要在祢的救恩中极其欢跃。}并且在那祢藉以赐万物生命的恩中，要极其欢跃。\par

\BibleQuote{祢已满足了祂心灵的愿望} \BibleRef{咏 20:2}。\BibleQuote{祢也没有拒绝祂口中的祈求}。祂渴望吃逾越节晚餐（见 \BibleRef{路 22:15}），并且甘愿舍命，也甘愿再取回来，祢就赐给了祂（见 \BibleRef{若 10:18}）。祂说：\BibleQuote{我的平安}，\BibleQuote{我留给你们}（见 \BibleRef{若 14:27}），这事就成就了。\par

\BibleQuote{因为祢以甘饴的祝福迎接了祂} \BibleRef{咏 20:3}。因为祂先饮了祢甘饴的祝福，我们罪恶的苦胆就没有伤害祂。\Quote{\OriginalTerm{细拉}{Diapsalma}。祢将宝石的冠冕戴在祂头上。}在祂开始宣讲时，贵重的宝石被带来，环绕着祂；这些宝石就是祂的门徒，祂的宣讲应从他们开始。\par

\Quote{祂求生命，祢就赐给了祂：}祂求复活，说：\BibleQuote{父啊，光荣祢的子}（见 \BibleRef{若 17:1}），祢就赐给了祂，\BibleQuote{永远的寿数，无穷的岁月}（见 \BibleRef{咏 20:4}）。这是指教会将在世上享有的漫长年代，以及此后无尽的永恒。\par

\BibleQuote{祂的荣耀在祢的救恩中何其伟大} \BibleRef{咏 20:5}。的确，祂的荣耀在祢使祂复活的救恩中是伟大的。\Quote{祢将光荣和尊大加在祂身上。}但当祢使祂坐在天上祢的右边时，祢还要为祂增添光荣和尊大。\par

\Quote{因为祢必将祝福赐给祂，直到永远。}这就是祢要赐给祂的永远祝福：\BibleQuote{祢要以祢的慈颜使他欢喜快乐}（见 \BibleRef{咏 20:6}）。按祂的人性，祢要以祢的慈颜使他欢喜，祂曾向祢举目仰望。\par

\Quote{因为君王寄望于上主。}因为这位君王不骄傲，而是心里谦卑，他寄望于上主。\BibleQuote{且因至高者的慈爱，他必不动摇}（见 \BibleRef{咏 20:7}）。因至高者的慈爱，祂顺从至死，且死在十字架上，这并不扰乱祂的谦卑。\par

\Quote{愿祢的手被祢所有的仇敌寻见。}哦君王，当祢来审判时，愿祢的威能被祢所有的仇敌寻见；他们在祢的卑微中未能认出它。\BibleQuote{愿祢的右手找出所有仇恨祢的人}（见 \BibleRef{咏 20:8}）。愿那荣耀，就是祢在圣父右边统治的荣耀，在审判之日找出所有仇恨祢的人，予以惩罚；因为如今他们还没有寻见它。\par

\Quote{祢必使他们如烈火之炉：}祢必使他们内心燃烧，因他们意识到自己的不虔敬：\Quote{在祢显现的时候：}在祢彰显的时候。\BibleQuote{上主必在震怒中扰乱他们，火要吞灭他们}（见 \BibleRef{咏 20:9}）。然后，他们被主的报复所扰乱，良心的控告之后，他们将被交付永火，被吞灭。\par

\Quote{祢必从地上消灭他们的果实。}他们的果实，因为是属地的，祢要从地上消灭。\BibleQuote{并从人子中消灭他们的后裔}（见 \BibleRef{咏 20:10}）。以及他们的作为；或者说，凡他们所诱惑的人，祢都不会将他们算在祢召叫承受永远基业的人子中。\par

\Quote{因为他们将恶事转向祢。}这惩罚必要报应在他们身上，因为他们所想象因祢的统治而临于他们身上的恶事，反而转向祢，致使祢受死。\BibleQuote{他们图谋恶计，却不能成就}（见 \BibleRef{咏 20:11}）。他们图谋计策，说：\BibleQuote{叫一个人替众人死，这是有益的}（见 \BibleRef{若 11:50}），这计策他们不能成就，因为他们不知道自己在说什么。\par

\Quote{因为祢必使他们卑微。}因为祢要把他们列在那些祢将从之转脸、轻视与厌恶的人中。\BibleQuote{在祢所遗留的物中，祢预备了他们的脸面}（见 \BibleRef{咏 20:12}）。并且在祢所遗留的这些事物，即属世王国的欲望中，祢为祢的受难预备了他们的厚颜无耻。\par

\BibleQuote{上主，求祢在祢的威能中高举}（见 \BibleRef{咏 20:13}）。主啊，祢在卑微中他们未能认出，愿祢在祢的威能中受高举，这威能他们曾以为是软弱。\Quote{我们要歌颂，赞美祢的大能。}我们将在心里和行动上庆祝并宣扬祢的奇事。\par

\SourceRecord{Translated by an anonymous scholar.；Nicene and Post-Nicene Fathers, First Series；Vol. 8.；Edited by Philip Schaff.；Buffalo, NY: Christian Literature Publishing Co.,；1888.；<http://www.newadvent.org/fathers/1801021.htm>.}

% structure: p0001 source=h1 target=section reason=page-title

\section{\WorkTitle{圣咏第二十二篇释义}}

\Emph{致终末，为清晨的摄取，达味圣咏。}\par

1. 主耶稣基督亲自发言，为祂自己的复活，\BibleQuote{致终末}。\BibleRef{若 20:1} 因为在七日的第一日清晨，祂复活了，被摄取进入永生，\BibleQuote{死亡不再统治祂}。\BibleRef{罗 6:9} 以下的话是被钉十字架者以第一人称说的。因为这圣咏的开头正是祂悬在十字架上所呼喊的话，祂也担当了旧人的位格，祂带走了旧人的可朽性。因为我们旧人已与祂同钉十字架。\BibleRef{罗 6:6}\par

2. \BibleQuote{天主，我的天主，求祢垂视我，祢为何舍弃了我，远离我的救恩？}\BibleRef{咏 21:1} 远离我的救恩：因为\BibleQuote{救恩远离罪人}。\BibleQuote{我罪恶的言语}。因为这并非义人的言语，而是我罪恶的言语。因为那被钉十字架的旧人说话，甚至不知天主为何舍弃了他；或者也可以这样理解：我罪恶的言语远离我的救恩。\par

3. 我的天主，我白天向祢呼求，祢不俯听。\BibleRef{咏 21:2} 我的天主，我在今世的顺境中向祢呼求，求这些顺境不要改变；祢却不俯听，因为我以罪恶的言语向祢呼求。\BibleQuote{夜间也是如此，但并非出于我的愚昧。} 同样，我在今世的逆境中向祢呼求顺境；祢照样不俯听。祢这样做并非出于我的愚昧，而是为使我获得智慧，知道该为什么事呼求祢——不是以罪恶的言语贪求现世生命，而是以归向祢的言语恳求永生。\par

4. \BibleQuote{但祢居于圣所，以色列的赞美。}\BibleRef{咏 21:3} 但祢居于圣所，因此不会俯听罪恶的不洁言语。\BibleQuote{赞美}属于看见祢的人；不属于那因尝禁果而寻求自己赞美的人，以致在他肉体的眼睛睁开后，试图躲避祢的视线。\par

5. \BibleQuote{我们的祖先曾寄望于祢。} 即所有的义人，他们不寻求自己的赞美，只寻求祢的赞美。\BibleQuote{他们寄望于祢，祢便拯救了他们。}\BibleRef{咏 21:4}\par

6. \BibleQuote{他们向祢呼求，便得了救。} 他们向祢呼求，并非以罪恶的言语——救恩远离罪恶的言语；因此他们得了救。\BibleQuote{他们寄望于祢，便没有蒙羞。}\BibleRef{咏 21:5} \BibleQuote{他们寄望于祢，} 他们的希望没有欺骗他们，因为他们没有把这希望放在自己身上。\par

7. \BibleQuote{但我是虫，不是人。}\BibleRef{咏 21:6} 但我——现在不是以\Person{亚当}的位格说话，而是以我自己的位格——耶稣基督，在肉身中生于人而无人的生育，为使我作为人超越众人；这样，人的骄傲至少可以效法我的谦卑。\BibleQuote{是人的耻笑，是百姓的鄙弃。} 在这谦卑中，我成了人的耻笑，以致他们辱骂说：\BibleQuote{你是祂的门徒。}\BibleRef{若 9:28} 并且百姓藐视我。\par

8. \BibleQuote{凡看见我的都嘲笑我。}\BibleRef{咏 21:7} 凡看见我的都讥笑我。\BibleQuote{他们开口说话，又摇着头。}\BibleRef{玛 27:39} 他们说话不是出于心，而是出于嘴唇。\par

9. 因为他们摇着头讥笑说：\BibleQuote{他倚靠上主，让上主救他吧；}\BibleRef{玛 27:43} \BibleQuote{上主既然喜悦他，就拯救他吧。}\BibleRef{咏 21:8} 这是他们的话，但只是\BibleQuote{用嘴唇}说的。\par

10. \BibleQuote{因为祢是从母胎中牵引我的。}\BibleRef{咏 21:9} 因为祢不仅从童贞女的胎中牵引我（因为这是一切人出生的规律，即从母胎中被牵引出来），而且从犹太民族的母胎中牵引我；凡将自己的救恩置于遵守安息日、割损礼等肉体习俗中的人，就被犹太民族的黑暗所遮盖，尚未进入基督的光明中。\BibleQuote{从我母的乳房，祢就是我的希望。} \BibleQuote{我的希望，} 天主啊，不是从我开始被童贞女奶汁喂养的时候（因为那时早已有希望），而是从我所说的会堂的乳房——即从母胎中，祢牵引了我，使我不致吮吸肉体的习俗。\par

11. \BibleQuote{我从母胎中就倚靠祢。}\BibleRef{咏 21:10} 这是会堂的母胎，它没有怀我，反而抛弃了我；但我没有跌倒，因为祢扶持了我。\BibleQuote{从我母亲的腹中，祢就是我的天主。} 从我母亲的腹中：我母亲的腹中没有使我像婴孩那样忘记祢。\par

12. \BibleQuote{祢是我的天主，} \BibleQuote{求祢不要远离我，因为患难临近了。}\BibleRef{咏 21:11} 因此，祢是我的天主，求祢不要远离我，因为患难临近了我；这患难是在我的身体里。\BibleQuote{因为无人帮助。} 因为祢若不帮助，谁能帮助呢？\par

13. \BibleQuote{许多牛犊围困了我。} 许多放荡的民众围困了我。\BibleQuote{肥壮的公牛包围了我。}\BibleRef{咏 21:12} 他们的首领因我的苦难而欢喜，\BibleQuote{包围了我。}\par

14. \BibleQuote{他们向我张口。}\BibleRef{咏 21:13} 他们向我张口，不是依据祢的圣经，而是出于自己的私欲。\BibleQuote{如同掠夺咆哮的狮子。} 如同狮子，其掠夺是我被捉拿、被牵引；其咆哮是\BibleQuote{钉死，钉死。}\BibleRef{若 19:6}\par

15. \BibleQuote{我好像水一般倾泻，我的一切骨骸都散开了。}\BibleRef{咏 21:14} \BibleQuote{我好像水一般倾泻，} 当我的迫害者跌倒时；由于恐惧，我身体的支柱——即教会，我的门徒们都离我而散了。\BibleRef{玛 26:56} \BibleQuote{我的心像蜡一般熔化，在我肺腑的中间。} 我的智慧——即圣经中关于我的记载——曾像是坚硬封闭、不为人所理解的；但当我受难的火焰临到时，它便像熔化的蜡，被彰显出来，并珍藏在我教会的记忆中。\par

16. \Quote{我的力气乾枯如瓦片} \BibleRef{咏 21:15}。我的力气因我的苦难而乾枯；不是如草，而是如瓦片，因火而变得更坚固。\Quote{我的舌头贴紧我的上颚。} 而那些我将藉以说话的人，将我的诫命存于心中。\Quote{祢将我带到死尘之中。} 祢将我带到那注定死亡、如风吹散地面之尘的不敬虔者中。\par

17. \Quote{因为许多犬类围着我} \BibleRef{咏 21:16}。许多围着我狂吠，不是为了真理，而是为了习俗。\Quote{恶人的议会围着我。} 恶人的议会围困了我。\Quote{他们刺透了我的手和脚。} 他们用钉子刺透了我的手和脚。\par

18. \Quote{他们数算了我所有的骨头} \BibleRef{咏 21:17}。当被悬在十字架的木材上时，他们数算了我所有的骨头。\Quote{确实，这些人观看，并注视着我。} 确实，这些相同的人，即未改变的人，观看并注视着我。\par

19. \Quote{他们分了我的外衣，并为我的内衣拈阄} \BibleRef{咏 21:18}。\par

20. \Quote{但祢，主啊，求祢不要远离我，将祢的帮助撤回} \BibleRef{咏 21:19}。但祢，主啊，使我再次复活，不像其他人在世界末日，而是立即复活。\Quote{看顾我的辩护。} \Quote{看顾}，使他们绝不能伤害我。\par

21. \Quote{救我的灵魂脱离刀剑。} \Quote{救我的灵魂} 脱离纷争的舌头。\Quote{并救我的独一者脱离犬类的爪} \BibleRef{咏 21:20}。并救我的教会脱离那按习俗狂吠的百姓的权势。\par

22. \Quote{救我脱离狮子的口：} 救我脱离这世界王国的口；\Quote{并救我的谦卑脱离独角兽的角} \BibleRef{咏 21:21}。并救我的谦卑脱离那骄傲者的高傲，他们自高自大，追求特殊优越，不容分享者，惟独我的谦卑例外。\par

23. \Quote{我要向我的弟兄们宣扬祢的圣名} \BibleRef{咏 21:22}。我要向谦卑的人，并向我的弟兄们宣扬祢的圣名，他们彼此相爱，正如他们蒙我所爱一样。\BibleRef{若 17:6} \Quote{在教会中我要歌颂祢。} 在教会中我要欢欣地宣讲祢。\par

24. \Quote{你们敬畏主的人，要赞美祂。} \Quote{你们敬畏主的人，} 不要寻求自己的赞美，而要 \Quote{赞美祂。} \Quote{雅各伯所有的后裔，要尊崇祂} \BibleRef{咏 21:23}。雅各伯所有的后裔——那为长者所服役的那一位——要尊崇祂。\par

25. \Quote{愿以色列所有的后裔都敬畏祂。} 愿所有重生、恢复见到天主的人 \Quote{敬畏祂。} \Quote{因为祂没有轻视，也没有忽略穷人的祈祷} \BibleRef{咏 21:24}。因为祂没有轻视那祈祷，不是那用罪恶之言向天主呼号、不愿越过虚浮生命之人的祈祷，而是穷人的祈祷，没有因短暂的浮华而膨胀。\Quote{祂也没有向我转脸。} 不像那说\Quote{我要向祢呼号，但祢不听}的人。\Quote{我向祂呼号时，祂就垂听了我。}\par

26. \Quote{我的赞美在于祢} \BibleRef{咏 21:25}。因为我不寻求自己的赞美，\BibleRef{若 8:50}因为祢是我的赞美，祢住在圣所；以色列的赞美，祢垂听现在向祢恳求的圣者。\Quote{在大会中我要向祢称谢。} 在全世界的教会中 \Quote{我要向祢称谢。} \Quote{我要在敬畏祂的人面前偿还我的誓愿。} 我要在敬畏祂的人面前献上我体血的圣事。\par

27. \Quote{穷人必会吃饱，且得满足} \BibleRef{咏 21:26}。谦卑的人和轻视世界的人必会吃饱，并效法我。因为他们既不贪求今世的富足，也不惧怕其匮乏。\Quote{寻求主的人必赞美祂。} 因为对主的赞美就是那饱足的涌流。\Quote{他们的心要永远活着。} 因为那食物是心之食物。\par

28. \Quote{大地四极的人都要记念自己，并归向主} \BibleRef{咏 21:27}。他们必要记念自己：因为外邦人生于死亡，专注于外在事物，天主已被他们遗忘；那时大地四极都要归向主。\Quote{万族的诸支派都要在祂面前朝拜。} 万族的诸支派要在自己的良心中朝拜。\par

29. \Quote{因为王国属于主，祂要治理万民} \BibleRef{咏 21:28}。因为王国属于主，不属于骄傲的人；祂要治理万民。\par

30. \Quote{地上所有的富足人都吃了，并且朝拜} \BibleRef{咏 21:29}。地上的富足人也吃了主卑微的体，虽然他们没有像穷人那样充满到效法的地步，但他们仍然朝拜。\Quote{凡降世之人都要在祂面前跌倒。} 因为只有祂看见，那些人如何跌倒，他们放弃天上的交往，甘愿在地上在人前显得幸福，而这些人是看不见他们跌倒的。\par

31. \Quote{我的灵魂要生活于祂。} 我的灵魂，在今世的轻蔑中在人看来仿佛死亡，却要生活，不是为自己，而是为祂。\Quote{我的后裔必要事奉祂} \BibleRef{咏 21:30}。我的作为，或那些藉着我信祂的人，必要事奉祂。\par

32. \Quote{未来的世代必将宣告给主} \BibleRef{咏 21:31}。新约的世代必将宣告，归于主的荣耀。\Quote{诸天要传扬祂的公义。} 福音作者们要传扬祂的公义。\Quote{给那将出生的百姓，那是主所造的。} 给那将因信而生为主之民的百姓。\par

\SourceRecord{Translated by an anonymous scholar.；Nicene and Post-Nicene Fathers, First Series；Vol. 8.；Edited by Philip Schaff.；Buffalo, NY: Christian Literature Publishing Co.,；1888.；<http://www.newadvent.org/fathers/1801022.htm>.}

% structure: p0001 source=h1 target=section reason=page-title

\section{圣咏第二十三篇释义}

达味自己的圣咏。\par

1. 教会向基督说：\Quote{主牧养我，我必一无所缺} \BibleRef{咏 22:1}。主耶稣基督是我的牧者，\Quote{我必一无所缺。}\par

2. \Quote{在青草地上祂使我躺卧} \BibleRef{咏 22:2}。在鲜嫩牧草之地，引我进入信德，祂安置我在那里受滋养。\Quote{祂领我到静谧的水旁。}借着圣洗之水，那些失去健康与力量的人得以复苏，祂领我上来。\par

3. \Quote{祂使我的灵魂苏醒，为了祂的名，引导我走上正义的途径} \BibleRef{咏 22:3}。祂领我走上祂正义的窄路，其中行走的人甚少；并非因我的功德，而是为了祂的名。\par

4. \Quote{纵使我走过死亡的幽谷} \BibleRef{咏 22:4}。纵使我走过此生，它不过是死亡的幽影。\Quote{我不怕凶险，因为祢与我同在。}我不怕凶险，因为祢借着信德居住在我心中；祢如今与我同在，使我经过死亡的幽影后，也能与祢同在。\Quote{祢的牧杖和棍杖，它们安慰了我。}祢的管教，如同羊群的牧杖，又如同渐长的孩童的棍杖，从属自然生命成长到属灵生命，它们对我不曾是重担；反而安慰了我，因为祢顾念我。\par

5. \Quote{在我敌人面前，祢为我摆设筵席} \BibleRef{咏 22:5}。自牧杖之后——当我仍是幼小者，过着属自然的生活，在牧场上被牧养于羊群中——我且说，自牧杖之后，当我开始处于棍杖之下时，祢在我面前摆设筵席，使我不要再如婴孩般吃奶，\BibleRef{格前 3:2}而是年长后吃干粮，得以坚强，对抗那困扰我的人。\Quote{祢用油膏了我的头。}祢以属灵的喜乐使我心神欢畅。\Quote{祢的杯爵何等甘美，令人陶醉！}祢的杯爵使人忘却从前虚妄的欢乐，何其卓越！\par

6. \Quote{祢的慈爱必追随我终生：}即在我活于此可朽生命的全部时日，并非祢的生命，而是我的。\Quote{使我得以长久住在主的殿宇中} \BibleRef{咏 22:6}。如今祢的慈爱不仅在此追随我，且使我永远住在主的殿宇中。\par

\SourceRecord{Translated by an anonymous scholar.；Nicene and Post-Nicene Fathers, First Series；Vol. 8.；Edited by Philip Schaff.；Buffalo, NY: Christian Literature Publishing Co.,；1888.；<http://www.newadvent.org/fathers/1801023.htm>.}

% structure: p0001 source=h1 target=section reason=page-title

\section{圣咏第二十四篇释义}

\Emph{达味自己的一首圣咏，关于一周的第一天。}\par

1. 达味自己的一首圣咏，论及主的荣耀与复活，此事发生在一周第一天的清晨，这日现今被称为主日。\par

\BibleQuote{大地属于主，及其中所充满的，世界的地极，及所有住在其中的}\BibleRef{咏 23:1}；当主受光荣，被传报给万民信仰时，整个世界的地极就成了祂的教会。\BibleQuote{祂将大地奠立在海上}。祂将教会最稳固地建立在这世界的所有波浪之上，使波浪被教会征服，不能伤害教会。\BibleQuote{祂预备大地在江河之上}\BibleRef{咏 23:2}。江河流入大海，贪欲之人堕入世界；教会也征服这些，当世俗欲念被天主的恩宠克服后，教会已藉爱预备好接受不朽。\par

\BibleQuote{谁能登上主的山？}谁能登上主的义的高处？\BibleQuote{或谁能站在祂的圣所？}\BibleRef{咏 23:3}。或者谁能居于那奠立在海上、预备在江河之上的地方？\par

\BibleQuote{手洁心净}\BibleRef{咏 23:4}。那么谁能登上那里并居住在那里呢？不就是行为无罪、思想纯洁的人吗？\BibleQuote{他没有徒然接受他的灵魂。}谁没有将他的灵魂算作转瞬即逝的东西，而是感到灵魂是不朽的，渴望永恒、坚定而不变。\BibleQuote{也没有以欺骗向邻人起誓。}因此，他毫不欺诈，因为永恒的事物是单纯而不欺人的，他对待邻人也是如此。\par

\BibleQuote{这样的人必领受主的祝福，和救他的天主的仁慈}\BibleRef{咏 23:5}。\par

\BibleQuote{这便是寻求主的一代}\BibleRef{咏 23:6}。因为寻求祂的人就是这样诞生的。\BibleQuote{寻求雅各伯天主的面容。\Emph{细拉}}如今，他们寻求天主的面容，就是那使年幼的居首的那一位。\BibleRef{罗 9:12}\par

\BibleQuote{撤去你们的门户，你们众位首领}\BibleRef{咏 23:7}。你们所有在人间寻求权柄的人，请撤除你们所制作的、出于欲望和恐惧的入口，免得它们妨碍。\BibleQuote{永恒的门户，请高举。}永恒生命的入口，即弃绝世俗、皈依天主，请高举。\BibleQuote{光荣的君王要进来。}那位我们可因之而夸耀却不骄傲的君王要进来了：祂战胜了死亡的阴门，为自己打开了天上的居所，成就了祂所说的：\BibleQuote{你们放心，因为我已战胜了世界。}\BibleRef{若 16:33}\par

\BibleQuote{谁是这位光荣的君王？}有死的人性惊惶而问：\BibleQuote{谁是这位光荣的君王？}\BibleQuote{是主，强壮而英勇。}就是你们认为软弱而被打败的那一位。\BibleQuote{是主，战争中的勇者}\BibleRef{咏 23:8}。触摸祂的伤痕，你们会发现它们已痊愈，人性的软弱已恢复为不朽。主的光荣，本是属于大地，在那里与死亡争战，已经还清了。\par

\BibleQuote{撤去你们的门户，你们众位首领。}让我们从这里直接升入天堂。先知再次高声吹号：\BibleQuote{你们空中的众首领，也要撤除你们在人心中所设的门户，就是那些\InnerQuote{敬拜天上万军}的人。}\BibleRef{列下 17:16}\BibleQuote{永恒的门户，请高举。}永恒公义、爱和贞洁的门户，请高举；通过这些门户，灵魂爱慕唯一真天主，不向那些被称为神的众多存在行淫。\BibleQuote{光荣的君王要进来}\BibleRef{咏 23:9}。\BibleQuote{光荣的君王要进来，}使祂在圣父右边为我们转求。\par

\BibleQuote{谁是这位光荣的君王？}怎么！你这空中权势的首领，\BibleRef{弗 2:2}也惊讶地问：\BibleQuote{谁是这位光荣的君王？}\BibleQuote{是万军的主，祂就是光荣的君王}\BibleRef{咏 23:10}。是的，祂的身体现在复活了，祂虽受过试探，却超越你而行进；那曾受天使、诱惑者试探的，超越众天使。你们中没有一个能把自己放在我们前面阻止我们的道路，使我们敬拜他为神；\BibleQuote{无论是权势、天使、或异能，都不能使我们与基督的爱隔绝。}\BibleRef{罗 8:39}\BibleQuote{信赖主，胜于信赖王侯；凡夸耀的，应夸耀主。}\BibleRef{格前 1:31}这些固然是治理此世的有力者，但是\BibleQuote{万军的主，祂就是光荣的君王。}\par

\SourceRecord{Translated by an anonymous scholar.；Nicene and Post-Nicene Fathers, First Series；Vol. 8.；Edited by Philip Schaff.；Buffalo, NY: Christian Literature Publishing Co.,；1888.；<http://www.newadvent.org/fathers/1801024.htm>.}

% structure: p0001 source=h1 target=section reason=page-title

\section{圣咏第廿五篇释义}

\Emph{致终点，达味自己的圣咏。}\par

基督说话，但以教会的身分发言；因为所说的话更多地指向已归向天主的基督徒子民。\par

\BibleQuote{主啊，我向祢举起了我的灵魂} \BibleRef{咏 24:1}：以属灵的渴望我举起了灵魂，这灵魂曾因肉欲的渴望被踏在地上。\BibleQuote{我的天主，我信赖祢，我不会蒙羞} \BibleRef{咏 24:2}。我的天主，由于信赖我自己，我甚至被带到这肉身的软弱中；而我，在离弃天主时曾想如天主一样，连最微小昆虫的死亡都害怕，因我的骄傲而被嘲笑蒙羞；因此，现在，\BibleQuote{我信赖祢，我不会蒙羞。}\par

\Quote{也不容我的仇敌嘲笑我。} 他们不可嘲笑我，他们曾用蛇一般的隐秘暗示诱捕我，并用 \Quote{做得好，做得好} 的话煽动我，使我陷于此。\BibleQuote{因为凡期待祢的，必不蒙羞} \BibleRef{咏 24:3}。\par

\Quote{愿那些不义地行虚妄之事的人蒙羞。} 愿那些为获取短暂之物而行不义的人蒙羞。\BibleQuote{主，求祢使我知道祢的道路，教导我祢的路径} \BibleRef{咏 24:4}：不是那些宽阔的、引多数人走向毁灭的道路；\BibleRef{玛 7:13} 而是祢的路径，窄狭而少人知晓，求祢教导我。\par

\BibleQuote{以祢的真理引导我：} 避免错误。\BibleQuote{教导我：} 因为靠我自己，我一无所知，只知道虚假。\BibleQuote{因为祢是我救恩的天主；我终日等候祢} \BibleRef{咏 24:5}。因为我被祢逐出乐园，远游到了远方，\BibleRef{路 15:13} 我无法靠自己返回，除非祢来迎回这流浪者；因为我的归返，在这整个世界的时间中，一直等候着祢的怜悯。\par

\BibleQuote{上主，求祢记起祢的怜悯} \BibleRef{咏 24:6}：上主，求祢记起祢仁慈的作为；因为人们以为祢似乎忘记了。\BibleQuote{并且祢的慈爱是永恒的。} 也请记住，祢的慈爱是永恒的。因为祢从未没有慈爱，祢曾使有罪的人也确实屈从于虚妄，但怀着望德；\BibleRef{罗 8:20} 并且没有剥夺他祢创造中如此众多而极大的安慰。\par

\BibleQuote{求祢不要记念我青年和愚昧的过犯} \BibleRef{咏 24:7}：我狂妄自大和愚昧的过犯，求祢不要留着报复，而愿它们被祢遗忘。\Quote{天主，求祢按祢的仁慈记念我。} 求祢确实记念我，不是按我所应得的忿怒，而是按祢相称的仁慈。\Quote{为了祢的良善，上主。} 不是为了我的功劳，而是为了祢的良善，上主。\par

\BibleQuote{上主是仁慈正直的} \BibleRef{咏 24:8}：上主是仁慈的，因为祂甚至如此怜悯罪人和不敬虔的人，以致饶恕了一切过去的事；但上主也是正直的，在召叫和赦免的仁慈之后（这仁慈是恩宠，不凭功德），祂将要求与最后审判相称的功绩。\Quote{因此，祂将为那些在道路上失败的人制定法律。} 因为祂首先赐下仁慈，以将他们带入道路。\par

\Quote{祂将引导温良者在审判中。} 祂将引导温良者，并将在审判中不使那些追随祂旨意、且不因抗拒而偏爱自己意愿的人蒙羞。\BibleQuote{祂将教导温顺的人祂的道路} \BibleRef{咏 24:9}。祂将教导祂的道路，不是给那些想要奔跑在前的人，好像他们更能管理自己；而是给那些当轻省的轭和容易的担子放在他们身上时，不昂首也不抬脚跟的人。\BibleRef{玛 11:30}\par

\BibleQuote{上主的一切道路都是仁慈和真理} \BibleRef{咏 24:10}：祂将教导他们什么道路呢？岂不是仁慈——在其中祂是宽恕的，以及真理——在其中祂是不朽的？祂已显出仁慈于赦罪，显出真理于审判功过。因此，\Quote{上主的一切道路}是指天主圣子的两次来临：一次在仁慈中，一次在审判中。那么，那看见自己毫无功绩却被释放，放下骄傲，并因体验过祂帮助的宽厚而从此警惕祂审判严厉的人，便持守祂的道路而到达祂。\Quote{给那些寻求祂的盟约和祂的见证的人。} 因为他们明白主在第一次来临时是仁慈的，在第二次来临时是审判者，他们以温良和柔和寻求祂的盟约——当祂用自己的血救赎我们获得新生时；并在先知和福音作者中寻求祂的见证。\par

\BibleQuote{上主，为了祢圣名的缘故，求祢垂怜我的罪，因为它是多重的} \BibleRef{咏 24:11}。祢不仅赦免了我信主之前所犯的罪；而且对于我多重的罪（因为即使在道路上也不乏失足），祢将因痛悔精神的祭献而变得仁慈。\par

\Quote{谁是那敬畏上主的人？} 从这种敬畏，他开始走向智慧。\BibleQuote{祂必为他在他所选择的道路上制定法律} \BibleRef{咏 24:12}。祂必为他在他所自由选择的道路上制定法律，使他现在不能逍遥犯罪。\par

\BibleQuote{他的灵魂必居于福乐中，他的后裔将继承土地} \BibleRef{咏 24:13}。他的工作将拥有更新身体的稳定继承。\par

\BibleQuote{上主是敬畏祂之人的依靠} \BibleRef{咏 24:14}：恐惧似乎属于软弱的人，但上主是敬畏祂之人的依靠。上主的名，在全世界被荣耀，也是敬畏祂之人的依靠。\Quote{并使祂的盟约向他们显明。} 祂使祂的盟约向他们显明，因为外邦人和地极都是基督的产业。\par

\BibleQuote{我的眼睛时常仰望上主，因为祂必把我的脚从罗网中拉出} \BibleRef{咏 24:15}。我也不畏惧地上的危险，只要我不注目地上：因为我所仰望的那一位，必把我的脚从罗网中拉出。\par

16. \Quote{求祢垂顾我，怜悯我，因为我孤独而又贫苦}\BibleRef{咏 24:16}。因为我是一个单一的百姓，持守着祢独一教会的卑微，这教会没有分裂或异端占据。\par

17. \Quote{我心中的困苦增多}\BibleRef{咏 24:17}。因不法之事增多，许多人的爱情冷淡，我心中的困苦也增多了。——\BibleRef{玛 24:12} \Quote{求祢救我脱离我的急难。}因为我必须忍受这事，藉着恒心忍耐到底而得救，求祢救我脱离我的急难。\par

18. \Quote{求祢看我的卑微和劳苦}\BibleRef{咏 24:18}。求祢看我的卑微，就是我从不因自夸义而断绝与教会的合一；也看我的劳苦，就是我在忍受那些混在我中间的不守规矩之人。\Quote{并赦免我的一切罪过。}并因这些祭献而蒙宽恕，赦免我的一切罪过，不仅是我在信主之前的青年时期和无知的罪，也包括我现在凭信德生活时，因软弱或此世黑暗而犯的罪。\par

19. \Quote{求祢看我的仇敌，他们何其众多}\BibleRef{咏 24:19}。因为不仅在外，甚至在内，在教会的共融中，他们也不缺乏。\Quote{并以不义的仇恨憎恨我。}他们恨我，而我是爱他们的。\par

20. \Quote{求祢守护我的灵魂，并拯救我。}求祢守护我的灵魂，使我不至转向模仿他们；并把我从他们与我混杂的混乱中拉出来。\Quote{愿我不致蒙羞，因为我投靠了祢}\BibleRef{咏 24:20}。愿我不致蒙羞，即便他们起来攻击我；因为我不是投靠自己，而是投靠了祢。\par

21. \Quote{纯朴与正直的人依附了我，因为我等候了祢，上主}\BibleRef{咏 24:21}。纯朴与正直的人，不仅像恶人那样在身体上与我混杂，而且以同心的纯朴与正直依附于我；因为我没有堕落去模仿恶人，却等候了祢，期待着祢末次收割时的簸扬。\par

22. \Quote{天主，求祢拯救以色列脱离他的一切困苦}\BibleRef{咏 24:22}。\Quote{天主，求祢拯救祢的百姓，}就是祢所预备来朝见祢的百姓，脱离他的困苦，不仅是他忍受在外的，也包括他在内所忍受的。\par

\SourceRecord{Translated by an anonymous scholar.；Nicene and Post-Nicene Fathers, First Series；Vol. 8.；Edited by Philip Schaff.；Buffalo, NY: Christian Literature Publishing Co.,；1888.；<http://www.newadvent.org/fathers/1801025.htm>.}

% structure: p0001 source=h1 target=section reason=page-title

\section{《圣咏集》第26篇释义}

\Person{达味}所作。\par

这可以归于\Person{达味}本人，但并非中保、\Person{耶稣基督}其人，而是指如今在基督内圆满建立的整个教会。\par

2. \BibleQuote{上主，请审判我，因为我行走在我的纯洁中} \BibleRef{咏 25:1}。上主，请审判我，因为在祢首先向我显示的慈悲之后，我因我的纯洁而有了一些功德，我持守了它的道路。\BibleQuote{并且信赖上主，我将不会被摇动。} 然而，即使如此，不是信赖我自己，而是信赖上主，我将住在祂内。\par

3. \BibleQuote{上主，请考验我，且试炼我} \BibleRef{咏 25:2}。然而，免得我任何隐秘的罪对我隐藏，上主，请考验我，且试炼我，使我知道，不是让祢（无所不知者）知道，而是让我自己和人知道。\BibleQuote{焚烧我的肾脏和我的心。} 将一种治疗性的净化，如同火一般，应用在我的快乐和思想上。\BibleQuote{因为祢的仁慈在我眼前} \BibleRef{咏 25:3}。因为，为使我不被那火吞噬，不是我的功德，而是祢的仁慈——祢借此引领我进入如此的生命——在我眼前。\BibleQuote{并且我因祢的真理而蒙祢悦纳。} 既然我自己的虚假令我不悦，而祢的真理令我喜悦，我也因它并在它内蒙悦纳。\par

4. \BibleQuote{我没有与虚妄的议会同坐} \BibleRef{咏 25:4}。我没有选择将我的心交给那些试图提供不可能之事的人——他们如何在享受短暂事物中蒙福。\BibleQuote{我也不会与作恶的人一同进入。} 而既然这正是所有邪恶的原因，因此我不会与作恶的人一同隐藏我的良心。\par

5. \BibleQuote{我憎恨作恶者的会众。} 但为了到达这个虚妄的议会，形成了作恶者的会众，这些我憎恨了。\BibleQuote{我也不会与不虔敬的人同坐} \BibleRef{咏 25:5}。因此，与这样的议会，与不虔敬的人，我不会同坐，也就是说，我不会给予我的同意。\BibleQuote{我也不会与不虔敬的人同坐。}\par

6. \BibleQuote{我要在无辜者中间洗净我的手} \BibleRef{咏 25:6}。我要在无辜者中间洁净我的作为；在无辜者中间我要洗净我的手，我将用这些手拥抱祢荣耀的恩赐。\BibleQuote{上主，我要环绕祢的祭台。}\par

7. \BibleQuote{为使我听到祢赞美的声音。} 为使我学习如何赞美祢。\BibleQuote{并使我宣扬祢一切奇妙的作为} \BibleRef{咏 25:7}。在我学习之后，我要陈明祢一切奇妙的作为。\par

8. \BibleQuote{上主，我爱慕祢圣殿的华美：} 即祢的教会。\BibleQuote{以及祢荣耀的居所} \BibleRef{咏 25:8}：祢居住并被荣耀的地方。\par

9. \BibleQuote{求祢不要将我的灵魂与不虔敬的人一同消灭} \BibleRef{咏 25:9}。那么，求祢不要将我的灵魂——它爱慕祢圣殿的华美——与那些憎恨祢的人一同消灭。\BibleQuote{也不要将我的生命与流血的人一同消灭。} 并与那些憎恨邻人的人一同消灭。因为祢的圣殿因两条诫命而华美。\par

10. \BibleQuote{在他们的手中是邪恶。} 那么，不要将我与不虔敬的人和流血的人一同消灭，他们的作为是邪恶的。\BibleQuote{他们的右手充满了礼物} \BibleRef{咏 25:10}。并且那赐予他们以获得永恒救恩的事物，他们转变成接受此世的礼物，\BibleQuote{以为敬虔是一种交易} \BibleRef{弟前 6:5}。\par

11. \BibleQuote{但我行走在我的纯洁中：求祢拯救我，并怜悯我} \BibleRef{咏 25:11}。让我主之血的如此伟大的代价为我成就完全的拯救；在这生命的危险中，求祢的仁慈不要离开我。\par

12. \BibleQuote{我的脚站在正直中。} 我的爱没有从祢的义中退缩。\BibleQuote{上主，在众教会中我要赞美祢} \BibleRef{咏 25:12}。上主，我不会向祢所召叫的人隐藏祢的祝福；因为紧接在爱祢之后，我加上爱我的邻人。\par

\SourceRecord{Translated by an anonymous scholar.；Nicene and Post-Nicene Fathers, First Series；Vol. 8.；Edited by Philip Schaff.；Buffalo, NY: Christian Literature Publishing Co.,；1888.；<http://www.newadvent.org/fathers/1801026.htm>.}

% structure: p0001 source=h1 target=section reason=page-title

\section{\WorkTitle{咏第二十七篇释义}}

\Emph{达味所作，在他受傅之前。}\par

1. 基督的年轻士兵在来到信德时发言说：\BibleQuote{上主是我的光明，我的救恩：我还怕谁呢？}\BibleRef{咏 26:1}。上主将赐给我对祂的认识和救恩：谁能把我从祂那里夺去？\BibleQuote{上主是我生命的保护者：我畏惧谁呢？}上主将击退我仇敌的一切攻击和陷阱：我不怕任何人。\par

2. \BibleQuote{当罪人走近我，要吃我的肉}\BibleRef{咏 26:2}。当罪人走近，要认识并侮辱我，好在我改善之后高抬自己；用他们辱骂的牙齿吞食的并非我，而是我的肉欲。\BibleQuote{我的仇敌，他们困扰我。}不仅是那些以友好意图责备我、想使我回心转意而困扰我的人，而且也是我的仇敌。\BibleQuote{他们软弱，跌倒了。}当他们这样做，为维护自己的意见时，他们就变得软弱，无法相信更好的事，并开始憎恨救恩的话语——正是这话语使我去做令他们不悦的事。\par

3. \BibleQuote{若军营列阵攻击我，我的心不畏惧。}但若反对者的众多人合谋列阵攻击我，我的心不畏惧，以致倒向他们。\BibleQuote{若战争起来攻击我，在此我仍要信赖}\BibleRef{咏 26:3}。若这世界的迫害起来攻击我，我要将希望寄托在这恳求中，就是我正在思想的恳求。\par

4. \BibleQuote{有一件事我求了上主，我要寻求这事。}因为我向上主求了一件事，我要寻求这事。\BibleQuote{就是愿我一生的岁月，常居住在上主的殿里}\BibleRef{咏 26:4}。就是只要我活在此生中，任何逆境都不能使我脱离那些在全世界持守上主的信仰之合一与真理者的行列。\BibleQuote{为能瞻仰上主的甘饴。}以此为目的，即坚忍于信德中，那甘饴的异象将显现给我，使我得以面对面地瞻仰。\BibleQuote{并受保护，成为祂的圣殿。}死亡被胜利吞灭后，我将穿上不朽，成为祂的圣殿。\par

5. \BibleQuote{因为在我遇祸之日，祂将我藏在祂的帐幕中}\BibleRef{咏 26:5}。因为在我遇祸之日，祂将我藏在祂降生成人的圣言的安排中，这祸患乃我的必死生命所暴露的试炼。\BibleQuote{祂将我藏在祂帐幕的隐密处。}祂保护了我，使我的心相信以得成义。\par

6. \BibleQuote{祂将我高举在磐石上。}并且为使我所信的彰显出来而得救恩，祂使我的宣认在祂自己的力量中显明。\BibleQuote{如今，看，祂已使我抬头高于我周围的仇敌}\BibleRef{咏 26:6}。当现在身体因罪而死亡时，祂为我保留了什么呢？看，我感到我的心思事奉天主的法律，并未被叛逆的罪恶法律所俘虏？\BibleQuote{我环绕而行，在祂的帐幕中献了欢庆的祭献。}我思考了世界巡回，相信基督；并且因为在这世界中，天主为我们而在时间中谦卑，我以欢庆赞美了祂：因为这样的祭献令祂喜悦。\BibleQuote{我要歌唱，赞美上主。}我要在心与行中因上主喜乐。\par

7. \BibleQuote{上主，请听我向你呼号的声音}\BibleRef{咏 26:7}。上主，请听我内在的声音，这声音以强烈的意向达到你的耳中。\BibleQuote{求你怜悯我，应允我。}求你怜悯我，并在此事中应允我。\par

8. \BibleQuote{我的心对你说：我曾寻求你的仪容}\BibleRef{咏 26:8}。因为我没有将自己显示给人；而是在隐秘中，惟有你听见之处，我的心对你说：我没有从你寻求任何除你之外的赏报，而是你的仪容。\BibleQuote{上主，我要求你的仪容。}我要恒久不懈地追求这寻求：因为不求寻常之物，而是你的仪容，上主，我要寻求，好能白白地爱你，因为我找不到更宝贵之物。\par

9. \BibleQuote{求你不要转面不顾我}\BibleRef{咏 26:9}：好使我找到我所寻求的。\BibleQuote{不要向你的仆人发怒而转开：}免得我在寻求你时，偶遇别的什么。因为对于爱慕并寻求你仪容真理的人，还有什么比这惩罚更痛苦呢？\BibleQuote{求你作我的助佑。}若你不助佑，我怎能找到？\BibleQuote{不要弃舍我，也不要轻视我，我的救主天主。}不要耻笑一个必死的人敢于寻求永恒者；因为你，天主，医治我罪的创伤。\par

10. \BibleQuote{因为我的父母抛弃了我}\BibleRef{咏 26:10}。因为这世界和这世俗之城的王国，我生于时间与必死之中，它们在我寻求你、并轻视它们所应许之物时抛弃了我，因为它们不能给予我所寻求的。\BibleQuote{但上主却收留了我。}但上主，那能将自己赐给我的，却收留了我。\par

11. \BibleQuote{上主，求你在你的道路上为我制定法律}\BibleRef{咏 26:11}。至于我，正要启程走向你，并宣告如此伟大的誓愿——从畏惧到达智慧，求你在你的道路上为我制定法律，免得我迷失时你的规则离弃我。\BibleQuote{并因我的仇敌，引导我走正路。}并因我的仇敌，引导我走其狭窄的正路。因为开始还不够，因为仇敌直到达到终点才停止。\par

12. \BibleQuote{不要将我交于困扰我者的灵魂}\BibleRef{咏 26:12}。不要容忍那些困扰我的人因我的恶而满足。\BibleQuote{因为不义的证人起来攻击我。}因为那些诬陷我的人起来攻击我，要使我离开你、唤回我，好似我寻求人的光荣。\BibleQuote{不义向自己说谎。}因此不义以其自己的谎言为乐。它并未动摇我，因为为此有更大的赏报在天上应许给我。\par

13. \BibleQuote{我相信在活人之地得见主的恩惠} \BibleRef{咏 26:13}。既然我的主先经历了这些，如果我也轻视垂死者的舌头（因为\BibleQuote{撒谎的口杀害灵魂} \BibleRef{智 1:11}），我相信在活人之地得见主的恩惠，那里没有虚伪的容身之地。\par

14. \BibleQuote{等候主，要有勇气；坚固你的心，等候主。} \BibleRef{咏 26:14}。但这要到什么时候呢？对凡人来说是艰难的，对爱人来说是缓慢的；但请听那位不会欺骗的声音说：\BibleQuote{等候主。}要勇敢地忍受腰间的灼热，刚毅地承受内心的煎熬。不要以为你尚未领受的就是被拒绝了。为了不因绝望而气馁，请看它是如何说的：\BibleQuote{等候主。}\par

\SourceRecord{Translated by an anonymous scholar.；Nicene and Post-Nicene Fathers, First Series；Vol. 8.；Edited by Philip Schaff.；Buffalo, NY: Christian Literature Publishing Co.,；1888.；<http://www.newadvent.org/fathers/1801027.htm>.}

% structure: p0001 source=h1 target=section reason=page-title

\section{《圣咏第廿八篇释义》}

属达味本人。\par

1. 这是中保亲自的声音，在苦难的争战中手强有力。祂似乎向敌人所愿的，并非恶意的愿望，而是宣告对他们的惩罚；正如福音中，\BibleRef{玛 11:20}论到那些城市，祂虽在其中行了奇迹，他们却仍不信祂，祂并非以任何邪恶的意愿说祂所说的话，而是预言将要降在他们身上的事。\par

2. \Quote{主啊，我向祢呼号；我的天主，求祢不要向我沉默}\BibleRef{咏 27:1}。\Quote{免得祢向我沉默，我就如同下到深渊的人一样。}因为正是由于祢圣言的永恒不断与我自己结合，我才不是像其余生于这世界深重苦难中的人那样；在那里，好像祢沉默了，祢的圣言不被认识。\Quote{主啊，我呼求的声音，求祢垂听，当我向祢祈祷，当我向祢的圣殿举起我的双手时}\BibleRef{咏 27:2}。当我为他们的得救而被钉十字架时，他们因信成为祢的圣殿。\par

3. \Quote{求祢不要将我的灵魂与罪人一同拉去，也不要与作恶的人一同毁灭我，就是那些与邻舍说和平}\BibleRef{咏 27:3}的人。那些对我说，\Quote{我们知道祢是从天主来的师傅}\BibleRef{若 3:2}\Quote{心里却存着恶意}。但他们心里说的是恶言。\par

4. \Quote{求祢按他们的作为报应他们}\BibleRef{咏 27:4}。求祢按他们的作为报应他们，因为这是公义的。\Quote{按他们恶意的情感。}因为他们图谋邪恶，不能发现善。\Quote{按他们手中的作为，求祢报应他们。}虽然他们所行的可能为他人有益于得救，但求祢按他们意志的作为报应他们。\Quote{给他们应得的报应。}因为他们为所听见的真理，愿用诡诈回报；愿他们自己的诡诈欺骗他们。\par

5. \Quote{因为他们不明白主的作为}\BibleRef{咏 27:5}。从何看出这事临到了他们？显然从这一点，\Quote{因为他们不明白主的作为。}事实上，这甚至已经是他们的报应，即他们怀着恶意试探祂如同一个人，却在祂身上认不出天主以何种计划派遣圣子降生成人。\Quote{也不明白祂手的作为。}也不为那些摆在眼前可见的作为所动。\Quote{祢要毁灭他们，不再建立他们。}愿他们不能加害于我，甚至在他们企图兴造机器攻击我的教会时，也一无所成。\par

6. \Quote{上主当受赞美，因祂垂听了我祈祷的声音}\BibleRef{咏 27:6}。\par

7. \Quote{上主是我的助佑，是我的保护者}\BibleRef{咏 27:7}。上主在如此大的苦难中帮助我，并以我的复活用不朽保护我。\Quote{我心依靠了祂，我就得了扶助。}\Quote{我的肉身得以繁荣：}就是说，我的肉身复活了。\Quote{我要以我的意志向祂颂谢。}因此，死亡的恐惧既已消灭，不再是因法律的恐惧所迫，而是以自由意志与法律一同，那些信我的人要向他颂谢；并且因为我住在他们内，我也要颂谢。\par

8. \Quote{上主是自己人民的力量}\BibleRef{咏 27:8}。不是那\Quote{不认识天主的正义，而想建立自己的正义}\BibleRef{罗 10:3}的人民。因为他们不以为自己靠自己能刚强：因为上主是自己人民的力量，他们在今生的艰难中与魔鬼争战。\Quote{也是祂的基督的得救的保护者。}为的是在战争的力量之后，借祂的基督拯救了他们，最后以和平的不朽保护他们。\par

9. \Quote{求祢拯救祢的人民，祝福祢的产业}\BibleRef{咏 27:9}。因此，在我的肉身复活之后，我代祷，因为祢曾说过：\Quote{你向我请求，我就将外邦人赐给你作产业；}\Quote{求祢拯救祢的人民，祝福祢的产业：}因为\Quote{我的一切都是祢的}\BibleRef{若 17:10}。\Quote{求祢牧养他们，并扶助他们直到永远。}求祢在此世牧养他们，并由此提拔他们进入永生。\par

\SourceRecord{Translated by an anonymous scholar.；Nicene and Post-Nicene Fathers, First Series；Vol. 8.；Edited by Philip Schaff.；Buffalo, NY: Christian Literature Publishing Co.,；1888.；<http://www.newadvent.org/fathers/1801028.htm>.}

% structure: p0001 source=h1 target=section reason=page-title

\section{圣咏第二十九篇释义}

达味本人的圣咏，为帐幕的完成。\par

1. 这是中保自己——手中有力量者——的圣咏，论及教会在此世战争中的成全，她在此世按时与魔鬼作战。\par

2. 先知说：\BibleQuote{天主的众子，请将上主的光荣和能力归给上主。}\BibleRef{咏 28:1} 你们这些宗徒所生的人，你们这些由群羊的领袖因福音所生的人，请将上主自己归给上主。\BibleRef{格前 4:15} \BibleQuote{将上主的光荣和威能归给上主。}\BibleRef{咏 28:2} \BibleQuote{将上主之名的光荣归给上主。}愿祂在普世得享荣耀的宣认。\BibleQuote{在祂的圣洁庭院中朝拜上主。}在你们那被扩大并圣化的心中朝拜上主，因为你们就是祂王权的圣洁居所。\par

3. \BibleQuote{上主的声音响在水面上。}\BibleRef{咏 28:3} 基督的声音响在众民之上。\BibleQuote{荣耀的天主发出雷声。}荣耀的天主从肉身的云中可怕地宣讲悔改。\BibleQuote{上主是在大水之上。}主耶稣自己，在向众民发出声音、使他们畏惧之后，就吸引他们归向自己，并住在他们内。\par

4. \BibleQuote{上主的声音具有威能。}\BibleRef{咏 28:4} 上主的声音如今就在他们自己里面，使他们有能力。\BibleQuote{上主的声音具有大能。}上主的声音在他们里面行大事。\par

5. \BibleQuote{上主的声音劈断香柏。}\BibleRef{咏 28:5} 上主的声音在心的破碎中谦卑骄傲者。\BibleQuote{上主必要劈断黎巴嫩的香柏。}上主必借悔改劈碎那些因世俗的尊贵而高傲自大的人，因为祂要\BibleQuote{拣选了世上卑贱的}\BibleRef{格前 1:28}来彰显祂的天主性。\par

6. \BibleQuote{并使祂们如同黎巴嫩的小牛。}\BibleRef{咏 28:6} 及至他们的骄傲被除去后，祂要使他们在模仿祂自己的谦卑中自卑，祂自己如同小牛犊被这世界的尊贵者牵去宰杀。\BibleRef{依 53:7} \BibleQuote{世上的君王起来，首领商议，一起反对上主和祂的受傅者。} \BibleQuote{祂是可爱的，如同独角兽的幼崽。}诚然，祂是可爱的、圣父的独生子，\BibleQuote{祂使自己空虚}\BibleRef{斐 2:7} 了祂的光荣，成为人，如同犹太人的孩子，那些\BibleQuote{不认识天主的正义}\BibleRef{罗 10:3} 并骄傲地夸耀自己的正义为独特的人。\par

7. \BibleQuote{上主的声音切断火焰。}\BibleRef{咏 28:7} 上主的声音无害地穿过所有迫害者的激动热情，或分开迫害者的狂怒，以致有人说：\BibleQuote{莫非这就是基督？}另有人说：\BibleQuote{不然，祂是欺骗民众的。}\BibleRef{若 7:41} 这样祂便切断了他们疯狂的喧嚷，使一些人转向祂的爱，使另一些人留在他们的恶意中。\par

8. \BibleQuote{上主的声音震动旷野。}\BibleRef{咏 28:8} 上主的声音震动那曾\BibleQuote{没有希望、在世上没有天主}\BibleRef{弗 2:12} 的外邦人，使他们归向信仰；那里没有先知，没有天主圣言的宣讲者，如同无人居住之地。\BibleQuote{上主也必震动卡德士的旷野。}那时主将使人完全认识祂圣经中的圣言，这圣言曾被不懂它的犹太人遗弃。\par

9. \BibleQuote{上主的声音使母鹿完美。}\BibleRef{咏 28:9} 因上主的声音首先完美了那些战胜并驱退毒舌之人。\BibleQuote{并要揭示树林。}然后祂要向他们启示神圣书籍的幽暗和奥秘的阴影深处，他们在那里自由放牧。\BibleQuote{在祂的圣殿中，人人在讲论祂的光荣。}在祂的教会中，所有重生而怀有永生希望的人，各因自己从圣神所领受的恩赐而赞美天主。\par

10. \BibleQuote{上主居住于洪水之中。}\BibleRef{咏 28:10} 因此上主首先居住在祂的圣人们所经历的这世界的洪水中，安全地保存在教会中，如同在方舟里。\BibleQuote{上主将作为君王永远坐着。}此后祂要住在他们里面永远为王。\par

11. \BibleQuote{上主将赐力量给祂的百姓。}\BibleRef{咏 28:11} 因为上主必赐力量给祂的百姓，使他们抵抗这世界的风暴和旋风，因为祂没有应许他们在世界上有平安。\BibleRef{若 16:33} \BibleQuote{上主必以平安祝福祂的百姓。}同一位主必祝福祂的百姓，赐予他们在祂自己的平安；因为祂说：\BibleQuote{我将我的平安赐给你们，我将我的平安留给你们。}\BibleRef{若 14:27}\par

\SourceRecord{Translated by an anonymous scholar.；Nicene and Post-Nicene Fathers, First Series；Vol. 8.；Edited by Philip Schaff.；Buffalo, NY: Christian Literature Publishing Co.,；1888.；<http://www.newadvent.org/fathers/1801029.htm>.}

% structure: p0001 source=h1 target=section reason=page-title

\section{论圣咏第三十篇}

\Emph{交与乐官。歌调。} \Emph{为献殿之用。达味作。}\par

1. 交与乐官。这是一篇关于复活、改变、身体更新为不朽状态之喜乐的圣咏，不仅关乎主，也关乎整个教会。因为在前一篇圣咏中，帐幕已经完成，我们在战争期间居其中；但现在这房屋被祝圣，它将在永远的和平中存留。\par

2. 那时，是整全的基督在说话：\Quote{上主，我称扬祢，因为祢提拔了我。}\BibleRef{咏 29:1} 我要赞美祢崇高的尊威，上主，因为祢提拔了我。\Quote{祢没有让我的仇敌因我而喜乐。}那些曾屡次企图以各种迫害压迫我于普世的人，祢没有让他们因我而喜乐。\par

3. \BibleQuote{上主，我的天主，我曾向祢呼号，祢就医治了我。}\BibleRef{咏 29:2} 上主，我的天主，我曾向祢呼号，我不再带着因必朽性而软弱多病的身体。\par

4. \BibleQuote{上主，祢从阴府救回我的灵魂，祢从下坑的人中拯救了我。}\BibleRef{咏 29:3} 祢从深沉的黑暗境况和可朽坏肉体的最低泥沼中拯救了我。\par

5. \Quote{上主的圣徒，请向上主歌咏。}先知预见了这些未来之事，便欢喜地说：\Quote{上主的圣徒，请向上主歌咏，并因记念祂的圣德而称谢。}\BibleRef{咏 29:4} 要向祂称谢，因为祂没有忘记祂所用来圣化你们的圣德，尽管这整个过渡时期都属于你们的愿望。\par

6. \BibleQuote{因为祂的震怒是片刻的。}\BibleRef{咏 29:5} 因为祂为了第一个罪报复了你们，你们为此以死亡偿付。\Quote{而生命却因祂的旨意而延续。}而永恒的生命，你们不能凭自己的力量返回，祂却赐予了，因为祂愿意如此。\Quote{晚间哭泣停留。}晚间开始了，当智慧之光从有罪的人撤离，当他被判受死时：从这晚间起，哭泣停留，只要天主的子民在劳苦和诱惑中，等待主的日子。\Quote{早晨的欢呼。}甚至到早晨，那时将有复活的欢呼，这欢呼在主早晨的复活中已预先闪耀。\par

7. \BibleQuote{但我在富足时说，我永远不会动摇。}\BibleRef{咏 29:6} 但我，那从起初就发言的子民，在我的富足中说——现在不再受任何匮乏之苦——\Quote{我永远不会动摇。}\par

8. \BibleQuote{上主，祢的恩宠使我稳固如山。}\BibleRef{咏 29:7} 但我的富足，上主，并非来自我自己，而是祢的恩宠使我稳固如山，我从这一点得知：\Quote{祢转面不顾我时，我便惊惶。}因为祢有时转面不顾罪人，而当我失去祢智识的光照时，我便惊惶。\par

9. \BibleQuote{上主，我向祢呼号，向我的天主哀祷。}\BibleRef{咏 29:8} 当我想起那困苦和悲惨的时期，仿佛置身其中，我听到祢的\OriginalTerm{首生者}{First-Begotten}，我的元首，正为我而将死的声音，祂说：\Quote{上主，我向祢呼号，向我的天主哀祷。}\par

10. \Quote{什么益处}如此？我流我的血，我下到腐败中？\Quote{灰尘岂能称谢祢？}如果我不立即复活，我的身体变得腐朽，\Quote{灰尘岂能称谢祢？}即那些不虔者的群众，我将通过我的复活使他们成义？\Quote{或宣扬祢的真理？}还是为其余的人得救而宣扬祢的真理？\par

11. 上主垂听了，并怜悯了我，上主成了我的助佑。祂也没有让\BibleQuote{祂的圣者见到腐朽。}\BibleRef{咏 29:10}\par

12. \BibleQuote{祢将我的哀号化为舞蹈。}\BibleRef{咏 29:11} 我，教会，接受了他——从死者中首生者（\BibleRef{默 1:5}），如今在祢房屋的祝圣中说：\Quote{祢将我的哀号化为舞蹈，祢除去了我的麻衣，给我披上喜乐。}祢撕去了我罪恶的遮盖、我必朽的悲伤，并以不朽的喜乐给我束上第一件长袍。\par

13. \BibleQuote{为让我的荣耀歌颂祢，使我不再缄默。}\BibleRef{咏 29:12} 但愿如今，不是我的卑微，而是我的荣耀不再哀悼，而要歌颂祢，因为祢从卑微中高举了我；愿我不再因罪恶的意识、对死亡的恐惧、对审判的恐惧而受刺伤。\Quote{上主，我的天主，我要永远称谢祢。}这就是我的荣耀，上主，我的天主，我要永远称谢祢，因为我本身一无所有，我的一切美善都来自祢，祢是\Quote{天主，万物中的万有。}\par

\SourceRecord{Translated by an anonymous scholar.；Nicene and Post-Nicene Fathers, First Series；Vol. 8.；Edited by Philip Schaff.；Buffalo, NY: Christian Literature Publishing Co.,；1888.；<http://www.newadvent.org/fathers/1801030.htm>.}

% structure: p0001 source=h1 target=section reason=page-title

\section{《\WorkTitle{圣咏第三十一篇讲疏}》}

《终曲，达味自己的一篇诗篇，\OriginalTerm{心醉神迷}{ecstasy}。》\par

1. 终曲，达味自己的一篇诗篇，那位在迫害中手强力的中保。因为加在标题上的\Quote{\OriginalTerm{心醉神迷}{ecstasy}}一词，表示心灵的超拔，这种超拔或由恐惧，或由某种启示而产生。但在这篇诗篇中，首先主要看到的是天主子民因所有外邦人的迫害和普世信仰的失落而惊慌失措。但首先是中保亲自发言；然后，被祂的血所救赎的子民献上感谢；最后，在患难中详细发言，这正属于心醉神迷；而先知本人的角色在接近末尾和末尾处两次插入。\par

2. \Quote{\BibleQuote{上主，我投靠了祢，求祢不要让我永久蒙羞}\BibleRef{咏 30:1}}。上主，我投靠了祢，求祢永远不让我蒙羞，当他们像嘲笑普通人一样嘲笑我的时候。\Quote{\BibleQuote{求祢按祢的公义拯救我，并解救我。}}并求祢按祢的公义从死亡的深渊拯救我，从他们的同伙中解救出来。\par

3. \Quote{\BibleQuote{求祢侧耳听我}\BibleRef{咏 30:2}}。在我卑微时，求祢就近听我。\Quote{\BibleQuote{快快来救我。}}不要拖延到世界终了，像对待所有信我的人一样，将我从罪人中分别出来。\Quote{\BibleQuote{求祢作我保护的天主。}}求祢作我的天主和保护者。\Quote{\BibleQuote{并作我避难的房屋，以拯救我。}}并像一所房屋，我躲在其中便可获救。\par

4. \Quote{\BibleQuote{因为祢是我的力量和避难所}\BibleRef{咏 30:3}}。因为祢是我忍受迫害者的力量，是逃避他们的避难所。\Quote{\BibleQuote{为祢圣名的缘故，祢必引导我，并养育我。}}并且为使我藉着祢被万民认识。我将在一切事上遵行祢的旨意；并且通过逐渐聚集圣人归我，祢将完成我的身体和我完美的身量。\par

5. \Quote{\BibleQuote{祢必使我脱离他们为我暗设的罗网}\BibleRef{咏 30:4}}。祢必使我脱离这些他们为我暗设的网罗。\Quote{\BibleQuote{因为祢是我的保护者。}}\par

6. \Quote{\BibleQuote{我把我灵魂交在祢手中}\BibleRef{咏 30:5}}。我把我的灵魂交托于祢的权能之下，很快就要取回。\Quote{\BibleQuote{祢救赎了我，上主，真理的天主。}}愿那被主苦难所救赎、因他们元首的得荣耀而欢欣的子民也说：\Quote{\BibleQuote{祢救赎了我，上主，真理的天主。}}\par

7. \Quote{\BibleQuote{祢憎恨那些徒然把持虚妄的人}\BibleRef{咏 30:6}}。祢憎恨那些把持世上虚假幸福的人。\Quote{\BibleQuote{而我却投靠了上主。}}\par

8. \Quote{\BibleQuote{我要因祢的仁慈而欢乐踊跃：}}因为这仁慈不会欺骗我。\Quote{\BibleQuote{因为祢垂顾了我的卑微：}}在其中祢以望德使我服从于虚妄之下。\BibleRef{罗 8:20}\Quote{\BibleQuote{祢救了我的灵魂脱离困厄}\BibleRef{咏 30:7}}。祢救了我的灵魂脱离恐惧的困厄，好使它以自由的爱侍奉祢。\par

9. \Quote{\BibleQuote{祢没有把我交在敌人手中}\BibleRef{咏 30:8}}。祢没有把我关闭起来，使我不得开口恢复自由，永远被交在魔鬼的权势下，它以现世的贪欲引诱我，以死亡恐吓我。\Quote{\BibleQuote{祢使我的脚站在宽阔之处。}}因我主复活已知，我的复活也已应许，我的爱脱离了恐惧的窄狭，继续在自由的开阔中行走。\par

10. \Quote{\BibleQuote{上主，求祢怜悯我，因为我困苦}\BibleRef{咏 30:9}}。但这是何等出乎意料的迫害者的残酷，在我心里投下如此恐惧？\Quote{\BibleQuote{上主，求祢怜悯我。}}因为我现在不再为死亡惊慌，而是为折磨和酷刑。\Quote{\BibleQuote{我的眼睛因忿怒而昏花。}}我曾注视祢，求祢不要舍弃我：祢发怒，使我的眼睛昏花。\Quote{\BibleQuote{我的灵魂和我的腹腔。}}因同样的忿怒，我的灵魂受到扰乱，我的记忆——我以此记住我的天主为我受了什么，并应许了我什么——也受到了扰乱。\par

11. \Quote{\BibleQuote{因为我的生命在痛苦中耗尽}\BibleRef{咏 30:10}}。因为我的生命是宣认祢，却因痛苦而耗尽，当敌人说：\Quote{让他们受折磨，直到他们否认祂。}\Quote{\BibleQuote{我的岁月在叹息中。}}我在这个世界度过的时光并没有被死亡夺走，而是存留，并在叹息中耗费。\Quote{\BibleQuote{我的力量因匮乏而衰弱。}}我缺乏这身体的健康，折磨的疼痛临到我；我缺乏身体的解体，死亡却迟迟不来；在这匮乏中，我的信心削弱了。\Quote{\BibleQuote{我的骨骸也动摇。}}我的坚定被动摇了。\par

12. \Quote{\BibleQuote{我成了我所有仇敌的羞辱}\BibleRef{咏 30:11}}。所有恶人都是我的仇敌；然而他们为自己的邪恶受折磨，直到他们招供为止；那么我，我的招供之后不是死亡，而是折磨的疼痛，我超过了他们的羞辱。\Quote{\BibleQuote{对我的邻人更是过分。}}这对那些已经亲近祢、并持守我所持守的信仰的人来说，似乎太过分了。\Quote{\BibleQuote{成了我相识者的恐惧。}}我因这可怕苦难的榜样，甚至在我的相识者中引起了恐惧。\Quote{\BibleQuote{看见我的人都从外面逃离我。}}因为他们不明白我内心无形可见的望德，就逃离我到外面有形可见的事物中。\par

13. \Quote{\BibleQuote{我被人遗忘，好像从心中死去的人}\BibleRef{咏 30:12}}。他们忘记了我，好像我从他们心中死了一样。\Quote{\BibleQuote{我好像一个失落的器皿。}}我自认为在主的一切服务中失落了，活在这世上，一无所获，当众人都害怕与我联合时。\par

14. \Quote{因为我已听见许多住在四周之人的斥责} \BibleRef{咏 31:13}。因为我已听见许多人在我附近，在这世界的旅途中，随从时间的循环，拒绝与我一同返回永恒的家乡，他们斥责我。\Quote{当他们聚集起来攻击我时，他们同谋要夺取我的灵魂。}他们设想了一个计谋，使我的灵魂——它本可藉死亡轻易逃脱他们的权势——竟同意他们，甚至不容我死去。\par

15. \Quote{但是我信赖祢，上主；我说：祢是我的天主} \BibleRef{咏 31:14}。因为祢并未改变，祢应拯救，祢必纠正。\par

16. \Quote{在祢手中} \Quote{是我的命运} \BibleRef{咏 31:15}。在祢的权能中是我的命运。因为我看不到任何功德，使祢从人类普遍的邪恶中特别拣选我获得救恩。虽然在我被拣选的事上祢有某种公正而隐秘的安排，但我——对此隐藏的事——已通过拈阄获得了我的主的衣袍。\BibleRef{若 19:24} \Quote{求祢救我脱离仇敌的手，和迫害我的人。}\par

17. \Quote{求祢使祢的容光照耀祢的仆人} \BibleRef{咏 31:16}。求祢让人们——他们不认为我属于祢——知道祢的容光照耀着我，我事奉祢。\Quote{以祢的仁慈拯救我。}\par

18. \Quote{上主，求祢不要让我蒙羞，因为我呼求了祢} \BibleRef{咏 31:17}。上主，求祢不要让我因那些侮辱我的人而羞愧，因为我呼求了祢。\Quote{愿不虔敬者蒙羞，并坠入阴府。}愿那些呼求石头的人反而蒙羞，并与黑暗同住。\par

19. \Quote{愿说谎的嘴唇变成哑巴} \BibleRef{咏 31:18}。在向万民显明祢在我身上所行的奥迹时，愿那对我编造谎言者的嘴唇因惊愕而哑然。\Quote{他们傲慢地轻视义人，口出恶言。}他们傲慢地轻视被钉十字架的基督，对祂口出恶言。\par

20. \Quote{上主，祢的慈祥何其伟大} \BibleRef{咏 31:19}。先知看见这一切，赞叹祢的慈祥何其丰盛，上主。\Quote{祢为敬畏祢的人所珍藏的。}祢深爱那些祢所纠正的人；但为免他们因安闲而疏忽，祢向他们隐藏了祢慈爱的甘饴——他们以敬畏祢为有益。\Quote{祢为仰望祢的人成全了它。}祢为仰望祢的人成全了这慈祥；因为祢不会从他们收回他们所坚忍持守到底的希望。\Quote{在世人面前。}因为这并不躲过世人的眼目——他们如今不再按亚当生活，而是按人子生活。\Quote{祢将他们隐藏在祢容颜的隐蔽处：}祢将把那座席永远保存在认识祢的隐蔽处，为那些仰望祢的人。\Quote{远离人的困扰。}如此，他们如今不再受人的困扰。\par

21. \Quote{祢在祢的帐幕中保护他们，远离口舌的争辩} \BibleRef{咏 31:20}。但在此世，当恶人的口舌反对他们，说：\InnerQuote{谁从那里来了？}时，祢将在帐幕中保护他们——即对主在时间中为我们所行和所受之事的信德帐幕。\par

22. \Quote{愿上主受赞美，因为祂在环围之城彰显了祂奇妙的仁慈} \BibleRef{咏 31:21}。愿上主受赞美，因为在最严厉的迫害的纠正之后，祂向普世——在人类社会的循环中——彰显了祂奇妙的仁慈。\par

23. \Quote{我在惊惶中说} \BibleRef{咏 31:22}。那时，那民族再次发言说：我在恐惧中说，当外邦人猛烈地攻击我时。\Quote{我从祢眼前被抛弃。}因为祢若垂顾我，就不会让我忍受这些。\Quote{因此，上主，当我向祢呼号时，祢听了我的祈祷声。}因此，祢限制了纠正，并显示我共享祢的眷顾；上主，当我从苦难中高声呼求时，祢听了我的祈祷声。\par

24. \Quote{热爱上主，你们所有祂的圣徒} \BibleRef{咏 31:23}。先知看见这些，再次劝勉说：\Quote{热爱上主，你们所有祂的圣徒；因为上主要求真理。}既然\Quote{义人仅仅得救，那犯罪和不虔敬的人将出现在哪里？} \BibleRef{伯前 4:18} \Quote{祂必报应那些极其骄傲的人。}祂必报应那些即使受挫仍不悔改的人，因为他们极其骄傲。\par

25. \Quote{你们要刚强，使你们的心坚强} \BibleRef{咏 31:24}：坚持行善而不灰心，以便按时收割。\Quote{所有信赖上主的人：}即那些合宜地敬畏并朝拜祂的人——你们要信赖上主。\par

\SourceRecord{Translated by J.E. Tweed.；Nicene and Post-Nicene Fathers, First Series；Vol. 8.；Edited by Philip Schaff.；Buffalo, NY: Christian Literature Publishing Co.,；1888.；<http://www.newadvent.org/fathers/1801031.htm>.}

% structure: p0001 source=h1 target=section reason=page-title

\section{圣咏第卅二篇释义}

达味所作，训诲诗。\par

1. 达味所作，训诲诗；由此使人明白：人得救不是因行为的功劳，而是因天主的恩宠，他承认自己的罪过。\par

2. \Quote{罪恶蒙赦免，过犯得遮掩的人，是有福的}\BibleRef{咏 31:1}：他们的罪过被遗忘。\Quote{主不归咎于他，口中也没有诡诈的人，是有福的}\BibleRef{咏 31:2}：他口中也没有自夸为义的话，尽管他的良心充满罪过。\par

3. \Quote{因为我闭口不言，我的骨头就衰老了：}因为我没有用口\Quote{宣认以得救}\BibleRef{罗 10:10}，我所有的坚固都因软弱而衰老。\Quote{我终日哀号}\BibleRef{咏 31:3}：当我不虔敬、亵渎时，呼喊着反对天主，好像为自己辩护、推诿自己的罪过。\par

4. \Quote{因为你昼夜以你的手重重压我：}因为通过你不断的鞭打惩罚，\Quote{我转侧不安，有刺扎透了我}\BibleRef{咏 31:4}：我因知道自己的痛苦而变得痛苦，被邪恶的良心刺痛。\par

5. \Quote{我承认了我的罪过，没有隐藏我的不义：}这就是说，我没有隐匿我的不义。\Quote{我说：我要向主承认我的不义，反对我自己：}我说，我要承认，不是反对天主（如我不虔敬时呼喊、闭口不言时那样），而是反对我自己，向主承认我的不义。\Quote{你就赦免了我心中的罪恶}\BibleRef{咏 31:5}；在未发之于声前，早已听见心中有认罪之言。\par

6. \Quote{为此，凡虔敬的人应在适当的时候向你祈祷：}凡正义的人都应为这心中的邪恶向你祈祷。因为他们并非凭自己的功德成为圣洁，而是藉着那适当的时候，即基督的来临，他救赎了我们脱离罪恶。\Quote{但纵然洪水泛滥，他们也不会接近他}\BibleRef{咏 31:6}：然而，不要以为当终局突然来临，如同诺厄的日子\BibleRef{玛 24:37}，还有认罪的地方，可以借此接近天主。\par

7. \Quote{你是我的避难所，救我脱离围绕我的压迫：}你是我的避难所，脱离围绕我心的罪过的压迫。\Quote{我的喜乐，求你救我脱离包围我的人}\BibleRef{咏 31:7}：我的喜乐在你内：求你救我脱离罪恶带给我的忧伤。\par

8. \Emph{细拉。}天主的回答：\Quote{我要教导你，指示你当行的路；}在你认罪之后，我要赐你理解，使你不偏离当行的路，免得你愿凭自己行事。\Quote{我要定睛看你}\BibleRef{咏 31:8}；借此我要向你坚固我的爱。\par

9. \Quote{不要像骡马一样无知：}因此想要自己治理。但先知说：\Quote{须用嚼环和辔头勒住他们的口。}天主啊，求你对那些\Quote{不愿亲近你的人}\BibleRef{咏 31:9}，如同人对骡马那样，用鞭打使他们服从你的治理。\par

10. \Quote{罪人的鞭打很多：}那不愿向主认罪、想自己做主的人，受鞭打很多。\Quote{但信赖主的人，慈爱环绕他}\BibleRef{咏 31:10}；但信赖主、服从他治理的人，慈爱必环绕他。\par

11. \Quote{你们义人，应当在主内喜乐欢欣：}你们义人哪，不要在自己内，要在主内喜乐欢欣。\Quote{心地正直的人，你们都要光荣}\BibleRef{咏 31:11}：你们都要光荣他，凡明白应当服从他，以便凌驾于万有之上的人。\par

\SourceRecord{Translated by J.E. Tweed.；Nicene and Post-Nicene Fathers, First Series；Vol. 8.；Edited by Philip Schaff.；Buffalo, NY: Christian Literature Publishing Co.,；1888.；<http://www.newadvent.org/fathers/1801032.htm>.}

% structure: p0001 source=h1 target=section reason=page-title

\section{《圣咏集第三十三篇》释义}

1. \BibleQuote{义人们，请在上主内欢欣：}义人们，不要在自己内欢欣，因为那并不安全；但要欢欣于上主。\BibleQuote{赞美在正直人原是美善。}\BibleRef{咏 32:1}：这些赞美上主的人，是顺从上主的人；否则他们就是扭曲和乖谬的。\par

2. \BibleQuote{弹着竖琴赞美上主：}赞美上主，将你们的身体奉献给祂，作为生活的祭献。\BibleRef{罗 12:1}\BibleQuote{弹着十弦瑟歌咏祂。}\BibleRef{咏 32:2}：让你们的肢体服侍于爱天主和爱近人，其中遵守了三诫和七诫。\par

3. \BibleQuote{你们应向上主唱新歌：}向祂咏唱信德恩宠之歌。\BibleQuote{在欢呼声中向祂弹奏吟咏。}\BibleRef{咏 32:3}：在喜乐中向祂弹奏吟咏。\par

4. \BibleQuote{因为上主的话是正直的：}因为上主的话是正直的，要使你们成为你们靠自己无法成为的。\BibleQuote{祂的一切工作都忠信可靠。}\BibleRef{咏 32:4}：免得有人以为凭行为的功绩达到了信德，其实天主所喜爱的一切行为都是在信德中完成的。\par

5. \BibleQuote{祂喜爱慈悲和公道：}因为祂喜爱慈悲——如今祂首先显示慈悲；也喜爱公道——祂以此要求祂所先显示的。\BibleQuote{上主的慈爱弥漫大地。}\BibleRef{咏 32:5}：普世之中，罪过因上主的慈悲而获得赦免。\par

6. \BibleQuote{藉上主的话，诸天得以建立：}因为义人不是靠他们自己，而是藉上主的话得以坚强。\BibleQuote{藉祂口中的气息，万军得以形成。}\BibleRef{咏 32:6}。他们所有的信德是藉着祂的圣神。\par

7. \BibleQuote{祂将海中的水聚集成堆：}祂将世界上的百姓聚集起来，使他们承认已死的罪过，免得他们因骄傲而随意流散。\BibleQuote{祂将深渊收藏在库房里。}\BibleRef{咏 32:7}：并将祂的奥秘保存在他们里面，作为宝藏。\par

8. \BibleQuote{愿全地都敬畏上主：}愿每一个罪人敬畏，好使他停止犯罪。\BibleQuote{愿世上的居民都畏惧祂。}\BibleRef{咏 32:8}：不要畏惧人的恐吓或任何受造之物，而要畏惧祂。\par

9. \BibleQuote{因为祂一发言，它们就形成：}因为那应当畏惧的事物并非由他人造成，而是祂一发言，它们就形成。\BibleQuote{祂一出命，它们就被创造。}\BibleRef{咏 32:9}：祂藉祂的圣言一命，它们就被创造。\par

10. \BibleQuote{上主使异邦人的计谋落空：}那些不寻求祂的国，而寻求自己国的人。\BibleQuote{祂使万民的计划失效：}那些贪求现世幸福的人。\BibleQuote{并挫败王侯的谋略。}\BibleRef{咏 32:10}：那些寻求统治这般百姓的人。\par

11. \BibleQuote{上主的谋略却永远常存；}然而上主的谋略——祂使那顺服于祂的人有福——永远常存。\BibleQuote{祂心中的思念流传万代。}\BibleRef{咏 32:11}：祂智慧的思念不可变更，延续至万代。\par

12. \BibleQuote{以天主为自己上主的民族是有福的：}有一个民族是有福的，她属于天上之城，只选择了上主为她的天主：\BibleQuote{祂为自己选定的子民。}\BibleRef{咏 32:12}：这子民不是出于自己，而是因天主的恩赐被拣选，以便祂拥有他们，不让他们无人照料而悲惨。\par

13. \BibleQuote{上主从天上俯视，看见所有的人类。}\BibleRef{咏 32:13}。上主从义人的灵魂中，慈悲地注视所有愿意更新生命的人。\par

14. \BibleQuote{从祂所预备的居所：}从祂所取的人性的居所，祂为自己预备的。\BibleQuote{祂注视地上所有的居民。}\BibleRef{咏 32:14}：祂慈悲地注视所有生活在肉身之内的人，为要治理他们。\par

15. \BibleQuote{祂逐个塑造他们的心：}祂以属灵的方式赐予他们的心脏适当的恩赐，使得全身不全是眼，也不全是耳，\BibleRef{格前 12:17}而是一人如此，另一人如彼，得以与基督结合。\BibleQuote{祂明了他们所有的作为。}\BibleRef{咏 32:15}。在祂面前，他们所有的作为都被明了。\par

16. \BibleQuote{君王不因兵力强大而获胜：}那统治自己肉身的人，若过度信赖自己的力量，将不得救。\BibleQuote{巨人也不因力量强大而自救。}\BibleRef{咏 32:16}：那与自身私欲的习惯或与魔鬼及其使者争战的人，若过度信赖自己的勇力，也将不得救。\par

17. \BibleQuote{马匹为得救是虚幻的：}那人受欺骗，以为藉着人可得着在人间获得的救恩，或以为靠自己的勇力之冲动能免于毁灭。\BibleQuote{它虽力大，却不能救人。}\BibleRef{咏 32:17}。\par

18. \BibleQuote{看，上主的眼目垂顾敬畏祂的人：}因为如果你寻求救恩，看，上主的慈爱临于敬畏祂的人。\BibleQuote{眷顾那些期望祂仁慈的人。}\BibleRef{咏 32:18}：那些不期望自己的力量，而期望祂仁慈的人。\par

19. \BibleQuote{为救他们的灵魂免于死亡，并在饥荒中养活他们。}\BibleRef{咏 32:19}。赐予他们圣言的滋养和永恒真理的滋养——他们因信赖自己的力量而失去了这滋养，因此连自己的力量也没有了，因为缺乏正义。\par

20. \BibleQuote{我的灵魂切望上主：}为叫它将来得以享受不朽坏的美味，同时在此停留期间，我的灵魂切望上主。\BibleQuote{因为祂是我们的助佑和盾牌。}\BibleRef{咏 32:20}：当我们努力趋向祂时，祂是我们的助佑；当我们抵抗仇敌时，祂是我们的盾牌。\par

21. \Quote{因为我们的心要因祂欢欣：}因为不在我们自己内，在其中没有祂就有极大需要；而是在祂自己内，我们的心要欢欣。\Quote{我们已信赖了祂的圣名} \BibleRef{咏 32:21}；因此我们信赖我们将来到天主面前，因为对我们这些远离的人，祂借着信德，派遣了祂自己的圣名。\par

22. \Quote{上主，求祢的仁慈临于我们，因为我们寄望于祢} \BibleRef{咏 32:22}：上主，求祢的仁慈临于我们，因为望德不会蒙羞，因为我们寄望于祢。\par

\SourceRecord{Translated by J.E. Tweed.；Nicene and Post-Nicene Fathers, First Series；Vol. 8.；Edited by Philip Schaff.；Buffalo, NY: Christian Literature Publishing Co.,；1888.；<http://www.newadvent.org/fathers/1801033.htm>.}

% structure: p0001 source=h1 target=section reason=page-title

\section{圣咏第34篇释义}

达味的一篇圣咏，当他在\Person{阿彼默肋客}面前改变面容，\Person{阿彼默肋客}就打发他走了，他便离去。\par

1. 因为那里有按亚郎品位的祭献，后来祂以自己的体和血设立了按默基瑟德品位的祭献；因此祂在司祭职上改变了面容，遣散了犹太人的王国，来到了外邦人中。那么，\Quote{祂动了情}是什么意思？祂充满了深情。有什么比我们的主耶稣基督的仁慈更充满深情呢？祂看见我们的软弱，为了将我们从永死中拯救出来，甘愿忍受如此大的伤害和侮辱而经受暂时的死亡。\Quote{祂击鼓：}因为鼓只有将皮张在木头上才能做成；达味击鼓，预示基督将被钉在十字架上。\Quote{祂在城门上击鼓：}什么是\Quote{城门}？不就是我们向基督关闭的心灵吗？祂借着十字架的鼓声打开了必死之人的心灵。\Quote{祂被自己的双手捧着：}怎么\Quote{被自己的双手捧着}？因为当祂把自己的体和血委托给人时，祂拿起那信徒所知道的，放在自己手中；并在某种意义上捧着自己，当祂说：\BibleQuote{这是我的身体。}\BibleRef{玛 26:26}\Quote{祂跌倒在门洞边：}意思是祂谦卑了自己。因为这就是跌倒在我们信仰的开端。门洞就是信仰的开端；教会从这里开始，最终达到目睹：好叫它相信未见之事，当它开始面对面看见时，得以享受它们。圣咏的标题如此；我们已简要听过了；现在让我们听听那动了情并在城门上击鼓者的话语。\par

2. \BibleQuote{我要时时赞美上主，对祂的赞颂常在我口中}\BibleRef{咏 33:1}。这是基督在说话，也让基督徒如此说；因为基督徒是在基督的身体内；因此基督成了人，好使基督徒能成为天使，他说：\BibleQuote{我要时时赞美上主。}我何时\BibleQuote{赞美上主}？当祂祝福你时？当现世的财物丰盈时？当你有大量的谷物、油、酒、金银、仆役和牲畜时？当这必死的健康毫无损伤、安然无恙时？当所有为你所生的人都长大成人，没有因夭折而失去什么，幸福完全统治你的家，所有事物都环绕你丰盈时？那时你赞美上主吗？不，而是\BibleQuote{时时。}因此，既在那时，也当按时间或按上主天主的鞭打，这些事物受到扰乱、被夺去、很少为你而生、或生了又逝去时。因为这些事发生了，随之而来的是贫穷、匮乏、劳苦、痛苦和诱惑。但你，曾唱过\BibleQuote{我要时时赞美上主，对祂的赞颂常在我口中}，当祂赐予时，赞美；当祂拿走时，也赞美。因为赐予的是祂，拿走的也是祂；但祂不会从赞美祂的人那里将自己拿走。\par

3. 然而，除了心中谦卑的人，谁能时时赞美上主呢？因为我们的主用自己的体和血教导了极深的谦卑：因为当祂托付自己的体和血时，祂托付了祂的谦卑，这在历史上所记载的、我们已略过的达味那看似疯狂的事迹中有所体现：\BibleQuote{他的唾沫流到胡须上。}\BibleRef{撒上 21:13}当宗徒被诵读时，你们听到了同样的唾沫，但流到胡须上。也许有人会说：我们听到了什么唾沫？刚才不是读到了宗徒说：\BibleQuote{犹太人要求神迹，希腊人寻求智慧}吗？但现在读的是：\BibleQuote{但我们宣扬，}他说，\BibleQuote{被钉在十字架上的基督}（那时祂击鼓），\BibleQuote{这对犹太人来说是绊脚石，对希腊人来说是愚妄；但对蒙召的，无论是犹太人还是希腊人，基督是天主的德能，天主的智慧。因为天主的愚妄比人智，天主的软弱比人强。}\BibleRef{格前 1:22}因为唾沫象征愚妄；唾沫象征软弱。但如果天主的愚妄比人智，天主的软弱比人强；那么唾沫不要冒犯你，而要注意它流到胡须上：因为如同唾沫象征软弱，胡须象征力量。祂用自己软弱之体遮盖了祂的力量，外表看起来好像唾沫般的软弱；但内在的神性力量如胡须般被遮盖。因此谦卑被托付给我们。你们要谦卑，才能时时赞美上主，使祂的赞颂常在你口中……\par

4. 但为什么人能时时赞美上主？因为他谦卑。什么是谦卑？不将赞美归于自己。愿意自己被赞美的是骄傲；不骄傲的就是谦卑。那么你就不骄傲吗？为了谦卑，请说这里所写的：\BibleQuote{我的心灵因上主而夸耀，谦卑的人听了，必会欢欣。}\BibleRef{咏 33:2}因此，那些不愿因上主被赞美的人，不是谦卑，而是凶暴、粗鲁、高傲、骄傲。上主愿意温顺的牲畜；你应是上主的\OriginalTerm{牲畜}{jumentum}；就是说，你要谦卑。祂骑在你身上，祂管理你；不要害怕你会失足、跌倒：那确实是你的软弱；但要思考谁骑在你身上。你是驴驹，但你背着基督。因为祂也骑着一头驴驹进了城；那牲畜是温顺的。……\BibleQuote{你们不要像骡马一样，无理智。}因为骡马有时扬起脖子，凭着自己的凶暴摔下骑手。它们被嚼环、辔头、鞭打驯服，直到学会顺从，驮着主人。但你，在颚被嚼环勒伤之前，要谦卑，驮着你的主；不要为自己求赞美，愿那骑在你身上的被赞美，你要说：\BibleQuote{我的心灵因上主而夸耀，谦卑的人听了，必会欢欣。}\BibleRef{咏 33:2}……\par

5. 现在接着的是：\BibleQuote{你们当与我一同称扬主} \BibleRef{咏 33:3}。是谁劝勉我们，要我们与他一同称扬主？弟兄们，凡是在基督身体内的，都应当为此努力，使主能与他一同被称扬。因为无论他是谁，他都爱主。他如何爱祂呢？是那样地爱，以致不嫉妒他的同爱者……让那些如此爱天主以致嫉妒他人的人羞愧吧。堕落的人爱一个战车御者，任何爱战车御者或猎人都希望全体人民与他一同爱，并劝勉说：与我一同爱这哑剧，与我一同爱这个或那个可耻之物。他在民众中呼唤，使可耻之物得以与他一同被爱；难道基督徒不也在教会中呼唤，使天主的真理得以与他一同被爱吗？弟兄们，你们要激发你们内的爱，并向你们每一位呼吁，说：\BibleQuote{你们当与我一同称扬主}。愿你们内有那样的热忱。为什么这些事被诵读和解释？如果你们爱天主，就要迅速地将所有与你们相连的人，以及你们家中所有的人，带入对天主的爱；如果基督的身体被你们所爱，也就是说，如果教会的合一被你们所爱，就迅速带他们去享受，并说：\BibleQuote{你们当与我一同称扬主}。\par

6. \BibleQuote{让我们齐声颂扬祂的名}。什么是\Quote{让我们齐声颂扬祂的名}？就是说，在合一中。因为许多抄本如此写：\BibleQuote{请你们与我一同赞扬主}，\BibleQuote{让我们齐声颂扬祂的名}。无论是说\Quote{一同}还是\Quote{在合一中}，都是一回事。因此，要迅速带领那些你能带领的人，通过劝勉、引导、恳求、辩论、说明理由，以柔和与温良。速速带领他们进入爱德；这样，如果他们颂扬主，他们就能在合一中颂扬祂……\par

7. \BibleQuote{我寻求了主，祂就应允了我} \BibleRef{咏 33:4}。主在哪里应允？在内里。祂在哪里赐予？在内里。你在那里祈祷，在那里蒙应允，在那里蒙祝福。你祈祷了，你蒙应允了，你蒙祝福了；而站在你旁边的人不知道：这一切都是在暗中进行的，正如主在福音中所说：\BibleQuote{当你祈祷时，要进入你的内室，关上门，向你在暗中的父祈祷；你的父在暗中看见，必要在明处报答你。} \BibleRef{玛 6:6} 因此，当你进入你的内室时，你进入了你的心。那些进入自己的心而欣喜，并在其中找不到任何邪恶的人是有福的……\par

8. \BibleQuote{我寻求了主，祂就应允了我。} 那么，谁没有被应允，是因为他没有寻求主。圣善的弟兄们，请注意：他没有说，我从主那里寻求了金子，祂就应允了我；我从主那里寻求了长寿，祂就应允了我；我从主那里寻求了这个或那个，祂就应允了我。寻求主以外的任何事物是一回事，寻求主自己则是另一回事。他说：\Quote{我寻求了}主，\Quote{祂就应允了我。} 但是你，当你祈祷说：杀死我的那个仇人，你寻求的不是主，而是，可以说，你使自己成为你仇人的审判官，并使你的天主成为刽子手。你怎么知道他不比你好呢？你所求其死的那个人，也许恰恰在这一点上比你更好，即他不寻求你的死。所以，不要从主那里寻求任何外在的东西，而要寻求主自己；祂就会应允你，而且当你还在说话时，祂会说：\BibleQuote{我在这里，我在这里。} \BibleRef{依 65:24} ……\par

9. 我已说过谁是这个劝勉者，即那不愿独享其所爱之物的爱者，他说：\BibleQuote{你们亲近祂，并要光亮} \BibleRef{咏 33:5}。因为他所说的，正是他自己所证明的。无论是一个属神的人在基督的身体内，还是我们的主耶稣基督自己按肉身而言，元首劝勉祂自己的肢体，他说什么？\BibleQuote{你们亲近祂，并要光亮}。或者更确切地说，某个属神的基督徒邀请我们亲近我们的主耶稣基督自己。但让我们亲近祂并得以光亮；不像犹太人那样亲近祂，以便被黑暗；因为他们亲近祂是为了钉死祂：让我们亲近祂，好领受祂的体血。他们因祂被钉而死而变黑暗；我们因吃喝被钉者而得以光亮。\BibleQuote{你们亲近祂，并要光亮}。看，这话是对外邦人说的。基督在犹太人发狂并观看时被钉；外邦人不在场；看，那些在黑暗中的已经亲近，那些没有看见的得以光亮。外邦人如何亲近？借着信德跟随，借着内心渴慕，借着爱德奔跑。你们的脚就是你们的爱德。要有两只脚，不要瘸腿。什么是你们的两只脚？爱的两条诫命：爱你的天主，并爱你的近人。用这双脚奔向天主，亲近祂，因为祂既劝勉你奔跑，又亲自倾泻祂自己的光，正如他慷慨而神圣地继续说的：\BibleQuote{你们的面容不再羞愧}。他说：\Quote{亲近祂}，\Quote{并要光亮，你们的面容不再羞愧。} 没有面容会羞愧，除非是骄傲人的面容。为什么？因为他想被高举，当他在这世上遭受侮辱、耻辱或不幸，或任何磨难时，他就羞愧了。但你不要害怕，亲近祂，你将不会羞愧……\par

10. 正如先知所证：\BibleQuote{穷人呼号，主俯允了他}\BibleRef{咏 33:6}。祂教导你如何可以得蒙俯允。所以你之所以不被俯允，是因为你富有。恐怕你说：你呼号了，却没有被俯允，请听原因；\Quote{穷人呼号，主俯允了他}。你要像穷人一样呼号，主必俯允。我怎么方能像穷人呼号呢？借着不倚靠你所拥有的，不自恃自己的力量；借着明白你是匮乏的；借着明白只要你没有那使你富足的一位，你就是贫穷的。主如何俯允了他？\Quote{并由他的一切困厄中拯救了他}。主如何从一切困厄中拯救人呢？\BibleQuote{主的天使将派遣环绕敬畏祂的人，并拯救他们}\BibleRef{咏 33:7}。圣经是这样写的，弟兄们，并非如某些坏抄本所写的那样：\Quote{主将派遣祂的天使环绕敬畏祂的人，祂必拯救他们}；而是这样：\BibleQuote{主的天使将派遣环绕敬畏祂的人，并拯救他们}。祂在这里称谁为主的天使，祂将派遣环绕敬畏祂的人，并拯救他们？我们的主耶稣基督自己，在先知书中被称为\OriginalTerm{大谋略的天使}{Angel of the great Counsel}、\OriginalTerm{大谋略的使者}{Messenger of the great Counsel}；先知们就是这样称呼祂的。那么，正是祂，大谋略的天使，即使者，将派遣给敬畏主的人，并拯救他们。所以不要害怕，免得你被隐藏：无论你在何处敬畏了主，那位天使就认识你，祂将派遣来援助你，并拯救你。\par

11. 现在祂要公开谈论同一件圣事：祂如何被自己双手所捧持。\BibleQuote{请你们体验并观看，主是何等美善}\BibleRef{咏 33:8}。圣咏如今岂不敞明自身，并向你展示那看似癫狂与持续疯癫、同一种\OriginalTerm{清醒的醉态}{sober inebriety}——那\Person{达味}在\Person{阿基士}王面前所象征的某种事物？当人对祂说，这是怎么回事？主说：\BibleQuote{除非人吃我的肉并喝我的血，在他内没有生命}\BibleRef{若 6:53}。而在他们心中掌权的\Person{阿基士}，即错误与无知，说了什么？\BibleQuote{这人怎能把自己的肉给我们吃呢？}\BibleRef{若 6:52}。如果你无知，\BibleQuote{请体验并观看，主是何等美善}；但若你不明白，你就是\Person{阿基士}王；\Person{达味}要改变容貌，离开你，舍弃你，离去。\par

12. \Quote{信赖祂的人有福了}。何需详述？凡不信赖主的人，是有祸的。谁是那不信赖主的人？就是信赖自己的人……\par

13. \BibleQuote{敬畏主吧，你们所有祂的圣者，因为敬畏祂的人不会缺乏}\BibleRef{咏 33:9}。因此，许多人因害怕遭受饥饿而不敬畏天主。有人对他们说：不要欺骗；他们说：我怎能养活自己？没有诡诈就没有技艺；没有欺诈就没有事业。但欺诈是天主惩罚的：你要敬畏天主。但如果我敬畏天主，我就没有生活的来源。祂向那战栗怀疑的人承诺富足，生怕如果他敬畏天主，他会失去多余之物。主曾在你看不起祂时喂养了你，岂会因你敬畏祂而抛弃你？要留心，不要说：某人富有，而我贫穷。我敬畏主；他却因不敬畏而获得了多少，而因敬畏我却一无所有！请看下文：\BibleQuote{富足的人缺乏而受饥饿，但寻求主的人不会缺少任何美物}\BibleRef{咏 33:10}。如果你按字面理解，祂似乎欺骗了你，因为你看见许多邪恶的富人在他们的财富中死去，活着时并不变穷；你看见他们变老，甚至在极大的富足中到达生命的终点。你看见他们的葬礼排场极其铺张，那人本身甚至到坟墓都是富有的，死在象牙床上，亲属围着他哭泣；于是你心里想，如果你碰巧知道他所犯的某些罪和恶行：我知道那人做了什么事；看，他老了，死在自己床上，朋友将他送到坟墓，他的葬礼如此隆重；我知道他所做的事；圣经欺骗了我，祂说假话，因为我听见并歌唱：富足的人缺乏而受饥饿。这人何时有过缺乏？他何时受过饥饿？但寻求主的人不会缺少任何美物。我每日起来上教堂，每日屈膝，每日寻求主，却没有一点好处；此人没有寻求主，却在这些好事中死去！这样想着，绊倒的网罗扼住他；因为他在地上寻求可朽的食物，而没有在天上寻求真正的赏报，所以他把头伸进了魔鬼的套索，他的咽喉被紧紧勒住，魔鬼把他牢牢抓住，使他行恶，好让他模仿那些在如此丰足中死去的恶人。\par

14. 因此，不要这样理解。……当你充满属灵财富时，你能贫穷吗？他富有是因为他有象牙床，而你贫穷是因为你心灵的内室充满了如此多的德行珍宝：正义、真理、爱德、信德、忍耐？展开你的财富，如果你拥有它们，并与富人的财富相比较。但某人在市场上找到了价值连城的骡子，买下了它们。如果你能在市场上找到信德可卖，你愿出多少钱买那信德？天主愿意你白白获取，而你却忘恩负义！那些富人因此缺乏，他们缺乏，而且更严重的是，他们缺乏的是食粮。……因为祂说过：\BibleQuote{我是从天上降下的活的食粮}\BibleRef{若 6:51}。又说：\Quote{饥渴慕义的人是有福的，因为他们必得饱足}。\Quote{但寻求主的人不会缺少任何美物：}但什么样的美物，我已经说过了。\par

15. \Quote{孩子们，你们来听从我：我要教导你们敬畏上主} \BibleRef{咏 33:11}。弟兄们，你们以为这是我说的话；你们要以为是达味说的；要以为是一位宗徒说的；不，要以为是我们主耶稣基督亲自说的：\Quote{孩子们，你们来听从我。} 让我们一同听从祂：你们要藉着我们听从祂。因为祂要教导我们；祂——谦卑者，祂——敲击者，祂——感动者，要教导我们……\par

16. \Quote{有谁愿意生命，喜爱看到幸福的日子？} \BibleRef{咏 33:12}。祂发问。你们中间哪一个人不回答：我？你们中间有谁不爱生命，即不愿意生命，不喜欢看到幸福的日子？你们不是每日这样抱怨，这样说：我们要忍受这些事到何时？一天比一天更糟：我们父辈的时代日子更欢乐，日子更好。哦，倘若你们能问那些人，你们的父辈，他们也会向你们抱怨他们自己的日子。我们的父辈曾幸福，我们却悲惨，我们遭遇邪恶的日子：某人统治我们，我们以为他死后会有些许安慰；结果更糟：天主啊，求祢向我们显示好日子！\Quote{有谁愿意生命，喜爱看到幸福的日子？} 让他不要在此寻求好日子。他寻求一件好事，但没有在正确的地方寻求。正如，如果你在一个国家寻找某个义人，而那人并未在那里生活过，就会对你说：你找的是一个好人，一个伟人，继续找吧，但不在这里；你在此处找是枉然，你永远找不到。你寻求好日子，让我们一起寻求，不要在此寻求……读圣经……\par

17. 所以，不要让基督徒抱怨，让他看他跟随谁的脚步：但如果他喜爱好日子，让他听从教导且说这话的祂：\Quote{孩子们，你们来听从我；我要教导你们敬畏上主。} 你愿意什么？生命和好日子。听从，并行动。\Quote{谨守你的口舌，远离邪恶} \BibleRef{咏 33:13}。去做这事。不，一个可怜人说，我不愿谨守口舌远离邪恶，却仍想要生命和好日子。假如你的一个工人对你说：我确实毁坏这葡萄园，但我向你要求我的报酬；你带我来葡萄园修剪整枝，我砍掉了所有有用的木头，我还会截断葡萄藤的主干，使你一无所获，做完这些后，你要付我工钱。你岂不说他疯了？你岂不在他动手拿刀之前就把他赶出家门？这样的人就是那些既作恶、又发假誓、又说亵渎天主的话、又抱怨、又欺诈、又醉酒、又争吵、又行奸淫、又用符咒、又求问占卜者，同时又希望看到好日子的人。对他们说：你们不能作恶而求善报。如果你不义，难道天主也不义吗？那么我该做什么？你渴望什么？我渴望生命，我渴望好日子。\Quote{谨守你的口舌，远离邪恶，你的嘴唇不说诡诈的话，} 即：不欺骗任何人，不对任何人说谎。\par

18. 但什么是\Quote{远离邪恶}？\BibleRef{咏 33:14}。你不伤害任何人、不杀人、不偷盗、不行奸淫、不作恶、不作假见证，这还不够；\Quote{远离邪恶。} 当你远离了，你说：现在我安全了，我做了所有事，我将拥有生命，我将看到好日子。祂不仅说\Quote{远离邪恶}，还说\Quote{并要行善。} 你不掠夺不算什么：要给赤身者衣服穿。如果你没有掠夺，你已远离了恶；但除非你接待陌生人进你家，你就不算行善。所以，要远离邪恶，以致行善。\Quote{寻求和平，并追随它。} 祂没有说：你在这里会有和平；要寻求它，并追随它。我往哪里追随它？它已先行。因为主是我们的和平，祂已复活，并升了天。\Quote{寻求和平，并追随它；} 因为当你也复活时，这必朽的将改变，你将在那里拥抱和平，没有人会打扰你。因为那里有完美的和平，你将不再饥饿……\par

19. \Quote{上主的双眼垂顾义人：} 所以不要害怕；要劳苦；上主的双眼在你身上。\Quote{祂的耳朵倾听他们的祈祷} \BibleRef{咏 33:15}。你还要什么？如果一个大户人家的家主不听仆人的抱怨，那仆人就会诉苦说：我们在这里受何等的苦，却没有人听我们。你岂能对天主说：我受何等的苦，却没有人听我？如果祂听了我的，或许你会说，祂就会拿走我的磨难：我向祂呼号，却仍有磨难。只要你紧守祂的道路，当你受磨难时，祂就垂听你。但祂是医生，你仍有腐烂的部分；你喊叫，但祂仍在割，不收回祂的手，直到祂割够了。因为那听人话而放过伤口和腐烂的医生是残忍的。母亲们如何为孩子的健康在浴池中搓揉他们。小孩子在她们手中岂不哭喊？她们岂是因为不放过、不听他们的哭声而残忍？她们岂不充满慈爱？然而孩子们哭喊，却得不到放过。同样，我们的天主也充满爱德，但正因如此，祂似乎不听，为的是饶恕我们、医治我们，直到永远。\par

20. 或许恶人说：我安全地作恶，因为上主的眼不在我身上：天主关注义人，却不看见我，我无论做什么都安全。圣神立刻看到人的心思，就说：\Quote{但上主的面容敌对作恶的人，为从地上剪除对他们的纪念} \BibleRef{咏 33:16}。\par

21. \BibleQuote{义人呼号，主垂听了他们，从他们的一切患难中拯救了他们}\BibleRef{咏 33:17}。义人是那三位青年；他们在火窑中呼号主，在赞美中火焰冷却了。火焰不能接近或伤害赞美天主的无罪义人，祂拯救他们脱离了烈火。\BibleRef{达 3:28}有人说：看，那些义人确实蒙垂听，正如所记载的：\BibleQuote{义人呼号，主垂听了他们，从他们的一切患难中拯救了他们。}但我呼号了，祂却不拯救我；要么我不是义人，要么我没有做祂命令的事，或者也许祂看不见我。不要怕：只管做祂所命令的；如果祂不拯救你的身体，祂将拯救你的灵魂。因为祂从火中救出了三位青年，难道祂没有从火中救出玛加伯吗？\BibleRef{加下 7:3}前一批人在火焰中歌唱，后一批人在火焰中断气，难道不是这样吗？三位青年的天主，难道不也是玛加伯的天主吗？祂拯救了前者，没有拯救后者。不，祂拯救了二者：但祂如此拯救三位青年，以致连属血肉的人都惭愧；但祂没有这样拯救玛加伯，好使那些迫害他们的人，因为他们以为他们战胜了天主的殉道者，而陷入更大的折磨。祂拯救了伯多禄，当天使来到狱中对他说话：\BibleQuote{起来，出去吧，}\BibleRef{宗 12:7}他的锁链突然脱落，他跟随天使，祂拯救了他。当伯多禄不被从十字架上拯救时，他就失去了义德吗？那时祂没有拯救他吗？那时祂也拯救了他。难道他长寿就使他成为不义的吗？也许祂在最后比在起初更垂听了他，那时祂确实拯救了他脱离一切患难。因为当祂初次拯救他时，他后来遭受了多少事！因为祂最后把他送到一个他不能再受任何邪恶的地方。\par

22. \BibleQuote{主亲近心灵破碎的人，拯救心神谦卑的人}\BibleRef{咏 33:18}。天主是至高的：让基督徒谦卑。如果他愿意至高的天主亲近他，就让他谦卑。一个伟大的奥秘，弟兄们。天主超乎万有：你抬高自己，就触不到祂；你谦卑自己，祂就降临到你身上。\BibleQuote{义人的灾难有许多}\BibleRef{咏 33:19}：祂是否说：因此让基督徒成为义人，因此让他们听我的圣言，好使他们不受任何磨难？祂没有这样应许；而是说：\BibleQuote{义人的灾难有许多。}相反，如果他们不义，他们的磨难较少；如果义人，磨难就多。但在少数的磨难之后，或甚至没有磨难，这些人将进入永久的磨难，永不得救赎；但义人在许多磨难之后，将进入永久的平安，在那里他们将永不受任何邪恶。\BibleQuote{义人的灾难有许多，但主拯救他脱离一切。}\par

23. \BibleQuote{主保全他们的一切骨骸，连一根也不容许折断}\BibleRef{咏 33:20}：这一点，弟兄们，我们不要按血肉来接受。骨骸是信友坚稳的支柱。因为如同在肉体中，骨骸给予我们坚固，同样在基督徒的心中，是信德给予坚固。那么，在信德中的忍耐，就如内在之人的骨骸：这就是那不可折断的。\BibleQuote{主保全他们的一切骨骸，连一根也不容许折断}。如果这是指着我们的主天主耶稣基督说的，\Quote{主保全祂儿子的一切骨骸，连一根也不容许折断}，正如在另一处也预表了祂，当说到应被宰杀的羔羊时，论到它说：\BibleQuote{一根骨头也不可折断}\BibleRef{出 12:46}，那么就在主身上应验了，因为当祂挂在十字架上时，祂在那些人来到十字架之前就断了气，发现祂的身体已经无生命，就不打断祂的腿，好使所记载的应验。\BibleRef{若 19:33}但祂也将这一应许给了其他基督徒：\BibleQuote{主保全他们的一切骨骸，连一根也不容许折断}。因此，弟兄们，如果我们看到任何圣人遭受磨难，并且或许被医生切割，或被某个迫害者肢解，以致他的骨头折断，我们不要说：这个人不是义人，因为主曾向祂的义人应许了这一点，祂说：\BibleQuote{主保全他们的一切骨骸，连一根也不容许折断}。你愿意看见祂说的是另一种骨头，就是我们称之为信德坚固支柱，即在一切磨难中的忍耐和坚忍吗？因为这些是不可折断的骨头。请听，并请你们在我们主的苦难中看到我所说的。主被钉在中间，两旁有两个强盗：一个讥笑，一个相信；一个被定罪，一个被称义；一个在今世和将来的世界都受罚，但主对另一个说：\BibleQuote{我实在告诉你，今天你就要同我一起在乐园里}\BibleRef{路 23:43}；然而那些来的人没有打断主的腿，却打断了强盗的腿：亵渎的强盗的腿被打断，正如相信的强盗的腿被打断。那么，那所说的\BibleQuote{主保全他们的一切骨骸，连一根也不容许折断}又在哪里呢？看，主对他说：\BibleQuote{今天你就要同我一起在乐园里}，祂能保全他所有的骨骸吗？主回答你：是的，我保全了它们：因为他的信德那坚固的支柱没有被那些打断他腿的打击所折断。\par

24. \BibleQuote{罪人的死是极恶的} \BibleRef{咏 33:21}。各位弟兄，请留意我所说的。主实在伟大，祂的仁慈也实在伟大；那赐给我们吃祂的身体——在其中祂曾受如此大的苦——并喝祂的血的那一位，何其伟大！祂如何看待那些心怀恶意并说：\Quote{某人死得不好，被野兽吞噬了；他不是义人，因此他死得不好；否则他不会死。}的人呢？那么，死在自己家中、自己床上的人就是义人吗？这就是我所惊奇的，因为我认得这人的罪过和恶行，然而他死得好：在自己家中，在自己门内，没有旅途的损伤，甚至没有在年老时。请听：\BibleQuote{罪人的死是极恶的。}在你看来是善终的，若你能看见内在，便是极恶。你看他在床上安卧，你能看见他被带到地狱里吗？请听，各位弟兄，并从福音中学习什么是罪人的\Quote{极恶之死}。在那个时代，不是有两个人吗？一个富人，身穿紫红袍和细麻衣，天天奢华宴乐；另一个穷人，满身疮痍，躺卧在富人的门口，狗来舔他的疮痍，他渴望从富人桌上掉下来的碎屑充饥。后来穷人死了（那穷人是义人），被天使带到\Person{亚巴郎}的怀抱里。那曾见他尸身躺在富人门口、无人埋葬的人，或许会说：\Quote{愿我的仇敌也这样死；凡迫害我的人，我愿看他这样死。}他的身体被人唾弃，伤口发臭；然而他在\Person{亚巴郎}的怀抱里安息。\BibleRef{路 16:19}若我们是基督徒，就让我们相信；若我们不信，各位弟兄，就不要假装自己是基督徒。信德引领我们至终点。主如此说了，事情便如此。难道星象家对你说的是真的，基督说的却是假的吗？但那个富人是怎么死的呢？穿紫红袍和细麻衣的人会怎么死呢？多么奢华，多么排场！那是怎样的葬礼！那尸体用多少香料埋葬！然而当他在地狱里受痛苦时，却从那被轻视的穷人的手指上，渴望一滴水滴在他燃烧的舌头上，却得不到。那么，请学习\BibleQuote{罪人的死是极恶的}是什么意思；不要再问那些铺着华贵衣物的床，用许多贵重物品包裹的肉体，朋友假装哀哭，家人捶胸，抬棺时前呼后拥的人群，大理石和金饰的墓碑。因为如果你问那些事，它们会回答你虚假的事：许多不是轻罪而是全然邪恶之人的死是最好的，他们配得如此哀悼、涂抹香料、包裹、抬出、埋葬。但请问福音，它会向你的信德显示那富人的灵魂在痛苦中燃烧，所有那些活人的虚荣给予他尸体的尊荣和丧礼，对他毫无益处。\par

25. 但因罪人有多种，而成为无罪者是困难的，或许在此生不可能，他立刻接着指明哪一类的罪人的死是极恶的。\BibleQuote{凡恨那义者的人}（他说）\BibleQuote{必要灭亡。}那义者是谁？不就是\BibleQuote{那使不虔敬者成义的一位}吗？\BibleRef{罗 4:5}是谁？不就是我们的主耶稣基督，祂也\BibleQuote{为我们作了赎罪祭}吗？\BibleRef{若一 2:2}那么，凡恨祂的人，便有了极恶的死；因为凡不藉着祂与我们的天主和好的人，都死在他们的罪恶中。\BibleQuote{因为主赎回祂仆人的灵魂。}但死的好坏应根据灵魂来理解，而非根据肉体可见的侮辱或荣誉。\BibleQuote{凡信赖祂的人，没有一个会灭亡。}\BibleRef{咏 33:22}；这就是人类成义的方式：此有死的人生，无论怎样进步，因不能无罪，在此不会灭亡，只要它信赖祂，在祂内有罪过的赦免。阿们。\par

\SourceRecord{Translated by J.E. Tweed.；Nicene and Post-Nicene Fathers, First Series；Vol. 8.；Edited by Philip Schaff.；Buffalo, NY: Christian Literature Publishing Co.,；1888.；<http://www.newadvent.org/fathers/1801034.htm>.}

% structure: p0001 source=h1 target=section reason=page-title

\section{对圣咏第三十五篇的诠释}

1. ……这篇圣咏的标题没有使我们耽搁，因为它既简短，又不难理解，尤其是在天主的教会中受培育的人。事实如此：\Quote{致达味自己。}因此这篇圣咏是致达味自己的：而\Person{达味}被解释为\Quote{手中有力的}或\Quote{可欲的}。因此这篇圣咏是致那位手中有力、可欲者，致那位为我们战胜了死亡、为我们应许了生命者：因为祂手中有力，在于祂为我们战胜了死亡；祂可欲，在于祂为我们应许了永生。有什么手比那只触摸过棺架、使死者复活的手更强呢？\BibleRef{路 7:14} 有什么手比那只战胜了世界的手更强呢——那只手不是以钢武装，而是被木头刺穿？或者有什么比祂更可欲呢？殉道者们虽未看见祂，却情愿死去，以便能堪当前往祂那里。因此这篇圣咏是致祂的：让我们的心、我们的舌向祂相称地歌唱——如果祂自己肯赐予一些可以歌唱的内容。……\par

2. \BibleQuote{主，求祢审判那些伤害我的人，求祢攻击那些攻击我的人。}\BibleRef{咏 34:1} \BibleQuote{如果天主与我们同在，谁能反对我们呢？}\BibleRef{罗 8:31} 天主怎能为我们做这事？\BibleQuote{拿起武器和盾牌，起来帮助我。}\BibleRef{咏 34:2} 看到天主为你武装，是多么壮观的景象。祂的盾牌是什么？祂的武器是什么？\Quote{主，}在另一处，这位在此说话的人说，\Quote{祢以祢善意的盾牌环绕了我们。}但祂的武器——藉此祂不仅能保护我们，也能打击祂的敌人——如果我们很好地得益，我们自身将成为这些武器。因为正如我们由祂获得武装，祂也由我们而获得武装。但祂由祂所造者而武装，我们则以我们从造物主所领受的事物而武装。这些我们的武器，宗徒在某处称之为\BibleQuote{信德的盾牌、救恩的头盔、圣神的宝剑，即天主的圣言。}\BibleRef{弗 6:16} 祂用你刚才听说的那些武器武装了我们，那些武器令人赞叹、不可征服、不可战胜、闪耀光辉；真正属神且不可见，因为我们也要与不可见的敌人作战。如果你看见你的敌人，让你的武器被看见。我们以信德武装，信德针对那些我们看不见的事物，我们推翻了那些我们看不见的敌人。……\par

3. \BibleQuote{倾出武器，阻止那迫害我者的道路。}\BibleRef{咏 34:3} 那些迫害你的人是谁？或许是你的邻人，或是你得罪过的人，或是你亏负过的人，或是想夺走你东西的人，或是你向他宣讲真理的人，或是你责备其罪恶的人，或是因你生活良好而冒犯了生活恶劣之人。确实，这些也是我们的仇敌，它们迫害我们：但我们被教导要认识另外的仇敌，就是那些我们与之无形作战的，宗徒警告我们说：\BibleQuote{我们不是对抗血肉之躯，}\BibleRef{弗 6:12} 也就是说，不是对抗人；不是对抗那些你看见的，而是对抗那些你看不见的；\BibleQuote{对抗率领者，对抗掌权者，对抗这个黑暗世界的统治者。}……\BibleQuote{整个世界都卧于邪恶中；}因此宗徒解释了他们是什么世界的统治者，他说：\Quote{这个黑暗的。}我说，这个世界的统治者，就是这些黑暗的统治者。……\par

4. \BibleQuote{愿那寻索我灵魂的人，蒙羞受辱。}\BibleRef{咏 34:4} 因为他们寻索它，是为了毁灭它。我但愿他们为善而寻索它！因为在另一篇圣咏中，他指责人们中间没有寻索他灵魂的人：\BibleQuote{我失去了避难所：没有寻索我灵魂的人。}这是谁说的\Quote{没有寻索我灵魂的人}？难道是祂吗？很久以前就预言了祂：\BibleQuote{他们刺穿了我的手和我的脚，数算了我的所有骨头，他们注视着我，看着我，他们分了我的衣服，为我的长衣拈阄。}如今这一切都在他们眼前发生了，却没有寻索祂灵魂的人……\par

5. ……许多人蒙受了羞辱，却得到了益处；许多人受辱之后，以虔诚的敬礼从迫害基督转为与祂肢体结合；若他们没有蒙羞受辱，就不会发生这事。因此他为他们祝福……让他们不要走在前面，而要跟随；让他们不要给人建议，而要接受建议。因为当主谈论祂未来的苦难时，\Person{伯多禄}想要走在主前面：他仿佛要给主健康方面的建议。病人给救主健康方面的建议！他对主肯定祂未来的苦难说了什么？\BibleQuote{主，请远祢。求祢怜悯自己。这事绝不会临到祢身上。}他想走在前面，好让主跟随；但主说了什么？\BibleQuote{撒殚，退到我后面去！}\BibleRef{玛 16:22} 走在前面，你就是撒殚；跟随，你就会成为门徒。对这些人也说同样的话：\BibleQuote{愿那些图谋恶事攻击我的人，转身退后，蒙受羞辱。}\par

6. 其他人呢？因为并非所有的人都如此被征服而悔改相信：许多人继续顽固执拗，许多人在心中保留着前行的精神，即使不发挥出来，却仍为此劳苦，一有机会就把它显露出来。关于这样的人，接下来如何呢？\BibleQuote{让他们像风前的尘埃}\BibleRef{咏 34:5}。\Quote{不虔敬者却不如此，绝不如此；他们像被风吹离地面的尘埃。}风是诱惑，尘埃是不虔敬者。当诱惑来到，尘埃被扬起，它既不能站立，也不能抵抗。\Quote{让他们像风前的尘埃，愿上主的天使使他们受困扰。}\BibleQuote{愿他们的道路变得黑暗而又滑溜}\BibleRef{咏 34:6}。这是一条可怕的道路！独自黑暗谁不害怕？独自滑路谁不避开？在一条黑暗而滑溜的路上，你如何行走？哪里落脚？这两种祸患是人的重大惩罚：黑暗是愚昧，滑路是纵欲。\Quote{愿上主的天使追赶他们，}使他们不能站立。因为任何一个在黑暗滑路上的人，当他看到自己一动脚就要跌倒，脚前又没有光亮时，或许会决心等到光明来到；但这里上主的天使正追赶他们。他预言这些事将要临到他们，并非他愿意它们发生。虽然先知在天主圣神内如此说这些话，正如天主也如此行，以确定的审判，以善良、正义、圣洁、恬静的审判，不是出于愤怒，不是出于苦毒的嫉妒，也不是出于报复仇敌的欲望，而是出于惩罚邪恶的正义；然而，这仍是一个预言。\par

7. 但是为何这些如此大的灾祸？凭什么报应？听凭什么报应。\BibleQuote{因为他们无缘无故为我暗设了网罗的败坏}\BibleRef{咏 34:7}。因为，请观察，犹太人对我们的元首就这样做了：他们暗设了网罗的败坏。他们为谁暗设罗网？为那看透暗设者心肠的那一位。然而，祂在他们中间如同一无所知的人，仿佛受了欺骗，而他们正是在那件事上受了欺骗，就是他们认为祂被欺骗了。因为祂所以好像受了欺骗而生活在他们中间，是因为我们将来要在如他们这样的人中间生活，以至于无疑地受欺骗。祂看见了出卖祂的人，反倒拣选他来从事一项必要的工作。藉着此人的恶，祂成就了极大的善；然而，在十二人中，他仍被拣选，免得这小小的十二人团体连一个恶人都没有。这对我们来说是忍耐的榜样，因为我们必须在恶人中生活，必须忍受恶人，或认识他们或不认识他们：祂给了你忍耐的榜样，免得当你开始在恶人中生活时退缩。既然基督的学校在十二人中并没有失败，那么，当在大教会中应验了有关恶人混杂的预言时，我们更应当坚定不移。...\par

8. 但还要做什么？\BibleQuote{他们无缘无故为我暗设了网罗的败坏。}\InnerQuote{无缘无故}是什么意思？我未曾对他们行恶，丝毫未伤害他们。\BibleQuote{他们无端诽谤了我的灵魂。}\InnerQuote{无端}是什么？即是妄言，毫无实据。\BibleQuote{愿他们不知道的罗网临到他们身上}\BibleRef{咏 34:8}。绝妙的报应，再公平不过了！他们暗设罗网，以为我不知道：愿他们所不知道的罗网临到他们身上。因为我知道他们的罗网。但临到他们的是什么样的罗网呢？是他们所不知道的。让我们听听，免得他说的不是那个。\Quote{愿他们不知道的罗网临到他们身上。}或许那是他们为祂暗设的罗网，而另外的罗网会临到他们自己。并非如此，而是什么呢？\BibleQuote{恶人被自己罪恶的绳索缠住。}\BibleRef{箴 5:22}。他们企图欺骗别人，自己却受了欺骗。他们竭力加害于人，祸患却从那里临到他们。因为接着有：\Quote{愿他们暗设的网罗缠住自己，愿他们陷入自己的圈套。}就如同一个人为另一个人准备了一杯毒药，却忘记而自己把它喝了下去；或如同一个人挖了一个坑，想让他的敌人在黑暗中掉进去，而自己却忘记了所挖的坑，先走了那条路，便掉进了坑里。...\par

9. 这是对那想伤害我的恶人而言：至于我呢？\newline{}\Quote{但我的灵魂要因主而欢跃}\BibleRef{咏 34:9}，正如从祂那里听见了\Quote{我是你的救恩}；不向外寻求别的财富，不寻求充盈于地上的享乐和善物，而是自由地爱那真正的净配，并非渴望从祂那里得到任何令人愉悦的东西，唯独以祂自身为目标，因祂而喜乐。因为有什么比天主更好的会赐给我呢？天主爱我：天主爱你。看，祂向你提议：\Quote{求什么吧！}\BibleRef{玛 7:7} 如果皇帝对你说：\Quote{求什么吧}，你会喷涌出什么命令、什么尊位！你会为自己设想多么伟大的事，既接受又施与！当天主对你说：\Quote{求什么吧}，你会求什么呢？倒空你的心思，施展你的贪欲，尽力伸展，扩大你的愿望：说这话的不是别人，而是全能的天主：\Quote{求什么吧}。如果你是个爱财产的人，你会渴望整个大地，使所有出生的人都成为你的农夫或奴隶。而当你拥有了整个大地之后呢？你会求海洋，然而你无法在其中生活。在这贪婪中，鱼会比你更优越。但或许你会占有岛屿。越过这些吧；求你无法飞翔的空气；将你的欲望伸向天空，把太阳、月亮、星辰据为己有，因为造物主说：\Quote{求什么吧}：然而你找不到比造万物者本身更珍贵、更好的东西。寻求那造万物的祂，在祂内并从祂那里，你将拥有祂所造的一切。万物都是珍贵的，因为都是美丽的；但有什么比祂更美丽呢？它们是强大的，但有什么比祂更强大呢？祂不会把别的赐给你胜过祂自己。如果你找到了更好的东西，就求它吧。如果你求别的，你便是对祂不敬，也损害自己，因为宁愿选择祂所造的而不选择祂，而祂本要赐给你祂自己……\par

\Quote{但我的灵魂要在主内喜乐，要因祂的救恩而欢跃。}天主的救恩就是基督：\Quote{因为我亲眼看见了祢的救恩。}\BibleRef{路 2:30}\par

10. \Quote{我所有的骨头都要说：主啊，谁能像祢？}\BibleRef{咏 34:10} 谁能说得出这些话相称的意义？我认为这些话只能诵念，不能解释。你为何寻求这个或那个？有什么像你的主？祂就在你面前。\Quote{不义的人向我述说快乐，但主啊，不合祢的法律！}曾有迫害者说：崇拜萨图尔努斯，崇拜墨丘利。我不崇拜偶像（他说）：\Quote{主啊，谁能像祢？他们有眼却不能看，有耳却不能听。}\Quote{主啊，谁能像祢？}祢造了眼睛能看，耳朵能听。但我（他说）不崇拜偶像，因为它们是工匠所造。崇拜一棵树或一座山？工匠也造了它们吗？这里也是，主啊，谁能像祢？他们向我展示地上之物，祢是地的创造者。或许他们由此转向更高的受造物，对我说：崇拜月亮，崇拜这太阳，它像天上的一盏大灯用光制造白昼。在这里我也清楚地说：\Quote{主啊，谁能像祢？}月亮和星辰是祢造的，太阳统治白昼是祢点燃的，诸天是祢布置的。还有许多不可见的更好的东西。但或许这里也有人对我说：崇拜天使，敬拜天使。在这里我也要说：\Quote{主啊，谁能像祢？}连天使也是祢造的。天使之所以为天使，不过是因看见祢。与他们一起拥有祢，比因崇拜他们而失落祢更好。\par

11. 基督的身体，圣教会啊，愿你所有的骨头都说：\Quote{主啊，谁能像祢？}如果在迫害下肉体倒下了，让骨头说：\Quote{主啊，谁能像祢？}因为关于义人经上说：\Quote{主保护他们所有的骨头，一根也不折断。}有多少义人的骨头在迫害下被折断了呢？最后，\Quote{义人因信德而生活}，并且\Quote{基督使不敬虔的人成义}\BibleRef{罗 4:5}。但祂如何使任何人成义呢？除非是信而宣认。\Quote{因为人心里相信，可成正义；口里宣认，可获救恩}\BibleRef{罗 10:10}。因此那盗贼，虽然因偷窃被带到审判官面前，又从审判官到十字架，却在十字架上成了义：他心中相信，口里宣认。因为主不会对一位不义且尚未成义的人说：\Quote{今天你就要同我在乐园里了}\BibleRef{路 23:43}，然而他的骨头却被折断。因为当人来取下尸体时，由于安息日临近，发现主已经死了，祂的骨头没有折断\BibleRef{若 19:33}。但对那些还活着、以便取下的人，腿被折断，使他们因这痛苦死去而被埋葬。那么，其中一个始终不虔敬地挂在十字架上的盗贼，骨头被折断，而另一个心中相信、口里宣认以致得救的盗贼，难道不是同样被折断了吗？那么所说\Quote{主保护他所有的骨头，一根也不折断}又在哪里呢？除非在主的身体中，所有义人都被称为骨头，那些心中坚固、坚强，不屈服于任何迫害、任何诱惑而同意作恶的人……\par

12. \Quote{你拯救弱小于强暴，拯救贫寒于盗匪}……那拯救者是谁？不就是那\OriginalTerm{强有力的手}{Strong in hand}吗？就是那\Person{达味}要拯救穷弱脱离强于他的人。因为魔鬼对你太强了，他控制了你，因为他战胜了你，当你同意他时。但那强有力的手做了什么？\BibleQuote{人怎能进入一个有力者的家，抢劫他的家具呢？除非先把那有力者捆住，然后才抢劫他的家。}\BibleRef{玛 12:29}凭他自己的力量，至圣、至荣，他捆住了魔鬼，倾倒武器阻挡他的道路，使他能拯救穷弱，他们本无援助。因为谁是你的援助者，除非是主，你对他说：\Quote{上主，我的力量，我的救主。}如果你倚仗自己的力量，你就会因此跌倒；如果倚仗别人的，那人会统治你，而不是帮助你。所以当寻求惟独他，他救赎了他们，使他们自由，并赐下自己的血来购买他们，从仆人使他们成为他的弟兄……\par

13. 让我们的头说：\BibleQuote{诬妄的见证者起来，质问我所不知道的事。}\BibleRef{咏 34:11}但我们对我们的头说：主，你有什么不知道的？你果真什么都不知道吗？难道你不知道那些指控你者的心吗？你没有预见他们的诡计吗？你不是明知故犯地把自己交在他们手中吗？你不是来为要受他们的苦吗？那么你不知道什么？他不知道罪，因此他不知道罪，不是因为没有判断，而是因为没有犯。日常用语中也有这类说法，比如你说某人不知道站立，意思是他不站立；他不知道行善，因为他不做善事；他不知道作恶，因为他不作恶。……\Person{基督}不知道的，有什么比亵渎更甚呢？为此，他被迫害者质询，因他说真话，就被判定说了亵渎的话。\BibleRef{玛 26:65}但被谁呢？被那些以下所说的人：\BibleQuote{他们以恶报善，使我的灵魂孤独。}\BibleRef{咏 34:12}我给了他们丰收，他们回报我荒芜；我给了生命，他们回报死亡；我给了尊荣，他们回报侮辱；我给了医药，他们回报创伤；在他们回报我的一切中，真是荒芜。这荒芜在树上他诅咒了，因为找果子时找不到。\BibleRef{玛 21:19}有叶子，却没有果子；有言语，却没有行为。看，言语丰盛，行为荒芜。\BibleQuote{你教导人不可偷盗，自己却偷盗；你说人不可奸淫，自己却奸淫。}\BibleRef{罗 2:21}这样的人就是那些以不知道的事指控\Person{基督}的人。\par

14. \BibleQuote{但我，当他们烦扰我时，我穿上苦衣，以禁食谦卑我的灵魂，我的祈祷要回到我的怀中。}\BibleRef{咏 34:13}……弟兄们，如果我们以虔诚的好奇暂时揭开面纱，以心灵的眼目细察这圣经的内在部分，我们会发现这也是主所做的。他或许称他的肉体为苦衣。为何是苦衣？因为具有罪肉的相似。因为宗徒说：\BibleQuote{天主派遣了自己的儿子，带着罪恶肉身的形状，为因罪而惩罚罪恶在肉身上。}\BibleRef{罗 8:3}也就是说，他给自己的圣子穿上\OriginalTerm{苦衣}{sackcloth}，为通过苦衣惩罚山羊。并非有罪，我说不在于天主的圣言，也不在于那圣洁的灵魂和人的心智，即圣言和天主的智慧如此与自己结合成为一位。不，甚至在他的身体里也没有罪，但主身上有罪肉的相似；因为死亡只来自罪，\BibleRef{罗 5:12}而且那身体确实是可死的。因为若非可死，就不会死；若非死，就不会复活；若非复活，就不会给我们显示永生的榜样。所以，因罪引起的死亡被称为罪；就像我们说希腊语、拉丁语，意指不是肉体的舌头本身，而是由舌头所做之事。因为舌头在肢体中之一，如眼睛、鼻子、耳朵等；但希腊语是希腊词语，并非舌头是词语，而是词语由舌头发出。……所以，主的罪是因罪而引起的；因为他取了同一团具有因罪而应得死亡的肉。简而言之，出自\Person{亚当}的\Person{玛利亚}为罪而死，\Person{亚当}为罪而死，而出自\Person{玛利亚}的主的肉体为除去罪而死。主穿上这苦衣，因此他不被认识，因为他隐藏在苦衣之下。他说：\Quote{当他们烦扰我时，我穿上苦衣。}也就是说，他们发怒，我隐藏。因为若他不愿意隐藏，他也不能死，因为在一刹那，他只要放出他力量的一滴（如果可称为一滴），当他们想要抓他时，在他一句话\BibleQuote{你们找谁？}\BibleRef{若 18:4}之下，他们都退后跌倒在地上。这样的力量，他本不能在受难中谦卑，除非他隐藏在苦衣之下。\par

15. 再者，若我们理解了苦衣，又怎能理解禁食？基督何时愿意进食？当祂在树上寻求果子时，\BibleRef{谷 11:13}，若祂找到了，祂会吃吗？基督何时愿意喝水？当祂对撒玛黎雅妇人说：\BibleQuote{给我水喝}时，\BibleRef{若 4:7}；当祂在十字架上说：\BibleQuote{我渴}时？\newline{}基督饥饿或口渴，岂非为了我们的善行？因为那些钉死并迫害祂的人中，祂找不到善行，祂便禁食；因为他们以荒芜回报祂的灵魂。祂的禁食是何等之大？祂在十字架上几乎只找到一个强盗，可以尝到一点味道！因为宗徒们逃走了，藏在群众中。甚至伯多禄，曾许诺坚持到底直到主的死亡，这时却三次否认了祂，痛哭流涕，仍藏在群众中，仍害怕被认出。最后，看见祂死了，他们都绝望于自己的安全，并绝望了，祂在复活后找到他们，并与他们谈话时，发现他们悲伤哀恸，不再盼望什么……主一直处于极大的禁食中，除非祂安慰他们，以便以他们为食。因为祂安慰他们，鼓励他们，坚固他们，并将他们转化为祂自己的身体。这样，我们的主也处于禁食中。\par

16. \BibleQuote{我的祈祷要返回我怀中。}\newline{}在这节经文的怀中，显然有极大的深度，愿主赐我们能够探测。因为\Quote{怀中}意味着秘密。我们自身，弟兄们，在此被善意地劝告，要祈祷于我们自己的怀中，在那里天主看见，天主听见，没有人的眼睛穿透，除了那援助者之外无人看见；在那里苏撒纳祈祷，她的声音虽未被听见，却被天主听见……我们也读到耶稣独自在山上祈祷，\BibleRef{玛 14:23}；我们读到祂彻夜祈祷，\BibleRef{路 6:12}，甚至在受难之时。那么这是什么？\BibleQuote{我的祈祷要返回我怀中。}我不知道如何更好地理解关于主的事：暂时采取现在出现的想法；或许以后会有更好的想法，或对我，或对某个更好的人：\BibleQuote{我的祈祷要返回我怀中}：我理解这是说，因为祂在自己怀中拥有圣父。\BibleQuote{因为天主在基督内使世界与自己和好。}\BibleRef{格后 5:19}祂在自己内有祂祈祷的对象。祂并不远离祂，因为祂自己曾说过：\BibleQuote{我在圣父内，圣父在我内。}\BibleRef{若 14:10}但由于祈祷更属于真人（因为按照基督是圣言，祂不祈祷，而是垂听祈祷；也不为自己寻求援助，而是与圣父一同援助众人）：那么\BibleQuote{我的祈祷要返回我怀中}是什么？岂不是在我内，我的人性在我内呼求我的天主性？\par

17. \BibleQuote{我如同邻人、如同弟兄般讨祂喜悦：如同哀恸悲伤者，我这样谦卑了自己。}\BibleRef{咏 34:14}\newline{}现在祂回顾祂自己的身体：让我们现在也注视这一点。当我们祈祷喜乐时，当我们的心灵得享平安，不是因世事的顺遂，而是因真理之光：谁觉察这光，就知道我在说什么，并看见且承认所说的话：\BibleQuote{我如同邻人、如同弟兄般讨祂喜悦}：那时我们的灵魂讨天主喜悦，并非远离，因为经上说：\BibleQuote{在祂内}，有人这样说，\BibleQuote{我们生活、行动、存在都在祂内}\BibleRef{宗 17:28}，而是如同弟兄、如同邻人、如同朋友。但若它并非如此，以致能如此喜乐、如此闪耀、如此亲近、如此依附于祂，却看到自己远离那里，那么让它做以下的事：\BibleQuote{如同哀恸悲伤者，我这样谦卑了自己。}\BibleQuote{如同弟兄，我这样讨祂喜悦}祂说，接近时；\BibleQuote{如同哀恸悲伤者，我这样谦卑了自己}祂说，被移除并远离之后……伯多禄岂不是在说：\BibleQuote{祢是基督，永生天主之子}时接近了？然而同一个人因说：\BibleQuote{主，这事千万不可临于祢}而远离了。最后，祂对他的\Quote{邻人}，如同是接近的，说了什么？\BibleQuote{约纳的儿子西满，你是有福的。}对那远离且不像的，祂说了什么？\BibleQuote{退到我后面去，撒殚。}\BibleRef{玛 16:16}对那接近的，祂说：\BibleQuote{不是血和肉启示了你，而是我在天之圣父。}祂的光照耀你，你在祂的光中闪耀。但当远离时，他反对主应为我们救恩的受难，祂说：\BibleQuote{你所体会的，不是天主的事，而是人的事。}有人正确地将这两者放在一起，在一篇圣咏中说：\BibleQuote{我在神魂超拔中说：我从祢眼前被抛弃。}若他不接近，他岂会在神魂超拔中说这些话？因为神魂超拔是心灵的提升。他将自己的灵魂倾注于自身，并接近天主；又因某种云彩和肉身的重量再次被抛回地上，回想曾在哪里，并看见自己在哪里，他说：\BibleQuote{我从祢眼前被抛弃。}那么，\BibleQuote{我如同邻人、如同弟兄般讨祂喜悦}，愿祂恩赐在我们内成就；但当这不成就时，但愿这也成就：\BibleQuote{如同哀恸悲伤者，我这样谦卑了自己。}\par

18. \BibleQuote{他们却对我欢喜，并聚集起来}\BibleRef{咏 34:15}，只针对我：他们欢喜，我悲伤。但我们刚才在福音中听见：\BibleQuote{哀恸的人是有福的。}\BibleRef{玛 5:5}若哀恸的人有福，则欢笑的人可悲。\BibleQuote{他们对我欢喜，并聚集起来：鞭子聚集攻击我，他们却不知道。}因为他们把我不认识的事归咎于我，他们也不认识他们指控的是谁。\par

19. \BibleQuote{他们试探了我，并以嘲笑嘲笑了我} \BibleRef{咏 34:16}。这就是说，他们讥笑我、侮辱我；这是论元首，也是论身体。弟兄们，请想想教会如今的荣耀，回想它昔日的屈辱，回想从前基督徒如何到处逃窜，无论在哪里被发现，都受嘲笑、被殴打、被杀害、被投给野兽、被焚烧，人们幸灾乐祸地攻击他们。正如元首所经历的，身体也经历了。因为正如主在十字架上所受的，他的身体在不久前的一切迫害中也同样受了；就连现在，对他的迫害也没有停止。无论在哪里，人们一遇到基督徒，就惯于侮辱、迫害、讥笑他，称他为迟钝、无知、无精神、无知识。任凭他们为所欲为，基督在天上；任凭他们为所欲为，他已荣耀了自己的刑罚，已将他的十字架铭刻在所有人的额上；不虔敬的人被允许侮辱，但不被允许发作；然而，从口中所吐之言，可知他心里所怀的：\Quote{他们向我咬牙切齿。}\par

20. \BibleQuote{主，你何时观看？救我的灵魂免于他们的诡诈，将我的独一者从狮子中救出} \BibleRef{咏 34:17}。因为在我们看来，时间缓慢；这是以我们的身份说的：\Quote{你何时观看？} 也就是说，何时我们才能看见那些侮辱我们的人受惩罚？何时审判者因厌倦而听从寡妇的请求？ \BibleRef{路 18:3} 但我们的审判者并非出于厌倦，而是出于爱，迟延了我们的救恩；出于理性，而非出于需要；并非他现在不能帮助我们，而是为了使我们众人的数目直到终结得以满足。然而出于我们的渴望，我们说什么呢？\BibleQuote{主，你何时观看？救我的灵魂免于他们的诡诈，将我的独一者从狮子中救出：} 也就是说，将我的教会从狂怒的权势中救出。\par

21. 最后，你想知道那独一者是什么吗？请读下面的话：\BibleQuote{主，我要在大会中颂谢你；在沉重的民族中我要赞美你} \BibleRef{咏 34:18}。他确实说：\Quote{我要颂谢你：} 因为颂谢是在全体群众中进行的，但并非在所有人中天主都被赞美；全体群众听到我们的颂谢，但并非在全体群众中都有天主的赞美。因为在全体群众中，也就是在遍及全世界的教会里，有糠秕和麦子：糠秕飞散，麦子存留；因此，\BibleQuote{在沉重的民族中我要赞美你。} 在沉重的民族中，就是那不被诱惑之风吹走的人，在这样的民族中，天主被赞美。因为在糠秕中，他常常被亵渎……\par

22. \Quote{不要让我那无理欢庆的仇敌得胜于我：} 因为他们因我的糠秕而欢庆于我。\Quote{那无故恨我的人，} 也就是说，我从未伤害过的人；\BibleQuote{以眼示意} \BibleRef{咏 34:19}：也就是说，伪装成假善人，\BibleQuote{因为他们确实对我讲和平；却以愤怒图谋诡诈。} \BibleRef{咏 34:20} \BibleQuote{他们向我大大张口} \BibleRef{咏 34:21}。起初以眼示意，那些狮子想要掠夺并吞食；起初谄媚地说和平，然后以愤怒图谋诡诈。他们讲了什么和平？\Quote{师傅，我们知道你不看人情面，诚实地教导天主的道路。给凯撒纳税是否可以，还是不可以？} 他们确实对我讲了和平。那么怎样？难道你不认识他们，他们以眼示意欺骗了你？他确实认识他们；因此他说：\BibleQuote{你们为什么试探我，假善人？} \BibleRef{玛 22:16} 后来，\Quote{他们向我大大张口，} 喊叫：\BibleQuote{钉死他，钉死他！} \BibleRef{路 23:21} 并说：\Quote{哈哈，哈哈，我们的眼睛看见了。} 这是当他们侮辱他时：\Quote{哈哈，哈哈，你这位基督，向我们说预言吧。} \BibleRef{玛 26:68} 正如他们关于钱试探他时，和平是伪装的，同样现在侮辱性的赞美也是一样。\BibleQuote{他们说：哈哈，哈哈，我们的眼睛看见了} \BibleRef{咏 34:21}：也就是说，你的作为，你的奇迹。这人是基督。\Quote{如果他是基督，让他从十字架上下来，我们就相信他。他救了别人，却不能救自己。} \Quote{我们的眼睛看见了。} 这就是他自夸的一切，当他\BibleQuote{自称是天主之子}时。 \BibleRef{若 19:7} 但主忍耐地悬挂在十字架上；他没有失去他的能力，而是显明了忍耐。因为他后来能从坟墓中复活，从十字架上下来又算得了什么大事呢？但他似乎向侮辱者屈服了；可爱的弟兄们，这是因为他复活后要显示给他自己的人，而不显示给他们，这是一个伟大的奥秘；因为他的复活象征了新生，但新生只有他的朋友知道，而非他的敌人。\par

23. \BibleQuote{主，你已看见，不要缄默} \BibleRef{咏 34:22}。什么是\Quote{不要缄默}？审判吧。因为论到审判，经上某处说：\Quote{我缄默了；难道我要永远缄默吗？} 论到审判的迟延，对罪人说：\Quote{你做了这些事，我缄默了；} \Quote{你想我完全像你一样。} 他如何缄默？他借着先知说话，在福音中亲口说话，借着福音作者说话，借着我们说话，当我们讲真理时。那么怎样？他在审判上缄默，不在训诫上，不在教导上。但先知仿佛呼求并预言他的审判：\BibleQuote{主，你已看见：不要缄默；} 也就是说，你不会缄默，你必然要审判。\Quote{主，不要远离我。} 在你的审判来临之前，不要远离我，正如你应许的：\Quote{看，我同你们天天在一起，直到今世的终结。}\par

24. \BibleQuote{主啊，兴起，顾及我的\OriginalTerm{审判}{judgment}}\BibleRef{咏 34:23}。是什么审判？就是你身处患难，被劳苦和痛苦折磨？难道许多恶人不也遭受同样的事吗？是什么审判？因此你是义的吗，因为你遭受这些事？不：而是什么？\BibleQuote{顾及我的审判}。接下来是什么？\BibleQuote{顾及我的\OriginalTerm{案件}{cause}，我的天主，我的主。}不是我的\OriginalTerm{惩罚}{punishment}，而是我的案件；不是那强盗与我共有的，而是那论及它的，\BibleQuote{为\OriginalTerm{义德}{righteousness}而受\OriginalTerm{迫害}{persecution}的人是有福的。}\BibleRef{玛 5:10}因为案件是区别的。因为惩罚对善人和恶人是同等的。因此，是案件，而非惩罚造就\OriginalTerm{殉道者}{Martyr}，因为如果惩罚造就殉道者，那么所有矿坑都将充满殉道者，每条锁链都将拖曳殉道者，所有被刀剑处决的都将被加冕。因此，让案件得到区别；让没有人说，因为我受苦，我就是义人。因为那首先受苦的，是为义德受苦，因此他加上了一个重大的例外：\BibleQuote{为义德而受迫害的人是有福的。}因为许多人有好的案件却施行迫害，许多人因坏的案件而遭受迫害。因为如果迫害不能正义地施行，就不会在某篇圣咏中说：\BibleQuote{暗中诽谤邻舍的人，我要迫害他。}……所以，让没有人说：我遭受迫害；让他不要细查惩罚，而要证明案件：免得他若不能证明案件，就被列在\OriginalTerm{不虔敬}{ungodliness}的人中。因此，这个人多么谨慎、多么卓越地推荐了自己：\BibleQuote{主啊，顾及我的审判，}不是我的惩罚；\BibleQuote{甚至我的案件，我的天主，我的主。}\par

25. \BibleQuote{主啊，请按我的\OriginalTerm{义德}{righteousness}审判我}\BibleRef{咏 34:24}；也就是说，顾及我的\OriginalTerm{案件}{cause}。不是按我的\OriginalTerm{惩罚}{punishment}，而是\BibleQuote{按我的义德，主，我的天主}，也就是说，按此审判我。\BibleQuote{不要让他们因我而欢喜；}也就是说，我的仇敌。\par

26. \BibleQuote{不要让他们心里说：啊哈，啊哈，我们正愿如此}\BibleRef{咏 34:25}；就是说，我们做了我们所能做的，我们杀了他，我们把他带走了。\BibleQuote{不要让他们说：}向他们表明他们什么也没做。\BibleQuote{不要让他们说：我们\OriginalTerm{吞下}{swallow}了他。}那些\OriginalTerm{殉道者}{Martyr}从何而说：\BibleQuote{若不是主在我们这边，他们早已活活吞下我们了。}\BibleQuote{吞下我们}是什么意思？就是吞入他们自己的\OriginalTerm{身体}{body}。因为你吞下什么，你就把它吞入你的身体。世界会吞下你；你要吞下世界，把它吞入你的身体：杀而吃。正如对伯多禄说的：\BibleQuote{宰了吃吧；}\BibleRef{宗 10:13}你要在他们身上杀死他们之所是，使他们成为你之所是。但反之，如果他们说服你走向\OriginalTerm{不虔敬}{ungodliness}，你就会被他们吞下。不是当他们\OriginalTerm{迫害}{persecution}你时你被他们吞下，而是当他们说服你成为他们那样时。\BibleQuote{不要让他们说：我们吞下了他。}你要吞下外邦人的身体。为什么是外邦人的身体？因为它会吞下你。你对它做它要对你做的事。因此，或许那\OriginalTerm{牛犊}{calf}被磨成粉末，投在水里，给以色列子民喝，\BibleRef{出 32:20}好使不虔敬的身体被以色列吞下。\BibleQuote{愿那些因我受害而欢喜的都蒙羞受辱；愿他们以羞耻和耻辱为衣}\BibleRef{咏 34:26}，好叫我们吞下那蒙羞受辱的。\BibleQuote{那些说恶言攻击我的：}愿他们蒙羞，愿他们受辱。\par

27. 现在，头与肢体说什么？\BibleQuote{愿那些喜爱我\OriginalTerm{义德}{righteousness}之\OriginalTerm{案件}{cause}的人欢呼喜乐：}就是那些依附我的\OriginalTerm{身体}{Body}的人。是的，让他们\BibleQuote{不断地说：愿主被尊崇，因祂喜悦祂\OriginalTerm{仆人}{servant}的顺利}\BibleRef{咏 34:27}。\BibleQuote{我的舌头要述说你的义德，和你的\OriginalTerm{赞美}{praise}}\BibleRef{咏 34:28}。谁的舌头能终日述说对天主的赞美？你看，我现在讲得稍长了些；你已疲倦。谁能终日赞美天主？我要提出一个补救方法，使你若愿意，可以终日赞美天主。无论你做什么，做得好，你就赞美了天主。当你唱圣歌时，你赞美天主，但你的舌头若不与你的心一同赞美祂，又有什么用？你停止唱圣歌，离开去休息？不要醉酒，你就赞美了天主。你睡觉去？起来不要作恶，你就赞美了天主。你经商？不要行不义，你就赞美了天主。你耕田？不要挑起纷争，你就赞美了天主。以你工作的纯洁，预备自己终日赞美天主。\par

\SourceRecord{Translated by J.E. Tweed.；Nicene and Post-Nicene Fathers, First Series；Vol. 8.；Edited by Philip Schaff.；Buffalo, NY: Christian Literature Publishing Co.,；1888.；<http://www.newadvent.org/fathers/1801035.htm>.}

% structure: p0001 source=h1 target=section reason=page-title

\section{\WorkTitle{圣咏第三十六篇释义}}

1. ……\BibleQuote{恶人在自己心里说：他要犯罪；在他眼前没有对天主的敬畏。} \BibleRef{咏 35:1} 他不是指一个人，而是指一群恶人，他们与自己作战，因为他们不理解，好能好好生活；不是因为他们不能，而是因为他们不愿意。因为，当一个人努力理解某事，却因肉身的软弱而不能时，是一回事——正如圣经某处说：\BibleQuote{可腐朽的身体压着灵魂，地上的帐幕拖累着深思许多事的心智；} 但人的心恶意地与自己作对，以致于它本可理解（只要存有善意）的事却不理解，不是因为难，而是因为意志作对，则是另一回事。但情况正是如此：人爱自己的罪，而恨天主的诫命。因为天主圣言是你的仇敌，如果你是你自己邪恶的朋友；但如果你是邪恶的仇敌，天主圣言就是你的朋友，也是你邪恶的仇敌。……\par

2. \BibleQuote{因为他在祂面前行诡诈。} \BibleRef{咏 35:2} 在谁面前？在他面前，那人的眼前没有对祂的敬畏，而行诡诈。\BibleQuote{为要查出不义而恨恶它。} 他行事如此，以致查不出来。因为有些人仿佛竭力寻找自己的不义，却害怕找到它；因为如果他们找到了，就会对他们说：离开它；这是你在不知道以前做的；你在无知中行了不义，天主赐予宽恕；现在你已经发现了，就放弃它，这样你的无知就易于得到宽恕；并且你能以清白的脸面向天主说：\BibleQuote{求你不要记念我幼年的罪愆和我的无知。} 因此他寻找它，因此他害怕找到它；因为他是在诡诈地寻找。当一个人说：我不知道那是罪？当他看见那是罪，并停止做那项只因为他无知而做的罪时，这样的人诚然会认识他的罪，好查出它并恨恶它。但现在许多人行诡诈来查出他们的不义：他们不是出于内心去查出不义并恨恶它。因为在寻找不义的过程中有诡诈，在找到它时就会有辩护。因为当一个人找到了他的不义，如今他已清楚知道那是不义。你说：不要做。而行诡诈去找它的人，如今他找到了，却不恨它；因为他会说什么？\Quote{有多少人这样做！谁不做呢？天主会毁灭他们所有人吗？或至少他说：如果天主不愿这些事发生，那些犯同样罪的人还能活吗？} 你看出你行诡诈去查出不义吗？因为如果你不是诡诈而是真诚地去做，你如今就会找到它并恨恶它；如今你找到了它，却为它辩护；所以当你寻找时，你是诡诈的。\par

3. \BibleQuote{他口中的言语是不义和诡诈：他不愿明白，以便行善。} \BibleRef{咏 35:3} 你看到他把这归因于意志：因为有些人愿意明白却不能，有些人则不愿意明白，因此不明白。\BibleQuote{他不愿明白，以便行善。}\par

4. \BibleQuote{他在床上图谋不义。} 祂说什么？\BibleQuote{在床上？} \BibleRef{咏 35:4} \BibleQuote{恶人在自己心里说：他要犯罪：} 上文所说的\BibleQuote{在自己心里}，在这里说\BibleQuote{在床上。} 我们的床就是我们的心：我们在那里遭受邪恶良心的颠簸；当我们的良心良好时，我们在那里安息。谁爱他心的床，就让他在这床内做些好事。那就是我们的床，主耶稣基督命令我们在那里祈祷。\BibleQuote{进入你的内室，关上门。} \BibleRef{玛 6:6} \BibleQuote{关上门}是什么意思？不要从天主那里期待外在的事物，而是内在的；\BibleQuote{你的父在暗中看见，必要公开报答你。} 谁不关门？谁向天主大量祈求这样的事，并以这样的方式引导他的一切祈祷，好领受今世之物。你的门是敞开的，众人看见你祈祷。关上门是什么意思？求那只有天主知道如何赐予的事。当你关上门祈祷时，你求的是什么？是\BibleQuote{眼所未见，耳所未闻，人心所未想到的。} 或许它还没有进入你的床，就是你的心。但天主知道祂要赐予什么；但什么时候呢？当主被启示时，当审判者显现时。……\par

5. \Quote{他在一切不善的道上立定了自己。} 什么是 \Quote{他立定了自己}？就是恒久地犯罪。因此论到某个虔诚而善良的人也说：\Quote{他没有站罪人的道路。} 如同这里的 \Quote{没有站}，那里的 \Quote{立定了自己}。\Quote{却不恨恶邪恶。} 这就是终结，这就是果实：若人不能没有邪恶，至少也当恨恶它。因为你一旦恨恶它，就几乎不会去行任何邪恶。因为罪在我们的可朽身体内，但宗徒怎么说？\Quote{不要让罪在你们可朽的身体内为王，使你们顺从它的私欲。} \BibleRef{罗 6:12} 什么时候它不再在其中？当他在我们身上所说的这句话应验时：\Quote{当这可朽坏的穿上不朽坏，这可死的穿上不死。} \BibleRef{格前 15:54} 在这事成就之前，身体中对罪有一种喜爱，但对智慧之言和天主的诫命有更大的喜爱和愉悦。要胜过罪和它的私欲。你要恨恶罪和邪恶，好使你能与天主结合，祂也像你一样恨恶它。现在既在心灵上与天主的律法结合，你就在心灵上事奉天主的律法。如果在肉身中你因此事奉罪的律法，\BibleRef{罗 7:25} 因为你里面有某些肉体的喜爱，那么当你不再争战时，这些就没有了。不争战而处于真正持久的和平是一回事；争战而得胜是另一回事；争战而被胜是另一回事；完全不争战而被带走是另一回事……\par

6. \Quote{上主，祢的仁慈达于高天，祢的真理上彻云表。} \BibleRef{咏 35:5} 我不明白祂所说的在天上的仁慈是什么。因为上主的仁慈也在地上。经上记载：\Quote{大地充满了上主的仁慈。} 那么，当祂说 \Quote{上主，祢的仁慈达于高天} 时，说的是哪一种仁慈呢？天主的恩赐一部分是暂时的、属地的，一部分是永恒的、属天的。谁为此敬拜天主，为获得那些向所有人开放的暂时和属地的福乐，他仍如同禽兽：他确实享受天主的仁慈，但不是那例外的仁慈，这仁慈只赐给义人、圣者、善人。那丰盈地赐给众人的恩赐是什么？\Quote{祂使太阳升起，光照恶人也光照善人；降雨给义人也给不义的人。} \BibleRef{玛 5:45} 谁没有拥有天主的这仁慈呢？首先，他有存在，他从禽兽中被分别出来，是有理性的受造物，从而能理解天主；其次，他享受这光明、空气、雨水、果实、四季的变迁，以及一切属地的舒适、身体的健康、朋友的关爱、家庭的平安？这一切都是美好的，都是天主的恩赐……\par

7. 但这人正确地理解了应从天主求得何种仁慈。\Quote{上主，祢的仁慈达于高天；祢的真理上彻云表。} 也就是说，祢赐给祢圣徒的仁慈，是属天的，不是属地的；是永恒的，不是暂时的。祢怎能向人宣告它呢？因为 \Quote{祢的真理上彻云表。} 谁能认识天主属天的仁慈，除非天主亲自向人宣告？祂如何宣告？藉着差遣祂的真理直达云表。云是什么？天主的圣言的宣讲者……真理上彻云表；因此，那在天上而不在地上的天主的仁慈才能向我们宣告。确实，弟兄们，云就是真理之言的宣讲者。当天主藉着祂的宣讲者发出警告时，祂藉着云发出雷声。当天主藉着祂的宣讲者行奇迹时，祂藉着云闪电，藉着云令人惊惧，并降下雨水。那么，那些宣讲天主福音的宣讲者，就是天主的云。让我们盼望仁慈，但盼望那在天上的仁慈。\par

8. \Quote{祢的公义好似天主的山岳，祢的断论如同深渊。} \BibleRef{咏 35:6} 谁是天主的山岳？那被称为云的，也就是天主的山岳。伟大的宣讲者是天主的山岳。正如太阳升起时，先以光明照亮山岳，然后光明降到大地的最低处：同样，我们的主耶稣基督来到时，首先光照了宗徒的高处，首先照亮了山岳，于是祂的光就降到世界的幽谷。因此，祂在某一圣咏中说：\Quote{我举目仰望山岳，我的救助从何而来？} 但不要以为山岳本身会给你救助：因为他们接受才能给予，并非凭自己施与。如果你停留在山岳上，你的希望就不会坚固；但你的希望和信赖应当在那光照山岳者身上。你的救助确实藉着山岳来到你那里，因为圣经藉着山岳、藉着伟大的真理宣讲者传给你们；但不要将希望寄托在他们身上。请听祂接着说：\Quote{我举目仰望山岳，我的救助从何而来？} 那么，山岳给你救助吗？不；请听下文：\Quote{我的救助来自上主，祂造了天地。} 救助藉着山岳而来，却不是来自山岳。那么来自谁呢？\Quote{来自那造了天地的上主。} ……\par

9. \Quote{祢的审判如同深渊。} 他称深渊为罪的深处，人因轻视天主而进入其中；正如在某处所说：\BibleQuote{天主任凭他们随从心中的情欲，行不合理的事。}\BibleRef{罗 1:28} ……因为他们骄傲、忘恩，所以理当被交给自己心中的情欲，成了极大的深渊，以致他们不仅犯罪，还诡诈行事，以免明白自己的罪孽并憎恨它。这就是邪恶的深处：不愿发现并憎恨它。但人如何进入那深处？请看：\Quote{祢的审判是深渊。} 正如山由天主的正义而成，并因祂的恩宠而变得高大；同样，也因祂的审判，那些沉到最低处的人进入深渊。因此，愿山使你喜悦，远离深渊，转向那所说的：\Quote{我的帮助来自主。} 但藉着什么？\Quote{我举目向山。} 这是什么意思？我要直说。在天主的教会里，你发现深渊，也发现山；你发现善人很少，因为山少，深渊广阔；也就是说，你发现许多人因天主的愤怒而生活邪恶，因为他们这样行事，以致被交给自己心中的情欲；如今他们甚至辩护自己的罪，不承认，却说：为什么？我做了什么？某人做了这事，某人做了那事。他们甚至辩护天主的话所谴责的事。这就是深渊。因此，在某处\BibleRef{箴 18:3}圣经说（要听这深渊）：\BibleQuote{罪人一旦进入罪之深处，就轻视。} 看，\Quote{祢的审判如同深渊。} 但你并非山，也未在深渊中；逃离深渊，趋向山；但不要停留在山上。\Quote{因为你的帮助来自主，祂创造了天和地。}\par

10. 因他说：祢的仁慈在天上，好使人知道也在地上，他说：\BibleQuote{主啊，祢拯救人和牲畜，祢的仁慈何其增多，天主啊}\BibleRef{咏 35:7}。祢的仁慈伟大，祢的仁慈多样，天主；祢向人和牲畜都显示。因为人的拯救来自谁？来自天主。牲畜的拯救不也来自天主吗？因为那造人的，也造了牲畜；那造两者的，拯救两者；但牲畜的拯救是暂世的。但有些人把这当作大事向天主祈求，而天主已赐给牲畜了。\Quote{祢的仁慈，天主啊，是增多的，} 以至于不仅赐给人，也赐给牲畜同样赐给人的拯救，即肉身的、暂世的拯救。\par

11. 难道人没有在天主那里保留一些牲畜不配、牲畜达不到的吗？显然有。那他们拥有的是什么？\Quote{人的子女投靠在祢翅膀的荫下。} 亲爱的，请留意这最愉快的句子：\Quote{祢拯救人和牲畜。} 首先，他提到\Quote{人和牲畜}，然后提到\Quote{人的子女}；好像\Quote{人}是一回事，\Quote{人的子女}是另一回事。有时在圣经中，\Quote{人的子女}泛指所有人，有时以某种特殊方式、带有某种特殊含义，以致并非指所有人；主要是在有所区分时。因为这里并非无故如此排列：主啊，祢拯救人和牲畜，但人的子女…… 仿佛将前者搁置一旁，他把\Quote{人的子女}分别出来。与谁分别？不仅与牲畜分别，也与那些向天主寻求牲畜的拯救、并把此当作大事的人分别。那么，谁是人的子女？就是那些投靠在祂翅膀荫下的人。因为那些人同牲畜一起以拥有为乐，但人的子女以望德为乐：那些人随同牲畜追求现世之物，这些人同天使一同盼望未来之物……\par

12. \Quote{他们因祢殿中的丰盈饱足} \BibleRef{咏 35:8}。祂应许我们一些伟大的事。祂要说出它，却未说出。难道祂不能，还是我们未能领受？我的弟兄们，我敢说，即使是那些向我们宣告真理的圣洁舌和心，他们所宣告的，既不能言说，甚至不能思想。因为这是一件伟大而不可言喻的事；就连他们也只是模糊地观看，如宗徒所说：\Quote{我们现在是藉着镜子观看，模糊不清，到那时就要面对面了。} \BibleRef{格前 13:12} 看，那些模糊观看的人，就这样迸发出来。那么，当我们面对面观看时，我们将成为什么？那是他们心中孕育，却不能用舌头说出来，使人能领受的。因为他说\Quote{他们因祢殿中的丰盈饱足}，有什么必要呢？他寻求一个词语，从人间事物中表达他想要说的；因为他看见人们沉醉在醉酒中，确实无节制地饮酒，却失去知觉，他看见了该说什么；因为当那不可言喻的喜乐被领受时，人的灵魂在某种程度上将失落，它将变为神圣，并因天主殿中的丰盈而饱足。因此在另一篇圣咏中也说：\Quote{祢令人沉醉的杯爵，何其美妙！} 殉道者正是因这杯爵而饱足，他们前往受难时，竟不认识自己的亲人。什么如此令人沉醉，以致不认识哭泣的妻子、儿女、父母？他们不认识他们，不认为他们在眼前。不要惊奇：他们醉了。他们是因什么而如此？看，他们领受了一个杯爵，因此饱足。因此他也感谢天主说：\Quote{我拿什么报答上主赐给我的一切恩惠？我要举起救恩的杯爵，呼号上主的名号。} 因此，弟兄们，让我们成为孩童，并信赖在祂翅翼的荫庇下，因祂殿中的丰盈而饱足。按我所能的，我已说了；按我所能的，我看见；而我看见的，我不能说出。\Quote{祢也必叫他们畅饮祢喜乐的川流。} 我们称水流如洪水为川流。将有一天，天主的仁慈如洪水泛滥，使那些现今信赖在祂翅翼荫庇下的人陶醉。那喜乐是什么？如同一条川流令口渴者陶醉。那么，现在口渴的人，当存希望；现在口渴的人，当怀希望；陶醉时，他必拥有；在拥有之前，他当在希望中口渴。\Quote{饥渴慕义的人是有福的，因为他们要得饱饫。} \BibleRef{玛 5:6}\par

13. 那么，你将被什么泉源涌溢，祂喜乐的川流又从何处涌出？\Quote{因为，} 他说，\Quote{在祢那里有生命的泉源。} 生命的泉源是什么？岂不是基督？祂曾以肉身来到你那里，为浸润你干渴的双唇；祂必满足你这位信赖者，祂曾浸润你这位口渴者。\Quote{因为，在祢那里有生命的泉源；藉祢的光明，我们才能看见光明。} \BibleRef{咏 35:9} 在这里，泉源是一回事，光明是另一回事；在那里却不是这样。因为那泉源，也是光明；无论你称它为什么，它都不是你所称呼的；因为你找不到合适的名字；因为它不限于一个名称。如果你说它只是光，就会有人对你说：那么，我被告知要饥渴就是无意义的，因为谁吃光呢？它明明直接对我说：\Quote{心里洁净的人是有福的，因为他们要看见天主。} \BibleRef{玛 5:6} 如果它是光，我必须预备我的眼睛。也要预备嘴唇；因为那光也是泉源：泉源，因为它满足口渴者；光，因为它光照盲者。在这里，有时光在一处，泉源在另一处。因为有时泉源甚至在黑暗中流淌；有时你在旷野忍受烈日，却找不到泉源：因此在这里这两者可能分离；在那里你不会疲倦，因为有泉源；在那里你不会昏暗，因为有光。\par

14. \Quote{求祢向认识祢的人显示祢的仁慈；向心地正直的人显示祢的正义。} \BibleRef{咏 35:10} 如我已说过的，心地正直的人是在此生遵行天主旨意的人。天主的旨意有时是你应健康，有时是应患病。如果你健康时天主的旨意是甘甜的，患病时天主的旨意是苦涩的，你就不是心地正直的人。为什么？因为你不想使你的意志符合天主的旨意，反而想使天主的旨意屈就于你的意志。那是正直的，但你是弯曲的；你的意志必须被调整为正直，以符合那旨意，而不是将那旨意扭曲以符合你；这样你就有一个正直的心。你在今世顺遂，愿天主受赞美，因祂安慰你；你在今世艰难，愿天主受赞美，因祂惩戒并考验你；这样你便成为心地正直的人，说：\Quote{我要时时赞美上主，对祂的赞颂常在我口中。}\par

15. \BibleQuote{让骄傲的脚不向我而来} \BibleRef{咏 35:11}。但现在他说：人子们将信任于祢翅膀的荫下；他们必因祢殿中的丰满而饱足。当一个人开始从那源泉丰盛地涌流时，让他谨慎，免得他变得骄傲。因为同样的泉源并不欠缺于第一个人亚当；但骄傲的脚向他而来，罪人的手挪移了他，即魔鬼骄傲的手。正如那诱惑者自称的：\BibleQuote{我要坐在北方的极处} \BibleRef{依 14:13}；他这样说服他，说：\BibleQuote{尝吧，你们将如同神} \BibleRef{创 3:5}。于是我们因骄傲如此堕落，以至于到达这必死的境地。又因骄傲伤害了我们，谦卑使我们痊愈。天主谦卑地来了，为从那如此巨大的骄傲创伤中医治人类。祂来了，因为\BibleQuote{圣言成了血肉，寄居在我们中间} \BibleRef{若 1:14}。祂被犹太人拿住，受他们辱骂。你们听见福音诵读时，他们说了什么，他们向谁说：\Quote{你附有魔鬼}；祂不说：\Quote{你附有魔鬼}，因为你们仍在罪中，魔鬼占据你们的心。祂没有说这话，若祂说了，祂说的是真话；但祂不宜说这话，免得祂似乎不是在传讲真理，而是回以恶言。祂放过了所听见的，如同没有听见。因为祂是医生，来医治狂人。正如医生不在意他从狂人口中听到什么，而在意狂人如何康复而变得清醒；即使他受到狂人的打击，他也不在意；反而当他给他造成新伤时，他治愈了他的旧热病：同样，主来到病人那里，来到狂人那里，以便无论祂听见什么，无论祂忍受什么，祂都轻视；借此教导我们谦卑，使我们藉由谦卑受教，从骄傲中被治愈；为此祂在这里祈祷被解救，说：\Quote{让骄傲的脚不向我而来；也不让罪人的手挪移我}。因为若骄傲的脚来了，罪人的手便挪移。什么是罪人的手？就是那出坏主意者的工作。你变得骄傲了吗？那出坏主意的人立刻败坏你。谦卑地将自己固定在天主内，不要太在意别人对你说的话。因此，在别处所说的：\Quote{求祢洗净我的隐密罪过；也求祢使祢的仆人免受他人的罪过}。什么是\Quote{我的隐密罪过}？\Quote{让骄傲的脚不向我而来}。什么是\Quote{也求祢使祢的仆人免受他人的罪过}？\Quote{不要容恶人的手挪移我}。守住内心，你必不怕外面。\par

16. 但你为何如此惧怕这事？因为经上说：\BibleQuote{所有行恶的都由此跌倒} \BibleRef{咏 35:12}；以致他们陷入了那深渊，关于此，经上说：\Quote{祢的判断如同巨大的深渊}；以致他们甚至进入了那深处，其中藐视的罪人已经跌倒。\Quote{跌倒}。他们起初如何跌倒？因着骄傲的脚。听听骄傲的脚。\Quote{他们虽然认识天主，却没有以祂为天主而光荣祂}。因此骄傲的脚向他们而来，他们由此进入深处。\BibleQuote{天主任凭他们随从心里的情欲，行那些不合理的事} \BibleRef{罗 1:21}。那说\Quote{让骄傲的脚不向我而来}的人所惧怕的，是罪根和罪首。他为何说\Quote{脚}？因为人因骄傲而行，离弃了天主，离开了祂。他称自己的脚为他的情感。\Quote{让骄傲的脚不向我而来；也不让恶人的手挪移我}：也就是说，不要让恶人的行为将我挪离祢，使我愿意效法他们。但他为何针对骄傲说\Quote{所有行恶的都由此跌倒}？因为那些现在不虔诚的人，已因骄傲跌倒。因此，主警告祂的教会时说：\BibleQuote{他要伤害你的头，你要伤害他的脚跟} \BibleRef{创 3:15}。蛇窥视着，等待骄傲的脚向你而来，当你跌倒时，好把你推下去。但你当注视它的头：一切罪的开端是骄傲。\BibleRef{德 10:13} \Quote{所有行恶的都由此跌倒：他们被赶出，不能站立}。首先，那在真理中站不住的一位，然后通过他，那些天主从乐园赶出的人。因此，那个谦卑的人——他曾说他不配解祂的鞋带——没有被赶出，而是站着听祂，并因新郎的声音而大大喜乐；不是因自己的声音，免得骄傲的脚向他而来，他就会被赶出，不能站立……\par

\SourceRecord{Translated by J.E. Tweed.；Nicene and Post-Nicene Fathers, First Series；Vol. 8.；Edited by Philip Schaff.；Buffalo, NY: Christian Literature Publishing Co.,；1888.；<http://www.newadvent.org/fathers/1801036.htm>.}

% structure: p0001 source=h1 target=section reason=page-title

\section{\WorkTitle{《圣咏第三十七篇释义》}}

论圣咏的第一部分。\par

1. 那些不愿因善而生、却愿长久作恶的人，听到末日的来临便恐惧战栗。然而，天主有意使那日子保持未知，这是有益的：为使人心常备不懈，期待那明知必来、却不知何时来临的事。但是，既然我们的主耶稣基督被派遣到我们中间，作我们的\Quote{师傅}，祂说，\BibleQuote{连人子也不知道那日子}，\BibleRef{谷 13:32}，因为知道那日子不属于祂作为我们师傅的职责，不应藉着祂而让我们知道。事实上，圣父所不知道的，圣子也不可能不知道；因为那正是圣父自己的知识，即祂的智慧；而祂的圣子、祂的圣言，就是\Quote{祂的智慧}。但是，因为知道那日子对我们没有益处——那日子被祂所知，但祂来确实是为教导我们，虽然并非教导我们那对我们无益的知识——所以祂不仅作为师傅教导我们一些事，也作为师傅留下一些未教的事。因为，作为师傅，祂知道该教导什么对我们有益，也知道不该教导什么对我们有害。这样，按照某种说话方式，圣子被说成不知道祂所不教导的：即如同我们日常所说的，祂使我们不知道的事，祂便说不知道……\par

2. 身为基督徒的你，正是为此而困扰：你看到行为恶劣的人事事亨通，拥有丰裕的这类事物；你看到他们身体健康，享有尊贵的荣誉；你看到他们的家庭没有遭遇不幸；你看到他们亲属的幸福、侍从的殷勤、极强的影响力、以及没有悲哀事件打扰的生活；你看到他们品行最放荡，而外在资源最丰富；于是你的心说：没有天主的审判；一切事物都随着偶然事件飘荡，在无序和不规则的变动中翻腾。因为你说，如果天主眷顾人间的事，那么恶人的昌盛和我的无辜受难又怎么解释呢？灵魂的每一种疾病在圣经中都有其适当的疗法。那么，如果一个人的疾病属于那种类型，他心里说出这样的话，让他饮下这首圣咏作为药方……\par

3. \BibleQuote{不要因作恶的人而嫉妒，也不要因行不义的人而嫉妒。} \BibleRef{咏 36:1} \BibleQuote{因为他们如草快要枯萎，又如青绿的禾草快要凋零。} \BibleRef{咏 36:2} 在你看来似乎长久的事，在天主眼中是\Quote{快要}的。你要使自己符合天主，那么它在你眼中也就成了\Quote{快要}。他在此称之为\Quote{草}的，我们理解为\Quote{青绿的禾草}。它们是一些没有价值的东西，只占据地表面，没有深厚的根基。在冬天，它们青绿；但当夏天的太阳开始炙烤时，它们就会枯萎。因为现在是冬天。你的荣耀尚未显现。但若你的爱有深厚的根，如同许多树在冬天那样，严寒过去，夏天（即审判之日）就会来临；那时，草的青绿将要枯萎。那时，树的荣耀将要显现。\BibleQuote{因为你们}（宗徒说）\BibleQuote{已经死了}，\BibleRef{哥 3:3} 正如树在冬天看起来如同死了、如同枯萎了一样。那么，如果我们死了，希望何在呢？根在里面；我们的根在哪里，我们的生命也在那里，因为我们的爱安放在那里。\BibleQuote{你们的生命与基督一同藏在天主内。} \BibleRef{哥 3:3} 这样扎根的人何时会枯萎呢？但我们的春天何时到来？我们的夏天何时到来？何时繁茂的枝叶环绕我们，丰硕的果实使我们富足？这一切何时发生？请听接下来的话：\BibleQuote{当基督，我们的生命显现时，那时你们也要与祂一同在光荣中出现。} 那么我们现在该做什么？\BibleQuote{不要因作恶的人而嫉妒，也不要因行不义的人而嫉妒。因为他们如草快要枯萎，又如青绿的禾草快要凋零。} \BibleRef{咏 36:1-2}\par

4. 那么你该做什么呢？\BibleQuote{你该信赖上主}，\BibleRef{咏 36:3} 因为他们也信赖，但不是\Quote{信赖上主}。他们的希望是易逝的，他们的希望是短暂的、脆弱的、飘忽的、易变的、没有根基的。\BibleQuote{你要信赖上主。} \BibleRef{咏 36:3} 你说：\Quote{看，我信赖；我该做什么？}\par

\BibleQuote{并要行善。} \BibleRef{咏 36:3} 不要行你看见那些在邪恶中亨通的人所行的恶事。\BibleQuote{行善，并安居在这地上。} \BibleRef{咏 36:3} 免得你行善却不\BibleQuote{安居在这地上}。因为教会就是主的土地。是她，天父——耕耘者——浇灌并栽培她。因为有许多人，好像行了善工，但由于他们不\BibleQuote{安居在这地上}，他们不属于农夫。因此，你行善，不要在这地之外，而要\BibleQuote{安居在这地上}。那么我将得到什么？\par

\Quote{你将在她的富饶中得到喂养。} 那地的富饶是什么？她的富饶就是她的主！她的富饶就是她的天主！正是祂，对祂说：\Quote{上主是我的产业和杯中的份。} 在最近的讲道中，亲爱的诸位，我们曾向你们提示：天主是我们的产业，同时我们也是天主的产业。请听，祂自己就是那地的富饶。\par

\BibleQuote{你要以上主为喜乐}，\BibleRef{咏 36:4} 好像你提出了问题，说：\Quote{请把那地的富饶指给我看，就是你叫我安居的那地。} 他说：\BibleQuote{你要以上主为喜乐。} \BibleRef{咏 36:4}\par

5.\BibleQuote{祂必将你心中的渴望赐给你。}要按正确的意义理解\BibleQuote{你心中的渴望}。要尽可能区分\BibleQuote{你心中的渴望}和肉身的渴望。在一篇圣咏中说：\BibleQuote{天主是我（心中的）力量。}这并非没有意义，因为下文接着说：\BibleQuote{天主永远是我的福分。}例如：有人遭受身体失明之苦。他祈求恢复视力。让他祈求吧；因为天主也做这事，也赐予这些恩惠。但连恶人也求这些事。这是肉身的渴望。有人生病，祈祷康复。从死亡边缘恢复健康。那也是肉身的渴望，凡属此类者皆是。什么是\BibleQuote{心中的渴望}？正如肉身的渴望是希望恢复视力，好能看见那种由眼睛可见的光；同样，\BibleQuote{心中的渴望}关乎另一种光。因为\BibleQuote{心里洁净的人是有福的，因为他们要看见天主。你们要因主而喜乐，祂必将你心中的渴望赐给你。}\par

6.\Quote{看}（你说），\Quote{我确实渴望它，我确实祈求它，我确实盼望它。那么我能完成它吗？}不。谁能？\BibleQuote{将你的路交托给主，信赖祂，祂必成就。}\BibleRef{咏 36:5}向祂诉说你所受的，向祂诉说你渴望的。因为你所受的是什么？\BibleQuote{肉身的私欲相反圣神，圣神的引导相反肉身。}\BibleRef{迦 5:17}那么你所渴望的是什么？\BibleQuote{我这个人真不幸！谁能救我脱离这该死的肉身？}\BibleRef{罗 7:24}因为是祂\Quote{自己}要\Quote{成就}，当你将你的路\Quote{交托给祂}之后；请听下文：\BibleQuote{藉我们的主耶稣基督的天主恩宠。}那么祂要成就的是什么呢？既然说：\Quote{将你的路交托给祂，祂必成就}？祂要成就什么呢？\par

\BibleQuote{祂必使你的正义如光出现}\BibleRef{咏 36:6}。因为现在，你的正义是隐藏的。现在属于信德的事，尚未看见。你相信一些事而实行。你尚未看见你所相信的。但当你开始看见先前所相信的时，\Quote{你的正义将出现在光中}，因为你的信德就是你的正义。因为\BibleQuote{义人因信德而生活。}\par

7.\BibleQuote{祂必使你的判断如同正午。}也就是说，\Quote{如同光明。}只说\Quote{如同光}还不够。因为我们称之为\Quote{光}，即使只是拂晓；我们称之为光，即使是日出之时。但正午的光最亮。因此，祂不仅\BibleQuote{使你的正义如光出现}，而且\BibleQuote{你的判断必如同正午}。因为现在你使你的判断跟随基督。这是你的目的：这是你的选择：这是你的判断。……\par

8.\Quote{那么我该做什么呢？}听你该做什么。\BibleQuote{顺服主，恳求祂}\BibleRef{咏 36:7}。让你的生活就是服从祂的诫命。因为这就是顺服祂；并且恳求祂，直到祂赐给你所应许的。让善行\Quote{持续}；让祈祷\Quote{持续}。因为\BibleQuote{人应当时常祈祷，不可灰心。}\BibleRef{路 18:1}你怎么证明你\Quote{顺服了祂}？在于做祂所命令的。但或许你还没有得到你的工价，因为你还不能。因为祂已经能够赐予，但你尚不能接受。在行动中操练自己。在葡萄园劳作；在一天结束时索要工钱。\BibleQuote{信实的是祂}那召唤你进入葡萄园的人。\Quote{顺服主，恳求祂。}\par

9.\Quote{看！我这样做；我确实\InnerQuote{顺服主，并且恳求祂。}但你怎么想？我的那个邻人是个恶人，生活堕落，却兴旺！他的偷窃、奸淫、抢劫，我都知道。他高抬自己，骄傲，因邪恶而上升，他不屑注意我。在这种情况下，我怎能忍耐得住？}这是一种病；喝下药方。\Quote{不要因那道路亨通者而烦躁。}他亨通，却是在\Quote{他的路上}；你受苦，却是在天主的路上！他的分是在他的路上兴盛，在路终时痛苦；你的分是在路上辛劳，在终点幸福。\BibleQuote{主认识义人的道路；恶人的道路必要灭亡。}你行走的是\BibleQuote{主认识}的道路，如果你在其中受苦辛劳，它们不会欺骗你。\Quote{恶人的道路}不过是暂时的幸福；在路的尽头幸福也结束了。为什么？因为那条路是宽路；它的终点通向地狱的火坑。现在，你的路是窄的；\BibleQuote{进去的人很少}\BibleRef{玛 7:13}；但最后它会进入何等广阔的田野，你应当思考。\BibleQuote{不要因那道路亨通者而烦躁，也不要因那图谋恶计的人而烦躁。}\par

\BibleQuote{息怒，弃绝忿怒}\BibleRef{咏 36:8}。你为何发怒？你为何因那激情与愤慨而亵渎，或几乎亵渎？反对\Quote{那成就恶谋的人，息怒，弃绝忿怒}。你岂不知那忿怒将你引向何处？你正要说天主不公义。它趋向于此。\Quote{看！为何那人亨通，而此人遭难？}思想它生出什么意念：扼杀那邪恶的念头。\Quote{息怒，弃绝忿怒}，使你如今恢复理智，可以说：\Quote{我的眼因忿怒而昏乱}。那是什么眼，岂不是信德之眼？我诉诸你的信德之眼。你信了基督：为何信？祂许诺了你什么？若基督许诺你的是今世福乐，那你就怨谤基督吧；是的！当你看见恶人兴隆时，就怨谤祂。祂许诺了什么福乐？什么福乐，除非在死人复活中？但在今世呢？是祂那份。祂的那份，我说！你，仆人和门徒，岂能嫌弃你的主、你的师傅所忍受的？……\par

\BibleQuote{因为作恶的必被剪除}\BibleRef{咏 36:9}。\Quote{但我看见他们亨通}。相信那位说的：\Quote{他们必被剪除}；祂看得比你更清楚，因为愤怒不能蒙蔽祂的眼。\Quote{因为作恶的必被剪除。但那等候主的，}——不是等候任何能欺骗他们的，而是实在是在于祂，那真理本身——\Quote{但那等候主的，他们必承受土地。}什么\Quote{土地，}岂非那耶路撒冷？凡被爱火点燃者，终将进入平安。\par

10. \Quote{但恶人要兴旺多久？我要忍耐多久？}你急躁不安；那在你看来似乎长久的事，很快就要过去。是软弱使得那实际上短暂的事显得长久，正如在病人的渴望中所见的那样。没有比为他调制解渴的药饮更显得漫长的了。尽管他身边的仆从竭力加快速度，免得病人发怒；\Quote{什么时候才能做好？（他喊道。）什么时候才能做好？什么时候才能端上来？}那些伺候你的人正在赶工，但你的软弱却以为那迅速完成的事是漫长的。因此，请看我们的医生如何迁就病人的软弱，说：\Quote{我要忍耐多久？还要多久？}\par

\BibleQuote{再过片刻，恶人就不存在了}\BibleRef{咏 36:10}。你果真是在罪人中，因罪人而抱怨吗？\Quote{再过片刻，他就不存在了。}免得因为我曾说：\Quote{那等候主的，他们必承受土地}，你就以为那等候必定非常长久。等候\Quote{片刻，}你将获得你所等候的无尽之物。片刻，一段中等的时间。回顾从亚当时代直到今日的岁月；通读圣经。他坠落出乐园几乎是昨天的事！如此多的世代已被度量并展开。过去那些世代如今在哪里？照样，剩下的少数世代也将过去。假如你从亚当被逐出乐园直到今天一直活着，你必定会看到那已飞逝的生命并不长久。但每个个体生命的长度是多少？任凭你加增年岁：将老年延长至极限：它是什么？岂不只是一阵晨风？然而，即使审判之日遥远，那时义人与不义之人的赏报将要来临：你的末日无论如何不会遥远。为此做好准备！因为当你离开此世时的样子，你将按你所是的样子被恢复到另一生命中。在那短促生命的终结，你还不会到达圣人们所在的地方，他们将要听到说：\BibleQuote{你们蒙我父祝福的，来承受为你们从创世以来预备的国度吧。}\BibleRef{玛 25:34}你还不会在那里吗？谁不知道呢？但你或许已经在那里，就是那个曾满身疮痍的乞丐，在远处安息，被那骄傲而不结果实的富翁在痛苦中看见的地方。确实隐藏在那安息中，你安心等待审判之日，那时你将重新获得身体，被改变以等同于天使。那么，我们所不耐烦并说\Quote{它何时来临？它会拖延很久吗？}的那事，到底有多久呢？我们的子孙后来说，我们的子子孙孙也会说；尽管他们一个接一个都会说同样的话，但那尚未来临的\Quote{片刻}也会过去，正如所有已过去的都已过去一样！哦，你这病人！\Quote{再过片刻，恶人就不存在了。是的，你要细察他的地方，却找不到他。}……\par

11. \BibleQuote{但谦卑的人必承受土地}\BibleRef{咏 36:11}。那土地就是我们多次讲过的圣城耶路撒冷，它将从她的这些朝圣之旅中得到解脱，永远与天主同在，并赖天主而生活。因此，\Quote{他们必承受土地。}他们的喜乐是什么？\Quote{他们必因丰富的平安而喜乐。}让那不虔敬的人在此地因他的大量金银、大量奴隶，最终因他的大量浴池、玫瑰、醉人的酒、最奢华丰盛的宴席而自得其乐。这是你所嫉妒的能力吗？这是令你喜悦的光荣吗？即使他永远如此，他的命运岂不令人哀叹吗？你的喜乐将是什么？\Quote{他们必因丰富的平安而喜乐。}平安将成为你的金银。平安将成为你的土地。平安将成为你的生命，你的天主——平安。平安将成为你所渴望的一切。……\par

\Emph{论圣咏的第二部分}\par

1. 接着是以下的话：\BibleQuote{恶人图谋攻击义人，向他咬牙切齿。}\BibleRef{咏 36:12}：\BibleQuote{但主要嗤笑他。}\BibleRef{咏 36:13}。嗤笑谁？当然是嗤笑那罪人，因为他\Quote{向他}……\Quote{咬牙切齿。}但主为什么\Quote{嗤笑他}？\Quote{因为主预见他的日子将要来到。}他看起来确实充满怒火，对自己即将面临的明天一无所知，正威胁着义人。但主却看见并\Quote{预见他的日子。}\Quote{什么日子？}就是\Quote{祂要照各人的行为报应各人}的日子。因为他正\BibleQuote{为自己积蓄忿怒，以致震怒的日子和天主公义审判的显示。}\BibleRef{罗 2:6}但预见这事的是主，你并未预见。是那预见者将此事启示给了你。你以前不知道\Quote{不义之人的日子，}就是他将受惩罚的日子。但那知道此事者已将它启示给你。知识的主要部分，就是使自己与那有知识者联合。祂有知识的眼目：愿你也有相信之心的眼目。天主\Quote{看见}的，愿你愿意相信。因为天主预见的不义之人的日子必要来到。那是什么日子？是执行一切报复的日子！因为报复必须加于不虔敬的人，报复必须加于不义的人，无论他悔改或不悔改。因为若他转离自己的道路，那正是他的\Quote{不义终结了}，这就是惩罚的施行……\par

2. \BibleQuote{恶人已经拔剑，张弓，要击倒穷困和贫乏的人，并杀害心正的人。}\BibleRef{咏 36:14}。\BibleQuote{他们的武器必进入自己的心。}\BibleRef{咏 36:15}。他的武器，即他的剑，要触及你的身体是容易的，正如迫害者的剑触及了殉道者的身体一样，但当身体被击中后，\Quote{心}却未受伤；但那\Quote{拔剑攻击}义人身体的人，他的心显然没有保持无伤。这正是这篇圣咏所见证的。它说，他们的武器，即\Quote{他们的剑要}，不是进入他们的身体，而是\Quote{他们的武器必进入自己的心。}他们本想杀害他的身体。让他们死于灵魂的死亡吧。因为对于那些他们企图杀害其身体的人，主已使他们免于忧虑，说：\BibleQuote{那杀身体而不能杀灵魂的，不要怕他们。}\BibleRef{玛 10:28}……\par

3. \Quote{他们的弓要被折断。}\Quote{他们的弓要被折断}是什么意思？他们的计谋将落空。因为上面祂曾说：\Quote{恶人已经拔剑，张弓。}通过\Quote{拔剑}他理解为公开的敌意；而通过\Quote{张弓}暗指秘密的阴谋。看！他的剑毁灭他自己，他的设陷阱也落空。落空是什么意思？即对义人没有造成损害。那么，例如（你问），它对那个人怎么没有损害呢？那人被这样剥夺了财物，因财产被夺而陷入困境？他仍有理由歌唱：\BibleQuote{义人所有的虽少，却胜过恶人的大量财富。}\BibleRef{咏 36:16}。\par

4. ……\BibleQuote{因为恶人的手臂要被折断，}\BibleRef{咏 36:17}。现在\Quote{他们的手臂}是指他们的权势。他在地狱里会做什么？会像那富豪必须做的吗？他在世上习惯于\Quote{奢华宴乐}，在地狱里却\Quote{受痛苦}？因此他们的手臂要被折断；\Quote{但主扶持义人。}祂如何\Quote{扶持}他们？祂对他们说什么？正如另一篇圣咏所说：\Quote{你要等待主，要勇敢，使你的心坚强；我说，你要等待主。}\Quote{等待主}是什么意思？你只受苦一时；你将永远安息；你的苦恼是短暂的；你的幸福是永久的。你只需\Quote{片刻}忧愁；你的喜乐将没有尽头。但在患难中，你的\Quote{脚}开始\Quote{滑跌}吗？基督受苦的榜样已摆在你面前。想想祂为你忍受了什么，而祂本身并无任何理由该忍受这一切。无论你的苦难有多大，你都不会遭受那些侮辱、那些鞭打、那件耻辱的袍子、那荆棘冠冕，以及最后祂所忍受的十字架；因为十字架现今已从人类刑罚中除去了。因为虽然在古人中罪犯被钉十字架，如今没有一个人被钉十字架。它被尊崇了，然后终止了。作为一种刑罚，它终止了；却在荣耀中继续存在。它从刑场移到了皇帝们的额上。那位使祂的苦难本身获得如此尊荣者，为祂忠信的仆人们还会保留什么？……\par

5. 但请注意，圣咏现在所说的话是否在他身上应验了。\Quote{主扶持义人。——不但如此}（同一位\Person{保禄}，虽忍受许多苦难，说：）\BibleQuote{我们也在患难中欢跃，因为知道患难生坚忍，坚忍生经验，经验生望德；望德不令人羞耻，因为天主的爱藉着所赐给我们的圣神，已倾注在我们心中。}\BibleRef{罗 5:3}现在他既是义人，又\Quote{被扶持}，这话说得恰当。因此，正如那些迫害他的人，在他如今\Quote{被扶持}时，不能伤害他，同样，他自己也没有伤害那些他所迫害的人。他说：\Quote{但主扶持义人。}……\par

6. 因此\Quote{主确实坚固义人。}祂如何坚固他们？\Quote{主知道无玷者的道路}（\BibleRef{咏 36:18}）。当他们受苦时，那些无知的人，那些没有分辨\Quote{无玷者的道路}知识的人，就以为他们在行恶道。那位\Quote{知道那些道路}的祂，知道用什么方式引导祂的\Quote{温良者}走正路。因此，在另一篇圣咏中祂说：\Quote{祂要引导温良者行正道，教训温良者走祂的路。}你以为那位躺在富人家门前、满身疮痍的乞丐（\BibleRef{路 16:20}），过路人是如何藐视他的！他们很可能掩鼻而唾弃他！然而，主知道如何为他保留乐园。另一方面，他们自己多么渴望那个\Quote{身穿紫红袍和细麻衣，天天奢华宴乐}（\BibleRef{路 16:19}）之人的生活！但主预见了那人的\Quote{日子来到}，知道等待他的是无尽的痛苦。因此\Quote{主知道义人的道路。}\par

7. \Quote{并且他们的产业直到永远}（\BibleRef{咏 36:18}）。这是我们所持守的信德。主也是凭信德知道吗？主知道这些事是那样清楚显明，甚至当我们成为与天使相等时，也无法言说。因为那将显明给我们的事，对我们而言不会像现在对那位不变者那样清楚。然而，关于我们本身，经上怎么说？\Quote{亲爱的，如今我们是天主的儿女，将来如何，还未显明；但我们知道，当祂显现时，我们要像祂，因为我们要看见祂实在怎样}（\BibleRef{若一 3:2}）。因此，确实有某种幸福的景象为我们保留；如果现在能\Quote{藉着镜子，模糊地}（\BibleRef{格前 13:12}）有所领略，那么我们就完全无法用语言表达那福乐的美妙，就是天主为敬畏祂的人所保留、在那些仰望祂的人身上所成全的。正是为了这个目的，我们的心在今世的一切烦恼和试探中受磨练。不要奇怪你是在患难中为此受磨练。你是为某种荣耀而受磨练。如今坚固的义人所说的话哪里来的：\Quote{现今的苦楚，与将来在我们身上要显示的光荣，是不能相比的}（\BibleRef{罗 8:18}）？那应许的荣耀是什么，不就是成为与天使相等并看见天主吗？祂赐给盲人多大的恩惠，就是使他的眼睛康复，能看见今生之光……那么，我们该给那位医治我们内眼、使其能看见永恒之光——祂自己——的医生什么报酬呢？……\par

8. \Quote{在凶恶之时，他们必不蒙羞}（\BibleRef{咏 36:19}）。在患难之日，在窘迫之日，他们不会\Quote{羞惭}，像那些希望落空的人那样羞愧。谁是那个\Quote{羞惭}的人？是那说\Quote{我没有得到我所希望的}的人。这也是罪有应得；因为你是从你自己或从你的朋友那里指望它。但\Quote{凡信赖世人者，是可咒骂的}（\BibleRef{耶 17:5}）。你羞愧，因为你的望德欺骗了你；你的望德建立在谎言上。因为\Quote{众人都是说谎者}。但如果你把望德寄托于你的天主，你就不会\Quote{羞惭}。因为你所信靠的那位，不会受骗。因此，我们刚才提到的那个人，那如今\Quote{坚固}的义人，当陷入凶恶之时，在患难之日，他说了什么来表示他没有\Quote{羞惭}？\Quote{我们以患难为荣，因为知道患难生忍耐，忍耐生老练，老练生望德，望德不叫人羞惭。}望德为什么不叫人\Quote{羞惭}？因为它建立在天主身上。因此紧接着说：\Quote{因为天主的爱，藉着所赐给我们的圣神，已倾注在我们心中了}（\BibleRef{罗 5:3}）。圣神已经赐给了我们；祂既已给了我们这样的\Quote{凭据}，又怎能欺骗我们？\Quote{在凶恶之时，他们必不蒙羞，在饥荒的日子，他们必得饱饫}。……\par

9. \Quote{因为恶人必丧亡。主的仇敌，当他们开始光荣和高傲时，就要立即消灭，如同烟雾一样}（\BibleRef{咏 36:20}）。从这比喻本身，你就认出他所暗示的事。烟雾从火源处升起，升到高处，并因升起而膨胀成一大团；但这一团越大，就越空虚；因为这团体积既无根基也无实质，只是松散、飘移、易逝，它就散入空中消失；所以你会看到它的庞大反而对它致命。因为它升得越高，扩展得越远，散布的范围越广，就越稀薄、越稀散、越易消失和消散。\Quote{主的仇敌，当他们开始光荣和高傲时，就要立即消灭，如同烟雾一样。}关于这样的人，经上说：\Quote{就如雅乃斯和杨布勒曾反对梅瑟，照样这些人也反对真理；他们是心思败落的人，在信德上是可弃绝的}（\BibleRef{弟后 3:8}）。但他们如何反对真理呢？还不是藉着他们骄傲自大的空虚膨胀，把自己抬高，好像伟大而正义的人，虽然即将化为虚空？但关于他们，他又说了什么？如同说到烟雾：\Quote{他们的愚昧将要显露给众人，如同那两个人的一样。}……\par

10. \BibleQuote{恶人借贷，而不偿还}\BibleRef{咏 36:20}。他领受了，却不归还。他不愿归还的是什么？是感恩。因为天主愿意从你得到什么，祂要求你什么，除非祂向你施恩？然而罪人领受了如此多的恩惠，却不愿归还！他领受了存在的恩赐，领受了成为人的恩赐，并且是远高于禽兽的存在；他领受了身体的形式，以及身体感官的区分——眼用以看，耳用以听，鼻用以嗅，舌用以尝，手用以触摸，脚用以行走；甚至还有身体的健康和强健。但到目前为止，这些我们甚至与禽兽共有；他还领受了更多：能够理解的理智，能够认识真理，能够分辨是非；能够寻求并渴望其造物主，赞美祂，并依附于祂。这一切，恶人也与其他人一样领受了；但因生活不善，他不偿还所欠的债务。因此，\BibleQuote{恶人借贷，而不偿还：}他不愿报答他领受自谁的那一位；他不愿感恩；甚至以恶报善，亵渎、抱怨天主、愤恨。因此他\BibleQuote{借贷，而不偿还；但义人却慈悲为怀，且施舍}\BibleRef{咏 36:21}。因此，前者一无所有；后者则富足。看，一边是匮乏；看，另一边是财富。前者领受而\BibleQuote{不偿还：}后者\BibleQuote{慈悲为怀，且施舍：}而且他绰绰有余。倘若他是穷人呢？即使如此，他仍是富足的：你只以信仰之眼来看他的财富。因为你看见空空的箱子；却看不见那充满天主的良心……\par

11. \BibleQuote{因为凡祝福祂的，必承受土地}\BibleRef{咏 36:23}，也就是说，他们必拥有那位义者：那唯一既是真正义人，又使人成义者；祂曾在世上贫穷，却带来巨大财富，用以使祂发现为贫穷的人富足。因为祂曾以圣神使穷人的心富足，并因他们承认罪过而倒空他们的灵魂，又以正义的富足充满他们：祂曾使渔夫富足，渔夫舍弃渔网，轻看已拥有的，却寻求捕捞那尚未拥有的。因为\BibleQuote{天主拣选了世上软弱的，以羞辱那强壮的}\BibleRef{格前 1:27}。祂不是藉演说家而赢得渔夫归向祂；而是藉渔夫赢得演说家归向祂；藉渔夫赢得元老；藉渔夫赢得皇帝。因为\BibleQuote{凡祝福祂的，必承受土地}；他们必与祂同为继承人，在那\BibleQuote{活人之地}——另一篇圣咏论及此地时说：\BibleQuote{祢是我的希望，我在活人之地的福分。}……\par

12. 注意下文：\BibleQuote{善人的脚步由主安排，他喜爱祂的道路}\BibleRef{咏 36:23}。为使那人自己\BibleQuote{喜爱主的道路，}他的脚步由主亲自安排。因为如果主不安排人的脚步，那么他天生如此扭曲，将永远行走在弯曲的小径上，因追逐弯曲的道路而无法回转。然而祂来了，呼召我们，救赎了我们，倾流了祂的血；祂付了这赎价；祂行了这善，忍受了这些恶。以祂所行的事来看，祂是天主！以祂所受的苦来看，祂是人！那位神人是谁？难道不是因为你，人啊，离弃了天主，天主才为你而成为人吗？因为祂使你成为人，这对你而言已是过于微小，不足以回报，或对祂而言过于微小，不足以赐予；除非祂自己也为你成为人。因为正是祂\BibleQuote{安排了我们的脚步；}使我们\BibleQuote{喜爱祂的道路。}……\par

13. 倘若人一生都在劳苦、受苦、疼痛、折磨、监禁、鞭打、饥饿、干渴中度过，每日每时，整个生命长度直到老年，然而人的整个生命不过寥寥数日。那劳苦过后，将临到永恒的王国；将临到无尽的幸福；将临到与天使同等；将临到基督的产业，并且基督，我们的\BibleQuote{同继承人}\BibleRef{罗 8:17}，将要来临。这劳苦何等重大，而你所领受的酬报何等伟大？那些在战争中服役、多年在创伤中辗转的退伍军人，从青年起入伍，到老年退役：为获得晚年几天的安息——此时年老本身就开始压垮那些未被战争摧毁的人——他们忍受何等巨大的艰辛！何等行军，何等霜冻，何等烈日，何等匮乏，何等创伤，何等危险！而忍受这一切时，他们所思虑的不过是晚年那几天的安息，他们并不知道能否活到那时。同样，\BibleQuote{善人的脚步由主安排，他喜爱祂的道路。}这正是我开始时所说的要点。如果你\BibleQuote{喜爱}基督的\BibleQuote{道路，}并且是真正的基督徒（因为真正的基督徒，是不轻视基督的道路，而\BibleQuote{喜爱}跟随基督\BibleQuote{的道路}历经祂的苦难），你就不要走任何其他道路，只走祂亲自走过的道路。那路看似痛苦，却正是得救之路；另一条路也许令人愉快，却布满强盗。\BibleQuote{他喜爱祂的道路。}\par

14. \Quote{他虽失足，却不致倾覆，因为主用手扶持他。}\BibleRef{咏 36:24}。请看什么叫做\Quote{喜爱}基督的\Quote{道路。}假若他遭受一些磨难，丧失名誉、痛苦、损失、侮辱，或人生中常遇到的其他种种不幸，他便把主放在自己眼前，想想祂曾受过怎样的试探！并且，\Quote{他虽失足，却不致倾覆，因为主用手扶持他，}因为祂已在他之前受过苦。人哪，你的脚步被安排得如此，使你要\Quote{喜爱主的道路}，你还有什么可怕的呢？你怕什么？痛苦？基督受过鞭打。你怕侮辱？祂曾被责骂说：\Quote{你附有魔鬼，}而祂自己正驱逐魔鬼。或许你害怕党派的阴谋和恶人的合谋？阴谋也曾针对祂。你在某次控告中无法显明你良心的纯洁，反而遭受冤屈和暴力，因为有人听信假见证反对你。假见证首先针对祂，不但在祂死前，甚至在祂复活之后……。\par

\Emph{论圣咏的第三部分。}\par

1. \Quote{我年幼时，如今已年老，从未见过义人被弃，亦未见他的后裔求乞食粮。}\BibleRef{咏 36:25}。\par

若这话仅由一人说出，人的一生有多长？又有什么可惊奇的？一个人固定在大地的某处，在他整个一生——如此短暂的人生中，从未见过\Quote{义人被弃，亦未见他的后裔求乞食粮，}尽管他已从幼年到了老年。这并不值得惊奇；因为可能在他生前，曾有过一些\Quote{义人求乞食粮；}也可能在大地的另一处，并非他所在之处，曾有过这样的人。再听另一件令我们印象深刻的事。你们中间任何一位现已年老的人，或许在回顾自己一生的经历时，翻阅他所认识的人，却发现不了任何义人求乞食粮、或他的后裔求乞食粮的实例；然而他转而查阅受启示的圣经，却发现正义的\Person{亚巴郎}曾在自己的家乡受困挨饿，并离开那地前往别处；他也发现同一位的儿子\Person{依撒格}，因同样的饥饿原因移居他乡寻找食粮。那么，怎能说：\Quote{我从未见过义人被弃，亦未见他的后裔求乞食粮？}若他在自己一生中确实如此，却发现圣经的记载与此不同，而圣经比人的生命更可信。\par

2. 那么，我们该怎么办？愿你们虔诚的关注帮助我们，使我们能辨别天主在这圣咏的经句中的旨意，祂要我们如何理解它们。因为有一种惧怕，恐怕任何不稳定的、不能属灵地理解圣经的人，会诉诸人的经验，并观察到天主忠实的仆人有时处于某种匮乏和缺乏中，以致被迫乞食；他会特别想起宗徒保禄，他说：\BibleQuote{受饥受渴；挨冻露体；}\BibleRef{格后 11:27} 并因此绊倒，心里说：\Quote{我所歌唱的难道是真的吗？我站在圣堂里以如此虔诚的声音所宣告的难道是真的吗？\InnerQuote{我从未见过义人被弃，亦未见他的后裔求乞食粮。}} 免得他心中说：\Quote{圣经欺骗我们；} 于是他全身对善工瘫痪了：当他在内的这些肢体，内在之人的肢体瘫痪了（这是更可怕的瘫痪），他从此便放弃善工，对自己说：\Quote{我为何行善？我为何把食物分给饥饿的人，给赤身露体者衣服穿，将无处安身的人领回自己家中，\BibleRef{依 58:7} 相信所写的经文呢？\InnerQuote{我从未见过义人被弃，亦未见他的后裔求乞食粮；} 而我却看到许多生活端正的人，大多却忍受饥饿。但若我或许弄错了，以为生活好的人和生活不好的人都生活得好，而天主知道并非如此；也就是说，祂知道我认为正直的人其实是不义的，那么我该如何看待亚巴郎呢？圣经本身称赞他为义人。我该如何看待宗徒保禄呢？他说：\InnerQuote{你们效法我，如同我效法基督一样。}\BibleRef{格前 11:1} 什么？难道我自己要遭受像他所忍受的那些苦难，\InnerQuote{受饥受渴；挨冻露体；}？}\BibleRef{格后 11:27}\par

3. 因此，当他这样想时，他全身对行善的能力瘫痪了，弟兄们，我们能像抬起瘫子一样，并像\Quote{揭开房顶}一样揭开这圣经，把他放到主面前。\BibleRef{路 5:19} 因为你们看出这是隐晦的。若隐晦，便是被遮盖的。我看到一个精神上的瘫子，我也看到这个房顶，并且相信基督就隐藏在房顶之下。让我尽量去做那些被称赞的人所做的事：他们揭开房顶，把瘫子放到基督面前；使祂能对他说：\BibleQuote{孩子，你放心，你的罪赦了。}\BibleRef{路 5:20} 因为正是这样，祂治好了他内在之人的瘫痪，解除了他的罪，坚固了他的信德……。\par

4. 但是谁是\Quote{义人}，谁\Quote{从未见过被遗弃，也未见过他的后裔讨饭}？如果你明白\Quote{饼}是什么意思，你就明白他是指谁。因为\Quote{饼}就是天主的圣言，它从不离开义人的口……现在看，如果\Quote{神圣的默想\InnerQuote{保守你}}在于反复咀嚼这饼，那么\Quote{你从未见过义人被遗弃，也未见过他的后裔讨饭}。\par

5. \Quote{他终日怜悯，借贷} \BibleRef{咏 36:26}。拉丁文中的\OriginalTerm{放债}{Fœneratur}一词，确实既指借出之人，也指借入之人。但在此段经文中，若用\Quote{\OriginalTerm{借贷}{fœnerat}}来表达，意义更为清晰。语法学家喜欢怎样规定，与我们何干？我们宁可用一个粗鄙之词，好让你们明白，也不愿因措辞文雅而使你们不得要领。因此，那\Quote{义人终日怜悯，并（fœnerat）借贷}。然而，放高利贷者不可欢喜。因为我们发现，这里所说的是某种特定的借贷，正如前面所说的某种特定的饼一样；好叫我们在所有段落中都\Quote{拆开屋顶}，找到通往基督的道路。我不愿你们放高利贷；我不愿你们这样做，因为天主也不愿你们这样做……从哪里可以看出天主不愿如此？在另一处经上说：\Quote{不放他的银钱取利}。而且，放高利贷是多么可憎、可厌、可咒之事，我相信连放高利贷者自己也知道。另一方面，我本人，不，我们的天主反而命令你做放债者，并向你说：\Quote{借贷给天主}。如果你借贷给人，你还有希望？如果你借贷给天主，难道竟没有希望？如果你把钱以高利贷借给人，也就是说，你把钱借给一个人，期望从他那里得到比你给出的更多的东西，不单是钱，而是任何东西，无论是麦子、酒、油，或其他什么，如果你期望收到比你给出的更多，你就是放高利贷者，在这件事上不值得称赞，只应受责备。\Quote{那么，}你说，\Quote{我该怎样做，才能\InnerQuote{有利地}借贷？}想一想放高利贷者做什么。他无疑想少给多收；你也这样做：少给多收。看你的本金如何增长！给出\Quote{暂时之物}，收到\Quote{永恒之物}：给出尘世，收到天国！也许你会说：\Quote{我给谁呢？}那曾命令你不要放高利贷的主，他自己就成为你放债的对象！听圣经如何说你可以\Quote{借贷给主}：\Quote{怜悯穷人的，就是借贷给主}。\BibleRef{箴 10:17}因为主不需要你任何东西。但你有一个需要你帮助的人：你伸手给他，他接受。因为穷人无法偿还你，他本想报答你，却找不到可还之物；他所能做的，只有为你祈祷的善意。当穷人向你祈祷时，他仿佛对天主说：\Quote{主，我借了这些，请为我作保。}于是，虽然你对穷人的偿还没有借据，却有一个可靠的担保。看，天主从祂自己的圣经中对你说：\Quote{给吧，不要怕；我偿还。你是给我的。}那些为别人作保的人怎么说？他们说：\Quote{我偿还；我承担。你是给我的。}那么，我们以为天主也这么说：\Quote{我承担；你是给我的}？确实，如果基督是天主——这是无疑的——祂亲自说过：\Quote{我饿了，你们给了我吃。}\BibleRef{玛 25:35}当他们对祂说：\Quote{主，我们几时见了你饥饿？}\BibleRef{玛 25:37}祂为表明自己是穷人的担保人，为祂的所有肢体负责，祂是头，他们是肢体，当肢体接受时，头也接受；祂说：\Quote{你们对我这些最小兄弟中的一个所做的，就是对我做的。}\BibleRef{玛 25:40}来，贪婪的放高利贷者，想想你给了什么，想想你将收到什么。如果你给了一小笔钱，而收钱的人因这小数给了你一座大庄园，价值远超过你所给的，你会多么感谢，多么欢喜！听那借给你的一位赐予什么产业：\Quote{我父所祝福的，你们来承受}\BibleRef{玛 25:34}——承受什么？他们所给的吗？断乎不是！你所给的是地上的东西，即使你不给，它们也会在地上朽坏。因为如果你不给，你能拿它们做什么？在地上会失去的，在天上被保存了。所以，我们要承受的就是那被保存的。你的功德被保存了，你的功德成了你的宝藏。因为想一想你要承受的是什么：承受——\Quote{从创世以来为你们预备的王国。}另一方面，那些不愿\Quote{借贷}的人将得到什么判决？\Quote{你们到为魔鬼和他的使者预备的永火里去。}\BibleRef{玛 25:41}我们所承受的王国叫什么？看下文：\Quote{这些人要进入永罚，而义人却要进入永生。}\BibleRef{玛 25:46}为这事取利；赚取这利。把你的钱放债以赚取这利。你有基督在天上坐于宝座，在地上乞讨。我们发现了义人如何借贷：\Quote{他终日怜悯，借贷}。\par

6. \Quote{他的后裔是有福的。}这里也不要让任何肉性的观念出现。我们看到许多义人的儿子死于饥饿；那么在何种意义上他的后裔将是有福的呢？他的后裔是他以后所留下的；是他在此播种、日后收割的。因为宗徒说：\Quote{我们行善不要厌倦；若不松懈，到了适当时节我们就要收割。所以我们一有机会，}他说，\Quote{就该向众人行善。}\BibleRef{迦 6:9}这就是你的那个\Quote{后裔}，它将\Quote{有福}。你将之交给大地，就能收获更多；难道你交给基督反而会失去它吗？请看宗徒在谈论施舍时特意称之为\Quote{种子}。因为他这样说：\Quote{那少种的，也要少收；那怀着祝福种的，也要在祝福中收获。}\BibleRef{格后 9:6}……\par

7. 因此，注意下文，不要懈怠。\Quote{远离恶，行善}\BibleRef{咏 36:27}。不要以为你不剥去已经穿着衣服的人的衣服就足够了。因为你不剥去已经穿着衣服的人的衣服，你确实\Quote{远离了恶}；但不要成为不结果实的，因而枯萎。因此，你宁可选择不剥去已经穿着衣服的人的衣服，也要去给赤身露体的人穿上衣服。因为这就是\Quote{远离恶，行善}。你会说：\Quote{我能从中得到什么好处？}你所借贷的那一位已经向你保证了他将给予你什么。他将赐予你永生。给予他吧，不要害怕！再听下文：\Quote{远离恶，行善，并永远居住。}当你施舍时，不要以为没有人看见你，或天主遗弃了你；当你施舍给穷人后，或许会有一些损失，或对所失之物的忧伤随之而来，你对自己说：\Quote{我行善有什么益处？我相信天主不爱行善的人。}那种嗡嗡声、那种低沉的抱怨声从何而来，岂不就是因为这些说法很常见吗？你们每个人此刻都认出了这些说法，要么是在自己的口里，要么是在朋友的口里。愿天主摧毁它们；愿祂从自己的田地里拔除荆棘；愿祂栽植\Quote{好种子}和\Quote{结果实之树}！因为人哪，你为何忧伤？你给了穷人一些东西，却失去了其他东西？难道你看不见你所失去的是你没有给予的吗？你为什么不听从你的天主的声音？你的信德在哪里？为什么它如此沉睡？在你心中唤醒它吧。想一想主自己对你所说的话，祂鼓励你做这样的善行：\Quote{要为自己预备永不腐朽的钱袋，在天上积蓄用不完的宝藏，那里没有贼偷窃。}\BibleRef{路 12:33}因此，当你为损失哀叹时，要记起这一点。你为什么哀叹，你这小信的人，更确切地说，你这心不健全的人？你为什么失去它？除非因为你没有把它借给我。你为什么失去它？谁把它拿走了？你会回答：\Quote{一个贼。}难道这不是我所预先警告你的吗？要你把它存在贼不能靠近的地方。那么，如果任何失去东西的人忧伤，让他为此忧伤：他没有把它存在那里，使它不能失去。\par

8. \Quote{因为主喜爱公义，不遗弃祂的圣徒}\BibleRef{咏 36:28}。当圣徒受苦时，不要以为天主不审判，或不公义地审判。那警告你要公义审判的祂，会不公义地审判吗？祂\Quote{喜爱公义，不遗弃祂的圣徒}。但是（要想到）圣徒的\Quote{生命}如何\Quote{与祂一同隐藏}，以致他们如今在地上受苦，如同冬天的树木，没有果实和叶子；当祂像新升的太阳显现时，那先前活在他们根中的生命，将在果实中显出来。祂确实\Quote{喜爱公义，不遗弃祂的圣徒}……\par

9. \Quote{但不义的人将受惩罚；恶人的后裔将被剪除。}正如前者的\Quote{后裔有福}，同样\Quote{恶人的后裔将被剪除}。因为恶人的\Quote{后裔}就是恶人的作为。另一方面，我们又看到恶人的儿子在世界中兴旺，有时也成为义人，并在基督内兴旺。因此，要小心你如何理解它；好让你揭开覆盖，直达基督。\BibleRef{路 5:19}不要以肉性的方式理解经文，否则你会受骗。但是\Quote{恶人的后裔}——恶人的一切作为——\Quote{将被剪除}：它们将没有果实。因为它们在短时间内确实有效；但日后他们将寻找它们的果实，却找不到他们所行所为的报偿。因为那节经文表达了那些失去他们所行所做之人的话：\Quote{骄傲给了我们什么好处？夸耀中的财富又带来了什么？这一切都像影子一样消逝了。}\BibleRef{智 5:8}那么，\Quote{恶人的后裔将被剪除}。\par

10. \BibleQuote{义人将承受土地} \BibleRef{咏 36:29}。这里再不要让贪欲悄悄地侵蚀你，也不要许诺你什么大的产业；不要希望在世上找到你所被命令轻视的东西。经文中的\Quote{土地}，是某种\Quote{生命之地}，即圣人之国。因此经上说：\BibleQuote{你是我的希望，在生命之地中我的份。}因为如果你的生命也是那里所说的同一生命，那么想一想你将承受的是何种\Quote{土地}。那是\Quote{生命之地}；这是那些即将死去之人的土地：当它们死去时，再次接纳那些在活着时滋养它们的东西。因此，正如那土地是怎样的，生命本身也将是怎样的：如果生命是永恒的，那么\Quote{土地}也将是你们的，直到永远。而\Quote{土地}如何能是你们所有，直到永远？\par

\BibleQuote{他们将在那里居住}（经上说）\BibleQuote{直到永远。}因此，那必定是另一块土地，他们在那里\BibleQuote{永远居住。}因为论及这土地（这个世界）经上说：\BibleQuote{天地要过去。} \BibleRef{玛 24:35}\par

11. \BibleQuote{义人的口讲论智慧} \BibleRef{咏 36:30}。看，这就是那\Quote{食粮}。请观察这位义人以何等满足来喂养它；他是如何在口中反复品尝智慧。\BibleQuote{他的舌头谈论审判。}\par

\BibleQuote{天主的法律在他心中} \BibleRef{咏 36:31}。免得你以为他只是嘴上说说，心里却没有；免得你把他算作那种人，如经上所说：\BibleQuote{这民族用嘴唇尊敬我，他们的心却远离我。} \BibleRef{依 29:13} 这对他有什么益处呢？\par

\BibleQuote{他的脚步必不滑倒。} \Quote{天主的话在心中}使人脱离陷阱；\Quote{天主的话在心中}使人脱离恶道；\Quote{天主的话在心中}使人脱离\Quote{滑脚之处。}那话不离弃你的主与你同在。现在，天主所保守的人会遭受什么恶事呢？你在葡萄园里安置一个看守，就以为安全不受盗贼侵扰；但那看守可能睡觉，可能自己跌倒，可能放进盗贼。但\BibleQuote{保护以色列的那位，既不打盹也不睡觉。} \BibleQuote{天主的法律在他心中，他的脚步必不滑倒。}因此，让他无所畏惧地生活；即使在不虔敬的人中间也无所畏惧；即使在不虔敬或不义的人中间也无所畏惧。因为不虔敬或不义的人能对义人做什么恶事呢？看哪！请看下文。\par

\BibleQuote{恶人窥伺义人，设法杀死他} \BibleRef{咏 36:32}。因为他所说的话，正如\BibleBook{智慧}{篇}中预言他将会说的：\BibleQuote{他实在使我们生厌，连看都不愿看他，因为他的生活不像别人一样。} \BibleRef{智 2:15}。因此他\Quote{设法杀死他。}怎么？那保护他、与他同住、不离弃他的嘴唇和心灵的主，难道会舍弃他吗？那么，先前所说的话又怎么样呢：\Quote{他决不遗弃他的圣人}？\par

12. \BibleQuote{恶人窥伺义人，设法杀死他} \BibleRef{咏 36:33}。那么，祂为什么把殉道者交在不虔敬的人手中呢？为什么他们对待他们\Quote{任意而行}呢？ \BibleRef{玛 17:12} 有的被刀杀；有的被钉十字架；有的被交给野兽；有的被火烧；有的被链子牵着，直到因长期衰弱而耗尽。的确，\Quote{主不遗弃他的圣人。}祂必不\Quote{把他交在他手中。}最后，祂为什么把自己的圣子交在\Quote{不虔敬的人手中}呢？这里也是如此，如果你要使你内在之人的所有肢体坚强起来，就掀开屋顶的覆盖物，找到通往主的道路。请听另一处圣经，它预见到我们的主将来要被不虔敬的人苦害，是怎么说的。它说：\BibleQuote{大地被交在恶人手中。} \BibleRef{约 9:24} \Quote{大地} \Quote{被交在不虔敬的人手中}是什么意思？是将肉身交到迫害者手中。但天主没有把他的\Quote{义者}留在那里：从被掳的肉身中，祂带领灵魂未受征服地出来……\par

\BibleQuote{主必不把他交在他手中，也不在审判他的时候定他的罪} \BibleRef{咏 36:33}。有些抄本作：\Quote{当他审判他时，审判将是为他的。}然而，\Quote{为他}是指当判决对他宣告时。因为我们表达时可以对人说：\Quote{请为我审判}，即\Quote{听我的诉讼。}因此，当天主开始听他义仆的诉讼时，既然\BibleQuote{我们都必须到基督的审判台前，站在它面前，各人按他本身所行的} \BibleRef{格后 5:10} 或善或恶领受报应，那么当他来到那审判时，祂必不定他的罪；尽管他今生可能被人定罪。即使总督对\Person{西彼廉}宣判了，但地上的审判席是一回事，天上的法庭是另一回事。从低级的法庭他接受死刑判决；从高级的法庭他接受冠冕。\Quote{在审判他的时候，他也不定他的罪。}\par

13. \BibleQuote{你应仰望上主} \BibleRef{咏 36:34}。当我在等候祂时，我该做什么？——\BibleQuote{并遵守祂的道。} 我若遵守，将获得什么？\Quote{祂必高举你，使你得地业。} \Quote{什么地业？} 请不要再让任何地产浮现在你心中：——那是经上所载的地业：\BibleQuote{你们蒙我父降福的，来承受自创世以来为你们预备的国度吧。} \BibleRef{玛 25:34} 至于那些困扰我们的人，我们曾在他们中间叹息，我们曾耐心忍受他们的绊脚石，当他们对我们发怒时，我们曾徒然为他们祈祷？他们将如何？下文是什么？\Quote{当恶人被剪除时，你必看见。}……\par

\BibleQuote{我见过恶人横跨高空，高耸如黎巴嫩的香柏} \BibleRef{咏 36:35}。假设他被\Quote{高举}；假设他高耸于\Quote{众人之上}；下文如何？\par

\BibleQuote{我经过，看啊，他已不在！我寻找他，却找不到他的踪影！} \BibleRef{咏 36:36}。为何他\Quote{已不在，踪影无处可寻}？因为你\Quote{经过}了。但你若仍属血肉，仍以为世俗的顺遂是真正的幸福，你便尚未\Quote{经过}他；你或是与他同类，或是在他之下。前进吧，超越他；当你进步，经过他之后，你以信德之眼观察他：你看到他的结局，便对自己说：\Quote{看啊！那个曾如此膨胀的人，如今不在了！} 就好像你经过一团烟云。因为这同样在这篇圣咏前面说过：\Quote{他们必如烟云消散。}……\par

14. \BibleQuote{持守纯全} \BibleRef{咏 36:37}；持守它，正如你从前贪婪时，持守你的钱袋；正如你曾紧握那钱袋，免得被贼人夺去；同样，\Quote{持守纯全}，免得它被魔鬼从你手中夺去。让它成为你可靠的基业，富人和穷人都能同样保有。\Quote{持守纯全。} 得金子而失纯全，对你有何益处？\par

\Quote{持守纯全，并注重正直的事。} 你当保持你的眼睛\Quote{正直}，好能看见\Quote{正直的事}；不要扭曲，不要用你注视着恶人的那种眼神；不要歪斜，以免天主在你眼中显得扭曲不公，似乎祂偏袒恶人，而用迫害来磨炼忠信者。你难道没看到自己的视线是如何歪曲的吗？扶正你的眼睛，\Quote{注视那正直的事。} 什么是\Quote{正直的事}？不要在意眼前的事。那么你将看见什么？\par

\Quote{因为为缔造和平的人，存有后备。} \Quote{存有后备}是什么意思？你死了，你并不会死。这就是\Quote{存有后备}的意思。他在此世之后仍将有某物存留给他，那就是\Quote{后裔}，它将\Quote{蒙福}。因此我们的主说：\BibleQuote{凡信我的，即使死了，仍要活着；} \BibleRef{若 11:25}——\Quote{看，为缔造和平的人存有后备。}\par

15. \BibleQuote{但悖逆的人在同一事上将被消灭} \BibleRef{咏 36:38}。\Quote{在同一事上}是什么意思？意思是永远；或是一同在同一个毁灭中。\par

\Quote{恶人的余孽必被剪除。} 如今\Quote{（有余地）为缔造和平的人}；因此，不缔造和平的人就是恶人。因为\BibleQuote{缔造和平的人是有福的，因为他们要称为天主的子女。} \BibleRef{玛 5:9}\par

16. \BibleQuote{但义人的救恩来自上主，祂是他们患难时的力量} \BibleRef{咏 36:39}。\BibleQuote{上主必帮助他们，拯救他们；祂必拯救他们脱离恶人} \BibleRef{咏 36:40}。因此，现在让义人容忍罪人；让麦子容忍莠子；让谷粒容忍糠秕：因为分离的时候将至，好种子将与火焚的杂质分开。\BibleRef{玛 13:30} 一种将被收进仓里，另一种则投入\Quote{永火}；因为义人与不义者起初放在一起，正是为让一方成为绊脚石，一方经受考验；而后一方受罚，一方得冠冕……\par

\SourceRecord{Translated by J.E. Tweed.；Nicene and Post-Nicene Fathers, First Series；Vol. 8.；Edited by Philip Schaff.；Buffalo, NY: Christian Literature Publishing Co.,；1888.；<http://www.newadvent.org/fathers/1801037.htm>.}

% structure: p0001 source=h1 target=section reason=page-title

\section{圣咏第三十八篇释义}

达味自己的一首圣咏，论及安息日的纪念。\par

这安息日的纪念是什么意思？这安息日是什么？因为他是在呻吟中\Quote{回想}安息日。你们在聆听圣咏时已经听见了，当我们诵读它时还要听见，他的呻吟、哀叹、眼泪和痛苦是多么深重。但以这种方式悲苦的人是有福的！因此主也在福音中称一些哀恸的人为有福的（\BibleRef{玛 5:4}）。\Quote{哀恸的人怎能有福？悲苦的人怎能有福？}不，正相反，他若不哀恸，才是悲苦的。那么，我们也应当如此理解这里的人，他记念安息日（即），某个哀恸的人；但愿我们自己也成为那样的\Quote{某个人}！因为这里有一个悲伤、呻吟、哀恸的人，他在记念安息日。安息日就是安息。无疑，他是在某种不安中，带着呻吟记念安息日……\par

2. \Quote{上主，求祢不要在祢的震怒中责斥我，也不要在祢的怒火中惩戒我}（\BibleRef{咏 37:1}）。因为将会有人在天主的\Quote{怒火}中被惩戒，在祂的\Quote{震怒}中被责斥。或许并非所有被\Quote{责斥}的人都会被\Quote{惩戒}；然而，在惩戒中得救的人也是有的。的确如此，因为称为\Quote{惩戒}，但却是\Quote{如同经过火一样}。但也有一些人会受\Quote{责斥}，却不会受\Quote{纠正}。因为祂必定要\Quote{责斥}那些祂将对其说：\Quote{我饿了，你们没有给我吃的}的人（\BibleRef{玛 25:42}）。……\Quote{也不要在祢的怒火中惩戒我}；好使祢在此生中洁净我，使我成为那样的，以后不再需要炼净之火，因为那些\Quote{将要得救，却如同经过火一样}的人（\BibleRef{格前 3:15}）。为什么？为什么呢？只因为他们\Quote{在根基上建造了木、草、禾秸}。他们原本应当在其上建造\Quote{金、银、宝石}（\BibleRef{格前 3:12}）；他们本应什么都不怕，无论是对那永远焚毁恶人的火，还是那炼净那些将要超越此火的人的火。因为经上说：\Quote{他自己将要得救，却如同经过火一样}。\par

3. 那么，此人凭什么理由祈求他不要\Quote{在震怒中被责斥，在怒火中被惩戒}呢？他仿佛对天主说：\Quote{既然我已承受的苦难众多，我求祢让这些足够吧；}于是他开始列举它们，以此令天主满足；献上他现在所受的，以免将来遭受更严重的痛苦。\par

4. \Quote{因为祢的箭射中了我，祢的手重重压在我身上}（\BibleRef{咏 37:2}）。\Quote{我肉体内外毫无完好之处，皆因祢的忿怒}（\BibleRef{咏 37:3}）。他现在开始述说这些在此世所受的苦难；然而，就连这已是出于上主的义怒，因为是出于上主的惩罚。\Quote{什么惩罚？}就是祂加于亚当的惩罚。不要以为惩罚没有加在他身上，或者天主编造了\Quote{你必定要死}这句话（\BibleRef{创 2:17}），或者我们在此生所受的任何事，不是出于那由原罪所赚得的死亡……那么，祂的\Quote{箭如何射中他}呢？正是那惩罚、那报应，以及在此生我们不得不承受的心灵和身体的痛苦，他正是用这些\Quote{箭}来描述。因为圣约伯也提到这些箭（\BibleRef{约 6:4}），并说上主的箭射中了他，当他忍受那些痛苦时。然而，我们也惯于称天主的话为箭；但他会因被这些箭射中而悲伤吗？天主的话是箭，仿佛点燃爱火，而非痛苦……那么，我们可以这样理解\Quote{箭射中}：祢的话牢牢钉在我心中；正是因这些话，我\Quote{记念了安息日}；而这安息日的记念，以及目前尚未拥有它，使我不能在此刻喜乐；并使我承认\Quote{我的肉体内毫无完好之处}，当我将这种完好与我在永恒安息中将拥有的完好相比时，就不该称之为完好；在那里\Quote{这可朽坏的将穿上不朽坏，这可死的将穿上不死}（\BibleRef{格前 15:53}），并且看出与那种完好相比，现今的这种不过是有病。\par

5. \BibleQuote{在我的骨头里也没有安息，因为我的罪过。} 通常人们会问，这是谁的话；有些人认为这是基督的话，因为这里有一些关于基督受难的话；我们很快就会谈到这些。我们自己也承认这是关于祂受难的话。但那没有罪过的祂，怎么能说：\BibleQuote{在我的骨头里也没有安息，因为我的罪过。}……因为如果我们说这不是基督的话，那么那些话：\BibleQuote{我的天主，我的天主，祢为什么舍弃了我？} 也就不会是基督的话。因为那里也有：\BibleQuote{我的天主，我的天主，祢为什么舍弃了我？} \BibleQuote{我的过犯之言远离我的救恩。} 正如这里你说：\BibleQuote{因为我的罪过，} 那里也有：\BibleQuote{我的过犯之言。} 如果基督完全没有 \Quote{罪} 和 \Quote{过犯}，我们就会开始认为那篇 \BibleBook{}{圣咏} 中的话也不是祂的。而如果那篇 \BibleBook{}{圣咏} 与基督无关，那是极其生硬和不一致的，因为在那篇 \BibleBook{}{圣咏} 中，祂的受难被清清楚楚地揭示出来，就好像从福音中读给我们听一样。因为那里有：\BibleQuote{他们分了我的外衣，为我的内衣拈阄。} 我何必提及主自己挂在十字架上时亲口说出了那篇 \BibleBook{}{圣咏} 的第一句：\BibleQuote{我的天主，我的天主，祢为什么舍弃了我？} 祂打算由此让人推断出什么，岂不是那整篇 \BibleBook{}{圣咏} 与祂有关，因为祂自己，身体的头，以祂自己的位格说出了它？当接下来说到：\BibleQuote{我的过犯之言，} 毫无疑问这是基督的话。那么 \Quote{罪} 从哪里来，岂不是从身体，即教会而来？因为基督的头和身体都在说话。为什么它们好像单一个人那样说话？因为正如祂说的：\BibleQuote{二人成为一体。}\BibleRef{创 2:24} \BibleQuote{这是}（宗徒说）\BibleQuote{极大的奥秘；但我是指着基督和教会说的。}……因为祂为什么不能说 \Quote{我的罪}，祂曾说过：\BibleQuote{我饿了，你们没有给我吃；我渴了，你们没有给我喝；我作客，你们没有收留我；我患病，我在监里，你们没有来探望我。}\BibleRef{玛 25:42} 主肯定没有在监里。祂为什么不能说这话呢？当有人对祂说：\BibleQuote{我们几时见了祢饥饿，或口渴，或在监里，而没有服事祢？} 祂回答说，祂是以祂身体的位格这样说的。\BibleQuote{你们没有给这些最小中的一个做，就是没有给我做。}\BibleRef{玛 25:44} 祂为什么不能说 \Quote{因为我的罪过}？祂曾对扫禄说：\BibleQuote{扫禄，扫禄，你为什么迫害我？}\BibleRef{宗 9:4} 然而，祂在天上，那时并未受迫害者之苦。但是，正如在那段经文中，头为身体说话，同样在这里，头说出了身体的话；同时你也听到了头本身的声音。然而，当你听到身体的声音时，不要把头与它分开；当你听到头的声音时，也不要把身体分开：因为 \BibleQuote{他们不再是两个，而是一体。}\BibleRef{玛 19:6}\par

6. \BibleQuote{在我的肉体上毫无健全，因为祢的忿怒。} 但是，亚当啊，天主对你发怒也许是不义的？人子啊，天主对你发怒是不义的？因为现在你已经承认你的惩罚，既然你是一个被安置在基督身体内的人，你说：\BibleQuote{在我的肉体上毫无健全，因为祢的忿怒。} 你要宣告天主发怒的正义：免得你似乎是在为自己辩解，而控告祂。继续说出主的 \Quote{忿怒} 从何而来。\BibleQuote{在我的肉体上毫无健全，因为祢的忿怒；在我的骨头里也没有安息。} 祂重复了前面所说的：\BibleQuote{在我的肉体上毫无健全；} 因为 \BibleQuote{在我的骨头里也没有安息} 与此相同。然而祂没有重复 \BibleQuote{因为祢的忿怒}，而是陈述了天主发怒的原因。\BibleQuote{在我的骨头里也没有安息，因为我的罪过。}\par

7. \BibleQuote{因为我的罪孽抬起了我的头，如同重担，沉重得令我无法承担。}\BibleRef{咏 37:4} 这里祂也是先说了原因，后说了结果。祂告诉我们发生了什么后果，以及从什么原因而来。\BibleQuote{我的罪孽抬起了我的头。} 因为只有不义的人才会骄傲，他的头被抬起。祂是\Quote{被抬起的}，因为他的\BibleQuote{头被高高抬起}敌对天主。你们在听 \WorkTitle{德训篇} 的诵读时听到：\BibleQuote{骄傲的开端是当人离开天主。} 那第一个拒绝听从诫命的人，\BibleQuote{他的罪孽抬起了他的头}敌对天主。因为他的罪孽抬起了他的头，天主对他做了什么？它们\BibleQuote{如同重担，沉重得令我无法承担}！轻浮的行为是抬起头，就好像抬起头的人没有什么可背负的。因此，既然那能被抬起的东西是轻的，它承受了重量，好把它压下去。因为\BibleQuote{他的恶行回到他自己的头上，他的强暴落到他自己的头顶。}\BibleQuote{它们如同重担，沉重得令我无法承担。}\par

8. \BibleQuote{我的创伤溃烂流脓} \BibleRef{咏 37:5}。如今，有创伤的人并非完全健康。此外，这些创伤\BibleQuote{溃烂流脓}。为何它们\Quote{溃烂}？因为它们\Quote{流脓}：这种情形如何适用于人的生命，谁不明白？只要一个人的灵魂的嗅觉正常，他就察觉罪是何等恶臭。与罪的恶臭相反的是那种馨香，关于此，宗徒说：\BibleQuote{我们在各处，在得救的人中，都是基督甘饴的馨香献给天主} \BibleRef{格后 2:15}。但这馨香从何而来，岂不是出于望德？从何而来，岂不是出于我们\Quote{追念安息日}？因为我们在此生哀叹的，与我们在来生所期待的，是不同的事物。我们哀叹的是\Emph{臭气}，我们寄望的是\Emph{芬芳}。因此，若没有那样的馨香吸引我们，我们绝不会追念安息日。但因着圣神，我们有了这样的馨香，以致我们对我们的新郎说：\BibleQuote{因你芳香的油膏，我们愿追随你} \BibleRef{歌 1:3}；我们便转离自身的不洁气味，转向祂，得着少许喘息。但确实，除非我们的恶行在我们的鼻孔中发出恶臭，我们绝不会以那些呻吟告明说：\BibleQuote{我的创伤溃烂流脓}。为什么呢？\Quote{因我的愚妄}。正如他先前所说的\Quote{因我的罪}，如今他也因同一原因说\Quote{因我的愚妄}。\par

9. \BibleQuote{我困苦，我极度弯腰} \BibleRef{咏 37:6}。他为何\Quote{弯腰}？因为他曾\Quote{高举}。如果你\Quote{谦卑，你将被高举}；如果你高举自己，你将\Quote{弯腰}；因为天主绝不缺少重担来压弯你……愿他为此事哀叹，好能获得另一事；愿他\Quote{追念安息日}，好配得抵达安息日。因为犹太人所庆祝的只是一个标记。那是何物的标记？是他所追念的，就是那说\Quote{我困苦，并极度弯腰}的。\Quote{直至终结}是什么意思？就是直至死亡。\par

\Quote{我终日哀痛。} \Quote{终日}，即\Quote{不停。} 所谓\Quote{终日}，他意指的是\Quote{我一生之久。} 但从何时他知道了？从他开始\Quote{追念安息日}之时。因为他\Quote{追念}他不再拥有的东西，你岂会要他\Quote{哀痛}？\Quote{我终日哀痛。}\par

10. \BibleQuote{我的灵魂饱受虚幻，我的肉身毫无安康} \BibleRef{咏 37:7}。哪里有完整的人，那里就有灵魂和肉身。\Quote{灵魂饱受虚幻}，\Emph{肉身}却\Quote{毫无安康}。还剩下什么能带来喜乐？难道人不应当\Quote{哀伤痛心}吗？\BibleQuote{我终日哀伤痛心}。愿哀伤成为我们的分，直到我们的灵魂脱去虚幻，身体披上安康。因为真正的安康无非是不朽。然而，灵魂的虚幻何其大，即使我试图表达，时间又岂够用？哪个灵魂不受它们支配？有一个简单的事例，我要提醒你们，以示我们的灵魂如何充满虚幻。这些虚幻有时几乎不让我们祈祷。我们不知道如何不借助形象思想物质事物，而那些我们不想要的形象却涌入脑海；我们想从这一个转到另一个，又放下那个另寻别的。有时你想回到之前想过的事情，放下现在所想；新的形象却呈现给你；你想唤起遗忘之事，它却不出现在脑海；反而出现了另一个你不想要的东西。同时，被你遗忘的那个哪里去了？为什么后来它不再被寻求时却出现了，而被寻求时，无数不想要的取而代之？我简要陈述了一个事实；我向你们，弟兄们，抛出了一个提示或建议，你们可以据此自行推断其余，并发现为我们的\Quote{虚幻}和\Quote{灵魂}哀伤意味着什么。因此他接受了虚幻的惩罚；他失去了真理。因为正如虚幻是灵魂的惩罚，真理就是它的赏报。但当我们置身于这些虚幻之中时，真理本身来到我们这里，发现我们被虚幻压倒，便亲自取了我们的肉身，或者说从我们取了肉身，即从人类取了肉身。祂将自己显现在肉身的眼前，以便祂能\Quote{因信德}治愈那些祂日后将向其启示真理的人，使真理能显现在已被治愈的眼上。因为祂自己就是\BibleQuote{真理} \BibleRef{若 14:6}，祂在当时应许给我们，那时祂的肉身要被肉眼看见，以便为那信德奠定基础，使真理成为其赏报。因为基督在地上显示的并不是祂自己，而是祂的肉身。因为如果祂显示了自己，犹太人就会看见并认识祂；但如果他们\BibleQuote{认识了祂，就决不会将光荣的主钉在十字架上} \BibleRef{格前 2:10}。然而，或许祂的门徒见过祂，当他们向祂说：\BibleQuote{将圣父显示给我们，我们就心满意足了} \BibleRef{若 14:8}；祂为表明他们所见不是祂自己，补充说：\BibleQuote{斐理伯！这么长久的时间，我和你们在一起，而你们还不认识我吗？谁看见了我，就看见了圣父} \BibleRef{若 14:9}。那么，如果他们看见了基督，为何还寻求圣父？因为如果他们看见的是基督，也就看见了圣父。所以他们还没有看见基督，他们渴望圣父显示给他们。为证明他们尚未看见祂，请注意，祂在另一处以赏报的方式应许说：\BibleQuote{那爱我的，必遵守我的命令；那爱我的，我父也必爱他；我也要爱他，并}（似乎有人对祂说：\Quote{祢爱他，祢将给他什么？}祂说）\BibleQuote{我要向他显示我自己} \BibleRef{若 14:21}。那么，如果祂以赏报的方式向爱祂的人应许此事，显然，应许给我们的真理的显现具有这样的性质：当我们看见它时，将不再说：\BibleQuote{我的灵魂饱受虚幻}。\par

11. \BibleQuote{我颓丧衰弱，过度悲哀} \BibleRef{咏 37:8}。那回想安息日超越高度的人，看见自己怎样被\Quote{深深压弯}。因为无法设想那安息之高度的人，看不到自己目前在哪里。因此另一篇圣咏说：\BibleQuote{我在出神中说：我从祢眼前被驱逐}。因为他的心神被提升到那里，看见某样崇高之物；但他尚未完全到达所见之物所在之处；好比一道永恒之光的闪光照耀了他，当他意识到自己尚未到达那他能以某种方式理解的事物时，他看见自己在哪里，以及他如何被人性软弱所束缚和\Quote{压弯}。他说：\BibleQuote{我在出神中说：我从祢眼前被驱逐}。这就是我在出神中看见的那某种东西，我由此感知自己离得有多远，因为我尚未到达那里。那说自己\Quote{被提到第三层天，在那里听到了不可言传的话，这是人不可以说出的}，已经是到达那里了。但他被召回到我们这里，为要使他因需要被成全，先哀伤自己的软弱，然后被披上能力。但因看见那些事而受到鼓励，为履行他的职事，他继续说：\BibleQuote{我听到了不可言传的话，这是人不可以说出的} \BibleRef{格后 12:4}。那么，你们问我或任何人那些人不可以说出的事又有何用？如果他不可以说出它们，谁又有权听到它们呢？然而，让我们在告解中哀伤叹息；让我们承认我们在哪里；让我们\Quote{记念安息日}，并耐心等待祂所应许的，祂以祂自己的位格也向我们显示了忍耐的榜样。\BibleQuote{我颓丧衰弱，过度悲哀}。\par

12. \BibleQuote{我以我心的呻吟哀号。} 你观察到天主的仆人们通常以呻吟转求；人们询问原因，表面上什么也没有，只有某位天主仆人的呻吟，如果这呻吟确实传到靠近他的人耳中。因为有一种隐秘的呻吟，是人听不到的。但如果某个强烈渴望的念头如此牢牢抓住了心，以至于内在之人的创伤以某种发出的呼喊表达出来，人们就会询问原因；一个人对自己说：\Quote{这或许是他呻吟的原因；} 以及 \Quote{这或许是他遭遇了这事或那事。} 谁能断定呢，除了那在他眼目和耳中呻吟的那一位？因此他说：\BibleQuote{我以我心的呻吟哀号；} 因为如果人们有时听到一个人的呻吟，他们大多只听到肉身的呻吟；他们听不到那以\Quote{他心的呻吟}呻吟的人。有人夺走了他的财物：他\Quote{哀号}，却不是\Quote{以他心的呻吟}：另一个因为埋葬了儿子，另一个埋葬了妻子；另一个因为葡萄园被冰雹毁坏；另一个因为酒桶变酸；另一个因为有人偷了他的牲口；另一个因为遭受损失；另一个因为害怕某个他的仇敌：所有这些人都以\Quote{肉身的呻吟}\Quote{哀号}。然而，天主的仆人，因为他因纪念安息日而\Quote{哀号}，那里是天主之国，是血肉不能承受的，他说：\BibleQuote{我以我心的呻吟哀号。}\par

13. 谁观察并注意到他呻吟的原因呢？\BibleQuote{我的一切渴望都在祢面前} \BibleRef{咏 37:9}。因为它不是在不能看见心的人面前，而是在祢面前，我的一切渴望都是敞开的！让你的渴望在祂面前；\BibleQuote{那位在暗中看见的圣父，必要报答你。} \BibleRef{玛 6:6} 因为是你心的渴望成为你的祈祷；如果你的渴望持续不断，你的祈祷也持续不断。因为宗徒并非没有意义地说：\BibleQuote{不断地祈祷。} \BibleRef{得前 5:17} 我们是否要\Quote{不断地}屈膝、俯伏身体，或举起双手，以致他说\Quote{不断地祈祷}？或者如果我们以这种意义说我们\Quote{祈祷}，我相信，我们不能\Quote{不断地。} 还有另一种内在的不断的祈祷，就是心的渴望。无论你做什么，只要你渴望那安息日，你就没有停止祈祷。如果你想从不停止祈祷，就不要停止渴望它。你渴望的持续就是你祈祷的持续。如果你停止渴望它，你就会停止说话。那些停止说话的是谁？就是那些关于他们所说的：\BibleQuote{因为罪恶增多，许多人的爱德变得冷淡。} \BibleRef{玛 24:12} 爱德的冷冻是心的沉默；爱德的燃烧是心的呼声。如果爱仍然持续，你就仍然在扬起你的声音；如果你总是在扬起你的声音，你总是在渴望某物；如果总是渴望某不在之物，你就是在\Quote{纪念安息日的安息。} 而且重要的是，你也应该明白，你的\Quote{你心的哀号}是在谁面前敞开的。现在那么思想，那些在天主眼目之前的渴望应当是什么样的。难道是对我们敌人死亡的渴望吗？这是人们自欺欺人认为可以合法祈求的事吗？因为我们有时祈求不应当的。让我们思想他们自欺欺人认为可以合法祈求的事！因为他们祈求某个人能死，他的遗产归他们。但是那些为敌人死亡而祈祷的人，也当听主说：\BibleQuote{要为你们的仇敌祈祷。} \BibleRef{玛 5:44} 让他们不要为此祈祷，让他们的敌人死；而是为此祈祷，让他们被挽回；那么他们的敌人就死了；因为从他们被挽回的时候起，从此他们就不再是敌人了。\Quote{我的一切渴望都在祢面前。} 如果我们设想我们的渴望在祂面前，而那个\Quote{呻吟}本身却不在祂面前，那会怎样？这怎么可能，既然我们的渴望本身在\Quote{呻吟}中表达出来？因此接着：\Quote{我的呻吟不向祢隐藏。}\par

从祢那里确实不隐藏；但向许多人隐藏。天主的仆人有时似乎在谦卑中说：\Quote{我的呻吟不向祢隐藏。} 有时他也似乎微笑。那么那渴望是否死在他心中？但如果内心有渴望，也就有\Quote{呻吟}。它并不总是传到人的耳中；但它从不停止在天主的耳中响起。\par

14. \Quote{我的心烦乱} \BibleRef{咏 37:10}。为何烦乱？\Quote{我的勇气也离开了我。} 一般有突然之事临到我们：\Quote{心烦乱；} 大地震动；天上打雷；可怕的袭击临到我们，或听见可怕的声音。也许在路上看见狮子；\Quote{心烦乱。} 也许有强盗埋伏等候我们；\Quote{心烦乱：} 我们充满惊恐；从各方都有令人焦虑之事。为何？因为 \Quote{我的勇气离开了我。} 因为如果那勇气仍然坚定不移，有什么可怕的呢？无论带来什么坏消息，无论什么威胁我们，无论听见什么声音，无论什么要坠落，无论什么显得可怕，都不会引起恐惧。但那烦乱从何而来？\Quote{我的勇气离开了我。} 为何我的勇气离开了我？\Quote{我眼睛的光也离开了我。} 同样，亚当也不能看见 \Quote{他眼睛的光}。因为 \Quote{他眼睛的光} 就是天主自己，当他得罪了祂时，就逃到荫蔽中，藏在乐园的树木间。\BibleRef{创 3:8} 他惊慌地躲避天主的容颜，寻求树木的遮蔽；从此在树木中他不再有 \Quote{他眼睛的光}，他曾习惯于在那光中喜乐。……\par

15. \Quote{我的爱人；} 我何必再提我的仇敌？\Quote{我的爱人和我的邻人走近前来，站在我的对面} \BibleRef{咏 37:11}。要明白他所说的：\Quote{站在我的对面。} 因为他们若站在我的对面，他们就跌倒在自己身上。\Quote{我的爱人和我的邻人走近前来，站在我的对面。} 现在我们当认出那说话之元首的话；现在让我们的元首在祂的受难中开始光照我们。然而当元首开始说话时，你不要将身体与祂分离。如果元首不愿意将自己与身体的话分开，难道身体竟敢将自己与元首的苦难分开吗？你要在基督的苦难中受苦：因为基督仿佛因你的软弱而犯了罪。只因刚才祂述说你的罪，仿佛以祂自己的位格说话，并称那些罪为祂自己的。……对于那些希望接近祂高举、却不思想祂谦卑的人，祂回答他们说：\Quote{你们能饮我将饮的杯吗？} \BibleRef{玛 20:22} 因此主的那些苦难也是我们的苦难：倘若每个人好好事奉天主，真正持守信德，各尽其责，在众人面前诚实行事，我倒要看看他是否不遭受基督在此详述的祂受难中的事。\Quote{我的爱人和我的邻人走近前来，站在我的对面。}\par

16. \Quote{我的邻人远远站着。} 哪些是走近的 \Quote{邻人}，哪些是远远站着的？犹太人是 \Quote{邻人}，因为 \Quote{近亲}，他们甚至在他们钉死祂时走近了；宗徒们也是祂的 \Quote{邻人}；他们也 \Quote{远远站着}，以免必须与祂同受苦。也可以这样理解：\Quote{我的朋友}，即那些假装是 \Quote{我的朋友}的人：因为他们假装是祂的朋友，当他们说：\Quote{我们知道你诚诚实实教导天主的道}；\BibleRef{玛 22:16} 他们想试探祂，是否该给凯撒纳税；当祂用他们自己的话驳倒他们时，他们想显得是祂的朋友。\Quote{但祂不需要任何人给人作证，因为祂知道人心里有什么}；\BibleRef{若 2:25} 所以当他们用朋友的话语对祂说话时，祂回答他们说：\Quote{假冒为善的人，为什么试探我？} \BibleRef{玛 22:18} \Quote{我的朋友和我的邻人} 于是 \Quote{走近前来，站在我的对面，我的邻人却远远站着。} 你明白我所说的：我称那些 \Quote{走近} 又同时 \Quote{远远站着} 的人为邻人。因为他们在身体上 \Quote{走近}，但在心里 \Quote{远远站着}。谁在身体上比那些将祂举在十字架上的人更靠近祂？谁在心里比那些亵渎祂的人更远离祂？请听先知依撒意亚描述这种远近；同时观察这种近和远。\Quote{这百姓用嘴唇尊敬我：} 看，他们用身体走近；\Quote{心却远离我。} \BibleRef{依 29:13} 同一些人同时 \Quote{近} 又 \Quote{远}：用嘴唇近，心里远。然而，因为宗徒们也因恐惧远远站着，我们更简单且更恰当地将这理解为指着他们；意思是有些人走近，另一些人远远站着；因为连伯多禄，他虽然比别人更大胆地跟随，仍然离得那么远，以致被质问和惊慌时，三次否认了曾应许 \Quote{准备同祂死} 的主。后来，为了使那从远处的人得以来近，复活后他听见了问题：\Quote{你爱我吗？} 并说：\Quote{主，我爱你；} \BibleRef{若 21:15} 这样说就使他 \Quote{近} 了，正如曾否认祂使他成为 \Quote{远}，直到用三次爱的宣认，他摆脱了三次否认。\Quote{我的邻人远远站着。}\par

17. \BibleQuote{他们也寻找我的灵魂，图谋施暴于我} \BibleRef{咏 37:12}。现在显明谁\Quote{寻找祂的灵魂}；即那些不在祂身体内、因而没有拥有祂灵魂的人。那些\Quote{寻找祂灵魂}的人，远离了祂的灵魂；然而他们\Quote{寻找它}是为要毁灭它。因为祂的灵魂也可以被正当地\Quote{寻找}。在另一处，祂责备某些人说：\Quote{没有人关心我的灵魂。}祂责备一些人没有寻找祂的灵魂；又责备另一些人寻找它。谁是那正当地寻找祂灵魂的人？那效法祂苦难的人。谁是那错误地寻找祂灵魂的人？就是那些\Quote{图谋暴力反对祂}并将祂钉十字架的人。\par

18. 祂接着说：\Quote{那些寻找我的过犯的，说了虛妄的话。}什么是\Quote{寻找我的过犯}？他们寻找许多事，却找不着。也许祂的意思是：\Quote{他们寻找我的罪状。}因为他们寻找可反驳祂的话，却\BibleQuote{找不着}\BibleRef{玛 26:60}。因为他们企图在\Quote{善者}身上找到恶行，在无辜者身上找到罪行。他们怎能在那位无罪者身上找到这些呢？但因为必须在无罪者身上寻找罪过，他们只好捏造他们找不着的东西。因此，\Quote{那些寻找我的过犯的，说了虛妄的话}，即不真实的话，\Quote{终日图谋诡诈}；也就是说，他们不断地策划背信弃义。你知道在祂受难之前，人们怎样以残暴的伪证反对主。你知道甚至在祂复活之后，人们怎样以残暴的伪证反对祂。那些看守祂坟墓的士兵们（依撒意亚论到他们说：\BibleQuote{我使恶人埋葬祂}（因为他们本是恶人，不愿说实话，收受贿赂就散布谎言）），想想他们说了何等\Quote{虚妄的话}。他们也受审问，说：\BibleQuote{我们睡觉的时候，祂的门徒来把祂偷去了} \BibleRef{玛 28:13}。这就是\Quote{说虚妄的话}。因为如果他们睡着了，又怎能知道所发生的事呢？\par

19. 祂接着说：\BibleQuote{但我像聋子，不听} \BibleRef{咏 37:13}。祂对所听到的不作回应，就好像不听他们一样。\BibleQuote{但我像聋子，不听；我又像哑巴，不开口。}祂又重复同样的话。\par

\BibleQuote{我成了一个不听的人，口中没有责备的话} \BibleRef{咏 37:14}。好像祂对他们无话可说，好像祂没有什么可指责他们的。难道祂之前没有多次指责他们吗？难道祂没有说许多话，并且说了\BibleQuote{祸哉，你们经师和法利塞人} \BibleRef{玛 23:13}以及许多别的话吗？然而当祂受苦时，祂一句也没有说这些；不是祂没有可说的，而是祂等待他们完成一切，并使关于祂的一切预言应验，正如论到祂所说的：\BibleQuote{又如羊在剪毛的人前默不作声，祂也一样不开口} \BibleRef{依 53:7}。祂必须在受难时保持沉默，尽管以后在审判时不会沉默。因为祂来是要受审判，而祂将要来审判；并且祂之所以要以大能来审判，正是因为祂曾在极大的谦卑中受了审判。\par

20. \BibleQuote{因为上主，我仰望祢；祢必应允我，我的天主} \BibleRef{咏 37:15}。好像有人对祂说：\Quote{祢为什么不开口？祢为什么不说\InnerQuote{停止}？祢为什么在十字架上不斥责不义的人？}祂接着说：\BibleQuote{因为上主，我仰望祢；祢必应允我，我的天主。}祂提醒你：若遭遇患难该当如何。因为你寻求为自己辩护，也许你的辩护无人倾听。那时你便受窘，仿佛败诉了；因为无人为你辩护或为你作证。只要你在内里保持无辜，在那里没有人能歪曲你的案件。在人前，伪证得胜了。那么在天主面前，当你的案件要受审时，伪证也能得胜吗？当天主审判时，除了你自己的良心，再无其他见证。在公义的审判官和你的良心面前，只要害怕你自己的案件。如果你没有恶案，你就不会害怕控告者，无需驳倒伪证，也无需寻找真理的见证。只要把良善的良心带上法庭，你就能说：\BibleQuote{因为上主，我仰望祢；祢必应允我，我的天主。}\par

21. \BibleQuote{因为我曾说：愿我的仇敌永不因我而欢跃；我的脚滑跌时，他们就对我妄自尊大。}\BibleRef{咏 37:16} 祂再次回到祂身体的软弱上；元首再次留意它的\Quote{脚}。元首并非如此在天上，以致抛弃祂在地上的；祂明明看见并观察我们。因为有时，按今世的方式，我们的脚\Quote{滑跌}，因陷入某些罪而失足；在那里，仇敌的舌头以最恶毒的恶意升起。由此我们看出他们真正的意图，即使在他们沉默时也是如此。然后他们以毫不留情的严厉说话；为他们本应悲伤的发现而欢喜。\Quote{我又说：以免我的仇敌一时欢跃于我。}我确实这样说了；然而，也许是为了纠正我，祢使他们\Quote{在我脚滑跌时对我妄自尊大}；也就是说，当我绊倒时，他们得意洋洋，说了许多话。因为对软弱者，他们本应同情而非侮辱；正如宗徒所说：\BibleQuote{弟兄们，若有人陷于过犯，你们属神的人要用温柔的心来挽回这样的人；也要当心自己，免得你也受试探。}\BibleRef{迦 6:1} 祂所说的那些人并非如此：\Quote{我的脚滑跌时，他们大大欢喜攻击我；}而是如同祂在别处所说的：\Quote{恨我的人若我跌倒就欢跃。}\par

22. \BibleQuote{因为我已为鞭打作了准备。}\BibleRef{咏 37:17} 这是极雄壮的表达；仿佛祂在说：\Quote{我为此而生，为要受苦。}因为祂不是要从别处出生，而是从亚当所生，亚当应受鞭打。但罪人在今生有时完全不受鞭打，或受的鞭打少于他们所应得的：因他们心中的邪恶已被视为毫无希望。然而，那些预备得永生的人，必须在今生受鞭打：因为那句断言是真实的：\BibleQuote{我儿，不要轻视主的管教，也不要厌倦祂的责备；}\BibleRef{箴 3:11}\BibleQuote{因为主所爱的祂必管教，又鞭打凡收纳为子的。}\BibleRef{希 12:6} 所以，让我的仇敌不要侮辱我；让他们\Quote{不要妄自尊大；}若我的父鞭打我，\Quote{我已为鞭打作了准备；}因为为我存留的产业。你不愿忍受鞭打：产业就不会赐予你。因为\Quote{每个儿子}都必须受鞭打。此话如此真实，没有罪的祂也未被豁免。因为\BibleQuote{我已为鞭打作了准备。}\par

23. \Quote{我的忧苦常在我眼前。} 这\Quote{忧苦}是什么呢？也许是对我鞭笞的忧苦。我各位弟兄，说实话，说实话，让我告诉你们：人们为自己的鞭笞哀哭，却不哀哭鞭笞的原因。这里的人不是这样。我各位弟兄，请听：若有人遭受任何损失，他更常说\Quote{我不该遭受这事}，而不思考他为何遭受这事；他为损失金钱而哀哭，却不哀哭正义的损失。若你犯了罪，就为你内里宝藏的损失而哀哭吧。你家里一无所有，但或许你心里更为空虚；但若你的心充满了它的美善——即你的天主——你为什么不说：\BibleQuote{主赐予的，主已收回；照主所喜悦的，就这样成就了。愿主的名受赞美。} \BibleRef{约 1:21} 那么，祂何以忧伤呢？是为了祂所受的\Quote{鞭笞}吗？天主不许。\Quote{我的忧苦}（祂说）\Quote{常在我眼前。} 仿佛我们问：\Quote{什么忧苦？这忧苦从何而来？} 祂说：\BibleQuote{因为我宣告我的过犯，我要关心我的罪。} \BibleRef{咏 37:18} 看，这就是忧苦的原因！这不是由鞭笞引起的忧苦，不是为药石，不是为创伤。因为鞭笞是对抗罪过的药石。我各位弟兄，请听：我们是基督徒，若某人的儿子死了，他为他哀哭；但若他犯罪，他却不为他哀哭。他看见他犯罪时，才应为他哀哭，为他悲叹。那时他应约束他，给他生活的规条；应给他施以纪律。或者他这样做了，而对方没有留意，那么，那时他本当为他哀哭；那时他生活在奢华之中，比因死亡而结束奢华更加致命地死去。那时，他在你家里行这些事，他不仅\BibleQuote{死了，而且是发臭的。} \BibleRef{若 11:39} 这些事值得哀悼，另一些事则可以忍受。我说，那些是可以容忍的，这些是值得哀哭的。他们应当像你听见这个人哀哭那样哀哭：\Quote{因为我宣告我的过犯。我要关心我的罪。} 你告明了罪，不要以为总能告明你的罪并再犯，就心安理得。你要这样\Quote{宣告你的过犯，如同关心你的罪。} \Quote{关心你的罪}是什么意思？就是关心你的创伤。若你说\Quote{我要关心我的创伤}，这是什么意思？岂不是我要努力使它痊愈吗？因为这就是\Quote{关心自己的罪}：不断地奋斗，不断地努力，全力以赴、热忱地治愈自己的创伤。看！你日复一日为你的罪哀哭；但也许你的眼泪确实流下，你的手却不动。要施舍，要赎你的罪，让穷人因你的慷慨而欢喜，好使你也因天主的恩宠而喜乐。他有所缺乏，你也有所缺乏；他需要你的帮助，你也需要天主的帮助。你藐视需要你帮助的人，那么当你需要天主帮助时，天主岂不藐视你吗？所以，你要供给需要你帮助的人，好使天主在你内满足你的需要。这就是\Quote{我要关心我的罪}的意思。我要做一切当做的事，以涂去并治愈我的罪。\Quote{我要关心我的罪。}\par

24. \BibleQuote{但我的仇敌却活着} \BibleRef{咏 37:19}。他们富裕，因世俗的繁荣而欢乐，而我却在受苦，\Quote{因我心里的叹息而咆哮}。祂的仇敌以什么方式\Quote{活着}呢？因为祂已经论到他们说，他们\Quote{说了虚妄的话}？请在另一篇圣咏中也听：\Quote{他们的儿子像幼小的植物，深深扎根。} 但上面祂曾说：\Quote{他们的口讲虚妄。他们的女儿修饰得如同圣殿的模样；他们的仓房满盈，涨溢更多；他们的牲畜肥壮，他们的羊群繁殖，在街上增多；没有篱笆倒塌；街上没有哭号。} 那么\Quote{我的仇敌}就\Quote{活着}。这是他们的生命；他们赞美这种生命；他们倾心于此；他们紧抓不放，以至自取灭亡。因为之后是什么？他们称\Quote{处于这种景况的人民}为有福的。但你怎么说呢？你这位\Quote{关心你的罪}的人？你这位\Quote{宣告你的过犯}的人？他说：\Quote{以天主为主的人民，真是有福。}\par

\Quote{但我的仇敌活着，且比我强盛，无理恨我的人增多。} \Quote{无理恨我}是什么意思？他们恨我——我愿他们好。若他们只是以恶报恶，他们不会不义；若他们不以善报我所行的善，他们就是忘恩负义。然而，那些\Quote{无理恨人}的，实际是以恶报善。犹太人就是这样。基督以善来到他们那里，他们却以恶报善。我各位弟兄，要提防这种恶；它很快就会潜入我们。你们中没有人因为我说了\Quote{犹太人就是这样}，就以为自己远离了危险。若一位兄弟为你好而责备你，你恨他，你就与他们一样。要留意，它是多么容易、多么快速地产生；要避免如此大的恶，如此容易犯的罪。\par

25. \Quote{他们也以恶报善，毁谤我，因为我追求正义} \BibleRef{咏 37:20}。正因如此，我以善报恶却得恶报。所谓\Quote{追求正义}是什么意思？并非放弃。为了使你不总是将\Emph{\OriginalTerm{迫害}{persecutio}}理解为贬义，他以\Emph{\OriginalTerm{追逼}{persecutus}}一词表示紧追不舍、彻底追随。\Quote{因为我追求正义。}请听我们的元首在受难中悲切呼喊：\Quote{他们弃绝了祢所爱的，如同弃绝令人厌恶的死人。}他\Quote{死了}还不够吗？为何还要\Quote{令人厌恶}？因为他是被钉死的。十字架之死在他们的眼中是极大的憎恶，因为他们没有察觉到经上预言说：\Quote{凡被挂在树上的，都是可咒骂的。} \BibleRef{申 21:23}因为他自己并未带来死亡，而是在此找到了从首人的诅咒中传下来的死亡；他承担了我们这源于罪恶的死亡，并被挂在树上。为避免有人（如某些异端者所想的）以为我们的主耶稣基督只有一个虚假的血肉之躯，以为他在十字架上所做的赎罪之死不是真实的死亡，先知注意到了这一点，说：\Quote{凡被挂在树上的，都是可咒骂的。}因此他表明天主之子真实地死了，是属血肉之躯应受的死亡；免得他若不被咒骂，你就以为他并未真实地死去。然而，那死亡并非虚幻，而是来自从诅咒衍生的原始根源，当时他说：\Quote{你必定要死。} \BibleRef{创 2:17}既然真实的死亡确实也临到他身上，为的是使真实的生命临到我们，死亡的诅咒也就临到他身上，为的是使生命的祝福也能临到我们。\Quote{他们弃绝了祢所爱的，如同弃绝令人厌恶的死人。}\par

26. \Quote{求祢不要舍弃我；上主，我的天主，求祢不要远离我} \BibleRef{咏 37:21}。让我们在他内说话，藉着他说话（因为他亲自为我们代祷），并且说：\Quote{上主我的天主，求祢不要舍弃我。}然而他曾说过：\Quote{我的天主，我的天主，祢为什么舍弃了我？} \BibleRef{玛 27:46}现在却说：\Quote{我的天主，求祢不要远离我。}若他不舍弃身体，难道舍弃了元首？那么，这些话语除了出自首人，还能是谁的呢？为了表明他携带了出自首人的真实血肉之躯，他说：\Quote{我的天主，我的天主，祢为什么舍弃了我？}天主并未舍弃他。若他不舍弃你——那信从他的人，难道圣父、圣子、圣神，独一天主，会舍弃基督吗？但他将首人的位格转移到了自己身上。我们由宗徒的话知道：\Quote{我们的旧人已与他同钉在十字架上。} \BibleRef{罗 6:6}然而，我们若不是因他\Quote{在软弱中}被钉，就不会脱去旧日本性。因为他为此而来，为使我们在他内得以更新，正是由于渴慕他、效法他受苦的榜样，我们得以更新。因此，那是软弱的呼声；我说的是那句：\Quote{祢为什么舍弃了我？}由此，在前面的经文中说到：\Quote{我过犯的话语。}仿佛在说：这些话语从罪人之位转到了我的位格上。\par

27. \Quote{求祢不要远离我。求祢快来帮助我，我救恩的主。} \BibleRef{咏 37:22}这就是那\Quote{救恩}，弟兄们，正如宗徒伯多禄所说：\Quote{先知们曾详细考究过} \BibleRef{伯前 1:10}，虽详细考究，却未寻见。他们探究并预言了它；而我们来了，找到了他们所寻求的。并且看，我们自己也尚未领受它；在我们之后还要生出他人，他们也将寻见，但同样尚未领受，就要离去，好让我们全体一同领受\Quote{在一天的终结的救恩银币}，与先知、圣祖和宗徒一起。因为你们知道，雇工或劳工在不同时刻被召进葡萄园，却都平等地领受了工钱。 \BibleRef{玛 20:9}所以，宗徒、先知、殉道者，以及我们自己，还有那些直到世界末日跟随我们的人，都要在终结本身领受永久的救恩；那时我们瞻仰天主的圣容，默观他的荣耀，永远赞美他，不再有缺陷，不再有罪恶的惩罚，不再有任何罪过的偏斜：赞美他，不再渴慕他，而是紧紧依附于我们曾一直渴慕、并曾因望德而喜乐的他。因为我们将在那城中，在那里，天主是我们的福乐，天主是我们的光明，天主是我们的食粮，天主是我们的生命；我们一切美好的事物——当我们远离时才如今悲伤的——都要在他内寻得。在他内将有那\Quote{安息}，当我们如今\Quote{记起}它时，不能不悲伤。因为那是我们所\Quote{记起}的\Quote{安息日}；对此回忆，已经说了如此伟大的事，我们也应当继续说如此伟大的事，并且永不停息地用我们的话语——不是用嘴唇，而是用内心——来述说；因为我们的嘴唇止息，正是为了让我们以心灵呼喊。\par

\SourceRecord{Translated by J.E. Tweed.；Nicene and Post-Nicene Fathers, First Series；Vol. 8.；Edited by Philip Schaff.；Buffalo, NY: Christian Literature Publishing Co.,；1888.；<http://www.newadvent.org/fathers/1801038.htm>.}

% structure: p0001 source=h1 target=section reason=page-title

\section{\WorkTitle{圣咏集第三十九篇释义}}

1. 这首圣咏的标题，我们刚刚咏唱并提议讨论的，是\Quote{论终末，为\Person{依杜通}，\Person{达味}自己的圣咏。}因此，我们必须寻找并留意一个被称为\Person{依杜通}之人的话语；如果我们每个人自己就是\Person{依杜通}，那么在他所咏唱的中，他认出自己，并听见自己说话。因为你们可以看到，按照人类的古老传承，谁被称为\Person{依杜通}；然而，让我们理解这个名字的翻译，并试图在词语的翻译中领悟真理。因此，根据我们从那些研究圣经的人那里，在从希伯来语译为拉丁语的名字中能够探询发现的，\Person{依杜通}被翻译为\Quote{超越他们}。那么，这个\Quote{超越他们}的人是谁？或者，他所\Quote{超越}的是谁？…因为有些人仍紧贴大地，仍俯伏于地，仍将心放在低下之处，仍将希望寄托在消逝的事物上，那个被称为\Quote{超越他们}的人已经\Quote{超越}了他们。\par

2. 你们知道有些圣咏的标题是\Quote{登阶之歌}；而在希腊文中，这个词是很明显的。\par

\SourceRecord{Translated by J.E. Tweed.；Nicene and Post-Nicene Fathers, First Series；Vol. 8.；Edited by Philip Schaff.；Buffalo, NY: Christian Literature Publishing Co.,；1888.；<http://www.newadvent.org/fathers/1801039.htm>.}

% structure: p0001 source=h1 target=section reason=page-title

\section{圣咏第四十篇释义}

1. 在我们的主耶稣基督所预言的这一切事中，我们知道一部分已经实现，另一部分我们希望在将来实现。然而，这一切都将应验，因为祂就是说话的那位\Quote{真理}，并且要求我们如同祂忠信地说话一样，成为\Quote{忠信}的……\par

2. 那么，让我们说说这篇圣咏所说的话。\Quote{我耐心等候了主}\BibleRef{咏 39:1}。我等候的，不是任何可欺人亦自欺的凡人的应许；我等候的，不是任何可因自身忧苦而在我得安慰前先被吞噬的凡人的安慰。难道一个同为凡人的兄弟，他自己也处于忧苦中时，还能安慰我吗？让我们一同哀悼，一同哭泣，让我们一同\Quote{耐心等候}，也一同同心祈祷。我所等候的，除了主，还有谁呢？主虽然延缓祂应许的实践，却从不收回它们。祂必会实现；确然，祂必会实现，因为祂已经实现了许多应许；对于天主的真理，我们不应有任何恐惧，即使祂至今尚未实现任何一个应许。看！让我们从此这样想：祂已应许了我们一切，尚未赐予我们任何拥有；祂是一位应许的保证者，一位忠信的支付者；你只须显出自己是应许的忠实的催讨者；如果你\Quote{软弱}，如果你是一个小孩子，就索取祂仁慈的应许。难道你没有看见柔嫩的羔羊用头撞击母羊的乳房，好使它们饱饮乳汁吗？……\Quote{祂留意了我，并垂听了我的呼号。}祂留意了这呼号，也垂听了它。你看，你没有白等。祂的眼目在你之上。祂的耳朵转向你。因为，\BibleQuote{上主的眼目看顾义人，祂的耳朵垂听他们的呼号。}那么如何？当你作恶并亵渎祂时，祂没有看见你吗？那么同篇圣咏所说的\BibleQuote{上主的面容敌对作恶的人}又作何解释？但目的是什么？\BibleQuote{好从世上剪除对他们的纪念。}因此，即使当你作恶时，祂\Quote{留意\Emph{于}你}；但祂\Quote{没有留意\Emph{到}你}。所以，对于\Quote{耐心等候了主}的人，单单说\Quote{祂留意了我}还不够，祂说：\InnerQuote{祂留意\Emph{到}了我}。也就是说，祂借着安慰我而留意了我，好使我得益。祂留意了什么？\BibleQuote{并垂听了我的呼号。}\par

3. \Person{祂}为你成就了什么？\Person{祂}为你做了什么？\BibleQuote{\Person{祂}把我从可怕的坑里拉上来，从泥沼中拉出来，把我的脚放在磐石上，使我的步伐稳当}\BibleRef{咏 39:2}。\Person{祂}已经赐给我们伟大的祝福；但\Person{祂}仍然是我们的债主；但让那些已经得到债务偿还的人，相信其余也将偿还，因为他甚至在得到任何东西之前就应该相信。\Person{上主}亲自用事实说服我们，\Person{祂}是忠实的应许者，慷慨的施与者。那么\Person{祂}已经做了什么？\BibleQuote{\Person{祂}把我从可怕的坑里拉出来}。那可怕的坑是什么？是罪孽的深渊，来自肉欲，这就是\BibleQuote{泥沼}\BibleRef{咏 39:2}的意思。\Person{祂}从哪里把你拉出来？从一个深渊中，在另一篇圣咏中你曾呼喊：\BibleQuote{\Person{上主}，我从深渊向你呼号}\BibleRef{咏 130:1}。而那些已经\BibleQuote{从深渊呼号}的人，并非完全处于最深的深渊；呼号的行为本身就在提升他们。有一些人在更深的深渊中，他们甚至没有察觉到自己在深渊中。这些人就是骄傲的蔑视者，不是谦卑的祈求赦免者，不是流泪哀求仁慈者；而是圣经这样描述的人：\BibleQuote{恶人陷入罪恶的深渊，便傲慢无礼}\BibleRef{箴 18:3}。因为他是在深渊中更深的人，他不满足于做一个罪人，除非他不认罪甚至辩护自己的罪。但那个已经\BibleQuote{从深渊呼号}的人，已经抬起头来以从深渊中呼号，已经被垂听，并且已经被\BibleQuote{从可怕的坑里，从泥沼中拉出来}。他已经拥有以前没有的信德；他已经拥有以前没有的望德；他现在在\Person{基督}里行走，而以前他在魔鬼里迷途。因此他说：\BibleQuote{\Person{祂}把我的脚放在磐石上，使我的步伐稳当}\BibleRef{咏 39:2}。现在\BibleQuote{那磐石就是\Person{基督}}\BibleRef{格前 10:4}。假设我们\BibleQuote{在磐石上}，并且我们的\BibleQuote{步伐已经安排妥当}，我们仍需继续行走；我们要向更高处前进。因为宗徒\Person{保禄}现在在磐石上，他的步伐已经稳当，他说了什么？\BibleQuote{这不是说我已经达到目标，或已是成全的；弟兄们，我不以为我已经抓住了}\BibleRef{斐 3:12}。那么，如果你没有抓住，为你做了什么？你凭什么感谢说：\BibleQuote{但我蒙受了怜悯}\BibleRef{弟前 1:13}？因为他的步伐现在稳当，因为他现在在磐石上行走？……因此，当他说：\BibleQuote{我向目标直奔，为获得天主在基督耶稣内召我向上争夺的奖品}\BibleRef{斐 3:14}时，因为\BibleQuote{他的脚现在放在磐石上}，并且\BibleQuote{他的步伐安排妥当}，因为他现在行走在正路上，他有理由感恩；也有理由继续祈求；为已经领受的感恩，同时要求仍然欠着的债。他为什么事情感恩？为罪过的赦免，为信德的光照，为望德的稳固支持，为爱德的火焰。但他在哪方面仍然向\Person{上主}要求欠债？\BibleQuote{从今以后，正义的冠冕已为我预备下了}\BibleRef{弟后 4:8}。因此还有东西欠我。欠什么？\BibleQuote{正义的冠冕，就是\Person{主}，正义的审判者，在那一天要赐给我的}\BibleRef{弟后 4:8}。\Person{祂}先是慈爱的父亲，\BibleQuote{把他从可怕的坑里带出来}；赦免他的罪，把他从\Quote{泥沼}中救出；此后\Person{祂}将是\Quote{正义的审判者}，向行走正直的人回报\Person{祂}所应许的；回报给（我说）\Person{祂}首先赐予行走正直能力的人。然后\Person{祂}作为\Quote{正义的审判者}将偿还；但偿还给谁？\BibleQuote{唯独坚持到底的，才可得救}\BibleRef{玛 10:22}。\par

4. \BibleQuote{\Person{祂}在我口里放了一首新歌}\BibleRef{咏 39:3}。什么是新歌？\BibleQuote{就是赞美我们的天主}\BibleRef{咏 39:3}。也许你以前向异神唱赞美诗；那是旧歌，因为是\Quote{旧人}唱出的，不是\Quote{新人}唱出的；让\Quote{新人}得以形成，让他唱\Quote{新歌}；他本身被造为\Quote{新}的，让他爱那些使他变\Quote{新}的\Quote{新}事物。因为有什么比\Person{天主}更古老？\Person{祂}居于万有之前，无始无终。当你归向\Person{祂}时，\Person{祂}对你成为\Quote{新}的；正是由于远离\Person{祂}，你才变得老旧；并且你说：\Quote{我因我所有仇敌而衰老}。我们因此向我们的\Person{天主}发出\Quote{赞美诗}；这赞美诗本身解放我们。\BibleQuote{因为我要呼求\Person{主}赞美\Person{祂}，我将安全脱离我所有仇敌}\BibleRef{咏 39:3}。因为赞美诗是赞美之歌。呼求\Person{天主}来\Quote{赞美}\Person{祂}，而非指责\Person{祂}……\par

5. 若有人问，在这篇圣咏中说话的是谁？我愿简要地说，\Quote{是基督。}但是，弟兄们，如同你们知道，也如同我们必须常说的，基督有时以祂自己的位格说话，有时以我们头的位格说话。因为祂自己就是\BibleQuote{身体的救主。}\BibleRef{弗 5:23}祂是我们的头；天主子，由童贞女诞生，为我们受苦，\BibleQuote{为使我们成义而复活，}坐在\BibleQuote{天主的右边，}以\BibleQuote{为我们转求：}\BibleRef{罗 8:34}祂也将要在审判时，按照恶人和善人所行的一切善恶，予以报应。祂屈尊成为我们的头；成为\Quote{身体的头，}由我们取了那祂将为我们而死的人身；祂又为我们使那肉身复活，好在那肉身中为我们树立复活的榜样；使我们学会盼望那过去所绝望的，从此将我们的脚立在盘石上，并在基督内行走。因此，祂有时以我们头的名义说话；有时也以我们这些祂的肢体的身份说话。因为当祂说\BibleQuote{我饿了，你们给了我吃的，}\BibleRef{玛 25:35}时，祂是代表祂的肢体说话，而非代表祂自己；而当祂说\BibleQuote{扫禄，扫禄，你为什么迫害我？}\BibleRef{宗 9:4}时，是头为它的肢体呐喊；但祂没有说\Quote{你为什么迫害我的肢体？}而是说\Quote{你为什么迫害我？}如果祂在我们内受苦，那么我们也将在祂内得到冠冕。基督的爱就是这样。有什么能与它相比？正是为此，\Quote{祂将一首赞歌放在我们口中，}而祂正是代表祂的肢体说这话。\par

6. \Quote{义人必将看见，并将敬畏，并将信赖主。}\Quote{义人将看见。}谁是义人？就是信者；因为\BibleQuote{义人因信德而生活。}\BibleRef{哈 2:4}因为在教会中有这样的秩序：有些人走在前头，另有些人跟随；走在前头的人使他们自己成为跟随者的\Quote{榜样}；跟随者则效法走在前头的人。但是，那些将自己作为后来者榜样的人，难道他们自己不跟从任何人吗？如果他们根本不跟从任何人，他们就会陷入错误。那么这些人自己也必须跟从某一位，就是基督自己……因此，\Quote{义人}将\Quote{看见，也将敬畏。}他们看见一边是窄路，另一边是\Quote{宽道}：这边看到少数，那边看到许多。但你是个义人；不要数算他们，而要衡量他们；拿出\Quote{公正的天平}，而不是\Quote{诡诈}的天平：因为你被称为义人。\Quote{义人将看见，并敬畏，}这句话适用于你。因此，不要数算那些挤满\Quote{宽路}的人群，他们明天将要挤满竞技场；他们用欢呼庆祝城市的周年纪念，同时却以邪恶的生活玷污城市本身。不要看他们；他们数目众多，谁能数算？但有一些人正在窄路上行走。拿出天平，我说。称量他们：看这一边你举起多少糠秕，那一边只有几粒麦子。让那些将要跟随的\Quote{义人}，即\Quote{信者}，来这样做。那么那些走在前头的人该做什么呢？让他们不要骄傲，不要\Quote{自高自大；}不要欺骗那些跟随他们的人。他们如何欺骗跟随者呢？就是向他们许诺在他们自身内有救恩。那么那些跟随者应当做什么呢？\Quote{义人将看见，并敬畏，并将信赖主；}而不是信赖那些走在前头的人。然而他们确实定睛注视那些走在前头的人，并跟随效法他们；但他们这样做，是因为他们意识到那些人是因谁而获得恩宠走在前头；而且因为他们信赖祂。因此，虽然他们以这些人为榜样，但他们仍信赖那一位，即其他人从祂领受了恩宠，从而成为他们所是的那一位。\Quote{义人将看见这一切，并敬畏，并将信赖主。}正如在另一篇圣咏中，\Quote{我举目向高山，}我们以高山来理解教会中所有杰出而伟大的属灵人物；他们因坚实而伟大，而非因膨胀的虚浮。正是借由他们，全部圣经被分发给我们；他们是先知，他们是福音作者，他们是纯正的导师：向他们我\Quote{举目，从那里我的帮助将临到。}免得你想到仅仅是人的帮助，他接着说：\Quote{我的帮助来自主，祂造了天地。义人将看见这一切，并敬畏，并将信赖主。}……\par

7. \BibleQuote{承蒙那以上主的名号为信赖的人，是有福的，他没有注目的是虚无和说谎的疯狂。} \BibleRef{咏 39:4} 看，那就是你本愿走的道路。看，那 \Quote{充满宽途的人群。} 那条路并非无缘无故通向竞技场。它并非无缘无故通向死亡。\Quote{宽途} 通向死亡，\BibleRef{玛 7:13} 它的宽广令人一时愉悦，终点却是永恒的狭隘。不错，但人群在窃窃私语；人群在共同欢乐；人群在匆忙赶路；人群在聚集！不要效仿他们，不要随从他们：他们是 \BibleQuote{虚无和说谎的疯狂。} 愿上主你的天主成为你的盼望。不要从上主你的天主那里盼望别的东西，而要让上主你的天主本身成为你的盼望。因为许多人盼望从天主手中获得财富，以及许多易朽和短暂的荣誉；总之，他们盼望从天主手中得到任何别的东西，唯独不是天主本身。但你要寻求你的天主本身；不，实际上，轻视一切其他事物，直奔向他！忘记其他，记念他。放下其他，\BibleQuote{直奔向他} \BibleRef{斐 3:14}。的确，正是他亲自纠正了你，当你偏离正路时；如今你已走上正路，他引导你正确前行，他引导你到达目的地。愿他成为你的盼望，他既引导你，也引导你到达目的地。世俗的贪欲将你引向何方？最终又将你带往何处？你起初想要一块田地，接着想拥有庄园；你想驱赶邻居；赶走了他们，你又觊觎其他邻居的土地；你的贪欲扩张，直到海边；到了海边，你又贪图海岛；当你把大地据为己有时，你或许还想占据天空。离弃你所爱的一切吧。创造天地的那一位比一切更美。\par

8. \BibleQuote{承蒙那以上主的名号为盼望的人，是有福的，他没有注目的是虚无和说谎的疯狂。} 为何 \Quote{疯狂} 被称为 \Quote{说谎的}？疯狂是虚假的，正如清醒才看见真理。因为你看到的好事其实是骗人的；你神志不清；一阵高烧使你癫狂；你所迷恋的并非真实事物。你为战车手喝彩，你为战车手欢呼，你疯狂地爱慕战车手。那是 \Quote{虚无，} 那是 \Quote{说谎的疯狂。} \Quote{'不是的'}（他喊道），\Quote{没有什么比这更好、更令人愉快的了。} 对一个高烧不退的人我能做什么？如果你们心中还有一点怜悯，请为这样的人祈祷吧。因为医生在绝望的情况下，通常转向房间里的那些人，他们围在床边哭泣，屏息等待他对病危患者的诊断。医生犹豫不决：他看不到任何好的前景；他不敢说出坏话，以免引起恐慌。他想出一个极其谨慎的说法：\Quote{良善的天主无所不能。请为他祈祷吧。} 那么，这些疯狂的人中，我该制止哪一个？哪一个会听我？哪一个不会认为我们可怜？因为他们以为我们失去了他们疯狂迷恋的许多伟大而多样的快乐，因为我们不像他们那样疯狂地迷恋它们；他们看不到这些快乐是 \Quote{说谎的。} ……\Quote{并且没有注目的是虚无和说谎的疯狂。} \Quote{某某赢了，} 他喊道，\Quote{他套上了某匹马，} 他大声宣告。他想要成为某种占卜者；他因离弃神性的源泉而向往占卜的荣誉；他经常发表意见，又经常出错。为何如此？正是因为他们 \Quote{说谎的疯狂。} 但为何他们说的有时应验？为了迷惑愚昧的人；使他们因喜爱那里真理的表象而陷入谎言的陷阱；让他们被撇下，被 \Quote{放弃，} 被 \Quote{砍掉。} 如果他们是我们身上的肢体，就必须被治死。\BibleQuote{治死你们在地上的肢体。} \BibleRef{哥 3:5} 愿我们的天主成为我们的盼望。创造万有的那位比万有更美！创造美丽的那位比一切美丽更美。创造大能的那位本身更大能。创造伟大的那位本身更伟大。他将成为你所爱的一切。在受造物中学习爱创造者；在作品中爱造物主。不要让被他所造的事物抓住你的感情，以致你失去他，因为你自己也是他所造的。\BibleQuote{有福的，是那以上主的名号为信赖的人，他没有注目的是虚无和说谎的疯狂。}……\par

9. 我们要用其他的景象来交换这些景象。那么，对于这位基督徒，我们想让他远离这些景象，该向他呈现什么景象呢？我感谢上主我们的天主；他在圣咏的下一节中指示我们，对于渴望观看的观众，我们应当呈现和提供什么景象。现在让我们假设他脱离了马戏团、剧院、竞技场；让他寻找，无论如何让他寻找一些可以观看的景象；我们不会让他没有可见之物。那么，我们该用什么来交换那些景象呢？且听下文。\par

\Quote{主、我的天主，祢所行的奇事何其众多}\BibleRef{咏 39:5}。他过去常常注视人的\Quote{奇事}；现在让他默观天主的奇事。天主\Quote{所行的奇事}何其众多。为什么这些在他眼中变得卑贱？他赞美驾驭四马战车的车夫，全都奔驰无瑕、毫不失足。或许主在属灵之事上没有行这样的\Quote{奇事}。让他克制情欲，克制怯懦，克制不义，克制不智——我说的是那些因过激而产生这些恶习的激情；让他克制这些并使之服从，让他掌握缰绳，不容自己被冲走；让他照自己意愿的方向引导它们；不要被迫走向他不愿去的地方。他过去常为车夫喝彩，他自己将因自己的御车术而受喝彩；他过去常喊叫车夫应得荣誉之衣，他自己将穿上不朽。这些是天主向我们展示的景观，这些是景象。祂从天上呼喊：\Quote{我的眼目注视着你。奋斗吧，我必帮助你；得胜吧，我必给你加冕。}\par

\Quote{在祢的思念中，无人能与你相比。}那么现在看看那演员！因为那人费了极大苦功学会了走绳索；悬在那里时，他也让你悬着心。转向那展示更奇妙景观的那一位。这人学会了走绳索；但他曾使另一人在海上行走吗？现在忘记你的剧场；请看我们的伯多禄；他不是走绳索者，而是，可以这样说，海上行者。你也要在别的水上行走（虽然不是伯多禄所行的水，以象征某种真理），因为这世界就是海。它有有害的苦涩；它有患难的波浪，诱惑的风暴；它里面有人，像鱼一样，以自毁为乐，并彼此吞食；你在这里行走，你在这上面踏足。你想看景象；你自己成为一个\Quote{景象}。为免你的精神沉没，请看那在你前面引领你的那一位，祂说：\Quote{我们成了一台戏，给世界、天使和世人看。}\BibleRef{格前 4:9}你踏在水上；不容你自己被淹在海中。你不会去那里，你不会\Quote{践踏它}，除非那是祂的命令，祂自己首先在海面上行走。因为伯多禄是这样说的：\Quote{如果是祢，请叫我在水面上步行到祢那里去。}\BibleRef{玛 14:28}因为\Quote{祂是}，祂在祈祷时听到了他；祂在表达渴望时允准了他的意愿；祂在沉没时扶起了他。这些就是\Quote{主所行的奇事}。以信德为眼目来观看它们。你也要同样行；因为尽管风惊扰你，浪涛向你汹涌，人性的软弱也许使你对自己的得救生出疑虑，你仍能\Quote{呼喊}，你可以说\Quote{主，我丧命了！}\BibleRef{玛 14:30}那命令你行在水上的，不容你丧亡。因为你既行在\Quote{磐石}上，就连在海上也不惧怕！你若没有\Quote{磐石}，必在海中沉没；因为你必须行走的磐石，是不沉在海中的。\par

10. 那么，请注意天主的\Quote{奇事}。\Quote{我传报了，我也说了；它们多得不可胜数。}有一个\Quote{数}，有些超出了这个数。有一个固定之数属于那天上的耶路撒冷。因为\Quote{主认识属于祂的人}\BibleRef{弟后 2:19}；敬畏祂的基督徒，相信的基督徒，遵守诫命的基督徒，行走在天主道路上的基督徒，远离罪恶的基督徒；若跌倒就告解的：他们属于\Quote{那数}。但只有他们吗？还有一些\Quote{超出那数}的。因为即使他们只是少数（与大多数群体的数量相比是少数），我们的圣堂却充满了多么大的数目，挤到墙壁边；他们彼此拥挤到何种程度，几乎因人数过多而窒息。但是，在这些相同的人中，当有公共表演时，又有成群的人涌向竞技场；这些是超出\Quote{那数}的。但我们之所以这样说，是为了使他们能进入\Quote{那数}。他们既然不在这里，就没有从我们这里听到这话；但你们离开这里后，让他们从你们口中听到。\Quote{我传报了，}他说，\Quote{我也说了。}这是基督在说话。\Quote{祂传报了，}以自己的位格，作为我们的元首。祂藉祂的肢体亲自传报了。祂亲自派遣了那些应当\Quote{传报}的人；祂亲自派遣了宗徒。\Quote{他们的声音传遍了天下，他们的言语传到了地极。}聚集起来的信徒何其众多；蜂拥而至的人群何其众多；其中许多真正皈依，许多只是外表如此：真正皈依的是少数；外表如此的是多数：因为\Quote{他们多得超出了那数。}\par

11. …这些是天主的\Quote{奇事}；这些是天主的\Quote{思念}，无人能与之相比；好观看者当戒除好奇，与我们一同寻求那些更卓越、更有益的事，当他得到它们时，必要喜乐……\par

12. 圣咏向天主说：\Quote{祭献和供物，祢不想要}\BibleRef{咏 39:6}。因为古时，当那真正的祭献——就是信徒所认识的——还在以预像预示时，人们曾举行那些仪式，作为将来实体的预像；其中许多人明白其意义，但更多的人对此一无所知。因为先知和圣祖们明白他们所举行的是什么；但那\Quote{顽固的百姓}\Person{}其余的成员是如此属肉体，以致他们所行的只是象征以后要来的事物；于是，当那最初的祭献被废除时，当\Quote{公绵羊、公山羊、牛犊}和其他牺牲的燔祭被废除时，\Quote{天主不想要它们}。为什么天主不想要它们？为什么祂起初要它们？因为所有那些东西，好比是许愿者的话语；当所许的承诺实现时，表达承诺的话语就不再被说出……那些祭献，作为承诺的表达，已被废除。什么作为其应验而被赐予？那就是\Quote{身体}，是你们所认识的，但不是你们所有人都认识；愿天主使那些认识它的人不致因此而被定罪。请注意这句话说出的时间；因为说话者是基督我们的主，有时为祂的肢体说话，有时以祂自己的身份说话。\Quote{祭献和供物}，祂说，\Quote{祢不想要}。那么，现在我们就没有祭献了吗？绝对不可！\par

\Quote{但祢为我预备了一个身体。}正因如此，祢不想要其他的祭献，好让祢\Quote{预备}这个；在祢\Quote{预备}这个之前，祢要其他的。承诺的应验废除了表达承诺的话语。因为如果它们仍在承诺，那么所承诺的尚未应验。这事曾以某些记号被应许；传达承诺的记号被废除了，因为所承诺的实体已经来到。我们是在这个\Quote{身体}内。我们分沾了这个\Quote{身体}。我们知道我们所领受的；而你们这些还不知道的人，不久将要认识它；当你们认识它时，愿你们不致因此而受审判。\Quote{因为那吃喝不相称的，就是吃喝自己的罪案。}\BibleRef{格前 11:29} 一个\Quote{身体}已为我们\Quote{预备}了；愿我们在这个身体内得以成全。\par

13. \Quote{赎罪的燔祭，祢也不要。}\Quote{于是我说：看，我来了！}\BibleRef{咏 39:7} 现在是时候了，\Quote{所承诺的应当来到}，因为那些藉以承诺的记号已被除去。的确，弟兄们，请看这些已被除去，那些已被应验。愿犹太民族此时若能，给我指出他们的司祭！他们的祭献在哪里？它们已被终结，已被除去。难道我们那时就该拒绝它们吗？我们现在拒绝它们，因为如果你如今选择举行它们，那是不合时宜的，不适当的，不协调的。你仍在许愿，而我已经领受了！他们留下某种可举行的事物，好让他们不完全没有记号……他们就像这样，如同带著标记的加音。然而，那些在那里举行的祭献已被除去；而那留给他们的，如同加音的标记，如今也已应验，他们却不知道。他们宰杀羔羊，吃无酵饼。\Quote{基督，我们的逾越节羔羊，已被祭杀。}\BibleRef{格前 5:7} 看，在基督的祭献中，我认出那被宰杀的羔羊！至于无酵饼呢？\Quote{所以，}他说，\Quote{我们过节，不可用旧酵母，也不可用奸恶的酵母}（他指明\Quote{旧}什么意思；那是\Quote{陈腐}的面粉，是酸的），\Quote{而要用真诚和正直的无酵饼。}\BibleRef{格前 5:8} 他们仍留在阴影中，不能承受荣耀的太阳。我们已在白日之中。我们有基督的\Quote{身体}，我们有基督的血。如果我们获得新生命，让我们\Quote{高唱新歌，歌颂我们的天主。}\Quote{赎罪的燔祭，祢也不要。于是我说：看，我来了！}\par

14. \Quote{书卷上关于我已有记载：我乐意承行祢的旨意；我的天主，是的，我乐意，祢的法律在我心中。}\BibleRef{咏 39:8} 看！祂转向祂的肢体。看！祂亲自\Quote{承行了}圣父的\Quote{旨意}。但这是写在什么\Quote{书卷的开端}呢？也许是在这本\WorkTitle{圣咏集}的开端。因为我们何必远求，或查考其他书卷呢？看！它写在这本\WorkTitle{圣咏集}的开端！\Quote{他的旨意是在上主的法律中}，就是说：\Quote{我的天主，我乐意，祢的法律在我心中}，这也就是\Quote{他昼夜默思祂的法律。}\par

15. \BibleQuote{我已在伟大的会众中宣扬了祢的义德} \BibleRef{咏 39:9}。祂现在向祂的肢体说话。祂是在劝勉他们去做祂已经做过的事。祂曾\Quote{宣扬}，让我们也宣扬。祂曾受苦；让我们\Quote{与祂同受苦}。祂已受光荣；我们将\Quote{与祂同受光荣}。\BibleRef{罗 8:17} \Quote{我已在伟大的会众中宣扬了祢的义德。}那是多么伟大的会众呢？就是在普世。有多伟大？甚至遍布万邦。为什么在万邦？因为祂是\BibleQuote{亚巴郎的后裔，在\Emph{祂}内万邦都将蒙受祝福。} \BibleRef{创 22:18}为什么在万邦？\Quote{因为他们的声音传遍了全地。} \Quote{看哪！上主，我不会管束我的口唇，祢也知道这一点。}我的口唇说话，我不会\Quote{管束}它们不出声。我的口唇确实在人的耳中发出可听见的声音；但祢\Quote{知道}我的心。\Quote{上主，我不会管束我的口唇；祢知道这一点。}人听到是一回事，天主\Quote{知道}是另一回事。为使这\Quote{宣扬}不只局限于口唇，免得我们被人说：\BibleQuote{凡他们对你们所说的，你们要行，但不要照他们的行为去做；} \BibleRef{玛 23:3}或者免得对百姓说：\Quote{用口唇赞美天主，心却不亲近祂，} \BibleQuote{这百姓用口唇尊敬我，他们的心却远离我；} \BibleRef{依 29:31}你当以口唇发出可听的宣认，也当以心亲近。\BibleQuote{因为人心里相信，以致成义；口里宣认，以致获得救恩。} \BibleRef{罗 10:10}正如那个盗贼的情形，他与主同挂在十字架上，在十字架上承认了主。别人在祂行奇迹时拒绝承认祂，这人却在悬挂十字架时承认了祂。那盗贼的其他肢体都被钉穿了；双手被钉子钉住，双足也被刺透；全身被钉在木架上；身体的其他部分都不能活动；只有心和舌头是自由的；他\Quote{以心} \Quote{相信}，\Quote{以口}作了\Quote{宣认}。他说道：\Quote{主，当祢进入祢的王国时，请记念我。}他盼望他的救恩在遥远的时候到来；他满足于在长久的等待之后接受它；他的希望寄托在遥远的目标上。然而，那日子并未推迟！回答是：\Quote{今天你就要同我在乐园里。}乐园有喜乐的树木！今天你与我在十字架的\Quote{木架}上；今天你将要与我在救恩的\Quote{木架}上。……\par

16. \BibleQuote{我没有将我的义德隐藏在心里} \BibleRef{咏 39:10}。所谓\Quote{我的义德}是什么意思？就是我的信德。因为\Quote{义人必因信德而生活。}假设迫害者以刑罚相威胁，就像他们曾经被允许做的那样，向你提问：\Quote{你是什么人？外教人还是基督徒？} \Quote{基督徒。}这就是他的\Quote{义德}。他相信，他\Quote{因信德而生活。}他并没有\Quote{将他的义德隐藏在心里。}他没有在心里说：\Quote{我确实相信基督，但我不会告诉这位向我发怒并威胁我的迫害者我所信的是什么。我的天主知道，我在内心里确实相信。祂知道我没有否认祂。}看哪！你说你内心有这样的事！你口唇上说的是什么？\Quote{我不是基督徒。}你的口唇作证反对你的心。\Quote{我没有将我的义德隐藏在心里。}……\par

17. \Quote{我宣扬了祢的真理和祢的救恩。}我宣扬了祢的基督。这就是\Quote{我宣扬了祢的真理和祢的救恩}的含义。\Quote{祢的真理}如何是基督？\BibleQuote{我是真理。} \BibleRef{若 14:6}基督如何是\Quote{祢的救恩}？西默盎在圣殿中认出了在祂母亲手中的婴孩，并说：\BibleQuote{因为我的眼睛看见了祢的救恩。} \BibleRef{路 2:30}这位老人认出了这婴孩；这位老人自己在那婴孩身上\Quote{成了小孩子}，因信德而更新。因为他曾从天主领受了一则神谕；它这样说：主曾对他说，在他未看见\Quote{天主的救恩}以前，他不得离开此生。这\Quote{天主的救恩}是一件向人显示的好事；但让他们呼喊：\Quote{上主，求祢向我们显示祢的仁慈，并把祢的救恩赐给我们。}……\par

18. \Quote{我没有向伟大的会众隐瞒祢的仁慈和祢的真理。}让我们也在那里；让我们也被列在这身体的肢体中：让我们不要保留主的\Quote{仁慈}和主的\Quote{真理}。你想听\Quote{主的仁慈}是什么吗？离开你的罪恶，祂会赦免你的罪恶。你想听\Quote{主的真理}是什么吗？持守义德。你的义德将获得冠冕。因为现在向你宣布了仁慈；此后\Quote{真理}将向你显示。因为天主并非只仁慈而不公义，也非只公义而不仁慈。那仁慈在你看来微不足道吗？祂不会将你从前的一切罪过归咎于你：你至今生活败坏，你仍然活着；今天好好生活，那么你就不会\Quote{隐瞒}这\Quote{仁慈。}如果这是\Quote{仁慈}的意思，那么\Quote{真理}又是什么意思呢？……\par

19. \Quote{主，求祢不要向远处收回祢的仁慈} \BibleRef{咏 39:11}。祂正在转向受伤的肢体。因为我没有\Quote{从大会中隐瞒祢的仁慈和祢的真理}，即从普世教会的合一中，求祢垂顾祢受苦的肢体，垂顾那些犯有忽略之罪和行为之罪的人：不要收回祢的仁慈。\Quote{祢的仁慈和祢的真理时常保护了我。}我不敢离弃我的恶道，除非我确知赦免；我无法坚持到底，除非我确知祢应许的实现……\par

\Quote{无数的祸患围绕着我} \BibleRef{咏 39:12}。谁能数清罪恶？谁能算清自己和别人的罪？这是那叹息之人的重担，他曾说：\Quote{求祢洗净我隐藏的过犯；也求祢从别人的过犯中宽恕祢的仆人，主。}我们自己的已经够多，还要加上\Quote{别人的}重担。我为自己害怕；我为有德行的弟兄害怕，我还要忍受邪恶的弟兄；在这样的重担下，若天主的仁慈缺失，我们将如何？\Quote{但是祢，主，不要远远离开。}求祢亲近我们！主亲近谁？\Quote{就是}那些\Quote{心碎的人}。祂远离骄傲的人；祂亲近谦卑的人。\Quote{因为主虽崇高，却垂顾卑微的人。}但骄傲的人不要以为自己不被看见：因为高处的，祂\Quote{从远处观看}。祂\Quote{从远处观看}那自夸的法利塞人；祂却就近帮助那做告解的税吏。\BibleRef{路 18:9} 一个夸耀自己的功德，隐藏自己的创伤；另一个不夸耀功德，却暴露创伤。他来到医生那里；他知道自己病了，需要痊愈；他\Quote{不敢举目望天，只是捶着胸膛}。他不宽恕自己，好让天主宽恕他；他承认自己有罪，好让天主\Quote{忽略}对他的指控。他惩罚自己，好让天主释放他免于惩罚……\par

20. \Quote{我的过犯抓住了我，以致我不能看见。}有什么事是我们\Quote{要看见}的？是什么阻止我们，使我们看不见？岂不是过犯吗？你的眼睛也许被某种渗入的眼液、或烟尘、或别的投入之物妨碍了，以致不能看见这光；你怎能举起受伤的眼睛凝视这白日之光呢？那么，你能举起受伤的心朝向天主吗？岂不需先被医治，才能看见？你说：\Quote{先让我看见，然后我才相信}，这不是显出你的骄傲吗？有谁这样说？因为那渴望看见的人，谁会说出\Quote{让我看见，然后我才相信}？我要向你显示光，或者更确切地说，光本身愿意向你显示自己！向谁？它不能向瞎子显示。他看不见。为什么看不见？因为眼睛被众多罪恶堵塞了……\par

21. \Quote{它们多过我头上的头发。}他把\Quote{头上头发}的数目置于计算之下。谁能数算自己头上的头发？更不用说他的罪，超过头上头发的数目。它们看似微小，却数量众多。你已防备了大罪；你现在不犯奸淫、不杀人、不抢夺他人财物、不亵渎、不作假见证；这些是较重的罪。你防备了大罪，对小罪你怎么办？你丢掉了重担，当心别被沙土淹没。\Quote{而我的心离弃了我。}你的心被你的天主离弃，有什么稀奇，既然连它自己都\Quote{离弃}自己？所谓\Quote{衰败}、\Quote{离弃}是什么意思？就是不能认识自己。祂的意思是：\Quote{我的心离弃了我。}我愿用心看见天主，却因众多罪恶不能；这还不够，我的心甚至不认识自己。因为无人完全认识自己：谁也不要凭自己的状态自恃。伯多禄能用心领悟自己心的状态吗？他曾说：\Quote{我愿与祢同死}。\BibleRef{路 22:33} 心中存着虚假的自信，同时隐藏着真实的恐惧；心不能领悟心的状态。这状态对生病的心自己也是未知，却对医生是显明的。预言在他身上应验了。天主知道他里面连他自己也不知道的，因为他的心离弃了他，他的心不为自己所知。\par

22. \Quote{主，求祢乐意拯救我} \BibleRef{咏 39:13}。好像祂在说：\Quote{\InnerQuote{祢若愿意，就能使我洁净。}\BibleRef{玛 8:2} 求祢乐意拯救我。主，求祢垂顾帮助我。}垂顾，就是垂顾痛悔的肢体，那些在痛苦中、在医生的手术刀下扭动的肢体；但仍然怀着希望。\par

23. \Quote{愿那些寻索我性命、要毁灭它的人，一同蒙羞受辱} \BibleRef{咏 39:14}。因为在另一处，祂提出指控说：\Quote{我向右观看，看见没有人寻找我的灵魂；}即没有人效法我的榜样。基督在受难中是发言者。\Quote{我向右观看，}即不是看那不虔敬的犹太人，而是看我右边的人，宗徒们——\Quote{并没有人寻找我的灵魂。}如此彻底地没有人\Quote{寻找我的灵魂}，以致那自负己力的人\Quote{否认了我的灵魂。}但人的灵魂被寻索有两种方式：或是为享受他的陪伴，或是为迫害他；因此，祂在这里论及另一等人，祂愿他们\Quote{蒙羞受辱}，这些人是\Quote{寻索祂灵魂}的。但为使你不至于理解为如同祂抱怨某些人不\Quote{寻索祂灵魂}时那样，祂加上\Quote{要毁灭它；}即他们寻索我的灵魂，为置我于死地……\par

24. \Quote{愿那些向我行恶的人，转身受辱。} \Quote{转身。}我们不要从坏的意义理解。祂是愿他们好；这是祂的声音，祂曾从十字架上说：\Quote{父啊，赦免他们，因为他们不知道他们做的是什么。} \BibleRef{路 23:34}。那么祂为什么对他们说，他们应当\Quote{转身}呢？因为他们以前骄傲，以致跌倒，现在变得谦卑，以便重新站起。因为当他们在前时，他们想要越过他们的主，比祂更好；但如果他们跟在祂后面，他们就承认祂比他们好；他们承认祂应当在前，祂领先，他们跟随。因此，祂这样斥责给祂出坏主意的伯多禄。因为主为了我们的救恩将要受苦时，也预言了这受难本身要发生的事；伯多禄说：\Quote{愿这事远离祢，} \BibleRef{玛 16:22} \Quote{天主不许！} \Quote{这事不可临到祢！}他原想走在主前面，给他的师傅出主意！但主为使他不是走在祂前面，而是跟在祂后面，就说：\Quote{撒殚，退到我后面去！}正因这个缘故，祂说\Quote{撒殚，}因为你试图走到祂前面，而你应该跟随祂；但如果你在后面，如果你跟随祂，你从此就不再是\Quote{撒殚。}那么怎样呢？\Quote{在这磐石上，我要建立我的教会。} \BibleRef{玛 16:18}……\par

25. \Quote{愿那些向我说\InnerQuote{好啊！好啊！}的人，快快蒙羞退去。} \BibleRef{咏 39:15}。他们无故称赞你。\Quote{伟人！好人！有教养有学问的人；但为什么是基督徒呢？}他们称赞你身上那些你本不愿被称赞的事；他们责怪你所欢喜的事。但如果你也许说：\Quote{人啊，你称赞我什么？我是一位有德行的人？一位义人？如果你这样想，是基督使我成为这样的；要赞美祂。}但对方说：\Quote{愿这事远离你。不要亏待自己！是你自己使自己这样的。} \Quote{愿那些向我说\InnerQuote{好啊！好啊！}的人蒙羞。}然后呢？\par

\Quote{愿所有寻求祢的人，主啊，都喜乐欢欣。} \BibleRef{咏 39:16}。那些不\Quote{寻求}我，而是\Quote{寻求祢}的人；他们不对我说\Quote{好啊！好啊！}而是看见我\Quote{以祢夸耀}，如果我有什么可夸耀的；因为\Quote{夸耀的，应当因主而夸耀。} \BibleRef{格前 1:31} \Quote{愿所有寻求祢的人，主啊，都喜乐欢欣。}\par

\Quote{并要不断地说：愿主受尊崇。}因为即使罪人成为义人，你也应将荣耀归于\Quote{那使不虔敬者成义的天主。} \BibleRef{罗 4:5}所以，无论是罪人，愿召他悔改得宽恕的天主受赞美；或是已经走在义德道路上的人，愿召他接受华冠的天主受赞美！愿主的圣名不断被\Quote{那些爱慕祢救恩的人}所尊崇。\par

\Quote{至于我，} \BibleRef{咏 39:17}。我，他们为我图谋恶事，他们\Quote{谋害我的性命，要把它夺去。}但请你转向另一类人。至于我，他们对我说：\Quote{好啊！好啊！} \Quote{我贫穷且困乏。}在我内没有什么可被夸为属于我自己的。愿祂撕裂我的麻衣，用祂的袍子遮盖我。因为\Quote{现在活着的，不再是我，而是基督在我内生活。} \BibleRef{迦 2:20}如果是基督\Quote{在你内生活}，你所有的一切都是基督的，你将来所有的一切也是基督的；那么你自己是什么？\Quote{我贫穷且困乏。}如今我不富足，因为我不骄傲。那说\Quote{天主，我感谢祢，因为我不像其他人}的人是富足的； \BibleRef{路 18:11}但税吏贫穷，他说：\Quote{天主，可怜我这个罪人吧！}前者饱足打嗝；后者因贫乏而哀哭：\Quote{我贫穷且困乏！}那么，贫穷困乏的人，你要做什么？在天主门前乞讨；\Quote{敲吧，就给你们开。} \BibleRef{玛 7:7}——\Quote{至于我，我贫穷且困乏。然而主眷顾我。}——\Quote{将你们的挂虑托付给主，祂必成就。}你靠自己操心能成就什么，能为自己提供什么？让那造你的主\Quote{眷顾你。}祂在你尚未存在时就眷顾了你，如今你成为祂愿意你成为的，祂岂会不眷顾你？因为如今你是一个信徒，如今你行走在\Quote{义德的道路上。}那\Quote{使太阳上升，照耀善人和恶人，降雨给义人和不义的人}的主，岂会不眷顾你？\BibleRef{玛 5:45}……\par

\BibleQuote{祢是我的助佑，我的拯救者；我的天主，求祢不要迟延。}\BibleRef{咏 39:17}。他正在呼求天主，恳求祂，害怕自己会跌倒：\BibleQuote{不要迟延。}\BibleQuote{不要迟延}是什么意思？我们最近读到关于苦难日子的经文：\BibleQuote{除非那些日子被缩短，凡有血肉的都不会得救。}\BibleRef{玛 24:22}基督的肢体——遍及各地的基督的身体——正如一个人，一个穷人，一个乞丐，在向天主祈求！因为祂也曾成为贫穷，祂\BibleQuote{原是富有的，为了你们却成了贫穷，好使你们因祂的贫穷而成为富有的。}\BibleRef{格后 8:9}是祂使真正的穷人富足，使假冒的富人贫穷。祂向祂呼号：\BibleQuote{我从地极呼求祢，当我心灰意冷时。}将会有苦难的日子，更大的苦难；它们会来到，正如圣经所说：随着时日推进，苦难也增多。谁也不该自行许诺福音所没有许诺的....\par

\SourceRecord{Translated by J.E. Tweed.；Nicene and Post-Nicene Fathers, First Series；Vol. 8.；Edited by Philip Schaff.；Buffalo, NY: Christian Literature Publishing Co.,；1888.；<http://www.newadvent.org/fathers/1801040.htm>.}

% structure: p0001 source=h1 target=section reason=page-title

\section{《圣咏集》第四十一篇诠释}

致民众，论殉道者庆节。\par

1. 殉道者的庄严日已破晓；因此，为光荣基督——殉道者的统帅——的苦难，祂没有顾惜自己，命令祂的士兵上阵；但祂首先战斗，首先得胜，为要以自己的榜样鼓励他们作战，以祂的尊威协助他们，以祂的应许给他们加冕：让我们从这篇圣咏中听一些关于祂苦难的教导。我时常向你们推荐，也不厌其烦地重复，这于你们有所裨益，要记住：我们的主耶稣基督常以祂自己——即我们的头——的位格发言，也常以祂的身体——即我们和祂的教会——的位格发言；但如此，话语如同出自一人之口，好使我们理解头和身体在完整合一中结合，彼此不离；正如在婚姻中所说的：\BibleQuote{二人成为一体}。如果我们承认在一肉体内有两者，那么也让我们承认在一声音内有两者。首先，我们应和读经员所咏唱的，虽然取自圣咏的中段，我将从那开始这篇讲道。\par

\BibleQuote{我的仇敌恶言毁谤我说：祂一死，祂的名字必被遗忘} \BibleRef{咏 40:5}。这是我们的主耶稣基督的位格；但请看，此处是否也包含了肢体？当我们的主自己在地上肉身行走时，这话也说过……他们看见众人跟随祂，就说：\BibleQuote{祂一死，祂的名字必被遗忘；}就是说，我们杀了祂，祂的名字在地上就不再存在，祂死了也不能再迷惑任何人；但因着杀害祂，人们会明白他们所跟随的不过是一个人，在祂内没有救恩的希望，就会弃绝祂的名字，它就不再存在。祂死了，祂的名字并未被遗忘，反而像种子一样被播撒：祂死了，但祂是一粒麦子，死了，却立刻发芽。\BibleRef{若 12:24} 当我们的主耶稣基督受光荣后，人们开始更多、更大量地信靠祂；于是祂的肢体开始听到头所听到的。如今我们的主耶稣基督在天上坐于右边，而祂自己仍在我们内在地上劳苦，敌人仍说：\BibleQuote{祂一死，祂的名字必被遗忘。}因为魔鬼在教会中挑起迫害，要消灭基督的名字。弟兄们，难道你们以为，那些异教徒向基督徒发怒时，他们彼此间没有说\Quote{要从地上抹去基督的名字}吗？为使基督再死一次，不在头内，而在祂的身体内，殉道者也被杀害。圣洁的鲜血倾流，有助于教会的增长；殉道者的死亡也帮助了它的播种。\BibleQuote{上主眼中看为宝贵的，是祂圣徒的死亡。}基督徒越发增多，敌人所说的\BibleQuote{祂一死，祂的名字必被遗忘}并未实现。如今这话仍被说。异教徒坐下计算年代，听他们的狂热者说：将有一个时候，基督徒不再存在，那些偶像将如从前一样受崇拜；他们仍说：\BibleQuote{祂一死，祂的名必被遗忘。}两次被征服，如今第三次当明智！基督死了，祂的名字并未消亡；殉道者死了，教会更增殖，基督的名字在万国中增长。祂预言了自己的死亡和复活，祂预言了殉道者的死亡和他们的冠冕，祂预言了祂教会未来之事，如果祂说了两次真话，难道第三次就撒谎吗？所以你们反对祂所信的是虚妄；更好的，是你们信祂，好使你们\BibleQuote{思念那贫苦的人}；祂\BibleQuote{本是富有的，却为了你们成了贫穷的，使你们因祂的贫穷而成为富有的。} \BibleRef{格后 8:9}……\par

2. \BibleQuote{那思念贫苦者的，是有福的：在邪恶的日子，主必拯救他} \BibleRef{咏 40:1}。因为邪恶的日子必要来到：无论你愿不愿意，它必来到；审判的日子必要临于你，若你\Quote{不思念贫苦者}，那便是邪恶的日子。因为你如今不愿相信的，终将显明。但你在它显明时也不能逃脱，因你在他隐藏时不信。你被邀请，要相信你所不见的，免得你看见时羞愧。\BibleQuote{那么，思念那贫苦者}，就是基督：在祂内认出隐藏的财富，你看见祂是贫穷的，\BibleQuote{在祂内蕴藏着一切智慧和知识的宝藏} \BibleRef{哥 2:3}。因此，在邪恶的日子，祂必拯救你，因为祂是天主；但作为人，祂将那属于人性的复活并改变为更好，祂把你提升到天上。但祂是天主，祂愿意在人身内并与人类共有一个位格，祂不能减少或增加，不能死或复活。祂因人的软弱而死，但天主不死……但正如我们正确地说：某人死了，虽然他的灵魂不死；同样我们正确地说：基督死了，虽然祂的天主性不死。为什么祂死了？因为贫苦者。不要让祂的死冒犯你，使你转离注视祂的天主性。\BibleQuote{那思念贫苦者的，是有福的。}也要思念穷人、贫乏者、饥饿口渴者、赤身露体者、患病者、被囚者；也思念这样的穷人，因为如果你思念他们，就是思念那说过的话：\BibleQuote{我饿了，我渴了，我作客，我赤身露体，我患病，我在监里} \BibleRef{玛 25:35}；如此，在邪恶的日子，主必拯救你……\par

3. \BibleQuote{不要把他交在敌人手中}\BibleRef{咏 40:2}。敌人就是魔鬼。愿没有人听见这些话时，把一个人想作是自己的敌人。或许有人想到自己的邻人，想到与自己有诉讼的人，想到要夺走他产业的人，想到强迫他出卖房屋的人。不要这样想；而要想到那敌人，就是主所说\Quote{\BibleQuote{敌人做了这事}}\BibleRef{玛 13:28}的那一位。因为正是他唆使人为了尘世的事物而崇拜他，因为这个敌人不能推翻基督徒的名号。因为他看见自己被基督的名声和赞美所征服，他看见，当他杀害基督的殉道者时，他们却得享冠冕，他反被战胜了。他开始无法说服人相信基督算不得什么；而且因为现在藉着辱骂基督很难欺骗人，他便藉着赞美基督来企图欺骗。在这之前他说了什么？你敬拜谁？一个犹太人，死了，钉在十字架上，一个无足轻重的人，甚至不能使自己摆脱死亡。当他在基督的名号之后看见全人类奔跑，看见在被钉十字架者的名号下，庙宇被推倒，偶像被打碎，祭献被废止；所有这些在先知书中预言的事被人类思量，人类惊讶不已，且现在封闭了他们的心，抵挡对基督的辱骂；他便用赞美基督来装扮自己，开始以另一种方式使人远离信德。基督的法律何其伟大，那法律何其有力、神圣、不可言喻！但谁能履行它？因我们救主的名，\Quote{践踏狮子和毒龙}。公开辱骂时，狮子咆哮；狡猾赞美时，毒龙潜伏。让那些怀疑的人来接受信德吧；不要说，谁能履行它？如果他们依靠自己的力量，他们就不能履行。让他们依靠天主的恩宠而相信，依靠（恩宠）而来；来是为得到帮助，不是为受审判。所有信友都这样生活在基督的名号下，每个人按自己的程度履行基督的命令，无论是已婚的，还是独身和贞女，他们尽天主所赐的生活；他们不依靠自己的力量，而知道他们应在祂内夸耀……\par

4. \BibleQuote{愿主扶持他}\BibleRef{咏 40:3}。但什么时候？或许在天上，或许在永生，以致仍然要为尘世的需要、为今生的急需而崇拜魔鬼。决不可！你拥有\BibleQuote{现今和将来生命的应许}\BibleRef{弟前 4:8}。那创造天地的主来到你面前，在地上。那么，请思量祂所说的：\Quote{愿主扶持他，在他的病榻上}。病榻就是肉身的软弱；免得你说：我不能管束、承担、约束我的肉身；你得到帮助，以便能够。愿主在你的病榻上扶持你。你的床榻背负着你，你并没有背负床榻，但你是内在的瘫子；祂来对你说：\BibleQuote{拿起你的床，回家去吧}\BibleRef{谷 2:11}。\Quote{愿主扶持他，在他的病榻上}。然后祂转向主自己，好像有人问：既然主帮助我们，为什么我们在今生要遭受如此巨大的痛苦、如此巨大的绊脚石、如此巨大的劳苦、来自肉身和世界的如此烦扰？祂转向天主，好像向我们解释祂医治的计划，祂说：\Quote{你在他软弱中，翻转了他一切的床榻}。床榻指一切属地的事物。每一个在这生命中软弱的灵魂，都为自己寻找某种可以安息的东西，因为劳苦的强度以及灵魂朝向天主延伸，它难以持久地忍受；它寻求地上某种可以安息的东西，并且仿佛以一种暂停的方式躺卧，如同无辜者所喜爱的事物……无辜者安息在他的房屋、他的家庭、他的妻子、他的儿女中；在他的贫穷中，他的小农场、他亲手种植的果园、他亲手建造的建筑物中；在这些中，无辜者安息。然而，天主不愿我们只爱永生，即使在这些纯洁的欢愉中，也掺入苦涩，使我们在这些中也遭受磨难，因此祂翻转我们软弱中一切的床榻。\Quote{你在他软弱中，翻转了他一切的床榻}。那么，当他因自己无辜拥有的事物而遭受磨难时，不要抱怨。他藉着更坏之事的苦涩，被教导去爱更好之物；免得他作为旅客前往家乡时，选择了旅店而非自己的家。\par

5. 但为何如此？因为祂\BibleQuote{鞭打祂所接纳的每一个儿子}\BibleRef{希 12:6}。为何如此？因为对犯罪的人说过：\BibleQuote{你必汗流满面，才有饭吃}\BibleRef{创 3:19}。因此，因为所有这些责罚——在我们的软弱中我们所有的床榻都被翻转了——人应当承认他因罪而受苦；让他转向自己，说出以下的话：\BibleQuote{我说：上主，求祢怜悯我，医治我的灵魂，因为我得罪了祢}\BibleRef{咏 40:4}。噢，上主，愿祢藉着磨难来操练我；祢判定要鞭打每一个祢将接纳的儿子，连独生子也未放过。祂确实没有罪而受鞭打；但我却说：\BibleQuote{我得罪了祢}\BibleRef{咏 40:4}……\par

6. \BibleQuote{我的仇敌恶言议论我说：他若死了，他的名字必归于无有}\BibleRef{咏 40:5}。关于这一点我们已经说过，并且由此开始。\par

7. \BibleQuote{进来观看}\BibleRef{咏 40:6}。基督所受苦的，教会也同样受苦；头所受苦的，肢体也同样受苦。\BibleQuote{因为门徒不能超过师傅，仆人不能大过主人}\BibleRef{玛 10:24}……\par

如果你属于基督的肢体，就进来，紧贴着头。\BibleQuote{如果你是麦子，就容忍莠子；如果你是谷粒，就容忍糠秕}\BibleRef{玛 13:30}。如果你是好鱼，就容忍网中的坏鱼。为什么在簸扬之前逃走？为什么在收割之前连自己也把麦子连根拔起？为什么在到达岸边之前就扯破网？\Quote{他们出去传扬}。\par

8. \BibleQuote{我的一切仇敌都齐声议论我} \BibleRef{咏 40:7}。他们都齐声反对我。与我同心合意，胜过反对我而\Quote{齐声}。什么是\Quote{齐声反对我}？就是同谋，一致筹划。基督于是对你们说：你们已同意反对我，现在请同意跟我一起吧：为何反对我？为何不与我一致？如果你们始终持守那同一件事，你们就不会分裂为教派。因为宗徒说：\BibleQuote{弟兄们，求你们大家言谈一致，在你们中不要有分裂。} \BibleRef{格前 1:10} \Quote{我的一切仇敌都齐声议论我：}他们竟\Quote{向我设计为恶}。其实是向他们自己，因为他们\Quote{为自己积蓄了不义}；但也是向我，因为他们的意图要受衡量：并非因为他们无力作恶；他们作恶是出于意志。魔鬼渴望消灭基督，犹达斯想要杀害基督；然而基督被杀并复活，我们得以存活，但魔鬼和犹达斯所得的，是他们恶意的报应，而非我们的救恩……记录他们曾对祂说恶言的那一位，所看的是他们说话时的意图，而非他们所说的话，\Quote{他们向我设计为恶。} 他们对基督行了什么恶？对殉道者行了什么恶？天主已把一切转为善。\par

9. \BibleQuote{他们向我口出恶言} \BibleRef{咏 40:8}。是什么样的恶言？聆听元首亲自说的。\BibleQuote{来吧，我们杀了他，产业就归我们了。} \BibleRef{谷 12:7} 愚人啊！产业怎能归你们呢？因为你们杀了他？看！你们甚至杀了他，产业却仍不是你们的。\Quote{那睡了的人岂不也会起来吗？}当你们因杀了他而得意时，他睡了；因为他在另一篇圣咏中说：\Quote{我睡了。}他们发怒要杀我；\Quote{我睡了。}如果我不愿意，我连睡也不会睡。\BibleQuote{我有权舍掉我的性命，也有权再取回它来。} \BibleRef{若 10:18} \Quote{我躺下睡了，又起来了。}任凭犹太人发怒；任凭\BibleQuote{大地被交在恶人手中，} \BibleRef{约 9:24} 任凭肉身落在迫害者手中，让他们把它悬在木架上，用钉子穿透，用长矛刺透。\Quote{那睡了的人，岂不也加上这事，就是起来吗？}为什么他睡了？因为\BibleQuote{亚当是那将来者的预像。} \BibleRef{罗 5:14} 亚当睡了，从他肋骨形成了厄娃。\BibleRef{创 2:21} 亚当预像基督，厄娃预像教会；因此厄娃被称为\BibleQuote{众生之母。} \BibleRef{创 3:20} 厄娃是在何时受造的？当亚当睡着时。基督肋旁流出的教会圣事是在何时？当他在十字架上睡了时……。\par

10. \BibleQuote{我素来信任的盟友，吃过我饭的人，竟抬脚踢我} \BibleRef{咏 40:9}：抬脚踢我，要践踏我。这我素来信任的盟友是谁？是犹达斯。基督信任他吗？竟说了\Quote{我所信任的}？祂难道从起初不认识他？难道在他出生之前不知道他会如此？祂不是曾对众门徒说：\BibleQuote{我拣选了你们十二人，但你们中有一个是魔鬼} \BibleRef{若 6:70}？那么基督怎会信任他呢？只因祂是在祂的肢体中，又因许多信徒信任了犹达斯，主便将这事归到自己身上……\Quote{我素来信任的盟友，吃过我饭的人。} 在受难时，如何指认出他来？借着自己预言的言语：用蘸饼指认了他，为应验经上论他所说的：\BibleQuote{吃过我饭的人。} \BibleRef{若 13:26} 再者，当他来出卖祂时，祂让他亲吻，\BibleRef{玛 26:49} 为应验经上论他所说的：\Quote{我素来信任的盟友。}\par

11. \BibleQuote{但是，上主，求你怜悯我} \BibleRef{咏 40:10}。这是僕人的身份，这是穷困贫乏者的身份：因为\Quote{眷顾穷困贫乏者的人是有福的。} 看，正如所说的：\Quote{求你怜悯我，使我起来，我要报复他们，} 事就这样成了。因为犹太人杀害基督，以免失去他们的地方。\BibleRef{若 11:48} 基督被杀，他们就失去了地方。他们被从王国中连根拔除，分散各处。祂复活后，用磨难报复了他们，但只是为警告，尚未到定罪。因为那城，就是百姓如咆哮的狮子，狂喊\BibleQuote{钉死祂，钉死祂} \BibleRef{路 23:21}的那城，犹太人从那里被连根拔除，现在拥有基督徒，已无一个犹太人居住。在那里栽种了基督的教会，而会堂的荆棘被拔除。因为这场火确实如\Quote{荆棘的火}般燃烧。但主却如绿树。祂自己曾说过，当某些妇女哀悼祂临终时……\Quote{因为如果他们对绿树这样做，对枯木又将如何？} 绿树何时能被荆棘的火烧尽？因为荆棘的火燃烧如刺丛中的火。火能烧尽荆棘，但绿树不易点燃。……但恐你们以为基督的天主父能复活基督，即祂圣子的肉身，而基督自己虽是等于圣父的圣言，却不能使自己的肉身复活；请听福音：\BibleQuote{拆毁这座圣殿，三天内我要把它重建起来。} \BibleRef{若 2:19} \Quote{但是}，圣史说（恐我们听后仍迟疑），\Quote{他所说的圣殿，是指祂的身体。} \Quote{请使我起来，我要报复他们。}\par

12. \BibleQuote{由此，我知道祢眷顾我，因为我的仇敌不得向我夸胜。}\BibleRef{咏 40:11} 因为犹太人当时确实夸胜了，当他们看见基督被钉十字架时；他们以为他们的恶意得逞了：他们亲眼看见他们残酷行为的果实，基督悬在十字架上：他们摇着头说，\BibleQuote{祢如果是天主子，就从十字架上下来吧！}\BibleRef{玛 27:39} 祂没有下来，虽然祂能够；祂没有显示祂的大能，而是教给了忍耐。因为如果他们这样说时，祂从十字架上下来，祂就会像是屈服于他们的侮辱，并因不能忍受羞辱而被认为失败了：祂在十字架上比他们侮辱时更坚定；祂稳固，他们动摇。所以他们没有依附于真正的元首。祂明白地教给了我们忍耐。因为祂所做的更为有力，祂没有做犹太人所挑战的事。因为从坟墓中复活远比从十字架上下来更为有力。\BibleQuote{我的仇敌不得向我夸胜。}他们那时夸胜了。基督复活了，基督受光荣了。如今他们看见在祂的名下人类归化：如今让他们侮辱，如今摇头吧；更确切说，如今让他们稳住头，或者如果他们摇头，让他们在惊奇和赞叹中摇头……\par

13. \BibleQuote{至于祢，祢支持我，因为我的无辜}\BibleRef{咏 40:12}。真是无辜；无玷的纯真，不求回报，无罪而受鞭笞。 \BibleQuote{祢因我的无辜支持我，并使我永远在祢面前坚强。}祢使我永远坚强，祢使我一时软弱；祢使我在祢面前坚强，祢使我在人面前软弱。那么，什么？赞美归于祂，光荣归于祂。\BibleQuote{上主以色列的天主是应当称颂的。}因为祂是以色列的天主，我们的天主，雅各伯的天主，幼子之天主，年幼民族之天主。没有人说，这是祂对犹太人说的，我不是以色列；反而犹太人不是以色列。因为长子，他是被弃的长子民族；幼子，是蒙爱的民族。\BibleQuote{大的要服事小的。}\BibleRef{创 25:13}如今已应验：如今，弟兄们，犹太人服事我们，他们如同我们的书童，我们研究，他们背着我们的书。请听犹太人如何服事我们，而且不无道理……在他们那里有法律和先知，在这法律和先知中，基督被宣讲。当我们面对外邦人时，展示这在基督教会中实现的事，这些事早先关于基督的名、关于基督头和身的预言已被预言，以免他们认为我们伪造了这些预言，并且从已经发生的事好像未来一样编造它们，我们就拿出犹太人的书。犹太人实在是我们的敌人，从敌人的书中我们说服对手……如果有敌人喧嚷说：\Quote{你们为自己伪造了预言；}就拿出犹太人的书，因为大的要服事小的。让他们在其中阅读那些我们如今看到应验的预言；并且让我们都说：\BibleQuote{上主以色列的天主是应当称颂的，从永远直到永远，全体人民要说：阿们！阿们！}\par

\SourceRecord{Translated by J.E. Tweed.；Nicene and Post-Nicene Fathers, First Series；Vol. 8.；Edited by Philip Schaff.；Buffalo, NY: Christian Literature Publishing Co.,；1888.；<http://www.newadvent.org/fathers/1801041.htm>.}

% structure: p0001 source=h1 target=section reason=page-title

\section{圣咏第四十二篇讲疏}

1. 我们承担了一篇与你们自己的\Quote{渴望}相应的圣咏的讲疏，打算向你们讲述。因为这圣咏本身以某种虔诚的\Quote{渴望}开始；这样歌唱的人说：\BibleQuote{天主，我的灵魂渴慕祢，有如牝鹿渴慕溪流。} \BibleRef{咏 41:1}。那么，说这话的是谁呢？正是我们自己，只要我们愿意！何必问是谁，除了你自己，因为成为你所问的事物就在你的能力之内？然而这不是一个人，而是\Quote{一个身体}；但\BibleQuote{基督的身体就是教会}。\BibleRef{哥 1:24}。这种\Quote{渴望}确实并非在所有进入教会的人身上都能找到；但愿所有尝过\Quote{主}的甘饴、并在基督内认明他们所喜爱之物的人，不要以为自己是唯一的；而是有这些种子散布在\Quote{田地}——主的整个大地——中；并且有一个基督的合一，其声音这样说：\BibleQuote{天主，我的灵魂渴慕祢，有如牝鹿渴慕溪流。} 而且，这也不难理解为那些尚未成为慕道者、正奔向圣洗圣事恩宠之人的呼声。为此，这篇圣咏通常在这种场合咏唱，使他们如同\Quote{牝鹿渴慕溪流}一样渴慕罪赦的泉源。让我们允许这一点；并且这个意义在教会中保持其位置；一个既真实又经习惯认可的位置。然而，我的弟兄们，在我看来，这种\Quote{渴望}即使在信徒的圣洗中也并未完全满足；而是如果她们知道自己寄居何处，以及要从这里迁往何处，她们的\Quote{渴望}也许会更加强烈地燃起。\par

2. 它的标题是：\Quote{论结局：科辣黑子孙的训诲诗。}我们在其他圣咏的标题中遇到过科辣黑的子孙；记得我们已经讨论并陈述过这个名称的意义。然而，我们现在仍须以这样的方式注意这个标题，使我们之前所说的不损害我们再次说它：因为并非所有人在我们说过的地方都在场。科辣黑可能是一个特定的人，也确实如此，并且他可能有儿子，可称为\Quote{科辣黑的子孙}；然而，让我们探寻这一圣事所隐藏的奥秘，使这名称揭示它所蕴含的玄义。因为基督徒被称为\Quote{科辣黑的子孙}，其中含有某种伟大的奥秘。为什么是\Quote{科辣黑的子孙}？他们是\Quote{新郎的子孙，基督的子孙}。那么为什么\Quote{科辣黑}代表基督？因为\Quote{科辣黑}等同于\Quote{加尔瓦略}。……因此，\Quote{新郎的子孙}，祂苦难的子孙，被祂宝血救赎的子孙，祂十字架的子孙，他们在额上带着祂仇敌在加尔瓦略所竖立之物，被称为\Quote{科辣黑的子孙}；为他们，这篇圣咏作为一篇\Quote{训诲}圣咏而咏唱。那么让我们的理解被唤醒吧：如果这圣咏向我们咏唱，就让我们以\Quote{理解}跟随它。……奔向溪流；渴慕溪流。\Quote{生命之泉与天主同在}；一个永不会干涸的\Quote{泉源}；在祂的\Quote{光}中，有永不会暗淡的光。愿你渴慕这光；一个你们肉体的眼睛所不知的泉源、一种光；为了看见这光，必须预备内在的眼；为了饮这泉源，必须燃起内在的渴。奔向泉源；渴慕泉源；但不要随随便便，不要像普通动物那样满足于奔跑；你要\Quote{像牡鹿一样}奔跑。\Quote{像牡鹿一样}是什么意思？愿你的奔跑没有懈怠；尽力奔跑；尽全力渴慕泉源。因为我们在\Quote{牡鹿}身上找到了敏捷的象征。\par

3. 但圣经也许要我们在鹿身上考虑的不仅这一点，还有另一点。听听鹿还有什么特性。它消灭蛇；消灭蛇之后，它被更强烈的渴所灼烧；消灭了蛇之后，它怀着比先前更锐利的渴奔向\Quote{溪流}。蛇是你们的恶习，消灭不义的蛇；然后你们会更加渴慕\Quote{真理之泉}。也许贪婪在你耳边低语某种黑暗的计谋，嘘声反对天主的话，反对天主的诫命。既然对你说：\Quote{轻视这个或那个事物}，如果你宁愿行不义而不愿轻视某些暂时的好处，你就选择被蛇咬，而不消灭它。因此，当你现在还纵容你的恶习、你的贪婪或你的食欲时，我何时能在你身上找到这种\Quote{渴望}，使你能奔向溪流？……\par

4. 在鹿身上还有另一个要点可以观察。据说牡鹿……当它们在群中漫游，或游泳到达大地的其他部分时，它们以这种方式互相支撑头部的负担：一只领头，其他跟随，将头靠在它身上；而跟随者后面的又将头靠在它们身上，如此一直到群末；但那只领头承受头部负担的，当疲倦时，退到后面，像其他的一样将头支撑着来休息；通过这样轮流支撑这负担，它们既完成了旅程，也不彼此抛弃。它们岂不就是宗徒所称呼的那种\Quote{鹿}吗？他说：\BibleQuote{你们要彼此担待重担，这样就满全了基督的法律。} \BibleRef{迦 6:2}……\par

5. \Quote{我的灵魂渴慕生活的天主} \BibleRef{咏 41:2}。我说：\Quote{正如牡鹿渴慕溪水，我的灵魂也同样渴慕祢，天主}，意思就是：\Quote{我的灵魂渴慕生活的天主}。它渴慕什么呢？\Quote{我何时能来朝见天主呢？}这就是我所渴慕的：\Quote{来朝见祂}。我在旅途中、奔跑中渴慕；抵达时我将得饱足。但\Quote{我何时能来？}这在天主眼中是很快的，但对我们的\Quote{渴慕}而言却是迟延的。\Quote{我何时能来朝见天主呢？}这也出自那\Quote{渴慕}，正如在另一处所呼喊的：\Quote{我向上主求得一事，我仍渴求：使我一生常住在上主的殿宇中。}为何如此？\Quote{使我得以仰望上主的美善。}\Quote{我何时能来朝见上主呢？}……\par

6. \Quote{他们天天向我说：\InnerQuote{你的天主在哪里？}而我的眼泪昼夜成为我的食物} \BibleRef{咏 41:3}。他说：我的眼泪不是苦涩的，而是\Quote{我的食物}。这些眼泪对我竟是甘甜的：我渴慕那活泉，既然尚不能饮，便急切地以眼泪为食物。因为他并没有说\Quote{我的眼泪成为我的饮料}，免得显得他渴慕眼泪如同渴慕\Quote{溪水}；而是仍存着那燃烧的渴慕，那催逼我奔向溪水的渴慕，他说：\Quote{我的眼泪成为我的食物}，因为我尚未到达那里。确实，他以眼泪为食物，反倒更加渴慕溪水。……\Quote{他们天天向我说：\InnerQuote{你的天主在哪里？}}\par

7. \Quote{我想起这些事，就在我以上倾注我的灵魂} \BibleRef{咏 41:4}。我的灵魂何时能达到那\Quote{在我灵魂以上}的追求对象呢？除非我的灵魂\Quote{在自我以上倾注自己}。因为若它停留在自身内，就看不到自身以外的任何事物；即使看见自己，也看不见天主。如今，让那侮辱我的仇敌说：\Quote{你的天主在哪里？}让他们说吧！我，只要尚未\Quote{看见}，只要我的幸福还在延迟，我就以眼泪作为我\Quote{昼夜的食物}。让他们仍然说：\Quote{你的天主在哪里？}我在一切有形体的自然物中，无论地上的还是天上的，寻找我的天主，却找不到祂；我在自己的灵魂中寻找祂的本体，也找不到；然而我仍想起这些事，并希望\Quote{借着所造之物，看见天主那看不见的事} \BibleRef{罗 1:20}，我便在我以上倾注我的灵魂，于是再没有任何存有可供我追求，除了我的天主。因为\Quote{那里}正是\Quote{我天主的居所}。祂的住所高于我的灵魂；从那里祂俯视我；从那里祂创造了我；从那里祂引导我、照顾我；从那里祂呼唤我、召叫我、指导我；在路上引领我，直到路途的终点。……\par

8. 因为当我\Quote{在我以上倾注我的灵魂}，为要到达我的天主时，我为何这样做？\Quote{因为我将进入祢帐幕的居所。}因为若我离开\Quote{祂帐幕的居所}去寻找我的天主，我便会迷失。\Quote{因为我将进入祢奇妙帐幕的居所，直抵天主的殿宇。}\par

\Quote{我要进入，}他说，\Quote{奇妙帐幕之所，直到天主的殿宇！}因为在\InnerQuote{帐幕}中，我已经有许多令我赞叹的事。看，我在帐幕中赞叹何等伟大的奇事！因为天主地上的帐幕就是信众；我在他们身上赞叹他们连肢体都顺服：在他们内\InnerQuote{罪恶不得为王，叫他们服从它的情欲；也不将他们的肢体献给罪恶作不义的武器，而是献给永生的天主行善。}\BibleRef{罗 6:12} 我赞叹在事奉天主的灵魂的服役中，肢体争战的景象……虽然帐幕奇妙，但当我来到\InnerQuote{天主的殿宇}时，我更是惊讶得无言。他在另一篇圣咏中提到那\InnerQuote{殿宇}，当时他对自己提出一个深奥而难解的问题（即：为何世上恶人常享顺遂，义人常遭不幸？），说：\Quote{我本想明了这事，但对我实在太难，直到我进入天主的圣所，明了他们的结局。}因为正是在那里，在天主的圣所，在天主的殿宇中，才有\InnerQuote{领悟}的泉源。在那里他\InnerQuote{明了了他们的结局}，并解决了恶人享福、义人受苦的问题。他如何解决？原来，恶人在这里暂缓受罚，是留给永无止境的刑罚；而善人在这里受苦，是为受考验，以在最后获得产业。正是在天主的圣所中，他明白了这一点，\InnerQuote{明了了他们的结局}……因为他告诉我们他的进展，以及他在那里的引导；仿佛我们问他：\Quote{你在地上赞赏帐幕，怎能到达天主的圣所呢？}他说：\Quote{在欢呼和赞美的声音中，在庆祝节期的音响中。}这里，当人们仅仅为自己的放纵而庆祝时，习惯在房屋前放置乐器，或安排一队歌手，以及任何助长荡逸的音乐。我们路过的人听到这些，就会问：\Quote{这里在做什么？}有人回答说是某个节庆。他们说：\Quote{正在庆祝生日，这里有婚宴。}这样，那些歌声才不显得不合时宜，而放纵的享乐也可因节庆而获得谅解。在\InnerQuote{天主的殿宇}中，有一个永不止息的节庆：因为那里不是一次庆祝、然后就过去的场合。天使的唱诗班举行永恒的\InnerQuote{节日}：天主面容的临在，永不枯竭的喜乐。这是一个既不由曙光开启、也不由黄昏结束的\InnerQuote{节日}。从那永恒不断的节庆中，有一种甜美悦耳的旋律触及心灵之耳，只要世界不淹没那声音。当他行在这个帐幕中，默观天主为救赎信众而行的奇妙工程时，那节庆的声音吸引他的耳朵，将\InnerQuote{牡鹿}带往\InnerQuote{水泉}。\par

9. 但是，弟兄们，只要我们\InnerQuote{住在这肉身内，就是与主远离}\BibleRef{格后 5:6}；而且\InnerQuote{可朽坏的肉身重压灵魂，这地上的帐幕，使那多虑的心神昏沉}\BibleRef{智 9:15}；即使我们以某种方式驱散云雾，随着\InnerQuote{渴望}引领前行，并短暂地触及那声音，以便努力从那\InnerQuote{天主的殿宇}中捕捉到一些东西，然而，可以说，因我们软弱的负担，我们又沉回平常的水平，返回到通常的状态。正如在那里我们有喜乐的原因，在这里也不乏悲伤的时机。因为那只以\InnerQuote{眼泪}为\InnerQuote{昼夜饮食}的牡鹿，被\InnerQuote{渴望}带往\InnerQuote{水泉}（即天主的属灵福乐），\InnerQuote{将自己的灵魂倾注于己之上}，以得到那\InnerQuote{超越}自己灵魂的，向着\InnerQuote{奇妙帐幕之所，天主的殿宇}前进，并被那内在属灵声音的甜蜜所引导，而轻视一切外在事物，被带向属灵的事物，但他终究仍是凡人；仍然在此叹息，仍然带着肉身的脆弱，仍然在这世界的\InnerQuote{绊脚石}\BibleRef{玛 18:7} 中处于危险。因此，他仿佛从那个世界回来，回顾自己；对自己说话，如今置身于这些悲伤之中，将这些与他进去所见的事物相比较，及看见后从那里出来所面对的事物；\par

\Quote{我的灵魂，你为何沮丧？为何在我内烦躁？}\BibleRef{咏 41:5} 看，我们刚才因某些内在的喜乐而喜乐；用心灵之眼，虽只是片刻一瞥，我们得以看见某种不变的事物：你为什么还要\InnerQuote{使我烦躁}，为什么还要\InnerQuote{沮丧}？因为你并不怀疑你的天主。现在，对于那些说\InnerQuote{你的天主在那里？}的人，你已不是没有话可说。我已经体验到某种不变的事物；你为什么还要使我烦躁？\par

\Quote{期望于天主。}仿佛他的灵魂默默回答他：\Quote{我为什么使你烦躁呢？不是因为我还没有到达那喜乐之处，就是那片刻被提去的地方？难道我已经喝着这\InnerQuote{泉水}而无后顾之忧吗？}……他回答那使他烦躁、并想以这世界充斥的邪恶来解释其烦躁的灵魂说：\Quote{期望于天主。}同时，你要生活在希望中：因为\InnerQuote{所见到的希望，不是希望；我们所希望若未见到的，我们就以坚忍等候}。\BibleRef{罗 8:24}\par

10. \\Quote{望天主}。为什么\\Quote{望德}？\\Quote{因为我将向祂宣认}。你要\\Quote{宣认}什么？\\Quote{我的天主是我面容的救恩}。我的\\Quote{救恩}（我的救恩）不能来自我自己；这就是我要说的，我要\\Quote{宣认}的。是我的天主是\\Quote{我面容的救恩}。因为他解释自己的恐惧，在那些他现在知道的事物中，在某种程度上已经获得了对它们的\\Quote{理解}之后，他又焦虑地回头看，恐怕敌人偷袭他：他还不能说\\Quote{我已经完全痊愈了}。因为我们只获得了\\Quote{圣神的最初果实}，我们内心叹息；等待义子名分，也就是身体的救赎。\\BibleRef{罗 8:23} 当那\\Quote{救恩}（那救恩）在我们内完美时，那时我们将永久居住在天主的房屋里，永久赞美祂，曾有人对祂说：\\Quote{住祢房屋的人是有福的，他们将永无止境地赞美祢。}现在还不是这样，因为所应许的救恩尚未实现；但我是\\Quote{在望德中}向天主宣认，说：\\Quote{我的天主是我面容的救恩。}因为我们是\\Quote{在望德中}\\Quote{得救的；但看得见的望德，就不是望德了。}...\par

11. \\BibleQuote{我的灵魂因我自己而烦乱} \\BibleRef{咏 41:6}。是因天主而烦乱吗？是因我自己而烦乱。我知道天主的义德保持不变；我自己的意志是否坚定不移，我不知道。因为我因宗徒的话而警醒：\\BibleQuote{那自以为站得稳的，要小心免得跌倒。} \\BibleRef{格前 10:12}。因此既然\\Quote{在我内没有健全，为了我自己，} 我对自己也毫无希望。\\Quote{我的灵魂因我自己而烦乱。}...\\Quote{因此，主，我记起了祢，从约旦\\OriginalTerm{约旦}{Jordan}地，从赫尔孟\\OriginalTerm{赫尔孟}{Hermon}小山。} 我从哪里记起了祢？从\\Quote{小山}，从\\Quote{约旦地}。也许从圣洗圣事中，那里赐予罪的赦免。因为没有人奔向罪的赦免，除非他对自己不满意；没有人奔向罪的赦免，除非他宣认自己是罪人；没有人宣认自己是罪人，除非在天主前谦卑自己。因此是从\\Quote{约旦地我记起了祢，从小山}；注意，不是\\Quote{大山}，好使你使\\Quote{小山}成为大山：因为\\Quote{凡高举自己的，将被贬抑；凡谦卑自己的，将受高举。} 你若要追问这些名称的含义，约旦意为\\Quote{他们的下降}。那么下降吧，使你可以\\Quote{被举升}；不要自高，免得你被摔下。\\Quote{和赫尔孟小山}。赫尔孟意为\\OriginalTerm{诅咒}{anathematizing}。诅咒你自己，因对自己不满；因为如果你对自己满意，天主将对你不满。因为那时天主赐给我们一切美物，是因为祂自己是美善的，不是因为我们配得；是因为祂是仁慈的，不是因我们在任何事上应得；是从\\Quote{约旦地，并从赫尔孟}，我记起了祢。并且因为他如此谦卑地记起，他将赢得他的举升以得果实，因为他不是\\Quote{自高}，那\\Quote{在主内夸耀}的人。\par

12. \BibleQuote{深渊以你瀑布之声呼唤深渊} \BibleRef{咏 41:7}。我或许能藉着你们的注意力完成这篇圣咏，我察觉到你们的热情。至于你们听讲的疲倦，我并不十分担心，因为你们看见我，这位讲话者，也在这番努力的炎热中劳作。无疑，是因你们看见我劳苦，你们才与我一同劳苦：因为我劳作不是为了自己，而是为了你们。\BibleQuote{深渊以你瀑布之声呼唤深渊}。他所呼求的是天主，是那\Quote{从约旦和赫尔孟之地记起他来}的天主。他怀着惊奇和赞叹说出这话：\BibleQuote{深渊以你瀑布之声呼唤深渊}。这呼唤的深渊是什么？它所呼唤的另一个深渊又是什么？确实，因为所说的\Quote{理解}是一个\Quote{深渊}。因为一个\Quote{深渊}是一种无法触及或领会的深处，主要应用于浩大的水体。因为有一种\Quote{深}，一种\Quote{深邃}，其底部无法用测锤触及。此外，在某处经文也说：\Quote{祢的判断是伟大的深渊}，圣经意思是暗示天主的判断是无法理解的。那么，这\Quote{深渊}是什么，它呼唤的另一个\Quote{深渊}又是什么？如果我们把\Quote{深渊}理解为极大的深度，难道人的心，你们不认为，不就是一个\Quote{深渊}吗？因为有什么比那个\Quote{深渊}更深远呢？人能说话，可从其肢体的动作被看见，可在交谈中被听见说话：但谁的意念能被看透？谁的心能被洞悉？他内在从事什么，内在能做什么，内在在做什么或打算什么，内在希望发生或不发生什么，谁能理解？我想，不无理由地把\Quote{深渊}理解为一个人，关于这人别处说：\Quote{人将到达深处的心，天主将被高举}。如果人是一个\Quote{深渊}，那么\Quote{深渊}如何呼唤\Quote{深渊}呢？难道人如天主那样被人呼求吗？不，但\Quote{呼求}等于\Quote{召唤他}。因为关于某个人说过：他呼求死亡；\BibleRef{智 1:16}也就是说，他生活的方式是邀请死亡；因为根本没有一个人会祈祷并明确呼求死亡：而是人因邪恶生活而邀请死亡。\Quote{深渊呼唤深渊}，因此是\Quote{人呼唤人}。因而智慧被学习，信德也是这样，当\Quote{人呼唤人}时。神圣的天主圣言的宣讲者呼唤\Quote{深渊}，他们自己不也是\Quote{深渊}吗？……\par

13. \BibleQuote{深渊以祢瀑布之声呼唤深渊}。我，全身颤抖，当我的灵魂因我自身而不安时，因祢的\Quote{判断}而大大恐惧……那些判断是轻微的吗？它们是伟大的、严厉的、难以承受的；但愿它们就是一切。\BibleQuote{深渊以祢瀑布之声呼唤深渊}，因为祢在警告，祢说，甚至在那些苦难之后还有另一个定罪积蓄。\BibleQuote{深渊以祢瀑布之声呼唤深渊}。\BibleQuote{我往哪里去，才能躲避祢的面容？我往哪里逃，才能躲避祢的神？} \BibleRef{咏 139:7}，因为看见深渊呼唤深渊，并且那些苦难之后更严厉的恐怖要来临。\par

14. \BibleQuote{祢所有的洪涛和巨浪都涌向我} \BibleRef{咏 42:7}。\Quote{巨浪}是我已经感受到的，\Quote{洪涛}是祢所宣告的。我所有的痛苦都是祢的巨浪；祢所有审判的宣告都是祢的\Quote{洪涛}。在那\Quote{巨浪}中，那深渊\Quote{呼唤}；在\Quote{洪涛}中，是它\Quote{呼唤}的另一个\Quote{深渊}。在我所受的这些痛苦中，有祢所有的巨浪；在祢所警告的更严厉的惩罚中，祢所有的\Quote{洪涛}都已临到我。因为那威胁者并不让祂的判断落在我们身上，而是将它们悬在我们头上。但因祢安然坐着，我对我的灵魂说：\BibleQuote{希望于天主，因为我将向祂宣认；我的天主是我面容的救援。} \BibleRef{咏 42:11}。我的痛苦越多，祢的仁慈就越甜美。\par

15. 因此接着是：\BibleQuote{主白昼必赐其慈爱，黑夜必歌颂祂}\BibleRef{咏 41:8}。在磨难中，无人有闲暇去听：要留心，当顺遂时；要听，当顺遂时；要学习，当平安时，学习智慧的纪律，像储存食物一样储存天主的话。因为在磨难中，每个人都必须从他在平安时听到的话中获得益处。因为在天主所赐的顺遂中，主\Quote{向你彰显祂的仁慈}，免得你忠心事奉祂，因祂使你脱离磨难；但只有在\Quote{黑夜}祂才\Quote{彰显}祂的仁慈，就是祂在白昼向你\Quote{彰显}的那份。当磨难实际来临时，祂不会让你失去祂的帮助；祂会向你证明，祂在白昼所彰显的是真实的。因为某处经上写着：\Quote{主的仁慈在患难时是合时的，如同旱天的雨云。}\Quote{主白昼必赐其慈爱，黑夜必歌颂祂。}除非磨难来临，祂不会显出祂是帮助者，因为你要从磨难中被祂所救，祂曾在\Quote{白昼}应许你。因此我们被警告要像\Quote{蚂蚁}。因为正如世俗的顺遂由\Quote{白昼}象征，逆境由黑夜象征；同样，另一种意义上，世俗的顺遂由\Quote{夏季}表达，逆境由\Quote{冬季}表达。蚂蚁做什么呢？它在夏季储存冬季有用的东西。因此，当还是夏季时，当顺遂时，当平安时，要听主的话。因为在这世界的风暴中，你怎能平安渡过整个大海而不受苦？怎么可能呢？有哪一位凡人曾幸免？即使有人曾幸免，那平静本身更令人畏惧。\Quote{主白昼必赐其慈爱，黑夜必歌颂祂。}……\Quote{我有祈祷，向我的生命的天主。}这是我在此的事业；我，这只\Quote{鹿，渴慕溪水}，追忆那甜美曲调，我经由会幕甚至被引到天主的殿宇；当这\BibleQuote{可朽坏的身体压迫灵魂}\BibleRef{智 9:15}时，我仍有\Quote{向我的生命的天主的祈祷}。因为为向上主祈祷，我不必从海外买什么；或为使祂俯听我，不必航海从远方带来乳香和香料，或带来\Quote{牛犊或公羊}。我有\Quote{向我的生命的天主的祈祷}。我内在有祭品可献；我内在有乳香可放于祭台上；我内在有祭献可平息我的天主。\Quote{天主所要的祭献是痛悔的精神。}我所拥有的\Quote{痛悔精神}的祭献，请听。\par

16. \BibleQuote{我要向天主说：祢是我的高举者。祢为何忘记了我？}\BibleRef{咏 41:9}。因为我在此受苦，仿佛祢忘记了我。但祢是在考验我，我知道祢只是推迟，并非全然收回祢所应许的。但，\Quote{祢为何忘记了我？}我们的元首也曾这样呼喊，仿佛以我们的名义说话：\Quote{我的天主，我的天主，祢为何舍弃了我？}我要向天主说：\Quote{祢是我的高举者；祢为何忘记了我？}\par

17. \Quote{祢为何弃绝了我？}\Quote{弃绝}我，也就是说，从理解不变真理的高度被抛弃。\Quote{祢为何弃绝了我？}为什么，当我已渴慕那些事物时，却因我罪恶的沉重负担而被抛下这些事物？这同一声音在另一处说：\Quote{我在恍惚中说}（即，在我心神超拔时，当他看见某种伟大事物时），\Quote{我在恍惚中说：我从祢眼前被抛弃。}因为他将自身所处的这些事物与所被提升至的那事物相比较，看见自己被远远抛离\Quote{天主眼前}，正如他在这里说：\Quote{祢为何弃绝了我？为何我悲哀地行走，仇敌折磨我，打断我的骨？}就是那仇敌，那诱惑者，魔鬼；当恶行随处增多，因\BibleQuote{许多人的爱又冷淡了}\BibleRef{玛 24:12}。当我们看到教会中坚强的肢体通常屈服于绊脚石的原因时，基督的身体岂不说：\Quote{仇敌打断我的骨}？因为坚强肢体就是\Quote{骨}；有时甚至坚强的人也屈服于诱惑。因为基督身体的任何成员想到这一点，难道不会以基督身体的声音呼喊：\Quote{祢为何弃绝了我？为何我悲哀地行走，仇敌折磨我，打断我的骨？}\par

你不仅看到我的血肉，甚至我的\Quote{骨}。看到那些被认为有些坚固的人，在诱惑下屈服，以致其余软弱的弟兄们看到强人跌倒时便绝望；那危险是多么大啊，我的弟兄们！\par

18. \BibleQuote{折磨我的人当面辱骂我。他们终日对我说：你的天主在哪里？}\BibleRef{咏 41:10}。尤其在教会的诱惑中，他们说：\Quote{你的天主在哪里？}这辱骂曾多少次投向殉道者！那些为基督之名而忍耐和勇敢的人，他们常被说：\Quote{你的天主在哪里？}\Quote{若祂能，让祂救你。}因为人们外在看见他们的苦难，却未内在地看见他们的冠冕！\Quote{折磨我的人当面辱骂我，他们终日对我说：你的天主在哪里？}正因如此，看到\Quote{我的灵魂因我自己而困扰}，我除了说这些话还能对它说什么呢？\par

\BibleQuote{我的灵魂，你为何忧闷？为何在我内烦躁？} \BibleRef{咏 41:11}。就如它似乎回答说：\Quote{你置身于如此巨大的邪恶中，难道不让我烦扰你吗？你向往美善、渴望并努力追求它，难道不让我烦扰你吗？}我该说什么呢，除非——\par

\Quote{你要盼望天主，因为我还要向祂称谢} \BibleRef{咏 41:11}。他说出了那称谢的话；他重复了坚固他盼望的根据。\Quote{祂是我面容的救援，是我的天主。}\par

\SourceRecord{Translated by J.E. Tweed.；Nicene and Post-Nicene Fathers, First Series；Vol. 8.；Edited by Philip Schaff.；Buffalo, NY: Christian Literature Publishing Co.,；1888.；<http://www.newadvent.org/fathers/1801042.htm>.}

% structure: p0001 source=h1 target=section reason=page-title

\section{圣咏第四十三篇释义}

1. 这是一首简短的圣咏；它满足听者心灵的需求，而不给守斋者的饥饿带来过重的考验。让我们的灵魂以它为食粮；我们的灵魂，就是那位在这首圣咏中歌唱者所说\Quote{沮丧}的灵魂；我想，这种沮丧或是由于某种拳头打击，或更可能是由于他所处的某种饥饿。因为守斋是自愿的行为；饥饿则是非自愿的。那处于饥饿状态的，是教会，是基督的身体；而那位延伸至全世界的\Quote{人}，其首在上，肢体在下：正是他的声音，在所有圣咏中，如今我们应当已完全熟悉并习以为常了；时而欢唱，时而忧伤；时而因希望而喜乐，时而因现实而叹息，仿佛就是我们自己的声音。因此，我们不必花时间向你们指出谁在此发言：愿我们每人成为基督身体的肢体；那么他就在此发言……\par

2. \BibleQuote{上主，求你审断我，并从不虔敬的国民中，为我的案件辩护}\BibleRef{咏 42:1}。我不畏惧你的审判，因我知道你的仁慈。他呼号说：\Quote{天主，求你审断我}。如今，在这朝圣的旅途中，你尚未将我的场所分开，因为我必须与\Quote{莠子}一起生活直到\Quote{收割}之时；你尚未将我的雨水与他们的分开；我的光明与他们的分开：\Quote{为我的案件辩护}。愿信你之人与不信你之人有所区别。我们的软弱相同，但我们的良心不同；我们的苦难相同，但我们的渴望不同。\Quote{不虔敬者的渴望必将消逝}，至于义人的渴望，我们本该怀疑，若那应许者不是\Quote{确实}的。我们所渴望的对象就是他本人，那应许者：他将把自己赐给我们，因为他已将自己赐给了我们；他将把自己不朽的赐给那时不朽的我们，正如他把自己可朽的赐给了可朽的我们……\par

3. 既然需要忍耐，以忍受直到收割之时的某种区别却不分离（如果我们可以这样说）（因为他们与我们同在，因此尚未分离；然而莠子仍是莠子，麦子仍是麦子，因此已有了区别）；既然需要某种力量，必须向那命令我们坚强、且没有他的使我们坚强我们就不能成为他所命令的那一位祈求；他对那说\BibleQuote{惟有忍耐到底的，必然得救}\BibleRef{玛 24:31}的那一位说，以免灵魂的力量因她将任何力量归于自己而受损，他便立刻补充说：\par

\BibleQuote{因为你是我的天主，我的力量。你为什么舍弃我？为什么我因仇敌的压迫而悲哀而行？}\BibleRef{咏 42:2}。我悲哀而行：仇敌以日常的诱惑压迫我；激起某种不当的爱恋，或某种无根据的恐惧；那与两者搏斗的灵魂，虽未被它们俘虏，但因处于危险中而紧缩于忧愁中，向天主说：\Quote{为什么？}\par

让她向他询问，并听见\Quote{为什么？}因为她在圣咏中追问自己沮丧的原因，说：\Quote{你为什么舍弃我？为什么我悲哀而行？}让她从依撒意亚那里聆听；让她想起刚才宣读的读经。\BibleQuote{神明从我这里出去，我造了每一口气。因着罪愆，我稍稍打击了他；我向他掩面，他就怀着忧伤，顺着自己内心的道路离开了我。}\BibleRef{依 57:16}。那么你为什么问：\Quote{你为什么舍弃我，为什么我悲哀而行？}你已经听见，是因着\Quote{罪愆}。\Quote{罪愆}是你悲哀的原因；让\Quote{正义}成为你喜乐的原因！你愿意犯罪，却又不愿受苦；这样，你自己不义还不够，还希望他也成为不义，因为你宁可不受他的惩罚。想想另一首圣咏中更好的话。\Quote{你谦卑我，对我有益，为使我学习你的律例。}由于被高举，我学会了我的不义；让我因着被\Quote{谦卑}，学习\Quote{你的律例}。\Quote{为什么我因仇敌的压迫而悲哀而行？}你抱怨仇敌。他确实压迫你，但却是你\BibleQuote{给魔鬼留了地步}\BibleRef{弗 4:27}。即使现在，你仍有出路；选择明智之路，迎接你的君王，将暴君关在门外。\par

4. 但为叫她这样做，请听她说什么，她恳求什么，她祈祷什么。你当为你所听的祈祷；当你听的时候要祈祷；让这些话成为我们众人的声音：\Quote{求祢发出祢的光明和真理。它们引导了我，带领我到了祢的圣山，进入祢的帐幕。}\BibleRef{咏 42:3}。因为那\Quote{光明}和\Quote{真理}在名称上确实是两个；所表达的实体却是一个。因为天主的\Quote{光明}若非天主的\Quote{真理}，还能是什么呢？或者天主的\Quote{真理}若非天主的\Quote{光明}，还能是什么呢？而基督的同一个位格就是这两者。\BibleQuote{我是世界的光；信我的，必不在黑暗中行走。}\BibleQuote{我是道路、真理、生命。}祂本身就是\Quote{光明}；祂本身就是\Quote{真理}。那么让祂来拯救我们，并\Quote{立刻将我们的案由从不敬神的民族中分别出来；让祂救我们脱离诡诈和不义的人，}让祂将麦子从莠子中分别出来，因为在收获时，祂将亲自派遣祂的天使，叫他们\Quote{从祂的国度里收集一切绊倒人的事，}\BibleRef{玛 13:41}并将它们扔进烈火中，同时将谷物收进仓里。祂将发出祂的\Quote{光明}和祂的\Quote{真理}；因为那些已经\Quote{带领我们到了祂的圣山，进入祂的帐幕。}我们拥有\Quote{抵押}；我们盼望奖赏。\Quote{祂的圣山}就是祂的圣教会。那座山，按达尼尔的神视，\BibleRef{达 2:35}从一颗很小的\Quote{石头}长起，直到粉碎了地上的列国；并且长到如此之大，以至\Quote{充满了地面。}这就是那座\Quote{山}，那说\BibleQuote{我以我的声音呼求上主，祂从祂的圣山应允了我}的人告诉我们他的祈祷是从那里被垂听的。愿那些在那山之外的人，不要希望被垂听而得永生。因为许多人在他们的祈祷中被垂听，是为许多事物。愿他们不因被垂听而自庆；魔鬼在它们的祈祷中被垂听，是为被遣入猪群。愿我们渴望因我们的渴望而被垂听得永生，藉此我们说：\Quote{求祢发出祢的光明和真理。}\BibleRef{玛 8:31}。那是需要心灵的眼睛的\Quote{光明}。因为祂说：\BibleQuote{心里纯洁的人是有福的，因为他们要看见天主。}\BibleRef{玛 5:8}。我们现在就在祂的圣山上，即在祂的教会里，在祂的帐幕内。\Quote{帐幕}是为旅居的人；房屋是为同住一团体的人。帐幕也是为那些离乡背井且在战斗状态中的人。当你听到帐幕时，要形成一个战争的概念；要提防敌人。但房屋是什么呢？\BibleQuote{住在祢房屋的人是有福的：他们要永远赞美祢。}\par

5. 那么既然我们已经甚至被引导到\Quote{帐幕}，并被安置在\Quote{祂的圣山}上，我们带着什么希望呢？\par

\Quote{那时我要走到天主的祭台前}\BibleRef{咏 42:4}。因为有一个看不见的高天祭台，不义的人不能接近。只有那毫无畏惧地接近这个祭台的人，才能接近那个祭台。他在此世\Quote{分离他的案由}的人，将在那里找到他的生命。\Quote{我要走到天主的祭台前。}从祂的圣山，从祂的帐幕，从祂的圣教会，我要走到高天的天主祭台前。那里有什么样的祭献呢？进去的人自己就被当作全燔祭。\Quote{我要走到天主的祭台前。}他说\Quote{我的天主的祭台}是什么意思？\par

\Quote{走向使我青春欢乐的天主。}青春意味着新：就好像他说：\Quote{走向使我新的欢乐的天主。}是祂使我的新欢乐，祂曾使我的旧状态充满哀伤。因为如今我在旧中\Quote{哀伤行走}，那时我将在新中\Quote{站立}，欢跃！\par

\Quote{是的，我要在竖琴上赞美祢，我的天主。}什么是\Quote{在竖琴上赞美}和\Quote{在琴瑟上赞美}的意思？因为他并不总是使用竖琴，也不总是使用琴瑟。音乐家的这两种乐器各有其独特的含义，值得我们思考和注意。它们都拿在手里，靠手指弹奏发声；它们象征我们某些身体的工作。两者都是好的，如果人懂得弹琴瑟，或弹竖琴。但因为琴瑟那种乐器，其共鸣箱（即那鼓，那中空的木头，琴弦张在其上而发声）位于上部，而竖琴的同一凹形共鸣板位于下部，所以我们的工作应当有所区分：何时是\Quote{在竖琴上}，何时是\Quote{在琴瑟上}；但两者都是天主所悦纳的，在祂耳中中听的。当我们按照天主的诫命行事，服从祂的命令并听从祂，以便完成祂的训示，当我们主动而非被动时，那是琴瑟在弹奏。因为天使也是如此：他们没有任何受苦。但当我们在此世上遭受任何磨难、试探、冒犯（正如我们只从我们自身较低的部分受苦；即由于我们是有死的，我们因原由而欠下一些磨难，也由于我们从那些不\Quote{在上}的人那里遭受许多痛苦）；这就是\Quote{竖琴}。因为从我们\Quote{在下}的部分升起甜美的旋律：我们\Quote{受苦}，我们弹奏琴瑟，或不如说我歌唱并弹奏竖琴……\par

6. 并且，为了从下方的共鸣板中汲取声音，他向自己的灵魂说话：他说，\Quote{我的灵魂，你为何忧伤？为何扰乱我？} \BibleRef{咏 42:5}。我身处困苦、疲惫、哀恸之中，\Quote{我的灵魂，你为何扰乱我？}说话者是谁？他是在向谁说话？他向灵魂说话，这是众所周知的：很明显，这话是直接向灵魂说的：\Quote{我的灵魂，你为何忧伤？为何扰乱我？}问题在于说话者。肯定不是肉身向灵魂说话，因为没有灵魂，肉身不能说话。因为灵魂向肉身说话，比肉身向灵魂说话更为合适……于是我们发现我们拥有某一部分，其中有\Quote{天主的肖像}，即理智和理性。正是这同一个理智曾祈求\Quote{天主的光}和\Quote{天主的真理}。正是这同一个理智，我们通过它辨别对错；我们通过它分辨真理与谬误。这就是我们称之为\Quote{理智}的；实际上，野兽缺乏这\Quote{理智}。凡是在自己内忽略这\Quote{理智}，将其看得比本性中其他部分更轻，并将其抛弃，就好像他根本没有一样，在圣咏中就对他说：\Quote{不要像马和骡子，它们没有理智。}因此，是我们的\Quote{理智}在向我们的灵魂说话。灵魂因困苦而枯萎，因痛苦而疲惫，在诱惑中\Quote{忧伤}，在劳苦中昏厥。理智瞥见上方的真理，渴望振奋精神，便说：\Quote{我的灵魂，你为何忧伤？}……\par

7. 弟兄们，这些说法是稳妥的：但你们也要在善工上警醒。借着遵守诫命去弹奏\Quote{瑟}；借着耐心忍受苦难去弹奏竖琴。你们从依撒意亚先知那里听到：\Quote{把你的饼分给饥饿的人} \BibleRef{依 58:7}，不要以为单靠禁食就足够了。禁食是克制你自己，并不能使他人得到饱饫。如果你能安慰别人，你自己的痛苦才会对你有益。看，你克苦了自己：你把从自己那里剥夺的东西给谁呢？你拒绝自己的东西，要施予何处？我们今天放弃的早餐可以填饱多少穷人？要这样禁食，使你能因别人吃了而欢喜，好像你吃了早餐一样；要因你的祈祷而禁食，使你的祈祷得以垂听。因为祂在那段经文中说：\Quote{你还在说话时，我就要说：我在这里}，只要你乐意以愉快的心情\Quote{把你的饼分给饥饿的人}。因为通常人们这样做是出于勉强和怨言，为了摆脱乞丐烦人的纠缠，而不是为了救济贫穷之人的饥肠。但\Quote{天主喜爱乐意捐献的人} \BibleRef{格后 9:7}。如果你勉强给饼，你就既失去了饼，也失去了善行的功劳。所以要发自内心去做：好使那\Quote{在暗中看见的} \BibleRef{玛 6:6}可以说：\Quote{你还在说话时，我在这里。}行义之人的祈祷何等迅速地蒙垂听！这就是人在今生的义：禁食、施舍和祈祷。你愿意你的祈祷飞升到天主面前吗？就给它加上施舍和禁食这两个翅膀。愿天主的光和天主的真理这样找到我们，使祂找到我们时我们无所畏惧，当祂来解救我们脱离死亡时——祂已经为我们而来接受死亡。阿们。\par

\SourceRecord{Translated by J.E. Tweed.；Nicene and Post-Nicene Fathers, First Series；Vol. 8.；Edited by Philip Schaff.；Buffalo, NY: Christian Literature Publishing Co.,；1888.；<http://www.newadvent.org/fathers/1801043.htm>.}

% structure: p0001 source=h1 target=section reason=page-title

\section{\WorkTitle{圣咏第四十四篇释义}}

1. 这篇圣咏是写给\Quote{科辣黑子孙}的，正如其标题所示。科辣黑相当于\Quote{秃头}一词；我们在福音中读到，我们的主耶稣基督被钉在\Quote{髑髅地}。显然，这篇圣咏是为祂\Quote{苦难}的子孙而唱的。关于这一点，我们从宗徒保禄那里得到了最确实、最明显的证据；因为当教会遭受外邦人迫害时，他曾引用本节作为安慰和鼓励忍耐的话。因为他在书信中所引用的，正是这里所说的：\BibleQuote{为了祢，我们整日被置于死地；我们被视作待宰的羊。}\BibleRef{罗 8:36} 那么，让我们在这篇圣咏中倾听殉道者的声音；并看看殉道者的声音所辩护的事业是何等美好，他们说：为了祢，等等……\par

2. 因此，标题不仅仅是\Quote{致科辣黑子孙}，而是\Quote{致科辣黑子孙：\Emph{有悟性的}}。同样，主自己在十字架上所喊出的那句圣咏的首节也是如此：\BibleQuote{我的天主，我的天主，祢为什么舍弃了我？}因为\BibleQuote{祂将我们喻作祂所说的话}\BibleRef{格前 4:6} 以及祂的身体（因为我们也是\Quote{祂的身体}，祂是我们的\Quote{头}），祂在十字架上所喊出的不是祂自己的呼声，而是我们的呼声。因为天主从未\Quote{舍弃}祂；祂自己也从未离开过圣父；但祂是为我们说了这话：\Quote{我的天主，我的天主，祢为什么舍弃了我？}因为下文接着说：\BibleQuote{我的罪恶的言语远离我的救恩}；这显明了祂是以谁的身份说这话；因为罪在祂内是找不着的……\par

3. \BibleQuote{天主啊，我们亲耳听见了，我们的祖先也给我们讲述了，祢在他们之日、在古时所作的事}\BibleRef{咏 43:1}。他们惊讶为何在这些日子里，祂似乎舍弃了那些祂愿意在苦难中磨练的人，他们回忆从祖先那里听到的往事；好像是在说：我们所受的这些苦，并不是祖先告诉我们的！因为在那另一篇圣咏中，祂也这样说：\BibleQuote{我们的祖先信赖祢；他们信赖了，祢就拯救了他们。但我是一条虫，不是一个人；是人的耻辱，是百姓的弃物。}他们信赖了，祢就拯救了他们；那么我寄望了，祢却\Quote{舍弃}了我？我白白地信仰了祢？我的名字被写在祢的册子上，祢的名字被铭刻在我身上，难道是徒然的吗？祖先告诉我们的是这样：\par

\BibleQuote{祢亲手摧毁了列邦，栽培了他们；祢打击了万民，赶走了他们}\BibleRef{咏 43:2}。也就是说：\Quote{祢赶走了\InnerQuote{万民}，使他们离开自己的土地，为叫祢能\InnerQuote{使他们}进入并栽培他们；并因祢的仁慈建立他们的王国。}这些是我们从祖先听来的。但也许是因为他们勇敢、善战、无敌、纪律严明、好斗，所以才能做成这些事。绝非如此。这不是祖先告诉我们的；圣经中也不是这样记载的。但圣经说的是什么呢？岂不是接下来所说的吗？\par

\BibleQuote{因为他们获得这土地，并非靠自己的刀剑，也不是靠自己的臂膀拯救了他们；而是靠祢的右手、祢的臂膀和祢慈颜的光照}\BibleRef{咏 43:3}。祢的\Quote{右手}是祢的能力；祢的\Quote{臂膀}是祢的圣子自己。以及\Quote{祢慈颜的光照}。这是什么意思呢？岂不是祢通过这样的奇迹与他们同在，使人能感知祢的同在？因为当天主通过任何奇迹与我们同在后，我们是否用肉眼看见祂的面容？不。是借奇迹的效果向人暗示祂的同在。事实上，所有对这类事迹感到惊奇的人都会说什么？\Quote{我看见天主临在。}\Quote{而是靠祢的右手、祢的臂膀和祢慈颜的光照；因为祢喜欢他们}：即这样对待他们，使祢在他们中喜悦；使任何考虑他们如何被对待的人，可以说：\Quote{天主实在与他们同在}；是天主推动他们。\par

4. \Quote{什么？难道祂现在不同于往昔了吗？}摒弃这种想法。因为接下来是什么呢？\par

\BibleQuote{你自身就是我的君王，我的天主。}\BibleRef{咏 43:4}. \Quote{你自身；}因为你是不改变的。我看见时代改变了，但时代的创造者是不改变的。\BibleQuote{你自身就是我的君王，我的天主。}你惯于引导我、治理我、拯救我。\BibleQuote{是你命令拯救雅各伯。}什么是\Quote{是你命令}？虽然在你自己的本体和本性中，无论你是什么，你对他们隐藏；虽然没有以你自身的所是与他们交流，使他们能\BibleQuote{面对面}看見你，但藉着任何受造物，\BibleQuote{是你命令拯救以色列。}因为那\BibleQuote{面对面}看见你，是保留给那些在复活中获得自由的人的。新约的\Quote{祖先们}也是如此，虽然他们看见你的奥秘启示，虽然他们宣讲那些启示给他们的隐密事，但他们说他们自己只是\BibleQuote{在镜中，模糊不清}地看见，但\BibleQuote{面对面看见}\BibleRef{格前 13:12}是保留给将来的时刻，那时宗徒自己所讲的将要来临。\BibleQuote{当基督，我们的生命，显现时，你们也要与祂一同在光荣中显现。}\BibleRef{哥 3:4}因此，那\BibleQuote{面对面}看见的景象是为你保留的，直到那时，若望也提到：\BibleQuote{亲爱的，我们已经是天主的子女，但将来如何，还没有显明；我们知道，主显现时，我们必要相似祂，因为我们要看见祂实在怎样。}\BibleRef{若一 3:2}虽然那时我们的祖先沒有看见你如你所是的\BibleQuote{面对面}，虽然那景象保留到复活，然而，即使是天使们显现，仍是\BibleQuote{是你命令拯救雅各伯。}你不仅以你自身临在；而且藉着任何受造物你显现时，是你藉着他们\Quote{命令}你为拯救你的仆人所做的；而那些你所\Quote{命令}的人所做的，是为了拯救你的仆人。既然你自身就是\BibleQuote{我的君王，我的天主，是你命令拯救雅各伯}，为什么我们还要受这些苦？\par

5. 但也许只是过去的事被描述给我们了，但这类事我们未来不该再盼望。不，其实仍是该盼望的。\BibleQuote{藉着你，我们将簸扬我们的仇敌}\BibleRef{咏 43:5}. 我们的祖先曾宣告你所作的工作，\Quote{在他们的时代，在古老的时代}，你的手消灭了外邦人：你\Quote{驱逐了万民，栽植了他们。}过去如此；但将来如何？\BibleQuote{藉着你，我们将簸扬我们的仇敌。}时候将到，所有基督徒的仇敌都要像糠秕一样被簸扬，像尘土一样被吹散，从地上被抛弃……因此未来的事如此。\BibleQuote{我不信靠我的弓，}正如我们的祖先不信靠\BibleQuote{他们的刀。我的刀也不能救我}\BibleRef{咏 43:6}.\par

6. \BibleQuote{因为你从我们的仇敌中拯救了我们}\BibleRef{咏 43:7}. 这也是以过去的形式讲述未来的事。之所以说成好像过去，是因为它确定如同过去。请注意，为什么先知们将许多事表达为好像过去，而预言的对象是未来的事，不是过去的事实。因为我们主自己将来的受难也被预言了，但经上说：\BibleQuote{他们刺穿了我的手和我的脚。我数过我的所有骨头；}而不是\Quote{他们要刺穿}和\Quote{他们要数说。}\BibleQuote{他们注视着我，盯着我；}而不是\Quote{他们要注视着我盯着我。}\BibleQuote{他们分了我的衣服。}经上没有说\Quote{他们将要分}它们。所有这些事都表达为好像过去，尽管它们尚未到来：因为对于天主，未来的事也像过去一样确定……因此，由于它们的确定性，那些尚未发生的事被说成好像已经过去。这就是我们所盼望的。因为经上说：\BibleQuote{你从我们的仇敌中拯救了我们，使恨我们的人蒙羞。}\par

7. \BibleQuote{我们终日要以天主夸耀}\BibleRef{咏 43:8}. 请注意他如何混入表示未来时间的词语，使你能看出之前用过去时所说的，实是对未来时间的预言。\BibleQuote{我们终日要以天主夸耀；我们要永远称谢你的名。}什么是\Quote{我们要夸耀}？什么是\Quote{我们要称谢}？就是\Quote{你从我们的仇敌中拯救了我们；}你要赐给我们永远的王国：在我们身上应验了这话：\BibleQuote{住在你殿中的有福了；他们要不断赞美你。}\par

8. 既然我们确信这些事将来必要成就，既然我们从祖先那里听说我们所说的那些事是过去发生的，那么我们目前的景况如何？\BibleQuote{但如今祢抛弃了我们，使我们受羞辱}\BibleRef{咏 43:9}。祢\Quote{使我们受羞辱}，不是在我们自己的良心面前，而是在人的眼中。因为曾有一个时期，基督徒受迫害；那时他们处处被遗弃，处处有人这样说：\Quote{他是基督徒！}仿佛那是一句侮辱和责骂。那么，那\Quote{我们的天主，我们的君王}、那\Quote{命令救恩归于雅各伯}的在哪里？那行了这一切事、\Quote{我们的祖先曾告诉我们}的在哪里？那将来要藉祂的圣神启示给我们的一切事的在哪里？祂改变了吗？不。这些事是为了\Quote{科辣黑子孙}的\Quote{智慧}而成就的。因为我们应当\Quote{明白}祂为何愿意我们在这期间忍受这一切事的道理。什么是\Quote{一切事}？\Quote{但如今祢抛弃了我们，使我们受羞辱，天主，祢不再同我们的军队出征。}我们出去迎敌，祢却不与我们同去。我们看见他们：他们非常强大，而我们毫无力量。祢的威能在哪里？祢的\Quote{右手}和祢的大能在哪里？那海水分开、埃及追兵被波浪淹没的地方在哪里？那阿玛肋克人的抵抗被十字架的记号征服的地方在哪里？\BibleRef{出 17:12}\Quote{天主，祢不再同我们的军队出征。}\par

9. \BibleQuote{祢使我们在敌人面前转身退后}\BibleRef{咏 43:10}，以致他们好像在前，我们在后；他们算为得胜者，我们算为失败者。\BibleQuote{仇恨我们的人任意抢夺。}他们\Quote{抢夺}了什么？不就是我们吗？\par

10. \BibleQuote{祢使我们如待宰的羊，把我们分散在列邦中}\BibleRef{咏 43:11}。我们被\Quote{列邦}所\Quote{吞吃}。这些人是说那些因受苦而逐渐同化、成为外邦世界\Quote{身体}一部分的人。因为教会为他们哀恸，如同为她的肢体中被吞吃的人哀恸。\par

11. \BibleQuote{祢将祢的子民廉价出售}\BibleRef{咏 43:12}。因为我们看见祢交付了谁，但祢得到了什么，我们却没有看见。\BibleQuote{在他们的喜乐庆祝中并无众多人群。}因为当基督徒逃避拜偶像的敌人的追逼时，那时还有举行任何聚集和荣耀天主的\Quote{喜乐庆祝}吗？那些在和平时期惯常在天主教会中齐声歌唱、以兄弟友爱的和谐音韵传入天主耳中的圣诗，那时还能歌唱吗？\par

12. \BibleQuote{祢使我们受邻邦的羞辱，被四周的人嘲笑讥讽}\BibleRef{咏 43:13}。\BibleQuote{祢使我们在异民中成了笑柄}\BibleRef{咏 43:14}。什么是\Quote{笑柄}？就是当人在咒骂时，用他们所憎恶之人的名字作为\Quote{笑柄}。\Quote{愿你如此死！} \Quote{愿你如此受罚！} \Quote{愿你如此被钉十字架！}即使在今天，基督的仇敌（就是那些犹太人自己）并不缺少；每当我们为基督辩护时，他们就对我们说：\Quote{愿你像祂那样死。}因为他们不会选择那种死法，除非他们对那样死去极其恐惧，或者他们无法理解其中所包含的奥秘。当药膏涂在盲人的眼睛上时，他看不见医生手中的眼药。因为十字架本身甚至是为了造福那些迫害者。后来他们藉此得了医治，并相信了他们自己所杀害的祂。\BibleQuote{祢使我们在异民中成了笑柄，在万民中成了摇头的对象}，一种侮辱性的\Quote{摇头}。\Quote{他们用嘴唇说话，摇头。}他们对主如此行，也对他们所能追捕、擒拿、嘲弄、出卖、折磨和杀害的一切圣徒如此行。\par

13. \BibleQuote{我的耻辱常在我面前，我面上的羞愧遮盖了我}\BibleRef{咏 43:15}。\BibleQuote{因那辱骂并亵渎者的声音}\BibleRef{咏 43:16}：这就是说，来自那些侮辱我、指控我敬拜祢、承认祢之人的声音！他们指控我背负那将涂去我一切指控的名字。\Quote{因那辱骂并亵渎者的声音}，即那反对我的人。\Quote{由于仇敌和迫害者。}此处传达的\OriginalTerm{领悟}{understanding}是什么？那些关于过去的事，不会发生在我们身上；那些所盼望的将来之事，尚未显明。过去的事，如祢以极大荣耀带领祢的子民出离埃及；拯救他们脱离迫害者；引导他们经过列邦，安置他们于被驱逐列邦的国度。将来之事是什么？当基督，我们的\Quote{首领}在祂的荣耀中显现时，带领子民出离这个世界的埃及：圣人被安置在祂右边，恶人在左边；恶人同魔鬼被定罪受永罚；与圣人一起从基督领受永恒的国度。这些是将来的事；前者是过去的事。在期间，我们的分是什么？磨难！\Quote{为何如此？}好显明那敬拜天主的灵魂，它对天主的敬拜到何等程度；它是否\Quote{白白地}敬拜那\Quote{白白地}赐予救恩的主。……你给了天主什么？你曾是恶人，却被赎回了！你给了天主什么？有什么不是你从祂\Quote{白白地}\Quote{领受}的？故有理称其为\OriginalTerm{恩宠}{grace}，因它是\OriginalTerm{白白地}{gratis}赐予的。\BibleRef{罗 11:6}向你要求的便是，你也应当\Quote{白白地}敬拜\Quote{祂}；不是因为祂赐你暂时之物，而是因为祂向你提供永恒之物。……\par

14. \BibleQuote{这一切临到了我们，我们却没有忘记祢}\BibleRef{咏 43:17}。\Quote{没有忘记祢}是什么意思？\Quote{也没有在祢的盟约中行事乖僻。}\par

\BibleQuote{我们的心没有退后，祢却使我们的脚步偏离祢的道路}\BibleRef{咏 43:18}。请看，这便是\OriginalTerm{领悟}{understanding}：\Quote{我们的心没有退后}；我们在极大的磨难和外邦人的迫害中，没有\Quote{忘记祢}，没有\Quote{在祢的盟约中行事乖僻}。\Quote{祢却使我们的脚步偏离祢的道路。}我们的\Quote{脚步}原在世间的享乐中；我们的\Quote{脚步}原在暂时的顺遂中。祢使\Quote{我们的脚步偏离祢的道路}；并指示我们\BibleQuote{通往生命的门是多么窄，路是多么狭}\BibleRef{玛 7:14}。\Quote{祢使我们的脚步偏离祢的道路}是什么意思？仿佛祂说：你处在磨难之中；你正遭受许多痛苦；你已失去了许多在今世所爱之物；但我并没有在道路上遗弃你，这狭窄的道路是我正在教你的。你曾寻求\Quote{宽路}。我对你说什么？这是我们通往永生的道路；照你所愿行走的路，你是走向死亡。\BibleQuote{通往毁灭的门是多么宽，路是多么广，找到它的人何其多！通往生命的门是多么窄，路是多么狭，找到它的人何其少}\BibleRef{玛 7:13}。那些少数人是谁？是那些耐心忍受磨难、耐心忍受试探的人；是那些在这一切患难中不\Quote{跌倒}的人；是那些不只为\Quote{一时}欢喜地接受圣言，而当太阳升起时在患难中枯萎的人；而是拥有\Quote{爱}的\Quote{根}，正如我们最近在福音中所读到的……\par

15. \BibleQuote{因为祢使我们屈辱在软弱之处}\BibleRef{咏 43:18}：因此祢必高举我们于刚强之处。\BibleQuote{死荫遮蔽了我们}\BibleRef{咏 43:19}。我们的这必朽性不过是死亡的\Quote{阴影}。真正的死亡是与魔鬼同受定罪。\par

16. \BibleQuote{若我们忘记了我们天主的圣名。}这里就是所谓\Quote{\Person{科辣黑}子孙}的\Quote{明悟}。\BibleQuote{并向外邦的神祇伸出了我们的手} \BibleRef{咏 43:20}。\BibleQuote{天主岂不查明此事？因为祂知晓人心的秘密} \BibleRef{咏 43:21}。祂\Quote{知晓}，却\Quote{查明它们}？若祂知晓人心的秘密，那么\BibleQuote{天主岂不查明此事？}这句话在那里做什么？祂在自己内\Quote{知晓}；祂为我们\Quote{查明它}。因为正是为此，天主有时\Quote{查明某事}，并谈论那对祂而言变得可知的事，而祂正使之对你成为可知。祂是在谈论祂自己的工程，而非祂的知识。我们常说：\Quote{一个愉快的日子}，当日光美好。但日子本身会体验到愉悦吗？不：我们说日子愉快，因为它使我们充满愉悦。我们也说\Quote{阴沉的天空}。并非云层中有任何这样的感觉，而是因为人们看到天空这样的景象时感到阴沉，所以称其为阴沉，因为它使我们阴沉。同样，当天主使我们知道时，祂被称为\Quote{知晓}。天主对\Person{亚巴郎}说：\BibleQuote{现在我晓得你敬畏天主了。} \BibleRef{创 22:12} 难道祂以前不知道吗？但\Person{亚巴郎}直到那时才认识自己：因为正是在那考验中，他认识了自己……而天主被称为\Quote{知晓}那祂曾使他知晓的事。\Person{伯多禄}认识自己吗？当他对那医生说：\BibleQuote{我愿与祢同死？} \BibleRef{路 22:33} 医生已经把脉，知道祂病人灵魂内的情况；病人却不知道。试探的关头来到；医生证明祂判断的正确：病人放弃了他的自负。因此，天主同时\Quote{知晓}并\Quote{查明它}。\Quote{祂早已知道。为什么祂要\InnerQuote{查明}？}为了你的原故：使你能认识你自己，并感谢那造你的主。\BibleQuote{天主岂不查明此事？}\par

17. \BibleQuote{为了祢的缘故，我们终日被杀；我们被看作待宰的群羊} \BibleRef{咏 43:22}。因为你可以看见一个人被处死；却不知道他为何被处死。天主知道这事。事情本身是隐藏的。但有人会对我说：\Quote{看，他因基督之名被囚在监里，他是为基督之名作宣认的人。}为什么异端者也宣认基督之名，却不为祂而死呢？不仅如此，让我说，在天主教会内，你以为就没有或曾缺少那些为了在人中的荣耀而受苦的人吗？若没有这样的人，宗徒就不会说：\BibleQuote{我若舍身被焚，却没有爱德，于我毫无益处。} \BibleRef{格前 13:3} 因此他知道，可能有某些人做这事并非出于\Quote{爱德}，而是出于虚荣。这在我们是隐藏的；唯有天主看见；我们不能看见。唯有祂能判断这事，因为祂\BibleQuote{知晓人心的秘密}。\Quote{因为}，为了祢的缘故\BibleQuote{我们终日被杀；我们被看作待宰的群羊}。我已经提过，\Person{保禄}宗徒曾从这里取材，以鼓励殉道者：使他们不致\Quote{在患难中丧胆}，这患难是他们为基督之名所受的。\par

18. \BibleQuote{醒来；主啊，祢为什么沉睡？} \BibleRef{咏 43:23}。谁被呼求，谁在说话？说这样的话的人，不更恰当被称为\Emph{沉睡}和昏睡吗？他回答你说：我知道我在说什么：我知道\Quote{那保护以色列者不打盹。}然而殉道者们呼喊：\BibleQuote{醒来；主啊，祢为什么沉睡？}主耶稣，祢曾被杀害；祢在受难中\Quote{沉睡}；对我们而言，祢如今已\Quote{醒来}。因为\Quote{我们}知道祢如今已再次\Quote{醒来}。祢醒来并复活又有何益？那迫害我们的外邦人以为祢死了；不相信祢已复活。那么，\Quote{求祢起来}对他们也是一样！\Quote{为什么沉睡，}虽不对于我们，却对于他们？因为他们若已相信祢复活了，岂能迫害我们这些信祢的人？但他们为什么迫害？\Quote{消灭、杀戮这人那人，凡是信了祢的，就是那死得不好的！}对他们而言，至今\Quote{祢仍在沉睡}；求祢在他们面前\Quote{起来}，使他们觉察到祢已\Quote{醒来}，并得安息。最后，当殉道者死去，说着这些话；当他们睡眠，并\Quote{唤醒}基督时，他们在睡眠中的的确确是死了，基督在某种意义上，在外邦人中复活了；即人们相信祂已复活，逐渐地，他们自己因信而皈依基督，聚集了一大群人；正是迫害者所惧怕的；于是迫害止息了。为什么？因为基督，从前对他们如同\Emph{睡着}，因他们不信，如今在外邦人中复活了。\Quote{起来，求祢不要永远抛弃我们！}\par

19. \BibleQuote{祢为什么掩起祢的面？}好像祢不在眼前；好像祢忘记了我们？\BibleQuote{并忘记我们的苦难和压迫？} \BibleRef{咏 43:24}。\par

20. \Quote{我們的靈魂屈服於塵土} \BibleRef{咏 43:25}。它屈服於何處？\Quote{塵土：}即\Emph{塵土}迫害我們。他們迫害我們，就是你曾說的：\Quote{惡人卻不是這樣；他們好像塵土，被風從地面吹散。} \Quote{我們的肚腹緊貼著地。}在我看來，他似乎表達了極度屈辱的懲罰，當人俯伏時，\Quote{他的肚腹緊貼著地。}因為凡是謙卑到屈膝地步的人，還有更低的屈辱程度可達；但若謙卑到\Quote{肚腹緊貼地面}，就再無更深的屈辱了。若有人想再進一步，那就不是使他俯伏，而是壓碎他了。或許他這話的意思是：我們在這塵土中\Quote{極其俯伏}；屈辱已無可復加。屈辱現已達到極點：那麼，願慈悲降臨……\par

21. \Quote{上主，求祢醒來，援助我們} \BibleRef{咏 43:26}。的確，親愛的諸位，祂已醒來並援助了我們。因為當祂醒來（即當祂復活，並被外邦人所認識）時，迫害止息，甚至那些曾經緊貼著地的人也從地上被舉起，並藉著補贖，回歸到基督的身體，儘管他們軟弱而不完全：這樣在他們身上應驗了經上的話：\Quote{祢的眼睛看見了我尚未成形的身體；在祢的書卷上，它們全部被記錄下來。}\par

\Quote{上主，求祢醒來，援助我們，並為了祢的聖名救贖我們；}這就是說，白白地；為了祢的聖名，不是為了我的功勞：因為祢願意這樣做，不是因為我配得祢為我這樣做。至於這點，即\Quote{我們沒有忘記祢；}我們的\Quote{心沒有後退；}我們\Quote{沒有向任何外邦神伸出雙手；}我們怎能靠自己做到呢？除非藉著祢的幫助？我們怎能有力氣做到呢？除非祢在我們內心呼喚、勸勉我們，並不捨棄我們？無論我們是在患難中受苦，還是在順境中歡樂，求祢救贖我們，不是因為我們的功勞，而是為了祢的聖名。\par

\SourceRecord{Translated by J.E. Tweed.；Nicene and Post-Nicene Fathers, First Series；Vol. 8.；Edited by Philip Schaff.；Buffalo, NY: Christian Literature Publishing Co.,；1888.；<http://www.newadvent.org/fathers/1801044.htm>.}

% structure: p0001 source=h1 target=section reason=page-title

\section{圣咏第四十五篇释义}

1. 这篇圣咏，正如我们曾与你们一同欢欣咏唱一样，我们恳求你们也以同样的方式与我们一同用心思考。因为它是关于神圣的婚宴而唱的；关于新郎和新娘；关于君王和他的子民；关于救主和那些将要得救的人……我们是祂的儿子，因为我们是\Quote{新郎的子女}；这篇圣咏正是向我们说的，其标题有这样的话：\Quote{科辣黑的子孙，关于将要改变的事。}\par

2. 何需我解释\Quote{关于将要改变的事}是什么意思呢？每一个自己\Quote{改变}了的人，都明白这话的意义。凡听见\Quote{关于将要改变的事}的人，请想想从前是什么，现在是什么。首先，让他看到世界本身已被改变：从前崇拜偶像，如今崇拜天主；从前服事自己所造之物，如今服事创造他们的那一位。请注意\Quote{关于将要改变的事}这话是在何时说的。如今，剩下的外教人已惧怕这种\Quote{改变}了的事态；那些不愿让自己被\Quote{改变}的人，看到教堂满座，神庙荒凉；看到这里人群聚集，那里孤寂冷清！他们对如此改变的事感到惊奇；让他们读一读这些事已预先被预言；让他们倾听那应许者的声音；让他们相信那履行应许者。但是，我们每位弟兄，也从\Quote{旧人}改变为\Quote{新人}：从无信者变为信徒；从盗贼变为施舍者；从奸夫变为贞洁者；从作恶者变为行善者。因此，我们当咏唱\Quote{关于将要改变的事}，并让那位藉以改变者开始被描述。\par

3. 接着经文说：\Quote{关于将要改变的事，科辣黑的子孙，为得理解；一首为所爱者的歌。}因为那位\Quote{所爱者}被祂的迫害者看见，但却不是为得\Quote{理解}。因为\Quote{他们若认识祂，就决不会把光荣的主钉在十字架上。}\BibleRef{格前 2:8}为获得这种\Quote{理解}，祂需要另一双眼睛，当祂说：\Quote{看见我的，就是看见我父。}\BibleRef{若 14:9}那么，现在让这篇圣咏论述祂；让我们在婚宴中欢喜，我们将会与那婚宴中的人在一起——那些被邀请赴婚宴的人；而被邀请的人正是新娘本身。因为教会是\Quote{新娘}，基督是新郎。通常，吟游诗人会为新郎新娘唱某些称为\OriginalTerm{婚歌}{Epithalamia}的诗歌。在那里所唱的一切，都是为赞美新娘和新郎而唱的。那么，在我们被邀请的婚宴中，难道没有洞房吗？那么，另一篇圣咏为何说：\Quote{祂在太阳中搭了帷幕；祂好像新郎走出洞房。}这婚配的结合是\Quote{圣言}与肉身的结合。这结合的新房，是童贞女的子宫。因为肉身本身与圣言结合；故此也说：\Quote{从此他们不是两个，而是一体。}教会从人类中被取为祂的净配，如此，肉身与圣言结合，成为教会之头；其余相信的人，成为那头的肢体……\par

4. \BibleQuote{我的心涌出优美的言辞}\BibleRef{咏 44:1}。说话的是谁？是圣父，还是先知？因为有些人认为这是圣父的位格在说话：\Quote{我的心涌出优美的言辞}，向我们暗示某种不可言喻的诞生。唯恐你或许以为祂取了某物，以之生出圣子（就如人取某物以生子，即婚姻的结合，没有它人不能生子），唯恐你因此以为天主生\Quote{圣子}需要婚姻结合，祂便说：\Quote{我的心涌出优美的言辞。}今天，你的心，人啊，生出一个计划，并不需要妻子：这计划生于你的心，你建造某些东西；在那建造物存在之前，设计就已存在；你将要生产的东西，已经存在于你用以生产它的那个东西之中；你赞美那尚未存在的建筑物——不是以其可见的形式，而是借设计的构思；没有人会赞美你的设计，除非你向他展示，或他看见你所做的。如果\Quote{万有是藉着圣言而造成的}\BibleRef{若 1:3}，而圣言出自天主，那么请你藉着圣言所创造的建筑来思考，并从那座建筑学习钦羡祂的智慧！那是怎样的一位圣言，藉着它创造了天地\BibleRef{希 11:3}，以及天上的一切光华，地上的各种丰饶，海洋的浩渺，空气的广被，星辰的璀璨，日月的光辉？这些是可见之物；还要超越这些；想想天使：\Quote{率领者、掌权者、宰制者和异能者}\BibleRef{哥 1:16}。万有都是藉着祂而造的。那么，这些美善之物是如何造成的呢？因为\Quote{涌出了一个优美的言辞}——藉着它，万有要被造成……\par

5. 接下来是：\Quote{我向君王讲述我所成就的事。}仍然是圣父在说话吗？如果圣父仍在说话，就让我们探究：我们如何能按照公教真信仰理解这句话：\Quote{我向君王讲述我所成就的事。}因为，如果是圣父向祂的圣子——我们的\Quote{君王}——讲述祂自己的工程，那么圣父向圣子讲述的是哪些工程呢？既然圣父的一切工程都是藉圣子而成就的？或者，在\Quote{我向君王讲述我的工程}这句话中，\Quote{我讲述}一词本身是否暗示了圣子的诞生？我担心这话能否被那些领悟迟钝的人理解；但我仍要说出来。让那些能跟上我的人跟上吧；免得因为我不说，连那些能跟上的人也跟不上。我们在另一篇\BibleBook{}{圣咏}中读到：\Quote{天主一次说了话。}祂藉先知如此多次说话，藉宗徒如此多次说话，在今日也藉祂的圣人们说话，祂却说\Quote{天主一次说了话}？除了指着祂的\Quote{圣言}，祂怎能\Quote{一次}说话呢？\BibleRef{希 1:1}然而，正如我们在上文另一句中理解了\Quote{我的心涌出美辞}是指圣子的诞生，似乎接下来的句子是一种重复，使得已经说过的\Quote{我的心涌出美辞}在现在所说的\Quote{我讲述}中再次被重复。因为\Quote{我讲述}是什么意思？\Quote{我发出圣言。}这圣言除了从祂的心、从祂的至深处，天主还能从哪里发出呢？你自己若不是从你的\Quote{心}中发出，就不会说出任何话；你这句一度发声、随即消逝的话，也只能从那里发出；而你对天主的\Quote{讲述}竟感到惊奇吗？然而天主的\Quote{讲述}是永恒的。你现在正在说一些话，因为你先前是沉默的；或者，看，你还没有发出你的话；但当你开始发出时，你仿佛\Quote{打破了沉默}；并产生了一个先前不存在的话。天主产生\Quote{圣言}并非如此。天主的\Quote{讲述}无始无终；然而祂所发出的\Quote{圣言}只有\Quote{一个}。如果祂所说的话消逝了，就让祂再发出另一个吧。但既然发出祂的那位是永存的，所发出的那位也是永存的；且只发出一次，没有终结；那么这\Quote{一次}本身也是无始的，并且没有第二次讲述，因为一次所说的话并不消逝。因此，\Quote{我的心涌出美辞}与\Quote{我向君王讲述我所成就的事}是同一回事。那么为什么说\Quote{我讲述我所成就的事}呢？因为在圣言自身中包含了天主的一切工程。凡天主在创造中所要造的一切，早已存在于\Quote{圣言}中；若非存在于圣言中，就不会存在于现实中——就像对你来说，若非存在于你的设计中，就不会存在于建筑物中；正如\BibleBook{若望}{福音}所说：\BibleQuote{凡受造的，在祂内是生命。}\BibleRef{若 1:3}那么，受造的东西已经存在；但它的存在是在圣言中；天主的一切工程都在那里存在，却尚未成为\Quote{工程}。然而\Quote{圣言}已经存在，因为\Quote{圣言是天主，与天主同在}；并且是天主之子，与圣父同为独一的天主。\Quote{我向君王讲述我所成就的事。}愿那领悟\Quote{圣言}的人听到祂\Quote{讲述}；愿他能与圣父一起看到永存的圣言；在祂内，那些尚未到来的事物也存在；在祂内，那些过去的事物也没有消逝。这些天主的\Quote{工程}存在于\Quote{圣言}中，如同在圣言中，如同在独生子中，如同在\Quote{天主的圣言}中。\par

6. 随后是什么呢？\Quote{我的舌头是快手笔者的笔。}我的兄弟，我问你们，天主的\Quote{舌头}与抄写员的笔有何相似之处呢？\Quote{磐石}与基督有什么相似呢？\newline{}\BibleRef{格前 10:4}\Quote{羔羊}与我们的救主有什么相似呢？\newline{}\BibleRef{若 1:29}或\Quote{狮子}与独生者的勇力有什么相似呢？\newline{}\BibleRef{默 5:5}然而，这些比拟已经作出；若非如此，我们就不能在一定程度上藉这些可见之物形成对\Quote{不可见者}的认识。因此，对于这卑微的笔的比喻，我们不要将它比作祂的卓越伟大，也不要轻蔑地拒绝它。因为我问，为何祂将祂的\Quote{舌头}比作\Quote{快手笔者的笔}？无论抄写员写得多么快，仍不能与另一篇圣咏所说的那种速度相比：\Quote{祂的话奔行极速。}但在我看来（如果人的悟性可以如此推测），这话也可以理解为是由圣父说的：\Quote{我的舌头是笔者的笔。}因为由\Quote{舌头}说出的话，发出声音就消逝了；而所写的则留存。既然天主发出\Quote{圣言}，而发出的圣言并不发出声音就消逝，而是发出并持续存在，天主宁可将其比作写下的文字，而非声音。但祂接着说\Quote{快手笔者的}，这是要激励心灵去\Quote{领悟}。然而，不要懒惰地停留于此，想着抄写员或某种快速的速记员；如果在此处所见的就是这些，那就会停留于此。要迅速地思索\Quote{快}这个词是什么意思。天主的\Quote{快}是无可超越的。因为书写时，字母接字母，音节接音节，话语接话语；只有当第一个写完时，我们才移到第二个。但在那\InnerQuote{快}中，没有任何可以超越速度，因为其中没有多个话语；然而，也没有任何事物被省略：因为万有都包含在唯一者内。\par

7. 看哪！现在，那如此发出的圣言，永恒的、永恒者的永恒之子，将作为\Quote{新郎}而来；\Quote{你比世人更美}\newline{}\BibleRef{咏 44:2}。\Quote{比世人}。我问，为什么不是比天使也美呢？为什么他说\Quote{比世人}，除非因为祂是人？免得你以为\Quote{基督耶稣是人}\newline{}\BibleRef{弟前 2:5}，他说：\Quote{你比世人更美}。即使祂自己\Quote{是人}，祂仍\Quote{比世人更美}；虽然在世人中，却\Quote{比世人更美}；虽然出于世人，却\Quote{比世人更美}。\Quote{你的嘴唇洋溢恩宠。}\Quote{法律是藉梅瑟传授的，恩宠和真理却由耶稣基督而来。}\newline{}\BibleRef{若 1:17}……\par

8. 也曾有人偏向将前面所有的段落理解为先知本人的话，甚至认为\Quote{我的心涌出美词}这一节也应理解为先知所说，仿佛他在唱一首赞美诗。因为无论谁向上主唱赞美诗，他的心就好像\Quote{涌出美词}，正如那亵渎上主的人，他的心是在涌出恶言。因此，接着所说的\Quote{我向君王倾诉我所创作的事}，他的意思是，人的主要工作就是赞美天主。归祂所有的是以祂的美善使你满足；归于你的是以感恩赞美祂。……\par

9. \Quote{我的舌头是快手笔者的笔。}曾有人以为先知是这样描述他所写的，因此将他的舌头比作\Quote{快手笔者的笔}；但他选择用\Quote{快手}一词来表达，是为了表明他所写的那些事即将\Quote{快速}来临；\Quote{快手}应被视为等同于\Quote{写快速的事}，即写那些不会长久迟延的事。因为天主并未长久迟延显示基督。凡是已经过去的，人们会多么迅速地意识到它已流逝！回想在你之前的世代，你会发现亚当的被造不过是昨天的事。我们读到，从起初一切就这样进行着：\newline{}\BibleRef{伯后 3:4}因此它们是\Quote{快速}完成的。审判之日也会\Quote{快速}来临。你要预先防备它的\Quote{快速}来临。它要\Quote{快速}来临；你更要\Quote{快速}皈依。审判者的面容将要显现；但你要留意先知的话：\Quote{让我们以忏悔来迎接祂的面容。}\par

10. \Quote{大能者！请将你的刀佩在腰间} \BibleRef{咏 44:3}。所谓\Quote{你的刀}，不就是\Quote{你的话}吗？他用这把刀驱散了他的敌人；用这把刀将儿子与父亲、\Quote{女儿与母亲、儿媳与婆婆}分开。我们在福音中读到这些话：\Quote{我来不是为送和平，而是送刀剑。} \BibleRef{玛 10:34} 并且\Quote{在一家之中，五个人将要分裂：三个反对两个，两个反对三个；} \BibleRef{路 12:52} 即\Quote{父亲反对儿子，女儿反对母亲，儿媳反对婆婆。} 这分裂是因哪把\Quote{刀}造成的，岂不是基督带来的那一把吗？的确，弟兄们，我们每天都看到这一点的体现。有些年轻人立志献身于天主的事业，他的父亲反对；他们\Quote{彼此分裂}：一个许诺世俗的遗产，另一个却爱慕天上的；一个许诺一样，另一个却选择更好的。父亲不应认为自己受了亏待：只是天主被置于他之上。然而，他与那愿意献身于天主服务的儿子争执。但属神的刀剑比血肉的纽带更有力地使他们分离。这同样发生在母亲反对女儿的事上；更甚的是儿媳反对婆婆。因为有时在一家中，婆婆和儿媳分别属于正统和异端。在那把刀强力作用的地方，我们并不害怕重复圣洗。难道女儿不能与母亲分裂，而儿媳不能与婆婆分裂吗？……\par

11. 他所说的\Quote{腰间}指的是什么？肉身。因此有这话：\Quote{权杖不离犹大，立法者不离他的腰间。} 亚巴郎本人（他得到了万民都要因他的后裔受祝福的应许）在他派仆人为儿子寻找并迎娶妻子时，凭着信心完全相信，在那个卑微的后裔中包含了伟大的名号，即天主子要从亚巴郎的后裔中诞生，从所有人中出来，他岂不就这样让仆人向他起誓，说：\Quote{请把你的手放在我的大腿下，} \BibleRef{创 24:2} 然后起誓；好像他说：\Quote{把你的手放在祭台上，或福音书上，或先知书上，或任何圣物上。} \Quote{请把你的手放在我的大腿下，} 他满怀信心，不以为耻，反而在其中领悟真理。\Quote{带着你的威能和荣耀。} 接受这义，你在这义中常常美丽而荣耀。\Quote{愿你驰骋，得胜前进，为王！} \BibleRef{咏 44:4} 我们岂不见到如此？岂不已经成就？他\Quote{驰骋，得胜前进，为王}；万民都屈服于他。在圣神中见到这些事是何等伟大！我们现在有能力在现实中经历同样的事！当这些话被说出时，基督尚未如此\Quote{为王}；尚未驰骋，\Quote{得胜前进}。那时他们正在传讲，如今已应验；在许多事上，我们已看到天主的应许应验；在少数事上，我们仍需恳求其应验。\par

12. \Quote{为了真理、温良和正义。} 真理已归还我们，当\Quote{真理从地上生出，正义由天上俯视。} 基督被呈现给人类的期待，使万民都在亚巴郎的后裔中蒙受祝福。福音已被宣讲。这就是\Quote{真理}。\Quote{温良}是何意？殉道者们受苦了，天主的国由此大大推进，传遍万民；因为殉道者们受苦，既不\Quote{跌倒}，也不抵抗；承认一切，不隐瞒任何；准备好面对一切，不退缩任何。奇妙的\Quote{温良}！这是基督的身体从它的头学来的。他首先\Quote{被牵去宰杀，像羔羊在剪毛者面前缄默，同样不开口。} \BibleRef{依 53:7} 如此温良，以至于在十字架上他说：\Quote{父啊，宽恕他们，因为他们不知道自己在做什么。} \BibleRef{路 23:34} 为何因为\Quote{正义}？他还要来审判，\Quote{按各人的行为予以报酬。} 他宣讲了\Quote{真理}；他耐心忍受了不义；他将来要带来\Quote{正义}。\par

13. \Quote{你的右手将奇妙地引导你。} 我们被他的右手引导；他藉他自己的右手。因为他是天主，我们是必死的人。他藉自己的右手被引导，即藉他自己的大能。因为父所有的大能，他也有；父的不朽，他也有；他有父的神性，父的永恒，父的能力。他的右手将奇妙地引导他，施行天主的作为，忍受人性的苦难，以他的至善推翻人邪恶的意志。即使现在，他仍在被引导到尚未有他临在的地方，而引导他的是他自己的右手。因为那引导他的是他自己赐予圣徒的。\Quote{你的右手将奇妙地引导你。}\par

14. \BibleQuote{你的箭是锋利的，极其有力}（见：\BibleRef{咏 44:5}）；这些话刺透人心，燃起爱火。因此，在\WorkTitle{雅歌}中记载：\BibleQuote{我为爱所受的伤}（见：\BibleRef{歌 2:5}）。因为她所说\Quote{我为爱所受的伤}，即是说，她陷入爱情，激情燃烧，为新郎叹息，她从他那里领受了圣言之箭。\Quote{你的箭是锋利的，极其有力}；既穿透，又有效；\Quote{锋利，极其有力}。\Quote{万民必仆倒在祢脚下}。谁是\Quote{仆倒}的？那些\Quote{受伤}的人也\Quote{仆倒}了。我们看见万国被基督征服；我们并没有看见他们\Quote{仆倒}。他解释他们在哪里\Quote{仆倒}，即\Quote{在心里}。正是在那里，他们曾高举自己反对基督，也正是在那里，他们\Quote{仆倒}在基督面前。\Person{扫禄}曾是基督的亵渎者：那时他高傲自大，他向基督祈祷，\Quote{他仆倒了}，他俯伏在他面前：基督的仇敌被击杀，好让基督的门徒得以生存！被一支从天射出的箭击中，\Person{扫禄}（还不是\Person{保禄}，仍是扫禄），仍高傲自大，尚未倒伏，在\Quote{心里}受了伤：他领受了那箭，他\Quote{在心}里仆倒了。因为他虽俯伏在地，却不是在那里心里仆倒：而是在他大声说\BibleQuote{主，祢要我做什么？}（见：\BibleRef{宗 9:6}）那里。但刚才你还要捆绑基督徒，将他们带去受罚：现在你却对基督说：\Quote{祢要我做什么？}啊，锋利而极其有力的箭，被它击中\Quote{扫禄}仆倒，从而成了\Quote{保禄}。正如他一样，\Quote{万民}也是这样；思量万国，观察他们如何服从基督。\Quote{万民}（于是）\Quote{必仆倒在祢脚下，在王的仇敌心中}；也就是，在祢仇敌的心中。因为他称他为王，承认他为王。\Quote{万民必仆倒在祢脚下，在王的仇敌心中}。他们曾是\Quote{仇敌}；他们被祢的箭击伤：他们仆倒在祢面前。从仇敌变成了朋友：仇敌死了，朋友存活。这就是\Quote{为那些将要转变的人}的意思。我们正寻求\Quote{理解}每一个单独的字，每一句独立的经文；但我们只应如此寻求它们的\Quote{理解}，以至于无人怀疑这些话是关于基督的。\par

15. \BibleQuote{天主，祢的宝座永远常存}（见：\BibleRef{咏 44:6}）。因为天主永远\Quote{\InnerQuote{祢受祝福}直到永远}，是由于\Quote{祢的唇倾注了恩宠}。然而，犹太王国的宝座是暂时的；属于那些在法律之下的人，不属于那些在\Quote{恩宠}之下的人。他来\Quote{赎回那些在法律之下的人}，并将他们置于\Quote{恩宠}之下。他的\Quote{宝座永远常存}。为什么？因为那第一个王国的宝座不过是暂时的：那么我们何以有\Quote{永远常存的宝座}呢？因为这是天主的宝座。啊，永恒的神性属性！因为天主不可能有暂时的宝座。\Quote{天主，祢的宝座永远常存——正直的权杖是祢王国的权杖}。\Quote{正直的权杖}是引领人类的权杖：他们以前是弯曲的、扭曲的；他们寻求为自己统治；他们爱自己，爱自己的恶行；他们不将自己的意志服从于天主；反而想要将天主的意志屈从于自己的私欲。因为罪人和不义的人通常对天主生气，因为天不下雨！却又希望天主不对他生气，因为他放荡。也正是出于这个原因，人每天坐着与天主争辩：\Quote{这是他应当做的：这件事他做得不好。}你确实看见自己所做的一切。他不知道他在做什么！是你弯曲了！他的道路是正直的。你何时能使弯曲的与正直的一致？它无法与之重合。好比你将一根弯曲的木棍放在平坦的路面上；它不能贴合上去；不能粘合；不能嵌入路面。路面处处平坦：但那根木棍是弯曲的；它无法与平整的路面吻合。天主的意志是\Quote{正直的}，你自己的是\Quote{弯曲的}：因为你不能符合它，它在你看来就是\Quote{弯曲的}：你当以它来治理自己；不要试图使之屈从于你的意志：因为你不能做到；它始终是\Quote{正直的}！你愿住在他内吗？\Quote{你当纠正自己}；这样，那统治你的他的权杖，对你而言将成为\Quote{正直的准则}。因此他也被称为王，源于\Quote{统治}。因为那不纠正的就不是\Quote{统治者}。为此，我们的王是\Quote{义人}的王。正如他是一位司铎（\OriginalTerm{司铎}{Sacerdos}）藉着圣化我们，同样，他藉着\Quote{统治}我们，是我们的王，我们的统治者……\par

16. \Quote{你喜爱正义，憎恨邪恶} \BibleRef{咏 44:7}。请看那里所描述的\Quote{正直的棍杖}。\Quote{你喜爱正义，憎恨邪恶。}靠近那\Quote{棍杖}；让基督作你的君王：让祂用那棍杖\Quote{统治}你，而不是击碎你。因为那棍杖是\BibleQuote{铁杖}，一根不可弯曲的棍杖。\BibleQuote{你必要用铁杖统治他们，像陶器一样打碎他们。}有些人祂统治，有些人祂\Quote{打碎}：祂\Quote{统治}属神的人，祂\Quote{打碎}属血肉的人……如果祂愿意打击你，祂会如此大声宣告祂要打击你吗？那么祂正在控制祂的手，不惩罚你的过犯；但你不可控制自己。你自己转向惩罚你的过犯：因为过犯不能不受惩罚；因此惩罚必须由你自己或由祂执行：那么你承认有罪，好让祂宽恕你。看看那忏悔圣咏中的例子：\BibleQuote{求祢转面不看我的罪过。}他意思是\Quote{转面不看我自己}？不是：因为在另一处他清楚地说：\BibleQuote{求祢不要转面不顾我。}那么\Quote{转面不看我的罪过。}我希望祢不看我的罪过。因为天主的\Quote{看}就是注意并惩罚。因此一位法官也被称为\Quote{注意}他所惩罚之事；即，把心思转向它，专注于它，甚至惩罚它，因为他是法官。同样，天主也是一位法官。\BibleQuote{求祢转面不看我的罪过。}但你自己，如果你希望天主转面不看它们，就不要转过你自己的脸不看它们。注意他如何在那篇圣咏中向天主提出这事：\BibleQuote{我承认，}他说，\BibleQuote{我的过犯，我的罪恶常在我的眼前。}他宁愿那他希望永远在自己眼前的事，不在天主的眼前。不要让任何人用对天主仁慈的虚假希望来恭维自己。祂的权杖是\Quote{正义的权杖}。我们能说天主不仁慈吗？有什么能超过祂的仁慈？祂对罪人显示了这样的忍耐，对一切归向祂的人都不计较过去。因此你要因祂的仁慈而爱祂，同时仍然希望祂是信实的。因为仁慈不能剥夺祂正义的属性，正义也不能剥夺仁慈的属性。与此同时，在祂延迟你的惩罚期间，你不可延迟惩罚自己。\par

17. \BibleQuote{因此，天主，祢的天主，用油傅了祢。}祂傅油给祢正是为此，使祢喜爱正义，憎恨邪恶。请注意他表达的方式。\BibleQuote{因此，天主，祢的天主，用油傅了祢：}即\Quote{天主傅了祢，哦，天主。}\Quote{天主}被天主\Quote{傅油}。因为在拉丁文中，这被认为是名词同一格的重复；但在希腊文中却有非常明显的区别：一个是受言者的名字，一个是说话者说\Quote{天主傅了祢。}\Quote{哦，天主，祢的天主用油傅了祢，}正如祂在说\Quote{因此，祢的天主，哦，天主，用油傅了祢。}照此理解；这在希腊文中最为明显。那么，被天主\Quote{傅油}的天主是谁？让犹太人告诉我们；这些圣经是我们和他们都共有的。是那被天主傅了油的天主：你听到\Quote{受傅者}；要理解为\Quote{基督。}因为\Quote{基督}这个名称来自\Quote{傅油}；祂被称为\Quote{基督}这个名称表达了\Quote{傅油}：除了在那个预言基督、傅油基督、基督之名将要来的王国里，任何其他地方的任何王国，任何其他地方的君王和先知都从未被傅过油。其他地方完全找不到：没有任何一个民族或王国。那么，天主被天主傅了油；祂被傅了什么油？不就是属神的油吗？因为可见的油是标记，不可见的油是奥秘；属神的油在内里。因此，\Quote{天主}被\Quote{傅油}是为了我们，并被派遣给我们；天主自己成了人，以便祂可以被\Quote{傅油}：但祂是那样地成了人，祂仍然是天主。祂是那样地是天主，以致不耻于成为人。\Quote{真人真天主；}毫无欺诈，毫无虚假，因为处处是真实的，处处是\Quote{真理}本身。那么天主就是人；正是因此，\Quote{天主}被\Quote{傅油}，因为天主成了人，并成了\Quote{基督}。\par

18. 这在雅各伯以石为枕而睡时预表了（\BibleRef{创 28:11}）。\newline{}先祖雅各伯将一块石头放在头下，枕着那石头睡觉，看见天开了，有梯子从天通到地，天使上去下来。这异象后他醒了，给石头傅了油，就起身走了。在那\Quote{石头}上，他认出了基督；为此他傅了油。请注意这所传的基督是什么。特别是先祖们只崇拜独一的天主，给石头傅油意味着什么？这却是一个象征性的举动；他就走了。因为他不是傅了油，就常在那里敬拜，在那里献祭。这是奥秘的表达，不是亵渎的开始。注意\Quote{石头}的含义。\BibleRef{咏 118:22}说：\BibleQuote{匠人弃而不用的石头，反而成了屋角的基石。}这里有一个大奥秘。\Quote{石头}就是基督。伯多禄称祂为\Quote{活石，固然被人弃绝，却是天主所拣选的}（\BibleRef{伯前 2:4}）。石头放在\Quote{头}下，因为\BibleQuote{基督是男人的头}（\BibleRef{格前 11:3}）。石头被傅油，因为基督的名字来源于祂的受傅。在基督的启示中，看见梯子从地上到天上，或从天上下到地上，天使上去下来。当我们从福音中引用主自己的证言时，会更清楚这含义。你们知道雅各伯就是以色列。当他与天使搏斗并\Quote{得了胜}（\BibleRef{路 23:23}），他所胜过的天使祝福了他，他的名字就改了，叫\Quote{以色列}；正如以色列子民\BibleQuote{得了胜}（\BibleRef{路 23:23}）反对基督，以致钉死祂，尽管如此，在信基督的人中，却被祂所胜过的祝福了。但许多人不信；因此雅各伯跛了。这里立即有祝福和跛行。祝福临到成为信徒的人；因为我们知道后来那民族中有许多人信了：跛行则临到不信的人。因为多数人不信，只有少数人信，所以为了产生跛行，他摸了他的\BibleQuote{大腿窝}（\BibleRef{创 32:25}）。大腿窝是什么意思？指他的后代众多……\par

19. \Quote{天主，你的天主，给你傅了油。}我们一直在讲那受傅的天主，即基督。基督的名字再清楚不过地藉此表达了：祂被称为\Quote{受傅的天主}。正如祂\Quote{在世人中最美丽}，同样祂\Quote{被喜油傅抹，胜过你的同伴}。那么谁是他的\Quote{同伴}？是世人；因为祂自己（作为人子）分享了他们的可朽性，为使他们分享祂的不朽性。\par

20. \Quote{你的衣服散发没药、沉香和肉桂的芬芳}（\BibleRef{咏 44:8}）。你的衣服散发出芬芳的香气。祂的衣服是指祂的圣人、祂的选民、祂的整个教会，祂展示教会如同祂的衣服；祂的长袍\BibleQuote{没有污点或皱纹}（\BibleRef{弗 5:27}），因着斑点祂用血\Quote{洗净}了它；因着\Quote{皱纹}祂钉在十字架上拉伸了它。因此有那里提到的某些香料所象征的甘饴芬芳。听保禄，那\Quote{宗徒中最小的}（那\Quote{衣穗}，患血漏的妇人摸了就得痊愈），听他说的：\BibleQuote{我们在各处，在得救的人中，也在丧亡的人中，都是基督的甘饴芬芳}（\BibleRef{格后 2:14}）。他没有说：\Quote{我们在得救的人中是甘饴芬芳，在丧亡的人中是恶臭；}但就我们自己而言，\Quote{我们在得救的人中，也在丧亡的人中，都是甘饴芬芳}。……爱他的人因\Quote{甘饴芬芳}的香气而得救；嫉妒他的人因那\Quote{甘饴芬芳}而丧亡。因此对丧亡的人，他不是恶臭，而是\Quote{甘饴芬芳}。正因为如此，他越是被那恩宠所充满，人们就越嫉妒他：因为没有人嫉妒不幸的人。他在传播天主圣言和按照\Quote{正直的棍杖}的规则管制自己的生活上，是荣耀的；那些在基督内爱他的人爱他，追随并追求甘饴芬芳的香气；他们爱新郎的朋友：这就是说，新娘自己，她在《雅歌》中说：\Quote{我们愿追随你香气的甘饴芬芳}。但其他人，越是看见他披戴着传福音的荣耀和无可指责的生活，就越被嫉妒所折磨，并发现那甘饴芬芳为他们成了死亡。\par

21. \Quote{祢从象牙宫殿出来，使君王们的女儿将祢欢喜。} 请随意选择，\Quote{象牙}宫殿，或\Quote{壮丽}、或\Quote{皇家}宫殿，正是从这些之中，君王们的女儿使基督欢喜。你愿理解\Quote{象牙宫殿}的属灵意义吗？要理解为天主的壮丽居所和帐幕，即圣人们的心；并通过这些\Quote{君王}本身，那些管束自己肉身的人；他们将人类情感中叛逆的平民置于自己的权下，他们惩戒身体，使之成为奴仆；因为正是从这些之中，君王们的女儿使祂欢喜。因为凡通过他们的宣讲和传福音所生的灵魂，都是\Quote{君王们的女儿}；而教会，作为宗徒们的女儿，乃是君王们的女儿。因为祂乃是\Quote{万王之王}；他们自己也是君王，关于他们曾说过：\BibleQuote{你们要坐在十二个宝座上，审判以色列十二支派。} \BibleRef{玛 19:28} 他们宣讲了\Quote{真理的圣言}，并生出教会，不是为自己，乃是为祂……因此正如\Quote{为他们的兄弟立嗣}，凡他们所生的，他们不给予\Quote{保禄派}或\Quote{伯多禄派}之名，而是\Quote{基督徒}。注意在这几节中，这层意义是否被清醒地保持。因为当他说\Quote{从象牙宫殿}时，他指的是王家宅邸，宽敞、尊贵、平安，如同圣人们的心；他又说：\Quote{因此君王们的女儿在祢的尊荣中使祢欢喜。} 她们的确是君王们的女儿，是你们宗徒的女儿，但仍然是\Quote{在祢的尊荣中}：因为她们为兄弟立嗣。因此保禄，当他看见那些他为其兄弟所立的人追随他自己的名字时，喊道：\BibleQuote{难道保禄为你们被钉死了吗？} \BibleRef{格前 1:13} ……不；因为他接着说：\Quote{或者你们是受洗归于保禄的名下吗？}\par

\Quote{君王们的女儿在祢的尊荣中使祢欢喜。} 要持守、紧紧抓住这个\Quote{在祢的尊荣中}。这就是拥有\Quote{婚宴礼服}的意思：寻求祂的尊荣，祂的荣耀。再者，将\Quote{君王们的女儿}理解为那些城市，它们由君王建立，并接受了信德：并从象牙宫殿（富裕、骄傲、高抬的宫殿）出来。\Quote{君王们的女儿在祢的尊荣中使祢欢喜；} 因为它们不寻求其创立者的尊荣，而是寻求祢的尊荣。请你在罗马指给我一座罗慕路斯的神庙，有如此大的尊荣，就像我可以指给你伯多禄的纪念碑一样。在伯多禄内，谁受到尊荣，难道不是那为我们而死的那一位吗？因为我们追随的是基督，不是追随伯多禄。即使我们是从已死之人的兄弟所生，我们却是以已死之人的名字被称呼。\BibleRef{申 25:26} 我们由一人所生，却是为另一人所生。看，罗马、迦太基以及其他几个城市都是君王们的女儿，然而它们\Quote{在祢的尊荣中使君王欢喜；} 所有这些构成了一位独一无二的新娘。\par

22. 何等的婚歌！看，在充满喜乐的歌声中，新娘亲自出来了。因为新郎正在来临。被描述的正是祂：我们的全部注意力都集中在祂身上。\par

\BibleQuote{王后站在祢的右边} \BibleRef{咏 44:9}。站在左边的那位不是王后。因为也有一个人要站在\Quote{左边}，将听到：\BibleQuote{进入永火里去！} \BibleRef{玛 25:41} 但将要站在右边的，将听到：\BibleQuote{来吧，我父所祝福的，承受自创世以来为你们预备的国度吧。} \BibleRef{玛 25:34} 在祢的右边站着王后，\BibleQuote{身穿金衣，披着各色绣衣}。这王后的衣服是什么？是既贵重又多色的：是各种不同语言中教义的奥秘：一种非洲语，一种叙利亚语，一种希腊语，一种希伯来语，一种这个，一种那个；正是这些语言产生了这件衣服的各种颜色。但正如衣服的各种颜色融合在同一件衣服中，同样，所有语言也在同一个信德中。在那件衣服里，要有多样性，但不要有裂痕。看，我们\Quote{理解了}语言多样性的各种颜色；而衣服指向合一：但在那个多样性本身中，\Quote{黄金}指的是什么？智慧本身。无论何种语言的多样性，但所宣讲的只有一种\Quote{黄金}：不是不同的黄金，而是同一黄金的不同形式。因为每一种语言所宣讲的是同一个智慧，同一个教义和纪律。语言中有多样性，黄金在思想中。\par

23. 先知向这位王后说话（因为他乐于向她歌唱），而且也是向我们每一个人说话，只要我们知道自己在哪里，并努力成为那一身体的一部分，实在在信德和望德中属于它，与基督的肢体结合。因为他是在对我们说话，说：\BibleQuote{听吧，女儿，你要看}\BibleRef{咏44:10}，如同是\Quote{祖先}中的一位（因为她们是\Quote{君王的女儿们}），尽管说话的是先知，或是宗徒；他向她说话，如同向女儿，因为我们惯于这样说：\Quote{我们的祖先先知们，我们的祖先宗徒们；}若我们称他们为\Quote{祖先}，他们也可以称我们为儿女：这是一个父亲向一个女儿发出的声音。\Quote{听吧，女儿，你要看。}\Quote{听}在先，然后\Quote{看}。因为他们带着福音来到我们这里；那福音已传给了我们，我们尚未看见，但听到时便相信了，因相信而将看见；正如新郎自己在先知中所说：\BibleQuote{一个我不认识的民族服事了我；我一听到耳中，它就服从了我。}\Quote{听到耳中}是什么意思？就是他们没有\Quote{看见}。犹太人看见了他，将他钉死；外邦人没有看见他，却相信了。让那来自外邦的王后穿着\Quote{金线衣，有各种颜色修饰}而来吧；让她从外邦人中穿着各种语言，在智慧的合一中而来吧；愿有人向她说：\Quote{听吧，女儿，你要看。}若你不听，你就不会\Quote{看}。\par

\BibleQuote{请你侧耳细听。}光是\Quote{听}还不够；要谦卑地听：低下你的耳朵。\BibleQuote{还要忘记你的民族和你的父家。}曾有一个\Quote{民族}，一个你的父家，你在那里出生，就是巴比伦民族，以魔鬼为你的君王。无论外邦人从哪里来，他们都是从他们的父魔鬼来的；但他们已弃绝了作魔鬼的儿女。\Quote{还要忘记你的民族和你的父家。}他使你成为罪人，将你生得可憎；另一位，因\BibleQuote{祂使不虔敬的人成义}\BibleRef{罗4:5}，将你重生得美丽。\par

24. \BibleQuote{因为君王恋慕你的美丽}\BibleRef{咏44:11}。那是怎样的\Quote{美丽}，岂非祂自己的化工？\Quote{恋慕美丽}——恋慕谁呢？恋慕那罪人、不义者、不虔敬者，正如她与她的\Quote{父}魔鬼同在，在她自己的\Quote{民族}中那样吗？不是，而是那被说\Quote{这位上来的是谁，被洗得洁白？}的人。她起初并不洁白，而是后来被\Quote{洗白}的。因为\BibleQuote{你们的罪虽似朱红，我将使它们洁白如雪}\BibleRef{依1:18}。\Quote{君王恋慕你的美丽。}这位君王是谁？\BibleQuote{因为祂是上主、你的天主。}现在你想，你岂不是应当舍弃你那个父亲和你自己的民族，来就这位君王，祂是你的天主？你的天主就是\Quote{你的君王}，你的\Quote{君王}也是你的新郎。你嫁给你的君王，祂是你的天主；由祂妆扮，由祂装饰；由祂救赎，由祂医治。你所有用以取悦于祂的一切，都是从祂而来。\par

25. \BibleQuote{提尔的女儿们要带着礼物朝拜祂}\BibleRef{咏44:12}。这位就是那位\Quote{君王，你的天主}，提尔的女儿们要带着礼物朝拜祂。提尔的女儿们是外邦人的女儿，以部分代整体。提尔，一座靠近这预言发出之地的城市，预表了将要信从基督的万民。那位客纳罕妇人就是从那里来的，她最初被叫做\Quote{狗}；因为你知道她是从那里来的，福音如此说：\BibleQuote{祂退到提洛和漆冬一带，看，有一个客纳罕妇人从那地区出来，}连同那里记载的一切。她起初在她\Quote{父家}，在她自己的\Quote{民族}中，不过是一条\Quote{狗}，但因为她来向那位\Quote{君王}呼求并追赶，因信祂而变得美丽，她听到了什么？\BibleQuote{妇人，你的信德真大！}\BibleRef{玛15:21}。\Quote{君王恋慕你的美丽。提尔的女儿们要带着礼物朝拜。}带着什么礼物？这位君王也愿被接近，并要充满祂的宝库；正是祂自己给了我们那可以充满礼物之物，并由你们充满。让她们来（祂说），\BibleQuote{带着礼物朝拜祂。}\BibleQuote{带着礼物}是什么意思？……\BibleQuote{施舍，一切对你们就都洁净了。}带着礼物来到祂面前，祂说：\BibleQuote{我喜欢仁慈胜过祭献。}对于那曾是未来事物的影子的古圣殿，他们曾带着公牛、公绵羊和山羊，以及各种不同的祭献动物而来；以那血成就一事，同时又预表另一事。如今那一切所预表的血已经来了：君王自己来了，祂自己要你们的\Quote{礼物}。什么礼物？施舍。因为祂自己将来要审判，并要亲自将\Quote{礼物}算给某些人。\BibleQuote{来吧，我父所祝福的。}\BibleRef{玛25:34} 为什么？\BibleQuote{我饿了，你们给了我吃的}\BibleRef{玛25:34}等等。这就是提尔的女儿们用来朝拜君王的礼物；因为当她们说：\BibleQuote{我们什么时候见过祢？}时，祂既在上又在下（因此有\Quote{上去下来}之说），祂说：\BibleQuote{你们对我这些最小兄弟中的一个所做的，就是对我做的。}\BibleRef{玛25:40}\par

26. ……\Quote{人民中的富人必恳求祢的面容。}而那些恳求那面容的人，以及他们所恳求的那位，都共同构成一个\OriginalTerm{净配}{Bride}，一个\OriginalTerm{王后}{Queen}，母亲和子女全都属于基督，属于他们的\OriginalTerm{元首}{Head}。……\par

27. \BibleQuote{君王女儿的一切光荣，都来自内在} \BibleRef{咏44:13}。不仅她外袍的表面是\BibleQuote{金穗，色彩缤纷}；而且那爱慕她美丽的那位，也知道她内在的美丽。什么是那些内在的魅力？是良心的魅力。在那里基督看见，在那里基督爱她：在那里祂向她说话，在那里惩罚，在那里加冕。所以，你们要隐密地施舍，因为\BibleQuote{君王女儿的一切光荣，都来自内在}。\BibleQuote{她穿着金穗，服饰多彩} \BibleRef{咏44:14}。她的美丽来自内在；然而在\Quote{金穗}之中，是各种语言的多样性：教义之美。如果它们没有那\Quote{来自内在}的美丽，又有何用？\Quote{贞女们将被带到君王之后。}这确实应验了。教会相信了；教会已在万民中形成。如今贞女们何等渴望在君王眼中蒙恩！她们为何如此感动？正是因为教会走在了她们前面。\Quote{贞女们将被带到君王之后，她的近亲女子也将被带到祢面前。}因为被带到祂面前的，不是外人，而是属于她的\Quote{近亲女子}。因为他曾说过\Quote{到君王那里}，他转而向祂说话：\Quote{她的近亲女子将被带到祢面前。}\par

28. \BibleQuote{她们必欢乐喜跃地被带来，被领进君王的圣殿} \BibleRef{咏44:15}。那\Quote{君王的圣殿}就是教会本身：正是教会进入\Quote{君王的圣殿}。那座圣殿由什么建造？由进入圣殿的人建造？除了天主的\Quote{忠信者}，谁是它的\BibleQuote{活石}？\BibleRef{伯前2:4} \BibleQuote{她们必被领进君王的圣殿。}因为有些贞女在君王的圣殿之外，就是异端者中的修女：她们确实是贞女；但这对她们有什么益处，除非她们被领进\Quote{君王的圣殿}？\Quote{君王的圣殿}在于合一：\Quote{君王的圣殿}不坍塌，不分裂，不割裂。那些活石的粘合剂就是\Quote{爱德}。\par

29. \BibleQuote{你的子孙将代替你的父亲} \BibleRef{咏44:16}。再明显不过了。如今请思考那\Quote{君王的圣殿}本身，因为祂代表它说话，为了那散布全世界的身体的合一：因为那些选择成为贞女的人，在君王眼中蒙恩，除非她们被领进君王的圣殿。\BibleQuote{你的子孙代替了你的父亲。}是宗徒生了你们：他们被\Quote{派遣}：他们是宣讲者：他们是\Quote{父亲们}。但他们可能永远在肉体中与我们同在吗？虽然他们中一位说：\Quote{我渴望离去，与基督同在，这更好得多；留在肉身中为你们是必要的。}他确实这样说了，但他能在这里留多久？能直到现在吗？能直到永远吗？难道教会因他们的离去而变得荒凉？绝对不是。\BibleQuote{你的子孙代替了你的父亲。}这是什么意思？宗徒被派给你们作为\Quote{父亲们}，代替宗徒的，是你们的子孙：就是被任命的主教们。因为如今，全世界的主教们从何而来？教会称他们为父亲；教会生了他们，并把他们安置在\Quote{父亲们}的宝座上。所以，不要以为你被遗弃了，因为你没有看到伯多禄，也没有看到保禄：没有看到那些生你们的人。从你们自己的后代中，兴起了一群\Quote{父亲}给你们。\BibleQuote{你的子孙代替了你的父亲。}请观察\Quote{君王的圣殿}多么广传，那些没有被领进\Quote{君王的圣殿}的贞女，当知道她们与那婚姻毫无关系。\BibleQuote{祢要使她们作全地的首领。}这是普世教会：她的子孙被立为\Quote{全地的首领}：她的子孙被任命代替\Quote{父亲们}。让那些被剪除的人承认这个真理，让他们来到那一个身体：让他们被领进君王的圣殿。天主到处建立了祂的圣殿：到处立定了\BibleQuote{先知和宗徒的基础。}\BibleRef{弗2:20}教会生了\Quote{子孙们}；使她\Quote{代替她的父亲}成为\Quote{全地的首领}。\par

30. \BibleQuote{他们世世代代必思念你的名，因此万民要向你宣认} \BibleRef{咏 44:17}。那么，实在\Quote{宣认}却是在\Quote{圣殿}之外宣认，又有何益？祈祷却不是住在山上的祈祷，又有何益？他说：\BibleQuote{我扬声呼号上主，祂从祂的圣山上应允了我。}从哪座\Quote{山}？就是经上所说：\BibleQuote{建在山上的城，是不能隐藏的。} \BibleRef{玛 5:14} 哪座\Quote{山}？就是达尼尔所见的那座山：\BibleQuote{由一块小石头长出来，打碎了世上所有的王国，并遮盖了全球。} \BibleRef{达 2:34} 希望领受的，就在那里祈祷；愿祈祷蒙垂听的，就在那里祈求；愿被宽恕的，就在那里宣认。\BibleQuote{因此，万民要永远向你宣认，世世无穷。}因为在永生中，的确不再有哀叹罪过之事；但是在天上那永存的城里，赞颂天主时，仍将不断宣认那幸福的伟大。因为正是这座城，另一篇圣咏曾对它歌唱：\BibleQuote{荣福之事，已经论及你，哦天主的城。}她正是基督的新娘，那位王后，一位\BibleQuote{君王的女儿，君王的配偶}……万民要为此原因甚至向她宣认；所有的心，如今为圆满的爱德所光照，赤露敞开，显明出来，使她可以最完全地认识自己的全部，而她在这里，在许多部分，对自己是未知的……\par

\SourceRecord{Translated by J.E. Tweed.；Nicene and Post-Nicene Fathers, First Series；Vol. 8.；Edited by Philip Schaff.；Buffalo, NY: Christian Literature Publishing Co.,；1888.；<http://www.newadvent.org/fathers/1801045.htm>.}

% structure: p0001 source=h1 target=section reason=page-title

\section{圣咏第四十六篇释义}

1. 它被称为 \Quote{一篇圣咏，直到终末，为\Person{科辣黑}之子，论及隐密之事}。那么，这是隐密的；但你们知道，在\Place{加尔瓦略}被钉十字架的那位自己，已撕裂了帐幔\BibleRef{玛 27:51}，使圣殿的隐密得以显露。再者，因为我们的主的十字架是一把钥匙，藉以开启封闭之物；让我们相信祂会与我们同在，使这些隐密得以启示。所说的 \Quote{直到终末}，常应理解为指基督。因为\BibleQuote{基督是法律的终末，使一切相信的人都获得义}。\BibleRef{罗 10:4}但祂被称为终末，不是因为祂消灭，而是因为祂完成。因为我们也称被吃的食物为结束，被编织的外衣也为结束；前者是为了消耗，后者是为了完成。因为当我们来到基督那里，就没有更远的地方可去，祂自己便被称为我们道路的终点。我们也不应以为，当我们来到祂那里，还当努力再来到父那里。因为\Person{斐理伯}也有过这种想法，他曾对祂说：\BibleQuote{主，请将父显示给我们，这为我们就够了}。他说\BibleQuote{这为我们就够了}时，他寻求的是满足和完成的终点。然后主说：\BibleQuote{\Person{斐理伯}，我与你们在一起这么久，你们还不认识我吗？看见了我，就是看见了父}。\BibleRef{若 14:8}所以，在祂内我们有了父，因为祂在父内，父也在祂内，祂与父原是一体。\par

2. \BibleQuote{天主是我们的避难所和力量}\BibleRef{咏 45:1}。有些避难所没有力量，人逃到那里，反而更软弱而非坚强。例如，你逃向世上某个更强大的人，为了使自己结交一位有力的朋友；这在你看来是一个避难所。然而，这世界是如此变幻无常，有权势者的崩溃一天比一天频繁，以致当你逃到这样的避难所时，你开始比以往更加恐惧……我们的避难所不是这样的，我们的避难所是力量。当我们逃到那里，我们将坚定不移。\par

3. \BibleQuote{在极重苦难中的助佑，这些苦难过于深重地找到了我们}。苦难很多，在每一种苦难中我们都必须逃到天主那里；无论是财产上的苦难、身体健康的苦难、亲人危险带来的苦难，还是维持生命所需的其他任何事，基督徒都不应有避难所，除了他的救主，除了他的天主；当他逃到那里，他就坚强。因为他不会在自己内坚强，也不会成为自己的力量，而是那成为他避难所的，将成为他的力量。但是，亲爱的诸位，在人类灵魂的一切苦难中，没有比罪恶的良知更大的苦难了。因为如果这一点没有受伤，人内里称之为良心的部分健全，那么无论他在其他地方遭受什么苦难，他都会逃到那里，并在那里找到天主……你们看，亲爱的诸位，当树木被砍下，被木匠检验时，有时表面上看起来好像损伤腐烂；但木匠查看树的内里，如果发现内里木质坚实，他就保证这树能在建筑中持久；他不会太在意表面损伤，只要内里坚实。同样，对人来说，没有比良心更内在的了；那么，如果外在健康而良心的精髓已经腐烂，又有何益？这些是紧密而极其剧烈的，正如这圣咏所说，极大的苦难；然而，即使在这些苦难中，主也藉着赦免罪恶成为助佑。因为不虔敬者的良心除了宽恕什么也不憎恨；因为若有人说他遭受极大苦难，作为国库的公认债务人，当他看见自己财产窘迫，无法偿清；若因讨债者年年压在他头上，他说自己遭受极大苦难，只在期望宽恕中才能喘息，且是在世俗之事上；那么，因罪恶浩繁而欠下刑罚的债务人，又当如何？从罪恶的良心中，他何时能偿还所欠的？若偿还，他必丧亡。所以，唯有藉着祂的宽恕，我们才能获得平安，但如此获得宽恕后，我们不可再回去负债……\par

4. 既然如此，获得了这样的平安，他们说什么？\BibleQuote{因此，当地被扰乱时，我们也不害怕}\BibleRef{咏 45:2}。刚才还焦虑，突然就平安了；从极大的苦难进入极大的安宁。因为在他们心中基督睡着了，所以他们颠簸；基督醒了（正如我们刚才从福音中听到的），祂命令风，风就平静了。\BibleRef{玛 8:24}因为基督藉着信德住在各人心里，这向我们表明，忘记信德的人，他的心像船在这世界的风浪中颠簸；好像基督睡着时它颠簸，但基督醒来，平静来临。不，主自己说了什么？\BibleQuote{你们的信德在哪里？}\BibleRef{路 8:25}基督被唤醒，信德被唤醒，使那曾在船上发生的事，也发生在他们心里。\BibleQuote{在极重苦难中的助佑，这些苦难过于深重地找到了我们}。祂使那里有了极大的安宁。\par

5. 看，何其宁静：\BibleQuote{因此，当地震动，山岳被投入海心时，我们决不畏惧。}那时我们将寻见不畏惧。让我们寻求被投入的山岳，若能寻见，显然这便是我们的保障。主曾真确地对祂的门徒说：\BibleQuote{如果你们有信德像芥子那样大，你们对这座山说：从这边移到那边去，它必会移去。}也许祂说\BibleQuote{对这座山}是指祂自己，因为祂被称为山：\BibleQuote{到末日，上主的山必被彰显。}\BibleRef{依 2:2}但这山位于众山之上，因为宗徒们也是山，支撑着这座山。因此接着说：\BibleQuote{到末日，上主的山必被彰显，屹立在众山之巅。}所以它超越众山之顶，被安置在众山之巅，因为众山宣讲这山。海象征这个世界，与之相比，犹太民族如同旱地。它未被拜偶像的苦涩所淹没，而是像干地，被外邦人的苦涩如海般包围。那时地要震动，即犹太民族；山要被投入海心，即首先那座屹立在众山之巅的大山。因为祂离开了犹太民族，来到了外邦人中间。祂从地被投入了海。谁投入了祂？宗徒们，祂曾对他们说：\BibleQuote{如果你们有信德像芥子那样大，你们对这座山说：从这边移到那边去，它必会移去。}即借着你们最忠实的宣讲，将成就此事：这座山，即我自身，要在外邦人中被宣讲，在外邦人中受荣耀，在外邦人中被承认，并且那关于我的预言也要应验：\BibleQuote{我未曾认识的民族要事奉我。}...\par

6. \BibleQuote{其中的水咆哮翻腾}\BibleRef{咏 45:3}：当福音被宣讲时，\BibleQuote{这是什么？他似乎是个讲论新奇鬼神的人}\BibleRef{宗 17:18}——这是雅典人说的；至于厄弗所人，他们以何等的骚动欲杀害宗徒们！在剧场中，为他们女神狄安娜，他们喧嚷起来，喊叫说：\BibleQuote{大哉！厄弗所人的狄安娜！}\BibleRef{宗 19:34}在那海浪和海的咆哮中，那些逃向这避难所的人并不畏惧。甚至宗徒保禄要进入剧场，却被门徒们拉住，因为他仍应为他们的缘故留在肉身内。然而，\BibleQuote{其中的水咆哮翻腾：山岳因它的威力而动摇。}谁的威力？海的？还是天主的？关于祂曾说过：\BibleQuote{护佑和力量，在困苦中极速的扶助。}因为山岳被动摇了，即这世界上的权势。天主的山是一回事，世上的山是另一回事：世上的山，其头是魔鬼；天主的山，其头是基督。但这些山被那些山摇动了。当山岳被动摇、海水咆哮时，他们便发出反对基督徒的声音；因为山岳被摇动，发生了大地震，海也震动。但这些是针对谁？针对那建立在磐石上的城。海水咆哮，山岳摇动，福音被宣讲。那么，天主的城如何？请听下面的话。\par

7. \BibleQuote{河流的支流使天主的城欢欣}\BibleRef{咏 45:4}。当山岳摇动，大海怒涛时，天主并未藉着河流的支流遗弃祂的城。这些河流的支流是什么？那圣神的丰沛流溢，主曾说过：\BibleQuote{谁若渴，让他到我这里来喝。信从我的人，从他的心中要流出活水的江河。}\BibleRef{若 7:37}这些江河于是从保禄、伯多禄、若望、其他宗徒、其他忠实的福音传播者的心中流出。既然这些江河从一条河流出，\BibleQuote{河流的支流使天主的城欢欣。}为使你知道这是指圣神说的，在同一福音中，圣史接着说道：\BibleQuote{他说这话，是指那信仰祂的人将要领受的圣神。那时圣神还没有赐下，因为耶稣尚未受到光荣。}\BibleRef{若 7:39}耶稣在复活后受到光荣，升天后受到光荣，在五旬节日圣神降临，充满了信徒们\BibleRef{宗 2:1}，他们说起各种语言，并开始向外邦人宣讲福音。因此天主的城欢欣，当时海因水声咆哮而扰动，山岳惊惶失措，询问该做什么，如何驱散这新教义，如何从地上根除基督徒的族类。他们反对谁？反对那\BibleQuote{河流的支流使天主的城欢欣}。因为祂借此表明祂所说的是哪条河；祂以\BibleQuote{河流的支流使天主的城欢欣}来指圣神。接着是什么？\BibleQuote{至高者圣化了祂的帐幕}：既然随后提到圣化，显然这些河流的支流应理解为圣神，每位信从基督的虔敬灵魂藉着祂被圣化，以成为天主的城的一位公民。\par

8. \BibleQuote{天主在她中间，她必不动摇}\BibleRef{咏 45:5}。让海翻腾，山摇动；\BibleQuote{天主在她中间，她必不动摇。}什么是\BibleQuote{在她中间}？是说天主站在一个地方，而信祂的人围绕祂吗？那么天主就被空间限制了；周围是宽的，被围绕的是窄的吗？绝对不可。不要这样想象天主，祂不受空间限制，祂的宝座是虔诚者的良心：天主的宝座在人心中，以至于如果人离开天主，祂仍安居于自身，并不像找不到地方的人那样跌倒。因为祂是提升你，使你在祂内，而不是依靠你，以至于你若离开，祂就跌倒。祂若离开，你必跌倒；你若离开，祂不会跌倒。那么，\BibleQuote{天主在她中间}是什么意思？表明天主对众人平等，不偏待人。因为如同中间之物与所有边界距离相等，天主被称为在中间，因为祂平等地眷顾众人。\BibleQuote{天主在她中间，她必不动摇。}她为何不动摇？因为天主在她中间。祂是\BibleQuote{在困苦中的帮助，困苦过于我们所能承担的。天主必以祂的面容帮助她。}什么是\BibleQuote{以祂的面容}？即以祂的显现。天主如何显现自己，使我们看见祂的面容？我已经告诉过你们；你们已经学到了天主的临在；我们通过祂的作为学到了。当我们从祂那里得到任何帮助，以至于我们毫不怀疑那是主赐予我们的，那时天主的面容就与我们同在。\par

9. \BibleQuote{外邦人喧嚷}\BibleRef{咏 45:6}。如何喧嚷？为何喧嚷？为了推倒天主的城，城中有天主在其中？为了推翻圣洁的帐幕，天主以祂的面容帮助它？不：现在外邦人是在有益的喧嚷中。因为接着是什么？\BibleQuote{邦国屈膝。}他说，邦国屈膝了；现在不昂首以便发怒，而是屈膝以便朝拜。邦国何时屈膝？当另一篇圣咏所预言的事成就时：\BibleQuote{诸王都要向祂俯伏，列国都要事奉祂。}什么原因使邦国屈膝？听其原因。\BibleQuote{至高者发出祂的声音，大地震动。}偶像崇拜的狂热者，像沼泽中的青蛙，越喧嚣、越污秽，在污秽和泥泞中聒噪。青蛙的呱呱声与云中的雷声相比如何？因为从云中\BibleQuote{至高者发出祂的声音，大地震动：}祂从祂的云中发出雷声。祂的云是什么？祂的宗徒，祂的宣讲者，藉着他们，祂在诫命中如雷，在奇迹中如电。这些云也是山：山因其高度和坚固，云因其雨水和丰饶。因为这些云浇灌了大地，经上论到它们说：\BibleQuote{至高者发出祂的声音，大地震动。}因为正是针对这些云，祂威胁某个不结果实的葡萄园，那里的山被投入海心；\BibleQuote{我要命令，}祂说，\BibleQuote{云不再降雨在其上。}\BibleRef{依 5:6}这应验在我刚才提到的事上，当山被投入海心时；当经上说：\BibleQuote{天主的道应当先讲给你们是必要的；但你们既然弃绝它，我们就转向外邦人。}\BibleRef{宗 13:46}那时就应验了：\BibleQuote{我要命令云不再降雨在其上。}犹太民族就像地上的干羊毛。因为你们知道，在某个奇迹中，地是干的，只有羊毛是湿的，但羊毛上的雨并未显现。\BibleRef{民 6:36}同样，新约的奥秘在犹太民族中并未显现。那里是羊毛，这里是帕子。因为奥秘隐藏在羊毛中。但在大地上，在所有民族中，基督的福音公开显露；雨水显明，基督的恩宠毫无遮蔽，因为它没有被帕子遮盖。但为了让雨从中而出，羊毛被挤压。因为他们压迫自己，把基督赶出，如今主从祂的云中向大地降雨，羊毛却保持干燥。但那时\BibleQuote{至高者发出祂的声音，}从那些云中；这声音使邦国屈膝并朝拜。\par

10. \BibleQuote{万军的上主与我们同在；雅各伯的天主是我们的扶持者}\BibleRef{咏 45:7}。不是任何人，不是任何权势，总之，不是天使，也不是任何受造之物，无论是地上的还是天上的，而是\BibleQuote{万军的上主与我们同在；雅各伯的天主是我们的扶持者。}祂派遣了天使，在天使之后来了，祂来是为使天使事奉祂，祂来是为使人能与天使同等。伟大的恩宠！如果天主与我们同在，谁能反对我们？\BibleQuote{万军的上主与我们同在。}与我们同在的是哪位万军的上主？\BibleQuote{如果天主与我们同在，谁能反对我们？祂没有怜惜自己的儿子，反而为我们众人把祂交出，怎么不把一切也白白赐给我们呢？}\BibleRef{罗 8:31}因此，让我们安全，在心灵的宁静中，以上主的食粮滋养一个良好的良心。\BibleQuote{万军的上主与我们同在；雅各伯的天主是我们的扶持者。}无论你的软弱多么大，请看谁扶助你。一个人病了，医生被叫到他那里。医生称病人为他的扶持者。谁扶助了他？正是祂。伟大的救恩希望；伟大的医生扶助了他。什么医生？除了祂，每个医生都是人：每个来给病人看病的医生，第二天也可能生病，唯有祂不会。\BibleQuote{雅各伯的天主是我们的扶持者。}你完全成为一个小孩子，像那些被父母抱起来的孩子。因为那些没有被抱起的，是被遗弃的；被抱起的，被养育。你认为天主这样抱起你，如同你母亲抱起婴儿时那样吗？不，而是到永远。因为你的声音在那篇圣咏中：\BibleQuote{我的父亲和母亲离弃了我，但上主收留了我。}\par

11. \BibleQuote{你们来看主的作为} \BibleRef{咏 45:8}。关于这个被提，主做了什么？思考整个世界，来看吧。因为如果你不来，你就看不见；如果你看不见，你就不相信；如果你不相信，你就远远站立；如果你相信，你就来；如果你相信，你就看见。因为我们是如何到达那座山的？不是步行吗？是靠船吗？是靠翅膀吗？是靠马吗？因为所有涉及空间和位置的事，不要担心，不要自寻烦恼，祂来找你。因为祂从一块小石头长大，变成一座大山，以至于充满了全地。那么，你为何还要通过陆地到祂那里去，祂已经充满了所有陆地？看，祂已经来了：你要留意。祂通过成长唤醒了沉睡的人；如果他们没有沉睡到如此之深，以至于即使面对来临的山也变得刚硬；但他们听到：\BibleQuote{你这睡着的人，当醒过来，从死者中起来，基督必光照你。} \BibleRef{弗 5:14} 因为对犹太人来说，看见这块石头是一件大事。因为石头还是小的：他们理所当然地轻视了它，轻视了就绊倒，绊倒就破碎了；剩下的就是被碾成粉末。因为关于这块石头是如此说的：\BibleQuote{凡掉在那石头上的，必被碎；那石头掉在谁身上，就要把谁压碎。} \BibleRef{路 20:18} 破碎是一回事，被碾成粉末是另一回事。破碎比碾成粉末轻；但祂在崇高降临时，只碾碎那些在卑微时已被破碎的人。因为在祂降临之前，祂在犹太人面前是卑微的，他们绊倒于祂，就破碎了；以后祂要在审判中来临，荣耀而崇高，伟大而有力，不再软弱受审判，而是刚强审判，并将那些绊倒于祂而破碎的人碾成粉末。因为\BibleQuote{一块绊脚石，一块使人跌倒的磐石} \BibleRef{伯前 2:8}是给那些不信的人的。因此，弟兄们，如果犹太人不认识祂，也就不奇怪了，因为他们轻视了那块放在他们脚前的小石头。更令人惊奇的是，那些现在连这样一座大山也不认识的人。犹太人在小石头上因看不见而绊倒；异端者却在大山上绊倒。因为现在那块石头已经长大，现在我们对他们说：看，达尼尔的预言已经应验：\BibleQuote{那块小石头变成一座大山，充满了全地。} \BibleRef{达 2:35} 你们为何要绊倒于祂，而不上去到祂那里？有谁如此盲目，竟会绊倒在大山上？祂到你这里来，是为了让你有绊脚的地方，而没有上去的地方吗？这是依撒意亚说的：\BibleQuote{来，我们到主的山上去} \BibleRef{依 2:3} \Quote{来，我们上去}是什么意思？\Quote{来}就是相信。\Quote{我们上去}就是进步。但他们既不来，也不上去，既不相信，也不进步。他们对着山狂吠。即使现在，他们屡次绊倒于祂而破碎，却不肯上去，总是选择绊倒。我们对他们说：\Quote{你们来看主的作为}，\Quote{祂在地上所设立的异兆} 是什么？它们被称为异兆，因为它们预告了某些事情，那些当世界相信时所行的奇迹的征兆。然后后来发生了什么，它们预告了什么？\par

12. \BibleQuote{祂使战争止息，直到地极}\BibleRef{咏 45:9}。这事我们尚未见成就：仍有战争，列国间为争夺主权而战；各教派间，犹太人、异教徒、基督徒、异端者之间，常有战争，有的为真理，有的为谬误而争。那么，\BibleQuote{祂使战争止息，直到地极} 尚未成就；但或许将来会成就。或者如今也已成就？在有些人身上成就了，在麦子中成就了，但在莠子中尚未成就。那么，\BibleQuote{祂使战争止息，直到地极} 是什么意思？祂称那些与天主争战的为战争。但谁与天主争战呢？不虔敬。而不虔敬对天主能做什么？什么也不能。陶器撞在磐石上，无论撞得多么猛烈，于己的损害岂不更大吗？这些战争是大的，是频繁的。不虔敬与天主争战，陶器被撞碎，就是人仗恃自己，过于依靠自己的力量而遭粉碎。这就是约伯论及一个不虔敬者时所说的盾牌。\BibleQuote{他挺着坚硬的盾牌，跑向天主。}\BibleRef{约 15:26} 所谓\BibleQuote{挺着坚硬的盾牌}是什么意思？就是过于信赖自己的保护。那些说\BibleQuote{天主是我们的避难所和力量，是我们在极大困苦中的帮助} 的人，是这样的吗？或在另一篇圣咏中：\BibleQuote{因为我不信赖我的弓，我的剑也不能拯救我。} 当人学到在自己内一无所能，在自己内毫无帮助时，他手中的武器就被折断，战争就止息了。那至高者从祂的圣云中发出的声音，使大地震动，万国屈膝，摧毁了这样的战争。祂已使这些战争止息直到地极。\BibleQuote{祂要折断弓，粉碎武器，用火焚烧盾牌。} 弓、武器、盾牌、火。弓是阴谋；武器是公开的战争；盾牌是徒然自恃的保护；焚烧它们的火，就是主所说的\BibleQuote{我来是把火投在地上}\BibleRef{路 12:49}，圣咏论及这火说：\BibleQuote{没有什么能躲避它的热力。} 这火燃烧起来，任何不虔敬的武器都不能存留于我们内，必须全部折断、粉碎、焚烧。你当安然无恙，不依靠自己的任何帮助；你越软弱，越没有自己的武器，那被称为\BibleQuote{雅各伯的天主是我们的扶持者} 的就越扶持你。……但当天主扶持我们时，祂会让我们徒手而去吗？祂武装我们，却是用另一种武器，福音的武器，真理、节制、救恩、信德、望德、爱德的武器。我们将拥有这些武器，但不是来自我们自己；我们原本拥有的武器被焚毁了；然而，如果我们被圣神的火点燃，就是那论及\BibleQuote{祂要用火焚烧盾牌} 的火，那么，你——曾想靠自己强大——天主使你软弱，为叫你在祂内坚强，因为你靠自己本是软弱的。\par

13. 接着是什么？\BibleQuote{你们要静止。} 为什幺？\BibleQuote{并要知道我是天主}\BibleRef{咏 45:10}。就是说，不是我，而是我是天主。我创造，我更新；我形成，我再形成；我造成，我重造。你若不能造你自己，又怎能更新你自己？人的灵魂的争闹骚动看不见这一点；对这争闹骚动说：\BibleQuote{你们要静止。} 就是说，抑制你们的灵魂，不要反驳。不要辩论，仿佛武装起来反抗天主。否则你们的武器仍然活着，尚未被火焚烧。但如果它们被焚烧了，\BibleQuote{你们要静止，} 因为你们没有可以用来战斗的东西。如果你们在自身内静止，从我寻求一切——你们从前曾仗恃自己——那么你们将要\BibleQuote{看到我是天主。} \BibleQuote{我要在异教徒中被尊崇，我要在地上被尊崇。} 刚才我以\Quote{地}的名称指犹太民族，以\Quote{海}的名称指其他民族。山岭被投入海心；万民骚动，万国屈膝；至高者发出声音，大地震动。\BibleQuote{万军的上主与我们同在，雅各伯的天主是我们的扶持者}\BibleRef{咏 45:11}。奇迹在异教徒中施行，异教徒的信德得以圆满；人自恃的武器被焚烧。他们在心灵的宁静中静止，承认天主是他们一切恩赐的根源。在这次荣耀之后，祂难道会抛弃犹太民族吗？宗徒论及这民族说：\BibleQuote{我对你们说，免得你们自以为聪明：以色列人一部分的固执已经发生，直到外邦人的数目满了}\BibleRef{罗 11:25}。就是说，直到山岭被移来，云彩在这里降雨，主用雷声在这里使万国屈膝，\BibleQuote{直到外邦人的数目满了。} 然后呢？\BibleQuote{这样，全以色列也必得救。} 因此，这里也遵循同样的顺序：\BibleQuote{我要被尊崇}（祂说）\BibleQuote{在异教徒中，我要在地上被尊崇；} 就是说，既在海中，也在地上，叫众人现在都能说以下的话：\BibleQuote{雅各伯的天主是我们的扶持者。}\par

\SourceRecord{Translated by J.E. Tweed.；Nicene and Post-Nicene Fathers, First Series；Vol. 8.；Edited by Philip Schaff.；Buffalo, NY: Christian Literature Publishing Co.,；1888.；<http://www.newadvent.org/fathers/1801046.htm>.}

% structure: p0001 source=h1 target=section reason=page-title

\section{\WorkTitle{圣咏第四十七篇释义}}

1. 本圣咏的标题如下：\Quote{致终末：为科辣黑之子，达味自己的一篇圣咏。}这些科辣黑之子也出现在其他圣咏的标题中，暗示着一个甜蜜的奥秘，隐含着伟大的圣事。我们乐意在其中有自知之明，并在标题中承认我们这些聆听、阅读并以明镜自照的人。科辣黑之子是谁？……或许是新郎之子。因为新郎在加耳瓦略山被钉十字架。回忆福音\BibleRef{玛 27:33}，那里他们钉死了主，你会发现祂被钉在加耳瓦略山。此外，那些嘲笑祂十字架的人，被魔鬼如同野兽般吞噬。这也由某处圣经所预示。当天主的先知厄里叟上坡时，有些孩童跟在他后面嘲笑说：\BibleQuote{上去，秃头！上去，秃头！}但他并非出于残忍，而是出于奥秘，让树林里的母熊出来吞噬了那些孩童\BibleRef{列下 2:23}。如果那些孩童没有被吞噬，他们会活到如今吗？或者他们既生而必死，难道不能因热病而丧命吗？但在他们身上若没有显示出奥秘，后代就无法受到警戒。因此，不要让任何人嘲笑基督的十字架。犹太人被魔鬼附身而遭吞噬，因为他们在加耳瓦略山钉死基督，并将祂举在十字架上，仿佛带着孩童般的理智，却不明所以地说：\BibleQuote{上去，秃头。}因为\Quote{上去}是什么？\BibleQuote{钉死祂，钉死祂。}\BibleRef{路 23:21}因为孩童被置于我们面前，作为谦卑的榜样，也被置于我们面前，作为愚昧的警戒。主将孩童领到跟前\BibleRef{玛 18:2}，并因为有人阻止他们而说：\BibleQuote{让小孩子到我这里来，因为天国正是属于这样的人。}\BibleRef{玛 19:14}孩童的榜样被置于我们面前，作为愚昧的警戒，正如宗徒所说：\BibleQuote{弟兄们，不要在心智上作小孩子}；但他又提出孩童作为效法的榜样：\BibleQuote{但在邪恶上却要作婴孩，好使心智成为成年人。}\BibleRef{格前 14:20}这篇圣咏是为科辣黑之子而唱的；因此是为基督徒而唱的。让我们作为新郎之子来聆听，那在加耳瓦略山被无知孩童钉死的新郎之子。因为他们配被野兽吞噬；我们却配被天使加冕。因为我们承认我们主的谦卑，并且不以此为耻。我们不以祂在奥秘中被称为\OriginalTerm{秃头}{Calvus}为耻，这称呼来自加耳瓦略山。因为就在祂被侮辱的十字架上，祂并没有让我们的额头变成秃头；相反，祂用自己的十字架标记了它。最后，要让你知道这些话是对我们说的，请看看所说的话。\par

2. \BibleQuote{万民哪，你们要鼓掌}\BibleRef{咏 46:1}。难道犹太人的民族就是万民吗？不，但以色列发生了部分的迟钝，致使无知的孩童喊叫：\BibleQuote{秃头！秃头！}于是主在加耳瓦略山被钉死，好让祂所流的血救赎外邦人，并应验宗徒所说：\BibleQuote{以色列发生了部分的迟钝，直到外邦人的数目全数进入。}\BibleRef{罗 11:25}因此，让那些虚妄、愚昧、无知的人嘲笑好了，让他们说\BibleQuote{秃头！秃头！}但你们是藉着祂在加耳瓦略山所流的血而蒙救赎的，你们要宣告：\BibleQuote{万民哪，你们要鼓掌。}因为天主的恩宠已降临到你们身上。\BibleQuote{你们要鼓掌。}什么是\Quote{鼓掌}？就是喜乐。但为什么用手？因为藉着善行。不要口里喜乐，双手却闲散。若你们喜乐，就当\Quote{鼓掌}。让祂看见万民的手，因为祂乐意赐予他们喜乐。万民的手是什么？是他们的善行。\BibleQuote{万民哪，你们要鼓掌，向天主欢呼胜利之声。}既要有声音，又要有手。若只有声音，不好，因为手是怠惰的；若只有手，也不好，因为舌头是沉默的。手和舌必须一致。愿这手承认，愿那手工作。\BibleQuote{向天主欢呼胜利之声。}\par

3. \BibleQuote{因为至高者上主是可畏的}\BibleRef{咏 46:2}。至高者降下，看似可笑；但升天以后，变为可畏。\BibleQuote{他是统治全地的大王。}不仅统治犹太人；因为祂也是他们的王。因为宗徒中有犹太人，他们之中成千上万的人卖了财产，把价银放在宗徒脚前\BibleRef{宗 4:34}，在他们身上应验了十字架上所写的名号：\BibleQuote{犹太人的君王。}\BibleRef{玛 27:37}因为祂也是犹太人的王。但\Quote{犹太人的王}是小的。\BibleQuote{万民哪，你们要鼓掌：因为天主是统治全地的大王。}因为祂不满足于只统治一个民族；因此祂从自己肋旁付出了如此巨大的代价，好买赎整个世界。\par

4. \Quote{祂将万民服在我们以下，将列邦放在我们脚下}\BibleRef{咏 46:3}。是谁征服的？征服了什么？谁在说话？或许是犹太人？当然，如果是宗徒，当然是；如果是圣人，当然是。因为天主将万民和列邦服在这些人的以下，以致今天他们在列邦中受尊敬，而那些被自己的同胞杀害的人：正如他们的主被自己的同胞杀害，却在列邦中受尊敬；被自己的人钉在十字架上，被外邦人朝拜，但那些人因被赎回而成为祂自己的人。因为祂赎买了我们，使我们不再是与祂隔绝的。那么，你是否认为这些话是宗徒说的：\Quote{祂将万民服在我们以下，将列邦放在我们脚下}？我不知道。宗徒们竟会如此骄傲地说，以列邦被放在他们脚下为乐——即基督徒在宗徒脚下——这令人奇怪。因为他们喜乐的是，我们与他们一同在那位为我们而死者的脚下。因为那些愿属于保禄的人，保禄对他们说过：\Quote{难道保禄为你们被钉在十字架上吗？}\BibleRef{格前 1:13}。那么这里，我们该如何理解？\Quote{祂将万民服在我们以下，将列邦放在我们脚下}。所有属于基督产业的人都在\Quote{万民}中，所有不属于基督产业的人也都在\Quote{万民}中；你看到基督的教会因基督的名如此被高举，以致所有尚未相信基督的人都伏在基督徒的脚下。现今有多少人奔向教会；他们尚未成为基督徒，却向教会求助；他们愿意在现世得到我们的援助，尽管还不愿与我们一同永远为王。当所有人都向教会求助，甚至那些尚未在教会内的人，难道祂没有\Quote{将万民服在我们以下，将列邦放在我们脚下}吗？\par

5. \Quote{祂为我们选择了产业，就是祂所爱的雅各伯的华美}\BibleRef{咏 46:4}。祂选择了雅各伯的某种美作为我们的产业。厄撒乌和雅各伯是两兄弟；他们在母腹中争斗，因此母亲的肠腹被搅动；当他们还在其中时，年幼的被拣选并优先于年长的，并且曾说：\Quote{两个民族在你腹中，年长的要服事年幼的}\BibleRef{创 25:23}。在万民中有年长的，在万民中有年幼的；但年幼的在好基督徒中，即蒙选、虔诚、忠信的；年长的在骄傲、不配、有罪、顽固、宁愿辩护而不承认自己罪过的人中；就如犹太人的百姓本身，\Quote{他们不知道天主的正义，企图建立自己的正义}\BibleRef{罗 10:3}。但正因为曾说：\Quote{年长的要服事年幼的}，所以很明显，不虔敬的被虔敬的所制服，骄傲的被谦卑的所制服。厄撒乌先出生，雅各伯后出生；但后出生的被优先于长子，那长子因贪吃而失去了\OriginalTerm{长子权}{birthright}。正如经上所写的：\BibleRef{创 25:30}他渴望扁豆羹，他兄弟对他说：你若愿意我把它给你，就把你的长子权给我。他更爱那按肉性所贪求的，胜过那因先出生而按灵性所赢得的；于是他放弃了长子权，好能喝那扁豆羹。但我们发现扁豆是埃及人的食物，因为那里盛产埃及。亚历山大的扁豆如此被重视，甚至传到我们这里，仿佛我们这里不长扁豆。因此由于贪求埃及食物，他失去了长子权。同样，犹太人百姓，关于他们曾说过：\Quote{心中转回埃及}\BibleRef{宗 7:39}。他们可以说是贪求扁豆，便失去了长子权。\par

6. \Quote{天主在欢呼声中上升}\BibleRef{咏 46:5}。就是祂，我们的天主，主基督，在欢呼声中上升；\Quote{上主在号角声中上升}。\Quote{上升}：到何处？除了我们所知的地方？犹太人没有跟随祂，甚至用眼睛也没有看见。因为当祂在十字架上被高举时，他们嘲笑祂；当祂升天时，他们没有看见祂。\Quote{天主在欢呼声中上升}。什么是欢呼？不就是无法用言语表达的喜乐赞叹吗？正如门徒们喜乐地赞叹，看见他们曾哀悼已死的那位进入天堂；真的，因那喜乐，言语不足以表达：只能以欢呼表达那无人能言说的事。也有号角的声音，天使的声音。因为经上说：\Quote{你要扬起你的声音像号角}。天使传报了主的升天：他们看见门徒们，他们的主上升，他们停留、赞叹、困惑、无言以对，但心中欢呼；那时号角的声音在天使清晰的话音中：\BibleQuote{你们加里肋亚人，你们为什么站着望天呢？这位就是耶稣}\BibleRef{宗 1:11}。好像他们不知道那就是同一耶稣。难道他们刚才没有看见祂在他们面前？没有听见祂与他们说话？不但如此，他们不仅看见祂显现的形象，还触摸了祂的肢体。难道他们自己那时不知道那就是同一耶稣？但他们因极度赞叹，因欢呼的喜乐，心灵仿佛被提，天使便说：\Quote{这位就是耶稣}。意思好像是：如果你们信祂，这就是那同一耶稣，祂被钉十字架，你们的脚曾绊倒；祂死了被埋葬，你们曾以为你们的希望失落了。看，这就是同一耶稣。祂已在你们之前上升，\BibleQuote{祂怎样往天上去，还要怎样来}\BibleRef{宗 1:11}。确实，祂的身体从你们的眼前移开了，但天主没有与你们的心分离：看见祂上升，相信祂不在，盼望祂回来；然而，借着祂隐秘的仁慈，感受祂的临在。因为那升天以便从你们眼前移开的那一位，曾应许你们说：\BibleQuote{看，我常与你们同在，直到世界的终结}\BibleRef{玛 28:20}。所以，宗徒正确地这样对我们说：\BibleQuote{主已经近了，应当一无挂虑}\BibleRef{斐 4:5}。基督坐在诸天之上；诸天遥远，但祂坐在那里，却近在咫尺……\par

7. \Quote{向我们的天主歌唱，歌唱}（\BibleRef{咏 46:6}）。他们——那些远离天主的人——曾嘲笑作为人的祂。\Quote{向我们的天主歌唱。}因为祂不但是人，而且是天主。祂是达味后裔的人（\BibleRef{罗 1:3}），是达味的主天主，按肉身来自犹太人。\Quote{他们的}（宗徒说）\Quote{是祖先，基督按肉身也是从他们来的}（\BibleRef{罗 9:5}）。如此，基督来自犹太人，但按肉身而言。然而这位按肉身来自犹太人的基督是谁？\Quote{祂是在万有之上，永远受赞美的天主。}天主在肉身之前，天主在肉身之内，天主带着肉身。不仅是肉身之前的天主，而且是创造肉身的土地之前的天主；不只是创造肉身的土地之前的天主，而且是首先创造的诸天之前的天主；是首先被造的日子的天主；是天使之前的天主；同一基督就是天主：因为\Quote{在起初已有圣言，圣言与天主同在，圣言就是天主}（\BibleRef{若 1:1}）。\par

8. \Quote{因为天主是全地的君王}（\BibleRef{咏 46:7}）。怎样？难道以前祂不是全地的天主吗？难道祂不是天地的主宰吗，因为万物确实是由祂造成的？谁能说祂不是他的天主呢？但并非所有人都承认祂为他们的天主；而只在祂被承认的地方，可以说，祂才是天主。\Quote{天主在犹大为人所知。}那时还未对科辣黑的子孙说：\Quote{啊，万民，你们要鼓掌。}因为那位在犹大为人所知的天主，是全地的君王：如今祂被众人承认，因为依撒意亚所说的应验了：\Quote{祂是拯救你的天主，祂必被称为全地的天主}（\BibleRef{依 54:5}）。\Quote{你们要明智地歌唱。}祂教训并警告我们要明智地歌唱，不要寻求耳朵的声音，而要寻求内心的光明。外邦人——你们从他们中被召叫成为基督徒——曾敬拜人手所造的神祇，为他们歌唱，但不是明智地。如果他们明智地歌唱，就不会敬拜石头。当一个有知觉的人向无知觉的石头歌唱时，他明智地歌唱了吗？但是现在，弟兄们，我们不是用眼睛看见我们所敬拜的，却正确地敬拜了。天主更受称赞，因为我们用眼睛看不见祂。如果我们用眼睛看见了祂，或许我们会轻视。因为即使是可见的基督，犹太人也轻视了；而不可见的主，外邦人却敬拜了。\par

9. \Quote{天主将为万民为王}（\BibleRef{咏 46:8}）。祂曾为一个国家为王，祂说：\Quote{将为万民为王}。当这话说出来时，天主只管治一个国家。这是一个预言，事情尚未显明。感谢天主，我们现在看到从前预言的事已经应验。天主在时间之前给我们写下了一个许诺，在时间圆满时祂偿还了我们。\Quote{天主将为万民为王}，这是一个许诺。\Quote{天主坐在祂的圣座上}。什么是祂的圣座？或许有人说：诸天，他理解得很好。因为基督已经上升（\BibleRef{宗 1:2}），正如我们所知，带着祂被钉在十字架上的身体，坐在圣父的右边；我们从那里期待祂来审判生者死者（\BibleRef{弟后 4:1}）。\Quote{天主坐在祂的圣座上}。诸天是祂的圣座。你也要成为祂的宝座吗？不要以为你不可能；要在你心中为祂预备一个地方。祂来了，并乐意坐下。同一基督确实就是\Quote{天主的德能，天主的智慧}（\BibleRef{格前 1:24}）；圣经论到智慧本身说什么？\Quote{义人的灵魂是智慧的宝座}（\BibleRef{智 7:27}）。所以，如果义人的灵魂是智慧的宝座，你的灵魂也当成为义人，你便成为智慧的王座。确实，弟兄们，所有生活良善、行为端正、在虔诚的爱德中往来的人，天主岂不住在他们内，并亲自命令吗？你的灵魂服从坐在它内的天主，它自己则命令肢体。因为你的灵魂命令你的肢体，使脚、手、眼、耳得以活动，它自己命令肢体如同仆人，但它自己却服侍坐在它内的主人。除非它不轻看服侍它的上主，否则它便不能好好管理它的下属。\par

10. \BibleQuote{万民的首领集合起来，归于亚巴郎的天主。}\BibleRef{咏 46:9}亚巴郎的天主，依撒格的天主，雅各伯的天主。\BibleRef{出 3:6}确实，天主说了这话，犹太人因而自夸，说：\BibleQuote{我们是亚巴郎的子孙。}\BibleRef{若 8:33}他们以祖先的名字自豪，带着他的肉身，却没有持守他的信德；由血统亲近他，在品行上堕落。但主对他们这些自夸的人说了什么？\BibleQuote{如果你们是亚巴郎的子孙，就该做亚巴郎所做的事。}\BibleRef{若 8:39}再者……\Quote{万民的首领：}万民的首领，不是一民的首领，而是万民的首领都\Quote{集合起来，归于亚巴郎的天主。}这些首领中也有一位百夫长，就是刚才你们在福音中听到的那位。因为他是一个百夫长，在人间享有尊荣和权柄，他是万民首领中的一位首领。基督到他那里去，他派朋友去迎接他，不，当基督真正经过他那里时，他派人去，请求主治好他病危的仆人。主愿意前往时，他派人传话给主：\BibleQuote{我不敢当主到我舍下来，只要你说一句话，我的仆人就会痊愈。因为我也是一个在权下的人，手下有兵。}\BibleRef{路 7:6}看他如何保持自己的身份！他先提到自己受别人管束，然后说别人受他管束。我受权管辖，我也掌权；我既在别人之下，又在别人之上。……仿佛他说：我既在权下，能命令那些在我之下的人，你既不受任何人管辖，难道不能命令你的受造物吗？因为万物都是藉着你造的，没有你，什么也没有造成。\Quote{你说一句话吧，}他说，\Quote{我的仆人就会痊愈。因为我实在不配你进我的屋。}……耶稣惊叹他的信德，斥责犹太人的不信。他们自以为是健康的，实则病势沉重，竟杀害了他们不认识的医生。因此，当他斥责并拒绝他们的骄傲时，说了什么？\BibleQuote{我告诉你们：将有许多人从东从西来，}不属于以色列家族的人；许多人将来到，他曾对他们说：\BibleQuote{万民啊，你们要鼓掌，}\BibleQuote{并且要与亚巴郎、依撒格、雅各伯一起坐席，在天国里。}亚巴郎虽然没有生他们出于自己的肉身，他们却要与他一起坐席，成为他的子孙。凭什么是他的子孙？不是出于他的肉身，而是追随他的信德。\BibleQuote{但本国的子民}即犹太\BibleQuote{将被驱逐到外面的黑暗中，在那里有哀号和切齿。}\BibleRef{玛 8:12}\par

11. 而那些属于亚巴郎的天主的人？\BibleQuote{因为地上有权势的神明大受高举。}那些是神明、天主子民、天主的葡萄园，正如经上所说：\BibleQuote{你们在我与我的葡萄园之间判断吧。}\BibleRef{依 5:3}他们却要进入外面的黑暗，不与亚巴郎、依撒格、雅各伯同席，不被集合到亚巴郎的天主面前。为什么？\BibleQuote{因为地上有权势的神明，}他们曾是地上的有权势的神明，倚靠于地。什么地？他们自己；因为每个人都是地。因为人对人说过：\BibleQuote{你原是土，还要归于土。}\BibleRef{创 3:19}但人应该信赖天主，从那里盼望帮助，而不是从自己。因为地不会为自己降雨，也不会为自己发光；但正如地从天期待雨和光，人也应从天主期待仁慈和真理。他们，这些\BibleQuote{地上有权势的神明，就大受高举，}即大显骄傲；他们自认为不需要医生，因此留在自己的病中，并因自己的病而堕入死亡。天然的枝条被折下来，好使谦卑的野橄榄枝被接上。\BibleRef{罗 11:17}因此，弟兄们，我们要持守谦卑、爱德和虔敬：既然我们蒙召，他们却显为无用，甚至因他们的榜样，我们当惧怕骄傲。\par

\SourceRecord{Translated by J.E. Tweed.；Nicene and Post-Nicene Fathers, First Series；Vol. 8.；Edited by Philip Schaff.；Buffalo, NY: Christian Literature Publishing Co.,；1888.；<http://www.newadvent.org/fathers/1801047.htm>.}

% structure: p0001 source=h1 target=section reason=page-title

\section{\WorkTitle{《圣咏第四十八篇讲疏》}}

1. 这首圣咏的标题是\BibleQuote{赞美之歌，属科辣黑子孙，安息日后第二日}。关于这事，愿主赐予你们如同穹苍之子般领受。因为在安息日后第二日，即我们称为主日的第一日之后的那一日，也称为第二工作日，穹苍被造成。\BibleRef{创 1:6}……然后，安息日后第二日我们不应理解为别的，而应理解为基督的教会：但基督的教会是在圣人们中，基督的教会是在那些名字被记录在天堂的人中，基督的教会是在那些不向此世的诱惑屈服的人中。因为他们配得上\Quote{穹苍}之名。因此，基督的教会，在那些强壮的人中，宗徒说：\BibleQuote{我们强壮者应担当软弱者的软弱}\BibleRef{罗 15:1}，被称为穹苍。这首圣咏就是歌唱这教会。让我们聆听、承认、联合、光荣、为王。因为那被称为穹苍的教会，在宗徒书信中也听到：\BibleQuote{真理的柱石和基础}\BibleRef{弟前 3:15}……\par

2. \BibleQuote{主伟大，应受赞美}\BibleRef{咏 47:1}……就是说，\BibleQuote{在我们天主的城内，在他的圣山上。}这是一座建在山上的城，不能隐藏；这是灯，不放在斗底下\BibleRef{玛 5:14}，为人所知，为人所宣扬。然而并非所有的人都在这城里，只有在那些\BibleQuote{主伟大，应受赞美}的人中。那么，那城是什么？让我们看看是否可能，既然经上说\BibleQuote{在我们天主的城内，在他的圣山上}，我们不应该查问这座山，在那里我们也能被垂听……那么，弟兄们，那山是什么？它需要极大的谨慎去查问，极大的热忱去探究，也需要劳力去占据和攀登。但如果它在任何部分的大地，我们该怎么办？我们要离乡背井，前往那山吗？不，当我们不在其中时，我们在外地。因为那是我们的城，如果我们是大王的肢体，他是同一城的头……因为有一块被人轻视的角石，犹太人绊跌在上面\BibleRef{罗 9:32}，从某座山非人手凿出，即来自犹太人的王国非人手所成，因为人性的运作没有与玛利亚结合，基督由她而生\BibleRef{玛 1:16}。但如果那石头，当犹太人绊跌其上时，仍留在那里，你们就没有可攀登之处。但发生了什么？达尼尔预言说什么？岂非那石头长大，成为一座大山？多大？以致充满全地\BibleRef{达 2:35}。因此，借着长大并充满全地，那山来到我们这里。为什么我们寻找那山如同不在，而不当作在眼前攀登它；好使主在我们中间\BibleQuote{伟大，应受赞美}？\par

3. 此外，……当他说了\BibleQuote{在我们天主的城内，在他的圣山上}之后，他又加上了什么？\BibleQuote{散布全地的喜乐，熙雍的山岭}\BibleRef{咏 47:2}。熙雍是一座山，为何说\Quote{山岭}？是否因为熙雍也属于那些从另一边来的，以便在角石上相遇，成为两堵墙，如同两座山，一座属于割损，另一座属于未割损；一座属于犹太人，另一座属于外邦人：不再敌对，虽然不同，因为来自不同的侧面，如今在角石上甚至不再不同。\BibleQuote{因为他是我们的和平，使两者合一}\BibleRef{弗 2:14}。同一角石\BibleQuote{匠人弃而不用的石头，反而成了屋角的基石}。这山在自己内联合了两座山；一座房屋，又是两座房屋；两座，因为来自不同侧面；一座，因为角石将两者联合在一起。也请听这话：\BibleQuote{熙雍的山岭：北方的边陲是大王的城。}……看外邦人；\BibleQuote{北方的边陲是大王的城}。北方常与熙雍相对：熙雍固然在南方，北方对着南方。那北方是谁，岂非那说\BibleQuote{我要坐在北方的边陲，我要与至高者相似}的那位？\BibleRef{依 14:13}魔鬼曾统治不敬神的人，占有敬拜偶像、崇拜魔鬼的万民；而世上无论何处所有的人类，因依附他，都成了北方。但既然那捆绑壮士的，夺去他的器具\BibleRef{玛 12:29}，并使它们成为自己的器具；从无信和魔鬼的迷信中得救的人，相信基督，被接合到那城，在角石上遇到了那来自割损的墙，而那原本是北方边陲的，成了大王的城。因此，在另一处经文中也说：\BibleQuote{从北方而来金黄色的云彩：全能者的光荣与尊贵是伟大的}\BibleRef{约 37:22}。因为医生的光荣是伟大的，当被绝望的病人康复时。\BibleQuote{从北方而来云彩}，不是黑云，不是暗云，不是阴沉的云，而是\BibleQuote{金黄色的}。这若不是借着基督的恩宠照亮，又如何能呢？看，\BibleQuote{北方的边陲是大王的城}。……\par

4. 让圣咏接着这样说：\Quote{天主必在它的宫殿中被认识。}现在在它的\Quote{宫殿中，是因为群山，因为两道墙，因为两个儿子。}\Quote{天主必在它的宫殿中被认识，}但祂推崇恩宠，因此祂加上\Quote{当祂收容她的时候。}因为那座城若非祂收容了她，那会是什么？它岂不立刻倒塌，除非它有那样的基础？因为\BibleQuote{没有人能奠立别的根基，那已奠立的根基就是耶稣基督。}\BibleRef{格前 3:11}所以，谁也不该夸耀自己的功德；但\BibleQuote{凡夸耀的，该在主内夸耀。}\BibleRef{格前 1:31}……主已收容了这座城，并在其中被认识，也就是说，祂的恩宠在那城中被认识：因为那座城所有的一切，凡是夸耀于主的，都不是来自自己。为此有话说：\BibleQuote{你有什么不是领受的呢？}\BibleRef{格前 4:7}\par

5. \BibleQuote{看哪，地上的众王聚集在一起。}\BibleRef{咏 47:3}看，那些北方之地，看他们如何前来，看他们如何说：\BibleQuote{来，我们上主的山去，祂必教导我们祂的道路，我们就在其中行走。}\BibleRef{依 2:3}\Quote{并合而为一。}向哪一个？不就是那\Quote{角石}吗？\BibleRef{弗 2:20}\BibleQuote{他们看见了，就惊奇。}\BibleRef{咏 47:4}在他们对基督的奇迹和荣耀惊奇之后，接着发生了什么？\BibleQuote{他们惊慌，他们摇动，}\BibleRef{咏 47:5}\Quote{战栗抓住了他们。}战栗何以抓住他们，不正是由于罪恶的自觉吗？那么，让他们奔跑，君王一个接一个；君王们，让他们承认这位君王。因此祂在别处说：\Quote{然而我已被祂立为君王，在祂神圣的熙雍山上。}……一位君王被听闻了，被立在熙雍，祂的产业被交到祂手中，直到地极。君王们应当害怕，免得失去王国，免得王国被夺走。如同可怜的\Person{黑落德}害怕，并为那婴孩而屠杀孩童。\BibleRef{玛 2:16}但因害怕失去王国，他配不上认识这位君王。但愿他也与贤士们一起朝拜这位君王：而不是因恶求王国而杀害无辜者，并罪恶地灭亡。因为就他而言，他毁灭了无辜者；但就基督而言，即使是一位婴孩，为祂而死的孩子们，祂却给他们加冕。因此，当话说：\Quote{然而我已被祂立为君王，在祂神圣的熙雍山上，}并且那立祂为王的，要将地极的产业赐给祂时，君王们应当害怕……因此也对他们说：\Quote{所以，你们君王应当明白；世上审判官，当受教。以敬畏事奉主，以战栗向祂欢乐。}他们做了什么？\Quote{有如同临产妇人的疼痛。}这\Quote{如同临产妇人的疼痛}是什么？岂不就是忏悔者的剧痛？看这痛苦和生产的同一概念：\Quote{因祢的敬畏}（\Person{依撒意亚}说）\Quote{我们怀孕，我们产生了救恩之神。}\BibleRef{依 26:17}这样，君王们从对基督的恐惧中怀孕，以致通过生产，他们因相信他们所惧怕的那一位而带来救恩。\Quote{有如同临产妇人的疼痛：}当你听到生产时，期待一个出生。旧人生产，但新人出生。\par

6. \BibleQuote{祢要以强风打破塔尔史士的船只。}\BibleRef{咏 47:6}简而言之，这是说：祢要倾覆万民的骄傲。但在这段历史中，哪里提到倾覆万民的骄傲？由于\Quote{塔尔史士的船只。}学者们曾探寻塔尔史士城，即这个名字指代的是哪座城：有人似乎认为基里基雅被称为塔尔史士，因为其首府称为塔尔索。宗徒\Person{保禄}就是这座城的人，他生于基里基雅的塔尔索。\BibleRef{宗 21:39}但有些人理解为迦太基，或许它有时如此称呼，或在某种语言中如此表示。因为在\Person{依撒意亚}先知书中如此写道：\Quote{哀号吧，迦太基的船只。}但在\Person{厄则克尔}书中，有些译者将这个词译为迦太基，有些译为塔尔史士：从这种差异可以理解，被称为迦太基的同一地方也称为塔尔索。但显然，迦太基在其统治初期因船只而繁荣，如此繁荣，以致在各国中，他们在贸易和航海方面卓越超群。因为当\Person{狄多}逃离她的兄弟，逃到非洲地区，在那里建造了迦太基时，她曾为航行带上了在她国家为商业准备的船只，当地的首领同意了；当迦太基建成后，同样的船只也没有停止贸易。因此那城变得过于骄傲，以致恰好可以通过它的船只理解万民的骄傲，他们信赖不确定的事物，如同信赖风的呼吸。现在，谁也不该信赖满帆，以及今生如大海般的表面美好状态。愿我们的根基在熙雍：我们应当在那里被建立，而不是被\BibleQuote{每一种教义之风吹来吹去。}\BibleRef{弗 4:14}因此，凡因今生不确定之事而自高自大的人，愿他们被倾覆，愿万民的一切骄傲都服从于基督，祂将\Quote{以强风打破塔尔史士所有的船只：}不是任何一座城的，而是\Quote{塔尔史士}的。\Quote{以强风}是什么意思？以极度的恐惧。因为如此，一切骄傲都惧怕那将要审判的那一位，以致谦卑地信靠祂，以免在祂被高举时自己害怕。\par

7. \Quote{我们听了，也看见了} \BibleRef{咏 47:7}。有福的教会！你一时听了，一时也看见了。她听在预许里，看见在实现中；听在预言里，看见在福音中。因为现今所实现的一切，从前都已预言了。那么，请抬起你的眼睛，遍观全世界；请看现今，他的\Quote{产业直到地极}已实现；请看现今，所说的\Quote{诸王都要向祂俯伏，万国都要事奉祂}已实现；请看现今，所说的\Quote{天主，请祢高居于诸天之上，愿祢的荣耀充满大地}已实现。请看祂，手足被钉子穿透，骨悬于木上被数点，衣被拈阄：\BibleRef{玛 27:35}，请看他们曾看到悬挂的那一位如今为王；请看曾在地上被轻视的那一位如今坐在天上：\BibleRef{玛 26:64}，请看因此实现了\Quote{大地四极必想起上主，而转向祂，万邦万族都要在祂面前叩拜}。看见这一切，便欢欣呼喊道：\Quote{我们听了，也看见了}。教会本身从外邦人中这样被召出来，是公正的……那些没有领受先知的人，倒先听了并理解了先知；那些先没有听的人，后来听到便惊讶。那些被派遣去的人反而留在后面，他们带着经卷却不明白真理；持有盟约的法版，却不继承产业。但我们……\Quote{我们听了，也看见了}。你在哪里听？在哪里看？\Quote{在万军之主的城里，在我们天主的城里。天主永远奠定了它}。愿异端者不要分派别而侮辱人，不要自高，他们常说：\Quote{看，基督在这里，或看，在那里}。\BibleRef{玛 24:23} 凡说\Quote{看，基督在这里，或看，在那里}的，是引人分派。天主所预许的是合一。君王们会聚在一处，而非因分裂而分散。但那个曾掌握世界的城邑，将来会倾覆吗？绝无此事！\Quote{天主永远奠定了它}。既然天主永远奠定了它，你为何害怕穹苍会坠落呢？\par

8. \Quote{我们在祢百姓当中领受了祢的仁慈，} \BibleRef{咏 47:8}。谁领受了，在哪里领受了？难道不是祢的百姓领受了祢的仁慈。如果祢的百姓领受了祢的仁慈，那又怎么是\Quote{在祢百姓当中}呢？好像领受的是一部分人，他们是在另一群人当中领受的。这是一个伟大的奥秘，但也是众所周知的。因此，从这些诗句中提取和引出你们所知道的，将不是更粗鲁，而是更甜美。如今，所有携带天主圣事的人都算作天主的百姓，但并非所有人都属于祂的仁慈。所有领受基督圣洗圣事的人都称为基督徒，但并非所有人都活得与那圣事相称。有一些人，宗徒论及他们说：\Quote{具有虔诚的外表，却否认其能力。} \BibleRef{弟后 3:5} 然而，因着这虔诚的外表，他们被列于天主的百姓中。如同禾场，在谷粒被打净之前，不仅有麦子，还有糠秕。但它也会属于谷仓吗？那么，在邪恶的百姓当中，有一个善良的百姓，他们领受了天主的仁慈。活得与天主的仁慈相称的人，听见并持守且实行宗徒所说的：\Quote{我们恳求你们，不要白白领受天主的恩宠。} \BibleRef{格后 6:1} 因此，谁若不白白领受天主的恩宠，他就不仅领受了圣事，也同时领受了天主的仁慈。……所以，那些有圣事却没有善行的人，既被说成属于天主，又被说成不属于天主；既被说成是祂的，又被说成是外人：是祂的，因着祂的圣事；是外人，因着他们自己的恶习。同样，也有陌生的女子：\BibleRef{歌 2:2}；她们是女子，因为具有虔诚的外表；陌生，因为失去了美德。愿百合花在那里；让它领受天主的仁慈：紧握好花的根，不要对从天降下的柔雨忘恩负义。愿荆棘忘恩负义，让他们在阵雨中生长：他们是为火而生，不是为谷仓。在祢没有领受祢仁慈的百姓当中，我们领受了祢的仁慈。因为\Quote{祂来到自己的地方，自己的人却没有接受祂，}然而，在他们中间，\Quote{凡接受祂的，祂就赐给他们权能，成为天主的儿女。} \BibleRef{若 1:11}……\par

9. 因为当他说了\Quote{我们在祢百姓当中领受了祢的仁慈，}他表明存在一个没有领受天主仁慈的百姓，在他们中间有一些人领受了天主的仁慈：然后，免得人们以为人数如此之少，几乎等于零，祂如何在接下来的话语中安慰他们呢？\Quote{天主，依照祢的名，祢的赞美也达于地极} \BibleRef{咏 47:9}。这是什么呢？……也就是说，正如祢被认识于普世，祢也同样被赞美于普世，并且现在不缺在普世赞美祢的人。但他们赞美祢，是那些生活得好的人。因为\Quote{天主，依照祢的名，祢的赞美也是如此，}不是在一部分地区，而是\Quote{达于地极。} \Quote{祢的右手充满正义。}\par

10. \Quote{让熙雍山欢乐，犹大的女儿们喜乐，因为祢的审判，上主} \BibleRef{咏 47:10}。哦，熙雍山，哦，犹大的女儿们，你们现在在莠子中劳作，在糠秕中劳作，在荆棘中劳作；却要因天主的审判而喜乐。天主的审判不错。你们要分别生活，虽然你们并非生而分别；从你们的口和心中发出声音并非徒然：\Quote{不要从罪人中夺去我的灵魂，也不要从流人血者那里夺去我的性命。} 他以这样的技艺簸扬，手中拿着簸箕，使得没有一粒麦子落入预备焚烧的糠秕堆中，也没有一根糠秕的芒刺落入要存入库中的堆子里。\BibleRef{玛 3:12} 喜乐吧，哦，犹大的女儿们，因为天主的审判不错，且不要过早判断。你们负责收集，他负责分离。但不要以为\Quote{犹大的女儿们}就是犹太人。犹大意为宣认；所有宣认的儿子都是犹大的儿子。因为\BibleQuote{救恩出自犹太人} \BibleRef{若 4:22}，无非是说基督出自犹太人。宗徒也说：\BibleQuote{外表上是犹太人的，不是犹太人；外表在肉身上的割损，也不是割损。而是内在是犹太人的，才是犹太人；割损是内心的，在于神，不在于文字；这样的人称赞不是来自人，而是来自天主。} \BibleRef{罗 2:28} 你们要作这样的犹太人；以内心的割损为荣，即使你们没有肉身的割损。让犹大的女儿们因祢的审判而喜乐，上主。\par

11. \Quote{你们要环绕熙雍，拥抱她} \BibleRef{咏 47:11}。要对那些生活在恶中的人说，在他们中间有人民，即得到天主怜悯的人民。在你们中间有生活善良的人民：\Quote{环绕熙雍。} 但要怎样？\Quote{拥抱她。} 不是以绊脚石，而是以爱围绕她；这样你们可以效法在你们中间生活善良的人，并通过效法他们，与基督结合，成为基督的肢体，他们是基督的肢体。\Quote{环绕熙雍，绕行她；在其中的塔楼上宣讲。} 在城防的高处，传扬她的赞美。\par

12. \Quote{将你们的心放在她的能力上} \BibleRef{咏 47:12}。不是叫你们有敬虔的外貌，却否认了它的能力 \BibleRef{弟后 3:5}，而是\Quote{将你们的心放在她的能力上。在其中的塔楼上宣讲。} 这城的能力是什么？谁要理解这城的能力，就让他理解爱的力量。这是一种无法征服的德行。爱火不能被世界的波涛、不能诱惑的溪流所熄灭。正如经上所载：\BibleQuote{爱情猛如死亡} \BibleRef{歌 8:6}。因为死亡来临时无法抵抗；无论你用尽什么技巧、什么药物迎对，凡生而可朽者都无法避免死亡的强力；同样，面对爱的强力，世界无能为力。因为从相反的角度作了死亡的比喻；死亡最具强力以夺取，爱最具强力以拯救。许多人因爱死于世界，而活于天主；这爱情点燃了殉道者——不是假装的，不是自高自大而膨胀的，不是如经上所写\BibleQuote{我若舍身被焚，却没有爱德，于我无益} \BibleRef{格前 13:3}——而是那些被基督之爱和真理真正引导到这一激情的人；对他们来说，折磨者的诱惑算得了什么？他们哭泣的朋友的眼神比仇敌的迫害更有暴力吗？因为有多少人被自己的子女抓住，以免他们受苦？有多少人的妻子跪倒在他们面前，以免她们成为寡妇？有多少人被父母禁止赴死；正如我们在真福佩尔佩图阿受难中知道和读到的！这一切都发生了；但无论眼泪多么大，多么有力地流淌，它们何时熄灭了爱情的炽热？这就是熙雍的能力，正如对别处所说：\BibleQuote{愿和平在你城垣之内，繁荣在你宫殿之中。}\par

13. 这里我们把什么理解为\Quote{将你们的心放在她的能力上，并分配她的房屋}？即，分开房屋与房屋。不要混淆。因为有一种房屋有敬虔的外貌，却没有敬虔；但有一种房屋既有外貌又有敬虔。你们要分配，不要混淆。但当你们\Quote{将你们的心放在她的能力上}，即因爱成为属神的人时，你们就分配而不混淆。那时你们不会作轻率的判断，那时你们会看到，只要我们在这一禾场上，恶者不能伤害善者。\Quote{分配她的房屋。} 也可有另一种理解：两座房屋，一出于割损，一出于未受割损，宗徒们奉命分配。因为扫禄被召叫，成为宗徒保禄，与同伴宗徒在合一中同意，他如此决定：他们应去割损者那里，他应去未受割损者那里。通过他们宗徒职务的这一分施，他们分配了大王城市的房屋；并在角石中相遇，在分施中分割了福音，在爱中联合了福音。确实这更应被理解；因为接下来的经文表明这里是对宣讲者说的：\Quote{分配她的房屋，以便你们传报给后来的世代。} 即使直到我们这些将后来的人，他们福音的分施也应传给我们。因为他们不只为他们在地上同住的人劳苦；主也不只为那些祂复活后屈尊显示给他们的宗徒，也为我们。因为祂对他们说话，当祂说\BibleQuote{看！我天天同你们在一起，直到今世的终结} \BibleRef{玛 28:20}时，祂也预示了我们。难道他们那时要在这里直到今世的终结吗？祂也说：\BibleQuote{我不但为他们祈求，也为那些因他们的话而信从我的人祈求} \BibleRef{若 17:20}。因此祂思念了我们，因为祂为了我们而受难。所以圣经说：\Quote{以便你们传报给后来的世代。}\par

14. 说什么呢？\BibleQuote{因为这是天主，是我们的天主。}\BibleRef{咏 47:13} 大地被看见了，大地的创造者却没有被看见；血肉被握住了，血肉中的天主却没有被承认。因为那血肉是由那些从之取了同一血肉的人所持有，因为童贞玛利亚出自亚巴郎的后裔。他们停留在血肉之上，却不懂得天主性。啊，宗徒们，啊，强大的城邑，你们要在城楼上宣讲，说：\BibleQuote{这是天主，是我们的天主。}同样，当祂被轻视时，当祂作为一块绊脚的石头放在绊跌者脚下时，为要谦逊忏悔者的心；同样，这话仍然真实：\BibleQuote{这是天主，是我们的天主。}当然，祂曾被看见，正如经上所说：\BibleQuote{此后祂在地上显现，与人同住。}\BibleRef{巴 3:37} \BibleQuote{这是天主，是我们的天主。}祂也是人，有谁能认识祂呢？\BibleQuote{这是天主，是我们的天主。}但或许如同假神一样暂时存在。因为它们虽被称为神，却不能真正是神，只是一时被如此称呼。因为先知说了什么，或主警告要对他们说什么？\BibleQuote{你们要这样对他们说：那些没有创造天地的神，必从地上、从天下消灭。}\BibleRef{耶 10:11} 祂不是这样的神：因为我们的天主超越众神。超越哪些众神？\BibleQuote{因为万民的神都是偶像，惟独主造了诸天。}同一的天主就是我们的天主。\BibleQuote{这是天主，是我们的天主。}到什么时候？\BibleQuote{直到永远：祂要永远统治我们。}如果祂是我们的天主，祂也是我们的君王。祂保护我们，作为我们的天主，免得我们死亡；祂统治我们，作为我们的君王，免得我们跌倒。但祂统治我们，并不击碎我们；因为凡祂不统治的，祂就击碎。祂说：\BibleQuote{你要用铁杖统治他们，如同打碎陶器一样粉碎他们。}但有一些祂不统治的；这些人祂并不宽恕，如同打碎陶器一样粉碎他们。因此，让我们愿意被祂统治并获得拯救，\BibleQuote{因为祂是我们的天主，直到永远，祂要永远统治我们。}\par

\SourceRecord{Translated by J.E. Tweed.；Nicene and Post-Nicene Fathers, First Series；Vol. 8.；Edited by Philip Schaff.；Buffalo, NY: Christian Literature Publishing Co.,；1888.；<http://www.newadvent.org/fathers/1801048.htm>.}

% structure: p0001 source=h1 target=section reason=page-title

\section{《圣咏第四十九篇释义》}

第一部分。\par

1. ……\Quote{万民，请听此言}\BibleRef{咏 48:1}。那么，不只是你们在这里的人。因为我们声音的能力怎能如此呼喊，以致万民都能听见呢？因为我们的主耶稣基督已借着宗徒们，并借着祂所派遣的众多语言宣讲了它；我们看到这篇圣咏，以前只在犹太人的会堂中、在一个民族中被诵念，如今却在整个世界、在所有教会中被诵念；并且这里所说的话应验了：\Quote{万民，请听此言。}……你们是属于谁的：\Quote{普世居民，请侧耳细听。}祂似乎第二次重复了这句话，以免先前只说\Quote{听}还不够。祂说：\Quote{我说，你们要听，要侧耳细听}，就是说，不要草率地听。什么是\Quote{侧耳细听}？这就是主所说的：\Quote{有耳听的，听罢！}\BibleRef{玛 11:15} 因为凡在祂面前的人都有耳朵，当祂说\Quote{有耳听的，听罢}时，祂要求的岂不是心灵的耳朵吗？这篇圣咏也触动了同样的耳朵。\Quote{普世居民，请侧耳细听。}也许这里有一些区别。我们固然不应局限自己的视野，但解释这种感官理解也无妨。也许\Quote{万民}和\Quote{普世居民}之间有些区别。因为也许祂要我们理解\Quote{居住}这个词有更深的意义，即把\Quote{万民}理解为所有恶人，而\Quote{普世居民}理解为所有义人。因为那不被束缚的人才是居住者；但被占据的人是居住的地方，而不是居住者。正如一个人是他财产的主人，但主人是不被贪婪之网束缚的人；而被贪婪束缚的人是被占有的，而不是占有者。……\par

2. 因此，让不敬神的人也听：\Quote{万民，请听此言。}也让义人听，他们并非徒然听闻，而是统治世界而非被世界统治：\Quote{普世居民，请侧耳细听。}\par

3. 祂又说：\Quote{无论贵贱，不分贫富}\BibleRef{咏 48:2}。\Quote{土生之子}指罪人；\Quote{人子}指信友和义人。你看，这个区别是被遵循的。谁是\Quote{土生之子}？大地之子。谁是大地之子？那些渴望地上产业的人。谁是\Quote{人子}？那些属于人子的人。我们之前已向你们的圣德解释了这一区别，并断定：亚当是人，但不是人子；基督是人子，但也是天主。因为凡属于亚当的，就是\Quote{土生之子}；凡属于基督的，就是\Quote{人子}。然而，让所有人都听，我不向任何人隐瞒我的言论。如果一个人是\Quote{土生之子}，让他因审判而听；另一个人是\Quote{人子}，让他因天国而听。\Quote{无论贵贱，不分贫富。}这里又重复了同样的话。\Quote{贵}指\Quote{土生之子}；但\Quote{贫}指\Quote{人子}。通过\Quote{贵}理解骄傲的人，通过\Quote{贫}理解谦卑的人。……祂在另一篇圣咏中说：\Quote{贫乏的人必将食而饱饫。}祂如何称赞穷人？\Quote{贫乏的人必将食而饱饫。}他们吃什么？就是信友所知道的那食物。他们如何饱足？通过效法他们主的苦难，并且并非无故获得他们的赏报。\Quote{贫乏的人必将食而饱饫，寻求上主的人必将赞美祂。}富人又如何？他们也吃。但他们如何吃？\Quote{地上所有富庶的人都吃了，并朝拜了。}祂没有说：\Quote{吃了，并饱足了}；而是说：\Quote{吃了，并朝拜了。}他们确实朝拜天主，但不会显示兄弟般的仁爱。这些人吃并朝拜；那些人吃并饱足：然而两者都吃。对于吃的人，他所吃的是被要求的：不要让分施者禁止他吃，但要劝诫他敬畏那要求他交账的主。那么，让罪人和义人、万民和居住在世界的人、\Quote{土生之子}和\Quote{人子}、\Quote{无论贵贱，不分贫富}，都一同听这些话：不分彼此，不被分开。那是收割时的事，簸扬者的手将成就这事。\BibleRef{玛 3:12} 现在，让富人和穷人一同听，让山羊和绵羊在同一牧场吃草，直到那一位来临，将绵羊置于右方，山羊置于左方。\BibleRef{玛 25:32} 让他们都一同听老师的话，免得彼此分离而听见审判者的声音。\par

4. 现在他们要听什么呢？\Quote{我的口要宣讲智慧，我的心所默思的是通达}\BibleRef{咏 48:3}。这个重复或许是为了避免：如果祂只说\Quote{我的口}，你会以为那对你说话的人只有嘴唇上的通达。因为许多人有嘴唇上的通达，心里却没有，正如经上所说：\Quote{这民族用嘴唇尊敬我，他们的心却远离我。}\BibleRef{依 29:13} 那么，那对你说话的人说什么呢？当祂说了\Quote{我的口要宣讲智慧}之后，为使你知道那从口里流出的，是从心底涌出的，祂又加上：\Quote{我的心所默思的是通达。}\par

5. \BibleQuote{我愿意侧耳听比喻，我要在竖琴上显示我的命题} \BibleRef{咏 48:4}……为何说\Quote{比喻}？因为\BibleQuote{我们现在是借着镜子观看，模糊不清，} \BibleRef{格前 13:12}如宗徒所说；\BibleQuote{我们住在肉身内，就是远离主} \BibleRef{格后 5:6}。因为我们的看见还不是面对面，那里不再有比喻，不再有谜语和比拟。我们现在所理解的，都是透过谜语看见。谜语是难以理解的暗喻。无论人如何修养他的心，致力于领悟奥秘，只要我们借着这朽坏的血肉观看，我们就只是部分看见……但是，正如祂在受审时，被那些相信祂的人和那些钉死祂的人看见；同样，当祂开始作审判者时，祂将被那些祂要定罪的人和那些祂要加冕的人看见。但是，祂所应许给爱祂的人的那种神性显现，当祂说：\BibleQuote{爱我的，必蒙我父爱他；爱我的，必遵守我的诫命；我也要爱他，并要向他显示我自己：} \BibleRef{若 14:21}这种显现，恶人是看不见的。这种显现是以某种方式亲切的：祂为祂自己的人保留，不会显示给恶人。这显现本身是怎样的？基督是怎样的？祂与圣父同等。基督是怎样的？\BibleQuote{在起初已有圣言，圣言与天主同在，圣言就是天主。} \BibleRef{若 1:1}为了这显现，我们现在叹息，在我们旅居此处时呻吟；最终我们将被带回这显现，这显现我们现在只能模糊看见。因此，如果我们现在模糊看见，就让我们\Quote{侧耳听比喻}，然后让我们\Quote{在竖琴上显示我们的命题}：让我们听我们所说的，做我们所吩咐的。\par

6. 祂说了什么？\BibleQuote{我为何要在邪恶的日子害怕？我脚跟的邪恶将围绕我} \BibleRef{咏 48:5}。祂开始有些隐晦。因此，如果他的脚跟的邪恶围绕他，他更应当害怕。不，祂说，那没有能力逃脱的人不要害怕。例如，害怕死亡的人，他怎样才能逃脱死亡？让他告诉我，他是如何逃脱亚当所欠的，这个从亚当出生的人。但他要想到，他从亚当出生，并且跟随了基督，应当偿还亚当所欠的，并获得基督所应许的。因此，害怕死亡的人根本无法逃脱；但害怕恶人所听到的诅咒\BibleQuote{你们进入永火里去吧，} \BibleRef{玛 25:41}的人，却有一条出路。那么，让他不要害怕。因为他为什么要害怕？他脚跟的邪恶会围绕他吗？如果他避开\Quote{他脚跟的邪恶}，并行走在天主的道上，他就不会来到邪恶的日子：邪恶的日子，最后的日子，对他不会是邪恶的。现在，当他们活着的时候，让他们警惕自己，让他们除去脚跟的邪恶：让他们行走在那条道上，行走在祂自己所说的道上：\BibleQuote{我是道路、真理、生命} \BibleRef{若 14:6}，并且让他们不要在邪恶的日子害怕，因为那成为\Quote{道路}的祂赐给他们安全。因此，让他们避开他们脚跟的邪恶。人用脚跟滑跌。请各位留意。天主对蛇说了什么？\Quote{她要踏碎你的头，你要伤害她的脚跟。}魔鬼图谋你的脚跟，为了在你滑跌时推倒你。他图谋你的脚跟，你要图谋他的头。他的头是什么？一个邪恶建议的开端。当他开始暗示邪恶思想时，你要在快感产生之前就把他赶走，从而避免同意；这样你就会避开他的头，他也不会抓住你的脚跟。但祂为何对厄娃说这话？因为人通过肉身滑跌。我们的肉身是我们内在的厄娃。\BibleQuote{爱自己妻子的，}他说，\BibleQuote{就是爱自己。}\BibleRef{弗 5:28}他说\Quote{自己}是什么意思？他继续说：\BibleQuote{从来没有人恨过自己的肉身}\BibleRef{弗 5:28}。因为魔鬼要通过肉身使我们滑跌，就像他通过厄娃使那个人亚当滑跌一样；厄娃被吩咐要图谋魔鬼的头，因为魔鬼图谋她的脚跟。\Quote{如果因此我们脚跟的邪恶会围绕我们，我们为何在邪恶的日子害怕，}既然我们皈依基督，就能不行邪恶；那么就没有什么会围绕我们，我们在最后的日子将喜乐而不忧愁？\par

7. 但是，那些\Quote{脚跟的邪恶必围绕}的人是谁？\BibleQuote{他们依靠自己的力量，并以财富丰裕而夸耀} \BibleRef{咏 48:6}。因此，我要避免这样的罪，\Quote{我脚跟的邪恶}就永不会围绕我。什么是避免这样的罪？我们不要依靠自己的力量，不要以自己财富的丰裕而夸耀，而要夸耀那应许给我们——谦卑的人得高举，并警告高傲的人受惩罚——的天主；这样，我们脚跟的邪恶就永不会围绕我们。\par

8. 有些人倚靠朋友，有些人倚靠德行，有些人倚靠财富。这就是不倚靠天主的世人的僭妄。他提到了德行，提到了财富，提到了朋友。\BibleQuote{弟兄不赎，难道人能赎吗？} \BibleRef{咏 48:7} 难道你指望人能赎你脱离将来的忿怒吗？如果弟兄不赎你，难道人能赎你吗？这弟兄是谁？若祂不赎你，就没有人能赎你。就是祂在复活后所说的：\BibleQuote{去，告诉我的弟兄们。} \BibleRef{玛 28:10} 祂愿意成为我们的弟兄：当我们对天主说\Quote{我们的天父}时，这就在我们身上显明了。因为那对天主说\Quote{我们的天父}的，就是对基督说\Quote{弟兄}。因此，让那以天主为父、以基督为弟兄的人，在邪恶的日子不要害怕。\BibleQuote{因为他脚后跟的不义必不围绕他，} 因为他既不倚靠自己的德行，也不夸耀财富丰裕，也不自夸有权势的朋友。让他倚靠那位为他而死、使他不再永远死去的主：祂为了他的缘故谦卑自下，好使他被高举；祂寻找那不虔敬的他，好使祂被这位信者寻求。因此，如果祂不赎，难道人能赎吗？如果人子不赎，难道有人能赎吗？如果基督不赎，难道亚当能赎吗？\BibleQuote{弟兄不赎，难道人能赎吗？}\par

9. \BibleQuote{他不将他的\OriginalTerm{赎罪}{propitiation}献给天主，也不付出\OriginalTerm{救赎}{redemption}他灵魂的代价} \BibleRef{咏 48:8}。那倚靠自己的德行、夸耀财富丰裕的人，就是\BibleQuote{不将他的赎罪献给天主}——即他用以在天主前为罪求得宽赦的满足——也不\BibleQuote{付出救赎他灵魂的代价}，他倚靠自己的德行、朋友和财富。但那些付出救赎自己灵魂代价的人是谁呢？就是主对他们说\BibleQuote{你们要藉着\OriginalTerm{不义的钱财}{Mammon of unrighteousness}结交朋友，好使他们在你们被接纳到永恒的居所时迎接你们}的那些人。\BibleRef{路 16:9} 那些不断行施舍的人，就是在付出救赎自己灵魂的代价。因此，宗徒藉弟茂德所嘱咐的那些人，他不愿他们骄傲，免得他们夸耀财富丰裕。最后，他不愿他们手中的财富闲置变旧，而要用来谋取救赎灵魂的代价。因为他说：\BibleQuote{你要嘱咐今世富足的人，不要心高气傲，也不要信赖无定的财富，而要信赖那厚赐百物给我们享用的永生天主。} \BibleRef{弟前 6:17} 好像他们曾说：\Quote{那么我们用我们的财富做什么呢？}他接着说道：\BibleQuote{要他们行善，富裕于善工，甘心施舍，乐意与人通财，} \BibleRef{弟前 6:18} 这样他们就不会失去那财富。我们如何知道？请听下文：\BibleQuote{为自己积蓄良好的根基，以备将来，好使他们把握住真正的生命。} \BibleRef{弟前 6:19} 这样，他们便付出救赎自己灵魂的代价。主也如此劝告：\BibleQuote{要为你们自己制作不朽坏的钱囊，在天上积攒永不匮乏的宝藏，那里没有贼偷，也没有虫蛀。} \BibleRef{路 12:33} 天主不愿你失去财富，而是给了你建议，让你改变它的位置。让你的爱德理解。假设你的朋友刚进你屋子，发现你把谷物堆放在潮湿地方，而他知道谷物易于霉烂，你或许不知道，他就会这样劝你说：\Quote{兄弟，你正在丧失辛苦积聚的财物，你把它放在潮湿地方，几天后这谷物就会霉烂。} 那我该做什么呢，兄弟？\Quote{把它移到高处。} 你听从朋友建议，把谷物从低处移到高处，难道不听从基督的命令，把你的宝藏从地上提升到天上吗？在那里，不是你所储存的会归还给你，而是你储存地上的，会得到天上的；你储存必朽的，会得到不朽的；这样，当你借给基督、在地上从你手中只得到小额贷款时，祂会在天上回报你巨大的赏报吗？然而，那些被\BibleQuote{脚后跟的不义围绕}的人，因为他们倚靠自己的德行，夸耀财富丰裕，倚靠无能相助的世人朋友，他们\BibleQuote{必不将赎罪和救赎灵魂的代价献给天主。}\par

10. 这样的人，他（圣咏作者）说了什么呢？\BibleQuote{是的，他永远劳苦，却要继续活到尽头} \BibleRef{咏 48:9}。他的劳苦将无止境，他的生命将有尽头。为什么他说\BibleQuote{他要活到尽头}？因为这样的人认为生命无非是每日的享乐。所以，当我们这时代的许多贫穷困苦、不稳定、不仰望天主对他们的劳苦所应许的人，看见富人在每日的筵席中，在金银的辉煌和闪烁中，他们会说什么呢？\Quote{只有这些人才是真正的人；他们真正活着！}这是一个说法，但愿不再被说：我们既警告你们，也仍然要警告你们，使说这话的人比如果不警告你们时要少。因为我们不敢以为我们这样说，就能使它不被说，而是要使它被更少的人说：因为它甚至要被说到世界终结。他说\BibleQuote{他活着}还是不够；他加上说，他雷鸣般地说，你以为只有他活着吗？让他活着吧！他的生命将要终结：因为他没有交付他灵魂的赎价，他的生命将要结束，他的劳苦将不会结束。\BibleQuote{他永远劳苦，却要继续活到尽头。}他怎能活到尽头呢？正如那个\BibleQuote{身穿紫袍和细麻衣，天天奢华宴乐} \BibleRef{路 16:19}的人所活的，他骄傲自大，轻视躺在他门前的满身疮痍的人，那人的疮被狗舔，并且渴望从富人桌子上掉下来的碎屑。那些财富对他有什么益处呢？两人的处境交换了：一个从富人的门口被带到亚巴郎的怀抱，另一个从他的丰盛筵席被投入火中；一个在平安中，另一个在焚烧；一个饱足，另一个口渴；一个劳苦到了尽头，却永远活着；另一个活到了尽头，却永远劳苦。那个在地狱痛苦中的富人，请求从拉匝禄的指尖滴一滴水在他的舌头上，说\BibleQuote{因为我在这火焰里很痛苦} \BibleRef{路 16:24}，却未得准许，这对他有什么益处呢？一个人渴望指尖的一滴水，正如另一个人渴望富人桌上的碎屑；但一个人的劳苦结束了，另一个人的生命结束了：这个的劳苦永远，那个的生命永远。我们或许在此尘世劳苦，但我们的生命不在这里：将来也不会被置于那样的境地，因为我们的生命将是基督直到永远；而那些\Quote{愿意}在此拥有他们生命的人，将永远劳苦，却活到尽头。\par

11. \BibleQuote{因为他必不见死亡，虽然他必看见智慧人死去} \BibleRef{咏 48:10}。那永远劳苦却活到尽头的人，\BibleQuote{必不见死亡，虽然他必看见智慧人死去。}这是什么意思呢？他必不明白死亡是什么，每当他看见智慧人死去。因为他对自己说：\Quote{这个人，尽管他聪明、与智慧同住、虔诚地敬拜天主，岂不是也死了吗？因此我要趁活着时享乐；因为如果那些在其他方面有智慧的人能有所作为，他们就不会死了。}正如犹太人看见基督挂在十字架上就藐视祂，说：\BibleQuote{如果这人是天主子，祂就该从十字架上下来} \BibleRef{玛 27:40}，不明白死亡是什么。如果他们明白了死亡是什么；如果他们明白了，我说。祂死了一时，好永远复活；他们活了一时，好永远死去。但因他们看见祂死去，他们就没有看见死亡，也就是说，他们没有理解什么是真正的死亡。甚至在\WorkTitle{智慧篇}中，他们说什么？\BibleQuote{我们用最耻辱的死刑定罪祂，因为按祂自己的话，祂将受尊重；}因为如果祂真是天主子，祂会从敌人的手中解救祂；如果祂真是祂的儿子，祂不会容许祂的儿子死去。但当他们看见自己在十字架上侮辱祂，而祂不从十字架上下来，他们就说，祂不过是一个人。话是这样说的：而且祂当然能从十字架上下来，祂能再次从坟墓中复活：但祂教导我们忍受那些侮辱我们的人；祂教导我们忍耐人的口舌，现在饮这杯苦杯，日后领受永久的救恩……\par

12. \Quote{鲁莽者和愚昧者将一同灭亡。}谁是\Quote{鲁莽者}？那不为未来打算的人。谁是\Quote{愚昧者}？那不明白自己处境邪恶的人。但你要明白你现在处于何等邪恶的处境，并要打算使你将来处于良好处境。藉着明白你现在的邪恶处境，你将不成为愚昧者；藉着为未来打算，你将不成为鲁莽者。那为自己打算的是谁？就是那个仆人，主人给他应花费的财物，后来对他说：\Quote{你不能作我的管家，把你管家的账目交代清楚}；他回答说：\Quote{我该做什么？我不能挖地，乞讨我又羞耻}（\BibleRef{路 16:1}），然而他仍用主人的财物为自己结交了朋友，好使他在被撤去管家职务时，能被这些朋友接待。他欺骗了主人，为使自己有朋友接待他；但你不要怕自己会欺骗，主亲自鼓励你这样做。祂亲自对你说：\Quote{用不义的钱财为自己结交朋友}（\BibleRef{路 16:9}）。也许你所拥有的，是从不义中得来的；或者这一事实本身就是不义：你拥有而别人没有，你富足而别人匮乏。用这不义的钱财，即不义之人称为财富的，为自己结交朋友，你便是有智慧的人；你是在为自己赚得利益，而不是在欺骗。因为现在你似乎失去了它。如果你把它存入宝库，你会失去它吗？因为孩子们，我的弟兄们，一旦找到一些钱用来买什么东西，就会把它放进存钱罐，过后才打开；他们因为看不到自己所拥有的，就因此失去它吗？不要怕：孩子们把钱放进存钱罐，就安心了；你把它放在基督手中，却害怕吗？要明智，为自己在天堂的未来做准备。因此要明智，效法蚂蚁，正如圣经所说：\Quote{夏天储存，免得冬天饥饿}；冬天是末日，是患难的日子；冬天是冒犯和痛苦的日子；要为你未来的需要积聚；但如果你不这样做，你就要既鲁莽又愚昧地灭亡。\par

13. 然而那个富人（\BibleRef{路 16:22}）也死了，并为他举行了类似的葬礼。看看人们将自己带到了何种地步：他们不看他在世时生活多么邪恶，而只关注他死时有何等排场！哦，有这么多人哀悼的他是有福的！但另一个人生活的方式却是很少有人哀悼。因为所有人都应当哀悼一个生活如此悲惨的人。但送葬队伍出现了；他被安放在昂贵的坟墓中，被包裹在昂贵的衣袍中，被埋葬在香料和香膏中。其次，他的纪念碑何等华丽！多么像大理石！他真住在那纪念碑里吗？他在里面是死的。人们认为这些是好事，就远离了天主，没有寻求真正的好事，被虚假的好事所欺骗。为此，请看下文。那不肯付出自己灵魂赎价的人，那不明白死亡的人，因为他看到有智慧的人死亡，他就变得鲁莽和愚昧，以便与他们一同死亡。他们将怎样灭亡，就是那些\Quote{把自己的财富留给外人}的人？\par

14. 但这些所谓的\Quote{外人}真的服事那些被称为\Quote{自己的}人吗？请听他们在什么事上服事他们，看看他们是如何被嘲弄的：为什么他说\Quote{对外人}？因为他们对他们毫无益处。然而，他们在什么事上自以为有益呢？\Quote{他们的坟墓将成为他们永久的房屋}（\BibleRef{咏 48:11}）。如今因为这些坟墓被建造，坟墓就是房屋。因为你常听到富人说：我有一座大理石房屋，我必须离开，我却不为自己考虑一座永久的房屋，在那里我将永远居住。当他想着为自己建造一座大理石或雕刻的纪念碑时，他仿佛在想一座永久的房屋：好像这个富人就会住在那里！如果他住在那里，就不会在地狱里焚烧了。我们必须考虑，恶人的灵魂所居之处，并不是他必死的尸身所安放之处；而是\Quote{他们的坟墓将成为他们永久的房屋。他们的住所代代相传。}\Quote{住所}是他们在世时暂时居住的地方；\Quote{房屋}是他们似乎永久居住的地方，即他们的坟墓。因此，他们离开住所（他们活着时居住的地方）留给家人，而他们则仿佛迁到永久的房屋，即他们的坟墓。\Quote{他们的住所代代相传}对他们有什么益处呢？假设一代又一代，将会有儿子、孙子、曾孙子；他们的住所对他们有什么益处呢？什么？听：\Quote{他们将在他们的土地上呼求他们的名字。}这是什么意思？他们将把饼和酒带到坟墓前，在那里呼求死者的名字。你想想看，那个富人死后，他的名字被多么大声地呼求，当人们在他的墓前喝得酩酊大醉时，却没有一滴酒落在他那焦灼的舌头上！人们是在满足自己的肚腹，而不是在服事他们朋友的灵魂。死者的灵魂什么也得不到，除非他们活着时自己所做的；但如果他们活着时自己什么也没做，那么死后什么也得不到。但活着的人做什么呢？他们只是\Quote{在他们的土地上呼求他们的名字。}\par

15. \\BibleQuote{人虽受尊荣，却不彻悟，竟与无理性的畜类相比，且变得与它们相似} \\BibleRef{咏 48:12} ……相反，他们应当藉善工为自己预备永恒的居所，应当为自己预备永生，应当先行赒济，紧随其后的是他们的善工，应当照顾有需要的同伴，应当施与那与他们同行的人，不应藐视躺卧在他们门前满身疮痍的基督，祂曾说过：\\BibleQuote{你们对我这些最小兄弟中的一个所做的，就是对我做的} \\BibleRef{玛 25:40}。然而，\\Quote{人虽受尊荣，却不彻悟}。什么是\\Quote{受尊荣}？人按天主的肖像和模样受造，故高于畜类。因为天主造人并非如造畜类那样；但天主造人，是要畜类为人服务：那么，是靠他的力量，而非靠他的理智吗？不。但他\\Quote{不彻悟}；那按天主的肖像所造的人，\\Quote{竟与无理性的畜类相比，且变得与它们相似}。因此在别处说：\\Quote{不要像马和骡子，它们没有理智}。\par

16. \\BibleQuote{他们自己的道路成为他们的绊脚石} \\BibleRef{咏 48:13}。这绊脚石是对他们，不是对你们。但何时它也会成为你们的绊脚石呢？如果你们以为这样的人是有福的。如果你们领悟到他们并非有福，他们自己的道路就是他们自己的绊脚石；不是对基督，不是对祂的身体，不是对祂的肢体。\\Quote{以后他们要用口赞美}。\\Quote{以后他们要用口赞美}是什么意思？虽然他们变得如此，只寻求现世的财物，但他们成了伪善者：当他们赞美天主时，是口唇赞美，而非内心赞美。像这样的基督徒，当有人向他们推荐永生，并告诉他们，因基督的名，他们应当轻视财富时，他们心中便做鬼脸：如果他们不敢公开做，怕脸红，或怕被人责备，但他们仍是在心里做，并加以嘲笑；于是，他们口中有祝福，心里却有咒骂。\par

\\Emph{第二部分}\par

1. \\BibleQuote{他们如羊被放在阴府中，死亡是他们的牧者} \\BibleRef{咏 48:14}。是谁的牧者？是那些其道路成为自己绊脚石之人的牧者。是谁的牧者？是那些只关心现在，而不思考未来之人的牧者；是那些不想任何生命，只想到那必称为死亡的生命之人的牧者。所以，他们如羊在阴府，有死亡作他们的牧者，并非无因。\\Quote{他们有死亡作他们的牧者}是什么意思？死亡究竟是一个事物，还是一种势力？是的，死亡或是灵魂与肉身的分离，或是灵魂与天主的分离；而人们所害怕的，是灵魂与肉身的分离；但真正的死亡，即人们不害怕的，是灵魂与天主的分离。人们往往在恐惧那使灵魂与肉身分离之事时，反而陷入那使灵魂与天主分离的境地。这就是死亡。但\\Quote{死亡是他们的牧者}又是如何呢？如果基督是生命，那么魔鬼就是死亡。我们在圣经许多地方读到，基督是生命。但魔鬼是死亡，并非因为他本身就是死亡，而是因为他带来了死亡。无论是亚当所陷入的死亡，那死亡是由他的劝说灌入人的，还是灵魂与肉身分离的死亡，人仍视他为肇始者：他首先因骄傲而堕落，嫉妒那站立着的人，以无形死亡推倒那站立者，好让他偿还可见的死亡。那些属于他的人，有死亡为他们的牧者；但我们这些思考将来不死的人，且并非无故在额上佩戴基督十字架的标记，除了生命之外没有其他牧者。对不信者而言，死亡是牧者；对信者而言，生命是牧者。如果阴府中的羊，其牧者是死亡，那么天上的羊，其牧者是生命。那么怎样？我们现在已在天上吗？我们因信德而在天上。因为若不在天上，那\\Quote{举起你们的心}又在哪里？若不在天上，那么与宗徒保禄一起说：\\BibleQuote{因为我们的家乡在天上} \\BibleRef{斐 3:20}又是从何而来？在身体上我们行走于地上，在心里我们居于天上。我们住在那里，如果我们把任何能抓住我们的事物送往那里。因为除了思想所在之处，无人居于心里；但他的思想在哪里，他的宝藏也在那里。他若积蓄在地，他的心便不离地；他若积蓄在天，他的心便从天下不坠；因为主明明说：\\BibleQuote{你的宝藏在哪里，你的心也必在哪里} \\BibleRef{玛 6:12}。\par

2. 那么，那些以死亡为牧者的人，似乎一时兴旺，而义人却劳苦：但为什么呢？因为仍是黑夜。什么是\\Quote{黑夜}？义人的功德不显，不义者的福乐仿佛有名无实。只要还是冬天，青草就比树木更显青翠。因为青草在冬天繁茂，树木在冬天则像枯干；当夏天来临，太阳发出更大的热力时，那在冬天看似枯干的树木就迸发绿叶，结出果实，而青草却枯干：你将会看到树的尊荣，草却干枯。同样，现在义人劳苦，是在夏天到来之前。根里有生命，却尚未显现在枝条上。但我们的根就是爱。宗徒怎么说？他说我们应当将根植于天上，好使生命成为我们的牧者，因为我们的居所不应离开天上，因为我们在地上应当如同已死那样行走；这样，我们生活在天上，在地上却是死的；而不是死在天上，活在地上。……我们的劳苦将在早晨显现，早晨必结果实：这样，现在劳苦的人，将来要统治；现在自夸骄傲的人，将来要被制服。因为下文是什么？\\Quote{他们如羊被放在阴府中，死亡是他们的牧者；义人要在清晨管辖他们}。\par

3. 忍耐黑夜，渴望清晨。不要以为黑夜有生命，清晨就没有生命。那么，睡觉的人活着，而醒来的人就不活吗？睡觉的人不是更像死亡吗？谁又是那些睡觉的人呢？就是那些宗徒\Person{保禄}所唤醒的人，只要他们愿意醒来。因为他曾对某些人说：\BibleQuote{睡着的人，醒来罢，从死者中起来，基督必要光照你。} \BibleRef{弗 5:14} 那些被基督光照的人，现在警醒，但他们警醒的果实尚未显现；到了清晨才会显现，就是当这世上一切不确定之事都过去的时候。因为这就是黑夜：它们在你眼中难道不像黑暗吗？……但是那些被人践踏、因信仰而受嘲笑的人，将要从他们作为牧者的生命本身听到：\Quote{你们蒙我父降福的，来承受自创世之初就为你们预备好的王国罢。} 因此义人\Quote{要统治他们}，但不是现在，而是\Quote{在清晨。} 不要有人说：我为什么是基督徒？我统治不了任何人，我倒想统治恶人。不要着急，你将统治，但要\Quote{在清晨。} \Quote{他们的帮助，从他们的荣耀中，将在地狱里衰老。} 现在他们拥有荣耀，在地狱里他们将要衰老。什么是\Quote{他们的帮助}？来自金钱的帮助，来自朋友的帮助，来自自身力量的帮助。但是当人一死，\Quote{那一天他的所有计划都要消失。} 他活着时在人间显得多么荣耀，当他死在地狱里时，就要承受同样巨大的衰老和刑罚的腐朽。\par

4. \BibleQuote{然而，天主要救赎我的灵魂} \BibleRef{咏 48:15}。请看一个希望未来之人的声音：\BibleQuote{然而，天主要救赎我的灵魂。} 也许这是一个仍希望摆脱压迫之人的声音。某人被囚在监，他说：\BibleQuote{天主要救赎我的灵魂；}某人被捆绑，他说：\BibleQuote{天主要救赎我的灵魂；}某人在海上遇险，被波浪和狂风抛掷，他说什么？\BibleQuote{天主要救赎我的灵魂。} 他们为了此生而希望得救。但这人的声音并非如此。请听下文：\BibleQuote{当祂收纳我时，天主要从地狱手中救赎我的灵魂。} 他所说的，正是基督如今在自己身上所显示的救赎。因为祂曾下降地狱，又升上天堂。我们在元首身上所见的，我们在身体上也找到了。因为关于我们在元首所信的，那些见过的人亲自告诉了我们，我们也透过他们看见了：\BibleQuote{因为我们都是一个身体。} \BibleRef{罗 12:5} 但那些听见的人比我们更好吗？我们这些被告知的人更差吗？生命本身，我们的牧者亲自说不是如此。因为祂责备祂的一位门徒，这门徒渴望触摸祂的伤痕，当他触摸了伤痕并喊出：\BibleQuote{我的主，我的天主！} \BibleRef{若 20:28} 看见祂的门徒怀疑，并顾念将要相信的全世界，祂说：\BibleQuote{因为你看见了我，你才相信；那些没有看见而相信的，才是有福的。} \BibleQuote{但是，当祂收纳我时，天主要从地狱之地救赎我的灵魂。} 那么这里有什么？劳苦、压迫、磨难、诱惑：别指望别的。喜乐在哪里？在未来的希望里……\par

5. ……或许你的心说：我真悲惨，我想我白信了，天主不关心人间的事。因此天主唤醒我们：祂说什么？\BibleQuote{不要害怕，纵使人致富} \BibleRef{咏 48:16}。因为你为何害怕？因为人致富了吗？你害怕你白白信了，也许你为了信仰所付出的辛劳和皈依的希望都落空了：因为也许在你面前出现了带罪过的利益，你本可以致富，如果你抓住那个带罪过的利益，就不必劳苦了；而你，记起天主所威胁的，便远离了罪过，鄙视了利益：你看到另一个人通过罪过获利，却未受损害；于是你害怕行善。\BibleQuote{不要害怕，}天主的圣神对你说，\BibleQuote{纵使人致富。} 难道你只能看见眼前的事物吗？那位复活的应许了将来的事物；祂并没有应许这世上的平安和今生的安息。每个人都寻求安息；他寻求一件好事，但不是在正确的地方寻求。这生命中没有平安；我们在地上所寻求的，在天上已被应许；我们在此世所寻求的，在来世已被应许。\par

6. \BibleQuote{不要害怕，即使人变得富有，即使他家的荣耀增加。} 为何\BibleQuote{不要害怕}？ \BibleQuote{因为他死时，什么也带不走} \BibleRef{咏 48:17}。 你见他活着，要想他死去。你注意他现世有什么，注意他带走什么。他带走什么？他存有黄金，存有白银，许多田产、奴隶：他死了，这些都留下，他不知道留给谁。因为尽管他留给想给的人，却不能为他想给的人保存。因为许多人甚至得到了未留给他们的东西，许多人失去了留给他们的东西。这一切都留下，他带走什么？也许有人说，他带走裹他的东西，以及为他花费的昂贵大理石坟墓，为了立纪念碑，他带走这个。我说，连这个也没有。因为这些东西呈现给他，他毫无感觉。如果你装饰一个沉睡未醒的人，他身边有装饰品在床上：也许装饰品放在他躺卧的身体上，而他在睡梦中看见自己衣衫褴褛。他所感觉到的比他所感觉不到的更真实。尽管当他醒来时，这也不存在：然而对他睡梦而言，他在梦中所见的比他所不觉的更真实。那么，弟兄们，人为什么要对自己说：让钱财在我死时花掉吧：我为什么留给我继承人财富？他们将有我的许多东西，让我也为自己的身体留点什么。死尸有什么？腐烂的肉有什么？没有感觉的肉有什么？如果那个舌头干渴的富人有什么，那么人就有自己的东西。我的弟兄们，我们在福音中读到，这个富人出现在火中，带着丝绸和细麻衣的包裹吗？他在地狱里也像在宴席上那样吗？当他渴了求一滴水时，那些东西都不在那里。所以人什么也带不走，死人带不走埋葬所取的东西。因为哪里有感觉，哪里就有那个人；哪里没有感觉，那个人就不在。那里躺倒的是装着人的器皿，是容纳人的房屋。我们称身体为房屋，称灵魂为房屋的居住者。灵魂在地狱中受折磨：身体躺在香料和香水里，裹在昂贵的细麻布中，对他有什么益处？就像房屋的主人被放逐，你却装饰他房屋的墙壁。他在放逐中穷困、饥饿，几乎找不到一个可以打盹的小屋，而你说：\Quote{他有福，因为他的房屋被装饰了。} 谁不会断定你是在开玩笑或者是疯了？你装饰身体，灵魂却在受折磨。给灵魂一些东西，你就给了死者一些东西。但当他求一滴水却得不到时，你能给他什么？因为他藐视事先送任何东西。为什么藐视？\Quote{因为他们的道路是他们的绊脚石。} 他不顾念别的，只顾及现世生命，只想着如何被埋葬，裹在昂贵的衣服里。他的灵魂被夺去，正如主所说：\BibleQuote{愚笨的人，今夜就要索回你的灵魂，你所预备的，将归谁呢？} \BibleRef{路 12:20} 这就应验了这篇圣咏所说的：\BibleQuote{不要害怕，即使人变得富有，即使他家的荣耀增加；因为他死时什么也带不走，他的荣耀不会随他下降。}\par

7. 愿你们的爱德留意：\BibleQuote{因为他的一生，他的灵魂受祝福} \BibleRef{咏 48:18}。 他活着时，为自己做得好。这是所有人说的，但说得虚假。这是祝福者心中的祝福，并非来自真理本身。因为你说什么？因为他吃喝了，因为他做了他所选的，因为他天天奢华宴乐，所以他为自己做得好。我说，他为自己做得不好。不是我说，而是基督说。他为自己做得不好。因为那个富人，当他天天奢华宴乐时，被认为为自己做得好；但当他开始在地狱里焚烧时，那时被认为好的就被发现是坏的。因为他与上面的人所吃的，他在下面的地狱里消化。我说的是不义，弟兄们，他惯常吃的是不义。他用肉体的口吃昂贵的宴席，用他心灵的口吃不义。他用心灵的口与上面的人所吃的，他在下面的惩罚中消化。确实，他吃了一阵子，却永远消化坏的东西。那么，不义可以被吃吗？也许有人说：他说什么？不义被吃？不是我说：听圣经：\BibleQuote{醋对牙，烟对眼，不义对使用它的人也一样} \BibleRef{箴 10:26}。因为那吃了不义的，即故意有不义的人，将不能吃正义。因为正义是饼。谁是饼？\BibleQuote{我是从天上降下的生活的饼} \BibleRef{若 6:51}。他自己是我们心灵的饼……那么，正义也能被吃吗？如果它不能被吃，主就不会说：\BibleQuote{饥渴慕义的人是有福的} \BibleRef{玛 5:6}。因此，\Quote{既然他的一生，他的灵魂受祝福}，在活着时它\Quote{将}受祝福，在死亡时它将受折磨……\par

8. \Quote{祂若使你顺利，你必向祂认罪。} 弟兄们，不要成为这样的人：请看，我们为此说话，为此歌唱，为此宣讲，为此劳苦——不要做这些事。你们的行为证明你们：有时在你们的行事中，你们听到真理，却亵渎。你们亵渎教会。为什么？因为你们是基督徒。\Quote{既然如此，我就投奔多纳特派，我要作个外教人。} 为什么？因为你们吃了饼，牙齿疼痛。当你们看见饼本身时，你们赞美；当你们开始吃时，牙齿疼痛；也就是说，当你们听到天主圣言时，你们赞美；当有人对你们说：\Quote{做这个}，你们却亵渎。不要如此行恶：要说：\Quote{这饼是好的，但我不能吃。} 但现在如果你们用眼睛看，你们就赞美；当你们开始合上牙齿时，你们说：\Quote{这饼是坏的，如同做它的人一样。} 于是，你们向天主认罪，当天主使你们顺利时，你们却撒谎，当你们歌唱时：\Quote{我要常常赞美天主，祂的赞颂常在我口中。} 如何常常？如果常常获利，祂常常受赞美；如果有时亏损，祂不受赞美，反而被亵渎。实在，你们常常赞美！实在，祂的赞颂常在你们口中！你们将如刚才他所描述的：\Quote{祂若使你顺利，你必向祂认罪。}\par

9. \BibleQuote{他必进入他祖先的世代} \BibleRef{咏 48:19}：就是说，他要效法他的祖先。因为现今的不义之人有兄弟，有祖先。古时的恶人是现今之人的祖先；而现今不义之人，是将来不义后代的祖先；同样，义人的祖先，古时的义人，是现今义人的祖先；而现今的义人，是将来义人的祖先。圣神愿意表明，当人抱怨正义时，正义并非邪恶；但这些人从起初就有他们的祖先，直到他们祖先的世代。\Person{亚当}生了两个人，一个里面有不义，一个里面有正义；不义在\Person{加音}，正义在\Person{亚伯尔}。\BibleRef{若一 3:12} 不义似乎胜过正义，因为不义的\Person{加音}杀了义人\Person{亚伯尔}\BibleRef{创 4:8}，是在夜间。到了早晨岂非如此？不，\Quote{但在早晨，义人要管辖他们。} 早晨将要来临，那时就必看见\Person{亚伯尔}在哪里，\Person{加音}在哪里。同样，所有属\Person{加音}的人，以及所有属\Person{亚伯尔}的人，直到世末。\Quote{他必进入他祖先的世代，直到永远他看不见光明。} 因为即使他在这里时，也处在黑暗中，喜爱虚假的美物，不爱真实的美物；同样，他将从那里进入地狱：从他梦境的黑暗，折磨的黑暗将要迎接他。因此，\Quote{直到永远他看不见光明。}\par

但是，这是为什么呢？他在诗篇中间所写的，也在结尾写道：\BibleQuote{人在尊贵中而不醒悟，就如无理性的畜类，变成与它们一样。} \BibleRef{咏 48:20} 但是你们，弟兄们，要思想你们是按天主的肖像和相似受造的人。天主的肖像\BibleRef{创 1:26}在内在，不在身体里；不在你们所看见的这些耳朵、眼睛、鼻孔、上颚、手和脚里；它却是如此构成的：那里有理智，有心神，有发现真理的能力，有信德，有望德，有爱德，在那里天主有祂的肖像。至少在那里你们能感知并看见这些事物都要过去；因为他在另一篇诗篇中说过：\Quote{人纵然在肖像中行走，却是徒然扰乱，积蓄财宝，不知将来谁将收取。} 不要扰乱，因为无论这些事物如何，都是短暂的，如果你们是在尊贵中而醒悟的人。因为如果在尊贵中你们不醒悟，你们就如无理性的畜类，变得与它们一样。\par

\SourceRecord{Translated by J.E. Tweed.；Nicene and Post-Nicene Fathers, First Series；Vol. 8.；Edited by Philip Schaff.；Buffalo, NY: Christian Literature Publishing Co.,；1888.；<http://www.newadvent.org/fathers/1801049.htm>.}

% structure: p0001 source=h1 target=section reason=page-title

\section{《圣咏第五十篇讲疏》}

1. 天主圣言于我们改过迁善、无论是关于应许的赏报还是当畏惧的惩罚，有多大的益处，各人当在自己内衡量；并当毫不欺瞒地将良心置于祂眼前，在如此大险中不要自欺：因为你们看到，连我们的主天主自己也不欺瞒任何人：祂虽以应许福乐并坚固我们的望德来安慰我们，但对那些生活邪恶并轻视祂圣言的人，祂必定毫不宽贷。各人当趁有时间省察自己，看看自己在何处，或是持守于善，或是从恶中转变。因为正如祂在这篇圣咏中所说，不是随便什么人，也不是随便什么天使，而是\BibleQuote{主，万神之神，已经发言}\BibleRef{咏 49:1}。但祂发言，做了什么事？\BibleQuote{祂从日出之地到日落之处呼唤了大地。}那\BibleQuote{从日出之地到日落之处呼唤了世界}的，是我们的主救主耶稣基督，\BibleQuote{圣言成了血肉}\BibleRef{若 1:14}，为能居住在我们中间。因此，我们的主耶稣基督就是\BibleQuote{万神之神}；因为万有是藉着祂造的，没有祂什么也没有造成。天主圣言，如果祂是天主，就真是万神之神；但祂是否是天主，福音回答说：\BibleQuote{在起初已有圣言，圣言与天主同在，圣言就是天主。}\BibleRef{若 1:1} 如果万有是藉着祂造的，如祂随后所说，那么若有任何神被造成，也是藉着祂造成的。因为那独一的天主不是受造的，祂自己就是唯一真实的天主。但这一位独一天主，圣父、圣子、圣神，是唯一的天主。\par

2. 但那些神是谁，或在哪里，天主是他们的真天主？另一篇圣咏说：\BibleQuote{天主站在众神的会中，在众神中要施行审判。}我们尚不知道是否或许有些神聚集在天上，并在他们的集会中——这就是\BibleQuote{会中}——天主站立审判。在同篇圣咏中，请看祂对他们所说的话：\BibleQuote{我曾说：你们是神，都是至高者的子民；但你们要像人一样死亡，像一位首领一样倒毙。}由此显然可见，祂称人为神，是因祂的恩宠而被神化，并非由祂的本体所生。因为那凭自己、而非凭别人而成义的一位，祂是义的；那凭自己、而非因分享别人而成神的一位，祂是天主。但那位使人成义者，自己也使人成神，因为藉着使人成义，祂使人成为天主的子女。\BibleQuote{祂赐给他们权能，成为天主的子女。}\BibleRef{若 1:12} 若我们成了天主的子女，我们也成了神：但这是恩宠的收养效果，而非本性的生育。因为天主的独生子、天主、与圣父同为一神、我们的主救主耶稣基督，在起初就是圣言，圣言与天主同在，圣言就是神。其余后来成为神的，是因祂自己的恩宠而成，并非由祂的本体所生，以致他们与祂相同，而是因恩宠得以接近祂，并成为基督的同承继者。因为在那位承继者内，爱如此伟大，以致祂愿意有同承继者。哪个贪心的人会愿意这样，有同承继者？但即使有人愿意这样，他与人分享遗产，分享者自己所得就会比独享时少；但我们在基督中同承继的遗产，并不因拥有者众多而减少，也不因同承继者人数众多而变窄：对多人如对少数一样大，对个人如对众人一样大。\BibleQuote{请看}，宗徒说，\BibleQuote{父赐给我们何等的爱，使我们得称为天主的子女，且真是如此。}\BibleRef{若一 3:1} 在另一处：\BibleQuote{亲爱的，我们现在是天主的子女，但将来如何，还未显明。}因此我们是在希望中，尚非在实质上。\BibleQuote{但我们知道}，他说，\BibleQuote{当祂显现时，我们要像祂，因为我们要看见祂实在怎样。}\BibleRef{若一 3:2} 独生子与祂相似是由于出生，我们相似是由于看见。因为我们与祂的相似，并不像祂那样，即祂与那生祂的相同：我们是相似，而非平等；祂因为是平等，所以相似。我们已听见那些被造成而被称为义的神是谁，因为他们被称为天主的子女；而那些不是神的神，万神之神对他们是可畏的？因为另一篇圣咏说：\BibleQuote{祂在众神之上是可畏的。}好像你问：什么神？祂说：\BibleQuote{因为列邦的神都是魔鬼。}对列邦的神，即对魔鬼，祂是可畏的；对祂自己所造成的神，即对子女，祂是可亲的。此外，我发现两类神都承认天主的尊威：魔鬼承认了基督，信友也承认了基督。\BibleQuote{你是基督，永生天主之子}，\BibleRef{玛 16:16} 伯多禄说。\BibleQuote{我们知道你是谁，你是天主子}，魔鬼说。我听到相似的承认，但找不到相似的爱；不，在这里是爱，在那里是恐惧。因此，对祂是可亲的人，就是子女；对祂是可畏的人，就不是子女；对祂是可亲的人，祂使他们成为神；对祂是可畏的人，祂证明他们不是神。因为这些是祂造成的神，那些是被视为神的神；这些是真理造成的神，那些是错误如此认为的。\par

3. 因此，\Quote{万神的天主，主说了话，} \BibleRef{咏 49:1}。祂说了话，以多种方式。祂亲自藉天使说话，祂亲自藉先知说话，祂亲自用祂自己的口说话，祂亲自藉祂的忠信者说话，藉我们的卑微，当我们说任何真理时，祂亲自说话。看，祂藉着不同的说话方式、藉着许多器皿、藉着许多工具，却亲自处处回响，藉着触摸、塑造、激励：看祂所做的事。因为 \Quote{祂说了话，并呼唤了世界。} 什么世界？或许是非洲！为了那些说基督的教会是多纳图地盘的人。确实，祂并非只呼唤了非洲，但祂也没有将非洲排除在外。因为那 \Quote{从日出之地到日落之处呼唤世界} 的，没有遗漏任何祂未呼唤的部分，在祂的呼唤中也找到了非洲。因此，愿她在合一中欢欣，而不在分裂中自傲。我们说得好，万神的天主的声音已进入非洲，并未停留在非洲。因为 \Quote{祂从日出之地到日落之处呼唤了世界。} 没有地方可以隐藏异端者的阴谋，他们没有地方可在谎言的阴影下藏身；因为 \Quote{没有人能躲避它的热气。} 那呼唤世界的，也呼唤了整个世界；那呼唤世界的，呼唤了祂所创造的全部。为何假基督和假先知起来反对我？为何他们试图用诡辩的话陷害我，说：\Quote{看！基督在这里，看！祂在那里！} \BibleRef{玛 24:23}。我不听那些指认部分的人：万神的天主已指出了全部：那 \Quote{从日出之地到日落之处呼唤世界} 的，已救赎了全部；却谴责了那些虚假主张部分的人。\par

4. 但我们已听到世界从日出之地到日落之处被呼唤：那呼唤者从何处开始呼唤呢？你们也要听这事：\Quote{出自熙雍，有祂美貌的形貌。} \BibleRef{咏 49:2}。显然，圣咏与福音一致，福音说：\Quote{从耶路撒冷开始，直到万邦。} \BibleRef{路 24:47}。听，\Quote{直到万邦：}祂从日出之地到日落之处呼唤了世界。听，\Quote{从耶路撒冷开始：} \Quote{出自熙雍，有祂美貌的形貌。} 因此，\Quote{祂从日出之地到日落之处呼唤了世界，} 与主的话相符，主说：\Quote{基督必须受难，第三天从死者中复活；并且必须因祂的名宣讲悔改及罪过的赦免，直到万邦。} \BibleRef{路 24:46}。因为万邦是从日出之地到日落之处。但那句\Quote{出自熙雍，有祂美貌的形貌，} 即祂福音的美貌由此开始，祂从此开始被宣讲，祂是\Quote{在世人中美丽无比，} 与主的话相符，主说：\Quote{从耶路撒冷开始。} 新事与旧事和谐，旧事与新事和谐：两位色辣芬彼此说：\Quote{圣，圣，圣，萨巴奥特的上主天主。} \BibleRef{依 6:3}。两约都和谐，两约有一声音：愿和谐的两约之声被听见，而非被剥夺继承权的冒名者之声。万神的天主于是做了这事：\Quote{祂从日出之地到日落之处呼唤了世界，祂的形貌从熙雍先行。} 因为在那地方有祂的门徒，\BibleRef{宗 1:4}，他们在祂复活后第五十天领受了从天上派遣的圣神。从那里传出福音，从那里开始宣讲，从那里充满整个世界，且是在信德的恩宠中。\par

5. 因为当主亲自来临时，祂来是为受苦，祂是隐藏的：虽然祂在自己内是强有力的，却在肉身上显得软弱。因为祂必须显现，以便不被察觉；必须被轻视，以便被杀害。在神性中有光荣的相似，但它隐藏在肉身内。\Quote{因为如果他们认识了，就决不会将光荣的主钉在十字架上。}\BibleRef{格前 2:8}所以，祂就这样隐藏在犹太人中间，在祂的仇敌中间，行奇迹，受苦害，直到祂被挂在木头上；犹太人看见祂悬在那里越发轻视祂，并在十字架前摇着头说：\Quote{如果祂是天主子，就从十字架上下来吧！}\BibleRef{玛 27:39}所以万神之神是隐藏的，祂说话更多的是出于对我们的怜悯，而非出于祂自己的尊严。因为若非从我们身上取来，那些话从何而来？\Quote{我的天主，我的天主，祢为什么舍弃了我？}但圣父何时舍弃了圣子，或圣子舍弃了圣父？圣父与圣子不是同一天主吗？那么，\Quote{我的天主，我的天主，祢为什么舍弃了我？}若非在软弱之肉中承认了罪人的声音，又是为何？因为祂取了罪恶肉身的相似体，\BibleRef{罗 8:3}为何祂不该取罪人的声音呢？所以，万神之神是隐藏的，无论祂行走在人间，还是饥饿、口渴、疲惫而坐、带着疲乏的身体睡眠，还是被捕、受鞭打、站在审判官面前，还是对骄傲的审判官回答说：\Quote{除非由上赐给你，你对我毫无权柄；}\BibleRef{若 19:11}还是像祭品一样被牵到\Quote{剪毛者前，祂不开口，}\BibleRef{依 53:7}还是被钉十字架、被埋葬，祂始终是隐藏的万神之神。祂复活后发生了什么事？门徒们惊奇，起初不信，直到他们触摸并摸到了。\BibleRef{路 24:37}但肉身复活了，因为肉身曾经死过；那不能死的神性，在复活的肉身中仍然隐藏。形象可见，肢体可握，伤痕可摸；那创造万物的圣言，谁能看见？谁能握住？谁能触摸？然而\Quote{圣言成了血肉，居住在我们中间。}\BibleRef{若 1:14}而多默，他抓住那人，尽可能地认识了天主。因为当他触摸了伤痕，他喊道：\Quote{我的主，我的天主！}但主所显示的形象和肉身，正是他们在十字架上所见、被安放在坟墓中的那一个。祂与他们一起四十天……但触摸的多默听到了什么？\Quote{因为你看见了我，才相信；那些没有看见而相信的，才是有福的。}\BibleRef{若 20:29}我们是预言的实现。那从日出之地到日落之处被召叫的世界，没有看见却相信了。所以，万神之神是隐藏的，对祂所行走其间的人、被祂所钉死的人、在祂眼前复活的人、以及我们这些相信祂在天上坐于右边却未曾在世上行走中见过祂的人，都是隐藏的。但即使我们看见了，难道我们看见的不正是犹太人看见并钉死的那位吗？我们未见而相信基督是天主，比他们见了却只以为祂是人，更为可贵。总之，他们因恶念而杀害，我们因善信而获得生命。\par

6. 那么，弟兄们，这是怎么回事？这位万神之神，往日隐藏，现今隐藏，祂会永远隐藏吗？显然不会：请听下文：\Quote{天主必显然而来}\BibleRef{咏 49:3}。那隐藏而来的，将显然而来。祂隐藏而来是为受审判，显然而来是为审判；隐藏而来是为站在审判官前，显然而来是为审判审判者：\Quote{祂必显然而来，绝不缄默。}但为什么？难道祂现在缄默吗？我们所说的一切话从何而来？那些诫命、警告、恐怖的号角从何而来？祂不缄默，却又缄默：祂不缄默是出于警告，缄默是出于报复；祂不缄默是出于训诫，缄默是出于审判。因为祂每日容忍罪人作恶，不关心天主，不在良心上，也不在地上：这一切都逃不过祂的眼目，祂普遍地劝诫所有人；而每当祂在地上惩罚任何人时，那是劝诫，而非定罪。所以祂在审判上是缄默的，祂隐藏在天上，至今仍为我们代祷；祂对罪人宽宏大量，不发出祂的忿怒，而是等待悔改。祂在另一处说：\Quote{我缄默不言，难道我永远缄默吗？}\BibleRef{依 42:14}当祂不再缄默时，\Quote{天主必显然而来。}哪一位天主？\Quote{我们的天主。}而且就是这位天主，祂是我们的天主：因为谁不是我们的天主，谁就不是天主。因为外邦人的神是魔鬼：基督徒的天主是真天主。祂必要来临，但\Quote{显然，}不再被戏弄，不再被掌掴和鞭打：祂必要来临，但\Quote{显然，}不再被芦苇击打头部，不再被钉十字架、杀死、埋葬：因为这一切是隐藏的天主意愿承受的。\Quote{祂必显然而来，绝不缄默。}\par

7. 但祂将要来审判，以下的话教导我们。\Quote{火要在祂面前燃烧。} 我们害怕吗？我们当改变，就不会害怕。让糠秕害怕火：火对金子能做什么？你所能做的，现在在于你的能力，因此你不会因不愿改正而经历那将要临到的事，甚至违背你的意愿。因为，弟兄们，如果我们可以使审判之日不来，我想即使那样我们也不该活得败坏。即使审判之日的火不来，而罪人面对的只是与天主面容的分离，无论他们在何等丰裕的享乐中，不见那创造他们的主，与祂不可言喻的容颜的甜美隔绝，无论罪恶多么长久和不受惩罚，他们也当哀哭自己。但我该说什么，又该对谁说？这是对爱慕者的惩罚，而非对轻蔑者的惩罚。那些已经开始多少感受到智慧与真理甜美的人，知道我说的是什么，单单与天主面容分离是多么大的惩罚；但那些没有尝过那甜美的人，若还不对天主的面容满怀渴望，就让他们甚至害怕火；让惩罚恐吓那些奖赏无法赢取的人。天主所应许的对你们没有价值，你们当畏惧祂所威胁的。祂同在的甜美将要来临；你没有悔改，你没有警醒，你不叹息，你不渴慕；你拥抱你的罪恶和肉体的享乐，你正为自己堆积碎秸，火将来临。\Quote{火要在祂面前燃烧。} 这火不像你家的炉火，然而如果你被迫把手伸进去，那威胁给你这个选择的人无论要你做什么，你都会做。如果他对你说：\Quote{写下反对你父亲性命的文书，写下反对你儿女性命的文书，因为你不做，我就把你的手伸进我的火里：} 你将会这样做，为使你的手不被烧，为使你的肢体暂时不被烧，虽然它不会永远痛苦。你的敌人不过威胁如此轻微的恶，你就作恶；天主威胁永久的恶，你却不从善吗？作恶，甚至不应该被威胁驱迫；行善，甚至不应该被威胁阻止。但因着天主的威胁，因着永火的威胁，你被劝阻离开邪恶，被邀请向善。它为何使你忧伤？岂不是因为你不相信？那么，让每个人省察自己的心，看看信德在那里持守什么。如果我们相信审判要来，弟兄们，让我们好好生活。现在是仁慈的时候，那时将是审判的时候。没有人会说：\Quote{把我召回到我过去的年岁。} 即使那时人们会后悔，但后悔也是徒然；现在要有悔改，趁悔改还有果实；现在，让一筐粪土施于树根，\BibleRef{路 13:8}，心中的忧伤和眼泪；免得祂来连根拔除。因为当祂拔出后，火就要被期待了。现在，即使枝条已经被折断，它们仍可以被接上：\BibleRef{罗 11:19}；那时，\BibleQuote{凡不结好果子的树，都要砍倒，投入火中。} \BibleRef{玛 3:10} \Quote{火要在祂面前燃烧。}\par

8. \Quote{有大风暴环绕祂} \BibleRef{咏 49:3}。\Quote{有大风暴，} 为簸扬这么大的禾场。在这风暴中，将有那样的簸扬，藉此从圣人中除去一切不洁，从信者中除去一切虚假，从敬神的人和敬畏天主圣言的人中除去一切嘲弄者和骄傲者。因为如今有一种混杂存在那里，从日出之地到日落之地。让我们看看那将要来的祂将如何行，祂将用那\Quote{成为围绕祂的大风暴}作什么。无疑，这风暴是为了进行一种分离。那就是他们不愿等待的分离，他们在来到岸边之前网就裂开了。\BibleRef{路 5:6} 但在这分离中，善人与恶人之间做出了一种区分。有些人现在轻省着肩膀跟随基督，没有世俗挂虑的重担，他们没有徒然听到：\BibleQuote{你若愿意是成全的，去变卖你所有的，施舍给穷人，你必有宝藏在天上；然后来跟随我。} \BibleRef{玛 19:21} 对于这样的人，曾说：\BibleQuote{你们必要坐在十二宝座上，审判以色列十二支派。} \BibleRef{玛 19:28} 因此有些人要与主一同审判，另一些人则被审判，但被放在右边。因为将来有一些人要同主一起审判，我们有最明确的证据，就是我刚才引用的：\Quote{你们必要坐在十二宝座上，审判以色列十二支派。}……\par

9. 主在复活后所行的，预示了我们复活后在天国中的情形，那里将没有恶人。……最后，那七千人——就是答覆厄里亚所说的：\BibleQuote{我为自己保留了七千人，从未向巴耳屈过膝的}\BibleRef{列上 19:18}——远远超过那鱼的数量。因此，那一百五十三条鱼\BibleRef{若 21:11}并非仅仅表达如此数目的圣人，而是圣经以这一特定数字表达了全部圣人和义人的数目；也就是说，在那一百五十三条中，可以理解为所有与复活进入永生相关的人。因为法律有十诫\BibleRef{申 4:13}；而恩宠之神——唯独藉着祂，法律得以成全\BibleRef{依 11:2}——被称为七重的。那么，必须探究这数字十和七的含义：十代表诫命，七代表圣神的恩宠；藉此恩宠，诫命得以满全。因此，十和七包含了所有属于复活、右边、天国、永生的人，即那些藉着圣神的恩宠——而非凭自己的行为或功绩——满全法律的人。但十和七，若你从一数到十七，依次累加所有数字：加二得三，加三得六，加四得十，加五得十五，加六得二十一，加七得二十八，加八得三十六，加九得四十五，加十得五十五，加十一得六十六，加十二得七十八，加十三得九十一，加十四得一百零五，加十五得一百二十，加十六得一百三十六，加十七得一百五十三——你会发现所有圣人的庞大数目都属于这少数鱼所代表的数字。同样，就如五个童女代表无数童女；就如地狱中受折磨之人的五个兄弟代表成千上万的犹太人；就如一百五十三条鱼代表千千万万的圣人：十二个宝座也不是指十二个人，而是指众多完美的人。\par

10. 但我看出接下来需要我们做什么；正如在五个童女的事例中，给出了为何许多属于五，为何那些五代表许多犹太人，以及为何一百五十三代表许多完美的人——同样，也要说明为何十二宝座不是指十二个人，而是指许多人。十二宝座意味着什么？它代表所有地方被赋予如此完美能力的人，就是那些完美的人，以致对他们说：\BibleQuote{你们要坐在十二宝座上，审判以色列十二支派}\BibleRef{玛 19:28}。为何所有地方的人都属于十二这个数字？因为我们所说的\Quote{普世}，是指整个世界；而大地的四极包含四个特定的方向：东、西、南、北。从所有这些方向，他们被召唤进入天主圣三中，并在天主圣三的信德和训诫中得以完美——既然三乘以四等于十二——你便明白为何圣人属于整个世界；他们将要坐在十二宝座上审判以色列十二支派，因为以色列十二支派也代表了整个以色列的十二支派。因为正如审判者来自普世，受审判者也来自普世。宗徒圣保禄自己在责备那些信奉教友的平信徒时，因为他们不将诉讼交给教会，却将与他们有纠纷的人拖到法庭面前，他说：\BibleQuote{你们不知道我们将要审判天使吗？}\BibleRef{格前 6:3}请看，祂使自己成了怎样的审判者：不仅他自己，还有所有在教会中按正义审判的人。\par

11. 既然显然，许多人要与主一同审判，但另一些人要受审判，然而不是平等的，而是按他们的功过；祂要带着祂所有的\OriginalTerm{天使}{Angels}来临，\BibleRef{玛 25:31}那时万民都要聚集在祂面前，在众天使中也要算那些已变得如此完美的人，他们坐在十二个宝座上审判以色列十二支派。因为人被称为天使：宗徒论到自己说：\Quote{你们接待我如同天主的使者。}\BibleRef{迦 4:14}关于若翰洗者，经上说：\Quote{看，我派遣我的使者在你面前，他要在你前面预备你的道路。}因此，带着众天使来临，祂也要同祂一起拥有圣人。因为\Person{依撒意亚}也明说：\Quote{他必与百姓的长老一同前来审判。}\BibleRef{依 3:14}那么，那些\Quote{百姓的长老}，那些刚才称为天使的，那些从普世而来的成千上万已完美的许多人，被称为上天。但另一些\Emph{则被称为}大地，然而是有果实的。什么是那有果实的大地？就是那要放在右边，将要听到说：\Quote{我饿了，你们给了我吃：}\BibleRef{玛 25:35}真正有果实的大地，宗徒在其中欢欣，当他们送给他所需之物时：\Quote{我并不求馈赠，}他说，\Quote{但只求果实。}\BibleRef{斐 4:17}他感谢说：\Quote{因为你们竟再开始关心我。}\BibleRef{斐 4:10}他说：\Quote{你们再开始，}如同树木曾因某种不结果实而枯萎了一样。因此主来审判（好让我们现在听圣咏，弟兄们），祂要做什么？\Quote{祂要从高天呼唤上天}\BibleRef{咏 49:4}。上天，即所有圣人，那些将成为审判者的完美之人，祂要从高处呼唤他们，使他们与祂同坐，审判以色列十二支派。因为当上天总是在高处时，祂如何\Quote{要从高天呼唤上天}？但祂在此称为上天的，在别处称为诸天。什么诸天？就是那些讲述天主荣耀的：因为\Quote{诸天述说天主的荣耀}，关于此经上说：\Quote{它们的声音传遍普世，它们的言语达于地极。}因为请看主在审判中分离：\Quote{祂要从高天呼唤上天和大地上来，分别祂的子民。}从谁中分别？不是从恶人中么？关于这些人，此后不再提，仿佛已被定罪受罚。请看这些善人，并分别。\Quote{祂要从高天呼唤上天和大地上来，分别祂的子民。}祂也呼唤大地，但不是为了联合，而是为了分开。因为起初祂召集他们在一起，\Quote{当万神之神发言，从日出之地到日落之地召集大地时，}祂尚未分开：那些仆人曾被派去邀请赴婚宴，\BibleRef{玛 22:3}他们聚集了善恶。但当万神之神显明地来临并不再沉默时，祂将如此呼唤\Quote{上天}使其与祂一同审判。因为上天是什么，诸天本身就是；正如大地是什么，诸地本身就是；正如教会是什么，众教会本身就是：\Quote{祂要从高天呼唤上天和大地上来，分别祂的子民。}现在，祂与上天一起分别大地，也就是说，上天与祂一同分别大地。祂如何分别大地？这样：把一些放在右边，另一些放在左边。但对已分别的大地，祂说什么？\Quote{我父所祝福的，你们来吧！承受自创世以来给你们预备了的国度吧！因为我饿了，你们给了我吃的，}等等。但他们说：\Quote{我们几时见了你饿了？}祂说：\Quote{你们对我这些最小兄弟中的一个所做的，就是对我做的。}\Quote{因此，祂要从高天呼唤上天和大地上来，分别祂的子民。}\par

12. \Quote{为祂聚集祂的\OriginalTerm{义人}{righteous}}\BibleRef{咏 49:5}。神圣而先知的声音，看见未来之事如同现在，它劝勉聚集的天使。因为祂要派遣祂的天使，万民都要聚集在祂面前。\BibleRef{玛 25:32}为祂聚集祂的义人。除了那些凭\OriginalTerm{信德}{faith}生活并施行\OriginalTerm{怜悯}{mercy}工作的人，还有什么义人？因为那些工作是正义的工作。你们有福音：\Quote{你们要小心，不要在人前施行你们的\OriginalTerm{正义}{righteousness}，为叫他们看见。}\BibleRef{玛 6:1}又好像有人问：什么正义？\Quote{所以，当你\OriginalTerm{施舍}{alms}时，}祂说。因此，祂指明施舍是正义的工作。那些人聚集给祂的义人：聚集那些怜悯贫困者的人，那些顾念贫困和贫穷的人：聚集他们，\Quote{愿主保护他们，使他们存活；}\Quote{为祂聚集祂的义人：那些将祂的\OriginalTerm{盟约}{covenant}置于\OriginalTerm{祭献}{sacrifice}之上的人：}也就是说，他们看重祂的应许超过他们所行的事。因为那些事是祭献，天主说：\Quote{我喜欢仁慈胜过祭献。}\Quote{谁遵守祂的盟约胜过祭献。}\par

13. \Quote{诸天要宣扬祂的正义} \BibleRef{咏 49:6}。这义真是天主赐予我们的，\Quote{诸天已经宣扬了}，福音传播者们已经预告了。通过他们，我们听说有些人将在右边，家主对他们说：\Quote{来吧，我父所祝福的，承受} \BibleRef{玛 25:34}。承受什么？\Quote{天国}。作为什么回报？\Quote{我饿了，你们给了我吃的}。有什么比给饥饿者擘饼更微不足道、更属世的呢？天国竟被如此估价。\Quote{擘开你的饼给饥饿的人，将无家可归的穷人领到你的屋里；若你看见赤身露体的，给他衣服穿} \BibleRef{依 58:7}。若你没有擘饼的资财，没有可以领人进来的屋子，没有可以遮盖的衣裳：给一杯凉水 \BibleRef{玛 10:42}，往银库里投两个小钱 \BibleRef{谷 12:42}。寡妇用两个小钱所买得的，与伯多禄舍弃渔网 \BibleRef{玛 4:20} 所买得的相同，与匝凯献出一半财产 \BibleRef{路 19:8} 所买得的相同。你所有的一切都值这么多。\Quote{诸天要宣扬祂的正义，因为天主是审判者}。确实是审判者，不是混淆而是分别。因为\Quote{主认识属于祂的人} \BibleRef{弟后 2:19}。即使麦粒藏在糠秕中，农夫也认识它们。没有人该害怕自己在糠秕中是一粒麦子；我们簸扬者的眼睛不会被欺骗。不要害怕那环绕祂的暴风会把你和糠秕混淆。暴风固然强大，但它不会把任何一粒麦子从麦粒的一侧扫到糠秕中去：因为审判者不是任何拿着三齿叉的乡下人，而是三位一体的天主。\Quote{诸天要宣扬祂的正义，因为天主是审判者}。让诸天去吧，让诸天传扬吧，让他们的声音传到每个地方，他们的话语直到地极；也让那身体说：\Quote{我从地极向你呼求，当我心沉重时}。如今它呻吟混合，将来分别时它要欢喜。让它呼求并说：\Quote{求你不要将我的灵魂与不敬虔的人一同消灭，不要将我的生命与流人血的人一同消灭}。祂不一同消灭，因为天主是审判者。让它向祂呼求并说：\Quote{上主，求你审判我，为我伸冤，脱离不圣洁的百姓}；让它说，祂将成就这事：祂的义人将聚集到祂那里。祂呼唤大地，为要分别祂的百姓。\par

14. \Quote{听吧，我的百姓，我要向你发言} \BibleRef{咏 49:7}。祂必来临，不再缄默；看，如今你若听，祂就不缄默。听吧，我的百姓，我要向你发言。因为若你不听，我就不向你发言。\Quote{听吧，我要向你发言}。因为若你不听，即使我发言，也不是对你说的。那么我何时向你发言呢？若你听？你何时听呢？若你是我的百姓。因为\Quote{听吧，我的百姓：}你若是一个异民，你就听不见。\Quote{听吧，我的百姓，我要向你发言：以色列，我要向你作证}……因为\Quote{你的天主}，这是恰当地对那一个人说的，天主更将他当作自己家里的人，如同在他家中，如同在他特有之中：\Quote{我是你的天主}。你还要什么？你向天主求赏报，使天主能给你些什么；使祂所赐给你的能成为你自己的？看，那将要赐予的天主本身就是你的。有谁比祂更富有？你曾渴望恩赐，你已得到了赐予者本身。\Quote{天主，你的天主，我是}。\par

15. 祂向人要求什么，让我们看；我们的天主，我们的君王和我们的皇帝向我们征什么税，既然祂愿意作我们的君王，并愿意我们作祂的省区？我们听祂的命令。愿穷人不要在天主的命令前战栗：天主命令人献给祂的，祂先赐予那命令的本身；你们只要虔诚。天主不索取祂未赐予的，而祂向所有人都赐予了祂所索取的。因为祂索取什么？现在让我们听：\Quote{我不会因你的祭献而责备你} \BibleRef{咏 49:8}。我不会对你说：为什么你没有为我宰杀肥牛？为什么你没有从你的羊群中选出最好的山羊？为什么那只公绵羊在你羊群中漫步，而没有放在我的祭台上？我不会说：查看你的田地、你的畜圈和你的围墙，寻找你可以给我的东西。\Quote{我不会因你的祭献而责备你}。那么：你不接受我的祭献吗？\Quote{但你的全燔祭常在我眼前} \BibleRef{咏 49:9}。圣咏另一处论及某些全燔祭说：\Quote{你若渴求祭献，我必献上，全燔祭你不喜悦}；然后他转向自己，\Quote{天主所要的祭献是忧伤的灵，破碎和谦卑的心，天主不轻视}。那么哪些是全燔祭祂不轻视？哪些全燔祭常在他眼前？\Quote{上主，请按你的善意恩待熙雍，建造耶路撒冷的城墙，那时你必悦纳正义的祭献、供物和全燔祭}。他说某些全燔祭天主会悦纳。但什么是全燔祭？全被火烧尽的：\OriginalTerm{燃烧}{causis} 是焚烧，\OriginalTerm{全部}{holon} 是全体：而 \Quote{全燔祭} 就是全被火烧尽的。有一种最炽热的爱火：让心灵被爱火点燃，让同样的爱驱使肢体为它所用，不容它们服务于贪欲，好让我们全然燃烧于天主之爱的火焰中，向天主献上全燔祭。这样的 \Quote{你的全燔祭常在我眼前}。\par

16. 那以色列或许还不明白祂常摆在眼前的那些\OriginalTerm{全燔祭}{holocausts}是什么，仍在思念牛、羊、公山羊：让它不要这样想：\Quote{我不接受你家的牛犊。} 我说全燔祭；你立刻在心思意念中奔向地上的羊群，从中为我选择一些肥美的：\Quote{我不接受你家的牛犊。} 祂是在预言新约，其中所有这些祭献都已停止。因为那时它们在预告将要来临的某一种祭献，藉其血我们得以洁净。\Quote{我不接受你家的牛犊，也不接受你羊群中的公山羊。}\par

17. \BibleQuote{因为林中的百兽都是我的} \BibleRef{咏 49:10}。我为何要向你求我所造的呢？它更是你的，因为是我给你拥有，还是我的，因为是我造的？\Quote{因为林中的百兽都是我的。} 但或许那以色列说：百兽是天主的，那些野生的我没有关在圈里，我没有拴在槽上；但这牛、羊和公山羊——这些都是我自己的。\Quote{山上的牲畜，和公牛。} 那些你没有拥有的，是我的；这些你拥有的，也是我的。因为如果你是我的仆人，你的全部财产都是我的。因为不可能仆人所获得的是主人的财产，而主人自己所创造的却反而不是主人的。因此，你没有捕获的林中的百兽是我的；山上的属于你的牲畜，以及你槽旁的公牛，也都是我的：它们全是我的，因为我造了它们。\par

18. \BibleQuote{我知道天空的一切飞鸟} \BibleRef{咏 49:11}。祂如何知道？祂称量了它们，数算了它们。我们中有谁知道天空的一切飞鸟？但即使天主赐一些人知道天空的一切飞鸟，祂自己的知道也不同于祂赐人知道的方式。天主的知道是一回事，人的知道是另一回事；同样，天主的拥有是一回事，人的拥有是另一回事。因为你所拥有的，并不完全在你的权力之下；或者你的牛，只要它活着，就在你的权力之下，以至于它要么不死，要么不需喂养。谁拥有至高权力，谁就拥有至高和最奥妙的认知。让我们在赞美天主时，将这归于天主。我们不要胆敢说：天主如何知道？我恳请你们，弟兄们，不要期望我向你们展开说明天主如何知道：我只说这一点：祂不像人那样知道，也不像天使那样知道；祂如何知道，我不敢说，因为我也无法知晓。但有一件事我知道：甚至在天上的一切飞鸟存在之前，天主已经知道祂要创造的是什么。那知道是什么？人啊，你开始看见，是在你被造之后，在你有了视觉之后。这些飞禽是藉天主的话语从水中生出的，祂说：\BibleQuote{让水生出飞禽} \BibleRef{创 1:20}。天主如何知道祂命令水生出的事物呢？现在当然祂知道祂所创造的，并且在创造之前祂就已经知道了。天主的知道如此伟大，以至于在它们被创造之前，它们就在祂内以一种不可言喻的方式存在：而祂却期望从你那里接受祂在创造之前就已拥有的东西？\Quote{我知道天空的一切飞鸟}，你不能把它们给我。你将要为我宰杀的东西，我都知道：不是因为我造了才认识，而是为了我才能造。\Quote{田野的华美与我同在}。田野的美丽，大地上所生的一切丰盛，\Quote{与我同在}，祂说。如何与祂同在？甚至在被造之前，它们就已经是这样吗？是的，因为所有将来的事都与祂同在，所有过去的事也与祂同在：将来的事是这样，以至于过去的一切并不离开祂。借由居于圣言中的天主不可言喻的智慧的一种认知，万有都与祂同在，而圣言本身就是万有。田野的华美岂不也以一种方式与祂同在，因为祂无所不在，并且祂自己说过：\BibleQuote{我充满天地} \BibleRef{耶 23:24}？什么不与祂同在呢？关于祂写道：\Quote{我若升到天上，祢在那里；我若下到地狱，祢也在那里}？与祂同在的是全体：但祂与它们同在的方式是，祂不会因祂所创造的任何事物而受沾染，或有所缺乏。因为与你同在的，或许有一根柱子，你站在它旁边，当你疲倦时，你倚靠它。你需要那与你同在的东西，但天主不需要与祂同在的田野。与祂同在有田野，有大地之美，有天空之美，有所有飞鸟，因为祂自己无所不在。何以万物都与祂相近？因为在万物存在或被造之前，万物都被祂所知道。\par

19. 谁能解释，谁能阐述在另一篇圣咏中向祂所说的：\Quote{你不需要我的美物}？祂已说过，祂不需要我们提供任何必需品。\BibleQuote{我若是饥饿，我不必告诉你}\BibleRef{咏 49:12}。那保护以色列的，既不饥饿，也不口渴，不疲乏，也不打盹。但是，看！我按你们的血肉说话：因为你们若不进食就会饥饿，或许你们以为天主也需要进食才会饿？即使祂饿了，祂也不告诉你们：万有都在祂面前，祂从何处愿意，就取用祂所需之物。这些话是为说服悟性薄弱的人，并非天主在宣告祂的饥饿。虽然为了我们的缘故，这位万神的天主甚至屈尊就饿。祂来，为饥饿，也为饱饫；祂来，为口渴，也为赐饮；祂来，为穿上可朽坏的身体，也为穿上不朽；祂来，为贫穷，也为使人富足。因为祂取了我们的贫穷，并未失去祂的财富，因为\BibleQuote{在祂内蕴藏着一切智慧和知识的宝库}\BibleRef{哥 2:3}。\BibleQuote{我若是饥饿，我不必告诉你；因为世界和其中的一切都是我的。}那么，不要再劳苦寻找能给我什么，因为我若没有你们，仍有我所要的。\par

20. 那么，你为什么仍思念你的羊群？\BibleQuote{我岂吃公牛之肉，岂喝公山羊的血吗？}\BibleRef{咏 49:13}。你已经听见祂不要求我们什么，祂愿意给我们一些命令。如果你曾想过这类事，现在请将你的思想从这些事上撤回：不要想给天主献上任何这类东西。若你有一头肥牛，杀给穷人吧：让他们吃公牛的肉，虽然他们不会喝公山羊的血。当你这样做时，祂会记在你身上，祂曾说过：\Quote{我若是饥饿，我不必告诉你：}并且祂要对你说：\BibleQuote{我饿了，你们给了我吃}\BibleRef{玛 25:35}。\Quote{我岂吃公牛之肉，岂喝公山羊的血吗？}\par

21. 那么，请说，我们的主天主，祢命令祢的子民、祢的以色列做什么？\BibleQuote{应向天主祭献赞颂的祭品}\BibleRef{咏 49:14}。我们也对祂说：\Quote{天主，我向祢所许的誓愿在我内，我将用散文向祢偿还。}我曾害怕祢可能命令一些非我能力所及的事，我以为那些事在我羽毛笔中，但或许刚才已被小偷取走。祢命令我什么？\Quote{应向天主祭献赞颂的祭品。}让我回到自我，在其中我能找到我可祭献的：让我回到自我；在我内我能否找到赞颂的祭献：愿祢的祭台成为我的良心。我们毫不忧虑，我们不前往阿拉伯寻找乳香；我们不筛选任何贪婪商人的口袋：天主要求我们的是赞颂的祭献。匝凯在产业中有赞颂的祭献；\BibleRef{路 19:8}那寡妇在钱袋中也有；\BibleRef{谷 12:42}某个穷主人或其他人在罐中也有；另一人既无产业，也无钱袋，也无罐，一无所有，全在心中：救恩临到了匝凯的家；并且这穷寡妇比那些富人投得更多：\BibleQuote{这人，即使奉献一杯凉水，也不会失掉他的赏报}\BibleRef{玛 10:42}而且还有\BibleQuote{平安归于地上善意之人}\BibleRef{路 2:14}。\BibleQuote{应向天主祭献赞颂的祭品，并向至高者奉献你们的祈祷}。这馨香使主悦乐。\par

22. \BibleQuote{在你困苦之日，你要呼求我；我必拯救你，你也要光荣我。}（\BibleRef{咏 49:15}）因为你不可依靠自己的能力，你的一切援助都是虚假的。\Quote{你在困苦之日要呼求我；我必拯救你，你也要光荣我。}因为为此我允许困苦之日临到你：或许若你不受困扰，你就不会呼求我；但当你受困扰时，你就呼求我；当你呼求我时，我必拯救你；当我拯救你时，你就要光荣我，好使你不再离开我。有一个人，在祈祷的热忱中变得迟钝冷淡，他说：\Quote{我找到了困苦和忧愁，我呼求了主的名。}他找到了困苦，如同某种有益之物；他曾腐烂在罪恶的泥沼中；如今他继续毫无感觉，他发觉困苦仿佛是一种腐蚀和切割。\Quote{我找到了，}他说，\Quote{困苦和忧愁，我呼求了主的名。}的确，弟兄们，困苦是人所共知的。请看遍布人间的那些苦难：一人因损失而痛苦哀哭；另一人因丧亲而悲伤；又一人被流放异地，悲叹思念家乡，视寄居为难以忍受；另一人的葡萄园遭冰雹袭击，他看到自己的劳作和一切辛劳都付诸东流。人何时能不悲伤呢？他在朋友中找到了敌人。人世间还有比这更大的不幸吗？这些事所有人都哀叹悲伤，这些就是困苦：在这一切中，他们都呼求主，他们做得对。让他们呼求天主，祂能教导如何忍受，或在忍受时治愈它。祂知道不让我们受试探超过我们所能承担的。（\BibleRef{格前 10:13}）让我们甚至在那些困苦中也呼求天主；但这些困苦确实找上了我们；正如另一篇圣咏所写：\Quote{在过于找上我们的困苦中，你是帮助者。}有一种困苦是我们应当找到的。愿这样的困苦找到我们：有一种困苦是我们应当寻求并找到的。那是什么？前面提到的今世的幸福、暂时事物的丰裕：那其实不是困苦，而是我们困苦的慰藉。什么困苦？我们旅居的困苦。因为我们尚未与天主同在，我们生活在试炼和困难中，我们不能没有恐惧，这本身就是困苦：因为还没有那应许给我们的平安。凡在旅居中没有发现这种困苦的人，就不会想到返回祖国。这就是困苦，弟兄们。诚然，当我们分饼给饥饿的人，接待旅客等善行时，这本身也是困苦。因为我们会遇到可怜的对象，从而对他们施以怜悯；而那些可怜对象的悲惨境遇也使我们生出同情。在那个地方，你找不到饥饿的人可喂，找不到旅客可接待，找不到赤身的人可遮盖，找不到病人可探望，找不到争执的人可调解，那对你将是多么好啊！因为那地方的一切都是至高的、真实的、圣洁的、永恒的。我们在那里的食粮是正义，饮料是智慧，衣服是不朽，房屋是天上的永恒居所，我们的坚定是不朽：疾病会临到吗？疲惫会使人昏昏欲睡吗？没有死亡，没有诉讼：那里有平安、宁静、喜乐、正义。没有敌人能进入，没有朋友会离去。那里的宁静是怎样的呢？如果我们思想并留意我们现在在哪里，以及那不能说谎者曾应许我们将在哪里，从祂的应许本身，我们就发现自己正处在何种困苦中。这种困苦，除非那寻求它的人，没有人能发现。你健康时，看看你是否可怜；因为病人容易发现自己可怜；当你健康时，看看你是否可怜：因为你尚未与天主同在。\Quote{我找到了困苦和忧愁，我呼求了主的名。}因此，\Quote{献上，} \Quote{颂赞的祭献给天主。}赞美祂的应许，赞美祂的召唤，赞美祂的劝勉，赞美祂的帮助：并要知道你处在何种困苦中。呼求祂，你必被拯救，你必光荣祂，必存留。\par

23. 但请看下文，我的弟兄们。因为现在某个人或其他人，由于天主曾对他说：\Quote{向天主献上赞美的祭献}，并且以某种方式命令了这贡品，他就在心里思量并说：我要每日起来，我要去圣堂，我要在晨祷唱一首赞美诗，在晚祷唱一首，在家唱三四首，我每日献上赞美的祭献，并祭献给我的天主。你这样做确实很好；但要小心，免得你现在因这样做而疏忽了：也许你的舌头赞美天主，而你的生活却诅咒天主。我的子民，万神之神，那发言的主，对你说：\Quote{从日出之地到日落之处召唤大地}，虽然你仍处于莠子之中，\BibleRef{玛 13:25}\BibleQuote{向你的天主献上赞美的祭献，并向祂呈上你的祈祷}：但要小心，免得你生活邪恶，却歌唱得好。为何如此？因为\BibleQuote{天主对罪人说：你为什么传述我的律例，将我的盟约放在你口中？}\BibleRef{咏 49:16}。你们看，弟兄们，我们是何等战栗地说这些话。我们将天主的盟约放在口中，我们说这些话。我们将天主的盟约放在口中，并向你们宣讲天主的训诲和判断。天主对罪人说什么？\Quote{你为什么？}那么，祂禁止犯罪的宣讲者吗？那又怎么说：\BibleQuote{他们所说的，你们要遵守；但他们所做的，你们不要学}？\BibleRef{玛 23:3}还有：\BibleQuote{无论出于真心还是假意，基督总被宣讲}？\BibleRef{斐 1:18}但这些话是说给听的人的，免得他们害怕从谁那里听到（无论从谁那里听到），而不是说给说好话却行恶事的人，让他们可以无忧无虑地生活。所以现在，弟兄们，你们无忧无虑：如果你们听到好话，你们就是听到天主，无论你们是从谁那里听到。但天主不会让那些说话的人不受责备：免得他们只凭说话，不顾自己，在邪恶的生活中沉睡，并对自己说：\Quote{因为天主不会把我们送到灭亡，虽然祂愿意通过我们的口对祂的子民说出这么多好话。}不，你要听你自己所说的话，无论你是谁，你说话的人：你希望别人听到你，你先听你自己；并且说出某人在另一篇圣咏中所说的话：\Quote{我要听在我内说话的上主天主，因为祂要向祂的子民讲和平。}那么，我是什么人，竟不听祂在我内所说的话，却希望别人听祂借我所说的话？我先听，我要听，并且首先我要听在我内说话的上主天主，因为祂要向祂的子民讲和平。让我听，并\BibleQuote{制伏我的身体，使它服从，免得我给别人宣讲，自己反而被淘汰。}\BibleRef{格前 9:27}\Quote{你为什么传述我的律例？}为什么传述对你无益的事？祂劝诫他听：不是放弃宣讲，而是承担服从。\Quote{但你，为什么将我的盟约放在你口中？}\par

24. \BibleQuote{但你憎恶管教}\BibleRef{咏 49:17}。你憎恶纪律。当我宽容时，你歌唱赞美我；当我责罚时，你抱怨，好像我宽容时我是你的天主，责罚时我就不是你的天主。\BibleQuote{凡我所爱的，我就责备和管教。}\BibleRef{默 3:19}\Quote{但你憎恶管教，将我的话抛在脑后。}那借你所说的话，你抛在脑后。\Quote{你将我的话抛在脑后：}丢到一个你看不见的地方，但却重压你。\Quote{你将我的话抛在脑后。}\par

25. \BibleQuote{你若见了贼，便与他同谋，又与奸夫分赃}\BibleRef{咏 49:18}。免得你或许说：我没有偷窃，我没有行奸淫。但是，若你喜悦那行奸淫的人，难道你因那喜悦不是与他同谋吗？难道你不是因赞同而与他分赃吗？因为，弟兄们，这就是与贼同谋，与奸夫分赃：即使你没有行，但若你赞许所行的，你就是同犯；因为\Quote{罪人因他灵魂的欲望而受赞美，行不义的人必受称许。}你没有行恶事，你却赞美恶人。难道这是小恶吗？\Quote{你与奸夫分赃。}\par

26. \BibleQuote{你的口充满恶意，你的舌头编造诡诈}\BibleRef{咏 49:19}。弟兄们，这是说某些人的恶意和诡诈，他们通过奉承，虽然知道所听到的是邪恶，但为了避免得罪说话的人，他们不仅不责备，反而以沉默表示赞同。这还不够，他们不说\Quote{你做了坏事}，反而说\Quote{你做得很好}，尽管他们知道那是邪恶：但他们的口充满恶意，他们的舌头编造诡诈。诡诈是一种言语上的欺骗，口是心非。祂没有说\Quote{你的舌头做出了诡诈}或\Quote{施行了诡诈}，而是为向你指出对作恶本身的一种快感，祂说\Quote{编造}。你作恶还不够，你还以此为乐；你公开赞美，内心暗笑。你使一个轻率暴露自己过失、不知道那是过失的人走向毁灭：你知道那是过失，却不说\Quote{你要往哪里去？}如果你看见他在黑暗中轻率行走，你知道那里有一口井，却保持沉默，你会是什么样的人？难道你不被视为他生命的敌人吗？然而如果他掉进井里，不是灵魂而是身体会死。他确实坠入恶习，在你面前暴露他的恶行：你知道那是恶行，却赞美并暗笑。哦，但愿他最终归向天主，那个被你嘲笑、你不愿责备的人，并且他说：\Quote{愿那些对我说\InnerQuote{好，好}的人蒙羞。}\par

27. \BibleQuote{你坐在你兄弟的对面，你曾诋毁他} \BibleRef{咏 49:20}。而这个 \Quote{坐} 正是属于他上面所说的 \Quote{拥抱}。因为人站立或路过时所做的事，并非出于乐意；但若他为此目的而坐，那该多么寻求闲暇去行此事！你正以勤勉行这邪恶的诋毁，你是坐着做的；你因此全神贯注；你是在拥抱你的邪恶，你在亲吻你的诡诈。\Quote{你又在你母亲的儿子面前设下了绊脚石。} 谁是 \Quote{母亲的儿子}？难道不是兄弟吗？那么他是在重复他上面所说的 \Quote{你的兄弟}。他是否暗示我们必须看出某种区别？显然，弟兄们，我认为必须区分。兄弟诋毁兄弟，例如，一个坚强、如今是博士和颇有分量的学者，诋毁他的兄弟，也许是一个善教善行的人；但另一个是软弱的，他因诋毁前者而为他设下绊脚石。因为当好人被那些看似有分量、有学识的人诋毁时，软弱的人就会绊倒在这绊脚石上，他们还不知道如何判断。因此，这个软弱的人被称为 \Quote{母亲的儿子}，还不是父亲的，仍然需要奶，依恋于乳房。他至今仍被他母亲教会的胸怀所怀抱，他不够坚强，不能接近他父亲餐桌上的固体食物，而是从母亲的乳房汲取营养，不善于判断，因为他仍是属血气的和属血肉的。\Quote{因为属神的人判断一切，} \BibleRef{格前 2:15} 但 \Quote{属血气的人不接受天主圣神的事，因为在他看来是愚妄。} \BibleRef{格前 2:14} 宗徒对这些人说：\BibleQuote{我不能向你们说话如同对属神的人，只能如同对属血肉的人，如同对在基督内的婴孩，我给你们喝奶，不是吃饭；因为你们还不能，就是现在还是不能。} \BibleRef{格前 3:1} 我曾作过你们的母亲：如在另一处所说：\BibleQuote{我在你们中间成了婴孩，如同乳母抚育自己的子女。} \BibleRef{得前 2:7} 不是乳母哺育别人的孩子，而是乳母抚育自己的孩子。因为有些母亲生养后交给乳母：那些生养的人并不抚育自己的孩子，因为他们交给别人乳养；但那些抚育的人，抚育的不是自己的孩子，而是别人的：但他自己生了，他自己抚育，没有将所生的托付给任何乳母；因为他说：\BibleQuote{我为你们再受产痛，直到基督在你们内形成。} \BibleRef{迦 4:19} 他抚育了他们，并给了奶。但有些似乎有学识和属神的人诋毁保禄。\BibleQuote{他的书信，他们说，是沉重有力的；但他亲临的软弱，言语是可轻视的。} \BibleRef{格后 10:10} 他本人在书信中说，某些诋毁他的人说了这些话。他们坐着，诋毁他们的兄弟，并为他们母亲的儿子（即那些要靠奶喂养的）设下了绊脚石。\BibleQuote{你又在你母亲的儿子面前设下了绊脚石。}\par

28. \BibleQuote{你行了这些事，我闭口不言} \BibleRef{咏 49:21}。因此，主我们的天主必要来临，且不再缄默。现在，\Quote{你行了这些事，我闭口不言。} 什么是 \Quote{我闭口不言}？我从报复中退让，我推迟了我的严厉，我延长了对你的忍耐，我长久等待你的悔改……\Quote{你曾设想不义，以为我像你一样；} 你曾设想我像你一样，而你不会像我。因为，他说：\BibleQuote{你们应当是成全的，如同你们的天父是成全的；他使太阳上升，光照恶人，也光照善人。} \BibleRef{玛 5:48} 你不肯效法他，他连恶人也赐予善事，而你呢，竟坐着诋毁善人！\BibleQuote{我要责斥你，} 当 \BibleQuote{天主显现来临，我们的天主，且不再缄默，} 时，\BibleQuote{我要责斥你。} 我责斥你时要对你做什么？要对你做什么？现在你不见自己，我要使你看见自己。因为如果你看见自己，而对自己不满，你就会使我喜悦；但因为你不见自己而自满，你必使我与你自己都不悦：在我审判你时使我不悦，在你焚烧时使你自己不悦。但我要对你做什么？他说。\Quote{我要把你的面摆在你面前。} 因为为什么你要逃避自己？你对自己是在背后，你不见自己：我使你看见自己：你把放在背后的，我将放在你面前；你将看见你的污秽，不是要你改正，而是要你羞愧……\par

29. 但是，\BibleQuote{你们这些忘记天主的，要明白这事} \BibleRef{咏 49:22}。看，他如何呼喊，并不缄默，不宽容。你忘记了天主，没有思想你的邪恶生活。要明白你是如何忘记了天主。\Quote{免得他终久像狮子一般撕碎，无人搭救。} 什么是 \Quote{像狮子一般}？像骁勇者，像大能者，像无人能抵抗者。他说 \Quote{狮子} 时就是指这一点。因为狮子可用于赞美，也可用于显示邪恶。魔鬼曾被称作狮子：\BibleQuote{你们的仇敌，}他说，\BibleQuote{如同咆哮的狮子巡游，寻找可吞食的人。} \BibleRef{伯前 5:8} 难道不是这样吗？魔鬼因凶猛残暴被称为狮子，而基督因神奇的大能被称为狮子？那里说：\BibleQuote{犹大支派的狮子已得胜？} \BibleRef{默 5:5}……\par

30. \Quote{赞颂之祭必要光荣我}\BibleRef{咏 49:23}。\Quote{赞颂之祭如何光荣我？}诚然，赞颂之祭绝不能有益于恶人，因为他们口中称扬你的盟约，却行你所不喜悦的可憎之事。他说，我立刻也对那些人说：\Quote{赞颂之祭必要光荣我。}因为如果你生活恶劣却口出善言，你尚未赞美；但另一方面，如果你开始活得好，却将行善归功于自己的功德，你仍未赞美……因此，税吏下去，成了正义的，而非法利塞人。所以，你们这些活得好的人听，你们这些活得坏的人也听：\Quote{赞颂之祭必要光荣我。}无人向我献此祭而仍是恶人。我不说：不可有恶人向我献此祭；而是说，献此祭的人，没有一个仍是恶人。因为赞美的人，是善的：因为他若赞美，他也行得善，因为若他赞美，不仅用舌头赞美，生活也与舌头一致。\par

31. \Quote{这就是那条路，我要以此向他显示天主的救恩。}在赞颂之祭中\Quote{就是那条路。}什么是\Quote{天主的救恩}？耶稣基督。而赞颂之祭如何向我们显示基督？因为基督带着恩宠来到我们这里。宗徒说了这话：\BibleQuote{我生活，已不是我生活，而是基督在我内生活；我现今在肉身内生活，是生活在对天主子的信德内，他爱了我，且为我舍弃了自己。}\BibleRef{迦 2:20}那么，罪人们应当承认，若是健康的，就不需要医生。\BibleRef{玛 9:12}因为基督为不虔敬的人死了。\BibleRef{罗 5:6}当他们承认自己的不虔敬，并首先效法那个税吏，说：\BibleQuote{主，可怜我这个罪人吧。}\BibleRef{路 18:13}展示伤口，恳求医师；并且由于他们不赞美自己，反而责备自己——\BibleQuote{使那夸耀的，不在自己内夸耀，而只夸耀主。}\BibleRef{格前 1:31}——他们承认基督来临的原因，因为他正是为此而来，为拯救罪人：因为\BibleQuote{耶稣基督来到这世界，是为拯救罪人；其中我是元凶。}\BibleRef{弟前 1:15}再者，那些犹太人夸耀自己的作为，宗徒正是这样斥责他们，说他们不属于恩宠，因为他们认为奖赏是应归于自己的功劳和善行。\BibleRef{迦 5:4}因此，知道自己属于恩宠的人，就知道基督是什么、基督的什么，因为他需要恩宠。恩宠之所以被称为恩宠，是\OriginalTerm{白白地}{gratis}赐予的；若是\OriginalTerm{白白地}{gratis}赐予，就没有任何先前的时期功德该当领受它……\par

\SourceRecord{Translated by J.E. Tweed.；Nicene and Post-Nicene Fathers, First Series；Vol. 8.；Edited by Philip Schaff.；Buffalo, NY: Christian Literature Publishing Co.,；1888.；<http://www.newadvent.org/fathers/1801050.htm>.}

% structure: p0001 source=h1 target=section reason=page-title

\section{圣咏第五十一篇释义}

1. 不可欺骗这群众的拥挤，也不可加重他们的软弱。我们请求安静和宁静，以便我们的声音在昨日的劳作之后，能稍有余力持续下去。应该相信，你们的爱德今天聚集得比往常更多，别无他事，只是为了你们能为那些被异己和乖僻的倾向所阻隔的人祈祷。因为我们既不是在谈论外邦人，也不是在谈论犹太人，而是在谈论基督徒；也不是在谈论那些尚是慕道者的人，而是在谈论许多甚至已受洗的人——你们与他们的圣洗池毫无不同，然而与他们的心却不相像。因为如今有多少我们的弟兄，我们想起他们，就哀叹他们走向虚妄和谎言的疯狂，而忽视了他们所受的召唤。他们若在竞技场中因任何原因偶然受惊，便立即画十字，并在额头上带着它站立——就在那地方，他们本已从中退去，若他们曾在心中带着它。必须恳求天主的仁慈，好赐予理解以谴责这些事，赐予逃避的意愿，并赐予宽恕的仁慈。幸好，今天咏唱了一篇关于补赎的圣咏。我们也对那些缺席者说话：对于他们，我们的声音就是你们的记忆。不要忽视受伤和软弱的人，但为了更易治愈，你们应当保持健全。通过责备纠正，通过讲话安慰，通过善行树立榜样。那曾与你们同在的，也将与他们同在。既然你们已越过这些危险，天主的仁慈之泉并未关闭。你们所到之处，他们也将来到；你们所经过之处，他们也将经过。实在是一件可悲的事，且极其危险，甚至是毁灭性的，并且确定是致命的事，那就是明知故犯。因为走向这些虚妄的人，一种是轻视基督之声的，另一种是知道自己从何逃离的。但是连这样的人我们也不该绝望，这篇圣咏就表明了这一点。\par

2. 因为其上写着它的标题：\Quote{达味自己的圣咏，那时纳堂先知来到他那里，当他进去见巴特舍巴的时候。}巴特舍巴是一个妇人，是另一个人的妻子。我们说话时确实带着悲痛和战栗，但天主不愿意那祂所愿意被写下的被沉默。我就要说，不是出于我的意愿，而是我不得不说的；我不是作为劝人效仿者，而是作为教导你们恐惧者。这位君王和先知达味，那将要按血统从其后裔而来的主\BibleRef{罗 1:3}，被这妇人的美貌所迷惑，与这别人的妻子犯了奸淫。这件事在这篇圣咏中没有被读到，但在其标题中出现；但在列王纪\BibleRef{撒下 11:2}中有更详细的叙述。两部圣经都是正经，基督徒必须毫无疑问地相信两者。犯罪已发生，并被记录下来。此外，他使她的丈夫在战争中被杀。在这事之后，主派遣先知纳堂到他那里\BibleRef{撒下 12:1}，去斥责如此重大的恶行。\par

3. 人应谨慎之事，我们已说过；但若跌倒，应效法什么，让我们听。因为许多人愿意与达味一同跌倒，却不愿与达味一同起来。故此，榜样并非为跌倒而设，而是为若跌倒，便得以复起。谨慎，以免跌倒。愿幼者不以长者的失足为乐，而愿长者的跌倒成为幼者的警戒。为此它被陈列，为此被写下，为此在教会中常被诵读和咏唱：愿未曾跌倒的人听见，以免跌倒；愿已跌倒的人听见，得以复起。如此伟大人物的罪过未被遮掩，却在教会中被宣告。那里有不良的听众，他们为自己犯罪寻找借口：他们寻求如何为自己已预备要犯的罪辩护，而非如何警戒自己未曾犯的罪；他们对自己说：若达味如此，为何我不如此？那灵魂更为不义，因它因达味做了而做，故比达味做得更恶。我要尽量更明白地说这同一件事。达味未曾为自己立下任何先例，如你一样：他因贪欲的失足而跌倒，并非因圣德的支持；你却将一位圣人置于眼前，以便犯罪：你不效法他的圣德，却效法他的跌倒。你爱达味身上达味自己所恨恶的；你预备犯罪，倾向犯罪：为了犯罪，你查阅天主的书；你听天主的圣经，为行天主所不喜悦的事。达味并未如此行；他受先知谴责，未曾因先知而绊倒。但其他人健康地听见，藉强者的跌倒衡量自己的软弱；并渴望避免天主所定罪的事，因而谨慎地约束自己的眼目。他们不将眼目固定在他人的肉体之美上，也不以乖谬的单纯而疏忽；他们不说：\Quote{我善意地观看，我善良地观看，我因爱德长久注视。}因为他们将达味的跌倒摆在眼前，并看见这位伟人为此目的而跌倒，为叫小人们不愿观看那能令他们跌倒的事。他们便约束眼目免于放纵，不轻易与人结交，不与陌生女子混杂，不将迎合的眼目抬向陌生的阳台、陌生的屋顶。因为达味从远处看见了她，那使他被迷住的女人。\BibleRef{撒下 11:2} 女子遥远，情欲近在咫尺。他所见的在别处，而使他跌倒的在他自己内。因此必须记住这肉身的软弱，回忆宗徒的话：\BibleQuote{所以不要让罪在你们必死的身体上为王。}\BibleRef{罗 6:12} 他没有说\BibleQuote{不要有罪}，而是说\BibleQuote{不要为王}。在你内有罪，当你喜悦时；它作王，若你同意了。肉体的快乐，尤其若走向不法与陌生之物，当被勒住，而非放松；当被治理驯服，而非被立为治理者。观看并无忧虑吧，若你没有任何可被激动的事。但你回答说：\Quote{我以坚强的决意克制。}你难道比达味更强吗？\par

4. 再者，他以这样的榜样劝诫，叫人在顺境中不可自高。因为许多人害怕逆境，却不害怕顺境。顺境对灵魂比逆境对肉体更为危险。首先，顺境使人腐化，好让逆境寻得可击碎之物。我的弟兄们，必须更加警惕福乐。因此，请看在我们自己的福乐中，天主的话如何夺去我们的安全感：祂说：\Quote{你们当以敬畏事奉主，又当以战兢向祂欢跃。}在欢跃中，为叫我们感恩；在战兢中，以免我们跌倒。达味在遭受撒乌耳迫害时，未曾犯这罪。当圣达味遭受仇敌撒乌耳之苦，当他被迫害所扰，当他四处逃窜以免落入其手时，他并未贪恋他人之妻，也未在犯奸后杀害丈夫。他在困苦的软弱中，愈显可怜，却愈与天主亲密。患难是有益的；外科医生的手术刀比魔鬼的诱惑更有益。当他的仇敌被推翻，压力解除，肿胀膨胀，他便变得自满。故此，这榜样有助于使我们惧怕福乐。他说：\Quote{我找到了患难和忧愁，并呼求了主的名。}\par

5. 然而这事已经发生了；我要向那些没有做过同样事的人说这些话，好使他们警醒，保持自己的纯洁，并当他们看到一位伟人如何跌倒时，卑微者应感到恐惧。但若有人已经跌倒，听到这些话，且他的良心中有任何邪恶之事，让他留意这篇圣咏的话；让他注意伤口的严重性，但不要对医师的伟大丧失希望。绝望中的罪是必死的。因此，没有人应该说：\Quote{若我已做了任何恶事，我便已被定罪；天主不赦免这样的恶事，我何不罪上加罪呢？我将在享乐、放荡、邪恶的贪欲中享受这个世界；既然改过的希望已经失去，至少让我拥有我所看见的，若我无法拥有我所相信的。} 因此，这篇圣咏使那些尚未相信的人警醒，同样也不愿让那些跌倒的人绝望。无论你是谁，曾犯了罪，却因绝望于自己的得救而犹豫要不要为你的罪作补赎，请听\Person{达味}的呻吟。\Person{纳堂}先知没有被派到你那里，但\Person{达味}本人已被派给你。听他呼喊，并与他一同呼喊；听他呻吟，并与他一同呻吟；听他哭泣，并与他同流眼泪；听他悔改，并与他一同喜乐。若罪不能从你身上被排除，也不要排除赦免的希望。当时\Person{纳堂}先知被派到那人那里，请注意这位君王的谦卑。他没有拒绝那劝诫者的话，他没有说：\Quote{你敢对我这君王说话吗？} 一位被尊崇的君王听了先知的话，愿祂卑微的子民听\Person{基督}的话。\par

6. 所以，请听这些话，并与他一同说：\Quote{天主，求祢按照祢的浩大仁慈，怜悯我。}\BibleRef{咏 50:1} 那恳求浩大仁慈的，承认巨大的痛苦。让那些因无知而犯罪的人向祢寻求一点仁慈：他说：\Quote{怜悯，} \Quote{我，按照祢的浩大仁慈。} 按照祢伟大的医治解除深重的创伤。我的创伤是深的，但我投靠于全能者。对我自己如此致命的伤口，我本应绝望，除非我能找到如此伟大的医师。\Quote{天主，求祢按照祢的浩大仁慈怜悯我；并按照祢丰厚的怜悯，涂抹我的罪愆。} 他所说的\Quote{涂抹我的罪愆}，就是\Quote{天主，求祢怜悯我。}而他所说的\Quote{按照祢丰厚的怜悯}，就是\Quote{按照祢的浩大仁慈。} 因为仁慈是浩大的，怜悯是众多的；因祢的浩大仁慈，祢的怜悯众多。祢垂顾嘲弄者以纠正他们，垂顾无知者以教导他们，垂顾忏悔者以赦免他们。他是无知中做了这事吗？有一个人做了某些，甚至许多恶事；他说：\Quote{怜悯，} \Quote{我获得了，因为我在不信中无知地做了这事。}\BibleRef{弟前 1:13} \Person{达味}不能这样说：\Quote{我是无知地做了这事。}因为他并非不知道触摸他人妻子是多么邪恶的事，也不知道杀害那毫不知情、甚至没有发怒的丈夫是多么邪恶的事。因此，那些无知中做了这事的人获得了\Person{主}的仁慈；而那些明知故犯的人，所得的或许并非任何一般的仁慈，而是\Quote{浩大的仁慈。}\par

7. \Quote{求祢多多洗净我的不义。}\BibleRef{咏 50:2} 什么是\Quote{多多洗净}？一个沾满污秽的人。多多洗净一个明知故犯者的罪。祢已经洗净了无知之人的罪。即便如此，也不应绝望于祢的仁慈。\Quote{并求祢清除我的过犯。} 按照祂是医师的方式，献上补偿。祂是天主，献上祭献。你能给出什么使你得以洁除？因为看哪，你呼求的是谁，你呼求的是公义者。祂恨恶罪恶，若祂是公义的；祂惩罚罪恶，若祂是公义的；你不能从\Person{主天主}那里夺去祂的公义：恳求仁慈，但要注意公义：有仁慈以赦免罪人，有公义以惩罚罪恶。那么呢？你求仁慈，难道罪可以不罚而安好吗？让\Person{达味}回答，让那些跌倒的人回答，与\Person{达味}一同回答说：\Quote{不，主，我的罪没有一个是未受惩罚的；我知道我所求仁慈的那位的公义：它不会不受惩罚，但为此我不愿祢惩罚我，因为我惩罚自己的罪；为此我恳求祢赦免，因为我承认。}\par

8. \Quote{因为我承认我的罪愆，我的过犯常在我面前。}\BibleRef{咏 50:3} 我没有把所做的事抛在背后，我不看顾别人而忘记自己，我不假装要拔掉弟兄眼中的草屑，而自己眼中却有梁木；\BibleRef{玛 7:5} 我的罪在我面前，不在我背后。因为当先知被派到我这里，把穷人的绵羊的比喻摆在我面前时，我的罪在我背后。因为\Person{纳堂}先知对\Person{达味}说：\Quote{有一个富人，拥有许多羊；但他的邻居穷人只有一只小母羊，他把它抱在怀里，用自己的食物喂养它；有一个旅客来到富人那里，他没有从自己的羊群中取羊，却贪图邻居穷人的那只小母羊，把它宰了给旅客：这人该受什么处罚？} 但那人发怒并宣判：那时，君王显然不知道自己在何处被逮住，宣称那富人应处死刑，并应四倍偿还那羊。极其严厉，极其公正。但他的罪尚不在他面前，他所做的在他背后；他尚未承认自己的罪愆，因此也不宽恕别人的。但是，为此被派来的先知，从他的背上取下那罪，放在他眼前，使他看到那如此严厉的判决正是针对他自己宣判的。为切开并医治他心上的伤口，他用自己的舌头作了解剖刀……\par

9. \Quote{我唯独得罪了祢，在祢面前行了恶}\BibleRef{咏 50:4}。这是什么？因为不是在人前玷污了他人的妻子并杀了她的丈夫吗？众人岂不知道达味所做的事吗？\BibleRef{撒下 11:4}什么是\Quote{我唯独得罪了祢，在祢面前行了恶}？因为唯有祢是无罪的。祂是正义的惩罚者，在祂内没有什么应受惩罚；祂是正义的谴责者，在祂内没有什么应受谴责。\BibleQuote{为叫祢在祢的言语上显示为义，并在祢受审判时得胜。}他向谁说话，弟兄们，他向谁说话，难以理解。他固然向天主说话，而显然天主父是不受审判的。\BibleQuote{并在祢受审判时得胜}是什么意思？祂预见那将受审判的审判者，一位义人被罪人审判，并在此得胜，因为在祂内没有什么应受审判。因为唯有在人间，那\Person{天主而人}能真正地说：\BibleQuote{你们中间谁能指证我有罪？}\BibleRef{若 8:46}但或许有某种事物逃过了人的眼目，他们并未找到那真正在那里却未显露的东西。在另一处\BibleRef{若 14:30}，祂说：\BibleQuote{看，世界的首领来了}——祂明察一切罪过——祂说：\BibleQuote{看，这世界的首领来了}，以死亡折磨罪人，掌管死亡；因为\BibleQuote{因魔鬼的嫉妒，死亡进入了世界}\BibleRef{智 2:24}祂说：\BibleQuote{看}，\BibleQuote{世界的首领来了}——祂说这话时临近受难：——\BibleQuote{在我身上他一无所获}，没有任何罪，没有任何该受死亡或该受定罪的。仿佛有人问祂：那么祢为何要死？祂接着说：\BibleQuote{但为叫世人知道我爱父，我怎样受命，我就怎样行；起来，我们走吧。}祂说，我并不当受，却为人当受的而受苦，好使他们能配得我的生命，而我是白白地为他们忍受死亡。于是，那毫无罪过的祂，达味先知在此刻说：\BibleQuote{我唯独得罪了祢，在祢面前行了恶，为叫祢在祢的言语上显示为义，并在祢受审判时得胜。}因为祢胜过所有人、所有审判者；凡自以为义的人，在祢面前都是不义的：唯独祢公正地审判，却受了不公正的审判；祢有权舍掉祢的生命，也有权再取回它。\BibleRef{若 10:18}祢在受审判时得胜。祢胜过所有人，因为祢超越人，是祢造了人。\par

10. \BibleQuote{看，我是在罪恶中受孕的}\BibleRef{咏 50:5}。好像他说：那些做了你所做之事的人，达味，都失败了；因为这不是微小的恶和小罪，即奸淫和杀人。那么，那些从母腹出生之日起从未做过这些事的人呢？就连他们，你也归咎一些罪过，好使祂开始受审判时胜过所有人。达味承担了人类的本性，看到了所有人的束缚，思考了死亡的后果，察觉了不义的起源，他说：\BibleQuote{看，我是在罪恶中受孕的。}难道达味是奸淫所生？他生于\Person{叶瑟}\BibleRef{撒上 16:18}，一个义人，和他的妻子？他说自己是在罪恶中受孕的，意思难道不是不义是从\Person{亚当}而来的吗？就连死亡的束缚本身，也随不义一同根植？没有人生来不带惩罚，不带应受的惩罚。另一位先知在别处也说：\Quote{在祢面前无人是洁净的，即使婴儿，在世只活了一天。}因为我们知道，借着基督的圣洗圣事，罪得以赦免，基督的圣洗圣事能赦罪。如果婴儿全然无罪，为什么母亲们在孩子生病时抱着他们跑向教会？那圣洗圣事，那赦免，除去了什么？我看见一个无辜者，他哭泣多于发怒。圣洗圣事洗去了什么？那恩宠解除了什么？解除了罪的产物。因为如果那个婴儿能对你说话，并且有达味那样的理解力，他会回答你：你为什么在意我，一个婴儿？你确实看不见我的行为，但我在罪恶中受孕，\BibleQuote{我的母亲在罪恶中孕育了我。}\par

除了这\\OriginalTerm{有死的情欲的束缚}{bond of mortal concupiscence}之外，基督没有经由男子，由童贞女因圣神受孕而生。不能说祂是在罪孽中受孕的，不能说：祂的母亲在罪过中将祂滋养在胎中，对那一位曾说过：\\BibleQuote{圣神将临于你，至高者的能力要庇荫你。}\\BibleRef{路 1:35} 因此，人之所以在罪孽中受孕、在罪过中被母亲滋养于胎中，并非因为与妻子交合是罪，而是因为那所生的确实是由应受惩罚的肉体所造。因为肉体的惩罚是死亡，其中确实有对死亡本身的承受力。因此宗徒没有谈论身体如同将要死，而是如同已死：他说：\\BibleQuote{身体固然因罪而死亡，但神魂却因义而生活。}\\BibleRef{罗 8:10} 那么，那由因罪而死的身体所孕育和播种的，怎能没有罪的束缚而生呢？在已婚者中，这种贞洁的行为没有罪，但罪的根源却带来\\OriginalTerm{应得的惩罚}{condign punishment}。因为没有哪个丈夫，因为他是丈夫，就不受死亡的约束，也没有任何其他原因使他受死亡的约束，除非是因为罪。甚至主也受死亡的约束，但不是因罪：祂承受了我们的惩罚，从而解除了我们的罪责。因此合理地：\\BibleQuote{在亚当内众人都死了，但在基督内众人都要复活。}\\BibleRef{格前 15:22} 因为宗徒说：\\BibleQuote{罪恶藉着一人进入了世界，死亡藉着罪恶也进入了世界，这样死亡就传给了众人，因为众人都犯了罪。}\\BibleRef{罗 5:12} 这是明确的论断：\\BibleQuote{在亚当内，}他说，\\BibleQuote{众人都犯了罪。}因此，只有那从未从亚当的作为所生的婴儿才可能是无辜的。\par

11. \\BibleQuote{因为，看，祢喜爱真理：祢将祢智慧的\\OriginalTerm{不确定和隐秘的事}{uncertain and hidden things}启示给我。}\\BibleRef{咏 50:6} 即，祢没有让那些祢所宽恕之人的罪不受惩罚。\\BibleQuote{祢喜爱真理：}因此祢先赐予了慈悲，以便祢也应保存真理。祢宽恕一个告明者，祢宽恕，但只当他惩罚自己：如此慈悲和真理都被保存：慈悲因为人被释放；真理因为罪被惩罚。\\BibleQuote{祢将祢智慧的\\OriginalTerm{不确定和隐秘的事}{uncertain and hidden things}启示给我。}什么是\\Quote{隐秘的事}？什么是\\Quote{不确定的事}？因为天主甚至宽恕这样的人。没有什么如此隐秘，没有什么如此不确定。因这不确定，尼尼微人忏悔了，因为他们虽然经过先知的警告，虽然经过那呼声：\\BibleQuote{三天之后，\\Place{尼尼微}将要倾覆。}\\BibleRef{纳 3:4} 他们对自己说，必须恳求慈悲；他们这样彼此推理说：\\BibleQuote{谁知道天主是否会改变祂的判决，而怜悯我们呢？}\\BibleRef{纳 3:9} 那是\\Quote{不确定的}，当人们说\\Quote{谁知道？}时，他们因不确定而忏悔，却赢得了确定的慈悲：他们俯伏在泪水中，在禁食中，在麻衣和灰烬中俯伏，呻吟，哭泣，天主宽恕了。\\Place{尼尼微}站立：\\Place{尼尼微}是否倾覆了？在人看来是一回事，在天主看来是另一回事。但我认为先知所预言的已经实现了。看看\\Place{尼尼微}是什么，看看它是如何倾覆的；在邪恶中倾覆，在良善中建立；正如迫害者扫罗被倾覆，宣讲者保禄被建立（\\BibleRef{宗 9:4}）。谁不会说，我们现在所在的这座城市，如果所有那些疯狂的人，放弃他们的琐事，带着痛悔的心一起跑到教会，并呼求天主的慈悲宽恕他们过去的行为，那真是幸福的倾覆呢？我们难道不会说，那个\\Place{迦太基}在哪里？因为那里没有了从前有的，它倾覆了；但若有了从前没有的，它建立了。因此对耶肋米亚说：\\BibleQuote{看，我指派你为要拔除、拆毁、毁灭、推翻，又要建设、栽植。}\\BibleRef{耶 1:10} 因此有主的声音：\\BibleQuote{我打击，我也医治。}\\BibleRef{申 32:39} 祂打击行为的腐朽，祂医治伤口的疼痛。医生们切割时就是这样做的；他们打击，也医治；他们拿起武器来击打，他们携带手术刀，前来治疗。但因为尼尼微人的罪过很大，他们说了\\Quote{谁知道？}这不确定的事，天主启示给了祂的仆人达味。因为当他在先知面前站着并受谴责时，他说了\\BibleQuote{我犯了罪。}立刻他从先知那里——即从先知内的天主圣神那里——听到：\\BibleQuote{你的罪已从你身上除去。}\\BibleRef{撒下 12:13} \\Quote{不确定和隐秘的事}，祂的智慧启示给了他。\par

12. \Quote{祢要洒我，}他说，\Quote{用牛膝草，我便洁净。}\BibleRef{咏 50:7} 我们知道牛膝草是一种卑微却具疗效的草药；据说其根部紧附于\OriginalTerm{盘石}{Rock}。因此，在奥秘中，它被用作洁净心灵的象征。你也要以你爱德的根紧紧抓住你的盘石；在你卑微的天主面前谦卑，以便在你受荣耀的天主面前被高举。你将被牛膝草洒过，基督的谦卑将洁净你。不要轻视这草药，要留意药效。我还要说一些我们从医生那里常听到、或在病人身上经历的事：他们说牛膝草适用于清肺。肺常被联想到骄傲，因为肺有膨胀，有呼吸。关于迫害者\Person{扫禄}，即骄傲的扫禄，经上说他要捆绑基督徒，\Quote{口吐杀气}\BibleRef{宗 9:1}；他口吐杀气，口吐血腥，他的肺还未洁净。也在此处听一位因牛膝草洁净而谦卑的人说：\Quote{祢洗涤我，}即洁净我：\Quote{我就比雪更白。}\Quote{虽然，}他说，\Quote{你们的罪虽似朱红，我要使之洁白如雪。}\BibleRef{依 1:18} 基督从这样的人中为自己预备了\Quote{没有瑕疵或皱纹的}衣裳。\BibleRef{弗 5:27} 此外，祂在山上发光如雪的衣裳，\BibleRef{玛 17:2} 象征着教会从一切罪污中被洁净。\par

13. 但牛膝草的谦卑何在？听听下文：\Quote{祢使我听到欢欣与喜乐，使谦卑的骨骸欢跃。}\BibleRef{咏 50:8} 我将在听从祢中喜乐，而不是在反对祢中喜乐。你犯了罪，为何为自己辩护？你要说话：忍受吧；听，顺服于天主的话语，免得你受窘迫，受更重的伤。罪已犯下，切勿辩护；让它归于\OriginalTerm{告解}{confession}，而非辩护。你以你罪之辩护者自居，你被打败了：你聘请了一位无罪的代理人，你的辩护对你不利。因为你是谁，竟为自己辩护？你当控告自己。不要说：\Quote{我什么也没做；}或\Quote{我做了多么大的事？}或\Quote{别人也做了。}如果你在犯罪时说没做什么，你将一无所是，你将一无所得：天主准备施恩，你却把自己关在门外；祂准备给予，不要以辩护的栏杆阻挡，而要敞开告解的怀抱。\Quote{祢使我听到欢欣与喜乐。}……\par

14. \Quote{求祢转面不看我的罪，并抹去我的一切过犯。}\BibleRef{咏 50:9} 如今谦卑的骨骸欢跃，如今被牛膝草洁净，我变得谦卑。\Quote{转面不看，}不是不看我自己，而是\Quote{不看我的罪。}因为他在另一处祈祷时说：\Quote{求祢不要转面不看我自己。}他不愿天主的面转离他，却愿天主的面转离他的罪。因为当天主不转面时，祂\OriginalTerm{注意}{adverts}罪；若祂注意，便\OriginalTerm{惩罚}{animadverts}。\Quote{并抹去我的一切过犯。}他正忙于那首要的罪；他提到更多，他愿自己的一切过犯都被抹去；他依靠医师的手，依靠那在诗篇开头他所呼求的\Quote{大仁慈}：\Quote{抹去我的一切过犯。}天主转面不看，如此便抹去；藉\Quote{转面不看}，祂抹去罪过。藉\Quote{转面看见}，祂写下它们。你已听见祂转面抹去，现在听祂转面看见，祂做什么？\Quote{上主的面容却注视作恶的人，为从地上消除对他们的纪念：}祂将消除对他们的纪念，但不是藉\Quote{抹去他们的罪过}。但这里他求什么？\Quote{求祢转面不看我的罪。}他求得好。因为他自己并不转面不看\Emph{自己的}罪，他说：\Quote{因为我承认了我的罪。}你有理由求，且求得好，即愿天主转面不看你的罪，如果你也不从那里转面不看；但如果你将你的罪置于背后，天主就在那里安置祂的面容。你若愿天主转面不看罪，就当将罪置于你的面前；然后你方能安心祈求，祂也垂听。\par

15. \Quote{天主，求祢在我内造一颗洁净的心。} \BibleRef{咏 50:10} \Quote{造}——他意在说，\Quote{好像\Emph{开始}一件新事。}然而，因为他悔罪而祈祷（犯了某罪，犯罪之前他更为无辜），他以何种方式说出\Quote{造}，就如此表明。\Quote{求祢在我内在部分更新正直之神。}他说，因我的所为，我灵的正直已变得老迈而弯曲。因为他在另一篇圣咏中说道：\Quote{他们使我的灵魂弯曲。}当一个人使自己屈从于世俗情欲时，他可以说是\Quote{弯曲}了；但当他被提升向天上的事物时，他的心便正直，为使天主善待他。因为经上说：\Quote{天主对心地正直的人何其美善！}此外，弟兄们，请听。有时天主在今世惩罚那在来世得赦免的罪人。就连对达味，先知已向他说过：\Quote{你的罪已得赦免，} \BibleRef{撒下 12:13} 却仍有天主为此罪所警告的某些事发生在他身上。因为他的儿子阿贝沙隆向他发动血战，多方羞辱他的父亲。\BibleRef{撒下 15:10} 他行走在忧伤中，处于屈辱的磨难里，如此顺服于天主，以至于将一切公义归于天主，承认自己所受的苦难无一不是应得的，如今拥有一颗正直的心，天主并不因此不悦。他耐心地听从一个诽谤之人、一个向他口吐恶咒的人（\BibleRef{撒下 16:10}），那人是敌方的一个士兵，与他的逆子在一起。当那人对君王咒骂不已时，达味的一个同伴怒不可遏，想要前去击打他；却被达味拦住。他为何拦住？因为他说：天主派他来咒骂我。他承认自己的罪过，拥抱补赎，不寻求自己的光荣，在所拥有的善上赞美主，在所受的苦中也赞美主，\Quote{时常赞美主，祂的赞美常在他口中。}所有心地正直的人都是如此：并非那些弯曲的人，他们自以为正直而天主弯曲；他们行恶时欢喜，受苦时亵渎；不，若置于磨难与鞭笞中，他们从扭曲的心中说：\Quote{天主，我对祢做了什么？}确实，因为他们对天主什么也没做，一切都做给了自己。\Quote{求祢在我内在部分更新正直之神。}\par

16. \Quote{求祢不要将我从祢面前抛弃} \BibleRef{咏 50:11}。转面不看我的罪恶：并\Quote{求祢不要将我从祢面前抛弃。}他畏惧谁的面，就呼求同一位的面。\Quote{求祢不要从我身上收回祢的圣神。}因为在一个告明的人内，有圣神。甚至现在，属于圣神的恩赐，就是你所做的让你不悦。不洁之神喜悦罪恶；圣神则不悦。虽然你仍恳求宽恕，但你在另一方面已与天主联合，因为你所行的恶使你厌恶：同一件事使你与祂都不悦。现在，为攻击你的热病，你们是二者：你与医师。因为一个人不可能单凭自己就有对罪的告明和对罪的惩罚：当一个人对自己发怒，且自憎自厌时，那便不是没有圣神的恩赐，他也不会说：求祢赐给我圣神，而是说：\Quote{求祢不要从我身上收回。}\par

17. \Quote{求祢归还我祢救恩的喜乐} \BibleRef{咏 50:12}。\Quote{归还}我原先所有的；因犯罪而失去的：即祢的基督。因为没有祂，谁能痊愈？因为甚至在祂成为玛利亚之子以前，\Quote{在起初已有圣言，圣言与天主同在，圣言就是天主；} \BibleRef{若 1:1} 如此，圣祖们期待将来取肉身的安排，正如我们相信已经成就。时代改变，信仰不变。\Quote{并以\OriginalTerm{主宰之神}{Principal Spirit}坚固我。}有人在此理解为天主圣三，即天主本身；唯独除去肉身的安排：因为经上记载：\Quote{天主是神。} \BibleRef{若 4:24} 因为那非物质而存在者，似乎以神的方式存在。因此有人在此理解为天主圣三的提及：\Quote{\OriginalTerm{正直之神}{upright Spirit}}指圣子；\Quote{\OriginalTerm{圣神}{Holy Spirit}}指圣神；\Quote{\OriginalTerm{主宰之神}{Principal Spirit}}指圣父。所以，这并非异端意见，不论是否如此，或\Quote{正直之神}应指人自身（当祂说：\Quote{在我内在部分更新正直之神}），就是我因犯罪而弯曲扭曲的，以便在此情况下，圣神就是主宰之神：他也不愿圣神被从他身上收回，从而在其中得以坚固。\par

18. 请看他又加上什么：\Quote{以圣神，}他说，\Quote{求祢坚定我。}何谓\Quote{坚定}？因为祢已赦免了我，因为我得以安稳，祢所赦免的不再被追究，因这安稳和这恩宠的坚定，我便不致忘恩。但我将做什么？\Quote{我要将祢的道路教导不义的人}\BibleRef{咏 50:13}。我自己\Emph{原是}不义之人（即我曾是不义的人，如今不再不义；圣神没有从我取走，我得以用圣神坚定）。\Quote{我要将祢的道路教导不义的人。}你要将什么道路教导不义的人？\Quote{不虔敬的人必归向祢。}若\Person{达味}的罪算为不虔敬，愿不虔敬的人不失望，因为天主宽恕了一个不虔敬的人；但让他们留心归向祂，学习祂的道路。若\Person{达味}的行为不算为不虔敬，而真正的不虔敬乃是离弃天主，不敬拜独一真神，或从未敬拜，或背弃所敬拜的，那么他说的话便有多余的效力：\Quote{不虔敬的人必归向祢。}祢的慈悲如此丰盈，以致归向祢的人，不仅是各种罪人，甚至不虔敬的人，都没有失望的理由。为何？因为相信那使不虔敬者成义的天主，这信德就被算为正义。\BibleRef{罗 4:5}\par

19. \Quote{求祢拯救我脱离血债，天主，我救恩的天主}\BibleRef{咏 50:14}。拉丁文译者用一个不合拉丁文法的词，却准确地表达了希腊文。因为我们都知道，拉丁文中不说\Emph{sanguines}（血，复数），也不说\Emph{sanguina}（中性复数），然而希腊译者如此使用复数，并非无因，而是因为希伯来原文如此，虔诚的译者宁可用不合拉丁文的词，也不愿不准确。那么，他为何用复数说\Quote{脱离血债}？在许多血债中，即罪过的根源——有罪的肉身中，他意指许多罪过。宗徒论及那些来自血肉败坏而产生的罪，说：\BibleQuote{血肉不能承受天主的国}\BibleRef{格前 15:50}。无疑，按同一位宗徒的真信德，那血肉将要复活，且要获得不朽，如祂所说：\BibleQuote{这可朽坏的必须穿上不朽坏的，这可死的必须穿上不死的}\BibleRef{格前 15:53}。因为败坏出于罪，故以败坏之名称呼罪。正如我们口中的那一小块肉和肢体，在发音时被称为舌头，而由舌头所发出的也称为舌头，我们说某人是希腊舌、拉丁舌；因为肉体并无不同，而是声音。同样，由舌头所发出的言语被称为舌头；由血所行的不义被称为血。因此，他顾念自己的众多不义，正如上文\Quote{涂抹我的一切罪过}，并将它们归咎于血肉的败坏，便说：\Quote{求祢救我脱离血债}，即救我脱离不义，洗净我的一切败坏……还未是实体，而是确定的希望。\Quote{我的舌头要歌颂祢的正义。}\par

20. \Quote{主，祢必开启我的嘴唇，我的口要宣扬祢的赞美}\BibleRef{咏 50:15}。\Quote{祢的赞美}，因为我被创造；\Quote{祢的赞美}，因为我犯罪时未被弃绝；\Quote{祢的赞美}，因为我受劝告去告解；\Quote{祢的赞美}，为了我得安稳而被洁净。\par

21. \BibleQuote{因为祢若乐意祭献，我必定献上}\BibleRef{咏 50:16}。\Person{达味}生活在献祭牲为祭的时期，他却预见了这未来的时代。我们岂不在这话中看见自己吗？那些祭献是预像，预示那唯一的救恩祭献。我们也不缺少可献给天主的祭献。请听他因罪而忧心、盼望恶行得赦时所说的话：\Quote{祢若乐意祭献，我必定献上。祢必不喜悦全燔祭。}那么我们便无物可献吗？我们就这样到天主面前吗？我们如何平息祂？献上吧；你身上确有可献的。不要从外面取乳香，而要说：\Quote{天主，我向祢所许的愿在我身上，我要向祢献上赞颂的祭。}不要从外面寻找可杀的牲畜，你身上就有可杀的。\BibleQuote{天主所要的祭献就是忧伤的痛悔，天主，祢不轻看痛悔和谦卑的心}\BibleRef{咏 50:17}。祂完全轻看公牛、公山羊、公绵羊；如今不是献这些的时候。它们曾在象征某事物、应许某事物时被献；当所应许的实现时，应许便废止了。\Quote{天主不轻看痛悔和谦卑的心。}你知道天主是至高者；你若自高，祂便远离你；你若自卑，祂必亲近你。\par

22. 请看这是谁：\Person{达味}似乎以一个人的身份恳求；你们要在这里看到我们的形象和教会的预表。\par

\BibleQuote{Deal kindly, O Lord, in Your good will with Sion}\BibleRef{咏 50:18} . With this \Place{熙雍} deal kindly. What is \Place{熙雍}? A city holy. What is a city holy? That which cannot be hidden, being upon a mountain established. \Place{熙雍} in \OriginalTerm{观望}{prospect}, because it has \OriginalTerm{观望}{prospect} of something which it hopes for. For \Place{熙雍} is interpreted \Quote{\OriginalTerm{观望}{prospect}}, and \Place{耶路撒冷}, \Quote{\OriginalTerm{和平的异象}{vision of peace}}. You perceive then yourselves to be in \Place{熙雍} and in \Place{耶路撒冷}, if being sure ye look for hope that is to be, and if you have peace with God. \Quote{And be the walls of \Place{耶路撒冷} built.} \Quote{Deal kindly, O Lord, in Your good will with \Place{熙雍}, and be the walls of \Place{耶路撒冷} built.} For not to herself let \Place{熙雍} ascribe her merits: do Thou with her deal kindly, \Quote{Be the walls of \Place{耶路撒冷} built:} be the battlements of our immortality laid, in \OriginalTerm{信德}{faith} and \OriginalTerm{望德}{hope} and \OriginalTerm{爱德}{charity}.\par

23. \Quote{Then You shall accept the sacrifice of righteousness} \BibleRef{咏 50:19} . But now sacrifice for iniquity, to wit, a spirit troubled, and a heart humbled; then the sacrifice of righteousness, praises alone. For, \Quote{Blessed they that dwell in Your house, for ever and ever they shall praise You:} for this is the sacrifice of righteousness. \Quote{\OriginalTerm{供物}{Oblations} and \OriginalTerm{全燔祭}{holocausts}.} What are \Quote{\OriginalTerm{全燔祭}{holocausts}}? A whole victim by fire consumed. When a whole beast was laid upon the \OriginalTerm{祭台}{altar} with fire to be consumed, it was called a \OriginalTerm{全燔祭}{holocaust}. May divine fire take us up whole, and that fervour catch us whole. What fervour? \Quote{Neither is there that hides himself from the heat thereof.} What fervour? That whereof speaks the \OriginalTerm{宗徒}{Apostle}: \Quote{In spirit fervent.} \BibleRef{罗 12:11} Be not merely our \OriginalTerm{灵魂}{soul} taken up by that divine fire of wisdom, but also our \OriginalTerm{身体}{body}; that it may earn their \OriginalTerm{不朽}{immortality}; so be it lifted up for a \OriginalTerm{全燔祭}{holocaust}, that death be swallowed into victory. \Quote{\OriginalTerm{供物}{Oblations} and \OriginalTerm{全燔祭}{holocausts}.} \Quote{Then shall they lay upon your \OriginalTerm{祭台}{altar} \OriginalTerm{牛犊}{calves}.} Whence \Quote{\OriginalTerm{牛犊}{calves}}? What shall He therein choose? Will it be the innocence of the new age, or necks freed from the yoke of the law?...\par

\SourceRecord{Translated by J.E. Tweed.；Nicene and Post-Nicene Fathers, First Series；Vol. 8.；Edited by Philip Schaff.；Buffalo, NY: Christian Literature Publishing Co.,；1888.；<http://www.newadvent.org/fathers/1801051.htm>.}

% structure: p0001 source=h1 target=section reason=page-title

\section{\WorkTitle{论圣咏第五十二篇}}

\Quote{结局，\Person{达味}的智慧，当\Person{厄东}人多厄格来告诉\Person{撒乌耳}说：\Person{达味}进了\Person{阿彼默肋客}的家。}但我们读到的是\Person{达味}进了\Person{阿希默肋客}的家。或许我们不应无理地认为，由于名字的相似和一个音节或一个字母的差异，标题因而有所不同。然而，在我们查阅过的圣咏抄本中，我们发现更多的是\Person{阿彼默肋客}而非\Person{阿希默肋客}。况且，在另一处你有一篇极其明显的圣咏，它暗示的并非名字的相似，而是完全不同的名字；例如，\Person{达味}在\Person{阿基士}王面前改变了面容，而不是在\Person{阿彼默肋客}王面前；他打发他走了，他就离开了：但该篇圣咏的标题却如此写着：\Quote{当他在\Person{阿彼默肋客}面前改变面容时}——这一名字的改变使我们更加专注于奥秘，以免你追求所谓的史实，而轻视神圣的遮蔽……\par

2. 你们要观察两类人：一类是劳苦的人，另一类是他们在其中劳苦的人；一类是思念地上之事的人，另一类是思念天上之事的人；一类是将自己的心沉重地压向深渊的人，另一类是将自己的心与天使结合起来的人；一类是信赖世上之物、即这个世界所充盈的东西，另一类是信赖那不说谎的天主所应许的天上之事。但这两类人是混杂在一起的。我们现在看到\Place{耶路撒冷}的公民、天国的公民在地上担负某些职务；例如，有人穿紫袍，是长官，是市政官，是总督，是皇帝，他治理地上的共和国；但他若是一个基督徒、一个信者、一个虔诚的人、一个轻看他所拥有的东西、并信赖他尚未拥有的事物的人，他的心就在天上。那位圣妇\Person{艾斯德尔}便属于这一类，她虽是国王的妻子，却为同胞代祷而甘冒危险；当她在天主面前祈祷时（她不能撒谎），她在祈祷中说，她的王后饰品对她来说不过是行经女子的布。\BibleRef{艾 14:16}那么，当我们看到天国的公民从事\Place{巴比伦}的事务、在尘世共和国中做某些尘世之事时，我们不要对他们失望；另一方面，我们也不要立刻祝贺所有我们看见做天上之事的人，因为毒恶之子有时也坐在\Person{梅瑟}的座位上，关于他们曾说过：\BibleQuote{凡他们所说的，你们要遵行；但不要效法他们的行为，因为他们能说不能做。}\BibleRef{玛 23:3}那些人，在尘世之物中，将心举向天堂；这些人，在天上言语中，却将心拖在地上。但必将有簸扬的时候，那时两者要被极其仔细地分开，以便没有一粒麦子落入将焚的糠秕堆中，也没有一根稻草落入要储入谷仓的禾捆中。\BibleRef{玛 3:12}只要现在还是混杂的，我们就要从中听出我们的声音，即天国公民的声音（因为这是我们应当向往的：在此忍受恶人，而不是被好人忍受）；让我们以耳、以舌、以心、以工与我们自己联合。如果我们这样做了，我们就在我们听到的这些事中说话。因此，让我们首先谈论尘世国度的邪恶身体。\par

3. \BibleQuote{有势力的人为何以恶事夸耀？}\BibleRef{咏 51:1}我的弟兄们，请注意恶毒的夸耀，恶人的夸耀。夸耀在哪里？\BibleQuote{有势力的人为何以恶事夸耀？}那就是，那在恶事上有势力的人，为何夸耀？人需要有势力，但在善事上，而非在恶事上。以恶事夸耀算得了什么大事？建造房屋属于少数人，随便一个无知的人都能拆毁。播种小麦、照料庄稼、等待成熟、并享受劳动成果，属于少数人；随便一个人用一点火星就能烧毁所有庄稼。养育婴儿、出生后喂养、教育、将他带到青年时期，是一项伟大的任务；随便一个人在一瞬间就能杀死他。因此，那些旨在毁灭的事情是最容易做到的。\BibleQuote{夸耀的，应当因主而夸耀。}\BibleRef{格前 1:31}夸耀的，应当因善事而夸耀。你夸耀，因为你以恶事有势力。你要做什么，有势力的人，你自夸甚多？你要杀人：这件事蝎子也能做，一场热病也能做，有毒的蘑菇也能做。你如此强大的势力竟沦为与有毒的蘑菇相等？因此，\Place{耶路撒冷}的善良公民们，他们不是以恶事，而是以善事夸耀：首先，他们不是在自己身上，而是在主内夸耀。其次，他们热切地做那些有利于建设的事情，并且做那些坚固持久的事情；而那些导致毁灭的事情，他们也可能为了进德之人的纪律而做，不是为了压迫无辜。那么，与这种势力相比，那地上的身体为何不能从这些话中听到：\Quote{有势力的人为何以恶事夸耀？}\par

4. \BibleQuote{在邪恶中，你的舌头终日思想了不义} \BibleRef{咏 51:2}：这就是说，在整个时间内，无厌倦、无间断、无停息。当你不在作恶时，你却在思想，以致当任何邪恶离开你的双手时，它并未离开你的心；你或是行一件恶事，或是当你不能行时，你说一句恶言，就是说，你口出恶言；或者连这也做不到时，你仍意愿并思想一件恶事。\Quote{终日}，就是说，不间断。我们期望这人受惩罚。他自己岂不是一个很小的惩罚吗？你恐吓他：你，当你恐吓他时，你要把他送到哪里？送到邪恶里去吗？把他送到他自己那里去吧。为了发泄你的怒气，你要把他交在野兽的权下；对他来说，他自己比野兽更坏。因为野兽只能损伤他的身体，他自己却不能完整地保存自己的心。在内部，他自己对自己发怒，而你却从外面寻求鞭打吗？不，为他祈求天主，使他能从他自身获得释放。然而，在这篇诗篇里，我的弟兄们，并没有为恶人祈求，也没有反对恶人，而是预言恶人将有什么结局。所以不要以为这篇诗篇是出于恶意说的：因为它是出于预言的神恩而说的。\par

5. 那么接下来是什么？你所有的大能和所有邪恶的思想终日不断，以及你舌头上不间断的恶意默想，成就了什么，做了什么？\BibleQuote{如同锋利的剃刀，你行了欺诈} \BibleRef{咏 51:3}。看恶人对圣人所做的是什么：他们刮掉他们的头发。我说的是什么？如果\Place{耶路撒冷}的那些子民，听从他们的主、他们的君王的声音说：\BibleQuote{不要害怕那些杀害身体，却不能杀害灵魂的}；他们听从刚才从福音中宣读的声音：\BibleQuote{人纵然赚得了全世界，却赔上自己的灵魂，为他有什么益处？} \BibleRef{玛 16:26}，他们就轻视一切现世的美好事物，尤其是生命本身。那么，\Person{多厄格}的剃刀，对一个在地上默想天国、即将进入天国、与天主同在、并将与天主常居的人，能做什么？那剃刀能做什么？它只能刮掉头发，使人成为秃头。这也属于基督，祂曾在\Place{髑髅地}被钉十字架。\BibleRef{玛 27:33} 它也造就了\Person{科辣黑}的儿子，这名字被解释为\Quote{秃头}。\BibleRef{编上 6:22} 因为这头发象征现世之物的多余。这些头发确实是天主在人的身体上并非多余地创造，而是为了一种装饰；然而，因为它们被剪掉时没有感觉，那些全心依恋主的人，对待这些现世之物就如对待头发一样。但有时甚至也借\Quote{头发}做出一些善事：\BibleQuote{当你擘饼给饥饿的人，将无家可归的穷人领进你的房屋；你若看见赤身露体的人，就给他遮盖} \BibleRef{依 58:7} 最后，殉道者们自己也效法主，为教会流血，听从那声音：\BibleQuote{如同基督为我们舍掉了性命，我们也应当为弟兄们舍掉性命}，\BibleRef{若一 3:16} 在某种程度上，他们用自己的头发为我们行了善，就是说，用那些剃刀可以砍掉或刮掉的东西。但是，因此甚至用头发也能行善，那个犯罪的女人也暗示了这一点：她曾在主的脚前哭泣，用她的头发擦干被泪水浸湿的脚。\BibleRef{路 7:38} 这表示什么？当你怜悯某人时，你也应在可能时救济他。因为当你怜悯时，你如同流泪；当你救济时，你如同用头发擦干。如果对任何人如此，更何况对主的脚呢？主的脚是什么？是圣史，关于他们曾说过：\BibleQuote{那传报和平、传报佳音的脚是多么美丽啊！} 因此，让\Person{多厄格}像剃刀一样磨利他的舌头，让他尽其所能磨利欺诈：他将除去多余的现世之物，但他能除去永恒的必要之物吗？\par

6. \BibleQuote{你爱恶意胜过良善}\BibleRef{咏 51:4}。在你面前原是良善；你本该爱她的。因为你不必花费什么，也不必远航去取你所爱的。良善在你面前，不义在你面前：比较并选择吧。但或许你有眼能见恶意，却没有眼能见良善。那恶毒的心有祸了。更糟的是，它自己转开，不去看它本可以看见的。因为关于这类事，在别处不是说了吗？\Quote{他不愿明白，好去行善。} 经上不是说，他不能；而是\Quote{他不愿}，他说，\Quote{明白，好去行善，} 他闭眼不看眼前的光明。接着呢？\Quote{他在床上图谋不义；} 即在他内心的隐秘处。这种责备就堆在这厄东人多厄格身上，他是一个恶毒的身体，一个地上的运动，不安定，不属于天上。\Quote{你爱恶意胜过良善。} 因为你想知道恶人如何能看见两者，却宁愿选择前者，转身远离后者吗？为何他在遭受不义时常喊叫？为何那时他极力夸大不义，称赞良善，指责那待他以恶意胜过良善的人？那么，让他自己成为看见的准则：他必凭自己受审判。此外，若他遵行所记载的：\BibleQuote{你当爱你的邻人如你自己；}\BibleRef{玛 22:39} 以及\BibleQuote{凡你们愿意人给你们做的，你们也要照样给人做；}\BibleRef{玛 7:12} 他在家里就有认识的途径，因为凡他不愿别人对他做的，他也不该对别人做。\Quote{你爱恶意胜过良善。} 不义地、混乱地、颠倒地，你竟要把水升到油之上：水会下沉，油仍在上。你竟要把黑暗置于光明之下：黑暗将被驱散，光明仍在。你竟要把地置于天之上：地因自身重量会落回原位。因此，因爱恶意胜过良善，你也要下沉。因为恶意永不能胜过良善。\Quote{你爱恶意胜过良善：不义胜过谈论公平。} 公平在你面前，不义在你面前：你只有一条舌头，你想往哪里转就转：那么为何转向不义而非公平呢？你不把苦毒的食物给你的肚腹，却把不义的食物给你的恶舌？正如你选择赖以生存的食物，也选择你能说的话。你偏爱不义胜过公平，偏爱恶意胜过良善；你确实偏爱，但你能胜过什么呢？唯有良善和公平。但你自己仿佛置于那些必须沉下去的事物之上，你并不会使它们高于善事，但你却与它们一同沉入恶事。\par

7. 因此诗篇接着写道：\BibleQuote{你爱一切沉沦的话语}\BibleRef{咏 51:5}。那么，如果你能从沉沦中逃脱，就救自己吧。你正逃避船难，却拥抱铅块！如果你不愿沉没，就抓住一块木板，被木头承载，让十字架带你通过。但现在因为你是厄东人多厄格，一个\Quote{运动}和\Quote{属土的}，你在做什么？\Quote{你爱一切沉沦的话语，诡诈的舌头。} 这已经先说了，沉沦的话语接着诡诈的舌头。什么是诡诈的舌头？诡诈的舌头是欺诈的仆役，心里怀着一样，口里说出另一样。而在这些之中是倾覆，在这些之中是沉沦。\par

8. \BibleQuote{因此天主最终必毁灭你}\BibleRef{咏 51:6}：尽管现在你像田野的草在太阳发热前似乎繁茂。因为\BibleQuote{凡有血肉的都是草，人的荣华如同草上的花：草枯干了，花凋谢了；但主的话却永远长存。}\BibleRef{依 40:6} 看，你可以把自己绑在什么上，绑在什么\Quote{永远长存}上。因为如果你把自己绑在草和草的花上，既然草要枯干，花要凋谢，\Quote{天主最终必毁灭你：}若不是现在，肯定在末了祂要毁灭，当那簸扬来到，糠秕堆与饱满的谷粒被分开的时候。饱满的谷粒不是要收进仓里，糠秕不是要丢进火里吗？那整个多厄格岂不站在左边，当主要说：\BibleQuote{你们进入永火里去，这永火是为魔鬼和他的天使预备的。}\BibleRef{玛 25:41} 因此\Quote{天主最终必毁灭：必把你拔出，从你的帐棚中除掉。} 那么这厄东人多厄格现在是在一个帐棚里：\BibleQuote{但仆人不永远住在家里。}\BibleRef{若 8:35} 他甚至也做点好事，即使不是以自己的行为，至少以天主的话语，以致在教会里，当他\BibleQuote{寻求自己的事}\BibleRef{斐 2:21} 时，他至少会说那些属于基督的事。\par

\Quote{但祂要把你从你的居所挪去。}\Quote{我实在告诉你们，他们已获得了他们的赏报。}\BibleRef{玛 6:2}\Quote{并从生活之地拔除你的根。}因此，我们必须在生活之地有根。让我们的根在那里。根是看不见的；果实可以看见，根却看不见。我们的根就是我们的爱，我们的果实就是我们的行为：你的行为必须出自爱，这样你的根就在生活之地。那时，那个\Person{多厄格}将被连根拔起，他绝不能再留在那里，因为他没有更深地将根扎在那里：\BibleRef{玛 13:5}但他的情形将像那些撒在石头上的种子一样，即使长出根来，因为没有水分，太阳一升起就立刻枯干了。但是，那些把根扎得更深的人，从宗徒那里听到了什么？\Quote{我屈膝为你们向我们的主耶稣基督的父祈求，使你们在爱中生根建基。}因为如今已经有了根，他说：\Quote{使你们能够领悟那高、阔、长、深，并认识基督那超越知识之爱，使你们充满天主的一切圆满。}如此伟大的果实是配得上这样伟大的根，这根是如此单一，如此萌芽，为如此深植的嫩芽。但这个\Person{多厄格}的根却要从生活之地被拔起。\par

9. \Quote{义人看见，必要恐惧；他们必要嘲笑他}\BibleRef{咏 51:7}。何时恐惧？何时嘲笑？因此我们要理解，并区分这恐惧和嘲笑的两个时间，它们各有各的用处。因为只要我们还在这个世界上，就不该嘲笑，免得以后哀哭。我们已经读到为这个\Person{多厄格}在终末所保存的是什么，我们已读到，因为我们明白并相信，我们看见便恐惧。因此，经上说了：\Quote{义人看见，必要恐惧。}我们看见恶人的结局是什么，为什么我们恐惧？因为宗徒说：\Quote{怀着恐惧和战栗，作成你们的得救}\BibleRef{斐 2:12}；因为圣咏上说：\Quote{你们当以敬畏事奉主，以战栗欢呼。}为什么\Quote{以敬畏}？\Quote{所以，那自以为站得稳的，要谨慎，免得跌倒。}\BibleRef{格前 10:12}为什么\Quote{以战栗}？因为他在另一处说：\Quote{弟兄们，如果一个人陷于某种过犯，你们属神的人，要以温柔的心神矫正他；但你自己要小心，免得也受引诱。}\BibleRef{迦 6:1}因此，现今的义人，即因信德而生活的人，看见这个\Person{多厄格}将有何结局，但他们也为自己恐惧：因为他们知道今天是什么样子，却不知道明天将会怎样。因此，如今\Quote{义人看见，必要恐惧。}但他们何时嘲笑？当邪恶过去之后；当它飞逝之后；如同现在这不确定的时间在很大程度上已经飞逝；当这世界的黑暗被驱散之后，我们现在行路若不靠圣经的灯盏，便会恐惧，如同在黑夜中。因为我们依靠预言行走；为此\Person{伯多禄}宗徒说：\Quote{我们有更确实的先知之言，你们要留意它，如同在黑暗处的一盏灯，直到天破晓，晨星在你们心中升起。}\BibleRef{伯后 1:19}只要我们还靠灯盏行走，就必须怀着恐惧生活。但是，当我们的日子来到，即基督的显现，同一位宗徒说：\Quote{当基督——你们的生命显现时，你们也要与祂一同出现在光荣中}\BibleRef{哥 3:4}，那时义人就要嘲笑那个\Person{多厄格}……\par

10. 但那些将要讥笑的人那时会说甚么呢？\BibleQuote{他们必讥笑他，且说：看，这人不以天主为他的助佑。}\BibleRef{咏 51:8} 请看这\OriginalTerm{属地的身体}{body earthly}！\Quote{你有多大，就看你拥有多少。}这是一句谚语，出自\OriginalTerm{贪婪}{covetousness}之人、攫取之人、压迫无辜之人、抢夺他人财物之人、背信弃义之人。这是什么谚语？\Quote{你有多大，就看你拥有多少。}也就是说，你拥有多少金钱，你获得多少，你就变得多强大。\BibleQuote{看，这人不以天主为他的助佑，却依靠他财富的众多。} 一个穷人，或许是一个恶人，切莫说：我不属于这身体。因为他听见先知说：\BibleQuote{他依靠他财富的众多。} 立刻，如果他贫穷，他留意自己的破衣，他或许看见在他附近有一个天主子民中穿着更华丽的人，他心里说：先知说的是这个人；难道说的是我吗？不要因此自外，不要分开自己，除非你看见了并畏惧，以便日后你可以讥笑。如果你缺少资财却燃烧着\OriginalTerm{贪欲}{cupidity}，这对你有何益处？当我们的主\Person{耶稣基督}对那个忧愁并离开祂的富人说了：\BibleQuote{去，变卖你所有的一切，给穷人，你将有宝藏在天上，然后来跟随我。}\BibleRef{玛 19:21} 并预言富人毫无希望，以致祂说：骆驼穿过针眼，比富人进入天国更容易。\BibleRef{玛 19:24} 门徒们不是立刻忧愁起来，心里说：\BibleQuote{谁能得救呢？} 因此当他们说：\BibleQuote{谁能得救呢？} 他们只想到少数富人吗？那么多穷人难道被他们忽略了吗？他们难道不能对自己说：如果富人进入天国是困难甚至是不可能的，如同骆驼穿过针眼一样不可能，那就让所有穷人进入天国，唯独富人被关在外面吧？因为富人何其少？但穷人多得难以计数。因为我们不能只看天国的外衣；每个人的衣裳要计算为义德的辉煌：所以穷人将等同于天主的使者，穿着不朽的长袍，在他们父的国里发光如太阳：我们为何要为少数富人忧虑或苦恼呢？宗徒们可没有这样想；但当主说了：\BibleQuote{骆驼穿过针眼，比富人进入天国更容易。} 他们心里说：\BibleQuote{谁能得救呢？} 这指什么呢？不是指财富，而是指欲望；因为他们看见穷人自己，即使没有金钱，却仍有贪婪。而且，为使你们知道，受谴责的不是富人的金钱，而是贪婪，请听我说：你看见站在你附近的那个富人，或许他有金钱，却没有贪婪；你虽然没有金钱，却有贪婪。一个浑身生疮、痛苦不堪、被狗舔舐、无人帮助、没有食物、甚至可能没有一件衣服的穷人，他被天使们送到\Person{亚巴郎}的怀抱里。\BibleRef{路 16:22} 啊！作为穷人，你现在欢喜吗？难道你也希望生疮？你的本钱不就是健康吗？这位\Person{拉匝禄}的功劳不在于贫穷，而在于虔诚。因为你看见谁被提升，却没有看见他被提升到何处。谁被天使提升？一个穷人，充满痛苦，满身生疮。他被提升到何处？到\Person{亚巴郎}的怀抱。阅读圣经，你会发现\Person{亚巴郎}是一个富人。\BibleRef{创 13:2} 为使你知道，受责备的不是财富；\Person{亚巴郎}有很多金、银、牲畜、家仆，他是一个富人，而穷人\Person{拉匝禄}却被提升到他的怀抱。富人到了富人的怀抱？难道不都是在天主面前富有的人，而在贪欲上都是穷人吗？……\par

11. 因此，那\Quote{只依恃自己的财富，以自己的虚妄而夸耀}的人已被定罪。因为还有什么比认为金钱比天主更有用更虚妄呢？因此，那说\Quote{那百姓是有福的，因为主是他们的天主}的人已被定罪。而你说：\Quote{但我是}；现在终于听到那身体说：\BibleQuote{但我好像天主殿中结实累累的橄榄树}\BibleRef{咏 51:9}。这不是一个人说的，而是那结实的橄榄树说的，骄傲的枝条被剪去，谦卑的野橄榄被接上。\BibleRef{罗 11:17}\BibleQuote{好像天主殿中结实累累的橄榄树，我信赖天主的仁慈。}他做了什么？\Quote{只依恃自己的财富。}所以他的根要从活人之地拔出。\Quote{但我}，因为\BibleQuote{好像天主殿中结实累累的橄榄树}，其根得滋养，不被拔出，\BibleQuote{我曾信赖天主的仁慈。}但或许只是现在？因为甚至在此，人有时也会犯错。他们确实敬拜天主，但现在不像那个\Person{多益}；尽管依靠天主，却是为了暂时的事物；所以他们对自己说：我敬拜我的天主，祂会使我在世上富足，会赐给我儿子、妻子。确实，唯独天主赐予这些事物，但天主不希望我们为这些事物而爱祂。因为祂也常将这些赐予恶人，为要使善人从祂那里寻求别的东西。那么你是如何说\BibleQuote{我曾信赖天主的仁慈}的呢？或许是为了获得暂时的事物？不，而是\BibleQuote{直到永远，永世无尽}。他通过加上\BibleQuote{永世无尽}来重复\BibleQuote{直到永远}，为要通过重复来确认他怎样扎根于对天国及永远幸福的爱与希望中。\par

12. \BibleQuote{我要永远称谢祢，因为祢行了这事}\BibleRef{咏 51:10}。\Quote{行了什么事？}\Person{多益}被定罪，\Person{达味}被加冕。\BibleQuote{我要永远称谢祢，因为祢行了这事。}伟大的宣认：\BibleQuote{因为祢行了这事！}\Quote{行了}什么？除了上面所说的那些事，即我好像天主殿中结实累累的橄榄树，信赖天主的仁慈，直到永远，永世无尽？是祢行了这事：不虔敬的人不能自己成义。但那使人成义的是谁？\BibleQuote{相信}，他说，\BibleQuote{在祂}那使\BibleQuote{不虔敬的人}成义者。\BibleRef{罗 4:5}\BibleQuote{你有什么不是领受的呢？既然是领受的，为什么你夸耀，好像不是领受的，好像是你自己的？}\BibleRef{格前 4:7} 那反对\Person{多益}、在地上容忍\Person{多益}直到他从居所被移走、从活人之地被拔出的人说：但愿我不这样夸耀，好像不是领受的，而是在天主内夸耀。\BibleQuote{我要永远称谢祢，因为祢行了这事}，即，因为祢所行的，不是按我的功绩，而是按祢的仁慈。但我做了什么？如果你记得，\BibleQuote{原先我是亵渎者、迫害者和凌辱者。}但你做了什么？\BibleQuote{但我蒙受了怜悯，因为我是在不信中做的。}\BibleRef{弟前 1:13}\BibleQuote{我要永远称谢祢，因为祢行了这事。}\par

13. \BibleQuote{我要仰望祢的圣名，因为它是甘饴的。}世界是苦涩的，但祢的圣名是甘饴的。即使世上有某些甘甜的事物，也是伴着苦涩被消化。祢的圣名不但因其伟大，也因其甘饴而被选。\Quote{因为不义的人向我讲述他们的快乐，但不像祢的法律，主。}因为如果殉道者们没有尝到甘甜，他们就不会如此平静地忍受极大苦难的苦涩。任谁都能体验到他们的苦涩，但很少人能尝到他们的甘甜。因此，对那爱天主超越一切甘甜的人来说，天主的圣名是甘饴的。\BibleQuote{我要仰望祢的圣名，因为它是甘饴的。}你如何证明它是甘饴的呢？请给我一个能品尝甘饴的味觉。你尽管赞美蜂蜜，尽量用你能说的言语夸耀它的甘甜：一个不知道蜂蜜是什么的人，除非他尝过，否则他不知道你在说什么。那么，那篇邀请你的圣咏说了什么？\BibleQuote{请你们体验，请你们观看，主是何等甘饴。}你不去体验，却说：它甘饴吗？什么是甘饴？如果你已体验，就在你的果实中找到它，而不只是在言语中，好像只是叶子，免得你因主的诅咒，像那棵无花果树\BibleRef{玛 21:19}一样枯萎。\BibleQuote{体验}，他说，\BibleQuote{并观看，主是何等甘饴。}体验并观看：如果你体验了，你便会看见。但对于一个不体验的人，你如何证明？通过赞美天主的圣名的甘饴，你所说的一切都是言语；另一样是味道。甚至连不敬虔的人也听到对祂赞美的言语，但除了圣人，没有人尝到它有多甘甜。此外，一个辨别出天主的圣名的甘饴，并想要展示它、想向人表明它，却找不到可以展示的人（因为对圣人无需展示，他们自己便体验并知道，而不敬虔的人却无法辨别他们不愿体验的东西），他做了什么？因为天主的圣名的甘饴，他立即将自己从不敬虔的人群中带走。\BibleQuote{我要仰望}，他说，\BibleQuote{祢的圣名，因为它甘饴，在祢的圣者面前。}祢的圣名是甘饴的，但不在不敬虔的人面前。我知道它是何等甘甜，但只对那体验过的人而言。\par

\SourceRecord{Translated by J.E. Tweed.；Nicene and Post-Nicene Fathers, First Series；Vol. 8.；Edited by Philip Schaff.；Buffalo, NY: Christian Literature Publishing Co.,；1888.；<http://www.newadvent.org/fathers/1801052.htm>.}

% structure: p0001 source=h1 target=section reason=page-title

\section{圣咏第五十三篇释义}

1. 我们将在主所供给的范围内，与你们一起解释这首圣咏。一位弟兄敦促我们，使我们有此意愿，并祈祷我们能有能力。若我因仓促而有所遗漏，那甚至屈尊赐予我们所能说出的话的那一位，必在你们内补足。它的标题是：\Quote{末后，为\OriginalTerm{玛厄勒特}{Maeleth}，属于达味自己的智慧。} \Quote{为\OriginalTerm{玛厄勒特}{Maeleth}，}按希伯来名称的解释，似乎是指\Quote{为一位正在生产或痛苦的人。}但世上谁是那生产并痛苦的人，信友都承认，因为他们正是这样的人。基督在此生产，基督在此痛苦：头在上，肢体在下。因为一个不生产也不痛苦的人不会说：\BibleQuote{扫禄，扫禄，你为什么迫害我？}（\BibleRef{宗 9:4}）祂与他一起，当祂在迫害中生产时，他悔改了，祂使他生产。因为他自己后来也被光照，被接枝在他从前迫害的那些肢体上；怀着同样的爱，他说：\BibleQuote{我的小孩子们，我为你们再次生产，直到基督在你们内形成。}（\BibleRef{迦 4:19}）因此，为基督的肢体，为祂的身体——教会（\BibleRef{哥 1:24}），为那同一个新人，即那同一个合一，其头在上，这首圣咏被歌唱……那么，我们是在谁中间生产并叹息呢？如果我们是在基督的身体内，如果在祂头之下生活，如果被算作祂的肢体，他们是谁？听吧。\par

2. \BibleQuote{愚妄人心里说：没有天主}（\BibleRef{咏 52:1}）。基督的身体就是在这样的人中间疼痛并叹息。如果这类人是这样的，我们并非为许多人生产；依我们的想法，似乎很少；很难遇到一个人心里说\Quote{没有天主}；然而，人数如此之少，以至于他们害怕在众人中间说这话，就在心里说，因为用口说他们不敢。那么，我们所受的并不大，很难找到：那种心里说\Quote{没有天主}的人很少见。但是，如果从另一层意义来考察，岂不是在更多人身上发现了我们原以为只在少数、罕见、几乎无人身上才有的情况吗？让那些行为邪恶的人来到当中，让我们看看放荡、大胆、邪恶之人的行为，他们人数众多；他们日复一日地助长自己的罪恶，他们的行为变成了习惯，甚至失去了羞耻感：如此众多的人，置身其间的基督的身体几乎不敢谴责那不受约束去犯的罪，并认为保持纯洁的完整，不去做那因习惯而不敢指责或一旦敢于指责就引发恶人谴责和反诉的事，是自身的一件大事，而善人的自由声音却更难发出。那些人就是心里说\Quote{没有天主}的人。我正在驳斥这样的人。凭什么驳斥？他们判断自己的行为中悦天主。因此，祂不是说\Quote{有人说}，而是\Quote{愚妄人心里说：没有天主。}这些人如此相信存在着天主，以致判断自己的行为中悦这位天主。但是，如果你有智慧，你就能看出\Quote{愚妄人心里说：没有天主}是怎么回事。如果你留意、理解、考察：那认为邪恶行为中悦天主的人，并不以为祂是天主。因为如果天主存在，祂就是公义的；如果祂是公义的，不义就使祂不悦，邪恶使祂不悦。但是，当你认为邪恶使祂喜悦时，你就是在否认天主。因为如果天主是不喜悦邪恶的那一位，但天主在你眼中似乎不是不喜悦邪恶的那一位，而除了不喜悦邪恶的那一位以外，别无天主；那么，当你心里说：天主纵容我的不义时，你无非是说：\Quote{没有天主。}\par

3. 让我们也注意那关乎我们主基督自己、我们元首自己的意义。因为当祂以奴仆的形体\BibleRef{斐 2:7}出现在地上时，那些钉死祂的人说：\Quote{祂不是天主。} 因为祂是天主子，祂确实是天主。但那些败坏、变成可憎恶的人说了什么？\Quote{祂不是天主：}我们杀了祂，\Quote{祂不是天主。} 你们在智慧篇中就有这些人的声音。因为在这节\BibleQuote{愚顽人心里说：没有天主}之前，仿佛需要理由说明愚顽人为何这样说，他又加上了：\BibleQuote{他们败坏，在邪恶中变为可憎} \BibleRef{咏 52:2}。你们要听那些败坏的人。\BibleQuote{因为他们私下议论，思想不正：} \BibleRef{智 2:1} 败坏始于邪恶的信仰，由此进到堕落的道德，再到最昭彰的恶行，这就是阶梯。但他们思想不正，私下说了什么？\BibleQuote{我们的生命短暂又愁烦。} \BibleRef{智 2:1} 从这邪恶的信仰就引出宗徒所说的：\BibleQuote{我们吃喝吧，因为明天我们要死了。} \BibleRef{格前 15:32} 但在前一段中，对放纵情欲本身描述得更详细：\BibleQuote{我们要用玫瑰花冠冕，趁它们还未枯萎；在每一处留下我们欢愉的记号。} \BibleRef{智 2:8} 在那放纵情欲的详细描述之后，接着是什么？\BibleQuote{我们要杀死贫穷的义人：} \BibleRef{智 2:10} 因此这就等于说：\Quote{祂不是天主。} 他们刚才似乎说出柔和的话：\BibleQuote{我们要用玫瑰花冠冕，趁它们还未枯萎。} 有什么比这更娇嫩、更柔和呢？你岂能从这柔和中期待十字架、刀剑吗？不要惊奇，荆棘的根也是柔嫩的；若有人触摸它们，不会被刺伤；但从中生出的东西会刺伤你。因此，那些人\BibleQuote{败坏，并在他们的邪恶中变为可憎。} 他们说：\BibleQuote{祂如果是天主子，就从十字架上下来吧。} \BibleRef{玛 27:40} 看，他们公然说：\Quote{祂不是天主。}……\par

4. \BibleQuote{上主从天垂看世人，要看有没有明智的人，寻求天主的人} \BibleRef{咏 52:3}。这是什么意思？\BibleQuote{他们败坏了，}所有这些说\BibleQuote{没有天主}的人？怎么？难道天主不知道他们成了这样？或者，若不是由祂告诉我们，我们的内心思想岂能被揭示？如果祂明白，如果祂知道，那么所说的\BibleQuote{祂要察看}是什么意思？因为这是询问者、不知情者的话。\BibleQuote{天主从天垂看，}等等。祂仿佛通过察看、从天俯视，找到了祂所寻求的，便作出判决：\BibleQuote{众人都偏离了正道，一同变为无用；没有一个行善的，连一个也没有} \BibleRef{咏 52:4}。有两个相当困难的问题：因为如果天主从天察看，为要看有没有明智或寻求天主的人；一个愚昧人会认为，天主不知道一切。这是一个问题。另一个是什么？如果没有一个行善的，连一个也没有；谁是在恶人中叹息的人呢？第一个问题这样解决：圣经常常这样说话，以致受造物因天主的恩赐所做的，天主被说成是所作的。……因此也有以下经文：\BibleQuote{因为圣神洞察一切，连天主的深奥也洞悉；} \BibleRef{格前 2:10} 不是因为那洞察万物的在搜寻，而是因为圣神被赐给了你们，使你们也能搜寻；并且那因祂的恩赐你们所作的，祂被说成是所作的；因为没有祂，你们就不会作。因此，当你们作时，天主被说成是所作的。正因如此，天主从天\BibleQuote{垂看世人。} 第一个问题，按我们的能力，就这样解决了。\par

5. 那通过垂看我们所承认的是什么？那通过垂看天主所承认的是什么？祂承认什么（因为在这里祂赐予它）？请听它是什么：就是\BibleQuote{众人都偏离了正道，一同变为无用；没有一个行善的，连一个也没有。} 那么，另一个问题是什么？不就是我先前所提的那个吗？如果\BibleQuote{没有一个行善的，连一个也没有，}就没有人在恶人中叹息。且住，主说：不要仓促下判断。我赐给了人行善的能力；但祂说，是出于我，不是出于他们自己：因为他们出于自己是恶的；当他们行恶时，他们是人子；当他们行善时，他们是我的儿子。因为天主所作的是：从人子中祂造就了天主子；因为从天主子，祂成了人子。请看这参与是什么：曾应许我们参与天主性；如果祂没有先成为有死者的参与者，那应许者就是说谎。因为天主子成了有死者的参与者，为使有死的人能成为天主性的参与者。那应许要与你们分享祂的善的那一位，先与你们分享了你们的恶；那应许给你们天主性的那一位，在你们身上显示了爱。因此，除去人是天主子，剩下的就是他们是人子：\BibleQuote{没有一个行善的，连一个也没有。}\par

6. \Quote{难道所有作恶的人，就是那些吞噬我的人民如同吃食面包的人，都不知道吗？}\BibleRef{咏 52:5}……所以这里有一个正被吞噬的天主的子民。不，\Quote{没有一个行善的，连一个也没有。}我们以上述规则作答。但这被吞噬的子民，这承受恶人之害的子民，这在一群恶人中呻吟哀叹的子民，如今已经从人子变成了天主的子女：因此他们被吞噬。因为，\Quote{你挫败了穷人的计谋，因为上主是他的希望。}因为常常，为了让天主的子民被吞噬，正是这一点——他们是天主的子民——在他们身上被藐视。\Quote{我要掠夺，我要抢夺；如果他是基督徒，他能对我怎样呢？}……但下文如何？\Quote{我要指证你，并要把你放在你面前。}你现在不会知道，以致你对自己不悦；你将要知晓，以致你哀哭。因为天主不能不向不义的人显示他们的不义。如果祂不显示，那么将来说\Quote{骄傲对我们有什么益处？财富的炫耀又给了我们什么？}的人又会是谁呢？\BibleRef{智 5:8}因为那时他们将知道，而现在他们不愿意知道。\Quote{难道所有……都不知道吗？}等等。为什么祂加上\Quote{如同吃食面包}？他们吞吃我的子民，就如同吃面包一样。因为所有其他我们吃的东西，我们可以时而吃这个，时而吃那个；并不总是吃这蔬菜，并不总是吃这肉，并不总是吃这水果：但总是吃面包。那么\Quote{吞噬我的人民如同吃食面包}是什么意思？他们不停地、无间断地吞噬。\par

7. \Quote{他们不曾呼号天主。}祂安慰那呻吟的人，并且主要是通过告诫，以免他们因效法那常常得福的恶人而喜爱作恶。那已应许给你们的东西为你们存留着：他们的希望是现世的，你们的希望是未来的；但他们的希望是短暂的，你们的希望是确定的；他们的虚假，你们的真实。因为他们\Quote{未曾呼号天主。}难道这些人不是每天祈求天主吗？他们确实\Quote{不}祈求天主。请留意，如果我能在天主亲自的帮助下这么说的话。天主愿意人白白地敬拜祂，白白地爱祂，就是纯洁地爱祂；不是因为祂赐予祂自己以外的任何东西而爱祂，而是因为祂赐予祂自己。那么，一个人呼号天主是为了使自己富有，他就不是呼号天主：因为他呼号的是他所愿意来临的自身之物……但现在你愿意箱子满，而良心空：天主不充满箱子，只充满胸怀。如果内在的需要压迫你，外在的财富对你有什么益处呢？因此，那些为了世俗的安慰，为了地上的美物，为了现世的生活和尘世的幸福而呼号天主的人，并没有呼号天主。\par

8. 因此，关于他们，接下来所言是什么呢？\Quote{他们在没有恐惧的地方，怀有恐惧} \BibleRef{咏 52:6}。因为一个人若丧失财富，那是恐惧吗？那里没有恐惧，然而在这种情况下，人们却害怕。但若一个人丧失智慧，那才是真正的恐惧，而在那种情况下，他倒不害怕……你们害怕不归还金钱，却情愿失去忠信。殉道者并未夺取他人的财物，反而连自己的也轻视，以免失去忠信；而且，当他们被剥夺公权时，失去金钱还算小事；他们在受苦时也舍弃了性命：他们失去生命，为要在永生中找到它。\BibleRef{玛 10:39} 因此，他们在那本应恐惧的地方恐惧了。然而那些论基督说\Quote{祂不是天主}的人，却在没有恐惧的地方恐惧了。因为他们说：\Quote{我们若让祂继续下去，罗马人就会来，夺去我们的土地和邦国。} \BibleRef{若 11:48} 啊，愚昧和鲁钝，心里说\Quote{祂不是天主}！你们害怕失去尘世，却失去了天堂；你们害怕罗马人前来，夺去你们的土地和邦国！他们能夺去你们的天主吗？那么还剩下什么？无非是你承认，你本想保住，却因保存不善而丧失了？因为你们因杀害基督而失去了土地、邦国和基督。你们宁愿杀害基督，也不愿失去土地；结果你们失去了土地、邦国和基督。他们因恐惧而杀害了基督：但这又是为什么呢？\Quote{因为天主驱散了那些取悦于人者的骨骸。} 他们因取悦于人而害怕失去自己的土地。但基督自己——他们论及祂曾说\Quote{祂不是天主}——却宁愿不取悦像他们那样的人：他们是人子，而非天主子；祂宁愿不取悦他们。于是他们的骨骸被驱散了，祂的骨骸却无人折断。\Quote{他们蒙受了羞辱，因为天主轻视了他们。} 事实上，弟兄们，就他们而言，极大的混乱已临到他们。在他们钉死主的那个地方——他们钉死主正是为了不失去土地和邦国——犹太人已不复存在。因此，\Quote{天主轻视了他们}；然而在轻视中，祂警告他们要悔改。现在让他们承认基督，并说：他们是称之为\Quote{祂不是天主}的那一位，就是天主。让他们回归祖先的产业，即亚巴郎、依撒格和雅各伯的产业，让他们与这些先人一同拥有永生：虽然他们失去了现世的生命。为何如此？因为人子已成了天主子。只要他们仍执迷不悟，就没有一个行善的，连一个也没有。\Quote{他们蒙受了羞辱，因为天主轻视了他们。} 祂仿佛转向这些人，说：\Quote{谁能从熙雍赐给以色列救恩？} \BibleRef{咏 52:7}。啊，愚人们，你们辱骂、侮辱、击打、唾污、荆棘冠冕、高举在十字架上，是谁呢？\Quote{谁能从熙雍赐给以色列救恩？} 岂不就是你们所说\Quote{祂不是天主}的那一位吗？\Quote{在天主转变祂子民的命运时。} 因为转变祂子民命运的，除祂以外别无他人，就是那位甘愿在你们手中成为俘虏的那一位。但谁能明白这事呢？\Quote{雅各伯要欢腾，以色列要喜乐。} \Quote{以色列；} 真正的雅各伯和真正的以色列，即那弟弟，哥哥曾服事他的那一位，\BibleRef{创 25:23} 他要亲自欢腾，因为他要亲自明白。\par

\SourceRecord{Translated by J.E. Tweed.；Nicene and Post-Nicene Fathers, First Series；Vol. 8.；Edited by Philip Schaff.；Buffalo, NY: Christian Literature Publishing Co.,；1888.；<http://www.newadvent.org/fathers/1801053.htm>.}

% structure: p0001 source=h1 target=section reason=page-title

\section{论圣咏第五十四篇}

这篇圣咏的题名若被理解，在它的冗长中便含有果实；正因为圣咏简短，让我们借由在题名上逗留，来弥补我们无需在圣咏上逗留。因为所唱的每一节都依赖于此。因此，若任何人留意那固定在房屋正面之物，他将会安全地进入；而一旦进入，就不会犯错。因为这在门柱上被显著地标明，即在内里他如何能不犯错。其题名如下：\Quote{在末尾，在诗歌中，致达味自己的理解，当齐弗人来到，对撒乌耳说：看，达味不是藏在我们中间吗？}撒乌耳是圣达味的迫害者，我们非常清楚：撒乌耳代表着属世的王国，属于死亡而非生命，这点我们也记得曾对你们的爱德讲述过。并且达味自己代表着基督或基督的身体，你们既已学习，就当知道并记起。那么齐弗人呢？有一个叫齐弗的村庄，居民是齐弗人，达味曾藏在他们的地域，那时撒乌耳要寻找并杀害他。这些齐弗人得知此事后，便向迫害他的君王出卖了他，说：\Quote{看，达味不是藏在我们中间吗？}他们的出卖对自己确实毫无益处，对达味自己也无损害。因为他们的恶意被显明；但撒乌耳即使在出卖后也不能捉住达味；反而在那同一地域的一个洞穴中，当撒乌耳被交到达味手中可杀之时，达味饶恕了他，且没有做那在他能力之内的事。\BibleRef{撒上 24:4}但另一人却在寻求做那不在他能力之内的事。让那些曾是齐弗人的人留心：让我们看看圣咏借着这些人的事件向我们呈现要理解的那些人。\par

2. 如果我们询问齐弗人由什么词翻译而来，我们发现，\Quote{兴盛的人。}那么某些敌人对圣达味是兴盛的，在他隐藏之前兴盛。我们可以在人类中找到他们，如果我们愿意理解圣咏。让我们这里首先找到隐藏的达味，我们就会找到他的对手兴盛。请注意隐藏的达味：\Quote{因为你们已经死了，}宗徒对基督的肢体说，\Quote{你们的生命与基督一同藏在天主内。}\BibleRef{哥 3:3}因此，这些隐藏的人，什么时候要兴盛呢？\Quote{当基督，}他说，\Quote{你们的生命，显现时，那时你们也要与祂一同在光荣中出现。}\BibleRef{哥 3:4}当这些人兴盛时，那些齐弗人就要枯萎。因为请看他们的光荣被比作什么花：\Quote{凡有血肉的如同草，人的光荣如同草上的花。}\BibleRef{依 40:6}结局是什么？\Quote{草枯干了，花也凋谢了。}那么达味在哪里？请看下文：\Quote{但上主的话却永远常存。}……\par

3. 这些人有时被软弱的\Quote{光明之子}观察到，他们的脚站立不稳，当他们看见恶人在顺境中兴盛时，就对自己说：\Quote{纯洁于我何益？我侍奉天主，遵守祂的诫命，不压迫任何人，不掠夺任何人，不伤害任何人，尽我所能施予，这对我有什么好处？看，我做了这一切，而他们兴盛，我劳苦。}但为什么呢？你也愿意成为一个齐弗人吗？他们在世间兴盛，在审判中枯萎，枯萎后将被投入永火：你也选择这个吗？你岂不知道那来到你这里的那一位向你应许了什么，祂在这里显示了什么在自己内？如果齐弗人的花是可羡慕的，难道你们的主自己不是也可以在这世界兴盛吗？或者祂缺乏兴盛的能力？不，相反，祂在这里宁可在齐弗人中隐藏，并对般雀比拉多说——好像他也是齐弗人的一朵花，对他的王国心存猜疑——\Quote{我的王国不属于这世界。}\BibleRef{若 18:36}因此，祂在这里是隐藏的：所有善人在这里也是隐藏的，因为他们的善是内在的，是隐藏的，是在心里的，那里有信德、爱德、望德，那里有他们的宝藏。这些善事在世上显现吗？这些善事和这些善事的报酬都是隐藏的。……\par

4. \BibleQuote{天主，因祢的名拯救我，以祢的德能审判我} \BibleRef{咏 53:1}. 让教会说这话，藏身于\Person{齐弗人}之中。让基督的身体说这话，暗中保守其道德之善，暗中期待其功行之赏报，让它说这话：\Quote{以祢的德能审判我。} 祢来了，\Person{基督}，祢谦卑地显现，祢受轻贱，祢被鞭打，祢被钉十字架，祢被杀害；然而，第三日祢复活了，第四十日祢升天了：祢坐在圣父的右边，无人看见；祢从那里派遣了祢的圣神，那些配得的人领受了；充满祢的爱，他们向普世万民宣讲了祢那谦卑的赞美：我看见祢的名为人类所尊崇，然而我们仍将祢视为软弱。因为连那位\OriginalTerm{外邦人的导师}{Teacher of the Gentiles}也没有说，他在我们中间知道什么，\BibleQuote{只知道耶稣基督，并祂被钉在十字架上} \BibleRef{格前 2:2}，为的是我们宁可选择祂的羞辱，而不是昌盛的\Person{齐弗人}的荣耀。然而，祂说了什么？\Quote{祂虽因软弱而死，却因天主的德能而活着。} 祂来是要因软弱而死，祂将要在天主的德能中审判：但通过十字架的软弱，祂的名声显赫。凡不相信这因软弱而显赫的名的人，当祂带着权能来时，必敬畏审判者。但唯恐那位曾经软弱的，当祂带着力量来时，用那簸箕把我们送到左边；愿祂\Quote{以祂的名拯救我们，以祂的德能审判我们。} 因为谁会如此鲁莽而渴望这个，比如对天主说\Quote{审判我}？不是常对人咒骂说\Quote{愿天主审判你}吗？所以显然那是咒骂，如果祂以祂的德能审判你，而没有以祂的名拯救你：但当祂先以名拯救你之后，祂将以随后的德能审判你，使你健康。你不要担心：那审判对你不会是惩罚，而是分开。因为在某篇圣咏中这样说：\Quote{天主，求祢审判我，把我的案件与不圣洁的民族分开。}……\par

5. \BibleQuote{天主，求你俯听我的祈祷，求你侧耳听我口中的言语} \BibleRef{咏 53:2}. 愿我的祈祷达到祢面前，它从祢永福的渴望中发出并射出：我把它送到祢的耳中，求祢助它到达，免得它中途坠落，如昏倒般掉下。但即使现在我所求的善事没有临到我，我仍然确信它们以后必会来到。因为即使是关于过犯，据说有一个人祈求了天主，却没有被垂听，这对他反而有益。因为这世界的匮乏激发了他祈祷，在现世的苦难中，他希望现世的苦难过去，草上的花重新回来；他说：\Quote{我的天主，我的天主，祢为什么舍弃了我？} 这真是基督的声音，但为了祂的肢体。他说：\Quote{我因我的过犯向祢终日呼求，祢没有垂听；夜间我也呼求，祢却没有垂听，但这于我不是愚昧：} 也就是说，\Quote{我夜间也呼求，祢却没有垂听；然而在这件事上，祢没有垂听，于我并非愚昧，反而是智慧，使我知道我该向祢求什么。因为我所祈求的，或许是我若得到反会受害的事。} 人啊，你求财富；有多少人被财富倾覆？你怎知道财富对你有益？不是许多穷人更安全地隐于无名；一旦变富，刚一开始显耀，便成了强者的猎物？他们若继续隐藏，继续不为人知，岂不更好？这些人被寻找，不是因为他们本身，而是因为他们拥有的！所以，弟兄们，在这些现世的事物上，我们在主内劝诫和敦促你们：不要如确定之事般祈求任何事物，而要求天主知道对你们有益的事物。因为你们完全不知道什么对你们有益。有时你们以为对你们有利的反而有害，以为有害的反而有利。因为你们有病，不要指定医生放在你们身边的药物。如果\OriginalTerm{外邦人的导师}{Teacher of the Gentiles}、\Person{保禄宗徒}说：\BibleQuote{我们不知如何祈求才对} \BibleRef{罗 8:26}，何况我们呢？然而，当他自认为明智地祈求，即求主除去他肉体上的刺、撒殚的使者来击打他，免得他因启示的浩大而自高时，他听到了主说什么？他所愿的事成就了吗？不，为了成就有益的事，他听到主说：\BibleQuote{我三次求主使它离开我，祂对我说：我的恩宠为你足够了，因为德能在软弱中才成全。} \BibleRef{格后 12:8} 我已把药膏敷在伤口上；我敷的时候知道，何时揭开也知道。不要让病人从医生的手中退缩，不要给医生出主意。所有现世的事都是如此。有苦难；如果你好好敬拜天主，你就会知道祂知道每个人什么有益；有顺境；要更加小心，免得这些东西腐蚀你的灵魂，使它离开那赐予这些事物的主。……\par

6. \BibleQuote{外人起来攻击我} \BibleRef{咏 53:3}。什么是\Quote{外人}？难道达味本人不属犹大支派的犹太人吗？但齐弗这个地方本身也属于犹大支派，是犹太人的地方。那么怎么说是\Quote{外人}呢？不是指城市，不是指支派，不是指亲属，而是指花……但请看齐弗人，看他们一时兴旺。有理由说是\Quote{外人的}儿子。你躲在齐弗人中间时说什么？\Quote{以主为天主的百姓，真是有福。} 因着这种情感，当说\Quote{外人起来攻击我}时，这祈祷就上达于主的耳中。\par

7. \BibleQuote{有权势的人寻索我的性命。} 兄弟，因为他们以一种新方式要毁灭圣洁之人的族类，以及那些不寄望于今世之人的族类——所有寄望于今世的人。他们确实混杂在一起，确实一同生活。这两种人彼此非常对立：一种人只寄望于世俗之事和暂时的福乐，另一种人坚定地寄望于主天主。尽管这些齐弗人同心合意，但不要过于相信他们的同心合意：考验尚缺；当任何考验来临时，以致一个人可能因世俗之花而受责备，我不对你说他会与主教争吵，甚至不会接近教会本身，以免有草枯落。兄弟，我为什么说这些话？因为现在你们都在基督名下乐意听到，并且按你们所理解的，你们在听到话语时欢呼；你们若不理解，就不会欢呼。你们的这种理解应当结出果实。但是否结果实，考验会检验；免得当你们被称为我们的人时，通过考验你们被发现是外人，而有人说：\BibleQuote{外人起来攻击我，有权势的人寻索我的性命。} 但愿不说后面的话：\BibleQuote{他们没有将天主放在眼前。} 因为一个人眼前只有世界，何时会将天主放在眼前呢？即他如何挣得钱上加钱，如何增加羊群，如何装满粮仓，如何对他的灵魂说：\Quote{你拥有许多好东西，快乐吧，欢宴吧，尽情享受吧。} 他会将天主放在眼前吗？那位向如此夸耀、如此绽放齐弗人之花的人说：\BibleQuote{愚昧人}（即\InnerQuote{不明智的人}，\InnerQuote{无智慧的人}），\BibleQuote{今夜就要索回你的灵魂；你所预备的一切，将归谁呢？}\BibleRef{路 12:20}\par

8. \BibleQuote{因为天主帮助我} \BibleRef{咏 53:4}。他们自己不知道自己在隐藏之中。但如果他们也把天主放在眼前，就会知道天主如何帮助我。因为所有圣人都由天主帮助，但在内里，无人看见。因为正如不虔敬之人的良心是极大的刑罚，同样，虔敬之人的良心也是极大的喜乐。\BibleQuote{我们的光荣}，宗徒说，\BibleQuote{就是良心的见证。}\BibleRef{格后 1:12} 在这内在之事上，而不是在外在的齐弗之花上，那现在说\Quote{因为天主帮助我}的人夸耀。虽然他所应许的尚在遥远，但今天我已有了甜美的、现成的帮助；今天我内心的喜乐使我发现，有人毫无理由地说：\Quote{谁向我们显示好事？因为你的面容的光照在我们身上，主，你将喜乐放在我心中。} 不是放在我的葡萄园，不是放入我的羊群，不是放入我的酒桶，不是放入我的餐桌，而是\Quote{放在我心中。} \Quote{因为天主帮助我。} 他如何帮助你？\Quote{主是提拔我灵魂的主。}\par

9. \BibleQuote{将灾祸转向我的仇敌} \BibleRef{咏 53:5}。因此，无论他们如何青翠，如何茂盛，他们是为火保留的。\BibleQuote{愿你以你的能力消灭他们。}因为他们现在茂盛，像草一样生长：你不要做不明智和愚昧的人，不要因思虑这些事而永远丧亡。因为，\Quote{将灾祸转向我的仇敌。} 如果你在达味本身的身体中有份，他必以自己的能力消灭他们。这些人兴旺于世间的福乐，却消灭于天主的能力。他们茂盛的样子与他们消灭的样子不同：他们茂盛一时，消灭永远；茂盛在虚假的美好事上，消灭在真实的痛苦中。\Quote{愿以你的能力消灭}那些你以你的软弱所忍受的人。\par

10. \BibleQuote{我甘愿向你献祭} \BibleRef{咏 53:6}。除非他亲身尝过，否则别人讲述时，谁能理解这内心的美善呢？\BibleQuote{我甘愿向你献祭}是什么意思？……因为在这里我要取什么祭献呢，弟兄们？或者，我当献上什么才堪当报答主的仁慈呢？我岂要从羊群中寻找牺牲，挑选公羊，在牛群中寻觅公牛，或者从\Place{舍巴}地带乳香来？我该做什么？该献什么？除了祂所说的：\BibleQuote{赞颂的祭献将光荣我}？那么为什么\Quote{甘愿}？因为我实在爱我所赞颂的。我赞美天主，并在这赞美中喜乐；我因赞美祂而喜乐，当我赞美祂时，我并不羞愧。因为祂的赞美不同于那些喜爱戏院荒唐事的人所赞美的——无论是战车御者、猎人还是任何演员，他们的赞美者邀请、劝勉其他赞美者一同呼喊；当众人都呼喊时，如果他们宠爱的角色被打败，他们往往都会脸红。我们的天主并非如此：愿祂被以意志赞美，以爱德爱慕；愿祂被爱和被赞美都是无偿的（即甘愿的）。什么是\Quote{无偿}？是为了祂本身，而非为了其他事物。因为你若赞美天主是为了祂给你别的东西，你就不再自由地爱天主了。如果你的妻子因财富而爱你，而一旦贫穷降临，她就开始想着通奸，你会羞愧。既然你愿意被配偶无偿地爱，难道你会因别的事物而爱天主吗？贪吝的人啊，你将从天主那里得到什么赏报？祂为你存留的不是大地，而是祂自己——那位造天地者。\BibleQuote{我甘愿向你献祭}：不要出于不得已而行此事。因为你若因别的事物赞美天主，你便是出于不得已而赞美。祂所赐的这一切，因施予者而成为美物。因为祂全然赐予，祂赐予这些暂时的事物：按照祂判断的高深，对某些人为益，对某些人为害……。\BibleQuote{我甘愿向你献祭}。为什么\Quote{甘愿}？因为是\OriginalTerm{无偿}{gratis}。什么是\OriginalTerm{无偿}{gratis}？\BibleQuote{主，我要称谢你的圣名，因为这是美善的}：别无他故，只因为这是\Quote{美善的}。他难道说：\BibleQuote{主，我要称谢你的圣名}，是因为你赐我肥沃的庄园，赐我金银，赐我广大的财富、丰厚的金钱、至高的尊位吗？不。而是什么？\BibleQuote{因为这是美善的}。我找不到比你的圣名更好的事物。\par

11. \BibleQuote{因为你从一切患难中救赎了我} \BibleRef{咏 53:7}。因此我领悟了你的圣名是何等美善：因为如果我在患难之前就能承认这事，或许我就不需要它们了。但患难被用来作为告诫，告诫归向了你的赞美。因为我若不因自己的软弱受告诫，就不知道我身在何处。因此，\BibleQuote{你从一切患难中救赎了我，我的眼也回看了我的仇敌}：是回看了那些齐弗人。是的，我以心灵的高超越过了他们的花，来到了你面前，并从那里回看他们，看见\BibleQuote{凡有血肉的都像草，凡人的光荣如同草上的花}；\BibleRef{依 40:6}另有一处也说：\BibleQuote{我看见恶人高升，如黎巴嫩的香柏树。我经过，看哪！他已不在了。}为什么说\Quote{他已不在了}？因为你经过了。什么是\Quote{因为你经过了}？因为你并非徒然地听见了\Quote{举起你们的心}；因为你没有留在地上——那会使你腐朽；因为你把灵魂高举向天主，超越黎巴嫩的香柏树，并从高处观察：\Quote{看哪！他已不在了}；你寻找他，却找不到他的地方。劳累不再在你面前；因为你进入了天主的圣所，明白了终局。同样，在这里他也如此结尾：\BibleQuote{我的眼也回看了我的仇敌。}所以，弟兄们，你们也要这样对待你们的灵魂：举起你们的心，磨砺你们的心思，学习真实地爱天主，学习轻视今世，学习甘愿献上赞颂的祭品；为的是，超越草上的花，你们可以回看你们的仇敌。\par

\SourceRecord{Translated by J.E. Tweed.；Nicene and Post-Nicene Fathers, First Series；Vol. 8.；Edited by Philip Schaff.；Buffalo, NY: Christian Literature Publishing Co.,；1888.；<http://www.newadvent.org/fathers/1801054.htm>.}

% structure: p0001 source=h1 target=section reason=page-title

\section{《第五十五篇圣咏释义》}

1. 这篇圣咏的标题是：\Quote{直到终结，在赞美诗中，达味自己的智慧篇}。何谓\Quote{终结}，我们愿简要提醒诸位，因你们已知道。\Quote{因为律法的终结就是基督，使所有信者获得正义。}（\BibleRef{罗 10:4}）因此，愿注意指向终结，指向基督。为何祂被称为终结？因为我们所做的一切，都指向祂，而当回归祂时，便不再有所求。因为有一个终结指消耗，有一个终结指成全。例如，我们在听到\Quote{食物吃完了}时想到一种含义，在听到\Quote{衣服织完了}时想到另一种含义：两种都说\Quote{完了}，但食物是归于无有，衣服则是得以完成。因此，我们的终结应当是我们的成全，而我们的成全就是基督。因为我们在祂内得以成全，祂是头，我们是肢体。祂被称为\Quote{律法的终结}，因为没有祂，无人能成全律法。因此，当你们在圣咏中听到\Quote{直到终结}——许多圣咏都有这样的标题——不要想到消耗，而要想到圆满。\par

2. \Quote{在赞美诗中}：在赞颂中。因为无论我们受困扰、受困苦，或是欢乐、踊跃，祂都应受赞颂，祂在磨难中训导，在喜乐中安慰。因为基督徒的心和口不应停止赞美天主；不是只在顺境中赞美，在逆境中诽谤，而是依本篇圣咏所指示的方式：\Quote{我要时时赞美上主，祂的赞颂常在我口中。}你欢乐时，承认圣父的抚慰；你受苦时，承认圣父的惩戒。无论祂抚慰或惩戒，祂都在教导那将要承受产业的人。\par

3. 那么，什么是\Quote{达味自己的智慧篇}？我们知道，达味确是一位圣先知、以色列王、叶瑟之子（\BibleRef{撒下 23:1}）。但因他的后裔，为了我们的救恩，按血肉主耶稣基督来到（\BibleRef{罗 1:3}），所以常在预象中以此名代表祂，达味被预表性地当作基督，因为同一血肉的起源。因为在某种形式上，祂是达味之子；在另一种形式上，祂是达味之主。按血肉是达味之子，按天主性是达味之主。因为万物都由祂造成（\BibleRef{若 1:3}），达味本人也是由祂造成，祂从达味的后裔来到人间。再者，当主询问犹太人，他们以为基督是谁的儿子时，他们回答：\Quote{达味的。}主斥责犹太人，因他们停留在血肉上，忽视了天主性；祂提出一个问题来责难他们：\Quote{那么，达味自己怎么在圣神内称祂为主，说：\InnerQuote{上主对我主说……}……如果达味在圣神内称祂为主，祂怎么会是达味的儿子？}（\BibleRef{玛 22:43}）祂提出了一个问题，但并未否认祂是儿子。你们听到了\Quote{主}，但要说祂怎么是达味的\Quote{儿子}；你们听到了\Quote{儿子}，但要说祂怎么是\Quote{主}。这个问题由公教信仰来解决。怎么是\Quote{主}？因为\Quote{在起初已有圣言，圣言与天主同在，圣言就是天主。}（\BibleRef{若 1:1}）怎么是\Quote{儿子}？因为\Quote{圣言成了血肉，居住在我们中间。}（\BibleRef{若 1:14}）因此，达味在预象中是基督，但基督——正如我们常提醒你们的爱德——是头和身体；我们不应说自己与基督无关，因为我们是祂的肢体，也不应将自己视为别物；因为\Quote{二人成为一体。}（\BibleRef{创 2:24}）宗徒说：\Quote{这是伟大的圣事，但我是指基督和教会说的。}（\BibleRef{弗 5:32}）所以，整个基督是\Quote{头和身体}；当我们听到\Quote{达味自己的智慧篇}时，也要在其中理解我们自己。愿基督的肢体理解，基督在祂的肢体中理解，基督的肢体在基督中理解：因为头和肢体是同一个基督。头在天上，曾呼喊：\Quote{你为什么迫害我？}（\BibleRef{宗 9:4}）我们藉着希望已与祂在天上，祂藉着爱在人间与我们同在。因此，\Quote{达味自己的智慧篇}。当我们听到时，愿我们受训诫，让教会理解：因为我们需十分留意，要明白我们现在处于何种邪恶中，并渴望从何种邪恶中被解救出来，要记住主的祈祷，其结尾说：\Quote{救我们免于凶恶。}（\BibleRef{玛 6:13}）因此，在这世界的许多磨难中，这篇圣咏以某种理解抱怨。没有理解的人不会与之同哀。此外，亲爱的诸位，我们应当记住：我们是按天主的肖像被造的，而且是在理解本身上。因为在许多事上，我们被野兽超越；但当人知道自己按天主的肖像被造时（\BibleRef{创 1:26}），他承认在自己内有某种超越畜生的东西。人思量一切所有，发现他与畜生的独特区别在于他有理解。因此，某些人轻视这从造物主领受的独特而特殊之物，造物主亲自斥责说：\Quote{不要像马骡那样无知。}……\par

4. \Quote{天主，求祢俯听我的恳求，不要轻忽我的祈祷：求祢垂顾我，应允我}\BibleRef{咏 54:1}。这些是一个诚恳、焦虑、处于困苦中的人的话。他正在祈祷，遭受许多苦难，渴望从邪恶中被解救。剩下的，我们须要听见他处于何种邪恶之中，当他开始说话时，让我们承认自己也在那里；好使我们同受困苦，共同祈祷。\Quote{我在我的\OriginalTerm{操练}{exercise}中悲伤，且被困扰}\BibleRef{咏 54:2}。在哪里悲伤，在哪里困扰？他说：\Quote{在我的操练中。}他提到了他所忍受的恶人，并将这种对恶人的忍受称为他的\Quote{操练。}你们不要以为世上有恶人是没有益处的，也不要以为天主不会利用他们行善。每个恶人活着，要么是为了被纠正，要么是为了通过他来操练善人。哦，但愿那些如今操练我们的人能够悔改，并与我们一同被操练！然而，只要他们仍是操练我们的人，我们就不要恨他们；因为我们不知道他们中哪个人会坚持到底作恶；常常当你以为自己恨的是一个仇敌时，你恨的却是一个兄弟，而你并不知道。魔鬼和他的使者在天主圣经中已被启示给我们，他们命定受永火。只有对他们，我们才应放弃改正的希望。……因此，既然这爱的规矩已为你定下，即效法圣父要爱仇敌，因为祂说：\Quote{爱你们的仇敌}\BibleRef{路 6:27}；若没有仇敌可忍受，你在这条诫命中如何被操练？你看到了，他在某种程度上对你是有益的；愿天主容忍恶人也有益于你，使你显出仁慈；因为也许你，如果你是个善人，恰是从恶人变为善人的；倘若天主不容忍恶人，连你也不会被发现应当感恩。愿祂因此容忍他人，因为祂也容忍了你。因为当你经过之后，关闭虔敬的道路是不合适的。\par

5. 那么，这个人置身恶人中，被他们的敌意所操练，他是如何祈祷的呢？为什么他说：\Quote{我在我的\OriginalTerm{操练}{exercise}中悲伤，且被困扰}？当他扩展他的爱以致爱仇敌时，他感到厌恶，被许多人的敌意四面围困，被许多人的疯狂所迫，在某种人性的软弱下沉沦了。他看到自己开始被魔鬼的邪恶暗示刺穿，从而对仇敌产生仇恨；在与仇恨搏斗以成全爱本身的过程中，在战斗中，在摔跤中，他被困扰了。因为他在另一篇圣咏中有这样的声音：\Quote{我的眼睛因愤怒而困扰。}那里接着说什么？\Quote{我在我所有的仇敌中衰老了。}仿佛在风暴和波浪中，他开始下沉，如同伯多禄一样。\BibleRef{玛 14:30} 因为他踏过这世界的波浪，就是爱仇敌者。基督在海面上无畏地行走，从祂心中绝不可能夺走爱仇敌的心，祂悬挂在十字架上时说：\Quote{父啊，宽恕他们，因为他们不知道他们所做的是什么。}\BibleRef{路 23:34} 伯多禄也想步行。祂为元首，伯多禄为身体：因为祂说：\Quote{在这磐石上，}\Quote{我要建立我的教会。}\BibleRef{玛 16:18} 他被命令行走，并依靠命令者的恩宠行走，而非自己的力量。但当看见风势猛烈时，他害怕了；然后他开始下沉，在他的\OriginalTerm{操练}{exercise}中被困扰。靠什么猛烈的风？\Quote{因仇敌的声音，因罪人的压迫。}\BibleRef{咏 54:3} 因此，正如他在波浪中呼喊：\Quote{主啊，我要丧亡了，救我！}\BibleRef{玛 14:30} 相似的声音已从这人发出：\Quote{求你应允我。}为何？你为何受苦？你为何呻吟？\Quote{我在我的操练中悲伤。}你确实将我安置在恶人中受操练，但他们过于兴起，超过我的力量：求你平静这被困扰的人，向这下沉的人伸出援手。\Quote{因为他们向我倾注不义，且以愤怒向我投下阴影。}你已听到波浪和风声：一个仿佛被他们侮辱的谦卑者，他正在祈祷；四面被侮辱的咆哮攻击，但他内在呼喊着他们所看不见的那一位……\par

6. 但此人虽烦恼忧伤，仍祈祷，他的眼因愤怒而似乎被扰乱。但若兄弟的愤怒成为积习，那就是仇恨。愤怒扰乱眼睛，仇恨却熄灭它：愤怒是一根稻草，仇恨是一根梁木。有时你恨一个人并责骂一个愤怒的人：在你内是仇恨，在你所责骂的人内是愤怒。因此，有理地对你说：\BibleQuote{先取出你眼中的梁木，然后你才能看清，取出你兄弟眼中的稻草。}\BibleRef{玛 7:5}为使你知道愤怒与仇恨之间的巨大差异：人每天都对自己的儿子发怒，但你能找出一个恨自己儿子的人吗？此人烦恼时仍祈祷，即使忧伤，也抵挡所有毁谤者的毁谤；不是为藉回击毁谤而胜过他们，而是为不恨他们中的任何人。因此他祈祷，因此恳求：\Quote{从仇敌的声音和罪人的磨难中。}\BibleQuote{我的心在我内烦乱。}\BibleRef{咏 54:4}这正如别处所说的：\Quote{我的眼因愤怒而烦乱。}若眼烦乱了，接着是什么？\Quote{死亡的恐惧落在我身上。}我们的生命是爱：若生命是爱，死亡就是仇恨。当人开始害怕自己会恨他所爱的人时，他害怕的正是死亡；一种更锋利、更内在的死亡，灵魂由此被杀，而非身体。你注意一个向你发怒的人；他能做什么？你的主已给你保证，说：\BibleQuote{你们不要害怕那杀害身体的。}\BibleRef{玛 10:28}他藉发怒杀害身体，你藉保持仇恨杀害了灵魂；他杀害别人的身体，你杀害了自己的灵魂。因此，\Quote{恐惧，}因此，\Quote{死亡的恐惧落在我身上。}\par

7. \BibleQuote{战栗和恐惧临到我，黑暗笼罩了我。}\BibleRef{咏 54:5} \Quote{于是我说：}\BibleQuote{凡恨自己弟兄的，就是在这黑暗中，直到现在。}\BibleRef{若一 2:9} 若爱是光，恨就是黑暗。那位处在此软弱中并在此操练中烦乱的人，对自己说什么？\Quote{谁能赐我鸽子般的翅膀？我要飞去安息。}\BibleRef{咏 54:6} 他或是盼望死亡，或是渴望独处。他说，只要我从事这工作，只要这命令给我，即我应爱仇敌，这些人的毁谤越来越多遮蔽我，就扰乱我的眼，干扰我的视线，穿透我的心，杀死我的灵魂。我情愿离去，但我软弱，以免因留下而罪上加罪；或者至少让我与人群分离片刻，以免我的伤口因频繁的打击而受苦，好使它痊愈后能重回操练。弟兄们，这就是所发生的事，天主的仆人心中时常产生对独处的渴望，唯一的原因就是由于多重的磨难和绊脚石，他就说：\Quote{谁能赐我翅膀？}他发现自己没有翅膀，还是翅膀被束缚？若缺少，就赐予；若被束缚，就解开；因为即使是松开鸟翅膀的人，也是给予或归还它的翅膀。因为它原先好像没有自己的能飞的翅膀。被束缚的翅膀是重担。\Quote{谁，}他说，\Quote{能赐我鸽子般的翅膀？我要飞去安息。}在哪里安息？我说这里有两个意义：要么，如宗徒所说，\BibleQuote{脱离肉身，与基督在一起，这实在是再好不过的。}\BibleRef{斐 1:23}……甚至那不能改变的，是属于你的，或是因人类共融，或是因教会的共融；他在内，你将怎么办？你将去哪里？你将怎样使自己分离，以免遭受这些？但到他那里去，说话，劝勉，安慰，恐吓，责备。我已做了一切，尽了我的能力，并已耗尽，我看我什么也没有成就；我所有的劳力都已耗尽，忧伤仍在。那么我的内心怎能从这样的人中得到安息，除非我说：\Quote{谁能赐我翅膀？} \Quote{如同鸽子，}然而，不是如同乌鸦。鸽子寻求从烦恼中飞走，但她不失去爱。因为鸽子被设定为爱的象征，在她的哀鸣中受人喜爱。没有什么像鸽子那样喜爱哀鸣：她日夜哀鸣，仿佛她在此被安置，正应哀鸣。那么这位爱者说什么？我无法忍受人的毁谤，他们咆哮，狂乱被带走，被忿恨点燃，在愤怒中遮蔽我；我不能为他们做好事；啊，但愿我能在某处安息，在身体上与他们分离，而不是在爱中；以免爱本身在我内被扰乱：用我的言语和言谈，我不能使他们受益，或许为他们祈祷，我就能受益。人们说这些话，但往往他们被这样束缚，以致不能飞去。因为他们或许不是被粘鸟胶束缚，而是被责任束缚。但若他们被关心和责任束缚，并且不能离开，就让他们说：\BibleQuote{我情愿脱离肉身，与基督在一起，这实在是再好不过的；但因你们的需要，留在肉身是必要的。}\BibleRef{斐 1:23}一只被情感拉住、而非被贪欲拉住的鸽子，因要履行的责任而不能飞去，并非因功德小。不过心中必须怀有渴望；没有人会忍受这种渴望，除非他开始了行走在那窄路上：\BibleRef{玛 7:14} 为使他知晓教会不乏迫害，甚至在此时期，当在教会中看到平静时，至少就我们的殉道者所受的迫害而言。但迫害并不缺少，因为这话是确实的：\BibleQuote{凡愿意在基督耶稣内热心生活的，都要遭受迫害。}\BibleRef{弟后 3:12}……\par

8. \Quote{看，我逃往远处，住在旷野} \BibleRef{咏 54:7}。在什么旷野？无论你在哪里，别人也会聚集到那里，他们会与你一同寻找旷野，依附你的生活，你无法摆脱弟兄们的团体：恶人也混杂在你中间；你的本分仍是操练。\Quote{看，我逃往远处，住在旷野。}在什么旷野？或许是良心中，那里无人进入，无人与你同在，只有你和天主。因为若是在某个地方的旷野，当人们聚集时，你该怎么办？只要你还生活在人中间，你就无法与人类分离。\par

9. \Quote{我期盼那救我脱离心灵软弱和风暴的} \BibleRef{咏 54:8}。那里有海，有风暴：你只剩下呼喊：\Quote{主啊，我快丧命了} \BibleRef{玛 14:30}。愿那无畏地踏浪者伸出援手，解除你的恐惧，使你在祂内坚定安稳，并在你内心对你说：\Quote{留意我，我所忍受的：}你或许在忍受一个恶毒的弟兄，或一个外在的仇敌；这些有哪一样我没有忍受过？外有犹太人咆哮，内有门徒出卖。因此风暴肆虐，但祂救人脱离心灵软弱和风暴。或许你的船正因祂在你内沉睡而受颠簸。海在翻腾，门徒们航行的小船正被抛掷；但基督在睡觉；他们终于意识到，在他们中间睡着的是风和浪的统治者与创造者；他们唤醒基督 \BibleRef{玛 8:24}；祂斥责风，便大大平静了。因此，你的心被搅扰，或许是因为你忘记了所信的那位；你忍无可忍，是因为你没有想起基督为你承受了什么。若基督没有进入你的心，祂就在睡觉；唤醒基督，重拾信德。因为如果你忘记了基督的苦难，基督就在你内睡着了；如果你记起了基督的苦难，基督就在你内醒着。但当你怀着全心默想祂所承受的，难道你也不会平静地忍受吗？或许还会喜乐，因为你发现自己与你的君王所受的苦难相似。因此，当你思考这些事而开始得到安慰和喜乐时，祂已起来，斥责了风，于是大大平静了。\Quote{我期盼那救我脱离心灵软弱和风暴的。}\par

10. \Quote{主啊，求你使他们沉没，并分裂他们的舌头} \BibleRef{咏 54:9}。他是指那些困扰和遮蔽他的人，弟兄们，他并非出于愤怒而希望此事。那些邪恶地自高之人，沉没对他们有益。那些邪恶地同谋之人，分裂他们的舌头对他们有益：愿他们同意善事，愿他们的舌头彼此一致。但如果所有仇敌都同谋暗中攻击我，他说，愿他们失去那邪恶的\Quote{同一心意}，愿他们的舌头被分裂，不要彼此同心。\Quote{主啊，求你使他们沉没，并分裂他们的舌头。}为何\Quote{沉没}？因为他们自高自大。为何\Quote{分裂}？因为他们为恶事联合。回想洪水之后那些骄傲的人建造的高塔：骄傲的人说：\Quote{免得我们被洪水消灭，让我们建造一座高塔} \BibleRef{创 11:4}。他们骄傲地以为自己得到坚固，建造了一座高塔，主就分裂了他们的舌头。于是他们开始彼此无法理解；由此产生了众多语言的起源。因为先前只有一种语言：但一种语言对合一的人是有益的，一种语言对谦卑的人是有益的；但当那一群人坠入骄傲的联盟时，天主怜悯了他们；祂分裂了舌头，免得他们因彼此理解而制造毁灭性的合一。通过骄傲的人，舌头被分裂；通过谦卑的宗徒，舌头被合一。骄傲之神分散语言，圣神合一语言。因为当圣神降临于门徒时，他们用各种语言说话 \BibleRef{宗 2:4}，被众人理解：分散的语言合而为一。因此，如果外邦人仍在咆哮，分裂他们的舌头对他们有益。他们若愿有同一舌头，让他们来到教会；因为在肉体语言的多样性中，心的信德中只有一种语言。\par

11. \Quote{因为我看见了城中的不义与悖逆。}此人寻求旷野是有理由的，因为他看见了城中的不义与悖逆。有一座骚动的城：就是那建造高塔的城，也是那被混乱并称为\Place{巴比伦}的城，它分散在无数民族中\BibleRef{创 11:9}；从那里，教会聚集到良心的旷野。因为他看见了城中的悖逆。\Quote{基督来了。}——\Quote{什么基督？}你悖逆——\Quote{天主子。}——\Quote{天主有儿子吗？}你悖逆——\Quote{祂由童贞女所生，受难，复活了。}——\Quote{这事怎能成就？}你悖逆。——至少留心十字架本身的荣耀。如今，那曾被敌人侮辱的十字架，已固定在君王的额上。效果证明了它的德能。它不是用铁，而是用木头征服了世界。十字架的木料曾被敌人视为可辱骂的，他们站在十字架前摇头说：\BibleQuote{如果祂是天主子，就让祂从十字架上下来吧。}\BibleRef{玛 27:40}祂向不信与悖逆的百姓伸出了双手。因为\BibleQuote{义人因信德而生活}\BibleRef{罗 1:17}，没有信德的人就是不义的。此处所说的\Quote{不义}，我理解为不信。主于是看见城中的不义与悖逆，向不信与悖逆的百姓伸出了双手；然而祂仍等候他们，说：\BibleQuote{父啊，宽恕他们，因为他们不知道自己在做什么。}\BibleRef{路 23:34}直到如今，那城的残余仍在狂怒，仍在悖逆。从所有人的额上，祂仍向不信与悖逆的残余伸着双手。\par

12. \Quote{不义与劳苦必昼夜环绕其城垣。}\Quote{在其城垣上}，即在其堡垒上，在其首领、尊贵之人身上。倘若那位尊贵之人成了基督徒，就不会再有异教徒了！人们常说：\Quote{如果他成了基督徒，就没有人会仍是异教徒。}又常说：\Quote{如果他也成了基督徒，谁还会是异教徒？}因为他们尚未成为基督徒，所以他们仿佛是那不信与悖逆之城的城垣。这些城垣能站立多久？不会永远站立。约柜正环绕\Place{耶里哥}的城墙；到了第七次绕城时，那不信与悖逆之城的全部城垣都要倒塌。\BibleRef{苏 6:5}在这一切发生之前，此人在试炼中感到困扰；他忍受着悖逆之人的余党，宁愿选择翅膀飞去，选择旷野的安歇。然而，让他继续留在悖逆之人中间，忍受威胁，喝尽辱骂，期待那拯救他脱离软弱心灵和风暴的那一位；让他仰望元首——他生活的模范，在望德中平静，即使事实上困扰。\Quote{不义与劳苦必昼夜环绕其城垣；其中是劳苦与不义。}因此那里有劳苦，因为那里有不义；因为那里有不义，所以也有劳苦。但让他们听到那伸出双手者的话：\BibleQuote{凡劳苦的，都到我这里来。}\BibleRef{玛 11:28}你们呼喊、悖逆、辱骂；祂却说：\BibleQuote{凡劳苦的，都到我这里来，}你们因骄傲而劳苦，要在我的谦卑中获得安息。祂说：\BibleQuote{你们跟我学习，因为我良善心谦，你们必会找到灵魂的安息。}\BibleRef{玛 11:29}他们为何劳苦？不正因为他们不善良、不心谦吗？天主成了谦卑的，人岂可骄傲而不脸红呢？\par

13. \BibleQuote{它的街市上没有断绝过高利贷和欺诈} \BibleRef{咏 54:11}。高利贷和欺诈至少没有隐藏，因为它们是坏事，但它们在公开场合肆虐。因为凡是在自己家里做坏事的人，至少会为自己的坏事脸红：\BibleQuote{它的街市上有高利贷和欺诈。}放贷甚至成为一种职业，放贷甚至被称为一门学问；人们谈论行会，一个仿佛对国家必要的行会，而其行业还缴纳赋税；所以那本应隐藏的，如今完全在大街上。还有另一种更坏的高利贷，就是你不宽恕欠你的债；你的眼睛在这句祷词中受到困扰：\BibleQuote{宽免我们的罪债，如同我们也宽免欠我们债的人。}因为当你要祈祷，来到这同一句祷词时，你在那里做什么？你听到了一句辱骂的话：你就想要判处谴责的惩罚。你只要同意索取和你所给的一样多，你这个伤害的放高利贷者！你被拳头打了，你想要杀人，邪恶的高利贷！你将如何祈祷？如果你停止祈祷，你怎能回转归向主？看，你会说：\BibleQuote{我们在天的父，愿你的名被尊为圣，愿你的国来临，愿你的旨意奉行在人间，如同在天上。}你要说：\BibleQuote{我们的日用粮，今天赐给我们。}你会来到：\BibleQuote{宽免我们的罪债，如同我们也宽免欠我们债的人。} \BibleRef{玛 6:9}即使在那邪恶的城市里，让这些高利贷盛行吧；但不要让它们进入那捶胸的城墙内！你要做什么？因为你和那节经文都在中间？一位天上的律师为你拟定了祈求。他知道那里常发生的事，便对你说：\Quote{否则你们不能获得。} \BibleQuote{我实实在在告诉你们：如果你们宽免人的罪过，他们也要被宽免；但如果你们不宽免人的罪过，你们的父也不会宽免你们。} \BibleRef{玛 6:14}这是谁说的？他知道你在那里做什么，在你站着祈求的地方。看，他如何甘愿成为你的辩护人；他自己是你的顾问，他自己是圣父的陪审官，你自己是你的审判官，曾说过：\Quote{否则你们不会领受。}你要做什么？你不会领受，除非你说话；如果你说谎，你也不会领受。因此，你必须要么做并说，否则你所求的将不会获得；因为不这样做的人，就在那些邪恶的高利贷中间。让那些仍然崇拜或渴望偶像的人从事这些吧；不要这样，天主的子民，不要这样，基督的子民，不要这样，你们这些作为头的身体！留意你们和平的纽带，留意你们生命的应许。因为你对所遭受的伤害索取赔偿，有什么益处呢？报复会使你振作吗？因此，你要因别人的恶而高兴吗？你遭受了恶；你宽恕吧；不要你们两人……\par

14. \BibleQuote{如果仇敌辱骂了我} \BibleRef{咏 54:12}。确实，上面他因仇敌的声音和罪人的磨难而\Quote{在操练中困扰}，或许是被放在那城中，那骄傲的、建造一座塔的城市，那城\Quote{沉没}了，致使舌头分散：留意他因假弟兄的危险而发出的内在呻吟。\BibleQuote{因为如果仇敌辱骂了我，我必定会忍受；如果恨我的人向我夸口，}即因骄傲践踏我，自高自大，用尽其权力威胁我：\BibleQuote{我必定会躲避他。}对于外面的人，你会在哪里躲藏？在那些里面的人中间。但现在看还有什么别的出路，除了你寻求孤独。\BibleQuote{但是你，}他说，\BibleQuote{同心合意的人，我的向导和我的朋友} \BibleRef{咏 54:13}。或许有时你给了我好的劝告，或许有时你走在我前面，并给了我一些有益的忠告：我们在天主的教会中曾在一起。\BibleQuote{但是你，……与我一同品尝甘美食物的} \BibleRef{咏 54:14}。这些甘美的食物是什么？并非所有在场的人都知道：但让知道的人不要变得酸涩，以便他们能够对尚不知道的人说：\BibleQuote{你们要尝尝并观看，主是何等的甘甜。} \BibleQuote{我们在天主的殿中同心同行。}那么分歧从何而来？你原是里面的，却成了外面的。他曾与我在天主的殿中同心同行：却另立了一座殿来对抗天主的殿。为什么我们同心同行的殿被离弃？为什么我们一同品尝甘美食物的殿被抛弃？\par

15. \Quote{愿死亡降至他们身上，愿他们活活地下入地狱} \BibleRef{咏 54:15} 。祂如何引述并使我们回想起分裂的最初开端，那时在那第一批犹太人中有些骄傲的人自高自大，宁愿在外面献祭？一种新的死亡降临他们：地裂开口，活活地吞没了他们。 \BibleRef{户 16:31} \Quote{愿死亡}，祂说，\Quote{降至他们身上，愿他们活活地下入地狱。} \Quote{活活}是什么意思？明知自己正在灭亡，却仍灭亡。请听活生生的人灭亡，被吞入地下的深渊，即被吞噬于尘世欲望的贪婪之中。你对一个人说：\Quote{弟兄，你有什么事？我们是兄弟，我们呼求同一天主，信仰同一基督，聆听同一福音，咏唱同一圣咏，回答同一阿们，高唱同一阿肋路亚，庆祝同一逾越节：你为什么在外面，而我却在里面？}有时一个人感到困窘，觉察到指控多么真实，便说：\Quote{愿天主追讨我们的祖先！}因此他活活地灭亡了。接着你继续这样警告：至少让分裂的恶独自存在，为什么还要加上重领圣洗的恶？请在我身上承认你所拥有的；如果你恨我，请在你的基督内宽恕我。这个恶事常常而且极大地令他们不悦……因为他们手中持有圣经，并且通过日常阅读深知天主教会遍布普世，以致一切反对都徒劳；他们深知找不到任何支持他们分裂的依据：因此他们活活地降入低处，因为他们所做的恶，他们明知是恶。但前者被天主义怒的火所吞噬。因为他们被争斗的欲望所点燃，不愿离开他们的恶领袖。火上浇油，纷争的热度加上吞噬的火。\Quote{因为邪恶在他们的住所中，在他们中间。} \Quote{在他们的住所中}，他们暂居并消逝之处。因为他们不会常在此处；然而他们为了保卫一时的敌意如此激烈地争斗。\Quote{在他们的住所中有不义；在他们中间有不义：}没有哪一部分像他们的心那样靠近他们的中心。\par

16. \Quote{因此我向主呼号} \BibleRef{咏 54:16} 。基督的身体与基督的一体性，在痛苦、疲惫、不安、磨炼的患难中，那一位人，一体性在同一个身体内，当祂在从地极呼号时使自己的灵魂疲惫时，说：\Quote{我从地极向你呼号，当我的心受困扰时。}祂自己是一位，但那一体性是一位！并且祂自己是一位，不是在一处地方的一位，而是从地极如同一人呼号。如何从地极有一人呼号，除非是一体？\Quote{我向主呼号。}你当向主呼号，不要向\Person{多纳图}呼号：免得他为你代替主成为主人，而在主下他不愿作同仆。\par

17. \Quote{在晚上、在早晨、在中午，我要述说传报，祂必垂听我的声音} \BibleRef{咏 54:18} 。你当传报喜讯，不要隐藏你所领受的，在\Quote{晚上}（论已往之事），在\Quote{早晨}（论将来之事），在\Quote{中午}（论永恒之事）。因此，祂所说\Quote{在晚上}属于祂所述说的；祂所说\Quote{在早晨}属于祂所传报的；祂所说\Quote{在中午}属于祂的声音被垂听的。因为终点在中午；也就是说，自此不再至于日落。因为在中午，光明极高，智慧的光辉，爱的热忱。\Quote{在晚上、在早晨、在中午。} \Quote{在晚上}，主在十字架上；\Quote{在早晨}，在复活中；\Quote{在中午}，在升天中。我要在晚上述说祂受死的忍耐，在早晨传报祂复活的生，我要祈求祂在中午俯听，祂坐在圣父的右边。祂必垂听我的声音，那为我们转求的。 \BibleRef{罗 8:34} 这个人是多么大的保障。多么大的安慰，多么大的避难所，\Quote{脱离心智的软弱和风暴}，抵挡恶人，抵挡不虔敬的人，无论他们是在外面还是里面，而那些在外面的人，虽然曾是在里面的。\par

18. 因此，我的弟兄们，你们在这些墙壁的集会中所见到的那些叛逆的人、骄傲的人、寻求自己利益的人、自高自大的人；他们没有对天主的纯洁、健全、安静的热忱，反而自我标榜；他们预备分裂，只是没有找到机会；这些人是主的禾场上的糠秕。\BibleRef{玛 3:12} 从那里，骄傲之风已将这几个人吹散；整个禾场不会飞走，除非祂在最后扬场。但我们能做什么呢？唯有与这人一同歌唱，与这人一同祈祷，与这人一同哀叹，并安全地说：\Quote{祂必在平安中救赎我的灵魂}\BibleRef{咏 54:18}。针对那些不爱和平的人：\Quote{祂必在平安中救赎我的灵魂。} \Quote{因为我与那些憎恨和平的人，曾是缔造和平的人。} \Quote{祂必在平安中救赎我的灵魂，脱离那些临近我的人。} 因为那些远离我的人，情况容易处理：他说\Quote{来，向偶像祈祷}，这不会很快欺骗我：他离我很远。你是基督徒吗？他说，是基督徒。从邻近的地方，他是我的对手，他就在眼前。\Quote{祂必在平安中救赎我的灵魂，脱离那些临近我的人：因为他们曾在许多事上与我同在。} 我为何说\Quote{临近我}？因为\Quote{他们曾在许多事上与我同在。} 在这节经文中出现两个命题。\Quote{他们曾在许多事上与我同在。} 圣洗圣事我们两者都有，在那事上他们与我同在；福音我们都读，在那事上他们与我同在；殉道者的庆节我们一同庆祝，他们在那事上与我同在；复活节的隆重礼仪我们参加，他们在那事上与我同在。但他们并非完全与我同在：在裂教中不与我同在，在异端中不与我同在。在许多事上与我同在，在少数事上不与我同在。但在这少数不与我同在的事上，那许多他们与我同在的事对他们毫无益处。请看，弟兄们，宗徒保禄列举了多少事：他说，如果缺少一样，那些事便都是枉然。他说：\Quote{我若能说人间的语言和天使的语言，} \Quote{若我有先知之恩、一切信德和一切知识；若能移山，若能把我所有的财物分施给穷人，若我牺牲自己的身体以至被火焚烧。}\BibleRef{格前 13:1} 他列举了多少事！在所有这些许多事中，只要缺少一样，就是爱德；前者在数量上更多，后者在分量上更重。因此，在所有圣事中，他们与我同在，但在爱德上不与我同在：\Quote{他们曾在许多事上与我同在。} 另一种表达方式：\Quote{因为他们曾在许多事上与我同在。} 那些从我自己分离出来的人，他们与我同在，不是在少数事上，而是在许多事上。因为在全世界，麦子少，糠秕多。因此他说什么？他们在糠秕中与我同在，在麦子中不与我同在。糠秕与麦子关系密切，从同一粒种子长出，在同一块田里扎根，被同一雨水滋养，遭受同一收割者，承受同一打场，等待同一扬场，\Emph{但}不会进入同一谷仓。\par

19. \Quote{天主必听见我，并要羞辱他们，祂是从永远存在的。}\BibleRef{咏 54:19} 因为他们依靠他们某个领袖，那领袖不过是昨天才开始的。\Quote{祂是从永远存在的，祂要羞辱他们。} 因为即使从时间上说，基督出自童贞玛利亚，但祂是从永远存在的：\Quote{在起初已有圣言，圣言与天主同在，圣言就是天主。}\BibleRef{若 1:1} \Quote{祂是从永远存在的，祂要羞辱他们。因为他们没有改变：} 我\Quote{指的是那些没有改变的人。} 祂知道有些人会坚持到底，并在自己的邪恶之坚持中死去。因为我们看见他们，他们没有改变：那些在同样的悖逆中、在同样的裂教中死去的人，他们没有改变。天主必羞辱他们，在惩罚中羞辱他们，因为他们因分裂而自高自大。他们没有改变，因为他们没有变得更好，反而变得更坏：无论在今生，还是在复活时。因为众人都要复活，但并非所有人都要改变。为什么？因为\Quote{他们没有改变，也不敬畏天主。}……\par

20. \Quote{祂伸手报复}\BibleRef{咏 54:20}。\Quote{他们亵渎了祂的盟约。} 请读他们所亵渎的盟约：\Quote{地上万民都要因你的后裔而蒙受祝福。} 而你对抗立约者的这些话，说什么？神圣的多纳图斯的阿非利加独得了这恩宠，在他那里存留着基督的教会。你至少说多纳图斯的教会。为何加上\Quote{基督的}？关于基督，经上说：\Quote{地上万民都要因你的后裔而蒙受祝福。} 你要跟随多纳图斯吗？那就撇开基督，然后分裂。请看接着的话：\Quote{他们亵渎了祂的盟约。} 什么盟约？对亚巴郎说了这许诺，以及对他的后裔。宗徒说：\Quote{然而，一个人的盟约既经确立，谁也不能取消或增加。对亚巴郎说了这许诺，以及对他的后裔。并没有说\InnerQuote{对后裔们}，好像为许多人；而是向一个说\InnerQuote{对你的后裔}，就是基督。}\BibleRef{迦 3:15} 因此，在这个基督里，有什么盟约被应许了？\Quote{地上万民都要因你的后裔而蒙受祝福。} 你放弃了所有民族的合一，只留在部分中，你亵渎了祂的盟约。……\par

21. \Quote{\BibleQuote{他的心已临近}} \BibleRef{咏 54:22}。我们理解为这是指谁呢，除非是指祂，因祂面容的怒火他们便被分裂？如何\Quote{他的心已临近}？如此，好使我们明白祂的旨意。因为藉着异端者，天主教会得到了彰显；藉着那些思想邪恶者，思想良善者得到了验证。因为许多事隐藏在圣经中；当异端者被剪除之后，他们以问题搅扰了天主的教会：于是那些隐藏的事便被开启，天主的旨意也被明白了。因此，在另一篇 \WorkTitle{圣咏} 中说道：\BibleQuote{好使那些用银炼净的人被驱逐出去}。因为祂说，让他们被驱逐，让他们出来，让他们显现。因此，在银匠中，有些被称为\OriginalTerm{排除者}{excluders}，即从混沌的团块中压出形状的人。因此，许多能够非常出色地理解和诠释圣经的人，曾隐藏在天主的子民中；但当没有毁谤者再逼迫他们时，他们并未宣告难题的解答。因为在亚略人对此狂吠之前，是否曾完美地论及天主圣三？在诺瓦提安人反对之前，是否曾完美地论及悔改？同样，在重洗者分离出去并反对之前，关于圣洗圣事的讨论并不完善；关于基督的合一性，那些已经陈述的教义也未被清楚地阐明，除非在这分裂开始压迫软弱者之后。为了那些知道如何讨论和解决这些问题的人（免得软弱者因不虔敬者的问难而困惑丧亡），藉着他们的讲论和辩论，将律法中幽暗的事带到光天化日之下……请看宗徒以何种方式将这隐晦的含义带进光明；他说：\BibleQuote{需要也有异端，好使经考验的人在你们中间显明出来}。\BibleRef{格前 11:19} 什么是\Quote{经考验的人}？用银考验，用圣言考验。什么是\Quote{显明出来}？就是被带出来。这是为什么？因为异端者。所以，这些人也因此\Quote{因祂面容的怒火而被分裂，祂的心已临近}。\par

22. \Quote{\BibleQuote{祂的言语比油还要柔和，它们本身就是利箭}} \BibleRef{咏 54:21}。因为圣经中的某些事看似艰难，因为它们隐晦；当被解释时，它们便变得柔和了。因为甚至在基督的门徒中，最初的分裂也仿佛是因祂话语的艰难而起的。因为当祂说：\BibleQuote{谁若不吃我的肉，不喝我的血，在他内便没有生命}，他们不明白，便彼此说：\Quote{这话生硬，谁能听得下去呢？}他们说\Quote{这话生硬}时，便离开了祂；祂与其余的十二人留在一起。当他们告诉祂，他们因祂的话而跌倒时，祂说：\Quote{你们也愿意走吗？}于是伯多禄说：\BibleQuote{你有永生的话，我们去投奔谁呢？}请注意，我们恳求你们，你们这些小孩子要学习虔敬。伯多禄那时是否理解了主话语的奥秘？他尚未理解；但他虔敬地相信他并不明白的话是美善的。因此，如果一句话是生硬的，且尚未被理解，让它在不虔敬的人面前是生硬的，但对你而言，让它藉着虔敬变得柔和；因为每当它被解决，它对你就会成为油，甚至渗入骨髓。\par

23. 此外，正如伯多禄在他们因主话语的艰难（如他们所想）而跌倒之后仍说：\BibleQuote{我们去投奔谁呢？}，同样，祂接着说：\BibleQuote{将你的忧虑托付于上主，祂必扶持你} \BibleRef{咏 54:22}。你是个小孩子，尚未明白话语的奥秘；或许面包向你隐藏，你还须用奶喂养：\BibleRef{格前 3:1} 不要对乳房发怒；它们会使你足以承受餐桌，如今你尚不足以承受它。看，藉着异端的分裂，许多艰难的事已被软化：祂的言语原是生硬的，现已比油柔和，它们本身就是利箭。它们武装了宣讲福音的人；这些话语本身就射向每个听者的胸膛，藉着那些及时和不合时宜地坚持的人；藉着那些话语，那些言词，犹如箭矢，人们的心被击向爱慕和平。它们原是生硬的，已被变得柔和。变软后并未失去它们的力量，反而转化成了利箭……请将自己投入主怀。看，你将投入主怀，不要让任何人取代主的位置。\Quote{将你的忧虑托付于上主}……\par

24. 但对于另一些人呢？\BibleQuote{但是天主，祢必使他们堕入灭亡的深坑} \BibleRef{咏 54:23}。灭亡的深坑就是沉沦的黑暗。盲人领盲人，两人都掉进坑里。\BibleRef{玛 15:14} 天主使他们堕入灭亡的深坑，并非因为祂是他们罪恶的肇始者，而是因为祂是他们不义的审判者。\BibleQuote{因为天主任凭他们随从心中的情欲}。\BibleRef{罗 1:24} 因为他们喜爱黑暗，而不喜爱光明；他们喜爱盲目，而不喜爱看见。因为请看，主耶稣已照耀整个世界，让他们与整个世界合一歌唱：\BibleQuote{没有一个人能躲避祂的热力}。但他们从整体移向部分，从身体移向伤口，从生命移向被砍断的肢体，除了堕入灭亡的深坑，还能遇到什么呢？\par

25. \Quote{流人血与诡诈的人}。称他们为流人血的人，是因为他们行杀害；但愿是肉体的杀害，而不是属灵的杀害。因为从肉体流出的血，是可见的，令人战栗；谁看见一个重洗者心中的血呢？那些死亡需要另一样的眼睛。虽然连对这可见的死亡，全副武装的\OriginalTerm{西库姆塞利奥内斯}{Circumcelliones}也各处不得安宁。如果我们想到这些可见的死亡，就有流人血的人。注意那武装的人，看他是否是一个和平的人，而非流人血的人。哪怕他仅携带一根棍棒，也还好；但他带着投石器，带着斧头，带着石头，带着长矛；带着这些武器，无论何处，他们横行，渴求无辜者的血。因此，连就这些可见的死亡而言，也有流人血的人。但即使对他们，我们也说：但愿他们只施行这样的死亡，而不杀害灵魂。这些流人血与诡诈的人，不要以为我们这样误解了流人血的人，认为他们是杀灵魂的；他们自己对他们自己的\OriginalTerm{马克西米安派}{Maximianists}就是这样理解的。因为当他们定罪他们时，在他们公会议的文书中，他们写下了这些话：\Quote{他们的脚迅速倾流（宣讲者的）血}，\Quote{他们的道路上有患难与祸害，和平的道路他们不认识。}这是他们关于马克西米安派所说的。但我问他们：马克西米安派何时曾流过肉体的血？不是因为他们若有一大群也不会流，而是因为他们在少数中的畏惧，反倒从别人那里受了一些苦，从未行过任何这样的事。因此我问那多纳特派人士说：在你们的公会议中，你们关于马克西米安派写下了\Quote{他们的脚迅速倾流人的血。}请指给我一个马克西米安派曾伤过的一根手指！他除了我所说的之外，还能回答什么？那些自合一中分离出来、用引人走入歧途而杀害灵魂的人，是属灵而非属肉体地流人血。你们解释得很好，但在你们的解释中，要承认他们自己的行为。\Quote{流人血与诡诈的人}。诡诈在于欺骗、掩饰、诱惑。那么，那些因祂面容的忿怒而分裂的人自身又如何呢？他们自己就是流人血与诡诈的人。\par

26. 但他论及他们说甚么？\Quote{他们不得终享天年。}甚么是\Quote{他们不得终享天年}？他们不能如他们所想的那样进展：在他们预期的时间内，他们必要灭亡。因为他就是那只鹧鸪，经上论它说：\BibleQuote{他在中年时他们必离弃他，在他末日他必成为愚拙的人。}\BibleRef{耶 17:11}他们进展，但只是一时。因为宗徒说：\BibleQuote{但恶人和骗子必越陷越深，他们自己迷惑，也迷惑别人。}\BibleRef{弟后 3:13}但\BibleQuote{瞎子领瞎子，两人必一同掉进坑里。}\BibleRef{玛 15:14}他们理当掉进\Quote{腐败的深坑。}因此他说甚么？他们必更坏地进展：但不会长久。因为不久前他说：\BibleQuote{但他们不会再进展了：}\BibleRef{弟后 3:9}也就是说，\Quote{不得终享天年。}让宗徒继续说明原因：\Quote{因为他们的疯狂必显露于众人，正如那另外一些人的疯狂一样。}\Quote{但是，主啊，我要仰望祢。}但他们不得终享天年，是理所当然的，因为他们寄望于人。但我从暂时的日子达到了永恒的日子。为甚么？因为主啊，我寄望于祢。\par

\SourceRecord{Translated by J.E. Tweed.；Nicene and Post-Nicene Fathers, First Series；Vol. 8.；Edited by Philip Schaff.；Buffalo, NY: Christian Literature Publishing Co.,；1888.；<http://www.newadvent.org/fathers/1801055.htm>.}

% structure: p0001 source=h1 target=section reason=page-title

\section{咏56释义}

1. 正如当我们将要进入任何房屋时，我们观看标题以了解它是谁的、属于谁，以免我们不适当地闯入不应进入的地方；同样，我们也不应因胆怯而退出我们应当进入的地方：简而言之，我们仿佛读到：这些产业属于某位或某位：因此，在这篇圣咏的门楣上，我们见到铭刻着：\BibleQuote{末后，为那些从圣民中被远放的人民，为达味本人，在标题的铭文上，当\OriginalTerm{异族人}{Allophyli}在\Place{Gath}拘禁了他时。} \BibleRef{撒上21:10} 因此，让我们认识那些在标题铭文上从圣民中被远放的人民。因为这确实属于那位如今你们已知道要从属灵上理解的达味。因为这里所推崇给我们的不是别的，正是那一位，关于祂曾说：\BibleQuote{律法的终结就是基督，使一切相信的人得称为义。} \BibleRef{罗10:4} 因此，当你们听到\BibleQuote{末后}时，要留意基督，免得在路上耽搁而达不到终点……\par

2. 那么，这些在标题铭文上从圣民中被远放的人民是谁呢？让标题本身向我们宣告这人民。因为在主的受难时，当主被钉十字架时，那里写了一个标题：用希伯来文、希腊文和拉丁文刻着：\BibleQuote{犹太人的君王；} \BibleRef{若19:19} 用三种语言，就如由三个证人证实了标题：因为\Quote{凭两三个证人的口，每一句话都得成立。} ……\par

3. 那么，那仍然属于标题本身的部分，即\BibleQuote{\OriginalTerm{异族人}{Allophyli}在\Place{Geth}拘禁了他}，是什么意思？\Place{Geth}是\Emph{\OriginalTerm{异族人}{Allophyli}}，即\Emph{外邦人}的一个城市，也就是远离圣民的人民。所有拒绝基督为王的人都成了外邦人。他们为何成了外邦人呢？因为即使那葡萄树，虽由祂所栽种，当它变酸时，听到了什么？\BibleQuote{你怎么变成了酸葡萄，成为外邦的葡萄树？} \BibleRef{耶2:21} 它未被称作\Quote{我的葡萄树}：因为如果是我的，就是甜的；如果酸了，就不是我的；如果不是我的，就必定是外邦的。\BibleQuote{拘禁了他}，于是，\BibleQuote{\OriginalTerm{异族人}{Allophyli}在\Place{Geth}}。我们确实发现，弟兄们，达味本人——\Person{叶瑟}的儿子、以色列的君王——在被\Person{撒乌耳}追赶时，曾在异邦人中间，在异乡地，并且在那城里与那城的君王在一起，\BibleRef{撒上21:10} 但我们没有读到他在那里被拘禁。因此，我们的达味——主耶稣基督，出自那达味的后裔——不仅他们拘禁了，而且\OriginalTerm{异族人}{Allophyli}在\Place{Geth}仍然拘禁着祂。关于\Place{Geth}，我们说过它是一个城市。但若问这名字的含义，它表示\Quote{榨酒池。}……那么，祂如何在这里被拘禁在\Place{Geth}呢？在榨酒池中被拘禁的是祂的身体，即祂的教会。在榨酒池中是什么意思？就是在受压中。但在榨酒池中，压榨是有果效的。葡萄在葡萄树上不受压榨，看起来完整，但从中流不出什么；它被投入榨酒池，被踩踏、被压榨；看似对葡萄造成了伤害，但这种伤害并非不结果实；相反，若未施加伤害，它将一直不结果实。\par

4. 因此，凡被那些远离圣人们的人所压迫的圣人们，都当留意这篇圣咏，在此看见自己，说出这里所说的话，承受这里所描写的苦楚。……所以，在聆听这篇圣咏的词句时，不要让任何人想到私人的敌意：\Quote{要知道我们的战斗不是对抗血肉，而是对抗首领和掌权者，以及邪恶的属神事物。}\BibleRef{弗 6:12}\newline{}也就是说，对抗魔鬼及其天使；因为即使当我们遭受那些骚扰我们的人时，是他正在煽动，正在点燃，仿佛驱动他的器皿。因此，让我们留意两个敌人：一个我们看得见，一个我们看不见；我们看见人，我们看不见魔鬼；让我们爱人，提防魔鬼；为人祈祷，对抗魔鬼祈祷，并向天主说：\Quote{主，求祢怜悯我，因为人践踏了我。}\BibleRef{咏 55:1}不要因为人践踏了你而害怕：你有酒，你已成了葡萄，好被压榨。\Quote{终日作战，他困扰了我。}——每一个远离圣人们的人。但为何此处不应理解为魔鬼本身呢？难道是因提到\Quote{人}？难道福音错了，因为它说：\Quote{一个仇人做了这事。}\BibleRef{玛 13:28}但藉着某种比喻，他也可以被称为人，却并非人。所以，无论是说这话的人所注视的是他，还是那远离圣人的民族和每个人——魔鬼藉着这类人困扰那些依附圣人、依附圣者、依附君王的天主子民；那君王的名号使他们愤怒，仿佛被击退并远离了——愿他说：\Quote{主，求祢怜悯我，因为人践踏了我。}愿他在此践踏中不昏厥，因他知道他所呼求的是谁，并因祂的榜样而坚强。酒榨中的第一串葡萄是基督。当那串葡萄藉着苦难被压榨时，\BibleRef{依 63:3}流出了那\Quote{令人沉醉的杯，何其美妙！}同样，愿祂的身体注视其元首说：\Quote{主，求祢怜悯我，因为人践踏了我；终日作战，他困扰了我。}\Quote{终日}——在一切时刻。不要让任何人对自己说：在我们父辈的时代有苦难，在我们时代却没有。如果你自以为没有苦难，那你还没有开始成为基督徒。宗徒的话在哪里？\Quote{凡在基督内虔敬生活的，都必受迫害。}\BibleRef{弟后 3:12}因此，如果你没有为基督受任何迫害，你要当心，恐怕你尚未开始在基督内虔敬地生活。但当你开始在基督内虔敬地生活时，你已进入了酒榨；准备好受压榨吧，但不要干枯，否则压榨时什么也流不出来。\par

5. \Quote{我的敌人终日践踏我。}\BibleRef{咏 55:2}那些远离圣人们的人，就是我的敌人。终日：已经说过\Quote{从白日的巅峰。}\Quote{从白日的巅峰}是什么意思？或许是一件高深的事。这并不奇怪，因为是白日的巅峰。或许他们因此远离圣人们，因为他们不能渗透白日的巅峰，那里的宗徒是十二个光辉的时辰。因此，那些如同对待人一样钉死祂的人，在白日犯了错误。但他们为何遭受黑暗，以致远离圣人们？因为白日正在高处照耀，他们不认识那隐藏在至高处的祂。\Quote{因为如果他们认识，就绝不会钉死光荣的主。}\BibleRef{格前 2:8}……\par

6. \Quote{因为许多与我作战的人要害怕。}\BibleRef{咏 55:3}何时害怕？当白日过去，他们高高在上之时。因为一时他们高高在上，当他们的高傲之时结束时，他们就会害怕。\Quote{但我，主，唯有在祢内寄望。}他不说：\Quote{但我不害怕。}却说：\Quote{许多与我作战的人要害怕。}当那审判之日来临时，\Quote{地上所有支派都要哀悼。}\BibleRef{玛 24:30}当人子的记号在天上出现时，那时所有圣人都将安然无恙。因为那他们所盼望、所渴望、所祈求到来之事，就要来到；但那些人不再有悔改的机会，因为在可结出果实的悔改之时，他们刚硬了心，不听天主的警告。他们也要在审判的天主面前筑起一道墙吗？你确实当承认这人的虔诚，并且如果你属于那身体，就效法他。当他说：\Quote{许多与我作战的人要害怕}之后，他没有接着说：\Quote{但我不害怕。}免得他将不害怕归于自己的能力，从而也处于暂时的高傲之中，因着暂时的骄傲而不能得到永远的安息；相反，他使你明白他为何不害怕。\Quote{但我，}他说，\Quote{主，唯有在祢内寄望。}他没有谈论自己的自信，而是谈论自信的原因。因为如果我不害怕，我也可能因心硬而不害怕，因为许多人因过于骄傲而什么都不怕。……\par

7. \BibleQuote{在天主内，我要赞美我的言论，在天主内我怀有希望：我不害怕血肉对我所做的事} \BibleRef{咏 55:4}。为何？因为在天主内，我要赞美我的\OriginalTerm{言论}{discourses}。若你在自己内赞美你的言论：我并非说你不必害怕；你不可能不害怕。因为你的言论，要么是假的，因此是你自己的，因为是假的；要么你的言论是真的，而你以为不是从天主领受，而是凭自己说的；它们将是真的，但你将是假的；但如果你知道，你凭天主的智慧、真理的信德，不能说任何真的，除非是从祂领受的，关于祂曾说：\BibleQuote{你所有的一切，哪一样不是领受的呢？} \BibleRef{格前 4:7} 那么，你在天主内赞美你的言论，为要你在天主内，因天主的言论而得赞美……\BibleQuote{在天主内我怀有希望，我不害怕血肉对我作什么。} 你岂不是不久前还在说：\BibleQuote{主啊，怜悯我，因为人已将我践踏；终日作战，他烦扰了我}？那么这里怎么又说：\BibleQuote{我不害怕血肉对我作什么}？他能对你作什么？你自己不久前说过：\Quote{他已将我践踏，他烦扰了我}。当他就作这些时，他什么也不作？他注视那从践踏流出的酒，并回答说：他确实践踏了，确实烦扰了；但对我能作什么？我本是葡萄，我将成为酒：\BibleQuote{在天主内我怀有希望，我不害怕血肉对我作什么。}\par

8. \BibleQuote{终日，他们憎恶我的言语} \BibleRef{咏 55:5}。你知道，他们就是这样。讲真理，宣讲真理，向外邦人宣讲基督，向异端者宣讲教会，向众人宣讲救恩：他们反驳，他们憎恶我的言语。但当他们憎恶我的言语时，你们以为他们憎恶谁，岂不是憎恶那一位，就是我在祂内要赞美我的言论的那一位？\BibleQuote{终日，他们憎恶我的言语。} 至少让这一点足够：让他们憎恶言语，不要再进一步的谴责、拒绝！愿他们远离！我为何说这话？当他们拒绝言语、当他们仇恨言语，那些从真理之源流出的言语时，他们会对传讲这些言语的人做什么？除了接下来的：\BibleQuote{他们的一切计谋都是为害我？} 如果他们仇恨面包本身，他们又怎能放过那盛面包的篮子呢？\BibleQuote{他们的一切计谋都是为害我。} 若是这样，甚至对主自身也是如此，那么，身体不要嫌弃那在头上已先行的事，好使身体能依附于头。你的主曾被轻视，你难道要自己受那些离圣人很远的人尊崇吗？你不要为自己要求那在祂身上未先行的事。\BibleQuote{门徒不会大过他的师傅；仆人不会大过他的主人。若他们称家主为贝耳则步，何况他的家人呢？} \BibleRef{玛 10:24} 他们的一切计谋都是为害我。\par

9. \Quote{他们要客居，且要隐藏。} \BibleRef{咏 55:6} 客居就是身在异乡。客居一词用来指那些生活在不属于自己国家的人。此世人人都为旅客：在此生，我们看到自己为肉体所覆盖，透过这肉体，人心无法被看见。因此宗徒说：\Quote{不等到时候，你们什么也不要判断，只等到主来，他要照亮暗中的隐情，并显明人心的意念；那时，各人才从天主那里获得称赞。} \BibleRef{格前 4:5} 在这事发生之前，在此世肉体生命的客居中，每人背负自己的心，每颗心对另一颗心都是关闭的。此外，那些计谋要加害这人的恶人，\Quote{要客居，并要隐藏：} 因为他们身在异乡，背着肉体，在心中隐藏诡计；无论他们想什么恶事，都隐藏起来。为什么呢？因为这生命仍是客旅。让他们隐藏吧；他们所隐藏的将要显露，他们自己也必不被隐藏。关于这隐藏之事，还有另一种解释，或许更可取。因为那些远离圣徒的人中，会混进一些假弟兄，他们给基督的身体带来更严重的磨难；因为他们不像完全的外人那样被避开……但就是那些人，弟兄们，我们也不要怕：\Quote{我必不怕肉体对我做什么。} 即使他们客居，即使他们混入，即使他们伪装，即使他们隐藏，他们只是肉体；你只管在主内盼望，肉体不能对你做什么。但他带来磨难，带来践踏。有酒添上，因为葡萄被压榨；你的磨难不会没有果子：别人看到你，效法你。因为你为了学习忍受这样的人，已仰望你的元首，那第一串葡萄，曾有一个人到祂那里来看祂，客居并隐藏，就是叛徒犹达斯。因此，所有带着假心混入的人，客居并隐藏的，你不要怕；这些人的父亲犹达斯曾与你的主同在：主确实知道他；虽然叛徒犹达斯客居并隐藏，但他的心向万有的主是敞开的。主明知故选了一个人，为安慰你们，使你们不知道应当避开谁。因为祂本可以不选犹达斯，因祂知道犹达斯：祂对门徒说：\Quote{我不是拣选了你们十二个吗？你们中却有一个是魔鬼。} \BibleRef{若 6:70} 所以连魔鬼也被选了。或者他若没有被选，那么祂如何选了十二个，而不是十一个呢？他确被选了，但为另一目的。十一个被选为考验的工作，一个被选为试探的工作。祂如何能给你一个榜样——你不知道应当避免哪些恶人，应当提防哪些假意做作、客居隐藏的人——除非祂对你说：看，我自己就曾有过那样一个人！先例已行，我忍受了我所知道的事，为给不知道的你们安慰。他向我所做的，也必向你们做；为使他能多行恶事，造成大破坏，他会控告，会捏造诬告……\par

10. \Quote{这些人要窥伺我的脚跟。} 因为他们要如此客居并隐藏，好窥看人在何处失足。他们专注于脚跟，要看何时可能失足；为要绊倒脚，或绊跌脚；当然，是为找到可指控的事。有谁行走而不失足呢？例如，在舌头上失足是多么迅速？正如经上所记：\Quote{谁在言语上不跌倒，就是个完人。} \BibleRef{雅 3:2} 请问，有谁胆敢称自己为完人或自以为完人呢？因此，人人都难免在言语上失足。但那些将要客居并隐藏的人，会挑剔一切言语，寻找机会设下网罗和难解的假控告，他们自己反被缠住，比他们要缠住的人更甚；好叫他们自己先被捉住而灭亡，而不是他们捉住别人来毁灭……凡我说过的善言、真言，都是出于天主，从天主说的；凡我不该说的其他话，我是作为人说的，但在天主之下说的。那扶持行走者的，警告迷途者的，宽恕认罪者的，召回舌头，召回失足者……你要留意你所责备之人的言论，或许他教你一些有益于你健康的事。他说：那在言语上失足的人，能教你什么有益健康的事呢？或许他正教你健康的事：叫你不要作挑剔言语的人，而要作收集训诫的人。\Quote{我的灵魂曾忍受了。} 我说我所忍受的事。他像一个有经验的人说话：\Quote{我的灵魂曾忍受了。他们要客居并隐藏。} 让我的灵魂忍受所有人，从外面来的不吠叫之人，在里面隐藏之人，让它忍受。从外面来的，如河水的试探：让你站在磐石上，它撞上来，却不把你冲倒；房子建在磐石上。\BibleRef{玛 7:25} 他在里面，他要客居并隐藏：假设糠秕靠近你，让牛群的践踏来到，让试探的碾轮来到；你被洁净，那其他东西被压碎。\par

11. \Quote{因为祢一无所有，祢必拯救他们} \BibleRef{咏 55:7} 。祂甚至教导我们为这些人祈祷。然而，\Quote{他们必寄居并隐藏}，无论他们如何诡诈，如何伪装，如何埋伏伺机；你要为他们祈祷，不要说：天主岂能改变这样的人，如此邪恶，如此悖逆？不要失望：留神你所求的那一位，而非你所求的对象。你看见疾病的严重，难道看不见医师的大能？\Quote{他们必寄居并隐藏：正如我的灵魂所经历。} 经历，祈祷：然后成就什么？\Quote{因为祢一无所有，祢必拯救他们。} 祢必拯救他们，使之对祢而言一无所是，即对祢毫无劳苦。在人看来他们已无望，但祢一语便可治愈；祢在医治中不会劳苦，尽管我们观看时惊愕不已。在这节经文中还有另一层含义：\Quote{因为祢一无所有，祢必拯救他们：} 没有他们任何先行的功绩，祢就拯救他们……他们不会给祢带来公山羊、公绵羊、公牛，不会在祢的圣殿中向祢献上礼物和香料，不会倾注任何良心准备的奠祭；他们身上一切都粗糙、污秽、可憎：尽管他们未给祢带来任何可使他们得救之物；\Quote{因为祢一无所有，祢必拯救他们，}即凭祢恩宠的白白赠予……\par

12. \Quote{在怒气中，祢将万民降下。} 祢发怒并降下，祢震怒并拯救，祢恐吓并召唤。祢用患难充满万物，好使人在患难中投奔祢，免得被享乐和虚假的安全所诱惑。从祢而来的怒气可见，但那是父亲般的怒气。父亲对儿子发怒，儿子藐视他的命令；父亲发怒，就责打他，击打他，拉耳朵，用手拖拽，领去学校。有多少人进入了，有多少人充满了上主的圣殿，是在祂的怒气中被降下，即被患难惊吓并满怀信德？因为患难正为此而激动，为要倒空那充满邪恶的器皿，好使它被恩宠充满。\par

13. \Quote{上主，我向祢倾吐了我的生命} \BibleRef{咏 55:8} 。因为我能活着是出于祢的作为，为此我向祢倾吐我的生命。但天主岂不知道祂所给予的？你所倾吐给祂的是什么？你岂要教导天主？断乎不可。那么，为何他说：\Quote{我向祢倾吐了我的生命}？莫非因为我倾吐我的生命对祢有益？这对天主有什么益处？是对天主的益处有益。我已向天主倾吐了我的生命，因为那生命是天主的作为。正如宗徒保禄倾吐他的生命，说：\Quote{我从前是亵渎者、迫害者和施暴者，}他要倾吐他的生命。\Quote{但是我蒙受了怜悯。} \BibleRef{弟前 1:13} 他倾吐他的生命，不是为自己，而是为祂：因为他如此倾吐，好使人在祂内相信，不是为了自己的益处，而是为了祂的益处……\Quote{上主，我向祢倾吐了我的生命。祢将我的眼泪放在祢面前。} 祢垂听了我的哀求。\Quote{正如祢的应许。} 因为祢曾许诺此事，便如此成就。祢说了要垂听哭泣者。我信了，我哭了，我被垂听了；我发现在应许中祢仁慈，在报偿中祢信实。\par

14. \BibleQuote{愿我的仇敌反身退后} \BibleRef{咏 55:9}。这件事对这些人大有益处，他并不是在诅咒他们。因为他们愿意走在前头，所以不愿意改正。你劝诫你的仇人生活美好，好使他改正自己；他却嗤之以鼻，拒绝你的话：\Quote{看那劝诫我的人，看那我要听他命令而生活的人！} 他愿意走在你前头，但走在前面却不能改正。他不注意你的话不是出于你自己，他也不注意你向天主陈明你的生活，而不是向你自己。因此，走在前头不能改正：使他反身退后，跟随他原先愿意走在前头的那位，这对他是一件好事。主向祂的门徒谈论即将来临的受难。伯多禄惊骇，说道：\BibleQuote{主，千万不可！} \BibleRef{玛 16:22} 他刚才还说：\BibleQuote{祢是默西亚，永生天主之子，} 既已宣认了天主，却怕祂死去，仿佛祂只不过是人。但主如此降临，正是为了受苦（因为我们除非靠祂的血被救赎，否则不能得救），祂刚才还赞扬了伯多禄的宣认……但主一开始谈论祂的受难，伯多禄就害怕，唯恐祂死于死亡，而实际上我们自己也本该灭亡，除非祂死去；于是他说：\BibleQuote{主，千万不可，这事绝不会临到祢！} 主对刚才还对他说：\BibleQuote{你是伯多禄，在这磐石上，我要建立我的教会} 的那位说：\BibleQuote{退到我后面去，撒殚！你是我的绊脚石。} 为什么刚刚还是\Quote{有福的}、\Quote{磐石}的那位，现在成了\Quote{撒殚}？因为祂说：\BibleQuote{你所体会的，不是天主的事，而是人的事。} \BibleRef{玛 16:23} 刚才他还\Emph{体会}天主的事，因为\BibleQuote{不是血肉启示了你，而是我在天之父。} 当他在天主内赞扬祂的话语时，他不是撒殚，而是伯多禄，来自\OriginalTerm{磐石}{petra}；但当出于他自己和人的软弱时，出于肉性的人爱，这爱会成为他自己和其余人得救的障碍，他就被称为撒殚。为什么？因为他愿意走在主前头，向天上的领袖提出世俗的劝告。\BibleQuote{主，千万不可，这事绝不会临到祢！} 你说\Quote{千万不可}，又说\Quote{主}：当然，如果祂是主，祂行事带着权能；如果祂是师傅，祂知道祂做什么，知道祂教什么。但你却想领导你的领袖，教导你的师傅，命令你的主，为天主选择：你走得太前了，退到后面去。这岂不也对这些仇敌有益？\BibleQuote{愿我的仇敌反身退后；} 但不要让他们停留在后面。为此，让他们反身退后，以免他们走在前头；而是要他们跟随，不是停留在原地。\par

15. \BibleQuote{在我呼求祢的每一日，看，我已知道祢是我的天主。} \BibleRef{咏 55:9} 这是一个伟大的知识。祂不是说：\Quote{我已知道祢是天主，} 而是：\Quote{祢是\Emph{我的}天主。} 因为祂是你的，当祂援助你时；祂是你的，当你并非祂的外人时。因此经上说：\BibleQuote{那主是天主的人民，是有福的。} 为什么说是\Quote{有福的人民}？因为祂不是谁的天主呢？祂确实是万有的天主；但祂特别被称为那些爱祂、持守祂、拥有祂、朝拜祂的人的天主，如同属于祂自己的家：他们就是祂的大家族，被独生子的宝血所赎买。天主赐给了我们何等伟大的恩典，使我们成为祂的，祂成为我们的！但事实上，远方的人已被从圣者中驱逐出去，他们是外邦之子。看另一篇圣咏如何论及他们：\BibleQuote{主，求祢拯救我，} 他说，\BibleQuote{脱离外邦之子的手，他们的口出虚言，他们的右手是行欺诈的右手。} ……\par

16. 所以，诸位弟兄，让我们纯洁而贞洁地爱天主。如果一个人为了报酬而敬拜天主，他的心就不是贞洁的。为何这样说？难道我们敬拜天主就没有报酬吗？显然我们有，但我们所敬拜的天主本身就是报酬。祂自己将是我们的报酬，因为\BibleQuote{我们必要看见祂实在怎样}\BibleRef{若一 3:2}。请注意，你将获得一份报酬……我要告诉你们，诸位弟兄：在这些人间的关系中，想一想一颗对天主贞洁的心是怎样的：当然，人间的关系是这样的：一个人如果因为妻子的嫁奁而爱她，他就不是真正爱妻子；一个女人如果因为丈夫给了她什么或给了很多而爱他，她就不是贞洁地爱丈夫。无论丈夫是富人还是穷人，他仍是丈夫。有多少被流放的男人，反而被贞洁的妻子更加深爱！许多贞洁的婚姻在丈夫的不幸中得到了证明：妻子们为了不让人以为她们爱任何其他东西胜过爱丈夫，不仅没有离弃，反而更加顺服。因此，如果肉体的丈夫是被自由地爱着，是被贞洁地爱着；肉体的妻子是被自由地爱着，是被贞洁地爱着；那么，天主——灵魂的真实而说实话的净配，使我们结出永生的果实，不让我们荒芜——又该如何被爱呢？因此，让我们这样爱祂，除了祂自己之外，不爱任何其他事物：这样，在我们身上就会发生我们所说过的、我们所歌颂的，因为在这里我们的声音也是：\Quote{在我呼号祢的那一天，看，我已知道祢是我的天主。}这就是呼求天主，自由地呼求祂。此外，关于某些人，经上说了什么？\Quote{他们没有呼号上主。}他们似乎是在呼求上主，但他们是向上主祈求关于遗产、增加财富、延长寿命以及其余暂时的事物：关于他们，圣经说了什么？\Quote{他们没有呼号上主。}因此，接下去是什么？\Quote{他们在没有恐惧的地方战栗恐惧。}什么是\Quote{没有恐惧的地方}？怕钱财被偷，怕家中有什么减少，最后，怕此生寿命比他们所期望的短少：但他们在没有恐惧的地方却颤抖恐惧。……\BibleQuote{我靠天主颂扬祂的许诺，我靠上主赞美祂的言语。}\BibleRef{咏 55:10}：\BibleQuote{我倚靠天主，必不惧怕，人能对我做什么？}\BibleRef{咏 55:11}。这正是前面所重复的同一个意思。\par

17. \BibleQuote{天主，在我身内有祢的誓愿，我要向祢献上赞颂的祭献。}\BibleRef{咏 55:12}\Quote{你们要向上主你们的天主许愿，并予以偿还。}什么誓愿？什么祭献？难道是以前在祭台上奉献的那些牲畜吗？你不要奉献那些东西：在你自身内就有你可以许愿和偿还的。从心灵的宝库中取出赞颂的馨香，从良心的库房内取出信德的祭献。无论你取出什么，都要用爱点燃。在你自身内要有誓愿，你可以向天主献上赞颂。是什么样的赞颂？祂赐给你什么？\BibleQuote{因为祢救了我的灵魂脱离死亡。}\BibleRef{咏 55:13}这正是他向祂陈诉的生命：\Quote{天主，我向祢陈述了我的生命。}因为我是什么？死亡。凭我自己，我是死的；藉着祢，我是什么？活的。因此，\Quote{天主，在我身内有祢的誓愿，我要向祢献上赞颂的祭献。}看，我爱我的天主：没有人能从我这里夺走祂；我能给祂的，没有人能从我这里夺走，因为它封闭在心灵内。因此，先前那充满信心的\Quote{人能对我做什么？}说得有理。任凭人发怒，任凭他们发怒，任凭他们成就他们所尝试的：他能夺走什么？黄金、白银、牲畜、男仆、婢女、田产、房屋，让他夺走一切吧：他难道能夺走在我身内的、我要向天主献上赞颂的誓愿吗？诱惑者被允许试探圣人约伯\BibleRef{约 1:12}，一瞬间他夺走了所有的一切：他享有的一切财物，他都夺去了；夺走了遗产，杀害了继承人；这并非一点点，而是一下子、一举，所有的都同时报来：当一切都被夺走时，只剩下约伯一人，但在他的身内却有赞颂的誓愿，他可以献给天主，显然在他身内就有：他圣洁胸怀的宝库，那偷窃的魔鬼没有搜刮，那里装满了可以献祭的资粮。听听他有什么，听听他拿出了什么：\BibleQuote{上主赐予，上主收回；上主所喜悦的，就这样成就了：愿上主的名受赞美。}\BibleRef{约 1:21}哦，内在的财富，小偷不能靠近！天主自己赐下了祂所接收的；祂自己用那些祂曾赐予的使他富有，使他可以献上祂所爱的。天主需要你的赞颂，天主需要你的宣认。但是，你能从你的田地里给出什么吗？祂自己降雨使你有所收获。你能从你的宝库里给出什么吗？祂自己放入你要给出的一切。你有什么不是领受的呢？\BibleQuote{你有什么不是领受的呢？}\BibleRef{格前 4:7}从心里你将要给出吗？祂也赐给了信德、望德、爱德：你必须取出这些，你必须献祭这些。但是，显然其他的东西，敌人能够违背你的意愿夺走；而这些，除非你愿意，他不能夺走。一个人甚至会违背自己的意愿失去这些：想要金子，会失去金子；想要房屋，会失去房屋；信德，没有人会失去，除非他轻视了她。\par

18. \Quote{因为祢救了我的灵魂脱离死亡，我的眼睛脱离眼泪，我的脚脱离滑跌，使我行走在天主面前，在活人的光明中} \BibleRef{咏 55:13}。外邦人不是天主所喜悦的人，他们被远离圣人，因为他们没有活人的光明，看不见什么令天主喜悦。\Quote{活人的光明}，就是不灭之光，圣人之光。不在黑暗中的人，在活人的光明中蒙悦纳。人被人观察，人所拥有的也被观察；但无人知道他是何种人：天主知道他是何种人。有时连魔鬼自己也逃脱了；除非他诱惑，否则他不会发现：正如我刚才提过的那个人……\Quote{约伯敬畏天主，岂是无缘无故的呢？} \BibleRef{约 1:9} 因为这就是真光，活人的光明，即\Emph{白白地}敬拜天主。天主在自己仆人的心中看到了他白白的崇拜。那颗心在主面前、在活人的光明中是蒙悦纳的：魔鬼的眼目逃脱了，因为他在黑暗中。天主容许那诱惑者进来，并非为使自己知道祂已知道的事，而是为使我们知晓并效法。诱惑者被容许进入；他夺走了一切，那人留下了，一无所有，失去了财产，失去了家人，失去了儿女，却充满了天主。的确留下了一个妻子。\BibleRef{约 2:9} 你们是否认为魔鬼仁慈，竟给他留下了一个妻子？他知道自己是通过谁欺骗了亚当……从头到脚受伤，但里面却完整无缺，他在这活人的光明中，在他内心的光明中，对那诱惑他的女人说：\Quote{你说话像那些愚昧的妇人一样} \BibleRef{约 2:10}，也就是说，像一个没有活人的光明的人。因为活人的光明就是智慧，而愚昧人的黑暗就是愚蠢。你说话像那些愚昧的妇人一样：你看见我的肉体，却看不见我内心的光明。因为那时她或许会更爱自己的丈夫，倘若她认出了他内在的美，看见他在天主眼中美丽的地方：因为他内心怀着可以向天主献上赞美的誓愿。仇敌何等完全地不敢侵犯那产业！他所拥有的以及他所盼望的、那在\Quote{德上加德}中前进的一切，多么完整！因此，弟兄们，愿这一切都为我们服务，使我们\Emph{白白地}爱天主，永远盼望祂，不怕任何人或魔鬼。无论那一个，除非被许可，否则什么也不能做；许可的原因只能是它对我们有益。让我们忍受恶人，让我们成为善人：因为我们自己也曾经是恶人。天主将拯救那些我们甚至不敢绝望的人，像什么都没有一样。所以，不要对任何人绝望，为我们所忍受的众人祈祷，永不离开天主。愿祂是我们的产业，我们的希望，我们的救恩。祂在这里是安慰者，在那里是酬报者，无处不是赐予生命者，生命的赋予者，不是另一生命，而是那曾说\BibleQuote{我就是道路、真理、生命} \BibleRef{若 14:6}的生命：好叫我们既在此世信德的光明中，也在来世眼见的光明中，如同在活人的光明里，在主面前蒙悦纳。\par

\SourceRecord{Translated by J.E. Tweed.；Nicene and Post-Nicene Fathers, First Series；Vol. 8.；Edited by Philip Schaff.；Buffalo, NY: Christian Literature Publishing Co.,；1888.；<http://www.newadvent.org/fathers/1801056.htm>.}

% structure: p0001 source=h1 target=section reason=page-title

\section{\WorkTitle{圣咏第五十七篇讲疏}}

诸位弟兄，我们刚才在福音中听见，我们的主、救主\Person{耶稣基督}多么爱我们：祂与圣父同为天主，与我们同为世人，出自我们自身，如今坐在圣父的右边；你们已听见祂多么爱我们……\par

因为这圣咏歌唱的是主的受难，请看它的标题是什么：\Quote{直到终末。}终末就是基督。\BibleRef{罗 10:4}为什么祂被称为终末？不是作为终结者，而是作为完成者。\par

\Quote{直到终末，不可败坏，属于达味自己，为标题的铭文；当他躲避撒乌耳的面，逃进一个山洞时。}我们查阅圣经，确实发现圣\Person{达味}——以色列君王，达味圣咏集也从他得名——曾遭受本国君王\Person{撒乌耳}的迫害，正如你们许多读过或听过圣经的人所知。\Person{达味}王当时以\Person{撒乌耳}为迫害者：一方极其温良，另一方极其凶残；一方柔和，另一方嫉妒；一方忍耐，另一方残忍；一方仁慈，另一方忘恩负义：\Person{达味}以如此的柔和忍受他，以至于当\Person{撒乌耳}落在他手中时，他没有碰他，也没有伤害他……这与\Person{基督}有何关系？如果当时所行的一切都是未来之事的预像，那么我们在其中找到\Person{基督}，而且远超过其他。因为这句\Quote{不可败坏，为标题的铭文}，我不知如何属于那位\Person{达味}。因为\Person{达味}自己没有任何\Quote{标题}是\Person{撒乌耳}要\Quote{败坏}的。但在主的受难中，我们看到写有一个标题：\Quote{犹太人的君王}，为使这标题羞辱那些人，因为他们对自己的君王仍不住手。因为在他们里面是\Person{撒乌耳}，在\Person{基督}里面是\Person{达味}。因为\Person{基督}，如宗徒福音所说，按肉身是\Person{达味}的后裔，正如我们知道并宣认的；按天主性，祂超越\Person{达味}，超越众人，超越天地，超越天使，超越一切有形无形之物……而且因为已通过圣神唱出：\Quote{直到终末，不可败坏，为标题的铭文。}\Person{比拉多}回答他们说：\BibleQuote{我所写的，我已经写了。}\BibleRef{若 19:22}为什么你们向我暗示虚假？我不败坏真理。\par

那么，\Quote{当他躲避撒乌耳的面，逃进一个山洞时}是什么意思？这件事先前的\Person{达味}也确实做了；但因为在\Person{达味}身上我们找不到标题的铭文，就让我们在\Person{基督}身上找到逃进山洞。\BibleRef{撒上 24:3}因为\Person{达味}藏身的那个山洞预表了某物。但他为什么藏身？是为了隐藏起来不被找到。藏在山洞里是什么意思？就是藏在地下。因为逃进山洞的人被泥土覆盖，使人看不见。但耶稣带着泥土，即祂从泥土领受的肉身，并把自己藏在其中，免得被犹太人认出是天主。\BibleQuote{因为如果他们知道了，决不会把光荣的主钉在十字架上。}\BibleRef{格前 2:8}所以为什么他们找不到光荣的主？因为祂把自己藏在一个山洞里，即向他们的眼目呈现肉身的软弱，而将天主性的威严藏在身体的遮盖下，如同藏在泥土的隐密处……但祂为什么愿意忍耐直到死亡？是为了躲避\Person{撒乌耳}的面而逃进一个山洞。因为山洞可理解为地下深处。而且，正如所有人明显和确信的，祂的身体被安放在一个岩石中凿成的坟墓里。这坟墓就是山洞；祂在那里躲避\Person{撒乌耳}的面。因为犹太人一直迫害祂，直到祂被安放在山洞里。我们凭什么证明他们这样迫害祂直到祂被安放？甚至当祂死在十字架上时，他们还用长矛刺祂。\BibleRef{若 19:34}但当祂被裹尸布包裹、举行葬礼后安放在山洞里时，他们再也不能对肉身做什么了。主因此从那个山洞中复活，安然无恙，没有朽坏，从祂躲避\Person{撒乌耳}之面的地方出来：向不虔敬的人隐藏自己，\Person{撒乌耳}预表了他们，却向祂的肢体显现。因为祂复活的肢体被祂的肢体触摸：祂的肢体，即宗徒们，触摸复活的主并相信了；\BibleRef{路 24:39}看，\Person{撒乌耳}的迫害毫无益处。现在我们听这圣咏吧；因为关于它的标题，我们已按主所赐的恩宠说了足够的话。\par

5. \BibleQuote{天主，求祢怜悯我，求祢怜悯我，因为我的灵魂曾信赖祢。} \BibleRef{咏 57:2} 基督在受难中说：\BibleQuote{天主，求祢怜悯我。} 天主对天主说：\BibleQuote{求祢怜悯我！} 那与圣父一同怜悯你的那一位，在你内呼喊：\BibleQuote{求祢怜悯我。} 因为祂那呼喊\BibleQuote{求祢怜悯我}的部分是你的：祂从你那里领受了这，为了你，好使你获得解救，祂穿上了血肉。血肉本身呼喊：\BibleQuote{天主，求祢怜悯我，求祢怜悯我：} 人本身，灵魂和血肉。因为圣言取了完整的人，圣言成为完整的人。因此，不要以为那里没有灵魂，因为福音作者这样说：\BibleQuote{圣言成了血肉，居住在我们中间。} \BibleRef{若 1:14} 因为人被称为血肉，正如圣经在另一处说：\BibleQuote{凡有血肉的，都要看见天主的救恩。} 难道血肉独自看见，而灵魂不在那里吗？…你听见师傅祈祷，学习你也祈祷。因为祂祈祷正是为此，为教导如何祈祷；因为祂受苦正是为此，为教导如何受苦；祂复活正是为此，为教导如何盼望复活。\BibleQuote{在祢翅膀的荫庇下，我必盼望，直到罪愆过去。} 这现在显然整全的基督在说：这里也有我们的声音。因为罪愆尚未过去，仍然猖獗。我们主自己最后说，罪愆将要增多：\BibleQuote{且由于罪愆增多，许多人的爱德必要冷却；但坚持到底的，这人才可得救。} \BibleRef{玛 24:12} 但谁能坚持到底，直到罪愆过去呢？那就是在基督身体内的人，在基督肢体中的人，并从元首学会了坚持的忍耐。你过去，看，你的诱惑过去了；你进入另一生命，圣人们已进入那里，如果你曾是圣人的话。殉道者们已进入另一生命；如果你曾是殉道者，你也进入另一生命。因为\Quote{你}从这里过去，但罪愆因此过去了吗？生出了其他不义的人，正如一些不义的人死去。同样，正如一些不义的人死去，另一些生出；一些义人离去，另一些生出。直到世界终结，罪愆不会缺少以压迫，正义也不会缺少以受苦…\par

6. \BibleQuote{我要呼求至高者天主} \BibleRef{咏 57:3} 。如果祂是至高者，祂如何听见你呼求？经验产生了信心：\BibleQuote{向天主，}他说，\BibleQuote{祂曾厚待了我。} 如果在我寻求祂之前，祂已厚待了我，当我呼求时，祂岂不垂听？因为主天主厚待了我们，给我们派遣了我们的救主耶稣基督，使祂为我们的过犯而死，并为我们的成义而复活。 \BibleRef{罗 4:25} 祂愿意祂的圣子为什么样的人死？为不敬神的人。但不敬神的人并不寻求天主，而是被天主所寻求。因为祂是至高者，然而我们的痛苦和呻吟并不远离祂：因为\BibleQuote{上主亲近心灵破碎的人。} \BibleQuote{天主曾厚待了我。}\par

7. \Quote{祂从天上派遣，拯救了我} \BibleRef{咏 56:3}。如今，那真人本身，那肉身本身，那天主子在取了我们的本性之后，显而易见祂如何得救：是父从天上派遣并拯救了祂，从天上派遣并使祂复活。但为使你们知道，主自己也使自己复活——两种说法都记载于圣经：既是父使祂复活，也是祂自己使自己复活。请听父如何使祂复活：宗徒说：\Quote{祂成了顺命至死，且死在十字架上。为此，天主极其举扬祂，赐给祂一个名号，超越一切名号。} \BibleRef{斐 2:8} 你们已听到父使子复活并举扬祂；再听祂自己如何使自己的肉身复活。祂以圣殿为比喻对犹太人说：\Quote{你们拆毁这座圣殿，三天之内我要把它重建起来。} \BibleRef{若 2:19} 而福音作者向我们解释了祂所说的话：\Quote{祂这话，是指祂身体的圣殿。} 因此，现在出于祈祷者的人位、出于人的人位、出于肉身的人位，祂说：\Quote{祂拯救了我。祂使践踏我的人蒙受羞辱。} 那些践踏祂的人，在祂死后侮辱祂的人，把祂当作人钉死的人（因为他们没有认出天主），祂使他们蒙受羞辱。请看这件事是否已经成就。我们并非相信此事尚未到来，而是承认它已经应验。犹太人对基督发怒，傲慢地对待基督。在哪里？在耶路撒冷城中。因为在他们统治的地方，他们骄傲自大，他们昂首挺胸。在主受难之后，他们从那里被连根拔除；他们失去了那国度，因为他们不愿承认基督为王。请看他们如何蒙受羞辱：他们分散到万国，无处定居，无处安身。但正因如此，他们仍是犹太人，为的是让他们携带我们的书卷，令自己蒙羞。每当我们想证明基督曾被预言，我们就向异教徒出示这些著作。恐怕一些难以相信的人说，这些书是我们基督徒编造的，以便连同我们所传的福音一起伪造了先知书，使我们所传的似乎早有预言；我们以此来说服他们：所有预言基督的著作都在犹太人那里，所有这一切著作犹太人都持有。我们从敌人那里拿出文献，去驳倒其他敌人。因此，犹太人是何等蒙羞！犹太人携带文献，基督徒借此得以相信。他们成了我们的图书管理员，正如奴隶通常跟在主人后面携带文书，一边因搬运而疲乏，一边因阅读而受益。犹太人就蒙受了这样的羞辱；那早已预言的话应验了：\Quote{祂使践踏我的人蒙受羞辱。} 弟兄们，这是多么大的羞辱：他们读到这节经文，自己却是瞎子，照镜子也看不见自己！因为犹太人在他们所携带的圣经中显现的情形，如同盲人的脸在镜中：别人看见，自己却看不见。\par

8. 当祂说\Quote{祂从天上派遣，拯救了我}时，你或许会问：祂从天上派遣了什么？祂从天上派遣了谁？难道派遣了一位天使去拯救基督，主子竟通过仆人得救？因为所有天使都是服侍基督的受造物。天使可能因服从而被派遣，因服务而被派遣，却不是为援助；正如经上所记：\Quote{天使服侍祂，} \BibleRef{玛 4:11} 并非像人怜悯一个贫困者，而是如同臣仆服侍全能者。那么，\Quote{祂从天上派遣，拯救了我}是什么意思？现在我们在另一节经文中听到祂从天上派遣了什么。\Quote{祂从天上派遣了祂的仁慈和祂的真理。} 为了什么目的？\Quote{并从幼狮群中拉出了我的灵魂。} \Quote{祂派遣了，} 祂说，\Quote{从天上派遣了祂的仁慈和祂的真理。} 而基督自己说：\Quote{我就是真理。} 因此，真理被派遣，为要从幼狮群中拉出我的灵魂；仁慈也被派遣。我们发现基督本身就是仁慈和真理：仁慈在于与我们一同受苦，真理在于报答我们……谁是幼狮？那是较小的百姓，被犹太首领欺骗、引入歧途，以致那些首领是狮子，这些百姓是幼狮。他们都咆哮，都杀害。因为我们也要在此处听到这些人的被杀，就在这圣咏的随后几节中。\par

9. \Quote{并从幼狮群中，} 祂说，\Quote{拉出了我的灵魂} \BibleRef{咏 56:4}。你为什么说\Quote{并拉出了我的灵魂}？你受了什么苦，以致你的灵魂需要被拉出？\Quote{我睡着时心烦意乱。} 基督暗示了祂的死亡……\par

10. \Quote{扰乱}从何而来？谁在扰乱？让我们看看他如何给犹太人烙上邪恶的良心，他们想要为自己杀害主的行为辩解。为此，正如福音所说，他们将祂交给了审判官，以免显得是他们杀了祂……让我们质问祂，并说：既然祢因扰乱而沉睡，谁曾迫害祢？谁杀了祢？难道是\Person{比拉多}，将祢交给士兵，挂在木头上，用钉子穿透吗？请听他们是谁——\BibleQuote{人子} \BibleRef{咏 56:5}。祂所说的是那些祂作为迫害者而受苦的人。但他们既未持刀，又如何杀害？他们未拔剑，未向祂攻击以杀祂；他们从何而杀？\BibleQuote{他们的牙齿是武器和箭矢，他们的舌头是利剑。} 不要看那未武装的手，而要看那武装的口；从那里出了剑，基督将被杀；同样，从基督的口中出了剑，犹太人将被杀。因为祂有一把双刃的剑：\BibleRef{默 1:16}，祂复活后击打了他们，并从他们中分离出那些祂要使其成为祂忠信子民的人。他们的是恶剑，祂的是善剑；他们的是恶箭，祂的是善箭。因为祂自己也有善箭，即善言，以此刺痛忠信的心，好使祂被爱。因此，他们的箭是一种，他们的剑是另一种。\BibleQuote{人子，他们的牙齿是武器和箭矢，他们的舌头是利剑。} 人子的舌头是利剑，他们的牙齿是武器和箭矢。那么，他们何时击打，岂非当他们呼喊：\Quote{钉死，钉死}的时候？\par

11. 主啊，他们对你做了什么？让先知在此欢欣！因为上面，所有这些诗句都是主在说话：确实是先知，但以主的身份，因为在先知中是有主的……他说：\Quote{请祢高升，}他说，\Quote{在诸天之上，天主啊。} 人在十字架上，而天主在诸天之上。让他们在地上继续狂怒，祢在天上审判。那些狂怒的人在哪里？他们的牙齿、武器和箭矢在哪里？难道\Quote{他们的鞭伤成了婴儿的箭矢}？因为诗篇在另一处这样说，想要证明他们徒然狂怒，徒然疯狂地冲向头：他们当时对钉在十字架上的基督无能为力，后来当祂复活并坐在天上时也是如此。婴儿如何为自己制造箭矢？用芦苇？但什么箭？什么力量？什么弓？什么伤口？\BibleQuote{天主啊，请祢高升于诸天之上，愿祢的荣耀充满全地} \BibleRef{咏 56:6}。那么，天主啊，为何高升于诸天之上？弟兄们，我们看不见高升于诸天之上的天主，但我们相信；但祂的荣耀充满全地，我们不仅相信，也看见。但异端者所受的是何种疯狂，我请你们注意。他们从基督教会之纽带中被割断，抓住一部分而失去整体，不愿与全地共融，而基督的荣耀正遍布全地。但我们公教徒在全地，因为我们与全世界共融，无论基督的荣耀散布何处。因为我们看到那时所唱的，如今已经实现。我们的天主已高升于诸天之上，祂的荣耀也充满全地。啊，异端的疯狂！你们所看不见的，你们与我一同相信；你们所看见的，你们却否认；你们与我一同相信基督高升于诸天之上，这是我们看不见的事；却否认祂的荣耀充满全地，这是我们看得见的事。\par

12. ……请诸位看主对我们说话，并以祂的榜样劝勉我们：\BibleQuote{他们为我的脚设下了网罗，屈抑了我的灵魂} \BibleRef{咏 56:7}。他们想要将牠从天上拉下来，并压到低处：\BibleQuote{他们屈抑了我的灵魂：在我面前掘了坑，自己反掉在其中。} 他们伤害了我，还是伤害了自己？看，祂已被高举于诸天之上，天主，并且看，祂的荣耀充满全地：我们看见基督的王国，犹太人的王国在哪里？既然他们做了不应该做的事，那么他们身上就发生了应该承受的事：他们自己挖了坑，自己掉了进去。因为他们迫害基督，对基督没有伤害，却伤害了自己。弟兄们，不要以为只有他们遭遇了这事。每一个为兄弟设坑的，必然自己掉进去……\par

13. 但善人的忍耐以准备好的心灵接受天主的旨意，并在磨难中夸耀，说出以下的话：\BibleQuote{天主啊，我的心已准备妥当，我要歌唱，我要弹奏} \BibleRef{咏 56:8}。他对我做了什么？他准备了坑，我的心已准备妥当。他准备了坑来欺骗，难道我不准备心灵去承受吗？他准备了坑来压迫，难道我不准备心灵去忍受吗？因此他将掉进去，但我要歌唱弹奏。请听在一位宗徒中心灵准备好的例子，因为他效法了他的主：\BibleQuote{我们在磨难中也夸耀，}他说，\BibleQuote{因为磨难生忍耐，忍耐生考验，考验生希望，而希望不令人羞惭，因为天主的爱借着所赐予我们的圣神，已倾注在我们心中。} \BibleRef{罗 5:3} 他曾在压迫、锁链、监牢、鞭打、饥渴、寒冷和赤身中，\BibleRef{格后 11:27}，在一切劳苦和痛苦的煎熬中，他却说：\BibleQuote{我们在磨难中夸耀。} 从何而来？岂不是因为他的心已准备好了？因此他歌唱并弹奏。\par

14. \Quote{我的荣耀，起来吧} \BibleRef{咏 56:9} 。那曾躲避撒乌耳面、逃进山洞的人说：\Quote{我的荣耀，起来吧：}耶稣在受难后受到光荣。\Quote{\OriginalTerm{琴瑟}{psaltery}和\OriginalTerm{竖琴}{harp}，起来吧。}祂呼唤什么起来？我见到两件乐器：但我见到一个基督的身体，一个肉身复活了，两件乐器也起来了。一件乐器是琴瑟，另一件是竖琴。\Quote{乐器}一词用于所有音乐家的工具。不仅那用风箱吹奏的大型乐器称为乐器，而且凡是适合演奏的、有形的、演奏者用作工具的都称为乐器。但这些乐器彼此有区别……那么这两件乐器为我们象征什么？因为主基督我们的天主正在唤醒祂的琴瑟和竖琴；祂说：\Quote{我将在黎明时起来。}我想你们现在看到主复活了。我们在福音中读到：\BibleRef{谷 16:2} 看到复活的时间。基督曾多久在阴影中被寻找？祂已光照，愿祂被承认；祂在黎明时复活了。但什么是琴瑟？什么是竖琴？主通过祂的肉身行了两种事：奇迹和苦难：奇迹来自上，苦难来自下。但祂所行的奇迹是神圣的；但祂通过身体行奇迹，通过肉身行奇迹。因此行神圣事的肉身是琴瑟：承受人间苦难的肉身是竖琴。让琴瑟奏响：瞎子得明，聋子得听，瘫痪者得力，跛子行走，病人痊愈，死人复活；这是琴瑟的声音。也让竖琴奏响：让祂饥饿、口渴、睡觉、被捉拿、被鞭打、被嘲笑、被钉十字架、被埋葬。因此，当你看到那肉身中有某些来自上、某些来自下的事，一个肉身复活了，我们在一肉身中认出琴瑟和竖琴。这两种完成的事成全了福音，并在万民中被宣讲：因为主的奇迹和苦难都被宣讲。\par

15. 因此，琴瑟和竖琴在黎明中起来，他向主称谢；并说什么？\Quote{主啊，我要在万民中称谢你，在列邦中歌颂你：因为你的慈爱高及诸天，你的信实上达云霄} \BibleRef{咏 57:10}。天高于云，云低于天：然而云属于这最近的天空。但有时云停留在山上，甚至如此远地在最近的气流中被卷起。但上面有一个天，是天使、宝座、宰制者、率领者、掌权者的居所。因此这可能似乎是应该说的：\Quote{你的信实高及诸天，你的慈爱上达云霄。}因为天使在天上赞美天主，看见信实的本体，毫无视觉的黑暗，没有虚假的混杂：他们看见、爱慕、赞美，不会疲倦。那里有信实：但在此处我们的苦难中确有慈爱。因为必须向可怜的人施行慈悲。上面没有可怜的人，所以不需要慈悲。我说这些，是因为似乎更合适说：\Quote{你的信实高及诸天，你的慈爱上达云霄。}因为我们理解\Quote{云彩}是真理的宣讲者，那些人承载着那肉身，在某种意义上黑暗，天主在其中透过奇迹闪耀，透过诫命雷响……光荣归于我们的主，也归于祂的慈爱和祂的信实，因为祂既没有因慈悲而放弃藉祂的恩宠使我们有福，也没有剥夺我们的信实：因为信实起初隐藏在肉身中来到我们这里，并藉祂的肉身医治我们内心的眼睛，以便日后我们能够面对面地看见祂。\BibleRef{格前 13:12} 因此让我们感谢祂，用同一篇圣咏的最后几节来说，我先前也曾说过：\Quote{天主，愿你崇高过于诸天，愿你的荣耀高过全地} \BibleRef{咏 57:11}。因为先知在多年前已向祂说了这话；现在我们看到这事；因此我们也要如此说。\par

\SourceRecord{Translated by J.E. Tweed.；Nicene and Post-Nicene Fathers, First Series；Vol. 8.；Edited by Philip Schaff.；Buffalo, NY: Christian Literature Publishing Co.,；1888.；<http://www.newadvent.org/fathers/1801057.htm>.}

% structure: p0001 source=h1 target=section reason=page-title

\section{论圣咏第五十八篇}

1. 我们所咏唱的这些话，与其说是宣讲，不如说是我们当聆听的。因为真理在人类集会上向众人呼喊：\BibleQuote{如果你们确实讲论正义，请判断正直的事，人子们。} \BibleRef{咏 57:1} 因为对哪一个不义的人而言，讲论正义不是一件容易的事呢？或者哪一个人，当被问及正义，且没有理由时，不会轻易回答正义的事呢？原来我们造物主的手在我们的心中已经写下了这真理：\BibleQuote{你不愿意别人对你做的事，你也不要对别人做。} \BibleRef{多 4:15} 关于这真理，甚至在法律颁赐之前，没有人被允许无知，为的是能有一个准则，用来判断那些甚至没有领受法律的人。但为了防止人们抱怨缺了什么，也在版上写下了他们心中所不读的。因为他们不是没有写，而是不愿意读。在他们眼前摆上了他们良心中所被迫看见的；并且好像天主的声音从外面带到他们那里，人就这样被赶回自己的内心，经上说：\BibleQuote{因为不虔敬的人的思想必将受审问。} \BibleRef{智 1:9} 哪里有审问，哪里就有法律。但因为人们贪图外在的事物，甚至从自身变成流亡者，所以也颁赐了成文法：不是因为心里没有写，而是因为你离弃了你的心，你被那无所不在的抓住了，并被唤回到你自身之内。因此，成文法对那些离弃了写在他们心中的法律的人喊叫什么呢？\BibleRef{罗 2:15} \BibleQuote{你们这些背信的人，回归你们的心吧。} \BibleRef{依 46:8} 因为谁教导你，不让别人接近你的妻子？谁教导你，不愿意别人对你行窃？谁教导你，不愿意受冤屈，以及其他任何普遍或特别可说的事？因为有许多事，若分别问人，他们会大声回答，他们不愿意受这些。来吧，如果你不愿意受这些事，你难道就是唯一的人吗？你难道不是生活在人类的团契中吗？那与你一同被造的，是你的同类；众人都是按照天主的肖像被造的，\BibleRef{创 1:26} 除非他们以尘世的遮盖抹去了祂所塑造的。所以，你自己不愿意被做的，也不要对别人做。因为你判断那件事中有恶，而你不愿意遭受；并且你被一条内在的法律约束着知道这事；它写在你心里。你在做某事时，在你的手中就有呼喊：当你从别人手中遭受这事时，你如何被迫使回到你的心呢？偷窃是好的吗？不！我问，奸淫是好的吗？众人都喊：不！杀人是好的吗？众人都喊，他们憎恶它。觊觎邻人的财产是好的吗？不！这是众人的声音。或者如果你还不承认，有一个觊觎你财产的人临近了：请愿意回答你愿意什么。所以众人，当被问及这些事时，都喊这些事不好。再者，关于行善，不仅是不伤害，而且还要施与和分施，任何饥饿的灵魂被这样质问：\Quote{你忍饥受饿，另一个人有饼，且超过所需，他知道你缺乏，却不给：当你饥饿时，这使你不悦；当你吃饱时，也让它使你不悦，当你知道了别人的饥饿。一个陌生人需要住所，来到你的地方，他不被接纳：于是他就喊叫那城市不人道，连在野蛮人中他也能找到家。他感受到不义，因为他遭受；你或许感觉不到，但你应该也把自己想象成一个陌生人；并且你要看到，那没有给予你所应给的东西的人，将如何使你厌恶，而在你的地方你也不会给陌生人。} 我问众人。这些事是真的吗？真的。这些事是公义的吗？公义的。但请听圣咏：\BibleQuote{如果你们确实讲论正义，请判断正直的事，人子们。} 不要只在唇间有正义，也要在行为上。因为如果你行事与说话相反，你说话好，判断却坏。\par

2. 现在，如果你们愿意，让我们来到当前的情况。因为那声音是甘甜的声音，如此熟悉于教会的耳朵，是我们主耶稣基督的声音，也是祂身体的声音，是教会劳苦、在地上旅居、生活在诽谤和谄媚之人的危险中的声音。你若不爱奉承者，就不会畏惧威胁者。因此，这声音所属的祂，已经观察并看见，所有的人都讲论正义。因为哪一个人不敢讲论正义，免得被称为不义呢？所以，当祂好像听见众人的声音，观察所有人的嘴唇时，祂向他们呼喊：\BibleQuote{如果你们确实讲论正义，}——如果不是虚假地讲论正义，如果嘴唇发出的不是一件事，而心里隐藏着另一件事——\BibleQuote{请判断正直的事，人子们。} 请从福音中听祂自己的声音，与这圣咏完全相同的：\BibleQuote{假善人，主对法利塞人说，你们既是恶人，怎能说出善事来？……你们或者使树好，它的果子也好；或者使树坏，它的果子也坏。} \BibleRef{玛 12:33} 你们为什么粉饰自己，泥墙？我知道你的内里，我不被你的覆盖所欺骗：我知道你出示什么，我知道你掩盖什么。\BibleQuote{因为祂不需要任何人给祂作证关于人：祂自己知道人心里有什么。} \BibleRef{若 2:25} 因为祂知道人心里有什么，祂造了人，并成为了人，为要寻找人……\par

3. 但如今你们做什么？我为何向你们说这些话？\BibleQuote{因为你们在心上作恶，在地上施行不义。}\BibleRef{咏 57:2} 也许只在心上作恶？请听下文：他们的心手相随，心手相从，事既思之，便行之；或者不行，并非因不愿，而是因不能，\Emph{凡你愿意而不能者，天主都算为已行}。\BibleQuote{因为你们在心上作恶，在地上施行不义。}然后呢？\BibleQuote{你们的手编织不义。}什么是\BibleQuote{编织}？从罪到罪，因罪而罪。这是什么？某人偷窃，是罪；他被看见，就想杀死看见他的人：罪与罪编织在一起。天主以祂隐藏的审判容许他杀死他所愿杀的人；他发觉事情泄露，又想杀第二人；他编织了第三罪。他计划这些时，恐怕被人发现或被证明是凶手，便去咨询占星家；第四罪加了进来；占星家或许给出严厉凶险的回答，他又跑去求巫师，求赎罪；巫师说他不能赎罪；又去找术士。谁能数算那些与罪编织在一起的罪？\BibleQuote{你们的手编织不义。}你们愈编织，愈将罪缚于罪上。从罪中解脱吧。但你说：我不能。向祂呼号吧。\Quote{我实在不幸！谁能救我脱离这该死的肉身？}\BibleRef{罗 7:24} 因为天主的恩宠必会来到，使正义成为你的喜乐，如同你曾以不义为乐；你，一个从束缚中被释放的人，要向天主呼喊：\Quote{祢解开了我的锁链。}\Quote{祢解开了我的锁链}不是别的，正是\Quote{祢赦免了我的罪}？听，为何它们是锁链：圣经回答说，\BibleQuote{各人被自己的罪孽束缚。}\BibleRef{箴 5:22} 不仅是锁链，而且是枷锁。枷锁是那些被扭结而成的：就是说，因为你将罪与罪编织在一起……\par

4. \BibleQuote{罪人一出母胎就与天主疏远，一离母腹就走错了路，说出虚谎的话。}\BibleRef{咏 57:3} 他们讲不义时，说的是虚谎的话，因为不义是欺骗的；他们讲正义时，说的也是虚谎的话，因为口里承认一事，心里却隐藏另一事。\Quote{罪人一出母胎就与天主疏远。}这是什么？让我们更仔细地探究：或许他这样说，是因为天主预知人即使在母腹中就将成为罪人。因为当\Person{黎贝加}尚怀孕，腹中怀着双胎时，何以说\Quote{我爱了\Person{雅各伯}，却恨了\Person{厄撒乌}}？因为曾说\Quote{年长的要服事年幼的}。那时天主的审判是隐藏的；但罪人从母胎，即从根源起，就已与天主疏远。与什么疏远？与真理疏远。与什么疏远？与福地、与福乐的生命疏远。也许他们从母胎就已疏远。那么，什么罪人从母胎疏远？因为人若不曾在母胎中被怀有，何能出生？或者今天活着听这些话而一无所获的人，若非出生，何能生存？因此，罪人或许是与某个母胎疏远了，那母胎中怀有那透过宗徒说话的慈爱：\BibleQuote{我再次为你们受产痛，直到基督在你们内形成为止。}\BibleRef{迦 4:19} 所以你要期待；你要被形成；不要为自己妄作论断，因为你或许不知道。你仍是属血肉的，你已被怀胎；从你领受基督之名的那一刻起，通过某种圣事，你已在母腹中诞生。因为人不仅从母腹出生，也在母腹中出生。首先他在母腹中出生，以便能从母腹出生。因此，甚至对\Person{玛利亚}说：\Quote{那在你内受生的，是出于圣神。}那时它尚未从她而生，但已在母腹受生。因此，在教会母腹中有些小者诞生，他们被形成后顺利娩出，而不致流产，这是好事。让母亲怀你，不要流产。如果你忍耐到底，直到你被形成，直到你拥有确定的真理教义，母腹应当保存你。但如果你因不耐烦而摇动你母亲的身体，她虽疼痛地把你们产出，但这对你损失更大，对她则损失较小。\par

5. 因此，他们之所以\Quote{一离母腹就走错了路}，是因为\Quote{他们说了虚谎的话}？或者，他们说了虚谎的话，岂不正是因为走错了路？因为在教会的母腹中，真理永存。谁若离开这个教会母腹，必然说虚谎的话；必然，我说，说虚谎的话；要么是不愿被怀胎，要么是被怀胎后母亲将其逐出。于是异端者攻击福音（我宁愿先谈那些被逐出而令我们哀悼的人）。我们向他们重申：看，基督曾说：\BibleQuote{基督必须受苦，第三天从死者中复活。}\BibleRef{路 24:46} 我在那里认出我们的元首，我在那里认出我们的新郎；你也与我一同认出新娘……\par

\BibleQuote{他们遭受愤恨，如同蛇一般} \BibleRef{咏 57:4}。你将要听到一件大事。\BibleQuote{他们遭受愤恨，如同蛇一般} \BibleRef{咏 57:4}。好像我们说过：你说的这是什么？接着又说：\BibleQuote{如同聋了的毒蛇}。何以聋？\BibleQuote{并塞住它的耳朵}。因此聋，因为它塞住耳朵。\BibleQuote{并塞住它的耳朵}。\BibleQuote{它不肯听行法术者的声音，不听贤者所配的药} \BibleRef{咏 57:5}。正如我们所听到的，连人们也以他们所能的研究说了这些，但无论如何，这是天主圣神比任何人更清楚知晓的事。因为这并不是徒然说的，而是因为关于毒蛇我们所听到的也可能是真的。当毒蛇开始受到马尔西亚念咒者的影响，用某种特殊咒语召唤它出来时，听听它做了什么……为了类比起见，请留意对你们所说的话，并注意要避免什么。因此，这里也给出了一个取自马尔西亚人的比喻，他念咒召唤毒蛇从黑暗的洞穴中出来；诚然，他要把毒蛇带到光明中，但毒蛇喜爱它的黑暗，在那里盘蜷隐藏自己，当它不愿出来时，它拒绝听见那些它感到约束它的话语，据说它把一只耳朵贴在地上，用尾巴堵住另一只耳朵，因此尽可能逃避那些话语，不向念咒者出来。圣神以此比喻提到某些人不听天主圣言，不仅不行，而且完全拒绝听，以免去行。\par

这事甚至在信仰初期就已发生。殉道者斯德望宣讲真理，向如同黑暗的心灵念咒，为将他们带到光明中；当他开始提到基督时，他们根本不愿听，圣经论到他们说什么？关于他们记载什么？它说：\BibleQuote{他们塞住耳朵} \BibleRef{宗 7:57}。但他们后来做了什么，斯德望的受难记叙述了。他们不是聋子，而是使自己聋了。……因为他们在基督被提名时做了这事。这些人的愤恨如同蛇的愤恨。你们为什么塞住耳朵？等一等，听一听，如果你能，就发怒吧。因为他们选择除了发怒什么也不做，他们不愿意听。但如果他们听了，也许他们会停止发怒。他们的愤恨如同蛇的愤恨……\par

\BibleQuote{天主已彻底打碎他们口中的牙齿} \BibleRef{咏 57:6}。指谁？指那些愤恨如同蛇一般的人，如同毒蛇塞住耳朵，不听行法术者的声音和不听贤者所配的药的人。主对他们做了什么？\BibleQuote{彻底打碎他们口中的牙齿}。这事已经成就，最初已经成就，现在仍在成就。但是弟兄们，只要说\BibleQuote{天主已彻底打碎他们的牙齿}就够了。法利塞人不愿听法律，不愿听基督关于真理的训诫，如同那蛇和毒蛇。因为他们喜爱过去的罪恶，不愿失去现在的生活，即不愿以地上的喜乐换取天上的喜乐……\Quote{在他们自己的口中}是什么意思？就是如此：他们用自己的口反对自己作出宣告：祂迫使他们用自己的口对自己宣判。他们曾因纳税的事要诽谤祂：\BibleRef{玛 22:17}祂没有说\Quote{纳税是合法的}或\Quote{纳税是不合法的}。祂愿意彻底打碎他们的牙齿，他们正张着嘴要咬人；但祂要在他们自己的口中行这事。如果祂说，应向凯撒纳税，他们就会诽谤祂，因为祂使犹太民族成为纳贡者，从而说了民族的坏话。因为他们因罪而纳税，受了屈辱，正如法律曾向他们预言的。他们说，如果我们祂命令我们纳税，我们就说祂是恶意中伤我们民族的人；但若祂说不纳税，我们就说祂教我们不服从凯撒。他们设下双重的圈套，如同要捉拿主。但他们来到谁面前？来到那位知道如何彻底打碎他们自己口中之牙齿的那位面前。祂说：\BibleQuote{把税币拿给我看} \BibleRef{玛 22:19}。你们为什么试探我，你们这些伪善的人？你们以为纳税的事吗？你们愿意行正义吗？你们寻求正义的劝告吗？\Quote{你们若果真谈论正义，就应判断正直的事，你们人子}。但现在因为你们一方面说，另一方面判断，你们是伪善的人：\Quote{你们为什么试探我，你们这些伪善的人}？现在我要彻底打碎你们口中的牙齿：\BibleQuote{把税币拿给我看}。他们就拿给祂看。祂不说，这是凯撒的；而是问这是谁的？为的是彻底打碎他们口中的牙齿。因为当祂问，这像和名号是谁的，他们说，是凯撒的。现在主就要彻底打碎他们口中的牙齿。你们现在已经回答了，现在你们口中的牙齿已被彻底打碎。\BibleQuote{凯撒的归凯撒，天主的归天主} \BibleRef{玛 22:21}。凯撒寻求他的像，你们就给他；天主寻求祂的像，你们也给祂。不要让凯撒从你们那里失去他的钱币，也不要让天主在你们内失去祂的像。他们就找不到可以回答的话。因为他们被派来诽谤祂；他们回去说，没有人能回答祂。为什么？因为他们口中的牙齿已被彻底打碎。同样的事也发生在下面：\Quote{祢凭什么权柄做这些事？我也问你们一个问题，回答我}。祂就问他们关于若翰的事，若翰的洗礼是从天上来的，还是从人来的？这样无论他们怎么回答都会对自己不利……\par

9. 主不喜悦那法利塞人，因为他请主吃饭，而一个罪妇走近主的脚前，他就向主抱怨说：\Quote{这人若是先知，必定知道这妇人是谁。}\BibleRef{路 7:39} 哦，你不是先知，你怎么知道主不知道那妇人是谁呢？原来因为他没有遵守犹太人外在的肉体上的洁净礼（其实这洁净远离内心），他就这样怀疑主。为免在这事上多说，主甚至愿意在他口中彻底打碎他的牙齿。主对他说：\Quote{一个债主有两个债户，一人欠五百银币，另一人欠五十银币；两人都无力偿还，债主就免了双方的债。这两人中，谁更爱他呢？}\BibleRef{路 7:41} 主这样提问，是为了让对方回答；对方这样回答，是为了让主在他口中打碎他的牙齿。……\par

10. \Quote{主打碎了狮子的颚骨。}不单是毒蛇的。毒蛇怎样呢？毒蛇诡诈地吐出毒液，散布出去，咝咝作响。外邦人公然发怒，像狮子咆哮。\Quote{外邦人为何怒吼，万民为何妄想？}当时他们窥伺着主。\Quote{给凯撒纳税，可以不可以？}\BibleRef{玛 22:17} 他们是毒蛇，是蛇；主在他们自己的口中彻底打碎了他们的牙齿。后来他们喊叫：\Quote{钉死他，钉死他！}现在不再是毒蛇的舌头，而是狮子的咆哮。但\Quote{主也打碎了狮子的颚骨。}或许这里不需要他在其中所没有加上的\Quote{他们的口中}。因为那些以诡诈问题窥伺的人，被迫用自己的回答被征服；但那些公然发怒的人，难道要用问题来驳斥吗？然而，连他们的颚骨也被打碎了：主被钉死之后，复活了，升了天，被尊为基督，被万国朝拜，被所有君王朝拜。犹太人如今尽管发怒吧，如果他们还能的话。我们也在异端者身上看到这个警告和先例，因为我们也发现他们像蛇一样因愤怒而变聋，不愿听从\Quote{智慧人所配的药}；主也在他们自己的口中彻底打碎了他们的牙齿。……\par

11. \Quote{他们必像流水一样被轻视}\BibleRef{咏 57:7}。弟兄们，不要被某些称为急流的溪水吓倒：它们充满了冬天的雨水，不要害怕；过一会儿它就流过去了，那水就流干了；它咆哮一时，很快会平息：它们不能持久。许多异端如今已经彻底死了：它们在自己的渠道里流了所能流的，流干了，渠道干涸了，几乎找不到它们的回忆，或者它们曾经存在过。\Quote{他们必像流水一样被轻视。}但不仅仅是他们；整个现世时代暂时咆哮着，寻找它可以拖走的人。让所有不虔敬的人，所有骄傲的人，像急流奔腾汇聚一样撞击他们骄傲的岩石，不要吓倒你们；它们是冬天的水，不能永远流淌；它们必然流到自己的地方，自己的结局。然而，主也喝了这世界的急流。因为他在这里受了苦，他喝了这急流，但只是在路上喝了，只是在经过时喝了，因为他没有停留在罪人的道路上。关于他，圣经说什么？\Quote{他在路上喝了急流，因此必抬起头来；}也就是说，他因此得荣耀，因为他死了；因此复活了，因为他受了苦。……\par

12. \BibleQuote{他们像蜡一样熔化，被带走}\BibleRef{咏 57:8}。因为你或许会说，并不是所有的人都像我这样被造得软弱，以使他们相信：许多人坚持作恶，执迷不悟。对于同样的恐惧，你不必担心：\BibleQuote{他们像蜡一样熔化，被带走。}他们不能与你对抗，不能持守：他们将在自己情欲之火中灭亡。因为这里有一种隐藏的惩罚，圣咏现在就要谈论它，直到结尾。只有寥寥数节；请留心。有一种将来的惩罚，地狱之火，永火。因为未来的惩罚有两种：一种是下界（阴府）的，就是那富人被焚烧的地方，他曾希望得到一滴水，从穷人的指尖滴到他的舌头上，那穷人他曾在他门前藐视过，他说：\BibleQuote{因为我在这火焰中受折磨。}\BibleRef{路 16:24}第二种是末日的惩罚，他们将听到：那些在左边的将被安置：\BibleQuote{你们进入那为魔鬼和他的使者所预备的永火里去。}\BibleRef{玛 25:41}这些惩罚将在那时显明，当我们离开此生时，或在世界末日人们来到死人复活时。那么现在就没有惩罚吗？天主竟容忍罪恶完全不受罚直到那一天吗？就在这里也有一种隐藏的惩罚，同样的惩罚他现在正在谈论……尽管如此，我们有时看到义人受这些惩罚之苦，而不义之人却与这些惩罚无缘：因此那后来欢欣地说的话，他的脚曾颤抖：\BibleQuote{天主对心地正直的人何其美善！但我的脚几乎动摇，因为我嫉妒恶人，看见罪人享平安。}因为他看见了恶人的福乐，并且很乐意成为一个恶人，看见恶人掌权，看见他们万事亨通，拥有一切现世之物的丰裕，而他自己当时还像一个婴孩，也向主渴望这些东西：于是他的脚动摇了，直到他看见在终末是应当盼望还是畏惧。因为他在同一篇圣咏中说：\BibleQuote{这事在我面前是劳苦，直到我进入天主的圣所，领悟到最后的结局。}因此，那并不是下界的惩罚，不是复活后永火的惩罚，也不是这在现世中义人和不义之人所共有的惩罚，而且常常义人的惩罚比不义之人更重；而是天主圣神要我们留意此生某种惩罚。请留意，听我即将讲述你们已知之事：但当它在圣咏中宣告出来时，就是更甜蜜的事，因为在未宣告之前，它被认为晦涩难解。因为看，我提出你们已经知道的事：但因为这事情是从你们从未见过的地方提出来的，所以即便是已知之事，也如同新事一般使你们喜悦。请听不敬虔之人的惩罚：他说：\BibleQuote{他们像蜡一样，}他说，\BibleQuote{熔化，被带走。}我说过，因着他们的情欲，这事就成就在他们身上。邪恶的情欲如同烧灼的火。火能烧毁衣服，难道奸淫的情欲不吞噬灵魂吗？当圣经论及预谋的奸淫时，它说：\BibleQuote{人若将火抱在怀中，他的衣服岂能不烧着呢？}\BibleRef{箴 6:27}你怀中抱着活炭；你的背心被烧透；你心中存着奸淫，你的灵魂还能安然无恙吗？但这些惩罚很少有人看见：因此天主圣神极力要我们留意它们。听宗徒说：\BibleQuote{天主就任凭他们随从心里的情欲。}\BibleRef{罗 1:24}看，那使他们像蜡一样熔化的火。因为他们放松了对贞洁的节制；因此这些人，随从他们的情欲，就被说成是松散和熔化的。何以熔化？何以松散？因情欲之火。\BibleQuote{天主就任凭他们随从心里的情欲，以致行那些不适宜的事，充满了各种不义。}……\par

13. \BibleQuote{火落在他们身上，他们却没有见到太阳。}你看出他是以何种方式谈论一种黑暗的惩罚。\BibleQuote{火落在他们身上}，傲慢之火，烟雾之火，情欲之火，愤怒之火。这火何等大？凡被这火落上的人，将见不到太阳。因此经上说：\BibleQuote{不可让太阳在你们的愤怒中落下。}\BibleRef{弗 4:26}所以，弟兄们，你们要害怕邪恶情欲之火，如果你们不想像蜡一样熔化，并从天主的面前灭亡。因为那火落在你们身上，你们就看不见太阳。什么太阳？不是那个与你们一同被走兽、昆虫以及善人与恶人所看见的太阳，因为\BibleQuote{祂使太阳升起，光照善人也光照恶人。}\BibleRef{玛 5:45}但有另一个太阳，那些人将要说：\BibleQuote{太阳没有为我们升起，一切都如影子般消逝。因此我们迷失了真理的道路，正义之光没有照耀我们，太阳也没有为我们升起。}\BibleRef{智 5:6}……\par

14. \BibleQuote{在荆棘生出你们的刺以前：祂要如同活人，如同发怒，将祂们吞灭} \BibleRef{咏 57:9}。什么是荆棘？它是一种有刺的植物，据说上面有某些最紧密的刺。起初它是一棵草；当它是草时，柔软而美丽；但上面仍有刺要长出来。所以现在罪是令人愉快的，似乎并不扎人。荆棘是草；然而现在仍有刺。\Quote{在荆棘生出刺以前：}是在悲惨的快乐和愉悦的明显折磨出现之前。让那些爱慕任何对象却得不到的人反省自己；让他们看看是否不被渴望折磨：当他们得到了他们非法渴望的东西时，让他们注意是否不被恐惧折磨。因此，让他们在这里看到自己的惩罚；在那种复活来临之前，那时复活在肉身中的人将不会被改变。\BibleQuote{因为我们都要复活，但并非我们都要被改变。} \BibleRef{格前 15:51}因为他们将有肉身的腐朽以便受苦，而不是在其中死去；否则那些痛苦也会结束。那时，那荆棘的刺，即所有痛苦的刺痛和折磨，将要生出。这样的刺，那些要说\BibleQuote{这些人就是有时我们曾讥笑的人} \BibleRef{智 5:3}的人将要承受，那是悔改的刺痛，但为时已晚，没有果实，如同荆棘的不毛。此时的悔改是医治的痛苦；那时的悔改是刑罚的痛苦。你不想受那些刺吗？在这里被悔改的刺刺透，好使你行那所说的事：\Quote{我在忧伤中转回，当刺扎入时：我认清了我的罪，没有掩盖我的过犯；我说，我要向主承认我的过犯，你就赦免了我心的不虔敬。}现在就这样做，现在就被刺透，不要在你身上成就那关于某些可憎之人的话：\Quote{他们被劈开，却没有被刺透。}观察那些被劈开却没有被刺透的人。你看见人被劈开，却看不见他们被刺透。看哪，他们在教会之外，并不悔改，以致他们应回到他们被劈开的地方。荆棘以后要生出它们的刺。他们现在不愿接受医治的刺透，以后将有刑罚的刺透。但即使在荆棘生出刺之前，已有火落在他们身上，使他们不见太阳，即天主的愤怒正在吞灭他们，他们仍活着：那是邪恶情欲、虚空荣誉、骄傲、贪婪的火；以及一切压住他们，使他们不认识真理，以致他们似乎不被征服，甚至不被真理本身制服的重担。因为弟兄们，还有什么比被真理制服并被真理胜过更光荣的事呢？让真理征服你，你愿意；因为即使你不愿意，真理本身也会征服你……\par

15. 下界的惩罚尚未到来，永恒的火尚未到来：让那在天主内成长的人现在将自己与不虔敬的人相比，将盲目的心与\OriginalTerm{受光照的}{enlightened}的心相比：比较两个人，一个在肉身中看见，一个看不见。肉身的看见有什么了不起呢？托彼特难道有肉身的眼睛吗？\BibleRef{多 4:3}他自己的儿子有，他没有；而一个盲人给一个看见的人指出了生命之路。因此，当你看到那种惩罚时，要喜乐，因为你不在其中。\par

因此，圣经说：\BibleQuote{义人见复仇时必喜乐} \BibleRef{咏 57:10}。不是那将来的惩罚；因为看看后面的话：\BibleQuote{他要在罪人的血中洗手。}这是什么意思？请你们的爱德注意。当杀人犯被击杀时，难道有无辜的人去到那里洗手吗？但\Quote{在罪人的血中他要洗手}是什么意思？当义人看到罪人的惩罚时，他自己成长；一个人的死是另一个人的生命。因为如果灵的血从那些内里已死的人流出，你看到这样的复仇，就在其中洗手；为将来更洁净地生活。他怎么洗手呢，如果他是义人？因为他手上有什么可洗的呢，如果他是义人？\BibleQuote{但义人因信德而生活。} \BibleRef{罗 1:17}因此，义人他称为信徒：从你相信的那一刻起，你立刻开始被称为义人。因为罪过已得到赦免。即使在你余生的那部分中，有些罪是你的，这些罪不能不停留，就像海水进入船舱一样；然而，因为你已相信，当你看见那完全转离天主的人在那盲目中被杀，那火落在他身上使他不见太阳——那时，你现在藉着信德看见基督，为了你能在实质上看见（因为义人因信德生活），观察那不虔敬之人死去，并洁净自己脱离罪恶。这样，你就在某种意义上在罪人的血中洗手。\par

16. \BibleQuote{有人要说：那么，义人果真有果实吗？}\BibleRef{咏 57:10}。看，在恩许的还没有来到以前，在永生还未赐下以前，在不虔敬的人被投入永火以前，在此生中义人已有果实。什么果实？\BibleQuote{因望德而喜乐，在患难中忍耐。}\BibleRef{罗 12:12}义人有什么果实？\BibleQuote{我们在患难中也夸耀，因为知道患难生出忍耐，忍耐生出考验，考验生出望德；望德不叫人蒙羞，因为天主的爱已借着赐给我们的圣神，倾注在我们心中了。}\BibleRef{罗 5:3}酒徒岂不喜乐？义人岂不喜乐？在爱中，义人有果实。一个可怜，即使他使自己沉醉；另一个有福，即使他饥渴。一个酗酒饱食，另一个以望德为粮。因此，让他看见他人的刑罚，自己的喜乐，并思念天主。那位如今已赐下如此信德、望德、爱德、祂圣经真理之喜乐的天主，在末后预备了何等的喜乐？在路上祂如此喂养，在家中祂将如何充满他？\BibleQuote{有人要说：那么，义人果真有果实吗？}让看见的人相信、看见并领悟。义人看见报应时将喜乐。但如果他没有眼睛可以看见报应，他将会悲伤，且不会因此改正。但如果他看见了，他就看到内心的眼被蒙蔽与内心的眼被光照之间的区别：贞洁的清凉与情欲的火焰之间的区别，望德的安全与犯罪中的恐惧之间的区别。当他看见这些时，让他与自己分离，并在同样的血中洗手。让他从比较中获益，并说：\BibleQuote{因此义人果真有果实，因此在地上有一位审判他们的天主。}不在那个生命里，不在永恒之火里，也不在阴府里，而是在这里，在地上……\par

17. 如果我们有点过于冗长，请谅解我们。我们奉基督的名劝勉你们，善加默想你们所听的事。因为即使宣讲真理也毫无价值，如果心与口不一致；而听到真理也无益处，如果一个人不把根基立在磐石上。那建立在磐石上的，是那听了而行的人：\BibleRef{玛 7:24}但听了而不行的人，是把根基立在沙土上；那既不听也不行的人，什么也不建立……\par

\SourceRecord{Translated by J.E. Tweed.；Nicene and Post-Nicene Fathers, First Series；Vol. 8.；Edited by Philip Schaff.；Buffalo, NY: Christian Literature Publishing Co.,；1888.；<http://www.newadvent.org/fathers/1801058.htm>.}

% structure: p0001 source=h1 target=section reason=page-title

\section{圣咏第五十九篇注释}

第一部分。\par

就如圣经惯于将圣咏的奥秘置于标题中，并以一个奥秘的崇高宣告来装饰圣咏的额，好让我们这些将要进入的人能够知道（仿佛在读门楣上的文字，知道里面在做什么）——这房子是属于谁的，或者那产业的主人是哪位：同样，在这篇圣咏中，也写有一个标题的标题。因为它写着：\Quote{直到终末，不要毁坏，为达味自己，直至题名题署。}这就是我所说的标题之标题。因为那题署是什么，他禁止毁坏，福音向我们指明了。因为当主被钉在十字架上时，比拉多写了一个题名并安放，\Quote{犹太人的君王，}\BibleRef{玛 27:37}用三种语言——希伯来文、希腊文和拉丁文\BibleRef{若 19:20}，这三种语言在全世界最为通行。……因此\Quote{不要毁坏}是最恰当和预言性的；因为甚至那些犹太人当时也向比拉多建议，说：\Quote{不要写\InnerQuote{犹太人的君王}，而要写：他自己说他是犹太人的君王。}\BibleRef{若 19:21}因为，他们说，这个题名已立他为王统治我们。比拉多说：\Quote{我所写的，我已经写了。}这就应验了：\Quote{不要毁坏。}\par

2. 并非只有这篇圣咏有这样的题署，即不要毁坏标题。有数篇圣咏这样标记在表面，但无论如何都是预言主的受难。因此，在这里也让我们领会主的受难，并让基督——头和身体——对我们说话。这样，总是或几乎总是，让我们从圣咏中听到基督的话，如同我们不仅注视那位头、天主与人之间的唯一中保、基督耶稣那个人。\BibleRef{弟前 2:5}……但让我们想到基督——头与整个身体，一个完整的人。因为对我们说了：\Quote{但你们是基督的身体和肢体，}\BibleRef{格前 12:27} 由宗徒保禄所说。因此，如果他是头，我们是身体；整个基督是头和身体。因为有时你会找到不适合于头的话，除非你将它们连接到身体，否则你的理解会摇摆不定；再者，你会找到仅适用于身体的话，然而基督却是在说话。在那里，我们不必害怕人会弄错：因为他很快就会调整到适应于头，即他看见不适合于身体的部分。……\par

3. 因此，让我们听听下面的话：\Quote{当撒乌耳派人看守房屋，为要杀他的时候。}这件事虽然不属于主的十字架，却属于主的受难。因为基督被钉死，死了，并被埋葬。那坟墓就好比房屋；看守它的犹太政府派人去看守基督的坟墓。\BibleRef{玛 27:66}在列王纪的圣经中确实有一个故事，关于撒乌耳派人看守房屋为要杀达味的事。\BibleRef{撒上 19:11}……但是正如撒乌耳没有实现他杀达味的意图一样，犹太政府也无法实现，即让睡觉的守卫的证词比醒着的宗徒的证词更有效。因为守卫被指示说什么？他们说：我们给你们任意多的钱；你们就说，当你们睡觉时，他的门徒来了，把他偷走了。看，何等虚假的证人反对真理和基督的复活，借着撒乌耳所预表的，他的敌人就提出了。不信的人，请询问睡着的证人，让他们回答你在坟墓里发生了什么事。如果他们睡着了，从哪里知道？如果醒着，为什么不抓住小偷？因此，让下面的话说吧。\par

4. \Quote{求你从我的仇敌中拯救我，我的天主，从起来攻击我的人中救赎我。}\BibleRef{咏 58:1}这事已经在基督的肉身上成就了，也在我们身上正在成就。因为我们的仇敌，即魔鬼及其使者，没有一天停止起来攻击我们，并因我们的软弱和脆弱而想要戏弄我们，用欺骗、建议、诱惑和各种网罗来缠住我们，当我们还在地上生活时。但让我们的声音警醒地朝向天主，并在基督的肢体中，在位于天上的头之下，呼喊：\Quote{求你从我的仇敌中拯救我，我的天主，从起来攻击我的人中救赎我。}\par

5.\BibleQuote{求你救我脱离作恶的人，脱离流人血的人，求你拯救我}\BibleRef{咏 58:2}。他们确实是流人血的人，他们杀害了那义者，在他身上找不到任何罪过；他们是流人血的人，因为当那外邦人洗手，想要释放基督时，他们却喊说：\BibleQuote{钉死他，钉死他}\BibleRef{玛 27:23}；他们是流人血的人，当基督的血的罪行被归咎于他们时，他们回答，并将这罪传给后代喝：\BibleQuote{他的血归到我们和我们的子孙身上}\BibleRef{玛 27:25}。但是，流人血的人也没有停止起来攻击他的身体；因为甚至在基督复活和升天之后，教会仍遭受迫害，首先是那从犹太人中生长出来的教会，我们的宗徒也属于那教会。起初，在那里，\Person{斯德望}（\OriginalTerm{冠冕}{Stephanus}）被石击（见：\BibleRef{宗 7:58}），并领受了那与他的名字相符的。因为斯德望确实意指冠冕。卑微地被石击，却崇高地被加冕。其次，在外邦人中，外邦的王国兴起，在它们身上应验了那预言：\Quote{地上的列王都要朝拜他，万民都要事奉他}之前，那王国的凶暴向基督的见证人怒吼：殉道者的血大量而频繁地倾流：当这血倾流时，如同播种一般，教会的田地更加多产地发出，并充满了整个世界，正如我们现在所看见的。因此，基督，不仅是头，也是身体，从这些流人血的人中被拯救出来。基督从流人血的人中被拯救出来，无论是过去的、现在的，还是将来的；基督被拯救出来，无论是那先行的、现在的，还是将来的。因为基督是整个基督的身体；凡现在、以前和以后的好基督徒，都是整个基督，他从流人血的人中被拯救出来；这声音并非虚空的：\Quote{求你救我脱离流人血的人}。\par

6.\BibleQuote{因为看哪，他们狩猎了我的灵魂……有强横的人冲向我}\BibleRef{咏 58:3}。然而，我们不可忽略这些强横的人：我们必须勤勉地探究这些起来的强横的人是谁。强横的人，起来攻击谁呢？岂不是攻击软弱的人、无力的人、不坚强的人吗？然而，软弱的人受赞美，强横的人受谴责。若要识别谁是强横的人，首先，主亲自称魔鬼为强横的人：\BibleQuote{没有人能进入强横人的家，夺取他的家具，除非先把那强横的人捆住}\BibleRef{玛 12:29}。因此，他用自己统治的锁链捆住了那强横的人；夺取了他的家具，并使他们成为自己的家具。因为所有不义的人都是魔鬼的家具……但在人类中间，有一些强横的人，具有应受责备和应受惩罚的力量，他们确实自信，但相信的是暂时的福乐。那人的事迹刚在福音（见：\BibleRef{路 12:16}）中读到，难道他不令你感到是强横的吗？他的田地出产丰富，他烦恼起来，想出了重建的计划，以便拆毁旧仓房，建造更宽敞的新仓房，完工后，他对自己的灵魂说：\Quote{你有许多财物，灵魂，吃喝享乐吧，尽情欢庆吧}……还有另一种强横的人，不是因财富，不是因身体力量，不是因任何地位上暂时卓越的权力，而是依靠自己的正义。这类强横的人必须警惕、畏惧、抵御，不可效仿：我说的是依靠自己的正义的人，而非依靠身体、财产、出身、荣誉的人；因为谁看不出这一切都是暂时的、易逝的、衰败的、飞逝的呢？但他们依靠自己的正义……\BibleQuote{你们的老师为什么同税吏和罪人一起吃饭呢？}\BibleRef{玛 9:11}啊，你们这些强横的人，你们不需要医生！这种强健不属于健康，而属于疯狂。因为连疯癫的人都比正常人更强壮、更有力；但他们的力量越大，死亡就越近。因此，愿天主使我们远离效仿这些强横的人……因此，这些强横的人就是攻击基督的人，他们称赞自己的正义。请听这些强横的人：当耶路撒冷的一些人，被他们派去捉拿基督，却不敢捉拿时（因为当他愿意时，他才被捉拿，那才是真正强横的人），他们说：\Quote{为什么你们不能捉拿他？}他们回答说：\Quote{从来没有一个人像他那样讲话。}而这些强横的人说：\BibleQuote{难道法利塞人中或经师中有人信了他吗？只有这些不明白法律的群众}\BibleRef{若 7:45}他们把自己置于那奔向医生的患病群众之上；这从何而来，不正是因为他们是强横的人吗？更糟的是，他们用他们的强横，把所有的群众也拉拢到自己一边，杀死了众人的医生……\par

7.接下来是什么？\BibleQuote{我没有不义，也没有罪恶，上主}\BibleRef{咏 58:4}。那些依靠自己正义的强横的人确实冲了上来，他们冲了上来，却没有在我身上找到罪恶。因为那些强横的人，即那些自以为义的人，凭什么能迫害基督呢？岂不是把他当作罪人吗？但是，无论如何，让他们看看自己多么强横，是在热病的狂躁中，而非在健康的活力中；让他们看看自己多么强横，如何像义人一样对不义的人发怒。然而，\Quote{我没有不义，也没有罪恶，上主。我无过地奔跑，我受到了引导。}因此，那些强横的人不能跟随我奔跑；所以他们把我当作罪人，因为他们没有看见我的脚步。\par

8. \Quote{没有不义，我奔跑，并被引导；起来迎接我，并观看。} 这是对天主说的\Emph{这话}。但为什么呢？如果祂不迎接，祂就不能观看吗？这就像你在路上行走，从远处不能被某人认出，你会呼唤他，说：迎接我，观看我如何行走；因为当你从远处看见我时，你不能看见我的脚步。同样，除非天主迎接，否则祂怎能看见他是如何没有不义地被引导，以及他是如何没有罪地奔跑呢？的确，这种解释我们也可以接受，即\Quote{起来迎接我}如同\Quote{帮助我}。但祂所加的\Quote{并观看}必须被理解为：使我的奔跑被看见，使我的被引导被看见；按照那个比喻，其中这话也对亚巴郎说过：\BibleQuote{现在我确知你敬畏天主。}\BibleRef{创 22:12} 天主说：\Quote{现在我确知：} 从何而知？岂不是因为我使你知？因为在试探的追问之前，每个人都对自己无知；正如伯多禄在自己的自信中无知（\BibleRef{玛 26:35}），并通过否认学会了拥有何种能力，在他自己的失足中，他意识到他的自信是虚假的：他哭了，在哭中他有益地学会了认识自己是什么，并成为他所不是的。因此，亚巴郎在被试探时，变得自知；天主说：\Quote{现在我确知}，也就是说，现在我使你知。同样，欢喜的日子因为它使人欢喜；悲哀的苦味因为它使尝者悲哀：同样，天主的观看是使人观看。\Quote{起来吧，}祂说，\Quote{迎接我，并观看}\BibleRef{咏 58:5}。什么是\Quote{并观看}？就是帮助我，即在这些人们中，使他们看见我的奔跑，跟随我；不要让那正直的在他们看来是弯曲的，不要让那遵守真理规则的在他们看来是扭曲的。\par

9. 我还被劝诫在此处谈论我们元首自身的崇高：因为祂甚至软弱至死，并取了肉身的软弱，为聚集耶路撒冷的子女到祂翅膀下，像一只带着幼雏显得软弱的母鸡。\BibleRef{玛 23:37} 我们难道没有在某时某地观察到这件事发生在某只鸟身上，甚至在我们眼前筑巢的鸟身上，例如家麻雀，燕子——如同我们每年的客人，鹳，以及各种鸟类——它们在我们眼前筑巢、孵蛋、喂养小鸡，甚至我们每天看到的鸽子？难道我们不知道、没有看见、没有观察过某只鸟对它的幼雏变得软弱吗？母鸡是如何经历这种软弱的？无疑，我所说的是众所周知的事，每天在我们眼前发生。她的声音如何变得嘶哑，她的全身如何变得虚弱？翅膀下垂，羽毛松散，你看到小鸡周围有些病态的东西，这就是母爱，表现为软弱。那么，主为什么愿意成为一只母鸡，在圣经中说：\BibleQuote{耶路撒冷，耶路撒冷，我多少次愿意聚集你的子女，像母鸡聚集小鸡在翅膀下，而你们不愿意。}\BibleRef{玛 23:37} 但祂已经聚集了万民，像母鸡聚集小鸡一样……\par

10. \Quote{祢，万军的上主天主，以色列的天主。} 祢是以色列的天主，人们以为祢只是那一个敬拜祢的民族的天主，当万民敬拜偶像时，祢，以色列的天主，\Quote{请眷顾所有民族。} 愿那预言应验：依撒意亚以祢的口吻对祢的教会，祢的圣城，那荒废者说：因为那被遗弃的妇人的儿子比那有丈夫的妇人的儿子更多。确实，对她说：\BibleQuote{不生育的，不生产的，欢呼吧！}\BibleRef{依 54:1} 等等，比那有丈夫的犹太民族更多，那民族获得了法律，有可见的君王。因为你们的君王是隐藏的，妳借着隐藏的新郎有更多的儿子……先知加说：\BibleQuote{扩展妳帐幕的地方，坚固妳的庭院：妳无需节省，加长妳的绳索，将坚固的橛子一次又一次地带往右边和左边。}\BibleRef{依 54:2} 在右边保守善人，在左边保守恶人（\BibleRef{玛 25:33}），直到簸箕到来（\BibleRef{玛 3:12}）；尽管如此，占据所有民族；被邀请参加婚筵的善人和恶人，婚筵坐满了客人（\BibleRef{玛 22:9}）；邀请是仆人的任务，分开是主的事。\BibleQuote{妳将居住在被遗弃的城市中。}\BibleRef{依 54:3} 被天主遗弃，被先知遗弃，被宗徒遗弃，被福音遗弃，充满魔鬼。因为妳将得胜；不要因妳可憎而羞愧。因此，虽有强人起来攻击妳，不要羞愧：当反对基督之名的法律颁布时，当成为基督徒是耻辱和恶名时。\Quote{不要因妳可憎而羞愧：因为妳将永远忘记耻辱，妳不再记念妳寡居的羞辱。}……\par

11. \Quote{不要怜恤一切作恶的人。}此处祂显然是在震慑。祂会不震慑谁呢？有哪个人反省自己的良心时不战栗呢？即使良心自知虔诚，若不在某种程度上自知有罪，也是奇怪的。因为凡犯罪的，也就是行不义。\BibleRef{若一 3:4}\Quote{主，祢若究察罪孽，谁能站得住呢？}然而，这话是真实的，并非无的放矢，也决不会落空：\Quote{不要怜恤一切作恶的人。}但祂甚至怜恤了保禄，他起初作为扫禄时曾作恶。因为他做了什么好事，凭此可受恩于天主？难道他不是恨祂的圣徒至于死地吗？\BibleRef{宗 9:1}难道他没有从大司祭那里领受文书，为要无论在哪里找到基督徒，就把他们押去受刑？当他决意此事，正在前行，呼吸并喘息着杀戮——如圣经为他作证的——难道不是从天上有大声音召唤他，把他击倒，又扶起他；使他眼瞎，又使他明亮；杀死他，又使他活；毁灭他，又恢复他？这是对什么功绩的回报？我们什么也不说；不如让他自己听吧：\Quote{我从前是}，他说，\Quote{亵渎者、迫害者和凌辱者，但我蒙受了怜悯。}\BibleRef{弟前 1:13}诚然\Quote{祢不会怜恤一切作恶的人}：这话可有两种理解：要么事实上天主不放过任何罪恶而不惩罚；要么有一种罪孽，天主确实不怜恤作它的人。\par

12. 一切不义，无论大小，都必定要受惩罚，或由人自己悔改而罚，或由天主报复而罚。因为连悔改的人也在惩罚自己。所以，弟兄们，让我们惩罚自己的罪恶吧，如果我们寻求天主的怜悯。天主不能怜恤一切作恶的人，如同纵容罪恶，或不铲除罪恶。总之，要么你惩罚，要么祂惩罚……\par

13. 现在让我们看看这句话的另一种理解方式。有一种罪孽，作它的人不能得到天主的怜悯。你或许要问，那是什么？就是辩护罪恶。当一个人辩护自己的罪恶时，他就行了极大的不义：他在辩护天主所恨恶的事。请看这是多么乖谬，多么邪恶。凡他所行的善，他愿归功于自己；凡所行的恶，他归咎于天主。因为人就是这样，将罪恶推诿于天主，这是一种更重的罪。……因此你辩护你的罪恶，竟把责任推给天主。这样罪人被开脱，好让审判官受责。然而，天主绝不怜恤一切作恶的人。\par

14. \Quote{愿他们在傍晚归正}\BibleRef{咏 58:6}。这是指某些人，他们曾是作恶的人，曾一度是黑暗，却在傍晚归正了。\Quote{在傍晚}是什么意思？在以后。\Quote{在傍晚}是什么意思？较晚的时候。因为先前，在他们钉死基督之前，他们本该认出他们的医师。因此，当祂被钉死——复活、升天之后——祂派遣了圣神，充满了在一间房子里的人，他们开始用万国的言语说话时，钉死基督的人害怕了；他们被良心刺痛，从宗徒那里寻求救恩的劝告，他们听到：\Quote{你们悔改，并要因我们主耶稣基督的名受洗，你们的罪就赦免了。}\BibleRef{宗 2:38}在杀害基督之后，在流基督的血之后，你们的罪被赦免了。……因此，\Quote{愿这些人归正}，他们也\Quote{在傍晚归正}。愿他们渴望天主的恩宠，认识自己是罪人；愿那些强壮的人变为软弱，富有的人变为贫穷，自以为义的人承认自己有罪，狮子变为狗。\Quote{愿他们在傍晚归正，像狗一样忍饥挨饿。他们要绕城而行。}什么城？那世界，圣经在某些地方称之为\Quote{环绕之城}；因为各处万民都环绕着犹太这一个民族，那里正说着这样的话，它被称为\Quote{环绕之城}。那些如今已成为饥饿的狗的人要绕这城而行。他们怎样绕行？通过宣讲。扫禄从一只狼在傍晚变成了狗，即晚期归正，藉着主掉落的碎屑，在他的恩宠中奔跑，绕城而行。\par

15. \BibleQuote{看，他们要用自己的口说话，他们的嘴唇上有剑}。\BibleRef{咏 58:7} 这里就是那把双刃剑，宗徒所说的：\BibleQuote{圣神的剑，就是天主的圣言}。\BibleRef{弗 6:17} 为何双刃？不是因为从两约出击吗？这把剑曾杀死了那些，对伯多禄说过的：\BibleQuote{起来，宰了吃}。\BibleRef{宗 10:13} \Quote{他们的嘴唇上有剑。因为谁曾听过？} 他们都在用自己的口说：\Quote{谁曾听过？} 那就是，他们对那些迟迟不相信的人感到愤怒。不久之前他们自己也不愿相信的人，现在看到别人不信就感到厌恶。确实，弟兄们，情况就是这样。你看一个人迟迟不肯成为基督徒；你每天都向他呼喊，他几乎不肯归化；假设他归化了，然后他就希望所有人都成为基督徒，并惊讶为什么别人还没有成为基督徒。他碰巧在晚上归化了；但因为像狗一样饥饿，他的嘴唇上也有剑；他说：\Quote{谁曾听过？} 什么是\Quote{谁曾听过？} \BibleQuote{有谁相信了我们的听闻？上主的臂膀向谁显示了呢？} \BibleRef{依 53:1} \Quote{因为谁曾听过？} 犹太人不信；他们转向了外邦人，并宣讲了。犹太人不信；然而，通过相信的犹太人，福音传遍了全城，他们说：\Quote{因为谁曾听过？} \BibleQuote{主，你却要讥笑他们}。\BibleRef{咏 58:8} 所有外邦人都将成为基督徒，而你说：\Quote{谁曾听过？} 什么是\BibleQuote{你将讥笑他们}？\BibleQuote{你将把万民视为虚无。} 对你来说，这是虚无；因为让万民都相信你是一件最容易的事。\par

16. \BibleQuote{我的力量，我要为你持守}。\BibleRef{咏 58:9} 那些强壮的人之所以跌倒，正是因为他们没有将力量持守于你；也就是说，那些起来攻击我并冲向我的人，依靠了他们自己。而\Quote{我的力量，我要为你持守}：因为如果我退后，我就会跌倒；如果我靠近，我会变得更坚强。请看，弟兄们，在人灵魂中有什么？它本身没有光，本身没有能力；灵魂中一切美好的东西，就是德行和智慧；但灵魂既不是为自己而有智慧，也不是为自己而强壮，它本身不是自己的光，也不是自己的德行。有一个德行的根源和泉源，有一个智慧的根，有一个——如果可以这么说的话——不变真理的领域。灵魂远离它就会变得黑暗，靠近它就会变得光明。\BibleQuote{你们亲近他，就会变得光明}：因为远离他，你们就会变得黑暗。因此，\Quote{我的力量，我要为你持守}：我不会远离你，我不会依靠我自己。\Quote{我的力量，我要为你持守，因为天主，你是我的提拔者。} 因为我曾在哪里，现在又在哪里？你从何处提拔了我？你赦免了我哪些罪过？我曾躺在哪里？我被提升到了何处？我本应记得这些事，因为在另一篇\WorkTitle{圣咏}中曾说：\BibleQuote{因为我的父亲和母亲舍弃了我，但上主却收留了我。}\par

17. \BibleQuote{我的天主，他的仁慈必在我以先}。\BibleRef{咏 58:10} 看这就是\Quote{我的力量，我要为你持守}：我绝不依靠自己。因为我带来了什么好处，以致你怜悯我、使我成义？你在我里面找到了什么，除了罪过之外？属于你的一切，只有你所创造的本性；其他是我自己的恶事，是你所涂抹的。我没有首先起来走向你，而是你到来唤醒了我：因为\Quote{他的仁慈必在我以先}。在我行任何善事之前，\Quote{他的仁慈必在我以先}。不幸的佩拉纠对此将如何回答？\BibleQuote{我的天主在仇敌中向我显示}。\BibleRef{咏 58:11} 他对我施予了多大的仁慈，他在仇敌中向我显示。让被聚集的人将自己与被遗弃的人比较，让被选的人将自己与被弃的人比较；让仁慈之器将自己与义怒之器比较；让他看见天主如何从同一团泥中造了一个器皿为尊贵，另一个器皿为卑贱。\par

\BibleQuote{因为天主愿意显示义怒，并彰显自己的能力，就以极大的忍耐包容了那些义怒之器，他们已为灭亡而预备好了。} \BibleRef{罗 9:22} 这是为什么？为了将他丰富的荣耀彰显在仁慈之器上。因此，如果天主包容了义怒之器，为要在仁慈之器上彰显他丰富的荣耀，那么最正确地说：\BibleQuote{他的仁慈必在我以先：我的天主在仇敌中向我显示。} 也就是说，他对我施予了多大的仁慈，他在这些人中向我显示，他没有怜悯他们。因为除非债务人被吊着，他对赦免债务的人就不那么感激。\BibleQuote{我的天主在仇敌中向我显示。}\par

18. 至于敌人本身又如何？\Quote{不要杀灭他们，恐怕他们忘记祢的法律。} 他是在为他的仇敌求情，他是在成全诫命。不要杀灭那些人，祢所杀灭的是他们的罪。但是被杀是什么意思？就是忘记主的法律。这是真正的死亡，落入罪恶的深渊；这确实也可以理解为针对犹太人。为何针对犹太人说\Quote{不要杀灭他们，恐怕他们忘记祢的法律}？那些杀害了我的仇敌，求祢不要杀灭他们。让犹太民族存留：诚然他们已被罗马人征服，诚然他们的城已荡然无存，犹太人不得进入他们的城，然而，犹太人仍然存在。因为所有那些省份都已被罗马人征服。如今谁能区分罗马帝国中的各个民族？因为他们都成了罗马人，都被称为罗马人。然而，犹太人带着一个标记存留下来；他们并没有被征服到被征服者吞没的地步。这并非无缘无故，正如加音杀了弟兄后，天主在他身上立了一个记号，免得任何人杀他。\BibleRef{创 4:15} 这就是犹太人所拥有的记号：他们坚守自己法律的残余，行割损，守安息日，献逾越节祭，吃无酵饼。所以他们确实是犹太人，没有被杀灭，他们对信主的民族是必要的。为什么？为的是使祂在我们仇敌中向我们显示祂的仁慈。\Quote{天主在我的仇敌中向我显示。} 祂向野橄榄枝显示仁慈，这野橄榄枝被接在因骄傲而被折下的枝条上。看，那些骄傲的枝条倒在那里；看，你这曾说谎的，被接在哪里；不可骄傲，免得你也被折下。\par

19. \BibleQuote{求祢以祢的德能把他们分散在外} \BibleRef{咏 58:11}。如今这事已经成就：犹太人被分散到万民中，成为他们自己的不义和我们真理的见证人。他们有自己的著作，其中预言了基督，而我们持守基督。倘若有时偶然有外教人怀疑，当我们告诉他基督的预言时，他因预言的清晰而惊讶，惊奇地以为这些是我们自己写的，那么我们便从犹太人的抄本中证明，这事在很久以前是如何被预言的。看，我们如何借着我们的仇敌来驳倒其他仇敌。\Quote{求祢以祢的德能把他们分散在外：} 从他们那里取走\Quote{德能}，取走他们的力量。\Quote{我的保护者，上主，求祢使他们降下。} \BibleQuote{他们口中的过犯，他们嘴唇的言语：愿他们因骄傲而被捕获；因咒骂和谎言，终局将被宣告，在终局的愤怒中，他们将不复存在。} \BibleRef{咏 58:12} 这些是晦涩的话，我担心它们未能很好地渗透……\par

\Emph{第二部分。}\par

1. 因为，看，犹太人是仇敌，这篇圣咏似乎暗示他们；他们持有天主的法律，因此论到他们说过：\Quote{不要杀灭他们，恐怕他们忘记祢的法律：} 为的是使犹太民族存留，并借着他们的存留，基督徒的人数得以增加。他们确实存留在万民中，他们是犹太人，并未停止成为他们原先的样子：就是说，这民族并未向罗马制度屈服到失去犹太形态的地步；而是被罗马人征服，却仍保留自己的法律，即天主的法律。但在他们身上做了什么？\BibleQuote{你们捐纳薄荷、茴香，却遗弃了法律上更重要的事：正义、怜悯和审判；滤出蠓虫，却吞下了骆驼。} \BibleRef{玛 23:23} 这是主对他们说的。事实上他们正是如此：他们持有法律，持有先知书；阅读一切，歌唱一切：却看不见其中先知的光，即耶稣基督。他们不仅现在看不见坐在天上的祂，就是当时祂谦卑地行走在他们中间时，他们也没有看见祂，并且因流这同一位的血而成为有罪；但并非所有人。今天我们也将此事提请你们的爱德注意。并非所有人：因为他们中有许多人归向了他们所杀害的那位，并因信祂而获得了甚至流祂血的赦免：他们为人类树立了榜样；人不应绝望，以为任何罪都不会被赦免，因为连杀害基督的罪，在他们告明时也被赦免了……\par

2. 祢要击杀他们里面的什么？乃是他们所喊的\Quote{钉死，钉死}，而非喊叫的人。因为他们想要涂抹、铲除、消灭基督；但祢却藉着使基督（他们想要消灭的那位）复活，击杀了\Quote{他们口中的过犯，他们嘴唇的言论}。因为那位他们喊着应被消灭的，如今活着，他们就惊骇；那位他们在地上藐视的，如今在天上被万民朝拜，他们就诧异：这样，他们的过犯和嘴唇的言论就被击杀了。所谓\Quote{让他们在自己的骄傲中被擒获}是什么意思？因为壮士徒然冲来，在他们身上仿佛发生了这样的事：他们自以为做了什么，并且得胜了主。他们能够钉死一个人，软弱似乎得胜，德行被杀害；他们自以为是什么，如同壮士，如同强人，如同得胜者，如同预备猎物的狮子，如同肥壮的公牛，如同他在另一处提及的：\Quote{肥壮的公牛围困了我。}但在基督的事上他们做了什么？他们杀害的不是生命，而是死亡……如今那些皈依的人身上发生了什么呢？因为有人告诉他们，他们所杀的那位复活了。他们相信祂复活了，因为他们看见祂在天上，从那里派遣了圣神，充满了那些信了祂的人；他们发现自己什么都没有定罪，什么都没有做成。他们所做的归于虚空，罪却留存。因此，既然所为落空，但罪仍留在行凶者身上；他们就在骄傲中被擒获，他们看见自己处于不义之下。所以他们只剩下认罪，而那把自己交给罪人的主，施予宽恕，并宽恕祂的死亡——祂曾被死人杀害，却使死人活过来。所以他们就在骄傲中被擒获了。\par

3. \Quote{并且从咒骂和谎言中，将宣告圆满，在圆满的怒中，他们将要归于无有。}这也难以理解，不知\Quote{他们将要归于无有}与什么相连。什么将要归于无有？因此让我们考察上文：当他们已在骄傲中被擒获时，\Quote{从咒骂和谎言中将宣告圆满。}什么是圆满？就是成全：因为成为圆满，就是被成全。成为圆满是一回事，被消灭是另一回事。一件事被圆满，是指它被完成以致被成全；一件事被消灭，是指它被完成以致不复存在。骄傲不容许人被成全，没有什么比骄傲更阻碍成全。所以，愿你们的爱德稍加留意我所说的；并见一种极度有害、极需提防的恶。你们以为这是怎样的恶？我能花多少时间细说骄傲中有多少恶？魔鬼单单因此就应受惩罚。诚然他是所有罪人的首领；诚然他是引诱人犯罪的诱惑者；没有人将奸淫、酗酒、淫乱、抢夺他人财物归咎于他；他仅仅因骄傲而堕落。既然骄傲的同伴是嫉妒，那么骄傲的人必定嫉妒……总之，一切罪恶行为中的恶习都可怕，但行善中的骄傲更可怕。因此，宗徒如此谦卑地说：\BibleQuote{当我软弱时，正是我刚强的时候。}\BibleRef{格后 12:10}因为他自己恐怕被此罪诱惑，那位知道自己在医治什么的医生，给他用了什么药来防止肿胀？他说：\BibleQuote{因了启示的伟大，免得我高举自己，就给了我一根刺入我肉身的刺，撒殚的使者来打击我；为此我三次求主，使它离开我；祂对我说：\InnerQuote{我的恩宠为你足够了，因为德能在软弱中才成全。}}\BibleRef{格后 12:7}看，这就是圆满。一位宗徒，外邦人的导师，藉着福音成为信徒之父，竟接受了肉身的刺而受打击。我们中谁敢说这话，除非他不以承认这事为耻？因为若我们说\Person{保禄}没有遭受这事，我们表面上给他尊荣，实际上却使他成为说谎者。但因为他诚实，且说了真话；我们理当相信，有一个撒殚的使者被赐给了他，免得他因启示的伟大而高举自己。看，骄傲的蛇多么可怕……\par

4. 所谓\Quote{在圆满的怒中将宣告圆满}是什么意思？有一种圆满的怒，也有一种消灭的怒。因为天主的每一种报复都称为怒：有时天主报复，为使人成全；有时祂报复，为使人受罚。祂如何报复使人成全？\BibleQuote{祂鞭打所收纳的每一个儿子。}\BibleRef{希 12:6}祂如何报复使人受罚？当祂把不虔敬的人放在左边，对他们说：\BibleQuote{你们进入那为魔鬼和他的使者预备的永火里去吧。}\BibleRef{玛 25:41}这是消灭的怒，而非圆满的怒。但是\Quote{在圆满的怒中将宣告圆满}；将由宗徒们宣讲：\BibleQuote{罪恶在哪里增多，恩宠在哪里也格外增多}\BibleRef{罗 5:20}，而人的软弱属于谦卑的治疗。那些人思想这事，发现并承认自己的不义，\Quote{将要归于无有。}\Quote{归于无有}指什么？指他们的骄傲。\par

5. \BibleQuote{他们要知道：天主怎样统御雅各伯，并统御地极}（\BibleRef{咏 58:13}）。因为他们先前自以为是义人，因为犹太民族领受了法律，遵守了天主的诫命；但证明给他们看，他们并没有遵守，因为在天主的诫命本身中，他们没有认出基督，因为\BibleQuote{部分盲目落在以色列身上}（\BibleRef{罗 11:25}）。连犹太人自己也看出，他们不该轻视外邦人，因为外邦人曾被他们视为狗和罪人。正如他们一样在过犯中被找到，同样他们也将获得救恩。\BibleQuote{不但给犹太人，}宗徒说，\BibleQuote{也给外邦人}（\BibleRef{罗 2:10}）。因为为此，匠人所弃的石头，竟成了屋角的基石，为使两者在其中结合：因为墙角能结合两面墙。犹太人自以为崇高伟大；他们视外邦人为软弱的、罪人、魔鬼的仆役、偶像崇拜者，然而两者都有过犯。连犹太人也已被证明是罪人；因为\Quote{没有行善的，连一个也没有}：他们放下了骄傲，没有嫉妒外邦人的得救，因为他们知道自己的软弱和他们的软弱是一样的；并在角石中结合，一起朝拜了主。……\par

6. \BibleQuote{他们将在傍晚归化}（\BibleRef{咏 58:14}）：那就是，即使迟了，即在我们的主耶稣基督被杀害之后：\BibleQuote{他们将在傍晚归化；以后他们必像狗一样饥饿}。但\Quote{像狗一样}，不像羊或牛犊：\Quote{像狗一样}，像外邦人，像罪人；因为他们也知道了自己的罪，他们曾自以为义……因此，罪人谦卑是好的；没有人比自以为健康的人更不可救药。\Quote{他们必环绕城市}。我们已经解释了\Quote{城市}；它是\Quote{环绕的城市}；属于万民。\par

7. \BibleQuote{他们必分散四方，为要觅食}（\BibleRef{咏 58:15}）；那就是，为要赢得他人，为要将信徒改变融入他们的身体。\Quote{但他们若不饱足，就必发怨言}。因为上面他也提到了他们的怨言，说：\Quote{谁曾听见？} \Quote{但你，上主，}他说，\Quote{必嗤笑他们，说：谁曾听见？}为什么？因为，你必轻视万民如同虚无。让这篇圣咏结束吧。看那角石欢跃，现在两面墙都欢喜（\BibleRef{弗 2:20}）。犹太人曾骄傲，却谦卑了；外邦人曾绝望，却被高举了：让他们来到角石，在那里相遇，在那里相聚，在那里找到平安之吻；从不同方向而来，但不要带着不同而来，那些受割礼的，这些未受割礼的。墙相隔很远，但来到角石之前；但在角石中，让他们持守自己，现在整个教会从两面墙说什么呢？\BibleQuote{但我要歌颂你的大能，我要在清晨颂扬你的仁慈}（\BibleRef{咏 58:16}）。在清晨，当诱惑被战胜；在清晨，当这世界的黑夜过去；在清晨，当我们不再惧怕强盗和魔鬼及其天使的埋伏；在清晨，当我们不再靠预言的光行走，而是默观天主圣言本身如同太阳。\Quote{我要在清晨颂扬你的仁慈}。因此另一篇圣咏说：\Quote{清晨我侍立在你面前，并默想}。同样，主自己复活在黎明，应验了另一篇圣咏所说的：\Quote{傍晚哭泣，清晨欢呼}。\Quote{因为你成了我的扶持者，在我患难之日的避难所。}\par

8. \BibleQuote{我的助佑者，我要向你弹奏，因为天主，你是我的扶持者}（\BibleRef{咏 58:17}）。我本来是什么？除非你救助。我有多么绝望？除非你医治。我躺在何处？除非你到我这里来。确实，我带着巨大的伤口身处危险，但那伤口需要一位全能的医生。对全能医生来说，没有不可医治的。……最后，思量所有可能的好事，无论是本性上的，还是意向中的，或者在归化本身、信德、望德、爱德、好行为、正义、敬畏天主中；所有这些都唯独来自他的恩赐，他这样总结：\Quote{我的天主是我的仁慈}：他充满了天主的恩赐，找不到什么可称他的天主为，除了\Quote{他的仁慈}。这名字，无人应在它面前绝望！如果你说，我的救恩，我看出他赐予救恩；如果你说，我的避难所，我看出你在他内避难；如果你说，我的力量，我看出他赐你力量：\Quote{我的仁慈}是什么意思？凡我所是，都是出于你的仁慈。……\par

\SourceRecord{Translated by J.E. Tweed.；Nicene and Post-Nicene Fathers, First Series；Vol. 8.；Edited by Philip Schaff.；Buffalo, NY: Christian Literature Publishing Co.,；1888.；<http://www.newadvent.org/fathers/1801059.htm>.}

% structure: p0001 source=h1 target=section reason=page-title

\section{《圣咏第六十篇释义》}

1. 达味王是一个人，但他所表征的并不止一个人：有时他表征由许多人组成、扩展至地极的教会；但有时他表征那唯一的人，即天主与人之间的中保、基督耶稣这个人。\BibleRef{弟前 2:5}因此，在这篇圣咏中，或更确切地说，在这篇圣咏的标题中，谈及了达味的某些胜利事迹：\Quote{交与乐官，调用\InnerQuote{为所更改的百合花}，作达味自己教诲所用，当他在叙利亚的美索不达米亚和叙利亚索巴尔焚烧，并转向约阿布，击溃厄东，在盐谷中击毙了一万二千人。}我们在列王纪中读到这些事：他所提的那些人，即叙利亚的美索不达米亚、叙利亚索巴尔、约阿布、厄东，都被达味击败了。这些事确实发生了，正如所记载的，所读到的：谁愿读就可以读。然而，圣神在圣咏标题中通常稍为偏离所发生的事件的表述，并说一些历史上找不到的话，因而更提醒我们，这类标题之所以被写下来，不是要我们了解已发生的事，而是要预表未来的事。\newline{}但这里特别插入这件事，是因为在历史上并没有记载他在叙利亚的美索不达米亚和叙利亚索巴尔焚烧。现在，让我们开始按照未来事物的意义来审视这些事，并将阴影的晦暗带到言语的光明中来。\par

2. 你们知道何为\Quote{\OriginalTerm{以至于结局}{to the end}}。因为\Quote{法律的结局是基督}。你们知道那些被改变的。因为除了那些从旧生活过渡到新生活的，还有谁呢？……\BibleQuote{因为你们从前是黑暗，但现在在主内是光明}\BibleRef{弗 5:8} 但他们被改变\Quote{\OriginalTerm{标题的铭文}{title's inscription}}，……那些从魔鬼的国度过渡到基督的国度的人。他们被改变归入这标题的铭文，这是好的。但他们被改变，正如随后所说的，\Quote{归于教导}。他又加上，\Quote{给\Person{达味}自己，归于教导:} 也就是说，他们改变不是为了自己，而是为了\Person{达味}自己，并且改变归于教导。……那么，基督会在何时改变我们，除非祂做了祂所说的事：\BibleQuote{我来是把火投到世界上}\BibleRef{路 12:49} 因此，如果基督来是把火投到世界上，为了它的健康与利益，我们不应问祂如何把世界投入火中，而应问如何把火投入世界。既然祂来是把火投到世界，让我们查问被烧的\OriginalTerm{美索不达米亚}{Mesopotamia}是什么，\OriginalTerm{叙利亚索巴尔}{Syria Sobal}是什么？因此，让我们按照希伯来文（这部圣经最初就是用希伯来文写成的）来考察这些名字的解释。他们说美索不达米亚解释为\Quote{崇高的召唤}。如今整个世界因召唤而崇高，叙利亚解释为\Quote{高峻}。但那曾是高峻的，已被烧毁而谦卑。索巴尔解释为\Quote{空虚的古老}。感谢基督烧毁了她。每当老灌木被烧毁，青翠之地便取而代之；并且当火先行烧毁老灌木时，新芽会更快、更丰盛、更完全地发出。因此，不要害怕基督的火，它只烧毁干草。\BibleQuote{因为凡有血肉的都是草，人的一切荣华如同草上的花}\BibleRef{依 40:6} 因此祂用那火焚烧那些东西。\Quote{\OriginalTerm{约阿布}{Joab}就转身}。约阿布解释为敌人。有一个敌人被转变，正如你们将理解的。如果被转为逃跑，那是魔鬼；如果皈依信仰，那是基督徒。如何转为逃跑？从基督徒的心中：\Quote{这世界的元首，祂说，已被赶出去}\BibleRef{若 12:31} 但一个转向主的基督徒如何能成为一个转变的敌人？因为他曾经是敌人，现在成了信徒。\Quote{他击打了\OriginalTerm{厄东}{Edom}}。厄东解释为\Quote{属地的}。那属地的应当被击打。因为一个人既然应当活在天上，为什么还要活在地上？因此，属地的生命已被杀死，让属天的生命活着。\BibleQuote{就如我们曾带有了那属土者的肖像，也要带有那属于天上者的肖像}\BibleRef{格前 15:49} 看，它被杀死了：\BibleQuote{你们要致死你们在地上的肢体}\BibleRef{哥 3:5} 但他击打厄东时，他击杀了\Quote{一万二千人在\OriginalTerm{盐坑之谷}{valley of salt-pits}}。一万二千是一个完全的数，十二位宗徒的数目也被归为这个完全的数；因为这不是无缘无故的，而是因为圣言要传遍整个世界。但天主的圣言\BibleRef{则 37:9}，就是基督，存在于云中，即在真理的宣讲者中。但世界由四部分组成。世界的这四部分对所有人都非常熟悉，在圣经中常被提及：它们与四风之名相同：东、西、北、南。圣言被传到了这四个部分，以便所有人都能因天主圣三而被召叫。十二这个数是由四乘以三得来的。因此，有理由说一万二千属地的事物被击杀，整个世界被击杀：因为教会是从全世界被拣选出来的，从属地的生命中被治死。为什么\Quote{在盐坑之谷}？谷是谦卑：盐坑象征滋味。因为许多人谦卑了，但却是虚空而愚昧地，在空虚的古老中谦卑。一个人为金钱受苦，为暂时的尊荣受苦，为今生的安慰受苦；他要受苦并被谦卑：为什么不为了天主？为什么不为了基督？为什么不为了盐的滋味？难道你们不知道，曾对你们说过：\BibleQuote{你们是地上的盐}，并且\BibleQuote{盐若失了味，便毫无用处，只好被丢在外面}\BibleRef{玛 5:13} 因此，明智地谦卑是件好事。看，现在异端者不正在被谦卑吗？难道不是连人间的法律也制定来定他们的罪，而神圣的法律（甚至先前已定罪了他们）统治着？看，他们被谦卑，看，他们被赶走，看，他们受迫害，但没有滋味；因为愚昧，因为虚空。因为盐已经失了味：因此被丢在外面，被人践踏。我们已经听了这圣咏的标题，现在也让我们听圣咏的话语。\par

3. \Quote{天主，祢驱退了我们，祢毁了我们} \BibleRef{咏 59:1}。那是\Person{达味}在说话吗？就是那击打、焚烧、击败的人，而不是祂对这些所做的事所临到的人，也就是说，他们被击打驱退，他们是恶人，而他们又得以复活、回归，为使他们成为善人？那\Person{达味}所做的毁灭，他手强有力，我们的基督，那人所预表的，祂行了那些事，祂用祂的剑和祂的火做了这毁灭：因为祂把两者都带到这世界来了。正如福音中你们读到：\Quote{我来是把火带到世界上} \BibleRef{路 12:49}，和\Quote{我来是把剑带到地上} \BibleRef{玛 10:34}。祂带来了火，用以烧毁\Place{叙利亚}的\Place{美索不达米亚}和\Place{叙利亚}的\Place{索巴尔}；祂带来了剑，用以击打\Place{厄东}。这毁灭又是为那些\Quote{改变归于标题题词}的人而做的。因此我们当听他们的声音：他们为得健康而被击打，让他们被举起来说话。因此，那些改变成更好状态、改变归于标题题词、改变归于为\Person{达味}本人教训的人，让他们说：\Quote{祢怜悯了我们。}祢毁灭了我们，为的是建造我们；祢毁灭了我们这些建造不善的，祢毁掉了虚空的旧性，为的是有一建筑归于新人，建筑存留到永远……。\par

4. \Quote{祢震动了大地，使地震动} \BibleRef{咏 59:2}。大地如何被震动？在罪人的良心中。我们往哪里去？我们往哪里逃，当这剑被挥舞，\Quote{悔改吧，因为天国近了}？\BibleRef{玛 3:2} \Quote{医治它的破碎，因为它已震动。}若不被震动，就不值得被医治；但是你说话、讲道、以天主恐吓我们，对将来的审判闭口不言，从天主的诫命你警告，从这些事你不停止；而那听见的人，若不惧怕、不震动，就不配被医治。另一个人听见了，被震动、被刺痛、捶胸、流泪……。\par

5. 第一项劳苦是，你应使自己不悦，应奋战出罪恶，应改变成更好的；第二项劳苦，为回报你的改变，是承受这世界的磨难和试探，并在其中坚持到底。因此当他谈论这些事，指出这些事情时，他加上了什么？\Quote{祢向祢的百姓显现了艰难的事} \BibleRef{咏 59:3}：向祢的百姓，如今在\Person{达味}胜利后成了进贡的。\Quote{祢向祢的百姓显现了艰难的事。}在哪里？在基督的教会所遭受的迫害中，当时流了这么多殉道者的血。\Quote{祢给我们喝了让人刺痛之酒。}\Quote{让人刺痛}是什么意思？不是杀戮。因为那不是毁灭性的杀戮，而是刺痛的药。\Quote{祢给我们喝了让人刺痛之酒。}\par

6. 为什么这样？\Quote{祢给了敬畏祢的人一个记号，叫他们逃避弓的面} \BibleRef{咏 59:4}。经由暂时的磨难，祂说，祢向祢自己的人预示要逃避永火之怒。因为宗徒\Person{伯多禄}说：\Quote{时候到了，审判要从天主的家开始。} \BibleRef{伯前 4:17} \Person{伯多禄}劝勉殉道者忍耐，当世界狂怒，当迫害者之手制造杀戮，当信徒的血远近流洒，当信徒在锁链、监牢、酷刑中遭受许多艰难之事时，在这些艰难中，我说，恐怕他们丧胆，\Person{伯多禄}对他们说：\Quote{时候到了，审判要从天主的家开始。}等等。那么审判中有什么呢？弓已弯曲，仍处于威胁姿态，尚未瞄准。看那弓中有什么：不是有箭要向前射吗？然而弦被向后拉，方向与将要射出的方向相反；弦拉得越向后，它向前弹出的速度就越快。我所说的是什么？审判延迟得越久，它就要以越大的速度来临。因此，我们要为暂时的磨难向天主感恩，因为祂给了祂的百姓一个记号，\Quote{叫他们逃避弓的面}：为的是祂的忠信者既在暂时的磨难中受锻炼，便堪以避免永火的惩罚，那火将找出一切不信这些事的人。\par

7. \Quote{为使祢所爱的人得救：求祢以右手拯救我，垂听我} \BibleRef{咏 59:5}。主，以祢的右手拯救我：这样拯救我，使我得以站在右边。我不求任何暂时的安全，在这事上愿祢的旨意成就。因为在某个时刻什么对我们有益，我们完全不知道：因为\Quote{我们不知道该如何祈祷} \BibleRef{罗 8:26}，但是\Quote{求祢以右手拯救我}，这样即使在这时候我遭受各种磨难，当所有磨难的夜晚过去，我可以在右边被找到，在绵羊中，而不是在左边，在山羊中。\BibleRef{玛 25:33} \Quote{且垂听我。}因为现在我配得上祢所愿意给予的；不是\Quote{以我过犯的话}我终日呼号，以致祢不垂听，也不是\Quote{在夜间以致祢不垂听}，而且这不是我的愚昧，实为我的警告，借着从盐坑谷加添味道，使我在磨难中知道该求什么；但我寻求永生，因此求祢垂听我，因为我求祢的右手……。\par

8. \BibleQuote{天主在自己的圣者中发言了}\BibleRef{咏 59:6}。在祂的哪一位圣者中？\BibleQuote{天主在基督内使世界与自己和好}\BibleRef{格后 5:19}。在那一位圣者中，你们在别处听到过：\Quote{天主，在圣者中就是你的道路。}\Quote{我要欢乐，要划分\Place{舍根}……并要测量\Place{史科特谷}。}\OriginalTerm{舍根}{Sichima}被解释为肩膀。但是根据历史，\Person{雅各伯}从岳父\Person{拉班}那里带着所有亲属回来时，把从\Place{叙利亚}带来的偶像藏在\Place{舍根}\BibleRef{创 35:4}，他在那里住了很久，最后从那里出来。但是他在那里为羊群和牲畜搭了帐棚，并称那地方为\OriginalTerm{史科特}{Tabernacles}。这些我要划分，教会说。这是什么？\Quote{我要划分\Place{舍根}}？如果指的是偶像被藏起来的故事，那它象征\OriginalTerm{外邦人}{Gentiles}；我划分外邦人。我划分是什么意思？\BibleQuote{因为并非人人都有信德。}\BibleRef{得后 3:2}我划分是什么意思？有些人会相信，有些人不会相信……肩膀被划分，好让某些人的罪使他们负重，而另一些人则背起基督的担子。因为祂要求虔敬的肩膀，当祂说：\BibleQuote{因为我的轭是柔和的，我的担子是轻省的。}\BibleRef{玛 11:30}另一种担子压迫并加给你重负，但基督的担子减轻你的负担：另一种担子有重量，基督的担子有翅膀。因为即使你从鸟身上拔掉翅膀，你除掉了一种重量；你拿走的重量越多，它就越停留在地上。你选择减轻其负担的那只鸟躺在地上：它不飞，因为你拿掉了重量：把重量还给它，它就飞了。基督的担子就是这样；让人背起它，不要闲懒：不理会那些不愿意背的人；让愿意的人背起它，他们就会发现它是多么轻省、多么甘饴、多么愉快、多么引人升天，并且多么使人脱离尘世……也许因为雅各伯的羊群，\Quote{\Place{史科特谷}}应被理解为犹太民族，同样被划分：因为相信的人从那里出来了，其余的人留在外面。\par

9. \BibleQuote{基肋阿得是我的}\BibleRef{咏 59:7}。这些名字在天主的圣经中读到。\OriginalTerm{基肋阿得}{Galaad}有自己的解释，包含一个伟大的奥秘：因为它被解释为\Quote{见证的堆集}。在殉道者中有多大的见证堆集啊！\Quote{基肋阿得是我的}，我的就是见证的堆集，我的是真正的殉道者……那时，教会在人中卑视，那时她被当作寡妇受到羞辱，因为她是属于基督的，因为她在额上佩戴着十字架的标记：那时还没有荣耀，只有指责；因此，当没有荣耀而只有指责时，就有了见证的堆集；通过见证的堆集，基督的爱被扩大；通过基督的爱的扩大，外邦人被占领。接着是\Quote{默纳协也是我的}；\OriginalTerm{默纳协}{Manasses}被解释为\OriginalTerm{被遗忘}{forgotten}。因为曾对她说：\BibleQuote{你必永远忘记羞耻，不再记念你寡居的羞辱。}\BibleRef{依 54:4}教会曾有一次羞耻，如今已被忘记：因为她不再记念她的羞耻和\Quote{寡居的耻辱}。因为当人们在羞耻时，就做成了见证的堆集。如今没有人再记得那羞耻，那时作基督徒是耻辱，现在无人记得，全都忘记了，现在\Quote{默纳协是我的，厄弗辣因是我头的力量。}\OriginalTerm{厄弗辣因}{Ephraim}被解释为\OriginalTerm{丰饶}{fruitfulness}。他说，丰饶是我的，这丰饶是我头的力量。因为我的头是基督。丰饶怎会是祂的力量？因为除非一粒麦子落在地里，否则它不会增多，独自存留。\BibleRef{若 12:24}基督在苦难中落在地上，复活时就结出了果实。他被悬挂，被人轻视：那粒麦子在里面，有能力吸引一切到它那里。一粒麦子中隐藏着多少种子，在眼中看似卑贱，但隐藏着一种将物质转化并结出果实的能力；同样，在基督的十字架上，隐藏着德能，显现出软弱。强大的麦粒啊！无疑，悬挂的是软弱的，无疑，民众在祂面前摇头，无疑，他们说：\BibleQuote{如果祂是天主子，让祂从十字架上下来吧。}\BibleRef{玛 27:40}请听祂的力量：天主的软弱比人更强壮。\BibleRef{格前 1:25}因此，如此大的丰收随之而来：它是我的，教会说。\par

10. \BibleQuote{犹大是我的君王；摩阿布是我希望的锅} \BibleRef{咏 59:7}。是哪个犹大？就是出自犹大支派的那一位。是哪个犹大？不就是雅各伯本人所说的那一位：\BibleQuote{犹大，你的兄弟要赞美你}（\BibleRef{创 49:8}）？那么，当我的君王犹大说：\BibleQuote{你们不要害怕那杀害身体的}（\BibleRef{玛 10:28}）时，我还怕什么呢？\BibleQuote{摩阿布是我希望的锅。}为何说\Quote{锅}？因为磨难。为何说\Quote{我希望的}？因为我的君王犹大走在了前面。摩阿布被看作是外邦人。因为那民族生于罪恶（\BibleRef{创 19:37}），那民族生于罗特的女儿们，她们在父亲醉酒时与他同寝，滥用了父亲。与其成为母亲，倒不如保持不生育更好。但这是一种形象，代表着那些滥用法律的人。因为不要在意拉丁语中法律一词是阴性，希腊语中是阳性；但无论是阴性还是阳性，其表达并不影响真理。因为法律具有阳刚之力，因为它统治，而不被统治。而且，宗徒保禄说了什么？\BibleQuote{法律原是好的，只要人用得合法}（\BibleRef{弟前 1:8}）。但罗特的女儿们不合法地用他们的父亲。同样，当人善用法律时，善行开始生长；当人恶用法律时，恶行便产生。此外，他们滥用父亲，即滥用法律，生出了摩阿布人，由此象征恶行。于是教会的磨难，锅在沸腾。关于这锅，在某处预言中说：\BibleQuote{一口被北风吹热的锅}（\BibleRef{耶 1:13}）。从哪里吹来？除了从魔鬼的地域，他曾说：\BibleQuote{我要在北方安设我的宝座}（\BibleRef{依 14:13}）。因此，最大的磨难临到教会，无非来自那些滥用法律的人。\par

11. \BibleQuote{我要向厄东投掷我的鞋} \BibleRef{咏 59:8}。教会说：\Quote{我要直抵厄东。}让磨难发怒，让世界因罪孽而沸腾，甚至及于那些过着世俗生活的人（因为厄东被解释为世俗），甚至及于那些人，\BibleQuote{我要向厄东投掷我的鞋。}这鞋是指什么？除了\WorkTitle{福音}。\BibleQuote{那传布和平、传布福音者的脚是多么美丽}（\BibleRef{罗 10:15}），并且\BibleQuote{用和平福音的准备穿在脚上}（\BibleRef{弗 6:15}）。弟兄们，在这些时代，我们看到多少世俗的人为了利益而行欺诈，为了欺诈而行伪誓，因恐惧而求问算命者和占星者；所有这些人都属于厄东族，是世俗的；然而所有这些人都朝拜基督，都在祂的鞋下；现在祂的鞋甚至伸展到了厄东。\BibleQuote{外方人都向我投降了。}\Quote{外方人}是谁？其他种族的人，不属于我的种族。他们\Quote{投降了}，因为许多人朝拜基督，却不能与基督一同为王。\par

12. \BibleQuote{谁领我进入设防城？} \BibleRef{咏 59:9}。设防城是什么？如果你们记得，我已在另一篇\WorkTitle{圣咏}中提及，其中说：\BibleQuote{他们必环绕这城。}因为设防城就是外邦人的包围圈；这外邦人的包围圈中间曾有唯一的犹太民族，崇拜唯一的天主；其余的外邦人包围圈向偶像祈祷，侍奉魔鬼。这城神秘地被称为设防城，因为外邦人从四面八方涌入，并围绕着那个崇拜唯一天主的民族。\BibleQuote{谁领我直抵厄东？}\par

13. \BibleQuote{天主，岂非祢将我们击退？天主，祢岂不率领我们的军队出征？} \BibleRef{咏 59:10}。难道不是祢领我们下去，祢却将我们击退？为何说\Quote{将我们击退}？因为祢消灭了我们。为何消灭了我们？因为祢曾发怒，又怜悯了我们。因此，是祢——那将我们击退的——领我们下去；是祢，天主，那不率领我们的军队出征的，领我们下去。\Quote{不率领我们的军队出征}是什么意思？世界要发怒，世界要践踏我们，将有一批证人，由殉道者所流的血建成，而愤怒的异教徒要说：\Quote{他们的天主在哪里？}那时\Quote{祢不率领我们的军队出征}：祢不向他们显示自己，不显示祢的大能，如同祢曾在达味、梅瑟、农的儿子若苏厄身上显示的，那时外邦人在他们的威力下屈服，在屠杀结束和大破坏修复后，祢领着祢的子民进入祢所应许的土地。因此这事祢不会做，\Quote{祢不率领我们的军队出征}，而是在内里工作。\Quote{不率领出征}是什么意思？不显示自己。因为当殉道者们被锁链牵着，被关进监狱，被带出去受凌辱，被抛给野兽，被刀剑砍杀，被火烧死时，他们岂不被人看作是遗弃的、无援助的吗？天主在内里怎样工作？祂在内里怎样安慰？祂怎样使他们盼望永生的甘甜？祂岂没有舍弃他们的心，其中人静静地居住，善则善，恶则恶？难道因为祂没有率领他们的军队出征，祂就舍弃了吗？正因为祂没有率领他们的军队出征，祂岂不更将教会领下到厄东，领下到设防城吗？因为如果教会选择争战并使用刀剑，她似乎是为现世生命而战；但她轻视现世生命，因此为那将有的生命建立了证人的堆。\par

14. 因此，祢，天主，既不愿意在我们的能力中前进，\BibleQuote{求祢从苦难中赐给我们援助，因为人的救恩是枉然的}\BibleRef{咏 59:11}。现在，那些没有盐的人去吧，为他们的朋友渴求暂时的安全，那是虚空的旧物。\BibleQuote{求祢赐给我们援助：}从祢似乎要舍弃的地方，从那里施以援手。\BibleQuote{在天主内我们要施展勇力，祂必将我们的仇敌归于无有}\BibleRef{咏 59:12}。我们不用刀剑，不用马匹，不用胸甲，不用盾牌，不用军队的威力，不在外作战。但在哪里？在内里，在我们不被看见的地方。哪里内里？\BibleQuote{在天主内我们要施展勇力}，并且好像卑贱的人，好像被践踏的人，好像无足轻重的人，但我们将是\BibleQuote{祂要将我们的仇敌归于无有}。总之，这事已经发生在我们仇敌身上。殉道者曾被践踏：借着受苦、忍耐、坚持到底，他们在天主内施展了勇力。祂自己也做了随后的事：祂将他们的仇敌归于无有。如今殉道者的仇敌在哪里？除非现在醉酒之人用他们的杯子迫害那些曾经被疯狂之人用石头迫害的人？\par

\SourceRecord{Translated by J.E. Tweed.；Nicene and Post-Nicene Fathers, First Series；Vol. 8.；Edited by Philip Schaff.；Buffalo, NY: Christian Literature Publishing Co.,；1888.；<http://www.newadvent.org/fathers/1801060.htm>.}

% structure: p0001 source=h1 target=section reason=page-title

\section{圣咏第六十一篇释义}

标题本身并不让我们耽搁。因为它原是\Quote{至于终结，在圣歌中，为\Person{达味}本人}。\Quote{在圣歌中}，即是在赞美中。\Quote{至于终结}，即是至于基督……但是在这篇圣咏中的声音（如果我们是在祂的肢体中，在身体内，正如因祂的劝勉我们敢信赖一样）我们应当承认是我们自己的，而非任何外人的。但我称它为我们自己的，并非好像只是那些现在在场的；而是为我们自己的，作为遍及普世的我们，从东方直到西方。为使你们知晓它因此是我们的声音，祂在这里说话如同一个人；但祂并非一个人；而是如同一个，合一在说话。但在基督内我们众人是一个身体：因为这个人，头在天上，肢体却仍在世上劳苦：因为他们劳苦，请看祂说什么。\par

2.\BibleQuote{天主，请俯听我的祈求，垂顾我的祈祷}\BibleRef{咏 60:1}。谁在说？祂，如同一个人。看是否是同一个：\BibleQuote{我从地极向你呼号，当我心中烦乱时}\BibleRef{咏 60:2}。所以现在并非一个：而是因此是一个，因为基督是一位，我们都是祂的肢体。因为哪一个人从地极呼喊呢？除了那继承产业，即曾对圣子本身说过：\BibleQuote{你向我请求，我就将万民给你作产业，将地极给你作领土。}所以这是基督的产业，这是基督的继承，这是基督的身体，这是基督的唯一教会，这是我们所是的合一，正从地极呼喊……但我为何呼号这事呢？\BibleQuote{当我心中烦乱时}。祂显示自己遍及万国，在整个圆球世界，在极大的光荣中，但也在极大的苦难中。因为我们在这种旅居中的生命不能没有诱惑：因为我们的进步是通过诱惑而实现的，人除非受诱惑，就不会认识自己，除非他得胜，就不能获得冠冕，除非他奋斗，就不能得胜，除非他经验仇敌和诱惑，就不能奋斗。因此这个人正在受苦，那从地极呼喊的，但祂并未被弃。因为祂愿意预表我们这些祂身体的成员，也在祂那已经死了又复活了，升了天的身体中，好使头所先到的地方，肢体可以确信他们将会跟随。因此当祂愿意受撒殚试探时，祂就用比喻将我们转移到祂自己身上。\par

3. 然而刚才在福音中读到，主耶稣基督在旷野中受魔鬼的试探。\BibleRef{玛 4:1} 基督完全受了魔鬼的试探。因为在基督内，你们被试探，因为基督从你们取了肉身，从祂自己给你们救恩；从你们取了死亡，从祂自己给你们生命；从你们取了辱骂，从祂自己给你们荣耀；因此从你们取了试探，从祂自己给你们胜利。如果我们在祂内被试探，我们就在祂内战胜魔鬼……\BibleQuote{你把我高举在磐石上。}所以现在在这里我们领悟了谁从地极呼喊。让我们回想福音：\BibleQuote{在这磐石上，我要建立我的教会}\BibleRef{玛 16:18}。所以她从地极呼喊，祂愿意她被建造在磐石上。但为使教会建在磐石上，谁成了磐石？听\Person{保禄}说：\BibleQuote{那磐石就是基督。}因此我们被建造在祂身上。为此，我们被建造在其上的那磐石\BibleRef{玛 7:24}首先被风、洪水、雨击打，当基督受魔鬼试探时。请看祂愿意将你们建立在何等稳固之上。因此我们的声音不是徒然的，而是被垂听：因为我们被置于伟大的希望中：\BibleQuote{你把我高举在磐石上}……\par

4.\BibleQuote{你引领我下来，因为你成了我的希望：来自敌人面前的坚固堡垒}\BibleRef{咏 60:3}。我的心烦乱，那从地极的合一说：我在诱惑和绊脚石中劳苦：外邦人嫉妒，因为他们被征服了；异端者潜伏，隐藏在基督徒名号的披风下：在教会内部，麦子遭受糠秕的欺凌：在这一切之中，当我的心烦乱时，我从地极呼喊。但不离弃我的是那将我高举在磐石上的那一位，为领我下来直到祂自己，因为即使我劳苦，当魔鬼通过如此多的地方、时机和境遇潜伏反对我时，祂为我是一座坚固堡垒，当我逃入其中避难时，我不仅逃脱敌人的武器，而且甚至我能安全地向敌人投掷我所愿的任何标枪。因为基督自己就是堡垒，祂自己为了我们从敌人面前成为堡垒，祂也是教会建于其上的磐石。你注意不要被魔鬼击打吗？逃入堡垒；魔鬼的标枪永远不会追上你到那堡垒：在那里你将受到保护，坚定不移。但你如何逃入堡垒呢？但愿一个人，或许在诱惑中，不要身体去寻找那堡垒，当找不到时，就疲倦或昏厥在诱惑中。你面前就是堡垒：记起基督，进入堡垒……\par

5. \Quote{我要在祢的帐幕内作客，直到永远} \BibleRef{咏 60:4}. 你看，我们所说的那一位，就是呼喊者。我们之中谁能作客直到永远？因为在这里我们只活几天，就过去了：我们在此是客，在天上才是居民。你在那里是客，在那里你将听到你主天主的声音说：\Quote{离开吧。}因为从那永久的家乡——天上，没有人会叫你离开。所以你在这里是客。因此另一圣咏也说：\Quote{我在祢面前是客，是旅人，如同我所有的祖先一样。}所以我们在这里是客；主将在那里赐给我们永久的居所：\Quote{在我父的家里，\InnerQuote{有许多住处}} \BibleRef{若 14:2}. 那些住处，祂不是像给客人那样赐予，而是像给公民一样，让他们永远居住。但在这里，弟兄们，因为教会在地上并非短暂停留，而是直到世界的终结都要存在，所以此处祂说：\Quote{我要在祢的帐幕内居住，直到永远。}……不错，你也许会希望诱惑只持续几天；但教会如何聚集她所有的儿子，除非她长久在此，直到终结？不要嫉妒以后将要来的其余人类：不要因为你已经渡过，就想砍断慈悲的桥梁：让它在这里直到永远。那么诱惑呢？随着罪恶增多，它们必然增多。因为祂自己说：\Quote{由于罪恶增多，许多人的爱德会冷淡。} \BibleRef{玛 24:12}. 但那从地极呼喊的教会，正是在祂接着所说的处境中：\Quote{唯独坚持到底的，才可得救。}但你如何坚持呢？……\Quote{我将在祢翅膀的荫蔽下被遮盖。}看，这就是何以我们在如此大的诱惑中仍然安全，直到世界终结，永世接纳我们；因为我们被遮盖在祂翅膀的荫蔽下。世上有炎热，但在天主的翅膀下有巨大的荫凉。\par

6. \Quote{因为天主，祢俯听了我的祈祷} \BibleRef{咏 60:5}. 什么祈祷？就是开头所说的：\Quote{天主，求祢俯听我的哀求。}……\Quote{祢将产业赐给了敬畏祢名的人。}因此让我们继续敬畏天主的圣名：永恒之父不会欺骗我们。儿子劳苦，为了继承父母的产业，父母死后他们继承；难道我们劳苦不就是为了从那位父那里承受产业吗？我们不是死后才继承祂，而是与祂一同在产业中永远生活。\par

7. \Quote{祢将日子叠日子加给君王的年岁} \BibleRef{咏 60:6}. 这位君王就是我们作为其肢体的那位。基督是君王，我们的头，我们的君王。祢将日子叠日子赐给祂；不仅是在那有限的时间中的日子，更是那些没有终结的日子。他说：\Quote{我要住在上主的殿中，直到悠长的时日。}为何说悠长的时日？因为现在是短暂的日子。凡有终结的都是短暂的；但这位君王的日子叠日子，所以不仅在这些日子过去时基督在祂的教会中为王，而且圣人们也要与祂一同在那些没有终结的日子里为王。……因为也曾论及天主的年岁：\Quote{祢却是一样的，祢的年岁没有终结。}如同年岁，也如同日子，也如同一天。关于永恒，你愿说什么都可以。因着这个理由，你愿意说什么都可以，因为无论你说什么，你所说的都太少了。因为你必须说点什么，以便有东西让你默想那不可言说的。\Quote{直到世世代代。}指这个世代和将来的世代：这个世代好比月亮，因为月亮由新而盈、由盈而亏、消失，这些有死的世代也是如此；还有那个世代，我们藉复活重生，与天主永远同在，那时我们不再像月亮，而像主所说的：\Quote{那时义人要在他们父的国里发光如同太阳。} \BibleRef{玛 13:43}. 因为在圣经中，月亮常象徵这有死状态的变幻无常。……\par

8. \BibleQuote{祂将在天主面前永远存留} \BibleRef{咏 60:7}；根据什么，或因何缘故？\Quote{祂的仁慈和真理，谁为祂寻求？}祂在另一处也说：\Quote{主的一切道路都是仁慈和真理，为寻求祂的盟约和祂的诫命的人们。}关于真理和仁慈的论述是广泛的，但我们承诺了简短。请简要听一听什么是真理和仁慈：因为所说的话不是小事：\Quote{主的一切道路都是仁慈和真理。}仁慈被提及，因为天主不看我们的功德，而看祂自己的良善，为能赦免我们所有的罪，并应许永生；但真理被提及，因为祂必不食言，必履行祂所应许的。让我们在此承认它，并实行它；这样，正如天主向我们显示了祂的仁慈和祂的真理，仁慈在于赦免我们的罪，真理在于彰显祂的应许；同样，我说，让我们也实行仁慈和真理：对软弱者、对贫困者、甚至对我们的仇敌施行仁慈；在真理上不犯罪，不罪上加罪……那么，谁做这事呢？除非是那些少数人中的一位，关于他们曾说：\Quote{那坚持到底的，必然得救。}因此这里也说：\Quote{祂的仁慈和真理，谁为祂寻求？}为什么是\Quote{为祂}？\Quote{谁寻求}就足够了。为什么祂加了\Quote{为祂}？只是因为许多人寻求在祂的书中学习祂的仁慈和真理；他们学到了，却为自己活，不为祂活；\BibleRef{格后 5:15}他们寻求自己的事，不寻求耶稣基督的事；\BibleRef{斐 2:21}他们宣讲仁慈和真理，却不实行仁慈和真理。但通过宣讲，他们认识它们；因为若不认识，就不能宣讲。但那爱天主和基督的人，在宣讲同一位的仁慈和真理时，他自己是为祂寻求之，不是为自己；也就是说，不是为了通过这宣讲获得暂时的好处，而是为了祂的肢体，即祂的忠信者们行善，以真理施行他所知道的；为叫那活着的人不再为自己活，而为主活，祂为众人死了。\BibleRef{格后 5:15}\par

9. \BibleQuote{因此我要向祢的名奏乐，日复一日还我的誓愿} \BibleRef{咏 60:8}。如果你向天主的名奏乐，就不要只奏一时。你要永远奏乐吗？你要永久奏乐吗？向祂日复一日还你的誓愿。什么是日复一日向祂还你的誓愿？从这一天直到那一天。继续在这日还誓愿，直到你到达那一天：那就是\BibleQuote{那坚持到底的，必然得救。}\BibleRef{玛 24:13}\par

\SourceRecord{Translated by J.E. Tweed.；Nicene and Post-Nicene Fathers, First Series；Vol. 8.；Edited by Philip Schaff.；Buffalo, NY: Christian Literature Publishing Co.,；1888.；<http://www.newadvent.org/fathers/1801061.htm>.}

% structure: p0001 source=h1 target=section reason=page-title

\section{论圣咏第六十二篇}

1. 它的标题是：\BibleQuote{直至终结，为耶杜通，达味自己的一篇圣咏。} 我记得已经向你们解释过耶杜通是什么……让我们看看他跳过了多远，他\Quote{跳过}了谁，以及他在什么地方——尽管他跳过了某些人，他却处于一种仿佛从某种属灵且安稳的位置往下看的境地……我说，他处于一个设防之地，就说：\BibleQuote{我的灵魂岂不归于天主？}\BibleRef{咏 61:1}。因为他曾听见：\BibleQuote{凡高举自己的，必被贬抑；凡谦卑自己的，必被高举。}\BibleRef{玛 23:12} 他害怕因跳过而骄傲，不为下面的事所自傲，却因那在上的而谦卑；于是，对那些仿佛以堕落威胁他、因他跳过而恼怒的嫉妒之人，他回答说：\BibleQuote{我的灵魂岂不归于天主？}……\BibleQuote{因为我的救恩来自祂。}\BibleQuote{因为祂是我的天主、我的救恩、我的扶持者，我必不再动摇。}\BibleRef{咏 61:2}。我知道谁在我之上，我知道谁向认识祂的人广施仁慈，我知道我应在谁的翅膀荫下盼望：\BibleQuote{我必不再动摇。}……\par

2. 因此，从那更高处、设防且受保护之处，那位以主为避难所、以天主为设防之地的人，注视着他所跳过的人，并如同从高塔上俯视着他们说话：因为关于祂也曾说过：\BibleQuote{是面对仇敌的力量之塔。} 因此他注意他们，并说：\BibleQuote{你们把重担加在人身上要到何时？}\BibleRef{咏 61:3}。你们以侮辱、投掷责骂、设伏、迫害，将重担加在人身上，这些是人所能承担的；但为了使人能承担，在他之下的是那创造人的天主。如果你们只看人，\BibleQuote{你们就杀吧，你们众人。} 看哪，加担子吧，发怒吧，\BibleQuote{你们众人就杀吧。}\BibleQuote{如同歪斜的墙，如同被冲撞的篱笆。} 要靠过去，冲撞过去，仿佛要推倒它。那\BibleQuote{我必不再动摇}在哪里？为什么呢？\BibleQuote{我必不再动摇。} 因为祂是我的天主、我的拯救者、我的扶持者，所以你们人能加担子给人；但你们能加担子给那保护人的天主吗？\BibleQuote{你们众人就杀吧。} 一个人有那么大的身躯，能使众人杀他吗？但我们应当意识到我们的位格，教会的位格，基督身体的位格。因为一个人——连同他的头和身体——就是耶稣基督，身体与肢体的救主：二人在同一肉体内，在同一声音里，在同一苦难中，当不义过去之后，又在同一安息中。因此，基督的苦难并不只在基督内；不，除了在基督内，就没有任何苦难。因为如果你理解基督为头和身体，那么基督的苦难就只在基督内；但如果你只理解头为基督，那么基督的苦难就不只在基督内。因为如果基督的苦难只在基督内，即只在头内，那么祂的一个肢体——宗徒保禄——怎么说：\BibleQuote{为补足基督苦难在我肉身上所欠缺的。}\BibleRef{哥 1:24}？因此，如果你是基督的肢体，无论你是哪一个人，听这些话的人，无论你是谁，只要你听（但如果你在基督的肢体中，你就听）：凡你从那些不属于基督肢体的人所受的一切，都曾是基督苦难所欠缺的。因此，因为欠缺，所以添加；你填补了度量，你使它不溢出：你所受的正是要从你的苦难中补充到基督整个苦难中的那一部分，这苦难已由我们的头承受，并在祂的肢体中——即我们自身——承受。对于这我们的共同国度，我们每人按自己的度量缴纳我们所欠的，并按我们所拥有的力量，仿佛贡献一份苦难的定额。所有人之苦难的仓库要等到世界终结时才能完全装满……因此，从义人亚伯尔的血直到则加黎雅的血，整个城都在说话。\BibleRef{玛 23:35} 从此以后，从若翰的血，经过宗徒的血，经过殉道者的血，经过基督信友的血，这一座城在说话，一个人说：\BibleQuote{你们把重担加在人身上要到何时？你们众人就杀吧。} 让我们看看你们是否能抹去，让我们看看你们是否能消灭，让我们看看你们是否能从地上除去祂的名，让我们看看你们外邦人是否不图谋虚妄，说：\BibleQuote{她何时死？她的名何时消灭？}\BibleQuote{如同她是一堵歪斜的墙，一面被冲撞的篱笆，} 你们要靠向她，冲撞她。从高处听：\BibleQuote{我的扶持者，我必不再动摇：} 因为我曾像一堆沙被冲撞而倒下，但主扶持了我。\par

3. \Quote{尽管如此，他们仍想击退我的荣耀} \BibleRef{咏 61:4}。他们杀害顺服的人，自己却被征服了，因为被杀者的鲜血增加了信徒；他们向这些人屈服，不能再杀害了；\Quote{尽管如此，他们仍想击退我的荣耀。}因为基督徒不能被杀害，人们便努力使基督徒蒙受羞辱。如今，不敬神的人的心因基督徒的荣耀而受折磨：如今，那属灵的若瑟，被他的兄弟出卖后，离开他的家进入埃及，仿佛进入外邦人中，经过监狱的羞辱，经过假见证的虚构故事，之后那关于他所预言的\Quote{铁剑刺透了他的灵魂}的事发生了：如今，他受到尊荣，不再服从于出卖他的兄弟，反而供应粮食给他们饥饿的人。\BibleRef{创 42:5}被他的谦卑和贞洁、廉洁、诱惑、苦难所征服，如今他们看见他受尊荣，却想阻止他的荣耀……这一切是针对一个人，还是一个人针对所有人；或是所有人针对所有人，或是一个人针对一个人？同时，当他说\Quote{你们加害于人}时，似乎是针对一个人；当他说\Quote{你们全都杀害}时，似乎是所有人针对一个人；但尽管如此，这也是所有人针对所有人，因为所有人都是基督徒，但在同一体内。但为什么那些敌基督的各种异端要一起被提及？它们也是一个吗？诚然，我敢说它们也是一个：因为有一座城和一座城，一个民族和一个民族，一位王和一位王。一座城和一座城是什么？一个巴比伦，一个耶路撒冷。无论她还被其他什么神秘的名字称呼，但有一座城和一座城；这个城由魔鬼为王，那个城由基督为王……\par

4. 请留心，弟兄们，请留心，我恳求你们。因为我仍乐于对你们说几句关于这座可爱的城的话。因为\Quote{关于你——天主的城——已经说了最荣耀的事。}并且，\Quote{若我忘记你，哦，耶路撒冷，愿我的右手忘记我。}因为只有一个故乡是宝贵的，真正只有一个故乡，唯一的故乡：除此之外我们所拥有的，都是在异乡的寄居。因此我将说你们可以承认、你们可以赞同的事：我将提醒你们所知道的事，我不会教导你们不知道的事。\Quote{不是首先的，}宗徒说，\Quote{那是属灵的，而是属血气的，然后是属灵的。}\BibleRef{格前 15:46}因此前者按时代更大，因为先出生的是\Person{加音}，后来是\Person{亚伯尔}：\BibleRef{创 4:1}但在此，年长的要服侍年幼的。\BibleRef{创 25:23}前者按时代更大，后者在尊荣上更大。为何前者按时代更大？因为\Quote{不是首先的，那是属灵的，而是属血气的。}\BibleRef{格前 15:46}为何后者在尊荣上更大？因为\Quote{年长的要服侍年幼的。}\BibleRef{创 25:23}……\Person{加音}首先建造了一座城，在无城之处建造。但当\Place{耶路撒冷}被建造时，并非在无城之处建造，而是起初已有一座城，称为\Place{耶步斯}，由此有耶步斯人。这城被攻取、征服、制伏后，如同旧城被拆毁，就建造了一座新城，称为\Place{耶路撒冷}，\BibleRef{苏 18:28}即和平的异象，天主的城。因此，每个由\Person{亚当}所生的人，还不属于\Place{耶路撒冷}，因为他带着不义的苗裔和罪的惩罚，已被交付于死亡，并在某种意义上属于一种旧的城。但如果他要在天主的人民中，他的旧我就要被推倒，并被建造成新的。因此，\Person{加音}在无城之处建了一座城。因为每个人都是从必死性和邪恶出发，以便日后能成为善的。\Quote{因为正如因一人的悖逆，众人都成了罪人，照样，因一人的顺服，众人都将成义。}\BibleRef{罗 5:19}并且我们在\Person{亚当}内都要死：\BibleRef{格前 15:22}我们每人都是\Person{亚当}所生。让他转向\Place{耶路撒冷}，他将被摧毁旧人，并建成新人。如同对被征服的\Place{耶步斯}人所说，为了建造\Place{耶路撒冷}，\Quote{你们要脱去旧人，穿上新人。}现在，对那些在\Place{耶路撒冷}被建造、并被恩宠之光照耀的人说：\Quote{你们从前曾是黑暗，但现在在主内是光明。}\BibleRef{弗 5:8}因此，邪恶的城从起初直到终点一直运行，而善的城通过恶人的改变被建造起来。而这两座城在此时是混合的，在终末将被分开；彼此相互冲突，一个支持不义，另一个支持真理。有时正是这种现世的混合导致某些属于\Place{巴比伦}城的人料理属于\Place{耶路撒冷}的事务，而又有某些属于\Place{耶路撒冷}的人料理属于\Place{巴比伦}的事务。我似乎提出了困难的问题。请忍耐，直到通过例子证明。\Quote{所有这些事}在旧百姓身上，如宗徒所写，\Quote{是以预表的方式发生在他们身上；但它们被写下来，是为了我们的劝诫，我们这些到了世界末了的人。}\BibleRef{格前 10:11}因此，要把那个民族看作也是预示着后来的民族；然后看我所说的。在\Place{耶路撒冷}有过伟大的君王：这是众所周知的事实，他们被列举，被命名。我说，他们全都是\Place{巴比伦}邪恶的公民，并且他们料理\Place{耶路撒冷}的事务：所有的人从那里在终末将被分离，不属于任何人，只属于魔鬼。我们再次发现\Place{耶路撒冷}的公民料理属于\Place{巴比伦}的某些事务。因为\Person{拿步高王}被一个奇迹折服，使那三个青年成为他王国的臣宰，派他们管理他的总督们；于是就有\Place{耶路撒冷}的公民料理\Place{巴比伦}的事务。现在观察这是如何在教会中、在这些时期得以实现和完成的……每一个地上的政体，某时必会灭亡，其国度将要消逝，当那国度来临时，我们祈祷\Quote{愿祢的国来临，}\BibleRef{玛 6:10}并且预言说\Quote{祂的王国将没有终结，}\BibleRef{路 1:33}我所说的地上的政体，有我们的公民管理它的事务。因为有多少信者、多少善人，既是他们城中的地方官，又是法官、将军、伯爵和君王？所有那些正义而善良的人，心中所装的只是关于你——天主的城——所说的最荣耀的事。并且，仿佛他们在那将要消逝的城中做奴役，即使在那里，圣城的导师们也嘱咐他们要对那些管辖他们的人保持忠诚，\Quote{或是向君王如元首，或向总督如奉天主差遣，为惩罚恶人、称赞善人，}\BibleRef{伯前 2:13}或作为仆役，应顺从他们的主人，\BibleRef{弗 6:5}甚至基督徒服从外教人，更好的应当对更坏的保持忠诚，暂时服侍，永远掌权。因为这些事发生，直到不义过去。仆役被命令要忍受不义且任性的主人：\Place{巴比伦}的公民被命令要由\Place{耶路撒冷}的公民忍受，甚至显出更多的关注，比他们若是同一\Place{巴比伦}的公民还多，如同履行诫命：\Quote{若有人强迫你走一里路，你就同他走两里。}\BibleRef{玛 5:41}……\par

5. \Quote{我怀着渴慕奔跑。}因为他们以恶报善：我正渴望他们；他们以为能阻止我的荣耀；我渴望将他们纳入我的身体。因为饮水时，我们不正是把外面的液体送入我们的肢体，吸入体内吗？梅瑟在牛犊的头颅上就是这样做的。\BibleRef{出 32:10} 牛犊的头颅是一个伟大的圣事。因为牛犊的头颅就是恶人的身体，形如吃草的牛犊，追求世俗之事：因为凡有血肉的都如同草。\BibleRef{依 40:6} ……如今还有什么比这更明显：那作为以色列民族预像的耶路撒冷城，藉着圣洗圣事，人们将得以进入其中？因此，水被洒出，为能作为饮料赐予。为此，直到终结，这人始终渴慕；他奔跑并渴慕。因为祂饮下许多人，却永不失去渴慕。因此有：\Quote{我渴，妇人，请给我水喝。} \BibleRef{若 4:7} 那位撒玛黎雅妇人在井边遇见了渴慕的主，并被祂的渴慕所充满：她先发现祂渴慕，为能因她的信德而饮下祂。当祂在十字架上时，祂说：\Quote{我渴。} \BibleRef{若 19:28} 尽管他们没有给祂祂所渴慕的。因为祂是为他们而渴慕：但他们给了醋，而非新酒，新酒应装在新皮囊里，而是旧酒，旧到其损失。\BibleRef{玛 9:17} 因为旧醋也指旧人，关于他们曾说：\Quote{因为他们没有改变}；即耶步斯人应被推翻，耶路撒冷得以建成。\BibleRef{撒下 5:9}\par

6. 同样，这身体的头从起初直到终结也怀着渴慕奔跑。好像有人对祂说：为何渴慕？祢缺少什么？基督的身体、基督的教会，祢享有如此尊荣、如此崇高、如此高位，甚至在此世建立，祢还缺少什么？那预言关于你的已应验：\Quote{地上的众王都要朝拜祂，万民都要事奉祂。} ……那些在耶路撒冷的节期充满教堂的人，在巴比伦的节期则充满剧场：尽管如此，他们全体事奉、尊敬、服从她——不仅是那些领受基督圣事却憎恨基督诫命的人，甚至那些连圣事也不领受的人，无论是外邦人还是犹太人——他们尊敬、赞美、宣扬，\Quote{但他们用口祝福。} 我不听口舌，那教导我的知道：\Quote{他们心里却咒骂。} 他们在那地方咒骂，即\Quote{他们以为能阻止我的荣耀。}\par

7. 你在做什么，耶杜通？基督的身体，你跃过他们？你在这一切中做什么？你要做什么？你会昏厥吗？你不坚持到底吗？你岂不听：\BibleQuote{那坚持到底的，必要得救。} \BibleRef{玛 10:22} 尽管罪恶增多，许多人的爱德会冷淡？\BibleRef{玛 24:12} 你跃过他们又在哪里？你的家乡在天上又在哪里？\BibleRef{斐 3:20} 但他们紧贴世俗之事，如同土生之人思念大地，且是蛇的食物。\BibleRef{创 3:14} 你在这些事中做什么？……\BibleQuote{然而，我的灵魂当顺服于天主。} \BibleRef{咏 61:5} 谁能忍受如此大事：公开的战争或暗中的埋伏？谁能忍受在公开的敌人和假弟兄中的如此大事？谁能忍受？是一个人吗？若是人，他能靠自己吗？我没有如此跃起以致高升而跌倒：\Quote{我的灵魂顺服于天主，因为我的忍耐来自祂。} 在如此多的绊脚石中，有何忍耐？除非\BibleQuote{如果我们盼望那看不见的，我们就耐心等待。} \BibleRef{罗 8:25} 我的痛苦来了，我的安息也将来临；我的磨难来了，我的炼净也将来临。难道金子会在炼金炉中闪光吗？它会在项链中闪光，在装饰中闪光：但让它忍受炉火，以便除去渣滓而进入光明。这就是火炉：其中有糠秕，有金子，有火，炼金者鼓风；火炉中糠秕燃烧，金子被炼净；糠秕化为灰烬，金子除去渣滓。火炉是世界，糠秕是不义之人，金子是义人；火是磨难，炼金者是天主：因此，炼金者所愿的，我便去做；造物主安置我在哪里，我便忍受。我奉命忍受，祂知道如何炼净。尽管糠秕燃烧为了点燃我，如同要吞灭我；那化为灰烬的，我却被除去渣滓。为何？因为\Quote{我的灵魂顺服于天主，因为我的忍耐来自祂。}\par

8. \BibleQuote{因为祂是我的天主、我的拯救者、我的提携者，我必不移动} \BibleRef{咏 61:6}。因为\Quote{祂是我的天主}，所以祂召叫我；\Quote{祂是我的拯救者}，所以祂使我成义；\Quote{祂是我的提携者}，所以祂使我受光荣。因为我在此蒙召并成义，但在那里我将受光荣；从那个我受光荣的地方，\Quote{我必不移动}。因为我是与祢同在地上作客旅，如同我所有的祖先一样。因此，我将从我的寓所移动，但从我的天上家乡我必不移动。\BibleQuote{在天主内有我的救恩和我的光荣} \BibleRef{咏 61:7}。我将在天主内得救，我将在天主内受光荣：因为不仅得救，而且受光荣；得救，是因为我从一个不敬虔的人被祂成义而成为义人；\BibleRef{罗 4:2} 但受光荣，是因为不仅成义，而且受尊荣。因为\BibleQuote{祂所预定的人，也召叫了他们} \BibleRef{罗 8:30}。召叫他们，祂在这里做了什么？\Quote{祂所召叫的人，也使他们成义；祂所使他们成义的人，也使他们受光荣。}因此，成义属于救恩，受光荣属于尊荣。受光荣如何属于尊荣，无需讨论。成义如何属于救恩，让我们寻求一些证据。看，从福音中想起：有一些人自以为义人，他们责怪主，因为主接纳罪人赴宴，并与税吏和罪人一同吃饭；因此面对这些骄傲自大、地上的强者极其高举、夸耀自己的健康（他们自以为有，实际却没有）的人，主回答了什么？\BibleQuote{健康的人不需要医生，而是有病的人} \BibleRef{玛 9:13}。祂称谁为健康，称谁为有病？祂接着说：\BibleQuote{我来不是召义人，而是召罪人悔改} \BibleRef{玛 9:12}。因此，祂称\Quote{健康}的义人，不是因为法利塞人真是如此，而是因为他们自己以为如此；并且因此骄傲，嫉妒病人有医生，而他们比那些人病得更重，竟杀害了医生。然而，祂称健康的人是义人，有病的人是罪人。因此，那个跳跃的人说：我的成义来自祂，我的受光荣来自祂：\Quote{因为天主是我的救恩和我的光荣。}\Quote{我的救恩}，使我得救；\Quote{我的光荣}，使我受尊荣。这是将来的事：现在呢？\Quote{天主是我的帮助，我的希望在于天主}，直到我达到完全的成义和救恩。\BibleQuote{因为我们得救是在于希望，但看得见的希望不是希望} \BibleRef{罗 8:24}……\par

9. \BibleQuote{你们在祂内盼望，众民的议会} \BibleRef{咏 61:8}。效法\Person{耶杜顿}，跃过你们的仇敌；那些攻击你们、拦阻你们道路、仇恨你们的人，你们要跃过：\Quote{在祂内盼望，众民的议会：在祂面前倾吐你们的心}……借着恳求、借着承认、借着盼望。不要把你们的心留在你们心中：\Quote{在祂面前倾吐你们的心}。你们所倾吐的不会失落。因为祂是我的提携者。如果祂提携，你们为何害怕倾吐？\Quote{将你们的忧虑交托给主，并在祂内盼望}。你们在那些耳语者、毁谤者、憎恶天主的人（\BibleRef{罗 1:29}）中间，他们能公开攻击，不能时则暗中埋伏，虚假地赞美，真实地敌视，在他们中间你们怕什么？\Quote{天主是我们的助佑}。他们能与天主同等吗？他们比祂更强吗？\Quote{天主是我们的助佑}，你们不必挂虑。\BibleQuote{如果天主为我们，谁能反对我们？} \BibleRef{罗 8:31} \Quote{在祂面前倾吐你们的心，}借着跃向祂，借着提升你们的灵魂：\Quote{天主是我们的助佑}。……\BibleQuote{然而，人的子孙是虚妄的，人的子孙在秤上是说谎者，为的是欺骗，他们因虚妄而合一} \BibleRef{咏 61:9}。确实有许多人：看，有那一个人，就是被从众多宾客中赶出去的那个人。\BibleRef{玛 22:11} 他们共谋，他们都寻求暂时之物，属血肉的人寻求属血肉的事，并且盼望将来得到它们，凡是盼望的人：即使因意见分歧而分裂，但因虚妄而合一。错误确实多样且多形式，凡一国自相纷争，必站不住：\BibleRef{玛 12:25} 但在所有人中，虚妄而说谎的意志都一样，属于一个君王，与之一同被投入永火 \BibleRef{玛 25:41}——\Quote{这些人因虚妄而合一。}并且请看祂如何渴望他们，请看祂如何在干渴中奔跑。\par

10. 祂于是转向他们，渴求他们：\BibleQuote{不要指望不义} \BibleRef{咏 61:10}。因为我的指望在天主。你们这些不愿走近并越过的人，\BibleQuote{不要指望不义。}因为我已越过，我的指望在天主；难道天主有不义吗？\BibleRef{罗 9:14} 我们做这事，做那事，想那事，如此调整我们的埋伏；\BibleQuote{因为虚荣合而为一。}你们渴求：那些图谋反对你们的人被你们所饮的人抛弃。\BibleQuote{不要指望虚妄。}不义是虚妄的，不义是虚无的，唯有正义是强大的。真理可能暂时隐藏，但不会被征服。不义可能暂时兴旺，但不能长存。\BibleQuote{不要指望不义：也不要贪图抢劫。}你不富，还要抢劫吗？你找到什么？你失去什么？啊，失利的收获！你找到金钱，你失去正义。\BibleQuote{不要贪图抢劫。}……因此，虚妄的人子，说谎的人子，既不要抢劫，也不要因财富涌流而倾心于它们：不要再爱虚妄，寻求谎言。因为\BibleQuote{那以天主为上主为指望的人有福了，他没有留意虚妄和说谎的愚昧。}你愿意欺骗，愿意行欺诈，你提出什么以便行骗。不诚实的秤。因为他说：\BibleQuote{人子在天秤上是说谎的，}以便他们通过拿出不诚实的秤来行骗。借着不诚实的秤，你们哄骗看着的人：难道你们不知道，有一个是称重者，另一个是衡量重量者？祂看不见你们为谁称重，但祂看见那称量你们和祂的那位。因此，不要再贪图欺诈或抢劫，也不要将你们的指望放在那些事上：我已劝诫，已预言，——这个\Person{耶杜通}说。\par

11. 接下来是什么？\BibleQuote{天主曾一次说了，这两件事我听见了，即能力属于天主} \BibleRef{咏 61:11} \BibleQuote{，主啊，仁慈也属于祢，因为祢要按各人的行为报应各人} \BibleRef{咏 61:12} ……\BibleQuote{天主曾一次说了。}\Person{耶杜通}，你说什么？如果你这个越过他们的人说：\BibleQuote{祂一次说了；}我转向另一段圣经，它向我说：\BibleQuote{天主在古时，曾多次并以多种方式，藉着先知向祖先说过话} \BibleRef{希 1:1}。那么，\BibleQuote{天主曾一次说了}是什么意思？难道不是那个在人类之初向亚当说话的天主吗？\BibleRef{创 3:17} 难道同一位天主没有向加音、诺厄、亚巴郎、依撒格、雅各伯、众先知以及梅瑟说话吗？梅瑟是一个人，天主向他说话多少次？看，甚至对一个人，天主不是一次而是多次说话。其次，祂曾对站在这里的圣子说：\BibleQuote{祢是我的爱子} \BibleRef{玛 3:17}。天主曾向宗徒们说话，向所有圣人说话，虽不是透过云中的响声，而是借着心中祂亲自为师。那么，\BibleQuote{天主曾一次说了}是什么意思？那人已越过许多，才能到达天主一次说话的地方。看，我已对你们的爱德简短地说了。在这里，在人中间，对人们多次、以多种方式、藉着多种形式的受造物，天主说了话；祂自己一次说了，因为天主生了唯一的圣言。……因为天主不可能不知道祂藉着圣言所造的：\BibleRef{若 1:3} 但如果祂造了祂所知道的，那么在祂内就有那被造物在受造之前就已存在。因为如果在祂内没有那受造之前就已存在的，祂又如何知道祂所造的？因为你不能说天主造了祂所不知道的。所以天主知道祂所造的。祂在造之前如何知道？除非藉着受造物，否则无从知道已存在之物；但对于你，即一个被造在较低处、被安置在较低处的人而言；但在所有这些受造物存在之前，它们已被造它们的那位知道，祂知道祂所造的。因此，在祂所造万物的圣言中，在它们存在之前，万物就已存在；在它们被造之后，万物也存在；但在这里是一样，在那里是另一样；在它们被造的本性中是一样，在造它们的技术中是另一样。谁能解释这事？我们可以努力：与\Person{耶杜通}一同前去，观看。\par

12. ……因为主也说：\BibleQuote{我还有许多事要告诉你们，但你们现在不能担负} \BibleRef{若 16:12}。那么，\BibleQuote{这两件事我听见了}是什么意思？这两件事我要告诉你们，不是凭我自己对你们说的，而是我所听见的，我就说。\BibleQuote{天主曾一次说了：}祂有唯一的圣言，独生子天主。在圣言内有一切，因为万有是藉着圣言造成的。祂有唯一的圣言，\BibleQuote{在祂内蕴藏着一切智慧和知识的宝藏} \BibleRef{哥 2:3}。祂有唯一的圣言，\BibleQuote{天主曾一次说了。} \BibleQuote{这两件事，}我要向你们说的，是我听见的：我不是凭自己说话，不是凭自己说：这属于\BibleQuote{我听见了} \BibleRef{若 8:26}。但新郎的朋友站着听祂，为要讲真理。因为他听祂，免得说谎言，讲自己的话：\BibleRef{若 8:44} 免得你说：你是谁，竟对我说这事？你从哪里向我说这事？我听见了这两件事，而我向你们说这两件事是我听见的，我也知道天主曾一次说了。不要轻视一个听者，他向你们说这两件对你们如此必要的事；我是说，那个藉着越过整个受造物而达到天主的独生圣言的人，在那里他学到了\BibleQuote{天主曾一次说了。}\par

13. 让祂如今说这两件事。因为这两件事对我们至关重要。\Quote{因为权能属于天主，主，仁慈也属于祢。}就是这两件事吗，权能和仁慈？显然是这两件：你们要认识天主的权能，认识天主的仁慈。几乎全部圣经都包含在这两件事里。因为这两件事，才有了先知，因为这两件，才有了圣祖，因为这两件，才有了法律，因为这两件，才有我们的主耶稣基督自己，因为这两件，才有了宗徒，因为这两件，才有了教会中天主圣言的一切宣讲与传播，就因为这两件，就因为天主的权能和祂的仁慈。你们要畏惧祂的权能，你们要爱慕祂的仁慈。但既不可如此信赖祂的仁慈，以致轻视祂的权能；也不可如此畏惧祂的权能，以致对祂的仁慈绝望。权能同祂一起，仁慈也同祂一起。祂贬抑这人，举扬那人；祂以权能贬抑这人，以仁慈举扬那人。\BibleQuote{因为天主愿意显示义怒并彰显祂的权能，遂以极大的容忍，宽容了那些义怒之器，注定要丧亡的} \BibleRef{罗 9:22}——你们听到了权能：去寻求仁慈吧——\BibleQuote{并要彰显，}祂说，\BibleQuote{祂的荣耀丰富归于仁慈之器。}因此，定不义之人的罪属于祂的权能。至于那要说\Quote{你做了什么事？}的人，\BibleQuote{因为人哪，你是谁，竟敢向天主抗辩？} \BibleRef{罗 9:20} 因此，你们要畏惧、战栗于祂的权能，但要盼望祂的仁慈。魔鬼也是一种权势；然而，他常常想要伤害却不能，因为那权势是在权能之下。因为如果魔鬼能随心所欲地伤害，就没有一个义人存留，世上也不会有任何一位信徒。他藉自己的器皿击打，如同击打一堵倾斜的墙；但他只能击打，直到他获得权能为止。但为使那墙不至倒塌，主必扶持：因为赐予诱惑者权能的，祂自己也向被诱惑者广施仁慈。因为魔鬼被允许诱惑是有定量的。并且，\Quote{祢以眼泪给我们量杯喝。}因此，不要怕那被允许做些什么的诱惑者，因为你有一位最仁慈的救主。他被允许诱惑你到什么程度，就对你有什么益处，好使你受磨练，受考验；使你藉自己认识你自己，因为你不认识自己。因为，除了靠着天主的权能和仁慈，我们哪里或从哪里能得安全呢？继宗徒的话之后，\BibleQuote{天主是信实的，祂不许你们受诱惑超过你们所能承受的。} \BibleRef{格前 10:13} ……不要怕仇敌：他做多少，就从他得了多少权能；要怕那拥有至高权能的那位：要怕祂，祂随意行事，祂从不行不义，凡祂所做的，都是公义的。我们或许会认为某事不义，但既然天主做了，就该相信它是公义的。\par

14. 因此，你说，如果有人杀害一个无辜的人，他是正义地还是不正义地？当然是不正义地。为什么天主允许这事？……我不能告诉你天主的计划，哦，人；但我可以说：那人杀死一个无辜者，做了不义之事，而此事若非天主允许，就不会发生；尽管那人行了不义，但天主允许这事并非不义。让这个理由隐藏在那个人的身上，无论他是谁，因为你的情感为此波动，他的无辜确实使你大大激动。因为我可以很快回答你：他若不是有罪的，就不会被杀；但你认为他无辜。我可以很快对你说这话。因为你不能查探他的心，筛察他的行为，衡量他的思想，从而可以对我说：他是不正义地被杀。因此我可以很容易地回答；但有一个无可争议的义者，毫无疑问的义者，没有罪，被罪人杀害，被一个罪人出卖；那就是基督自己，我们不能说祂有任何不义，因为\Quote{祂没有抢夺的，祂偿还了}，这成了对我回答的反驳。我为什么要提到基督呢？你说：\Quote{我在与你交涉。}我也与你交涉。关于祂你提出了问题，关于祂我在解答问题。因为在那里我们知道了天主的计划，若非祂亲自启示，我们就不知道；这样，当你发现了那个天主的计划，祂借此允许祂的无辜圣子被不义的人杀害，而这个计划令你满意，如果你正义，它也不能令你不满意，你就可以相信，在其他事上，天主也按祂的计划行事，只是你未曾察觉。啊！弟兄们，需要义者的血来涂抹罪债；需要忍耐的榜样，谦卑的榜样；需要十字架的记号来击败魔鬼和他的天使；我们需要主的苦难；因为主的苦难，世界被救赎。主的苦难成就了多少善事！然而，这位义者的苦难不会发生，除非不义的人杀害了主。那么，这藉主的苦难赐予我们的善事，要归功于基督的不义杀害者吗？绝不可。他们愿意，天主允许。即使他们只是愿意，他们也是有罪的；但天主若不允许，除非这事是正义的，否则祂不会允许……因此，我的弟兄们，犹大，那对基督的污秽叛徒，以及基督的迫害者，都是恶毒的，不敬的，不义的，都该受谴责；尽管如此，圣父没有放过自己的圣子，反而为了我们众人把祂交出。\BibleRef{罗 8:32}如果你能，就排序；如果你能，就区分：向天主偿还你的誓言，就是你的嘴唇所说出的。看不义者在这里做了什么，义者做了什么。一个愿意，另一位允许：一个不义地愿意，另一位正义地允许。让不义的意愿受谴责，让正义的允许受光荣。因为基督遭遇了什么恶事？基督死了。那些恶意行恶的人都是邪恶的，然而，祂在他们身上所行的，并未遭受任何恶。被杀的是必死的肉身，用死亡杀死死亡，给出忍耐的教训，预先送出复活的榜样。义者藉不义者的恶行成就了多少善事！这是天主伟大的奥秘：甚至你所做的善事，也是祂亲自赐予你的，而祂亲自藉你的恶行成就善事。因此，不要惊奇，天主允许，且在审判中允许；祂在度量、数目、重量中允许。在祂那里没有不义：\BibleRef{罗 9:14}你只要属于祂；将你的希望寄托于祂，让祂成为你的助佑，你的救恩；让祂成为堡垒，力量的塔楼，让祂成为你的避难所，祂必不让你受试探超过你所能承载的，反而在试探中给你一条出路，使你能够支持：\BibleRef{格前 10:13}这样，祂让你承受试探，是祂的能力；祂不允许在你身上加增超过你所能承受的，是祂的仁慈：\Quote{因为权力属于天主，上主，慈爱也归于祢，因为祢将按各人的行为施行报应。}\par

15. 教会的这种渴望，也切愿汲取你所看见的那个人。同时，为使你们知道在混杂的基督徒群众中，有多少人用口祝福，心中却咒骂，这个人曾是基督徒和信徒，如今以忏悔者身份归来，因主的大能而恐惧，转向主寻求仁慈。因为他作信徒时被仇敌引入歧途，长期作占星术士，受迷惑，迷惑人，受骗，骗人，引诱、欺哄，说了许多反对天主的谎话——祂赐给人行善和不行恶的能力。他常说：行奸淫的不是自己的意志，而是维纳斯；杀人的不是自己的意志，而是马尔斯；天主不行正义，而是朱庇特；还有许多其他亵渎的话，且不是轻微的。你们想，他从多少基督徒那里骗取了钱财？多少人从他那里买了谎言，我们曾对他们说：\BibleQuote{人子们，你们心钝要到几时呢？为什么喜爱虚妄，寻求谎言？} 如今，如他必须被相信的那样，他因自己的谎言而战栗，曾是许多人的诱惑者，终于察觉自己是被魔鬼诱惑了，便以忏悔者身份转向天主。弟兄们，我们认为，这是由于心中极大的恐惧而发生的。我们该说什么呢？如果一个外教人中的占星术士归正，那真是极大的喜乐；但无论如何，可能会显得，如果他归正了，他是渴望在教会中获得圣职。他是一位忏悔者，除了仁慈之外别无所求。所以，他必须被推荐给你们的目光和心灵。你们用心灵爱他，用目光守护他。你们看见他，留意他，无论他走到哪里，向那些此刻不在这里的其余弟兄指出他；这样的勤勉就是仁慈，以免那诱惑者夺回他的心并强占它。守护他，不要让他的言行、他的道路逃过你们的眼睛；这样，凭你们的见证，就可以向我们证明他确实转向了主。因为关于他的生活，报告不会沉默，当他这样被呈现给你们观看并怜悯时。你们知道在\WorkTitle{宗徒大事录}中是如何记载的：许多迷失的人，即从事此类技艺、追随邪僻教义的人，把他们的书籍带到宗徒们面前；被焚烧的卷册如此之多，以致作者需要估价并记下总价。\BibleRef{宗 19:19} 这确实是为了天主的荣耀，好让那些迷失的人也不会被那位知道如何寻找迷失者的人所绝望。所以，这个人曾经迷失，如今被寻找，被找到，\BibleRef{路 15:32} 被领到这里，他带来要焚烧的书籍，这些书曾要焚烧他，好使它们被投入火中后，他本人能转入安息之地。弟兄们，你们知道，他在复活节前曾敲过教会的大门：因为他在复活节前就开始向教会寻求基督的良药。但由于他所操练的技艺是那种被怀疑为谎言和欺骗的，他被推迟了，以免他试探；但最终他被接纳了，以免他更危险地被试探。请你们为他在基督内祈祷。今日的祈祷立刻为你们倾流，向我们的主天主为他祈祷。因为我们知道并确信，你们的祈祷抹去他所有的不敬。愿主与你们同在。\par

\SourceRecord{Translated by J.E. Tweed.；Nicene and Post-Nicene Fathers, First Series；Vol. 8.；Edited by Philip Schaff.；Buffalo, NY: Christian Literature Publishing Co.,；1888.；<http://www.newadvent.org/fathers/1801062.htm>.}

% structure: p0001 source=h1 target=section reason=page-title

\section{圣咏第六十三篇释义}

1. 这篇圣咏的标题是：\Quote{达味自己，当他身在厄东旷野时。} 厄东这个名字象征这个世界。因为厄东原是一个迷失的民族，在那里崇拜偶像。厄东在这里没有任何好的意义。既然不是好的意义，就必须理解为：我们今生所承受的如此巨大的劳苦，所受的如此重大的需要，正是以厄东之名来表示。这里是一片旷野，其中有许多干渴，而你将听到那位如今在旷野中干渴者的声音。但如果我们承认自己是干渴的，我们也当承认自己将要喝饮。因为在这世界干渴的人，在来世将得饱足，正如主所说的：\BibleQuote{饥渴慕义的人是有福的，因为他们要得饱饫。} \BibleRef{玛 5:6} 因此，在这世界里我们不应爱慕满足。在此我们必须干渴，在别处我们将得饱足。但如今，为了使我们不致在这旷野中昏厥，祂将祂话语的露珠洒在我们身上，不让我们完全枯干，以致我们不应被再次寻找，而是使我们干渴到能够喝饮。为了使我们可以喝饮，我们被洒了些许祂的恩宠：然而我们仍然干渴。那么，我们的灵魂对天主说什么呢？\par

2. \BibleQuote{天主，我的天主，我从光明中醒寤向着祢} \BibleRef{咏 62:1}。什么是醒寤？就是不要沉睡。什么是沉睡？有灵魂的沉睡，有身体的沉睡。身体的睡眠我们都应当有：因为若不睡眠，人就会昏厥，身体本身就会昏厥。我们脆弱的身体无法长期保持灵魂醒寤并紧张地从事活跃的工作；若灵魂长时间专注于活跃的追求，身体脆弱而属土，无法承载她，无法永远支持她活动，就会昏厥倒下。因此天主赐给了身体睡眠，借此身体各肢体得以恢复，以便能够支持灵魂醒寤。但我们要留心，即我们的灵魂自己不要沉睡：因为灵魂的睡眠是邪恶的。身体的睡眠是好的，借此身体的健康得以恢复。但灵魂的睡眠是忘记她的天主。凡灵魂忘记了她的天主，就睡着了。因此宗徒对某些忘记天主、仿佛在沉睡中做出崇拜偶像之愚行的人说：\BibleQuote{你这睡着的人，醒来吧，从死者中起来，基督必要光照你。} \BibleRef{弗 5:14} 宗徒岂是唤醒一个身体睡着的人吗？不，他是在唤醒一个沉睡的灵魂，因为他唤醒她，使她可以被基督光照。因此，针对这些醒寤之事，这个人说：\BibleQuote{天主，我的天主，我从光明中醒寤向着祢。} 因为你不会自己醒寤，除非你的光明兴起，将你从睡眠中唤醒。因为基督光照灵魂，使它们醒寤；但若祂收回祂的光，它们就打盹。为此，在另一篇圣咏中向祂说：\BibleQuote{求祢光照我的眼睛，免得我沉睡于死亡。} ...\par

3. \BibleQuote{我的灵魂渴慕祢} \BibleRef{咏 62:2}。请看那厄东的旷野。请看他是如何在此干渴：但要看这里的好事：\BibleQuote{渴慕了祢。} 因为有人干渴，却不是为天主。因为每一个希望得到某物的人，都处于渴望的热度中；渴望本身就是灵魂的干渴。请看人心中有何等的渴望：有人渴求金子，有人渴求银子，有人渴求财产，有人渴求遗产，有人渴求大量钱财，有人渴求许多牲畜，有人渴求妻子，有人渴求尊荣，有人渴求儿子。你看这些渴望，如何在人心里。众人都被渴望所燃烧，却很难找到一个人说：\BibleQuote{我的灵魂渴慕了祢。} 因为人渴慕世界，却不觉察自己正处在厄东的旷野中，在那里，他们的灵魂本当渴慕天主。……\par

4. 因此必须渴慕智慧，必须渴慕义德。我们不得满足于它，不得因它而饱足，除非此生命结束，我们到达天主所应许之处。因为天主应许了与天使同等：\BibleRef{路 20:36} 而如今天使并不像我们一样渴，也不像我们一样饿；但他们拥有真理、光明、不朽智慧的全部。因此他们是有福的，并且出于如此巨大的福乐，因为他们是在那城，即天上耶路撒冷，我们如今远离它，在异地漂泊，他们注视我们这些旅人，并怜悯我们，并因主的命令帮助我们，为使我们终有一日能回到这共同的家乡，并在那里与他们一起，因主真理与永恒之泉而饱足。如今愿我们的灵魂渴慕：我们的肉体也同样渴慕，并以多种方式？他说：\Quote{以多种方式为了祢，} \Quote{我的肉体亦然。} 因为我们的肉体也蒙应许了复活。正如我们的灵魂蒙应许了福乐，同样我们的肉体也蒙应许了复活。……因为如果天主创造了我们这些原本不存在的，祂再造我们这些曾经存在的，岂是大事？因此，不要认为这是不可信的，因为你们看见死人仿佛腐朽，化为灰烬与尘土。或者若有死人被焚烧，或被狗撕裂，你们以为他就不能由此复活吗？所有被肢解、化为某种尘土的事物，在天主那里都是完整的。因为它们进入世界的元素，这些元素就是我们当初被造时所来自的；我们看不见它们；但天主会将其带出，祂知道从何处，因为甚至在我们存在之前，祂就从祂所知道之处创造了我们。因此，我们被应许了这样的肉体复活，即，尽管复活的是我们现在所携带的同一肉体，但它不再具有现在所有的腐朽。因为现在由于脆弱的腐朽，如果我们不吃，就会虚弱饥饿；如果不喝，就会虚弱口渴；如果长时间不睡，就会虚弱困倦；如果长时间睡，就会虚弱，因此我们醒着……其次，看我们的肉体如何没有固定的状态：婴儿期进入儿童期，你寻找婴儿期，婴儿期已不存在，因为现在代替婴儿期的是儿童期；这同样又进入青年期，你寻找儿童期，却找不着；青年人成为中年人，你寻找青年人，他已不在；中年人成为老年人，你寻找中年人，却找不着；老年人死去，你寻找老年人，却找不着；因此我们的年龄不停滞：处处是疲劳，处处是虚弱，处处是腐朽。注视到天主向我们应许的复活的希望，在我们所有这些多样的疲弱中，我们渴慕那种不朽：因此我们的肉体以多种方式渴慕天主。\par

5. 然而，我的弟兄们，良善基督徒和信徒的肉体，即使在此世，也确实渴慕天主：因为如果肉体需要面包，如果需要水，如果需要酒，如果需要金钱，如果这肉体需要牲畜，它应当从天主那里寻求，而不是从魔鬼、偶像和我不知道什么现世的能力那里。因为有些人，当他们在此世遭受饥饿时，离开天主，向墨丘利或朱庇特祈求，请他们赐予，或向他们称为\Quote{天上的}的那位，或任何类似的魔鬼祈求：他们的肉体不是为天主而渴。但那渴慕天主的人，应当在各处渴慕祂，既灵魂也肉体：因为对灵魂而言，天主也赐予祂的饼，即真理之言；对肉体而言，天主赐予必需之物，因为天主创造了灵魂和肉体。为了你们的肉体，你们向魔鬼祈求：难道天主创造了灵魂，而魔鬼创造了肉体吗？那创造了灵魂的，也创造了肉体；那创造了二者的，也喂养二者。愿我们每一部分都渴慕天主，并在多种劳苦之后，愿每一部分都单纯地得到饱足。\par

6. 然而，我们的灵魂和我们的肉体以多种方式渴慕，不是为其他任何一位，而是为祢，主，即我们的天主？它在何处渴慕？\Quote{在干旱、无路、无水之地}。我们谈论了此世，它就是\Place{厄东}，这是\Place{厄东}的旷野，圣咏由此得名。\Quote{在干旱之地}。只说\Quote{不毛之地}是不够的，那里无人居住；它还是\Quote{无路，且无水}。啊，愿那沙漠也有一条路：啊，愿进入其中的人，即使奔跑，也知道如何从那里出去！……沙漠是邪恶的、可怕的、可惧的：然而天主怜悯了我们，并在沙漠中为我们开辟了一条路，即主\Person{耶稣基督}自己：\BibleRef{若 14:6} 并在沙漠中为我们制造了安慰，即派遣祂圣言的宣讲者给我们：并赐给我们沙漠中的水，即以圣神充满祂的宣讲者，为在他们内创造一口涌到永生的水泉。\BibleRef{若 4:14} 看哪！我们在这里拥有一切，但它们不属于这沙漠……\par

7. \Quote{如此，我在圣所中朝拜了祢，为能看见祢的能力和祢的荣耀} \BibleRef{咏 62:3}。……人若不先在旷野中渴慕，即在他所处的邪恶中渴慕，决不会到达那善，即是天主。但他说：\Quote{我在圣所中朝拜了祢。} 在圣所中有极大的安慰。\Quote{我在圣所中朝拜了祢}是什么意思？是为了祢能看见我：祢看见了我，是为了我能看见祢。\Quote{我朝拜了祢，为能看见。} 他没有说：\Quote{我朝拜了祢，为叫祢看见，} 而是说：\Quote{我朝拜了祢，为能看见祢的能力和祢的荣耀。} 因此宗徒也说：\Quote{但现在，}他说，\Quote{认识天主，更确切地说，是被天主所认识。} \BibleRef{迦 4:9} 因为你首先朝拜了天主，为能使天主得以显现给你。\Quote{为能看见祢的能力和祢的荣耀。} 实在，在那被遗弃的地方，即那旷野中，若有人仿佛从旷野中竭力获得足够的给养，他决不会看见主的能力和主的荣耀，而会渴死，也找不到路、安慰或水，使他在旷野中得以存活。但是，当他提升自己到天主面前，以至从内心深处对祂说：\Quote{我的灵魂渴慕祢；我的肉身也多么渴慕祢！} 免得他或许也为别人肉身的必需品而求，而不是求于天主，或不再渴望天主所应许我们的肉身复活；我说，当他提升自己时，他将得到不小的安慰。\par

8. ……但你们刚听到福音宣读时，祂以何等言辞宣告了祂的尊威：\BibleQuote{我与父原是一体} \BibleRef{若 10:30}。看，多么伟大的尊威，与父多么伟大的平等，因着我们的软弱而降到了肉身。我们被爱得何其深，竟在我们爱天主之前。如果在我们爱天主之前，我们已被祂所爱，以致祂那与己平等的圣子，为了我们的缘故成了人，那么如今我们爱祂，祂又为我们预备了什么呢？因此许多人认为天主子显现于世上是一件微不足道的事；因为他们不在圣者内，他们看不见那同一者的能力和那同一者的荣耀：就是说，他们还没有一颗圣洁的心，从而能觉察那德性的卓越，并向天主感恩，也看不见那一位为他们自身而来的如此伟大者，祂降生到何等境地，经历何等苦难，祂的荣耀和祂的能力，他们都看不见。\par

9. \Quote{因为祢的仁慈比生命更美好。} 人的生命有许多，但天主应许的只有一种生命：祂赐予我们，并非因我们的功德，而是因祂的仁慈……有什么比罪人应受惩罚更公正的事呢？虽然罪人受罚是公正的，但祂的仁慈却是不惩罚罪人，反而使他成义，使罪人成为义人，使不敬虔的人成为敬虔的人。因此\Quote{祂的仁慈比生命更美好。} 是什么生命？是那些人们为自己选择的生活。有人选择了经商的生活，有人选择了田园生活，有人选择了放贷的生活，有人选择了军旅生活；这个一种，那个另一种。生活多种多样，但\Quote{祢的生命比我们的生命更美好。} ……\Quote{我的嘴唇要赞美祢。} 我的嘴唇不会赞美祢，除非祢的仁慈走在我前头。因祢的恩赐我赞美祢，藉祢的仁慈我赞美祢。因为我不能赞美天主，除非祂赐给我能力赞美祂。\par

10. \Quote{为此，我在一生中要称颂祢，因祢的名我要举起双手} \BibleRef{咏 62:5}。如今在我的一生中，即祢所赐予我的，不是我在世上随同众人从众多生活中所选择的，而是祢藉着祢的仁慈赐予我的，使我赞美祢。\Quote{为此，我在一生中要称颂祢。} 什么是\Quote{为此}？即我将我赞美祢的生命归于祢的仁慈，而非我的功德。\Quote{因祢的名我要举起双手。} 因此，在祈祷中举起双手吧。我们的主在十字架上为我们举起了双手，祂的双手为我们伸展，正是为了我们的双手能伸向善工：因为祂的十字架带给我们仁慈。看，祂举起了双手，将祂自己作为祭献献给了天主，藉着那祭献，我们所有的罪都被抹去了。让我们也在祈祷中向上主举起双手：若我们的双手行善工，它们向上主举起时便不致羞愧。因为举起双手的人要做什么？为何曾命令我们举手祈祷？因为宗徒说：\BibleQuote{举起圣洁的手，没有愤怒和争论} \BibleRef{弟前 2:8}。这是为了使当你向上主举手时，你的善工会浮现于你脑海中。因为那些手举起是为了获得你所愿的，你也想用那些手行善工，使它们不致因向上主举起而羞赧。\Quote{因祢的名我要举起双手。} 这就是我们在厄东，在这旷野，在这无水无路之地的祈祷；在这里，基督为我们是道路，\BibleRef{若 14:6}，但不是今世之路。\par

11. …我们的祖先已死，但天主活着：在此我们不能永远有父，但在那里，当我们拥有我们的父家时，我们将永远有一位生活的圣父。……那是怎样的地方？但你在此爱财富。天主自身将成为你的财富。但你爱一口好泉。有什么比那智慧更清澈？有什么更光明？凡在此可爱之物，你将拥有那创造万有的祂来替代一切，\BibleQuote{有如我灵魂饱沾脂膏：我的口唇要以欢乐的唇舌赞美你的圣名。}在这旷野里，我要因你的名举起双手：愿我的灵魂饱沾脂膏，\BibleQuote{我的口唇要以欢乐赞美你的圣名。}因为现在是祈祷，只要仍有渴慕；当渴慕过去，祈祷过去，赞美接续。\BibleQuote{我的口唇要以欢乐赞美你的圣名。}\par

12. \BibleQuote{若我在床上想起了你，在破晓时我默想了你\BibleRef{咏 62:7}：因为你成了我的助佑}\BibleRef{咏 62:8}。他称\Quote{床}为安息。当人安息时，让他记念天主；当人安息时，让他不要因安息而松懈，忘记天主：若他在安息时记念天主，他在行动中便默想天主。因为他称破晓为行动，因为每人在破晓时开始做某事。因此他说了什么？因此，若我在床上没有记念你，在破晓时我也没有默想你。难道那在闲适时不想天主的人，在行动中会想天主？但那在安息时记念祂的人，在行事时默想同一位，免得在行动中跌倒。因此他又加了什么？\BibleQuote{因为你成了我的助佑。}因为除非天主帮助我们的善工，它们无法由我们完成。我们应当做有价值的事：即如同在光明中，因为基督指明了道路，我们工作。凡行恶事的人，他在黑夜工作，不在破晓；按照宗徒所说：\BibleQuote{醉酒的人，是在夜间醉酒；睡觉的人，是在夜间睡觉；我们这些属于白昼的，应当清醒。}\BibleRef{得前 5:7}他劝勉我们，要在白昼后正直行走：\BibleQuote{如同在白昼，我们要正直行走。}\BibleRef{罗 13:13}又说：\BibleQuote{你们}，他说，\BibleQuote{是光明之子，白昼之子；我们不属于黑夜，也不属于黑暗。}\BibleRef{得前 5:5}谁是黑夜之子、黑暗之子？就是那些行一切恶事的人。他们如此是黑夜之子，以至于害怕他们所行的事被看见。……因此没有人在破晓工作，除非那在基督内工作的人。但那在闲适时记念基督的人，在一切行动中默想同一位，祂在善工中成为他的助佑，免得因他的软弱而失败。\BibleQuote{在你的翅膀荫庇下，我将欢跃。}我在善工中喜乐，因为你的翅膀荫庇在我之上。若你不保护我，我不过是小鸡，鹞鹰会攫取我。因为我们的主曾向耶路撒冷——祂被钉死的那座城——说：\BibleQuote{耶路撒冷，}祂说，\BibleQuote{耶路撒冷，我多少次愿意聚集你的子女，如同母鸡聚集小鸡，而你却不愿意。}\BibleRef{玛 23:37}我们是幼小的；因此愿天主在祂翅膀的庇荫下保护我们。当我们长大了呢？对我们有益的事是，即使那时祂仍保护我们，使我们在祂这位更大者之下，常是小鸡。因为无论我们长得多大，祂总是更大。没有人说：让我还是幼小时就受祂保护；仿佛有一天他会长到自足的地步。没有天主的保护，你什么也不是。我们常愿受祂保护；这样，我们若常在祂面前微小，便常能在祂内成为伟大。\BibleQuote{在你的翅膀荫庇下，我将欢跃。}\par

13. \BibleQuote{我的灵魂紧贴着你}\BibleRef{咏 62:9}。看一个渴望的人，看一个饥渴的人，看他如何紧贴天主。愿这情感在你内生发。若已萌芽，愿它被浇灌而成长；愿它达到如此力量，以致你也能全心说：\BibleQuote{我的灵魂紧贴着你。}那黏胶在哪里？那黏胶就是爱。拥有爱，你的灵魂便以它如胶般紧贴在天主之后。不是与天主一起，而是在天主之后；让祂前行，你跟随。因为那愿意行在天主之前的人，会按自己的主意生活，不遵从天主的诫命。为此，伯多禄也曾受责，当他想要给那将要为我们受苦的基督出主意时，……\BibleQuote{主，请免了祢吧，求祢怜悯自己。}主说：\BibleQuote{退到我后面去，撒殚：因为你所体会的，不是天主的事，而是人的事。}\BibleRef{玛 16:22}为何是人的事？因为你想要行在我前面，退到我后面去，好使你跟随我：这样，如今跟随基督的他能说：\BibleQuote{我的灵魂紧贴着你。}他合理加上：\BibleQuote{你的右手扶持了我。}这是基督在我们内所说的：即在祂为我们承担、为我们奉献的那人性中，祂说了这话。教会也在基督内说了这话，她在她的头内说了这话；因为她也在此地受了极大的迫害，并且甚至现在她仍通过她的各个肢体遭受迫害……\par

14. \BibleQuote{他们徒然寻索我的灵魂；他们必往大地深处去。} \BibleRef{咏 62:9}。他们不愿失去大地，当他们钉死基督时，他们已进入大地深处。大地深处是什么？是地上的情欲。行走在地上，比因情欲而进入地下更好。因为凡是为自己的得救而贪求地上事物的人，就在地下：因为他把大地放在自己面前，把大地压在自己身上，而把自己放在底下。因此，他们害怕失去大地，当看见大批群众跟随主耶稣基督，因为祂正行奇事时，他们说了什么呢？\BibleQuote{我们若让祂这样活着，罗马人必要来，夺去我们的地方和民族。} \BibleRef{若 11:48} 他们害怕失去大地，却到了地下：甚至他们所害怕的事也临到了他们。因为他们愿意杀死基督，以免失去大地；而他们正因为杀了基督，反而失去了大地。因为当基督被杀时，主自己曾对他们说：\BibleQuote{天主要从你们手中夺去天国，交给一个结果子的民族。} \BibleRef{玛 21:43} 随后，巨大的迫害灾难降临到他们身上；罗马皇帝和外邦的君王征服了他们；他们被赶出了他们钉死基督的那个地方，如今那地方充满了赞美基督的人；那里没有犹太人了，那些基督的敌人被清除了，那里充满了赞美基督的人。看，他们因罗马人的手失去了那地方，因为他们杀了基督——他们杀基督正是为了不因罗马人而失去那地方。因此，\BibleQuote{他们必进入大地深处。}\par

15. \BibleQuote{他们将被交在刀剑的手下。} \BibleRef{咏 62:10} 事实上，这事已清楚地临到他们身上：他们被攻破的敌人掳掠了。\BibleQuote{他们将成为狐狸的分。}他把当时统治世界的君王称为狐狸，那时犹大被征服。请听，好使你知道并明白，他称那些人为狐狸。主自己曾称黑落德王为狐狸：\BibleQuote{你们去告诉那只狐狸。} \BibleRef{路 13:32} 你们看，我的弟兄们：他们不想有基督作君王，于是成了狐狸的分。因为当总督比拉多在犹大按犹太人的声音杀死基督时，他对犹太人说：\BibleQuote{我要钉死你们的君王吗？} \BibleRef{若 19:15} 因为祂被称为犹太人的君王，祂确实是真正的君王。而他们拒绝基督说：\BibleQuote{我们没有君王，只有凯撒。} 他们拒绝了羔羊，选择了狐狸；他们理应成了狐狸的分。\par

16. \Quote{真正的君王}是这样写的，因为他们选择了一只狐狸，而不要真正的君王。\Quote{真正的君王：} 就是那真正的君王，祂受苦时，祂的称号被写在了上面。因为比拉多写了一个牌子，放在祂头上：\BibleQuote{犹太人的君王}，用希伯来、希腊和拉丁三种文字，好让所有经过的人都能读到君王的荣耀，以及犹太人自己的耻辱，因为他们拒绝了真正的君王，选择了狐狸凯撒。\Quote{真正的君王必因天主喜乐。} 他们成了狐狸的分……\Quote{说非义之言的口已被封住。} 现在没有人敢公开反对基督，现在所有人都敬畏基督。\Quote{因为说非义之言的口已被封住。} 当羔羊在软弱中时，连狐狸也敢对羔羊发怒。但犹大支派的狮子得胜了，\BibleRef{默 5:5} 狐狸便沉默了。\par

\SourceRecord{Translated by J.E. Tweed.；Nicene and Post-Nicene Fathers, First Series；Vol. 8.；Edited by Philip Schaff.；Buffalo, NY: Christian Literature Publishing Co.,；1888.；<http://www.newadvent.org/fathers/1801063.htm>.}

% structure: p0001 source=h1 target=section reason=page-title

\section{圣咏第六十四篇释义}

虽然这圣咏主要论及主的受难，但殉道者若非先看见那位首先受苦的祂，也不能坚强；若非盼望在复活中看见祂所彰显的，他们也不会忍受像祂所受的苦难。但你们的圣善知道，我们的头是我们的主耶稣基督，凡紧附于祂的，都是祂的肢体……。谁也不要以为，现今我们不在苦难的考验中。你们常听到这事：从前那整个教会仿佛一同被击打，如今却通过个人受到试探。魔鬼确实被束缚，以致不能为所欲为，不能任意而行；然而，它被允许试探，只要有利于人进步。对我们而言，没有试探是不利的；我们不应恳求天主使我们不受试探，而应求我们不被\BibleQuote{引入诱惑}。\BibleRef{玛 6:13}\par

因此，我们也当这样说：\BibleQuote{天主，求你俯听我忧闷时的祈祷，从仇敌的恐吓中拯救我的灵魂}\BibleRef{咏 63:1}。仇敌曾向殉道者发怒；基督身体的这个祈祷是在求什么？它是在求脱离仇敌，求仇敌无权杀害他们。难道因为他们被杀害，祈祷就没有蒙垂听？难道天主抛弃了祂怀着痛悔之心的仆人，轻视了仰望祂的人？绝不可能。因为\BibleQuote{谁呼求了天主而遭受抛弃？谁信赖了祂而被遗弃？}\BibleRef{德 2:10}。所以他们的祈祷蒙垂听，他们被杀，却仍从仇敌那里得了拯救。另一些人因害怕而妥协，保全了性命，却反被仇敌吞噬。被杀的反得拯救，活着的反被吞没。因此也有那感恩的呼声：\Quote{活活地吞了我们。}……因此，殉道者的声音祈求：\BibleQuote{从仇敌的恐嚇中拯救我的灵魂}\BibleRef{咏 63:1}，不是求仇敌不杀我，而是求我不怕那杀人的仇敌。因为仆人在圣咏中祈求的，正是主刚才在福音中所命令的。主刚才命令什么？\BibleQuote{你们不要害怕那杀害肉体，而不能杀害灵魂的；但更要害怕那能使灵魂和肉身都陷入地狱中的。}\BibleRef{玛 10:28}祂又重复说：\BibleQuote{我告诉你们，你们要害怕祂。}\BibleRef{路 12:5}谁是那杀害肉体的？就是仇敌。主命令什么？不要怕他们。因此，祈求祂赐予祂所命令的。\Quote{从仇敌的恐嚇中拯救我的灵魂。}拯救我脱离对仇敌的恐惧，使我臣服于对祢的敬畏。我不怕那杀身体的，却怕那有权柄将灵魂和肉身都投入地狱的。因为我并非渴望没有恐惧，而是摆脱对仇敌的恐惧，成为上主敬畏的仆人。\par

\BibleQuote{你曾保护我脱离了恶人的集会，脱离了作孽者的群众}\BibleRef{咏 63:2}。现在让我们注视我们的头基督本身。许多殉道者遭受了类似的事，但没有什么比殉道者之首更显耀；在他内，我们当看见他们所经历的一切。祂被保护脱离恶人的集会，是借着天主保护自己，是圣子和祂所携带的人性保护祂的肉体：因为祂是人子，也是天主圣子；按天主的形体是天主子，按仆人的形体是人子；祂有权柄舍弃性命，也有权柄再取回来。\BibleRef{若 10:18}仇敌能对祂做什么？他们杀害了肉体，却不能杀害灵魂。看清楚了。因此，主仅仅用言语劝勉殉道者是不够的，除非祂以榜样来加强劝勉。你们知道恶毒的犹太人如何聚集，作孽者的群众如何众多。那是什么罪孽？就是他们愿意杀害主耶稣基督的罪孽。\BibleQuote{我向你们显示了许多善事，你们为了哪一件要杀死我？}\BibleRef{若 10:32}祂忍受了他们所有的软弱，医治了他们所有的病人，宣讲天国，不缄默指责他们的恶习，为使他们厌恶这些恶习，而非厌恶那医治他们的医生；然而他们对于祂所有的救治忘恩负义，像发高烧而胡言乱语的人对医生狂怒，图谋杀害那来医治他们的一位，仿佛由此可以证明祂是否真的是一个会死的人，或者是超越凡人且不允许自己被杀的一位。我们在\BibleBook{}{智慧篇}中觉察到这些人的话：\Quote{我们用最可耻的死刑定祂的罪，且要考验祂，因为在祂的言论中必有关怀；如果祂真是天主子，愿天主解救祂。}让我们看看发生了什么事。\par

\BibleQuote{因为他们磨快了自己的舌头有如刀剑}\BibleRef{咏 63:3}。另一篇圣咏也说：\Quote{人子们的牙齿是武器和箭矢，他们的舌头是锋利的刀剑。}不要让犹太人说：我们没有杀害基督。因为他们正是为了这个目的把祂交给判官比拉多，好让自己表面上免于祂死的罪责……但如果那因不情愿而行了此事的人有罪，那迫使他去做的人岂不无辜？绝不然。但那判了祂死刑，并下令把祂钉在十字架上的人，在某种程度上也杀了祂；你们犹太人，也杀了祂。你们用什么杀的？用舌头的刀剑：因为你们磨快了舌头。你们什么时候击打的？除了你们喊叫\BibleQuote{钉死他，钉死他}\BibleRef{路 23:21}的时候，还能是什么时候？\par

5. 但我们不可忽略心中所想起的事，免得有人因阅读圣经而感到不安。一位圣史说主在第六时辰被钉十字架\BibleRef{若 19:14}，另一位说在第三时辰\BibleRef{谷 15:25}；除非我们理解这一点，否则便会不安。当第六时辰开始时，比拉多据说已坐上审判席；实际上，当主被高举在木头上时，正是第六时辰。但另一位圣史着眼于犹太人的心思，他们如何希望自己显得对主的死无罪，而通过他的记述证明他们有罪，说主在第三时辰被钉十字架。但考虑到历史的一切情况，在主受比拉多控告、为要钉死祂之前，可能发生了多少事？我们发现，当他们喊叫\Quote{钉死，钉死}时，可能是第三时辰。因此，更真实地说，他们是在喊叫的时刻杀了祂。官吏的差役在第六时辰钉死了祂，违法者在第三时辰喊叫：那些人在第六时辰用手所做的，这些人在第三时辰用舌头做了。那些喊着发怒的比那些顺从服侍的更有罪。这就是犹太人全部的精明，这就是他们视为大事的东西：让我们杀祂，又不杀祂；如此杀祂，以致我们不被认为杀了祂。\par

6. \Quote{他们弯了弓，一件苦毒的事，为要在暗中射那无玷的} \BibleRef{咏 63:4}。弓称为埋伏。因为用剑肉搏的人公开作战；射箭的人欺骗以便击中。因为箭击中之前，人未能预见它要来伤人。但人心的埋伏能逃脱谁呢？难道能逃脱我们的主耶稣基督吗？祂不需要任何人向祂作证关于人，\Quote{因为祂自己知道人里面有什么}，\BibleRef{若 2:25}正如圣史所证。然而，让我们听听他们，并观看他们的行为，仿佛主不知道他们的谋划。他用的措辞\Quote{他们弯了弓}等同于\Quote{暗中}：仿佛他们用埋伏欺骗。因为你知道他们用何种诡计做这事：他们用钱贿赂亲近祂的门徒，好让祂被交给他们\BibleRef{玛 26:14}，他们找来假证人；他们用何等的埋伏和诡计行事，\Quote{为要在暗中射那无玷的}。极大的不义！看，从暗处射来一枝箭，击中了无玷的那一位，祂甚至没有斑点可被箭穿透。祂确实是只无玷的羔羊，完全无玷，永远无玷；不是那斑点被移除的，而是从未染上任何斑点的。因为祂赦免罪过使许多人成为无玷的，而祂自己因无罪而无玷。\Quote{他们突然射祂，并不惧怕}。刚硬的心啊，竟想杀死一位曾复活死人的那一位！\Quote{突然}：即阴险地，仿佛出其不意，仿佛未曾预见。因为主似乎像一个不知道的人，在不知道的人中间，祂不知道什么不知道、什么知道：是的，不知道祂什么都不知道，而且祂知道一切，并为此而来，好让他们做他们以为凭自己能力所做的事。\par

7. \Quote{他们坚定了邪恶的言论} \BibleRef{咏 63:5}。行了如此伟大的奇迹，他们不被感动，他们坚持邪恶言论的计谋。祂被交给审判官：审判官战栗，而将祂交给审判官的他们不战栗；权力战栗，而凶暴不战栗；他要洗手，他们却玷污自己的舌头。但为何如此？\Quote{他们坚定了邪恶的言论}。比拉多做了多少事，多少事为了阻止他们！他说了什么？他做了什么？但\Quote{他们坚定了邪恶的言论：钉死，钉死}。\BibleRef{路 23:21}重复就是\Quote{邪恶言论}的坚定。让我们看看他们如何\Quote{坚定了邪恶的言论}。\Quote{我要钉死你们的王吗？}他们说：\Quote{除了凯撒，我们没有王}。\BibleRef{若 19:15}祂将天主子作为王献出；他们投靠一个人：他们配得拥有一个，而没有另一位。\Quote{我在这个人身上查不出什么该死的事}，审判官说。而那些\Quote{坚定了邪恶的言论}的人说：\Quote{祂的血归在我们和我们子孙身上}。\BibleRef{玛 27:25}\Quote{他们坚定了邪恶的言论}，不是对主，而是对\Quote{自己}。因为当他们说\Quote{归在我们和我们子孙身上}时，如何不是对自己呢？因此他们所坚定的，是对自己坚定的：因为同样的声音在别处说：\Quote{他们在我面前挖了一个坑，却掉进了里面}。死亡没有杀死主，而是祂杀死了死亡；但不义杀死了他们，因为他们不愿杀死不义……\par

8. \Quote{他们讲述，为要隐藏陷阱：他们说：谁看见他们？} \BibleRef{咏 63:5}。他们以为他们会逃脱祂，他们正在杀害的那一位；他们会逃脱天主。看，假设基督是一个人，像其余人一样，不知道为祂所设的计谋：难道天主也不知道吗？人心啊！你为何对自己说：谁看见我？当祂看见你时，祂已造了你？\Quote{他们说：谁看见他们？}天主看见了，基督也看见了：因为基督也是天主。但他们为何认为祂看不见？听听后面的话。\par

9. \Quote{他们搜寻了不义，他们失败了，搜寻着搜寻} \BibleRef{咏 63:6}：即致命的和敏锐的计谋。不要让祂被我们出卖，而是被祂的门徒；不要让祂被我们杀死，而是被审判官；让我们做一切，却似乎什么也没做……\par

10. 但他们遭遇了什么？\Quote{他们寻求的寻求失败了。}为何？因为他说：\Quote{谁看见他们？}也就是说，没有人看见他们。他们这样说，他们在自己中间这样想，以为没有人看见他们。请看邪恶的灵魂遭遇什么：它离开真理之光，因自身不见天主，便以为自己不被天主看见……\par

11. 接着是什么？\Quote{将有一人并一颗深邃的心走近。}他们说：谁看见我们？他们寻求的寻求失败了，邪恶的计谋失败了。有一个人走近了这些计谋，祂让自己被当作一个人抓住。因为若非祂是人，就不会被抓住，若非祂是人，就不会被看见，若非祂是人，就不会被击打，若非祂是人，就不会被钉十字架或死亡。因此，有一个人走近了所有这些苦难，这些苦难在祂内若祂不是人则毫无功效。但若祂不是人，就无人可得救赎。有一人走近，\Quote{并一颗深邃的心，}即一颗隐秘的心：在人面前显示为人，在内里保持天主性；隐藏了\Quote{天主的形像，}在其中祂与圣父平等（见：\BibleRef{斐 2:6}），并显示了仆人的形像，在其中祂小于圣父。因为祂亲自说了两者：但有一件事是祂以天主的形像说的，另一件事是以仆人的形像说的。祂以天主的形像说：\Quote{我与父原是一体。}（见：\BibleRef{若 10:30}）祂以仆人的形像说：\Quote{父比我大。}（见：\BibleRef{若 14:28}）那么，以天主的形像，祂说\Quote{我与父原是一体}？……\par

12. \BibleQuote{他们的击打成了婴儿的箭矢}（见：\BibleRef{咏 63:7}）。野蛮何在？那狮子般的咆哮何在？那人民咆哮说：\Quote{钉死，钉死}？那些弯弓埋伏的人何在？难道\Quote{他们的击打不是成了婴儿的箭矢}吗？你们知道婴儿用细竹竿给自己做箭矢的方式。他们击打什么？从何处击打？是什么手？是什么武器？是什么臂膀？是什么肢体？\par

13. \BibleQuote{他们的舌头在他们身上变得软弱}（见：\BibleRef{咏 63:8}）。让他们现在磨利他们的舌头像剑一样，让他们为自己确认恶毒的言论。他们理当为自己确认了这一点，因为\Quote{他们的舌头在他们身上变得软弱。}这能对抗天主强有力吗？他说：\Quote{不义向自己说了谎}；\Quote{他们的舌头在他们身上变得软弱。}看，主已复活了，那被杀的……你对那从十字架上没有下来、却从坟墓中复活的那一位作何感想？他们因此成就了什么？即使主没有复活，他们又能成就什么，除了殉道者的迫害者也成就的？因为殉道者尚未复活，然而他们并未成就任何事；我们如今正在庆祝他们尚未复活的诞辰。他们狂怒的疯狂何在？他们把他们那些寻求带到何处？在那些寻求中他们失败了，以致甚至当主已死并被埋葬时，他们仍在坟墓设置守卫？因为他们对比拉多说：\Quote{那个骗子}；主耶稣基督被称为此名，为安慰祂的仆人当被称为骗子时；因此他们对比拉多说：\Quote{那个骗子还活着时曾说，三天后我要复活}（见：\BibleRef{玛 27:63}）……他们在坟墓设置士兵守卫。大地震动时，主复活了：在坟墓周围行了如此奇迹，以至于那些来守卫的士兵本身也成为见证人，如果他们选择说实话的话；但同样贪婪曾俘虏了一个门徒，基督的同伴，也俘虏了那个守卫坟墓的士兵。他们说：\Quote{我们给你钱}（见：\BibleRef{玛 28:12}），你们却说，当你们自己睡觉时，祂的门徒来了，把祂偷走了……你们引用睡觉的证人：确实你们自己睡着了，在寻求这样的计谋时失败了。如果他们睡着了，他们能看见什么？如果什么也没看见，他们如何是证人？但是\Quote{他们寻求的寻求失败了}：失去了天主的光，在他们计划的完成中失败了：当他们所愿的无法完成时，他们确实失败了。这是为什么？因为\Quote{将有一人并一颗深邃的心走近，而天主被高举了。}……\par

14. \BibleQuote{人人都敬畏}（见：\BibleRef{咏 63:9}）。不敬畏的人，甚至不算是人。\Quote{人人都敬畏；}即每个人都运用理智察觉所发生的事。因此那些不敬畏的，宁可称为牲畜，宁可称为野蛮残忍的野兽。那人民至今尚是跳跃咆哮的狮子。但实际上人人都敬畏：即那些将要相信、在将来的审判前战栗的人。\Quote{人人都敬畏：他们传扬了天主的作为。}……\Quote{人人都敬畏了：他们传扬了天主的作为，并察觉了祂的所做。}什么是\Quote{他们察觉了祂的所做}？哦，主耶稣基督，是祢默默无声，如同被牵去祭献的羊羔，在剪毛者前不开口（见：\BibleRef{依 53:7}），我们以为祢是受击打、在痛苦中（见：\BibleRef{依 53:4}），并知道如何担当软弱？（见：\BibleRef{依 53:3}）是祢隐藏了祢的美，哦，祢在人子前形貌俊美？是祢看似没有美丽和优雅？（见：\BibleRef{依 53:2}）祢在十字架上忍受了人们辱骂并说：\Quote{如果他是天主子，让他从十字架上下来。}（见：\BibleRef{玛 27:40}）……这件事，那些希望他从十字架上下来的人没有察觉；但当祂复活、受光荣升天后，他们察觉了天主的作为。\par

15. \Quote{义人必因主而喜乐} \BibleRef{咏 63:10} 。义人并不忧愁。因为当主被钉十字架时，门徒们忧愁；他们被忧愁压倒，悲伤地离去，以为失去了希望。祂复活了，即使向他们显现时，祂发现他们仍忧愁。祂蒙住了在路上行走的两个人的眼睛，使他们认不出祂，祂发现他们在叹息呻吟，祂就留住他们，直到为他们解释了圣经，并借着同一圣经指明，事情必须如此成就。因为祂在圣经中指明，主在三天后必须复活。\BibleRef{路 24:46} 如果祂从十字架上下来了，又怎能在第三天复活呢？……因此，让我们众人都在主内喜乐，让我们众人因信德成为\Emph{一个义人}，让我们众人共属一个身体，共有一个头，并让我们在主内，而非在自己内喜乐：因为我们的善不是我们自己对我们自己，而是那创造了我们的那一位。祂本身就是使我们喜乐的善。愿没有人因自己喜乐，没有人依靠自己，没有人对自己绝望：愿没有人依靠任何人，人应当引导他成为自己希望的同伴，而非希望给予者。\par

16. 如今因为主已经复活，如今因为祂已升到天上，如今因为祂已表明有另一种生命，如今因为祂隐藏于内心深处的计谋显然并非空虚，因为那宝血正是为此而流，作为被赎者的赎价；如今因为一切明明可知，因为一切已被宣讲，因为一切已被相信，普天之下，\Quote{义人必因主而喜乐，并要仰望祂；所有心正的人都要得称赞。}……天主不悦纳你，你却悦纳自己，你有一颗乖僻弯曲的心：更糟的是，你竟想用你的心纠正天主的心，使祂做你所愿的，而你本该做祂所愿的。那么怎样？你要以你内心的败坏，使永远正直的天主之心变得弯曲吗？用天主的正直纠正你的心，岂不更好？难道你的主没有教导你这一点吗？我们刚才不就在谈论祂的苦难吗？当祂说：\Quote{我的心灵忧闷得要死}时，祂不是在承担你的软弱吗？\BibleRef{玛 26:38} 当祂说：\Quote{父啊！如果可以，请让这杯离开我}时，祂不是在祂自身内预表你吗？\BibleRef{玛 26:39} 因为圣父与圣子的心并非两个不同的心：而是在仆人的形像内，祂肩负了你的心，为要以祂的榜样教导它。现在看，烦恼仿佛找出了你的另一颗心，那心愿意那即将来临的事过去；但天主不愿意。天主不依从你的心，你要顺从天主的心。\par

17. 然后呢？如果\Quote{所有心正的人都要得称赞}，那么心歪的人就要被定罪。两件事摆在你面前，趁有时间选择吧……如果你成了心歪的人，那审判就会来到，天主做这一切事的理由都会显明：你在这生命中没有用天主的正直纠正你的心，没有为你自己预备右边，那里\Quote{所有心正的人都要得称赞}，你将站在左边，在那里你将要听见：\Quote{你们到那为魔鬼及其使者预备的永火里去。} \BibleRef{玛 25:41} 到那时还有时间纠正心灵吗？所以现在纠正吧，弟兄们，现在纠正吧。谁阻止呢？圣咏被咏唱，福音被宣读，读经者呼喊，讲道者呼喊；主是长久忍耐的；你犯罪，祂宽容；你仍犯罪，祂仍宽容，而你仍罪上加罪。天主忍耐要到几时呢？你也会发现天主是公义的。我们恐吓，因为我们害怕；教导我们不害怕，我们就不再恐吓。但天主教导我们害怕，比任何人教导我们不害怕更好……你出产谷粒，就等待谷仓；你生出荆棘，就等待火焰。但无论是谷仓还是火焰的时刻都还未到来：现在当预备，这样就不会害怕。我们这些说话的，和你们这些听我们说话的，都因基督的名而活着：要改变我们的计划，将邪恶的生活转变为善的生活，难道没有地方，没有时间吗？如果你愿意，今天不能做成吗？如果你愿意，现在不能做成吗？要行这事，需买什么，需寻求什么药材？要航行到哪个印度？预备哪艘船？看，我说话的时候，就改变心灵；那如此经常、如此长久被呼吁要成就的事，就成了；若不做成，则带来永罚。\par

\SourceRecord{Translated by J.E. Tweed.；Nicene and Post-Nicene Fathers, First Series；Vol. 8.；Edited by Philip Schaff.；Buffalo, NY: Christian Literature Publishing Co.,；1888.；<http://www.newadvent.org/fathers/1801064.htm>.}

% structure: p0001 source=h1 target=section reason=page-title

\section{\WorkTitle{圣咏第六十五篇释义}}

1. 在这篇圣咏的标题中，我们应当承认圣先知的神谕。标题写着：\Quote{交与乐官，达味圣咏，耶肋米亚和厄则克耳之歌，关于流徙之民开始出发之时。}我们的祖先在流徙巴比伦时期的情况，并非所有人都知道，只有那些勤于聆听或阅读圣经的人才了解。因为以色列被掳的百姓从耶路撒冷城被掳到巴比伦为奴。\BibleRef{列下 24:14} 但圣耶肋米亚预言，七十年后百姓将从被掳中回归，并重建他们因被敌人摧毁而哀悼的耶路撒冷城。那时，在巴比伦被掳的百姓中有先知，其中包括厄则克耳先知。然而，那百姓一直等待，直到依照耶肋米亚的预言，七十年期满。当七十年届满时，被拆毁的圣殿得以重建；大部分百姓从被掳中回归。然而，正如宗徒所说：\BibleQuote{这些事发生在他们身上，作为预像，并且记载下来，是为了劝戒我们这些生活在末世的人。}\BibleRef{格前 10:11} 因此，我们也应当首先认识我们的被掳，然后认识我们的解救；我们应当认识我们被掳的巴比伦，以及我们叹息渴望回归的耶路撒冷。因为这两座城，按字面意义，确实是两座城。但昔日的耶路撒冷如今已不被犹太人居住。因为主被钉十字架后，他们遭受了严厉的惩罚，从那块他们以不敬的放纵疯狂攻击自己医生的地方被连根拔起，分散到万民中，那片土地已赐给基督徒。主曾对他们说的预言便应验了：\BibleQuote{因此，天主的国要从你们中夺去，并要交给一个结果实的民族。}\BibleRef{玛 21:43} 但当时耶路撒冷的统治者们看见许多群众跟随主，宣讲天国，行了许多奇事，便说：\BibleQuote{如果我们让他这样，众人都会信从他，罗马人必要来，夺去我们的地方和民族。}\BibleRef{若 11:48} 他们为了不失去地方，杀了主；但正因杀了主，他们反而失去了地方。因此，那座地上的耶路撒冷是天上永恒之城的一个预像。但当所预表的事更清晰地被宣讲时，预像的阴影便被废弃了；因此，如今那里已没有圣殿，那圣殿原是为未来主的身体的肖像而建造的。我们有了光明，阴影已经过去；然而，我们仍处于某种被掳之中：\BibleQuote{我们几时住在肉身内，就是离主远行。}\BibleRef{格后 5:6}\par

2. 请看这两座城的名字：巴比伦和耶路撒冷。巴比伦意为混乱，耶路撒冷意为和平的异像。现在请观察混乱之城，以便你领悟和平的异像；使你忍受前者，叹息向往后者。如何区分这两座城？我们现在能将它们彼此分开吗？它们混合在一起，从人类之初起就混合着一直延续到世界终结。耶路撒冷始于亚伯尔，巴比伦始于加音；后来才建造了实际的城。耶路撒冷建在耶步斯人的土地上，起初称为耶步斯（\BibleRef{苏 18:28}），当天主的子民从埃及被拯救出来，领进应许之地时，耶步斯人被驱逐出去。巴比伦则建于波斯最内陆的区域，它在诸民族中长期昂首。这两座城在特定时期被建造，以显示两座城从古时便开始，并在这世界中持续到终末，但终将分离。那么，我们现在如何区分它们，既然它们混合在一起？那时主会显示，当他把一些人放在右边，另一些人放在左边时。右边将是耶路撒冷，左边将是巴比伦……两种爱构成了这两座城：对天主的爱构成耶路撒冷，对世界的爱构成巴比伦。因此，让每个人反问自己爱什么，他就会发现自己属于哪座城。如果他发现自己属于巴比伦，让他根除贪欲，种植爱德；如果他发现自己属于耶路撒冷，让他忍受被掳，盼望自由。……所以现在，弟兄们，让我们聆听、歌颂，并渴望我们身为公民的那座城。我们所歌唱的喜乐是什么？我们心中如何重新形成对我们城市的爱，这爱因长期流徙而被遗忘？但我们的父亲从那里给我们寄来了书信，天主赐给了我们圣经，藉着这些书信，应在我们心中激起回归的渴望：因为我们因热爱流徙，将脸转向敌人，而将背朝向祖国。那么，这里歌唱的是什么呢？\par

3. \BibleQuote{天主，歌咏于熙雍，理应归于祢}\BibleRef{咏 64:1}。那祖国就是\Place{熙雍}：\Place{耶路撒冷}与熙雍是同一的；你们应当知道这名字的解释。正如耶路撒冷被解释为\Quote{和平的异象}，同样，熙雍是\Emph{注视}，即异象与观想。一些伟大而不可言说的景象被应许给我们：而这就是天主自己，他建造了这城。这城美丽而优雅，但它的建造者是何等更加美丽！他说：\BibleQuote{天主，歌咏理应归于祢}。但在哪里？\BibleQuote{在熙雍}，在\Place{巴比伦}则不相宜。因为当一个人开始更新，他已经在心中于耶路撒冷歌唱，与宗徒一同说：\BibleQuote{我们的家乡是在天上}\BibleRef{斐 3:20}。因为\BibleQuote{我们固然在血肉中行事}他说，\BibleQuote{却不是按照血肉而作战}\BibleRef{格后 10:3}。我们已在渴望中到达那里，已将希望如锚一般预先投入那土地，以免在这海上颠簸而遭船难。同样，正如一艘下了锚的船，我们虽可正确地说它已到达陆地，因它仍在摇晃，但它已在某种程度上安全地抵达陆地，抵御了狂风暴雨；同样，面对这寄居生活中的诱惑，我们的希望因扎根于那\Place{耶路撒冷}城而使我们不被冲撞在岩石上。因此，凡按此希望歌唱的，就在那城中歌唱：所以让他说：\BibleQuote{天主，歌咏于熙雍，理应归于祢}……\par

4. \BibleQuote{在\Place{耶路撒冷}，也必有人向祢偿还誓愿。}在此我们立誓，而我们将来在那里偿还，这是一件好事。但那些在此立誓而不偿还的人是谁呢？就是那些在他们所立誓的事情上没有坚忍到底\BibleRef{玛 24:13}的人。因此另一篇圣咏说：\BibleQuote{你们当向主你们的天主许愿，并偿还；}又说：\BibleQuote{在\Place{耶路撒冷}，有人向祢偿还。}因为在那里我们要成为完整的，即义人复活中的完整：我们的全部誓愿将得到偿还，不仅是灵魂，还有肉身本身，不再朽坏，因为不再在\Place{巴比伦}，而是成为属天的、改变了的身体。应许的是怎样的改变呢？宗徒说：\BibleQuote{我们众人都要复活，但不都改变……死亡啊，你的刺在哪里？}因为现在当使用\OriginalTerm{心思的初果}{first-fruits of the mind}时，从而生出对\Place{耶路撒冷}的渴望，可朽坏血肉的许多事物攻击我们，但当死亡被胜利吞没时，它们就不再攻击。和平将得胜，战争将结束。但当和平得胜时，那被称为和平异象的城也将得胜。因此，死亡方面将不再有争战。如今我们是在何等大的死亡中争战！因为肉欲的快乐由此而来，甚至非法地给我们暗示许多事物：我们不予同意，但尽管如此，在不同意中我们仍在争战……\par

5. 他说：\BibleQuote{请听，我的祈祷，一切血肉都要来到祢面前}\BibleRef{咏 64:2}。我们又听到主说，那赐给祂的\BibleQuote{权柄于一切血肉}\BibleRef{若 17:2}。因此，当这么说时，那君王甚至现在就开始显现：\BibleQuote{一切血肉都要来到祢面前}。他说：\BibleQuote{到祢面前}，\BibleQuote{一切血肉都要来到}。为什么\Emph{一切}血肉要来到祂面前？因为祂取了血肉。一切血肉要到哪里去？祂取了血肉的\OriginalTerm{初果}{first-fruits}，取自\Person{童贞玛利亚}的胎；既然初果已被取去，其余的要跟随，以便\OriginalTerm{全燔祭}{holocaust}得以完成。那么，\Emph{一切血肉}指什么？指一切人。而一切人又指什么？难道所有人都被预言将要相信基督吗？难道不是也有许多不敬虔的人被预言将要受审判吗？难道不是每天都有不信的人死于自己的不信吗？那么我们怎样理解\BibleQuote{一切血肉都要来到祢面前}呢？他用\Emph{一切血肉}来表示\Quote{各种血肉}：从各种血肉中，他们都要来到祢面前。各种血肉是什么意思？难道穷人来了，富人就没有来？难道谦卑的人来了，高傲的人就没有来？难道没有学问的人来了，有学问的人就没有来？难道男人来了，女人就没有来？难道主人来了，仆人就没有来？难道老人来了，年轻人就没有来？或者年轻人来了，青少年就没有来？或者青少年来了，男孩就没有来？或者男孩来了，婴孩就没有被带来？简而言之，难道犹太人来了（因为宗徒是从他们中来的，起初背叛而后相信的成千上万的人\BibleRef{宗 2:41}），而希腊人就没有来？或者希腊人来了，罗马人就没有来？或者罗马人来了，蛮族就没有来？谁能数算所有来到祂面前的民族呢？对祂曾说过：\BibleQuote{一切血肉都要来到祢面前}。\par

6. \Quote{不义之人的言语压过我们，我们的不义祢将宽免}\BibleRef{咏 64:3}。每一个人，无论他生在何处，都学习那地、那区或那城的语言，并习惯那地的风俗和生活。一个生在外教人中的男孩，当他的父母已建议崇拜石头时，他该如何避免崇拜一块石头？他从他们那里听到的最初话语，那错误随着乳汁被他吸入；因为说话的是长者，而学话的男孩是个婴儿，这小家伙除了跟随长者的权威、认为他们所推荐的便是好的之外，还能做什么？因此那些后来归向基督的民族，将父母的邪恶铭刻在心，现在说出耶肋米亚先知自己所说的话：\Quote{我们的祖先真正崇拜了虚谎，那是无益的虚幻}\BibleRef{耶 16:19}——我说，当他们现在这样说时，他们弃绝了不义父母的意见和亵渎。……有些人教导邪恶之事，将我们引走；他们使我们成为巴比伦的公民；我们离开了造物主，崇拜了受造物；离开了创造我们的那一位，崇拜了我们自己所制造的。因为\Quote{不义之人的言语压过我们}，然而他们并未压碎我们。为何？\Quote{我们的不敬祢将宽免}，这话只对某位司铎说，他献上某种祭品，使不敬得以赎免和宽息。因为当天主对不敬转为仁慈时，不敬便被称为得到宽免。天主对不敬转为仁慈是什么意思？就是祂变为宽恕，赐予赦免。但为获得天主的赦免，须藉某种牺牲来行赎罪。因此，那一位我们的司铎，由主天主派遣，自我们中取了祂能向主奉献的；我们说的正是那从童贞母胎中取来的血肉初果。祂向天主献上了这全燔祭。祂在十字架上伸展双手，为能说：\Quote{愿我的祈祷如馨香上达于祢面前，愿我举手作为晚祭。}如你们所知，主在黄昏时分被钉在十字架上：\BibleRef{玛 27:46}，我们的不敬得到了宽免；否则它们早已吞没了我们：不义之人的言语压过了我们；朱庇特、萨图尔努斯和墨丘利的宣讲者曾使我们迷失：\Quote{不虔之人的言语压过了我们。}但祢要做什么？\Quote{我们的不敬祢将宽免。}祢是司铎，祢是祭品；祢是奉献者，祢是所献之物。\BibleRef{希 9:7}……\par

7. \Quote{你所拣选并接纳于你的人，真是有福。}\BibleRef{咏 64:4}谁是那被祂拣选并接纳的人？难道是某个人被我们的救主耶稣基督拣选，还是祂自己按肉身（因为祂是人）被拣选并接纳？……或者，难道不是基督自己接纳了一位有福者，而祂所接纳的这位并不以复数形式提及，而是以单数形式？因为祂接纳了一个人，因为祂接纳了合一。祂没有接纳分裂，没有接纳异端：他们自成一堆，没有一个能被接纳到祂那里。然而，那些持守基督的纽带并成为祂肢体的人，在某种程度上构成了一个人，宗徒论到这人说：\Quote{直到我们众人对天主子达成一致的认识，成为完人，达到基督圆满年龄的尺度。}\BibleRef{弗 4:13}因此，一个人被接纳到祂那里，其首是基督，因为\Quote{男人的首是基督。}\BibleRef{格前 11:3}同样，那位有福的人就是\Quote{未随从不虔之人的计谋}等等所描述的那一位；他就是被接纳的人。祂并不离弃我们，我们在祂的肢体中，受同一首治理，同属一圣神而生活，渴慕同一祖国。……祂将赐给我们什么？他说：\Quote{祂将居住在}——\Quote{祢的庭院中。}耶路撒冷，就是那些开始从巴比伦出发的人所歌唱的：\Quote{祂将居住在祢的庭院中：我们必因祢房屋的美物而心满意足。}什么是天主房屋的美物？弟兄们，让我们设想一座富户的房屋，它拥簇着何等众多的美物，装饰得何其丰盛，有多少金银器皿；仆役队伍何等庞大，马匹牲畜何其多；一句话，房屋本身以图画、大理石、天花板、柱子、壁龛、房间使我们愉悦。所有这些确实是人渴望的对象，但它们仍属于巴比伦的混乱。耶路撒冷的公民啊，断绝一切这样的渴望，断绝它们！如果你要回去，别让俘虏生活取悦你。但你若已开始出发，不要回头，不要在路上徘徊。仍有敌人不缺少向你推荐俘虏和寄居生活。愿不虔之人的言语不再压制你。你该渴慕天主的房屋，渴慕那房屋的美物；但不要渴慕你惯常在你家中、邻居家中、或恩主家中渴慕的那些东西……\par

8. \Quote{你的圣殿令人惊异，在于义德} \BibleRef{咏 64:5}。这是那房屋的福祉。祂没有说：你的圣殿令人惊异，在于柱子，在于大理石，在于镀金的天花板；而是说\Quote{令人惊异，在于义德}。若你没有眼睛来看大理石和金子，你心中有眼睛能看见义德之美。若义德中没有美，为何人们爱一位公义的老人？他身体上有什么悦目之处？四肢弯曲，额头皱纹，头发灰白，老态龙钟，满是抱怨。但或许这衰弱的老人不悦你的眼，却悦你的耳：用何言语？用何歌声？即便他年轻时唱得好，如今也因年迈丧失了一切。或许他牙齿脱落，难以完整发音，他的言语之声能悦你的耳吗？然而，若他是公义的，不贪图他人财物，将自己所有分给穷人，给出良好的忠告，公正地判断，持守圆满的信德，并愿为信德献出甚至那残破的肢体——因为许多殉道者正是老人——我们为何爱他？用肉体的眼睛我们在他身上看到什么好事？什么也没有。因此，在义德中有一种美，我们用心灵之眼看见，并爱慕，且燃起热情：人们在那些殉道者身上找到了多少可爱之处，尽管野兽撕裂了他们的肢体！当鲜血染遍各处，当脏腑从怪兽的牙齿间涌出，眼睛岂非只有可怖之物？有什么可爱之处，除非在那肢体撕裂的可怕景象中，义德之美仍是完整的？这是天主之家的福祉，你当以这些准备自己得以饱足……\Quote{饥渴慕义的人是有福的，因为他们要得饱饫。} \BibleRef{玛 5:6} \Quote{你的圣殿令人惊异，在于义德。}弟兄们，不要以为这圣殿是别的什么，而就是你们自己。你们爱慕义德，你们就是天主的圣殿。\par

9. \Quote{天主，我们的救主，求祢俯听我们} \BibleRef{咏 64:5}。如今祂已揭示祂所称为天主的是谁。\Quote{救主}尤其指主耶稣基督。现在更明显地显明了祂所说的\Quote{凡有血肉的，都要求到祢面前}。那被接入天主圣殿中的那个人，既是多人，也是一人。以一人之身份祂说过：\Quote{天主，求祢俯听，即我的饥饿}；因这同一人由多人组成，如今祂说：\Quote{天主，我们的救主，求祢俯听我们。}如今更公开地宣讲祂：\Quote{天主，我们的救主，求祢俯听我们，祢是地极和远方海上万民的希望。}看，为何曾说\Quote{凡有血肉的，都要求到祢面前}。他们从四方而来。\Quote{地极的希望}，不是一角的希望，不是仅\Place{犹太}的希望，不是仅\Place{阿非利加}的希望，不是仅\Place{潘诺尼亚}的希望，不是东方或西方的希望；而是\Quote{地极的希望，以及远方海上}的希望。\Quote{远方海上}，因为是在海上，所以是远方。因为\Quote{海}象征这世界，咸涩如苦，风暴动荡；其中悖逆和堕落欲望的人，如同彼此吞噬的鱼。看这邪恶之海、苦涩之海、波涛汹涌之海，看它充满了怎样的人。谁不是因他人死亡而贪求遗产？谁不是因他人损失而贪求收益？多少人因他人跌倒而希望高升？多少人为了购买，希望他人变卖货物？他们如何互相压榨，能力强的如何吞噬！当一条鱼吞噬了另一条，更大的又吞噬它……因为那些邪恶的鱼被网住，他们说自己不能忍受；但他们自己变得比他们所说不能忍受的人更邪恶。因为那网确实网住了好鱼和坏鱼。主说：\Quote{天国好比撒在海里的网，网罗各种鱼，网满了，人就把网拉上岸，坐在岸上，把好的收在器皿里，将坏的抛弃。} \BibleRef{玛 13:47} 祂指出什么是岸，什么是海的尽头。\Quote{天使要出来，把恶人从义人中分开，把他们扔在火窑里，那里要有哀号和切齿。}啊，你们这些在网中的\Place{耶路撒冷}居民，是好鱼；要忍受恶鱼，不要挣破网：你们与他们一同在海里，却不会一同在器皿里。因为\Quote{祂是地极的希望}，祂也是\Quote{远方海上}的希望。远方，因为也在海上。\par

10. \Quote{祂以大能预备群山} \BibleRef{咏 64:6}。不是靠它们自身的力量。因为祂预备了伟大的宣讲者，并称他们为群山；他们自卑，在祂内被高举。\Quote{祂以大能预备群山。}那些群山中的一位说过什么？\Quote{我们自己也断定是必死的，为叫我们不要依靠自己，而依靠那使死人复活的天主。} \BibleRef{格后 1:9} 那依靠自己、不依靠基督的，不属于祂以自己大能所预备的群山。\Quote{祂以大能预备群山，以能力束腰。}\Quote{能力}我明白；\Quote{以能力束腰}是什么意思？他们以基督为中间，\Quote{束腰}指祂，即从各方面被围绕。我们都共同拥有祂，因此祂在中间；所有信祂的人都围绕祂：因我们的信德不是出于自己的力量，而是出于祂的能力；因此祂以自己能力束腰，而非我们的力量。\par

11. \BibleQuote{你搅动了海的深处} \BibleRef{咏 64:7}。祂已行了这事：可见祂所行的。因祂以大能预备了山，派遣它们去宣讲；祂被信徒以大能束腰；海被搅动，世界被搅动，并开始迫害祂的圣人。\Quote{以能力束腰：你搅动了海的深处}。祂没有说\Quote{你搅动了海}，而是海的深处。海的深处是不虔敬之人的心。正如从深处一切都被更彻底地搅动，且深处坚固一切：凡由舌头、手、各种势力为迫害教会而出的一切，都是从深处发出的。若心中没有不义的根，这一切就不会针对基督发出。祂搅动了深处，或许是为使深处也被倒空：因为对于某些恶人，祂倒空了海的深处，使海成为荒芜之地。另一篇圣咏这样说：\Quote{祂使海变为旱地}。所有不虔敬的外邦人，凡相信的，曾是海，已成了陆地；起初是咸涩的波浪，后来结出正义的果实。\BibleQuote{你搅动海的深处：海浪的响声谁能忍受？} \Quote{谁能忍受}是什么意思？谁人能忍受海浪的响声，即世间高位权势的命令？但如何忍受呢？因为祂以大能预备了山。因此，祂所说的\Quote{谁能}忍受，如此说：我们自己凭自己本不能忍受那些迫害，除非祂赐予力量。\par

12. \BibleQuote{万民将被扰乱} \BibleRef{咏 64:8}。起初它们将被扰乱：但那些以基督的大能预备的山，它们也被扰乱吗？海被搅动，冲击着山：海破碎了，山却屹立不动。\BibleQuote{万民将被扰乱，凡人都要惧怕}。看，如今众人都惧怕：那些先前被搅乱的人，如今都惧怕了。基督徒未曾惧怕，如今基督徒被惧怕。所有从前迫害的人，如今都惧怕了。因为那位以大能束腰者已得胜，万民都来到祂面前，以致其余人因人数甚少而惧怕。\BibleQuote{凡住在地极的人，因祢的奇迹都要惧怕}。因为宗徒们行了奇迹，从此地极的人都惧怕并相信了。\BibleQuote{祢必使早晨与晚上的出离成为喜乐}，即祢使之喜悦。在此生中，已有什么应许给我们？有早晨的出离，有晚上的出离。早晨象征世间的顺境，晚上象征世间的逆境……起初，当福音应许利益时，对你来说是早晨；但现在傍晚临近，你变得忧愁。但那赐给你早晨出离的，也必赐给你晚上的出离。正如你因主的光照而藐视世界的早晨，同样也因主的苦难而藐视世界的晚上，对你的灵魂说：这个人还能对我做什么，比我的主为我所受的更多？愿我持守正义，不认同不义。让他对肉体发怒，网罗必被打破，我将飞向我的主，祂对我说：\BibleQuote{不要害怕那杀害身体，却不能杀害灵魂的} \BibleRef{玛 10:28}。至于身体本身，祂也给了保障，说：\BibleQuote{你们连一根头发也不会失落} \BibleRef{路 21:18}。此处高妙地写下：\Quote{祢必使早晨与晚上的出离成为喜乐}。因为若你不以那出离本身为乐，你就不会努力从那里出去。你若不被救主的应许所吸引，就不会将头冲向应许的利益。同样，你若在受试探和恐吓者面前屈服，是因为你没有在那位在你之前受苦、为给你打开出路的主身上找到喜乐。\par

13. \BibleQuote{祢眷顾了大地，并使它沉醉} \BibleRef{咏 64:9}。祢如何灌醉了大地？\BibleQuote{祢的杯爵令人沉醉，何其荣耀！} \BibleQuote{祢眷顾了大地，并使它沉醉。} 祢派遣了祢的云彩，它们降下真理的宣讲，大地被灌醉了。\BibleQuote{祢丰富地增加了它，以使它富足。} 从何而来？\BibleQuote{天主的河充满了水。} 天主的河是什么？是天主的子民。最初的子民充满了水，用以浇灌其余的大地。请听祂应许水：\BibleQuote{谁若渴，到我这里来喝！信我的人，从他腹中要流出活水的江河} \BibleRef{若 7:37}。若是河流，也是一条河；因为就合一而言，许多成为一。许多教会而一个教会，许多信徒而一个基督的新娘：同样，许多河流而一条河。许多以色列人相信了，并充满了圣神；从那里他们被分散到万民中，开始宣讲真理，从充满水的天主之河，整个大地被浇灌。\BibleQuote{祢为他们预备了食物，因为这样是祢的预备。} 并非因为他们对祢有功，祢已赦免了他们的罪：他们的功过是恶的，但祢因祢的仁慈，\BibleQuote{因为这样是祢的预备，} 就这样\BibleQuote{祢为他们预备了食物。}\par

14. \BibleQuote{它的犁沟使祢沉醉}\BibleRef{咏 64:10}。因此，首先要开出犁沟以使它们沉醉：愿我们心硬的土壤被天主圣言之犁翻开。\BibleQuote{它的犁沟使祢沉醉：祢增多它的出产。}我们看见，他们相信，因他们的相信，其他人也相信，并因那些人而相信；一个人仅仅自己成为信徒，赢得另一个人，是不够的。种子也是如此增多：几粒种子撒下，田野便长起来。\BibleQuote{在它的水滴中，当它升起时，它要喜乐。}也就是说，或许在它长成大河之前，\BibleQuote{当它升起时，在它的水滴中}，即在其适于它的人中，\BibleQuote{它要喜乐。}因为对于那些尚是婴儿和软弱的人，圣事的一些部分是被点滴施予的，因为他们不能接受圆满的真理。请听他是如何向那些正在升起、即初升而能力尚小的婴儿点滴的：宗徒说：\BibleQuote{我无法对你们如同对属神的人讲话，只能如同对属血肉的人，如同对基督内的婴儿。}\BibleRef{格前 3:1}当他说\BibleQuote{对在基督内的婴儿}时，他称他们为已升起的人，但尚不宜接受那丰盛的智慧，关于这智慧他说：\BibleQuote{我们在成全者中讲论智慧。}\BibleRef{格前 2:6}愿它在自己的水滴中喜乐，当它正在升起并成长，当它强壮起来时，也要接受智慧：正如婴儿被喂以奶，然后变得适于吃干粮，然而起初那干粮正是它不适于吃的，因为奶是从干粮而来。\par

15. \BibleQuote{祢要祝福祢的恩惠之年的冠冕}\BibleRef{咏 64:11}。如今种子正在撒播，所种的在生长，也将有收获。如今仇敌已在种子之上撒下了莠子；恶者在善者中兴起，虚假的基督徒，有同样的叶子，却没有同样的果实。因为那些像麦子一样生长的东西被恰当地称为莠子，例如毒麦，例如野燕麦，以及所有初叶相同的东西。因此关于撒莠子，主这样说：\BibleQuote{有仇敌来了，在麦子中间撒上了莠子；}\BibleRef{玛 13:25}但他对麦子做了什么？麦子没有被莠子窒息，反而因忍耐莠子使麦子的果实增加。因为主亲自对某些想要拔除莠子的工人说：\BibleQuote{让两者一同长到收割的时候。}\BibleRef{玛 13:30}……征服魔鬼，你将有冠冕。\BibleQuote{祢要祝福祢的恩惠之年的冠冕。}他再次归因于天主的恩惠，免得有人夸耀自己的功绩。\BibleQuote{祢的平原必充满丰盈。}\par

16. \BibleQuote{旷野的极处将要肥沃，山丘要以欢跃环绕}\BibleRef{咏 64:12}。平原、山丘、旷野的极处，所指的也是人。平原，因为平等：由于平等，我说，公义的民族被称为平原。山丘，因为高举：因为天主使谦卑的人在自己内升高。旷野的极处是所有的民族。为何是旷野的极处？他们是被遗弃的，没有先知被派遣到他们那里：他们就像无人经过的旷野一样。天主圣言没有被派遣到外邦人：只有以色列民族才有先知宣讲。我们来到主前：在犹太人那同一民族中麦子相信了。因为那时他对门徒说：\BibleQuote{你们说，收割的时候还远；举目观看，田地已经发白，可以收割了。}\BibleRef{若 4:35}因此有第一次收割，在最后时期将有第二次收割。第一次收割是来自犹太人，因为先知被派遣到他们那里，宣告即将来临的救主。因此主对他的门徒说：\BibleQuote{看，田地已经发白，可以收割了：}\BibleRef{若 4:35}他说：\BibleQuote{别人劳苦，而你们进入了他们的劳苦。}\BibleRef{若 4:38}先知劳苦撒种，而你们带着镰刀进入了他们的劳苦。因此第一次收割完成了，并且从那里，用那时被净化的麦子，播种了整个世界；这样产生了另一场收割，在终结时将被收取。在第二次收割中，莠子被撒下，现在这里有了劳苦。正如在第一次收割中，先知劳苦直到主来临：同样在第二次收割中，宗徒劳苦，所有真理的宣讲者劳苦，直到最后主派遣他的天使来收割。之前，我说，那是一片旷野，\BibleQuote{但旷野的极处将要肥沃。}看，先知未曾发声的地方，先知的主被接受了，\BibleQuote{旷野的极处将要肥沃，山丘要以欢跃环绕。}\par

17. \BibleQuote{羊群中的公绵羊都披上}\BibleRef{咏 64:13}：必须理解为\Quote{以欢跃}。因为山丘以何欢跃环绕，羊群中的公绵羊也以同样的欢跃披上。公绵羊正是这些山丘。因为他们是山丘，由于更卓越的恩宠；是公绵羊，因为他们是羊群的领袖……\BibleQuote{他们要欢呼：}因此他们充满麦子，因为他们要欢呼。他们要欢呼什么？\BibleQuote{因为要唱一首赞美诗。}因为向天主欢呼是一回事，唱赞美诗是另一回事；欢呼不义是一回事，欢呼天主的赞美是另一回事。如果你欢呼以亵渎，你已生出荆棘；如果你欢呼以赞美诗，你便丰盈于麦子。\par

\SourceRecord{Translated by J.E. Tweed.；Nicene and Post-Nicene Fathers, First Series；Vol. 8.；Edited by Philip Schaff.；Buffalo, NY: Christian Literature Publishing Co.,；1888.；<http://www.newadvent.org/fathers/1801065.htm>.}

% structure: p0001 source=h1 target=section reason=page-title

\section{咏第六十六篇释义}

1. 这篇圣咏在标题上题写着：\Quote{为终末，一首关于复活的圣咏之歌。}当你们听到\Quote{为终末，}无论何时圣咏重复出现，都要理解为\Quote{为基督：}宗徒说：\BibleQuote{因为律法的终向就是基督，使所有信者成义。}\BibleRef{罗 10:4}因此，这里如何歌唱复活，以及是谁的复活，你们将听到，只要祂愿意赐予和启示。因为基督徒已经知道复活在我们元首内成就，并要在肢体中成就。教会的元首是基督，\BibleRef{哥 1:18}基督的肢体就是教会。那在元首中预先实现的，将在身体上跟随。这是我们的希望；我们为此相信，为此忍耐并坚持在这世界如此多的悖逆中，希望安慰我们，直到那希望成为现实……犹太人曾持有死人复活的希望：他们盼望只有自己因法律的工价和圣经的义行而复活进入永生，因为只有犹太人拥有这些，而外邦人没有。基督被钉十字架，\BibleQuote{以色列的一部分成了盲目，直到外邦人的数目满盈，}\BibleRef{罗 11:25}正如宗徒所说。死人复活也开始应许给那些相信耶稣基督、相信祂复活了的外邦人。因此，这篇圣咏是针对犹太人的狂妄和骄傲，为安慰那蒙召分享同样复活希望的外邦人。\par

2. ……因此他开始说：\Quote{应在天主内欢欣。}谁？\Quote{普世大地}\BibleRef{咏 65:1}。所以不是犹太地独享。请看，弟兄们，教会的普世性如何展现于普世广传：你们不要只哀悼那些嫉妒外邦人获得恩宠的犹太人，更要为异端者们哀哭。因为那些未曾被聚集的尚且值得哀悼，何况那些已被聚集却又分裂的呢？\Quote{应在天主内欢呼，普世大地。}什么是 \Quote{\OriginalTerm{欢呼}{jubilate}}？倘若你无法表达于言语，就迸发为欢欣之声。因为 \Quote{\OriginalTerm{欢呼}{jubilation}} 不出于言语，而是人欢欣的声音被发出，如同心在劳动，将所想象之物无法表达的喜乐化为声音。\Quote{应在天主内欢欣，普世大地：}愿无人只在局部欢呼；愿普世大地都欢欣，愿天主教会欢呼。天主教会怀抱全体：凡抓住局部而与整体分离的，该哀号，而非欢呼。\par

3. \Quote{但你们要演奏以光荣祂的名}\BibleRef{咏 65:2}。他说了什么？通过你们的\Quote{演奏}，愿祂的名受赞美。但什么是\Quote{演奏}？演奏也是拿起一种称为 \OriginalTerm{索尔特里}{psaltery} 的乐器，通过手的拨动和动作配合歌声。因此，如果你们欢呼以致天主听见，也要演奏一些东西使人们都能看见和听见：但不可归功于你们自己的名……因为如果你们为了自己被光荣而行善事，我们就要重复主对某些人所说的话：\BibleQuote{我实在告诉你们，他们已经获得了他们的赏报。}\BibleRef{玛 6:2}又说：\BibleQuote{否则，你们在天上的父那里将没有赏报。}\BibleRef{玛 6:1}你会说：那么，我该隐藏我的善工，不在人前做吗？不。但祂说什么？\Quote{让你们的善工在人前照耀。}于是我就处在疑惑中了。一方面祢对我说：\Quote{你们要小心，不要在人前做你们的义德；}另一方面祢对我说：\Quote{让你们的善工在人前照耀；}我该持守什么？做什么？不做什么？一个人不能侍奉两个主子，正如不能听从命令不同的两个主任一样。主说，我命令的不是不同的事。注意终点，因为终点歌唱：看你以什么终点做它。如果你为了被光荣而做，我已禁止；但如果为了天主被光荣而做，我已命令。因此，演奏不可归功于你们自己的名，而要归于你们天主主的名。演奏吧，愿祂被赞美；生活善好，愿祂被光荣。因为那善好生活从哪里来？如果你们从永远就拥有它，你们绝不会生活不善；如果你们从自己拥有它，你们就永远只能生活善好。\Quote{将光荣归于祂的赞美。}他将我们的全部注意力都引向天主的赞美，为我们不留下任何可被赞美之处。我们因此更要光荣喜乐：让我们紧依于祂，在祂内被赞美。你们听到宗徒被宣读时的话：\BibleQuote{弟兄们，看看你们的蒙召，按肉身来说，有智慧的不多，有权势的不多，有尊贵的不多，但天主拣选了世上愚拙的，为使有智慧的蒙羞。}\BibleRef{格前 1:26}……但主后来也拣选了雄辩家；但他们本会骄傲，如果祂没有先拣选渔夫；祂拣选了富人；但他们本会说因为他们的财富而被拣选，除非祂先拣选了穷人：祂后来拣选了皇帝；但更好的是，当一个皇帝来到罗马，他应摘下皇冠，在渔夫的墓前哭泣，而不是渔夫在皇帝的墓前哭泣。\BibleQuote{天主拣选了世上软弱的，为使强壮的蒙羞，}等等。\BibleRef{格前 1:27}……接着呢？宗徒总结说：\BibleQuote{为使任何血肉之躯在天主前不可夸耀。}你们看，祂如何从我们身上拿走，为赐予光荣：拿走了我们的，为赐予祂自己的；拿走了空虚的，为赐予充实的；拿走了不稳妥的，为赐予稳固的……\par

4. \BibleQuote{你们要向天主说：祢的作为何等可畏！}\BibleRef{咏 65:3}。为何是可畏而非可爱？再听另一篇圣咏的话：\BibleQuote{当以敬畏事奉主，并在他面前战兢欢乐。}这是什么意思？再听宗徒的话：\BibleQuote{怀着恐惧，}他说，\BibleQuote{并战战兢兢地，成就你们的救恩。}\BibleRef{斐 2:12}为何怀着恐惧战兢？他随后说明了原因：\BibleQuote{因为是天主在你们内运作，使你们愿意并实行祂的美意。}\BibleRef{斐 2:13}因此，如果天主在你们内运作，你们是因天主的恩宠而善行，而非凭自己的力量。所以，你们若喜乐，也要畏惧：免得那赐给谦卑之人的，因骄傲而被夺去。……弟兄们，如果对古时从圣祖的根上被剪下的犹太人，我们不应自高，反而要畏惧，并对天主说：\BibleQuote{祢的作为何等可畏！}那么，对于那些刚被剪下的新伤，我们岂不更不应自高？从前犹太人被剪下，外邦人被接上；从这接枝上，又有异端者被剪下；但我们也对他们不应自高，免得那乐于辱骂被剪下之人的人，自己也被剪下。我的弟兄们，一位主教的呼声，尽管不配，已向你们发出：我们劝你们警醒，无论你们在教会内是谁，不要辱骂那些不在教会内的人；反而要祈祷，愿他们也能进入。因为天主能再次把他们接上。\BibleRef{罗 11:23}宗徒正是对这同一批犹太人说了这话，并且在他们身上应验了。主复活了，许多人相信了：他们钉死祂时并不明白，但后来却信了祂，如此大过也被赦免了。主血的倾流被赦免了杀人者，且不说弑神者：\BibleQuote{因为他们若知道，决不会钉死荣耀的主。}\BibleRef{格前 2:8}如今，对杀人者，无辜者之血的倾流已被赦免；那因疯狂而倾流的同一血，又因恩宠而被他们喝下了。……哦，外邦人的圆满，你要向天主说：\BibleQuote{祢的作为何等可畏！}并要如此喜乐，以致你畏惧，不要自高于被剪下的枝子。\par

5. \BibleQuote{祢的力量浩大，祢的仇敌必向祢说谎。}为此他说：\BibleQuote{祢的仇敌必向祢说谎，}为使祢的力量彰显为大。这是什么意思？请更留心听。我们的主耶稣基督的力量在复活中最为彰显，这圣咏的标题即由此而来。复活后，祂显现给门徒。\BibleRef{宗 10:40}祂没有显现给仇敌，而是显现给门徒。被钉死时，祂显现给众人；复活后，显现给信徒：为叫后来凡愿意的都可相信，而相信的人得到复活的许诺。许多圣人行了诸多奇迹，但没有一个死后复活：因为即便是他们从死者中复活的人，也复活了以便死去。……因此，当主行奇迹时，犹太人或许会说：梅瑟也行了这些事，厄里亚也行了，厄里叟也行了；他们为自己说这话，因为那些人也曾使死人复活，并行了诸多奇迹。因此，当有人向祂要求一个记号时，祂提到那独特的、唯独在祂身上将成就的记号，说：\BibleQuote{这邪恶淫乱的世代寻求记号，除了约纳先知的记号外，再没有记号给它：因为正如约纳在大鱼腹中三天三夜，人子也要在地心里三天三夜。}\BibleRef{玛 12:39}约纳如何在大鱼腹中？岂不是后来活着被吐出来吗？地狱之于主，正如大鱼之于约纳。祂提到这独特的记号，这是最强大的记号。死后复活比未曾死更有力量。主为人时所显示的力量的伟大，正在于复活的德能。……\par

6. 也要留意福音中假见证的谎言，看它是如何关乎复活的。因为当有人对主说：\Quote{你行这些事，给我们显什么征兆呢？}\BibleRef{若 2:18}；此外，祂也借着另一关于约纳的比喻说了同样的事，\BibleRef{玛 12:39}，好叫你明白这独特的征兆特别被指明了。祂说：\Quote{拆毁这座圣殿，}又说\Quote{三天之内，我要把它重建起来。}他们就问：\Quote{这圣殿建筑了四十六年，你三天之内就要重建它吗？}\BibleRef{若 2:19}。福音作者解释这话的意思说：\Quote{但耶稣所说的是指祂身体的圣殿。}\BibleRef{若 2:21}。看，祂曾说过要在这同一件事上——即祂以圣殿作比喻的同一来源——向人显示祂的能力，因为祂的肉体是隐藏于内的\OriginalTerm{天主性}{Divinity}的圣殿。犹太人外表看见圣殿，却看不见内住的\OriginalTerm{神性}{Deity}。假见证就是利用\Emph{主}的这些话——祂提及未来复活时所论的圣殿——编造谎言来反对祂。当假见证被问及曾听见祂说什么时，他们控告祂说：\Quote{我们听见祂说：我要拆毁这座圣殿，三天之内我要把它重建起来。}\Quote{三天之内我要重建，}他们听见了；\Quote{我要拆毁，}他们却没有听见；所听见的是\Quote{你们要拆毁。}他们改了一个词和几个字母，以支持他们的假见证。但你为了谁而改词呢？人世的虚妄，人世的软弱啊！你要改变那不变的\OriginalTerm{圣言}{Word}吗？你改你的词，你能改天主的圣言吗？……那么，他们为什么说祢曾说过\Quote{我要拆毁，}而不说祢所说的\Quote{你们要拆毁}呢？这似乎是为了替他们自己辩护，以免被控无故拆毁圣殿。因为基督是自愿死去的；尽管如此，你们还是杀死了祂。看，我们承认，你们这些说谎者，是祂自己拆毁了圣殿。因为宗徒曾说过：\Quote{祂爱了我，且为我舍弃了自己。}\BibleRef{迦 2:20}关于父也曾说：\Quote{祂没有怜惜自己的儿子，反而为我们众人把祂交出来。}\BibleRef{罗 8:32}……无论如何，是祂自己拆毁了圣殿，是祂自己——祂曾说：\Quote{我有权舍掉我的灵魂，也有权再取回它；谁也不能从我夺去，而是我自愿舍掉它；我有权舍掉，也有权再取回。}\BibleRef{若 10:18}——拆毁了圣殿。是祂自己因祂的恩宠，因你们的恶意，拆毁了圣殿。\Quote{祢的敌人必因祢大能的威势向祢说谎。}看，他们说谎，看，他们被相信，看，祢受压迫，看，祢被钉十字架，看，祢受侮辱，看，有人向祢摇头：\Quote{祂如果是天主子，就从十字架上下来吧。}\BibleRef{玛 27:49}看，祢愿意时，舍掉生命，并被长矛刺透肋旁，\OriginalTerm{圣事}{Sacraments}从祢肋旁流出；\BibleRef{若 19:34}祢被从木架上取下，裹上殓布，安放在坟墓里，有卫兵把守，免得门徒把祢偷去；祢复活的时候到了，地动山摇，坟墓裂开，祢秘密地复活，公开显现。那些说谎的人在哪里？那邪恶意志的假见证在哪里？祢的敌人岂没有因祢大能的威势向祢说谎吗？\par

7. 也要把看守坟墓的卫兵给他们，让他们述说所见的事，让他们拿钱也去说谎。\BibleRef{玛 28:12}……他们也加入了敌人的谎言：说谎的人数增加了，为使信者的赏报增加。因此他们说谎，\Quote{因祢大能的威势}说谎：为叫说谎者惭愧，祢向诚实的人显现了，祢向那些祢所使之成为诚实的人显现了。\par

8. 让犹太人固守他们的谎言吧；至于祢，因为他们因祢大能的威势说谎，愿成就以下经文：\Quote{普世都要朝拜祢，向祢歌颂，歌颂祢的圣名，至高者。}\BibleRef{咏 65:4}。片刻之前，至卑微者，如今是至高者：至卑微者在说谎的敌人手中；至高者在赞美天使的头上。哦，你们外邦人，哦，最遥远的万民，离开说谎的犹太人，前来朝拜。\Quote{你们来，观看主的作为：在世人子以上，祂的计谋可畏。}\BibleRef{咏 65:5}。祂确实也被称为人子，也确实成了人子：祂是天主子，具有天主的形体；\BibleRef{斐 2:6}又是人子，取了奴仆的形体：但不要因祂的形体而判断祂如同其他相似的人：祂是\Quote{可畏的}，\Quote{在世人子以上的计谋中。}世人子商议要钉死基督，被钉的基督却使钉死祂的人瞎眼。那么，世人子啊，你们用敏锐的计谋反对你们的主——祂内隐藏着尊严，外貌却显出软弱——你们做了什么？你们商议要毁灭，祂却要使人瞎眼并拯救；要叫骄傲的人瞎眼，拯救谦卑的人；但要叫那些骄傲的人瞎眼，好使他们因瞎眼而自卑，因自卑而认罪，认罪后得蒙光照。\Quote{在世人子以上的计谋中可畏。}确实可畏。看，以色列的一部分硬心已发生了：\BibleRef{罗 11:25}看，基督所出的犹太人被排斥在外；看，原本反对犹太的外邦人，在基督内却在其中。\Quote{在世人子以上的计谋中可畏。}\par

9. 因此，祂藉祂计谋的恐怖做了什么呢？祂将海变成了旱地。因为紧随此的是：\Quote{祂将海变成了旱地} \BibleRef{咏 65:6}。海曾是这个世界，因咸苦而苦涩，因风暴而翻腾，因迫害的波浪而狂怒——它曾是海：确实，海已变成了旱地，如今那曾被咸水充满的世界渴求甘甜的水。谁做了这事？是祂，\Quote{祂将海变成了旱地}。如今所有外邦人的灵魂说什么呢？\Quote{我的灵魂对祢如同无水之地。} \Quote{祂将海变成了旱地；他们要在河中徒步经过。} 那些已变成旱地的人，虽然他们从前是海，\Quote{要在河中徒步经过。} 什么是河？这条河是世间一切的必朽性。观察一条河：有些东西来而逝去，另一些即将逝去的东西接踵而至。河的水不也是从地湧出而流淌的吗？每一个出生的人必须让位于即将出生的人：这整个流转的秩序就是一种河。不要让灵魂贪婪地投进这条河，让她不要投进去，而是要站定不动。她将如何越过注定灭亡之物的享乐呢？让她信从基督，她将徒步经过：她以祂为首领而经过，徒步经过。\par

10. \Quote{在那里，我们要在祂内喜乐。} 哦，你们犹太人，你们以自己的行为夸口：放下自夸的骄傲，领受在基督内喜乐的恩宠。因为在那里我们将在祂内喜乐，而非在自己内：\Quote{在那里，我们要在祂内喜乐。} 我们何时喜乐？当我们已徒步越过河之后。永生被应许，复活被应许，在那里，我们的肉身将不再是河流：因为它现在是一条河流，只要它是必朽的。观察是否有任何年岁站定不动。孩童渴望长大；他们不知道随着年岁的更替，他们生命的长度在减少。因为随着他们成长，岁月不是被增加，而是被取走：正如河的水总是靠近，但从源头退去。孩童渴望长大，好逃脱长辈的束缚；看，他们长大了，这事很快发生，他们到达了青年：让那些从童年出来的人，如果可能，保留他们的青年：那也过去了。老年接踵而至：即使让老年成为永恒，它也被死亡移除。因此，有一条由所生之肉身的河流。这条必朽的河流，为免因对必朽之物的贪欲而侵蚀并冲走人，他能轻易地经过，他谦卑地，即徒步地经过，以那首先经过者为首领，祂在路上的急流中饮了，因此抬起了头。因此，徒步经过那条河，即轻易地经过那滑逝的必朽性，\Quote{在那里，我们要在祂内喜乐。} 但现在，除了在祂内，或在祂的希望中，还能在哪？因为即使现在我们喜乐，我们是在希望中喜乐；但那时我们要在祂内喜乐。现在也在祂内，但通过希望：\Quote{但那时将要面对面。} \BibleRef{格前 13:12} \Quote{在那里，我们要在祂内喜乐。}\par

11. 在谁内？\Quote{在祂以内，祂以祂的德能永远为王} \BibleRef{咏 65:7}。因为我们有什么德能，并且是永远的呢？如果我们的德能是永远的，我们就不会滑倒，不会陷入罪，不会应得刑罚性的必朽性。祂，因祂的善意，担当了我们的所应得把我们推落之处。\Quote{祂以祂的德能永远为王。} 让我们成为祂的有分者，在祂的德能中我们将坚强，但祂在自己的德能中。我们受光照，祂是光照的光：我们转离祂，就在黑暗中；祂不能转离自己。藉祂的热力我们得温暖；从那里退离时我们变冷，再次亲近祂时我们得温暖。因此，让我们向祂说话，求祂保守我们在祂的德能中，因为\Quote{在祂内我们要喜乐，祂以祂的德能永远为王。}\par

12. 但这恩典并非只赐予相信的犹太人……\Quote{祂的眼目眷顾外邦人。} 我们做什么呢？犹太人将要抱怨；犹太人会说：\Quote{祂赐给我们的，也赐给了他们；赐给我们福音，也赐给他们福音；赐给我们复活的恩宠，也赐给他们复活的恩宠；我们领受了法律，并在法律的义中生活，遵守了祖传的诫命，这难道对我们毫无益处吗？难道毫无功效吗？他们与我们一样。} 让他们不要争斗，不要争辩。\Quote{不要让那些心怀苦毒的人在自己内受举扬。} 哦，可怜而耗损的肉身，你不是有罪的吗？你的舌头为何叫嚷？让良心被听从。\Quote{因为所有人都犯了罪，需要天主的荣耀。} \BibleRef{罗 3:23} 认识你自己，人性的软弱。你领受了法律，为叫你成为法律的违反者：\BibleRef{罗 5:20} 因为你没有遵守和履行你所领受的。因着法律，临于你的不是法律所命令的成义，而是你所行的过犯。因此，如果罪恶增多了，你为什么嫉妒那更丰盛的恩宠呢？不要心怀苦毒，因为\Quote{不要让那些心怀苦毒的人在自己内受举扬。} 祂似乎在某种意义上发出了诅咒：\Quote{不要让那些心怀苦毒的人受举扬；} 是的，让他们受举扬，但不要\Quote{在自己内。} 让他们在自己内受贬抑，在基督内受举扬。因为\Quote{凡自谦卑的，必被举扬；凡自举扬的，必被贬抑。} \BibleRef{玛 23:12} \Quote{不要让那些心怀苦毒的人在自己内受举扬。}\par

13. \BibleQuote{你们万民，请赞美我们的天主} \BibleRef{咏 65:8}。看，那心怀怨恨的人已被击退，与他们算了账：有些人归正了，有些人仍然傲慢。不要让那些嫉妒外邦人得到福音恩宠的人恐吓你们；如今亚巴郎的子孙已来到，万民都要因他蒙受祝福。\BibleRef{创 12:3} 你们当祝福那使你们蒙福的那一位，\BibleQuote{你们万民，请赞美我们的天主，请听祂赞美的声音。}不要赞美自己，而要赞美祂。什么是祂赞美的声音？那就是：我们所有的一切美好，都是出于祂的恩宠。\BibleQuote{祂使我的灵魂得以存活} \BibleRef{咏 65:9}。看，这就是祂赞美的声音：\BibleQuote{祂使我的灵魂得以存活。}因此，灵魂原是处于死亡之中：在你自身内，它原是处于死亡之中。因此，你们不应当自高自大。所以，在你自身内，灵魂原是处于死亡之中：它将在哪里获得生命呢？除非是在那一位说\BibleQuote{我就是道路、真理、生命} \BibleRef{若 14:6} 的那位里面。正如宗徒对某些信徒所说：\BibleQuote{你们从前原是黑暗，但现在在主内却是光明}。\BibleRef{弗 5:8} ……\BibleQuote{且未使我的脚动摇。}祂使我的灵魂得以存活，祂引导我的双脚，使它们不致失足，不被移动和摇摆；祂使我们活着，祂使我们坚持到底，为使我们得以永远活着。……\par

14. \BibleQuote{因为你曾考验了我们，天主！你曾锻炼了我们，如同银子被锻炼} \BibleRef{咏 65:10}。你并未像锻炼干草那样锻炼我们，而是像锻炼银子：你将火加于我们身上，并未将我们化为灰烬，而是洗净了我们的不洁，\BibleQuote{你曾锻炼了我们，如同银子被锻炼}。且看天主如何对待那些祂已使其灵魂得以存活的人。\BibleQuote{你曾引领我们进入网罗：}并非使我们被捕获而丧亡，而是使我们受到考验并从中得救。\BibleQuote{你把重担放在我们的背上：}因为我们曾不当地自高自大，傲慢不已；曾不当地自高自大，我们被压低，为使我们被压低后，能合宜地被高举。\BibleQuote{你把重担放在我们的背上：} \BibleQuote{你使人骑在我们的头上} \BibleRef{咏 65:11}。这一切，教会都在各种不同的迫害中经历过；她在她的个别肢体中也曾经历，如今仍在经历。因为在今生，没有一个人能说自己免于这些考验。因此，甚至有人被置于我们头上：我们忍受那些我们所不愿忍受的人，我们为了那些我们认为比自己更坏的人而受苦。但是，如果罪过不存在，一个人就是正当地居于上位；然而，罪过越多，他就越是卑下。将自己视为罪人，并因此忍受那些被置于我们头上的人，这是好的：为使我们也能向天主承认，我们受苦是应得的。因为你为何要愤怒地忍受那位公义者所行的事呢？\BibleQuote{你把重担放在我们的背上：你使人骑在我们的头上。}当天主这样做时，他似乎是发怒了：不要害怕，因为他是父亲，他永远不会愤怒到毁灭的地步。当你生活邪恶时，如果他宽容，那才是更愤怒。总之，这些磨难是他管教所用的棍杖，免得他施以惩罚的判决。……\par

15. \Quote{我们经过了火和水。}火和水在今生都是危险的。当然，水似乎能灭火，而火似乎能烤干水。同样，这些就是考验，其中充满此生。火灼烧，水腐蚀：两者都须畏惧，既怕磨难的灼烧，也怕腐败的水。每当有逆境，以及世上所谓的不幸，那就像火；每当有顺境，世间的丰收环绕，那就像水。要当心火不烧伤你，水不腐蚀你……不要急于进入水：先经过火，再到达水，这样你也能渡过水。因此，在奥迹礼仪、慕道和驱逐礼中，首先使用火。因为不洁之神常常喊叫：\Quote{我燃烧，}若非那火，又是什么呢？但在驱逐礼之火后，我们来到圣洗圣事：这样从火到水，从水到安息。然而，正如在圣事中，同样在世界的试探中也是如此：首先是恐惧的紧迫临近，代替火；之后恐惧被除去，我们应当担心世俗的幸福腐蚀我们。但当火没有使你爆裂，当你没有沉入水中，却游了出来；那么，透过纪律你到达安息，经过火和水，你被领到一处清凉之所。因为那些在圣事中作为标记的事物，其真实本体就在永生的完美境界中……但我们在那里并非麻木，而是安息；虽然被称为热，我们在那里不会感到热，而是要在心灵中炽热。请看另一篇圣咏中同样的热：\Quote{没有一个人能躲避它的热。}宗徒又说了什么？\BibleQuote{在心灵中要热切。} \BibleRef{罗 12:11} 因此，\Quote{我们经过了火和水：而你引领我们进入凉爽之所。}\par

16. 请注意，他既未提及阴凉之处，也未提及那令人渴慕的火：\Quote{我要带着全燔祭进入祢的殿宇}\BibleRef{咏 65:13}。何为全燔祭？即整个祭品被焚烧，但用的是天主的火。因为当整个祭品被焚烧时，这牺牲便称为全燔祭：一部分牺牲是一回事，全燔祭是另一回事：当整个被焚烧，整个被天主的火吞噬时，称为全燔祭；当只烧一部分时，则称为牺牲。的确，每个全燔祭都是牺牲，但并非每个牺牲都是全燔祭。因此，基督的身体、基督的合一正许诺这全燔祭：\Quote{我要带着全燔祭进入祢的殿宇}。愿祢的火吞噬我的一切，不给我留下任何属于我的东西，一切都归于祢。但这将在义人的复活时实现，\Quote{当这可朽坏的穿上了不可朽坏的，这可死的穿上了不死的时，经上所记载的话就应验了：\InnerQuote{死亡被胜利吞灭了。}}\BibleRef{格前 15:54}。胜利，就好比是天主的火：当它连我们的死亡也吞灭时，就成了全燔祭。肉身不再有任何可朽的，精神不再有任何可责的：整个必死的生命将被耗尽，以便在永生中达于圆满，使我们从死亡中被保存于生命。这些便是全燔祭。那么，\Quote{在全燔祭中}还会有什么呢？\par

17. \Quote{我要向祢偿还我的誓愿，就是我嘴唇所分别的}\BibleRef{咏 65:14}。誓愿中的分别是什么？这分别就是：你责备自己，赞美祂；认清自己是受造物，祂是造物主；自己是黑暗，祂是光明者，你应向祂说：\Quote{上主，我的天主，祢必点亮我的灯，祢必照亮我的黑暗}。因为每当你，灵魂啊，说光来自你自己时，你便没有分别。你若不分清，便不能偿还分明的誓愿。要还清誓愿：承认自己是可变的，祂是不变的；承认离了祂你便一无所是，而祂离了你仍是完美的；你需要祂，祂却不需要你。要向祂呼喊：\Quote{我曾对上主说：祢是我的天主，因为祢不需要我的善物}。虽然天主接纳你为使你成为全燔祭，但祂不增长，不增加，不变得更富有，不变得更齐全：祂为你所成就的一切，对你是更好，对成就者则否。你若能分清这些，便能向你的天主偿还你嘴唇所分别的誓愿。\par

18. \Quote{我的口在我困苦中曾发言。}困苦是何等甘甜，何等必要！在此，那同一口在困苦中说了什么？\Quote{我要向祢献上带有骨髓的全燔祭}\BibleRef{咏 65:15}。什么是\Quote{带有骨髓的}？愿我在内心保持祢的爱，它不浮于表面，而是在我的骨髓里我爱祢。因为没有什么比我们的骨髓更内在：骨头比肉更内在，骨髓比骨头更内在。因此，谁若在表面上爱天主，更愿取悦于人，而内心却存有别的情感，他便不献上带有骨髓的全燔祭；但主注视谁的骨髓，便收纳谁的全人。\Quote{连同馨香和公绵羊。}公绵羊是教会的首领：整个基督的身体在说话：这是他们献给天主的祭品。馨香是什么？祈祷。\Quote{连同馨香和公绵羊。}因为公绵羊尤其为羊群祈祷。\Quote{我要向祢献上公牛连同公山羊。}我们看到公牛践踏谷粒，它们也被献给天主。宗徒说过，对福音的宣讲者应理解为经上所记载的：\Quote{不可笼住正在打谷的公牛的嘴。难道天主关心牛吗？}因此，那些公绵羊是伟大的，公牛是伟大的。其余的人呢？他们或许意识到某些罪过，或许在途中滑倒受伤，正借着补赎得到医治？他们也会继续存在，难道不属于全燔祭吗？让他们不要害怕，他也加上了公山羊。\Quote{我要向祢献上公牛连同公山羊。}借着同轭，公山羊得救；它们自身没有力量，与公牛同轭便蒙收纳。因为他们用不义的钱财结交了朋友，好使那些人接他们进入永恒的帐幕。\BibleRef{路 16:9}因此，那些公山羊不会在左边，因为他们用不义的钱财为自己结交了朋友。但哪些公山羊会在左边呢？就是那些被说\Quote{我饿了，你们没有给我吃的}的人，\BibleRef{玛 25:42}而不是那些用施舍赎罪的人。\par

19. \BibleQuote{你们来，听，我要述说，凡敬畏天主的你们}\BibleRef{咏 65:16}。我们来，我们来听，他即将述说什么，\InnerQuote{你们来，听，我要述说。} 但是对谁，\InnerQuote{你们来，听}？\InnerQuote{你们凡敬畏天主的。} 如果你们不敬畏天主，我就不述说。不可能向任何没有天主敬畏之处述说。愿敬畏天主开启耳朵，好让有些东西进入，并有一条道路，使我即将述说的可以进入。但他将述说什么？\BibleQuote{祂为我的灵魂行了何等大事。} 看，他想要述说：但他将述说什么？难道是大地如何广阔，天空如何延伸，星辰多少，太阳和月亮的变化如何？这受造物完成它的历程：但是他们极其好奇地探究了它，却不认识它的造主。\BibleRef{智 13:1} 这事要听，这事要接受，\BibleQuote{你们敬畏天主的，祂为我的灵魂行了何等大事：} 如果你们愿意，也为你们的灵魂。\BibleQuote{祂为我的灵魂行了何等大事。} \BibleQuote{我用我的口呼求了祂}\BibleRef{咏 65:17}。而正是这事，他说，已成就于他的灵魂；即他应用口呼求祂，他说，已成就于他的灵魂。看，弟兄们，我们曾是外邦人，即使不在我们自己，也在我们的父母。宗徒说什么？\BibleQuote{你们知道，当你们是外邦人时，你们怎样被牵引，走向无言的偶像。}\BibleRef{格前 12:2} 愿教会如今说，\BibleQuote{祂为我的灵魂行了何等大事。} \BibleQuote{我用我的口呼求了祂。} 我这个人曾向石头呼号，向聋的木偶呼号，向聋哑的偶像说话：如今天主的肖像已转向它的造主。我曾\BibleQuote{对木偶说：你是我的父亲；对石头说：你生了我：}\BibleRef{耶 2:27} 如今我说：\BibleQuote{我们在天之父}\BibleRef{玛 6:9}。…\BibleQuote{我用我的口呼求了祂，并在我的舌下高举了祂。} 看，那献上骨髓的全燔祭的，在隐密中是如何纯洁。你们要这样做，弟兄们，要效法，使你们能说：\BibleQuote{你们来，看祂为我的灵魂行了何等大事。} 因为凡他所说的一切，都是藉祂的恩宠在我们灵魂中成就的。看他所说的其他事物。\par

20. \BibleQuote{若我心中看到了不义，愿主不要俯听}\BibleRef{咏 65:18}。弟兄们，现在请思考，人因怕人而脸红，多么轻易地、每日地责备不义；他做了坏事，他做了卑劣的事，这家伙是个恶棍：这也许是出于对人的顾虑而说的。看看你是否在你心中看到不义，是否你所责难的别人的事，你自己正想要做，因此你对他叫嚷，不是因为他做了，而是因为他被发现了。回到你自己，在内心做你自己的审判者。看在你隐秘的室内，在心灵的最深处，那里只有你和那看见者独处，在那里让不义令你不悦，好使你能取悦于天主。不要看重它，即不要爱它，而要鄙视它，即藐视它，并远离它。无论它许诺了什么令人愉悦的事来引诱你犯罪；无论它恐吓了什么令人痛苦的事来驱使你作恶；一切都是虚无，一切都过去：它值得被鄙视，好被践踏；不要顾盼，以免被接纳。……\par

21. \BibleQuote{因此天主俯听了我}\BibleRef{咏 65:19}。因为我心中没有看到不义。\BibleQuote{祂垂听了我祈祷的声音。} \BibleQuote{愿我的天主受赞美，因祂没有推开我的恳求和祂的仁慈}\BibleRef{咏 65:20}。从那里领会意思，在那里祂说：\BibleQuote{你们来，听，我要告诉你们，你们凡敬畏天主的，祂为我的灵魂行了何等大事：} 祂已经说了你们所听到的话，并且在结束时这样结论：\BibleQuote{愿我的天主受赞美，因祂没有推开我的恳求和祂的仁慈。} 因为说话者就这样到达了复活，在那里我们已藉望德也在其中：是的，我们自身就是那个声音，这声音是我们的。因此，只要我们在此世，就让我们向天主祈求这一点：使祂不推开我们的恳求和祂的仁慈，即我们要不断祈祷，祂要不断怜悯。因为许多人在祈祷中变得软弱，在自己悔改的新鲜中热心祈祷，之后软弱，之后冷淡，之后疏忽：好像他们已经安全了。仇敌警醒：你却在睡觉。主自己在福音中给出了命令：\BibleQuote{人应当常常祈祷，不要灰心。} 祂用那个不义的审判官作比喻，他既不敬畏天主，也不尊重人，那寡妇天天缠着他要听他申诉；他因厌烦而让步，不是出于怜悯：那个邪恶的审判官对自己说：\BibleQuote{我虽不敬畏天主，也不尊重人，但因这寡妇天天烦扰我，我要给她申冤，给她报仇。} 主说：\BibleQuote{如果那不义的审判官这样做了，难道你们的天父不更要给祂所拣选的人报仇，就是那些昼夜向祂呼号的人吗？我告诉你们，祂要快快为他们伸冤。} 因此，让我们在祈祷中不要灰心。虽然祂延迟祂要赐予的，但祂并没有丢弃：既然确信祂的应许，让我们在祈祷中不要灰心，这是出于祂的良善。因此祂说：\BibleQuote{愿我的天主受赞美，因祂没有推开我的恳求和祂的仁慈。} 当你看到你的恳求\BibleQuote{没有从你那里被推开}，你就确信：祂的仁慈没有从你那里被推开。\par

\SourceRecord{Translated by J.E. Tweed.；Nicene and Post-Nicene Fathers, First Series；Vol. 8.；Edited by Philip Schaff.；Buffalo, NY: Christian Literature Publishing Co.,；1888.；<http://www.newadvent.org/fathers/1801066.htm>.}

% structure: p0001 source=h1 target=section reason=page-title

\section{圣咏第六十七篇释义}

1. 你们的爱德还记得，在已经讲解过的两篇圣咏中，我们曾激励我们的灵魂赞颂主，并以虔诚的歌唱说：\BibleQuote{我的灵魂，你要赞颂主。}因此，如果我们在那两篇圣咏中激励了我们的灵魂去赞颂主，那么在这篇圣咏中，便恰当地说：\BibleQuote{愿天主怜悯我们，并祝福我们}\BibleRef{咏 66:1}。愿我们的灵魂赞颂主，并愿天主祝福我们。当天主祝福我们时，我们成长；当我们赞颂主时，我们也成长。两者对我们都有益处。祂不会因我们的祝福而增加，也不会因我们的诅咒而减少。诅咒主的人，自己减少；祝福主的人，自己增加。首先，在我们内有主的祝福，结果我们也祝福主。那是雨水，这是果实。因此，就好像对天主这位农夫献上果实，祂降雨浇灌并耕耘我们。让我们以不贫乏的虔诚、不空洞的声音，而是以真诚的心来咏唱这些话。因为很明显，天主圣父被称为农夫\BibleRef{若 15:1}。宗徒说：\BibleQuote{你们是天主的田地，是天主的建筑物}\BibleRef{格前 3:9}。在这世界可见的事物中，葡萄树不是建筑物，建筑物也不是葡萄园：但我们是主的葡萄园，因为祂耕耘我们以结果实；我们是天主的建筑物，因为耕耘我们的祂，居住在我们内。而同一宗徒说了什么？\BibleQuote{我栽种，阿颇罗浇灌，然而使之生长的，是天主。所以，栽种的不算什么，浇灌的也不算什么，只在乎那使之生长的，就是天主}\BibleRef{格前 3:6}。因此，是祂赋予生长。那些人难道是农夫吗？因为栽种、浇灌的人被称为农夫：但宗徒说：\BibleQuote{我栽种，阿颇罗浇灌。}我们要问，他自己从何处做了这事？宗徒回答说：\BibleQuote{然而不是我，而是天主的恩宠偕同我}\BibleRef{格前 15:10}。因此，无论你转向何处，或通过天使，你将发现天主是你的农夫；或通过先知，同一主是你的农夫；或通过宗徒，你要承认同一位是你的农夫。那么，我们呢？我们也许是那位农夫的工人，而且这工作也是由祂自己赐予的能力、由祂自己恩赐的恩宠而做的……\par

2. \BibleQuote{愿祂的慈颜光照我们。}你或许会问，\Quote{祝福我们}是什么意思？人们渴望从天主那里得到祝福，方式多种多样：一个人希望自己被祝福，使他房屋充满今生的必需品；另一个人希望自己被祝福，使他获得毫无瑕疵的身体健康；又一个人希望自己被祝福，如果他正生病，则使他恢复健康；另一个人切望儿子，且因无子而悲伤，希望自己得祝福，使他有后裔。谁能数算渴望自己从天主那里得到祝福的人的各种愿望呢？但我们中谁会认为，如果田地为他结果实，或任何人的房屋充满现世富足，或身体本身健康得以保持而不丧失，或丧失后得以恢复，这不是天主的祝福呢？……\par

3. \\Quote{每一个蒙福的灵魂都是纯朴的，}不依附于世俗之事，也不以粘住的翅膀匍匐，而是因美德的荣光而闪耀，以双爱之双翼跃入自由的空气；并看到从她那里被撤去的是她所践踏的，而非她所安息的，她安然地说：\\BibleQuote{上主赐予，上主收回；上主所喜悦的，就这样成就了：愿上主的名受赞美。}……但不要让任何软弱的人说：\\Quote{我何时才能有像圣\\Person{约伯}那样的伟大德行？}你惊叹于树的雄伟，因为你刚刚出生；这棵你惊叹的大树，你在它的枝荫下乘凉，曾经只是一根枝条。但你害怕当你变成如此时，这些会被夺走吗？观察它们也会被恶人夺去。那么你为什么延迟悔改？你害怕在良善时失去的东西，也许在邪恶时你仍会失去。如果你在良善时失去了它们，那么在你身边有安慰者取走了它们：宝箱空了黄金；心中充满了信德：外表你贫穷，但内里你富足：你随身携带你的财富，即使你赤身从船难中逃出，也不会失去。为什么那可能发生在恶人身上的损失，不在你良善时发现你？因为你看到恶人也遭受损失？但他们遭受更大的损失：房屋空了，良心更空。任何恶人失去了这些东西，外在没有什么可倚靠的，内在也没有什么可安息的。他在遭受损失后逃离了他曾在众人面前炫耀财富的地方；现在不能在众人面前炫耀了：他不回到自己的内心，因为他一无所有。他没有效法蚂蚁，在夏天没有为自己采集谷粒。我所说的夏天是什么意思？当他生活安宁，当他拥有今世的繁荣，当他有闲暇，当所有人都称他幸福时，那就是他的夏天。他本应效法蚂蚁，聆听天主的圣言，采集谷粒，并储存在内。考验来了，麻木的冬天、恐惧的暴风雨、悲伤的寒冷降临到他身上，无论是损失、安危的危险、家人的丧亡、或者耻辱和羞辱；那是冬天；蚂蚁依赖它在夏天所采集的；在它秘密的储藏所里，无人看见，它因夏天的劳作而恢复。当它在夏天为自己储存这些时，所有人都看到它；当它在冬天以这些为食时，无人看见。这是什么？看天主的蚂蚁，它日复一日起来，急忙前往天主的教会，祈祷，聆听读经，咏唱圣歌，消化所听到的，在心中思量，将打谷场上采集的谷粒储存在内。那些谨慎聆听此刻正在谈论之事的人，就这样做，所有人都看见他们前往教会，从教会回来，听讲道，听读经，选择一本书，打开并阅读：当这些事情发生时，所有人都看见。那只蚂蚁正在行走它的路，在看见它的人面前搬运和储存。冬天总会来临，因为谁不会遇到呢？发生损失，发生丧亡：其他人也许认为它可怜，不知道蚂蚁内部有什么可吃，他们说：\\Quote{可怜啊，遭遇这事的人，或者你认为，他遭遇这事后心情如何？他多么痛苦？}他以自己衡量，凭自己的力量同情；因此他被欺骗：因为他用来衡量自己的尺度，他会用在他不认识的人身上……懒惰的人哪！趁你能的时候在夏天采集吧；冬天不会容许你采集，只容许你吃你所采集的。因为有多少人遭受如此考验，以致没有机会阅读或聆听任何东西，甚至可能无法接近安慰他们的人。蚂蚁留在它的巢里，让它看看是否在夏天采集了什么，以便在冬天可以恢复自己。\par

4. ……有两种诠释，都须给出：\Quote{光照，}他说，\Quote{祢的容颜对我们，}向我们显示祢的面容。因为天主并非曾经照耀祂的容颜，如同它曾是无光的；而是祂对我们照耀它，使得那对我们隐藏的，为我们开启，而那本是存在却对我们隐藏的，为我们揭露，也就是被光照。或者确实是这样：\Quote{祢的肖像光照我们：}所以他说这话，在\Quote{祢的容颜光照我们}中：祢将祢的容颜印在我们身上；祢按祢的肖像和模样造了我们，\BibleRef{创 1:26}祢造了我们作祢的钱币；但祢的肖像不该留在黑暗中：求祢发出祢智慧的光芒，驱散我们的黑暗，让祢的肖像在我们内闪耀；让我们知道我们是祢的肖像，让我们听见在\WorkTitle{雅歌}中所说的：\BibleQuote{你若不认识自己，女人中最美的。}\BibleRef{歌 1:8}因为祂对教会说：\BibleQuote{你若不认识自己。}\BibleRef{歌 1:8}这是什么？你若不知道你自己是按天主的肖像被造的。教会的灵魂，宝贵的，以无玷羔羊的血赎回来的，注意你是何等贵重，思想为你付出的是什么。因此，让我们说，并渴望祂\Quote{光照祢的容颜对我们}。我们佩戴祂的容颜：如同皇帝的面容被提及，的确一种神圣的面容是天主在祂自己的肖像中：但不义的人不认识他们自身的天主肖像。为了使天主的容颜对他们被光照，他们该说什么？\BibleQuote{主，我的天主，祢必照亮我的烛台，祢必照亮我的黑暗。}\BibleRef{咏 18:28}我处在罪的黑暗中，但通过祢智慧的光芒，愿我的黑暗被驱散，愿祢的容颜显现；若它透过我显得有些变形，求祢使它重新成型，那由祢所成型的。\par

5. \Quote{我们在地上要认识祢的道路}\BibleRef{咏 66:2}。\Quote{在地上，}在此，此生中，\Quote{我们要认识祢的道路。}什么是\Quote{祢的道路}？那引向祢的道路。愿我们承认我们往哪里去，承认我们在行路时身在何处；这在黑暗中我们无法做到。祢远离寄居的人，祢给我们呈现了一条道路，我们必须藉以回归祢。让我们\Quote{在地上承认祢的道路。}祂的道路是什么，我们曾渴望\Quote{我们在地上要认识祢的道路}？我们自己去探究这事，而不是靠自己学习。我们可以从福音中学习：\BibleQuote{我是道路，}\BibleRef{若 14:6}主说：基督曾说：\BibleQuote{我是道路。}但你害怕走错吗？祂又加了：\BibleQuote{并且是真理。}\BibleRef{若 14:6}谁在真理中走错？离开真理的人就走错。真理是基督，道路是基督：在其中行走。你害怕在到达祂之前死去吗？\BibleQuote{我是生命：我是，}祂说，\BibleQuote{道路、真理和生命。}\BibleRef{若 14:6}好像祂在说：\Quote{你怕什么？藉着我你行走，向我你行走，在我内你安息。}所以\Quote{我们在地上要认识祢的道路}是什么意思，不就是\Quote{我们在地上要认识祢的基督}？但让圣咏自己回答：免得你们认为必须从其他经文中援引见证，而这里或许缺乏；祂通过重复已表明了什么意义，\Quote{我们在地上要认识祢的道路：}好像你在问，\Quote{在地上，什么道路？}\BibleQuote{在万民中祢的救恩。}\BibleRef{咏 66:2}你问在地上？听：\BibleQuote{在万民中。}\BibleRef{咏 66:2}你找什么道路？听：\BibleQuote{祢的救恩。}\BibleRef{咏 66:2}基督难道不是他的救恩吗？而老年西默盎所说的，那位在福音中的老人，我说，保存年岁直到圣言之婴儿时期的话是什么？\BibleRef{路 2:30}因为那老人将天主的婴儿圣言抱在手中。难道那在母胎中屈尊俯就的，会不屑于在老人手中？同一位曾在童贞女的胎中，也在老人手中，一个软弱的婴孩，既在腹中，也在老人手中，赐给我们力量，万有藉祂而造；如果万有，甚至祂的母亲。祂来谦卑，祂来软弱，但披上了软弱以改变为刚强，因为\BibleQuote{祂虽因软弱被钉死，却因天主的德能活着，}\BibleRef{格后 13:4}宗徒说。祂当时就在老人手中。而那老人说什么？他欢喜自己现在必须从这世界被释放，看见在他手中拿着那一位，藉祂并在他内他的救恩被扶持；他说什么？\BibleQuote{现在，祢放，}他说，\BibleQuote{主，祢的仆人平安离去，因为我亲眼看见了祢的救恩。}\BibleRef{路 2:29}因此，\BibleQuote{愿天主祝福我们，并怜悯我们；愿祂光照祂的容颜对我们，使我们在地上认识祢的道路！}\BibleRef{咏 66:2}在地上？\BibleQuote{在万民中。}\BibleRef{咏 66:2}什么道路？\BibleQuote{祢的救恩。}\BibleRef{咏 66:2}\par

6. 既然万邦都已知晓天主的救恩，接下来是什么呢？\Quote{愿万民都称谢你，天主}\BibleRef{咏 66:3}；他说，\Quote{称谢你，}\Quote{万民。}有一个异端者站出来说：我在非洲有子民；另一个从别处说：我在迦拉达有子民。你在非洲，他在迦拉达：因此我需要一个到处都有子民的人。你听到那声音时，确实敢于欢欣：\Quote{愿万民都称谢你，天主。}听听接下来的诗句，他并非只说一部分：\Quote{愿万民都称谢你。}你们要与万邦一同走在路上；要与万民一同走在路上，和平之子们，独一天主教会的儿女们，你们要走在路上，边行边看。行路人这样做以消解旅途的劳苦。你们要在这路上歌唱；我以这同一条路恳求你们，在这路上歌唱：唱一首新歌，不要让任何人唱旧歌：唱你们家乡的颂歌，不要让任何人唱旧歌。新路，新行路人，新歌。请听宗徒劝勉你们唱新歌：\Quote{所以凡在基督内的，就是新的受造物；旧事已过，看，都成了新的。}你们要在路上唱新歌，就是你们所学的\Quote{在地上}。在什么地上？\Quote{在万邦中。}因此，新歌也不属于一部分。凡是在一部分里唱歌的，唱的是旧歌：无论他愿意唱什么，他唱的是旧歌，旧人唱：他是分裂的，他是属血肉的。的确，他越是属血肉，就越是旧的；他越是属神，就越是新的。请看宗徒说什么：\Quote{我曾不能对你们说话，如同对属神的人，只如同对属血肉的人。}\BibleRef{格前 3:1} 他如何证明他们是属血肉的呢？\Quote{因为有人说：我是属保禄的；另一人说：我是属阿颇罗的：你们岂不是}他说，\Quote{属血肉的吗？}\BibleRef{格前 3:4} 因此，在圣神内，你要在安全之路上唱新歌。正如行路人歌唱，常常在夜间歌唱。周围一切都是令人敬畏的声音，或者不如说没有声音，而是寂静；越寂静越可怕；然而，甚至惧怕强盗的人也唱歌。你在基督内唱歌是何等安全！那条路没有强盗，除非你离弃了路，落入强盗手中……你们为何害怕承认，并在承认中与整个大地一同唱新歌？在整个大地，在至公的平安中，你们害怕向天主承认，怕祂定你们承认者的罪？如果你们不承认而隐藏，承认了反会被定罪。你们害怕承认，但即使不承认也不能隐藏：如果你们保持缄默，将以承认之事被定罪，而你们本可以因承认而得救。\Quote{天主，万民都要向祢承认。}\par

7. 由于这种承认不会导致刑罚，他接着说：\Quote{愿万民欢欣踊跃}\BibleRef{咏 66:4}。如果强盗在向人承认后哀哭，那么信徒在天主面前承认后应当喜乐。如果人是法官，酷刑和恐惧迫使强盗承认：有时恐惧逼出承认，痛苦强要承认；那个在酷刑中哀哭却害怕承认后被杀的人，尽其所能忍受酷刑；如果他被痛苦征服，他就开口招供死亡。无论怎样，他都不欣喜，绝不欢跃：在承认之前，铁爪撕裂他；承认之后，刽子手将他作为定罪的囚犯带走：无论如何都是悲惨的。但是\Quote{愿万民欢欣踊跃。}从何而来？藉着同样的承认。为什么？因为他们所承认的那位是善的：祂要求承认，为的是释放谦卑的人；祂定不承认者的罪，为的是惩罚骄傲的人。因此，你在承认之前要忧伤，承认之后要欢跃，现在你要被治愈。你的良心曾积聚恶液，肿胀如疖，折磨你，使你不得安息：医生施用药水浸泡，有时切开它，用磨难的惩戒施以手术刀：你要认出医生的手，承认吧，让一切恶液在承认中流出流尽：现在欢跃，现在喜乐，剩下的将易于愈合……\Quote{愿万民欢欣踊跃，因为你以正义审判万民。}为了使不义的人不害怕，他又加上：\Quote{你在地上引导万民。}万民曾经败坏，万民曾经弯曲，万民曾经悖逆；因为他们败坏、弯曲、悖逆的恶劣应得，他们害怕审判者的来临：那同一审判者的手伸出来，仁慈地施予万民，他们被引导，以便行走正路；既然他们先已认出祂是修正者，为何还要害怕审判者的来临呢？让他们把自己交在祂手中，祂亲自在地上引导万民。但被引导的万民行走在真理中，在祂内欢跃，行善工；如果偶尔有水渗入（因为他们航行在海上），通过极小的孔洞，通过缝隙进入船舱，他们用善工将水抽出去，以免水越积越多，沉没船只，每天抽水，禁食、祈祷、施舍，以纯洁的心说：\Quote{宽免我们的罪债，如同我们也宽免我们的罪人}\BibleRef{玛 6:12}——说着这样的话，你行走安稳，在路中欢跃，在路中歌唱。不要害怕审判者：在你成为信徒之前，你已找到一位救主。祂寻找你这不敬之人，为的是救赎你；已经得赎的你，祂会弃绝而毁灭吗？\Quote{你在世上引导万民。}\par

8. 祂欢跃、喜乐、劝勉；在劝勉中，祂重复了同样的诗句。\Quote{大地已给出了她的果实}\BibleRef{咏 66:6}。什么果实？\Quote{愿万民都向祢承认。}大地原先是荆棘丛生；来了那拔除荆棘者的手，来了祂尊威与慈悲的召唤，大地开始承认；如今大地给出了她的果实。若非先被雨水浇灌，她会给出果实吗？若非天主的慈悲从天而降，她会给出果实吗？你们说，让我读给他们听：大地如何因被雨水浇灌而给出果实。请听那降在大地上的主的话：\Quote{你们悔改吧，因为天国临近了。}\BibleRef{玛 3:2} 祂降雨，而同样的雨也是雷声；它令人恐惧：你当畏惧那打雷的祂，并领受那降雨的祂。看，在那雷声与降雨的天主的声音之后，在那声音之后，让我们从福音本身看一些事。看，那城中的声名狼藉的妓女闯入了一所陌生的房屋，她并非被主人邀请，而是被那一位被邀请者所召叫；\BibleRef{路 7:37} 不是用舌头召叫，而是藉着恩宠。那患病妇人知道她在那里有一个地方，她知道她的医生正在那里坐席。她进去了，她是一个罪妇；她不敢靠近，只到脚前；她在祂脚前哭泣，用眼泪洗濯，用头发擦干，用香膏涂抹。你为何惊讶？大地已给出了她的果实。我说，这事是因主藉着祂自己的口在那里降雨而发生的；我们在福音中读到的事发生了；而藉着祂藉着云彩降雨，藉着派遣宗徒以及他们宣讲真理，大地更加丰盛地给出了她的果实，那庄稼如今已充满了世界。\par

9. 大地的果实首先在耶路撒冷。因为教会从那里开始：圣神降临在那里，充满了聚集在一处的圣人们；奇迹发生了，他们用众人的舌音说话。他们充满了天主之神，当地的人民悔改，畏惧并领受神圣的雨水，藉着告解结出了如此多的果实，以致他们把所有的财物聚集在一起，分给穷人，为使人不称任何东西为自己的，而一切在他们当中成为公有，并且他们一心一意归向天主。\BibleRef{宗 4:32} 因为他们所流的血已被赦免，被宽恕的主赦免了，为使他们如今甚至学会喝他们所流出的血。那地方果实丰盛：大地已给出了她的果实，既是大果实，又是最美好的果实。难道只有那大地该给出她的果实吗？\Quote{愿天主，我们的天主，祝福我们；愿天主祝福我们}\BibleRef{咏 66:7}。愿祂仍祝福我们：因为祝福在于增多，通常最首要且最恰当地被理解。让我们在\BibleBook{}{创世纪}中证明这一点；看看天主的工程：天主造了光，\BibleRef{创 1:3} 天主分开光与黑暗：祂称光为昼，称黑暗为夜。并没有说祂祝福了光。因为同样的光因昼夜而返回并变化。祂称天空为穹苍，在诸水之间；并没有说祂祝福了天空：祂把海与干地分开，并命名了二者，称干地为地，称水的聚集为海：这里也没有说天主祝福……\par

10. 我们应如何愿意祂\Emph{来到我们这里}？藉着善生，藉着善行。不要让过去的事取悦我们；不要让现在的事抓住我们；我们不要像用尾巴\Quote{堵住耳朵}一样，不要将耳朵压在地上；免得因过去的事我们被阻止聆听，免得因现在的事我们被缠住而无法默想未来的事；让我们伸向前面的事，忘记背后的事。\BibleRef{斐 3:13} 我们如今劳苦、叹息、渴望、谈论的事，我们稍许（无论多么微小）领悟却无法领受的，将在义人的复活中领受，并完全享受。我们的青春将如鹰一般更新，只要我们愿意将我们的旧人在基督磐石上撞碎。弟兄们，关于蛇所说的那些话，或关于鹰所说的那些话，究竟是真实的，还是更似人的传说而非真理，然而圣经中却有真理，圣经并非无缘无故提及此事：让我们行其所指，而不必费力探究它有多真实。你要成为这样的人，使你的青春能如鹰般更新。并且要知道，除非你的旧人在磐石上被撞碎，否则它无法被更新；也就是说，除非藉着磐石的助佑，藉着基督的助佑，你不可能被更新。你不要因过去生活的愉快而对天主圣言充耳不闻；不要因现在的事而被如此控制缠住，以致说：我没有时间阅读，我没有时间聆听。这就是把耳朵压在地上。因此，你不要成为这样的人；而要成为你在另一端所发现的那样的人，即忘记背后，努力向前，以便将你的旧人在磐石上撞碎。如果为你作了一些比喻，如果你在圣经中找到了它们，就相信；如果你没有找到它们，只是听人传说，就不要太相信它们。事情本身或许如此，或许并非如此。你当从中获益，让那比喻有助于你的救恩。你不愿从这个比喻获益，却从另一个比喻获益，这无关紧要，只要你去做；并且，安心等候天主的国，免得你的祈祷与你争吵。因为，基督徒啊，当你说\Quote{愿祢的国来临}时，你如何说\Quote{愿祢的国来临}？\BibleRef{玛 6:10} 省察你的心：看，请注意，\Quote{愿祢的国来临}：祂向你呼喊：\Quote{我来了}，你难道不害怕吗？我们曾多次告诉你们的爱德：若心与舌相违，宣讲真理便毫无用处；若听了真理而无果效，听真理也毫无用处。我们仿佛从这崇高之处向你们说话：但我们在恐惧中多么伏在你们脚下，天主知道，祂对谦卑的人仁慈；因为人的赞美之声并不如人忏悔的虔诚和如今义人的行为那样令我们喜悦。我们只以你们的进步为乐，然而这些赞美给我们带来多大危险，祂知道，我们祈求祂从一切危险中解救我们，并愿祂屈尊在祂的国内，在每一次考验中拯救我们，认识我们并与我们一起加冕。\par

\SourceRecord{Translated by J.E. Tweed.；Nicene and Post-Nicene Fathers, First Series；Vol. 8.；Edited by Philip Schaff.；Buffalo, NY: Christian Literature Publishing Co.,；1888.；<http://www.newadvent.org/fathers/1801067.htm>.}

% structure: p0001 source=h1 target=section reason=page-title

\section{圣咏第六十八篇释义}

1. 论这篇圣咏，标题似乎不需要费力的讨论：因为它显得简单而容易。因为是这样的：\Quote{为终向，为达味自己，一首诗歌的圣咏。}但我们在许多圣咏中已经提醒你们什么是\Quote{终向：因为法律的终向是基督，使所有相信的人成义。}\BibleRef{罗 10:4}祂是使人成全的终向，而非消耗或毁灭的终向。然而，若有人试图探究\Quote{诗歌的圣咏}是什么意思，为何不单是\Quote{圣咏}或\Quote{诗歌}，而是两者；或者诗歌的圣咏与圣咏的诗歌有何区别，因为即使这样，有些圣咏的标题也是这样题写的：他或许会发现一些东西，我们留给更敏锐、更有空闲的人去研究……\par

2. \BibleQuote{愿天主兴起，使祂的仇敌四散}\BibleRef{咏 67:1}。这事已经发生，基督已经兴起，\BibleQuote{祂是在万有之上，永远可赞美的天主}\BibleRef{罗 9:5}，祂的仇敌，即犹太人，已分散到万邦；就在他们施敌意的地方，他们在战争中被推翻，从此分散到各处：现在他们仇恨，但害怕，并且就在这害怕中，他们做了随后的事：\Quote{愿恨祂的人从祂面前逃跑。}心灵的逃跑确实是害怕。因为就肉体的逃跑而言，他们能从那位处处显示祂临在效能者的面前逃到哪里呢？\Quote{我往哪里去躲避你的神，}他说，\Quote{从你的面前我往哪里逃？}因此，他们是用心灵，而不是用身体逃跑；也就是说，由于害怕，而不是由于隐藏；并且不是从他们看不见的那张脸，而是从他们被迫看见的那张脸。因为祂的面就是祂在教会中的临在被称为……\par

3. \BibleQuote{如烟消散，愿他们消散}\BibleRef{咏 67:2}。因为他们从仇恨的火焰中升腾到骄傲的烟气，口出狂言攻击上天，呼喊说：\BibleQuote{钉死，钉死}\BibleRef{若 19:6}。他们戏弄被俘的祂，嘲笑悬挂的祂：但很快就败在他们曾趾高气扬反对的那位之下，他们消失了。\Quote{如蜡在火前熔化，愿罪人在天主面前灭亡。}虽然在本段中他可能指那些人的硬心在忏悔的眼泪中溶解；然而这也可理解为他在警告未来的审判；因为虽然在今世，他们像烟一样自高（即骄傲）而消散，但最终他们要受永罚，以致从祂面前永远灭亡，当祂在自己的荣耀中显现时，如同火，为惩罚不敬虔的人，并为义人的光明。\par

4. 最后，接着是：\BibleQuote{愿义人喜乐，并因天主面前而欢欣，愿他们因喜乐而愉悦}\BibleRef{咏 67:3}。因为那时他们将听到：\Quote{来吧，我父所祝福的，你们领受国度吧。}\Quote{愿他们喜乐，}因此，那些劳苦的人，\Quote{并因天主面前而欢欣。}因为在这欢欣中，不像在人前，没有虚空的夸耀；而是在那位无误地看顾祂所赐予者的面前。\Quote{愿他们因喜乐而愉悦：}不再像今世那样战战兢兢地欢欣，只要\Quote{人生在世是一场考验。}其次，他转向那些他赐予如此大希望的人，并在他们活在世上的时候对他们说话和劝勉：\BibleQuote{你们要向天主歌唱，歌颂祂的名}\BibleRef{咏 67:4}。关于这一主题，我们在标题的释义中已经说过合适的话。他向天主歌唱，是向天主而活；他歌颂祂的名，是为祂的荣耀而工作。如此歌唱，如此歌颂，也就是如此生活、如此工作，\Quote{你们要为祂修路，}他说，\Quote{那一位已升到没落之上。}你们要为基督修路：这样通过传报喜讯者的佳美脚踪，\BibleRef{依 52:7}许多相信者的心就有一条路向祂敞开。因为那位已升到\Quote{没落}之上的，就是祂：要么因为归向祂的新生命若不先让旧生命因弃绝世俗而没落，就不能接受祂；要么因为祂在复活时征服了身体的堕落，从而升到没落之上。\Quote{因为主是祂的名。}他们若认识祂，就绝不会钉死光荣的主。\BibleRef{格前 2:8}\par

5. \Quote{在祂面前欢乐。} 你们这些人，曾被告知：\Quote{向天主歌唱，歌颂祂的名，为那上升至日落之上者修路}，也\Quote{在祂面前欢乐}：好像\BibleQuote{似乎忧愁，却常常喜乐}（\BibleRef{格后 6:10}）。因为当你们为祂修路，为祂预备道路，使祂来并占有万民时，你们要在人前忍受许多忧愁的事。但不仅不要丧胆，反而要欢乐，不是在人前，而是在天主面前。\BibleQuote{在指望中喜乐，在患难中忍耐}（\BibleRef{罗 12:12}）\Quote{在祂面前欢乐。} 因为那些在人前烦扰你们的人，\BibleQuote{他们必因祂的面容而惊惶，祂是孤儿的慈父，是寡妇的审判者}（\BibleRef{咏 67:5}）。因为他们以为这些人被遗弃了，这些人常常因天主圣言的剑（\BibleRef{玛 10:34}）而父母与子女、丈夫与妻子被分离；但那些孤苦和寡居的人却得到\Quote{孤儿的慈父和寡妇的审判者}的安慰；他们得到祂的安慰，就是那向祂说：\Quote{因我的父母离弃了我，但主却收留了我}；以及那些仰望主、日夜专务祈祷的人（\BibleRef{弟前 5:5}）。当那些人看见自己无法得逞，因为举世都去跟随祂时（\BibleRef{若 12:19}），他们将因祂的面容而惊惶。因为从这些孤儿寡妇中，即那些在这世界盼望中无分的人，主为自己建造了一座圣殿；为此祂接着说：\Quote{主在祂的圣所里。}\par

6. 因为祂在说\BibleQuote{天主使同心合意的人同住一室}（\BibleRef{咏 67:6}）时，已揭示了祂的地方：同心合意、同感一灵的人，这就是主的圣所。因为当祂说\Quote{主在祂的圣所里}时，好像我们询问在什么地方，既然祂充满万有，没有空间能容纳祂，祂立刻补充了一些话，使我们不要在自己之外寻找祂，而要同心合意同住一室，好配得祂也亲自住在我们中间。这就是主的圣所，是大多数人寻求拥有的地方，一个在祈祷中蒙垂听的地方。……因为如同在人的大宅中，主人并非住在每一处，而是住在更私密、更尊贵的地方；同样，天主不是住在祂房屋中的所有人之内（因为祂不住在卑贱的器皿中），但祂的圣所是那些人，就是祂\Quote{使同心合意}或\Quote{使习惯相同}的人同住一室。因为希腊文称为\OriginalTerm{方式}{\GreekText{τρ}\'{\GreekText{ο}}\GreekText{ποι}}的，在拉丁文既可译为\Emph{modi}（方式），也可译为\Emph{mores}（习惯）。所以\Quote{主在祂的圣所里。}\par

7. 但是为证明祂以祂的恩宠为自己建造这地方，并非因为那些被建造之人的先功德，请看下文：\Quote{祂以大能领出被捆绑的人。} 因为祂解开罪的沉重枷锁，那些人被捆绑以致不能行走在诫命的道路上；但祂以大能领出他们，这大能是他们在祂的恩宠之前所没有的。\Quote{同样，惹怒祂、住在坟墓中的人。} 即完全死去，从事死工的人。因为这些人抗拒正义，惹祂发怒：那些被捆绑的人也许愿意行走却不能，向天主祈求能行，对祂说：\Quote{从我的患难中领我出来。} 他们蒙垂听，就感谢说：\Quote{祢解开了我的捆绑。} 但这些惹怒祂、住在坟墓中的人，是另一种，圣经在另一处指出说：\BibleQuote{从死人，如同从不存在的人那里，忏悔就消失了。}（\BibleRef{德 17:28}）因此有话说：\BibleQuote{当罪人陷入邪恶深处时，他就藐视。}（\BibleRef{箴 18:3}）因为一方面渴望，另一方面抗拒正义，是两回事：一方面想从恶中得救，另一方面宁愿辩护自己的恶行而不愿忏悔：然而这两种人都蒙基督的恩宠以大能领出来。有什么大能？不就是那使人抵抗罪恶直到流血的能力吗？因为从每种人中，都有适合的人用来建造祂的圣所：那些被释放的，这些被复活过来的。因为就连那妇人，撒殚捆绑她十八年，祂一声令下就解开了捆绑（\BibleRef{路 13:16}）；拉匝禄的死，祂一声呼唤就战胜了（\BibleRef{若 11:43}）。那在身体上做了这些事的主，更能奇妙地在性格中行事，并使人同心合意同住一室：\Quote{祂以大能领出被捆绑的人，同样也领出惹怒祂、住在坟墓中的人。}\par

8. \BibleQuote{天主，当祢在祢百姓面前出行时} \BibleRef{咏 67:7}。祂的出行被感知，当祂在祂的工程中显现时。但祂并不显现给所有人，只显现给那些懂得观察祂工程的人。因为我现在不谈那些对所有人都显而易见的工程——天、地、海和其中所有的一切；而是那些工程，藉此祂领出被力量束缚的人，同样也领出那些住在坟墓中挑衅的人，并使他们同心合意地住在一间房屋里。这样，祂就在祂的百姓面前出行，也就是说，在那些领悟祂这恩宠的人面前出行。最后，紧接着的是：\BibleQuote{当祢在旷野中经过时，大地震动} \BibleRef{咏 67:8}。外邦人曾是旷野，他们不认识天主；他们是旷野，在那里天主自己从未赐予法律，没有先知居住，也没有预言主要来临。\Quote{当祢，}然后，\Quote{在旷野中经过，}当祢在外邦中被宣讲时；\Quote{大地震动，}属地上的人被激动而相信。但它为何震动？\Quote{因为诸天从天主的面前滴下。}也许有人会想起那时，当在旷野中天主在祂的百姓面前经过，在以色列子民面前，白天在云柱中，夜间在火柱的光辉中；\BibleRef{出 13:21}并断定正因如此，\Quote{诸天从天主的面前滴下，}因为祂为祂的百姓降下了玛纳：\BibleRef{出 16:15}而同一件事也出现在随后的话中：\Quote{西乃山从以色列天主的面前，} \Quote{祢以自愿的雨为祢的产业分开天主} \BibleRef{咏 67:9}，即那位在西乃山上向\Person{梅瑟}说话的天主，当祂赐下法律时，以致玛纳就是那自愿的雨，是天主为祂的产业，即为祂的百姓所分开的；因为祂只这样喂养了他们，并没有喂养其他民族：所以接下来所说的\Quote{它被削弱了，}被理解为产业本身被削弱了；因为他们发怨言，厌弃玛纳，贪图肉食和他们从前在\Place{埃及}所习惯的那些食物。\BibleRef{户 11:5}…最后，在旷野中所有那些人都被击倒，除了两人之外，没有一人被认为配进入预许之地。\BibleRef{户 14:23}即使有人说在他们的儿子中那产业得以完成，我们仍更应坚持属灵的意义。因为那些事都是以预像的方式发生在他们身上，\BibleRef{格前 10:11}直到破晓，阴影消散。\BibleRef{歌 2:17}\par

9. 愿主开启我们这些敲门者；并愿祂奥迹的隐秘事理，按祂所乐意赐予的，得以显明。因为为了使大地被感动而趋向真理，当福音传向外邦人的旷野时，\Quote{诸天从天主的面前滴下。}这些就是诸天，在另一篇圣咏中歌颂道：\Quote{诸天述说天主的荣耀。}…同样在此，\Quote{诸天滴下；}但是\Quote{从天主的面前。}因为甚至这些人本身，他们\BibleQuote{藉着信德得救，并且这不是出于自己，而是天主的恩赐，不是出于行为，免得有人自夸。因为我们原是祂的化工，} \BibleRef{弗 2:8} \Quote{使人心神一致地住在屋子里。}\par

10. 但接着的话，\BibleQuote{从以色列的天主面前，西乃山}是什么意思？是否应理解为\Quote{滴落}，使祂以天的名义称呼的，也以西乃山的名义表达？正如我们说那些山被称为天？这个意思也不应使我们困扰，因为祂说的是\Quote{山}，而不是\Quote{群山}，而在那里被称作\Quote{天}，而不是\Quote{诸天}；因为在另一篇圣咏中，在说了\BibleQuote{诸天宣扬天主的光荣}之后，按照圣经用不同词语重复同一意义的习惯，随后又说\BibleQuote{穹苍述说祂手的化工}。祂先说\Quote{诸天}，而非\Quote{天}；但后来不是\Quote{诸穹苍}，而是\Quote{穹苍}。因为天主称穹苍为天，\BibleRef{创 1:8}，正如创世记所记载的。这样，诸天和天、群山和山，并不是不同的东西，而是完全相同的东西；就像许多教会和唯一的教会，不是不同的东西，而是完全相同的东西。那么，为什么\BibleQuote{西乃山，生子为奴}\BibleRef{迦 4:24}，如宗徒所说？是否西乃山应当理解为法律本身，即\BibleQuote{诸天从天主面前滴落}，使大地震动？而这正是大地的震动，当人们因无法履行法律而困扰时？若是如此，这就是自愿的雨，为证实这一点祂说：\BibleQuote{天主，祢为祢的产业降下自愿的雨}；因为\BibleQuote{祂没有这样对待任何民族，也没有向他们显示祂的律法}。因此天主将这自愿的雨分别出来赐给祂的产业，因为祂赐下了法律。而\BibleQuote{变得软弱}，或者是法律，或者是产业。法律可以被理解为变得软弱，因为它没有被履行；并非本身软弱，而是因为它通过威胁惩罚使人软弱，并不通过恩宠帮助。因为宗徒也用了同样的词，他说：\BibleQuote{因为法律所不能做的，因它因肉体而软弱}\BibleRef{罗 8:3}，意欲暗示通过圣神它得以成全；然而，他说法律本身变得软弱，因为软弱的人无法履行它。但产业，即百姓，无疑被理解为因赐下法律而变得软弱。因为\BibleQuote{法律介入，使过犯增多}\BibleRef{罗 5:20}。但接着的话\BibleQuote{但祢使它成全了}，是指法律而言，因为当主在福音中说\BibleQuote{我来不是为废除法律，而是为成全}\BibleRef{玛 5:17}时，法律得以成全。……在这些话中还有另一种意义，在我看来更合理。因为根据上下文，恩宠本身被理解为自愿的雨，因为它不是因先前的功绩而赐予，而是白白地赐予。\BibleQuote{因为如果是恩宠，就不再是由于行为；否则恩宠就不再是恩宠}\BibleRef{罗 11:6}。……\BibleQuote{但祂赐恩宠给谦卑的人}\BibleRef{雅 4:6}。它变得软弱，但祢使它成全了，因为\BibleQuote{德能在软弱中得以成全}\BibleRef{格后 12:9}。确实，有些拉丁文和希腊文抄本没有写\Quote{西乃山}，而是\BibleQuote{从西乃的天主面前，从以色列的天主面前}。也就是说，\BibleQuote{诸天从天主面前滴落}；并且，好像要问是哪位天主，祂说\BibleQuote{从西乃的天主面前，从以色列的天主面前}，即从赐法律给以色列百姓的天主面前。那么，为什么\BibleQuote{诸天从天主面前滴落}，是从这位天主面前，不正是因为应验了所预言的\Quote{祂赐福，祂赐下了法律}吗？法律用以恐吓依赖人力的人；祝福用以拯救仰望天主的人。因此，哦天主，祢成全了祢的产业；因为它在自身中变得软弱，好由祢使之成全。\par

11. \BibleQuote{祢的牲畜住在其中}\BibleRef{咏 67:10}。\Quote{祢的}，不是牠们自己的；隶属于祢，不为自身自由；为祢匮乏，不为自足。最后，祂接着说：\BibleQuote{天主，祢以祢的甘甜为穷人预备了。}\Quote{以祢的甘甜}，不是以他的相称。因为他是穷人，他变得软弱，好使他得以成全；他承认自己贫乏，好被充满。这就是那甘甜，在另一处说：\Quote{主必赐甘甜，我们的地也要出产果实}；好使善行不是出于恐惧，而是出于爱；不是出于对惩罚的畏惧，而是出于对正义的爱。因为这才是真实且健全的自由。但主为匮乏者预备了这甘甜，不是为富足者，后者的贫穷是一种耻辱；正如在另一处所说：\Quote{这是对富足之人的羞辱，对骄傲之人的轻视。}因为祂称那些骄傲的人为富足者。\par

12. \\Quote{主将赐予圣言} \\BibleRef{咏 67:11}：即作为祂将要居住其中的动物的食物。但这些承受圣言的动物将做什么呢？岂非那随后所说的事？\\Quote{给他们以丰厚的德能宣讲福音。} 是什么样的德能？岂非那带领被捆绑之人出来的力量？此处他或许也指那德能，即宣讲福音时他们以之行了奇事。那么，谁\\Quote{将赐予圣言给以丰厚的德能宣讲福音的人}？\\Quote{是君王，}他说，\\Quote{\\OriginalTerm{爱子}{Beloved}之德能的君王} \\BibleRef{咏 67:12}。因此，圣父是圣子之德能的君王。因为当没有指定任何特定的人为爱子时，通过名称的替代，理解为独生子。难道圣子自己不就是祂德能的君王吗？即那些服侍祂的德能的君王？因为德能的君王将以丰厚的德能赐予圣言给宣讲福音的人，关于祂曾如此说：\\Quote{万军的上主，祂是光荣的君王？} 但圣咏作者没有说\\Quote{德能的君王}，而是说\\Quote{爱子之德能的君王}，这是圣经中极常见的表达，只要人留意；这尤其在那些即使已经表达了人自己的名字的情况下出现，以至于毫无疑问说的是同一个人。此类表达在梅瑟五书中多处可见：\\Quote{梅瑟便这样做了，上主怎样吩咐梅瑟，他就怎样做了。} 他没有使用我们通常的表达，即\\Quote{梅瑟便这样做了，上主怎样吩咐他}，而是说\\Quote{梅瑟便这样做了，上主怎样吩咐梅瑟}，仿佛受吩咐的梅瑟是一个，而遵行的梅瑟是另一个，然而却是同一个人。在新约中极难找到这类表达。\\BibleRef{罗 1:3} ...\\Quote{因此，君王，} \\Quote{爱子之德能的，}可如此理解，仿佛是说，祂自己的德能的君王，因为基督既是德能的君王，而爱子就是同一基督。然而，这个理解并非具有如此大的紧迫性，以至于不能接受其他解释：因为圣父也可理解为祂的爱子之德能的君王，爱子自己曾对祂说：\\Quote{凡我所有的，都是你的；你所有的，也是我的。} \\BibleRef{若 17:10} 但若有人问，主耶稣基督的天主父是否也可被称为君王，我不知道有谁敢在宗徒说\\Quote{愿尊崇和光荣归于万世的君王，那不死不灭，不可见的，惟一天主}之处拒绝给祂这个名号呢？\\BibleRef{弟前 1:17} 因为即使这是对天主圣三本身说的，其中也有天主父。但若我们不按血肉理解\\Quote{天主，求你将审判权赐给君王，将你的正义赐给君王的儿子：}我不知道除了\\Quote{你的儿子}外还说了什么。因此，圣父也是君王。因此，这首圣咏中的诗句\\Quote{爱子之德能的君王；}两方面都可理解。所以，当他说了\\Quote{主将赐予圣言给以丰厚的德能宣讲福音的人：}因为德能本身由祂统辖，并服侍赐予它给祂的那一位；祂自己，他说，祂将赐予圣言给以丰厚的德能宣讲福音的人，就是爱子之德能的君王。\par

13. 接下来有：\\Quote{关于爱子，以及关于房屋的美丽以分配战利品。} 这个重复属于赞美之词.. ..但无论它是重复出现，还是被视为只说一次，已经写下的词，即\\Quote{爱子}，我猜想必须如此理解那随后的话：\\Quote{以及关于房屋的美丽以分配战利品；}就好像是说了\\Quote{甚至是为分配房屋之美丽的战利品而被拣选的}，那就是，甚至是为分配战利品而被拣选的。因为基督借着将战利品分给祂的房屋，使祂的房屋，即教会，变得美丽；正如身体因肢体的分布而美丽一样。此外，那些从被击败的敌人身上剥下的东西被称为战利品。这一点，福音在以下经文提示我们：\\Quote{没有人能进入壮士的家，抢夺他的器具，除非先捆绑住那壮士。} \\BibleRef{玛 12:29} 因此，基督以属灵的束缚捆绑了魔鬼，借克服死亡，并借从地狱升到诸天之上：祂借祂降生成人的圣事捆绑了他，因为虽然在祂身上找不到任何该死的理由，却仍被允许杀害祂；并从被他如此捆绑的魔鬼那里，取走了他的器具如同战利品。因为他曾在悖逆之子身上运行，\\BibleRef{弗 2:2} 以他们的不信来成就自己的意愿。这些器具，主借着赦免罪过予以洁净，圣化这些从被击倒并捆绑的仇敌手中夺来的战利品，并将它们分给祂房屋的美丽；有的立为宗徒，有的立为先知，有的立为牧者和教师，\\BibleRef{弗 4:11} 为圣职的工作，为建立基督的身体。因为身体是一个，却有许多肢体；虽然身体的肢体很多，身体却只是一个；基督也是这样。\\BibleRef{格前 12:12} \\Quote{众人岂都是宗徒？岂都是先知？岂都是有奇能者？岂都享有治病的恩宠？岂都说各种语言？岂都解释？} \\BibleRef{格前 12:29} \\Quote{然而这一切都是同一圣神所作，祂随自己的心意，个别分配给各人。} \\BibleRef{格前 12:11} 这就是房屋的美丽，战利品分给之处，以致一位因这美丽而被点燃的喜爱者，呼喊说：\\Quote{上主，我喜爱祢房屋的恩宠。}\par

14. 现在，在接下来的经文中，他转而直接对组成圣殿之美的肢体说话，说：\\Quote{你们若在份中安睡，银白的鸽翼，与在其两肩之间的金辉} \\BibleRef{咏 67:13}。首先，我们必须在这里考察词句的顺序，句子如何结束；当说到\\Quote{你们若安睡}时，这当然有待确定；其次，当他说的\\Quote{银白的鸽翼}，是应以单数理解为\\Quote{这只翅膀}，还是以复数理解为\\Quote{这些翅膀}。但希腊文排除了单数，我们在那里读到总是以复数写成。但仍然不确定是这些翅膀，还是\\Quote{哦，你们翅膀}，如此他好像是在对翅膀本身说话。因此，无论是通过前面的词句，句子结束，顺序是\\Quote{主赐予圣言给以大能力传福音的人，你们若在份中安睡，哦，银白的鸽翼：}或是通过后面的，顺序是\\Quote{你们若在份中安睡，银白的鸽翼在\\Place{塞尔孟}将被雪洗白：}也就是说，翅膀本身将被洗白，如果你们在\\Quote{份中}安睡：因此他可以被理解为对分赐给圣殿之美的人，如同战利品，即，如果你们在\\Quote{份中}安睡，哦，你们这些分赐给圣殿之美的人，\\Quote{藉着圣神显示于益处} \\BibleRef{格前 12:7}，这样\\Quote{确实藉圣神赐予一人智慧之言，另一人知识之言}等等，如果你们在份中安睡，那么银白的鸽翼将在\\Place{塞尔孟}被雪洗白。也可以这样：\\Quote{如果你们作为银白的鸽翼，在份中安睡，它们将在\\Place{塞尔孟}被雪洗白}，如此那些藉恩宠获得罪赦的人应当被理解。因此，在\\WorkTitle{雅歌}中也对教会本身说：\\Quote{那上来洗白的是谁？}因为天主的这项应许通过先知宣告，说：\\Quote{如果你们的罪如朱红，我必使其白如雪。}也可以这样理解，以至于在所说的\\Quote{银白的鸽翼}中，应理解为你们将成为，这样意思是：哦，你们这些如同战利品分赐给圣殿之美的人，如果你们在\\Quote{份中}安睡，你们将成为银白的鸽翼：也就是说，你们将被提升到更高之处，然而依附于教会的纽带。因为我认为在此所能更好了解的银白鸽子，无过于那曾说的\\Quote{我的鸽子是唯一的} \\BibleRef{歌 6:9}。但她是银白的，因为她被神圣的言教所教导：因为在另一处，主的言被称为\\Quote{练过的银子，七次被火炼净}。因此，在份中安睡是一件伟大的善事，有些人认为这指的是两部盟约，如此\\Quote{在份中安睡}就是安息于那些盟约的权威，即，同意任一盟约的见证：这样，每当从中提出并证明任何事物时，一切争论便在和平的同意中结束……\par

15. \\Quote{在两肩之间}。这确实是身体的一部分，是心脏区域的部分，然而是在后部，即背部：他说那银白鸽子的这部分是\\Quote{在黄金的鲜嫩中}，即在智慧的力量中，这力量我认为没有比爱更好的理解。但为什么在背部，而不是在胸前？虽然我奇怪这句话在另一篇圣咏中是如何放置的，那里说：\\Quote{祂必在两肩之间荫庇你，你当在祂翼翅下寄望：}因为翼翅下除了在胸前之外，没有东西能被荫庇。而在拉丁文中，确实\\Quote{在两肩之间}，或许在一定程度上可理解为两部分，前和后，我们可以认为肩膀是头在两旁的部分；而在希伯来文中，或许这个词是歧义的，也可以以此方式理解：但在希腊文中，这个词是\\OriginalTerm{在两肩之间}{\GreekText{μετ}}\"\par

\SourceRecord{Translated by J.E. Tweed.；Nicene and Post-Nicene Fathers, First Series；Vol. 8.；Edited by Philip Schaff.；Buffalo, NY: Christian Literature Publishing Co.,；1888.；<http://www.newadvent.org/fathers/1801068.htm>.}

% structure: p0001 source=h1 target=section reason=page-title

\section{圣咏第六十九篇释义}

我们生于这个世界，并被增添到天主之民中，正值那芥菜籽的嫩芽已伸展其枝的时刻；正值那起初微不足道的酵母已发酵了三升面——即由诺厄的三个儿子重聚的整个世界的时候，\BibleRef{创 9:19}因为从东方、西方、北方和南方，将有人来与圣祖们同席，\BibleRef{玛 8:11}而那些由他们的肉身所生却未效法他们信德的人，将被驱逐在外。因此，我们已睁眼看见基督教会的光荣；那曾被告知、预言要因有更多儿子而喜乐的荒胎妇人，我们已发现正是如此，她已忘却了寡妇的羞辱和耻辱：所以我们读到任何预言中基督受辱或我们自己的话时，或许会感到惊讶。而且，我们可能对它们感触不深，因为我们没有生活在那些苦难重重、这些话语被热切诵读的时代。但若我们再思量苦难的众多，并观察我们所行之路（如果我们确实行在其中），它何等狭窄，如何通过艰辛和苦难通向永恒的安息，\BibleRef{玛 7:14}以及在人事中所谓的幸福比不幸更可怕，因为不幸常从苦难中结出善果，而幸福却以邪恶的安全感腐蚀灵魂，为诱惑者魔鬼留出余地——当我们审慎正确地判断，如同那咸的牺牲所做的那样，\Quote{地上的人生是试炼}，并且没有人是绝对安全的，也不应该安全，直到他抵达那家园，从那里没有朋友离开，也没有敌人进入；即使在教会的光荣中，我们也承认我们苦难的声音：作为基督的肢体，在爱的联系中服从我们的元首，互相支持，我们将从圣咏中说出我们在这里从先我们殉道者所说的话；那苦难从起初直到终结是众人共同的……\par

2. 这篇圣咏的标题是：\Quote{直到终结，为那些将要改变的人，为达味本人。}现在你听到的改变是指向更好的；因为改变或是朝向更坏，或是朝向更好……那么，我们朝向更坏的改变，归于我们自己；我们朝向更好的改变，让我们赞美天主。\Quote{为那些}因此\Quote{将要改变的人}，这篇圣咏是为他们。但这改变从何而来，岂非藉着基督的苦难？"Pascha"一词在拉丁文中被解释为"逾越"。因为Pascha不是希腊文，而是希伯来文。它确实在希腊语中听起来像"苦难"，因为\GreekText{π}\par

\SourceRecord{Translated by J.E. Tweed.；Nicene and Post-Nicene Fathers, First Series；Vol. 8.；Edited by Philip Schaff.；Buffalo, NY: Christian Literature Publishing Co.,；1888.；<http://www.newadvent.org/fathers/1801069.htm>.}

% structure: p0001 source=h1 target=section reason=page-title

\section{圣咏第七十篇释义}

1. 感谢那\Quote{麦粒}，因为祂甘愿死去，为能繁衍增多；感谢天主的独生子、我们的主与救主耶稣基督，祂不屑于经历我们的死亡，为使我们堪当祂的生命。看，祂曾是独一的，直到祂离开此世；正如祂在另一篇圣咏中所说：\Quote{我独自行走，直到我离开这里。}因为祂是一粒麦子，却怀有极大的繁衍能力；当我们庆祝殉道者的诞辰时，有多少麦粒因效法祂的苦难而欢欣！因此，祂的许多肢体，联同我们的救主元首，因爱与和平的纽带（正如你们所知晓的，因为你们常听见）而结合，成为一人；如同一个人的声音常在圣咏中被听到，如此一人呼喊，仿佛众人呼喊，因为众人在一人内合而为一……\par

2. 因此，在这篇圣咏中，是受苦之人的声音，确实是殉道者在危难中的声音，但他们信赖自己的元首。让我们怀着同情之心倾听他们，并与他们交谈，尽管我们并未遭受同样的苦难。因为他们已获得冠冕，我们仍在危险中；并非同样的迫害折磨我们，而是更甚者或许是在如此巨大的恶表中间。因为我们这个时代更充满那主所哀叹的灾祸：\BibleQuote{祸哉，世界因了恶表！} \BibleRef{玛 18:7}且\BibleQuote{由于罪恶的增多，人的爱情必要冷淡。} \BibleRef{玛 24:12}因为即使在索多玛，圣洁的罗特也没有受到任何人的身体迫害，也没有人告诉他不应居住在那里；\BibleRef{创 19:19}他的迫害乃是索多玛人的恶行。如今基督坐在天上，如今祂受光荣，如今君王的颈项已屈服于祂的轭下，他们的前额已置于祂的标记之下，如今不再有人敢公开践踏基督徒；然而，我们仍在这种种乐器和歌声中呻吟，那些殉道者的仇敌，因言语和刀剑已无能为力，便以他们的放荡来迫害他们。但愿我们只为外邦人忧伤！这将是某种安慰：等待那些尚未被基督十字架所标记的人，等待他们被标记，因祂的权威而皈依，不再疯狂。此外，我们看见有些人额上佩戴着祂的标记，同时在那同一额上显露着放荡的无耻，在殉道者的纪念日与庆祝中，不是欢欣，而是侮辱。在这一切中我们呻吟，这就是我们的迫害，倘若我们心中有那爱说：\BibleQuote{有谁软弱，我不软弱呢？有谁跌倒，我不心焦呢？} \BibleRef{格后 11:29}那么，没有天主的仆人是没有迫害的。宗徒的话是真实的：\BibleQuote{凡是愿意在基督内虔诚生活的，都必要遭受迫害。} \BibleRef{弟后 3:12}……\par

3. \BibleQuote{天主，求你快来援助我。} \BibleRef{咏 69:1}因为我们在这世界上需要永久的援助。但我们何时不需要呢？如今在困苦中，让我们尤其要说：\Quote{天主，求你快来援助我。}\Quote{愿那些寻求我性命的，受辱而惊惶。}基督在说话：无论元首说话还是身体说话；是那位说过\BibleQuote{你为什么迫害我？} \BibleRef{宗 9:4}的祂在说话；是那位说过\BibleQuote{你们对我这些最小兄弟中的一个所做的，就是对我做的。} \BibleRef{玛 25:40}的祂在说话。因此，这人的声音，即整个人的声音，包括元首和身体，是为人所知的；这无需常提，因为已众所周知。\Quote{愿他们受辱，}他说，\Quote{并惊惶，那些寻求我性命的。}在另一篇圣咏中祂说：\Quote{我向右观望，却没有人认识我；逃路已为我断绝，也没有人寻找我的性命。}在那里祂说到迫害者，没有人寻找祂的性命；但此处却说：\Quote{愿那些寻求我性命的，受辱而惊惶。}……而你们从主那里听说的，\BibleQuote{你们要爱你们的仇敌，善待恨你们的人，为迫害你们的人祈祷} \BibleRef{玛 5:44}，又在哪里呢？看，你遭受迫害，却诅咒那些迫害你的人，你怎么效法你们的主先行的苦难呢？祂挂在十字架上说：\BibleQuote{父啊，宽恕他们，因为他们不知道他们所做的是什么。} \BibleRef{路 23:34}对说这样话的人，殉道者回答说：你将主摆在我面前，说：\Quote{父啊，宽恕他们，因为他们不知道他们所做的是什么。}你也理解我的声音，为使它也成为你的。因为我关于我的仇敌说了什么？\Quote{愿他们受辱而惊惶。}对殉道者的仇敌已施了这样的报复。那曾迫害斯德望的扫禄，他受辱而惊惶了。他正口吐杀气，寻找一些人拖去杀害；从天上有声音说：\BibleQuote{扫禄，扫禄，你为什么迫害我？} \BibleRef{宗 9:4}他就受辱而仆倒在地，并被提升至服从，他原是受迫害之火的煽动。因此，殉道者愿他们的仇敌\Quote{受辱而惊惶。}因为他们若不受辱而惊惶，就必为自己的行为辩护；他们认为自己荣耀，因为他们扣押、捆绑、鞭打、杀戮、跳舞、侮辱，并且因这一切行为终有一日受辱而惊惶。因为若他们受辱，他们也必皈依；因为他们若不先受辱而惊惶，就不能皈依。那么，让我们愿这些事临到我们的仇敌，让我们无所畏惧地愿这些事。看，我已经说了，也让我与你们同说：愿所有仍跳舞、歌唱、侮辱殉道者的人\Quote{受辱而惊惶}；最后在这围墙内，受辱的他们捶胸吧！\par

4. \BibleQuote{他们必转后退而且惭愧，因为对我思想邪恶的事}\BibleRef{咏 69:2} 起初是迫害者的攻击，现在却留下了恶意思想的他们。事实上，在教会中，有不同的迫害时期依次相随。当君王迫害时，教会受到了攻击；因为关于君王曾预言他们会迫害也会相信，当一件应验后，另一件就要随之而来。随之而来的事也发生了；君王相信了，教会得到了平安，教会开始被置于最高的尊严地位，甚至在这世上，在这生命中：但迫害者的咆哮并未停息，他们已将攻击转为思想。在这些思想中，正如在无底坑中，魔鬼已被捆绑，他咆哮却冲不出来。因为论到教会的这些时期，经上说：\Quote{罪人看见，便要恼怒。}他要做什么？他起初所做的？拖曳、捆绑、打击？他不这样做。然后呢？\Quote{他要咬牙切齿，而消瘦下去。}对这些人的行为，殉道者仿佛是愤怒的，然而他却为这些人祈祷。因为正如他对那些人曾表示好意，关于他们他曾说：\Quote{愿那寻找我灵魂的，蒙羞惧怕。}同样现在：\Quote{愿那对我思想邪恶的，转后退而且惭愧。}为什么？为使他们不超前，而是跟随。因为那批判基督宗教、并愿意按自己体系生活的人，他好像愿意超前于基督，仿佛基督确实犯了错误，软弱无能，因为祂要么愿意受苦，要么能在犹太人手中受苦；而他却聪明地防范这一切，逃避死亡，甚至卑鄙地说谎以逃避死亡，杀害自己的灵魂以便身体存活，他自认为是一个具有非凡而谨慎计谋的人。他在批判基督时超前，在某种意义上他超越了基督：让他相信基督，跟随基督。因为那正是为了思想邪恶的迫害者所愿望的，主自己曾对伯多禄说过。现在在某处，伯多禄愿意走在主前面……之前不久：\Quote{西满·巴尔约纳，你是有福的，因为不是血肉启示了你，而是我在天之父。}转瞬之间：\BibleQuote{退到我后面去，撒殚。}\BibleRef{玛 16:23} 什么叫做\Quote{退到我后面去}？跟随我。你愿意走在我前面，你愿意给我出主意，最好你跟随我的主意：这就是\Quote{退后}，退到我后面去。祂制止那超前的人，为使他退后；并且称他为撒殚，因为他愿意走在主前面。此前不久，\Quote{有福的}；现在，\Quote{撒殚}。为什么此前\Quote{有福的}？因为，祂说：\BibleQuote{给你启示的不是血肉，而是我在天之父。}为什么现在\Quote{撒殚}？因为祂说：\BibleQuote{你所体会的不是天主的事，而是人的事。}那么，我们这些愿意适当庆祝殉道者诞辰的人，要渴慕效法殉道者；不要愿意走在殉道者前面，自以为比他们更有见识，因为我们躲避了为正义和信德而受苦，而他们却没有躲避。因此，愿那些思想邪恶、在放荡中喂养自己心灵的人，\Quote{转后惭愧}。让他们后来从宗徒那里听到：\Quote{但那时你们在那些事上得了什么果实？就是你们现在所惭愧的。}\par

5. 接下来是什么？\BibleQuote{他们立刻转后退而惭愧，因为对我说：好啦，好啦}\BibleRef{咏 69:3} 迫害者有两种：辱骂者和谄媚者。谄媚者的舌头比杀人者的手更迫害人：因为圣经也称此为熔炉。的确，当圣经论到迫害时，它说：\BibleQuote{如同在熔炉中试炼了他们}（论及被杀的殉道者），\BibleQuote{并如同全燔祭的牺牲接纳了他们。}\BibleRef{智 3:6} 听听连谄媚者的舌头又是何等情形：他说：\BibleQuote{证明，}他说，\BibleQuote{银子金子的是火；但人由赞美他之人的舌头受到试验。}\BibleRef{箴 27:21} 那是火，这也是火：你应从两者中平安出来。批评者已破碎了你，你在熔炉中像瓦器一样被打碎。圣言塑造了你，然后苦难的考验来到：那已形成的，必须被坚固；如果塑造得好，就有火来增强。因此祂在受难中说：\Quote{我的力量像瓦片一样枯干了。}\par

6. 当他们都转面退避并羞愧时，无论是那些寻求我灵魂的人，还是那些对我图谋恶事的人，或是那些以邪僻虚伪的善意用言语缓和所施加打击的人，当他们自己被转离并困惑时，将会发生什么呢？\Quote{愿他们在祢内欢欣踊跃}，不是在我内，不是在这个人或那个人内，而是在祂内，他们原是黑暗，却已成了光明。\BibleQuote{愿寻求祢的所有人都欢欣踊跃于祢内} \BibleRef{咏 69:4}。寻求天主是一回事，寻求人是另一回事。\Quote{愿寻求祢的人欢乐}。那时，寻求自己的人将不会欢乐，因为祢在他們寻求祢之前已先寻求了他们。那羊尚未寻求牧人，它已离群，祂便下到它那里；\BibleRef{路 15:4} 祂寻求了它，并把它扛在肩上背回来。难道祂会轻视你吗？羊啊，你寻求祂，而祂在你轻视祂、不寻求祂时先寻求了你。现在，你该开始寻求那首先寻求了你、并把你扛在肩上背回来的祂。你要做祂所说的：\BibleQuote{我的羊听我的声音，我也认识他们，他们跟随我} \BibleRef{若 10:27}。如果你寻求那首先寻求了你的祂，并成了祂的羊，你听到你牧人的声音，跟随祂；请看祂向你展示了关于祂自己的什么，关于祂身体的什么，为使你关于祂不致犯错，关于教会不致犯错，以致无人对你说：那不是基督的却是基督，或那不是教会的却是教会。因为许多人说基督没有肉身，基督未以肉身复活：你不要跟随他们的声音。你要听牧人亲自的声音，祂披上了肉身，为寻求失落的肉身。祂已复活，并说：\BibleQuote{你们摸摸看，因为鬼神没有肉和骨，如同你们看到我有的} \BibleRef{路 24:39}。祂把自己显示给你，你要跟随祂的声音。祂也把教会显示出来，为使人不因教会之名欺骗你。\Quote{祂应当}，祂说，\BibleQuote{基督必须受苦，第三天从死者中复活，并要因祂的名向万邦宣讲悔改，以得罪之赦免，从耶路撒冷开始} \BibleRef{路 24:46}。你已得到你牧人的声音，不要跟随陌生人的声音；\BibleRef{若 10:5} 你若跟随了牧人的声音，就不必惧怕盗贼。但你如何跟随呢？如果你不对任何人说，好像凭他自己的功劳，\Quote{好，好}；也不以同样的言语为乐，以免你的头被罪人的油脂养肥。\Quote{愿所有寻求祢的人都在祢内欢欣踊跃，并愿他们说}——愿那些欢欣踊跃的人说什么？\Quote{愿主常常受显扬！}愿所有欢欣踊跃并寻求祢的人都说这话。什么？\Quote{愿主常常受显扬；是的，他们爱祢的救恩。}不仅是\Quote{愿主受显扬}，而且是\Quote{常常}... 你是罪人，愿祂受显扬，为使祂召唤；你认罪，愿祂受显扬，为使祂赦免；现在你生活正义，愿祂受显扬，为使祂引导；你坚持到底，愿祂受显扬，为使祂光荣。\Quote{愿主}，因此，\Quote{常常受显扬；是的，他们爱祂的救恩。}因为他们从祂那里获得救恩，而非从自己。我主天主的救恩，就是我们的主耶稣基督救主；凡爱救主的，就承认自己被医治；凡承认自己被医治的，就承认自己曾患病。不是他们自己的救恩，好像他们能自救；也不是人的救恩，好像能靠人得救。\Quote{不要}，他说，\Quote{信赖王侯，不要信赖人子，在他们那里没有救援。}为何如此？\Quote{救援属于上主，愿祢的祝福临于祢的百姓。}\par

7. 请看，\Quote{愿主受显扬}：你将永不吗？你将无处吗？在祂内有某种东西，在我内则无；但若在我内的一切都在祂内，那么是祂，而不是我。但你呢？\BibleQuote{但我困苦贫穷} \BibleRef{咏 69:5}。祂是富有的，祂是丰盈的，祂一无所缺。请看我的光明，请看我从何处被光照亮；因为我呼求：\Quote{上主，祢必要点亮我的灯}。那么你呢？\Quote{但我困苦贫穷}。我像一个孤儿，我的灵魂像一个孀居无倚无靠的寡妇：我寻求帮助，我一直承认我的软弱。我的罪已被赦免，现在我已开始遵从天主的诫命：然而，我仍然是困苦贫穷的。为什么仍然困苦贫穷？因为\BibleQuote{我发觉在我的肢体内另有一条法律，与我理智的法律交战} \BibleRef{罗 7:23}。为什么困苦贫穷？因为\BibleQuote{饥渴慕义的人是有福的} \BibleRef{玛 5:6}。我仍然饥饿，仍然口渴：我的满盈被延迟，并未被夺去。\Quote{天主啊，求祢援助我}。最合适的是，拉匝禄被解释为\Quote{受助者}：那个困苦贫穷的人，被送到亚巴郎的怀中；\BibleRef{路 16:23} 并预表教会，教会应当常常承认自己需要援助。这是真实的，这是虔敬的。\Quote{我曾对上主说：祢是我的天主}。为什么？\Quote{因为我不需要祢的财物}。祂不需要我们，我们需要祂：因此祂是真正的主。因为你并不是你仆人真正的主：两人都是人，都需要天主。但如果你以为你的仆人需要你，好使你给他面包；你也需要你的仆人，好使他助力你的劳作。你们彼此都需要对方。因此，你们中谁也不是真正的主，谁也不是真正的仆人。你要听真正的主，你是祂真正的仆人：\Quote{我曾对上主说：祢是我的天主}。为什么祢是主？\Quote{因为祢不需要我的财物}？但你呢？\Quote{但我困苦贫穷}。请看那困苦贫穷的人：愿天主喂养，愿天主减轻，愿天主援助：\Quote{天主啊，}他说，\Quote{求祢援助我}。\par

8. \BibleQuote{我的助佑者和拯救者是你；主，不要迟延。}你是助佑者和拯救者：我需要救助，求你帮助；我被缠住，求你拯救。因为除了你，没有人能从缠累中拯救。在我们周围，各种忧虑的罗网，从这边和那边，我们好像被荆棘和蒺藜撕扯，我们行走一条窄路，也许我们已陷入荆棘中：让我们向天主说：\BibleQuote{你是我的拯救者。}那给我们显示窄路的那一位，\BibleRef{玛 7:14}教导我们跟随它....\par

9. 什么是\BibleQuote{不要迟延}？因为许多人说，基督来临还有很长时间。那么，因为我们说\BibleQuote{不要迟延，}他就会在他决定来临之前来临吗？这个祈祷\BibleQuote{不要迟延}是什么意思？愿不要让你的来临在我看来来得太迟。因为在你看来似乎很久，对天主却不久，对他而言千年如一日，或如同更夫的三小时。但如果你没有忍耐，对你来说就是迟了；而当你觉得迟了，你就会偏离他，变得像那些在旷野中疲倦的人，急忙向天主求那些他在福地为你们保留的美好事物；而当在旅途中没有赐下这些美好事物（或许他们本会被这些事物败坏），他们就抱怨天主，内心回到埃及：那个身体已经离开的地方，内心却回去了。那么，你不可这样，不可这样：要敬畏主的话，他说：\BibleQuote{要记住\Person{罗特}的妻子。}\BibleRef{路 17:32}她也在路上，但已从索多玛人中被救出，回头看了；在她回头的地方，她就留在那里：她变成了一根盐柱，为了调教你。因为她被赐给你作为榜样，好使你明智，不要在迷途上止步。留意她的驻足而前行：留意她的回头，而你应向着前面的事物伸展开去，如同\Person{保禄}一样。\BibleRef{斐 3:13}什么是不要回头？\BibleQuote{忘记背后的事，}他说。因此你跟随，被召向天上的赏报，将来你要以此夸耀。因为同一宗徒说：\BibleQuote{为我保留了正义的冠冕，在那一天，主，公义的审判者，必将赐给我。}\BibleRef{弟后 4:8}\par

\SourceRecord{Translated by J.E. Tweed.；Nicene and Post-Nicene Fathers, First Series；Vol. 8.；Edited by Philip Schaff.；Buffalo, NY: Christian Literature Publishing Co.,；1888.；<http://www.newadvent.org/fathers/1801070.htm>.}

% structure: p0001 source=h1 target=section reason=page-title

\section{圣咏第七十一篇释义}

1. 在全部圣经中，那拯救我们的天主的恩宠将自己推荐给我们，为使我们被推荐。这首圣咏所歌颂的，正是我们打算讨论的……宗徒赞扬这恩宠：因这恩宠，他成了犹太人的敌人，他们夸耀法律的文字和自身的义德。他在所读的经文中赞扬这恩宠时，这样说：\Quote{因为我原是宗徒中最小的，不配称为宗徒，因为我迫害过天主的教会。}\BibleRef{格前 15:9}\Quote{但是，我蒙了怜悯，}他说，\Quote{因为我是在不信中无知地做的。}\BibleRef{弟前 1:13}稍后又说：\Quote{这话是确实的，值得完全接纳：就是基督耶稣来到世上，为拯救罪人，而我是其中的魁首。}\BibleRef{弟前 1:15}难道在他之前没有罪人吗？那么，他是魁首吗？是的，他在众人之前，不是在时间上，而是在邪恶的倾向上。\Quote{但是我蒙了怜悯，}他说，\Quote{是为叫基督耶稣在我这魁首身上显示他的全部坚忍，为给那些将要信从他而得永生的人作榜样。}就是说，每个罪人和不义的人，已经对自己绝望，已有角斗士的心态，以为既然必须被定罪，就可以为所欲为，但他看到宗徒保禄，他那么大的残忍和极恶的倾向竟被天主赦免，于是不再对自己失望，而转向天主。天主也在本圣咏中把这恩宠推荐给我们……\par

2. 这首圣咏的标题，一如往常，是一个标题，在门口就暗示了屋内所发生的事：\Quote{为耶曷纳达布的儿子们，并为那些首先被掳去的人，记载于达味本人。}耶曷纳达布（在耶肋米亚的预言中他向我们推荐）是某个人，他命令他的儿子们不喝酒，不住房屋，而住在帐幕里。儿子们遵守并持守了父亲的诫命，因此从上主那里得到了祝福。上主并没有命令这事，而是他们自己的父亲。但他们接受它，就好像它是上主天主的诫命一样；因为即使上主没有命令他们不喝酒、住在帐幕里，但上主曾命令儿子们应服从父亲。只有在一种情况下，儿子不应服从父亲：如果父亲命令了任何违背上主天主的事。因为父亲实在不应生气，当天主被置于他之上时。但当父亲命令的并不违背天主时，他必须如同天主一样被听从：因为服从父亲是天主所命令的。天主因耶曷纳达布的儿子们的服从而祝福了他们，并将他们丢在他不服从的百姓面前，责斥他们，因为耶曷纳达布的儿子们服从他们的父亲，而他们却不服从他们的天主。但耶肋米亚在谈论这些主题时，他对以色列百姓抱有这个目的：使他们预备自己要被掳到巴比伦去，不要指望别的，只指望他们将成为俘虏。于是这首圣咏的标题似乎由此染上了色彩，以致当他说\Quote{为耶曷纳达布的儿子们}之后，他又加上\Quote{并为那些首先被掳去的人}：并不是说耶曷纳达布的儿子们被掳去了，而是因为与那些将被掳去的人相对的是耶曷纳达布的儿子们，因为他们服从了他们的父亲；为使他们明白他们被掳，是因为他们没有服从天主。此外，耶曷纳达布被解释为\Quote{上主自愿的人}。这是什麼意思？上主自愿的人？自由地以意志事奉天主。什麼是上主自愿的人？\Quote{天主啊，我所许的愿都在你内，我将向你献上赞美的祭。}什麼是上主自愿的人？\Quote{我要自愿向你献祭。}因为如果宗徒的教导劝勉奴隶以善意而非出于勉强地事奉人间的奴隶主，并因自由地事奉而在心里成为自由人；那么，事奉天主更当以完全、圆满和自由的意志，因为祂能看见你的意志本身……第一个人使我们成了俘虏，第二个人已从俘虏中解救了我们。\Quote{因为就如所有的人在亚当内死了，同样所有的人在基督内也都要复活。}但在亚当内他们是因肉身的出生而死，在基督内他们是因心中的信德而得解救。你无权选择不从亚当出生：但你有权选择相信基督。所以，无论你多么愿意属于第一个人，你都会属于俘虏。而什麼叫愿意属于？或者什麼叫属于？你已经属于了，请呼喊：\Quote{谁能救我脱离这死亡的身体呢？}\BibleRef{罗 7:24}那么，让我们来听这个正在呼喊的人。\par

3. \Quote{天主，我寄望于你，上主，我永远不会蒙羞。}\BibleRef{咏 70:1}我已经蒙羞了，但不是永远。因为那被说\Quote{你们那时有什么果实？那些事现在使你们羞惭。}\BibleRef{罗 6:21}的人如何不蒙羞呢？那么要做什么，才能使我们永不蒙羞？\Quote{你们亲近他，就得以光明，你们的脸上就不会羞惭。}你在亚当内蒙羞了，离开亚当，亲近基督，那么你就不会蒙羞。\Quote{上主，我寄望于你，我永远不会蒙羞。}如果我在自己内现在蒙羞，在你内我永远不会蒙羞。\par

4. \\BibleQuote{以祢的义德解救我，拯救我} \\BibleRef{咏 70:2}。不是凭我的，而是凭祢的：因为若凭我的，我便成了祂所说的那种人：\\BibleQuote{他们不认识天主的义德，企图建立自己的义德，便不服从天主的义德} \\BibleRef{罗 10:3}。因此，\\Quote{以祢的义德}，不是我的。因为我的义德是什么？罪孽已经先去了。当我成为义人时，那将是祢的义德：因为我因祢赐给我的义德而成为义人；这义德将既是我的，也是祢的，即由祢赐给我的。因为我相信那使不虔敬的人成义的那一位，以致我的信德被算为义德 \\BibleRef{罗 4:5}。即使如此，义德将成为我的，然而却不是我的，仿佛是凭我自己赐给我自己的；如同那些凭文字自夸、并拒绝恩宠的人所想的……因此，你承认那在你内的美善来自天主，这是一件小事，除非你因此不高举自己超过那尚未拥有的人，他或许在领受后会超越你。因为当\\Person{扫禄}用石头砸\\Person{斯德望}时 \\BibleRef{宗 7:59}，他迫害的基督徒有多少！然而，当他被归化后，他超越了所有之前的人。因此，你要对天主说你在圣咏中所听到的：\\BibleQuote{上主，我仰望祢，我永远不会羞愧：以祢的义德}，不是我的义德，\\BibleQuote{解救我，拯救我}。\\BibleQuote{求祢侧耳听我}，这也是谦卑的宣认。那说\\BibleQuote{侧耳听我}的人，是在承认自己如同病人躺在站着的医生脚前。最后，请注意说话的是一位病人：\\BibleQuote{求祢侧耳听我，拯救我}。\par

5. \\BibleQuote{求祢作我保护的天主} \\BibleRef{咏 70:3}。不要让敌人的箭射到我：因为我不能保护自己。\\Quote{保护}还算小事；祂又加了：\\Quote{并作坚固的堡垒，使祢拯救我}。\\Quote{作坚固的堡垒}，求祢作我的堡垒……看，天主自己成了你逃往的地方，而你起初是害怕逃离祂的。\\Quote{作坚固的堡垒}，祂说，\\Quote{使祢拯救我}。除了在祢内，我不能安全；除非祢成为我的安息，我的病痛无法痊愈。求祢把我从地面举起；我要躺在祢身上，以便能上升到坚固的堡垒。有什么比这更坚固？当你逃往那地方作为避难所时，告诉我你还会惧怕什么敌人？谁会埋伏来攻击你？据说有一个人，当皇帝经过时，从山顶上喊叫：\\Quote{我说的不是你}；另一个人回头说：\\Quote{我也不是说你}。他藐视了一个带着闪亮兵器、强大军队的皇帝。从何处？从坚固之处。如果他在大地的高处是安全的，你在那创造天地的主内是多么安全？我若为自己选择了别的地方，就不能安全。人啊，你若找到了一个更好的堡垒，尽管选择吧。除了逃到祂那里，没有地方可躲避祂。你若想躲避祂的愤怒，就逃到祂的和解中。\\BibleQuote{因为祢是我的坚固和我的避难所}。\\Quote{我的坚固}是什么？藉着祢我坚固，因祢我坚定。\\BibleQuote{因为祢是我的坚固和我的避难所}：为使我因祢而坚固，无论我在自己内有何软弱，我都要逃到祢那里避难。因为基督的恩宠使你坚固，使你屹立不摇，抵抗敌人的一切诱惑。但那里仍有人性的脆弱，仍有最初的俘虏，仍有在肢体中与心神之法律交战、并企图在罪恶之法律中俘虏人的法律 \\BibleRef{罗 7:23}；仍有一个腐败的身体压迫着灵魂 \\BibleRef{智 9:15}。不论你因天主的恩宠多么坚固，只要你仍带着一个土制的器皿，其中盛有天主的宝藏，就必须从这同一泥器有所畏惧 \\BibleRef{格后 4:7}。因此，\\Quote{祢是我的坚固}，使我在今世能坚定面对一切诱惑。但如果诱惑众多，它们使我烦扰：\\Quote{祢是我的避难所}。因为我要承认自己的软弱，好使我像\\Quote{野兔}一样胆小，因为我像\\Quote{刺猬}一样充满荆棘。正如在另一篇圣咏中所说的：\\Quote{盘石是刺猬和野兔的避难所}；但盘石是基督 \\BibleRef{格前 10:4}。\par

6. \BibleQuote{天主，求祢从罪人的手中拯救我。}\BibleRef{咏 70:4} 总的说来，罪人，身处其中的那现在要从被掳中得释放的人在劳苦：他现在呼喊：\BibleQuote{我真不幸！谁能救我脱离这死亡的身体呢？天主的恩宠，借我们的主耶稣基督。}\BibleRef{罗 7:24} 内在有一个仇敌，即那在肢体中的法律；外在也有敌人：你向谁呼喊？向那已有人向他呼喊的，\Quote{求祢洗净我隐秘的罪，主，并从陌生的罪中宽免祢的仆人。}……但这些罪人分为两种：有些领受了法律，有些没有领受：所有外邦人都没有领受法律，所有犹太人和基督徒都领受了法律。因此，罪人是一般用语；或是违法者，如果他领受了法律；或只是不义者，没有法律，如果他未领受法律。关于这两种人，宗徒说：\BibleQuote{凡在法律之外犯了罪的，也必在法律之外丧亡；凡在法律之内犯了罪的，也必按法律受审判。}\BibleRef{罗 2:12} 但你身处两种人中间呻吟，要向天主说出你在圣咏中所听见的：\Quote{我的天主，求祢从罪人的手中拯救我。}是哪一个罪人？\Quote{从那违法者和不义之人的手中。}违法者确实也是不义之人；因为违法者没有不是不义之人的；但每一个违法者都是不义之人，并非每个不义之人都违法。因为，\Quote{没有法律，}宗徒说，\Quote{就没有违犯。}\BibleRef{罗 4:15} 这样那些没有领受法律的，可以称为不义之人，但不能称为违法者。两者都按其所应得受审判。但我，那要从被掳中借祢的恩宠得释放的，向你呼喊：\Quote{求祢从罪人的手中拯救我。}从他的手是什么意思？从他的权势，当他狂怒时，他不至于引我同意他；当他埋伏时，他不至于引诱我犯罪。\Quote{从罪人和不义之人的手中。}……\par

7. 最后，我说这话的理由是：\BibleQuote{因为祢是我的忍耐。}\BibleRef{咏 70:5} 如果祂是忍耐，正确的，祂也就是接着所说的：\Quote{主，我自幼的希望。}我的忍耐，因为我的希望；更确切地说，我的希望，因为我的忍耐。宗徒说：\BibleQuote{磨难生忍耐，忍耐生练达，练达生希望，希望不叫人蒙羞。}\BibleRef{罗 5:3} 我有理由寄希望于祢，主，我必永远不蒙羞。\Quote{主，我自幼的希望。}从你幼年起，天主就是你的希望吗？祂不也是从你的童年和婴儿期开始吗？确实，他说。因为看接着的话，免得你以为我说\Quote{我自幼的希望，}好像天主从未有益于我的婴儿或童年；听接着的话：\Quote{我从母胎中就受祢扶持。}再听：\BibleQuote{从我母腹中，祢就是我的保护者。}\BibleRef{咏 70:6} 那么，为什么说\Quote{从我幼年}，除非那是我开始寄希望于祢的时期？因为以前在祢内我没有希望，虽然祢是我的保护者，引领我安全直到我学会希望于祢之时。但从我幼年起，我开始在祢内希望，从祢武装我抵抗魔鬼的时候起，以致在祢军队的装备中，以祢的信德、爱德、望德以及其他恩赐武装起来，我与祢无形的敌人作战，并从宗徒那里听见：\BibleQuote{我们不是与血和肉搏斗，而是与率领者、掌权者，}等等。\BibleRef{弗 5:12} 在那里，是个年轻人与这些东西搏斗；但尽管他是年轻人，他也会跌倒，除非祂是他所呼喊者的希望：\Quote{主，我自幼的希望。}\Quote{在祢内，我的歌唱常存。}难道只是从我刚开始希望于祢直到如今吗？不，而是\Quote{常存。}\Quote{常存}是什么意思？不仅在信德之时，也在亲眼目睹之时。因为现在，\BibleQuote{我们住在肉身内，便远离主，因为我们凭信德往来，并非凭目睹。}\BibleRef{格后 5:6} 将有一个时候，我们将看见那未曾看见而相信的；但当所相信的被看见时，我们将喜乐；但当他们所不信的被看见时，不敬神的人将蒙羞。然后那现在作为希望的实体将来到。但是，\BibleQuote{可见的希望不是希望。但我们若希望那未看见的，必须忍耐等待。}\BibleRef{罗 8:24} 如今你呻吟，如今你奔往避难所，为能得救；如今你带着软弱恳求医生：当你得了完全的康复，当你被造得\BibleQuote{相似天主的众天使}\BibleRef{玛 22:30}时，你难道会忘记那使你获救的恩宠吗？绝不。\par

8. \BibleQuote{我好似成了众人中的怪物} \BibleRef{咏 70:7}。这里在希望之时，在呻吟之时，在卑微之时，在忧愁之时，在软弱之时，在带着锁链之声之时——那么这里如何？\BibleQuote{我好似成了众人中的怪物}。为什么，\BibleQuote{我好似成了怪物}？为什么那些以为我是怪物的人侮辱我？因为我相信我所未见之事。因为他们乐享所见之事，他们纵情于饮酒、荒淫、暗室淫乐、贪婪、财富、抢劫、世俗的尊荣、粉饰的泥墙，在这些事上他们欢喜；但我却行不同的路，轻视眼前之事，甚至惧怕世上的顺境，除了天主的应许之外无所信托。而他们说：\BibleQuote{我们吃喝吧，因为明天我们就要死了。} \BibleRef{格前 15:32}。你说什么？请重复：\BibleQuote{我们吃吧}，他说，\BibleQuote{喝吧}。来吧，你后来说了什么？\BibleQuote{因为明天我们就要死了}。你使我恐惧，却没有诱骗我。当然，正是你后来所说的那句话，使我害怕与你同流合污。\BibleQuote{因为明天我们就要死了}，你说；而那前面已有：\BibleQuote{我们吃喝吧}。因为当你说\BibleQuote{我们吃喝吧}时，你加了\BibleQuote{因为明天我们就要死了}。请听我另一边的说法：\Quote{是，我们禁食祈祷，\InnerQuote{因为明天我们就要死了。}} 我持守这条窄狭的路，\BibleQuote{我好似成了众人中的怪物，但你是大能的援助者。} 与我同在，哦主耶稣，向我说：不要在这窄路上气馁，我已先走过了它，我就是道路本身，\BibleRef{若 14:6} 我引导，在我内引导，引归我自己。因此，虽然\BibleQuote{我成了众人中的怪物}，我仍不害怕，因为\BibleQuote{你是大能的援助者。}\par

9. \BibleQuote{愿我的口充满赞美，好使我能以诗歌述说祢的荣耀，终日述说祢的伟岸} \BibleRef{咏 70:8}。什么是\Quote{终日}？无间断。在顺境中，因为祢安慰；在逆境中，因为祢纠正；在我存在之前，因为祢造了我；当我在存在时，因为祢赐予健康；当我犯罪时，因为祢宽恕；当我悔改时，因为祢帮助；当我坚持时，因为祢加冕。\par

10. 我自幼的希望，\BibleQuote{在我年老时，不要抛弃我} \BibleRef{咏 70:9}。这年老之时是什么？\BibleQuote{当我的力量衰竭时，不要离弃我。} 这里天主这样回答你：是的，让你的力量衰竭，好使我的力量在你内存留：为使你与宗徒一同说：\BibleQuote{当我软弱时，正是我有力量时。} \BibleRef{格后 12:10} 不要害怕，你将在那软弱中、在那年老时被抛弃。但为何？难道你的主在十字架上不也曾软弱吗？难道那些大能的人和肥壮的公牛在祂面前，像一个无力的人被俘虏和压迫，摇着头说：\BibleQuote{如果祂是天主子，就让祂从十字架上下来吧！}？祂因软弱而舍弃了吗？祂宁愿不从十字架上下来，免得祂似乎不是显示能力，而是屈服于辱骂祂的人。那悬挂在十字架上而不愿下来的祂，教导了你什么？那不是在辱骂中忍耐，并且你应在你的天主内刚强吗？或许也是在祂的人性中曾说过：\BibleQuote{我好似成了众人中的怪物，但你是大能的援助者。} 按照祂的软弱，而非按照祂的能力；按照祂将我们转化为祂自己的那一方面，而非按照祂亲自降临的那一方面。因为祂成了众人中的怪物。或许这也是祂的老年；因为由于它的老旧，称之为老年并无不妥，宗徒也说：\BibleQuote{我们的旧人已与祂同钉十字架。} \BibleRef{罗 6:6} 如果那里有我们的旧人，那么老年就在那里；因为旧，就是老年。然而，因为一句真话是：\BibleQuote{你的青春将如鹰更新；} 祂自己第三日复活，并应许在世界末日时的复活。元首已经先行，肢体将要跟随。你为什么害怕祂会离弃你，在你年老时抛弃你，当你的力量衰竭时？是的，那时在你内将是祂的力量，当你的力量衰竭时。\par

11. 我为何说这话？\BibleQuote{因为我的仇敌议论我，那些窥伺我灵魂的人已共同商议\BibleRef{咏 70:10}：说，天主舍弃了祂，追赶并捉拿祂，因为没有人拯救祂\BibleRef{咏 70:11}}这是指着基督说的。因为那具有伟大神能、与圣父同等、曾使死人复活的祂，忽然在敌人手中变得软弱，好像无力一般，被捉拿了。若非他们先在心中说\Quote{天主舍弃了祂}，祂何时会被捉拿呢？因此，十字架上才有那呼声：\BibleQuote{我的天主，我的天主，祢为什么舍弃了我？}那么，难道天主真的舍弃了基督，尽管\BibleQuote{天主在基督内使世界与自己和好} \BibleRef{格后 5:19}，尽管基督也是天主——按血统来说出自犹太人，\BibleQuote{祂是在万有之上，永远受赞美的天主} \BibleRef{罗 9:5}——难道天主舍弃了祂？绝对不可。但那是我们旧人的声音，因为我们的旧人与祂一同被钉在十字架上\BibleRef{罗 6:6}；祂取了那同一个旧人的身体，因为玛利亚出自亚当。因此，他们心里所想的，祂在十字架上说了出来：\BibleQuote{祢为什么舍弃了我？} \BibleRef{玛 27:46}为什么这些人认为我在他们的恶事中被撇弃了？\Quote{因为他们若认识，决不至于把荣耀的主钉在十字架上。\BibleRef{格前 2:8} 追赶并捉拿祂。}然而，弟兄们，让我们更亲切地将此理解为基督肢体的声音，并在这些话中认出我们自己的声音：因为祂是替我们说话，并非按祂自己的权能与威严，而是按祂为我们所成为的，并非按祂所是的、创造我们的那一位。\par

12. \BibleQuote{上主，我的天主，求祢不要远离我}\BibleRef{咏 70:12}诚然如此，主绝不远离。因为\BibleQuote{上主亲近那些心灵破碎的人。}\Quote{我的天主，求祢垂顾我的援助。}\BibleQuote{愿那些谋害我性命的人，蒙受羞辱而失败}\BibleRef{咏 70:13}祂愿求什么？\Quote{愿他们蒙受羞辱而失败。}祂为何愿求？\Quote{那些谋害我性命的人？}什么是\Quote{谋害我性命？}犹如挑起某种争闹。因为被挑战参与争吵的人被称为被谋害。若果真如此，让我们提防那些谋害我们灵魂的人。什么是\Quote{谋害我们的灵魂？}首先是挑唆我们违抗天主，好使我们在恶事中不悦于天主。因为何时你正直，使以色列的天主善待你——对心地正直的人善？何时你正直？你愿意听吗？当你所行的善事中，天主令你喜悦；而你所受的恶事中，天主不令你不悦。请看，弟兄们，我所说的，你们要提防那些谋害你们灵魂的人。因为所有与你们交往、要使你们在忧愁与磨难中疲惫的人，其目的就是使你们在所遭遇的事上不悦于天主，并从你们口中发出：\Quote{这是怎么回事？我做了什么？}现在你当真没做坏事，你是义人，祂不义？我是个罪人，你说；我承认，我不自称为义。但罪人啊，你难道作了与那享福者一样多的恶事？像盖乌塞乌斯那样？我知道他的恶行，知道他的不义，我虽是个罪人，却远不及他；然而我看见他享有一切福乐，而我却遭受如此大的恶事。那么我不说：天主啊，\Quote{我对祢做了什么}，因为我根本没做恶事；而是因为我没做那么多以致该受这些。再者，你是义人，祂不义？醒醒吧，可怜的人，你的灵魂已被谋害！他说：我没有自称为义。那你说什么？我是罪人，但我没犯那么大罪，以致该受这些。那么你不向天主说：我是义人，祢是不义；而是说：我不义，祢更不义。看，你的灵魂已被谋害，看，现在你的灵魂在作战。与谁？你的灵魂对抗天主，对抗那创造它的那一位。你既存在却向祂呼喊，真是忘恩负义。那么，回到你疾病的告明中，祈求医生的医治之手。不要以为那些暂时昌盛的人有福。你受管教，他们被暂缓；或许为你——受管教而改过——保留了一份基业……最后，请看下文：\Quote{愿那些想我恶事的人，披上羞愧与耻辱。}\Quote{羞愧与耻辱}，羞愧因良心不安，耻辱因谦逊。愿这事临到他们，他们就会变好……\par

13. \BibleQuote{但我始终仰望祢，并要加增祢一切的赞美}\BibleRef{咏 70:14}这是什么意思？\Quote{我要加增祢一切的赞美}，应当引起我们的注意。你要使天主的赞美更完全吗？还有可以加添的吗？如果那已经是全部的赞美，你还能加添什么？天主在祂一切善工中受赞美，在祂的每个受造物中，在整个万有的建立中，在世代的管理与安排中，在季节的秩序中，在穹苍的高处，在大地的丰饶中，在海洋的环抱中，在受造物各处显现的卓越中，在人类子孙本身中，在颁赐法律中，在将祂的子民从埃及人的奴役中解救出来，以及祂其余一切奇妙作为中——但尚未因祂使血肉之躯复活获得永生而受赞美。那么，愿这赞美因我们的主耶稣基督的复活而被加添：为使我们在此认出祂的声音超越过去一切的赞美；这样，我们也正确理解这首诗……\par

14. \BibleQuote{我的口要传述祢的公义}\BibleRef{咏 70:15}：不是我的公义。从此我要增添对祢的一切赞美：因为我若为义，是因祢在我内的义，而非我自己的义，因为祢使不虔敬的人成义。\BibleRef{罗 4:5}\Quote{终日述说祢的救恩。}什么是\Quote{祢的救恩}？愿无人自以为是救自己的，\Quote{救恩属于上主。}没有人能靠自己救自己，\Quote{人的救恩是虚幻的。}\Quote{终日述说祢的救恩：}时时如此。逆境临到时，宣讲主的救恩；顺境临到时，宣讲主的救恩。不要在顺境中宣讲，在逆境中闭口不言；否则就没有经上所说的\Quote{终日}。因为终日包括它自己的黑夜。当我们说，例如三十天过去了，难道不也提到夜晚吗？我们不是用\Quote{天}这个词语也包括夜晚吗？在\BibleBook{}{创世纪}中是怎么说的？\Quote{过了晚上，过了早晨，这是一天。}\BibleRef{创 1:5}因此一整天就是白天连同它的黑夜：因为黑夜为白天服务，而非白天为黑夜服务。你在有死肉体中所做的一切，都应为义服务；凡你因天主的诫命而做的，不要为肉体的利益而行，免得黑夜为白天服务。因此要终日述说天主的赞美，也就是说，在顺境和逆境中：在顺境中，犹如在白昼；在逆境中，犹如在黑夜：但仍要终日述说天主的赞美，这样你才不会徒然地歌唱：\Quote{我要时时赞美天主，对他的赞美常在我口中。}……\par

15. 因此，他说：\Quote{因为我不知买卖。}这些买卖是什么？让商人们听见并改变他们的生活；如果他们曾是那样，就不要再那样；让他们不认识他们曾是什么，让他们遗忘；最后，让他们不赞同，不称赞；让他们不赞成，谴责，悔改，如果买卖是一种罪。因为为此，商人啊，由于某种贪求获得，每当你遭受损失，你就会亵渎；你里面就不会有经上所说的\Quote{终日赞美祢。}但每当你为了出售货物的价格，不仅说谎，而且甚至发假誓，你的口中怎能终日有天主的赞美？如果你是一个基督徒，甚至因你的口，天主的名被亵渎，以致人们说：看，基督徒是什么样的人！因此，如果这人为此终日赞美天主，因为他不知买卖，那么让基督徒改正自己，不要做生意。但一个商人会对我说：看，我确实从远方把货物带到这些地方，这些地方没有我带来的东西，借此我可以谋生；我只求我劳力的报酬，即高价卖出我低价买进的；因为我如何能生活，经上不是写着\BibleQuote{工人配得他的工资}吗？\BibleRef{路 10:7}但他所论的是说谎、发假誓。这是我的过错，不是买卖的过错：因为如果我不愿意，我本可避免这个过错。那么，我，商人，不把自己的过错推给买卖：但如果我说谎，是我说谎，不是买卖。因为我可以这样说：我花这么多钱买进，我要卖这么多钱；你愿意，就买。买主听到这个真理，不会生气，所有的人也不会因此而少来找我：因为他们爱真理胜过爱金钱。因此，他说，你要告诫我：不说谎，不发假誓；而不是放弃我赖以谋生的生意。因为当你让我放弃这个，你让我去做什么？也许去学一门手艺？我将成为一个鞋匠，为人做鞋。难道他们不说谎？他们不发假誓？当他们与人约定做鞋，从另一个人那里收了钱，就放弃正在做的，又接另一个人的活，并欺骗那个他们承诺尽快完成的人？他们不是常说：今天我就做，今天就能完成？其次，在缝制中，他们不也犯了许多欺诈？这些是他们的行为，这些是他们的言语：但他们是恶人，而非他们所从事的职业。因此，所有邪恶的工匠，不敬畏天主，或为获利，或为怕损失或匮乏，都说谎、发假誓；在他们里面没有持续的天主的赞美。那么，你如何让我退出买卖？你愿意我成为一个农夫，对打雷的天主抱怨，以致害怕冰雹而去求问巫师，好知道如何保护自己免受天气之害；以致我盼望穷人遭饥荒，好让我能高价卖出我囤积的粮食？你把我引向这个吗？但你说，好农夫不做这样的事。好商人也不做那些事。然而，为什么，甚至生儿子也是一件恶事，因为当他们头痛时，邪恶而不信的母亲就求助于不虔敬的咒符和咒语？这些是人的罪，而非事物的罪。一个商人可能会对我说——那么，主教啊，请看你如何在\BibleBook{}{圣咏集}中读到你所理解的买卖：免得你理解错了，却禁止我做买卖。那么，请告诉我该如何生活；如果好，那就好；但我知道一点：如果我成为恶人，那不是买卖使我如此，而是我的不义。当真理被说出时，无人可以反驳它。\par

16. 那么让我们探究他所说的\OriginalTerm{交易}{tradings}，确实那不认识交易的人，终日赞美天主。在希腊语中，\InnerQuote{交易}一词源于行动，在拉丁语中源于不活动的缺失：但无论是源于行动还是不活动的缺失，让我们考察它是什么。因为那些活跃的交易人，仿佛依靠自己的行为，他们赞美自己的行为，却得不到天主的恩宠。因此，交易人反对这篇圣咏所颂扬的恩宠。因为圣咏颂扬那恩宠，为使人不得夸耀自己的行为。因为在某处经上说：\Quote{医生必不能使人复活}，难道人应当放弃医药吗？但这是什么？在此名称下所理解的是骄傲的人，向人许以救恩，而\Quote{救恩属于上主}……因此，主有理将那些人逐出圣殿，向他们说：\Quote{经上记载：我的殿宇将称为祈祷之所，但你们却使它成了交易之所；}那即是，夸耀你们的行为，追求活动，也不听圣经反对你们的不安和交易的话：\Quote{你们要安静，要知道我是天主。}……\par

17. 但在一些抄本中有：\Quote{因为我未曾认识学识。} 有些书用\Quote{交易}，另一些用\Quote{学识}：如何一致是件难事；然而译者的歧异或许显示意义，并不导致错误。那么让我们也探究如何理解学识，以免我们冒犯文法学家，如同刚才我们冒犯了交易人一样：因为文法学家也可以光荣地履行他的职务，既不发假誓也不说谎。那么让我们考察那未曾被认识、而其口中终日有赞美天主的学识。有一种犹太人的学识：因为让我们把这归于他们；在那里我们将找到所说的：正如我们查问交易人，关于行为和行为的问题，我们发现那被称为可憎的交易，正是宗徒所指责的，说：\Quote{因为他们不认识天主的正义，而想建立自己的正义，就不服从天主的正义。} \BibleRef{罗 10:3}……正如我们发现了先前对交易人的指控，即那些夸耀行为，因不容不活动的事务而高抬自己的人，是不安的人而非好工人；因为好工人是天主在其内工作的人；同样我们也发现了一种犹太人的学识……\Person{梅瑟}写了五卷书：但在围绕水池的五条门廊中，\BibleRef{若 5:2}有病的人躺着，却不能痊愈。看字句如何留下，定罪有罪者，而不拯救不义者。因为在那些五条门廊中，五卷书的预象，病人被放弃而不是得到痊愈。那么在那个地方什么使病人痊愈？水的扰动。当那水池被扰动时，一个病人下去，就痊愈了一个，一个为了合一：无论什么人下去到那同一扰动中，都没有痊愈。那么从地极呼喊的奥体之合一如何被颂扬？另一个人除非水池再次被扰动，就不会痊愈。水池的扰动于是预表了主耶稣基督来时犹太人的骚动。因为天使来到时，池中的水被感知到扰动。那么被五条门廊环绕的水，是被法律环绕的犹太民族。在门廊里病人躺着，只有在水中被搅动时才痊愈。主来了，水被扰动；他被钉十字架，愿他下来，使病人痊愈。愿他下来是什么意思？愿他谦卑自己。因此，无论你们是谁，爱慕字句而无恩宠，你们将留在门廊里，你们将生病，躺着，不痊愈……因为同样的预象，\Person{厄里叟}首先派他的仆人拿着棍杖去复活死去的孩子。他的女东家寡妇的儿子死了；有人报告给他，他把棍杖给了仆人：他去，他说，放在死孩子身上。先知不知道他在做什么？仆人先去了，他把棍杖放在死者身上，死者没有起来。\Quote{因为如果颁赐的法律能赋与人生命，那么正义就确实出自法律了。} \BibleRef{迦 3:21} 由仆人送去的法律没有叫人活；然而他派仆人拿着棍杖，自己随后来了，并使孩子活了。\BibleRef{列下 4:20} 因为那婴儿没有起来，\Person{厄里叟}自己来了，现在带着主的预象，主曾先派他的仆人带着棍杖，好像带着法律：他来到躺着的死孩子面前，把他的肢体放在孩子身上。一个是婴儿，另一个是成人：他缩紧和缩短了他全长的尺寸，为能适应那死孩子。死者于是活了过来，当活者使自己适应死者：主人做了棍杖没有做的事；恩宠做了字句没有做的事。那么那些留在棍杖里的人，以字句自夸，因此没有得生命。但我要因你的恩宠夸耀……在那同一恩宠中我夸耀：\Quote{我未曾认识学识}，就是说，那依靠字句而远离恩宠的人，我全心拒绝了。\par

18. 有理由接续说：\Quote{我要进入主的能力中：}不是我的，而是主的。因为他们以自己的文字能力为荣，因此，他们不知道恩宠与文字结合……但因\Quote{文字叫人死，但圣神却叫人活：}\BibleRef{格后 3:6}\Quote{我不曾认识\OriginalTerm{学问}{literature}，我要进入主的能力中。}因此，接下来的这句经文巩固和完善了这一意义，以便扎根在人心中，不让任何其他的解释从任何地方潜入。\Quote{主啊，我要记念惟独你的公义}\BibleRef{咏 70:16}。啊！\Quote{惟独。}为什么他加上\Quote{惟独，}我问你？只需说\Quote{我要记念你的公义}就足够了。\Quote{惟独}，他说，全然如此：我不思想属于自己的。\Quote{你有什么不是你领受的呢？既然你也领受了，为什么还夸耀，好像没有领受呢？}\BibleRef{格前 4:7}惟独你的公义拯救我，我自己的只有罪恶。那么，我不可夸耀自己的力量，不可停留在文字中；我要弃绝\Quote{学问}，即那些以文字自夸的人，并凭自己的力量乖僻地、像狂人一样依靠的人：我要弃绝这样的人，我要进入主的能力中，好使我软弱时，就变得刚强；为使你在里面刚强，因为\Quote{我要记念惟独你的公义。}\par

19. \Quote{天主，你从我的青年时代教导了我}\BibleRef{咏 70:17}。你教导了我什么？就是我只应记念惟独你的公义。因为回顾我过去的生活，我看见什么是我所欠的，什么是我没有欠的替代所领受的。所欠的是惩罚，所付清的是恩宠；所欠的是地狱，所赐予的是永生。\Quote{天主，你从我的青年时代教导了我。}从我信德的最初，你以它更新了我，你教导我，在我里面没有任何先行的东西，能使我可以说你所赐予的是我所欠的。因为谁转向天主不是脱离不义呢？谁被救赎不是脱离囚禁呢？但谁能说他的囚禁是不义的，当他抛弃了他的统帅而投靠了叛徒？天主是我们的统帅，魔鬼是叛徒：统帅颁布了命令，叛徒提出了诡计：在命令和欺骗之间，你的耳朵在哪里？难道魔鬼比天主更好？难道背叛者比造你的更好？你相信了魔鬼所应许的，经历了天主所威胁的。现在，从囚禁中被解救出来，仍然处于希望中，尚未在实质上，凭信德行走，非凭眼见，他说：\Quote{天主，} \Quote{你从我的青年时代教导了我。}从我转向你那时起——被你更新，我原是你所造；再被创造，我原被造；再被塑造，我原被塑造：从我归依那时起，我学会了：在我之前没有任何功绩，而是你的恩宠白白地临到我，\OriginalTerm{白白地}{gratis}，为使我记念惟独你的公义。\par

20. 青年时期之后是什么？因为他说：\Quote{祢教导了我，} \Quote{从我的青年时期起。} 青年时期之后又是什么？因为在你的第一次皈依中，你学到了：在皈依之前，你并不是义人，而是不义在先，为的是驱逐不义后，爱德可以接续；并且你被更新为新人，但只是在望德中，尚非在实质上；你学到了你的任何善行都不是先行的，你是因天主的恩宠而皈依了天主。如今，自从你皈依之后，难道你就拥有自己的东西，应当依靠自己的力量吗？正如人们习惯说的：现在离开我，你需要为我指明道路；够了，我将走在这条路上。而那为你指明道路的人说：\Quote{难道你不愿我引导你到那地方吗？} 但如果你自高自大，就说：\Quote{别管我，够了，我将走在这条路上。} 他们撇下了你，而因你的软弱，你会再次迷路。当初那引领你走上道路的那位，若继续引导你，对你是好的。但若他不再引领你，你又会迷失。那么你就对他说：\Quote{主啊，求祢引领我走在祢的道路上，我将在祢的真理中行走。} 但你踏上道路，却是青年时期，正是信德的更新和开始。因为以前你是流浪者，行走在自己的道路上；穿越林地和崎岖之地，四肢受伤，寻找一个家，即你灵魂的安息之所，在那里你可以说：一切都好；并在安全中如此说，脱离一切不安、一切试探、总之脱离一切囚禁；但你却没有找到。我该说什么？难道有一个人来给你指路吗？是道路本身来到了你这里，而你被安置在其中，并非由于你先前的功德，因为你显然是在迷路。那么，自从你踏入这条道路以来，你现在能自己引导自己吗？那教导你道路的，如今会离弃你吗？不，他说：\Quote{祢从我的青年时期就教导了我；直到现在，我还要宣扬祢奇妙的作为。} 因为祢仍在行奇妙的事：祢引导我，那已将我放在道路上的那一位；这些就是祢奇妙的作为。你认为天主的奇妙作为是什么？在天主的奇妙作为中，还有什么比使死者复活更奇妙呢？但你说：难道我是死人吗？除非你曾经是死者，就不会对你说：\Quote{你这睡着的人，醒来吧，从死者中起来，基督必要光照你。} \BibleRef{弗 5:14} 所有不信者、所有不义之人都是死的；他们在身体上活着，但在心中已经熄灭。但那使死人按身体复活的人，使他回来看见这光、呼吸这空气；但那使他复活的人本身并非他的光和空气；他开始看见，如同他以前看见一样。灵魂并非如此复活。因为灵魂是由天主复活的；尽管身体也是由天主复活的；但当天主使身体复活时，是使它回到世界；当天主使灵魂复活时，是使它回到自己那里。若这世界的空气被抽走，身体就死亡；若天主被抽走，灵魂就死亡。因此当天主使灵魂复活时，除非那位使她复活者与她同在，否则她即使复活也不能生活。因为他并不先使她复活，然后留下她让她自己生活：如同拉匝禄，在死后四天被复活，是因主的肉身临在而复活……主离开那城或那地方，拉匝禄就不再活了吗？灵魂的复活并非如此：天主使她复活，若天主离去，她就死亡。因为我将大胆地说，弟兄们，但这是真理。有两种生命：身体的生命和灵魂的生命；如同身体的生命是灵魂，同样灵魂的生命是天主；同样，若灵魂离去，身体死亡；若天主离去，灵魂死亡。这就是他的恩宠，即他复活我们并与我们同在。因为他从我们过去的死亡中复活了我们，并在某种意义上更新了我们的生命，我们对他说：\Quote{天主啊，祢从我的青年时期就教导了我。} 但因为他并不从他所复活的人那里撤退，唯恐当他撤退时他们死亡，我们对他说：\Quote{直到现在，我还要宣扬祢奇妙的作为。} 因为当你与我同在时，我活着；我的灵魂的生命就是祢，若她自我放任，就会死亡。因此当我的生命临在，即我的天主，\Quote{直到现在，} 接下来是什么？\par

21. \Quote{直到年老和暮年} \BibleRef{咏 70:18}。这是老年的两个术语，被希腊人加以区分。因为青年期之后的年长阶段在希腊人中有另一个名称，而同年长阶段之后来临的最后年龄又有另一个名称；因为\OriginalTerm{年长}{\GreekText{πρεσβύτης}}意指年长，而\OriginalTerm{}{\GreekText{γ}}\par

\SourceRecord{Translated by J.E. Tweed.；Nicene and Post-Nicene Fathers, First Series；Vol. 8.；Edited by Philip Schaff.；Buffalo, NY: Christian Literature Publishing Co.,；1888.；<http://www.newadvent.org/fathers/1801071.htm>.}

% structure: p0001 source=h1 target=section reason=page-title

\section{论圣咏第七十二篇}

1. \Quote{为撒罗满} 确实，这篇圣咏的标题预先注明了。但其中所说的事，不能适用于那按肉体为以色列王的撒罗满，按圣经关于他所说的那些事而言；但却最切合地适用于主基督。由此可见，\Quote{撒罗满}这个词是比喻性的用法，在他身上应理解为基督。因为\Quote{撒罗满}解释为\Quote{和平缔造者}；因此，这个词最真实且卓越地适用于祂——藉着祂这位中保，我们这些仇敌得以与天主和好，获得了罪的赦免。因为\BibleQuote{当我们还是仇敌的时候，藉着祂的圣子的死，我们与天主和好了} \BibleRef{罗 5:10}。祂本身就是那位和平缔造者……既然我们找到了真正的撒罗满，即真正的和平缔造者；接下来让我们观察这篇圣咏关于祂教导了什么。\par

2. \Quote{天主，求祢将祢的审判赐给君王，将祢的正义赐给君王的儿子} \BibleRef{咏 71:1}。主自己在福音中说：\BibleQuote{圣父不审判任何人，但将一切审判交给了圣子} \BibleRef{若 5:22}。这就是\Quote{天主，求祢将祢的审判赐给君王}。那位君王也是君王的儿子，因为圣父当然也是君王。圣经上正是这样记载的：\BibleQuote{君王为自己的儿子摆设婚宴} \BibleRef{玛 22:2}。但按圣经的惯用方式，同一件事被重复了。因为他在\Quote{祢的审判}中所说的，又在\Quote{祢的正义}中以不同方式表达了；他在\Quote{君王}中所说的，又在\Quote{君王的儿子}中以不同方式表达了……但这些重复大大地强调了神圣的言语，无论是同一词汇还是以其他词汇重复同一意义；它们主要出现在圣咏中，也出现在那些意在唤醒心灵情感的论述中。\par

3. 接下来是：\Quote{为秉公审判祢的百姓，秉义判断祢的穷苦人} \BibleRef{咏 71:2}。圣父君王将祂的审判和正义赐给祂的圣子君王的目的，在祂说\Quote{为秉公审判祢的百姓}时已充分表明，即为了审判祢的百姓。类似的表达方式见于撒罗满：\Quote{撒罗满的箴言，达味的儿子，为使人认识智慧和训诲} \BibleRef{箴 1:1}，即撒罗满的箴言，为了认识智慧和训诲。同样，\Quote{求祢将祢的审判赐给祂，为审判祢的百姓}：即求祢将祢的审判赐给祂，为了审判祢的百姓。但祂先前在\Quote{祢的百姓}中所说的，随后在\Quote{祢的穷苦人}中重复了；先前在\Quote{秉公}中所说的，随后在\Quote{秉义}中重复了——按照那种重复的方式。这确实表明，天主的百姓应当是穷苦的，即不骄傲，而是谦卑。因为\Quote{神贫的人是有福的，因为天国是他们的} \BibleRef{玛 5:3}。连有福的约伯在失去那些巨大的地上财富之前，也安处于这种贫穷中。我之所以认为应提及此事，是因为有些人更容易将一切财物分给穷人，却不愿自己成为天主的穷人。因为他们被骄傲所膨胀，认为自己生活良好应归功于自己，而非天主的恩宠；因此，即使他们似乎行了极大的善事，现在实际上也生活得不好……\par

4. 但是，既然他改变了词语的顺序（尽管他一开始说\Quote{天主，求祢将祢的审判赐给君王，将祢的正义赐给君王的儿子}，先提审判后提正义），现在又先提正义后提审判，说\Quote{为秉公审判祢的百姓，秉义判断祢的穷苦人}，他就更清楚地表明，他把审判称为正义，证明语序并不造成差异，因为两者指同一事物。因为通常我们说\Quote{错误审判}指不义的审判；但\Quote{不义或不公的正义}我们通常不说。因为若是不义或不公，就不再称为正义了。此外，通过先用审判再用正义重复，或先用正义再用审判重复，他清楚地表明，他特别命名的那种审判通常代替正义使用，即不能理解为给出恶审判的那种审判。因为在祂说\Quote{不要按外貌判断，但要以公义的审判判断} \BibleRef{若 7:24}的地方，祂表明可能存在错误审判；当祂说\Quote{要以公义的审判判断}时，祂禁止一种，命令另一种。但当祂不加任何修饰地说\Quote{审判}时，祂即刻要人理解为公义的审判；如祂所说的：\Quote{你们忽略了法律上更重要的事：怜悯和审判} \BibleRef{玛 23:23}。另如耶肋米亚所说：\Quote{不以正义积累财富} \BibleRef{耶 17:11}。祂不是说\Quote{以错误或不公的审判积累财富}，或\Quote{不以正义或公义的审判}，而是说\Quote{不以审判}；祂只把正义和公义的称为审判。\par

5. \BibleQuote{愿群山带给人民和平，丘陵带给正义}\BibleRef{咏 71:3}。群山是较大的，丘陵是较小的。这无疑就是另一篇圣咏所说的\BibleQuote{小的与大的}。因为当以色列离开埃及，即天主的子民脱离这世界的奴役时，那些群山曾像公羊一样欢跃，丘陵像羊羔一样雀跃。那么，在教会中因卓越的圣德而显赫的人，就是群山，他们适宜教导其他人，\BibleRef{弟后 2:2} 藉着说话使人忠实地受教，藉着生活使人效法而获益；而丘陵则是那些藉着自身的服从跟随前者卓越的人。那么，为何说\BibleQuote{群山带来和平，丘陵带来正义}呢？也许即使这样说也没有什么不同：\BibleQuote{愿群山带给人民正义，丘陵带来和平}？因为二者都需要正义，二者都需要和平：或许正义本身在另一个名字下被称为和平。因为这才是真正的和平，不是不义之人之间所缔造的那种。或者，必须以一种不容忽视的区分来理解他所说的\BibleQuote{群山带来和平，丘陵带来正义}？因为在教会中卓越的人应该以警醒的关怀谋求和平，免得为了自身的优越而骄傲行事，制造分裂，破坏合一的纽带。但让丘陵藉着效法和服从如此跟随他们，以致他们更喜爱基督过于他们：免得被邪恶的群山（因为他们似乎卓越）的空虚权威引入歧途，从而脱离基督的合一……\par

6. 因此，以下的理解也最为贴切：\BibleQuote{愿群山带给人民和平}，即我们理解和平在于那使我们与天主和好的和好：因为群山为祂的子民领受这和平……\BibleQuote{因此，愿群山为百姓领受和平，丘陵领受正义}，这样，两者一致，就会发生经上所记载的：\BibleQuote{正义与和平彼此相拥}。而其他抄本所说的\BibleQuote{愿群山为百姓领受和平，丘陵}我认为应理解为各种宣讲福音和平的方式，无论是先行的还是后来的。但在这些抄本中，接着说到：\BibleQuote{在正义中祂将审判百姓中的穷人}。但那些抄本是更受认可的，它们有我们上面所解释的：\BibleQuote{愿群山带给人民和平，丘陵带给正义}。但有些抄本有\Quote{祢的百姓}，有些则没有\Quote{祢的}，只有\Quote{百姓}。\par

7. \BibleQuote{祂要审判百姓中的穷人，拯救穷人的子孙}\BibleRef{咏 71:4}。在我看来，穷人和穷人的子孙似乎是同一回事，正如同一个熙雍和熙雍女子。但如果要区分理解，我们把穷人理解为群山，而穷人的子孙理解为丘陵：例如，先知和宗徒是穷人，而他们的子孙，即那些在他们的权威下受益的人，是穷人的子孙。而前面所说的\Quote{审判}，以及后面的\Quote{拯救}，似乎是一种解释，说明祂如何审判。因为祂审判是为了拯救，即从那些将被毁灭和定罪的人中分别出来那些祂赐予\BibleQuote{那在末次时期将要显现的救恩}的人。\BibleRef{伯前 1:5} 因为藉着这样的人，向祂说：\BibleQuote{不要将我的灵魂与不敬虔的人一同消灭}；以及\BibleQuote{天主，求祢审判我，为我分辨我的案件脱离不圣洁的民族}。我们也要注意，祂不是说\Quote{审判穷百姓}，而是\Quote{审判百姓中的穷人}。因为上面当祂说\Quote{以正义审判祢的百姓，以公义审判祢的穷人}时，祂称天主的子民为祂的穷人，即只有善人和属于右边的人。但是因为在这个世界上，右边和左边的人一起牧养，他们像绵羊和山羊一样在最后要被分开，\BibleRef{玛 25:32} 整个混杂在一起的群体，祂用\Quote{百姓}这一名称来称呼。而且因为在这里祂以善义使用审判，即为了拯救的目的：因此祂说\Quote{祂要审判百姓中的穷人}，即在百姓中为救恩分别那些穷人。\Quote{祂要压制诬告者。}在这里没有比魔鬼更适合被认定为诬告者了。诬告是他的专长。\BibleQuote{约伯敬畏天主岂是无缘无故呢？}\BibleRef{约 1:9} 但是主耶稣藉着祂的恩宠帮助祂自己的人来压制他，使他们无缘无故地敬畏天主，即以主为乐。祂也这样压制了他：因为魔鬼，即这世界的王，在祂里面一无所获，\BibleRef{若 14:30} 就藉着犹太人的诬告杀害了祂，这些诬告者被诬告者用作他的器皿，在悖逆之子身上运行。\BibleRef{弗 2:2}……\par

8. \Quote{祂将存续到太阳}，或\Quote{祂将与太阳一同存续}\BibleRef{咏 71:5}。因为我们的某些作者认为，希腊文中的\OriginalTerm{一同存续}{\GreekText{συμπαραμενεῖ}}应当更准确地翻译。但如果用拉丁文一个词来表达，就必须用\Emph{compermanebit}来表达；然而，因为拉丁文无法表达这个词，为了至少把意思译出来，便译成了\Quote{祂将与太阳一同存续}。因为\Quote{祂与太阳同存续}不是别的，就是\Quote{祂将与太阳一同存续}。但祂与太阳一同存续，这有何重大意义呢？万物藉祂而造，没有祂什么也不能造成，\BibleRef{若 1:3}但这预言预先发出，为的是那些以为基督信仰的名称在这世界上只会存续到特定时期、之后便不复存在的人。因此，\Quote{祂将存续}，\Quote{与太阳一同}，只要太阳升起落下，即只要这些时间轮转，天主的教会——即基督在世上的身体——就不会缺乏。但祂加上\Quote{在月亮之前，世世代代}；祂本来可以用\Quote{在太阳之前}来表达，即既与太阳一同又在太阳之前；这将被理解为既与时间一同又在时间之前。那在时间之前的便是永恒的；那不被时间改变的才真正被认为是永恒的，正如\Quote{在起初已有圣言}，\BibleRef{若 1:1}。但祂选择用月亮来暗示可朽事物的盈亏变化。最后，当祂说了\Quote{在月亮之前}，似乎想解释为何插入月亮，祂说\Quote{世代}，即\Quote{的世代}。仿佛祂在说：在月亮之前，即在那世世代代之前，这些世代在可朽之物的离去与接续中消逝，如同月亮的盈亏。因此，关于祂在月亮之前存续，还有什么比祂以不朽性超越所有可朽之物更恰当的理解呢？以下理解也不算离题：如今祂已挫败了诬告者，坐在圣父的右边，这就是与太阳一同存续。因为永恒荣耀的光辉被理解为圣子，\BibleRef{希 1:3}仿佛太阳是圣父，而祂的光耀是圣子。但这些都是就创造者不可见的实体而言的，并非指那可见的受造界，其中虽有天体，而发光体中太阳最为卓越，这个比喻正是从其中取来的；正如比喻也从地上的事物取来，例如石头、狮子、羔羊、有两个儿子的父亲等。因此，祂挫败了诬告者，与太阳一同存续，因为祂藉复活战胜了魔鬼，坐在圣父的右边，\BibleRef{谷 16:19}在那里祂不再死亡，死亡再也不能管辖祂，\BibleRef{罗 6:9}。这也在月亮之前，如同从死者中首生者走在教会之前，教会正在可朽之物的离去与接续中前行。这就是\Quote{世世代代}。或者也许因为\Quote{世代}是指我们藉以受生为可朽者的世代；而\Quote{世世代代}是指我们藉以重生为不朽者的世代。正是如此，祂走在教会之前，为要在月亮之前存续，作为从死者中首生者。诚然，希腊文\OriginalTerm{世代世代}{\GreekText{γενεας} \GreekText{γενεῶν}}，有些人解释为不是\Quote{世代}，而是\Quote{一代的世代}，因为\OriginalTerm{\GreekText{γενεας}}{\GreekText{γενεας}}在希腊文中格位模糊，不知是单数属格\OriginalTerm{\GreekText{της} \GreekText{γενε}}{\GreekText{της} \GreekText{γενε}}…\par

\SourceRecord{Translated by J.E. Tweed.；Nicene and Post-Nicene Fathers, First Series；Vol. 8.；Edited by Philip Schaff.；Buffalo, NY: Christian Literature Publishing Co.,；1888.；<http://www.newadvent.org/fathers/1801072.htm>.}

% structure: p0001 source=h1 target=section reason=page-title

\section{圣咏第七十三篇释义}

1. 这篇圣咏有一个题名，即一个标题：\Quote{达味之子耶西的咏歌已终止。阿撒夫的圣咏。}我们在许多圣咏的题名上看到\Person{达味}的名字，但从未加上\Quote{耶西之子，}唯独在此处是这样。我们应当相信，这不是无意义或随意而为的。因为天主处处向我们暗示，并邀请我们以虔敬的爱研究以理解之。那么，\Quote{达味之子耶西的咏歌已终止}是什么意思？咏歌是伴随着歌唱的天主赞词：咏歌是包含赞美天主的歌曲。如果赞美，但不是赞美天主，就不是咏歌：如果赞美，而且是赞美天主，但未歌唱，就不是咏歌。所以，如果它是一首咏歌，必须兼具这三件事：赞美，赞美天主，以及歌唱。那么，\Quote{咏歌已终止}是什么意思？赞美天主的歌唱已终止。他似乎在讲述一件令人痛苦的事，可以说是可悲的事。因为颂赞者不仅赞美，而且欢喜地赞美；颂赞者不仅歌唱，而且爱他所歌唱的对象。在赞美中，有宣认的言说；在歌唱中，有爱慕的情感。他说，因此\Quote{达味的咏歌已终止}，并加上了\Quote{耶西之子}。因为\Person{达味}是以色列王，\Person{耶西}之子\BibleRef{撒上 16:19}，是在\WorkTitle{旧约}的某个时期，那时\WorkTitle{新约}隐藏在其中，如同果实藏在根中。因为如果你在根中寻找果实，你不会找到；然而，难道你在枝条中找不到任何果实，除了那从根中发出的吗？…同样，\Person{基督}自己要按肉身诞生，也隐藏在根中，即圣祖们的种子中，并在某个时期必须被揭示，如同果实的出现，正如所记载的：\BibleQuote{从\Person{耶西}的根上将发出嫩枝}\BibleRef{依 11:1}。同样，\Person{基督}内的\WorkTitle{新约}本身，在那些早期时代是隐藏的，只有先知们和极少数虔敬的人才知道，不是通过现实事物的显明，而是通过未来事物的启示。兄弟们，这是什么意思？（只提一件事）\Person{亚巴郎}派遣他忠信的仆人为他的独生子娶妻，使他起誓，并在誓言中对他说：\Quote{请把你的手放在我大腿下，并起誓}。\Person{亚巴郎}的大腿下有什么，使他把手放在那里起誓？那里有什么，除了那时对他所应许的：\BibleQuote{地上万民都要因你的后裔蒙受祝福}\BibleRef{创 22:18}？在大腿的名义下，肉身被象征。从\Person{亚巴郎}的肉身，通过\Person{依撒格}和\Person{雅各伯}，且不提及许多名字，通过\Person{玛利亚}，我们的主\Person{耶稣基督}降生。\par

2. 但是，如何证明根在圣祖们中呢？让我们询问\Person{保禄}。\Person{外邦人}现今相信\Person{基督}，并仿佛想要与钉死\Person{基督}的\Person{犹太人}夸耀；尽管从同一民族也来了另一面墙，在角石，即\Person{基督}本身处相遇，即未受割礼的面墙，也就是从不同方向来的\Person{外邦人}。我说，当各民族抬举自己时，祂就如此压制他们。\Quote{因为如果你，}祂说，\BibleQuote{从自然的野橄榄上被砍下来，却逆性接在它们中间，就不要对枝条夸口：因为如果你夸口，不是你托着根，而是根托着你。}\BibleRef{罗 11:17}因此，祂谈到一些因不信而从圣祖的根上折断的枝条，以及接在其中的野橄榄，使之分享橄榄的肥汁，即由\Person{外邦人}而来的教会。是谁把野橄榄接在橄榄上呢？橄榄通常接在野橄榄上；我们从未见过野橄榄接在橄榄上。因为如果有人这样做，会发现只结野橄榄的果子。因为被接上的枝条生长，所结的果子属于该品种。不是根的果子，而是接枝的果子。宗徒显示天主以祂的全能做了这事，即野橄榄被接入橄榄的根上，却不结野橄榄，而是结橄榄——将此归于天主的全能，宗徒说：\Quote{如果你从自然的野橄榄上被砍下来，且逆性被接入好橄榄上，就不要}他说，\BibleQuote{对枝条夸口。}\BibleRef{罗 11:17}…\par

3. 在旧约时代，弟兄们，我们的天主向那肉性子民的应许是属地和属时间的。应许了一个地上的王国，应许了他们脱离埃及后被带领进入的那片土地：由\Person{若苏厄}（农之子）带领他们进入应许之地，在那里建造了地上的耶路撒冷，\Person{达味}在那里为王：他们脱离埃及后，渡过红海，得到了那土地……这样的应许也是那些不能持久的，然而通过它们，预示了将要持久的未来应许，以致所有那些时间性的应许是未来事物的预象和一种预言。因此，当那王国衰落时，\Person{达味}（\Person{叶瑟}之子）在那里为王，他是一个人，虽然是一位先知，虽然是圣的，因为他看见并预见了基督要来，按肉体基督也要从他的后裔出生：但仍是人，还不是基督，还不是我们的王、天主之子，而是达味王、叶瑟之子：因为那时那王国将要衰落，通过接受那王国，当时天主被肉性的人赞美；他们只把这件事看为大事，即他们暂时从压迫他们的人手中被解救，通过红海逃脱了逼迫的仇敌，被带领经过旷野，找到了土地和王国：他们只为此赞美天主，尚未领悟天主在那些预象中预先设计并应许的事。因此，当那些肉性子民（达味统治他们）赞美天主的事物衰落时，\Quote{达味的圣咏}就止息了——不是天主之子的，而是\Quote{叶瑟之子}……\par

4. 这篇圣咏是谁的声音？\Quote{是阿撒夫的。}阿撒夫是什么？正如我们从希伯来文译成希腊文，再从那希腊文译成拉丁文的解释中所见，阿撒夫被解释为\OriginalTerm{会堂}{Synagogue}。所以这是会堂的声音。但你们一听到\Quote{会堂}，不要立刻憎恶它，好像它是谋杀主的人。那会堂确实是谋杀主的人，无人怀疑这一点；但要记住，从会堂中出了公羊，我们是那些公羊的后代。因此在一篇圣咏中说：\Quote{请将公羊的子孙献给上主。}那些公羊是谁？伯多禄、若望、雅各伯、安德肋、巴尔多禄茂和其余宗徒。因此他（扫禄）起初也是扫禄，后来成了保禄：即起初骄傲，后来谦卑……所以连保禄也是从会堂来到我们这里，伯多禄和其他宗徒也是从会堂来的。因此，当你们听到会堂的声音时，不要看它的功过，而要观察它的后代。所以在这篇圣咏中，是会堂在说话，是在达味的圣咏——即叶瑟之子的圣咏——止息之后，也就是在时间性的事物（天主曾通过它们被肉性子民赞美）止息之后。但这些为什么会止息呢？岂不是为了使其他事物被寻求？要寻求什么呢？难道是那里没有的东西吗？不，而是那些以预象隐藏在那里的事物：不是那些还不存在的事物，而是那些那里某种程度上隐藏在奥秘的隐秘之物中的事物。是什么事物？宗徒本人说：\BibleQuote{这些事是我们的预像。} \BibleRef{格前 10:6}……\par

5. 所以这是会堂，即那些以虔诚方式敬拜天主的人，但却是为了属地的、现世的事物（因为有不敬虔的人向魔鬼寻求现世的福乐；但这一民族在这点上比外邦人更好，因为虽然他们寻求的是现世和暂时的福乐，却是从唯一的天主那里寻求，他是万有——包括属神和属肉体的——的创造者）。因此，当那些按肉体虔诚的人观察时——即那由善人、当时的好人（不是属神的人，如其中的先知，如那些少数领悟属天、永恒国度的人）组成的会堂——那会堂，我说，观察到它从天主那里接受了什么，以及天主向那一民族应许了什么：丰盛的属地事物、土地、和平、属地的幸福；但在这一切中都有预象，而他们没有觉察到那里在预象中隐藏着什么，就认为天主给了这些作为大事，并且没有更好的东西给他所爱和事奉的人了；他们注意到并看见某些罪人、不敬虔者、亵渎者、魔鬼的仆人、魔鬼的儿子，生活在极大的邪恶和骄傲中，却拥有丰盛的此类属地、暂时的事物——正是他们自己为这类事物事奉天主的；于是心中生出一个极邪恶的念头，使脚蹒跚，几乎从天主的道上滑落。看，这个念头就在旧约的子民中：但愿它不在我们肉性的弟兄中，如今新约的幸福已公开宣告……\par

6. \Quote{以色列的天主何其良善！}但向谁呢？\Quote{向心地正直的人。} \BibleRef{咏 72:1}向乖僻的人呢？他显得乖僻。同样，在另一篇圣咏中他说：\Quote{你与圣洁的人成为圣洁，与无辜的人成为无辜，与乖僻的人成为乖僻。}\Quote{与乖僻的人成为乖僻}是什么意思？乖僻的人会认为你乖僻。绝不是天主变成了乖僻。绝无可能：他是什么，就是什么。但就像太阳对拥有清晰、健全、健康、强壮眼睛的人显得柔和，但对脆弱的眼睛却刺出硬矛——可以说是这样；前者注视它就得力量，后者则受折磨，但太阳本身并未改变，而是人改变了：同样，当你开始变得乖僻时，天主在你看来也会显得乖僻，是你改变了，不是他。因此，那对善人来说是喜乐的事物，对你来说将是惩罚。他记起这事，就说：\Quote{以色列的天主对心地正直的人何其良善！}\par

7. 但对你呢？\BibleQuote{我的脚几乎要滑倒}（\BibleRef{咏 72:2}）。脚何时滑倒？若不是心不正？心为何不正？听：\BibleQuote{我的脚步几乎跌倒}。他所说的\Quote{几乎}，与\Quote{险些}同义：他所说的\Quote{我的脚几乎要滑倒}，也就是\Quote{我的脚步跌倒}。我的脚几乎要滑倒，我的脚步几乎跌倒。脚滑动了，但脚为何滑动，脚步为何跌倒？脚滑动到走偏，脚步跌倒到坠落：不是完全，而是\Quote{几乎}。但这是什么？我已经要迷失，但尚未迷失：我已经要跌倒，但尚未跌倒。\par

8. 但为何甚至这样？\BibleQuote{因为我嫉妒罪人}，他说，\BibleQuote{看到罪人的平安}（\BibleRef{咏 72:3}）。我观察罪人，见他们享有平安。什么平安？暂时的、短暂的、易逝的、属地的：但这样的平安也是我渴望从天主那里得到的。我看见那些不事奉天主的人拥有我所渴望的东西，为的是我能事奉天主：于是我的脚滑动了，我的脚步几乎跌倒。但罪人为何拥有这些，他简要地说：\BibleQuote{因为他们的死亡无法逃避，他们的鞭笞是\OriginalTerm{坚定}{firmament}}（\BibleRef{咏 72:4}）。现在我已经明白，他说，为何他们享有平安，在地上兴旺；因为他们的死亡无法逃避，因为确实而永恒的死亡正在等待他们，这死亡既不会逃避他们，他们也无法逃避它，\BibleQuote{因为他们的死亡无法逃避，他们的鞭笞是坚定的}。而且他们的鞭笞是坚定的。因为他们的鞭笞不是暂时的，而是永恒坚定的。由于这些邪恶之事将来对他们来说是永恒的，现在呢？\BibleQuote{他们不在人的劳苦中，也不与人一同受鞭笞}（\BibleRef{咏 72:5}）。难道魔鬼本身不也逃避与人的鞭笞吗？尽管永恒惩罚正在为他预备？\par

9. 因此，就这一点而言，这些人行些什么，当他们不受鞭笞，不与人一同劳苦时？\BibleQuote{因此}，他说，\BibleQuote{骄傲抓住了他们}（\BibleRef{咏 72:6}）。观察这些骄傲、无纪律的人；观察那被献为祭物的\OriginalTerm{公牛}{bull}，被容许自由游荡；并破坏它所能破坏的，直到它被宰杀的日子。如今，弟兄们，这是一件好事，我们应在先知的话语中听到关于这头公牛的事，就像我所说的。因为圣经在另一处这样提到他们：他说他们好像被预备为祭物，并因邪恶的自由而被饶恕。\BibleRef{箴 7:22}\BibleQuote{因此}，他说，\BibleQuote{骄傲抓住了他们。}\BibleQuote{什么是\InnerQuote{骄傲抓住了他们}？}\BibleQuote{他们被自己的不义和不敬虔所披戴。}他没有说\Quote{遮盖}，而是\Quote{披戴}，从各方被他们的不敬虔完全包裹。他们确实可怜，既看不见也看不见，因为他们被包裹了；他们的内在部分看不见。因为凡能看见暂时似乎幸福的恶人的内在部分的，凡能看见他们折磨良心的，凡能审视他们被如此巨大的欲望和恐惧的扰动折磨的灵魂的，都会看到即使是当他们被称为幸福时，他们也是可怜的。但因为\BibleQuote{他们被自己的不义和不敬虔所披戴}，他们看不见；但也不被看见。圣神认识他们，那说这些话关于他们的：我们应当以同样的眼睛来审视这样的人，正如我们知道我们看见一样，如果从我们眼中除去不敬虔的遮盖……\par

10. 首先，描述的是这些人。\BibleQuote{他们的不义好像从肥膘中发出} \BibleRef{咏 72:7} ……一个贫穷的乞丐犯了偷窃罪；不义从瘦弱中发出；但当一个富人拥有如此多的财富时，他为什么还要抢夺别人的东西？前者的不义从瘦弱中发出，后者的不义从肥膘中发出。因此，当你对那瘦弱的人说：\Quote{你为什么做了这事？}他谦卑困苦地回答：\Quote{需要逼迫了我。} \Quote{你为什么没有敬畏天主？} \Quote{穷乏紧迫。}对一个富人说：\Quote{你为什么做这些事，而不敬畏天主？}——假设你足够重要以至于能说这话——看他是否甚至屈尊听闻；看是否甚至连你，也有不义从他的肥膘中向你发出。因为他们如今向他们的师长和责备者宣战，成为说实话之人的仇敌，长久以来习惯了被谄媚者的话语哄骗，耳朵柔软，心不健全。谁敢对一个富人说：\Quote{你抢夺别人的财物，做了恶事？}或者，若有人胆敢说话，且他是这样一个不能抵抗的人，他如何回应？他所说的一切都是蔑视天主。为什么？因为他骄傲。为什么？因为他肥。为什么？因为他被献为祭品。\Quote{他们逾越到了心中的计谋。}在这里，他们逾越了。什么是\Quote{他们逾越了}？他们越过了道路。什么是\Quote{他们逾越了}？他们超越了人类的界限，不认为自己和其余人一样。我说，他们逾越了人类的界限。当你对这样的人说：\Quote{这个乞丐是你的弟兄；}当你对这样的人说：\Quote{这个穷人是你的弟兄；}你们有相同的父母，\Person{亚当}和\Person{厄娃}：不要在意你的高傲，不要在意你被高举的虚浮；尽管有仆役围绕你，尽管有数不尽的金银，尽管有大理石房屋容纳你，尽管有雕花天花板遮蔽你，你和穷人共同拥有那宇宙的屋顶——天空；但你和穷人的不同在于那不属于你、从外部附加在你身上的东西：你在它们中看见你自己，而不是它们在你们中。观察你自己，你与穷人的关系；是你自己，而不是你所拥有的。为什么你藐视你的弟兄？在你们母亲的腹中，你们都是赤身露体的。诚然，即使当你离开此生，这些身体腐烂，灵魂呼出后，让富人和穷人的骨头区分出来吧！我所说的，是状况的平等，是人类共同的命运，人人都在其中出生：因为在此一个人成为富人，而穷人不会永远在此；正如富人并非带着财富而来，他也不会带着财富而去；两者相同的进入，相似的离去。我再加上一句，也许你们会交换处境。如今福音处处被传讲：观察一个满身疮痍的穷人，正躺在富人家门前，\BibleRef{路 16:19}渴望用富人桌子上掉下来的碎渣充饥；也观察那个像你的人，身穿紫红袍和细麻衣，天天奢华宴乐。我说，那穷人死了，被天使们带到\Person{亚巴郎}的怀抱里；但那个富人死了，并被埋葬了；也许没有人关心他的埋葬……弟兄们，那穷人的辛劳何等巨大！那富人的奢华何等长久！但他们所换得的境况却是永久的……为时已晚，他配得地说：\BibleQuote{请派遣\Person{拉匝禄}，} \BibleRef{路 16:27} \BibleQuote{也叫他去告诉我的弟兄们。}因为自己没有获得悔改的果实。不是因为悔改没有被给予，而是悔改将是永远的，悔改之后没有救恩。因此这些人\Quote{他们逾越到了心中的计谋。}\par

11. \BibleQuote{他们图谋并说出恶毒} \BibleRef{咏 72:8}。但人们甚至带着恐惧说出恶毒：但这些人如何？\BibleQuote{他们高谈不义。}他们不仅说不义，而且公开地，在众人耳中，骄傲地：\Quote{我要这样做；} \Quote{我要给你点颜色看看；} \Quote{你会知道你要和谁打交道；} \Quote{我绝不让你活。}你本可以只思想这些事，不将其说出来！至少在思想的密室中，邪恶的欲望可以被限制，他至少可以在思想中克制它。为什么？难道他瘦弱吗？\BibleQuote{他们的不义好像从肥膘中发出。} \BibleQuote{他们高谈不义。}\par

12. \BibleQuote{他们以口对抗上天，他们的舌头凌驾于大地之上} \BibleRef{咏 72:9}。因为这里\Quote{凌驾于大地之上}是指，他们越过一切世俗之事？越过一切世俗之事是什么意思？他说话时，不把自己当作一个会突然死去的人；他好像要永远活着一样威胁：他的思想超越世俗的脆弱，他不知道自己是包裹在何等样的器皿里；他不知道在别处针对这样的人所写的话：\Quote{他的灵魂要出去，他要归于他的尘土，在那一天，他的一切思想都要消亡。}但这些人不想念他们的末日，说出骄傲之言，以口对抗上天，他们凌驾于大地之上。如果一个强盗不想他的末日，即他受审的末日，当他被送进监狱时，没有什么比他更可怕的了：然而他可能逃脱。你往哪里逃能逃脱死亡？那一天是确定的。你还有多长的时间可活？那有尽头的长时有多长，即使它很长？再加上它是无：那被称为长时的东西本身并不是长时，而且是不可靠的。为什么他不思想这些？因为他以口对抗上天，他的舌头凌驾于大地之上。\Quote{岁月满盈必在他们身上找到。}\par

13. \BibleQuote{因此，我的百姓必要归回这里} \BibleRef{咏 72:10}。如今，阿撒夫自己也归回这里。因为他看见这些事在恶人身上丰盛，在骄傲人身上丰盛：他正归向上主，开始询问和探讨。但何时呢？\BibleQuote{当在他们内寻得满全的日子时。} 什么是\BibleQuote{满全的日子}？\BibleQuote{但时期一满，天主就派遣了他的儿子。} \BibleRef{迦 4:4} 这就是那时期的圆满本身，当他来教导人轻视暂时的事物，不把恶人所贪求的任何对象看作大事，承受恶人所恐惧的一切。他成了道路，他召叫我们回归内心的思想，告诫我们应从天主寻求什么。请看，他从一种自我反作用于自身的思想，并且仿佛召回其冲动的波浪，进而过渡到选择真实的事物。\par

14. \BibleQuote{他们便说：天主怎能知晓？在至高者那里岂有知识？} \BibleRef{咏 72:11}。请看，他们经过怎样的思想。看，不义之人有福，天主不关心人间的事。他果真知道我们所做的事吗？请看人们正在说什么。我们正在探讨，弟兄们，\BibleQuote{天主怎能知晓，}等等（愿基督徒不再这样说）。因为你看来天主不知道，至高者那里没有知识，这是怎样的呢？他回答说：\BibleQuote{看哪，他们自身是罪人，在世上却获得了丰厚的财富。} \BibleRef{咏 72:12} 他们既是罪人，在世上却获得了丰厚的财富。他承认他不愿成为罪人，以便能拥有财富。一个属血肉的灵魂，为了可见和属世的事物，竟会出卖自己的义德。那种为了黄金而保留的义德是什么样的义德？仿佛黄金比义德本身更宝贵，或者，当一个人否认他人财物的寄存物时，被告否认的人所受的损失，比否认的人更大。前者失去一件衣服，后者失去忠诚。\BibleQuote{看哪，他们自身是罪人，在世上却获得了丰厚的财富。} 因此，天主不知晓，因此，在至高者那里没有知识。\par

15. \BibleQuote{我便说：因此我白白地洁净了我的心} \BibleRef{咏 72:13}。因为我事奉天主，却没有这些；他们不事奉他，却富有这些：\BibleQuote{因此我白白地洁净了我的心，在无辜者中洗净了我的手。} 我白白地做了这事。我善生的酬报在哪里？我事奉的报酬在哪里？我生活良善，却处于贫困；而不义之人却富有。\BibleQuote{我在无辜者中洗净了我的手，我终日受鞭打。} \BibleRef{咏 72:14} 天主的鞭打并未离开我。我事奉得好，却受鞭打；他不事奉，却受尊荣。他给自己提出了一个大问题。灵魂不安，灵魂将过去那些将要逝去的事物，转变成轻看属地、渴慕永恒的事物。在这思想中，有一个灵魂自身的过渡；她在某种风暴中颠簸，终将抵达港口。这有如病人，当康复尚远时，病情较轻；当康复在即，他们发更高的热；医生称之为\Quote{\OriginalTerm{危象}{critical accession}}，借此他们过渡到健康：更高的热，却导向健康；更高的温度，但康复在即。此人也是如此在发烧。因为，弟兄们，这些是危险的话，令人反感，几乎亵渎：\Quote{天主怎能知晓？} 这就是为什么我说\Quote{几乎}；他没有说天主不知道；他没有说至高者没有知识；而是仿佛在询问、犹豫、怀疑。这与他在前面所说的\Quote{我的脚步几乎滑倒}是一样的。他没有肯定，但怀疑本身是危险的。他正在从危险中过渡到健康。现在请听健康之言：\BibleQuote{因此我白白地洁净了我的心，在无辜者中洗净了我的手：我终日受鞭打，我的责罚在早晨。} 责罚就是纠正。受责罚者正在被纠正。什么是\Quote{在早晨}？它没有被延迟。恶人的责罚被延迟了，我的没有被延迟：前者太迟了，或根本不来；我的在早晨。\par

16. \\BibleQuote{若我说：我要这样述说；看，我遗弃了你众子的世代} \\BibleRef{咏 72:15}：这就是说，我要这样教训。你如何教训？说至高者没有知识，天主不知道吗？你将要提出这种意见，说人度正义的生活是没有原因的，义人丧失了他的劳苦，因为天主更偏袒恶人，或者祂谁也不关心？你将这样说、这样宣告吗？祂用一种权威抑制自己。什么权威？人有时想要迸发这种情绪，但被圣经召回，圣经时常指导我们度善生，说天主关心人间事，祂区分虔敬的人和不虔敬的人。因此，这个人也想要提出这种情绪，他克制自己。而他说什么？\\BibleQuote{我遗弃了你众子的世代。}正如有些抄本所载：\\BibleQuote{看，你众子的世代，我与它和谐一致：}这就是说，与你众子的世代，我与之和谐一致；也就是说，我同意它，我与之相符：若我这样教训，我就与所有人都不合拍。因为凡和谐歌唱的人，是与人一同奏调；但不同奏调的人，就不是和谐歌唱。难道我要说与亚巴郎所说、与依撒格所说、与雅各伯所说、与众先知所说不同的东西吗？因为众人都说天主关心人间事，我难道要说祂不关心吗？难道我比他们更有智慧？比他们更明白？一种最有益处的权威将他的思想从不敬神中唤回。然后呢？为使他不再遗弃，他做了什么？\\BibleQuote{我决意要明白} \\BibleRef{咏 72:16}。愿天主与他同在，使他明白。同时，弟兄们，他被阻止从大堕落中，当他不以为自己已经知道，而是决意明白他所不知道的。因为刚才他愿意显得好像知道，并宣告天主不关心人间事。因为这已成为不义之人最邪恶、最不敬神的教义。要知道，弟兄们，许多人争论说天主不关心人间事，说万事由偶然统治，或我们的意志受制于星宿，每人不是按功德对待，而是由他星宿的必然性——一种邪恶的教义，一种不敬神的教义。那脚步几乎滑倒、步伐几乎倾倒的人正要走向这些思想，进入这个错误；但因为他与天主子民的世代不和谐，他决意要明白，并定断那与他不同意于天主义人的知识。让我们听他说什么：他如何决意要明白，得到帮助，学到一些东西，并宣告给我们。他说：\\BibleQuote{我决意} \\BibleQuote{要明白。} \\BibleQuote{这劳苦在我面前。}真是巨大的劳苦：要明白天主如何既关心人间事，而恶人享福，善人却受苦。这问题意义重大，因此，\\BibleQuote{这劳苦在我面前。}仿佛有一堵墙立在我面前，但你有圣咏的声音：\\BibleQuote{靠着我的天主，我要越过这墙。}\par

17. ……他这样做了；因为他说明劳苦在他面前多久；\\BibleQuote{直到我进入天主的圣所，并明白终末之事} \\BibleRef{咏 72:17}。这是一件大事，弟兄们：他说，现在我长期劳苦，在我面前我看到一种难以克服的劳苦，要明白天主如何既是公义的，又关心人间事，并且并非因为罪人作恶而在地上享福就成为不义；但虔诚人和事奉天主的人却常消磨于试炼和劳苦中；要知道这事是极大的困难，但只有\\BibleQuote{直到我进入天主的圣所。}因为在圣所里，为了让你解决这问题，呈现给你的是什么？他说：\\BibleQuote{我明白了} \\BibleQuote{终末之事：}不是现今之事。他说，我从天主的圣所极目到终局，我越过现今之事。所有被称为人类的，所有那必死的群体，都要来到天平前，都要来到秤前，在那里人的行为将被称量。如今万物被云遮蔽；但天主知道各人各自的功过。他说：\\BibleQuote{我明白了} \\BibleQuote{终末之事：}但不是靠我自己；因为劳苦在我面前。我怎么\\BibleQuote{明白终末之事}？让我进入天主的圣所。在那里他也明白了这些人如今享福的原因。\par

18. 也就是，\Quote{因奸诈祢安放在他们身上}\BibleRef{咏 72:18}。因为他们本奸诈，即是欺诈；因他们本奸诈，他们反受欺诈。这是什么意思？因他们是欺诈的，他们反受欺诈？他们欲在一切恶行中欺骗人类，他们自己却也受欺骗，即选择现世之善而离弃永恒之善。因此，弟兄们，在他们行骗之时，他们自己反受欺骗。就我刚才所说，弟兄们，\Quote{那为得一件衣服却失去自己信实的人，有何智慧呢？}那衣服被夺的人受了欺骗，还是那蒙受如此大损失的人受了欺骗？若衣服比信实更珍贵，则前者所受损失更大；但若信德无可比拟地超越全世界，则后者似乎只损失了衣服；然而对前者说：\Quote{人纵然赚得了全世界，却赔上了自己的灵魂，为他有什么益处？}\BibleRef{玛 16:26}。因此他们遭遇了什么？\Quote{因奸诈祢为他们安排：在他们被高举之时，祢把他们推倒。}他没有说：祢因他们被举而推倒他们；并非他们被举之后祢才推倒他们，而是在他们被举之时，他们就已经被推倒了。因为如此被举即是跌倒。\par

19. \Quote{他们怎样忽然成了荒凉？}\BibleRef{咏 72:19}。他是在惊奇他们，并领悟到末后的结局。\Quote{他们已消失}。诚如烟，它上升时就消散，他们也如此消失了。他说\Quote{他们已消失}是什么意思？如同一个领悟末后结局的人所说的：\Quote{他们因自己的不义而灭亡}。\Quote{如人睡醒后的梦}\BibleRef{咏 72:20}。他们如何消失？如同一个睡醒之人的梦消失。设想一个人在睡梦中看见自己找到了宝藏；他成了富翁，但只到醒来为止。\Quote{如人睡醒后的梦}：他们就这样消失了，如同一个醒来之人的梦。于是人们寻找它，却找不到；手里空空，床上空空。一个穷人睡去，在梦中成了富翁：他若不醒，便是个富翁；他醒了，便找回了梦中失去的忧虑。这些人也必将找到他们为自己准备的苦难。当他们从这生命苏醒时，那像在睡梦中紧握的东西就消失了。\Quote{如人睡醒后的梦}。免得有人问：\Quote{那么，他们的荣耀在你眼中算小事吗？他们的地位在你眼中算小事吗？碑文、肖像、雕像、荣誉、成群的食客在你眼中算小事吗？}\Quote{哦，主啊，}他说，\Quote{在你的城中，你将使他们的肖像归于虚无。}……他夺去富人的骄傲；他给出劝告。好像他们说：我们是富人，你禁止我们骄傲，禁止我们夸耀财富的排场：那么我们该怎样处理这些财富？难道到了这样的地步，以致什么也不能做吗？他说：\Quote{让他们在善工上富足，慷慨施舍，乐意通财。}\BibleRef{弟前 6:18} 这有什么益处呢？\Quote{让他们为自己积蓄美好的基础，以备将来，使他们能把握真正的生命。}\BibleRef{弟前 6:19} 他们该在哪里为自己积蓄财宝？就在他进入天主的圣所时所注目的那地方。愿我们所有富足的弟兄们，拥有大量金钱、金银、房子、荣誉的，都因刚才所说的话而战栗：\Quote{你将使他们的肖像归于虚无。}他们难道不配受这些吗？即是天主在他的城中使他们的肖像归于虚无，因为他们自己已经在他们的世俗城中使天主的肖像归于虚无了。\par

20. \Quote{因为我的心喜悦}\BibleRef{咏 72:21}。他是在说出自己被什么所诱惑：\Quote{因为我的心喜悦，}他说，\Quote{我的腰子\OriginalTerm{腰子}{reins}也改变了。}当那些现世之事令我喜悦时，我的腰子改变了。也可以这样理解：\Quote{因为我的心喜悦}于天主，我的腰子也改变了，就是说，我的私欲改变了，我变得完全贞洁。\Quote{我的腰子改变了。}且听如何改变。\Quote{我成了虚无，竟不自知}\BibleRef{咏 72:22}。我，正是现在谈论富人这些事的这个人，也曾向往过诸如此类的事：因此，当我的脚步几乎滑倒时，\Quote{连我也成了虚无}。\Quote{我成了虚无，竟不自知。}所以，即使是他们，就是我所指责的那些人，我们也不该对他们绝望。\par

21. \Quote{不自知}是什么意思？\Quote{我好似走兽\OriginalTerm{走兽}{beast}成为于祢面前，但我常与祢同在}\BibleRef{咏 72:23}。这人与他人之间有很大区别。他因渴求现世之物而变得好似走兽，当他成为虚无而不知永恒之事时：但他并没有离弃他的天主，因为他并没有从魔鬼、从恶魔那里渴求这些事。因我已向你们指出过这一点。这声音来自会众，即来自那不敬拜偶像的百姓。我确实成了走兽，当我从天主那里祈求现世之物时：但我从未离弃我的那位天主。\par

22. 因为那时，虽然我已成了走兽，却没有离开我的天主，所以接下去说：\Quote{祢握住了我右手的右手。}祂没有说我的右手，而是说\Quote{我右手的右手。}如果它是右手的右手，那么一只手就有一只手。\Quote{祢所握住的我右手的右手，}为的是祢能引导我。因为祂把手放在哪里？为了权能。因为我们常说，一个人手中拥有什么，就是他权能之下有什么：正如魔鬼对天主论及约伯时所说的：\BibleQuote{祢伸手打击他，夺取他所有的一切。}\BibleRef{约 1:11}什么是\Quote{伸手}？就是施展权能。他把天主的\Quote{手}称为天主的权能：正如在另一处所记载的：\BibleQuote{死亡和生命，都在舌头的权下。}\BibleRef{箴 18:21}舌头有手吗？但什么是\Quote{在舌头的权下}？在舌头的权能之下。什么是\Quote{在舌头的权能之下}？\BibleQuote{凭你的口，你将成义；凭你的口，你将受罚。}\BibleRef{玛 12:37}因此，\Quote{祢握住了我右手的右手，}即我右手的权能。我的右手是什么？就是我常与祢同在。我偏向左，因为我成了走兽，也就是说，因为我内心有属地的贪欲；但右手是我的，因为我常与祢同在。祢握住了我这右手的\Quote{手}，就是说，祢引导了那权能。什么权能？\BibleQuote{祂给了他们权能，成为天主的子女。}\BibleRef{若 1:12}他现在开始成为天主的子女，属于新约。请看，他右手的右手是如何被握住的。\Quote{祢以祢的旨意引导了我。}什么是\Quote{以祢的旨意}？不是凭我的功德。什么是\Quote{以祢的旨意}？请听那位曾一度是走兽、贪恋属地事物、按旧约生活而生活的宗徒。他怎么说？\BibleQuote{我原先是个亵渎者、迫害者和施暴者，但我蒙受了怜悯。}\BibleRef{弟前 1:13}什么是\Quote{以祢的旨意}？\BibleQuote{靠天主的恩宠，我成为今日的我。}\BibleRef{格前 15:10}\Quote{祢以光荣接引了我。}现在他被接到什么光荣里，又是在什么光荣里，谁能解释、谁能述说？让我们期待吧，因为那将是在复活的时候，在末日的事上。\par

23. 他开始思想同一天上的福乐，并责备自己，因为他曾成了走兽，贪恋属地的事物。\BibleQuote{除祢以外，在天上我还有谁？在地上我也别无眷恋。}\BibleRef{咏 72:25}。从你们的声音中，我看得出你们已经明白了。他将自己属地的意愿与将要领受的天上赏报作了比较；他看见那里为他所保留的是什么；在思想并热切思念某种不可言喻的事物时——这事物是眼所未见，耳所未闻，人心也未曾想到的，\BibleRef{格前 2:9}——他没有说\Quote{我在天上有这个或那个}，而是说：\Quote{我在天上有什么？}那我在天上所拥有的是什么呢？它是什么？它是多么伟大？属于哪一类？\Quote{并且，}因为我在天上所拥有的不会消逝，\Quote{我从祢那里在地上又意愿了什么呢？}……祢为我保留了不朽的天上财富，就是祢自己，而我却曾从祢那里意愿了地上那不敬神的人也能拥有、恶人也能拥有、甚至堕落之人也能拥有的东西：金钱、黄金、白银、珠宝、家业，这些许多坏人也拥有，许多放荡的男女也拥有：我曾把这些当作大事向我的天主在地上求取，尽管我的天主在天上亲自为我保留了自己！\par

24. \BibleQuote{我的心肠和我的肉体衰残，我心中的天主，} \BibleRef{咏 72:26} 这就是为我所保留在天上的：\BibleQuote{我心中的天主，我的福分就是天主。}这不是暂时的\Quote{我的福分}，而是\BibleQuote{我的福分就是天主，直到永远。}因为请看祂是怎样爱了他：祂使他心怀洁德：\BibleQuote{我心中的天主，我的福分就是天主，直到永远。}他的心成了洁净的，因为如今天主被白白地爱着，不向祂寻求任何别的赏报。谁若向天主寻求别的赏报，并因此愿意事奉天主，他就把他愿意领受的看为比赐予者的天主更珍贵。那么，难道天主就没有赏报吗？没有除了祂自己以外的赏报。属于天主的赏报，就是天主自己。祂爱这个，珍视这个；若爱任何别的东西，这爱就不是纯洁的。你正远离不朽的烈火，你将变冷，将被腐化。不要远离。不要远离，远离就是你的败坏，就是你的淫乱。如今他回来了，如今他悔改了，如今他选择了悔改，如今他说：\BibleQuote{我的福分就是天主。}他是怎样因他所选择为福分的那一位而喜悦呢。\par

25. \BibleQuote{看，那自远于祢的人必要丧亡} \BibleRef{咏 72:27}。他因此离开了天主，但未远离；因为他说：\BibleQuote{我几乎成了畜类}，并且\BibleQuote{我时常与祢同在}。但他们却远离了，因为他们不仅贪求世俗之物，更从魔鬼和恶魔那里寻求。\BibleQuote{那自远于祢的人必要丧亡}。那么，远离天主是什么意思？\BibleQuote{祢消灭了所有背离祢而行淫的人}。与此种淫乱相对的是贞洁的爱。什么是贞洁的爱？灵魂热爱她的新郎：她向所爱的新郎要求什么？或许如同女子为自己选择女婿或新郎：她可能选择财富，爱他的金子、地产、银子、牛马、家仆等。绝非如此。她单爱祂本人，爱祂毫无所求：因为在祂内拥有一切，因为\BibleQuote{万物是藉着祂造的} \BibleRef{若 1:3}。\par

26. 但你们做了什么？\BibleQuote{至于我，亲近天主为我是有益的} \BibleRef{咏 72:28}。这就是完全的益处。你们还要更多吗？我因你们的意愿而忧心。弟兄们，你们还要什么？除了亲近天主，没有更好的，\BibleQuote{那时我们将面对面看见祂} \BibleRef{格前 13:12}。但现在呢？我还如客旅在说话：他说：\BibleQuote{亲近}，\BibleQuote{天主为我是有益的}；但现在在我旅居之时（因为本体尚未到来），我\BibleQuote{将我的希望寄托于天主}。因此，只要你们尚未亲近，就应将希望寄托于那里。你们在摇摆，将锚抛向陆地。你们尚未借着临在亲近，却要借着希望牢牢亲近。\BibleQuote{将我的希望寄托于天主}。在此你们将如何寄托你们的希望？你们的事业是什么，岂非赞美你们所爱的祂，并使他人与你们一同爱祂？看，若你们爱一个车夫，难道不会拉拢其他人一同爱他吗？爱车夫的人无论去哪里都谈论他，以便他人也能爱他。难道被遗弃的人被白白地爱着，而爱天主却要求回报吗？你们要白白地爱天主，不将天主据为己有。...因为接下来是什么？\BibleQuote{为使我能在熙雍女子的一切庭院中宣扬祢的一切赞美}。\BibleQuote{在庭院中}：因为在教会之外宣讲天主是徒劳的。赞美天主并宣示祂的一切赞美是一桩小事。你们要在熙雍女子的庭院中宣扬。要联合一体，不要分裂人民；而是将他们引向一体，使他们成为一体。我忘记了自己说了多久。现在这篇圣咏结束了，即使根据这个接近度，我想我做了长篇讲道，但这不足以满足你们的热情；你们过于急躁。但愿你们以这种急躁夺取天国。\par

\SourceRecord{Translated by J.E. Tweed.；Nicene and Post-Nicene Fathers, First Series；Vol. 8.；Edited by Philip Schaff.；Buffalo, NY: Christian Literature Publishing Co.,；1888.；<http://www.newadvent.org/fathers/1801073.htm>.}

% structure: p0001 source=h1 target=section reason=page-title

\section{圣咏第七十四篇释义}

1. 这篇圣咏的标题是：\Quote{关于阿撒夫的领悟}。\Person{阿撒夫}在拉丁文中译为\Quote{集会}，在希腊文中译为\OriginalTerm{会堂}{Synagogue}。让我们看看这个会堂领悟了什么。但首先让我们理解会堂；由此我们将理解会堂所领悟的。通常，任何集会都可称为会堂：既有兽类的集会，也有人类的集会；但当我们听到\Quote{领悟}时，这里便没有兽类的集会……因为圣咏的标题这样规定说：\Quote{关于阿撒夫的领悟}。因此，我们将要听见的是某个具有领悟的集会的声音。但由于会堂通常特别指以色列人民的集会，以致无论我们在何处听到会堂，通常都理解为犹太民族；让我们看看，或许这篇圣咏中的声音不是来自那个民族。但那是什么样的犹太人和什么样的以色列人民呢？因为他们不是糠秕，而或许是谷粒；\BibleRef{玛 3:12}；不是折断的树枝，而或许是那些坚固的枝子。\BibleQuote{因为不是所有以色列人都是真以色列人}\BibleRef{罗 9:6}……因此有些真正的以色列人，其中有一位曾被称为：\BibleQuote{看，这确是一个以色列人，在他内毫无诡诈}\BibleRef{若 1:47}。我不是说我们也是以色列人的方式，因为我们是亚巴郎的后裔。因为宗徒曾对外邦人说：\BibleQuote{所以你们是亚巴郎的后裔，按照恩许是承继人}\BibleRef{迦 3:29}。因此，按照这个，凡追随我们祖先亚巴郎的信德之足迹的，我们都是以色列人。但在这里，让我们像宗徒所说的那样理解以色列人的声音：\BibleQuote{因为我自己也是以色列人，亚巴郎的后裔，本雅明支派}\BibleRef{罗 11:1}。因此，在这里让我们理解先知们所说的：\BibleQuote{残余者将得救}\BibleRef{罗 9:27}。因此，让我们在此处听见那残余得救者的声音；这样，那接受了旧约并专注于现世许诺的会堂就能发言；正因如此，他们的脚才动摇。因为在另一篇标题也是阿撒夫的圣咏中，说了什么？\Quote{天主待正直的人何其美善。但我的脚几乎滑跌。}好像我们要问：你的脚为什么滑跌？他说：\Quote{几乎，我的脚步滑倒，因为我嫉妒恶人的事，见恶人的享乐}。因为他依照旧约中天主的许诺期待现世的幸福，却看见这些幸福在恶人身上满盈；那些不敬拜天主的人，竟拥有他从天主那里期待的这些东西；仿佛他白白侍奉了天主，他的脚就摇动了……但恰巧，不是出于我们自己，而是出于天主的安排，我们刚才从福音中听到：\BibleQuote{法律是藉梅瑟传授的，恩宠和真理却是由耶稣基督而来的}\BibleRef{若 1:17}。因为如果我们分辨两个约，旧约和新约，它们的圣事不同，许诺也不同；然而，诫命大多相同……仔细审视时，要么全都相同，要么福音中几乎没有什么不是先知们说过的。诫命相同，圣事不同，许诺不同。让我们看看为什么诫命相同；因为我们必须依照这些来侍奉天主。圣事不同，因为有些圣事赐予救恩，另一些则许诺救主。新约的圣事赐予救恩，旧约的圣事许诺救主。因此，当你现在已拥有所许诺的事物，为何还寻求那许诺的事物？既然已拥有救主？……天主通过新约从祂儿子们手中拿走了那些如同孩童玩具的东西，以便在他们成长时赐予更有益的东西；因此决不能认为那些先前的东西不是祂所赐的。两者都是祂所赐的。但法律本身是藉梅瑟传授的，恩宠和真理却是由耶稣基督而来的：\BibleRef{若 1:17} 恩宠是因为那藉文字命令的，藉爱德得以满全；真理是因为所许诺的得以实现。因此，这就是阿撒夫所领悟的。总而言之，所有曾向犹太人许诺的事物都被夺走了。他们的王国在哪里？圣殿在哪里？傅油在哪里？司祭在哪里？他们中间的先知如今在哪里？从那一位由先知预言者来临的时刻起，那个民族中现在已没有这些事物；如今她失去了现世的事物，却尚未寻求天上的事物。\par

2. 因此，你们不应紧抓住尘世的事物，虽然天主确实赐予它们……你们看，犹太人在害怕失去尘世事物时，怎样杀害了天上之王。他们得到了什么？他们甚至失去了那些尘世事物本身：在他们杀害基督的地方，他们自己被杀害；当不愿失去土地而杀害了生命的赐予者时，他们失去了那同一片被杀害的土地；就在他们杀害祂的那一刻，正是为了使他们在那一刻得以明白他们遭受这些事的原因。因为当犹太人的城市被推翻时，他们正在庆祝逾越节，成千上万的人，整个民族都聚集在一起庆祝那个节日。在那里，天主（确实是通过恶人，但祂自己是良善的；通过不义的人，但祂自己是公义的，并且公义地）这样报复了他们，以至于成千上万的人被杀害，城市本身被摧毁。在这篇圣咏中，就是关于这件事，\Quote{阿撒夫的训诲}在抱怨，而在这抱怨中，训诲仿佛区分了尘世事物与天上事物，区分了旧约与新约：为叫你们看见你们正经历什么，你们应寻求什么，应舍弃什么，应依恋什么。他就是这样开始的。\par

3. \BibleQuote{天主，祢为何永远抛弃了我们？}\BibleRef{咏 73:1}。\Quote{永远抛弃了}，以那称为会堂的集会之身份。\BibleQuote{天主，祢为何永远抛弃了我们？}祂并非责备，而是追问\Quote{为何}，为了什么目的，因为什么祢做了这事？祢做了什么？\Quote{祢永远抛弃了我们。}什么是\Quote{永远}？或许直到世界的终结。祢抛弃了我们直到基督，祂是每个相信者的终结？\BibleRef{罗 10:4}。因为\BibleQuote{天主，祢为何永远抛弃了我们？}\BibleQuote{祢的灵向祢羊群的羊发怒。}祢为何向祢羊群的羊发怒？岂不是因为我们紧抓尘世的事物，却不认识牧人？\par

4. \BibleQuote{求祢记得祢从起初就拥有的会众}\BibleRef{咏 73:2}。这难道是外邦人的声音吗？祂从起初就拥有了外邦人吗？不，祂拥有了亚巴郎的后裔，即按肉身而生的以色列民族，由我们的祖先圣祖们所生：我们成为了他们的儿子，不是由他们的肉身所生，而是因效法他们的信德。但那些从起初就被天主拥有的人，遭遇了什么？\BibleQuote{求祢记得祢从起初就拥有的会众。祢赎回了祢产业的棍杖。}那同一会众，作为祢产业的棍杖，祢赎回了。这同一会众，他称为\Quote{产业的棍杖}。让我们回顾最初所做的事，当祂愿意拥有那同一会众，从埃及拯救它时，祂给了梅瑟什么记号，当梅瑟对祂说：\BibleQuote{我该给他们什么记号，使他们相信是祢派遣了我？天主对他说：你手里拿的是什么？一根棍杖。把它丢在地上，等等。}\BibleRef{出 4:1}这暗示什么？因为这不是徒然做的。让我们查考天主的著述。蛇诱惑人什么？死亡。因此死亡来自蛇。如果死亡来自蛇，那么蛇中的棍杖就是死亡中的基督。因此，当他们在旷野中被蛇咬而致死时，主命令梅瑟在旷野中举起一条铜蛇，并告诫百姓，凡被蛇咬的，仰望它就必痊愈。\BibleRef{户 21:8}事情就这样成了：同样，被蛇咬的人因仰望蛇而中了毒得痊愈。\BibleRef{若 3:14}因蛇而得痊愈是一件伟大的圣事。因仰望一条蛇而从蛇痊愈是什么？那就是因相信一位死者而从死亡痊愈。然而梅瑟害怕并逃跑了。\BibleRef{出 4:3}梅瑟从那蛇逃跑是什么？弟兄们，岂不是我们已知在福音中发生的事吗？\BibleRef{路 24:21}基督死了，门徒们害怕，并从那曾有的盼望中退缩了。……但是，当时有数千犹太人，就是钉死基督的人，相信了：因为他们被发现就在眼前，他们如此相信，以至于变卖了一切所有，把货价放在宗徒脚前。\BibleRef{宗 4:34}因为这事原是隐藏的，而天主棍杖的救赎在外邦人中更为显著：他解释他所说的话，就是他曾说的：\Quote{祢赎回了祢产业的棍杖。}他这话不是论及外邦人，因为在他们身上是明显的。但论的是什么？\Quote{熙雍山。}然而就连熙雍山也可以有别的理解。\Quote{就是祢曾居住的那一座。}在那人民从前所在的地方，那里立了圣殿，那里举行了祭献，那里当时有这一切预表基督的必要事物。一个许诺，当所许诺的事物已被赐予时，现在就成了多余的了……\par

5. \BibleQuote{举祢的手于他们骄傲的终结} \BibleRef{咏 73:3}。如同祢在终结时击退了我们，同样地\Quote{举祢的手于他们骄傲的终结}。谁的骄傲？那些推翻耶路撒冷者的骄傲。但被谁推翻，岂非由外邦人的君王？祢的手举起来很公义，因为他们如今也已认识基督。\BibleQuote{因为法律的终结就是基督，使一切相信的人成义。} \BibleRef{罗 10:4}他多么祝愿他们！仿佛在愤怒中说话，看似在说恶事：但愿他说的恶事发生！不，如今在基督之名下，我们应欢喜它正在发生。如今执掌权杖者正在服从十字架圣言；如今先知预言的正在应验：\Quote{大地众君王都要崇拜祢，万国都要事奉祢。}如今在君王的额上，十字架圣号比冠冕上的宝石更珍贵。\Quote{举祢的手于他们骄傲的终结。仇敌的恶意在祢的圣地行了何等的大事！}在那些曾是祢圣地的地方，即圣殿、司祭职、以及当时的一切圣事中。的确，那时仇敌行了恶事。因为那时行这些事的外邦人，敬拜假神，崇拜偶像，服事魔鬼；尽管如此，他们在天主的圣者身上行了许多恶事。他们何时能行，除非被允许？但何时会被允许，除非那些起初应许的圣物不再需要，因为应许者本身已被抓住？因此，\Quote{仇敌的恶意在祢的圣地行了何等的大事！}\par

6. \BibleQuote{凡恨祢的人都自夸了} \BibleRef{咏 73:4}。看看这些魔鬼的仆人，偶像的仆人：例如那时的外邦人，当他们推翻天主的圣殿和城时，\Quote{他们自夸了。} \Quote{在祢的庆节中。}记住我说过的，耶路撒冷正是在庆祝那个节期时被推翻的：就是他们钉死主的节期。他们聚在一起发怒，聚在一起灭亡。\BibleQuote{他们竖立了他们的记号，自己的记号，却不晓得} \BibleRef{咏 73:5}。他们有记号要竖在那里，他们的军旗、他们的鹰、他们的龙、罗马的记号；或者甚至是他们起初放置在圣殿里的雕像；或许\Quote{他们的记号}是他们从自己先知的魔鬼那里听到的事。\Quote{却不晓得。}不晓得什么？\BibleQuote{你们对我毫无权柄，除非从上头赐给你们。} \BibleRef{若 19:11}他们不晓得，那加给他们的尊荣不是要他们去折磨、夺取或推翻城，而是他们的不敬虔成了天主的斧子。他们成了祂发怒的工具，却并非祂平息后的国度。因为天主也做一个人常做的事。有时一个人在怒中捡起路上一根棍子，随便什么枝条，用它打他的儿子，然后将棍子丢进火里，却为儿子保留产业：同样，天主有时通过恶人来训导善人，并借着那些将被定罪者暂时的权势，使那些将得救者受到管教。因为你们为何以为，弟兄们，管教竟如此加于那民族，以致完全灭亡？这个民族中后来有多少人相信了，还有多少人将要相信？有些是糠秕，有些是麦子；然而两者都经受打谷机的碾压；但在同一打谷机下，一个被压碎，另一个被净化。天主藉卖主的犹大之恶赐给我们多大的善！藉犹太人的残忍，相信的外邦人获得了多大的善！基督被钉死，为的是让被蛇咬的人能仰望十字架上的那位……\par

7. 现在我们应快速掠过耶路撒冷毁灭之后的诗节，因为它们既明了，我也不愿在敌人的惩罚上停留。\BibleQuote{彷佛在森林中用斧头砍树，他们一举砍断了殿门；用铁锤和斧头将她推倒} \BibleRef{咏 73:6}。他们合谋，坚决地\Quote{用铁锤和斧头}将她推倒。\BibleQuote{他们用火焚烧祢的圣所，将祢圣名的居所亵渎于地} \BibleRef{咏 73:7}。\par

8. \Quote{他们在心里说（他们的亲族原为一体）}——他们说什么？\BibleQuote{来吧，我们把这主的庆节从地上消灭} \BibleRef{咏 73:8}。\Quote{主}一词是加在这人，即加在阿撒夫口中的。因为狂怒的他们不会称那他们正在拆毁其圣殿者为\Quote{主}。\Quote{来吧，我们把主的全部庆节从地上消灭。}阿撒夫怎样呢？在这些话中阿撒夫有什么领悟呢？他连所受的管教也不能得益吗？心灵之曲直不仍未恢复吗？起初的一切都被推翻了：犹太人的司铎不在了，祭台没有了，牺牲没有了，圣殿也没有了。那么，难道没有别的东西可以承认接替这离去之物吗？或者，除非那所应许的已经来到，这预兆性的记号难道会被除去吗？因此，让我们现在看阿撒夫在这处的领悟，看他是否从苦难中获益。注意他所说的话：\BibleQuote{我们的记号我们未见，不再有先知，而我们，他尚且还不认识} \BibleRef{咏 73:9}。看哪，那些犹太人说他们还不被认识，即他们还在充军，尚未被释放，仍在期待基督。基督会来，但祂要来作审判者：第一次来是召唤，以后是分别。祂会来，因为祂已经来过，而且祂会再来是显然的；但此后祂将从高处而来。祂在你面前，以色列啊。你跌倒了，因为你绊倒在那躺卧的石头上了：为了不被碾碎，你要留意祂从高处而来。因为先知早已预告：\Quote{凡跌在这石头上的，必被摔碎；这石头落在谁身上，就要把谁压得粉碎。}祂小的时候叫你摔碎，大的时候叫你碾碎。如今你的记号你不见，如今没有先知：而你说，\Quote{而我们，他尚且还不认识：}因为你们自己还不认识祂。\Quote{不再有先知；而我们，他尚且还不认识。}\par

9. \BibleQuote{天主，敌人辱骂要到几时？} \BibleRef{咏 73:10}。痛哭吧，如同被遗弃、被撇下似的；哭吧，如同一个病人，你宁愿击打医生而不愿痊愈：祂还不认识你。看祂所行的，这还不认识你的那一位。因为那些未曾有人向他们宣讲的，将要看见；那些未曾听见的，将要明白：而你还在哭喊：\Quote{不再有先知，而我们，他尚且还不认识。}你的领悟在哪里？\Quote{仇敌终如此羞辱你的名。}仇敌终如此羞辱你的名，是为了受辱后你责备，责备后你在终时认识他们；或者无疑地，\Quote{终如此}意即直到终了。\par

10. \BibleQuote{你为何撤回你的手，将你的右手从你怀中抽回直到终了？} \BibleRef{咏 73:11}。这又是给梅瑟的另一个记号。因为正如上面从棍杖得了记号，如今也从右手得了。当那件关于棍杖的事完成后，天主给了第二个记号：\BibleQuote{把你的手放在怀里，}祂说，\BibleQuote{他就放了；抽出来，他就抽出来；见手上有癞疾，} \BibleRef{出 4:6}，即不洁。因为皮上的白就是癞病 \BibleRef{肋 13:25}，并非肤色白皙。因为天主的基业本身，即祂的子民，被赶出去就成了不洁。但祂对他说什么？\Quote{再把手放回怀里。}他放回去，就恢复了原来的颜色。阿撒夫说，祢何时做这事呢？祢将祢的右手从祢怀中抽回多久，以致它不洁地留在外面呢？放回去吧，让它恢复颜色，让它承认救主。\Quote{你为何撤回你的手，将你的右手从你怀中抽回直到终了？}这些话他盲目地呼喊，不明白，而天主做祂所做的事。因为基督来为了什么？\Quote{部分的愚昧发生在以色列身上，直到外邦人的数目满了，然后全以色列都要得救。} \BibleRef{罗 11:25}。因此现在，阿撒夫啊，承认那先前的事，好叫你至少能够跟随，如果你不能走在前头。因为基督不是徒然来临，或徒然被杀，或徒然那麦粒落在地里；而是落在地里好结出许多子粒来。 \BibleRef{若 12:24} 旷野里高举了蛇，为医治那被蛇咬中的人免于中毒。 \BibleRef{户 21:9} 注意所发生的事。不要以为祂来临是虚妄的：免得祂第二次来临时，发现你成为恶者。\par

11. 阿撒夫明白了，因为圣咏标题上有\Quote{阿撒夫的领悟。}他说什么呢？\BibleQuote{但是，天主，我们的君王，在万世之前，已在世间施行了救恩} \BibleRef{咏 73:12}。一方面我们呼喊：\Quote{不再有先知，而我们，他尚且还不认识：}但另一方面，\Quote{我们的天主，我们的君王，在万世之前}（因为祂就是在起初的圣言 \BibleRef{若 1:1}，万世藉祂造成），\Quote{已在世间施行了救恩。}\Quote{因此，天主，我们的君王，在万世之前，}行了什么事？\Quote{已在世间施行了救恩：}而我还如同被遗弃者一样呼喊！……如今外邦人醒了，我们却打鼾，仿佛天主遗弃了我们，我们在梦中胡说。\Quote{祂已在世间施行了救恩。}\par

12. 所以现在，阿撒夫啊，按你的悟性修正你自己，告诉我们天主在大地中间施行了怎样的救恩。当你那属地的救恩被推翻时，祂做了什么，祂应许了什么？\BibleQuote{祢以祢的能力巩固了海洋}\BibleRef{咏 73:13}。就好像犹太民族是脱离波浪的旱地，外邦人在他们的苦毒中就是海洋，从四面冲刷着那地；看，\BibleQuote{祢以祢的能力巩固了海洋}，那地却仍然渴望祢的雨露。\BibleQuote{祢以祢的能力巩固了海洋，祢在水中击碎了龙的头。}龙的头，就是魔鬼的骄傲，外邦人被其所附，祢在水中击碎了它们；因为祢藉着圣洗释放了那些被它们所附的人。\par

13. 在龙的头之后还有什么？因为这些龙有自己的首领，他本身就是第一条大龙。关于他，那在人间施行救恩的天主做了什么？听：\BibleQuote{祢击碎了龙的头}\BibleRef{咏 73:14}。是哪条龙？我们藉着龙理解所有在魔鬼手下作战的恶魔：那么，我们应当理解哪一条独一的龙，它的头被击碎，不就是魔鬼本身吗？天主对他做了什么？\BibleQuote{祢击碎了龙的头。}那就是罪的开端。那个头是受到诅咒的部分，即厄娃的后裔要伤害蛇的头。\BibleRef{创 3:15}因为教会受到训诫，要躲避罪的开端。那罪的开端像蛇的头是什么？一切罪的开端是骄傲。\BibleRef{训 10:13}因此，龙的头已被击碎，魔鬼的骄傲已被击碎。天主对他做了什么，在人间施行救恩？\BibleQuote{祢把他给了厄提约丕雅人民作为食物。}这是什么？我如何理解厄提约丕雅人民？难道不是指所有外邦民族吗？而且特别指黑人：因为厄提约丕雅人是黑色的。他们本身就是被召叫到信德中的，曾为黑暗的人；的确如此，所以对他们说：\BibleQuote{因为你们从前是黑暗，但现在在主内是光明。}...因此，那个被人民崇拜的牛犊也是如此，他们不信、背道，寻求埃及人的神，离弃了那把他们从埃及奴役中拯救出来的天主；因此，那伟大的圣事得以举行。因为当梅瑟看到他们崇拜和朝拜偶像时，\BibleRef{出 32:19}被对天主的热情所激励，暂时惩罚他们，以便使他们畏惧而逃避永死；然而，他把牛犊的头丢进火里，磨成粉末，毁灭了，撒在水上，给人民喝：这样，一个伟大的圣事得以举行。哦，先知的愤怒，心灵没有扰动而是受到光照！他做了什么？把它丢进火里，以便首先消灭那形状本身；一点一点地磨碎，以便逐渐消耗：丢进水里，给人民喝！这岂不是说明魔鬼的崇拜者已成了它自己的身体吗？同样，人们宣认基督，成了基督的身体；因此对他们说：\BibleQuote{但你们是基督的身体和肢体。}\BibleRef{格前 12:27}魔鬼的身体必须被消耗，而且也要被以色列人消耗。因为宗徒们出自那个民族，最初的教会也出自那个民族。……因此，魔鬼随着它肢体的丧失而被消耗。这在梅瑟的蛇中也得到了预演。因为术士们也照样做了，扔下他们的棍子，变成了蛇；但梅瑟的蛇吞了所有那些术士的棍子。\BibleRef{出 7:12}所以，现在甚至要看出魔鬼的身体：这就是正在发生的事，它正被信了的外邦人吞食，它成了厄提约丕雅人民的食物。这一点再次可见于：\BibleQuote{祢把他给了厄提约丕雅人民作为食物}，如今所有的人都咬他。什么是咬他？就是责备、责怪、控告。正如经上所说的，虽然是禁止的方式，但想法表达出来了：\BibleQuote{但如果你们彼此咬噬、吞食，要小心，免得你们彼此消灭。}\BibleRef{迦 5:15}什么是彼此咬噬、吞食？你们彼此诉讼，彼此诋毁，彼此辱骂。因此，现在要留意，这些咬噬如何使魔鬼被消耗。有哪一个人，即使是一个异教徒，当他向仆人发怒时，不会对他说：撒殚？看，魔鬼成了食物。基督徒这样说，犹太人这样说，异教徒也这样说：他崇拜魔鬼，却用魔鬼来诅咒！\par

14. \Quote{祢劈开了泉源和洪流} \BibleRef{咏 73:15}：为使它们涌流智慧之河，涌流信德之宝藏，灌溉外邦人的咸地，藉着它们的灌溉使所有不信者归向信德的甘甜……在某些人身上，天主的圣言成为涌到永生的水泉 \BibleRef{若 4:14}；但另一些人听了圣言，却不这样持守好使他们生活美满，却又口舌不停，他们就成为洪流。因为那些不是长流的才恰当地称为洪流：有时洪流也用作河水的次义，正如经上所说：\Quote{祢以祢快乐的洪流给他们畅饮。} 因为那洪流永不干涸。但洪流严格来说是指那些夏季干涸、冬季因雨水泛滥而奔流的河流。因此你看到一个信德健全的人，他将坚持到底，在任何考验中不背弃天主；为了真理，不为虚谎和错误，忍受一切艰难。这人何以如此刚强？不是因为圣言在他内成为涌到永生的水泉吗？ \BibleRef{若 4:14} 但另一个人接受了圣言，他宣讲，他不沉默，他奔流：但夏季证明了他是泉源还是洪流。然而，二者都灌溉了大地，由那在大地中间施行了救恩者所灌溉：愿泉源涌溢，愿洪流奔涌。\par

15. \Quote{祢晒干了厄堂的河流} \BibleRef{咏 73:15}……什么是厄堂？因为这是希伯来文。厄堂是什么意思？坚固、强壮。这坚固强壮者是谁，天主晒干了他的河流？不就是那巨龙吗？因为\Quote{没有人能进入壮士的家，抢夺他的家具，除非先捆住那壮士。} \BibleRef{谷 12:29}这就是那壮士，仗恃自己的德行，离弃了天主；这就是那壮士，他说：\Quote{我要将我的宝座安置在北方，我要相似至高者。} \BibleRef{依 14:13}从这悖逆能力的杯中，他给人喝了。他们想要强壮，以为藉着禁果就能成为神。亚当成了强壮者，经上责备他说：\Quote{看，亚当成了我们中的一个。} \BibleRef{创 3:22}……好像他们强壮似的，\Quote{没有服从天主的正义。} \BibleRef{罗 10:3}你们看，一个人抛弃了自己的力量，变得软弱、贫乏、远远站着，甚至不敢举目望天，只捶着胸膛说：\Quote{主啊，可怜我这个罪人罢！} \BibleRef{路 18:13}现在他软弱了，现在他承认自己的软弱，他不强壮了：他是旱地，愿他得泉源和洪流的灌溉。那些仗恃自己德行的人至今仍是壮士。愿他们的河流干涸，愿外邦人、巫师、占星家、魔术的学说不再盛行：因为壮士的河流已被晒干：\Quote{祢晒干了厄堂的河流。} 愿那学说干涸；愿人的心智被真理的福音所淹没。\par

16. \Quote{白昼属于祢，黑夜也属于祢} \BibleRef{咏 73:16}。谁不知道这事呢？因为祂亲自创造了这一切；因为万物是藉着圣言造成的。 \BibleRef{若 1:3}正是对那在大地中间施行了救恩的同一者说：\Quote{黑夜属于祢。} 在此我们应领悟一些属于祂所施行之救恩的事。\Quote{白昼属于祢。} 这些人是谁？属神的人。\Quote{黑夜也属于祢。} 这些人是谁？属血气的人……\Quote{祢造了太阳和月亮：} 太阳，属神的人；月亮，属血气的人。他还属血气，愿他不被遗弃，也得以成全。太阳仿佛是智慧人，月亮仿佛是愚笨人：然而祢没有遗弃。因为经上这样记载：\Quote{智慧人如太阳恒久不变，愚笨人却如月亮变化无常。} \BibleRef{德 27:11}那么，因为太阳恒久不变，即智慧人如太阳恒久不变，愚笨人如月亮变化无常，那还属血气、还愚笨的人就该被遗弃吗？宗徒所说的话又在哪里：\Quote{对智慧人和愚笨人，我都欠了债}？ \BibleRef{罗 1:14}\par

17. \Quote{祢造成了大地的四极} \BibleRef{咏 73:17}……看，那在大地中间施行了救恩者是怎样造成了大地四极的。\Quote{祢造成了大地的四极。夏天和春天，祢造成了它们。} 在圣神内火热的人是夏天。我说，祢造了在圣神内火热的人；祢也造了在信德上作初学的人，他们是\Quote{春天。} \Quote{夏天和春天，祢都造成了。} 他们不得夸耀好像没有领受：\Quote{祢造成了他们。}\par

18. \BibleQuote{求祢记念这个受造物} \BibleRef{咏 73:18}。祢的什么受造物？\Quote{仇敌辱骂了主。} \Person{阿撒夫}，哀伤你昔日在理解上的盲目吧：\Quote{仇敌辱骂了主。} 曾有人在本国对基督说：\Quote{这人是罪人，我们不知道祂从哪里来；}我们知道\Person{梅瑟}，天主曾对他说话；这人是撒玛黎雅人。\Quote{愚昧的人民也亵渎了祢的名。} \Person{阿撒夫}那时正是愚昧的人民，但不是\Person{阿撒夫}那时的理解。在先前的圣咏中怎么说？\Quote{我如同畜牲在祢面前，但我常与祢同在；}因为祂没有去外邦人的神祇和偶像那里。虽然他像畜牲一样无知，却又像人一样认识。因为他说：\Quote{我常与祢同在，如同畜牲；}后来在同一篇圣咏中，\Person{阿撒夫}在哪里？\Quote{祢握住了我的右手，以祢的旨意引导了我，并以荣耀提携了我。} 以祢的旨意，而非我的义德；以祢的恩赐，而非我的功行。因此，这里也如此：\Quote{仇敌辱骂了主，愚昧的人民也亵渎了祢的名。} 那么他们都灭亡了吗？断乎不是……因为连宗徒\Person{保禄}也曾因不信而被折断，又因信德而被接回根源上。所以显然\Quote{愚昧的人民亵渎了祢的名}，那时他们曾说：\BibleQuote{如果祂是天主子，就让他从十字架上下来吧。} \BibleRef{玛 27:40}\par

19. 但你现在怎么说，\Person{阿撒夫}，有了理解？\BibleQuote{不要将向祢告明的灵魂交付给野兽} \BibleRef{咏 73:19}。……交付给哪些野兽？除了那些在水面上被击碎了头颅的野兽？因为同一魔鬼被称为野兽、狮子、龙。他说，不要将向祢告明的灵魂交付给魔鬼和他的天使。如果我还思念地上的事，如果我还渴求地上的事，如果在新约启示后我仍停留在旧约的应许中，就让蛇吞吃吧。但是因为如今我已放下了骄傲，不再承认自己的义德，只承认祢的恩宠；愿骄傲的野兽无权加害于我。\Quote{求祢不要永远忘记祢穷人的灵魂。} 我们曾富有，我们曾强壮；但祢已经使\Place{厄堂}的河流干涸：我们不再建立自己的义德，而是承认祢的恩宠；我们是穷人，求祢垂听祢的乞丐。如今我们不敢举目望天，只是捶着胸膛说：\BibleQuote{主啊，可怜我这个罪人吧。} \BibleRef{路 18:13}\par

20. \BibleQuote{求祢顾念祢的盟约} \BibleRef{咏 73:20}。实现祢所应许的：我们拥有约版，盼望着继承。\Quote{求祢顾念祢的盟约，} 不是那旧的：我不是为了\Place{迦南}地求，不是为了暂时的制伏仇敌求，不是为了儿女肉身的丰裕求，不是为了地上的财富求，不是为了暂时的福乐求：\Quote{求祢顾念祢的盟约，} 在其中祢应许了天国。如今我承认祢的盟约：如今\Person{阿撒夫}有了理解，\Person{阿撒夫}不再是野兽，如今他看见了那曾被论及的：\BibleQuote{看，日子将到，上主说，我必与以色列家和犹大家订立新约，不像我昔日与他们祖先所立的盟约。} \BibleRef{耶 31:31} \Quote{求祢顾念祢的盟约，因为被黑暗遮蔽的人已充满了不义之家的尘土：} 因为他们有不义的心。我们的\Quote{家}就是我们的心：心地纯洁的人喜欢住在其中。 \BibleRef{玛 5:8} 所以，\Quote{求祢顾念} \Quote{祢的盟约：} 让残余者得救： \BibleRef{罗 9:27} 因为许多注意尘土的人被黑暗遮蔽，充满了尘土。因为尘土进入了他们的眼，使他们瞎了，他们成了风从地面上吹散的尘土。\Quote{被黑暗遮蔽的人已充满了不义之家的尘土。} 因为由于注意尘土，他们被黑暗遮蔽，关于他们有另一篇圣咏说：\Quote{愿他们的眼被蒙蔽不得看见，愿祢时常使他们的腰弯曲。} 因此，\Quote{被黑暗遮蔽的人已充满了不义之家的尘土，} 因为他们有不义的心……\par

21. \BibleQuote{不要让谦卑的人羞愧地离开} \BibleRef{咏 73:21}。因为骄傲使他们羞愧。\Quote{穷乏无助的人必要赞美祢的名。} 你们看，弟兄们，贫穷是多么甘甜：你们看见穷乏无助的人属于天主，但\BibleQuote{神贫的人是有福的，因为天国是他们的。} \BibleRef{玛 5:3} 谁是神贫的人？是谦卑的人，在天主的话语前战栗，承认自己的罪，不依靠自己的功德，也不依靠自己的义德。谁是神贫的人？是那些行善时赞美天主，行恶时责备自己的人。\BibleQuote{我的灵要安息在谁身上？} 先知说，\BibleQuote{但在谦卑、温良、敬畏我言语的人身上。} \BibleRef{依 66:2} 因此如今\Person{阿撒夫}已经理解，如今他不再依附尘土，如今他不要求旧约中地上的应许……\par

22. \Quote{上主，求祢起来，为我判断我的案件}\BibleRef{咏 73:22}……因为我不能显示我的天主，好像我在追随一个空虚的东西，他们就辱骂我。不只是外邦人、犹太人或异端者，甚至有时一位公教弟兄，当宣讲天主的许诺、预言将来的复活时，也会做鬼脸。尽管如此，他虽已用永生救恩的水洗净，带着基督的圣事，也许会说：\Quote{有谁复活过？}以及\Quote{我未听过我父亲从坟墓里说话，自从我埋葬了他！}\Quote{天主赐给祂的仆人以暂时的法律，让他们遵循：有谁能从阴间回来？}我该如何对待这样的人？我能否向他们显示他们看不见的？我不能：因为天主不应为他们的缘故而成为可见的……你说：\Quote{我看不见，我该信什么？}那么，你的灵魂是可见的吗？愚昧的人！你的身体是可见的，你的灵魂谁看见了？既然只有你的身体可见，为什么不把你埋葬？他惊讶我说：如果只有身体可见，为什么不把你埋葬？他回答（因为他知道这一点）：因为我活着。我如何知道你是活着的，而我却看不见你的灵魂？我如何知道？你会回答：因为我说话，因为我行走，因为我工作。愚昧的人！通过身体的动作我知道你活着，通过受造物的工程你不能认识造物主吗？也许那个说\Quote{我死后便不复存在}的人，既识字，又从\Person{伊壁鸠鲁}那里学了这学说——他是一种昏聩的哲学家，更确切说是愚昧的爱好者而非智慧的，连哲学家们自己也称他为猪；他说\Quote{至善}是身体的快乐；他们称这位哲学家为猪，在肉体的泥沼中打滚。也许这位有学问的人从他那里学会了说：我死后便不复存在。愿\Place{厄堂}的河流干涸！愿外邦人的学说消亡，愿耶路撒冷的种植繁荣！让他们看见所能见的，在心里相信所不能见的！的确，现今普世所见的一切，当时天主在大地中心施行救恩时，那些事尚未发生，但那时已被预言，如今已应验；然而愚昧的人仍心中说：\Quote{没有天主。}祸哉，悖逆的心！因为剩余的事必会成就，正如那时尚未发生而被预言要发生的事。天主是否已经实现了祂所许诺的一切，唯独在审判之日欺骗了我们？基督尚未在地上，祂许诺了，祂成就了；没有童贞女怀孕，祂许诺了，祂成就了；宝贵的血尚未流出，以消除我们死亡的罪状，祂许诺了，祂成就了；肉体尚未复活得永生，祂许诺了，祂成就了；外邦人尚未相信，祂许诺了，祂成就了；异端者尚未以基督之名武装起来反对基督，祂预言了，祂成就了；外邦的偶像尚未从地上消除，祂预言了，祂成就了；当祂预言并成就了这一切时，唯独在审判之日撒谎了吗？它必定会来，正如这些事来到；因为这些事在发生之前也是未来的，先被预言为未来，后来才应验。它必会来临，我的弟兄们。没有人可以说：它不会来临；或者说：它会来临，但还很遥远。但对你而言，离开此世是近在咫尺的……如果你做了魔鬼所暗示的事，轻视了天主的命令，审判之日必会来到，你将发现天主所威胁的是真实的，魔鬼所许诺的是虚假的。\Quote{求祢记念祢所受的辱骂，就是愚昧之人终日对祢的辱骂。}因为基督仍然被辱骂；义怒之器终日不缺，直到时间的终结。仍然有人说：\Quote{基督徒宣讲的是虚空的事；}仍然有人说：\Quote{复活是愚拙的事。}\Quote{求祢记念祢所受的辱骂。}但哪些辱骂呢？岂不是\Quote{愚昧之人终日对祢的辱骂？}有智慧的人会这样说吗？不，因为智慧的人被称为有远见的人。如果智慧的人是有远见的人，他因信德而远见；因为肉眼几乎只能看见脚前的事。\par

23. \Quote{求祢不要忘记向祢呼求者的声音}\BibleRef{咏 73:23}。当他们叹息期盼祢从\WorkTitle{新约}所应许的，并凭信德行走时，\Quote{求祢不要忘记向祢呼求者的声音。}但那些人仍说：\Quote{祢的天主在哪里？愿恨祢者的骄傲时常上达于祢。}不要忘记他们的骄傲。祂不会忘记：无疑，祂或是惩罚，或是纠正。\par

\SourceRecord{Translated by J.E. Tweed.；Nicene and Post-Nicene Fathers, First Series；Vol. 8.；Edited by Philip Schaff.；Buffalo, NY: Christian Literature Publishing Co.,；1888.；<http://www.newadvent.org/fathers/1801074.htm>.}

% structure: p0001 source=h1 target=section reason=page-title

\section{\WorkTitle{圣咏第七十五篇释义}}

1. ……圣咏的标题如此说：\Quote{至终，不可败坏。}什么是\Quote{不可败坏？}即你所应许的，要实行。但何时？\Quote{至终。}所以，让心灵的目光朝向\Quote{终末。}让一切途中发生的事都过去，好使我们达到终末。让骄傲的人因现世的福乐而欢跃，让他们因尊荣而自大，在黄金中闪耀，家中满溢奴仆，被门客的侍从环绕：这些事都要过去，如影子般消逝。当那终末来临，当所有现今仰望主的人欢喜时，他们却要遭受无尽的忧伤。当温良的人领受骄傲者所讥笑之物时，骄傲者的夸耀要变为悲哀。那时将有我们在\WorkTitle{智慧篇}中所知的声音：因为当他们看见圣人的荣耀时，他们要这样说，这些圣人在卑微时忍耐了他们，在受高举时没有顺从——那时他们要说：\BibleQuote{这些人就是我们昔日所讥笑的。}\BibleRef{智 5:3}他们也要说：\BibleQuote{骄傲为我们带来了什么？财富的夸耀又给了我们什么？}万事如影子般消逝。因为他们依靠可朽坏之物，他们的望德必被败坏；但我们的望德在那时却是实体。为使天主的应许对我们保持完整、可靠、确定，我们以信德之心说了：\Quote{至终不可败坏。}所以不要害怕，免得任何有权势的人败坏天主的应许。祂不败坏，因为祂是信实的；没有比祂更强大者能使祂的应许被败坏：让我们因此对天主的应许确定无疑，现在就从圣咏开始之处歌唱。\par

2. \BibleQuote{我们要称谢祢，上主，我们要称谢祢，并呼求祢的圣名}\BibleRef{咏 74:1}。不要呼求，在你称谢之前：先称谢，再呼求。因为你所呼求的，是你呼到你自己那里的。呼求是什么，岂不就是呼到你自己那里吗？如果祂被你呼求，也就是说，如果你把祂呼到你那里，祂要到谁那里去呢？祂不会到骄傲的人那里。祂确是崇高的，自高的人达不到祂。为了达到崇高的事物，我们抬高自己，如果达不到，就寻找一些工具或梯子，以便抬高后达到崇高：相反，天主既崇高，却能被谦卑的人达到。经上记载：\Quote{主亲近那些心碎的人。}心碎就是虔敬、谦卑。那破碎自己的人，是对自己发怒。让他对自己发怒，好使祂对他仁慈；让他做自己的判官，好使祂成为他的辩护者。所以，天主确实会在被呼求时来到。祂来到谁那里？祂不会到骄傲的人那里。……因此，要留意你们所作的：因为如果祂知道，祂并非不留意。所以，最好祂不留意，而不被知道。因为那\Quote{不留意}是什么？岂不就是不知道？不知道是什么？不惩诫。因为正如惩诫的行为常被提及。这里有人祈求祂不留意：\Quote{求你转面不看我的罪过。}如果祂转面不看你了，你会怎样？那是可怕的事，应当害怕，免得祂抛弃你。反之，如果祂不转面不看，祂就惩诫。天主知道这件事，天主能做到这件事，即：既转面不看犯罪者，又不转面不看认罪者。……所以，称谢并呼求。因为藉着称谢，你洁净圣殿，当祂被呼求时，祂可以进入其中。称谢并呼求。愿祂转面不看你的罪过，但不转面不看你自己；转面不看你所做的事，但不转面不看祂自己所做的事。因为祂创造了你为人，而你的罪是你自己创造的。……\par

3. 但是，重复能够加强语气，这是我们从圣经许多章节中学到的。因此，主说：\BibleQuote{实实在在。}\BibleRef{若 1:51}同样，在某些圣咏中，有\Quote{但愿如此，但愿如此。}为表示事情，一次\Quote{但愿如此}就足够了；为表示确认，又加了一次\Quote{但愿如此}。……圣经中到处都有无数这样的例子。我们提这些，足以引起你们注意这种表达方式，以便你们在所有类似情形中观察：现在请注意实质内容：他说：\Quote{我们要称谢祢，}\Quote{并且我们要呼求。}我已经说过，为什么在呼求之前先要有称谢：因为你所呼求的，你邀请祂。但如果你自高自大，祂在被呼求时就不愿来：如果你自高，你就不能称谢。你不能向天主隐瞒任何祂不知道的事。因此，你的称谢并非教导祂，而是洁净你自己。\par

4. ……你们现在要听基督的话。因为这些话看来不像是祂的，却说：\Quote{我们必要向祢——天主——认罪，必要向祢认罪，且要呼求祢的名。}现在开始了以头的身份所作的讲论。但无论是头说话，或是肢体说话，都是基督在说话：祂以头的身份说话，祂以身体的身份说话。然而，经上说了什么？\Quote{二人要成为一体。}（\BibleRef{创 2:24}）\Quote{这是伟大的圣事：}他说，\Quote{我——是在基督和教会内说话。}（\BibleRef{弗 5:32}）祂自己也在福音中说：\Quote{因此，不再是两个，而是一体。}（\BibleRef{玛 19:6}）为使你们知道这在某种意义上乃是两个人，却又因婚姻的纽带而成为一体，祂在依撒意亚中如同一个人说话，说：\Quote{就如给新郎佩上冠冕，又如给新娘披戴装饰。}（\BibleRef{依 61:10}）祂称自己为头中的新郎，身体中的新娘。因此，祂如同一位在说话，让我们听祂，并在祂内也让我们说话。让我们成为祂的肢体，好使这声音也能成为我们的。他说：\Quote{我要传述祢一切奇妙的事。}基督正在宣讲祂自己，甚至在祂现今存在的肢体中宣讲祂自己，为引导他人归向祂，使那些原本不在的人得以亲近，并与祂的那些肢体联合，福音正是通过祂的这些肢体被宣讲；并且使在一个头之下、在一个圣神内、在一个生命中，成为一个身体。\par

5. 祂又说了什么？\Quote{当我到了约定的时期，}他说，\Quote{我要按公义审判。}（\BibleRef{咏 74:2}）祂何时要按公义审判？当祂到达约定的时期。现在还不是精确的时期。感谢祂的仁慈：祂先宣讲公义，然后才审判公义。因为如果祂愿意在宣讲之前审判，谁能被找到而得以解脱？谁能遇见祂而得以赦免？所以现在是宣讲的时期：\Quote{我要传述，}他说，\Quote{祢一切奇妙的作为。}你们要听祂传述，听祂宣讲：因为如果你们轻视祂，\Quote{当我到了约定的时期，}祂说，\Quote{我要按公义审判。}祂说：现在，我赦免一个告罪者的罪；将来，我不会宽恕一个轻视者……祂作为人子接受了时期；祂作为天主之子治理时期。请听祂作为人子如何接受了审判的时期。祂在福音中说：\Quote{祂给了祂审判的权柄，因为祂是人子。}（\BibleRef{若 5:27}）按照祂作为天主之子的本性，祂从未领受过审判的权柄，因为祂从不缺乏审判的权柄；按照祂作为人子的本性，祂接受了一个时期，就如出生、受苦、死亡、复活、升天，同样也来审判。在祂内，祂的身体也说这些话，因为祂不会没有他们而审判。因为祂在福音中说：\Quote{你们要坐在十二个宝座上，审判以色列十二支派。}（\BibleRef{玛 19:28}）因此，整个基督——即在圣徒中的头和身体——说：\Quote{当我到了约定的时期，我要按公义审判。}\par

6. 但现在呢？\Quote{大地已经流泻}（\BibleRef{咏 74:3}）。如果大地已经流泻，它因什么流泻？岂不是因罪恶？因此它们也被称为过犯。所谓犯罪，就好比因某种流动性而从稳固的美德与公义中滑落。因为每个人犯罪都是由于贪爱低级事物：正如他被对高级事物的爱所坚固，他也因对低级事物的欲望而堕落，仿佛融化一般。慈悲的赦罪者观察到世事的这种流逝，祂是罪恶的慈悲赦免者，尚不是刑罚的追讨者；祂观察着说：大地本身确实因住在她内的居民而流泻了。以下的话是解释，而非添加。就好像你在说：大地怎样流泻了？根基被抽走了吗？其中的什么东西被吞进了某种深渊？我所谓的大地是指所有住在她内的人。他说：我发现了大地有罪。我做了什么？\Quote{我坚固了它的支柱。}祂坚固的支柱是什么？祂称宗徒为支柱。所以保禄宗徒论到他的同宗徒们说：\Quote{那被视为柱石的。}（\BibleRef{迦 2:9}）那些支柱若不是被祂坚固，又会怎样呢？因为在一场地震中，就连这些支柱也动摇了：在主受难时，所有宗徒都绝望了。因此，那些在主受难时动摇的支柱，藉着复活被坚固了。建筑的起点通过其支柱呼喊，而在所有支柱中，建筑师自己呼喊。因为保禄宗徒是其中之一，当他说：\Quote{你们既然要寻求基督在我内说话的凭证——}（\BibleRef{格后 13:3}）因此，他说：\Quote{我坚固了它的支柱。}我复活了，我表明死亡不可怕，我向那些惧怕的人表明，连身体本身在死亡中也不会消亡。创伤曾使他们恐惧，伤痕却坚固了他们。主耶稣本可以没有任何伤痕地复活：因为对于那种能力，将身体恢复到如此完美的健康，以致过去的伤口不留任何痕迹，又算什么呢？祂有能力使之即使无伤痕也完好；但祂愿意拥有可以坚固那动摇的支柱的东西。\par

7. 弟兄们，现在我们已听到了那每日不隐藏的事；让我们现在来听听祂藉着这些柱石所呼喊的……祂呼喊什么？\BibleQuote{我曾对不义的人说：不要行不义} \BibleRef{咏 74:4} ……但他们已经行了，并是有罪的；大地已经流泄，一切住在地上的人也如此。那些钉死基督的人感到刺心，\BibleRef{宗 2:37} 他们承认了自己的罪，从宗徒那里学到了一些，以免对宣讲者的宽恕绝望。因为祂作为医生而来，因此没有来到健全的人那里。祂说：\BibleQuote{因为不需要医生，而是需要病人。我来不是召义人，而是召罪人悔改。} \BibleRef{玛 9:12} 因此，\BibleQuote{我曾对不义的人说：不要行不义。} 他们没有听从。因为从前对我们说过：我们不听，我们跌倒了，成了可朽的，生为可朽：大地流泄了。现在让他们甚至那位来到病人那里的医生，他们在健全时不愿听祂以免跌倒，让他们在躺下时听祂以便起来……\BibleQuote{我曾对不义的人说：不要行不义；对犯过者说：不要高举你们的角。} 如果你们的角不被高举，基督的角必在你们内被高举。你们的角是不义的角，基督的角是威严的角。\par

8. \BibleQuote{因此不要自高：不要对天主说不义的话} \BibleRef{咏 74:5} ……祂在另一篇圣咏中说什么？\BibleQuote{你做了这些事，} 列举了某些罪过。祂说：\BibleQuote{你做了这些事，我又沉默了。} 什么是\BibleQuote{我沉默了}？祂从不沉默于诫命，但暂时沉默于惩罚：祂停止复仇，不对被定罪者宣告判决。但这人这样说：我做了如此这般的事，天主没有报复；看，我安然无恙，没有灾祸临到我。\BibleQuote{你做了这些事，我又沉默了；你猜疑不义，以为我像你一样。} 什么是\BibleQuote{以为我像你一样}？因为你是不义的，你甚至以为我也是不义的；好像我是你恶行的赞同者，不是反对者，不是报复者。然后祂又对你说什么？\BibleQuote{我要责备你，将你摆在你自己的面前}？这是什么？因为现在你暗中犯罪，你摆出自己，看不见自己，不省察自己；我要将你摆在你自己面前，并从你自身带来惩罚。同样在这里，\BibleQuote{不要对天主说不义的话。} 注意。许多人说这不义的话，但不敢公开，免得作为亵渎者被虔诚的人厌恶；他们在心里啃噬这些事，在内里吞食这种不敬虔的食物；他们乐于出言反对天主，如果舌头未爆发，在心里也不沉默。因此另一篇圣咏说：\BibleQuote{愚妄人心里说：没有天主。} 愚妄人说了，但惧怕人：他不愿在人们可能听见的地方说；他在那个地方说，而祂（指天主）自己会听见他说的话。因此同样在这篇圣咏中（亲爱的弟兄们，请留意），当祂说\BibleQuote{不要对天主说不义的话}时，祂看见许多人在心里这样做，祂又加上：\BibleQuote{因为既不从东，也不从西，也不从山岭的沙漠 \BibleRef{咏 74:6} ，因为天主是审判者} \BibleRef{咏 74:7}。你们的不义，天主是审判者。如果祂是天主，处处都在。你要从哪里逃避天主的眼目，以至在某处说你祂听不到的话？如果天主从东方审判，你退到西方，任意反对天主；如果从西方，你就到东方去说；如果祂从山岭的沙漠审判，你就到人群中去，在那里喃喃自语。祂不从任何地方审判，因为祂处处隐秘，处处公开；没有人能够按祂所是认识祂，也没有人能够不认识祂。小心你所做的。你正对天主说不义的话。\BibleQuote{上主的神充满了世界}（另一圣经经文这样说），\BibleQuote{包含万物的认知声音；因此说坏话的人不能隐藏。} \BibleRef{智 1:7} 所以不要以为天主在地方里：祂与你同在，正如你将成的样子。什么是你将成的样子？你若作善，祂就是善的；你若作恶，祂在你眼中就是恶的；你若作善，祂就是帮助者；你若作恶，祂就是复仇者。你在暗中有了一位审判者。想要作恶时，你从公开场合退入家中，没有仇敌看见；从家中那些公开和在人眼前的地方，你躲进卧室；你甚至在卧室里也害怕别的见证者，你退入内心，在那里思考。祂比你的心更内在。所以无论你逃到哪里，祂就在那里。你要从哪里逃避你自己？无论你逃到哪里，你岂不跟着自己？但既然有一位比你更内在，你无法逃避发怒的天主，只能逃向和好的天主。根本没有你逃避的地方。你要逃避祂吗？逃向祂吧……那么我们现在该做什么？\Quote{让我们来到祂面前，} \OriginalTerm{在忏悔中}{\GreekText{ἐν} \GreekText{ἐξομολογήσει}}，在忏悔中来到祂面前：祂将温柔地来临，你曾使祂发怒。\BibleQuote{既不从山岭的沙漠，因为天主是审判者：} 不从东，不从西，不从山岭的沙漠。为什么？\BibleQuote{因为天主是审判者。} 如果祂在任何地方，祂就不是天主；但因为天主是审判者，不是人，所以不要指望祂出自地方。你将成为祂的地方，如果你为善，如果你在忏悔后呼求了祂。\par

9. \BibleQuote{一位祂贬抑，另一位祂高举} \BibleRef{咏 74:7}。审判者天主做了什么？请观察圣殿中的那两个人，你便看见祂贬抑谁，高举谁。\BibleQuote{他们上圣殿去祈祷，}祂说，一个是法利塞人，另一个是税吏。\BibleQuote{我实在告诉你们：那个税吏下去，成了义，胜过那个法利塞人；因为凡高举自己的，必被贬抑；凡贬抑自己的，必被高举。}如此便解释了本篇圣咏的一节。审判者天主做了什么？\BibleQuote{一位祂贬抑，另一位祂高举}：祂贬抑骄傲的人，高举谦卑的人。\par

10. \BibleQuote{因为上主手中的杯，盛满了纯酒，搀和了酒糟} \BibleRef{咏 74:8}。确实如此。\BibleQuote{祂将这杯倒出，给这个人；然而，酒糟仍未倒空；地上所有的罪人都要喝。}让我们稍作振奋；这里有些晦暗之处……首先我们遇到的问题是：\Quote{盛满了纯酒，搀和了酒糟}。既是\Quote{纯酒}，怎又会\Quote{搀和}呢？但祂说\Quote{上主手中的杯}（我是对基督教会内受过教育的人说话），你们切不可在心中把天主描绘成具有人形的有限存在，免得即使圣殿关闭，你们却在心中塑造了偶像。因此，这个杯确实象征某种事物。我们将会明白这一点。而\Quote{在上主的手中}，意即在天主的权能中。因为天主的\Quote{手}常指天主的权能。甚至论及世人时也常说：\Quote{他手中拥有它}，即他拥有它在他权能中，他愿意时便去做。\Quote{盛满了纯酒，搀和了酒糟}。接着祂自己解释说：\BibleQuote{祂倾斜了，祂说，从这杯给这个人}；然而，酒糟仍未倒空。看，它如何盛满了搀和的酒。因此，不要因它既是纯酒又是搀和而惊惶：纯酒因其纯正，搀和因其酒糟。那么，那里的酒是什么？酒糟又是什么？\BibleQuote{祂从这杯给这个人倾斜，}以致酒糟仍未倒空，这又是什么意思？\par

11. 请回忆祂从何处说到这一点：\BibleQuote{一位祂贬抑，另一位祂高举。}那在福音中通过两个人——法利塞人和税吏——所预表的，\BibleRef{路 18:10}让我们更广泛地理解为两个民族：犹太民族和异邦民族：犹太民族是那个法利塞人，异邦民族是那个税吏。……正如前者因骄傲而退后，后者因认罪而亲近。因此，上主手中盛满纯酒的杯，就主让我领悟而言，……那盛满搀和纯酒的杯，似乎是赐给犹太人的法律，以及所有被称为旧约的圣经；那里有天主各种话语的份量。因为新约隐藏在其中，仿佛隐藏在有形圣事的酒糟里。肉身的割礼是一件伟大的奥秘，且由此理解心灵的割礼。耶路撒冷圣殿是一件伟大的奥秘，且由此理解主的身体。应许之地被理解为天国。祭献牺牲和牲畜有伟大的奥秘：但在所有那些祭献中，人们理解的是十字架上那唯一的祭献和唯一的牺牲，主，祂代替了所有这些祭献，因为我们只有一个；因为即使那些预表了这些，也就是说，通过这些预表了那些。那个民族接受了法律，他们接受了公正而良善的诫命。\BibleRef{出 20:1}有什么比\BibleQuote{你不可杀人，你不可奸淫，你不可偷盗，你不可作假见证，孝敬父母，你不可贪恋邻人的财物，应敬拜唯一的天主，并只事奉祂}更公正呢？这些都属于酒。但那些属肉体的东西仿佛沉下去了，好让它们留在他们那里，而所有的属神理解则从那里倾注出来。但\Quote{上主手中的杯}，即在上主的权能中：\Quote{纯酒}，即纯粹的法律：\Quote{搀和了酒糟}，即连同肉体圣事的酒糟。因为祂贬抑那骄傲的犹太人，高举那忏悔的异邦人；\BibleQuote{祂从这杯给这个倾斜，}即从犹太民族倾斜给异邦民族。倾斜了什么？法律。从那里蒸馏出属灵的意义。\BibleQuote{然而，酒糟仍未倒空}，因为一切肉体的圣事都留给了犹太人。\BibleQuote{地上所有的罪人都要喝。}谁要喝？\BibleQuote{地上所有的罪人。}谁是地上的罪人？犹太人固然是罪人，却是骄傲的；同样，异邦人是罪人，却是谦卑的。所有的罪人都要喝，但要注意，谁喝的是酒糟，谁喝的是酒。因为那些喝酒糟的人归于虚无；那些喝酒的人却成了义。我甚至敢称他们为酩酊大醉的，且我不害怕：但愿你们全都如此陶醉。请回忆：\BibleQuote{你的杯何等令人陶醉，何其优雅！}但为什么？我的弟兄们，你们以为所有因承认基督而甘愿死去的人都清醒吗？他们如此沉醉，以致不认识自己的朋友。所有他们的亲属，用尘世的诱惑力图使他们脱离天上赏报的希望，都不被这些沉醉的人承认，也不被他们听见。他们岂非心被改变而沉醉？他们岂非心思远离了此世而沉醉？\BibleQuote{他们要喝，}祂说，\BibleQuote{地上所有的罪人。}但谁要喝酒？罪人要喝，但为使他们不再停留于罪中；为使他们得以成义，为使他们不致受罚。\par

12. \Quote{但我}——因为众人皆饮，但我是独特的，即基督与祂的身体—— \Quote{我要永远欢跃，向雅各伯的天主歌唱} \BibleRef{咏 74:9}，在那应许的终结之处，正如经上所说：\Quote{不可败坏。} \Quote{我要折断一切恶人的角，义人的角必被高举} \BibleRef{咏 74:10}。这就是说，祂贬抑一个，高举另一个。恶人不愿自己的角被折断，但末了它们必被折断。你那时不愿祂折断它们，不如今日就折断它们。因为你已在上文听见，切勿轻视：\Quote{我曾对妄为的人说：不要妄为，对罪犯说：不要高举你们的角。}当你听见\Quote{不要高举你们的角}时，你轻视了并高举了你的角：你终必来到那样的结局，那时将发生：\Quote{我要折断一切恶人的角，义人的角必被高举。}恶人的角是骄傲者的高位；义人的角是基督的恩赐。因为角被理解为高举。你在世上憎恨属世的高举，好使你得到天上的；你若爱属世的，祂就不接纳你到天上的；你那被折断的角将归于羞耻，正如那被高举的角将归于光荣。因此，现在是选择的时候，那时却没有。你不得说：让我离去，我再来选择。因为已经有话在前：\Quote{我曾对妄为的人说。}若我没有说，你就可预备借口，预备辩护；但若我已经说了，你就当首先抓住认罪，免得你陷入永罚；因为那时认罪已太迟，且没有辩护。\par

\SourceRecord{Translated by J.E. Tweed.；Nicene and Post-Nicene Fathers, First Series；Vol. 8.；Edited by Philip Schaff.；Buffalo, NY: Christian Literature Publishing Co.,；1888.；<http://www.newadvent.org/fathers/1801075.htm>.}

% structure: p0001 source=h1 target=section reason=page-title

\section{\WorkTitle{圣咏第七十六篇 释义}}

犹太人惯常在这篇我们咏唱的圣咏中自夸，说：\Quote{天主在犹达亚为人所知，在以色列他的名为大}，并且辱骂不认识天主的外邦人，说只有他们自己认识天主，因为先知说：\Quote{天主在犹达亚为人所知}。所以在其他地方，他是不被认识的。但天主确实在犹达亚为人所知，只要他们明白什么是犹达亚。因为确实，除非在犹达亚，天主是无法被认识的。看，连我们也这样说：除非一个人曾在犹达亚，认识天主于他是不可可能的。但宗徒说的是什么呢？\BibleQuote{那在暗中的犹太人，才是真犹太人，那出自心割损、而非文字割损的，才符合精神。}\BibleRef{罗 2:29}因此有属于肉身割损的犹太人，也有属于心割损的犹太人。我们许多圣祖既有属于肉身割损的，作为信德印记；也有属于心割损的，为了信德本身。这些人在名字上自夸，却失去了他们的行为，他们是从这些圣祖堕落的；我说，他们是从这些圣祖堕落的，因此在肉身上仍是犹太人，在内心却是外邦人。因为这些犹太人出自亚巴郎，从他生了依撒格，从他生了雅各伯，从雅各伯生了十二位圣祖，从十二位圣祖产生了全体犹太人民。但他们通常被称为犹太人，是因为犹大是雅各伯十二子之一，是十二位圣祖之一，从他这一脉，王权来到了犹太人中间。因为全体人民按雅各伯十二子的数目，有十二支派。我们所谓支派，就好比是分开的家族和人群的集合。我说，那人民有十二支派，在这十二支派中，一支是犹大，国王出自这支派；另一支是肋未，司铎出自这支派。但因为侍奉圣殿的司铎没有分得土地，\BibleRef{户 18:20}而必须让十二支派分享所有应许之地：因此，既然抽出了一支地位更高的支派，即司铎所在的肋未支派，本应剩下十一支派，除非通过若瑟二子的收养，使十二之数得以补全。\par

观察这是怎么回事。雅各伯的十二个儿子之一就是若瑟……。这若瑟有两个儿子：厄弗辣因和默纳协。雅各伯临终时，仿佛立遗嘱般，将这两个孙子收为儿子，并对儿子若瑟说：\BibleQuote{以后所生的将归于你；但这两个归于我，他们要同他们的兄弟分地。}\BibleRef{创 48:5} 那时应许之地还未被赐下或分定，但他是在圣神内说话，预言。因此，加上若瑟的两个儿子，仍然构成十二支派，因为现在有十三个。因为代替若瑟一个支派，加了两个，就成为十三个。然后除去肋未支派，就是侍奉圣殿、靠所有其他分得土地之人的什一税而生活的司铎支派，剩下十二个。在这十二支派中有犹大支派，君王由此而出。因为起初从另一支派赐下了撒乌耳王，\BibleRef{撒上 9:1} 但他因是恶王而被弃绝；之后从犹大支派赐下达味王，从他起，君王都出自犹大支派。\BibleRef{撒上 16:12} 雅各伯祝福他儿子们的时候曾预言这事说：\BibleQuote{权杖必不离犹大，领导杖不离他两腿之间，直到那应许者来到。}\BibleRef{创 49:10} 但从犹大支派出来了我们的主耶稣基督。因为祂如圣经所说，正如你们刚才听到的，出自达味的后裔，由玛利亚所生。\BibleRef{弟后 2:8} 但就我们的主耶稣基督的神性而言，祂与圣父同等，不仅比犹太人更先，也比亚巴郎本身更先；\BibleRef{若 8:58} 不仅比亚巴郎更先，也比亚当更先；不仅比亚当更先，也比天地和万世更先：因为万物是藉着祂造的，没有祂，什么也没有造。\BibleRef{若 1:3} 所以，因为预言中说：\BibleQuote{权杖必不离犹大等等}\BibleRef{创 49:10} 将以前的时代查考，我们发现犹太人总是有出自犹大支派的君王，并且在主诞生时做王的黑落德以前，他们没有外族君王。外族君王从那时期开始，从黑落德开始。\BibleRef{路 3:1} 在黑落德以前，所有人都是犹大支派的，但仅仅是直到那应许者来到。因此，当主亲自来到时，犹太人的王国被推翻，从犹太人那里夺去。现在他们没有君王，因为他们不承认真正的君王。现在你们看，他们是否还配称为犹太人。现在你们确实看到，他们不配称为犹太人。他们自己用他们的声音放弃了那名号，以致他们不配被称为犹太人，仅在肉身上算是。他们什么时候从这名号断绝？他们说：\BibleQuote{我们除了凯撒没有君王。}\BibleRef{若 19:15} 噢，你们这被称为犹太人而其实不是的，如果你们除了凯撒没有君王，那么犹大的权杖已经失落；那么那应许者已经来了。那么，那些由犹太人而成为基督徒的，才是更真实的犹太人；其余不信基督的犹太人，已经连这名字本身也失去了。所以，真正的犹大就是基督的教会，她相信那位出于犹大支派、经由童贞玛利亚而生的君王；相信那位宗徒刚才在写给弟茂德的书信中所说的：\BibleQuote{你要记得，耶稣基督从死者中复活了，出于达味的后裔，正如我所传的福音。}\BibleRef{弟后 2:8} 因为犹大有达味，出于达味的就是主耶稣基督。我们相信基督的，属于犹大；我们承认基督。我们这些没有看见却用信德持守祂的人。不要让那些不再是犹太人的犹太人辱骂。他们自己说：\BibleQuote{我们除了凯撒没有君王。}\BibleRef{若 19:15} 对他们来说，他们的君王是基督、达味的后裔、犹大支派，岂不更好？然而，因为基督自己按肉身是达味的后裔，但祂是在万有之上永远可赞美的天主，\BibleRef{罗 9:5} 祂就是我们的君王和我们的天主；祂是我们的君王，因为按肉身生于犹大支派，是主基督、救主；但祂是我们的天主，祂在犹大之前，在天地之前，万物是藉着祂造的，\BibleRef{若 1:3} 包括精神和物质的。因为如果万物是藉着祂造的，那么连祂由之而生的玛利亚本人，也是由祂所造……\par

2. \BibleQuote{在犹大地天主是已知的，在以色列祂的名为大。}\BibleRef{咏 75:1} 关于以色列，我们也应当像对待犹大那样理解：就如他们不是真正的犹太人，同样那也不是真正的以色列。因为以色列被说成是什么？是看见天主的人。他们如何看见了天主？祂在肉身中行走在他们中间，而他们以为祂是人，就把祂杀了……？\BibleQuote{在以色列祂的名为大。} 你愿意成为以色列吗？请看主所说的那个人：\BibleQuote{看，这确是一个以色列人，在他内毫无诡诈。}\BibleRef{若 1:47} 如果真正的以色列人是内心毫无诡诈的人，那么诡诈和说谎的人就不是真正的以色列人。所以让他们不要说，天主与他们同在，祂的名在以色列为大。让他们证明自己是以色列人，我就承认\BibleQuote{在以色列祂的名为大。}\par

3. \BibleQuote{在平安中为祂预备了居所，祂的住处就在熙雍} \BibleRef{咏75:2}。熙雍算是犹太人的地域；真正的熙雍是基督徒的教会。但是希伯来名称的解释如下：犹太解释为\Quote{宣认}；以色列解释为\Quote{看见天主者}。犹太之后是以色列。你愿意看见天主吗？首先你应当宣认，然后在你内为天主预备居所；因为\BibleQuote{在平安中为祂预备了居所}。因此，只要你还没有宣认自己的罪，你就在某种程度上与天主争吵。你怎么不与祂争辩呢？你赞美祂所不喜悦的。祂惩罚窃贼，你却赞美偷盗；祂惩罚醉汉，你却赞美醉酒。你与天主争辩，你没有在内心为祂预备地方：因为平安是祂的居所。你如何开始与天主和好？你从宣认开始与祂和好。有一篇圣咏的声音说：\BibleQuote{你们当以宣认向主开始}。\BibleQuote{你们当以宣认向主开始}是什么意思？就是开始与主结合。怎样做？使你所不喜悦的，也正是祂所不喜悦的。祂不喜悦你的恶行；如果这使你喜悦，你就与祂分离；如果它使你不悦，你就通过宣认与祂结合……\par

4. \BibleQuote{祂在那里折断了弓的力量、盾牌、刀剑和战争} \BibleRef{咏75:3}。祂在哪里折断的？在永久的平安里，在完美的平安里。现在，我的弟兄们，那些正确相信的人看到他们不应该依靠自己：他们自己一切威胁的力量，以及他们内心一切磨砺用于恶事的，祂都折断了；还有他们视为大德以暂时保护自己的一切，以及他们因护守自己的罪而与天主进行的战争，祂都在那里折断了。\par

5. \BibleQuote{祢从永久的山岭发出奇妙的光} \BibleRef{咏75:4}。何为永久的山岭？就是祂自己所使之永恒的那些；就是大山，真理的宣讲者。祢光照，却从永久的山岭发出：大山先领受祢的光，而从大山所领受的祢的光，大地也因此得着衣裳。但那座大山，宗徒们领受了；宗徒们领受了，如同初升之光的曙光……因此，在另一处，圣咏说什么？\BibleQuote{我举目望向群山，我的救助从那里来}。那么，你的盼望就在山岭，从那里救助临到你？你停留在山岭吗？留意你做什么。山岭之上还有一位：山岭之上是祂，山岭在祂面前颤抖。他说：\BibleQuote{我举目} \BibleQuote{望群山，我的救助从那里来}。但接着说什么？\BibleQuote{我的救助}，他说，\BibleQuote{来自上主，祂创造了天地}。我确实向群山举目，因为通过群山圣经向我展示；但我的心在祂那里，祂光照一切山岭……\par

6. \BibleQuote{所有心中愚妄的人都惊惶} \BibleRef{咏75:5}……他们怎会惊惶？当福音被宣讲时。何为永生？谁是从死者中复活的？当宗徒保禄谈论死者复活时，雅典人感到惊奇，以为他说的尽是神话。但因他说还有另一种生命，是眼睛未曾看见，耳朵未曾听见，人心也未曾想到的（\BibleRef{格前2:9}），所以心中愚妄的人都惊惶了。但他们遭遇了什么？\BibleQuote{他们睡了他们的觉，一切富人手中一无所获}。他们贪爱眼前的事物，在眼前的事物中沉睡；于是这些眼前的事物成了他们的快乐：如同一个人在梦中看见自己找到了宝藏，只要他不醒来，他就是富足的。梦使他富足，醒来使他贫穷。或许睡眠使他躺在地上，躺在硬地上，贫穷或是个乞丐；在梦中他看见自己躺在象牙或金床上，堆积着羽毛；只要他睡着，他就睡得舒服；醒来发现自己在硬地上，睡眠就是在此地抓住了他。这些也是如此：他们进入此生，通过暂时的欲望，仿佛在此打盹；财富和虚浮的荣华飞逝而去，抓住了他们，他们便消逝了：他们不明白能用这些财富行多少善。因为他们若知道有另一种生命，就会在那里为自己积蓄必朽的财富；如同税吏长匝凯，当他迎接主耶稣到自己家里时，看见了那善，他说：\BibleQuote{我把一半财产给穷人，我若欺骗了谁，我就四倍偿还} \BibleRef{路19:8}。这人不在做梦者的空虚中，而是在清醒者的信德中……\par

7. \BibleQuote{雅各伯的天主，你一斥责，所有骑马的人都沉睡了}\BibleRef{咏 75:6}。那些骑马的是谁？就是那些不愿谦卑的人。骑马本不是罪，但向天主扬起权力的颈项，自以为有什么了不起，这才是罪。因为你富有，你就骑上了马；天主斥责，你却沉睡。祂斥责的愤怒何其大，愤怒何其大。请你们的爱德留意这可怕的事。斥责有响声，响声常使人惊醒。天主斥责的力量如此之大，以至于祂说：\BibleQuote{雅各伯的天主，你一斥责，所有骑马的人都沉睡了}\BibleRef{咏 75:6}。看，那骑马的法老睡得何等沉。因为他心里没有苏醒，因为面对斥责，他硬了心\BibleRef{出 14:8}。因为心硬就是沉睡。我问你们，我的弟兄们，那些在福音、阿们、阿肋路亚响彻全世界的时候，仍然不肯谴责自己的旧生活、醒来进入新生活的人，他们睡得如何？天主的圣经从前只在犹太，如今却响彻全世界。在那一民族中，创造万物的唯一真神被传扬，当受钦崇和朝拜；如今哪里还有人不谈论祂？基督虽然曾在十字架上受讥笑，却复活了；那曾被祂受讥笑的十字架，如今祂印在了君王的额上：而人却仍在沉睡……\par

8. \BibleQuote{你是可畏的，你发怒的时候，谁能站立在你面前？}\BibleRef{咏 75:7} 如今他们沉睡，感觉不到你发怒；但祂发怒正是为了让他们沉睡。他们沉睡时所感觉不到的，到了最后必将感觉到。因为审判活人死人的审判者将要出现。\BibleQuote{你发怒的时候，谁能站立在你面前？} 因为如今他们随心所欲地说话，与天主争辩说：基督徒是谁？基督是谁？那些相信看不见之事、放弃看得见的快乐、跟随未见之事信德的人，是何等愚昧！你沉睡打鼾，尽你所能地说反对天主的话。\Quote{罪人哪，何时，主啊，罪人何时自夸？他们回答并说出不义的话？} 但何时没有人应答，没有人说话呢？除非当人转向反对自己的时候？……\par

9. \BibleQuote{你从天降下审判，大地战栗，并且安息}\BibleRef{咏 75:8}。那现在自扰的，那现在说话的，末后必恐惧而安息。她若现在安息，末后倒可以欢乐。安息？何时？\BibleQuote{当天主起来审判，要拯救所有心里谦卑的人}\BibleRef{咏 75:9}。谁是心里谦卑的人？就是那些没有骑上喷鼻息的马，却在谦卑中承认自己罪过的人。\BibleQuote{人的思念要向你承认，思念的残余要向你举行庆典}\BibleRef{咏 75:10}。前者是思念，后者是思念的残余。什么是初次的思念？那是我们开始的起点，那使你开始宣认的善念。宣认使我们与基督结合。但宣认本身，即初次的思念，在我们内产生思念的残余；而这些\Quote{思念的残余要向你举行庆典}。那宣认的思念是什么？就是谴责过去的生活，对以往的所是不再满意，以便成为所不是的，这本身是初次的思念。因为你应当这样从罪恶中退出：在向天主宣认之后，带着初次的思念，以免忘记你曾是个罪人；正因你曾是个罪人，你就向天主举行庆典。此外，也可以这样理解：初次的思念包含宣认和脱离旧生活。但如果你忘记了你从哪些罪恶中被解救出来，你就不会向拯救者感恩，也不会向你的天主举行庆典。看，宗徒扫禄——后来成为保禄，起初名叫扫禄——他初次的宣认思念，就是当他听到从天而来的声音时……他发出了顺服的初次思念：当他听见\BibleQuote{我就是你所迫害的纳匝肋人耶稣}\BibleRef{宗 9:5}，他说：\BibleQuote{主啊，你吩咐我做什么？}\BibleRef{宗 9:5} 这就是宣认的思念：如今他呼求他所迫害的主。至于思念的残余如何举行庆典，你在保禄身上已经听到，当时宗徒本人被诵读：\BibleQuote{你要记得，耶稣基督从死者中复活了，出自达味的后裔，照我所传的福音}\BibleRef{弟后 2:8}。什么是\Quote{你要记得}？即使你初次宣认的思念从记忆中被抹去，愿思念的残余留在记忆中……\par

10. 基督曾一次为我们祭献，在我们相信的时候；那时是思念。但现在有思念的残余，当我们记念那来到我们中间的那一位，以及祂宽恕了我们什么；借着这些思念的残余，即记忆本身，祂每日为我们祭献，仿佛每日更新我们，那曾藉祂初次恩宠更新我们的。因为如今主在圣洗圣事中更新了我们，我们成了新人，在望德中喜乐，好使我们在患难中忍耐\BibleRef{罗 12:12}；然而，那已赐予我们的，不应从我们记忆中消失。如果你现在的思念不再是当初的——因为初次的思念是离开罪恶：但如今你不再离开，而是那时离开了——愿思念的残余存在，免得那使我们痊愈者从记忆中消失……\par

11. \BibleQuote{你们许愿，当向主、我们的天主还愿。}\BibleRef{咏 75:11}。各人要量力许愿，并且还愿。不要许愿而不还；但要每人都许愿，尽其所能。不要延迟许愿，因为你将不是凭自己的力量完成所许的愿。你若依靠自己，就会失败；但若依靠你所许愿的那一位，你必将安全地还愿。\BibleQuote{你们许愿，当向主、我们的天主还愿。} 我们众人应当共同许什么愿呢？就是信靠他，从他那里希望获得永生，并按照一个对众人皆普遍的尺度虔诚生活。因为确有某种对一切人共同的尺度。不偷窃并\Emph{不仅仅}是对守贞者的要求，而是对已婚妇女也要求的；不犯奸淫是对所有人要求的；不贪杯（饮酒使灵魂沉溺，在自身内腐化天主的殿宇）是对所有人同样要求的；不骄傲是对所有人同样要求的；不杀人、不恨弟兄、不谋害任何人，是对所有人共同要求的。这一切我们都应当许愿。此外还有个人特有的许愿：有人向天主许愿忠贞于婚配，除了妻子之外不认识别的女人；同样，女人除了丈夫之外不认识别的男人。另有些人，即使已有了这样的婚姻，也许愿今后不再有类似的事，既不愿也不接受相同的事；这样的人所许的愿比前者更大。还有人从生命开始就许愿守童贞，甚至不愿经历那些已有经验而放弃的人所经历的事；这样的人所许的愿是最大的。有人许愿自己的家成为所有来到的圣人的接待之所；他们许了一个大愿。另有人许愿放弃所有财物，分施给穷人，进入一个团体，进入圣人的团体；他确实许了一个大愿。\BibleQuote{你们许愿，当向主、我们的天主还愿。} 让各人按自己所愿意的许愿；他要注意还他所许的愿。若有人对他向天主所许的愿观望回头，便是邪恶。有些守贞的妇女愿意结婚：她愿意什么呢？和任何一位处女一样。她愿意什么呢？和她的母亲一样。她愿意了什么坏事吗？当然是坏事。为什么？因为她已经向主、她的天主许了愿。因为宗徒保禄关于这样的人说了什么呢？尽管他说年轻的寡妇若愿意可以结婚（\BibleRef{弟前 5:14}），但他在另一处说：\BibleQuote{但她如此存留，按我的意见，将更有福。}（\BibleRef{格前 7:40}）他指出，若她如此存留，更为有福；但若她愿意结婚，也不应受谴责。但关于某些许愿而不还的人，他说了什么？\BibleQuote{他们被判罪，因为他们废弃了起初的信德。}（\BibleRef{弟前 5:12}）什么叫\BibleQuote{废弃了起初的信德}？就是他们许了愿，却没有还。因此，被安置在隐修院里的弟兄，不要说：我要离开隐修院；因为并不是只有隐修院里的人才能进入天国，也不是不在那里的人就不属于天主。我们回答他：但他们没有许愿；你许了愿，你却回头了。主在警告审判之日时，他说什么？\BibleQuote{你们要记得罗特的妻子。}（\BibleRef{路 17:32}）他对所有人说话。因为罗特的妻子做了什么？她从索多玛被救出，在路上她回头了。在她回头的地方，她就停在那里。因为她成了一尊盐柱（\BibleRef{创 19:26}），好让人因注视她而得调味，得智慧，不被迷惑，不回头，免得因立恶表而自己停留，却调味了别人。因为我们现在正对我们的某些弟兄说这话，我们也许看见他们在所立善志上似乎软弱。你也要像那人一样吗？我们把他们中有过回头的人摆在他们面前。他们自身失了味，却调味了别人，因为他们被提及，好使人因害怕他们的榜样而不回头。\Quote{你们许愿，当还愿。} 因为那个罗特的妻子是关乎所有人的。已婚妇女起了意要犯奸淫；她从她已达到的地方回头了。寡妇许愿守寡而愿意结婚，她愿意了对于已婚者合法的事，但对自身却不合法，因为她从她的地方回头了。有守贞的童女，已奉献给天主；但愿她也有其他真正修饰童贞本身的恩赐，没有这些恩赐，那童贞便是不洁的。因为如果她身体无玷而心灵腐化，那算什么呢？他说了什么？如果有人没有触碰她的身体，但她或许醉酒、骄傲、好争辩、多言多语？这一切天主都谴责。若在她许愿之前，她结了婚，她就不会被谴责；她选择了更好的事，超越了她合法的事；她却骄傲，行了许多违法的事。我说，在她许愿之前，结婚是合法的，但骄傲从不合法。啊，天主的童女，你愿意不结婚，这是合法的；你却高抬自己，这是不合法的。谦卑的童女胜过谦卑的已婚妇女，但谦卑的已婚妇女胜过骄傲的童女。但回头望结婚的女子被谴责，不是因为她愿意结婚，而是因为她已经前行，又因回头而成了罗特的妻子。不要迟延，你们这些有能力的人，被天主感动去抓住更高召叫的人；因为我们说这些话，不是要你们不许愿，而是要你们既许愿也要还愿。因为我们谈论了这些事，你或许本来愿意许愿，现在却不愿意许愿了。但看看圣咏对你说的话。他没有说\Quote{不要许愿}，而是说\Quote{许愿并还愿}。因为你听到了\Quote{还愿}，就不愿许愿了吗？难道你愿意许愿而不还愿吗？不，两者都做。一件事由你的宣认完成，另一件事将靠天主的助佑成全。仰望那引导你的，你就不会回头看你被领出的地方。那引导你的在你前面行走；他领你出来的地方在你后面。爱他引导你，他必不谴责你回头。\par

12. \Quote{凡在祂周围的人都必献上礼物。}谁是在祂周围呢？……凡是众人共有的，都在中间。为什么说是在中间？因为它对所有人距离相同，对所有人近在咫尺。不在中间的，就好比是私有的。公共的放在中间，是为了使所有来的人都可使用它，都得光照。愿无人说：这是我的；免得他缺乏对那放在中间为众人的一份作出自己的贡献。那么，\Quote{凡在祂周围的人都必献上礼物}，这话是什么意思？所有明白真理是众人共有、且不因骄傲而将其据为己有的人，他们必献上礼物，因为他们有谦卑；但那些把公共的如同放在中间的据为己有的人，试图引诱人结党，这些人必不献上礼物……\Quote{向那可畏的主。}因此，凡在祂周围的人都当恐惧。为此他们应当恐惧，且以战栗赞美；因为他们是在祂周围，好使众人能达到祂，祂也公开迎接众人，公开光照众人。这就是与众人一同敬畏。当你好像把祂当作自己的，不再视为公共时，你就被高举到骄傲；尽管经上记着：\Quote{你们当以恐惧事奉主，以战栗向祂欢跃。}因此，在祂周围的人都必献上礼物。因为知道真理是众人共有的，他们是谦卑的。\par

13. 他们向谁献上礼物呢？\Quote{向那可畏的主，向那夺去王侯之灵的主} \BibleRef{咏 75:12}。因为王侯之灵是骄傲的灵。他们便不是祂的灵；因为他们若知道什么，就想要据为己有，而非公共；但那位向众人显示自己平等、将自己置于中间，好使众人尽可能、尽所能地获取，不是取自任何人，而是取自天主，也因此成为他们自己的，因为他们已属于祂。所以他们必须谦卑：他们已失去了自己的灵，而有了天主的灵……因为如果你承认自己是尘土，天主从尘土中造人。凡在祂周围的人都献上礼物。所有谦卑的人都向祂忏悔，朝拜祂。\Quote{他们向那可畏的主献上礼物。}为此，你们当以战栗向那可畏的主欢跃：\Quote{向那夺去王侯之灵的主：}即夺去骄傲人傲慢的那位。\Quote{向那在世上君王中可畏的主。}世上的君王是可畏的，但祂超越一切，祂使世上的君王畏惧。你若作世上的君王，天主对你就将是可畏的。你或许会说：我怎能成为世上的君王？管辖大地，你便是世上的君王。因此，不要因贪恋权势而放眼极广大的省份，以扩展你的王国；你要管辖你所承载的大地。请听那管辖大地的宗徒：\Quote{我这样打仗，不像打空气；我痛击我的身体，使它成为奴隶，免得我给别人传报了，自己反而落选。} \BibleRef{格前 9:26}……\par

\SourceRecord{Translated by J.E. Tweed.；Nicene and Post-Nicene Fathers, First Series；Vol. 8.；Edited by Philip Schaff.；Buffalo, NY: Christian Literature Publishing Co.,；1888.；<http://www.newadvent.org/fathers/1801076.htm>.}

% structure: p0001 source=h1 target=section reason=page-title

\section{\WorkTitle{圣咏第七十七篇释义}}

1. 这篇圣咏的标题如此写道：\Quote{交与乐官，依耶杜通调，阿撒夫的诗歌。}\newline{}所谓\Quote{交与乐官}，你们知道是何意。\OriginalTerm{耶杜通}{Idithun}被解释为\Quote{越过那些人}，\OriginalTerm{阿撒夫}{Asaph}被解释为\Quote{集会}。因此，这里发言的是\Quote{一个越过}的\Quote{集会}，为能到达\OriginalTerm{终向}{End}，即耶稣基督。……\par

2. \Quote{我向主呼求}，他说\BibleRef{咏 76:1}。但许多人向主呼求是为了获得财富、避免损失、求得朋友的安全、家宅的安稳、现世的福乐、世俗的尊荣，甚至仅仅为了身体的健康——这是穷人的遗产。为了诸如此类的事，许多人向主呼求；但几乎没有人为了主本身而呼求。因为一个人很容易向主祈求什么，而不是祈求主自己；仿佛祂所赐的比祂自己更甘饴。因此，凡是为了任何其他事向主呼求的人，还不是那\OriginalTerm{越过}{leaping over}的人……祂确实垂听你，当你寻求祂本身的时候，而不是当你通过祂寻求任何其他事物的时候。关于有些人，经上记载：\Quote{他们呼求，却无人拯救；向主呼求，祂却不垂听。}为什么呢？因为他们的声音不是向主的。圣经在另一处说明了这样的人，说：\Quote{他们并未呼求主。}他们不停地向祂呼求，却并未呼求主。什么是\Quote{并未呼求主}？他们没有呼唤主到自己身边：他们没有邀请主进入自己的心，他们不愿被主居住。因此，他们遭遇了什么？\Quote{他们无故战栗恐惧。}他们为失去眼前之物而战栗，因为他们没有充满他们所不呼求的那一位。他们没有白白地爱，以至于失去现世之物后能说：\Quote{上主所喜悦的，就这样成就了，愿上主的名受赞美。}\BibleRef{约 1:21}因此，这人说：\Quote{我的声音向主，祂必垂听我。}让他告诉我们这如何成就。\par

3. \Quote{我在患难之日寻求天主}\BibleRef{咏 76:2}。你是谁，做这样的事？在你患难之日，留意你所寻求的。如果监狱是患难的缘由，你寻求从监狱出来；如果发烧是患难的缘由，你寻求健康；如果饥饿是患难的缘由，你寻求饱足；如果损失是患难的缘由，你寻求增益；如果流放是患难的缘由，你寻求肉身的家园。我何必一一列举，何时能列举一切？你愿意做一个\Quote{越过}的人吗？在你患难之日，寻求天主：不是通过天主寻求别的，而是从患难中寻求天主，使天主除去患难，使你无忧无虑地依附于天主。\Quote{我在患难之日寻求天主：}没有别的，而是\Quote{我寻求了天主。}你是如何寻求的？\Quote{我在夜间在祂面前举手。}……\par

4. 不可认为磨难仅指某些特定事物。因为每一个尚未跃越之人，都认为除非是此生遭遇某种悲哀之事，否则便不是磨难；但那跃越之人，却将整个此生视为自己的磨难。他如此热爱天上的乡国，以至地上的旅程本身就成了最大的磨难。因为此生若非磨难，还能是什么呢？我请问你们，整个被称为试探的生命，又怎能不是磨难呢？你们在约伯传上读到：\BibleRef{约 7:1} \BibleQuote{人在世上岂不是一个试探吗？}他说，人的生命在地上被试探吗？不，生命本身就是试探。因此如果是试探，它必然就是磨难。在这个磨难中，也就是说，在这生命中，那个跃越之人寻求了天主。如何？\Quote{用我的手}，他说。\Quote{用我的手}是什么呢？就是用我的作为。因为他寻求的并非任何有形之物，为能找到并触摸所失之物，用双手寻求金钱、金子、银子、衣服，总而言之，一切可以用手握住之物。然而，我们的主耶稣基督自己也曾愿意被以手寻求，当祂向怀疑的门徒显示钉痕时。\BibleRef{若 20:27} ……那么，用手寻求难道不也属于我们吗？正如我所说，用作为寻求是属我们的。何时呢？\Quote{在夜间}。\Quote{在夜间}是什么呢？就是此世。因为这是黑夜，直到在我们的主耶稣基督荣耀的来临中白昼显现。你们愿意看这是如何的黑夜吗？若我们在此没有一盏灯，我们必留在黑暗中。因为伯多禄说：\BibleQuote{我们有更确实的先知之言，你们应留意它，如同在暗处照明的灯，直到白昼破晓，晨星在你们心中升起。}\BibleRef{伯后 1:19} 因此在这黑夜之后必有白昼到来，与此同时，在这黑夜中并不缺乏灯。也许这正是我们现在所做的：通过解释这些段落，我们正带来一盏灯，为能在黑夜中喜乐。这灯本该常在你们的家中燃烧。因为有人曾说：\BibleQuote{不要熄灭圣神。}\BibleRef{得前 5:19} 他接着又说，仿佛在解释自己所说的：\Quote{不要轻视先知之恩。}也就是说，让这灯在你们内常燃。甚至这光与某种不可言说的白昼相比，也被称为黑夜。因为信徒的生命与外教人的生命相比是白昼，但与天使相比却是黑夜。因为天使拥有我们尚未拥有的白昼。我们已经拥有外教人所没有的白昼，但信徒尚未拥有天使所拥有的白昼；但他们将在复活时，当他们与天使相等时，获得所应许的。\BibleRef{玛 22:30} 因此，在这现时既是白昼又是黑夜中——与我们所渴慕的未来白昼相比是黑夜，与我们所弃绝的过去黑夜相比是白昼——我就在这黑夜中，让我们用手寻求天主。不要让作为停止，让我们寻求天主，不要有闲散的渴望。若我们在路上，就让我们善用资源，以便能到达终点。让我们用手寻求天主……\Quote{我用双手在夜间在祂面前，我没有受骗。}\par

5. ……\BibleQuote{我的灵魂拒绝受安慰}\BibleRef{咏 76:2}。巨大的疲惫占据了我，以致我的灵魂向一切安慰关闭了门。这疲惫从何而来？也许他的葡萄园被冰雹打了，或他的橄榄树不结果，或收成被雨水打断。这疲惫从何而来？从另一篇圣咏中聆听。因为其中有同样声音：\Quote{疲惫压垮了我，因为罪人离弃了你的法律。}于是他说，他被如此巨大的疲惫征服，因为这种恶事；以至于他的灵魂拒绝受安慰。疲惫几乎吞噬了他，忧伤完全将他淹没，无法医治，他拒绝受安慰。那么还剩下什么？首先，看他从哪里得到安慰。他难道没有等待一个同情他的人吗？……\BibleQuote{我记念了天主，就喜乐}\BibleRef{咏 76:3}。我的双手并非徒劳，他们找到了伟大的安慰者。在不懈怠时，\BibleQuote{我记念了天主，就喜乐。}因此，天主应受赞美，此人记念了祂，就喜乐了，在忧伤中得了安慰，在似乎绝望时得了振奋：因此，天主应受赞美。最后，因为他得了安慰，接着他说：\Quote{我饶舌了。}在这同一安慰中，因记念天主，我喜乐了，并\Quote{饶舌了。}\Quote{我饶舌了}是什么意思？就是我说着喜乐的话，欢欣鼓舞。因为那些喜极而语、不能也不愿沉默的人，正是通常所说的饶舌者。此人成了这样的人。他又说什么？\Quote{我的灵衰微了。}\par

6. 由于疲倦，他憔悴不堪；因思念天主，他曾欣喜；但又因妄语而昏厥。下文是什么？\BibleQuote{我的所有仇敌都抢先守夜} \BibleRef{咏 76:4}。我的所有仇敌都看守着我；他们加倍看守我；他们已在我之先守了夜。他们在哪里不设陷阱？难道我的仇敌不是已抢先于一切守夜吗？因为这些仇敌是谁？不就是宗徒所说的那班：\BibleQuote{你们不是对抗血肉} \BibleRef{弗 6:12}……与魔鬼及其天使交战，我们正进行敌对。他称他们为世界的统治者，因为他们统治爱世界的人。他们并不统治世界，如同天地的主宰；而是他称罪人为世界。……与魔鬼及其天使无和好可言。他们自己嫉妒我们进天国。他们对我们绝不能被安抚：因为\BibleQuote{我的所有仇敌都抢先守夜}。他们比我更警醒以欺骗我，比我更守护自己。他们怎能不抢先守夜？他们在各处安放绊脚石，各处设下陷阱。疲倦侵袭心灵，我们须惧怕忧愁吞噬我们；在喜乐中要害怕心神因妄语而昏厥：\BibleQuote{我的所有仇敌都抢先守夜}。最终，就在那同一妄语中，当你说话，且无所畏惧地说话时，有多少次敌人抓住并指责，甚至构陷诽谤——\Quote{他这样说，他这样想，他这样讲！}人该做什么，除非下文？\Quote{我惊惶不安，遂沉默无言。}因此，当他惊惶时，免得在妄语中抢先守夜的敌人寻找并发现诽谤，他便沉默不言。……\par

7. \BibleQuote{我想起古昔的日子} \BibleRef{咏 76:5}。如今他，仿佛一个被赶出门外的人，躲进了避难所：他在自己内心的秘密处交谈。且让他告诉我们他在那里做什么。他安好。请观察他在想什么。他是在里面，在自己家中思想古昔的日子。没有人对他说：你说话不当；没有人对他说：你话说多了；没有人对他说：你思想乖僻。愿他安好，愿天主助佑他：让他思想古昔的日子，并告诉我们他在那内室中做了什么，他到达了什么，越过了什么，停留在了何处。\BibleQuote{我想起古昔的日子，也思念了永恒的年岁。}什么是永恒的年岁？这是一个伟大的思想。你看这思想除了深沉的寂静外还需要什么？远离一切外在的喧哗，一切人间纷扰，让他内在安静，以思想那些永恒年岁。我们所处的年岁是永恒的吗？或者我们祖先所处的？或者我们后裔将要处的？断不可认为它们是永恒的。因为这些年岁中有什么留存？我们说话并说：\Quote{在这年中：}而我们从这年得到了什么，除了我们所在的这一天？因为今年的已往日子已经过去，不再拥有；但未来的日子尚未到来。我们只在这一天，却说\Quote{在今年}；不如你说\Quote{今天}，如果你希望说任何现在的事物。因为整年中你有什么现在的东西？凡是过去的不再存在；凡是未来的尚未存在：那么怎么是这个\Quote{今年}？改正说法：说\Quote{今天}。你说的是真话，从今以后我要说\Quote{今天}。再看这个，今天早晨的时辰已过，未来的时辰尚未到来。因此再改正：说\Quote{在这一时辰}。而在这时辰中你得到了什么？它的一些瞬间已过去，那些未来的尚未到来。说\Quote{在这一瞬间}。在哪个瞬间？当我发出音节时，如果我说两个音节，后者在前者过去后才发音；简言之，同一音节中，如果有两个字母，后者在前字母过去后才发音。那么我们从这些年岁中得到了什么？这些岁月是变化的：必须思想永恒的年岁，那些站立的年岁，不是由来来去去的日子构成的；在别处圣经论到天主说：\BibleQuote{但是你永存，你的年岁无穷尽。}这些岁月，这个人跨越过去，不是在外妄语，而是在寂静中思想了。\par

8. \BibleQuote{我夜间用心默想} \BibleRef{咏 76:6}。没有诽谤者在他言语中找陷阱，他在心中默想。\Quote{我妄语。}看，这就是先前的妄语。要当心，免得你的心神昏厥。他说，我并非像在外那样妄语：现在是另一种方式。怎么样？\Quote{我妄语，并探察我的心神。}如果他探察大地寻找金矿脉，没有人会说他是愚笨的；相反，许多人会称他聪明，因他渴望得到金子：人内心蕴藏着何等财富，而他却不去挖掘！这人正在检查他的心神，并与那同一心神说话，在说话中他正在妄语。他自问，自省，审判自己。他继续说：\Quote{我探察了我的心神。}他必须害怕停留在自己的心神内：因为他曾在外面妄语；并且因为他的所有仇敌都抢先守夜，他在那里发现了忧伤，心神便昏厥了。那在外面妄语的人，看啊，现在开始在安全中内在妄语，在那里独自秘密，思想永恒的年岁……\par

9. 而你找到了什么？\Quote{天主必不永远拒绝} \BibleRef{咏 76:7}。他在此生中找到了厌倦；在任何地方都找不到可靠、无所畏惧的安慰。无论他投靠什么人，都在他们身上找到丑闻，或惧怕它。因此，他没有一处能免于忧虑。保持沉默对他来说是坏事，恐怕他因好言而沉默；在外面说话和喋喋不休令他痛苦，恐怕他所有的敌人，预先警戒，在他话语中寻找诽谤。他在此生中极其困窘，便多想到另一生命，那里没有这样的考验。而他要何时到达那里？因为我们受苦在此处显然是天主的愤怒。这便是在依撒意亚中所说的：\Quote{我必不永远作你们的报应者，也不永久对你们发怒。} \BibleRef{依 57:16}……这愤怒要永远存留吗？此人在沉默中没有找到这个。因为他说什么？\Quote{天主必不永远拒绝，祂必不再加添，使祂仍然喜悦。}也就是说，使祂仍然喜悦拒绝，祂不再永远拒绝。祂必须召回祂的仆人，祂必须接纳回归于主的逃亡者，祂必须垂听被束缚之人的声音。\Quote{或是祂将永远断绝仁慈，从这一代到那一代？} \BibleRef{咏 76:8}。\par

10. \Quote{或是天主将忘记仁慈？} \BibleRef{咏 76:9}。在你内，从你到他人，除非天主赐给你，否则就没有仁慈：难道天主自己会忘记仁慈吗？溪水奔流，源头会干涸吗？\Quote{或是天主将忘记仁慈，或是在愤怒中收住祂的仁慈？}也就是说，祂会如此愤怒，以至于不再施恩？祂更容易收住愤怒而非仁慈。\par

11. \Quote{于是我说。}如今他越过自己说了什么？\Quote{如今我已开始：} \BibleRef{咏 76:10} ，当我甚至走出自己时。从此以后，不再有危险：因为即使停留在我自己内，也是危险。\Quote{于是我说，如今我已开始：这是至高者右手的转变。}\par

12. \Quote{我记念了主的作为} \BibleRef{咏 76:11}。如今请看他在主的作为中漫游。因为他先前在外面喋喋不休，并因此忧伤，他的精神消沉：他在内心与自己的心、与自己的精神喋喋不休，并且当查考那同一精神时，他记起了永恒的年岁，记起了主的仁慈，即天主必不永远拒绝；于是他开始无所畏惧地因祂的作为喜乐，无所畏惧地在其中欢腾。让我们现在也倾听那些作为，让我们也欢腾。但让我们也在情感中跃升，不因短暂之事而喜乐。因为我们也有我们的床。为何我们不进入其中？为何我们不守在沉默中？为何我们不查究我们的精神？为何我们不思想永恒的年岁？为何我们不在天主的作为中喜乐？让我们现在就如此倾听，并以祂亲自说话为乐，好使我们离开这里后，能做我们在祂说话时所惯常做的事；只要我们正在做祂所提到的开始：\Quote{如今我已开始。}以天主的作为为乐，就是甚至忘记你自己，如果你能唯独以祂为乐。因为有什么比祂更好的事呢？难道你没有看见，当你转向你自己时，你转向了更糟的事？\Quote{因为我必从起初记念祢奇妙的作为。}\par

13. \Quote{我要默想祢的一切作为，并述说祢的情感} \BibleRef{咏 76:12}。看，第三次的言语！他在外言语，当他暗示时；他在内心精神中言语，当他前进时；他在天主的作为上言语，当他到达自己所前进的地方时。\Quote{并述说祢的情感：}不是任何情感。哪一个人没有情感而活着？弟兄们，难道你们以为，敬畏天主、崇拜天主、爱天主的人，没有情感吗？你们确实会以为并敢以为，绘画、剧场、打猎、驯鹰、钓鱼会牵动情感，而默想天主不会牵动某种内在的情感吗？当我们凝视宇宙，把自然界的景象置于眼前，并在其中努力发现造物主，发现祂无处不令人喜悦，但超乎万有之上时？\par

14. \Quote{天主，祢的道路在圣者中} \BibleRef{咏 76:13}。他现在默想天主在我们周围的仁慈作为，他在这些作为中言语，并在这些情感中欢腾。他首先从那里开始：\Quote{祢的道路在圣者中？}祢那在圣者中的道路是什么？祂说：\Quote{我就是道路、真理和生命。} \BibleRef{若 14:6} 所以，人哪，从你们的情感中回转吧……\Quote{有哪一位大神像我们的天主？}外邦人有关于他们神明的情感，他们崇拜偶像，它们有眼却看不见；有耳却听不见；有脚却不能行走。你为什么向一位不行走的神行走？他说，我不崇拜这样的东西，那你崇拜什么？那里面的神性。那么你崇拜那东西，关于它曾在别处说过：\Quote{因为外邦人的神都是魔鬼。}你要么崇拜偶像，要么崇拜魔鬼。他说，既不是偶像，也不是魔鬼。那你崇拜什么？星宿、太阳、月亮，那些天上的东西。那么，造了地上和天上万物的那位，要好得多。\par

15. \BibleQuote{祢是独自施行奇迹的天主} \BibleRef{咏 76:14}。祢的确是伟大的天主，在肉身、在灵魂中独自施行奇迹。聋子听见，瞎子看见，软弱者康复，死人复活，瘫痪者得坚振。但这些奇迹那时是在身体上施行的，让我们看看那些在灵魂上施行的。那先前醉酒的人清醒了，那先前崇拜偶像的人成为了信徒：他们把自己的财物施舍给穷人，而先前他们抢夺别人的财物……\BibleQuote{独自施行奇迹}。梅瑟也行了，但不是独自；厄里亚也行了，甚至厄里叟也行了，宗徒们也行了，但没有一人是独自的。为使有能力行这些，祢与他们同在：当祢行这些时，他们并不与祢同在。因为他们并不与祢同在，正如祢甚至造了这些人本身。\BibleQuote{独自}是什么意思？难道是父，而不是子？或子，而不是父？不，而是父、子和圣神。因为不是三位神，而是一位天主独自施行奇迹，甚至在这位\OriginalTerm{踊越者}{leaper-over}本身也是如此。因为他的踊越以及达到这些事，是天主的奇迹：当他与自己的心神喃喃自语时，为了使他能踊越那同一天心神，并喜乐于天主的工程，他自己也行了奇迹。但天主行了什么？\BibleQuote{祢在万民中显明了祢的德能。}由此这阿撒夫的会众踊越；因为祂在万民中显明了祂的\OriginalTerm{德能}{virtue}。祂在万民中显明了祂的什么德能？\BibleQuote{我们却宣讲被钉的基督，……基督是天主的德能，天主的智慧。}\BibleRef{格前 1:23}。如果天主的德能就是基督，那么祂就在万民中显明了基督。我们难道还没有感知到这一点吗？我们如此愚昧，如此低伏，如此毫无踊越，以致看不到这一点吗？\par

16. \BibleQuote{祢以祢的手臂救赎了祢的百姓} \BibleRef{咏 76:15}。\BibleQuote{以祢的手臂}，\OriginalTerm{手臂}{arm}，即祢的德能。\BibleQuote{主的手臂向谁显示了呢？}\BibleRef{依 53:1}。\BibleQuote{祢以祢的手臂救赎了祢的百姓，以色列和若瑟的子孙。}怎么如同两个民族，\Quote{以色列和若瑟的子孙}？难道若瑟的子孙不在以色列的子孙中吗？……祂提醒我们要作出某种区分。让我们省察自己的心神，或许天主在那里放置了什么——我们甚至应在夜间用手寻求的天主，以免受骗——我们或许会在这\Quote{以色列和若瑟的子孙}的区分中发现我们自己。藉着若瑟，祂愿意另一个民族被理解，愿意外邦人的民族被理解。为什么藉着若瑟来理解外邦人的民族？因为若瑟被他的弟兄们卖到了埃及。\BibleRef{创 37:28}。那个若瑟被兄长嫉妒，被卖到埃及，在埃及劳苦、卑微；被显扬后，兴盛、统治。这一切预示了什么？无非是基督被祂的弟兄们出卖，被逐出本乡，如同进入外邦人的埃及？在那里起初卑微，当殉道者忍受迫害时；如今显扬，如我们所看到的；因为在祂身上应验了：\BibleQuote{大地的一切种类都要朝拜祂，万民都要事奉祂。}因此若瑟代表外邦人的民族，而以色列代表希伯来民族。天主救赎了祂的百姓，\BibleQuote{以色列和若瑟的子孙}。藉着什么？藉着\OriginalTerm{角石}{corner stone}，\BibleRef{弗 2:20}在其中两堵墙被连接在一起。\par

17. 他继续如何？\BibleQuote{大水一见祢，天主，就惊惶，深渊也震荡} \BibleRef{咏 76:16}。\OriginalTerm{众水}{waters}是什么？是万民。这众水曾在《默示录》中被问及，\BibleRef{默 17:15}回答是：万民。在那里我们清楚地看到水是万民的象征。但上文曾说：\BibleQuote{祢在万民中显明了祢的德能。}因此有理地，\BibleQuote{大水一见祢，就惊惶。}它们因惊惶而改变。\OriginalTerm{深渊}{abysses}是什么？是众水的深处。当\OriginalTerm{良心}{conscience}被击打时，万民中哪一个人不震荡？你寻求海的深处，有什么比人的良心更深？那就是当深渊震荡时的深处，那时天主以祂的手臂救赎了祂的百姓。深渊是如何震荡的？当众人在宣认中倾泻他们的良心时。\par

18. 在赞美天主时，在告明罪过时，在圣诗与圣歌中，在祈祷中，\BibleQuote{有许多大水的声音。云彩发出了声音} \BibleRef{咏 76:17}。那水的声音，那深渊的翻腾，就因为\BibleQuote{云彩发出了声音}。什么云彩？就是宣讲真理之言的云彩。什么云彩？就是天主威胁某座葡萄园——那园子不结葡萄却长出荆棘——时所说的：\BibleQuote{我要命令我的云彩，不再降雨在其上} \BibleRef{依 5:6}。总之，宗徒们离开了犹太人，前往外邦人那里；当他们在万民中宣讲基督时，\BibleQuote{云彩发出了声音}。\BibleQuote{因为你的箭已经射透。}云彩的这些声音，他又称为箭。因为福音作者的话就是箭。这些都是比喻。严格来说，箭不是雨，雨也不是箭；但天主的圣言既是箭，因为它击打；又是雨，因为它浇灌。所以，当\BibleQuote{你的箭已经射透}时，没有人再对深渊的翻腾感到惊奇。\BibleQuote{你的箭已经射透}是什么意思呢？它们没有停留在耳中，而是刺透了心。\BibleQuote{你雷声的响声是在轮中} \BibleRef{咏 76:18}。这是什么意思？我们该如何理解？愿主帮助。小时候，每当听到天上的雷声，我们常想象好像从车房里有马车出来。雷声确实发出像马车滚动的声音。我们是否要回到这种童稚的念头，去理解\BibleQuote{你雷声的响声是在轮中}，好像天主在云中有某些马车，马车的行驶发出那声音？断乎不是。这是幼稚、虚妄、无意义的。那么，\BibleQuote{你雷声的响声是在轮中}是什么意思呢？你的声音在滚动。这我仍不理解。我们该怎么办？让我们问问\Person{耶杜通}本人，或许他能解释自己所说的：\BibleQuote{你雷声的响声是在轮中。}我不明白。我要听你说：\BibleQuote{你的电光已经照亮世界。}那么说吧，我本不明白。世界是一个轮子。因为圆形世界的轨道也合乎情理地被称为\OriginalTerm{圆轮}{orb}；因此一个小轮子叫做\OriginalTerm{小圆轮}{orbiculus}。\BibleQuote{你雷声的响声是在轮中：}\BibleQuote{你的电光已经照亮世界}。那些云彩在轮中周游世界，带着雷声和闪电周游，震动了深渊，以诫命打雷，以奇迹闪电。\Quote{他们的声音传遍大地，他们的言语直达地极。} \Quote{大地震动战栗：}即所有住在大地上的人。但以比喻来说，大地本身是海。为什么？因为万民被称为海，因为人类生活是苦涩的，且暴露于风暴和暴风雨中。此外，你若观察这事，人们如何像鱼一样彼此吞噬，强者如何吞灭弱者——那便是海，福音作者们往那里去了。\par

19. \BibleQuote{你的道路是在海中} \BibleRef{咏 76:19}。但现在你的道路是在圣者内，如今\BibleQuote{你的道路是在海中}，因为圣者自己就在海中，他甚至有理由行走在海的水面上 \BibleRef{玛 14:25}。\BibleQuote{你的道路是在海中}，即你的基督在外邦人中被宣讲。\BibleQuote{你的道路是在海中，你的路径在多水之处}，即在多民之中。\BibleQuote{你的脚踪将无人知道。}他触摸了一些人，若不是那些犹太人，那才奇怪。看，基督的怜悯如今如此广泛地传向外邦人，以致\BibleQuote{你的道路是在海中。你的脚踪将无人知道。}怎么无人知道呢？除了那些仍说基督尚未降临的人？他们为什么说基督尚未降临？因为他们尚未认出他行走在海面上。\par

20. \BibleQuote{你曾带领你的百姓如同羊群，藉梅瑟和亚郎的手} \BibleRef{咏 76:20}。为什么添加了这一句，有点难以发现。……他们驱逐了基督；他们患病，不愿他作他们的救主；但他开始在外邦人中，在万民中，在多民中。然而，那民族的一支余民得救了。忘恩负义的群众留在外面，就是\Person{雅各伯}大腿的瘸踌 \BibleRef{创 32:32}。大腿的宽度被理解为后裔的众多，而在以色列子民中，大部分变得空虚愚昧，以致不认识基督在水面上的脚步。\BibleQuote{你曾带领你的百姓如同羊群，}他们却不认识你。虽然你赐予他们如此大的恩惠，分开大海，使他们从水间干地经过，在波涛中淹没了追赶的敌人，在旷野中为他们的饥饿降下玛纳，藉\BibleQuote{梅瑟和亚郎的手}带领他们回家：他们仍将你推开，以致在海中有了你的道路，而你的脚踪他们却不认识。\par

\SourceRecord{Translated by J.E. Tweed.；Nicene and Post-Nicene Fathers, First Series；Vol. 8.；Edited by Philip Schaff.；Buffalo, NY: Christian Literature Publishing Co.,；1888.；<http://www.newadvent.org/fathers/1801077.htm>.}

% structure: p0001 source=h1 target=section reason=page-title

\section{圣咏第七十八篇释义}

1. 这篇圣咏确实包含那些据说是发生在古老民族中的事迹；但新近的后代民族受到劝诫，要警惕不可对天主的恩惠忘恩负义，从而激起祂的愤怒反对他们，而他们本应领受祂的恩宠。……其标题首先引起并抓住我们的注意。因为它并非无故题作\Quote{\OriginalTerm{阿撒夫}{Asaph}的领悟:}；也许是因为这些话需要一位读者，他不只感知表面发出的声音，而是某种内在的意义。其次，当要叙述并提到这一切似乎需要听者多于解释者的事情时，他说：\BibleQuote{我要开口说比喻，我要从起初陈明命题。}谁不会在此从沉睡中醒来？谁敢匆匆读过这些比喻和命题，仿佛它们不言自明，而它们的名字本身就表明应当以更深的眼光去探寻？因为比喻表面上具有某种事物的相似：尽管它是希腊词，但现在已被用作拉丁词。值得注意的是，在比喻中，那些被称为事物相似的东西，与我们密切相关的事物进行比较。而命题，希腊文称为\OriginalTerm{命题}{\GreekText{προβήλματα}}，是其中包含某些须通过辩论来解决的问题。那么，谁会草率地阅读比喻和命题？谁在听到这些话语时不以警觉的心留意，为能通过理解获得其中的果实？\par

2. 祂说：\Quote{\BibleQuote{我的百姓，请侧耳听我的法律}}\BibleRef{咏 77:1}。我们可能假设在此说话的是谁，除了天主？因为正是祂将祂的子民从埃及拯救出来并聚集在一起，这聚集恰当地称为\OriginalTerm{会堂}{Synagogue}，\Person{阿撒夫}一词被解释为指代这个。那么，是否\Quote{\OriginalTerm{阿撒夫}{Asaph}的领悟}是说阿撒夫本人领悟了，还是必须比喻地理解为同一位会堂，即同一子民领悟了，他们对之说道：\Quote{\Quote{我的百姓，请侧耳听我的法律}}？那么，为什么祂通过先知的口责备同一子民说：\Quote{\BibleQuote{但以色列不认识我，我的百姓不明白}}\BibleRef{依 1:3}？然而，事实上，在那子民中甚至也有明白的人，他们拥有后来启示的信德，这与法律的文字无关，而是与圣神的恩宠有关。因为他们不可能没有同样的信德，他们能够预见并预言这信德在基督内的启示，因为甚至那些旧礼也是未来之事的预表。难道只有先知们有这信德，而百姓没有？不，确实，就连那些忠信地听从先知的人，也蒙受同样的恩宠以明白他们所听到的。但毫无疑问，天国的奥秘在旧约中是隐藏的，在时期满全时应在新约中揭开。因为宗徒说：\Quote{\BibleQuote{他们都饮了那随行的神磐石，而那磐石就是基督}}\BibleRef{格前 10:4}。因此，在奥秘中，他们的食粮和饮料与我们的相同，但在预表上相同，而不是在形式上；因为同样的基督自身在磐石中为他们所预表，在肉身中向我们显现。但他却说：\Quote{\BibleQuote{但天主并非欢喜了所有的人}}\BibleRef{格前 10:5}。诚然，所有的人都吃了同样的神粮，喝了同样的神饮，也就是说，预表某种属神的事物；但并非在所有人中天主都喜悦。当他说\Quote{\Quote{并非在所有}}时，显然其中有一些是天主所喜悦的；虽然所有的圣事是共同的，恩宠——圣事的效力——并非对所有人都相同。正如在我们的时代，现在信仰已经启示，那时是隐藏的，对一切因父、及子、及圣神之名受洗的人，\BibleRef{玛 28:19}重生的水泉是共同的；但正是那恩宠——这些水泉的圣事所指向的，藉此基督身体的肢体要与他们的元首一同为王——并非对所有人都相同。因为甚至异端者也有同样的圣洗，并且假弟兄也在公教名义的共融之中。\par

3. 然而，不论那时或现在，那说\Quote{我的子民哪，你们要听从我的法律}的声音并非没有益处。这一表达在所有圣经中都值得注意，他说的是\Quote{你们要听从}，而不是\Quote{你要听从}。因为一个子民是由许多人组成的；对这许多人而言，随后的内容以复数形式说出。\Quote{倾耳于我口中的言语}。\Quote{你们要听从}与\Quote{倾耳}相同；他在那里说\Quote{我的法律}，在这里说\Quote{我口中的言语}。因为人若虔诚地听从天主的法律和祂口中的言语，必是谦卑倾耳的人，而非骄傲昂首的人。因为凡倾倒之物，必被谦卑的凹面所接纳，而被傲慢的凸面所弹开。因此，在另一处他说：\Quote{倾耳，接受明智的言语}。\BibleRef{箴 5:1}所以我们已受到充分劝诫，要以倾耳，即谦卑的虔诚，接受这篇属于\Person{阿撒夫}这种理解的圣咏。而且，它并非被说成是\Person{阿撒夫}自己的，而是对\Person{阿撒夫}说的。这一点从希腊语冠词可以明显看出，并且在某些拉丁语抄本中也有发现。因此，这些言语属于理解，即智能，它被赐予了\Person{阿撒夫}本人；我们最好不要将其理解为针对一个人，而是针对天主子民的集会；因此我们决不应脱离这集会。因为虽然我们通常将犹太人的集会称为\Quote{会堂}，而将基督徒的集会称为\Quote{教会}，因为\Quote{集合}通常被理解为更适用于兽类，而\Quote{召集}则更适用于人类；然而我们也发现那也被称为教会，并且也许更适合我们说：\Quote{主，我们的天主，求祢拯救我们，从异民中召集我们，使我们能称谢祢的圣名}。我们也不应不屑于成为，不，我们应表达说不尽的感恩，因为我们是祂手中的羊，祂在说\Quote{我还有别的羊，不属于这一羊栈，我也必须领它们来，好使只有一个羊群和一位牧者}（\BibleRef{若 10:16}）时已预见了我们，也就是说，通过把忠信的外邦子民与忠信的以色列子民联合起来，关于他们祂曾说过：\Quote{我受派遣，无非是为了以色列家迷失的羊}（\BibleRef{玛 15:24}）。因为所有民族都要聚集在祂面前，祂要把他们分开，好像牧人把绵羊和山羊分开一样（\BibleRef{玛 25:32}）。因此，让我们这样来听那被说的话：\Quote{我的子民哪，你们要听从我的法律，倾耳于我口中的言语}。不是像对犹太人说的，而是像对我们自己说的，或者至少好像这些话也同样对我们（如同对他们）说了。因为当宗徒说\Quote{但他们并不是所有的人都蒙天主喜悦}时，由此表明也有蒙天主喜悦的人；他立刻加上\Quote{因为他们在旷野中被消灭了}；接着他又说\Quote{这些事成了我们的预像}。……因此，这些言语更是为我们而唱的。因此在这篇圣咏中，与其他事物一起，也说了\Quote{好让后代的子民知道，将要出生的儿子们，他们要兴起}。此外，如果那被蛇咬死、被毁灭者毁灭、被剑杀死的事都是预像，正如宗徒明确宣告的那样，因为所有那些事都确实发生了：因为他不是说\Quote{以预像的方式说话}或\Quote{以预像的方式写作}，而是说\Quote{这些事以预像的方式发生在他们身上}；那么，我们岂不应以更大的虔诚勤奋躲避那些预像所象征的惩罚吗？因为毫无疑问，正如在好事中，预像所象征的善远胜于预像本身；同样，在坏事中，预像所象征的恶也远为更甚，而作为预像的恶事本身已是如此之大。因为正如应许之地，那子民被引往之地，与天国相比微不足道，基督徒子民正被引往天国；同样，那些作为预像的惩罚，尽管如此严厉，与它们所象征的惩罚相比也微不足道。但是，宗徒称为预像的，这篇圣咏，据我们判断，称为比喻和议论：它们的目的并非在于它们曾经发生的事实，而在于通过合理的比较所指涉的那些事。因此，让我们—祂的子民—听从天主的法律，并倾耳于祂口中的言语。\par

4. 他说：\Quote{我要开口说比喻，我要说出古时的谜语}。\BibleRef{咏 77:2}他所说的\Quote{从起初}是什么意思，从后续的话中看得很清楚。因为那起初不是天地被造之时，也不是人类在第一个人中受造之时，而是那被领出埃及的集会之时；好使这意义属于\Person{阿撒夫}，\Person{阿撒夫}被解释为\Quote{集会}。但愿那说了\Quote{我要开口说比喻}的那位，也赐给我们通达的理解力来理解它们！因为如果祂开口说比喻，祂也照样揭开比喻本身；如果祂说出\Quote{议论}，祂也照样说出它们的解释，我们就不必在此劳苦了；但现在所有的事都如此隐晦封闭，以至于即使我们能靠祂的帮助触及任何可资健康滋养的东西，我们仍必须汗流满面地吃饼，并且不只以身体的劳动，也以心灵的劳动来偿付古老判决的惩罚（\BibleRef{创 3:19}）。那么，让他发言吧，让我们听这些比喻和议论。\par

5. \BibleQuote{我们听到何等大事，也知道了它们，我们的父亲们曾告诉我们}\BibleRef{咏 77:3}。主在上文说话。因为这些话还能被认为是出自其他什么人呢？\BibleQuote{我的子民哪，你们要侧耳听我的法律}？那么为什么现在突然有一个人在说话？因为我们在这里有一个人说的话：\BibleQuote{我们的父亲们曾告诉我们}。毫无疑问，天主如今要藉着一个人的职务说话，正如宗徒所说：\BibleQuote{你们愿意证实那在我内说话的基督吗？}\BibleRef{格后 13:3}。祂起初愿意亲自说出这些话，免得一个说祂话的人因是凡人而被轻视。因为天主的那些通过我们身体感官到达我们的话语也是如此。造物主以一种不可见的工作感动受造物，不是使本体变成任何有形体、暂时之物；当祂藉着有形和暂时的标记，无论是属于眼或耳的，按人所能领受的程度，使祂的旨意被人知晓时。因为如果一个天使能使用空气、雾气、云彩、火和任何其他自然物质或有形样式；而人能使用面容、舌头、手、笔、文字或任何其他表意符号，以暗示自己心灵的隐密：总之，如果他是一个人，他派遣人的使者，对一个人说：\BibleQuote{你去，他就去；对另一个人说，你来，他就来；对他的仆人说，你做这事，他就做这事；}\BibleRef{路 7:8}。那么天主，作为主，万物都服从祂，该会用多么更大、更有力的权能使用同一个天使和人，以宣告祂所喜悦的一切呢？……因为那些在旧约中所听到的，是在新约中所知道的：听到时它们正在被预言，知道时它们正在被应验。哪里应许得以实现，听觉就不会被欺骗。\BibleQuote{而我们的父亲们，}\Person{梅瑟}和先知们，\BibleQuote{曾告诉我们。}\par

6. \Quote{他们没有向另一世代的后代隐藏}。\BibleRef{咏 77:4}。这是我们的世代，在这一世代中我们获得了重生。\Quote{传扬主的赞美和祂的权能，以及祂所行的奇妙化工。}这段话的次序是：\Quote{我们的祖先也告诉了我们，传扬主的赞美。}主受到赞美，为的是祂可以被爱。因为有什么事物比祂更有利于我们的救恩而值得被爱呢？\Quote{祂在雅各伯中立了见证，在雅各伯中设了法律} \BibleRef{咏 77:5}。这就是上面所说的开端：\Quote{我要从起初宣布比喻。}所以，开端是旧约，终结是新约。因为恐惧在法律中盛行。\Quote{但是法律的终向是基督，使一切相信的人都成义} \BibleRef{罗 10:4}由于祂的恩赐，\Quote{爱已倾注在我们心中，藉着所赐予我们的圣神} \BibleRef{罗 5:5}而圆满的爱驱除恐惧，\BibleRef{若一 4:18}因为如今没有法律，天主的正义已显示出来。但由于祂藉法律和先知所获得的见证，\BibleRef{罗 3:21}因此，\Quote{祂在雅各伯中立了见证。}因为即使那个用如此奇妙和充满奥义的工作所搭建的会幕，也被称为\Quote{见证的会幕}，其中有遮盖约柜的幔子，如同遮盖法律服役者面容的帕子；因为在那个时代有\Quote{比喻和箴言。}因为那些正在宣讲并即将成就的事，隐藏在隐晦的意义中，并未在显露的显现中被看见。但宗徒说：\Quote{你们几时转向基督，帕子就会除去。} \BibleRef{格后 3:16}因为\Quote{天主的恩许在祂内都是\InnerQuote{是}，并藉着祂说\InnerQuote{阿们}。} \BibleRef{格后 1:20}所以，凡紧贴基督的，就拥有法律字句中所未能觉察的全部美善；但凡远离基督的，既不能觉察，也不拥有。\Quote{祂在以色列中设了法律。}按照他惯常的方式，他在这里重复。因为\Quote{祂立了见证}与\Quote{祂设了法律}相同，而\Quote{在雅各伯}与\Quote{在以色列}相同。正如这是同一个人有两个名字，法律和见证也是同一事物的两个名称。有人问：\Quote{立了}与\Quote{设了}之间有区别吗？的确，正如\Quote{雅各伯}与\Quote{以色列}之间的区别：不是因为他们是两个人，而是这两个名字因不同的理由赐给了同一个人；雅各伯是因为抓夺，因为他出生时抓住了他兄弟的脚跟，\BibleRef{创 25:26}但以色列是因为见了天主，\BibleRef{创 32:28}所以\Quote{立了}是一回事，\Quote{设了}是另一回事。因为，据我所判断，\Quote{祂立了见证}之所以这样说，是因为藉着它某事物被举扬了；宗徒说：\Quote{没有法律，罪是死的；但我从前没有法律时，曾活着；但诫命一来，罪便活了。} \BibleRef{罗 7:8}请看，那被见证（即法律）所举扬的东西，好叫隐藏的事显现出来，正如他稍后所说：\Quote{但罪为了显示它是罪，藉着善事给我产生了死亡。} \BibleRef{罗 7:13}但\Quote{祂设了法律}这句话，仿佛是加在罪人身上的轭，因此又说了：\Quote{因为法律不是为义人设立的。} \BibleRef{弟前 1:9}所以，就其证明某事而言，它是见证；就其命令某事而言，它是法律；虽然它是同一事物。因此，正如基督是一块石头，对于相信的人是屋角的基石，对于不信的人则是绊脚石和使人跌倒的磐石；同样，法律为那些不按合法方式使用法律的人（\BibleRef{弟前 1:8}），是见证，用以定罪罪人，使其当受惩罚；但对于那些正当使用法律的人，是见证，指出罪人应当投奔谁以得解救……\par

7. \Quote{何等伟大的事，}他说，\Quote{祂曾命令我们的祖先，使他们传给他们的儿子？}\BibleRef{咏 77:5}。 \Quote{为使另一世代知晓，将出生的儿子们，他们要起来，并传给他们的儿子}\BibleRef{咏 77:6}。 \Quote{为使他们寄望于天主，不忘记天主的作为，并寻求祂的诫命}\BibleRef{咏 77:7}。 \Quote{为使他们不致像他们的祖先，成为弯曲和苦毒的世代：这世代未曾引导自己的心，其精神也未曾信赖天主}\BibleRef{咏 77:8}。这些话语确是指出两个民族，一个属于旧约，另一个属于新约：因为当他说\Quote{使他们传给他们的儿子}时，他暗示了他们接受了诫命，但他们没有知道或实行它们；但他们自己接受了诫命，为的是\Quote{另一世代可能知晓}，即前者所不知晓的。\Quote{将要出生并兴起的儿子们。}因为那些已经出生的并没有兴起：因为他们没有把心放在上面，而是放在地上。因为兴起是与基督一同：因此经上说过：\BibleQuote{如果你们与基督一同兴起，就该思念上面的事。}\BibleRef{哥 3:1}他说，\Quote{并且他们可以告诉他们的儿子，为使他们寄望于天主。}……\Quote{并且不忘记天主的作为：}这就是说，不可夸大和夸耀自己的行为，好像是自己做的；而\Quote{是天主在工作}，在他们内行善事，\BibleQuote{既有意愿，也有能力按善意行事}\BibleRef{斐 2:13}。\Quote{并且要寻求祂的诫命。}……祂所命令的诫命。那么，他们既然已经知道了，为何还要寻求呢？除非藉着寄望于天主，他们寻求祂的诫命，为的是藉着祂的帮助，使这些诫命得以实现？他立刻加上理由说：\Quote{并且其精神未曾信赖天主}，这就是说，因为它没有信德，而信德能获得法律所命令的。因为当人的精神与工作的天主圣神合作时，天主所命令的就被成就了：这除非因着相信那使不虔敬者成义的天主，是不会发生的。\BibleRef{罗 4:5}这个信德，弯曲和苦毒的世代没有；因此论到同一世代曾说：\Quote{其精神未曾信赖天主。}因为这句话更准确地指出了天主的恩宠，它不仅赦免罪过，而且使人的精神与其合作行善事，好像他在说，他的精神没有相信天主。因为信赖天主，不是相信他的精神能离开天主而行正义，而是与天主一起。这就是相信天主：这当然比相信天主更多。因为有时我们也必须相信一个人，尽管我们不可相信他。因此，相信天主就是这个：在相信中依附于行善事的天主，以便与祂一起行善事……\par

8. 最后，\Quote{厄弗辣因的子孙，他们弯弓射箭，却在战争之日转身逃跑}\BibleRef{咏 77:9}。他们追求正义的法律，却没有达到正义的法律。\BibleRef{罗 9:31}为什么？因为他们不属于信德。因为他们就是那世代，其精神未曾信赖天主：但可以说，他们是属于行为的：因为他们没有像他们弯弓射箭（这些是外在行为，如同法律的行为）那样，也引导自己的心，而义人是因信德生活，信德藉着爱德行事；人藉此依附于天主，天主在人内按善意既有意愿也有能力行事。因为弯弓射箭并在战争之日转身逃跑，不就是在听道之日留意并打算，而在诱惑之日离弃吗？事先挥舞兵器，而在行动时刻拒绝作战？但他说\Quote{弯弓射箭}，似乎他应该说弯弓射箭……据说有些希腊抄本有\Quote{弯弓射箭}，所以我们无疑应理解为箭。但是，藉着厄弗辣因的子孙，他愿意理解为整个苦毒的世代，这是以部分代表全体的表达方式。也许选择这一部分来代表全体，是因为从这些人身上特别期待某种好处……虽然\Person{雅各伯}因儿子年幼而将他放在左手边，但他仍用右手祝福，并以隐晦意义的祝福使他超过兄长。\BibleRef{创 48:14}……因为当时预表着那后来的将成为最先的，最先的将成为最后的\BibleRef{玛 20:16}，藉着救主的来临，关于祂曾说过：\Quote{那在我以后来的，成了在我以前的。}\BibleRef{若 1:27}同样，正义的\Person{亚伯尔}被喜爱胜过兄长；对\Person{依市玛耳}也是如此，\Person{依撒格}；对\Person{厄撒乌}也是如此，虽比他早生的孪生兄弟\Person{雅各伯}；还有\Person{培勒兹}本人甚至也在出生上超过他的孪生兄弟，那兄弟先伸手出母腹，开始出生；\Person{达味}也被喜爱胜过他的兄长：\BibleRef{撒上 16:12}正如所有这些比喻和类似者预先发生，不仅言语，而且行为，同样，基督徒民族被喜爱胜过犹太民族，如同\Person{亚伯尔}被\Person{加音}杀害\BibleRef{创 4:8}，为救赎那民族，基督被犹太人杀害。这件事甚至在\Person{雅各伯}交叉伸开双手时预表了：他用右手触摸站在左边的\Person{厄弗辣因}，把他放在站在右边、他自己用左手触摸的\Person{默纳协}之前。\BibleRef{创 48:14}\par

9. 但他说\BibleQuote{他们在战争之日退后了}是什么意思，下面的经文教导了我们，其中他清楚地解释了这一点：\BibleQuote{他们没有遵守天主的盟约，也不愿行走在祂的法律中} \BibleRef{咏 77:10}。这就是\BibleQuote{他们在战争之日退后了}：他们没有遵守天主的盟约。当他们弯腰射弓时，他们也说出了最勇敢承诺的话，说：\BibleQuote{凡主我们的天主所说的一切，我们都要做，我们都要听从} \BibleRef{出 19:8}。\BibleQuote{他们在战争之日退后了}：因为服从的承诺不是听，而是试探来证明。但谁的心神与天主相托付，就紧握那信实的天主，祂\BibleQuote{不容许他受试探超过他所能承担的，而且在试探中也要开一条出路} \BibleRef{格前 10:13}，使他能忍受，不致在战争之日退后。……因此，这些人被如此标记：他说，\BibleQuote{这是一个没有纠正自己心的世代}。说的不是行为，而是心。因为当心被纠正时，行为就正确；但当心没有被纠正时，行为就不正确，即使它们看似正确。而弯曲的世代如何没有纠正自己的心，已充分显明，当他说：\BibleQuote{他们的心神没有与天主相符}。因为天主是正直的：因此，通过依附于正直，如同依附于不变的规则，人的心可以得到纠正，而它本身原是弯曲的。……\par

10. \BibleQuote{他们忘记了祂的恩惠，和祂向他们显示的奇妙作为；在他们先祖面前，祂所行的奇事} \BibleRef{咏 77:11}。这是什么，是一个不容忽视的问题。关于那些先祖，他刚谈论过，他们曾是弯曲悖逆的一代。……哪些先祖？既然这些正是那些他不愿后代像他们的先祖？如果我们认为他们是那些后裔所由生的人，例如亚巴郎、依撒格、雅各伯，那时他们早已长眠，当天主在埃及显示奇妙之事时。因为接着有：\BibleQuote{在埃及地，在塔尼斯平原} \BibleRef{咏 77:12}；那里说天主在他们先祖面前向他们显示了奇妙之事。他们莫非以心神临在？因为主在福音中论及同样的事说：\BibleQuote{因为众人都是为祂而活} \BibleRef{路 20:38}。或者我们更合适地理解为先祖是指梅瑟和亚郎，以及其他在圣经中记载也领受了圣神的长老们，梅瑟也领受了圣神，为使他们帮助他管理和承载同一民族？\BibleRef{户 11:17} 为什么他们不应被称为先祖？这不与天主是唯一的父方式相同，祂以祂的圣神重生那些祂使之成为儿子以承受永远产业的人；而是为了表示尊敬，因他们的年岁和慈爱的关怀：正如年长的保禄所说：\BibleQuote{我写这些事，不是要叫你们羞愧，而是劝告你们，如同我所亲爱的儿子} \BibleRef{格前 4:14}，尽管他深知主曾说过：\BibleQuote{不要称地上的人为父，因为只有一位是你们的父，就是天主} \BibleRef{玛 23:9}。这话并非是要从我们通常的说话方式中抹去这种人间尊称；而是免得我们借此重生得永生的天主的恩宠被归于任何人的能力或圣德。因此当他说：\Quote{我生了你们}；\Quote{在基督内}，和\Quote{借着福音}；免得那属于天主的被以为是出于他。……因此，埃及地必须被理解为这个世界的预象。\BibleQuote{塔尼斯平原}是谦逊诫命的平坦表面。因为谦逊诫命就是塔尼斯的解释。因此，在这个世界上，让我们接受谦逊的诫命，以便在另一个世界，我们堪当接受祂所许诺的举扬，祂为了我们的缘故在此成为卑微的。\par

\BibleQuote{祂劈開大海，使他們通過，將水立起如瓶} \BibleRef{詠 77:13}，為了使水起初如同被關閉一般站立，祂能藉祂的恩寵遏制肉性慾望的漲落潮汐，當我們棄絕這世界時，使所有罪惡被徹底洗淨，如同仇敵一般，使信友之民藉\OriginalTerm{聖洗聖事}{Sacrament of Baptism}通過。祂 \BibleQuote{在白天以雲彩帶領他們，在整夜以火光照明} \BibleRef{詠 77:14}，也能在屬靈上引導人的行徑，若信德向祂呼喊：\Quote{請照禰的言語引導我的行徑。} 在另一處論及祂說：\Quote{因為祂必使你的道路正直，並使你的行徑在平安中延長}，藉我們的主耶穌基督，祂在這世界的聖事，如同在白天，在肉身中顯明，好似在雲彩中；但在審判時，祂將顯明，如同在夜間的恐怖中；因為那時將有大災難，如同火一般，為義人發光，為不義者焚燒。祂 \BibleQuote{在曠野劈開磐石，給他們水喝，如同深淵} \BibleRef{詠 77:15}；\BibleQuote{從磐石中引出水來，使水如江河降下} \BibleRef{詠 77:16}，祂必然能將聖神的恩賜傾注於乾渴的信德（那件事的實行所屬靈象徵的恩賜），我說是從那隨後的\OriginalTerm{屬靈磐石}{Spiritual Rock}，就是基督，傾注出來：祂站著呼喊說：\BibleQuote{誰若渴，讓祂到我這裡來} \BibleRef{若 7:37}；又說：\BibleQuote{誰若喝了我所賜的水，從他胸中要流出活水的江河} \BibleRef{若 4:14}。因為祂說了這話，正如福音所載，\BibleRef{若 7:39}，是指那將要領受祂恩賜的聖神，就是信祂的人要領受的，如同杖接近了苦難的木頭，好使恩寵為信徒湧流出來。\par

\Quote{他們，}如同乖謬和惹怒的世代，\Quote{仍然犯罪得罪祂} \BibleRef{詠 77:17}：就是不信仰。因為這罪，就是聖神要責備世界的，如主說：\Quote{論到罪，是因為他們沒有信我} \BibleRef{若 16:9}。\Quote{他們在乾旱中激怒了至高者，}其他抄本作\Quote{在無水之地，}這是從希臘文更準確的翻譯，所指無非是乾旱。是在那曠野的乾旱中，還是在他們自己的乾旱中？因為他們雖曾飲於磐石，但使他們的口腹乾燥，心靈卻乾枯，沒有正義的果實使之清新。在那乾旱中，他們本該更忠信地向天主懇求，好使那位已使他們口腹滿足的天主，也賜予他們行為的正直。因為信友的靈魂向祂呼喊：\Quote{願我的眼睛看到正直。}\par

\Quote{他們在心裡試探天主，為自己的靈魂求食物} \BibleRef{詠 77:18}。憑信德求問是一回事，試探是另一回事。最後接著說：\Quote{他們毀謗天主說：天主豈能在曠野擺設筵席？} \BibleRef{詠 77:19}。\Quote{因為祂擊打磐石，水就湧出，溪流氾濫：祂還能賜下糧食嗎？或為祂的百姓預備筵席嗎？} \BibleRef{詠 77:20}。因此，他們不信，就為自己的靈魂求食物。宗徒雅各伯並非這樣吩咐為心靈求食物，而是勸告信友去求，而不是像那些試探並毀謗天主的人。他說：\Quote{你們中若有人缺乏智慧，就該向天主求，祂厚賜於眾人，也不斥責，必會賜給他；但該憑信德求，一點不疑惑} \BibleRef{雅 1:5}。這信德，那世代卻沒有，\Quote{他們的心沒有歸向祂，他們的心靈也沒有對祂保持忠信。}\par

\BibleQuote{因此主聽見，就延遲，在雅各伯中點火，怒火上達以色列} \BibleRef{詠 77:21}。祂解釋了祂所稱的火。祂稱怒火為火：雖然嚴格來說，火也確實燒死了許多人。那麼祂所說的\Quote{主聽見，就延遲}是什麼意思呢？難道拖延了領他們進入應許之地的時間嗎？他們正被領往那裡，本可在幾日內完成，但因罪惡的緣故，他們不得不在曠野中耗盡，在那裡他們也確實耗盡了四十年。若是如此，祂是拖延了那百姓，而不是那些試探並毀謗天主的本人；因為他們都死在曠野，他們的子女才進入了應許之地。或者，祂延遲了懲罰，為了先滿足不信的私慾，免得祂被認為發怒，因為他們向祂求祂所不能做的事？所以，\Quote{祂聽見了，}就\Quote{延遲報復；}在祂做了他們認為祂不能做的事之後，\Quote{怒火才上達以色列。}\par

15. 最后，在简要提及这两件事之后，他显然接着叙述的顺序。\BibleQuote{因为他们不信从天主，也不指望祂的拯救} \BibleRef{咏 77:22}。因为当他说到为何在雅各伯中燃起烈火，愤怒上升于以色列时，也就是说，\BibleQuote{因为他们不信从天主，也不指望祂的拯救}；他立即列举了他们忘恩负义的明显恩惠，说：\BibleQuote{祂命令天上的云彩，开启了天门} \BibleRef{咏 77:23}。\BibleQuote{祂给他们降下玛纳作食物，赐给他们天上的食粮} \BibleRef{咏 77:24}。\BibleQuote{人吃了天使的食粮；祂赐给他们丰盛的珍馐} \BibleRef{咏 77:25}。祂从天上带来南风，用祂的能力引导西南风 \BibleRef{咏 77:26}。\BibleQuote{祂给他们降下鲜肉如尘土，有翅的飞鸟如海沙} \BibleRef{咏 77:27}。\BibleQuote{它们落在他们的营中，在他们帐幕的四周} \BibleRef{咏 77:28}。\BibleQuote{他们吃了，且极其饱足；祂满足了他们的欲望，他们并未失去所愿} \BibleRef{咏 77:29}。看，这就是祂延迟的原因。但祂延迟了什么呢？让我们听。\BibleQuote{食物还在他们口中的时候，天主的愤怒就降在他们身上} \BibleRef{咏 77:30}。看，这就是祂所延迟的。因为在\BibleQuote{祂延迟}之前，之后\BibleQuote{在雅各伯中燃起烈火，愤怒上升于以色列}。因此，祂延迟，为要首先做他们以为祂不能做的事，然后把他们应受的惩罚降在他们身上。因为如果他们寄望于天主，不仅他们肉体的欲望，连精神的欲望也会得到满足。因为那位\BibleQuote{开启了天门，给他们降下玛纳作食物}以填饱不信者的大能者，并不缺乏能力将祂自己——那真正从天上来的食粮，即玛纳所预表的——赐给信徒；这确实是天使的食粮，天主的圣言以不朽的方式喂养不朽的天使；为了使人能吃到这食粮，祂成了血肉，并居住在我们中间。\BibleRef{若 1:44}因为祂自己就是那食粮，藉着福音的云彩降在整个世界上，并且（如同打开了天堂之门）藉着宣讲者的心被传扬，不是向一个抱怨和试探的会堂，而是向一个相信并寄望于祂的教会。祂也能以肉体所说、迅速穿过空气的话语的标记——如同飞鸟——喂养那些不试探而相信之人的软弱信德；但不是从北方（那里寒冷和迷雾盛行）而来的，即取悦于这个世界的雄辩，而是从天上带来南风；去往何处？除了大地之外？为使信德软弱的人，藉着听见尘世之事，得滋养以领受天上的事……\par

16. 但至于不信者，那弯曲和悖逆的世代，仿佛食物还在他们口中的时候，\BibleQuote{天主的愤怒上升于他们，击杀了他们中最多的} \BibleRef{咏 77:31}；就是说，他们中的多数，或者如一些抄本所写：\BibleQuote{他们中的肥壮者}，然而在我们所拥有的希腊抄本中，我们并没有找到这个读法。但如果这个读法更真实，那么\BibleQuote{他们中的肥壮者}除了指骄傲中的强壮者之外，还能指什么呢？关于这些人，经上说：\BibleQuote{他们的不义如脂肪流出}。\BibleQuote{祂捆绑了以色列的选民。}即使在那个世代中也有选民，弯曲和悖逆的世代并没有与他们混杂。但他们被捆绑，以致他们丝毫不能有益于那些他们出于父爱想要帮助的人。因为当天主向人发怒时，人的怜悯能带来什么益处呢？或者祂愿意人们明白，连选民也同时与他们一同被捆绑，为要使那些在心思和生命上不同的人，与他们一同受苦，作为义德和忍耐的榜样？因为我们知道，圣人们被与罪人一同掳去，没有别的原因；因为在希腊抄本中，我们读到的不是\OriginalTerm{捆绑}{\GreekText{ἐνεπόδισεν}}（意为\BibleQuote{捆绑}），而是\OriginalTerm{同被捆绑}{\GreekText{συνεπόδισεν}}，这更确切地说是\BibleQuote{一同被捆绑}。\par

17. 但那弯曲和悖逆的世代，\BibleQuote{在这一切事上仍旧犯罪，不信祂奇妙的作为} \BibleRef{咏 77:32}。\BibleQuote{他们的日子在空虚中消逝} \BibleRef{咏 77:33}。尽管如果他们相信，他们本可与那位经上对祂说\BibleQuote{你的年岁没有穷尽}的主一同拥有不消逝的真实日子。因此，\BibleQuote{他们的日子在空虚中消逝，他们的岁月急速而过。}因为整个人类的生命都是急促的，那看似较长的也不过是稍长片刻的蒸汽。\par

18. 尽管如此，\BibleQuote{当祂击杀他们时，他们寻求祂：}不是为了永生，而是害怕过早结束那蒸汽。那时寻求祂的，其实不是那些被祂击杀的人，而是那些因他们的榜样而害怕被击杀的人。但圣经这样说，好像被击杀的人寻求天主，因为他们是一个民族，说话如同一个身体：\BibleQuote{他们回转，清早来到天主面前} \BibleRef{咏 77:34}。\BibleQuote{他们记起天主是他们的帮助者，至高的天主是他们的救赎者} \BibleRef{咏 77:35}。但这一切都是为了获得暂时的好处，避免暂时的祸患。因为那些为暂时的福乐而寻求天主的人，其实不是在寻求天主，而是寻求事物。因此，他们是以奴隶般的恐惧敬拜天主，而非自由的爱情。这样，天主没有被敬拜，因为被敬拜的是他们所爱的事物。因此，既然天主被发现比万有更伟大、更美好，祂必须被爱超过万有，以便祂能被敬拜。\par

最后，让我们在此看以下经文：\Quote{他们用口爱慕祂，} 他说，\Quote{用舌头向祂说谎。} \BibleRef{咏 77:36} \Quote{但他们的心不正向祂，他们在祂的盟约上也不忠信。} \BibleRef{咏 77:37} 祂发现他们口里一套，心里一套，因为人的隐秘在祂面前是赤露敞开的，祂毫无阻碍地看见他们所更喜爱的是什么。因此，人心在天主面前是正直的，当它单单为天主的缘故寻求天主时。因为有一件事他曾向主祈求，他仍要寻求，使他能永远住在主的殿中，默想祂的甘饴。忠信者的心对祂说：我必得饱足，不是埃及的肉锅，也不是瓜、韭菜、葱、蒜等弯曲苦毒的世代宁愿选择甚至超过天上食粮的东西（\BibleRef{出 16:3} ），也不是可见的玛纳和那些飞禽；而是：\Quote{我得了满足，当祢的荣耀显现时。} 因为这就是新约的产业，他们在这产业上不忠信；然而那信仰，即使在它还隐藏时，是在蒙选者中；现在它已经显明了，却不在许多蒙召者中。\Quote{因为蒙召的多，蒙选的少。} \BibleRef{玛 20:16} 因此，那弯曲和苦毒的世代就是如此，甚至当他们似乎寻求天主时，口里爱慕，舌头说谎；但心里不正向天主，因为他们更喜爱那些他们为了寻求天主帮助的东西。\par

\Quote{但是祂，因为祂是仁慈的，将为他们的罪恶而宽恕，并不毁灭他们。并且祂将多次转离祂的怒气，不会燃起祂全部的怒火。} \BibleRef{咏 77:38} 这些话使许多人因天主的仁慈而自认为他们的不义可以不受惩罚，即使他们坚持成为像那世代那样，即被描述为\Quote{弯曲和苦毒的；这世代没有正直他们的心，其心神没有与天主相契。} 与他们一致是没有益处的。因为，用他们的话来说，如果天主或许不会毁灭即使是恶人，那么毫无疑问祂不会毁灭善人。那么，我们为什么不选择那毫无疑问的善呢？因为那些用舌头向祂说谎的人，尽管心里另有打算，确实认为并愿意连天主也成为说谎者，当祂警告这样的人永罚时。但他们用谎言欺骗不了祂，祂用真理说话也不欺骗他们。因此，关于这些神圣的话语，弯曲的世代自欺欺人，但愿它不要像自己的心一样把它们扭曲；因为即使心被扭曲了，它们仍是正直的。首先，它们可以按照福音中的话理解：\Quote{好使你们成为你们在天之父的子女，因为祂使太阳升起，照耀恶人也照耀善人，降雨给义人也给不义的人。} \BibleRef{玛 5:45} 谁能看不见祂宽容恶人的仁慈是何等伟大？但在审判之前，祂如此宽容那个民族，并没有燃起祂全部的怒火，彻底根除并毁灭他们；这事在祂的话语和祂的仆人梅瑟为他们的罪所作的代祷中明显可见，那时天主说：\Quote{你且由我灭绝他们，我要使你成为一个大民族。} \BibleRef{出 32:10} 他代祷，更愿意自己被除灭而不是他们被除灭；他知道他是在一位仁慈者面前行事，因为祂绝不会除灭他，反而会为了他的缘故宽恕他们。让我们看看祂是多么宽容他们，而且至今仍在宽容……\par

其次，为了不让我们似乎对神圣话语施加暴力，也避免在说过\Quote{祂不毁灭他们}的地方，我们应当说：\Quote{但以后祂要毁灭他们。}关于这首圣咏本身，让我们转向圣经中一个很常见的短语，藉此可以更仔细、更真实地解决这个问题。在稍往下提及这些同一批人时，当祂提到埃及人因他们而遭受的事后，祂说：……\Quote{祂领他们到了祂的圣山，就是祂右手所获得的山。祂从他们面前驱逐了外邦，用绳索将土地分给他们为产业。} 如果有人用这些话质问我们，并对我们说：祂怎能将这一切作为恩赐给予他们呢？因为同那些人并没有被领进应许之地，正如那些从埃及被拯救出来的人一样，他们死了？除了回答：他们因子孙的继承而成为同一民族，我们还能回答什么？……\par

\Quote{祂记起他们不过是血肉，是一阵去而不返的风。} \BibleRef{咏 77:39} 因此，祂藉着恩宠呼唤并怜悯他们，祂亲自召唤他们回来，因为他们凭自己不能回来。因为血肉如何能回来，\Quote{是一阵行走而不回转的风，}当邪恶的罪孽之重压着它坠入最低、最远的邪恶之处，若非藉着恩宠的拣选？……因为这样也解决了一个不小的问题，即箴言中如何记载，当圣经论及不义之道时说：\Quote{凡行在这道上的，都不能返回。} \BibleRef{箴 2:19} 这话说得好像所有不虔敬的人都绝望了；但圣经只赞美恩宠；因为人凭自己能走在那道上，却不能凭自己回来，除非被恩宠召回。\par

23. 那么我要说说这些乖僻而苦毒的人，\Quote{他们多少次在旷野中激怒祂，在无水之地惹动祂的义怒！} \BibleRef{咏 78:40}。\Quote{他们转而试探天主，激怒了以色列的圣者} \BibleRef{咏 78:41}。他是在重复他们那同样的不信，就是他在前面已经提及的。但重复的原因，是为了也提及祂为他们的缘故加在埃及人身上的灾祸：所有这些事他们当然应该记念，而不该忘恩负义。最后，接下来是什么？\Quote{他们不记念祂的手，就是在祂救赎他们脱离那搅扰者的手的日子} \BibleRef{咏 78:42}。于是他开始讲述祂对埃及人所行的事：\Quote{祂在埃及设立了他的奇事，在塔尼斯平原上设立了他的异兆} \BibleRef{咏 78:43}：\Quote{祂使他们的江河变为血，使他们的雨水不能喝} \BibleRef{咏 78:44}，或者更确切地说，\Quote{水的涌流}，正如一些人由希腊文\OriginalTerm{涌流之水}{\GreekText{τὰ} \GreekText{ὀμβρήματα}}所理解的，我们在拉丁文中称之为\OriginalTerm{泉水}{scaturigines}，即从地下涌出的水。\Quote{祂使狗蝇来到他们中间，吞吃了他们；又使青蛙毁灭他们} \BibleRef{咏 78:45}。\Quote{祂将他们的果实交给霉病，将他们的劳苦交给蝗虫} \BibleRef{咏 78:46}。\Quote{祂用冰雹打坏了他们的葡萄树，用霜冻打坏了他们的桑树} \BibleRef{咏 78:47}。\Quote{祂将他们的牲畜交给冰雹，将他们的财产交给火} \BibleRef{咏 78:48}。\Quote{祂将祂的愤恨之怒降在他们身上，愤恨、恼怒和患难，藉恶天使降下的惩罚} \BibleRef{咏 78:49}。祂为祂的怒气的路径铺平了道路，将他们的牲畜关闭在死亡中 \BibleRef{咏 78:50}。\Quote{祂击杀了埃及地一切头生的，在含的帐棚中，是他们劳力的初果} \BibleRef{咏 78:51}。\par

24. 所有这些对埃及人的惩罚都可以按照各人选择的解释方式，通过寓意解释来理解，并将它们与它们所指涉的事物相比较。我们也将努力这样做；而且我们越得到天主的援助，就会做得越恰当。因为要做这事，这篇圣咏中的那些话约束了我们，其中说：\Quote{我要开口说比喻，我要从起初宣告谜语。} 为此缘故，这里甚至提到了一些事，我们读不到它们曾发生在那时埃及人身上，尽管在\WorkTitle{出谷纪}中所有的灾祸都按顺序被非常仔细地叙述，以至于我们确信，那在\WorkTitle{出谷纪}中没有提到的事并非在圣咏中无故提及，而只能以象征方式解释；同时我们也能理解，即使是其余明显发生的事情，也是为了一些象征的意义而做或描述的。因为圣经在许多先知言论的段落中都是这样做的。……因此，在埃及人的灾祸中，即在被称为\WorkTitle{出谷纪}的书中，圣经特别谨慎地按顺序叙述了所有使他们受苦的事，但并没有发现圣咏中所说的\Quote{祂将他们的果实交给霉病}。同样，当他说\Quote{祂将他们的牲畜交给冰雹}时，又加上了\Quote{以及他们的财产交给火}：关于牲畜被冰雹打死的事，在\WorkTitle{出谷纪}中读到了（\BibleRef{出 9:25}），但是财产被火烧的事却完全读不到。尽管声音和火与冰雹一同降临，正如雷声通常伴随着闪电一样；然而，没有记载任何东西被交给火去焚烧。最后，那些冰雹不能伤害的柔软之物，据说没有被击打，即被硬物击伤；这些东西后来被蝗虫吞噬了。此外，这里所说的\Quote{以及他们的桑树被霜冻打坏}，也不在\WorkTitle{出谷纪}中。因为霜冻与冰雹大不相同；在晴朗的冬夜，大地因霜冻而变白。\par

25. 那么，这些事物究竟象征什么，让诠释者尽其所能地解释，让读者和听者公正地判断。水变成血，在我看来似乎象征对事物起因的一种肉性看法。狗蝇是狗的行为，它们在刚出生时甚至不认自己的父母。青蛙是极其多言的虚妄。霉病暗中为害，有些人将其解释为锈，另一些人则解释为黑霉；这种邪恶之物与什么恶习相比更为恰当呢？岂不是与最不易显露的恶习（例如过度自信）相比吗？因为这是一种有害的气息，暗中在果实中起作用；正如在道德中，暗自骄傲，当人自以为了不起，其实什么都不是。\BibleRef{迦 6:3} 蝗虫是以口伤害的恶意，即不忠的见证。冰雹是不义夺取他人之物；由此产生偷窃、抢劫和掠夺；但掠夺者自己反被其邪恶掠夺。霜冻象征这样的过犯：对邻人的爱因愚昧的黑暗（如同夜的寒冷）而被冻结。至于火，如果这里不是指从闪电云中的冰雹中所提及的火，因为此处他说：\Quote{他把他们的财产交给了火，}这里暗示某物被烧毁，但我们读到并未因那火而发生——在我看来，我说，它象征愤怒的野蛮，甚至可能犯下杀人罪。至于野兽的死亡，依我判断，是象征失贞。因为产生后裔的欲望，我们与野兽共有。因此驯服并规范这种欲望，便是贞洁之德。首生者的死亡，是剥去那使人能与人类联合的正义本身。但这些事物的比喻性意义是否如此，或者是否以另一种方式理解更好，谁不会被这事实所触动：埃及人遭受十灾，而刻有十诫的石板被赐予，好让天主的子民受其统治？关于两者的比较，既然我们已在别处讲过，无需在此加重这篇圣咏的阐释；我们只提醒诸位：在这里，虽然顺序不同，却也纪念了埃及的十灾，因为代替\WorkTitle{出谷纪}中有而此处没有的三灾（即虱子、脓疮、黑暗），此处纪念了另外三灾（即霉病、霜冻和火）；这火不是闪电之火，而是他们的财产被交付于其中的火，那在\WorkTitle{出谷纪}中并未读到。\par

26. 但已经足够清楚地暗示，这些事是由于天主的审判，藉邪恶天使的媒介而临到他们身上，在这邪恶的世界里，如同在埃及和塔尼斯平原一般，我们应在那里谦卑，直到那世界来到，我们得以从这卑微中被高举。因为在希伯来语中，埃及本身便表示黑暗或磨难，而塔尼斯，如我所指出的，在这语言中被理解为谦卑的命令。因此，关于邪恶天使，在这篇圣咏中，当他正谈论那些灾祸时，插入了一些不容草率放过的话：他说：\Quote{祂把他们交于，藉邪恶天使之折磨}。众所周知，魔鬼及其天使是如此邪恶，以致为他们预备了永火，没有信者不知道这一点；但是，竟有来自主天主的折磨由他们被派往某些祂判定应得此罚的人身上，这对那些不惯于思考天主完美的公义如何善用甚至邪恶事物的人，似乎是件难事。因为就这些事物的本质而言，除了祂自己，谁创造了它们？但祂并没有将它们造成邪恶；然而，祂既为善，便善用它们，即适宜而公正地使用；正如另一方面，不义之人以邪恶的方式使用祂美好的受造物。因此，天主不仅用邪恶天使来惩罚恶人，如同这篇圣咏所谈论的那些人，也如同阿哈布王的事例，一个说谎的灵按天主的旨意欺骗了他，使他阵亡于战争：\BibleRef{列上 22:20}，而且也用它们来考验并彰显善人，如祂在约伯身上所做的。但就可见元素的物质方面而言，我认为善与恶的天使都能按照赋予各者的能力加以使用：正如善人与恶人也会尽其所能，按照人性的尺度使用这类事物。因为我们使用土、水、气和火，不仅用于维持生计的必需品，也用于许多多余、嬉戏和奇妙精巧的操作。因为无数被称为\OriginalTerm{机械装置}{\GreekText{μηχανήματα}}的东西，是由这些元素科学地运用而塑造的。但天使对这些事物有远为广大的能力，包括善与恶的，尽管善者拥有的更大；但仅限于天主意愿和眷顾所命令或准许的范围内；我们也以此条件拥有能力。因为即使在这些情况下，我们也并非能为所欲为。但在最无误的书中，我们读到魔鬼甚至能从天上降火，以神奇而可怖的猛烈烧毁了一位圣人的大量牲畜：\BibleRef{约 1:16}；倘若不是基于圣经的权威读到，恐怕没有信者敢将此归于魔鬼。但那人，因天主的恩赐而正义、坚定且拥有虔诚的知识，并没有说：主赐予，魔鬼取走了；而是说：\Quote{主赐予，主收回：}\BibleRef{约 1:21}；他深知，即使魔鬼能对这些元素所做的，若非其主的意愿和许可，仍不会对天主的仆人做出；他挫败了魔鬼的恶意，因为他知道是谁在利用魔鬼来考验他。因此，在不信从之子中，如同在他自己的奴仆中，魔鬼运行工作，\BibleRef{弗 2:2}，如同人使用牲畜，且仅在天主公义审判所容许的范围内。但有一回事，是他的力量被更大的力量阻止，甚至不能随心所欲地对待他自己的；另一回事，是给予他力量甚至对付那些不属于他的人。正如人对待自己的牲畜，按人的理解，他随意而行，然而若被更大的力量阻止，他就不行；但要动别人的牲畜，他必须等到主人给予许可。在前一种情况下，他已拥有的力量受到限制；在后一种情况下，未曾拥有的力量被赐予。\par

27. 如果情况如此，如果天主藉邪恶天使把这些灾祸降在埃及人身上，我们岂敢说，水也是藉同样的天使变成血，青蛙也是藉它们创造的，甚至连法郎的术士也能藉他们的法术造出类似之物；\BibleRef{出 8:7} 以致于邪恶天使站在双方，一方折磨他们，另一方欺骗他们，这是按照最公义且全能的天主的判断和安排，祂公正地甚至利用不义之人的恶行？我不敢如此说。因为法郎的术士何以绝不能造出虱子？\BibleRef{出 8:18} 岂不是因为连这些邪恶天使也不被允许这样做吗？或者，更确切地说，原因隐藏着，且超出我们探究的能力？因为如果我们假定天主藉邪恶天使行那些事，是因为正在施行惩罚，而非赐予祝福，好象天主从不藉善天使施行惩罚，而只藉那些如同天上义怒的执行者；那么结论将是，我们必须相信连索多玛也是被邪恶天使倾覆的，而亚巴郎和罗特似乎曾在他们的屋檐下接待了邪恶天使；这与最清楚的圣经相悖，我们万万不可这样想。那么显然，这些事可能藉善天使或恶天使行在人身上。该做什么或何时做，我不明了；但那行事者，祂明了，以及祂愿意启示的人，那人也明了。然而，就神圣圣经向我们应用所允许的而言，我们读到，恶人受到惩罚，既藉善天使，如索多玛人，也藉邪恶天使，如埃及人；但善人藉善天使以肉体的补赎被考验和证明，我未想到。\par

28. 但就本圣咏的这段经文而言，若我们不敢将那些由受造物奇妙形成的事归诸恶天使，我们确实有一样东西可以无疑地归诸他们：那就是牲畜之死、首生者之死，尤其是这一切的根源，即心硬，使他们不肯放走天主的子民。\BibleRef{出 4:21} 因为当天主被说成是造成这种最邪恶恶毒的顽固时，祂不是藉着暗示和激励造成的，而是藉着离弃，使悖逆之子行天主适当而公义地准许的事。\BibleRef{弗 2:2} ...此外，那些我们说过由这些肉体灾祸所象征的恶习，因先前所说：\Quote{我要开口说比喻，} 极宜认为是由恶天使在那些因天主公义而受制于他们的人中造成的。因为即使当宗徒所说的事发生时：\Quote{天主遂放任他们随从心中的情欲，行不合宜的事，} 那些恶天使怎能不住在其中并以此为乐，仿佛那是他们自己的工作呢？人的骄傲最公义地受制于他们，除了那些被恩宠拯救的人以外。\BibleQuote{对这些事，谁够资格呢？} \BibleRef{格后 2:16} 因此，当他曾说：\BibleQuote{祂向他们发怒，忿怒、恼怒、困苦，藉恶天使降罚；} 至于他所加的：\BibleQuote{祂为祂的怒气的路径开了一条路} \BibleRef{咏 77:50} ，我请问，谁的眼目足以看透，以致能理解并领会藏在如此深奥中的意义呢？因为天主怒气的路径是祂以隐藏的公义惩罚埃及人的不虔敬；但祂为同一条路径开了一条路，仿佛藉恶天使将他们从隐秘处拉出，进入明显的恶行，从而最明显地惩罚了那些最明显不虔敬的人。从恶天使的这种权势中，惟有天主的恩宠能拯救人，宗徒论及此恩宠说：\BibleQuote{祂拯救了我们脱离黑暗的权势，并将我们迁到祂爱子的国里；} \BibleRef{哥 1:13} 这些事是那民族的预象，当他们从埃及人的权势下被拯救，并被迁到流奶流蜜的许地时，这预象正象征恩宠的甘甜。\par

29. 圣咏在追忆埃及人的灾祸之后继续说：\BibleRef{咏 77:51} 又说：\BibleQuote{祂领出自己的子民，如同羊群，领他们像羊群在旷野中通过} \BibleRef{咏 77:52} 。\BibleQuote{祂怀有望德带领他们前行，他们不畏惧，而他们的敌人被海淹没} \BibleRef{咏 77:53} 。这事以更大的益处应验了，因为这是更内在的事：当我们脱离黑暗的权势，在心灵中被迁到天主的国，并在属灵的牧场中成为天主的羊群，行走在这世界仿佛在旷野中，因为我们的信德无人能见。因此宗徒说：\BibleQuote{你们的生命已与基督一同藏在天主内。} \BibleRef{哥 3:3} 但我们正怀有望德被引领回家，\BibleQuote{因为我们得救是因望德。} \BibleRef{罗 8:24} 我们不必害怕。因为：\BibleQuote{天主若为我们，谁能反对我们呢？} \BibleRef{罗 8:31} 我们的敌人被海淹没了，祂藉着圣洗中的罪赦，抹去了他们。\par

30. 接着经文说：\BibleQuote{祂领他们进入祂的圣山} \BibleRef{咏 77:54} 。这何等更好地进入了圣教会！\Quote{祂右手所夺得的山。} 基督所获得的教会何等崇高，关于祂曾说过：\BibleQuote{主的臂膀向谁显示了呢？} \BibleRef{依 53:1} \BibleRef{咏 77:55} 。\Quote{祂从他们面前驱逐了列邦。} 也从祂信徒的面前。因为列邦在某种意义上就是外邦错谬的恶灵。\Quote{祂用测量绳给他们分地。} 在我们内，\BibleQuote{这一切都是同一圣神所运行，随祂的旨意分给各人。} \BibleRef{格前 12:11}\par

31. \Quote{祂使以色列的支派住在他们的帐幕中。}他说，祂使以色列的支派住在外邦人的帐幕中；我认为这更能以属灵的方式来解释，因为借助基督的恩宠，我们正被提升至天上的光荣，犯罪的天使正是从那光荣中被驱逐、被摔落的。因为那弯曲和苦毒的一代，虽为这些肉身的祝福，仍未脱去旧人的外衣，\Quote{竟仍然试探}，\Quote{并激怒了至高者天主，不遵守祂的法度}\BibleRef{咏 77:56}；\Quote{他们背离，不守盟约，如同他们的祖先}\BibleRef{咏 77:57}。因为他们曾在一项盟约和诫命之下说：\Quote{凡主我们的天主所说的一切，我们都要遵行，我们都要听从。}\BibleRef{出 19:8}他所说的\Quote{如同他们的祖先}确实值得注意；尽管整篇圣咏的文本似乎都是针对同一群人，但现在看来，那些话涉及的是已经身在应许之地的人，而所说的祖先则是那些在旷野中激怒天主的人。\Quote{他们被扭曲，}他说，\Quote{变成弯曲的弓，}或如一些抄本所说，\Quote{变成悖逆的弓}\BibleRef{咏 77:58}。但这意思在接下来的话中更为清晰地显明，他说：\Quote{他们用他们的丘坛激怒了祂}\BibleRef{咏 77:59}。这表明他们堕入了崇拜偶像。因此弓被扭曲，不是为了主的名，而是反对主的名；主曾对同一子民说：\Quote{除我之外，你不可有别的神。}\BibleRef{出 20:3}祂用弓来比喻心灵的意向。最后，这同一观念更为清晰地展开，他说：\Quote{他们以他们的偶像惹动祂的怒气。}\par

32. \Quote{天主听见，就厌恶：}即祂留意并施行了报复。\Quote{并使以色列极大衰微}\BibleRef{咏 77:60}。因为当天主厌恶时，那些凭天主帮助才成为他们本来面目的人，又算什么呢？但他无疑是在追忆那件事，即司铎厄里时代他们被培肋舍特人击败，主的约柜被掳，他们被大量屠戮。他说的就是这事。\Quote{祂弃绝了史罗的帐幕，祂在人间居住的帐幕}\BibleRef{咏 77:61}。祂巧妙地说出了为何弃绝祂的帐幕，当他说：\Quote{祂在人间居住的帐幕。}因此，当他们不配祂在他们中间居住时，祂为何不弃绝这帐幕呢？这帐幕并非为祂自己而设立，而是为了他们的缘故，而如今祂断定他们不配祂在他们中间居住。\Quote{祂将他们的力量交与俘虏，将他们的荣美交在敌人手中。}那使他们自以为不可战胜并引以为傲的约柜，他称为他们的\Quote{德能}和\Quote{荣美}。最后，后来当他们行为邪恶并夸耀主的圣殿时，主藉先知恐吓他们，说：\Quote{你们看我对史罗所行的，那里是我的帐幕。}\BibleRef{耶 7:12}\Quote{祂用刀剑结束了祂的人民，并厌弃了祂的产业}\BibleRef{咏 77:62}。\Quote{他们的少年人被火吞灭：}即愤怒。\Quote{他们的贞女不举哀}\BibleRef{咏 77:63}。因为即使在恐惧敌人时，也无暇为此。\Quote{他们的司铎倒在刀下，他们的寡妇也不哀哭}\BibleRef{咏 77:64}。因为厄里的两个儿子倒在刀下，其中一人的妻子成为寡妇，随即在分娩中死去\BibleRef{撒上 4:19}，由于同样的混乱，她不能以葬礼的仪式被哀悼。\Quote{主如睡醒一般}\BibleRef{咏 77:65}。因为当祂将祂的子民交在祂所恨恶的人手中，当有人对他们说：\Quote{你的天主在哪里？}时，祂似乎睡了。\Quote{祂于是醒起，如睡醒的人，又如酒醉的勇士。}除了主的圣神，无人敢对天主说这话。因为他是以亵渎的恶人看来似乎的方式说的；好像一个醉酒的人沉睡很久，当祂不像人想的那样迅速救助时。\par

33. \Quote{祂从后面击打祂的仇敌}\BibleRef{咏 77:66}，即是那些因能夺走祂的约柜而欢喜的人；因为他们被击打在后部。\BibleRef{撒上 5:6}在我看来，这似乎是那惩罚的记号，人若回顾背后之事，就要受此折磨；正如宗徒所说，他应视这些事为废物。\BibleRef{斐 3:8}因为那些这样领受天主盟约的人，并不脱去旧日的虚妄，他们就像那外邦民族，将所夺的盟约之柜放在自己的偶像旁边。然而，即使他们不愿意，那些旧事仍要消逝：因为\Quote{凡有血肉的如同草，人的光荣如同草上的花。草枯干了，花也凋谢了：}\BibleRef{依 40:6}但主的约柜\Quote{却永远长存}，即天国隐秘的盟约，那里有天主永恒的圣言。但那些喜爱背后之事的人，正因这些事本身而最公义地受刑罚。因为\Quote{祂给了他们永久的羞辱。}\BibleRef{咏 77:67}\par

34. \BibleQuote{祂弃绝了}，他说，\BibleQuote{若瑟的帐幕，厄弗辣因支派祂没有拣选} \BibleRef{咏 77:68}。\BibleQuote{祂却拣选了犹大支派} \BibleRef{咏 77:69}。祂没有说，祂弃绝了\Person{勒乌本}的帐幕，他是\Person{雅各伯}的长子；\BibleRef{创 49:3}也没有说他们中的后来者，在出身次序中先于犹大的人；因此他们被弃绝、不被拣选，而犹大支派被拣选。因为或许可以说他们理当被弃绝；因为在\Person{雅各伯}祝福他儿子们的祝福中，他提到了他们的罪过，\BibleRef{创 49:5}并深深憎恶他们；虽然在它们之中，\Person{肋未}支派堪当成为司祭支派，\Person{梅瑟}也出自那里。\BibleRef{出 2:1}祂也没有说，祂弃绝了\Person{本雅明}的帐幕，或没有拣选\Person{本雅明}支派，从这支派中已经有了一位君王开始；因为从那里拣选了\Person{撒乌耳}；\BibleRef{撒上 9:1}因此由于那时间的非常接近，当他被弃绝和拒绝，而\Person{达味}被拣选时，\BibleRef{撒上 16:1}这本来很方便地说出来；但却没有说：而是他特别提到了那些似乎因更卓越的功绩而卓越的人。因为\Person{若瑟}在\Place{埃及}供养了他的父亲和兄弟们，并因他的虔诚、贞洁、智慧而被不虔敬地出卖后，被最公正地高举；\BibleRef{创 41:40}而\Person{厄弗辣因}因他祖父\Person{雅各伯}的祝福被偏爱超过他的长兄：然而天主\BibleQuote{弃绝了\Person{若瑟}的帐幕，\Person{厄弗辣因}支派祂没有拣选。}在此处，藉着这些著名功绩的名字，我们还能理解什么，除非是那整个以旧的贪求，向主索求地上赏报的民族，被弃绝和拒绝，而犹大支派被拣选，并非为了那同一个\Person{犹大}的功绩？因为\Person{若瑟}的功绩远为更大，但因犹大支派，既然基督按肉身由此而出，圣经确证了基督的新民族被偏爱超过那旧民族，主开口说比喻。此外，因此在那随后的话中，\BibleQuote{祂所拣选的\Place{熙雍}山}，我们更好地理解基督的教会，不为了现世肉体的祝福而敬拜天主，而是以信德的眼睛从远处盼望未来和永恒的赏报：因为\Place{熙雍}也被解释为\Quote{瞭望}。\par

35. 最后随后，\BibleQuote{祂像独角兽一样建造了祂的圣所} \BibleRef{咏 77:70}：或者，如一些诠释者由此造出一个新词，\Quote{祂的圣化}。独角兽被正确地理解为那些坚毅希望被提升到那唯一之事物的人，关于这事另一首圣咏说：\BibleQuote{我寻求了主一件事，我必要求它}。但是天主的圣化，按照宗徒\Person{伯多禄}，被理解为圣洁的子民和王家的司祭。\BibleRef{伯前 2:9}但随后的话：\BibleQuote{在祂为永远所建立的土地上：}希腊抄本有\OriginalTerm{归于永世}{\GreekText{εἰς} \GreekText{τὸν} \GreekText{αἰῶνα}}，不论我们称为\Quote{归于永远}或\Quote{归于一个时代}，取决于拉丁译者的喜好；因为它意指两者：因此后者出现在一些拉丁抄本中，前者出现在另一些中。有些还有复数形式，即\Quote{归于诸时代}：这在我们手中的希腊抄本中没有找到。但是哪位信徒会怀疑，教会，尽管有人离去，有人前来，她以有死的方式离开此生，却仍为永远所建立？\par

36. \BibleQuote{祂拣选了祂的仆人达味} \BibleRef{咏 77:71}。我说，犹大支派，为了达味；但达味为了基督；因此犹大支派为了基督。在祂经过时，盲人呼喊：\BibleQuote{可怜我们吧，达味之子！} \BibleRef{玛 20:30}，并立刻因祂的怜悯而获得光明，因为他们呼喊的是真实的。因此宗徒并不是粗略地提及，而是留意到，写信给\Person{弟茂德}：\BibleQuote{你要记念耶稣基督从死人中复活了，祂是达味的后裔}等等 \BibleRef{弟后 2:8}。因此救主自己，按肉身是达味的后裔，在这段经文下以达味的名义被预示，主开口说比喻。不要让我们不安，当祂说\BibleQuote{祂拣选了达味}，在这名字下祂预意基督，祂加上了\BibleQuote{祂的仆人}，而不是祂的儿子。是的，甚至从这里我们可以看出，不是与圣父同永恒的独生子的本体，而是\OriginalTerm{奴仆的形体}{form of a servant}是从达味的后裔所取的。\par

37. \Quote{祂从羊群里选取了他，从怀胎的母羊中召选了他：牧养雅各伯祂的仆人，以色列祂的产业。}\BibleRef{咏 78:72}。这\Person{达味}，的确，\Person{基督}的肉身来自他的后裔，从牧放牲畜转到了统治人：但我们的\Person{达味}，即\Person{耶稣}自己，从人到人，从\Place{犹太}人到外邦人，却按照比喻从羊到羊被取去并转移。因为现在在那地方没有\Quote{在基督内的犹太地区的教会}，在吾主近日的苦难与复活后，它们属于受割损的人；关于她们，宗徒说：\Quote{但我在面容上为犹太地区的在基督内的教会所不认识}等等。\BibleRef{迦 1:22}从此那些受割损之民的教会已经消失：因此在地上现今存在的\Place{犹太}，现在没有\Person{基督}。祂已从那里被移去，现在祂在牧养外邦人的羊群。确实，祂是从\OriginalTerm{怀胎的母羊}{teeming sheep}后面被取去的。因为那些先前的教会是这样的，正如在\BibleBook{}{雅歌}中论到它们所说：\Quote{你的牙齿——像一群剪了毛的母羊从洗池中上来，全都生双胞胎，其中没有不生育的。}\BibleRef{歌 4:2}因为他们当时像脱去羊毛一样卸下了世界的重担，\BibleRef{宗 2:45}，当时他们把变卖物品的价钱放在宗徒脚前，\BibleRef{宗 4:34}，从那\OriginalTerm{洗池}{Laver}上来，关于那洗池，\Person{伯多禄}宗徒在他们因流了\Person{基督}的血而忧苦时劝诫他们，说：\Quote{你们悔改吧，你们每人要因主耶稣基督的名受洗，你们的罪必得赦免。}\BibleRef{宗 2:38}但他们生了双胞胎，即\OriginalTerm{双重的爱}{twin love}的两条诫命所产生的工作，就是爱天主和爱近人：因此他们中间没有一个不生育的。我们的\Person{达味}从这些怀胎的母羊后面被取去，如今在外邦人中牧养别的羊群，那些羊群也是\Quote{雅各伯}和\Quote{以色列}。因为经上说：\Quote{牧养雅各伯祂的仆人，以色列祂的产业。}……除非有人愿意作这样的区分：即在这时间，\Person{雅各伯}服务；但他将是天主的永恒产业，到那时他将面对面看见\Person{天主}，因此他得了以色列的名字。\BibleRef{创 32:28}\par

38. \Quote{祂牧养了他们，}他说，\Quote{以祂心中的纯真}\BibleRef{咏 78:73}。有谁比祂更纯真？祂不仅没有任何罪可以被征服，甚至也没有任何罪可以征服？\Quote{以祂双手的理解引导他们回家：}或如一些抄本所作，\Quote{以祂双手的理解力。}任何人或许会认为，这样说更好：\Quote{以双手的纯真和心中的理解；}但祂比其他人更知道祂所说的，祂宁愿将\OriginalTerm{纯真}{innocence}与心相连，将\OriginalTerm{理解}{understanding}与双手相连。据我判断，原因是：因为许多人自以为纯真，他们不行恶是因为害怕若行了就会受苦；但他们仍然存有行恶的意愿，只要能够不受惩罚。这样的人似乎有\OriginalTerm{双手的纯真}{innocence of hands}，却没有\OriginalTerm{心中的纯真}{innocence of heart}。请问，如果纯真不在心中——人就是按\OriginalTerm{天主的肖像}{image of God}被造的，那纯真算什么？\BibleRef{创 1:27}但在此，当祂说\Quote{以祂双手的理解（或理智）引导他们回家}，依我看来，祂所说的理解是祂自己在信徒中所产生的；因此\Quote{祂双手的}——因为造作属于双手，但这里是指可以理解的天主之手的含义；因为甚至\Person{基督}也是人，但同时也是\Person{天主}……。\par

\SourceRecord{Translated by J.E. Tweed.；Nicene and Post-Nicene Fathers, First Series；Vol. 8.；Edited by Philip Schaff.；Buffalo, NY: Christian Literature Publishing Co.,；1888.；<http://www.newadvent.org/fathers/1801078.htm>.}

% structure: p0001 source=h1 target=section reason=page-title

\section{圣咏第七十九篇释义}

1. 这首诗篇的标题既简短又简单，我认为我们无需赘言。但我们在这里读到的预言已经明显应验。因为当这些事在大卫王时代被歌唱时，外邦人的敌意尚未对耶路撒冷城或天主的圣殿——那时甚至尚未建造——造成任何此类事件。因为大卫死后，他的儿子撒罗满为天主建造了圣殿，谁不知道呢？因此，那些在圣神中被视为未来的事，被说成好像已经过去。\par

\Quote{天主，外邦人进入了你的基业}\BibleRef{咏 78:1}。在这种表达形式下，其他将要发生的事也被说成已经完成。不必对此感到惊奇，这些话是对天主说的。因为并非向那不知道的天主陈述，而是由祂的启示预知；而是灵魂以那种天主所知的虔诚情感与天主交谈。因为即使是天使向人宣告的事，也是向不知道的人宣告；但他们向天主宣告的事，是向知道的天主宣告，当他们献上我们的祈祷，并以不可言喻的方式咨询永恒的真理——作为不变的法则——关于他们的行动。因此，这位天主的人将他将从天主学到的告知天主，如同学生向老师，并非不知而是判断；所以要么认可他所教导的，要么谴责他所未教导的：尤其因为先知在祈祷的外表下，将那些在这些事将要发生之时的人转化为他自己。但在祈祷中，习惯上将天主在报复中所做的事告知天主，并加上恳求，求祂今后怜悯和宽恕。同样，在此处，预言者说出审判，仿佛这些审判是由那些遭遇它们的人所说的，而这哀叹和祈祷本身就是一种预言。\par

2. \Quote{他们玷污了你的圣殿，使耶路撒冷成为苹果的贮藏所。}\Quote{他们将你仆人的尸首给天空的飞鸟作食物，将你圣徒的肉给地上的野兽}\BibleRef{咏 78:2}。\Quote{他们像水一样在耶路撒冷周围倾流他们的血，无人埋葬他们}\BibleRef{咏 78:3}。如果我们中有任何人认为在这预言中必须理解为提图斯罗马皇帝所造成的那次耶路撒冷毁灭——当时主耶稣基督在复活升天后已在万民中传扬——我不明白那个民族如何还能被称为天主的基业，因为他们不信从基督，拒绝并杀害了祂，那个民族成了被弃绝的，甚至在祂复活后也不信祂，还杀害了祂的殉道者。因为以色列民族中凡信仰基督的——基督的恩赐曾赐给他们，在某种程度上是应许的健康而多产的实现——关于他们，主自己说：\Quote{我奉差遣，只是到以色列家迷失的羊那里去}\BibleRef{玛 15:24}，这些人就是从他们中出来的应许之子；他们被算为后裔；\BibleRef{罗 9:8}他们属于天主的基业。从他们中来的有义人若瑟，以及生基督的童贞玛利亚：\BibleRef{玛 1:16}有新郎的朋友若翰洗者及其父母匝加利亚和依撒伯尔：\BibleRef{路 1:5}有西默盎老人\BibleRef{路 2:25}和亚纳寡妇，他们不是靠肉体的感觉听基督说话；而是当祂还是婴孩不说话时，靠圣神认出了祂：有蒙福的宗徒：有纳塔乃耳，在他内毫无诡诈：\BibleRef{若 1:47}有另一位若瑟，他自己也等候天主的国：有那大批群众在祂的牲畜前后来，喊着说：\Quote{奉主名来的，当受赞美：}\BibleRef{玛 21:9}其中也有那群孩童，祂宣称在他们身上应验了：\Quote{从婴孩和吃奶者的口中，你完备了赞美。}\有祂复活后的那些，一天三千人，另一天五千人受洗，被爱的火熔为一心一意，其中无人说任何东西是自己的，而是大家共有。\BibleRef{宗 4:32}有圣执事，其中斯德望在宗徒之前殉道。\BibleRef{宗 7:59}有犹大省那么多的教会，他们在基督内，保禄对他们只是闻名不见面，\BibleRef{迦 1:22}却因臭名远扬的残暴而出名，更因基督最仁慈的恩宠而出名。甚至有他本人，按照从前关于他的预言：\Quote{一只掠夺的狼，早晨抢掠，晚上分食；}\BibleRef{创 49:27}即起初作为迫害者抢掠致死，后来作为宣道者喂养致生命。这些人是从那个民族中天主的基业……因此即使在此时，因恩宠的拣选，也有残余得救。这个从那个民族中出来的残余属于天主的基业；而不是那些他在稍后说\Quote{其余的人被蒙蔽了}的人。因为他这样说：\Quote{那么，以色列所追求的，没有获得；但蒙拣选的人获得了；其余的人被蒙蔽了。}\BibleRef{罗 11:7}因此这蒙拣选的人，这残余，这天主的子民——天主没有弃绝的——被称为祂的基业。但在那个没有得到这事的以色列中，在那些被蒙蔽的其余人中，不再是天主的基业，关于他们，在基督在天上受荣耀之后，在提图斯皇帝的时代，有可能说出：\Quote{天主，外邦人进入了你的基业}，以及在这篇圣咏中似乎预言了关于那个民族之圣殿和城市毁灭的其他事。\par

3. 此外，在此我们应当或理解为在基督降生成人之前由其他敌人所行的事；那时甚至有圣先知，当巴比伦俘掳发生时，\BibleRef{列下 24:14}，那一民族遭受严重的苦难，以及在安提约古时期，玛加伯们忍受了可怕的痛苦，最荣耀地获得了冠冕。\EditorialNote{玛加伯下 第七章}或者，若在主复活升天之后，天主的产业必须被理解为在这里所谈论的；那么在此必须理解为：在拜偶像者和基督圣名之敌手中，祂的教会，在如此众多的殉道者中所忍受的事……。这个教会，这个天主的产业，是出于受割礼的和未受割礼的聚集而成的，也就是说，出于以色列人和其余的外邦人，借着那块被匠人弃而不用的石头，它成了屋角的基石；在这角石上，仿佛两堵从不同方向而来的墙联合在一起。因为祂是我们的和平，祂使两者合而为一，为要将两者在祂自己内造成一个新人，缔造和平，并将两者在一个身体内与天主联合：在这个身体内，我们是天主的子女，\BibleQuote{呼号：阿爸，父啊！}\BibleRef{罗 8:15}\OriginalTerm{阿爸}{Abba}是因他们的语言；\OriginalTerm{父啊}{Father}是因我们的语言。因为\OriginalTerm{阿爸}{Abba}与\OriginalTerm{父啊}{Father}相同……。\par

4. 但在紧接着的经文中，\BibleQuote{他们使耶路撒冷成为苹果的储藏所；}连教会本身也正确地以此名称被理解，即那自由的耶路撒冷，我们的母亲，\BibleRef{迦 4:26}关于她曾有记载：\BibleQuote{被遗弃者的儿子远多于有夫之妇的儿子。}\BibleRef{依 54:1}\Quote{成为苹果的储藏所，}这个说法，我认为必须理解为因迫害的荒废所造成的遗弃：就是说，如同苹果的储藏所；因为当苹果过去之后，苹果的储藏所就被废弃了。而且，当教会似乎因迫害她的外邦人而被遗弃时，殉道者的灵魂，如同许多极其甘甜的苹果从主的园中，被送往天上的宴席。\par

5. \BibleQuote{他们使，}他说，\BibleQuote{祢仆人的尸体成为天上飞鸟的食物，祢圣徒的肉体成为地上野兽的食物}\BibleRef{咏 78:2}。\Quote{尸体}一词在\Quote{肉体}中得到了重复；\Quote{祢的仆人}一词在\Quote{祢的圣徒}中得到了重复。唯有这一点有所不同，即\Quote{给天上的飞鸟，和地上的野兽}。那些写作\Quote{尸体}的人比那些写作\Quote{可朽的}的人诠释得更好。因为\Quote{尸体}只用于已死的人；而\Quote{可朽的}甚至用于活的身体。那么，正如我所说，当殉道者的灵魂如苹果般已往他们的农夫那里去时，他们将他们的尸首和肉体摆在天上飞鸟和地上野兽面前：仿佛其中任何一部分会从复活中失落，而祂，连我们的头发都被数过的，却要从自然界的隐秘之处恢复全部。\BibleRef{玛 10:30}\par

6. \BibleQuote{他们像水一样倾流了他们的血，}就是说，大量肆意地，\BibleQuote{在耶路撒冷的周围}\BibleRef{咏 78:3}。如果我们在此理解为地上的耶路撒冷城，我们就看见他们的血在城周围被倾倒，这些人是仇敌能在城墙外找到的。但如果我们理解为那座耶路撒冷，关于她曾有记载：\BibleQuote{被遗弃者的儿子远多于有夫之妇的儿子，}\BibleRef{依 54:1}，她的周围就是普世大地。因为在先知的那段经文中写有：\BibleQuote{被遗弃者的儿子远多于有夫之妇的儿子：}稍后对同一处说：\BibleQuote{那拯救你的，要被称为普世大地的以色列的天主。}\BibleRef{依 54:5}因此，这耶路撒冷在本圣咏中的周围必须理解如下：即那时教会已扩展，在普世世界结果实、成长，当时每一部分都有迫害在肆虐，并残害殉道者，他们的血像水一样被倾流，为天上的宝藏带来极大的增益。但至于附加的\BibleQuote{且没有人埋葬}，它要么不应看作是难以置信的事，即某些地方有如此大的恐慌，以至于没有任何圣洁尸体的埋葬者出现；要么，在许多地方，未埋葬的尸体可能躺卧很长时间，直到被虔诚者以某种方式偷走而得以埋葬。\par

7. \BibleQuote{我们成了，}他说，\BibleQuote{我们邻邦的羞辱}\BibleRef{咏 78:4}。因此，不在于人眼中宝贵的，这羞辱来自人，而是\BibleQuote{在主的眼中，祂的圣者的死亡是宝贵的。}\BibleQuote{讥讽和嘲笑：}或者，如一些人解释的，\BibleQuote{对我们周围之人的戏弄。}这是对前一句的重复。因为上面被称为\Quote{羞辱，}同样在\Quote{讥讽和嘲笑：}中得到了重复；上面所说的\Quote{对我们的邻邦，}同样在\Quote{对我们周围之人。}中得到了重复。此外，就地上的耶路撒冷而言，邻邦和那个民族的周围之人当然被理解为其他民族。但就那自由的耶路撒冷我们的母亲而言，\BibleRef{迦 4:26}在她的周围也有邻邦，在这些人中，作为她的仇敌，教会居住在普世世界的周围。\par

8. 其次，现在发出了一个明显的祈祷，从中可以看出，回忆先前的苦难并非为了告知，而是祈祷：\Quote{主啊，祢要愤怒到几时，直到终结吗？祢的嫉妒要如火焚烧吗？} \BibleRef{咏 78:5}。显然，他在求天主不要愤怒到终结，即如此大的压迫、苦难和蹂躏不要持续到终结；而是让祂节制祂的惩戒，正如另一篇圣咏所说：\Quote{祢用眼泪的饼喂养我们，给我们的眼泪有量有度。} 因为\Quote{主啊，祢要愤怒到几时，直到终结吗？}这句话的意思，如同说：主啊，求祢不要愤怒到终结。在接下来的\Quote{祢的嫉妒要如火焚烧吗？}中，两个词都必须理解，即\Quote{到几时}和\Quote{直到终结}：就好像是说：祢的嫉妒要如火焚烧直到终结吗？因为这两个词必须像前面稍高之处所用的那个词\Quote{他们做了}一样来理解。前一句有\Quote{他们使祢仆人的尸体成了天上飞鸟的食物}，而后一句说\Quote{祢圣者的肉体成了地上野兽的食物}，但前一句的\Quote{他们做了}显然应当理解为隐含的。\par

此外，天主的愤怒和嫉妒并非天主的情绪，如一些不理解圣经的人所指责的那样：但在愤怒的名义下，应当理解为惩罚不义；在嫉妒的名义下，应当理解为要求贞洁，以免灵魂藐视其主的律法，因离弃主而行淫而丧亡。这些在人的苦难的实际作用中是猛烈的；但在天主的安排中却是平静的，正如对祂所说的：\Quote{但是，主啊，祢以从容施行审判。} \BibleRef{智 12:18} 这些话语清楚表明，这些苦难是因罪恶而降临到人身上，即使他们是信者：尽管由此因他们的忍耐而绽放殉道者的荣耀，以及因虔敬地承受主的鞭策而显出纪律的轭。关于这一点，玛加伯人在剧烈的痛苦中，三青年在无害的火中 \BibleRef{达 3:21}，圣先知们在被掳中，都作了见证。虽然他们最勇敢、最虔敬地承受了父亲的管教，但他们并不隐瞒这些事是由于他们应得的罪过而临到他们的……\par

9. 但他接着说的\Quote{求祢将愤怒倾注在不认识祢的异邦，和没有呼求祢名的王国} \BibleRef{咏 78:6}，这也是预言，而非愿望。这些话不是出于恶意的诅咒，而是被圣神预见并预言：正如叛徒犹大的情况，将要降临在他身上的恶事，被预言得仿佛是愿望一样。因为正如先知以祈使语气说出\Quote{大能者啊，将祢的剑佩在腰间：在祢的尊荣和威严中，奋勇前进，胜利为王}，却不是命令基督，同样，说\Quote{求祢将愤怒倾注在不认识祢的异邦}的人，也不是愿望，而是预言。他照常重复说：\Quote{和没有呼求祢名的王国。} 因为异邦在王国中被重复了；不认祂在未呼求祂名中被重复了。那么，主在福音中关于\Quote{多和少}的鞭笞的话 \BibleRef{路 12:47} 该如何理解呢？难道是愤怒更强烈地倾注在不认识主的异邦吗？因为他说\Quote{倾注愤怒}，用这个词已经清楚表明他愿意理解愤怒有多大。随后他说：\Quote{向我们的邻人报七倍。} 难道这不表明在不知道主人旨意但仍呼求祂名的仆人，与远离如此伟大主人家庭的陌生人之间有很大区别吗？这些人如此不认识天主，以至于连天主也不呼求；反而呼求偶像、魔鬼或他们选择的任何受造物，而不是永远可称颂的造物主。他预言这些事涉及的那些人，他甚至没有暗示他们如此不知道他们天主的旨意，以至于仍敬畏主自己；而是如此不认识主自己，以至于不呼求祂，并且成为祂名的敌人。因此，在不知道他们天主旨意、但仍生活在祂家中和祂屋内的仆人，与不仅存心反对认识主自己、而且不呼求祂名、并且甚至在祂的仆人身上攻击祂的名的敌人之间，有很大的区别。\par

10. 最后，经文接着说：\Quote{因为他们吞吃了雅各伯，使他的地方荒凉}\BibleRef{咏 78:7}。……我们应如何理解雅各伯的\Quote{地方}，必须加以说明。因为雅各伯的地方可能是指那座城，那里也有圣殿，主曾命令整个民族聚集到那里献祭、敬拜并庆祝逾越节。但如果基督徒的聚会因迫害者而受阻碍、被压制，就是先知所要表达的意思，那么他似乎应该说\Quote{地方}（复数）荒凉，而不是\Quote{地方}（单数）。不过，我们可以把单数理解为复数，就像衣服指代服装，军队指代士兵，牲畜指代野兽：许多词语通常这样使用，不仅出现在俗语中，也出现在最受认可的权威的雄辩中。神圣的圣经本身也不排斥这种表达方式。因为就连她也把青蛙（单数）当作青蛙（复数），蝗虫（单数）当作蝗虫（复数），以及无数类似的表达。至于\Quote{他们吞吃了雅各伯}这句话，也很好理解，因为许多人被强行转入他们邪恶的身体，也就是他们的团体中。\par

11. ……他接着说：\Quote{求你不要记念我们祖辈的旧恶}\BibleRef{咏 78:8}。他没有说\Quote{往昔的}，那可能甚至是最近的；而是说\Quote{旧恶}，即从祖先而来的。因为这种邪恶应受审判，而非矫正。\Quote{愿你的仁慈迅速向我们迎面而来。}迎面而来，即在你审判时。因为\Quote{仁慈胜过审判}\BibleRef{雅 2:13}。而\Quote{不怜悯人的，审判也必不施怜悯}。但他又加上\Quote{因为我们已变得极其贫穷}：为此，他愿天主的仁慈被理解为迎面而来，好让我们的贫穷，即软弱，因祂的怜悯而得到扶助，以遵行祂的诫命，使我们不致受审判而被定罪。\par

12. 因此，接下来是：\Quote{天主，我们的医治者，求你援助我们}\BibleRef{咏 78:9}。藉着\Quote{我们的医治者}这个词语，他充分解释了他在前面所说的\Quote{因为我们已变得极其贫穷}所意指的贫穷是什么。因为正是那种疾病需要医治者。但当他希望我们得到援助时，他既不对恩宠忘恩负义，也不取消自由意志。因为得到援助的人，他自己也做了某事。他又加上：\Quote{为你的圣名的光荣，主啊，求你拯救我们}，好叫那夸耀的，不是夸耀自己，而是夸耀主\BibleRef{格前 1:31}。\Quote{求你宽恕我们的罪过}，他说，\Quote{为你的圣名}，而不是为我们自己。因为我们的罪过除了应得和相当的惩罚之外，还配得什么？但\Quote{求你为你的圣名宽恕我们的罪过}。因此，你拯救我们，即你从恶事中解救我们，既扶助我们行正义，又怜悯我们的罪过——没有罪过，我们在今生无法存活。因为\Quote{在你面前，凡有血肉的都不能成义}。但罪就是不义。\Quote{你若究察罪孽，谁能站得住呢？}\par

13. 但他所添加的：\Quote{恐怕外邦人常说：他们的天主在哪里？}\BibleRef{咏 78:10}，必须理解为对外邦人本身说的。因为对真天主绝望的人，认为祂不存在或不顾念祂的子民、不怜悯他们，会有不好的结局。至于接下来这句话：\Quote{愿你在我们眼前，在万民中为你仆人被流之血伸冤}，可以理解为当那些曾迫害祂产业的人相信真天主的时候；因为这也是复仇，即他们凶残的不义被天主之言的剑所击杀，正如经上所说：\Quote{束上你的剑}；或者当顽固的仇敌最终受罚的时候。因为他们在今世所遭受的身体上的痛苦，可能与好人共享。还有另一种复仇：当教会经过如此多的迫害——迫害者以为她会彻底灭亡——之后，在世上扩展和结果，罪人、不信者和仇敌看见了，就愤怒；\Quote{他必咬牙切齿，渐渐消灭}。谁敢否认这甚至是最重的惩罚呢？但我不知道他所说的\Quote{在我们眼前}，如果藉着这种惩罚理解为在内心隐秘处进行的、甚至折磨那些对我们微笑的人的事，是否足够优雅；因为我们看不见他们在内里所受的苦。但无论是在他们相信时他们的不义被消灭，还是在他们坚持邪恶时最后惩罚临到他们，这句话\Quote{愿我们在万民中看见复仇}的意思都不难无疑地理解。\par

14. 诚然，如我们所说，这是预言，而非祈愿……主在福音中\BibleRef{路 18:3}为我们立了那位寡妇的榜样，她渴望伸冤，便向不义的法官求情；法官最终听了她，并非出于正义，而是因为厌倦。但主以此教导我们，表明那位公义的天主更会速速为他的选民伸冤，因为他们昼夜向他呼号。同样，在祭台下的殉道者也呼喊\BibleRef{默 6:9}，求天主在审判中为他们伸冤。那么，\Quote{爱你们的仇敌，善待恨你们的人，为迫害你们的人祈祷}\BibleRef{玛 5:44}又在哪里？\Quote{不以恶报恶，不以辱骂还辱骂}\BibleRef{伯前 3:9}和\Quote{不以恶报恶}\BibleRef{罗 12:17}又在哪里？……因为主劝我们爱仇敌时，以我们在天之父的榜样教导我们，\Quote{他使太阳上升，光照善人，也光照恶人；降雨给义人，也给不义的人}\BibleRef{玛 5:45}。难道因此他就不以暂时的惩罚来管教，或最终不审判那些顽固不化的人吗？因此，爱仇敌应如此：他受罚时天主的公义不应使我们不悦，而应使我们喜悦，但这喜悦不是因他的灾祸，而是因良善的审判者。而恶毒的灵魂会因仇敌改正而逃脱惩罚而悲伤；当他见仇敌受罚时，便因伸冤而高兴，并非因爱慕天主的公义而喜悦，而是因他所恨之人的痛苦而高兴；当他将审判交托天主时，他希望天主施以比他所能施的更重的惩罚；当他给饥饿的仇敌食物、给口渴的仇敌水喝时，他怀着恶意理解经上所载：\Quote{你这样行，就是把炭火堆在他头上}\BibleRef{罗 12:20}。……因此，在这篇圣咏中，祈求者表面上是请求，实际是在预言对恶人的将来惩罚；我们应当明白，天主圣人们爱了他们的仇敌，不愿任何人遭受不善，即今生的敬虔和来世的永生；但在恶人的惩罚中，他们喜悦的不是恶人之祸，而是天主的善判；并且在圣经中凡读到他们对人的憎恨，那都是对恶习的憎恨，而每个人若爱自己，就必须在自己身上憎恨这些恶习。\par

15. 至于随后的话，\Quote{愿被囚者的哀叹上达于你面前}，或如有些抄本所记，\Quote{在你眼前，被囚者的哀叹}：很少有人能发现圣人们曾被迫害者投入锁链；如果在如此众多、多样的刑罚中发生此事，也极为罕见，因此不能相信先知在此节特意指此。实际上，锁链指身体的软弱和可朽性，它压沉重灵魂。由于肉身的脆弱，如同某种痛苦和困扰的材料，迫害者能迫使她趋向不敬。宗徒渴望脱离这些锁链，与基督同在\BibleRef{斐 1:23}；但为了他所传福音之人的缘故，他仍需留在肉身内。直到这必朽坏的穿上不朽坏的，这必死的穿上不死的\BibleRef{格前 15:54}，软弱肉身如同锁链般阻碍了情愿的精神\BibleRef{玛 26:41}。这些锁链并非人人都感觉到，只有那些在自身内叹息负重、渴望穿上来自天上的居所的人\BibleRef{格后 5:4}才会感到，因为死亡令人恐惧，而现世生命充满忧伤。先知为这些叹息之人加倍叹息，愿他们的叹息\Quote{上达于主面前}。也可以理解为那些被智慧诫命束缚的人，他们耐心承受，这些链便化为装饰；因此经上说：\Quote{把你的脚放在她的锁链中}\BibleRef{德 6:24}。他说：\Quote{按你手臂的伟大，收纳被处死者的子孙为义子}，或如某些抄本所载：\Quote{藉受罚者的死，拥有儿子}。在我看来，圣经在这里充分显示了被囚者的哀叹是什么——他们为基督之名忍受了最残酷的迫害，这些迫害在此圣咏中清楚地预言了。他们遭受各种苦难时，常为教会祈祷，愿他们的血不致于为后代徒劳；好使仇敌以为教会会因此灭亡，而主的庄稼反因而更加丰盛。因为\Quote{被处死者的子孙}指那些不仅不为先前的苦难所吓倒，反而因知晓他们为主之名受难而受其荣耀激励，以庞大的人群信从了主的人。因此他说：\Quote{按你手臂的伟大}。因为在基督子民身上发生了如此大的奇迹，以致那些以为迫害教会便能得利的人，绝不相信会有这样的事。\par

16. 他说：\Quote{归还}，\Quote{给我们的近人七倍于他们怀里} \BibleRef{咏 78:13}。他不是在祈愿任何恶事，而是在预言并宣告将要发生的公道。至于数字七，即七倍的报复，他要人看到惩罚的圆满，因为通常以这个数字来表示圆满。故此也有对善人的这句话：\Quote{他在此世将得到七倍}，这代表了全部。\BibleQuote{似乎一无所有，却拥有万有。} \BibleRef{格后 6:10}。他说的是近人，因为教会居住在他们中间，直到分离的日子，因为现在还没有进行肉体的分离。他说\Quote{入他们怀里}，意思是现在在隐秘中，以致于在此生隐秘中执行的报复，日后要在列国面前显明于我们眼前。因为当一个人被交付于邪僻的心思时，他在内心的深处正在接受他将来应得的惩罚。\Quote{他们凌辱了祢的羞辱，上主啊。}祢就将七倍报应于他们的怀里，也就是说，为这羞辱，祢在他们的隐秘处最完全地斥责他们。因为他们如此羞辱了祢的圣名，以为可以通过祢的仆人将祢从地上抹去。\par

17. \BibleQuote{但我们——祢的子民} \BibleRef{咏 78:14}，必须普遍理解为一切虔诚的真正基督徒的族类。\Quote{我们}，就是那些他们认为有能力消灭的人，\Quote{祢的子民，祢牧场的羊群}，为使夸耀的，在主内夸耀 \BibleRef{格前 1:31}，\BibleQuote{将向祢称谢，直到一代。}但有些抄本作\BibleQuote{将向祢称谢，直到永远。}这一分歧源于希腊文的歧义。因为希腊文作\OriginalTerm{希腊语}{\GreekText{εἰς} \GreekText{τὸν} \GreekText{αἰῶνα}}，既可解释为\Quote{直到永远}，也可解释为\Quote{直到一代}；但我们必须根据上下文来理解哪一个解释更好。依我看来，这段经文的意思显示，我们应当说\Quote{直到一代}，即直至今世的终结。但接下来的诗句，依照圣经、尤其是圣咏的惯例，是前一句的重复，只是顺序改变，把在前句中置后的放在前面，而将置前的放在后面。因为前一节曾说\BibleQuote{我们向祢称谢}，此处却说\BibleQuote{我们将宣扬祢的赞美}。同样，前一节曾说\Quote{直到一代}，此处却说\Quote{直到世世代代}。因为世代的重复表明永恒；或者，如一些人理解的，是因为有两个世代，旧的和新的……但在圣经的许多地方，我们已经向你们指出，\Quote{宣认}也用来表示赞美：正如在此处经文，\BibleQuote{你应当在宣认中说这些话语：\InnerQuote{上主的作为极其美妙。}} \BibleRef{德 39:33} 尤其是救主亲自所说的，祂毫无罪过，无需悔改去认罪：\BibleQuote{父啊，天地的主宰，我称谢祢，因为祢将这些事瞒住了智慧及明达的人，而启示给了小孩子。} \BibleRef{玛 11:25} 我说这些，是为了更清楚地看出，在\BibleQuote{我们将宣扬祢的赞美}这一表达中，重复了上文所说的\BibleQuote{我们向祢称谢}。\par

\SourceRecord{Translated by J.E. Tweed.；Nicene and Post-Nicene Fathers, First Series；Vol. 8.；Edited by Philip Schaff.；Buffalo, NY: Christian Literature Publishing Co.,；1888.；<http://www.newadvent.org/fathers/1801079.htm>.}

% structure: p0001 source=h1 target=section reason=page-title

\section{圣咏第八十篇释义}

1. 若偶然有隐晦之事需要解释者的职责，那么那些明显的事理应当要求我履行读者的职责。此歌是关于主的来临、我们的救主耶稣基督以及祂的葡萄园的。但此歌的歌唱者是那位\Person{阿撒夫}，就表面看来，他已蒙光照并悔改，由他的名字你们可知会堂被象征。最后，圣咏的标题是：\Quote{为终结，关乎那将改变的人}；意即，为那更好者。因为基督是法律的终结，\BibleRef{罗 10:4}祂特意来临，为要改变人，使成为更好。祂又加上\Quote{为阿撒夫自己的证言}。这是真理的良善证言。最后，这证言既承认基督，也承认葡萄园；即元首与身体，君王与百姓，牧人与羊群，以及全部圣经的整个奥秘——基督与教会。但圣咏的标题以\Quote{为亚述人}作结。亚述人被解释为\Quote{引导的人}。因此，这不再是那心尚未指引的后代，而是现今引导的后代。因此，我们当听他在这证言中所说的。\par

2. 什么是\BibleQuote{以色列的牧者，请留心听，引领若瑟如羊群的祢}？\BibleRef{咏 79:1}祂被呼求来临，祂被期待直到来临，祂被切望直到来临。因此愿祂找到\Quote{引导的人}：祂说：\Quote{引领若瑟如羊群的祢}。若瑟自己如羊群，若瑟自己也是一只羊。请注意若瑟；因为即使他的名字的诠释也大大帮助我们，因为它意指增加；祂确实来临，为要使那落入土中的子粒得以复活，结出许多子粒来；\BibleRef{若 12:24}即，为使天主的子民增加。⋯⋯\BibleQuote{坐在革鲁宾之上的祢}。革鲁宾是天主光荣的宝座，被诠释为知识的圆满。天主坐在知识的圆满之中。虽然我们理解革鲁宾是天上崇高的权能与德能；然而，若你愿意，你将成革鲁宾。因为若革鲁宾是天主的宝座，请听圣经所说：\Quote{义人的灵魂是智慧的宝座。}你如何说：我将成为知识的圆满？谁来完成这事？你具有完成它的方法：\BibleQuote{法律的圆满是爱德。}\BibleRef{罗 13:10}不要追求许多事而费心劳神。枝干的宽大令你害怕：抓住根，不要想树的高大。你们中若有爱德，知识的圆满必然随之而来。因为那认识爱德的人有何不知道呢？既然经上说过：\BibleQuote{天主是爱。}\BibleRef{若一 4:8}\BibleQuote{显现罢。}\BibleQuote{在厄弗辣因、本雅明和默纳协面前。}\BibleRef{咏 79:2}我说，在犹太民族面前，在以色列百姓面前显现。因为有厄弗辣因，有默纳协，有本雅明。但让我们看诠释：厄弗辣因是结果实的，本雅明是右手之子，默纳协是遗忘者。求祢在结出果实者面前显现，在右手之子面前显现；求祢在遗忘者面前显现，使他不再遗忘，而祢—曾解救他者—进入他的思念。⋯⋯因为祢软弱，当有人说：\BibleQuote{如果祂是天主子，就让他从十字架上下来吧。}\BibleRef{玛 27:40}祢似乎没有能力；迫害者凌驾于祢之上；祢早前已显示这事，因为雅各伯本人也在角力中获胜，一个人与天使角力。他何时能获胜，除非天使愿意？人获胜了，天使被战胜了；胜利的人抓住天使，说：\BibleQuote{除非祢祝福了我，我不放祢走。}\BibleRef{创 32:26}伟大的圣事！祂被战胜，却祝福了战胜者。被战胜，因为祂愿意；在肉身上软弱，在尊威中刚强。⋯⋯祢在软弱中被钉死，在能力中复活：\BibleRef{格后 13:4}\Quote{振起祢的能力，来拯救我们。}\par

3. \BibleQuote{天主，求祢使我们回转。}因为我们背弃了祢，除非祢使我们回转，我们将不会回转。\BibleQuote{求祢彰显祢的慈颜，我们便会得救。}\BibleRef{咏 79:3}祂的面容有任何暗晦吗？祂没有暗晦的面容，但祂在它前面放置了肉身的云彩，如同软弱的帕子；当祂悬在树上时，祂并未被认作祂以后坐在天上时所将被认出的那一位。因为事情是这样发生的。基督活在世上，行奇迹时，阿撒夫不认识祂；但当祂死了，复活了，升天之后，他认识祂了。他因心被刺透，他可能也在诉说关于祂的这证言，即我们现在在这圣咏中所承认的。祢遮蔽了祢的面容，我们便病了：求祢彰显同一面容，我们便会痊愈。\par

4. \Quote{万军的上主，祢向祢仆人的祈祷恼怒要到几时？}\BibleRef{咏 79:4}。现在是祢的仆人。祢曾恼怒仇人的祈祷，还要向祢仆人的祈祷发怒吗？我们已经归化，认识祢，祢还向祢仆人的祈祷发怒吗？祢显然会发怒，其实，如同父亲纠正，而非审判者定罪。祢显然会这样发怒，因为经上记载：\Quote{我儿，当你前去服事上主时，应立在正义与敬畏中，预备你的灵魂接受试探。}\BibleRef{德 2:1}。不要以为现在天主的忿怒已经过去，因为你已经归化。天主的忿怒已离开你，但只是不永远定罪。然而祂鞭打，并不宽容，因为祂鞭打所收纳的每一个儿子。\BibleRef{希 12:6}。你若不愿受鞭打，为何渴望被收纳？祂鞭打所收纳的每一个儿子。祂连自己的独生子也没有宽容，却鞭打每一个人。但即便如此，\Quote{祢向祢仆人的祈祷恼怒要到几时？}不再是你的仇人：而是\Quote{祢向祢仆人的祈祷恼怒}，要到几时？下文是：\Quote{祢用眼泪作食物喂养我们，又用大量泪水给我们解渴}\BibleRef{咏 79:5}。何为\Quote{大量}？请听宗徒：\Quote{天主是忠信的，祂不许你们受试探超过你们所能承受的。}\BibleRef{格前 10:13}。这量是按你们的能力；这量是为训练你们，而不是压垮你们。\par

5. \Quote{祢使我们成为邻国争论的对象}\BibleRef{咏 79:6}。这事显然发生了：因为从亚萨中被拣选的人，正是那些去向外邦人宣讲基督的人，他们曾听见人说：\Quote{这宣讲新鬼神的是谁？}\Quote{祢使我们成为邻国争论的对象。}因为他们宣讲那位本身就是争论对象的基督。他们宣讲什么？宣讲基督死后复活。谁会听？谁会知道？这是新事。但神迹相随，奇迹使难以置信的事变得可信。祂遭到反驳，但反驳者被征服，从反驳者变为信徒。那里，却有大火焰：殉道者们以眼泪作食物，大量饮用泪水，但有限量，不超过他们所能承受的；为的是在眼泪的量之后，有喜乐的冠冕。\Quote{我们的仇敌曾讥笑我们。}那些讥笑的人如今在哪里？长久以来人们说：谁在崇拜死者？谁在敬拜被钉者？长久如此。讥笑者的鼻子在哪里？如今，责难者岂不躲进洞穴，以免被看见？但你们看见下文：\Quote{万军的上主，求祢转变我们，显示祢的慈颜，我们便会得救}\BibleRef{咏 79:7}。\Quote{祢从埃及移来一棵葡萄树，祢驱逐了外邦，把她栽种了}\BibleRef{咏 79:8}。这事成就了，我们知道。多少外邦被驱逐？阿摩黎人、赫特人、耶步斯人、基尔加斯人、希威人。驱逐并推翻他们之后，从埃及被拯救的子民被领进应许之地。葡萄树从何处被驱逐，又栽种在何处，我们已经听见。让我们看看接下来发生了什么事，她如何相信，如何生长，覆盖了多大的土地。\par

6. \Quote{祢在她眼前开了一条路，栽种了她的根，她就充满了大地} \BibleRef{咏 79:9}。除非在她眼前开了一条路，她怎能充满大地呢？在她眼前所开的路是什么？祂说：\Quote{我是}，祂说，\Quote{道路、真理、生命} \BibleRef{若 14:6}。她充满大地，这话说得有理。这葡萄园如今已经实现了末后的事。但在这期间呢？\Quote{她的荫影遮蔽了群山，她的枝条遮蔽了天主的香柏} \BibleRef{咏 79:10}。\Quote{祢伸展了她的枝子直到海，她的嫩枝直到河} \BibleRef{咏 79:11}。这需要注释者的职责，读者和赞美者是不够的：请以注意帮助我；因为在这篇圣咏中提及这葡萄园，常使疏忽大意者陷入黑暗……但无论如何，最初的犹太民族就是这葡萄树。但犹太民族统治直到海和河。直到海：在圣经中 \BibleRef{户 34:5} 显示海在其附近。直到约但河。因为约但河那边有一部分犹太人定居，但约但河内是整个民族。因此，\Quote{直到海直到河}是犹太人的国度，以色列的国度：但不是\Quote{从海直到海，从河直到地极}；这是葡萄园未来的完美，他在这里预言了。我说，当他向你预言了完美，他回到开始，完美由此而出。你要听开始吗？\Quote{直到河。}你要听终结吗？\Quote{祂将统治从海直到海：}即\Quote{她充满了大地。}让我们看看阿撒夫的证言，关于第一个葡萄园发生了什么，以及第二个葡萄园必须期待什么，不，同一个葡萄园……那么，在她眼前开了路以便充满大地的葡萄园，起初在哪里？\Quote{她的荫影遮蔽了群山。}群山是谁？先知。为什么她的荫影遮蔽了他们？因为他们隐晦地讲论预言将要来的事。你从先知听到：守安息日，第八天给孩子行割损，献公羊、牛犊、山羊的祭。不要困扰，她的荫影遮蔽了天主的群山；在荫影之后将有显现。\Quote{她的灌木遮蔽了天主的香柏}，即她遮蔽了天主的香柏；非常崇高，但属于天主。因为香柏是骄傲者的预表，必须被推翻。\Quote{黎巴嫩的香柏}，世界的崇高，这葡萄树在生长中遮蔽了它们，以及天主的群山，所有圣先知和圣祖。\par

7. 然后呢？\Quote{祢为何拆毁了她的篱笆？} \BibleRef{咏 79:12}。现在你们看见那个犹太民族的覆亡：你们已经从另一篇圣咏中听到：\Quote{用斧与锤他们击倒了她。}这怎么可能发生，除非她的篱笆已被拆毁？她的篱笆是什么？她的防御。因为她向她的栽植者自高自大。被派到她那里的仆人，向她索取报酬，园丁们就鞭打、击打、杀害他们；独生子也来了，他们说：\Quote{这是继承人；来，我们杀他，产业就成为我们的了。}他们杀了他，把他扔出葡萄园。他被扔出去后，却更圆满地拥有了他被扔出的地方。因为主通过依撒意亚如此威胁她：\Quote{我要拆毁她的篱笆。}为什么？\Quote{因为我指望她结葡萄，她却结出荆棘。} \BibleRef{依 5:2} 我从那里指望果实，却找到了罪。那么，阿撒夫，你为什么问：\Quote{祢为何拆毁了她的篱笆？}难道你不知道为什么？我指望她行正义，她却行不义。难道她的篱笆不该被拆毁吗？当篱笆被拆毁时，外邦人来了，葡萄园受到攻击，犹太王国被抹去。起初他悲叹这事，但并非没有希望。因为他现在是在论及心灵的导向，即为了\Quote{亚述人}，为了\Quote{导向的人}，这篇圣咏。\Quote{祢为何拆毁了她的篱笆，使所有过路的人都摘取她的葡萄？}什么是\Quote{过路的人}？暂时掌权的人。\par

8. \Quote{来自树林的野猪使她荒凉} \BibleRef{咏 79:13}。在来自树林的野猪中，我们理解什么？对犹太人来说，猪是不洁之物，在猪身上他们仿佛想象外邦人的不洁。但推翻犹太民族的是外邦人；但那位推翻的君王，不仅是不洁的猪，而且还是野猪。因为野猪是什么？不过是凶猛的猪，狂怒的猪。\Quote{来自树林的野猪使她荒凉。}\Quote{来自树林，}即来自外邦人。因为她本是葡萄园，而外邦人是树林。但当外邦人信了时，说了什么？\Quote{那时，林中所有的树木都要欢呼。}\Quote{来自树林的野猪使她荒凉；一只\OriginalTerm{独特的野兽}{singular wild beast}吞噬了她。}\Quote{一只独特的野兽}是什么？就是那使她荒凉的野猪，就是那独特的野兽。独特，因为骄傲。因为每个骄傲者这样说：是我，是我，没有别人。\par

9. 但这有什么益处呢？\BibleQuote{万军的天主，求你常以慈悲的目光垂视}\BibleRef{咏 79:14}。虽然这些事已经发生，\BibleQuote{但求你垂视}。\BibleQuote{求你从天上垂看，眷顾这葡萄园}。\BibleQuote{完善你右手所栽种的她}\BibleRef{咏 79:15}。不要栽种别的，只使她成全。因为她就是亚巴郎的后裔，就是那使万民蒙受祝福的后裔：\BibleRef{创 22:18}，那里是根，野橄榄枝被接在其上。\BibleQuote{完善你右手所栽种的这葡萄园}。但他在哪方面使她完善呢？\BibleQuote{也完善你为自己所坚固的人子}。还有什么比这更清楚的吗？你们为什么还期待我们用言语向你们解释，不如和你们一同赞叹说：\Quote{完善你右手所栽种的这葡萄园，并完善人子}使她完善？哪一个人子？就是\BibleQuote{你为自己所坚固的那一位}。一座坚固的堡垒：尽力建造吧。\BibleQuote{因为除已奠立的基础，即耶稣基督外，任何人不能再奠立别的}\BibleRef{格前 3:11}\par

10. \BibleQuote{被火烧过、又被掘起的事物，因你面容的斥责而灭亡}\BibleRef{咏 79:16}。让我们看看并思考什么是被火烧过和掘起的事物。基督斥责了什么？罪：因他面容的斥责，罪已灭亡。那么，为什么罪是被火烧过和掘起呢？一切罪，在人内有二个原因：欲望和恐惧。想一想，省察你们的心，筛你们的良心，看看是否可能有罪，除非要么出于欲望，要么出于恐惧。摆在你面前一个报酬引诱你犯罪，即一件令你喜欢的事物；你做了，因为你欲望它。但或许你不会被贿赂诱惑；你用威胁恐吓，你做了，因为你恐惧。例如，有一个人贿赂你做假见证。例子无数，但我摆出更清楚的例子，使你们能想象其余的。你听从了天主，且在心中说：\BibleQuote{人纵然赚得了全世界，却赔上了自己的灵魂，为他有什么益处？}\BibleRef{玛 16:26}。我不被贿赂诱惑，为了钱财而丧失灵魂。他转而激起你内心的恐惧，那不能用贿赂腐蚀你的人，开始威胁损失、流放、屠杀、或许还有死亡。那么，如果欲望没有得逞，或许恐惧会使你犯罪。……邪恶的恐惧做了什么？它仿佛掘起了什么。因为爱燃烧，恐惧谦卑：因此，邪恶之爱的罪被火烧起；邪恶之恐惧的罪被掘起。一方面，邪恶的恐惧谦卑，而善的爱光照；但各自以不同的方式。因为农夫为树求情，不要砍倒它，说：\BibleQuote{我遂要挖松它周围的土，再施上粪}\BibleRef{路 13:8}。挖的沟代表一个恐惧者的虔诚谦卑，粪篮代表一个忏悔者的有益卑污状态。但关于善爱之火，主说：\BibleQuote{我来是把火投在地上}\BibleRef{路 12:49}。愿那火点燃心灵炽热的人，也为那被天主和邻人之爱所燃烧的人。如此，正如一切善工由善的恐惧和善的爱而成，同样，一切罪由恶的恐惧和恶的爱而犯。因此，\Quote{被火烧过和掘起的事物}，即一切罪，\Quote{因你面容的斥责而灭亡}。\par

11. \BibleQuote{愿祢的手扶持祢右手的人，并扶持祢所坚固的人子。} \BibleRef{咏 79:17} \BibleQuote{我们决不离开祢……祢将使我们复生，我们必呼求祢的名。} \BibleRef{咏 79:18} 祢将对我们成为甘饴，\BibleQuote{祢将使我们复生。} 因为以前我们爱的是大地，不是祢；但祢已使我们在地上的肢体死去了。 \BibleRef{哥 3:5} 因为旧约带有地上的应许，似乎劝勉人不要白白地爱天主，而应因天主赐予地上某物而爱祂。你爱什么，以至于不爱天主？告诉我。你若能爱任何祂所没有造的东西，就爱吧。环视整个受造界，看看是否你在何处被欲望的粘胶所缠住，而受阻于爱造物主，除非是因你所轻视的祂自己所造之物。但你为何爱那些东西？不正是因为它们美丽吗？它们能像创造它们的那位那样美丽吗？你因看不见祂而欣赏这些事物；但透过你所欣赏的这些事物，爱那你看不见的祂。审视受造物：若是出于自己，就停留在其中；但若是出于祂，那么它对于爱者而言之所以有害，无非是因为它被置于造物主之上。我为什么说这些？是为了这句圣咏，弟兄们。我说，那些为肉身的利益而敬拜天主的人是死的：\BibleQuote{因为随从肉性的思念是死亡。} \BibleRef{罗 8:6} 而不无偿地敬拜天主，即因祂本身是善、而非因祂赐予某些善物（甚至赐予不善的人），这些人也是死的。你愿从天主得到金钱？连强盗也有。妻子、子女众多、身体健康、世俗尊荣，且看多少恶人拥有这些。难道你就是为了这些而敬拜祂？你的脚将动摇，当你看见那些不敬拜祂的人也拥有这些时，你会以为自己敬拜毫无理由。我再说，这一切祂甚至赐予恶人，唯独祂自己为善人保留。\BibleQuote{祢将使我们复生；} 因为我们原是死的，当我们依恋地上的事物时；我们是死的，当我们带着属土的人肖像时。\BibleQuote{祢将使我们复生；} 祢将更新我们，赐给我们内在之人的生命。\BibleQuote{我们必呼求祢的名；} 也就是说，我们将爱祢。祢将是我们罪过的甜蜜赦免者，祢将是成义者的全部赏报。\BibleQuote{上主，万军的天主，求祢使我们复兴，显示祢的慈颜，我们便会得救。} \BibleRef{咏 79:20}\par

\SourceRecord{Translated by J.E. Tweed.；Nicene and Post-Nicene Fathers, First Series；Vol. 8.；Edited by Philip Schaff.；Buffalo, NY: Christian Literature Publishing Co.,；1888.；<http://www.newadvent.org/fathers/1801080.htm>.}

% structure: p0001 source=h1 target=section reason=page-title

\section{圣咏第八十一篇注解}

1. 这首圣咏的标题是：\Quote{致终末，为\OriginalTerm{压榨机}{presses}，在\OriginalTerm{安息日第五日}{fifth of the sabbath}，\Person{阿撒夫}本人的圣咏。}在这一个标题中，许多奥秘被堆积在一起，以至于圣咏的门楣指明了其中所包含的事。既然我们要谈论压榨机，就不要期望我们会谈论油槽、压榨机、橄榄筐；因为圣咏中并没有这些，因此它指明了更大的奥秘……\par

你们在诵读时并没有听到这类事。因此，把压榨机当作教会现在正在进行的奥秘。在压榨机中，我们观察到三样东西：压力，以及压力的两个结果：一个要收藏起来，另一个要丢弃。于是，在压榨机中发生了踩踏、压碎、重压：借此，油秘密地滤入油槽，渣滓公开地流入街道。\par

请专注地观看这伟大的景象。因为天主的作为从未停止向我们展示那些我们可以以极大喜悦观看的事，竞技场的疯狂也不能与这景象相比。那属于渣滓，这属于油。因此，当你们听到亵渎者无耻地胡言乱语，说苦难在基督徒时代增多——因为你们知道他们爱说这话——这是一句古老的谚语，却是从基督徒时代开始的：\Quote{天主不降雨；算在基督徒头上！}虽然那是古代人这样说的。但现在他们也说：\Quote{天主降雨，算在基督徒头上！天主不降雨，我们不播种；天主降雨，我们不收割！}他们故意将此作为骄傲的借口，而本应使他们更恳切地祈祷，他们却宁愿亵渎而不愿祈祷。\par

因此，当他们谈论这些事，当他们如此夸口，当他们说这些话，而且说得狂妄，不是出于恐惧，而是出于傲慢时，不要让他们扰乱你。因为假定压力增多了；你要作油。让那因无知黑暗而发黑的渣滓傲慢无礼；让它像被扔在街道上一样，公开地嘲弄：但你独自在你心中，在那里那在暗中观看的必会报答你，你滤入油槽中。\par

……只举一例，连那些制造苦难的人也在抱怨的事：他们说，在我们这个时代，有多少掠夺，有多少无辜者的苦难，有多少抢夺他人财物的事！这样，你们确实注意到了渣滓，即他人的财物被夺取；你们却不留意油，即连自己的财物也给了穷人。古代没有这样的抢夺他人财物的人：但古代也没有这样的施舍自己财物的人……\par

2. 为什么又是\Quote{在安息日第五日}？这是什么意思？让我们回到天主的起初作为，或许我们会在那里找到一些东西，也能理解奥秘。因为安息日是第七天，在这一天主\Quote{停止了祂一切的工作} \BibleRef{创 2:2}，预示着我们将来停止一切工作的伟大奥秘。因此，安息日的第一日被称为第一日，我们也称之为主日；安息日的第二日是第二日；……而安息日本身是第七日。所以请看，这首圣咏是对谁说的。因为在我看来，它是对领受圣洗者说的。因为在第五天，天主从水中创造了动物：在第五天，也就是\Quote{安息日第五日}，天主说：\BibleQuote{让水生出有生命灵魂的爬虫。} \BibleRef{创 1:20} 所以请看，你们这些水已生出有生命灵魂的爬虫的人。因为你们属于压榨机，在你们身上，水所生出的，一样被滤出，另一样被丢弃。因为有许多人生活不配他们所领受的圣洗。因为有多少领受圣洗者选择去填满竞技场而不是这座大殿！有多少领受圣洗者要么在街上搭棚，要么抱怨他们没有搭棚！\par

但是，这首圣咏，\Quote{为\OriginalTerm{压榨机}{presses}}，和\Quote{在安息日第五日}，是\Quote{向\Person{阿撒夫}}歌唱的。\Person{阿撒夫}是一个名叫此名的人，如同\Person{耶杜通}、\Person{科辣黑}，以及我们在圣咏标题中发现的其他名字：然而，这名字的解释暗示了一个隐藏真理的奥秘。事实上，阿撒夫在拉丁文中被解释为\Quote{聚会}。因此，\Quote{为压榨机，在安息日第五日}，圣咏\Quote{向\Person{阿撒夫}}歌唱，即，为一种区分性的压力，对由水重生的领受圣洗者，这首圣咏是向主的聚会歌唱的。我们已经阅读了门楣上的标题，并理解了这些\Quote{压榨机}的意义。现在，如果你们愿意，让我们看看这作品的内部，即压榨机的内部。让我们进入，观察，欢喜，恐惧，渴望，躲避。因为当我们将开始诵读，并在主的帮助下，讲论祂所赐予我们的时，你们将在这一内部房屋中，即在圣咏本身的经文中，找到所有这一切。\par

3. 看，你们自己，\Person{阿撒夫}，主的聚会。\BibleQuote{要向天主我们的助佑者欢跃} \BibleRef{咏 80:1}。你们今天聚集在一起，你们今天就是主的聚会，如果这首圣咏确实是向你们唱的：\BibleQuote{要向天主我们的助佑者欢跃}。别人向竞技场欢跃，你们向天主欢跃；别人向他们的欺骗者欢跃，你们向你们的助佑者欢跃；别人向他们的神——肚腹欢跃，你们向你们的天主——你们的助佑者欢跃。\BibleQuote{要向雅各伯的天主欢呼。} \BibleRef{创 25:23} \BibleQuote{要向雅各伯的天主欢呼。}因为你们也属于雅各伯：是的，你们就是雅各伯，年幼的、年长的服事他的民族。\BibleQuote{要向雅各伯的天主欢呼。} \BibleRef{创 25:23}\par

4. \Quote{拿起圣咏，将手鼓交给祂} \BibleRef{咏 80:2}。既是\Quote{拿起}，又是\Quote{交给}。什么是\Quote{拿起}？什么是\Quote{交给}？\Quote{拿起圣咏，将手鼓交给祂。}宗徒保禄在某处 \BibleRef{斐 4:15} 责备并痛心地说，在给予和收取的事上，没有人与他相通。什么是\Quote{在给予和收取的事上}？岂不是他在另一处 \BibleRef{格前 9:11} 所明示的：\Quote{我们若把属神的事撒给你们，若从你们收取属血肉的事，还算大事么？}确实，手鼓因用皮制成，属于血肉。因此，圣咏是属神的，手鼓是属血肉的。所以，天主的百姓，天主的会众，\Quote{请你们拿起圣咏，将手鼓交给祂}：拿起属神的事，交给属血肉的。这也是我们在那位蒙福殉道者的墓前所劝勉你们的：领受属神的事，就该交出属血肉的事。因为现今所建的这些，无论为收存活人或死者的遗体，都是暂时的，而时间正在逝去。我们难道在审判之后能把这些建筑物带到天堂吗？然而没有这些，我们此时便不能完成那些属于拥有天堂的事。所以，若你们切望得到属神的事，就当虔诚地支出属血肉的事。\Quote{拿起圣咏，将手鼓交给祂}：接纳我们的声音，回应你们的双手。\par

5. \Quote{悦耳的瑟琴，伴着竖琴。} 我记得我们曾向你们的爱德说明过瑟琴与竖琴的区别。……因为天主的圣言宣讲是属天的。但若我们等候属天的事，就不可在从事属地的事上懈怠；因为\Quote{瑟琴是悦耳的}，但是\Quote{伴着竖琴}。这正以前面另一种方式表达：\Quote{拿起圣咏，将手鼓交给祂}：此处以\Quote{瑟琴}代替\Quote{圣咏}，以\Quote{竖琴}代替\Quote{手鼓}。然而，这劝告我们：要以身体的工作回应天主圣言的宣讲。\par

6. \Quote{吹响号角} \BibleRef{咏 80:3}。这就是：大声并勇敢地宣讲，不要畏惧！正如先知在某处说：\Quote{高声呼喊，像号角一样扬起你的声音。} \BibleRef{依 58:1} \Quote{在月朔吹号角。} 曾命令在月朔吹号角：如今犹太人仍按肉体这样做，却按心神不明白。因为月朔就是新月；新月就是新生命。什么是新月？\Quote{所以，谁若在基督内，他就是一个新受造物。} \BibleRef{格后 5:17} 什么是\Quote{在月朔吹号角}？要满怀信心地宣讲新生命，不要畏惧旧生命的喧嚷。\par

7. \Quote{因为这是为以色列定的诫命，为雅各伯的天主定的典章} \BibleRef{咏 80:4}。有诫命之处，便有审判。因为\Quote{凡在律法之下犯了罪的，将按律法受审判。} \BibleRef{罗 2:12} 而那颁赐诫命者，主基督——圣言成了血肉——说：\Quote{我为审判来到了这世界，叫那些看不见的看见，那些看见的变成瞎子。} \BibleRef{若 9:39} 什么是\Quote{叫那些看不见的看见，那些看见的变成瞎子}？岂不是使谦卑者被高举，骄傲者被推翻？因为不是那些看见的要变成瞎子，而是那些自以为看见的人要被揭穿为瞎眼。这在压榨的奥迹中成就，叫那些看见的看不见，那些看不见的却看见。\par

8. \Quote{祂为若瑟立了这证据} \BibleRef{咏 80:5}。请注意，弟兄们，这是什么？若瑟意为增加。你们记得，你们知道若瑟被卖到埃及：若瑟被卖到埃及，预表基督转向外邦人。在那里，若瑟经历患难后被高举；在这里，基督在殉道者受苦之后被光荣。从此，外邦人更多地归属若瑟，从此即增加；因为\Quote{不怀孕、不生养的，你要欢呼；因为被弃者的儿女比有夫之妇的儿女更多。} \BibleRef{依 54:1} \Quote{祂立了它，直到他离开埃及地。} 请注意，这里也预指\Quote{安息日的第五日}：当若瑟离开埃及地时，即百姓借着若瑟繁衍，得以穿越红海。因此，那时水也生出有生命的爬虫。 \BibleRef{创 1:20} 那在预象中百姓穿越红海所预示的，无非就是信友借圣洗圣事的穿越；宗徒作证：\Quote{弟兄们，我不愿意你们不知道，我们的祖先都曾在云柱下，都从海中走过，都曾在云中和海中受洗归于梅瑟。} \BibleRef{格前 10:1} 所以穿越红海所象征的，无非是受洗者的圣事；追赶的埃及人所象征的，无非是过去罪过的众多。你们看最明显的奥迹：埃及人紧追、逼迫；罪也如此紧追，但只到水边为止。那么，你们这些尚未前来受洗的人，为什么害怕前来接受基督的圣洗、穿越红海呢？什么是\Quote{红}？以主的血祝圣的。你为什么不敢前来？或许是某些大罪的良心在折磨并刺痛你的心神，告诉你所犯的罪如此重大，以致你绝望它会被赦免。要害怕，若还有任何一个埃及人活着，你的罪就会留下什么！ \BibleRef{出 14:29}\par

但你们渡过红海，当你们被\Quote{以大能的手和伸出的臂}带出你们的过犯时，你们将领悟你们所不知道的奥秘：因为\Person{若瑟}自己也曾，\Quote{当他走出埃及地时，听见了一种他不认识的语言}。你们将听见一种你们不认识的语言：那些现在认识的人听见并认出，作证并知道。你们将听见你们应当把心放在哪里：刚才当我说这话时，许多人明白并高声回应，其余的人沉默不语，因为他们没有听见那种不认识的语言。那么，让他们赶快，让他们渡过，让他们学习。\par

9. \BibleQuote{他从重担中转回了他的背}\BibleRef{咏 80:6}。谁\Quote{他从重担中转回了他的背}，除了那呼喊\BibleQuote{凡劳苦和负重担的，你们都到我这里来}\BibleRef{玛 11:28}的那位？这同一件事以另一种方式被表明。埃及人追赶所做的事，罪恶的重担也做同样的事。就好比说：从什么重担？\Quote{他的手在篮子中服役}。篮子象征奴役的工作；清洗、施肥、搬运泥土，都是用篮子做的，这些工作是奴役性的，因为\BibleQuote{凡犯罪的，就是罪恶的奴隶；并且如果子使你们自由，你们就真正自由了}\BibleRef{若 8:34}。同样，世上被遗弃的事物也被视为篮子，但连篮子天主也用碎块填满了；\BibleQuote{十二筐}\BibleRef{玛 14:20}他用碎块填满了；因为\BibleQuote{他选择了世上被遗弃的事物，来使那些强有力的事物蒙羞}\BibleRef{格前 1:27}。而且当\Person{若瑟}用篮子服役时，他搬运泥土，因为他做砖。\Quote{他的手在篮子中服役}。\par

10. \BibleQuote{在患难中你呼求了我，我便救了你}\BibleRef{咏 80:8}。让每个基督徒的良心认识自己，如果它虔诚地渡过了红海，如果它以相信和遵守的信德听见了一种它不认识的语言，就让它认识到自己在患难中被俯听了。因为那是一个极大的患难，被罪恶的重担压垮。良心怎从地上被提升而喜乐。看，你受洗了，你昨天还沉重不堪的良心，今日使你喜乐。你在患难中被俯听了，记住你的患难。在你来到水边之前，你背负了多大的焦虑！你实践了多少斋戒！你在心中背负了多少患难！多少内在、虔诚、热心的祈祷！你的仇敌被杀死了，你所有的罪都被涂去。在患难中你呼求了我，我便救了你。\par

11. \Quote{我在暴风的隐秘处听到了你。}不是在海上的暴风，而是在心里的暴风。\Quote{我在\OriginalTerm{反叛之水}{water of contradiction}中考验了你。}的确，弟兄们，的确，那在暴风的隐秘处被俯听的人，应当在反叛之水中被考验。因为当他相信了，当他受了洗，当他开始走在天主的道上，当他努力被压入酒缸，并把自己从那跑在街上的酒糟中拉出来时，他将有许多扰乱者、许多侮辱者、许多诽谤者、许多沮丧者、许多甚至在他们能的地方威胁、阻止、压制的人。这一切就是全部的\Quote{反叛之水}。我猜想今天在这里有一些人，例如，我认为很可能有一些人，他们的朋友想把他们匆匆拉去竞技场，以及今天节日我不知道什么无聊的事：也许他们带那些人一起去了教堂。但无论他们带了那些人，还是他们自己没有允许自己被拉去竞技场，他们都在\Quote{反叛之水}中被试炼了。那么，不要羞于宣告你们所知道的，甚至在亵渎者中间捍卫你们所相信的……无论外邦的恶人如何肆虐，但愿我们自己内部的恶人不帮助他们！\par

你们回忆起来基督所说的话，祂如此诞生，是为了\BibleQuote{许多人的跌倒和许多人的兴起，并成为一个反对的标记}\BibleRef{路 2:34}。我们知道，我们看见：十字架的标记已经被树立，并且它被人反对。有人反对十字架的荣耀：但在十字架上有一个不能被腐蚀的标题。因为圣咏中有一个标题：\Quote{为标题的题词，你不要破坏。}它是一个反对的标记：因为犹太人说：\BibleQuote{不要写\InnerQuote{犹太人的君王}，而要写\InnerQuote{祂说：我是犹太人的君王}}\BibleRef{若 19:21}。反对被挫败了；回答是：\Quote{我所写的，我已经写上了。}\par

12. 所有这一切，从圣咏开始直到这节经文，我们听到了关于\OriginalTerm{压榨的油}{oil of the press}。剩下的是更令人悲伤和警告的：因为它属于\OriginalTerm{压榨的酒糟}{lees of the press}，甚至直到末了；也许在那插入的\OriginalTerm{细拉}{Diapsalma}也不是没有意义。但听这些话也是有益的，好叫那看见自己已经属于油的人可以欢喜；那有危险流入酒糟的人可以警惕。留心两者，选择其中一个，害怕另一个。\par

\BibleQuote{听啊，我的百姓，我要说话，我要向你们作证}\BibleRef{咏 80:8}。因为不是对一个外邦百姓，不是对一个不属于压榨的百姓：\BibleQuote{你们判断}，祂说，\BibleQuote{在我和我的葡萄园之间}\BibleRef{依 5:3}。\par

13. \BibleQuote{以色列，如果你们听了我，在你们中间必不会有新的神}。\BibleRef{咏 80:9} \Quote{新的神}是为了时间而造的：但我们的天主不是新的，从永远到永远。而我们的基督，或许作为人，是新的，但作为天主是永远的。因为在起初之前有什么呢？确实，\BibleQuote{在起初已有圣言，圣言与天主同在，圣言就是天主}。\BibleRef{若 1:1} 而我们的基督自己就是那成了血肉的圣言，以便居住在我们中间。\BibleRef{若 1:14} 所以，断不容任何人里面有新的神。新的神要么是石头，要么是幻影。有人说，不是石头，我有银的和金的。那说\Quote{外邦人的偶像是银子和金子}的人，恰当地选出了极贵重的东西。它们是伟大的，因为是金银的；是贵重的，是闪亮的；但是，\Quote{有眼却不能看}！这些神是新的。还有什么比作坊里出来的神更新呢？是的，尽管那些旧神已被蜘蛛网覆盖，不是永久的就是新的。关于外邦人，就说到这里。\par

14. 因为如果你们内有错误，你们就不会敬拜一个陌生的神。如果你们不想一个假神，你们就不会敬拜一个手工制造的神：因为\Quote{决不会}有陌生的神在你们中间。\Quote{因为我就是}。为什么你们要崇拜那不存在的东西？\BibleQuote{因为我上主是你的天主}。\BibleRef{咏 80:10} 因为\Quote{我是我所是}，而且的确祂说\Quote{我是}，那自有的，超越一切受造物：但我在时间中给了你们什么好处呢？\Quote{是我领你出了埃及地}。这不是只对那个民族说的。因为我们都是被领出埃及地的，我们都经过了红海；追赶我们的敌人都淹死在水里。我们不要辜负我们的天主；不要忘记那永存的天主，而在我们内制造一个新的神。\Quote{我领你出了埃及地}，天主说。\Quote{你要大大张口，我就使它饱足}。你们因为心中立了新的神而自陷于狭窄；打碎那虚妄的像，从良心中推倒那虚假的偶像：\Quote{大大张口}，在宣认中，在爱中：\Quote{我就使它饱足}，因为生命的泉源在我这里。\par

15. \BibleQuote{但我的百姓没有听从我的声音}。\BibleRef{咏 80:11} 因为祂不会对这些话，除非是对祂自己的百姓。因为\BibleQuote{我们知道：凡法律所说的，是对在法律之下的人说的}。\BibleRef{罗 3:19} \Quote{以色列没有听从我}。谁？对谁？以色列对我。啊，忘恩负义的灵魂！藉着我而存在的灵魂，被我召叫的灵魂，被我带回到希望中的灵魂，被我洗净了罪恶的灵魂！\Quote{以色列没有听从我}！因为他们受了洗并经过了红海：但在路上他们抱怨、争辩、诉苦、骚动不安，辜负了那从追赶的敌人中拯救他们、引导他们走干地、经过旷野、且有食物和饮料、夜间有光、白天有荫的那一位。\par

16. \BibleQuote{我任凭他们随从心里的情欲而行}。\BibleRef{咏 80:12} 看那压榨机：开口敞开着，渣滓流出。\Quote{我任凭他们，}不是随从我的命令的健康；而是，随从他们心里的情欲：我把他们交给自己。宗徒也说：\BibleQuote{天主任凭他们随从心里的情欲}。\BibleRef{罗 1:24} \Quote{我任凭他们随从自己心里的情欲，他们必行在自己的情欲中。} 如果你们至少是在筛入主隐藏的酒池中，如果你们至少对祂的仓库有了全心的爱，就有你们所厌恶的东西。有些人拥护竞技场，有些人拥护露天剧场，有些人拥护街头摊位，有些人拥护剧院，有些拥护这个，有些拥护那个，有些最终拥护他们的\Quote{新神}；\Quote{他们必行在自己的情欲中}。\par

17. \BibleQuote{但愿我的百姓肯听我的话，但愿以色列肯遵行我的道}。\BibleRef{咏 80:13} 因为或许那个以色列说：看，我犯罪了，这是明显的，我随从自己心里的情欲：但我能做什么呢？这是魔鬼做的。这是恶魔做的。什么是魔鬼？谁是恶魔？当然是你们的仇敌。\BibleQuote{我必使他们的仇敌都屈服，必伸手攻击那些迫害他们的人}。\BibleRef{咏 80:14} 但现在他们抱怨仇敌有什么用？他们自己成了更坏的仇敌。因为怎样？下文是什么？你们抱怨仇敌，你们自己是什么？\par

18. \BibleQuote{天主的仇敌向祂说了谎}。\BibleRef{咏 80:15} 你弃绝吗？我弃绝。而他又回到他所弃绝的。事实上，你弃绝了什么？无非是恶行、魔鬼的行为、该受天主谴责的行为：偷窃、抢夺、伪证、杀人、奸淫、亵渎圣事、可憎的仪式、邪术……\par

19. 因此，如果所有这些行为\Quote{不得承受天主的国}（是的，不是行为，而是\Quote{行这样事的人}；\BibleRef{迦 5:21} 因为这样的行为在火中必不存在：因为当他们在火中焚烧时，不会再去偷窃或奸淫；但是\Quote{行这样事的人不能承受天主的国}）；所以他们不会站在右边，与那些将听到\BibleQuote{来吧，我父所祝福的，承受国度}的人一起，因为\Quote{行这样事的人不能承受天主的国}。因此，如果他们不在右边，就只能是左边。对左边的人，祂将说什么？\BibleQuote{你们到永火里去}。因为\Quote{他们的时候将是永远的}。\par

\Quote{那么，请向我们解释，}有人说，那些在根基上建造木、草、禾秸的人，如何不丧亡，却\BibleQuote{要得救，却像是从火中经过的一样}？这确实是一个隐晦的问题，但我按自己的能力简要地告诉你们。弟兄们，有些人是完全轻视这世界的，他们对时间流逝中所流动的一切都不感到愉悦，他们不以爱恋依附于任何地上的工作，他们是圣洁、贞洁、节制、公义的，甚至可能变卖所有财产分给穷人，或者\BibleQuote{拥有如同不拥有，享用这世界如同不享用}。\BibleRef{格前 7:30}但也有另一些人，以某种情感依附于允许软弱之人的事物。他不抢夺别人的产业，但如此爱自己的产业，以致若失去它就会不安。他不贪恋别人的妻子，但如此依恋自己的妻子，如此与自己的妻子同居，以至于在此事上没有遵守为了生育子女而定的法律尺度。他不夺取别人的东西，但追讨自己的东西，并与自己的弟兄打官司。因为对他们曾说：\BibleQuote{你们彼此有诉讼，这已经是你们的大错。}\BibleRef{格前 6:7}但他命令这些诉讼应在教会内审理，不要拖到法庭，然而他说这是错误。因为一个基督徒为地上的事物争竞，与那应许了天国的人不相称。他没有把整个心高举向上，而是有一部分拖在地上……因此，若你爱你的财产，却不因它行强暴，不因它作假见证，不因它杀人，不因它发假誓，不因它否认基督：因你不愿为它行这些事，你就有基督为根基。但因你爱它，若失去便忧伤，你在根基上安放的，不是金、银、宝石，而是木、草、禾秸。因此你将得救，当你所建造的开始燃烧时，却像是从火中经过一样。因为任何人若在这根基上建造奸淫、亵渎、亵圣、偶像崇拜、伪誓，不要以为他将\BibleQuote{经由火得救}，好像它们是\BibleQuote{木、草、禾秸}；但那在天国根基上（即在基督上）建造对地上事物的爱，他那对暂时之物的爱将被焚烧，而他本人将藉由正确的根基而得救。\par

...\BibleQuote{祂用上等麦子喂养他们，又从磐石中使他们饱享蜂蜜}\BibleRef{咏 80:16}。在旷野中，祂从磐石中引出了水\BibleRef{出 17:6}，不是蜂蜜。\Quote{蜂蜜}是智慧，在心灵的食粮中因甘甜而居首位。那么，有多少主的仇敌，向主说谎，不但用上等麦子喂养，又从磐石中用蜂蜜，即从基督的智慧中喂养？有多少人被祂的话、被对祂圣事的认识、被祂比喻的阐释所喜悦？有多少人喜悦，欢呼喝彩！而这蜂蜜并非来自任何普通人，而是\BibleQuote{从磐石}。但\BibleQuote{那磐石就是基督}\BibleRef{格前 10:4}。那么，有多少人被那蜂蜜满足，呼喊说：这是甘甜的；说：无法想象或说出任何更好、更甘甜的了！然而主的仇敌却向祂说了谎。我不愿再停留于令人忧闷的事；虽然圣咏以此为目的以恐怖结束，但请你们从它的结尾转向标题：\Quote{欢跃于天主我们的助佑。}转向天主。\par

\SourceRecord{Translated by J.E. Tweed.；Nicene and Post-Nicene Fathers, First Series；Vol. 8.；Edited by Philip Schaff.；Buffalo, NY: Christian Literature Publishing Co.,；1888.；<http://www.newadvent.org/fathers/1801081.htm>.}

% structure: p0001 source=h1 target=section reason=page-title

\section{圣咏第八十二篇释义}

1. 这篇圣咏与其他同名圣咏一样，如此命名，或出于写作之人的名字，或出于同一名字的解释，以便在含义上指称亚萨所代表的会堂；尤其在第一句中暗示了这一点。它的开头是：\BibleQuote{天主站在众神的会集中}\BibleRef{咏 82:1}。然而，我们切不可将这些神理解为外邦人的神、偶像，或天上地下的任何受造物，除非是人；因为在这句之后不久，同篇圣咏便说明并解释它所指的、天主站在其会集中的神是什么，它说：\Quote{我曾说：你们是神，你们都是至高者的儿子；但你们要像人一样死去，如一位首领般跌倒。}在这些至高者的儿子的会集中——关于他们，同一至高者曾以依撒意亚的口说：\BibleQuote{我生了儿子，抚养他们，他们却轻视了我}\BibleRef{依 1:2}——天主站着。所谓会堂，我们理解为以色列人民，因为会堂一词是专用于他们的，尽管他们也被称为教会。相反，宗徒们从未称我们的集会为会堂，而总是称为\OriginalTerm{教会}{Ecclesia}；无论是为了区别，还是因为会堂得名于聚集，而教会得名于召集，二者有所不同：因为\Quote{聚集}（或\Quote{聚拢}）一词用于牲畜，尤其用于称为\Quote{羊群}的那一类；而\Quote{召集}（或\Quote{召唤在一起}）则更用于有理智的受造物，如人。……所以，我认为天主站在其中的众神的会集，那是清楚的。\par

2. 接下来的问题是，我们应当理解是圣父、圣子、圣神，还是天主圣三\Quote{站在众神的会集中，并在众神中间施行判断}；因为每一位都是天主，而天主圣三本身是一位天主。要说明这一点确实不容易，因为不能否认天主对受造物有一种非形体性的、符合其本性的属神临在，以奇妙的方式存在，但只有少数人明白，且明白得不完全；正如对天主所说：\Quote{我若升到天上，祢在那里；我若下到阴府，祢也在那里。}因此，恰当的说法是，天主无形地站在人的会集中，正如祂充满天地，祂曾借先知的口这样论及自己\BibleRef{耶 23:24}；并且，祂不仅被说成，而且在某种意义上，被知道站在祂所创造的万物中——按人所能理解的，如果人也站着并听从祂，且因祂内在的声音而大大喜悦。但我想，这篇圣咏暗示了在某一特定时间发生的事，即天主站在众神的会集中。因为那种充满天地的站立，既不专属于会堂，也不随时间而变化。所以，\Quote{天主站在众神的会集中}，即那一位曾说\BibleQuote{我被派遣，只为以色列家迷失的羊}\BibleRef{玛 15:24}。原因也被提到：\Quote{但在中间，要判断众神。}……\par

3. \BibleQuote{你们要审判不义要到几时，且袒护恶人？}\BibleRef{咏 82:2}；正如在另一处所说：\Quote{你们心怀愚顽要到几时？}直到那心灵之光来临为止？我赐了法律，你们固执反抗；我派了先知，你们虐待他们，或杀害他们，或容忍那些如此行的人。但如果那些杀害了被派往他们那里的天主的仆人的人甚至不值得被对他们说话，那么你们——当这些事情发生时保持沉默，即你们会效仿那些当时沉默的人，仿佛他们是无辜的——\Quote{你们要审判不义要到几时，且袒护恶人？}如果继承人现在来了，祂要被杀吗？祂不是为你们的缘故愿意成为一个在监护人下的孩子吗？祂不是为你们的缘故像有需要的人一样饥饿和口渴吗？祂不是向你们呼喊说：\BibleQuote{你们跟我学，因为我是良善心谦的}\BibleRef{玛 11:29}？祂不是\BibleQuote{成了贫穷的，当时祂是富有的，好使你们因祂的贫穷而成为富有的}\BibleRef{格后 8:9}？那么，\BibleQuote{你们当为孤儿和穷人申辩，应为谦卑和贫乏的人伸冤}\BibleRef{咏 82:3}。不要相信那为自己而富足骄傲的人，而要相信那为你们而谦卑贫穷的祂；你们要宣告祂是正义的。但他们将嫉妒祂，且毫不留情，说：\Quote{这是继承人，来，我们杀祂，产业就成为我们的了。}那么，\BibleQuote{要拯救穷人和贫乏的人，从恶人的手中救出他们}\BibleRef{咏 82:4}。这是为了让人知道，在基督诞生并被处死的那个民族中，那些人数众多——正如福音所说，犹太人怕他们，因而不敢下手害基督——后来却同意并允许基督被恶毒和嫉妒的犹太首领杀害的人，对于如此重大的罪行也并非无罪；因为如果他们愿意，他们仍会被惧怕，以致恶人的手永远无法胜过祂。因为关于他们，别处说：\Quote{哑巴狗，不知道吠叫。}关于他们也有这话：\Quote{看，义人灭亡，无人放在心上。}就他们这些希望祂灭亡的人而言，祂灭亡了；但祂怎能因死亡而灭亡呢？祂反而藉此找回那已灭亡的。所以，如果那些因默许而让如此邪恶的行为发生的人应受责备和斥责，那么那些故意并恶意地行这事的，又当如何受责备，甚至不仅仅受责备，而应受何等严厉的惩罚呢？\par

4. 诚然，以下的话最适合他们全体：\BibleQuote{他们不知道，也不明白，在黑暗中行走}\BibleRef{咏 81:5}。\BibleQuote{因为如果他们知道了，决不会将光荣的主钉在十字架上}\BibleRef{格前 2:8}。而那些其他人，如果他们知道了，就决不会同意要求释放巴拉巴，而将基督钉在十字架上。但是，正如上述的盲目发生在以色列身上局部地，直到外邦人的满全进来；这个人民的盲目导致了基督的被钉十字架，\Quote{大地的一切根基将会动摇}。它们已经动摇了，并将继续动摇，直到外邦人的预定满全进来。因为在主实际死亡时，大地震动，磐石裂开。\BibleRef{玛 27:51}如果我们以大地的基础来理解那些拥有丰富属地财富的人，那么的确预言他们应当被动摇，或者因惊奇于谦卑、贫穷、死亡竟在基督内被如此爱戴和尊崇，而在他们看来却是极大的痛苦；或者甚至他们自己应当爱慕并追随它，而轻视此世的虚幻幸福。所有大地的根基就这样被动摇，一部分人赞叹，一部分人甚至被改变。因为，我们毫无荒谬地称那些在其上天国在圣人及信友身上被建立的人为天堂的基础；他们的首要基础是基督自己，生于童贞女，关于他宗徒说：\BibleQuote{除了那已奠基的，即基督耶稣，没有人能奠立别的根基}\BibleRef{格前 3:11}；其次是宗徒和先知们自己，因他们的权威天乡被拣选，好使我们因服从他们而与他们一同被建造；因此他致厄弗所人说：\BibleQuote{你们已被建立在宗徒和先知的基础上，基督耶稣自己为角石}\BibleRef{弗 2:20}。……然而属地幸福的王国是骄傲，对此基督的谦卑来对抗，斥责那些他愿意藉谦卑使之成为至高者子女的人，并责备他们：\BibleQuote{我曾说：你们是神，你们都是至高者的子女}\BibleRef{咏 81:6}。\BibleQuote{但你们要像人一样死亡，并像一位王子一样跌倒}\BibleRef{咏 81:7}。他是对那些人说了这话，\BibleQuote{我曾说：你们是神，}以及特别对那些未预定得永生的人；而对另一些人，\BibleQuote{但你们要像人一样死亡，}等等，\BibleQuote{并要像一位王子一样跌倒，}以此也区分诸神；或者他是责备全体，为了区分听从者和接受纠正者：\BibleQuote{我曾说：你们是神，你们都是至高者的子女：}也就是说，我对你们全体允诺了天国的幸福，\BibleQuote{但你们}，因肉身的软弱，\BibleQuote{要像人一样死亡，}并因灵魂的傲慢，\BibleQuote{像一位王子一样，}即魔鬼，不得被高举，而\BibleQuote{要跌倒。}仿佛他说道：虽然你生命的时日如此之少，使你如人速速死亡，但这无助于你的改正；但你像魔鬼一样，他的时日在此世很多，因为他不在肉身死亡，你被高举以致跌倒。因为由魔鬼般的骄傲，那乖僻而盲目的犹太人首领嫉妒基督的光荣：由此发生，并且至今仍发生，即基督被钉至死的谦卑在那些爱慕此世高贵的人眼中被轻视。\par

5. 因此，为治愈这一恶习，以先知本人的口吻说：\BibleQuote{天主，请起来，审判大地}\BibleRef{咏 81:8}；因为大地在钉死你时高高昂起：从死者中起来，审判大地。\BibleQuote{因为你将在万民中毁灭。}消灭什么？岂不是大地吗？即消灭那些贪恋属地之事的人，或在信友中消灭属地情欲和骄傲本身的感觉；或分离那些不信的人，如同大地被践踏而灭亡。如此，他的肢体（他们的生活在天上）审判大地，并在万民中毁灭它。但我不该忽略提及，有些抄本写作：\BibleQuote{因为你将在万民中承受为产业。}这也符合意义，且没有任何事阻止两种含义同时并存。他的产业由爱而来，而他藉其命令和仁慈恩宠培养此爱，他毁灭了属地的情欲。\par

\SourceRecord{Translated by J.E. Tweed.；Nicene and Post-Nicene Fathers, First Series；Vol. 8.；Edited by Philip Schaff.；Buffalo, NY: Christian Literature Publishing Co.,；1888.；<http://www.newadvent.org/fathers/1801082.htm>.}

% structure: p0001 source=h1 target=section reason=page-title

\section{圣咏第八十三篇释义}

1. 这篇圣咏的标题是\Quote{阿撒夫圣咏的诗歌}。我们已经多次说过，阿撒夫的意思是\Quote{会众}。因此，那名叫\Person{阿撒夫}的人，在许多圣咏的标题中代表天主之会众。但在希腊语中，会众被称为\Quote{会堂}，这已成为犹太人的一种专有名称，即称为\Quote{会堂}；同样，基督徒民族更常被称为\Quote{教会}，因为它也是聚集而成的。\par

2. 因此，天主的子民在这篇圣咏中说：\Quote{天主，谁能相似祢？}（见：\BibleRef{咏 82:1}）。我认为这更适合用于基督，因为祂虽取了人的样式（见：\BibleRef{斐 2:7}），却被那些轻视祂的人认为与常人无异；祂甚至被\Quote{列于不义之人中}（见：\BibleRef{依 53:12}），但目的是为了受审判。但当祂来审判时，将要实现这里所说的：\Quote{天主，谁能相似祢？}因为如果圣咏不常对主基督说话，那么每一位信徒都确信是对基督所说的那句话也不会被说出：\Quote{祢的御座，天主，永永远远，祢王国的权杖是正直的权杖。}因此，现在也对祂说：\Quote{天主，谁能相似祢？}因为藉着祢的贬抑，祢曾屈尊与许多人相似，甚至与同祢一起被钉的强盗相似；但当祢带着荣耀来临时，\Quote{谁能相似祢？}……\par

3. \Quote{因为看，祢的仇敌喧嚷，恨祢的人抬起头来}（见：\BibleRef{咏 82:2}）。在我看来，这预示末世，那时这些现在因畏惧而被压抑的事物将爆发为自由言论，但完全无理性，以致它与其称为言论或谈话，不如称为\Quote{喧嚷}。因此，他们那时不是开始仇恨，而是\Quote{恨祢的人}将\Quote{抬起头来}。不是\Quote{诸头}，而是\Quote{头}；因为他们甚至到了那种地步，他们将有那\Quote{高抬自己超过一切被称为神和受崇拜者}的头（见：\BibleRef{得后 2:4}）；这样，在他身上尤其要应验：\Quote{凡高举自己的，必被贬抑}（见：\BibleRef{路 14:11}）；当那位曾被告知：\Quote{天主，不要缄默，也不要温和}的祂，将\Quote{用祂口中的气息杀死他，并以祂来临的光辉消灭他}（见：\BibleRef{得后 2:8}）。\Quote{他们恶毒地谋害祢的百姓}（见：\BibleRef{咏 82:3}）。或者，如其他抄本所载：\Quote{他们狡猾地谋划计策，谋害祢的圣人。}这是以轻蔑的态度说的。因为他们怎能够伤害天主的民族或百姓，或祂的圣人？他们知道如何说：\Quote{若是天主与我们同在，谁能反对我们？}（见：\BibleRef{罗 8:31}）\par

4. \Quote{他们说：来吧，让我们从一国中消灭他们}（见：\BibleRef{咏 82:4}）。他用单数代替复数：如\Quote{这是谁的牲畜？}，尽管所问的是一群，意思却是\Quote{这些牲畜}。最后，其他抄本作\Quote{从列国中}，译者们更侧重意思而非字词。\Quote{来吧，让我们从一国中消灭他们。}这就是那种\Quote{喧嚷}，而非说话，因为他们以虚妄的言语徒然地发出噪音。\Quote{不要再提以色列的名字。}另一些人更清楚地表达为\Quote{愿以色列的名字不再被纪念。}因为\OriginalTerm{提及名字}{memoretur nominis}是拉丁语中不常见的说法；更常见的说法是\OriginalTerm{名字被提及}{memoretur nomen}；但意思相同。因为说\Quote{提及名字}的人翻译的是希腊短语。但这里以色列必须理解为亚巴郎的后裔，关于这后裔，宗徒说：\Quote{因此，你们是亚巴郎的后裔，按照恩许的承继人}（见：\BibleRef{迦 3:29}）。不是按照血肉的以色列，关于它，他说：\Quote{看，按血肉的以色列。}\par

5. \BibleQuote{他们同心合意，勾结起来，同你作对，缔结盟约}\BibleRef{咏 82:5}：好像他们能更强一样。事实上，圣经中称为\Quote{盟约}的，不仅指立约者死后才生效的遗嘱，而且他们称一切盟约和法令为盟约。因为拉班和雅各伯立了一个盟约\BibleRef{创 31:44}，这当然是在活人之间有效；在圣经中这类例子不胜枚举。然后他开始以某些民族专名提及基督的敌人；这些名字的解释充分表明他想要被理解的是什么。因为用这样的名字最恰当地比喻真理的敌人。例如\Quote{厄东人}解释为\InnerQuote{血人}或\InnerQuote{地土之人}\Quote{依市玛耳人}是\InnerQuote{顺从自己的人}，因此不是顺从天主，而是顺从自己。\Quote{摩阿布}是\InnerQuote{来自父亲}；从坏的方面理解，没有比将其与历史实际相联系更好的解释：罗特，一位父亲，通过其女儿所促成的非法结合生了摩阿布；正是因为这个缘故他得了这个名字。\BibleRef{创 19:36}但他的父亲是好的，正如\BibleQuote{法律原是好的，只要人用的合法}\BibleRef{弟前 1:8}，而不是不洁和非法地使用。\Quote{哈盖尔人}是归化者，即外方人，在反对天主圣民的敌人中，这个名字也指那些人——不是那些成为公民的人，而是那些坚持异己和敌对心态、当有作恶机会时就表现出来的人。\Quote{革巴耳}是\InnerQuote{空虚的山谷}，即假谦卑。\Quote{阿孟}是\InnerQuote{不安定的民族}或\InnerQuote{忧伤的民族}。\Quote{阿玛肋克}是\InnerQuote{舔地的民族}；因此别处说：\BibleQuote{他的仇敌要舔尘土。}\Quote{异族}，虽然按其拉丁名字本身充分表明他们是异己，因此自然是敌人，但在希伯来语中称为\Quote{培肋舍特人}，解释为\InnerQuote{因醉酒跌倒}，如同被世俗奢侈灌醉的人。\Quote{提洛}在希伯来语中称为Sor；无论解释为狭窄还是患难，在这些天主圣民的敌人身上，必须按宗徒所说：\BibleQuote{患难和困苦加于一切作恶的人}\BibleRef{罗 2:9}。所有这些在圣咏中这样列举：\BibleQuote{厄东人的帐幕，依市玛耳人，摩阿布，哈盖尔人，革巴耳，阿孟，阿玛肋克，培肋舍特人，以及住在提洛的人。}\par

6. 好像要指出他们成为天主圣民的敌人的原因，他接着说：\BibleQuote{因为亚述与他们一同来了。}如今亚述常用来比喻魔鬼，\BibleQuote{那在悖逆之子身上运行}\BibleRef{弗 2:2}的，好像在自己的器皿中，使他们攻击天主圣民。\BibleQuote{他们帮助了罗特的子孙}，他说：因为所有敌人，由于魔鬼——他们的首领在他们内运行，\BibleQuote{帮助了罗特的子孙}，罗特解释为\Quote{堕落者}。但背弃的天使很好地解释为堕落的子孙，因为他们因背离真理而转向成为魔鬼的随从。这些人就是宗徒所说的：\BibleQuote{你们战斗不是对抗血肉，而是对抗率领者，对抗掌权者，对抗这黑暗世界的霸主，对抗天界里邪恶的鬼神}\BibleRef{弗 6:12}。那些无形的敌人得到了不信之人的帮助，因魔鬼在他们内运行，以便攻击天主圣民。\par

7. 现在让我们看看先知的神恩祈求什么降在他们身上，与其说是诅咒，不如说是预言。\BibleQuote{求你待他们如同米德杨、息色辣，如同雅宾在克雄河边}\BibleRef{咏 82:9}\BibleQuote{他们都在恩多尔灭亡，成了地上的粪土}\BibleRef{咏 82:10}。历史记载，所有这些都被以色列——当时的天主圣民——征服和制伏了，正如他接着提到的那些人一样：\BibleQuote{使他们的首领好像敖勒布和则布，好像则巴黑和匝耳慕纳}\BibleRef{咏 82:11}。这些名字的意义如下：米德杨解释为扭曲的审判；息色辣为快乐被排除；雅宾为智者。\BibleRef{民 4:7}但在这些被天主圣民征服的敌人中，应理解为宗徒所说的那位智者：\BibleQuote{智者在哪里？经师在哪里？今世的辩士在哪里？}\BibleRef{格前 1:20}敖勒布是干燥，则布是狼，则巴黑是牺牲（即狼的牺牲），因为他也有他的牺牲；匝耳慕纳是骚动的阴影。所有这些都符合天主圣民以善胜恶的恶事。此外，克雄——他们被击败的河流——解释为他们的顽固。恩多尔——他们灭亡之处——解释为诸世代之泉源，即肉性的世代，他们沉溺于此，因而灭亡，他们不关心那引入生命的重生，因为在那里\BibleQuote{人们不娶不嫁}\BibleRef{路 20:35}，因为他们不再死亡。所以关于他们正说：\BibleQuote{他们成了地上的粪土}，因为从他们生出的只是地上的果实。因此，正如所有这些都作为预像被天主圣民征服，他祈求那些其他敌人也能在真理中被征服。\par

8. \BibleQuote{他们的所有首领曾说：我们要占据天主的圣所作为产业}\BibleRef{咏 82:12}。这就是那虚妄的喧嚣，正如上文所说，你的敌人曾用它发出嘟囔。但所谓\Quote{天主的圣所}，除非指天主的殿，还能指什么？正如宗徒所说：\BibleQuote{天主的殿是圣的，这殿就是你们}\BibleRef{格前 3:17}。敌人除了企图占据——即使之服从自己——天主的殿，使它屈从于他们不敬神的意愿之外，还想要什么？\par

9. 接着是什么呢？\Quote{我的天主，使他们像轮子一样}\BibleRef{咏 82:13}。将这话理解为他们认为的事物毫无恒心，是很恰当的；但我认为也可以正确地解释为\Quote{使他们像轮子一样}，因为轮子在后面的部分被抬起，在前面的部分被抛下；天主子民的一切仇敌也遭遇如此。因为这不是一个愿望，而是一个预言。他接着又说：\Quote{像风面前的碎秸}。\Quote{面前}是指临在；因为风有什么脸面呢？它既无形体特征，只是一种运动，乃是空气的一种波动而已。但这里是用它来指诱惑，轻浮虚荣之心因此被卷走。\par

10. 这种轻浮，由于轻易同意邪恶，随之而来的是严厉的折磨；因此他继续写道：\par

\Quote{好像火烧树林，火焰焚烧山岭}\BibleRef{咏 82:14}：\Quote{求祢以祢的暴风追赶他们，以祢的愤怒扰乱他们}\BibleRef{咏 82:15}。他说\Quote{树林}是因为不结果实，\Quote{山岭}是因为高傲；天主子民的仇敌正是如此：缺乏正义，充满骄傲。当他说\Quote{火}和\Quote{火焰}时，他是用另一个词重复天主审判和惩罚的观念。但说\Quote{以祢的暴风}，他接着解释说，那就是\Quote{祢的愤怒}；而前面的\Quote{祢要追赶}与\Quote{祢要扰乱}相对应。然而，我们必须注意，天主的愤怒没有任何躁动的情绪；因为祂的愤怒是祂正义惩罚方式的表达：正如法律可以说是愤怒的，当执行者因法律的制裁而被激发去惩罚时。\par

11. \Quote{求祢使他们的脸面充满羞耻，叫他们寻求祢的名，上主}\BibleRef{咏 82:16}。他所预言于他们的是美好而可欲的；除非在那些天主子民的仇敌团体中，有一些人如此蒙恩，以致在最后审判之前先得到这恩赐，否则他不会这样预言。因为如今他们混杂在一起，这就是仇敌的团体，就嫉妒而言，他们与天主子民争竞。现在，他们能作声就作声，抬头就抬头；但只是个别地，而不是像世界末日时普遍地那样，那时最后的审判即将降临。但这同一个团体，既包括那些从这数目中将要相信并转入另一个身体的人（因为他们的脸面充满羞耻，是为了寻求上主的名），也包括那些在同样的邪恶中坚持到底的人，他们如同风前的碎秸，像树林和不结果实的山岭一样被烧尽。他再次回到这些人的话题，说：\Quote{他们要惭愧，永远受苦恼}\BibleRef{咏 82:17}。因为寻求上主之名的人不会永远受苦恼，而是因他们的罪恶感到羞耻，他们为此而苦恼，以便寻求上主的名，通过这名他们不再受苦恼。\par

12. 他又回到这些后者，他们在同一敌人团体中，为此目的而蒙羞，以便他们不再永远蒙羞；并为此目的而被消灭，就他们是邪恶而言，以便他们成为良善，永远存活。因为论到他们，他说了\Quote{让他们惭愧而灭亡}，他立刻加上\Quote{让他们知道祢的名是上主，惟有祢是普天下至高者}\BibleRef{咏 82:18}。他们得知这事，就应如此惶恐，以致中悦天主；他们应如此灭亡，以致得以存留。他说：\Quote{让他们知道，祢的名是上主}，仿佛其他被称为主的，不是真正被称为主，而是虚假的，因为他们作为仆人统治，与真正的主相比并不是主；正如经上说的：\Quote{我是自有者}\BibleRef{出 3:14}；仿佛万物与造物主相比，就不是存在的。他又加上：\Quote{惟有祢是普天下至高者}，或其他抄本作\Quote{普天之上}，如同可说是在全天之上或全天之上。但他更愿意用后者，以压制尘世的骄傲。因为尘世不再骄傲，即人不再骄傲，曾对他说：\Quote{你原是灰土}\BibleRef{创 3:19}，以及\Quote{灰土尘埃，怎么还敢傲慢？}\BibleRef{德 10:9}——当他说上主是普天下至高者时，意思是任何人的想法都不能对抗\Quote{按祂旨意蒙召的人}，以及关于这些人的话：\Quote{天主若帮助我们，谁能反对我们呢？}\par

\SourceRecord{Translated by J.E. Tweed.；Nicene and Post-Nicene Fathers, First Series；Vol. 8.；Edited by Philip Schaff.；Buffalo, NY: Christian Literature Publishing Co.,；1888.；<http://www.newadvent.org/fathers/1801083.htm>.}

% structure: p0001 source=h1 target=section reason=page-title

\section{《圣咏第八十四篇释义》}

1. 这篇圣咏的标题是\Quote{为榨酒池}。而且，正如你们与我一同观察到的，我的挚爱（因为我看见你们听得非常专注），在它的正文中既没有提到任何榨酒池，也没有提到酒篮、酒缸，或任何榨酒池的器具或建筑；我们没有听到任何这类内容被诵读；因此，这个题写在它上面的标题\Quote{为榨酒池}是什么意思，并非一个容易的问题。因为确实，如果标题之后提到了我所列举的那些东西，血肉之人可能会相信这是关于那些可见的榨酒池的诗歌；但既然它有这个标题，后来却没有提到那些我们所熟知的榨酒池，我毫不怀疑，还有另一些榨酒池，是天主之神要我们在这里寻找和理解的。因此，让我们回想在那些可见的榨酒池中发生的事，并看看这如何属神地发生在教会中。葡萄挂在葡萄藤上，橄榄挂在树上。因为这两种果实通常要预备榨酒池；当它们挂在枝头时，似乎享受着自由的空气；葡萄还不是酒，橄榄也不是油，除非它们被压榨。同样，那些天主在创世之前就预定要与他独生子的肖像相符的人（\BibleRef{罗 8:29}），那伟大的葡萄串首先并在特别的意义上在他的苦难中被压榨了。因此，这类人在他们接近天主的服务之前，在世上享受着某种甜美的自由，就像挂着的葡萄或橄榄；但是，正如经上所说：\BibleQuote{我儿，你前去服事天主时，当在正义和敬畏中站立，预备你的灵魂接受试探}（\BibleRef{德 2:1}），所以每个人，当他接近天主的服务时，就发现自己来到了榨酒池；他将遭受苦难，将被压碎，将被挤压，不是为使他在此世灭亡，而是为使他流入天主的仓库。他身上情欲的遮盖被剥去，就像葡萄皮；这在情欲中已在他身上发生，宗徒曾说：\BibleQuote{脱去旧人，穿上新人}。这一切不作压榨是无法完成的：因此，这时期的各教会被称为榨酒池。\par

2. 但我们这些被安置在榨酒池中的人是谁呢？\Quote{科辣黑的子孙}。因为接下来是：\Quote{为榨酒池，交与科辣黑的子孙}。\OriginalTerm{科辣黑的子孙}{Sons of Core}已被解释为\Quote{秃头者之子}，正如那些懂那语言的人，根据他们对天主应尽的服务，所能向我们解释的那样……\par

3. 但是，被置于压榨之下，我们被压碎，是为了这个目的：使我们对于那曾将我们带向世俗、现世、暂时、不稳定和可朽坏之物的爱，在这些事物中在此生遭受痛苦、压榨的磨难和丰富的试探之后，我们开始寻求那不属于此生、不属于此地的安息；并且主成为，如经上所写，\BibleQuote{穷人的避难所}。什么是\Quote{穷人}？就是那仿佛一无所有、无助、无靠、在地上没有任何可安息之物的人。因为对于这样的穷人，天主是亲近的。因为即使人在世上富有金钱，……他们充满了恐惧而非享受。因为有什么像滚动的东西那样不稳定呢？金钱本身被铸成圆形，不是不恰当地，因为它不会静止。因此，这样的人虽然拥有一些东西，仍是穷人。但那些没有这种财富却只渴望它的人，也被算在将被弃绝的富人中；因为天主不看能力，而看意愿。所以穷人是在这世界的一切财物上匮乏的，因为即使财物在他们周围丰裕，他们也知道它是多么短暂；他们呼号天主，在世上没有任何事物可以令他们喜悦并使他们沉溺，被置于丰富的压榨和试探中，如同在榨酒池中，他们流出来，成了油或酒。这些是什么？不过是善的渴望。因为天主是他们唯一的渴望；他们不再爱地上的事物。因为他们爱那创造了天地的主；他们爱他，却尚未与他同在。他们的渴望被延迟，是为了使它增长；它增长，是为了它能领受。因为对于渴望者，天主将给予的并非任何微小之物，而他也不需被轻微锻炼，以便配得领受如此巨大的善；天主将给予的不是他所创造的任何东西，而是那创造万物的他自己。锻炼自己以领受天主：那件你将永远拥有的事物，你要长久地渴望……\par

4. 因此，至亲爱者，各人当尽力向主天主许愿并实践，尽力而为；愿无人回头，无人以从前的兴趣自娱，无人从面前转向背后；愿他奔跑直到抵达：因为我们不是用脚奔跑，而是用渴望。但在此生无人可以说他已经抵达。因为谁能如保禄那样完美呢？（\BibleRef{斐 3:13}）然而他说：\BibleQuote{弟兄们，我不认为我已经得到了。}\par

5. 所以，若你即使幸福时也感受到这世界的激情，如今你明白你就在榨酒池中……\newline{}因此，若世界以幸福对你微笑，你要想象自己在榨酒池中，并说：\Quote{我找到了困苦与忧伤，我便呼号了上主的名。}祂没有说\Quote{我找到了困苦}而无意义——这一种困苦是本被隐藏的：因为有些人的困苦在这世界是隐藏的，他们以为远离上主而幸福。\BibleQuote{因为我们几时住在这肉身内，就是远离了主。}\BibleRef{格后 5:6} 若你远离了父亲，你会不幸福；你远离了上主，还会幸福吗？所以有人认为他们处境良好。但那些领悟到：无论财富与享乐多么丰裕，万物都听命于他们，没有烦恼侵入，没有逆境恐吓，但只要远离了上主，他们便处于恶劣境况的人——用最锐利的眼光，这些人找到了困苦与忧伤，并呼号了上主的名。这就是在这篇圣咏中歌唱的那一位。祂是谁？基督的奥体。那又是谁？你，若你愿意：我们众人，若我们愿意：因为基督的奥体只有一个。\par

\BibleQuote{万军的上主，祢的居所是多么可爱！}\BibleRef{咏 83:1} 他曾在某些居所中，即在榨酒池中；但他向往另一些居所，那里没有压榨；在此他叹息向往那些居所，从这里他仿佛因渴望而流注到那里。\par

6. 随后呢？\BibleQuote{我的灵魂渴慕且渴慕上主的庭院。}\BibleRef{咏 83:2} 说它\Quote{渴慕且渴慕}还不够：它渴慕什么？\Quote{上主的庭院。}葡萄被压榨时已耗尽：但为了什么？为了变成酒，流进酒缸，进入库房的其余部分，在那里安享极大的宁静。在这里它被渴慕，在那里它被接纳：在这里是叹息，在那里是喜乐：在这里是祈祷，在那里是赞美：在这里是呻吟，在那里是欢悦。我提的这一切，愿无人在此因羞愧而转离；愿无人不愿受苦。有危险：恐怕葡萄怕了榨酒池，反而被飞鸟或野兽吞吃。\par

7. 你已听到榨酒池中的呻吟：\BibleQuote{我的灵魂渴慕且渴慕上主的庭院；}\BibleRef{咏 83:2}聆听它如何在希望中坚持、欢跃：\BibleQuote{我的心灵和我的肉身因生活的天主而欢跃。}它们为何欢跃？\BibleQuote{因生活的天主。}你内什么欢跃？\BibleQuote{我的心灵和我的肉身。}它们为何欢跃？他说：\BibleQuote{因为，麻雀找到了房屋，斑鸠筑了窝，在何处产雏。}\BibleRef{咏 83:3} 这是什么？他提到了两样，又加了两种与它们相配的鸟形：他曾说他的心灵和肉身欢跃，而他用麻雀和斑鸠来对应这两者：心灵如麻雀，肉身如斑鸠。麻雀找到了房屋：我的心灵找到了房屋。它在现世的美德中——信仰、希望、爱德——尝试振翅，藉此飞向它的房屋；当它到了那里，便安息下来；此处麻雀哀鸣的声音将不再在那里。因为正是那哀鸣的麻雀，在另一篇圣咏中他说：\BibleQuote{像屋顶上孤零零的麻雀。}从屋顶它飞回房屋。现在让它待在屋顶，践踏它的血肉房屋：它将得到天上的房屋，永久的居所：那只麻雀将止息哀鸣。但斑鸠得了幼雏，即肉身：\BibleQuote{斑鸠筑了窝，在何处产雏。}麻雀得房屋，斑鸠得窝，而且是产雏的窝。房屋是作为永远的选择，窝是暂时构架的：我们用心灵思想天主，仿佛麻雀飞回它的房屋；我们用肉身行善工。因为你们看见圣人们的肉身行了多少善工；我们藉此实行所命我们做的事，藉此在今生得到帮助。\BibleQuote{把你饼分给饥饿的人，领穷苦无家之人到你屋里；若你见人赤身露体，就给他衣服穿；}\BibleRef{依 58:7} 以及其他命令我们的这类事，我们都只藉肉身而行……\newline{}我们说的，弟兄们，是你们知道的：有多少人似乎在教会之外行善工？甚至有多少外邦人供饥者食、衣裸者衣、接待旅客、探望病患、安慰囚犯？多少人这样做？斑鸠似乎也产雏：但找不到窝。异端者在教会外能行多少工？他们不把幼雏放在窝里。它们会被践踏、压碎：它们不能被保存、不能被保护……\newline{}在信仰中摆放你的幼雏：在那窝中作你的工。因为窝是什么，那窝是什么，紧接着就说：在说了\BibleQuote{斑鸠找到了窝，在何处产雏}\BibleRef{咏 83:3}之后，好像你问：什么窝？\BibleQuote{上主、万军的天主，我的君王，我的天主，祢的祭台。}\BibleQuote{我的君王，我的天主}是什么意思？祢统治我，祢造了我。\par

8.……\BibleQuote{居住在你殿里的人是有福的}\BibleRef{咏 83:4}……如果你有自己的房屋，你是贫穷的；如果拥有天主的房屋，你是富有的。在你自己的房屋里，你会畏惧盗贼；天主的房屋，祂自己就是墙垣。因此，\BibleQuote{居住在你殿里的人是有福的}。他们拥有天上的耶路撒冷，没有限制，没有压力，没有界限的分歧和划分；所有人都拥有它，每人都有全部。这些财富是伟大的。弟兄不拥挤弟兄：那里没有匮乏。接着，他们将在那里做什么？因为对人而言，需要乃是一切工作的母亲。我简要地说，弟兄们，你们在心里思考任何一种职业，看看是否独有需要产生了它们。那些似乎能给予他人很大帮助的卓越技艺，如辩护之词或医治之药，因为这些是此生中最优秀的职业；拿走诉讼者，辩护人还能帮助谁？拿走伤口和疾病？医生还能治愈什么？我们日常生活的所有那些必需的工作，都源于需要。耕田、播种、犁地、航海；是什么产生这一切工作，不正是需要和匮乏吗？拿走饥饿、口渴、赤身；谁还需要这些东西？……例如，命令：\BibleQuote{把你的饼分给饥饿的人}。如果没有人饥饿，你能分饼给谁？\BibleQuote{接纳无家可归的穷人到你家里}\BibleRef{依 58:7}。如果所有人都住在自己的家乡，还能接纳谁？如果他们都享有永久的健康，还能探望谁？如果那里有永久的和平，还能和解谁？如果有永生，还能埋葬谁？那时，没有一件世人共有的高尚行为会是你的事业，也没有任何这些善工；幼燕将飞离它的巢穴。那么怎样？你已说过我们将拥有什么：\BibleQuote{居住在你殿里的人是有福的}。现在说说他们将做什么，因为我看不到那时还有任何需要驱使我行动。即使我现在正在说的和争论的，也源于某种需要。在那里，会有这样的辩论来教导无知者或提醒健忘者吗？或者，福音会在那默观天主圣言本身的地方被诵读吗？……\BibleQuote{他们要不断赞美你}。这将是我们全部的本分，一场不停的阿肋路亚。弟兄们，不要以为那里会有任何厌倦：如果你们在这里不能长久忍受说这番话，那是因为某种需要将你们从那享受中拉走。如果未见之物在此不能给予如此大的喜乐，如果在肉体的压迫和软弱下，我们如此热切地赞美我们所相信的，那么我们将如何赞美我们所看见的？\BibleQuote{当死亡在胜利中被吞没，当这必死的穿上不朽时}\BibleRef{格前 15:54}，没有人会说：\Quote{我已经站了很久；}没有人会说：\Quote{我已经禁食很久了，}\Quote{我已经看守很久了。}因为那里将有大忍耐，我们不朽的身体将在对天主的默观中得支柱。如果我们现在分施给你们的圣言能支持你们软弱的肉体站立这么久，那喜乐将产生什么效果？它将如何改变我们？\BibleQuote{因为我们必要像祂，因为我们要看见祂实在怎样}\BibleRef{若壹 3:2}。变得像祂，我们何时会昏厥？什么会拉走我们？弟兄们，我们永不会厌倦赞美天主、爱天主。如果爱会衰竭，赞美也会衰竭。但若爱是永恒的，因为那里将有不可穷尽的美，不要害怕你不能永远赞美那你能永远爱的那位。让我们为这生命叹息吧。\par

9. 但我们如何能到那里？\BibleQuote{那以你为坚强的人是有福的}\BibleRef{咏 83:5}。他知道自己在哪里，并且由于肉体的脆弱，他不能飞到那福乐之境：他思虑自己的重担，正如别处所说：\BibleQuote{因为可朽坏的身体压着灵魂，这地上的帐幕阻挠那多虑的思想}\BibleRef{智 9:15}。圣神向上呼唤，肉体的重量向下拉回：在提升和压下的双重努力之间，产生了一种挣扎：这挣扎朝向酒榨的压力。请听宗徒如何描述这同一酒榨的挣扎，因为他自己曾在那里受折磨，在那里被压榨……\BibleQuote{我这个人真不幸呀！谁能救我脱离这死亡的身体呢？天主的恩宠，借着我们的主耶稣基督}\BibleRef{罗 7:24}……\Quote{因为我按内心，喜爱天主的法律。}但我该做什么？我如何飞升？我如何到达那里？\Quote{我看见在我的肢体中另有法律}等等……正如宗徒的话，那困难与几乎无法解脱的挣扎因添加\Quote{借着我们的主耶稣基督的天主恩宠}而得以缓解；同样，当他因热切渴望天主的房屋和那些对天主的赞美而叹息，并且当因感到身体的重担和肉体的重量而产生一种绝望时，他再次唤醒希望，说道：\BibleRef{咏 83:5}\BibleQuote{那以你为依靠的人是有福的}。\par

10. 那么，天主藉祂的恩宠，向祂所抓住以引导前进的人，供应什么呢？祂接着说：\Quote{他在他的心中设置了阶梯。}…这阶梯设在何处？\Quote{在他的心中，在涕泣之谷} \BibleRef{咏 83:6}。因此，在此你们有一个酒榨，就是涕泣之谷，磨难中极虔诚的眼泪乃是爱之人的新酒。…他们出去时\Quote{流泪，}他说，\Quote{撒种。}因此，愿藉天主的恩宠，上行的阶梯被安置在你们心中。藉着爱而上升。因此，这首\Quote{登阶}圣咏被称为…\Quote{他设置了上升的阶梯，到他所指定的地方} \BibleRef{咏 83:7}。如今我们哀叹；我们的哀叹从何而来，岂非来自那放置我们上升阶梯的地方？我们的哀叹从何而来，岂非来自那原因，宗徒因此呼唤自己是一个可怜的人，因为他看到在自己的肢体中存在另一条法律，与心中法律的交战？\BibleRef{罗 7:23} 这又从何而来？来自罪的惩罚。我们曾以为，在我们领受诫命之前，我们凭自己的力量就能轻易成为义人；\Quote{但诫命一来，罪复活了；我却死了，}\BibleRef{罗 7:9}宗徒说。因为法律赐给人，并非能立刻拯救他们，而是要向他们显示他们正处在何等严重的疾病之中。…但当藉着赐下的法律显明了罪时，罪反增大了，因为它既是罪，又是违背法律；\Quote{罪，}他说，\Quote{藉着诫命寻得机会，在我内生出各种贪欲。}\BibleRef{罗 7:8} 他所说的\Quote{藉着法律寻得机会}是什么意思？人领受了诫命，便试图凭自己的力量去遵守；被情欲征服，他们便也成了违犯这诫命的罪人。但宗徒说什么？\Quote{罪恶在哪里增多，恩宠在哪里就更丰富地增多；}\BibleRef{罗 5:20} 也就是说，疾病加重，药物变得更为有效。因此，我的弟兄们，撒罗满那五道走廊，中间有水池，曾治好病人吗？病人，圣史说，躺在五道走廊里。\BibleRef{若 5:3} 在福音中我们有读到。那五道走廊就是梅瑟五书中的法律。为此，病人从家里被抬出来，好躺在走廊里。所以法律带出了病人，却没有治好他们；但藉着天主的祝福，水被搅动，如一位天使降临其中。看见水搅动，那个能下去的人就下去而痊愈了。那被五道走廊围住的水，就是被圈在法律中的犹太人。主来了，搅动了这百姓，以致祂自己被杀害。因为如果主没有降临下来搅动犹太人，祂会被钉十字架吗？所以，被搅动的水象征主的受难，这受难源于祂搅动犹太百姓。病人相信这受难，如同下到搅动的水中的人，因此而痊愈。法律不能治愈的人，即当他躺在走廊里时，藉着恩宠，藉着对我们主耶稣基督受难的信德而痊愈。\par

11. \Quote{祂要赐福，}他说，\Quote{那赐下法律的。}…恩宠将在法律之后来到，恩宠本身就是福。那恩宠和福赐给我们什么？\Quote{他们将由能力走向能力。}因为在这里藉着恩宠赐下许多能力。\Quote{因为藉着圣神，赐给这人智慧的话，赐给那人知识的言语，同一圣神；赐给另一人信德，另一人治病的恩赐，另一人行异语的能力，另一人解释异语，另一人先知话。}\BibleRef{格前 12:8} 许多能力，但对今生是必要的；从这些能力我们走向\Quote{一个能力。}哪一个\Quote{能力}？\Quote{基督，天主的德能和天主的智慧。}\BibleRef{格前 1:24} 祂在此地赐下不同的能力，祂将用自己取代所有在涕泣之谷中必要和有用的能力。因为在圣经和许多作者中，描述了四种对生活有用的德性：明智，藉以辨别善恶；正义，藉以给予每人所应得的，\Quote{不欠任何人什么，}\BibleRef{罗 13:8} 但爱所有的人；节制，藉以抑制情欲；刚毅，藉以忍受一切困苦。这些德性如今藉着天主的恩宠在涕泣之谷中赐给我们；从这些德性我们攀登到那另一种德性。那是什么，岂不是唯独默观天主的德性？…接着那里说：\Quote{他们将由能力走向能力。}什么能力？默观的能力。什么是默观？\Quote{万神之神将在熙雍显现。}万神之神，基督徒的基督。…当一切都完成，那必朽性所需要的一切，祂要向心地纯洁的人显现，正如祂所是的，\Quote{天主与天主同在，}圣言与圣父同在，\Quote{万物藉祂而被造。}\par

12. 再次，从对那喜乐的思念，他回到自己的叹息中。他看见在希望中已来临的事物，和他实际所处的位置。…因此，回到这地方特有的呻吟，他说：\Quote{万军的上主天主，求你俯听我的祈祷：雅各伯的天主，求你侧耳倾听}\BibleRef{咏 83:8}：因为正是你，从雅各伯中使他成为以色列。因为天主向他显现，他被称为以色列，\BibleRef{创 32:28} 即看见天主者。因此，雅各伯的天主，求你俯听我，使我成为以色列。我何时成为以色列？当万神之神在熙雍显现之时。\par

13. \Quote{天主，我们的护盾，求你垂顾。并垂视你基督的面容}\BibleRef{咏 83:9}。因为天主何时不垂视祂基督的面容？这是什么意思，\Quote{垂视你基督的面容}？藉着面容我们被认识。那么，垂视你基督的面容是什么意思？使你的基督被所有人认识。垂视你基督的面容：愿基督被所有人认识，好使我们能由能力走向能力，好使恩宠增多，因为罪恶增多了。\par

14. \Quote{在祢的宫院里住一日，胜过千日}\BibleRef{咏 83:10}。他叹息、他渴慕的正是这些宫院。\Quote{我的灵魂渴慕并憔悴于上主的宫院：}在那里一日胜过千日。人们渴慕千日，盼望在此世长寿：让他们轻视这千日，渴慕那一日吧，那一日既无升起也无落下：一日，永恒之日，无昨日可让位，无明日可压迫。让我们渴慕这一日。我们要与千日有何相干？我们从千日走向那一日；让我们奔向那一日，从力量走向力量。\par

15. \Quote{我宁愿被遗弃在上主的殿宇里，也不愿住在恶人的帐幕中}\BibleRef{咏 83:11}。因为他找到了流泪谷，找到了使他升起的谦卑：他知道若自高必跌倒，若自谦必被高举；他选择了被遗弃，好被提升。在这上主压酒池的帐幕旁，即在天主教会旁，有多少人渴望被高举，喜爱尊荣，拒绝认识真理。若这节经文常在他们心中，他们岂不抛弃尊荣，奔向流泪谷，从而在心里找到上升之路，从德行走向德行，将他们的希望寄托于基督，而不寄托于某个人？这话真好，是可喜乐的话，是可选择的话。他自己选择了在上主的殿宇中被遗弃；但那邀请他赴宴的，在他选择了低座时却叫他到高座，对他说：\Quote{请你上座。}\BibleRef{路 14:10} 然而他只选择在上主的殿宇中，无论在殿中何处，只要不在门槛之外。\par

16. 他为何选择？……\Quote{因为上主喜爱仁慈和真理}\BibleRef{咏 83:12}。上主喜爱仁慈，祂首先以此仁慈来帮助我：祂喜爱真理，以此将所许诺的赐予相信祂的人。\BibleRef{罗 11:29} 听听宗徒保禄的身上所显示的仁慈和真理吧：保禄先前是扫禄，是迫害者。他需要仁慈，并说仁慈已向他显示：\Quote{我先前是亵渎者、迫害者和施暴者；但我蒙受了仁慈，为使基督耶稣在我身上显示全部的忍耐，给那些将信祂而得永生的人。}如此，当保禄如此大罪获赦时，无人应绝望于任何罪过的赦免。看！你已有了仁慈。……看，我们看见保禄视祂为负债者，既受了仁慈，便要求真理。他说：上主在那日必要偿还。祂要偿还你什么？岂不是祂所欠你的？祂如何欠你？你给了祂什么？\Quote{谁先给了祂，而使祂偿还呢？}\BibleRef{罗 11:29} 上主使自己成为负债者，并非因接受，而是因许诺：并不是对祂说：还清你所接受的，而是：还清你所许诺的。他说：祂向我显示了仁慈，使我成为无辜的；因为我先前是亵渎者和施暴者，但藉祂的恩宠我成了无辜的。然而，那首先显示仁慈的，岂能否认祂的债？\Quote{祂喜爱仁慈和真理。祂要赐下恩宠和光荣。}何种恩宠？即那位所说：\Quote{因天主的恩宠，我成为今日的我}\BibleRef{格前 15:10}。何种光荣？即他所说：\Quote{为我预备了正义的冠冕}\BibleRef{弟后 4:8}。\par

17. 因此\Quote{上主不会将福乐赐予行为正直的人}\BibleRef{咏 83:12}。那么，人哪，你们为何不愿保持正直，除非为了获得福乐？……你们看见财富在强盗、不敬者、邪恶者、卑贱者手中；在丑恶和罪恶的人手中你们看见财富：天主因他们同属人类，因祂丰沛的恩泽，将这些赐予他们；祂也\Quote{使太阳升起，光照善人，也光照恶人；降雨给义人，也给不义的人}\BibleRef{玛 5:45}。祂赐给恶人如此之多，难道不为你留下什么？祂留下了；放心吧！祂曾怜悯你，当你还不敬时，难道在你成为虔诚者后遗弃你？祂曾将祂圣子死亡的无偿恩赐给了罪人，难道不为那藉此死亡得救的人保留什么？所以放心。你因相信祂的许诺而视祂为负债者。那么，在这里——在压酒池中，在苦难中，在困厄中，在这危险的现世生活中——我们还有什么留下？我们剩下什么才能到达那里？\Quote{上主，万军的天主，那全心依靠祢的人，真是有福。}\par

\SourceRecord{Translated by J.E. Tweed.；Nicene and Post-Nicene Fathers, First Series；Vol. 8.；Edited by Philip Schaff.；Buffalo, NY: Christian Literature Publishing Co.,；1888.；<http://www.newadvent.org/fathers/1801084.htm>.}

% structure: p0001 source=h1 target=section reason=page-title

\section{圣咏第八十五篇释义}

1. …它的标题是：\Quote{为终的圣咏，属科辣黑子孙。} 我们应明白，除了宗徒所说的终以外，再没有别的终：因为\Quote{基督是法律的终向。}\BibleRef{罗 10:4} 因此，当他在圣咏标题的开头写下\Quote{为终}时，便将我们的心引向基督。如果我们定睛注视祂，就不会迷路：因为祂就是我们所渴望到达的真理本身，也是我们所奔跑的道路。\BibleRef{若 14:6}\par

2. 先知向祂歌唱未来，却使用仿佛过去时的言辞：他讲述未来的事如同已经成就，因为在天主眼中，未来之事已经发生。……\Quote{主，祢已恩待祢的土地。}\BibleRef{咏 84:1} 仿佛祂已经这样做了。\Quote{祢已扭转了\Person{雅各伯}的被掳。} 祂的古老子民\Person{雅各伯}，即\Place{以色列}子民，由\Person{亚巴辣罕}的子孙所生，在应许中注定有一天成为天主的继承人。那确实是一个真实的民族，\WorkTitle{旧约}赐给了他们；但在\WorkTitle{旧约}中，\WorkTitle{新约}被预表：那是预像，这是表达的真理。在那个预像中，借着某种对未来的预言，应许之地赐给了那个民族，位于犹太人所居住的地区；那里也有\Place{耶路撒冷}城，我们大家都听说过它的名字。当这个民族获得这块土地后，他们遭受了周围许多仇敌的困扰；当他们得罪了自己的天主时，他们被交到被掳中，不是为毁灭，而是为惩戒；他们的父不是定他们的罪，而是鞭打他们。被掳之后，他们获得释放，多次既被掳又获释；他们如今仍处于被掳之中，那是因为一个大罪，就是他们钉死了自己的主。那么，我们该如何理解他们所说的\Quote{祢已扭转了\Person{雅各伯}的被掳}呢？……这篇圣咏以诗歌预言了。\Quote{祢已扭转了\Person{雅各伯}的被掳。} 这是对谁说的？对基督；因为经上说：\Quote{为终，属科辣黑子孙：} 因为祂已扭转了\Person{雅各伯}的被掳。请听\Person{保禄}自己承认：\Quote{我这个人真不幸！谁能救我脱离这该死的身体呢？} 他问是谁能，立刻想到：\Quote{藉着我们的主\Person{耶稣基督}，天主的恩宠。}\BibleRef{罗 7:24} 先知论到这种天主的恩宠，对我们的主\Person{耶稣基督}说：\Quote{祢已扭转了\Person{雅各伯}的被掳。} 请留意\Person{雅各伯}的被掳，留意，并看这就是：祢扭转了我们的被掳，不是将我们从并未遭遇的蛮族人中释放，而是将我们从恶行中、从我们的罪恶中释放，\Person{撒殚}曾藉以控制我们。因为若有人从罪恶中被释放，罪人之王便没有地方可以控制他了。\par

3. 因为祂如何扭转了\Person{雅各伯}的被掳？看，那释放是属神的，看，它是如何在内在完成的。\Quote{祢已赦免，}他说，\Quote{祢子民的罪孽：祢已遮盖了他们所有的罪过。}\BibleRef{咏 84:2} 你看，祂如何扭转了他们的被掳，因为祂赦免了罪孽：罪孽曾俘虏他们；你的罪孽被赦免，你就获得自由。因此，你要承认你是在被掳中，好配得释放：因为不认识自己仇敌的人，怎能呼求释放者？\Quote{祢已遮盖了他们所有的罪过。} 什么是\Quote{祢已遮盖}？就是不去看它们。祢如何不看它们？就是不报复它们。祢不愿看我们的罪过：因此祢不看它们，因为祢不愿看它们：\Quote{祢已遮盖了他们所有的罪过。} \Quote{祢已平息了祢所有的愤怒：祢已从祢的震怒中转离。}\BibleRef{咏 84:3}\par

4. 虽然这些事是关乎未来说的，尽管言辞的时态是过去时，接着是：\Quote{天主，我们的救恩，求祢扭转我们。}\BibleRef{咏 84:4} 他刚才叙述仿佛是已经成就的，为何又祈求实现呢？岂不是因为他想表明，他先前在预言中是作为过去时说的？但他所说已成就的尚未成就，他借此表明，他祈求它实现：\Quote{天主，我们的救恩，求祢扭转我们，并求祢从我们身上转离祢的愤怒。} 你先前不是说过：\Quote{祢已平息了祢所有的愤怒，祢已从祢的震怒中转离}吗？怎么现在又说\Quote{求祢从我们身上转离祢的愤怒}？先知回答说：我之所以说这些事是已经成就的，因为我看它们将要成就；但因为它们尚未成就，我祈求它们来临，而我已经看见了。\par

5. \BibleQuote{求祢不要向我们永远发怒} \BibleRef{咏 84:5}。因为因天主的愤怒，我们必死；因天主的愤怒，我们在世上面临匮乏、汗流满面地吃饼。\BibleRef{创 3:19} 这是亚当犯罪时所受的判决：那亚当就是我们每个人，因为\BibleQuote{在亚当内众人都死了；}\BibleRef{格前 15:22} 对他的判决也延及我们。因为我们那时尚不存在，但我们已在亚当内：所以凡发生在亚当身上的事也发生在我们的身上，以致我们都要死：因为我们都曾在祂内....至此，你父亲的罪不会伤害你，若你改变了自己；同样，若他改变了自己，也不会伤害你。但我们的本性所承受的死亡轭，是从亚当而来。它承袭了什么？肉身的软弱、痛苦的折磨、贫穷的居所、死亡的锁链、诱惑的罗网：这一切我们都带在肉身中；这就是天主的愤怒，因为这是天主的报复。但由于我们应当重生，藉信德成为新人，一切必死性将在复活中被除去，整个人要在新生命中恢复；\BibleQuote{就如众人都在亚当内死，照样众人也都在基督内活过来；}\BibleRef{格前 15:22} 有鉴于此，先知说：\BibleQuote{求祢不要向我们永远发怒，也不要把祢的忿怒从这一代伸展到另一代。}第一代因祢的忿怒成了必死的；第二代因祢的仁慈将是不死的……\par

6. \BibleQuote{天主啊，祢必将使我们回转，并使我们活着} \BibleRef{咏 84:6}。并非好像我们凭自己自愿、没有祢的仁慈就转向祢，然后祢使我们活着；而是这样：不仅我们被活着是从祢而来，而且我们的皈依本身，也是为了使我们活着。\BibleQuote{祢的子民要因祢而欢乐。}他们为自己的恶而因自己欢乐；为自身的善而因祢欢乐。因为他们曾想从自身获得欢乐，就在自身中找到了祸；但现在因为天主是我们的一切欢乐，凡要稳妥地欢乐的，让他因那不可朽坏者而欢乐。因为，我的弟兄们，你们为何要因银子欢乐呢？或你的银子朽坏，或你朽坏：无人知道哪个在先；然而这是确定的，两者都将朽坏；哪个在先，是不确定的。因为人不能常存于此，银子也不能常存于此：金子、衣服、房屋、金钱、广阔的土地，最后，连这光本身也一样。因此你不要愿意因这些而欢乐；而要因那永远不落的光而欢乐：因那没有昨日在前的黎明而欢乐，也没有明天接续的黎明。那是什么光？\BibleQuote{我}，祂说，\BibleQuote{是世界的光。}\BibleRef{若 8:12} 那对你们说\BibleQuote{我是世界的光}的那一位，将你们召唤到祂身边。当祂召唤你们时，祂使你们皈依；当祂使你们皈依时，祂医治你们；当祂医治了你们，你们就要看见你们的皈依者，向祂曾说：\BibleQuote{求祢向我们显示祢的仁慈，上主，并赐给我们祢的救恩}\BibleRef{咏 84:7}：祢的救恩，就是祢的基督。有福的，是天主显示仁慈于他的人。天主显示仁慈于他的人，不能沉溺于骄傲。因为通过向他显示祂的救恩，祂说服他：人所拥有的一切善，并非出于自身，而是出于那作为我们一切善的祂。当人看见他所拥有的一切善并非出于自己，而是出于他的天主时；他就看见在他内一切可赞美的事都是出于天主的仁慈，而非自己的功劳；看见这事，他就不骄傲；不骄傲，就不被高举；不举高自己，就不堕落；不堕落，就站立；站立，就紧紧依附；紧紧依附，就存留；存留，就享受并因他的主、他的天主而欢乐。那造了他的将成为他的喜乐：这喜乐无人能破坏，无人能打断，无人能夺去....因此，若望在他的书信中说了什么？\BibleQuote{亲爱的，现在我们已是天主的子女，将来如何还未显明。}\BibleRef{若壹 3:2} 谁不会欢喜呢？假如他正在外漂泊，不知自己的出身，忍受匮乏，处于痛苦和辛劳中，忽然有人宣布说：你是元老的儿子；你父亲在祖业上拥有丰厚的家产；我命令你回到你父亲那里：如果有一个他能信任其承诺的人对他说这话，他岂不欢喜？一位我们能信任的人，基督的宗徒，来对我们说：你们有父亲，你们有祖国，你们有产业。那父亲是谁？\BibleQuote{亲爱的，我们是天主的子女。}\BibleRef{若壹 3:2} ...因此祂应许向我们显示祂自己。想想吧，我的弟兄们，祂的美是怎样的。你们所见所爱的一切美好事物，都是祂造的。如果这些是美好的，祂自己该是怎样的？如果这些是伟大的，祂有多么伟大？因此，从我们在此所爱的这些事物，让我们更加渴慕祂；轻看这些事物，让我们爱祂；好使我们藉着这爱因信德净化我们的心，当祂的显现来临时，能发现我们的心已被净化。那将要显示于我们的光应当发现我们是健全的：这就是现在信德的工作。这就是我们在此所说的：\BibleQuote{并赐给我们祢的救恩：}赐给我们祢的基督，使我们认识祢的基督，看见祢的基督；不像犹太人看见祂并钉死祂，而要像天使看见祂而欢乐。\par

7. \Quote{我要倾听} \BibleRef{咏 85:8}。先知发言：天主在他内发言，而世界在外喧嚷。因此，他稍从世界的喧嚷中退出，转回自身，又从自身转向他在内里所听到声音的那一位；他仿佛封住耳朵，对抗今生的纷扰不安，对抗那被可朽坏的身体所重压的灵魂，对抗那因尘世的帐幕压迫而思虑万端的想象 \BibleRef{智 9:15}，他说：\Quote{我要倾耳听从我内里说话的主天主；}他听到了什么？\Quote{因为祂要向祂的百姓说和平。}因此，基督的声音，天主的声音，就是和平；它召唤人归向和平。听！它说：凡尚未在和平中的人，你们要爱和平：因为从我这里，你们能找到什么比和平更好的吗？什么是和平？就是没有战争的地方。这是什么？是没有战争的地方？就是没有矛盾，没有抗拒，没有反对之处。请想想我们是否已在那里：请想想现在是否没有与魔鬼的争战，是否所有圣人及信者不与群魔之首角力。他们怎能与看不到的他角力呢？他们与自己的私欲角力，魔鬼藉此向他们暗示罪恶：而由于不同意他所暗示的，他们虽未被征服，却仍在战斗。因此，有战斗之处就没有和平。……我们为振奋自己所安排的一切，又在那里找到疲乏。你饿了吗？有人问你：你回答，我饿了。他摆上食物供你振奋；你继续享用，因为你需要它；然而在继续享用你需要振奋之物时，你从中找到疲乏。久坐使你疲倦；你站起来，藉行走振奋自己；继续那振奋，多行又使你疲倦；你又要坐下。请给我找出任何能振奋你、且你若持续在其中便不再疲倦的东西。那么，人在这里所拥有的和平，被如此多的烦恼、欲望、匮乏、疲惫所对抗，又是什么呢？这不是真正的、完全的和平。完全的和平是什么？\Quote{这必朽坏的必须穿上不朽坏的，这必死的必须穿上不死的。} \BibleRef{格前 15:53}……多食不停：这本身就会杀死你；多禁食不停：你会因此而死；一直坐着，决心不起，你会因此而死；总在行走从不休息，你会因此而死；一直醒着不睡觉，你会因此而死；一直睡觉从不醒着，你也会因此而死。因此，当死亡在胜利中被吞灭时，这些事将不再存在；那时将有圆满而永恒的和平。我们将置身于一座城，弟兄们，当我论及这城时，我难以住口，尤其当恶事频发时。谁不渴望那城呢？那里没有朋友出去，没有敌人进入，没有诱惑者，没有煽动叛乱者，没有分裂天主百姓的人，没有为服事魔鬼而劳苦教会的人；因为这一切的首领本身已被投入永火，同那些同意他、且不愿离开他的人一起。那时，在天主子女中将有纯全的和平，彼此相爱，彼此看见充满天主，因为天主将成为万物中的万有。\BibleRef{格前 15:28} 我们将拥有天主作为我们共同注视的对象，天主作为我们共同的财产，天主作为我们共同的和平。因为凡祂现在赐予我们的一切，祂自己将成为我们，以代替祂的恩赐；这将是圆满而完全的和平。这就是祂向祂的百姓所说的和平；这也就是那位曾说\Quote{我要倾听主天主将对我说什么：因为祂要向祂的百姓、向祂的圣人、并向那些将心转向祂的人说和平}的人所要倾听的。看，我的弟兄们，你们愿意那和平属于你们吗？请将你们的心转向祂：不是转向我，或转向某一个人，或转向任何人。因为无论何人想把人的心转向自己，他必与他们一同跌倒。哪样更好呢？是与你所转向的人一同跌倒，还是与那你要转向祂的那一位一同站立？我们的喜乐、我们的和平、我们的安息、一切烦恼的终结，除了天主以外别无他物：\Quote{将心转向祂的人是有福的。}\par

8. \BibleQuote{然而，祂的救恩临近那敬畏祂的人}\BibleRef{咏 84:9}。在犹太民族中，当时就有一些敬畏祂的人。普世各地都崇拜偶像：人们敬畏魔鬼，而不敬畏天主；但在那个民族中，天主受到敬畏。但为何敬畏祂？在旧约中，人们敬畏祂，免得祂将他们交于俘虏，免得祂夺走他们的土地，免得祂用冰雹摧毁他们的葡萄园，免得祂使他们的妻子不生育，免得祂从他们那里夺走儿女。因为天主的这些属尘世的应许俘获了他们的心，他们的心尚微小，为了这些事而敬畏天主；但祂临近那些即使为了这些事而敬畏祂的人。外邦人为土地向魔鬼祈祷：犹太人为土地向天主祈祷：他们祈求的是同一件事，但所祈求的对象不同。犹太人虽然寻求外邦人所寻求的，却与外邦人不同；因为他向那创造万有的主祈求。天主远离外邦人，却临近犹太人；然而祂也眷顾那些远处的人，也眷顾那些近处的人，正如宗徒所说：\BibleQuote{祂来向你们远处的人，也向那近处的人，传报了和平。}\BibleRef{弗 2:17} 祂说的近处的人是谁？犹太人，因为他们崇拜独一真神。远处的人是谁？外邦人，因为他们离开了创造他们的主，去崇拜自己所造之物。因为任何人远离天主，不在于空间，而在于情感。你爱天主，你就临近祂；你恨天主，你就远离祂。你站在同一地方，既可以近也可以远。我的弟兄们，先知所关注的就是这一点：虽然他看见天主的慈悲普及众人，但他仍看到对犹太人显出的某种特殊和独特的恩惠，于是他说：\Quote{我要听天主上主说的话：祂要向祂的子民宣讲和平；}祂的子民将不仅限于犹太，而是从万国聚集而来：\Quote{因为祂要向祂的圣徒和那些全心归向祂的人宣讲和平}，也向普世所有全心归向祂的人宣讲和平。\Quote{然而，祂的救恩临近那敬畏祂的人，为使光荣居住在我们的地上：}就是说，在先知出生的那块土地上，更大的光荣将居住，因为基督开始从那里被宣讲。宗徒们来自那里，首先被派遣到那里；先知们来自那里，圣殿首先在那里，祭献在那里呈献给天主，圣祖们在那里，祂自己从亚巴郎的后裔中来到那里，基督在那里显现，基督在那里显现；因为童贞玛利亚从那里出生，她孕育了基督。祂在那里用脚行走，在那里行奇迹。第三，祂赐予那民族如此大的尊荣，以致当一位客纳罕妇人拦住祂，祈求医治她的女儿时，祂对她说：\BibleQuote{我被派遣，无非是为了以色列家亡失的羊。}\BibleRef{玛 15:24} 先知看到这一点，就说：\Quote{为使光荣居住在我们的地上。}\par

9. \BibleQuote{慈爱与真理彼此相遇}\BibleRef{咏 84:10}。\Quote{真理在我们的地上}，在犹太人身上，\Quote{慈爱}在外邦人的地上。因为真理在哪里？就在天主的话语所在之处。慈爱在哪里？就在那些离开了自己的天主而转向魔鬼的人身上。祂也垂顾了他们吗？是的，祂似乎说：召唤那些远处逃亡的人，他们远离了我：召唤他们，让他们找到我这位寻找他们的人，因为他们自己不愿寻求我。因此，\Quote{慈爱与真理彼此相遇：正义与和平互相亲吻}。行正义，你就能得和平；好让正义与和平互相亲吻。因为如果你不爱正义，你就不会得和平；因为这两者，正义与和平，彼此相爱，互相亲吻：这样，行了正义的人就能找到和平亲吻正义。它们两者是朋友：你或许只愿得到一样，而不愿得到另一样。因为没有人不愿得到和平。但所有的人都愿意行正义吗？问问所有的人，你愿得和平吗？全人类会异口同声地回答你：我愿意，我渴望，我情愿，我爱和平。也要爱正义：因为这两者，正义与和平，是朋友；它们互相亲吻：如果你不爱和平的朋友，和平本身也不会爱你，也不会来到你这里。因为渴望和平有什么了不起呢？每个恶人都渴望和平。因为和平是好事。但你要行正义，因为正义与和平互相亲吻，它们不会彼此争吵……\par

10. \BibleQuote{真理从地生出，正义由天俯视}\BibleRef{咏 84:11}。\Quote{真理从地生出：}基督由女人所生。天主子从肉身中出来。什么是真理？天主子。什么是地？肉身。问基督从哪里出生，你就看到\Quote{真理从地生出}。但那从地生出的真理，早在天地之先，并且天地是藉着祂造的；但为使正义由天俯视，即为使人类因天主的恩宠而成义，真理由童贞玛利亚诞生；这样祂就能奉献祭献以使人成义，就是苦难的祭献，十字架的祭献。祂怎能为我们献赎罪祭，除非祂死亡？祂怎能死亡，除非从我们接受了祂能藉以死亡的身体？也就是说，除非祂从我们接受了可朽的肉身，基督就不能死亡：因为天主的圣言不会死，天主性不会死，天主的德能和智慧不会死。祂怎么能奉献祭献，一个医治的牺牲，如果祂不死？祂怎么能死，除非祂穿上肉身？祂怎么能穿上肉身，除非真理从地生出？\par

11. 关于同一经文，我们可以提及另一层含义。\Quote{真理从地上生出}：即人的宣认。因为，人啊，你曾是个罪人。大地啊，当你犯罪时听到了判决：\Quote{你是土，还要归于土}，\BibleRef{创 3:19}，愿真理从你生出，好使正义从天俯视。真理如何从你生出？你既是罪人，既不义？宣认你的罪，真理便从你生出。因为，你若不义却自称义，真理怎能从你生出？但若不义者宣认自己是不义的，\Quote{真理便从地上生出}。……\Quote{正义从天俯视}是什么？是天主的正义，好似祂说：我们宽恕这人吧，因为他没有宽恕自己；我们赦免他吧，因为他亲自承认。他改变自己以惩罚自己的罪，我也要改变以释放他。\par

12. \Quote{因为上主必赐甘甜，我们的地也要出产} \BibleRef{咏 84:12}。……祂将赐给你行义的甘甜，以致从前不义令你喜悦的，如今义德开始令你喜悦；以致你起初以醉酒为乐的，如今以节饮为乐；你起初以偷窃为乐的，夺取他人所无的，如今却寻求将你所有的给予那没有的；你以抢劫为乐的，如今以施舍为乐；你以观看为乐的，如今以祈祷为乐；你以轻浮淫荡的歌曲为乐的，如今以咏唱圣诗为乐；你起初跑向戏院的，如今跑向教会。这甘甜从何而生？岂不是因为 \Quote{天主赐下甘甜}？因为，看，你明白我的意思；看，我已向你宣讲了天主的话，我将种子撒在你虔诚的心田，发现你的灵魂已因告解的犁沟而松软；你以虔诚的专注接受了种子；现在请思考你所听的话，如同那些捣碎云彩的人，免得飞鸟把种子叼走，使所撒的种子能在那里发芽；除非天主降雨在上面，撒种有何益处？这就是 \Quote{我们的地也要出产} 的意思。愿祂以祂的眷顾——闲暇、事务、家中、床上、用餐时、交谈中、路上——眷顾你的心，即使我们不在身旁。愿天主的雨降临，使播种的种子发芽；当我们不在身边、安静休息或忙于他事时，愿天主赐予我们所撒种子的成长，以致日后看到你们品格的改善，我们也因你们的果实而喜乐。\par

13. \Quote{因为正义要在祂面前行走，祂要在路上引导自己的脚步} \BibleRef{咏 84:14}：这正义就是在于宣认罪过：因为这就是真理本身。因为你应当对自己正义，惩罚自己：因为这是人正义的开端，即你应当惩罚自己这个恶人，而天主使你成为善人。因此，既然这是人正义的开端，这便成为天主来临的道路，使天主能来到你面前：在那里，在罪过的宣认中，为祂预备一条路。因此，当若翰以悔改的水施洗，并要人们来到他那里悔改先前行为时，他如此说：\Quote{预备主的道路，修直祂的途径。} 你曾在自己罪中自娱，人啊：愿你所是的令你不悦，以便你能成为你所不是的。预备主的道路：愿那宣认罪过的正义先行，祂必来临并眷顾你，因为如今祂有了安放脚步之处，有了来到你面前的道路。在你宣认罪过之前，你关闭了天主的路：没有路使祂能来到你面前。宣认你过往的生活，你就打开了一条路；基督就来到你面前，\Quote{在路上安置祂的脚步}，好以祂自己的脚步引导你。\par

\SourceRecord{Translated by J.E. Tweed.；Nicene and Post-Nicene Fathers, First Series；Vol. 8.；Edited by Philip Schaff.；Buffalo, NY: Christian Literature Publishing Co.,；1888.；<http://www.newadvent.org/fathers/1801085.htm>.}

% structure: p0001 source=h1 target=section reason=page-title

\section{圣咏第八十六篇释义}

1. 天主赐予人最大的恩赐，莫过于使祂的圣言——祂藉以创造万物的圣言——成为他们的元首，并将他们作为肢体与祂联合：使天主子也成为人子，与圣父同一天主，与人同一人；如此，当我们向天主祈求怜悯时，我们并不将圣子与祂分离；而当圣子的身体祈祷时，它并不将它的元首与自己分离；而我们的主耶稣基督，天主子，是祂身体的唯一救主，祂为我们祈祷，在我们内祈祷，也受我们祈祷。祂作为我们的司铎为我们祈祷；祂作为我们的元首在我们内祈祷；祂作为我们的天主受我们祈祷。因此，让我们在祂内认出我们的话语，并在我们内认出祂的话语。当论及我们的主耶稣基督，尤其在预言中，若含有低于天主尊威的谦卑程度时，我们不要迟疑将其归于祂，因为祂并不迟疑将自己与我们联合。……祂以天主的形态受祈祷，以仆人的形态祈祷；在那里是创造者，在这里是受造者；祂未改变地取了受造者，使之得以改变，并使我们与祂自己成为同一人，即头和身体。因此，我们向祂祈祷，藉着祂，在祂内；我们与祂说话，祂与我们说话；我们在这篇圣咏的祷词内，在祂内说话，祂在我们内说话；这篇圣咏的标题是\Quote{达味的祈祷}。因为我们的主按肉身是达味之子；但按祂的天主性，是达味的主，也是他的创造者。……因此，当任何人听到这些话时，不要说\Quote{基督不说}；也不要说\Quote{我不说}；反之，如果他承认自己在基督的身体内，让他说：基督说，我也说。你不可愿意没有祂而说什么，祂也不没有你而说什么。……\par

2. \BibleQuote{主，求祢侧耳，俯听我}\BibleRef{咏 85:1}。祂以仆人的形态说话；仆人，你在你主的形式中说话吧：\Quote{主，求祢侧耳。}祂侧耳，你若不抬高颈项；因为祂亲近谦卑的人；远离自高的人，除非祂自己抬举那本为谦卑的人。天主向我们侧耳。因为祂在上，我们在下；祂在高处，我们在低处，却未被遗弃。\BibleQuote{因为当我们还是罪人的时候，基督已为我们死了。为义人死，是罕有的；为好人，或许有人敢死。}\BibleRef{罗 5:8}但我们的主为恶人死了。因为我们没有任何先行功绩，使天主子应为我们而死；正因没有功绩，祂的仁慈更大。那么，祂为义人保存生命的应许是多么确定、坚固，祂曾为恶人舍去自己的生命！\Quote{因为我贫穷而可怜。}祂不侧耳听富足的人；而是侧耳听贫穷可怜的人，就是谦卑和忏悔的人，需要怜悯的人；而不是那自满、自高自夸、仿佛一无所缺的人，他说：\BibleQuote{我感谢祢，因为我不像这个税吏。}\BibleRef{路 18:11}因为富有的法利塞人夸耀自己的功德；贫穷的税吏承认自己的罪过。\par

3. 然而，我的弟兄们，不要按我所讲的以为天主不垂听那些拥有金银、仆役和田产的人——如果他们是继承这些家业或拥有这样的世位；只要他们记住宗徒的话：\BibleQuote{劝戒今世富足的人，不要心高气傲。}\BibleRef{弟前 6:17}因为那些不心高气傲的人，在天主眼中是贫穷的，祂侧耳倾听贫穷、可怜和匮乏的人。因为他们知道自己的盼望不在于金银，也不在于他们一时似乎富足的事物。财富不毁灭他们，就足够了；财富不伤害他们，就足够了；因为财富不会对他们有益。真正有益的是仁爱的工作，由富人或穷人完成：富人用意愿和行为，穷人单用意愿。因此，当一个人如此，鄙视自己心中一切惯于使人傲慢的事，他就是天主的穷人；天主向他侧耳，因为祂知道他的心痛悔。……难道那穷人因贫穷的功劳被天使带去，\BibleRef{路 16:19}而那富人因财富的罪过被送受折磨吗？在那穷人身上，是谦卑所受的尊荣；在那富人身上，是骄傲所等待的惩罚。我马上要证明，那富人受折磨的不是财富，而是骄傲。可以肯定，那穷人被抱到亚巴郎的怀中；关于亚巴郎自己，圣经说他在地上有许多金银，是富有的。\BibleRef{创 13:2}如果每个富人都被迅速带去受折磨，亚巴郎怎能比那穷人先到，以致准备好接纳他被抱到怀中呢？但亚巴郎在他的财富中是贫穷的、谦卑的、敬畏一切诫命并服从的。他视所有财富为虚无，以至于因天主的命令，他准备祭献自己的儿子\BibleRef{创 22:10}，他本为儿子保存财富。因此，你们要学习成为贫穷可怜的，无论你们拥有世物与否。……\par

4. \BibleQuote{求祢保全我的灵魂，因为我是圣洁的} \BibleRef{咏 85:2}。我不知道是否有人能说\Quote{我是圣洁的}，但祂曾存在于世上而无罪，祂并未犯罪而是赦免一切罪。我们承认这是祂的声音说：\BibleQuote{求祢保全我的灵魂，因为我是圣洁的}；当然是在祂所取仆人的形象中。因为在其中是肉身，在其中也是灵魂。因为祂并非如某些人所言，仅是肉身与圣言：而是肉身与灵魂，以及圣言，而这一切是唯一天主子、唯一基督、唯一救主；在天主形象中与圣父同等，在仆人形象中为教会之首。因此，当我听见\Quote{因为我是圣洁的}，我认出祂的声音：难道我排除自己的声音吗？当然祂如此说话时是与祂的身体不可分离地说话。那么，我敢说\Quote{因为我是圣洁的}吗？如果作为使人圣洁的圣洁，且无需他人圣化，那便是骄傲和虚假；但如果是作为被圣化的圣洁，正如经上所记：\BibleQuote{你们应是圣洁的，因为我是圣洁的}，\BibleRef{肋 19:2}，那么基督的身体可以冒险，那个\Quote{从地极呼号}的人可以与他的首领同在，并在他的首领之下，说\Quote{因为我是圣洁的}。因为他已领受了圣洁的恩宠，圣洗圣事的恩宠，以及罪过的赦免。\BibleRef{格前 6:11}……要对你的天主说：我是圣洁的，因为祢圣化了我；因为我领受了，并非因为我有；因为祢赐予了，并非因为我堪当。因为从另一方面讲，你开始损害我们的主耶稣基督自己。因为如果所有在祂内因信德而受洗的基督徒都穿上了祂，正如宗徒所说：\BibleQuote{凡受洗归于基督的，就是穿上了基督}，\BibleRef{迦 3:27}如果他们已经成了祂身体的肢体，却说他们不是圣洁的，那就损害了他们的首领，他们作为肢体却仍不圣洁。你看你在哪里，并从你的首领领受尊严。因为你曾处于黑暗，\BibleQuote{但如今在主内成为光明}。\BibleRef{弗 5:8}\Quote{你们从前是黑暗}，他说；但你们仍是黑暗吗？难道光照者来临是为了你们仍留在黑暗，或是为了你们在祂内成为光明？因此，每个基督徒独自地，以及整个基督的身体，可以说起，可以在各处呼号，当它忍受苦难、各种诱惑和诽谤时，可以说起：\Quote{求祢保全我的灵魂，因为我是圣洁的：我的天主，求祢拯救祢的仆人，因为他信赖祢。}看，那圣洁的人并不骄傲，因为他信赖天主。\par

5. \BibleQuote{上主，求祢怜悯我，因为我终日向祢呼求} \BibleRef{咏 85:3}。不是\Quote{一天}，要理解\Quote{终日}意味着不断：从基督的身体在苦难中呻吟开始，直到世界的终结，当苦难过去时，那人呻吟并呼求天主：而我们各人按自己的限度，在整个身体的呼喊中有一份。你在你的日子里呼喊，而你的日子已经过去：另一个在你之后来临，在他的日子里呼喊：而你在这里，他在那里，另一个在别处：基督的身体终日呼喊，它的肢体交替离去和接续。那唯一的人延至世界的尽头：基督的同一肢体呼喊，有些肢体已安息在祂内，有些仍在呼喊，有些当我们安息时将呼喊，在他们之后其他人将呼喊。这是整个基督的身体，祂垂听它的声音，说：\BibleQuote{我终日向祢呼求}。我们的首领在天主圣父右边为我们转求：祂恢复一些肢体，鞭打另一些，洁净一些，安慰一些，创造一些，召叫一些，召回一些，纠正一些，复原一些。\par

6. \BibleQuote{求祢使祢仆人的灵魂喜乐，因为上主，我向祢举起了我的灵魂} \BibleRef{咏 85:4}。求祢使它喜乐，因为我向祢举起了它。因为它曾在地上，从地上感受了苦楚：免得它在苦楚中枯萎，免得它失去祢恩宠的一切甘甜，我向祢举起了它：求祢以祢自己使它喜乐。因为唯有祢是喜乐：整个世界充满苦楚。确实，祂有理由劝诫祂的肢体举起自己的心。愿他们听见并实行：愿他们向祂举起在地上软弱的事物。那里心不会朽坏，如果它被举向天主。如果你有谷物在楼下房间里，你会把它搬到高处，免得它腐烂。你愿意移动你的谷物，却任凭你的心在地上腐烂吗？你会把你的谷物搬到高处：将你的心举向天堂。而我怎能做到，你说？需要什么绳索？什么机器？什么梯子？你的情感是阶梯：你的意愿是道路。借着爱，你上升；借着忽视，你下降。站在地上，你就在天上，如果你爱天主。因为心不是像身体那样被举起：身体被举起会改变地方：心被举起会改变意志。\par

7. \BibleQuote{因为祢，上主，是良善宽仁的}\BibleRef{咏 85:5}……甚至祈祷也常被虚妄的思想所阻碍，以致心很难停留在天主身上；它想要持守，却仿佛从自身逃离，找不到可以围住自己的框架，没有栅栏可以拦住它的飞驰和游荡，好让它静止下来，在它的天主内享乐。在众多祈祷中，几乎难以找到一次这样的祈祷。每个人都可以说这事发生在自己身上，但不会说它没有发生在别人身上，除非我们在圣经中读到达味在一个地方祈祷说：\BibleQuote{上主，我既找到了我的心，好能向祢祈祷。}\BibleRef{撒下 7:27}他说他找到了自己的心，好像它习惯从他那里逃走，而他像追捕逃犯一样追赶它，却抓不住它，于是向天主呼喊：\Quote{因为我的心离弃了我。}因此，我的弟兄们，思考他在这里说的话，我想我明白了他所说的\Quote{宽仁}的意思。我似乎感到，他称天主为宽仁，是因为祂容忍我们的那些软弱，却仍期待我们祈祷，好使我们成全；当我们把祈祷献给祂时，祂感恩地接纳，垂听，不记念我们漫不经心倾泻出的众多祈祷，而接纳那几乎难以找到的一个。因为，我的弟兄们，有哪一个人，当他的朋友向他说话，他正要回答时，却看到朋友转身与别人交谈，谁能容忍这事？或者如果你向法官申诉，请他听你，而当你正对他说话时，突然离开他，开始与你的朋友交谈，谁能容忍这事？然而天主容忍了这么多人的心，他们祈祷时却想着不同的事……那么怎样呢？我们该对人类绝望，说每个在祈祷中渗入游荡思想并打断祈祷的人都已经定罪了吗？我的弟兄们，若我们这样说，我不知道还有什么希望存留。因此，因为在天主面前有希望，因为祂的仁慈伟大，让我们向祂说：\BibleQuote{因为上主，我曾向祢举起我的灵魂。}我是怎样举起的呢？按我所能的，按祢赐我力量的，按我能抓住它逃离时那样……我因软弱而沉沦：求祢医治我，我就站立；求祢坚固我，我就强壮。但在祢成就这事之前，祢容忍我：\BibleQuote{因为祢，上主，是良善宽仁的，且有大仁慈。}这就是说，不仅是\Quote{仁慈}，而且是\Quote{大仁慈}：因为正如我们的不义增多，祢的仁慈也同样增多。\BibleQuote{向凡呼求祢的人施大仁慈。}那么圣经在多处所说的\BibleQuote{他们呼求，我却不俯听}\BibleRef{箴 1:28}是什么意思呢？然而祢确实向凡呼求祢的人施仁慈；但有些人呼求，却不是呼求祂，正如经上所说：\Quote{他们没有呼求天主。}他们呼求，但不是呼求天主。你呼求你爱的；你呼求你召向自己的、你希望来到你这里的。因此，如果你呼求天主是为了这个理由，好使钱财来到你这里，好使遗产来到你这里，好使世俗地位来到你这里，你呼求那些你渴望来到你这里的事物：但你使天主成为你欲望的帮助者，而不是你需要的倾听者。如果天主赐给你所愿的，祂就是良善的吗？如果你愿的是恶，祂不赐予，岂不更显仁慈？然后，如果祂不赐予，天主对你而言就什么也不是；你说：我祈祷了这么多，祈祷了这么多次，却没有被俯听！为什么，你求了什么呢？也许是你仇敌的死亡。假如他同时也在祈祷你死呢？创造你的那一位，也创造了他：你是人，他也是人；但天主是审判者：祂俯听双方，却不允准任何一方的祈祷。你难过，因为你祈祷反对他时没有被俯听；要高兴，因为他的祈祷也没有被俯听来反对你。但你说：我求的不是这个；我求的不是我仇敌的死，而是我孩子的生命；我求了什么恶呢？你认为你没有求恶。如果\BibleQuote{他被接去，免得邪恶改变他的心意}\BibleRef{智 4:11}呢？但你说他是个罪人，所以我希望他活着，好能改正。你希望他活着，好能变得更好；如果天主知道，如果他活着会变得更坏呢？……因此，如果你呼求天主作为天主，要确信你必被俯听：你就有份于这句经文：\BibleQuote{向凡呼求祢的人施大仁慈。}……\par

8. 弟兄们，请思考并反省天主赐给罪人何等美善的事物；由此学习祂为祂自己的仆人存留了什么。祂每日赐给亵渎祂的罪人天空与大地、泉水、果实、健康、儿女、财富、丰裕：这一切美善唯有天主赐予。凡将这般恩惠赐予罪人的，你想想祂为祂忠信者存留了什么？岂能相信，祂将如此美好之物赐予恶人，却不为善人存留任何事物吗？不，祂确实存留着，不是大地，而是天堂为他们存留。也许我言及天堂，是过于寻常；祂自己，那创造诸天的，更为美好。诸天虽美，但创造诸天者更美。但我见诸天，却不见祂。因为你用眼目能见诸天，但尚未有心目能见创造诸天者；因此祂从天堂降临地上，为洁净人心，使人能看见那创造天地的。但你当以全然的忍耐等候救恩。祂知道用何种疗法来医治你：祂知道该用何种切割、何种灼烧。你因犯罪而自招疾病；祂不仅来护理，也来切割与灼烧。你岂不见人在医生手下所受的痛苦？医生向他们许诺一个不确定的希望。医生告诉你：\Quote{我会治好你，只要我切割。} 说话的是人，听话的也是人；说话者不确定，听者也不确定，因为那向人说话的人并未造人，也不完全知道人内在的运作；然而，一个不了解人内在之人的话，人却相信了，他顺服其手足，任由自己被捆绑，常常在未被捆绑的情况下被切割或灼烧；他或许获得几天的健康，但刚痊愈时，不知何时会死；或许在治疗中死去；或许无法治愈。但天主曾向谁许诺过什么而欺骗过他呢？\par

9. \BibleQuote{上主，求祢使我的祈祷稳立祢的耳中}\BibleRef{咏 85:6}。祈祷者极大的恳切！也就是说，不要让我的祈祷离开祢的耳中，要把它固定在其中。他是怎样劳苦，为了把他的祈祷固定在天主的耳中？让天主回答并对我们说：你愿意我把你的祈祷固定在我的耳中吗？将我的法律固定在你心中；\BibleQuote{并垂听我祈祷的声音。}\par

10. \BibleQuote{在困苦之日，我向祢呼号，因为祢应允了我}\BibleRef{咏 85:7}。稍前他曾说：我终日呼号，终日受困苦。所以没有一个基督徒可以说哪一天他没有受困苦。通过\Quote{终日}我们理解为整个时间。那么，即使我们境遇顺利，也有困苦吗？正是这样，有困苦。怎么会有困苦呢？因为\BibleQuote{我们几时住在肉身内，就是离开主。}\BibleRef{格后 5:6} 无论这里如何富足，我们尚未到达我们急于返回的家乡。喜欢旅居的人，不爱自己的家乡；若家乡可爱，旅居便是苦涩的；若旅居苦涩，则终日有困苦。何时没有困苦？当人在家乡有喜乐时。\Quote{祢的右手有永远的福乐。} \Quote{祢要以喜乐充满我，}他说，\Quote{在祢面前，使我得见主的喜乐。} 那里，劳苦和叹息都要过去：那里将不再是祈祷，而是赞美；那里有\OriginalTerm{阿肋路亚}{Alleluia}，那里有\OriginalTerm{阿们}{Amen}，与天使和谐的声音；那里有不失落的注视和不知疲倦的爱。因此，只要我们不在那里，你就看见我们不在那美善之中。但万物都丰裕吗？如果万物都丰裕，看看你是否确信万物不会朽坏？但我有了我以前没有的：更多的钱财来了，是我从前没有的。也许更多的恐惧也来了，是你从前没有的：也许你越是贫穷，就越安全。总之，即使你拥有财富，拥有这世界充裕的富足，你得到保证这一切不会朽坏；此外，天主对你说：你将在这些事物中永远存留，它们将永远与你同在，但我的面容你不得看见。谁也不要求问肉身的意见：你们当求问圣神的意见：让你的心回答你；让已经开始在你内的望德、信德、爱德回答。如果我们得到保证我们将永远处于世俗财富的丰裕中，如果天主对我们说：我的面容你不得看见，你们会因这些财富而欢欣吗？也许有人会选择欢欣，说：这些事物丰裕于我，我很好，别无所求。他还没有开始成为爱天主的人：他还没有开始像远离家乡的人一样叹息。绝不可，绝不可这样：让那些诱惑都退去：让那些虚假的谄媚退去：让那些他们每日对我们说的话消失：\Quote{你的天主在哪里？} 让我们在我们之上倾注我们的灵魂，让我们在眼泪中忏悔，让我们在忏悔中叹息，让我们在痛苦中呻吟。凡与我们同在而非我们的天主的，就不是甘甜的：如果祂不赐下祂自己，那赐下万有的，我们也不愿要祂所赐的一切。\par

11. \Quote{主，在诸神中没有像祢的} \BibleRef{咏 85:8}。他说了什么？\Quote{在诸神中}等等。让外邦人按他们的意愿造神；让他们带来银匠和金匠、磨光工、雕刻匠；让他们造神。那是怎样的神？有眼却不能看，以及圣咏随后提到的其他事情。但我们不崇拜这些，他说；我们不崇拜它们，这些都是象征。那么你们崇拜什么？更糟的东西：因为外邦人的神就是\OriginalTerm{魔鬼}{devils}。那么怎样？他们说，我们也不崇拜魔鬼。在你们的庙宇里肯定没有别的，除了魔鬼之外也没有别的激励你们的先知。但你们说什么？我们崇拜天使，我们以天使为神。你们完全不知道天使是什么。天使崇拜独一的天主，并且不赞成那些希望崇拜天使而不崇拜天主的人。因为我们发现高位天使禁止人朝拜他们，并命令他们朝拜真天主。\BibleRef{默 19:10} 但当他们说天使时，假设他们指的是人，因为经上说：\Quote{我曾说：你们是神，都是至高者的儿子。} 无论人怎么想，受造物不像造物主。除了天主，宇宙中任何其他事物都是天主造的。造物主与受造物之间有多大差别，谁能适当想象？因此这个人说：\Quote{主，没有像祢的；没有谁能像祢行事。} 但他没有说天主如何不像它们，因为无法言说。愿你们的\OriginalTerm{爱德}{Charity}留意：天主是\OriginalTerm{不可言喻}{ineffable}：我们更容易说祂不是什么，而不是祂是什么。你想到大地：这不是天主；你想到海洋：这不是天主；想到地上的一切，人和动物：这不是天主；想到海洋中的一切，空中飞翔的：这不是天主；想到天上发光的，星辰、太阳和月亮：这不是天主；天空本身：这不是天主；想到天使、异能、权势、总领天使、宝座、座位、宰制者：这不是天主。那么祂是什么？我只能告诉你祂不是什么。你问祂是什么？\Quote{眼所未见，耳所未闻，人心所未曾想到的。} \BibleRef{格前 2:9} …\par

12. \Quote{主，祢所造的万民都要来朝拜祢} \BibleRef{咏 85:9}。他预告了教会：\Quote{万民。} 如果有一个民族是天主没有造的，它就不会朝拜祂；但没有一个民族是天主没有造的，因为天主造了\Person{亚当}和\Person{厄娃}，万民的根源，从此万民生出。所以万民都是天主造的。这话是何时说的？当时在祂面前朝拜的只有希伯来一个民族中的少数圣人，那时这话就说了：现在请看所说的是什么：\Quote{祢所造的万民}等等。当这些话被说出时，它们未被看见，却被人相信；现在它们已被看见，为什么还要否认？\Quote{祢所造的万民都要来朝拜祢，主，并要光荣祢的圣名。}\par

13. \BibleQuote{因为祢是伟大的，且行了奇妙的事；惟有祢是伟大的天主。}\BibleRef{咏 85:10}。不要让任何人自称伟大。将来有些人要自称为伟大；针对这些人，经上说：\BibleQuote{惟有祢是伟大的天主。}因为当天主说惟有祂是伟大的天主时，有什么伟大的事归于天主呢？谁不知道祂是伟大的天主呢？但因将来有些人要自称为伟大，并使天主显得渺小，针对这些人，经上说：\BibleQuote{惟有祢是伟大的天主。}因为祢所说的必定应验，而不是那些自称为伟大的人所说的。天主藉祂的圣神说了什么？\BibleQuote{万民。}那个自称为伟大的人，无论他是谁，说了什么？\Quote{绝非如此：天主并没有在万民中受敬拜；万民都已消失，惟独\Place{阿非利加}存留。}这是你，自称为伟大的人所说的；另一位，惟有祂是伟大的天主，说了别的话。那位惟有祂是伟大的天主说了什么？\BibleQuote{万民。}我看见惟一的伟大天主所说的话：让那虚假伟大的人缄默；他只是在表面上伟大，因他鄙弃成为渺小。谁鄙弃成为渺小？说这话的人。主说：\BibleQuote{你们中谁愿成为伟大的，就应作你们的仆人。}\BibleRef{玛 20:26}如果那人曾愿作他兄弟们的仆人，他就不会将他们与他们的母亲分离；但当他希求伟大，却不愿成为渺小（这本是为他有益的）时，天主却相反骄傲的人，赐恩宠给谦卑的人\BibleRef{雅 4:6}，因为惟有祂是伟大的，祂成就祂所预知的一切，并斥责那些亵渎的人。因为这样的人亵渎基督，他们说教会已从全世界消失，只留在\Place{阿非利加}。如果你对他说：\Quote{你将失去你的别墅}，他或许几乎忍不住要动手打你；然而他却说，基督失去了祂用自己宝血所赎买的产业！我的弟兄们，请看这是多么大的冤屈。圣经说：\BibleQuote{君王的光荣在于民众多；王子受辱在于百姓稀少。}\BibleRef{箴 14:28}那么，你对基督行了这种冤屈，竟说祂的百姓减少到那微小数目。你生来就是为了这个吗？你自称基督徒就是为了这个吗？好使你嫉妒基督的荣耀，你说你在额头上带着祂的标记，心里却失去了祂？君王的光荣在于民众多：承认你的君王，给祂光荣，给祂一个广大的民族。你或许说：\Quote{我给祂什么广大的民族？}不要从你自己的心里给祂，你就会正确地给。你也许会说：\Quote{我从哪里给呢？}看，你从这给：\BibleQuote{万民都要来朝拜祢，上主。}说这个，承认这个，你就给了广大的民族：因为万民在一位内成为一体；这就是真正的一体。因为正如有一个教会和众教会，那些众教会也是一个教会，同样，那曾是众民族的现在是一个民族：先前是众民族，许多民族，如今是一个民族。为什么是一个民族？因为一个信德、一个望德、一个爱德、一个期待。最后，若是一个家乡，为何不是一个民族呢？我们的家乡是天上的，我们的家乡是耶路撒冷；凡不是这家乡的公民，就不属于那个民族；但凡是其公民的，就在那唯一的天主民族内。这个民族，从东到西，从北到海，遍及全世界的四方。天主说：从东和西，从北和海，归光荣于天主。祂预知了这一切，祂成就了这一切，惟有祂是伟大的。因此，那不愿成为渺小的人，应当停止对那惟一伟大者说这些话；因为不能有两个伟大者：天主和\Person{多纳图}。\par

14. \BibleQuote{上主，求祢引导我走祢的道路，我将在祢的真理中行走。}\BibleRef{咏 85:11} 祢的道路、祢的真理、祢的生命，就是基督。因此，身体属于祂，而身体是出于祂。\BibleQuote{我是道路、真理、生命。}\BibleRef{若 14:6} \BibleQuote{上主，求祢引导我走祢的道路。}哪条道路？\BibleQuote{我将在祢的真理中行走。}引导到道路是一回事，在道路中引导是另一回事。请看人处处贫穷，处处需要帮助。那些在道路之外的人不是基督徒，或尚未成为天主教徒：让他们被引导到道路；但当他们已被引领到道路，并在基督内成了天主教徒时，他们必须由祂在道路本身中引导，免得跌倒。现在他们确已在道路中行走。\BibleQuote{上主，求祢引导我走祢的道路：}我已确实在祢的道路中，求祢引导我到那里。\BibleQuote{我将在祢的真理中行走：}当祢引导时，我必不迷途；若祢容我离去，我必迷途。那么，要祈祷祂不让你离去，而是引导你直到终点。祂如何引导你？常常劝诫，常常向你伸出祂手。上主的手臂，向谁启示了呢？\BibleRef{依 53:1}因为赐下祂的基督，祂就伸出祂手；伸出祂手，祂就赐下祂的基督。祂引导到道路，藉着引导到祂的基督；祂在道路中引导，藉着在祂的基督内引导，而基督就是真理。\BibleQuote{求祢引导我，}因此，\BibleQuote{上主，走祢的道路，我将在祢的真理中行走：}确实是在祂内，祂曾说：\BibleQuote{我是道路、真理、生命。}\BibleRef{若 14:6}因为祢在道路和真理中引导，祢引导到哪里？岂不是引到生命？那么，在祂内，祢引向祂。\par

15. \Quote{愿我的心喜乐，以致敬畏你的名。}喜乐中存有敬畏。怎能有喜乐，若是敬畏？敬畏岂不常是痛苦的？将来必有毫无畏惧的喜乐，如今是带着敬畏的喜乐；因为尚未有完全的安稳，也非完全的喜乐。若无喜乐，我们便丧胆；若全然安稳，我们便妄喜。因此愿他既将喜乐洒在我们身上，又将敬畏击入我们心中，好藉喜乐的甘甜引导我们进入安稳的居所；藉赐予我们敬畏，使我们不致妄喜，且不从正路偏离。因此圣咏说：\Quote{当以敬畏事奉主，且战战兢兢向他欢欣。}宗徒\Person{保禄}也同样说：\Quote{当带着恐惧战兢，作成你们的救恩；因为是天主在你们内工作。}\BibleRef{斐 2:12} 诸位弟兄，无论何种顺境来临，反倒更当畏惧：你们以为顺利的事，实则是诱惑。遗产临到，财富临到，某种幸福的丰盈洋溢：这些都是诱惑；当心不要让它们败坏你们。此外，任何符合基督和基督真爱的顺境：或许你赢得了你的妻子，她原是\Person{多纳图}派的人；或许你的儿子原为外教人，如今成了信徒；或许你赢得了你的朋友，他曾想拉你到戏院，而你却拉他到教会；或许你某个仇敌曾对你狂怒，如今放下愤怒，变得温和，承认了天主，不再向你吠叫，却同你一同斥责邪恶：这些都是令人愉快的事。若非为此，我们为何喜乐？我们的喜乐若非这些，又是什么呢？但既然苦难、诱惑、纷争、分裂及其他邪恶也增多，没有这些世界便不能存在，直到邪恶逝去；愿那喜乐不使我们自满，但愿我们的心得喜乐，以致敬畏主的名，免得它在某方面喜乐，却在另一方面受打击。不要在旅途中期待安稳：若我们在此地企望它，它将成为身体的粘胶，而非人的安全。\Quote{愿我的心喜乐，以致敬畏你的名。}\par

16. \Quote{上主我的天主，我要全心称谢你，永远荣耀你的名。}\BibleRef{咏 85:12} \Quote{因你对我的慈爱浩大，你救了我的灵魂，免于极深的阴府。}\BibleRef{咏 85:13} 弟兄们，不要恼怒，如果我并非像有把握那样解释我所说的。因为我是人，关于神圣圣经所赐予我的，我才敢讲论：没有自己的东西。我未曾见过\OriginalTerm{阴府}{Hades}，你们也没有；或许我们另有道路，而不经过阴府。这些事是不确定的。但因那不可反驳的圣经说：\Quote{你救了我的灵魂，免于极深的阴府。}所以我们知道似乎有两个阴府，一个上层，一个下层：若没有上层，怎会有下层？只有与上层比较，下层才被称为下层。由此看来，弟兄们，似乎有某个天使的天上居所：那里有难以言喻的喜乐的生命，那里有不朽和不腐，那里一切事物按照天主的恩赐和恩宠常存。创造的那一部分是上方。若那是上方，而此尘世部分，有血肉，有可朽坏性，有出生和死亡，离去和继承，变化和不恒，有恐惧、欲望、恐怖、不确定的喜乐、脆弱的希望、易逝的存在；我想这一切无法与我刚提到的那个上天相比；若这部分无法与那部分相比，则一个在上，一个在下。死后我们去往何处？除非有一个比我们在肉体和可朽状态中所处的这个深渊更深的地方？因为宗徒说：\Quote{身体因罪而死。}\BibleRef{罗 8:10} 所以连此处也是死者；你们不必因它被称为\Emph{\OriginalTerm{阴府}{infernum}}而惊讶，若它充满死者。因为他不是说身体\Quote{将要死}，而是说：\Quote{身体是死的。}即便现在我们的身体确有生命；但与将来如天使身体的那个身体相比，人的身体被发现是死的，尽管仍有生命。但从这个\Emph{\OriginalTerm{阴府}{infernum}}，即从阴府的这一部分，另有一个更低的，死者去那里；天主愿从那里拯救我们的灵魂，甚至派遣他的独生子到那里。因为正是为了这两个阴府，弟兄们，天主子被派遣，从各方释放。到这一个阴府，他藉出生被派遣；到那一个，他藉死亡被派遣。因此那圣咏中的声音，不是凭人的猜测，而是宗徒解释时说的：\Quote{你必不将我的灵魂遗留在阴府。}所以这也是他的声音：\Quote{你从极深的阴府救了我的灵魂。}或是我们的声音，经由主耶稣基督自己：因他正是为此下到阴府，好使我们不留在阴府。\par

17. 我还要提另一种意见。因为或许就在地狱本身中，有一个较低的地方，那些犯了最重罪的恶人被扔在那里。至于地狱中是否有一些地方\Person{亚巴郎}所在之处，我们不能充分确定。因为主还没有来到地狱，以便解救所有先去的圣人的灵魂，而\Person{亚巴郎}在那里安息。\BibleRef{路 16:22} 有一个富翁在地狱受痛苦时，他举目望见\Person{亚巴郎}。他若举目，就不可能看见他，除非一个在上，一个在下。当他说：\BibleQuote{打发\Person{拉匝禄}来}。\Person{亚巴郎}回答说：\BibleQuote{我儿，你应记得你一生享了福，而\Person{拉匝禄}同样受了苦；但如今他在这里享安息，而你却受痛苦。此外，在你我之间，有一个大的深渊固定，以致我们这边不能到你那边去，也没有人能从那边过来。}\BibleRef{路 16:24} 因此，或许在这两个地狱之间，一个里面义人的灵魂得了安息，另一个里面恶人的灵魂受折磨，有一个在这里等候和祈祷，安置在基督的身体内，以基督的声音祈祷，说天主已从他的灵魂得救，脱离\OriginalTerm{最下层的地狱}{nethermost hell}，因为天主把他从那些可能成为他被拖入最下层的地狱之痛苦的罪中解救出来……。某人因一件麻烦的案子要被送进监狱；另一个人前来为他辩护；他感谢时说什么？你已救我的灵魂出狱。一个欠债者要被吊起来；他的债务被偿还了，他便被说是从被吊中得救了。他们并不在这些灾祸中；但因他们正朝这些灾祸前进，若没有援助，他们就会在其中，他们便正确地说，他们从那里被救出来，因为他们的解救者不容许他们被带去。所以，弟兄们，无论是这样或是那样，请视我为天主圣言的探究者，而非鲁莽的断言者。\par

18. \Quote{天主，\OriginalTerm{违法者}{transgressors of the law}起来攻击我}\BibleRef{咏 85:14}。他称谁为违法者？不是那些没有接受法律的外邦人，因为没有人会违反他没有接受的东西；宗徒清楚地说：\BibleQuote{因为没有法律，就没有作恶}\BibleRef{罗 4:15}。他称违法者为\OriginalTerm{作恶者}{prevaricators}。那么，弟兄们，我们理解的是谁？如果我们从主自己领受这话，那么违法者就是犹太人……他们不遵守法律，却控告基督好像他违法。我们知道主受了什么苦。你以为他的身体现在不受这样的苦吗？这怎么可能？\BibleQuote{如果人称家主为\Person{贝耳则步}，何况他的家人呢？门徒不能超过师傅，仆人也不能超过主人。}身体也受违法者的苦，他们起来攻击基督的身体。谁是违法者？难道犹太人胆敢起来攻击基督吗？不，因为不是他们给我们带来许多麻烦。因为他们还没有相信；他们还没有承认自己的救恩。起来攻击基督身体的是不良基督徒，基督的身体天天受他们的苦。所有分裂、所有异端、所有内里邪恶生活的人，将他们自己的品性嫁接到生活良善的人身上，把他们拉到自己一边，用不良的交际败坏善行；这些人\BibleQuote{违法起来攻击我}\BibleRef{格前 15:33}。让每个虔诚的灵魂发言，让每个基督徒的灵魂发言。那个不受这苦的灵魂，让它不要发言。但如果它是基督徒的灵魂，它知道自己受苦；如果它承认自己的苦难，让它在此承认自己的声音；但如果它没有受苦，就让它也没有声音；但为了它不至于没有受苦，让它行走窄路，\BibleRef{玛 7:14}并开始在基督内虔敬地生活：它必须受这种迫害。因为宗徒说：\BibleQuote{凡愿意在基督耶稣内虔敬生活的，都必受迫害}\BibleRef{弟后 3:12}。\par

\BibleQuote{\OriginalTerm{强大者的会堂}{synagogue of the powerful}寻索我的灵魂。}强大者的会堂就是骄傲者的集会。强大者的会堂起来攻击头，即我们的主耶稣基督，异口同声地喊道：\BibleQuote{钉死他，钉死他：}\BibleRef{若 19:6}关于这等人曾说过：\BibleQuote{人子，他们的牙齿是枪和箭，他们的舌头是快刀。}他们没有动手，只是喊叫；凭喊叫他们击打，凭喊叫他们钉死了他。那些喊叫者的意愿成就了，当主被钉十字架时：\BibleQuote{他们就没有把你放在他们眼前。}他们怎么不把他放在眼前？他们不知道他是天主。他们本该把他当作人而饶恕：他们所见到的，他们应当依此行事。假定他不是天主，他是人：难道他就该被杀吗？饶恕他是人，而承认他是天主。\par

19. \Quote{而你，主天主，是慈悲和宽仁的，\OriginalTerm{含忍}{longsuffering}，极富怜悯，并是真实的}\BibleRef{咏 85:15}。为何含忍、极富怜悯，并是慈悲？因为当他挂在十字架上时说：\BibleQuote{父啊，宽赦他们，因为他们不知道他们做的是什么}\BibleRef{路 23:34}。他向谁祈祷？他为谁祈祷？谁在祈祷？他在哪里祈祷？圣子向圣父祈祷，为不虔敬者被钉，在极大的侮辱中，不是言语的侮辱，而是被施加死刑，挂在十字架上；仿佛为此他伸开双手，好能这样为他们祈祷，使他的\BibleQuote{祈祷像馨香一样呈到父面前，举起双手像晚祭一样。}\par

20. 所以，如果祢是\Quote{真实的}，就\Quote{求祢垂顾我，怜悯我：将能力赐给祢的仆人}。因为祢是\Quote{真实的}，所以\Quote{求祢将能力赐给祢的仆人}\BibleRef{咏 85:16}。让忍耐的时期过去，审判的时期来临。如何\Quote{赐予能力}？\Quote{父不审判任何人，但将一切审判交给了子}\BibleRef{若 5:22}。祂复活后，将亲自来到大地进行审判：祂曾显得可鄙，将显现可畏。祂将显示祂的能力，祂曾显示祂的忍耐；在十字架上是忍耐，在审判中将是能力。因为祂将作为人审判，但却是带着荣耀：\Quote{就如你们看见祂升天}，天使说，\Quote{祂也要这样降来}\BibleRef{宗 1:11}。祂的形像本身将前来审判，因此不虔敬的人也必看见祂：因为他们不会看见天主的形像。因为\Quote{心里洁净的人是有福的，因为他们要看见天主}\BibleRef{玛 5:8}……在看见父的异象中，也有看见子的异象；在看见子的异象中，也有看见父的异象。因此祂加上一句推论，并说：\Quote{你不信我在父内，父在我内吗？}\BibleRef{若 14:10}也就是说，既在看见我时父被看见，也在看见父时子也被看见。父和子的异象是不可分的：在本性本体不可分离之处，异象也不可分离。为使你们知道，应当预备好心灵去看见父、子和圣神的天主性——虽然尚未看见，我们却相信，并且藉着相信洁净心灵，好能有朝一日眼见——主自己在另一处说：\Quote{那接受我的命令并遵守的，就是爱我的人；爱我的必蒙我父爱他，我也要爱他，并将我自己显示给他}\BibleRef{若 14:21}。难道当时与祂交谈的人没有看见祂吗？他们既看见祂，又没有看见祂？他们看见了一些东西，相信了一些东西：他们看见了人，相信了天主。但在审判中，恶人将与义人一同看见同一位主耶稣基督作为人；审判之后，他们将看见天主，远离恶人。\par

21. \Quote{求祢拯救祢婢女的儿子。}主就是那位婢女的儿子。是哪一位婢女？就是当祂被预报将由她诞生时，她回答说：\Quote{看，上主的婢女，愿照祢的话成就于我吧}\BibleRef{路 1:38}。祂拯救了祂婢女的儿子，也就是祂自己的儿子：祂自己的儿子，是在天主的形象内\BibleRef{斐 1:6}；婢女的儿子，是在仆人的形象内。因此，天主婢女的儿子以仆人的形象降生；祂说：\Quote{求祢拯救祢婢女的儿子。}祂从死亡中被拯救，如你们所知，祂的肉躯虽曾死过，却复活了……每一个在基督身体内的基督徒都可以说：\Quote{求祢拯救祢婢女的儿子。}或许他不能说：\Quote{将能力赐给祢的仆人}，因为获得能力的是那位子。然而，为何他也不说这话？难道不是对仆人说：\Quote{你们要坐在十二个宝座上，审判以色列十二支派}吗？\BibleRef{玛 19:28}而仆人们说：\Quote{你们不知道我们要审判天使吗？}\BibleRef{格前 6:3}所以每一位圣人也获得能力，每一位圣人也都是祂婢女的儿子。如果他是由异教母亲所生，却成了基督徒，那又怎样？异教徒的儿子怎能是祂婢女的儿子？他按肉身确实是异教母亲的儿子，但按圣神却是教会的儿子。\par

22. \Quote{求祢向我显示一个善待我的征兆}\BibleRef{咏 85:17}。什么征兆呢？不就是复活的征兆吗？主说：\Quote{这邪恶淫乱的世代寻求征兆，除了约纳先知的征兆外，再没有征兆给他们}\BibleRef{玛 12:39}。因此，在我们的元首身上已经显示了一个善待我们的征兆；我们每个人也可以说：\Quote{求祢向我显示一个善待我的征兆}，因为在末次号角响起、主来临时，\Quote{死人要复活成为不朽坏的，我们也要被改变}\BibleRef{格前 15:52}。这将是善待我们的征兆。\Quote{让恨我的人看见，并感到羞愧。}在审判中，他们将因自己的毁灭而羞愧，他们现在不愿因自己的痊愈而羞愧。所以现在让他们羞愧：让他们指责自己的道路，让他们持守善道：因为没有人能不经羞愧而生活，除非他先羞愧并重新做人。现在天主赐给他们健康羞愧的契机，如果他们不藐视告解的良药；但如果他们现在不愿羞愧，那么到时他们将羞愧，那时\Quote{他们的邪恶要在他们面前作证}。他们将如何羞愧？当他们说：\Quote{这些人是我们曾一度讥讽和侮辱的对象。我们愚妄人曾以为他们的生命是疯狂：现在他们怎么被列在天主的儿女中！骄傲给了我们什么益处？}那时他们将说这些话：让他们现在就说，他们是为了自己的健康而说的。因为让每个人谦卑地转向天主，现在就说：\Quote{我的骄傲给了我什么益处？}并从宗徒那里听到：\Quote{那时你们以那些事为羞耻，它们有什么光荣呢？}\BibleRef{罗 6:21}你们看，即使现在也有一种有益的羞愧，当仍有补赎的机会；但到那时，羞愧将是迟来的、无用的、徒然的……\par

23. \Quote{因为祢，上主，帮助了我，安慰了我。} \Quote{帮助了我，}是在斗争中；\Quote{安慰了我，}是在忧苦中。因为没有人寻求安慰，除非他处于痛苦中。难道你们不愿受安慰吗？就说你们是快乐的，你们听到：\Quote{我的子民}（现在你们回答，我听到一阵喃喃声，像是记得圣经的人的声音。愿天主，就是那将这些写在你们心中的，在你们的作为中证实它。你们看，弟兄们，那些对你们说 \InnerQuote{你们是快乐的} 的人，是在诱惑你们），\Quote{哦，我的子民，那些称你们为有福的人迷惑了你们，扰乱了你们脚下的路。} 同样从宗徒 \Person{雅各伯} 的书信：\BibleQuote{你们要受苦，要哀恸：让你们的欢笑变成哀恸。} \BibleRef{雅 4:9} 你们看到你们所听到的诵读：什么时候会有人在安全之地对我们说这样的话？这确实是冒犯和诱惑以及一切邪恶之地，为使我们在这里叹息，配得在那里喜乐；在这里受苦，在那里受安慰，并说：\Quote{因为祢解救了你的眼睛免于流泪，我的脚免于跌倒：我将在活人之地中悦乐上主。} 这是死人之地。死人之地过去，活人之地来临。在死人之地有劳苦、忧愁、恐惧、磨难、诱惑、呻吟、叹息：这里有虚假的快乐者，真正的痛苦者，因为快乐是虚假的，痛苦是真实的。但那承认自己处于真实痛苦中的人，也将处于真实的快乐中：然而现在，因为你们是痛苦的，请听主说：\BibleQuote{哀恸的人是有福的。} \BibleRef{玛 5:4} 哦，哀恸的人是有福的！没有什么比哀恸更接近痛苦：没有什么比幸福更远离和相反于痛苦：祢说到那些哀恸的人，祢称他们为有福的！要理解，祂说，我说的话：我称那些哀恸的人为有福的。为什么有福？在希望中。为什么哀恸？在行动中。因为他们在这个死亡中，在这些磨难中，在他们的漂泊中哀恸：并且因为他们承认自己处于这个痛苦中，并哀恸，他们是有福的。他们为什么哀恸？真福 \Person{西彼廉} 在他的苦难中陷入忧愁：现在他与他的冠冕一同受安慰；现在即使受安慰，他仍是忧愁的。因为我们的主耶稣基督仍在为我们转求：所有与他同在的殉道者都为我们转求。他们的转求不会过去，除非我们的哀恸过去：但当我们的哀恸过去后，我们所有人都将异口同声，在一个子民中，在一个国度中，得到安慰，成千上万的人与弹奏竖琴的天使、与生活在同一城中的天上掌权者的歌咏团联合。谁在那里哀恸？谁在那里叹息？谁在那里劳苦？谁在那里需要？谁在那里死亡？谁在那里施予仁慈？谁在那里给饥饿者擘饼，那里所有人都被正义的食粮所满足？没有人对你说：接待陌生人；那里没有人会成为你的陌生人：所有人都住在自己的家乡。没有人对你说：使你们争论的朋友和解；他们在永久的平安中享见天主的圣容。没有人对你说：探望病人；健康和永生永远常存。没有人对你说：埋葬死者；所有人都将得到永生。仁慈的工作停止了，因为痛苦不再存在。我们在那里要做什么？也许我们会睡觉吗？如果我们现在与自己对阵，虽然我们携带一栋睡眠之屋，即我们的肉身，并且保持警醒，借着这些灯光，而这个庄严的庆节给我们警醒的心思；那么那一天会赐予我们何等醒觉！因此我们将醒着，不会睡觉。我们要做什么？将不会有仁慈的工作，因为没有痛苦。也许会有现在这里这些必要的工作，如播种、耕田、烹饪、磨面、织布？没有这些，因为没有匮乏。因此，将没有仁慈的工作，因为痛苦已过去：在没有匮乏和痛苦的地方，将既没有必需的工作，也没有仁慈的工作。那里将有什么？我们将有什么事务？什么行动？难道没有行动，因为有安息？我们将坐在那里，呆滞不动，什么都不做吗？如果我们的爱变得冷淡，我们的行动也将变得冷淡。那么，那安息在天主圣容面前的爱，我们现今渴望、为之叹息的，当我们到达祂那里时，将如何点燃我们？祂，我们尚未看见时就如此叹息的，当我们到达祂那里时，将如何光照我们？祂将如何改变我们？祂将使我们成为什么？那么我们要做什么，弟兄们？让 \BibleBook{}{圣咏} 告诉我们：\Quote{居住在你殿中的人是有福的。} 为什么？\Quote{他们要赞美你，直到永远。} 这将是我们的事业：赞美天主。你爱并赞美。如果你停止爱，你就会停止赞美。但你不停止爱，因为你所看见的那一位是如此的，不会因任何疲倦而冒犯你：祂既满足你，又不满足你。我说的是奇妙的事。如果我说祂满足你，我害怕，因你好像满足了，想要离开，如同从午餐或晚餐离开一样。那么我说什么？难道祂不满足你吗？我又害怕，如果我说祂不满足你，你似乎处于匮乏中：并且几乎是空的，你里面有些空虚需要被充满。那么我要说什么，除了那可以说，但难以想到的？祂既满足你，又不满足你：因为我在圣经中两者都找到。因为当祂说：\BibleQuote{饥饿的人是有福的，因为他们将得饱足；} \BibleRef{玛 5:6} 又论及智慧说：\BibleQuote{那些吃祢的仍会饥饿，那些喝祢的仍会口渴。} \BibleRef{德 24:21} 不，但祂没有说\Quote{再}，而是说\Quote{仍然}：因为\Quote{再会口渴}就像是一次被充满后离开，消化了，再回来喝。所以是：\Quote{那些吃祢的仍会饥饿：} 因此他们吃的时候仍饥饿：那些喝祢的，即使在喝的时候仍然口渴。什么是喝的时候口渴？永不疲倦。那么，如果那里有那不可言喻的永恒甘甜，祂现在向我们要求什么，弟兄们，除了真诚的信德、坚固的望德、纯洁的爱德？人可以在主所赐的道路上行走，忍受困难，并接受安慰。\par

\SourceRecord{Translated by J.E. Tweed.；Nicene and Post-Nicene Fathers, First Series；Vol. 8.；Edited by Philip Schaff.；Buffalo, NY: Christian Literature Publishing Co.,；1888.；<http://www.newadvent.org/fathers/1801086.htm>.}

% structure: p0001 source=h1 target=section reason=page-title

\section{圣咏第八十七篇释义}

1. 刚才所咏唱的圣咏，若按其字数来看，是简短的，但其思想却深具意义.. ..那圣咏歌唱赞美的主题是一座城，其市民就是我们——就我们是基督徒而言；我们远离该城，因我们仍有死；我们正趋向这城；该城的君王，经由那些荆棘与蒺藜缠绕而不可发现的道路，亲自为我们开辟了一条通往该城的途径。我们如此在基督内行走，作为朝圣者直到抵达，因渴慕那城中不可思议的安息而叹息，这安息正如所应许的，是\BibleQuote{人的眼睛未曾看见如此，耳朵未曾听见，人心也未曾想到的}；让我们咏唱一颗渴望之心的诗歌：因为那真正渴望的人，即使舌头沉默，也在灵魂内如此歌唱；而那并不渴望的人，无论在人间如何喧嚷，在天主面前却是无声的。请看，那些写下并传给我们这些话的人，是多么热爱那城啊！他们咏唱时怀有何等深情！那深情是由对那城的爱在他们内产生的；那爱是由圣神所激发的；他说：\BibleQuote{天主的爱，藉着所赐予我们的圣神，已倾注在我们心中}。让我们怀着圣神的热情，聆听关于那城的话。\par

2. \BibleQuote{她的基础在圣山上。} \BibleRef{咏 86:1} 圣咏此前并未提及那城；它就这样开始，说：\BibleQuote{她的基础在圣山上。} 谁的基础？毫无疑问，基础——尤其是山上的基础——属于某座城。因此，那位市民充满了圣神，怀着对那城的许多爱与渴望，仿佛经过长时间的内省后，脱口而出：\BibleQuote{她的基础在圣山上。} 似乎他已谈论过那城。他怎能对一个在心中从未缄默的话题闭口不言呢？因为\BibleQuote{她的基础}这一说法，若先前从未提及，又怎能写下来？但正如我所说，他在心中对那城长久静默地沉思，向天主呼号后，向人耳这样脱口而出：\BibleQuote{她的基础在圣山上。} 而且，假设听到的人在询问他所谈的是哪一座城，他就加上：\BibleQuote{主喜爱熙雍的门户。} 看，那就是一座城，其基础在圣山上，一座名为熙雍的城，主喜爱她的门户，如他接着所说：\BibleQuote{超过雅各伯所有的帐幕。} 但这\BibleQuote{她的基础在圣山上}是什么意思？这座城所建筑的圣山是什么？另一位市民——\Person{保禄}——更明确地告诉我们：先知也是这城的市民，宗徒也是；他们说话鼓励其他市民。但这些人——我指的是先知和宗徒——如何是市民呢？或许在此意义上：他们自己就是那些山，这城的基础在其上，而主喜爱这城的门户。那么，让另一位市民清楚说明这一点，以免我像是在猜测。他对外邦人说话，告诉他们如何回归，并被建造在神圣结构之中，他说：\BibleQuote{被建造在宗徒和先知的基础上。} 因为宗徒和先知——那城的基础停靠在他们之上——不能凭自己的力量站立，他加上：\BibleQuote{耶稣基督自己为角石。} \BibleRef{弗 2:20} 因此，外邦人可能认为自己与熙雍无关，因为熙雍是这世界的一座城，作为预表与其影相似，指向他随后所说的那座熙雍——天上的耶路撒冷，宗徒说：\BibleQuote{它是我们众人的母亲。} \BibleRef{迦 4:26} 他们可能被认为与熙雍无关，因为他们不属于犹太百姓，但他这样对他们说：\BibleQuote{所以你们不再是外人和旅客，而是与圣徒同国，是天主家里的人，并且被建筑在宗徒和先知的基础上。} \BibleRef{弗 2:19} 你们看到如此宏伟之城的结构：但整个建筑倚靠在哪里呢？它安歇在哪里，才永不倒塌？他说：\BibleQuote{耶稣基督自己为角石。}\par

……但为使你知道基督既是最早也是最高的根基，宗徒说：\BibleQuote{奠立别的根基，没有人能奠立，除了那已奠立的根基，即耶稣基督。}\BibleRef{格前 3:11}那么，先知和宗徒如何也是根基，而基督却是如此，无人能高于祂呢？我想，除了正如祂明明被称为\OriginalTerm{圣中之圣}{Sanctus sanctorum}，以象征的方式也称为\OriginalTerm{根基中的根基}{Fundamentum fundamentorum}，还能如何呢？因此，你若思想奥秘，基督是圣中之圣；你若思想受牧养的羊群，祂是牧者中的牧者；你若思想建筑物，祂是柱石中的柱石。在物质建筑中，同一块石头不能既在下面又在上面；若在底部，就不能在顶部；\Emph{反之亦然}：因为几乎所有物体都受空间限制，不能同时遍在或永在；但天主性既在每一处，便可以从每一处取象征来表现它；它并非外在属性中的任何一物，却能在象征中成为一切。基督是门，如同我们所见木匠做的门那样吗？当然不是；然而祂说：\BibleQuote{我就是门。}或者祂是牧者，如同看守羊群的人那样吗？尽管祂说：\BibleQuote{我是善牧。}这两个名字出现在同一段经文中：在福音中，祂说牧者由门进入；原话是：\BibleQuote{我是善牧；}同段中又说：\BibleQuote{我就是门。}那么，由门进入的牧人是谁呢？\BibleQuote{我是善牧；}而祢，善牧，所进入的门是什么？祢如何是万有呢？其意义是：万有都藉着我。要解释：当\Person{保禄}由门进入时，难道基督不进入吗？为什么呢？不是因为\Person{保禄}是基督，而是因为基督在\Person{保禄}内，\Person{保禄}藉基督行事。宗徒说：\BibleQuote{你们既然寻求基督在我内说话的凭证？}\BibleRef{格后 13:3}当祂的圣徒和忠信的门徒由门进入时，难道基督不也由门进入吗？我们如何证明这点呢？因为\Person{扫禄}（尚未被称为\Person{保禄}）正在迫害那些圣徒时，祂从天上呼叫他：\BibleQuote{扫禄，扫禄，你为什么迫害我？}\BibleRef{宗 9:4}所以祂自己就是根基和\OriginalTerm{角石}{angularis lapis}；从底部升起——如果确实是从底部升起的话，因为这根基的基础正是建筑物最高的高升。正如物质建筑物的支撑安放在地上，精神建筑物的支撑却安放在高处。如果我们在地上建造自己，我们就会把根基放在最低处；但既然我们的建筑物是天上的，我们的根基已先我们升到天上。因此，我们的救主、角石、宗徒和强有力的先知，这些承载城市建筑物的山丘，构成了一种活的建筑。这建筑如今从你们心中呼喊；为使你们被建造进入它的结构，天主的手——如同工匠的手——甚至藉我的舌头工作。诺厄的方舟用\OriginalTerm{方木}{quadratis lignis}制成，并非没有意义，这些方木预表了教会的形状。因为成为方形是什么意思呢？请听方石的象征：基督徒应有这样的品质：在一切考验中永不跌倒；虽受推挤，仿佛被翻转，却仍不倒下；同样，一块方石无论怎样转动，都直立不倒。……在尘世的城市中，建筑物是一回事，居住其中的市民是另一回事；那座城市是由其居民建筑起来的，他们自己就是构成城市的石块，因为那些石头是活的：\BibleQuote{你们也就像活石，被建造成一座属神的殿宇，}\BibleRef{伯前 2:5}这话是对我们说的。让我们继续默想那座城市。\par

4. \BibleQuote{主喜爱熙雍的门，胜过雅各伯所有的帐幕} \BibleRef{咏 86:2}。我方才说了这些，好叫你们不以为门是一样东西，基础是另一样。为什么宗徒和先知是基础？因为他们的权威支撑着我们的软弱。为什么他们是门？因为通过他们，我们进入天主的国：因为他们向我们宣讲它；当我们藉着他们进入时，我们也藉着基督进入，祂自己就是门。经上也提到耶路撒冷有十二座门，\BibleRef{默 21:12}，而唯一的门是基督，十二座门也是基督，因为基督住在十二座门内，因此宗徒的数目是十二。在这十二的数字中有一个深奥的奥秘：我们的救主说：\BibleQuote{你们要坐在十二宝座上，审判以色列十二支派} \BibleRef{玛 19:28}。如果那里有十二个宝座，就没有保禄宗徒（第十三位宗徒）审判座的位置，尽管他说他不仅要审判人，还要审判天使——但除堕落的天使外还有谁呢？\BibleQuote{你们不知道我们要审判天使吗？} \BibleRef{格前 6:3}，他写道。世人会回答：你为什么夸耀你将作审判官？你的宝座在哪里？我们的主讲论了十二个宝座给十二位宗徒：一位，犹达斯，堕落了，他的位置由玛弟亚取代，十二个宝座的数目就补满了：\BibleRef{宗 1:15} 那么，首先为你的审判座找到位置，然后才恐吓说你将审判。因此，让我们思考十二宝座的意义。这句话典型地代表着一种普遍性，正如教会注定要传遍全世界；因此这建筑物被称为一同建造成基督：并且因为审判官从各方而来，所以提及十二宝座，正如十二扇门，从各方进入那城。因此，不仅有那十二位和保禄宗徒对十二宝座拥有权利，而且由于普遍的意义，所有将要坐席审判的人也都拥有；正如所有进入城的人，都从十二扇门中的某一扇进入。地球有四个方向：东、西、北、南；在圣经中经常提到它们。我们的主在福音中宣布，祂要从四方召唤祂的羊\Quote{从四方}；\BibleRef{谷 13:27} 因此教会从四方被召来。如何被召？在圣三的每一方面都被召：没有别的召法，只有因父、及子、及圣神之名的圣洗：如此，四乘以三，得十二。因此，用你们全心敲这些门：让基督在你们内呼喊：\Quote{给我打开正义的门}。因为祂作为元首走在我们前面；祂在祂的身体内跟随自己……。\par

5. \BibleQuote{人们论及你，天主的城，说了极其荣耀的事} \BibleRef{咏 86:3}。他仿佛默观了那地上的耶路撒冷城：因为请想想，他指的是哪座城，论及它说了某些极其荣耀的事。如今地上的城已被毁灭：遭受敌人的愤怒后，它倒塌在地；它不再是原先的样子：它曾展示预像，而影子已经消逝。那么，\Quote{人们论及你，天主的城，说了极其荣耀的事} 是从何说起呢？请听从何说起：\BibleQuote{我要提拔辣哈布和巴比伦，归于认识我的人} \BibleRef{咏 86:4}。在那城里，先知以天主的位格说：\Quote{我要提拔辣哈布和巴比伦}。辣哈布不属于犹太民族；巴比伦不属于犹太民族；如下一节经文所明示：\BibleQuote{连培肋舍特、提洛，以及雇士人，也在那里}。因此，的确，\Quote{人们论及你，天主的城，说了极其荣耀的事}：因为不仅从亚巴郎血肉而生的犹太民族被包括在内，而且所有的民族也都包括，其中提到了一些，以使人理解全部。他说：\Quote{我要提拔}不单是\InnerQuote{辣哈布}，还有\InnerQuote{巴比伦}，但和谁在一起呢？\InnerQuote{归于认识我的人}。谁是那个妓女？是耶里哥的那个妓女，她接待了探子，并引导他们从另一条路出城：她预先相信了应许，敬畏天主，被吩咐在窗口挂一条朱红线绳，即在她额头上佩戴基督血的标记。她在那里得救，并如此代表了外邦人的教会：因此我们的主对骄傲的法利塞人说：\BibleQuote{我实在告诉你们：税吏和娼妓要在你们之先进入天国} \BibleRef{玛 21:31}。他们走在前头，因为他们使用暴力：他们以信德冲出一条路，为信德开了一条路，无人能抵抗，因为强暴的人夺取它。因为经上记载：\BibleQuote{天国是以猛力夺取的，以猛力夺取的人就攫取了它} \BibleRef{玛 11:12}。那就是强盗的行径，在十字架上比在埋伏处更勇敢。\Quote{我要提拔辣哈布和巴比伦}。巴比伦意指这世界的城：正如有一座圣城耶路撒冷，一座不圣的城巴比伦；所有不圣的都属巴比伦，正如所有圣的都属耶路撒冷。但他从巴比伦滑向耶路撒冷。如何做到？惟有藉着那成就不虔敬者为义的那一位：耶路撒冷是圣徒的城，巴比伦是恶人的城；但祂来了，使不虔敬者成义：因为经上说：\Quote{我要提拔}，不单是\InnerQuote{辣哈布}，还有\InnerQuote{巴比伦}，但和谁在一起呢？\Quote{归于认识我的人}……。\par

6. 请听一个深奥的奥秘。\Place{辣哈布}在那里，藉着祂，\Place{巴比伦}也藉着祂，现在不再是\Place{巴比伦}，而是开始成为\Place{耶路撒冷}。女儿与母亲分离，并将成为那位王后的成员，对她说：\Quote{忘记你的民族和你的父家，那么国王将喜爱你的美貌。}因为\Place{巴比伦}如何能向往\Place{耶路撒冷}？\Place{辣哈布}如何能到达那些根基？\Place{培肋舍特人}、\Place{提洛人}、\Place{厄提约丕雅人}又如何能？请听这句：\BibleQuote{\Place{熙雍}，我的母亲，一个人要说。}因此有一个人说这话：藉着他，我提到的所有人都得以接近。这个人是谁？只要我们听、只要我们明白，它就会告诉我们。接着，好像有人提出了一个问题：\Place{辣哈布}、\Place{巴比伦}、\Place{培肋舍特人}、\Place{提洛人}和\OriginalTerm{摩黎人}{Morians}是藉着谁的帮助获得了进入？看，他们是藉着谁来的：\BibleQuote{\Place{熙雍}，我的母亲，一个人要说；有一个人在她内诞生，而至高者亲自建立了她。}\BibleRef{咏 86:5}我的弟兄们，有什么比这更清楚的呢？的确，因为\BibleQuote{关于你，天主的城，人们谈论了极其卓越的事。}看，\BibleQuote{\Place{熙雍}啊，母亲，一个人要说：}什么人？就是在她内诞生的那一位。因此，就是在她内诞生的那个人，而祂亲自建立了她。然而，祂如何能在祂亲自建立的城中诞生呢？城早已建立，为的是祂可以在此诞生。你们若能明白，就这样理解吧。\BibleQuote{母亲\Place{熙雍}，他要说；但说\InnerQuote{母亲\Place{熙雍}}的是一个人；确实，有一个人在她内诞生；然而建立她的（不是一个人，而是）至高者。}正如祂创造了祂将由此出生的母亲，同样，祂建立了祂将在此诞生的城。我们的希望是何等大啊，弟兄们！为我们，至高者——祂建立了这座城——称那座城为母亲；祂在她内诞生，而至高者建立了她。\par

7. 好像有人说：\Quote{你们怎么知道这个？}我们所有人都吟唱过这些圣咏；而基督——为我们而降生的人，在我们之前就存在的天主——在我们众人内吟唱。但说祂在我们之前，这对那在天、地和时间之前就已存在的祂来说，难道算得上什么吗？祂，为我们降生成人，在那城中（当祂是至高者时）也建立了她。但我们如何确定这事呢？\BibleQuote{上主在记录万民时，将述说这事。}\BibleRef{咏 86:6}正如下一节经文所说。\BibleQuote{上主在记录万民和他们的首领时，将宣告。}什么首领？就是在城内诞生的人；那些在她城墙内诞生、在她内成为首领的人：因为他们能在她内成为首领之前，天主拣选了世上卑贱的来羞辱那强健的。渔夫、税吏，曾是首领吗？他们确实是首领，但因为他们是在她内成为这样的。他们是怎样的首领呢？来自\Place{巴比伦}的首领，这世界的信主君王，来到\Place{罗马}城，如同来到\Place{巴比伦}的头部：他们不去皇帝的神庙，却去了渔夫的坟墓。他们确实从何处成为首领呢？\BibleQuote{天主拣选了世上懦弱的，来羞辱那强健的；祂也拣选了世上愚拙的，以及那无有的，为的要使那有的归于无有。}\BibleRef{格前 1:26}这是祂所做的，祂\Quote{从尘埃里提拔弱小者，从粪土中高举穷乏人。}为了什么目的？\Quote{为使祂与他的人民的首领同坐。}这是伟大的作为，是喜乐和欢跃的深邃源泉。演说家后来进入那城，但若渔夫未曾先到，他们绝不可能如此。这些事确实荣耀，但它们能在哪里发生呢？除非在城里，关于天主之城，人们谈论了极其卓越的事。\par

8. 因此，在汇聚和融合了一切喜乐欢跃的源头之后，他是如何总结的呢？\BibleQuote{仿佛一切将要喜乐之人的居所，就在祢内}见：\BibleRef{咏 86:7}。仿佛一切得成喜乐、一切欢欣的人，都要在那城里居住。在我们此世的旅途中，我们遭受创伤；我们最终的家乡，将是唯独喜乐的家乡。劳苦与叹息将要消逝；祈祷要过去，赞美的诗歌将接续而来。那里将是幸福之人的居所；不再有渴望者的叹息，只有享有者的欢欣。因为祂必将亲临，我们为之叹息的那一位；\BibleQuote{我们将相似祂，因为我们将看见祂实在怎样}见：\BibleRef{若一 3:2}；在那里，我们全部的工作就是赞美和享有天主的临在；此外我们还能求什么，当祂——万有由祂而造——独自满足我们时？我们将居住于祂内，祂也居住于我们内；我们将隶属于祂，\BibleQuote{好使天主成为万物之中的万有}见：\BibleRef{格前 15:28} \BibleQuote{有福的是居住在祢殿中的人}。如何有福？是拥有金银、众多奴仆、子孙满堂吗？\BibleQuote{有福的是居住在祢殿中的人：他们将永永远远赞美祢。}这福分在于唯一的工作，而这工作就是安息！我的弟兄们，当我们到达这番境地时，就让这成为我们唯一的渴望吧。让我们预备自己，好能在天主内喜乐，赞美祂。引导我们到达那里的善工，在那里将不再需要。我昨天曾尽我所能描述了我们在那里的境况：爱德的工作将不复存在，因为那里不再有苦难；你不会发现一个匮乏的人，一个赤身露体的人，没有人会遭遇干渴之苦，没有外乡人，没有需要探望的病人，没有需要埋葬的死者，没有需要调停的争论者。那么你将会做什么？我们将种植新葡萄树、耕地、经商、航行，以供养身体的需要吗？那里是深沉的宁静；一切因需要而产生的劳苦工作都将止息；需要既死，其工作也必消亡。那么我们的状态将是什么？凡人的口舌曾这样告诉我们，尽其所能地说：\Quote{仿佛一切将要完善之人的居所，就在祢内。}他为什么说\Quote{仿佛}？因为那里的喜乐，是我们在此所不识的。我在此看到许多欢乐，许多人在此世欢乐，有人以此事为乐，有人以彼事为乐；但没有什么能与那喜乐相比，但它将\Quote{仿佛}是喜乐。因为若我说喜乐，人们立刻会想到人们在酒、宴乐、贪婪和世俗的荣华中那种喜乐。人被这些事兴奋，带着一种疯狂的喜乐；但\BibleQuote{没有喜乐——主说——归于恶人。}见：\BibleRef{依 48:22}。\BibleQuote{有一种喜乐，是人耳未曾听见，人眼未曾看见，人心也未曾想到的。}见：\BibleRef{格前 2:9} \Quote{仿佛一切将要喜乐之人的居所，就在祢内。}让我们预备寻求另一种欢愉：因为在此找到的只是某种阴影，而非实体；免得我们期望在那里享受如同在此所喜好的事物；否则我们的克己就会变成贪婪。有些人，当被邀请参加丰盛的宴席，还有许多昂贵菜肴尚未上桌时，他们不先开斋；如果你问原因，他们会告诉你在斋戒：这确实是一项伟大的工作，基督徒的工作。但不要急于称赞他们：要察验他们的动机；他们是在为肚腹，而非为宗教考虑。为了不让普通的菜肴败坏他们的胃口，他们一直等到更精致的食物摆上。那么这种斋戒是为了口腹之欲。斋戒无疑很重要：它对抗肚腹和口舌；但有时它却为它们效劳。因此，我的弟兄们，如果你们想象我们将从那天国的号角催促我们前往的境域里找到任何相似的享乐，并为此而戒绝现世的享受，以便在那里更丰盛地领受，你们就在模仿我所描述的那些人，他们只为了更大的宴乐而斋戒，只为了更大的放纵而禁欲。你们不要这样：当为自己预备一种不可言喻的喜乐：从所有世俗和现世的情欲中洁净你们的心。我们将看见某样事物，看见它就会使我们有福；单单那样就足以满足我们。那么，我们就不吃了吗？不，我们吃：但那是我们的食物，它将不断更新，永不止息。\Quote{在祢内，是一切将要——仿佛是——喜乐之人的居所。}祂已经告诉我们，我们将如何成为喜乐。\BibleQuote{有福的是居住在祢殿中的人：他们将永永远远赞美祢。}让我们尽我们所能赞美主，但要掺杂着哀叹：因为当我们赞美时，我们渴望祂，却尚未拥有祂。当我们拥有祂时，我们的一切忧伤将被除去，只剩下赞美，纯洁而永恒。现在，让我们祈祷。\par

\SourceRecord{Translated by J.E. Tweed.；Nicene and Post-Nicene Fathers, First Series；Vol. 8.；Edited by Philip Schaff.；Buffalo, NY: Christian Literature Publishing Co.,；1888.；<http://www.newadvent.org/fathers/1801087.htm>.}

% structure: p0001 source=h1 target=section reason=page-title

\section{圣咏第八十八篇释义}

1. 这篇第八十七篇圣咏的标题包含一个值得探究的新主题：这里出现的\Quote{为\OriginalTerm{梅肋客}{Melech}应唱}一词，在其他地方未曾发现。我们已经对\Quote{\OriginalTerm{诗篇之歌}{Psalmus Cantici}}和\Quote{\OriginalTerm{歌之诗篇}{Canticum Psalmi}}这两个标题的含义发表了看法；而\Quote{\OriginalTerm{科辣黑之子}{sons of Core}}一词不断重复出现，也常被解释；\Quote{直到终结}亦然；但此标题中的下一部分却是独特的。因为\Quote{梅肋客}我们可以译为拉丁文\Quote{为合唱团}，因为合唱团是希伯来词梅肋客的意思。……此处预言了吾主的受难。宗徒伯多禄说：\BibleQuote{基督也为你们受了苦，给你们留下了榜样，叫你们追随他的足迹。}\BibleRef{伯前 2:21}这就是\Quote{应唱}的意思。宗徒若望也说：\BibleQuote{正如基督为我们舍命，我们也应当为弟兄舍命。}\BibleRef{若一 3:16}这也是应唱。但合唱团象征和谐，即在于爱德：因此无论谁，若效法吾主的受苦而将自己的身体交出被焚，却没有爱德，就不在合唱团中应唱，因而对他毫无益处。\BibleRef{格前 13:3}此外，正如拉丁文中使用\Quote{\OriginalTerm{领唱者}{Precentor}}和\Quote{\OriginalTerm{应唱者}{Succentor}}这两个术语来指代音乐中首先演唱者和接续者；同样，在这首受难之歌中，基督在前，殉道者的合唱团跟随，直到获得天堂的荣冠。这就是由\Quote{科辣黑之子}所唱的，即基督受难的效仿者：正如基督在哥耳哥达被钉十字架，而哥耳哥达一词正是希伯来语科辣黑的解释。\BibleRef{玛 27:33}这也是\Quote{以色列人\OriginalTerm{厄曼}{Æman}的智慧}：是此标题结尾处的词语。厄曼据说意为\Quote{他的兄弟}：因为基督屈尊使那些理解祂十字架奥迹的人成为祂的兄弟，他们不仅不以此为耻，反而忠实地以此夸耀，不因自己的功德而自夸，而是感恩于祂的恩宠：以致可以对其中每人说：\BibleQuote{看，这确是一个以色列人，在他内毫无诡诈。}\BibleRef{若 1:47}正如圣经论及以色列本人所说，他毫无诡诈。\BibleRef{创 25:27}\par

2. \BibleQuote{上主，我救恩的天主，我昼夜在你面前呼号。}\BibleRef{咏 87:1}因此，我们现在要听基督在预言中在我们面前歌唱，祂自己的合唱团应以效仿或感恩来应唱。\par

\BibleQuote{愿我的祈祷上达于你面前，求你侧耳俯听我的呼号。}\BibleRef{咏 87:2}因为连吾主也祈祷了，不是在天主性上，而是在仆人的形态上；因为祂在此也受苦。祂既在顺境时祈祷，即\Quote{白日}，也在患难时祈祷，我想这指的是\Quote{黑夜}。祈祷上达于天主面前，就是蒙悦纳；祂侧耳，就是祂怜悯地垂听：因为天主并没有我们这样的身体器官。然而，这段经文照例是重复的。\par

3. \BibleQuote{因为我的灵魂充满了灾祸，我的生命临近了地狱。}\BibleRef{咏 87:3}我们敢说基督的灵魂\Quote{充满了灾祸}么？因为苦难的能力只及于身体……因此灵魂没有身体也能感受痛苦，但身体没有灵魂则不能。那么，我们为何不能说基督的灵魂充满了人性的灾祸，尽管没有人的罪？另一位先知论及祂说，祂为我们忧伤：\BibleRef{依 53:4}而福音作者说：\BibleQuote{他带着伯多禄和载伯德的两个儿子，开始忧闷恐怖起来；}吾主也亲口对他们说：\BibleQuote{我的心灵忧闷得要死。}\BibleRef{玛 26:37}作此圣咏的先知预见此事将要发生，就介绍祂说：\Quote{我的灵魂充满了灾祸，我的生命临近了地狱。}因为这里用其他话语表达了相同的意思，正如祂所说：\Quote{我的心灵忧闷得要死。}\Quote{我的心灵忧闷}这句话，相当于\Quote{我的灵魂充满了灾祸}；而后面那句\Quote{要死}，相当于\Quote{我的生命临近了地狱。}吾主承担了这些人性的软弱情感，正如祂承担了人性软弱的肉身和人性肉身的死亡一样，不是出于祂本性的必然，而是出于祂慈悲的自由意志，为能转化祂的身体——即教会（祂屈尊成为其元首）——也就是祂在圣洁忠信的门徒中的肢体：这样，若有人身处试探中而感到忧闷痛苦，就不会因此认为自己与祂的恩宠隔绝；使身体如同跟随领唱者的合唱团，从其元首学习到这些忧苦并不是罪，而是人性软弱的证明。我们读到宗徒保禄，他是这身体的主要肢体，我们听他承认自己的灵魂充满了这样的灾祸，他说：\Quote{为我的弟兄、我血统的同胞，我心中常有极大的忧愁和不断的痛苦。}如果我们说，吾主在祂受难临近时也因他们而忧伤，因为他们将犯下最可怕的罪，我想这并不过分。最后，我们的救主在十字架上所说的：\BibleQuote{父啊，宽赦他们吧！因为他们不知道他们做的是什么。}\BibleRef{路 23:34}在这篇圣咏下文表达为：\BibleQuote{我算被列于下坑的人中，}\BibleRef{咏 87:4}这是指那些不知他们在做什么的人，他们以为祂像其他人一样死去，受制于必然并为其所胜。\Quote{坑}一词用于指极度痛苦或地狱。\BibleQuote{我成了一个无人帮助的人。}\BibleRef{咏 87:4}\par

4. \BibleQuote{在死者中自由} \BibleRef{咏 87:5}。这些话最清楚地显示了我们的主的位格：因为除了祂，谁在死者中自由呢？祂虽有罪肉的形状，却在罪人中独独无罪？\BibleRef{罗 8:3}……因此，祂，\BibleQuote{在死者中自由}，有能力舍弃自己的性命，也能再取回；无人能从祂取走，而是祂自愿舍弃；祂能按自己的意愿复活自己的肉身，如同被他们毁坏的圣殿；当所有人在祂受难前夕都离弃祂时，祂并不孤单，因为祂作证：祂的圣父没有舍弃祂；\BibleRef{若 8:29}然而，祂仍被祂为之祈祷的仇敌——他们不知道自己所做的是什么——视为\BibleQuote{如同无援助的人；如同受伤而卧在坟墓中的人。}但祂接着说：\BibleQuote{你尚不纪念的人。}在这句话中，要注意基督与其他死者之间的区别。因为虽然祂受伤了，死后被放在坟墓里，\BibleRef{玛 27:50}但那些不知道自己在做什么、也不知道祂是谁的人，把祂看作是那些因受伤而丧亡、卧在坟墓中的人，这些人尚不被天主纪念，也就是说，他们的复活时辰尚未到来。因为圣经这样称呼死者为\Quote{安眠}，是希望他们被视为将要醒来，即复活。但祂受伤并在坟墓中安眠，在第三日醒来，成为\Quote{如同屋顶上独居的麻雀}，即在天堂圣父的右边；现在祂\BibleQuote{不再死，死亡再不能统治祂。}\BibleRef{罗 6:9}因此，祂与那些天主尚未如此纪念而使其复活的人大不相同：因为那要先在元首身上发生的事，是为身体预留在末日的。天主被称为\Quote{纪念}，是当祂行事时；被称为\Quote{忘记}，是当祂不行事时：因为天主既不改变，故不能忘记，也不能回忆，因祂永不会忘记。因此，我被那些不知自己所做的人视为\BibleQuote{一个无援助的人}；而我虽是\BibleQuote{在死者中自由}，他们却把我当作\BibleQuote{如同受伤而卧在坟墓中的人。}然而，这些如此看待我的人，进一步被称为\BibleQuote{从祢手中被割去}，也就是说，当我这样被他们对待时，\BibleQuote{他们就从祢手中被割去}；那些相信我没有援助的人，被剥夺了祢手的援助：因为他们正如他在另一篇圣咏中所说的，在我面前挖了坑，自己却掉进其中。我倾向于这种解释，而不是将\BibleQuote{从祢手中被割去}这句话应用于那些卧在坟墓中、天主尚未纪念的人：因为义人也属于后者，尽管天主尚未召唤他们复活，但经上说他们的\BibleQuote{灵魂在天主手中}，\BibleRef{智 3:1}也就是说，他们\Quote{居于至高者的保护之下，并住在天上天主的荫庇下。}但那些从天主手中被割去的人，是那些相信基督从祂手中被割去，并如此把祂当作恶人来对待而敢于杀害祂的人。\par

5. \BibleQuote{他们将我放在极深的坑中} \BibleRef{咏 87:6}，即最深之坑。因为希腊文如此。但极深的坑若不是极深的苦难，还有什么比它更深呢？为此，在另一篇圣咏中说：\Quote{祢也将我从祸坑中拉出。}\BibleQuote{在黑暗之处和死荫中}，当他们不知道自己所做之时，他们这样判断祂，把祂放在那里；他们不认识祂，\BibleQuote{这世界上的众王侯没有一个认识祂。}\BibleRef{格前 2:8}至于\Quote{死荫}，我不知道应理解为身体的死亡，还是如经上所写的：\BibleQuote{那行走在黑暗和死荫之地的人，已有光照耀他们。}\BibleRef{依 9:2}因为借着信仰，他们从罪的黑暗和死亡中被带入光明和生命。那些不知道自己所做的人，这样看待我们的主，并在他们的无知中，把祂算作祂来要帮助的那些人之一，好使他们自己不再如此。\par

6. \Quote{您的义怒重重地压在我的身上} \BibleRef{咏 87:7}，或者，如其他抄本所载，\Quote{您的愤怒；} 或者，又如其他抄本，\Quote{您的烈怒：} 希腊文 \OriginalTerm{怒气}{\GreekText{θυμὸς}} 经历了不同的翻译。因为凡希腊文抄本有 \OriginalTerm{愤怒}{\GreekText{ὀργὴ}} 之处，译者都毫不犹豫地用拉丁文 \Emph{ira} 来表达；但若出现 \OriginalTerm{怒气}{\GreekText{θυμὸς}} 一词，大多数人反对将其译为 \Emph{ira}，尽管许多最优拉丁文体的作者，在翻译希腊哲学作品时，正是这样将该词译为拉丁文的。不过我不想在此问题上作进一步讨论；只是若我也要提出另一个术语，我认为 \Quote{义怒} 比 \Quote{烈怒} 更可接受，因为在拉丁文中，后者不用于理智清醒的人。那么，\Quote{您的义怒重重地压在我的身上} 是什么意思，除非是那些不认识荣耀之主的 \BibleRef{格前 2:8} 的人所相信的，他们以为天主的愤怒不仅被激起，而且重重地压在那一位身上，他们竟敢将祂置于死地，而且不只是死，还是那种他们视为最可咒骂的死，即十字架的死；因此宗徒说：\Quote{基督救赎了我们脱离法律的诅咒，为我们成了诅咒，因为经上记载：凡挂在树上的，都是可咒骂的。} \BibleRef{迦 3:13} 为此，为了赞扬祂那以极度谦卑而完成的服从，他说：\Quote{祂贬抑自己，听命至死，} 而因为这似乎不够，他又加上：\Quote{且死在十字架上；} \BibleRef{斐 2:8} 并且，据我看来，出于同样的观点，他在本篇圣咏中说：\Quote{您的一切悬垂，} 或如某些人译作 \Quote{波浪，} 另一些人译作 \Quote{抛掷，} \Quote{都加在我的身上。} 我们也在另一篇圣咏中读到：\Quote{您的一切悬垂和波浪都临到我身上，} 或如某些人所译的更好：\Quote{都经过了我身上：} 因为希腊文是 \OriginalTerm{经过}{\GreekText{διῆλθον}}，而不是 \OriginalTerm{进入}{\GreekText{εἰσῆλθον}}；并且当两词 \Quote{波浪} 和 \Quote{悬垂} 同时使用时，二者不能互相等同。在那一段中，我们将 \Quote{悬垂} 解释为威胁，\Quote{波浪} 为实际的苦难：两者都是由天主的审判所施行的；但在那处经文中说：\Quote{一切全经过了我身上，} 此处却说：\Quote{您将这一切全加在我身上。} 在另一种情形，即，虽然某些恶事发生了，但他说，这里所提及的一切都经过了；但在本处，\Quote{您将它们加在我身上。} 恶事作为悬而未决之事，当它们不触及人时，就经过；而当它们触及人时，就如波浪。但当使用 \Quote{悬垂} 一词时，他并未说它们经过了，而是说：\Quote{您将它们加在我身上，} 意思是所有悬而未决之事都已发生。关于祂苦难所预言的一切，只要还存于预言中待将来应验，都是悬而未决的。\par

7. \Quote{您使我的相识远离了我} \BibleRef{咏 87:8}。若将\Quote{相识}理解为祂所认识的人，那就包括所有人；因为祂有谁不认识呢？但祂称那些人为\Quote{相识}，是指那些自己认识祂的人，按他们在那时所能认识的程度而言；至少他们知道祂是无辜的，虽然他们只将祂看作人，而不看为天主。尽管祂也可能称祂所认可的义人为\Quote{相识}，正如祂称恶人为\Quote{不认识的人}，至终祂将对后者说：\Quote{我不认识你们。} \BibleRef{玛 7:23} 在接下来的话中，\Quote{他们使我成了他们憎恶的对象；} 祂之前所称的\Quote{相识}，也可能指这些人，因为他们连对那种死亡的方式也感到恐惧；但更好的解释是指祂前面所说的那些迫害祂的人。\Quote{我被交出，却没有脱离。} 这是因为当祂在屋内受审时，祂的门徒却在外边？\BibleRef{玛 26:56} 还是我们应当对 \Quote{我不能脱离} 作更深的解释，即祂说：\Quote{我仍隐藏在我的隐秘计划中，我未显明我是谁，我没有启示自己，没有显现自己}？于是接下来——\par

\Quote{我的眼睛因匮乏而变得昏花} \BibleRef{咏 87:9}。这里我们应当理解什么眼睛？如果是指祂受苦时所受的肉身的眼睛，我们并未读到祂的眼睛因匮乏（即饥饿）而昏花，就像通常所发生的那样；因为祂在晚餐后被出卖，并在同一天被钉十字架：如果是指内在的眼睛，那么在祂内有一个永不熄灭的光，它们又怎能因匮乏而变得昏花呢？但祂是通过祂的眼睛来指代那些肢体，祂自己是这些肢体的头，祂爱这些肢体如同更明亮、更卓越、超越其余的部分。宗徒在写作时，借用我们身体的比喻，曾论及这个身体：\Quote{若全身是眼睛，听觉在哪里？} 等 \BibleRef{格前 12:17}。他藉这些话所愿表达的，更清楚地说明，补充说：\Quote{你们却是基督的身体，各自作肢体。} \BibleRef{格前 12:27} 因此，当那些眼睛，即圣宗徒们——启示他们的不是血肉，而是那在天上的父，使伯多禄说：\Quote{祢是基督，永生天主之子} \BibleRef{玛 16:16}——看见祂被出卖，遭受如此苦难，看见祂不是他们所希望的那样，因为祂没有出来，没有在德能和权能中显示自己，仍隐藏于祂的隐秘中，像是一个被征服的、软弱的人承受一切，他们就因匮乏而变得昏花，好像他们的食粮、他们的光被撤走了。\par

8. 祂继续说：\Quote{我呼求了祢。} 这确实是在十字架上时最清楚地实行的。但下文：\Quote{我终日向祢伸开我的双手，} 必须考查应如何理解。因为如果我们在此表达中理解为十字架的木头，如何能将其与\Quote{终日}调和呢？能否说祂在十字架上悬了整个一日，将夜晚算作白日的一部分？但如果此表达意指与夜相对的白日，即使在祂被钉的那一日，第一段并非短暂的时间也已经过去。但如果我们取\Quote{日}为时间之意义（尤其因为此词用阴性，这个性在拉丁语中限于此义，虽然在希腊语中并非如此，在那里它总是用阴性，我认为这是为何在我们自己的译本中同样译作阴性的原因），问题的结就会拉得更紧：因为如果祂甚至在十字架上伸开双手还不足一日，如何能意指整个时间段呢？再者，如果我们以整体代部分，正如圣经有时如此使用这个表达，我不记得当明确加上\Quote{整}字时，有以整体代部分的例子。因为在福音中主说：\BibleQuote{人子要在地心里三天三夜，} \BibleRef{玛 12:40} 以部分代整体并非特别牵强，因为表达不是三个\Quote{整}日和三个整夜：因为中间的一个日是整日，另外两个是部分，最后一个日是第一日的一部分，第一夜是最后一夜的一部分。但如果这里不是指十字架，而是指祈祷，我们在福音中读到祂以仆人的形态向天主圣父倾心祈祷，据说祂在受难前长久祈祷，在受难前夕，以及在十字架上时，我们并未读到祂终日如此行。因此，终日伸开的双手，可理解为祂从未停止努力的善工之持续。\par

9. 但正如祂的善行仅有益于那些预定得永生的，并非有益于所有人，甚至连祂行善所在的人群中也不是所有人都得益，祂接着说：\Quote{难道祢要在死人中显示奇迹吗？}\BibleRef{咏 87:10}。若我们以为这是指肉身生命已离去的之人，那么在死人中的确行了伟大的奇迹，因为他们中有一些人复活了：\BibleRef{玛 27:52}；并且在我主下降阴府、作为死亡的征服者上升时，在死人中行了伟大的奇迹。祂在这话中，\Quote{难道祢要在死人中显示奇迹吗？}是指那些心中如此死寂的人，以致基督如此伟大的工作都不能唤醒他们进入信德的生命：因为祂不是说奇迹不显示给他们，因为他们看不见，而是说奇迹不使他们得益。因为，正如祂在此段所说：\Quote{我终日向祢伸开我的双手：}因为祂常将祂的一切工作归因于祂圣父的旨意，不断宣告祂来是为承行祂圣父的旨意：\BibleRef{若 6:38}。同样，一个不信的民族看到同样的工作，另一位先知说：\Quote{我终日向悖逆的民族伸开双手，他们不信，反而抗拒。}\BibleRef{依 65:2}。所以，那些是死人，奇迹没有显示给他们，并非因为他们没有看见，而是因为他们没有藉此重获生命。下一节\Quote{医生岂能使他们复活，使他们赞美祢？}意指：死人不会藉这样的方法复活，使他们赞美祢。在希伯来文中，据说有一个不同的表达：此处用的是\OriginalTerm{巨人}{giants}，而这里用的是\OriginalTerm{医生}{physicians}；但七十贤士译者，他们的权威如此之大，以至于可以当之无愧地说他们是因天主圣神的默感而翻译的，因为他们奇妙的一致同意，他们并非因错误，而是利用希伯来文中表达这两种意思的词在发音上的相似性，得出使用这个词是表示巨人一词所意指的含义的迹象。因为若你们以为巨人是指骄傲的人，宗徒所说的\Quote{智者在哪里？经师在哪里？今世的辩士在哪里？}\BibleRef{格前 1:20}，那么称他们为医生并无不妥，仿佛他们凭自己的独门技巧应许灵魂的救恩；对此有话说：\Quote{救恩属于上主。}但若我们按好的意义来理解巨人一词，如同论及我主时所说的：\Quote{祂如巨人般欢跃奔跑祂的赛程；}即巨人之巨人，在那些最大最强、在祂教会中具有超性力量的人中为首。正如祂是群山之山；正如经上所载：\Quote{到末日，上主殿宇的山必要矗立在群山之上。}\BibleRef{依 2:2}；以及圣中之圣：称这些伟大而有能力的人为医生也无不妥。为此宗徒说：\Quote{或者能激动我的骨肉之亲发愤，好救他们一些人。}\BibleRef{罗 11:14}。但即使是这样的医生，即使他们不凭自己的能力治愈（如肉身医生也不是凭自己），然而就他们藉忠信的职务协助得救而论，他们能治愈活人，却不能使死人复活：正如所说的：\Quote{难道祢要在死人中显示奇迹吗？}因为天主的恩宠，被用来在某种方式上使人的心灵获得新生命，以便能够从任何传道者那里听到救恩的教训，这是极为隐藏而奥秘的。这种恩宠在福音中这样论及：\Quote{除非那派遣我的父吸引他，没有人能到我这里来；}\BibleRef{若 6:44}……为表明，灵魂藉以相信、并从其旧有情欲的死亡中跃入新生命的那信德本身，是天主赐予我们的。所以，无论最好的圣言宣道者和通过奇迹说服真理者对人做了何种努力，都如同大医生一样：然而，若他们死了，且未因祢的恩宠获得第二次生命，\Quote{难道祢要在死人中显示奇迹，或是医生能使他们复活？而他们}——这些被医生复活的人——\Quote{怎能赞美祢？}因为这一宣认表明他们活着；不像别处所载的：\Quote{感谢从死人中消失，如同从无有之人中消失。}\BibleRef{德 27:26}\par

10. \Quote{岂有人在坟墓中述说祢的仁慈，在毁灭中述说祢的信实？}\BibleRef{咏 87:11}。当然，\Quote{述说}一词可理解为重复使用：岂有人在毁灭中述说祢的信实？圣经喜欢将仁慈与信实相连，尤其在圣咏中。\Quote{毁灭}也是对\Quote{坟墓}的复述，指的是在坟墓中的人，即上面在\Quote{难道祢要在死人中显示奇迹吗？}一节中所称的\Quote{死人}：因为身体是死魂灵的坟墓；因此我主在福音中说：\Quote{你们好像刷白的坟墓，外面看来美丽，里面却充满死人的骨骸及各种污秽。同样，你们外面叫人看来好像是义人，里面却满是伪善和不法。}\BibleRef{玛 23:27}\par

11. \BibleQuote{祢的奇事岂能在黑暗中被知道，祢的公义岂能在被遗忘之地被知道？} \BibleRef{咏 87:12}，黑暗对应那遗忘之地：因为黑暗是指不信的人，正如宗徒所说：\BibleQuote{因为你们从前是黑暗；} \BibleRef{弗 5:8}而那被遗忘之地，就是那忘记了天主的人；因为那不信的灵魂能陷入如此浓重的黑暗，\BibleQuote{以致愚妄人心里说：没有天主。} 因此，整个段落的含义可以这样联系上下文得出：\BibleQuote{主，我呼求了祢，}在我的苦难中；\BibleQuote{我整天向祢伸出了我的手} \BibleRef{咏 87:13}。我从未停止伸出我的手来光荣祢。那么，恶人为何对我狂怒，岂不是因为\BibleQuote{祢不在死者中显奇事}？因为这些奇事没有感动他们相信，也没有医生能使他们复活好赞美祢，因为祢隐藏的恩宠没有在他们内运作，吸引他们相信：因为没有人能到我这里来，除非是祢所吸引的。那么，\BibleQuote{祢的慈爱岂能在坟墓中显扬？}这就是说，那死于灵魂的坟墓，灵魂因肉体的重量而死：\BibleQuote{或祢的信实在灭亡中？}就是说，在这样一种不能相信或感觉任何这些事的死亡中。\Quote{因为在这死亡的黑暗中}，也就是说，在那个因忘记祢而失去生命之光的人中，\BibleQuote{祢的奇事和祢的公义怎能被知道。}……\par

12. 但是，为了使那些超越言语的祝福的祈祷更加热切和持久，那将持续到永恒的恩赐被延迟，而短暂的苦难却被允许增多。因此接着说道：\BibleQuote{主，祢为何弃绝了我的祈祷？} \BibleRef{咏 87:14}，这可以与另一篇圣咏相比较：\BibleQuote{我的天主，我的天主，祢为何舍弃了我？} 这理由被当作问题提出，并非好像天主的智慧被指责为无故这样做；这里也是如此。\BibleQuote{主，祢为何弃绝了我的祈祷？}但如果仔细留意这个原因，就会发现它已在上面指明了；因为正是为了使圣人们的祈祷因如此大福的延迟和他们在生活苦难中所遭遇的逆境而仿佛被拒绝，从而使被煽动的火焰燃烧得更加明亮。\par

13. 为此，他在下文中简要描绘了基督身体的苦难。因为这些苦难不仅发生在头身上，因为也对扫罗说：\BibleQuote{你为什么迫害我？} \BibleRef{宗 9:4}而保禄自己，作为同一身体中的一个蒙拣选的肢体，说：\BibleQuote{在我身上，为基督的身体——教会，补充基督苦难所欠缺的。} \BibleRef{哥 1:24} \BibleQuote{那么，主，祢为何弃绝了我的灵魂？为何向我掩面？}\par

\BibleQuote{我自幼贫穷，在劳苦中；我虽被高举，却被摔倒，惊慌失措} \BibleRef{咏 87:15}。\par

\BibleQuote{祢的愤怒从我身上经过，祢的威吓使我惊惶} \BibleRef{咏 87:16}。\par

\BibleQuote{它们终日如水环绕我，一起围困我} \BibleRef{咏 87:17}。\par

\BibleQuote{祢使我的朋友远离我，使我的熟人因我的痛苦而疏远} \BibleRef{咏 87:18}。所有这些灾祸都已在基督身体的肢体中发生，并且正在发生，天主向他们的祈祷转面，不按他们所愿望的垂听，因为他们不知道他们的愿望若实现并不对他们有益。教会是\BibleQuote{贫穷的}，因为她在旅途中饥渴慕义，渴望那将在本乡充满她的食粮；她\BibleQuote{自幼在劳苦中}，正如基督的身体在另一篇圣咏中所说：\BibleQuote{从我幼年以来，敌人屡次迫害我。}也正因如此，她的一些肢体甚至在此世被高举，好使他们在其中更加谦卑。那构成圣人和信徒合一、以基督为头的身体，受到天主的愤怒；但并未存留，因为只有关于不信者才记载说：\BibleQuote{天主的愤怒常在他身上。} \BibleRef{若 3:36}天主的威吓惊扰信友的软弱，因为一切可能发生的，即使实际上未发生，也宜谨慎恐惧；有时这些威吓如此扰乱那反省的灵魂，因着周围临头的灾祸，以致它们如水四面环绕我们，并在我们的恐惧中包围我们。由于教会在此世旅途中从未免于这些灾祸，它们时而发生在这一个肢体上，时而在另一个上，他加上\BibleQuote{终日}，表示持续的时间，直到世界的终结。常常，朋友们和熟人们，为了他们的世俗利益，在恐惧中离弃圣人们；关于此，宗徒说：\BibleQuote{众人都离弃了我，但愿这罪不归于他们。} \BibleRef{弟后 4:16}但这一切有何目的呢？无非是为了在清晨，即在不信的夜晚之后，这圣身体的祈祷得以在信德的光中先于天主，直到那救恩来临，我们目前是因希望而得救，并非因拥有，我们以耐心和信德期待它。那时，主将不会拒绝我们的祈祷，因为再也没有什么可求的，而一切合理所求的都将获得；祂也不会向我们转面，因为我们将看见他实在怎样：\BibleQuote{我们将看见祂实在怎样} \BibleRef{若一 3:2}我们也不再贫穷，因为天主将成为我们的富足，成为万有中的万有：\BibleQuote{天主将成为我们的富足，成为万有中的万有} \BibleRef{格前 15:28}我们将不再受苦，因为不再有软弱；被高举之后，也不会遇到卑辱和混乱，因为那里没有逆境；甚至不会承受天主的瞬间愤怒，因为我们将在祂不变的爱中安息；祂的威吓不会惊扰我们，因为祂应许的实现将祝福我们；我们的朋友和熟人，也不会因恐惧而远离我们，因为那里没有可怕的仇敌。\par

\SourceRecord{Translated by J.E. Tweed.；Nicene and Post-Nicene Fathers, First Series；Vol. 8.；Edited by Philip Schaff.；Buffalo, NY: Christian Literature Publishing Co.,；1888.；<http://www.newadvent.org/fathers/1801088.htm>.}

% structure: p0001 source=h1 target=section reason=page-title

\section{圣咏第八十九篇释义}

1. 亲爱的，请因天主的恩宠理解这篇我即将讲解的圣咏，这是关于我们对主耶稣基督的望德，并要欢欣鼓舞，因为那应许者必会全部实现，正如祂已经实现了很多；因为给我们对祂信赖的不是我们自己的功德，而是祂的仁慈。按我的理解，标题所指的正是\Quote{以色列人\Person{厄堂}的领悟}，这篇圣咏也因此得名。你们看，厄堂指的是谁。但这个词的意思是\Quote{坚强}。这世上没有一个人是坚强的，除非怀着对天主诸许的望德；因为按我们自己应得的，我们是软弱的，但在祂的仁慈中，我们是坚强的。那么，圣咏作者在自己内软弱，在天主的仁慈中坚强，就这样开始：\BibleQuote{我要永远歌颂祢的仁慈，主；我的口要传扬祢的信实直到万代}（\BibleRef{咏 88:1}）。\par

2. 让我的肢体服侍主，他说：我说话，但我讲的是祢的事。\Quote{我的口要传扬祢的信实}：若我不服从祢，我就不是祢的仆人；若我凭自己说话，我就是说谎者。那么，讲论祢的事和凭我自己的身份说话，是两回事：一是我的，一是祢的；信实是祢的，言语是我的。让我们听听他要传扬的是怎样的信实，他要歌颂的是怎样的仁慈。\par

3. \BibleQuote{因为祢说过：仁慈要永远被建立}（\BibleRef{咏 88:2}）。这就是我所歌颂的：这是祢的信实，我的口正是为传扬它而服务。祢这样说：我建立，是为了不摧毁；因为有些人祢摧毁而不建立；有些人祢摧毁后又重新建立。因为若不是有些人被摧毁是为了重新被建立，耶肋米亚就不会写道：\BibleQuote{你看，我今天立了你来拆毁和建立。}（\BibleRef{耶 1:10}）的确，所有从前敬拜偶像和石头的人，若不先摧毁他们旧日的谬误，就不能在基督内被建立。然而，除非有些人被摧毁而不再被建立，就不会有这记载：\Quote{他要摧毁他们，而不建立他们。}……在接下来的经文中，他把这两个词——仁慈和信实——合并在一起；\BibleQuote{因为祢说过：仁慈要永远被建立；祢的信实要在天上坚立}——其中仁慈和信实重复出现，\BibleQuote{因为主的一切道路都是仁慈和信实}，因为应许应验中的信实，除非先有罪过赦免中的仁慈，就无法显明。其次，如许多事情在预言中应许给那按血统由亚巴郎的子孙而来以色列人民，那人民也得以增长，使天主的诸许能在他们中实现；同时天主甚至没有对外邦人关闭祂善意的源泉，虽然祂将他们置于天使的管辖之下，却保留了以色列人民作为自己的一份；宗徒明确提到主的仁慈和信实，是分别针对这两方。因为他称基督是\BibleQuote{为天主的信实而作割礼的仆役，为证实对祖先的恩许}（\BibleRef{罗 15:8}）。你看天主如何未行欺骗；你看祂如何未弃绝祂所预知的人民。因为当宗徒处理犹太人的跌倒时，为防止有人以为他们已被天主如此抛弃，以致那场面上簸扬的麦子一粒都不能到达粮仓，他说：\BibleQuote{天主并没有弃绝祂所预知的人民；因为我也是以色列人。}（\BibleRef{罗 11:1}）如果那整个民族都是荆棘，那我对你们说话的我怎么会是麦子呢？如此，天主的信实在那些信了的以色列人中得以应验，来自割礼的一面墙就这样被带来与角石相遇。但这石头不会形成一处角，除非它接纳来自外邦人的另一面墙；因此，前一面墙特别关乎信实，后一面关乎天主的仁慈。\BibleQuote{但是我说，耶稣基督为了天主的信实成了割礼的仆役，为证实对祖先的恩许，并使外邦人因祂的仁慈而光荣天主。}（\BibleRef{罗 15:8}）因此，\BibleQuote{祢的信实祢要在天上竖立}这句话被加进来，是合理的；因为所有那些被召成为宗徒的以色列人，都变成了宣讲天主光荣的诸天；正如他们笔下所写：\BibleQuote{诸天述说天主的光荣，穹苍显示祂的化工。}……因为，虽然他们在教会充满普世之前就被提升到此处，但如同\Quote{他们的话语传遍了世界地极}，我们有理由认为我们刚才所读的\BibleQuote{祢的信实祢要在天上竖立}在他们身上应验了。\par

4. \BibleQuote{祢说过：我与我的被选者立了盟约}（\BibleRef{咏 88:3}）。什么盟约？岂不就是新约，藉此我们被更新为新的产业，因着我们对它的渴望和爱，我们唱一首新歌。\BibleQuote{我与我的被选者立了盟约，}圣咏作者说，\BibleQuote{我向我的仆人达味起了誓。}那领悟的人，他的口传扬信实，说得何等自信！我毫不畏惧地说话，因为\Quote{祢已经说过}。如果祢因祢说了而使我不怕，那么当祢起誓时，祢岂不更使我如此！因为天主的誓言是恩许的保证。人正当地被禁止发誓（\BibleRef{玛 5:34}），免得因发誓的习惯，由于人可能受欺骗，而陷于妄证。\par

5. 那么，让我们看看天主所誓许的是什么。祂说：\BibleQuote{我已誓许，}祂说，\BibleQuote{对\Person{达味}我的仆人：我要永远立定你的后裔} \BibleRef{咏 88:4}。但\Person{达味}的后裔不就是\Person{亚巴郎}的后裔吗？\Person{亚巴郎}的后裔是什么？祂说：\BibleQuote{而且是给你的后裔，}祂说，\BibleQuote{就是基督} \BibleRef{迦 3:16}。但或许那基督——教会的元首、身体的救主 \BibleRef{弗 5:23}——是\Person{亚巴郎}的后裔，因此也是\Person{达味}的后裔；但我们不是\Person{亚巴郎}的后裔吗？我们当然是；正如宗徒所说：\BibleQuote{如果你们属于基督，那么你们就是\Person{亚巴郎}的后裔，按照恩许的继承人} \BibleRef{迦 3:29}。因此，弟兄们，让我们这样理解这句话：\BibleQuote{我要永远立定你的后裔}，不仅指基督由\Person{童贞玛利亚}所生的肉身，也指我们所有相信基督的人，因为我们是那元首的肢体。这身体不能失去它的元首：如果元首永远在光荣中，肢体也是如此，这样基督便永远保持完整。\BibleQuote{我要永远立定你的后裔：并设立你的宝座，直到万代。}我们认为祂说\BibleQuote{永远}，是因为\BibleQuote{直到万代}；因为祂在上面说过：\BibleQuote{我要用我的口永远向万代宣扬你的真理。}\BibleQuote{直到万代}是什么意思？到每一个世代：因为这个词不需要像若干世代来临与过往那样多次重复。通过重复，世代的增多被表明和展示出来。是否应该理解为两个世代，正如你们所知，我亲爱的弟兄们，以及我之前所解释的？因为现在有一个血肉的世代；将来在死者复活中有一个未来的世代。基督在这里被宣讲；祂在那里也将被宣讲：在这里祂被宣讲，好使人相信祂；在那里，祂将被迎接，好使人看见祂。\BibleQuote{我要从这一代到另一代设立你的宝座。}基督如今在我们内有一个宝座，祂的宝座设立在我们内：因为除非祂坐在我们内统治我们，祂就不会管理我们；但如果我们不受祂管理，我们就会自我倾覆。因此祂坐在我们内，统治我们；祂也坐在另一个世代中，就是将来从死者复活而来的世代。基督将永远统治祂的圣徒。天主已应许了这事；祂已说了：如果这还不够，天主已誓许了。既然这应许是确定的，不是因我们的功德，而是因祂的怜悯，没有人应当在宣讲那他所不能怀疑的事时感到害怕。那么，让那力量激励我们的心，\Person{厄堂}的名字正是来源于此——\Quote{坚强的心}：让我们宣讲天主的真理、天主的言语、祂的应许、祂的誓言；让我们通过这些方法在各处得到坚固，光荣天主，并通过将祂携带在我们内，成为诸天。\par

6. \BibleQuote{上主，诸天要赞颂祢的奇事} \BibleRef{咏 88:5}。诸天不会赞颂自己的功绩，而是赞颂祢的奇事，上主。因为在每一个对流失者的怜悯、对不义者的成义行为中，我们所赞颂的，岂不是天主的奇事吗？祢赞颂祂，因为死者已复活；更要赞颂祂，因为流失者被救赎。何等恩宠，何等天主的慈悲！你看见一个人，昨天是酗酒的漩涡，今天是节制的装饰；昨天是奢侈的污水池，今天是贞洁的美德；昨天是亵渎天主的人，今天是赞美祂的人；昨天是受造物的奴隶，今天是造物主的崇拜者。从所有这些绝望的状态中，人们如此被转变：让他们不要看自己的功德，让他们成为诸天，并赞颂那使他们成为诸天的天主的奇事……\par

7. \Quote{因为有谁在天上的云彩中，能与主相比？} \BibleRef{咏 88:6}。这是诸天的赞美吗？这是他们的雨水吗？什么？传道者是否因\Quote{云彩中无人能与主相比}而自信？兄弟们，这看来是赞美的一个崇高理由，即云彩不能与它们的造物主相比吗？若按字面意义而非奥秘意义理解，并非如此：什么？在云彩之上的星辰能与主相比吗？什么？太阳、月亮、天使、诸天，能与主相比吗？那么，他说这话时，仿佛意指某种崇高的赞美：\Quote{因为有谁在天上的云彩中？}等等。我的兄弟们，我们理解这些云彩，如同诸天，乃是真理的传道者：先知、宗徒、天主圣言的宣布者……因此，如果云彩是真理的传道者，让我们首先探问他们为何是云彩。因为同一些人既是诸天又是云彩：因真理的光辉而成为诸天，因肉身的隐晦而成为云彩；因为所有的云彩都因他们的可朽性而朦胧；且它们来来去去。正是由于肉身的这些隐晦，也就是云彩，宗徒说：\Quote{所以，时候未到，你们什么也不要判断，直到主来到，祂要照亮黑暗中的隐密事。} \BibleRef{格前 4:5}此刻你看见一个人正在说什么；但他心里有什么，你却看不见：从云彩中逼出的，你看见了；云彩内隐藏的，你见不着。因为谁的眼睛能穿透云彩？因此，云彩就是肉身中的真理传道者。万物的造物主亲自在肉身中来临……我们因肉身而被称为云彩，因云彩的雨水而成为真理的传道者；但我们的肉身以一种方式来临，祂的则以另一种。我们也被称为天主的儿子，但祂是天主之子，具有另一种意义。祂的云彩来自童贞女，祂是从永远就是圣子，与圣父同永恒。\Quote{那么，在天上的云彩中，有谁能与主相比？在众子中，有谁能与主相似？}让主亲自说，祂能否找到一个与自己相似者。\Quote{人们说人子是谁？}因为我显现，因为我被看见，因为我行走在你们中间，或许此刻我已变得平常：请说：\Quote{人们说我是谁？}他们以人的传言回答祂：\Quote{有的说你是洗者若翰；有的说是厄里亚；有的说是耶肋米亚，或先知中的一位。}这里提到了许多云彩和天主的众子：因为他们都是义人且圣洁，作为天主的众子，耶肋米亚、厄里亚和若翰也被称为天主的众子；作为天主的传道者，他们被称为云彩。你们说了人们以为我是谁的云彩：那么你们也说：\Quote{你们说我是谁？}伯多禄以众人之名代答，一人代表那合一者，回答说：\Quote{你是基督，永生天主之子；} \BibleRef{玛 16:13}不像那些不能与你相比的天主众子：你已在肉身中来临：但不像那不能与你相比的云彩。\par

8. ……\Quote{天主在义人的会议中，极其可畏，在祂四周所有人中，应受敬畏。} \BibleRef{咏 88:7}天主无所不在；那么，谁是祂四周的人呢？祂遍及各处。因为如果祂有一些人在祂四周，那祂就被描绘成各面有限的了。此外，若真对天主并论及天主说：\Quote{祂的伟大，无穷无尽；}谁存留，谁是祂四周的人？除非因为那位无所不在者，选择在一个地方由肉身而生，在一个民族中居住，在一个地方被钉十字架，从一个地方复活并升天。祂在那里做了这一切，外邦人在祂四周。若祂留在祂做这些事的地方，祂就不会\Quote{伟大，并在祂四周所有人中应受敬畏}；但既然祂在那里宣讲时，以这样的方式派遣祂自己名字的传道者到全世界所有民族中；藉着在祂仆人中行奇迹，祂变得\Quote{伟大，并在祂四周所有人中应受敬畏}。\par

9. \Quote{万军的上主，谁能相似祢？祢的真理，至大之主，四面环绕。} \BibleRef{咏 88:8}祢的权能伟大：祢造了天地及其中万物；但祢的慈爱更伟大，将祢的真理显示给祢四周所有人。因为若祢只在祢曾屈尊降生、受苦、复活、升天的地方被宣讲，那么天主应许的真理就会实现，以坚定对祖先所作的许诺；但\Quote{使外邦人因天主的怜悯而光荣天主} \BibleRef{罗 15:9}的许诺就不会实现，若非那真理从祢曾屈尊显现的地方被解释并传播到祢四周的人。在那个地方，祢从祢自己的云彩中发出雷声；但为将雨水撒播给周围的外邦人，祢派遣了其他的云彩。确实，在祢的权能中，祢实现了祢所说的：\Quote{从此以后，你们要看见人子坐在权能者的右边，乘着天上的云彩降来。} \BibleRef{玛 26:64}\par

10. …因為你們如同習慣受到天主雲雨灌溉的人，聽到了\Quote{你的真理}遂即在\Quote{你的周圍}。但是，沒有迫害、沒有反對的時候，何時說過\Quote{他生來是為成為一個被反對的記號}？\BibleRef{路 2:34} 既然那民族，你曾屈尊降生並居住在其中，如同一塊與外邦人波浪隔開的土地，顯得乾旱，準備好接受雨水，而其餘的民族如同大海，處於他們荒蕪的苦澀中；那麼，你的宣講者們，在你周圍散播你的真理，當那大海的波浪猛烈翻騰時，他們又當如何？\Quote{你管轄了海中的高傲}\BibleRef{詠 88:9}。因為大海如此翻騰的結果是什麼？不就是我們今天所慶祝的節日嗎？它殺死了殉道者，播下了血的種子，教會的豐收便生長起來。那麼，讓雲彩安全地出發吧：讓它們在你周圍散播你的真理，不要害怕野蠻的波浪。\Quote{你管轄了海中的高傲。}大海膨脹、衝擊、咆哮；但是\Quote{天主是忠信的，他不許你們受試探超過你們所能受的}\BibleRef{格前 10:13}；因此，\Quote{你平定了它的波浪，當它們漲起時。}\par

11. 最後，你自己在海上做了什麼，以平息它的狂怒並削弱它？\Quote{你打傷了驕傲者，如同受傷的人}\BibleRef{詠 88:10}。在海中有一條驕傲的蛇，聖經的另一處經文說：\Quote{我必吩咐蛇，牠必咬他}\BibleRef{亞 9:3}；又說：\Quote{那裡有鱷魚，你所造來戲弄他的}。他的頭在水面上被打碎。他說：\Quote{你打傷了驕傲者，如同受傷的人。}你自己謙抑了自己，驕傲者就被打傷了；因為驕傲者藉著驕傲抓住了驕傲的人；但那偉大的謙抑了自己，藉著信賴他而變為微小。當弱小者因那從偉大降下到謙卑的榜樣而得滋養時，魔鬼便失去了他所掌握的；因為驕傲的只能抓住驕傲的人。當這樣的謙卑榜樣擺在他們面前時，人學會了譴責自己的驕傲，並效法天主的謙卑。同樣，魔鬼也因失去那些在他權下的人而被打傷；不是受懲戒，而是被推翻。\Quote{你打傷了驕傲者，如同受傷的人。}你自己受了傷，也打傷了別人：你的血流出，為了塗抹罪債的記錄\BibleRef{哥 2:14}，不能不傷害他。因為他驕傲的根據是什麼？不就是他所持反對我們的罪狀嗎？這罪狀、這記錄，你用你的血塗抹了；因此你打傷了他，從他那裡救出了這麼多犧牲者。你該明白魔鬼被打傷了，不是由於肉體的刺傷（他沒有肉體），而是由於他驕傲的心被擊碎。\Quote{你用大能的手臂驅散你的仇敵。}\par

12. \Quote{諸天是你的，大地也是你的}\BibleRef{詠 88:11}。從你而來，在你的大地上降雨。諸天是你的，他們在你的周圍宣講你的真理；\Quote{大地是你的}，它在你周圍領受了你的真理；這雨水產生了什麼結果？\Quote{你奠定了世界和其中所充滿的。}\Quote{你創造了北方和海洋}\BibleRef{詠 88:12}。因為沒有任何事物有能力反抗你，反抗他的創造者。世界固然可以因自己的惡意和意志的乖戾而狂怒，但它能越過創造者所設的界限嗎？那創造了萬有的。那麼，我為何害怕北風？我為何害怕海洋？是的，在北方有魔鬼，他曾說：\Quote{我要坐在北方的極處；我要與至高者相似}\BibleRef{依 14:13}；但你打傷了驕傲者，如同受傷的人。如此，你在他們身上所做的，對於你的統治比他們自己的意志對於他們的邪惡更有力量。\Quote{你創造了北方和海洋。}\par

13. \Quote{大博爾和赫爾孟因你的名而歡樂。}這些山在此被理解，但它們有含意。\Quote{大博爾和赫爾孟因你的名而歡樂。}大博爾，解釋為\Quote{臨近的光}。但那光從何而來，關於它曾說：\Quote{你們是世界的光}\BibleRef{瑪 5:14}，若不是出自那記載的：\Quote{那普照每人的真光正在進入這世界}\BibleRef{若 1:9}？那麼，這光，即世界的光，來自那光，它不是從任何其他源頭點燃的，所以不必擔心它會熄滅。這樣，光就從他那裡來，他是那盞燈，不是放在斗底下，而是放在燈臺上；大博爾是\Quote{臨近的光}。赫爾孟的意思是\Quote{他的咒罵}。正義的光來臨，並成為對他的咒罵。對誰？除了那魔鬼，那受傷的、驕傲的？那麼，我們的照明是從你而來；他被我們視為可咒罵的，因為他曾以他自己的錯誤和驕傲抓住我們，這也是從你而來。\Quote{因此，大博爾和赫爾孟要歡樂}，不是靠自己的功勞，\Quote{而是因你的名}。因為他們要說：\Quote{主啊，榮耀不要歸於我們，不要歸於我們，但要歸於你的名}，因為那翻騰的大海，免得\Quote{外邦人說：他們的天主在哪裡？}\par

14. \Quote{你有大能的手臂}\BibleRef{詠 88:13}。願沒有人將任何事歸於自己。\Quote{你有大能的手臂}：我們由你所造，我們由你所保護。\Quote{你有大能的手臂：願你的手強而有力，願你的右手被高舉。}\par

15. \BibleQuote{义德与审判是祢宝座的预备} \BibleRef{咏 88:14}。\newline{}祢的义德与审判将在终结时显明：它们现今是隐藏的。另有一篇圣咏论及祢的义德，\Quote{论及圣子的隐藏之事}。那时，祢的义德与审判必彰显：有些人将被置于右边，另一些在左边；\BibleRef{玛 25:33}而不信者将战栗，当他们看见现今所嘲笑且不信的事；义人将欢欣，当他们看见现今虽未见却相信之事。\BibleQuote{义德与审判是祢宝座的预备}：尤其是在审判之日。那么现今如何呢？\Quote{仁慈和真理走在祢面前。}我本当畏惧祢宝座的预备、祢的公义和将要来临的审判，若非仁慈和真理行在祢面前：为何我应在终结时畏惧祢的义德，既然祢的仁慈行在祢面前，涂抹了我的罪，并藉显示真理实现了祢的应许？\Quote{仁慈和真理走在祢面前。}因为\Quote{上主的一切道路都是仁慈和真理。}\par

16. 在这一切中，我们岂不欢欣？或我们岂能抑制我们的喜乐？或言语岂足以表达我们的喜乐？或舌头岂能说出我们的欢跃？因此，若言语不足，\BibleQuote{上主，认识欢跃呼喊的百姓，真是有福。} \BibleRef{咏 88:15}。有福的百姓！你正确地领会吗？你理解欢跃的呼喊吗？因为除非你理解欢跃的呼喊，你便不能有福。我所谓理解欢跃的呼喊是何意？你是否知道那无法言喻的喜乐的根源？因为这喜乐并非出于你自己，因为\BibleQuote{夸耀的，应因主而夸耀。} \BibleRef{格前 1:31}。所以不要因自己的骄傲而欢欣，而要因天主的恩宠而欢欣。看，那恩宠是如此伟大，以至舌头无法表达其伟大，然后你便理解欢跃的呼喊。……上主啊，\BibleQuote{他们要在祢面容的光明中行走。他们要因祢的名终日欢欣} \BibleRef{咏 88:16}。那座大博尔和赫尔孟山要因祢的名欢欣：他们若愿意，全日都要因祢的名欢欣；但他们若因自己的名欢欣，就不会终日欢欣：因为当他们以自己为乐而因骄傲跌倒时，他们就不会继续在喜乐中。因此，为能终日欢欣，\Quote{他们要因祢的名欢欣，并因祢的义德被高举。}不是因自己的义德，而是因祢的义德：免得他们有天主的热诚，却不按认识。因为宗徒指出有些人，他们有天主的热诚，却不按认识，\BibleQuote{不明了天主的义德，企图建立自己的义德}，且不因祢的光明而欢欣，因而\BibleQuote{没有服从天主的义德。} \BibleRef{罗 10:2}。为什么？因为\BibleQuote{他们有天主的热诚，却不按认识。}但那认识欢跃呼喊的百姓（因为前者因缺乏认识而犯错，但认识欢跃呼喊的百姓是有福的），他们应从何处呼喊，从何处欢欣，岂不是因祢的名，在祢面容的光明中行走？他们理应被高举，但是在祢的义德中：愿人人完全抛开自己的义德，战栗不已；天主的义德必来临，他必被高举，\Quote{并因祢的义德被高举。}\par

17. \BibleQuote{因为祢是他们力量的荣耀；并因祢的善意，祢必高举我们的角} \BibleRef{咏 88:17}：因为祢看为善，并非因我们配受。\par

18. \BibleQuote{因为我们的被举起是出于上主} \BibleRef{咏 88:18}。因为我曾像一堆沙，行将倒塌；我本已倒塌，若非上主将我举起。\BibleQuote{因为我们的被举起是出于上主，并出于以色列的圣者我们的君王。}祂就是你的举起，祂就是你的光照：在祂的光中你安全，在祂的光中你行走，在祂的义德中你被高举。祂举起你，祂保守你的软弱：祂赐给你力量，出于祂自己，而非出于你。\par

19. \BibleQuote{祢曾在异象中对祢的儿子们说话，说} \BibleRef{咏 88:19}。祢在祢的异象中说话。祢向祢的先知启示了这事。为此，祢在异象中说话，即在启示中说话；先知因此被称为先见。他们在内里看见某事，要在外面宣讲；他们秘密地听见，公开地宣讲。然后\BibleQuote{祢曾在异象中对祢的儿子们说话，说：我已将救助放在一位大能者身上。}你理解这大能者是谁？\BibleQuote{我从百姓中高举了一位特选者。}这特选者是谁？那位你已欢欣祂已被高举的那位。\par

20. \BibleQuote{我找到了我的仆人大卫：}那从大卫的后裔而出的大卫：\BibleQuote{我用我的圣油傅了祂} \BibleRef{咏 88:20}：因为论及祂说：\BibleQuote{天主，祢的天主，以喜乐的油傅了祢，胜过祢的同伴。}\par

21. \BibleQuote{我的手必扶持祂，我的臂膀必坚固祂} \BibleRef{咏 88:21}：因为人的被举起发生了；因为肉身由童贞玛利亚的子宫所摄取，\BibleRef{路 1:31}因为那位在神性形式上与圣父平等的，取了仆人的形式，并服从至死，且死在十字架上。\BibleRef{斐 2:6}\par

22. \BibleQuote{仇敌不能对祂施行强暴} \BibleRef{咏 88:22}。仇敌确实狂怒，但他不能对祂施行强暴：他惯常伤害，但他不能伤害。那么他如何折磨祂？他必锻炼祂，但必不伤害祂。他的狂怒必有裨益；因为他所狂怒的对象，将在得胜中被加冕。因为，若他不向我们发怒，我们如何战胜他？或者，若我们不战斗，天主我们的助佑在哪里？因此，仇敌必尽其所能；但\BibleQuote{他不能对祂施行强暴：邪恶之子不能前来伤害祂。}\par

23. \BibleQuote{我必在祂面前粉碎祂的仇敌}\BibleRef{咏 88:23}。他们因自己的阴谋而粉碎，并且因他们相信而粉碎；因为他们逐渐相信；正如牛犊的头被磨碎，他们将变为天主百姓的饮料。因为梅瑟把牛犊的头磨碎，撒在水面上，叫以色列子民喝。\BibleRef{出 32:20} 所有不信的人都磨碎了：他们逐渐相信；他们被天主的百姓喝下，并进入基督的身体。\BibleQuote{我必在祂面前粉碎祂的敌人，并赶跑恨祂的人。}\par

24. \BibleQuote{我的忠信和慈爱也与祂同在}\BibleRef{咏 88:24}。上主的一切道路都是慈爱和忠信。要记住，尽你所能，这两项德性多么常被强调，使我们回报给天主。因为祂向我们显示慈爱以消除我们的罪，显示忠信以实践祂的许诺；同样，我们行走在祂的道路上，也该回报祂以慈爱和忠信：慈爱，在于怜悯不幸的人；忠信，在于不偏私地判断。不要让忠信夺去你的慈爱，也不要让慈爱妨碍忠信：因为如果你因慈爱而作出了违反忠信的判断，或因严格的忠信而忘记了慈爱，你便没有行走在天主的道路上，那里\BibleQuote{慈爱和忠信彼此相迎。}\BibleQuote{因我的名，祂的角必被高举。}我何必多说呢？你们是基督徒，请认出基督。\par

25. \BibleQuote{我要将祂的手安置在海上}\BibleRef{咏 88:25}：就是说，祂将统治外邦人；\BibleQuote{祂的右手在江河中。}河流注入大海：贪婪的人们滚滚坠入这个世界的苦海中；然而所有这些人都将服从于基督。\par

26. \BibleQuote{祂要称呼我说：祢是我的父亲，是我救恩的举起者}\BibleRef{咏 88:26}。\par

\BibleQuote{我也要立祂为长子，为地上诸王中最高的}\BibleRef{咏 88:27}。我们正在庆生日的殉道者们，曾为这些虽未见却相信的事而流血；如今我们既见到他们所信的，岂不更应勇敢？因为他们那时尚未见到基督在诸王中受高举；那时诸王还在商议敌对上主和祂的受傅者；同一篇圣咏后面的话尚未应验：\BibleQuote{所以，你们君王，应当明智；地上的审判官，应受教。}如今，基督确实已在诸王中受高举。\par

27. \BibleQuote{我要为祂永远保存我的慈爱，我与祂的盟约忠实可靠}\BibleRef{咏 88:28}。因着祂，盟约是忠实的；在祂内，盟约得以订立；祂是盟约的封印者、中保、保证、见证、遗产、共同继承人。\par

28. \BibleQuote{我要使祂的后裔永远存续}\BibleRef{咏 88:29}。不仅在此世，而且直到永无穷尽之世：祂的后裔，即祂的产业，亚巴郎的后裔，就是基督，将要进入那里。但如果你属于基督，你也就是亚巴郎的后裔；如果你注定作祂永远的继承人，\BibleQuote{祂要立定祂的后裔直到永远，祂的宝座如诸天之日。}地上君王的宝座如同大地之日；诸天之日与大地之日不同。诸天之日就是那些岁月，经上说：\BibleQuote{祢却一样，祢的岁月没有穷尽。}大地之日很快被接替者超越：前面的与我们隔绝；后面的也不停留；它们来了又去，几乎未到已去。大地之日就是如此。但诸天之日，也就是诸天的\Quote{一日}，那永不穷尽的岁月，既无起始也无终结；那里没有一日被夹在昨天和明天之间；无人盼望未来，也不失去过去；诸天之日永远是现在，那里祂的宝座将永永远远……\par

29. 这是天主许诺的坚强保证。这位达味的儿子们是新郎的子女；因此所有基督徒都称为祂的儿子。但天主的许诺确实重大：如果基督徒，即\BibleQuote{如果祂的儿子们抛弃我的法律，不遵行我的判断}\BibleRef{咏 88:30}；\BibleQuote{如果他们亵渎我的规条，不遵守我的诫命}\BibleRef{咏 88:31}；我不会轻弃他们，也不会把他们从我面前逐出以致灭亡；但我要做什么？\BibleQuote{我必用杖惩罚他们的过犯，用鞭责罚他们的罪}\BibleRef{咏 88:32}。这不仅是一位呼唤者的仁慈；也是那惩罚和鞭打者的仁慈。因此，让你父亲的手临到你；如果你是善子，就不要拒绝管教；因为\BibleQuote{哪个儿子，父亲不惩戒他呢？}\BibleRef{希 12:7} 愿祂惩戒他，只要祂不收回祂的慈爱；愿祂打这执拗的儿子，只要祂不剥夺他的继承权。如果你明白了你父亲的许诺，不要害怕受鞭打，而要害怕被剥夺继承权：\BibleQuote{因为主所爱的，祂必惩戒；祂鞭打祂所收纳的每一个儿子。}\BibleRef{希 12:6} 有罪的儿子岂能拒绝管教，当他看见无罪的独生子被鞭打？\BibleQuote{我必用杖惩罚他们的过犯。}宗徒也同样警告：\BibleQuote{你们愿意怎么样？要我带着棍棒到你们那里去吗？}\BibleRef{格前 4:21} 不要让虔诚的儿女说：如果你带着棍棒来，那就根本不要来。因为用父亲的杖受教，总比在强盗的抚爱中灭亡更好。\par

\Quote{然而，我的慈爱仍不离开他}\BibleRef{咏 88:33} 从谁呢？从那位达味，我曾给他这些许诺，我曾用喜乐之圣油膏他，使他超越同伴。你们认出这位天主决不收回他慈爱的那一位了吗？免得有人焦虑地说，既然他讲论的是基督——他不会收回对他的慈爱，那么罪人将如何呢？他曾说过这样的话吗？\Quote{我决不将我的慈爱完全从他们收回}？他说：\Quote{我要用刑杖}，他说，\Quote{责罚他们的过犯，用鞭子责罚他们的罪过}。你们为了自己的安全，指望\Quote{我决不将我的慈爱完全从他们收回}。事实上，有些抄本是这样读的，但并非最准确的；不过，凡是这样读的地方，这种读法也绝不与原意相矛盾。因为怎能说天主不会完全收回他对基督的慈爱呢？这位身体的救主在天上或地上犯了什么罪吗？他\Quote{坐在天主的右边，也为我们代求}？\BibleRef{罗 8:34} 然而这是就基督而言：但指他的肢体，他的身体就是教会。因为在这种意义上，他说他不会收回对他的慈爱是一件大事，好像我们不认识那在父怀里的独生子；\BibleRef{若 1:18} 因为在那里，人并不算为他的位格，而是那一个位格是天主也是人。因此，当他没有从他身体、他肢体上收回他的慈爱时，他也就没有完全收回他对他的慈爱；在这些肢体身上，当他坐在天上时，他仍然在地上受迫害；当他从天上呼喊说：\Quote{扫禄，扫禄}，不是说你为什么迫害我的仆人，也不是你为什么迫害我的圣徒，也不是我的门徒，而是\Quote{你为什么迫害我？}\BibleRef{宗 9:4} 正如那时，当他坐在天上时没有人迫害他，他却呼喊说：\Quote{你为什么迫害我？}因为头认出了它的肢体，他的爱不允许头将自己与身体的结合分离；同样，当他不收回对他的慈爱时，那必然是他不收回对我们这些肢体和身体的慈爱。然而，我们不应因此无忧无虑地犯罪，并乖谬地保证自己不会灭亡，无论我们的行为如何。因为有些罪过和过犯，要界定并讨论它们，或者对我来说不可能，或者即使可能，在我们目前的时间内也太冗长。因为没有人能说自己无罪；因为如果他这样说，他是在说谎；\Quote{如果我们说我们没有罪，我们就是欺骗自己，真理也不在我们内。}\BibleRef{若一 1:8} 因此，每个人都必要因自己的罪受责罚；但如果他是个基督徒，天主的慈爱不会从他身上收回。当然，如果你犯了这样的过犯，以致拒绝管教者的手、责打你的杖，并对天主的纠正生气，当他责打你时逃避你的父亲，并且不容许他作你的父亲，因为当你犯罪时他不饶恕你；你就使自己与你的产业疏远了，他没有抛弃你；因为如果你接受责打，你就不会被剥夺继承权。\Quote{我也不会在我的信实中造成伤害。}因为他在释放时的慈爱不会被收回，免得他在报复时的信实造成伤害。\par

\Quote{我决不亵渎我的盟约，也不放弃我口中所出的话}\BibleRef{咏 88:34} 因为他的儿子们犯罪，我也不会因此被证明是虚谎的：我应许了，我必成就。假设他们选择继续犯罪，甚至毫无希望，以至于陷入罪恶，冒犯他们父亲的慈颜，并应被剥夺继承权；那不仍然是天主自己吗？经上论他说：\Quote{从这些石头}他\Quote{给亚巴郎兴起子孙来}？\BibleRef{玛 3:9} 因此，弟兄们，我告诉你们：许多基督徒犯小罪，许多人因罪受鞭打、受管教、受惩戒、得医治；许多人完全背离，硬着颈项抗拒父的管教，甚至完全拒绝天主作他们的父，虽然他们有基督的标记，因而陷入这样的罪中，以致只能宣告他们\Quote{行这样事的人不能承受天主的国。}\BibleRef{迦 5:21} 然而，基督不会因他们而失去继承产业：不会因糠秕的缘故，麦子也一起灭亡：\BibleRef{玛 3:12} 也不会因坏鱼的缘故，从那个网中就没有什么被收进器皿里。\BibleRef{玛 13:47} \Quote{主认识那些属于他的人}\BibleRef{弟后 2:19} 因为那在我们出生以前就预定我们的，毫无疑问地应许说：\Quote{凡他所预定的，他也召叫了他们；凡他所召叫的，他也使他们成义；凡他所成义的，他也使他们得荣耀。}\BibleRef{罗 8:29} 让绝望的罪人们随意犯罪吧！让基督的肢体回答说：\Quote{如果天主与我们同在，谁能反对我们呢？}因此，天主不会在他的信实中造成伤害，也不会\Quote{亵渎他的盟约。}他的盟约坚定不移，因为在他的预知中，他预定了他的继承者们；并且\Quote{他不会放弃他口中所出的话。}\par

32. 请你倾听，好使你在望德中得到确证，使你在安全中得享平安——若你自知属于基督的肢体。\BibleQuote{我凭我的圣洁一次起誓：我决不向达味说谎} \BibleRef{咏 88:35}。难道你要等天主第二次起誓吗？若祂一次起誓便失信，祂要起誓多少次呢？祂为了我们的生命起了一次誓，派遣祂的独生子为我们而死。\BibleQuote{我凭我的圣洁一次起誓，我决不向达味说谎。} \BibleQuote{他的后裔将永存不变} \BibleRef{咏 88:36}。他的后裔永存不变，因为主认识属于祂的人。\BibleQuote{他的宝座在我面前有如太阳，} \BibleQuote{又如永存不变的月亮；那天上的忠信见证} \BibleRef{咏 88:37}。他们是祂的宝座，祂在其上坐而为王。但如果他们是祂的宝座，那么也是祂的肢体；因为甚至我们的肢体也是我们头脑的座位。请看，我们所有的其他肢体如何支撑我们的头：但头并不支撑任何高于自身的东西，反而由我们其余肢体支撑，就如整个人体是头脑的座位。因此，祂的宝座——即所有天主在其上为王的人——\BibleQuote{在我面前将如同太阳}，祂说：因为义人在我父的国里\BibleQuote{将发光如太阳}。\BibleRef{玛 13:43}但这里的太阳是指属灵的太阳，而非肉体的太阳，即那从天上照耀、祂使之上升照善人亦照恶人的太阳。\BibleRef{玛 5:45}最后，那个太阳不仅在人眼前，也出现在牲畜和最小的昆虫眼前；因为最卑贱的动物中哪一个看不见那太阳？祂说什么来区分此处所指的太阳呢？\BibleQuote{在我面前如同太阳。}不是在人面前，不是在肉体面前，不是在必死的动物面前，而是\BibleQuote{在我面前，又如同月亮。}但什么样的月亮？是\BibleQuote{永存不变的月亮。}因为我们所认识的月亮，虽然变得圆满，但满月后的第二天便开始亏损。\BibleQuote{祂将如同永存不变的月亮}，祂说。祂的宝座将如月亮一样完美，但那是永存不变的月亮。如果如同太阳，为什么也如同月亮？圣经通常用月亮象征这肉身的可朽性，因其盈亏变化，因其短暂性。月亮也被解释为\Place{耶里哥}：一个人从\Place{耶路撒冷}下到\Place{耶里哥}，落在强盗手中： \BibleRef{路 10:30}因为他从不朽堕入可朽。因此，肉身类似于月亮，每月经历盈亏；但我们的肉身将在复活中臻于完美：\BibleQuote{并成为天上的忠信见证。}因此，如果只有我们的心智得以完美，祂只会把我们比作太阳；如果只有我们的身体，则比作月亮；但既然天主将使我们身心皆完美，就心智而言，经上说\BibleQuote{在我面前如同太阳}，因为唯有天主看见心智；而\BibleQuote{如同月亮}，则指肉身，这肉身\BibleQuote{在死者复活中将永存不变，}并成为\BibleQuote{天上的忠信见证}，因为凡关于死者复活所宣告的都是真实的。我恳求你们再更清楚地听一次并记住它：因为我知道有些人理解，有些人或许还在探究我的意思。基督徒信仰中，没有任何一条教义像肉身复活那样遭遇到如此多的反对。最后，祂——因被反对而成为标记的那位——死后重新取回祂的肉身，以面对诘难者；祂本可完全治愈祂的伤口，使伤疤完全消失，却将那些伤疤留在祂的身体上，以治愈心中怀疑的伤口。的确，在基督徒信仰中，没有什么像肉身复活那样遭到如此顽固、好辩的反对。关于灵魂不朽，许多外邦哲学家曾长篇大论地争论，并在许多著作中留下记载说灵魂是不朽的；但当他们论及肉身复活时，他们确实不怀疑，却最公开地否认，宣称这属土的肉身升天绝对不可能。因此，那月亮将永存不变，并成为天上的忠信见证，反对所有反驳者。\par

33. 这些应许是关乎基督的，如此确定、如此坚定、如此显明、如此无可置疑。因为虽有一些神秘地隐晦，但另一些却如此清楚，以至于一切隐晦之处都能轻易地藉着它们显明。既然如此，请看以下经文：\BibleQuote{但是祢已经弃绝、羞辱、厌弃了祢的受傅者}\BibleRef{咏88:38}。\BibleQuote{祢推翻了祢仆人的盟约，将祂的圣所玷污在地上}\BibleRef{咏88:39}。\BibleQuote{祢拆毁了他所有的篱笆，使他的堡垒成为恐怖}\BibleRef{咏88:40}……这是怎么回事呢？祢应许了这一切，却使它们反其道而行。那些不久以前还使我们满心喜悦的应许如今在哪里呢？我们曾何等欢愉地喝彩，何等无畏地夸耀！这如同一个人应许，另一个人毁坏。而这就是奥秘：因为经文不是说\Quote{另一个人}，而是说\Quote{祢}，那应许者，甚至屈尊俯就人的疑惑而起誓的祢。祢如此应许，又如此行事！我从哪里得到祢的誓言？哪里找到祢应许的实现？难道天主应许，或如此起誓，竟是虚假的吗？然而，这些应许和这些作为又是为何呢？我回答说，祂这样做正是为了成就那些应许。但我是谁，竟敢说这话？因此，让我们看看这是否是真理之言；那么我所说的便不是毫无根据的了。\Person{达味}就是那位，这应许要在他后裔，即\Person{基督}身上实现，是向\Person{达味}应许的；既然这些应许是对\Person{达味}说的，人们便期待它们在\Person{达味}身上应验。此外，以免任何基督徒声称这些应许是指着\Person{基督}说的，另一个人却将它们应用于\Person{达味}，因为他描述所有这些应许都在\Person{达味}身上应验，从而可能出错；祂便在\Person{达味}身上取消了它们，从而迫使我们在看到它们未在\Person{达味}身上应验时，转向别处寻求它们的应验。同样，在\Person{厄撒乌}和\Person{雅各伯}的案例中，我们看到年长的被年幼的朝拜，尽管经上写着：\BibleQuote{年长的要服事年幼的}\BibleRef{创25:23}；所以当你看到这在两兄弟身上未应验时，你便寻求两个民族，在其中发现天主按其真理屈尊应许之事如何完成。\BibleQuote{从你身所生的}，主对\Person{达味}说，\BibleQuote{我将放在你的海上}。祂应许从\Person{达味}后裔中要生出永远的事物；而\Person{撒罗满}作为他的儿子出生，获得了如此大的智慧，以至人们相信天主关于\Person{达味}身所生的应许在他身上实现了；但\Person{撒罗满}堕落了，为盼望\Person{基督}留下了余地；因为天主既不能被欺骗，也不欺骗人，祂必不会让祂的应许停留在一个祂知道将要堕落的人身上，而是使你可以在\Person{撒罗满}堕落后回顾天主，并索求祂的应许。主啊，难道祢欺骗了？祢没有实现祢的应许吗？祢不展示祢所起誓的吗？或许天主会回答说：我起誓并应许了；但是\Person{撒罗满}不会坚持到底。那么如何？主天主，祢不是预先知道他不会坚持到底吗？祢确实知道。那么祢为何应许我，赐予那不会坚持到底的人以永恒的事物呢？祢不是已经回答说：\BibleQuote{但是如果他的子女离弃我的法律，不遵行我的典章；如果他们不遵守我的律例，亵渎我的盟约}，然而我的应许仍然有效，我的誓言必将实现：\BibleQuote{我一次凭我的圣洁起了誓}，在内心，在某种奥秘中，在先知们饮用的泉源里，他们从那里向我们涌出这些事，\BibleQuote{我一次起了誓}，我必不辜负\Person{达味}。那么，请显明祢所起誓的，赐给我们祢所应许的。应验从那一个\Person{达味}身上被取走，为的是不让人在那一个\Person{达味}身上寻求；因此，等候我所应许的。\par

34. 连\Person{达味}自己也知道这事。请听他的话：\BibleQuote{祢已经弃绝他，使他归于无有}。祢的应许又在哪里呢？\BibleQuote{祢延迟了祢的受傅者}。这句话在众多悲伤之中令我们欣慰：因为天主的应许仍然有效；祢延迟了祢的受傅者，并没有除掉祂。请看那无知者所寄望于从\Person{达味}身上应验天主应许的\Person{达味}的命运，为的是使那些应许更可靠地在他人的应验中被信赖。\BibleQuote{祢延迟了祢的受傅者：祢推翻了祢仆人的盟约}。因为犹太人的旧约在哪里呢？那应许之地在哪里？他们曾居住其中并犯罪，在其倾覆后他们流落远方。你求问犹太人的王国，它不存在；你求问犹太人的祭台，它不存在；你求问犹太人的祭献，它不存在；你求问犹太人的司祭职，它不存在。\BibleQuote{祢推翻了祢仆人的盟约，将他的圣所玷污在地上}。祢显明他们以为神圣的，原是属地的。\BibleQuote{祢拆毁了他所有的篱笆}，祢曾用这些篱笆围护他：因为若不是篱笆被拆毁，他怎能被掠夺呢？\BibleQuote{祢使他的堡垒成为恐怖}。为什么恐怖？为的是能对罪人说：\BibleQuote{因为天主既然没有怜惜了原来的枝条，你也当小心，免得祂也不怜惜你}\BibleRef{罗11:21}。\par

\BibleQuote{行路的人都掠夺了他：}即所有行路的外邦人，意即所有经过此生的人，都掠夺了以色列，掠夺了达味。首先，看他在万民中的碎片：因为关于犹太人，经上记载：\BibleQuote{他们必成为狐狸的份。}因为圣经称邪恶、狡猾、懦弱的君王（他人的勇气令其恐惧）为狐狸。因此，当我们主亲自论及恐吓的黑落德时，祂说：\BibleQuote{你们去告诉那只狐狸。}\BibleRef{路 13:32}那无所畏惧的君王不是狐狸，而是犹大的狮子，经上论及祂说：\BibleQuote{你俯伏，又起来，如狮子睡卧。}\BibleRef{创 49:9}祢按祢的旨意俯伏，按祢的旨意起来；因为祢愿意，祢便睡了。因此，在另一篇圣咏中祂说：\BibleQuote{我睡了。}那句子不完整吗：\BibleQuote{我睡了，且安息，又再起来，因为主必扶持我}？为什么加上\Quote{我}这个词？因此，当我强烈强调\Quote{我}这个词时，他们对我发怒，他们扰乱我；但若我不愿意，我就不睡。那么，那些被预言要成为狐狸之份的人，现在被这样说：\BibleQuote{行路的人都掠夺了他；他成了邻舍的羞辱。}\BibleRef{咏 88:41}\BibleQuote{祢高举了他敌人的右手，使他的所有仇敌都欢乐。}\BibleRef{咏 88:42}看看犹太人，便看见所有预言的事都成就了。\BibleQuote{祢收回了他的刀剑之助。}他们过去如何以少数人作战，击倒许多人。\BibleQuote{祢收回了他的刀剑之助，在战争中祢不赐他胜利。}\BibleRef{咏 88:43}所以他自然被征服，自然被俘，自然被逐出他的王国，自然四散；因为他失去了那地为己有，他曾为此杀害了主。\BibleQuote{祢使他不再受洁净。}\BibleRef{咏 88:44}这是什么？在一切恶事中，这是极可畏的事；因为无论天主如何打击，如何发怒，如何鞭打责罚，但愿祂鞭打那被捆绑的、祂要洁净的人，而不是\BibleQuote{使他不再受洁净。}因为若祂使他不再受洁净，他便不能受洁净，必被驱逐。那么，犹太人不再受什么洁净？从信德；因为我们因信德生活；\BibleRef{迦 3:11}论到信德，经上说：\BibleQuote{因信德洁净他们的心。}\BibleRef{宗 15:9}而且只有基督的信德能洁净；因不信基督，他们不再受洁净。\BibleQuote{祢使他不再受洁净，并将他的宝座推倒在地。}所以祢把它打碎了。\BibleQuote{祢缩短了他宝座的日子。}\BibleRef{咏 88:45}他们以为会永远统治。\BibleQuote{祢以羞辱遮盖他。}这一切都发生在犹太人身上，基督并没有被夺走，只是祂的来临被延迟了。\par

35. 因此，让我们看看天主是否履行祂的应许。在这些严厉的惩罚之后（这些惩罚已被记载为施加于这个民族和王国），为了使天主不被认为在其中履行了祂的应许，从而不另赐一个在基督内的王国（这王国将永无穷尽）；先知这样对祂说：\BibleQuote{主，祢要隐藏到几时直到永远？}\BibleRef{咏 88:46} 因为可能不是向他们隐藏直到永远；因为\BibleQuote{以色列人一部分的硬心，直到外邦人的数目全满了，这样，全以色列都要得救。}\BibleRef{罗 11:25} 但在此期间，\BibleQuote{祢的忿怒如火焚烧。}\par

36. \BibleQuote{求祢记念我的本体是什么}\BibleRef{咏 88:47}。那位达味，按肉身安置在犹太人中间，按希望安置在基督内的，说：\BibleQuote{记念我的本体是什么。}因为并非因犹太人堕落，我的本体就失败；因为从那个民族来了童贞玛利亚，从她来了基督的肉身；那肉身不犯罪，反而洁净罪恶；在那里，达味说，是我的本体。\BibleQuote{求祢记念我的本体是什么。}因为根并未完全枯萎；那蒙应许的后裔必要来到，是由天使借中保之手所规定的。\BibleRef{迦 3:19}\BibleQuote{祢没有使世人全是虚无}\BibleRef{咏 88:47}。看！所有世人都归于虚无；然而祢没有使他们成为虚无。如果所有归于虚无的，都是祢没有使之成为虚无的，难道祢没有保留一些工具来净化他们的虚无吗？这个祢为自己保留以洁净人脱离虚无的，就是祢的圣者，我的本体在祂内；因为从祂而来的人，都是祢没有使之成为虚无的，且从他们自己的虚无中被洁净。对他们说：\BibleQuote{人子们哪，你们的心几时沉重？你们为何喜爱虚无，寻求虚谎？}\BibleRef{咏 4:2}也许他们会焦虑，并转离他们的虚无，当他们发现自己被虚无玷污时，会寻求洁净：那么帮助她们，使他们稳妥。\BibleQuote{你们也要知道，主使祂的圣者成为奇妙。}祂使祂的圣者受人钦佩：从此祂洁净了所有人的虚无：在那里，达味说，是我的本体：求祢记念它！\BibleQuote{因为祢没有使世人全是虚无。}所以祢有所保留来洁净他们：祢所保留的是谁？\BibleQuote{有哪个人活着，而不见死亡？}那活着而不见死亡的人，将净化他们脱离虚无。因为祂没有使所有人成为虚无，也不能因祂如此创造而轻视自己的受造物，而不去转化和洁净他们。\par

37. \BibleQuote{什么人生存而看不到死亡？}\BibleRef{咏 88:48}。因为祂从死者中复活后，不再死亡，死亡不再统治祂。\BibleRef{罗 6:9}正如在另一篇圣咏中所说：\BibleQuote{祢决不会将我的灵魂遗弃在地狱，也不让祢的圣者见到腐朽。}宗徒的教导承接了这一证言，并在宗徒大事录中\BibleRef{宗 2:29}如此反驳不信者：诸位弟兄，我们知道圣祖达味死了、埋葬了，他的肉身也见了腐朽。因此不能指着他说：\BibleQuote{祢也不让祢的圣者见到腐朽。}那么这话是指谁说的呢？\BibleQuote{什么人生存而看不到死亡？}也许没有这样的人。不然，但\BibleQuote{是谁？}这句话是要让你探求，而非绝望。但或许有一个人\BibleQuote{生存而看不到死亡}，然而或许他讲的不是那死了的基督？没有一个人\BibleQuote{生存而看不到死亡}，除了那为凡人而死的基督。为使你确信这话是指着祂说的，请看下文：\BibleQuote{什么人生存而看不到死亡？}难道祂从未死过吗？祂死过。那么祂如何生存而永不见死亡？\BibleQuote{祂要从地狱的手中将祂自己的灵魂救出。}的确，唯独祂被论及，因祂在众人中唯一\BibleQuote{生存而看不到死亡：祂要从地狱的手中将祂自己的灵魂救出}，因为尽管其他信徒将从死者中复活，他们自己也将永世生存，不再见死亡，但他们不能自己将灵魂从地狱手中救出。那将自己的灵魂从地狱手中救出的，也亲自拯救祂信徒的灵魂；他们无法自救。证明祂拯救自己的灵魂：\BibleQuote{我有权舍去我的生命，也有权再取回它。没有人夺去它从我这里；因为我亲自睡眠，但由我自己舍下，再取回。}因为就是祂亲自将自己的灵魂从地狱手中救出。\par

38. 但在对基督的信德中发生了巨大的困难，外邦人在他们的愤恨中长期说：\Quote{他何时死亡，他的名字何时消灭？}为了那些早已相信基督、但注定要怀疑一段时间的人，下面的话接续：\BibleQuote{主，祢从前的慈爱在哪里？}\BibleRef{咏 88:49}。我们现在已承认基督是我们的净化者，我们现在拥有了祂，祢的应许要在祂身上实现；求祢在祂身上显示祢所应许的。正是祂自己要生存，不见死亡：正是祂将自己的灵魂从地狱手中救出：然而我们仍在受苦。殉道者们就是这样说的，我们正在庆祝他们的诞辰。祂要生存，不见死亡：祂将自己的灵魂从地狱手中救出：然而\BibleQuote{为了祢的缘故，我们终日被杀；被视为待宰的羊群。}\BibleQuote{主，祢从前的慈爱在哪里，就是祢凭祢的真理向达味所起的誓？}\par

39. \BibleQuote{主，求祢记念祢仆人所受的羞辱}\BibleRef{咏 88:50}。即使当基督活着，并坐在圣父右边时，辱骂仍落在基督徒身上：他们长久因基督的名受辱骂。那位生育的寡妇，她的子女比有夫之妇的子女更多，听到了恶名，听到了辱骂；但教会，她如此繁衍，向右向左扩展，不再记念她寡居的羞辱。\BibleQuote{主，求祢记念}；在记念祂中有丰盛的甘甜。\BibleQuote{记念}，不要忘记。记念什么？\BibleQuote{祢仆人所受的羞辱：以及我如何将许多人的羞辱怀揣在怀中。}他说：我去传讲祢，听到了辱骂，并将它们怀揣在怀中，因为我在实践预言。\BibleQuote{被人毁谤，我们就劝勉；我们成了世上污秽，万物中的渣滓，直到今日。}\BibleRef{格前 4:13}长久以来，基督徒将辱骂怀揣在他们怀中，在心中：不敢对抗辱骂者；从前，回答外邦人算是犯罪：现在，仍作外邦人则是犯罪。感谢归于主！祂记念了我们的辱骂：祂高举了祂受傅者的角，使祂成为地上君王中的奇妙者。如今没有人侮辱基督徒，即使有人侮辱，也不在公开场合：他说话时，仿佛更怕被人听见，而非急于被人相信。\BibleQuote{我将许多人的羞辱怀揣在怀中。}\par

40. \BibleQuote{主，祢的敌人以何处亵渎了祢}\BibleRef{咏 88:51}，无论是犹太人还是外邦人。\BibleQuote{他们以何处亵渎了祢。}他们以何处亵渎了祢？\BibleQuote{以祢受傅者的改变。}他们反对说基督死了，被钉十字架。疯子们，你们的辱骂是什么？虽然如今没有人再用它：但假设还有一些人如此说，你们的辱骂是什么？基督死了？祂没有被消灭，而是改变了。因那三日，祂被称为\Quote{死}。那么，祢的敌人用何处亵渎了祢？不是用祢受傅者的丧失，不是用祂的毁灭，而是用祂的\Quote{改变}。祂从暂时生命被改变为永恒生命；祂从犹太人被改变为外邦人；祂从地上被改变到天上。让祢虚妄的敌人仍因祢受傅者的改变而亵渎祢吧。但愿他们能被改变：那样他们就不会亵渎基督的改变，因他们自己不改变，故对改变不悦。\BibleQuote{在他们当中没有改变，他们也不敬畏天主。}\par

41. 他们亵渎了基督的改变；但你如何回答？\Quote{主的祝福，直到永远。阿们，阿们。} \BibleRef{咏 88:52}。感谢祂的仁慈，感谢祂的恩宠。我们表达感谢：我们不是给予、也不是归还、更不是回报感谢：我们用言语表达感谢，而实际上我们内心保留着感恩之情。祂白白拯救了我们，未曾顾念我们的不虔；我们未寻找祂时，祂已寻找我们；祂找到了我们，救赎了我们，从魔鬼的奴役和其邪恶天使的权势下解放了我们；祂吸引我们归向祂，藉着那信德净化我们，只有那些不信的人，祂才将他们从中释放出来，而他们正因为不信而不能被净化。让那些仍然不信的人每天随意说什么吧；他们存留的人将一天比一天少；让他们辱骂、嘲笑、指控——不是基督的死亡，而是基督的改变。难道他们没有看到，当他们说这些话时，他们要么因相信而失败，要么因死亡而失败？因为他们的诅咒是暂时的；而主的祝福却\Quote{直到永远。}为了确认这祝福，还加上了\Quote{阿们，阿们。}这是天主契约的印章。因此，既已确信祂的许诺，让我们相信过去，认出当下，盼望未来。不要让敌人引诱我们偏离正路，愿那如同母鸡聚集小鸡在翅膀下的那一位养育我们；免得我们离开祂的翅膀，尚未长成羽毛就被空中的鹰叼走。因为基督徒不应在自己身上寄望：他若希望强壮，就当在母亲的温暖中成长。这就是那聚集小鸡的母鸡；为此我们的救主曾责备不信的耶路撒冷：\Quote{看，你们的房屋必给你们留下一片荒凉。} \BibleRef{玛 23:38}。因此经上说：\Quote{你使他的堡垒成了恐怖。}既然他们不愿被聚集在这母鸡的翅膀下，并给了我们一个警告，教导我们畏惧那空中飞翔、每日找寻可吞噬之物的不洁之神；让我们聚集在这母鸡的翅膀下，即天主的智慧，因为祂甚至为小鸡而虚弱至死。让我们爱我们的主天主，让我们爱祂的教会：爱祂如父，爱她如母；爱祂如主，爱她如婢女——正如我们自己是婢女的儿子。但这婚姻是由伟大的爱情结合在一起的：无人能冒犯了一方，而赢得另一方的欢心。没有人可以说：\Quote{我确实去拜偶像，我求问附身者和算命人，但我没有背弃天主的教会；我是天主教徒。}当你依附你的母亲时，你已冒犯了你的父亲。另一个人说：\Quote{我绝不如此；我不求问巫术，不寻找附身者，绝不借渎神的占卜求问，我不去拜偶像，不在石头前叩首；尽管我属于多纳徒派。}你虽未冒犯你的父亲，但若他惩罚你冒犯母亲的行为，于你有何益处？你承认主、尊崇天主、宣讲祂的名、承认祂的儿子、宣认祂坐在祂的右边，却亵渎祂的教会，这于你有何用？难道人间的婚姻类比不能说服你吗？假设你有一位庇护人，你每日向他献殷勤，踏破他的门槛，每日不仅问候，甚至敬拜他，向他致以最忠诚的礼遇；若你中伤他的妻子一句，你还能再进他的家门吗？所以，最亲爱的，请同心合意地持守天主父和我们的母亲教会。以节制庆祝圣人们的诞辰，以便我们效法那些先我们而去的人，并且那些为你祈祷的人能因你而喜乐；如此，\Quote{主的祝福必常与你同在，直到永远。阿们，阿们。}\par

\SourceRecord{Translated by J.E. Tweed.；Nicene and Post-Nicene Fathers, First Series；Vol. 8.；Edited by Philip Schaff.；Buffalo, NY: Christian Literature Publishing Co.,；1888.；<http://www.newadvent.org/fathers/1801089.htm>.}

% structure: p0001 source=h1 target=section reason=page-title

\section{论第九十篇圣咏}

1. 这篇圣咏的标题是\Quote{天主的人梅瑟的祈祷}。天主藉着祂的这个人——梅瑟，将法律赐给祂的子民，藉着他将他们从为奴之家解救出来，并领他们在旷野中四十年。所以梅瑟是旧约的仆役，新约的先知。正如宗徒所说：\BibleQuote{这一切事发生在他们身上，是为作预像；并写下来，为警戒我们这些活在末世的人。} \BibleRef{格前 10:11} 因此，照此赐予梅瑟的神恩安排，这篇圣咏应被检视，因它从其祈祷得名。\par

2. 他说：\BibleQuote{主，从这一代到那一代，祢是我们的避难所。} \BibleRef{咏 89:1} 意思是：或在每一代，或在两代，即旧的一代和新的一代。因为，如我所说，他（梅瑟）是与旧的一代有关的旧约的仆役，又是与新的一代有关的新约的先知。耶稣基督，那盟约的担保者，以及在那一代进入婚姻的新郎，说：\BibleQuote{如果你们相信梅瑟，也必会相信我，因为他所写的，就是指着我。} \BibleRef{若 5:46} 现在，不可相信这篇圣咏完全是那位梅瑟所写，因为它没有以他诗歌中使用的任何表达为特征；但是，这位伟大的天主仆人的名字被用来作某种指示，以引导读者或听者的注意力。他说，\BibleQuote{主，祢是我们的避难所，从一代到另一代。}\par

3. 他接着说明祂如何成为我们的避难所，因为祂开始成为我们的避难所，这是祂以前所没有的；并不是说祂在成为我们的避难所之前不存在：\BibleQuote{山岳尚未生出，大地和世界尚未造成，从亘古直到永远，祢是天主。} \BibleRef{咏 89:2} 所以，祢是永远的，在我们存在之前，在世界存在之前，自从我们转向祢，祢就成了我们的避难所。但是，\Quote{山岳尚未生出}等语，在我看来，含有特殊的意义；因为山岳是地上的高处，如果天主甚至在土地形成之前就已存在（或者，如一些书卷所说，从同一个希腊词，即\Quote{造成}），既然土地是由祂形成的，那么祂说祂在众山之前，或在土地的某些部分之前，有何必要呢？因为天主不仅是在土地之前，而且是在天地之前，甚至在整个有形和无形的受造物之前。但很可能是，整个理性受造物由这种区别标示：山岳象征天使的高贵，而土地象征人的卑微。因此，虽然一切受造之物都无不可称为被\Quote{造}或被\Quote{形成}；但若这些词语有恰当的含义，则天使是\Quote{被造的}；因为当它们被列在祂天上的工作中时，总结本身是这样说的：\Quote{祂说了话，它们就造成了；祂命令了，它们就创造了。} 而土地是\Quote{被形成的}，为使人能从中在身体上被创造。因为圣经使用这个词，我们读到：天主创造了，或\BibleQuote{天主用地上的尘土形成了人。} \BibleRef{创 2:7} 所以，在受造物中最崇高的部分（因为有什么比天上受造物中的理性部分更高呢）被造成之前，在土地被造成之前，为的是祢能在世上拥有敬拜祢的人；而这还是小事，因为所有这些都在时间中或在时间的同时开始；但\BibleQuote{从亘古直到永远，祢是天主。} 若说\Quote{从永远到永远}，就更好；因为天主，在万世之前，祂的存在不是从某个世代开始，也非到某个世代终止，这世代有终，而祂无终。但在圣经中常发生，因希腊词歧义，导致拉丁译者将世代译作永恒，将永恒译作世代。但他（圣咏作者）非常正确地没有说：\Quote{祢从世代而来，祢要到世代而去}，而是使用现在时动词，暗示天主的本质是完全不变的。祂不是\Quote{曾有}和\Quote{将有}，而只是\Quote{是}。所以有\Emph{我是自有者}这句话；以及\BibleQuote{\Emph{我是}派遣我到你们这里来。} \BibleRef{出 3:14} 还有：\BibleQuote{祢要将它们改变，它们就被改变；但是祢是同一的，祢的岁月不会穷尽。} 看，这就是那作为我们避难所的永恒，好让我们从时间的可变性逃往那里，永远安居。\par

4. 但是，既然我们在此世的生活面临许多巨大的诱惑，且恐怕我们可能因此被引离那避难所，就让我们看看，天主的人的祈祷随之寻求什么。\BibleQuote{求祢不要将人转向卑贱。} \BibleRef{咏 89:3} 意思是：不要让人从祢永恒崇高的事物转向，贪恋时间的事物，品味世俗的事物。这个祈祷正是天主自己在天主经中命令我们祈求的：\BibleQuote{不要让我们陷于诱惑。} \BibleRef{玛 6:13} 他接着说：\Quote{祢又说：回来吧，你们这些人类之子。} 仿佛他说：我向祢求祢命令我所求的：将荣耀归于祂的恩宠，当照经上所记：\BibleQuote{凡要夸耀的，应当夸耀主。} \BibleRef{格前 1:31} 没有祂的帮助，我们无法靠己意的努力克服此生的诱惑。\Quote{求祢不要将人转向卑贱；祢又说：回来吧，你们这些人类之子。} 但求祢赐予祢所命令的，俯听那至少会祈祷者的祈求，并帮助那愿者的信德。\par

5. \BibleQuote{因为在祢眼中，千年如已逝的昨日}\BibleRef{咏 89:4}：因此我们应当转向祢的避难所，祢是永恒不变的，离开我们周围流逝的景象；因为无论人对此生期望多长久，\BibleQuote{在祢眼中，千年如昨日}，而非如同明日，因为明日尚未来到：一切有限的时间都被视为已经过去。因此宗徒宁愿\BibleQuote{向着前面}\BibleRef{斐 3:13}，即渴慕永恒之事，忘记背后之事，应以此理解暂时的事物。但为避免有人以为天主将千年算作一日，仿佛在天主那里日子如此之长，而这样说其实只是为了贬低时间的长度；他便加上：\BibleQuote{又如夜间的一更}，那只有三个小时。然而竟有人胆敢声称自己知道时间，主对这些冒充者说：\BibleQuote{圣父以自己的权柄所定的时候、日期，不是你们可以知道的}\BibleRef{宗 1:7}。他们断言这段时间可以定为六千年，如同六日一样。他们也没有注意这句话：\BibleQuote{如同一日，已经过去}；因为说这话时，还不到一千年，而\BibleQuote{如同夜间一更}的说法，本应警告他们不要被季节的不确定性所欺骗：因为即使起初六日天主完成祂的工程似乎给他们的看法提供了某种合理性，但六更（共计十八小时）却与那种看法不符。\par

6. 天主的人，更确切地说是先知之神，似乎正在诵念一条写在天主隐秘智慧中的法律，天主藉此决定了有罪之人的生命期限，并定下了肉身的苦难，其言如下：\BibleQuote{他们的年岁如同虚无；在早晨，愿它如草枯萎}\BibleRef{咏 89:5}。因此，旧盟约的继承者所追求的福乐，他们曾向主、他们的天主祈求作为大恩，在祂神秘的眷顾中得以领受这条法律。梅瑟似乎正在诵念：\Quote{他们的年岁将被视为虚无。}这些事物就是这样：在它们来临之前不存在；一旦来临，也很快不复存在；因为它们来到世上不是为了存留，而是为了消逝。\BibleQuote{在早晨，}即在它们来临之前，\BibleQuote{像炎热一样让它过去；}但\BibleQuote{在傍晚，}意指它们来临之后，\BibleQuote{让它倒下，干涸，枯萎}\BibleRef{咏 89:6}。这就是在死亡中\BibleQuote{倒下}，在尸首中\BibleQuote{干涸}，在尘土中\BibleQuote{枯萎}。这不就是有血肉的身体吗？其中可诅咒的肉欲存于其中。\BibleQuote{凡有血肉的都如草，一切人的美丽都像田野的花；草枯花谢，但是主的言语永远常存}\BibleRef{依 40:6}。\par

7. 他毫不隐瞒这命运是因罪而受的惩罚，立刻加上：\BibleQuote{我们因祢的震怒而消耗，因祢的愠怒而惊慌}\BibleRef{咏 89:7}：我们在软弱中消耗，因怕死而惊慌；因为我们变得软弱，却又害怕结束这软弱。祂说：\BibleQuote{别人要给你束上腰，带你往你不愿意去的地方}\BibleRef{若 21:18}，虽然这不是要受惩罚，而是要通过殉道获得冠冕；主的灵魂将我们转化到祂内，忧伤至死：因为\Quote{主的出离}无非就是\Quote{死亡}。\par

8. \BibleQuote{祢将我们的过犯摆在祢面前}\BibleRef{咏 89:8}：也就是说，祢没有隐瞒祢的愤怒：\BibleQuote{并将我们的年岁放在祢面容的光中。}\BibleQuote{祢面容的光}与\BibleQuote{在祢面前}和\BibleQuote{我们的过犯}如前文一样相互对应。\par

9. \BibleQuote{我们的日子都消逝了，我们在祢的震怒中衰竭}\BibleRef{咏 89:9}。这些话充分证明我们受制于死亡是一种惩罚。他谈到我们的日子消逝，或因人们在其中因热爱流逝的事物而衰败，或因它们减少到如此之少；这正如他在以下经文中宣称的：\BibleQuote{我们的岁月好像蜘蛛般的思虑。}\BibleQuote{我们的年岁是七十岁，若强壮可到八十岁，但其中所矜夸的不过是劳苦愁烦}\BibleRef{咏 89:10}。这些话似乎表达了此生的短暂与痛苦：因为那些到达七十岁的人被称为老人。然而，到八十岁，他们似乎还有一些力量；但如果活过八十，他们的存在因多重愁苦而艰辛。但许多甚至未满七十岁的人却经历最衰残悲惨的老年；而有时老人即使超过八十岁也被发现精神矍铄。因此，最好在这些数字中寻找某种属灵的意义。因为天主的愤怒并未因亚当的罪而更大——藉着亚当一人，\BibleQuote{罪恶进入了世界，死亡藉着罪恶，于是死亡就蔓延到所有人身上}\BibleRef{罗 5:12}——因为他们活得远不如古人长久；即使他们日子的长度也被嗤笑，相比千年如已逝的昨日和三个小时，尤其在那他们激起天主愤怒降下洪水灭绝他们的时候，他们的寿命正是最长的。\par

10. 此外，七十岁与八十岁合计为一百五十；圣咏集明确暗示此数为一神圣数目。一百五十与十五具有同样的相对意义，后者由七与八合成：前者藉守安息日指向旧约；后者藉主的复活指向新约。故此\Place{圣殿}有十五级台阶。故此圣咏中有十五首\Quote{登阶歌}。洪水的水位超出最高的山岭十五肘：\BibleRef{创 7:20}以及其他许多同类事例。\Quote{我们的岁月如蜘蛛般在思虑中度过。}我们劳碌于可朽坏之物，编织可朽坏之工；正如\Person{依撒意亚}先知所说，它们丝毫不能遮蔽我们。\BibleRef{依 59:6}\Quote{我们一生的年日在于本身}等等。此处区分了\Quote{本身}和\Quote{力量}：\Quote{在于本身}，即在于年日本身，可能意指旧约所应许的暂时事物，由七十这个数字所象征；\Quote{但若}不在于本身，而在于\Quote{他们的力量}，则指非暂时事物，而是永恒事物，\Quote{八十岁}，因为新约包含永生的新生命与复活的希望；接着又说，若他们度过这后一段时期，\Quote{他们的力量便是劳苦与愁烦}，暗示那超出此信仰、寻求更多者之命运。也可如此理解：虽然我们已建立于新约——即八十所象征——但我们的生命仍是劳苦与愁烦，因为我们\Quote{在自己内叹息，等待义子之身，即我们身体的救赎；因为我们因希望而得救；若我们希望那未见之事，则我们耐心等候。}\BibleRef{罗 8:23}这与天主的仁慈有关，他接着说：\Quote{因你的仁慈临于我们，我们必受管教}；因为\Quote{主管教他所爱的，鞭打他所接纳的每个儿子}\BibleRef{希 12:6}，又给一些大能者一根刺在肉中，击打他们，使他们不致因启示的伟大而过于高傲，好使软弱中成全能力。\BibleRef{格后 12:7}有些抄本作\Quote{受教导}而非\Quote{受管教}，这同样表达天主的仁慈；因为无人能不受劳苦愁烦而受教导，因为能力在软弱中得成全。\par

11. \Quote{因为谁能知道你怒气的威力？谁能因敬畏你而数算你的忿怒？}\BibleRef{咏 89:11}。他说，很少有人能知道你怒气的威力；因为当你宽恕时，你的怒气对大多数人反而更重；好使我们知道劳苦与愁烦不属于怒气，而属于你的仁慈，当你管教并教导你所爱的人，救他们脱离永罚的痛苦：正如另一篇圣咏所说：\Quote{罪人激怒了上主：他必不按他怒气的威势向他追讨。}与这句话一起理解的还有：\Quote{谁知道？}寻找一个知道如何因敬畏你而数算你忿怒的人如此困难，以致他加上这句话，意思是：你似乎宽恕一些人，但你对他们更为愤怒，为使罪人在其道路上亨通，最终受到更重的刑罚。因为人怒气的威力杀死身体后，便无能为力；但天主有能力在此惩罚，并在身体死亡后将人投入地狱；借着那些如此被教导的少数人，恶人虚妄而诱人的亨通被断定是天主更大的忿怒……\par

12. \Quote{显扬你的右手}\BibleRef{咏 89:12}。这是大多数希腊抄本的读法；而有些拉丁抄本则作\Quote{把你的右手显给我看。}什么是\Quote{你的右手}？不就是你的基督吗？关于他，经上说：\Quote{上主的臂膀向谁显示了呢？}\BibleRef{依 53:1}求你如此显扬他，使你的忠信者在他内学会向你祈求并盼望那些作为他们信德奖赏的事物，这些事物不在旧约中出现，而在新约中启示出来：使他们不致认为从尘世与短暂福乐而来的快乐应受高度重视、渴求或喜爱，以致当他们看见不尊敬你的人竟有这些时，他们的脚滑倒；使他们的步履不致动摇，因为他们不知道如何数算你的忿怒。最后，依照这属于他的人的祈祷，基督已被如此显扬，以致藉由他的苦难表明：所应渴求的并非旧约中看似极为宝贵的那些赏报——它们只是未来之事的影像——而是永恒的事物。天主的右手也可理解为他将用这只手把圣徒与恶人分开；因为当这只手鞭打他所接纳的每个儿子，不容他在更大的怒气中在罪中亨通，反而仁慈地用左手鞭打他，好使他在净化后被置于右手时，这只手就变得为人所知。\BibleRef{玛 25:33}大多数抄本的读法\Quote{把你的右手显给我看}可指向基督，也可指向永恒的福乐：因为天主并没有形体的右手，正如他也没有激起强烈情感的忿怒。\par

13. 但他所添加的，\Quote{那些在智慧中心受束缚的}；其他抄本读作\Quote{受教的}，而非\Quote{受字句的}：希腊动词表达两种含义，只差一个音节。然而，正如经上所说，这些人也将自己的\Quote{脚放在智慧的桎梏}中，受教于智慧（他指的是心灵的脚，而非身体的），并被其金链束缚\BibleRef{德 6:24}，不偏离天主的道路，也不成为祂的逃兵；无论我们采用哪种读法，含义中的真理都是安全的。天主使这些在智慧中心受字句或受教的人，在新约中如此显扬，以致他们为信德轻视了犹太人和外邦人的不虔所憎恶的一切；并容许自己被剥夺那些在旧约中被按肉体判断者视为崇高应许的事物。\par

14. 当他们如此显扬，以致轻视这些事物，并因向往永恒之物而通过苦难作证（故在希腊语中被称为见证人或殉道者）时，他们长期忍受了许多痛苦的暂时磨难。这位天主之人留意到这一点，以梅瑟之名继续的先知之神如此说：\Quote{上主，求祢回来，要到何时？求祢对祢的仆人们心软}\BibleRef{咏 89:13}。这是那些人在那逼迫的时代忍受许多苦难，因心被智慧的锁链牢牢束缚，以致连这样的艰难也不能诱使他们逃离他们的主转向今世的美好事物，所说的话。另一篇圣咏中与这句\Quote{上主，求祢回来，要到何时？}一致的话是：\Quote{上主，祢要掩面不顾我到何时？}而那些以最肉性的精神将人身体的形式归于天主的人，要知道祂面容的\Quote{转开}和\Quote{回转}并不像我们身体的动作那样，让他们回想同篇圣咏上文的话：\Quote{祢将我们的过犯摆在祢面前，将我们的隐恶放在祢面容的光中。}那么，他如何在此处说\Quote{回来}，使天主垂顾，好像祂在愤怒中转开了脸？而前者却以这样的方式谈到天主的愤怒，暗示祂并没有转开祂的面容不看祂所恼怒之人的过犯和生活，反而将他们摆在祂面前和在祂面容的光中？\Quote{要到何时}属于祈求正义，而非愤慨的不耐烦。\Quote{心软}，有些人用动词\Quote{软化}来翻译。但\Quote{心软}避免了歧义；因为软化是一个常用动词：因为可以说他软化，即倾吐祈祷者，也可以说他被软化，即承受祈祷者：因为我们说：我软化你，以及我向你软化。\par

15. 接着，预见到未来的祝福，他如同已蒙恩赐般说道：\Quote{我们早晨因祢的仁慈而心满意足}\BibleRef{咏 89:14}。预言就这样在我们夜间的劳苦和悲伤中，如同黑暗中的灯，被点燃给我们，直到天破晓，晨星在我们心中升起。\BibleRef{伯后 1:19}因为心里洁净的人是有福的，他们要看见天主；那时义人将因他们现今所饥渴的福分而饱足\BibleRef{玛 5:8}，同时，他们凭着信德行走，是远离主的。\BibleRef{格后 5:6}因此有这话：\Quote{在祢面前有圆满的喜乐}；以及：\Quote{清晨他们必就近，且瞻仰}；正如其他译者所说：\Quote{我们早晨因祢的仁慈而心满意足}；那时他们将得满足。正如他别处所言：\Quote{当祢的荣耀显现时，我必心满意足。}因此说：\Quote{主，将父显示给我们，我们就满足了}；我们的主亲自回答：\Quote{我要将自己显示给熙雍}\BibleRef{若 14:8}；在这应许实现之前，没有祝福能使我们满足，或应当使我们满足，免得我们的渴望在半途止步，而当它们应当增强直到获得目标时却止步。\Quote{我们一生所有的日子都欢喜快乐。}那些日子是无尽的日子：它们共同存在：它们就是这样满足我们：因为它们不给予后继的日子：因为那里没有尚未存在的东西，因为它尚未到达我们，也没有因已过去而不再存在的东西；所有的一切都在一起：因为只有一天，它存留而不逝去：这就是永恒本身。这些日子就是经上所记载的：\Quote{谁愿生活，爱慕好日子？}这些日子在另一处被称为年：对此，对天主说：\Quote{唯有祢永存，祢的年数没有穷尽}：因为这些不是被视为虚无的年，或如影消逝的日子；而是实际存在的日子，那位如此说的人——\Quote{上主，求祢让我知道我的结局}（即，到达哪一终点后我将不再改变，不再需要任何祝福），\Quote{以及我日子的数目是什么}（是什么，而不是不是什么）——曾祈求知道其数。他将它们与今生的日子区分开来，关于后者他接着说：\Quote{看，祢使我的日子如同手掌之宽}，这些日子不存留，因为它们不站立，不存留，而是迅速更替：其中没有一刻我们的存在不是一部分已过去，另一部分将来临，没有一刻保持不变。但那些年日和日子，在其中我们也将永不匮乏，反而永远被更新，将永不匮乏。让我们的灵魂热切渴望那些日子，让它们炽热地渴望它们，使我们在那里得充满、满足，并说出我们现在所预说的话：\Quote{我们已心满意足}等等。\Quote{在祢使我们降卑的时候之后，现在祢又安慰了我们，并在我们见过邪恶的年岁中}\BibleRef{咏 89:15}。\par

16. 但在现今仍是邪恶的日子里，让我们这样说。\Quote{求祢垂顾祢的仆人和祢的作为} \BibleRef{咏 89:16}。因为祢的仆人本身就是祢的作为，不仅在于他们是人，更在于他们是祢的仆人，即服从祢命令的人。因为我们原是祂的化工，不只是在亚当内，而是在基督耶稣内，为行善事而受造的，就是天主所预先预备，使我们应去行事的：\BibleRef{弗 2:10}\Quote{因为是天主在你们内运作，使你们愿意并实行祂的美意。} \BibleRef{斐 2:13}\Quote{并指引他们的子孙：}使他们心里正直，因为天主对这样的人是慷慨的；因为\Quote{天主对以色列，对心里正直的人，原是慷慨的。}……\par

17. \Quote{愿上主我们天主的灿烂光辉照在我们身上} \BibleRef{咏 89:17}；因此有句话说：\Quote{上主，祢面容的光辉已印在我们身上。}又说：\Quote{求祢在我们身上，使我们双手的作为正直：}使我们行事不为追求地上的赏报：因为那样就不正直，而是歪曲的。在许多抄本中，圣咏到此为止，但有些抄本末尾还附有一节，如下：\Quote{并使我们双手的工作正直。}学者们在这句话前加了一颗星，称为星号，表示这些字见于希伯来文，或见于某些希腊文译本，但不在\WorkTitle{七十贤士译本}中。如果我们要解释这一节，在我看来，它的意思是：我们所有的善行都是爱的一项工作；因为爱是法律的满全。\BibleRef{罗 13:10}因为在上一节他说道：\Quote{求祢在我们身上，使我们双手的作为正直，}这里他说\Quote{工作}，而非\Quote{作为}，仿佛急于在最后一节表明，我们所有的作为都是一，即都指向一项工作。因为那时作为是正义的，当它们指向这一个目的时：\Quote{诫命的目的就是爱德，来自纯洁的心、良好的良心和真诚的信德。} \BibleRef{弟前 1:5}因此，有一项工作，在其中包含了一切，\Quote{那藉着爱德运作的信德}：\BibleRef{迦 5:6}因此主在福音中说：\Quote{天主要求的工作，就是要你们信祂所派遣来的那一位。} \BibleRef{若 6:29}因此，在这篇圣咏中，旧的和新的生命、必死的和永生的生命、被算为空虚的岁月和充满慈爱与真正喜乐的岁月，即第一人的惩罚和第二人的统治，被如此清晰地标明；我想，\WorkTitle{圣咏}的标题用了\Quote{天主的人梅瑟}之名，是为了让虔敬正直的圣经读者得到一个暗示：梅瑟的法律——在其中天主似乎只或几乎只应许现世报偿给善行——无疑在面纱下包含着像这篇圣咏所展示的某些希望。但当任何人转向基督时，那面纱就会被除去：\BibleRef{格后 3:15}他的眼睛将被揭开，使他得以默思天主法律中的奇妙之事，这是藉着祂的恩赐——我们向祂祈祷：\Quote{求祢开启我的眼睛，使我看到祢法律中的奇妙。}\par

\SourceRecord{Translated by J.E. Tweed.；Nicene and Post-Nicene Fathers, First Series；Vol. 8.；Edited by Philip Schaff.；Buffalo, NY: Christian Literature Publishing Co.,；1888.；<http://www.newadvent.org/fathers/1801090.htm>.}

% structure: p0001 source=h1 target=section reason=page-title

\section{\WorkTitle{圣咏第九十一篇释义}}

1. 这篇圣咏是魔鬼敢于用来试探我们的主耶稣基督的：因此我们当留意它，以便这样武装起来，能够抵抗那试探者，不依靠自己，而依靠那在我们之前受过试探的主，使我们受试探时不至被击败。试探对祂并非必要；基督的试探是我们的教训，但如果我们听从祂对魔鬼的回答，以便我们自己受试探时也同样回答，那么我们就正从门进入，正如你在福音中所听到的那样。因为从门进入是什么意思？就是从基督进入，祂自己说过：\Quote{我就是门：}\BibleRef{若 10:7}……祂敦促我们在那些祂若非成为人便不能做的事上效法祂；因为祂怎能受苦，除非祂成了人？否则祂怎能死、被钉十字架、受屈辱？因此，当你遭受这个世界的磨难时（这些磨难是魔鬼或公开藉着人，或隐秘地，如\Person{约伯}的情形所施加的），你要勇敢，要坚忍；\Quote{你必住在至高者的护佑下，}正如这圣咏所说的：因为如果你离开至高者的帮助，无力自助，你就会跌倒。\par

2. 因为许多人勇敢，当他们遭受人的迫害时，看见人们公开向自己发怒：他们以为那时他们是在效法基督的苦难，如果人们公开迫害他们；但如果受到魔鬼隐秘的攻击，他们相信他们不会从基督那里获得冠冕。当你效法基督时，永远不要害怕。因为当魔鬼试探我们的主时，旷野里没有人；他隐秘地试探祂；但他被征服了，当他公开攻击时也被征服了。你要这样做，如果你愿意从门进入，当敌人隐秘地攻击你时，当他要求一个人以便通过身体上的痛苦、热病、疾病或任何其他身体上的痛苦（如同\Person{约伯}所受的）来伤害他时。他没有看见魔鬼，但他承认天主的能力。他知道除非从全能万有的统治者那里得到那能力，魔鬼对他没有任何能力：他将全部光荣归于天主，没有将能力归于魔鬼……\par

3. 那么，那如此效法基督的人，以对天主的希望忍受这世界的一切苦难，以至于不陷入任何罗网，不被任何恐惧击溃，他就是\Quote{住在至高者的护佑下，将安居在天主的保护中}\BibleRef{咏 90:1}，用你们所听到和歌唱的这篇圣咏开头的话。当我们读到这些经文时，你们会认出这些非常熟悉的话，魔鬼正是用它们来试探我们的主。\Quote{他要向主说：你是我的提携者，我的避难所：我的天主}\BibleRef{咏 90:2}。谁这样对主说话？\Quote{那住在至高者护佑下的人：}而不是住在他自己的护佑下。这人是怎样的？他住在至高者的护佑下，他不骄傲，不像那些吃了果子想成为神的人，失去了他们被造时的不朽。因为他们选择住在自己的护佑下，而不是住在至高者的护佑下：因此他们听从了蛇的暗示，\BibleRef{创 3:5}并藐视了天主的诫命；并最终发现，在他们身上实现的，是天主所威胁的，而不是魔鬼所应许的。\par

4. 那么你也这样说：\Quote{我必信靠祂。因为祂必亲自救我}\BibleRef{咏 90:3}，并非我自己。请注意，祂所教导的，不正是我们所有的信靠应在天主内，不在人内吗？祂将从哪里救你？\Quote{从猎人的罗网，和从严厉的话语。}从猎人的网罗中得救确实是一种巨大的祝福；但如何从严厉的话语中得救？许多人因严厉的话语而陷入猎人的罗网。我所说的意思是什么？魔鬼和他的天使布置罗网，如同猎人一样；那些在基督内行走的人远离那些罗网，因为他不敢在基督内布置他的网；他把它设在路的边缘，不在路上。那么，让基督作你的道路，你就不会陷入魔鬼的罗网……\par

然而，所谓\Quote{来自严厉的话}是什么意思？魔鬼曾用严厉的话陷害了许多人：例如，在外邦人中承认基督信仰的人，遭受外邦人的侮辱；他们听到责难时便脸红，因而从自己的路上退缩，结果落入猎人的网罗。严厉的话能对你做什么呢？什么也不能。敌人藉着辱骂来陷害你的网罗，对你毫无作用吗？网罗通常是在篱笆尽头为捕鸟而设的，有人往篱笆里扔石头：那些石头不会伤害鸟。谁曾因为往篱笆里扔石头而打中鸟呢？但是鸟被无害的响声惊吓，落入网中；同样，那些害怕诽谤者虚妄的责难、因无端的侮辱而脸红的人，就落入猎人的网罗，被魔鬼掳去……正如在外邦人中，害怕他们责难的基督徒落入猎人的网罗；同样，在基督徒中，那些努力比其余人更勤勉、更好的人，注定要忍受来自基督徒本身的侮辱。那么，我的弟兄，如果你偶尔找到一个没有外邦人的城市，这有什么益处呢？那里没有人因他是基督徒而侮辱他，因为那里没有外邦人；但是有许多生活败坏的基督徒，在他们中间，那些决心度正直生活、在醉酒中保持清醒、在放荡中保持贞洁、在占星问卜者中真诚敬拜天主、不询问这类事、在轻浮表演的观众中只去教堂的人，正遭受这些基督徒的辱骂和严厉的话，他们对这样的人说：\Quote{你是大能的、正义的，你是厄里亚，你是伯多禄；你从天上来了。}他们侮辱他：无论他向哪边转，他每边都听到严厉的言词；如果他害怕了，放弃基督的道路，他就落入猎人的网罗。但当他听到这样的话时，不偏离正路是什么意思？听到这些话时，他有什么安慰，使他不理会它们，并能从门进入？让他说：我算什么人，竟被称为仆人、罪人？他们对我的主耶稣说：\Quote{你附有魔鬼。}（\BibleRef{若 8:48}）你刚刚听到针对我们主的严厉的话：我主不必受此苦，但祂这样做是为了警告你防备严厉的话，免得你落入猎人的网罗。\par

5. \Quote{祂必在祂的双肩之间保护你，你必在祂的翅膀底下投靠}（\BibleRef{咏 90:4}）。祂这样说，为叫你的保障不是出于你自己，为叫你不要以为自己能自卫；祂要保护你，使你脱离猎人的网罗和严厉的话。\Quote{在祂的双肩之间}这个说法，既可以理解为前面，也可以理解为后面；因为肩膀是在头部附近；但在\Quote{你必在祂的翅膀底下投靠}这句话中，显然，天主展开的翅膀的保护把你放在祂的双肩之间，因此天主的翅膀在这边和那边把你放在中间，在那里你必不害怕有人伤害你：只是你要当心，永远不要离开那个地方，因为没有仇敌敢靠近。如果母鸡在翅膀下保护小鸡，那么你在天主的翅膀下，甚至对抗魔鬼和他的天使——那些像鹰一样在空中飞翔、企图掳走软弱幼小者的权势——岂不更加安全？因为把母鸡比作天主的智慧并非没有根据；因为基督自己，我们的主和救主，说自己好比一只母鸡：\Quote{我多少次愿意聚集你的子女，……}等等（\BibleRef{玛 23:37}）。那耶路撒冷不愿意：但愿我们愿意……如果你们观察别的飞禽，弟兄们，你们会发现许多孵蛋并温暖幼雏的；但没有一个像母鸡那样因同情小鸡而衰弱自己。我们看到燕子、麻雀、鹳鸟在巢外，无法判断它们是否有幼雏；但我们知道母鸡是母亲，因她声音的衰弱和羽毛的松弛：她因爱小鸡而完全改变；她使自己软弱，因为它们软弱。同样，因为我们软弱，天主的智慧使自己软弱，当\OriginalTerm{圣言}{Word}成了血肉，居住在我们中间（\BibleRef{若 1:14}），叫我们能在祂的翅膀底下投靠。\par

6. \Quote{祂的真理必以盾牌四面护卫你}（\BibleRef{咏 90:5}）。所谓\Quote{翅膀}就是\Quote{盾牌}；因为既没有翅膀也没有盾牌。如果两者都是字面的，那么一个怎能与另一个相同呢？翅膀能是盾牌，或盾牌能是翅膀吗？然而，所有这些说法都是藉着相似来比喻的。如果基督真是\OriginalTerm{石头}{Stone}（\BibleRef{宗 4:10}），祂就不能是\OriginalTerm{狮子}{Lion}；如果祂是狮子（\BibleRef{默 5:5}），就不能是\OriginalTerm{羔羊}{Lamb}；但祂既被称为狮子，又被称为羔羊（\BibleRef{若 1:29}），以及石头、牛犊，以及其他这类比喻，因为祂既不是石头，也不是狮子，也不是羔羊，也不是牛犊，而是耶稣基督，我们众人的救主，因为这些是比喻，不是字面的名称。经上说：\Quote{祂的真理必是你的盾牌}：这盾牌向我们保证，祂不会把那些信赖自己的人与那些希望于天主的人混为一谈。一个是罪人，另一个也是罪人；但假定一个自恃的人是藐视者，不承认自己的罪，他会说：如果我的罪使天主不悦，祂就不会让我活着。但另一个不敢举目，只捶着胸膛说：\Quote{天主，可怜我这个罪人罢！}（\BibleRef{路 18:13}）两者都是罪人：但这个讥笑，那个哀悼；这个藐视自己的罪，那个承认自己的罪。但天主的真理不看情面，辨别忏悔者与否认自己罪的人，谦卑者与骄傲者，自恃者与信赖天主者。\Quote{你不必害怕黑夜的恐怖。}\par

7. \BibleQuote{也不怕白日飞的箭，也不怕黑暗中行走的瘟疫，也不怕正午使人毁灭的毒病} \BibleRef{咏 90:6}。以上两句对应以下两句；\BibleQuote{你不必害怕}指的是\BibleQuote{黑夜的惊惶，白日飞的箭}：既因\BibleQuote{黑夜的惊惶}，来自\BibleQuote{黑暗中行走的瘟疫}；也因\BibleQuote{白日飞的箭}，来自\BibleQuote{正午魔鬼的毁灭}。什么是黑夜所当怕的，什么是白日所当怕的？当一个人无知地犯罪时，他仿佛在黑夜中犯罪；当他明知故犯时，就是在白日中犯罪。前两种罪较轻，后两种则重得多；但这一点隐晦难明，若蒙天主的降福，我能解释清楚使你们理解，那便值得你们留意。他把无知者所屈从的轻微诱惑称为\BibleQuote{黑夜的惊惶}；把明知故犯者所受的轻微诱惑称为\BibleQuote{白日飞的箭}。什么是轻微诱惑？那些并不紧迫逼迫我们以致胜过我们，但若拒绝便能迅速过去的诱惑。假设这些是重诱。如果迫害者威胁并严重惊吓无知者——我的意思是那些信德尚不稳固、不知道他们成为基督徒是为了盼望来生的人——他们一旦因现世灾祸而惊慌，便以为基督抛弃了他们，认为作基督徒没有益处；他们不明白，他们成为基督徒正是为了胜过当下，盼望将来。那黑暗中行走的瘟疫逮住了他们。但也有些人知道他们是蒙召为将来的望德；知道天主所应许的不属于今生今世；知道必须忍受这一切诱惑，才能领受天主所应许的永恒事物；这一切他们都知晓。然而，当迫害者更猛烈地逼迫他们，以威胁、刑罚、酷刑相逼，他们终于屈服了，尽管深知自己的罪，却如同在白日跌倒一般。\par

8. 但他为什么说\Quote{正午}？因为迫害极其炎热；所以正午象征过度的炎热……那\Quote{正午}的魔鬼代表猛烈迫害的热度：因为我们的主说过：\BibleQuote{太阳升起，因为没有根，就枯干了}；当解释时，祂将这话应用于那些因迫害而跌倒的人，说：\BibleQuote{因为他们在自己内没有根}。因此，我们正确地将那在正午毁灭的魔鬼理解为强烈的迫害。亲爱的，请听我描述主从其中拯救了祂的教会的迫害。起初，当世俗的皇帝和君王以为能通过迫害从地上铲除基督徒的名号时，他们宣告：凡承认自己是基督徒的，应被击打。凡不愿被击打的人，就否认自己是基督徒，明知自己在犯罪；白日飞的箭击中了他。但那些不看重现世生命、却坚信来生的人，通过承认自己是基督徒而避开了那箭；肉身被击打，灵魂获得释放；与主同在，开始平安地等待身体在死人复活中的救赎；他逃脱了那诱惑，逃脱了白日飞的箭。\BibleQuote{凡自认是基督徒的，应被斩首}，这如同白日飞的箭。那\Quote{正午的魔鬼}尚未出现，它用可怕的迫害灼烧，甚至以极大的热度煎熬强壮者。请听后来发生的事；当敌人看到许多人急忙赴殉道，新皈依的人数随着受苦者的增加而增长时，他们彼此说：我们要消灭人类，因为相信祂名的人如此成千上万；如果我们杀尽他们，地上将几乎无人幸存。于是太阳开始炽燃，发出可怕的热度。他们的第一道法令是：凡承认自己是基督徒的，应被击打。第二道法令是：凡承认自己是基督徒的，应受折磨，甚至折磨到否认自己是基督徒为止……因此，许多不否认的人却在折磨中失败了；因为他们被折磨直到否认。但那些坚持承认基督的人，刀剑一击杀死身体，将灵魂送到天主那里，又能做什么呢？这也是长期折磨的结果：然而，谁能找到能够抵抗如此残酷而持续的折磨的人呢？许多人失败了：我相信，是那些自恃的人，那些不居住于至高者的庇护下、不居住于天上天主的荫庇下的人；那些没有对主说\BibleQuote{你是我的举起者}的人；那些不信任祂翅膀的荫庇，却过分依赖自己力量的人。天主推倒他们，以向他们显示：是祂保护他们，祂掌管他们的诱惑，祂只允许每个人所能承受的分量临到他们。\par

9. 那时有许多人跌倒在午间魔鬼面前。你们想知道有多少吗？他继续说，并说：\Quote{有千人跌倒在你的旁边，万人倒在你右边；但这灾祸不会临近你。}\BibleRef{咏 90:7} 弟兄们，这话是对谁说的？岂不是对耶稣基督说的吗？……因为肢体、身体和头是彼此不分离的：身体和头就是教会和她的救主。那么，怎么说\Quote{有千人跌倒在你的旁边，万人倒在你右边}？因为他们要跌倒在毁灭人的午间魔鬼面前。弟兄们，从基督旁边、从祂右边堕落，是一件可怕的事；但他们怎么会从祂旁边堕落呢？为什么有的在祂旁边，有的在祂右边？为什么千人跌倒在祂旁边，万人跌倒在祂右边？为什么千人跌倒在祂旁边？因为千人比那跌倒在祂右边的万人少。这些人是谁，借着基督的名很快会清楚；因为祂曾应许一些人要与祂一同审判，即那些舍弃一切跟随祂的宗徒……那些审判者就是教会的首领，成全的人。祂对他们说：\Quote{你若愿意成为成全的，去变卖你所有的，施舍给穷人。}\BibleRef{玛 19:21} 这句话\Quote{你若愿意成为成全的}是什么意思？意思是：你若愿意与我一同审判，而不被审判……当时有许多这样的人，他们已经把一切分施给穷人，并已应许自己在基督旁边审判万民的席位，却在迫害的烈焰折磨中，如同在午间魔鬼面前，失败了，否认了基督。这些就是跌倒在祂\Quote{旁边}的人：当他们正要与基督一同审判世界时，他们跌倒了。\par

10. 我现在要解释谁是在基督右边跌倒的人……因为许多人从作审判者的希望中跌倒了，然而还有更多的人从在祂右边的希望中跌倒了，因此圣咏作者这样对基督说：\Quote{有千人跌倒在祢旁边，万人倒在祢右边。}既然有许多人并不顾念这一切，基督与他们仿佛同为一身，祂便加上说：\Quote{但这灾祸不会临近祢。}这话是仅仅对头说的吗？当然不是；当然也不临近保禄、伯多禄、所有宗徒和所有在折磨中没有失败的殉道者。那么\Quote{不会临近}是什么意思？为什么他们受这样的折磨？折磨临近了肉身，却没有到达信德的范围。他们的信德远远超出了折磨者所威胁的恐怖所能触及的范围。任凭他们折磨，恐怖不会临近；任凭他们折磨，他们却要嘲笑折磨，依靠那在他们之前得胜的那一位，使其余的人也能得胜。谁得胜呢？岂不是那些不依靠自己的人吗？……谁不害怕？那不依靠自己而依靠基督的人。但那些依靠自己的人，虽然他们甚至希望坐在基督旁边审判，虽然他们希望自己是在祂右边，仿佛祂对他们说：\Quote{你们这些蒙我父祝福的，来吧！}等等；然而那午间魔鬼追上了他们，迫害的炽热以暴力恐吓；许多人从审判席位的希望中跌倒了，关于他们经上说：\Quote{有千人跌倒在祢旁边；}许多人也从对他们职责的奖赏的希望中跌倒了，关于他们经上说：\Quote{万人倒在祢右边。}但这跌倒和午间魔鬼\Quote{不会临近祢}，即头和身体；因为主认识谁是属于祂的。\BibleRef{弟后 2:19}\par

11. \Quote{不过，你将亲眼观看，见到恶人的报应。}\BibleRef{咏 90:8} 这是什么意思？为什么说\Quote{不过}？因为恶人被允许压迫祢的仆人，迫害他们。那么他们迫害祢的仆人就可以不受惩罚吗？不是不受惩罚，因为虽然祢允许了他们，祢自己的人也因此获得了更明亮的荣冠，\Quote{不过}，等等。因为他们所愿意的恶，而不是他们无意识所行的善，将得到报应。所缺少的只是信德的眼睛，借此我们可以看到他们只是暂时被高举，而他们要永远哀号；那些被赋予权柄统治天主的仆人的，将要听到：\Quote{可咒骂的，离开我，到那给魔鬼和他的使者预备的永火里去吧！}\BibleRef{玛 25:41} 但如果每个人都有眼睛，正如经上所说\Quote{你将亲眼观看}，那么看到恶人在今生亨通，并用眼注视他，思考如果他不能改过自新，最终会怎样，这不是一件小事；因为那些现在想要雷击别人的人，日后自己将遭受雷击。\par

12. \Quote{因为祢，上主，是我的希望。}\BibleRef{咏 90:9} 现在他来到了那能力面前，这能力拯救他免于因\Quote{跌倒和午间魔鬼}而堕落。\Quote{因为祢，上主，是我的希望；祢将祢的避难所安置得极高。}\Quote{极高}是什么意思？因为许多人把他们的避难所建于天主内，作为逃避暂时迫害的庇护；但天主的庇护是高的，极其隐秘，你可以逃到那里躲避将来的忿怒。里面\Quote{祢将祢的避难所安置得极高。灾祸不会临到祢，祸患也不挨近祢的帐棚。}\BibleRef{咏 90:10}\par

13. 圣城不仅是一个地区的教会，也是全世界的教会；不仅是这个时代的，而且从亚伯尔本人直到那些将要诞生并信从基督直到终结的人，全部圣人的集合，属于一座城；这座城是基督的身体，基督是这身体的元首。在那里，也居住着天使，他们是我们的同胞；我们劳苦，因为我们还是旅人；而他们已经在那城中等待我们的到来。从我们流离失所的那座城，也有书信送达我们：那些书信就是圣经，它们劝勉我们善度生活。我为什么只说书信呢？君王亲自降临，成为我们在流浪中的道路：好使我们行走在他内，既不迷路，也不疲乏，也不落入强盗手中，也不被设在路旁的网罗捕获。因此，我们认识到这正是基督的整个位格连同教会……他自己是元首，是天主，与圣父同等，是天主的圣言，万物借着他而造成的（\BibleRef{若 1:3}）。但作为天主，他是为了创造；作为人，他是为了更新；作为天主，他是为了创造；作为人，他是为了恢复。那么，让我们定睛在他身上，聆听圣咏。亲爱的，请听。这就是这学堂的教学和道理，它能使你们不仅理解这一篇圣咏，而且能理解许多篇，只要你们谨记这一规则。有时一篇圣咏，以及所有先知预言，在论及基督时，只赞扬元首；有时从元首转向身体，即教会，却未明显地改变所谈论的位格：因为元首与身体并不分离，两者被当作一体来谈论……\par

14. 那么，我的弟兄们，关于我们的元首说了什么？\Quote{因为主，你是我的希望，}等等。关于这一点我们已经说过，\BibleQuote{因为他为你吩咐了他的天使，在你行走的一切道路上保护你}（\BibleRef{咏 90:11}）。这些话你们刚才在宣读福音时已经听到；因此请留意。我们的主受洗后，就禁食。他为什么受洗？好使我们不轻看受洗。因为当若翰对我们的主说：\BibleQuote{你到我这里来受洗吗？我本该受你的洗；}我们的主回答说：\BibleQuote{你暂且容许吧，因为我们应当这样完成全义}（\BibleRef{玛 3:14}）。他愿意完成一切谦卑，好使那没有玷污的被洗涤。…我们的主于是受了洗，受洗后他受了试探；他禁食四十天，这个数字，正如我常提到的，有深意。一切事物不能同时解释，以免占用必要的时间。四十天后，他饿了。他本可以禁食而不觉饥饿；但那样他怎能受试探呢？或者，如果他没有胜过诱惑者，你怎能学习与他争战呢？他饿了；那时诱惑者说：\BibleQuote{你若是天主子，就命令这些石头变成饼。}我们的主耶稣基督用五个饼使数千人吃饱，难道用石头变出饼是件大事吗？他从无中造出饼。因为使这么多人吃饱的那些食物从哪里来的？那饼的来源在主的手中。这并不稀奇；因为他自己用五个饼造出足够数千人吃的饼，他也每天用几粒种子在地上长出丰盛的收成。这些是我们的主的奇迹：但由于它们不断运作，人们便忽视了。那么，我的弟兄们，主难道不能用石头造出饼吗？他甚至用石头造出人来，正如洗者若翰自己所说的：\BibleQuote{天主能从这些石头给亚巴郎兴起子孙来}（\BibleRef{玛 3:9}）。那么他为什么没有这样作呢？为要教导你如何回答诱惑者，好使你若陷入困境，诱惑者暗示说：如果你是一位基督徒，属于基督，他现在会丢弃你吗？…请听我们的主：\BibleQuote{人生活不只靠饼，也靠天主口中所发的一切言语。}你认为天主的圣言是饼吗？如果天主的圣言——万物借他而造成——不是饼，他就不会说：\BibleQuote{我是从天上降下来的生活的饼}（\BibleRef{若 6:41}）。因此，当你在饥饿时，你已学会回答诱惑者。\par

15. 如果他这样试探你：如果你是基督徒，你就会行奇迹，如同许多基督徒所做的那样？你被恶毒的暗示所欺骗，就会试探主你的天主，向他说：如果我是基督徒，在你眼前，你把我算在你自己的数目中，也让我做些如同你的圣人们所做的许多事吧。你试探了天主，仿佛若不是如此，你就不是基督徒。许多渴望这类事的人都跌倒了。因为那个行邪术的西满曾渴望宗徒们的这种恩赐，他想用钱购买圣神（\BibleRef{宗 8:18}）。他喜爱行奇迹的能力，却不喜爱谦卑的仿效。…那么，如果他这样试探你，\Quote{行奇迹}？为使你不试探天主，你该怎样回答？像我们主所回答的一样。魔鬼对他说：\Quote{你跳下去，因为经上记着：\BibleQuote{他为你吩咐了他的天使}，}等等。如果你跳下去，天使会接住你。我的弟兄们，如果我们的主跳下去，护卫的天使确实会接住我们主的肉身；但他对他怎么说呢？\Quote{经上又记着：\BibleQuote{你不可试探主你的天主}}（\BibleRef{申 6:16}）。你把我当作一个人。因为魔鬼来见他，带着这样的意图，要试探他是否是天主子。他看见他的肉身；但他的大能显现在他的作为中：天使已经见证过。他看见他是可朽的，好试探他，使基督徒因基督的试探而受教。那么经上记着什么？\BibleQuote{你不可试探主你的天主。}因此，我们不要试探主，以致说：如果我们属于你，就让我们行个奇迹吧。\par

16. 让我们回到圣咏的话。\newline{}\BibleQuote{他们要用手托着你，免得你的脚碰在石头上。}\BibleRef{咏 90:12}\newline{}\newline{}\newline{}基督被接升天时，是由天使的手托起的；这并不是说，若天使不托住他，他就会跌倒，而是因为他们正服侍他们的君王。不要说，那些托起他的人比那被托起的更好。难道牲畜因为托起人的软弱，就比人更好吗？我们也不该这样说；因为若牲畜撤去扶持，骑乘者就会跌倒。但我们应该怎样说呢？因为连天主也有话说：\BibleQuote{天是我的宝座。}\newline{}既然天托住天主，天主坐在其上，难道天就更好吗？同样，在这圣咏中，我们也可以理解为天使的服务：这不涉及我们主的任何软弱，而是关乎他们给予的尊荣和他们的服务……\newline{}什么是天主的手指，福音向我们解释；因为天主的手指就是圣神。我们如何证明这点？我们的主回答那些用贝耳则步的名号驱魔来控告他的人时说：\BibleQuote{我若靠着天主的神驱魔，}\BibleRef{玛 12:28}\newline{}另一位福音作者在叙述同样的话时说：\BibleQuote{我若靠着天主的手指驱魔。}\BibleRef{路 11:20}\newline{}因此，在一个经文中清楚声明的事，在另一个经文中却是隐晦地表达。你先前不知道什么是天主的手指，但另一位福音作者通过称它为天主的神解释了它。\newline{}那么，由天主的手指所写的法律，是在羔羊被宰杀后的第五十天赐下的；而圣神在我们主耶稣基督受难后的第五十天降临。羔羊被宰杀，逾越节被庆祝，五十天满了，法律被赐下。但那法律是要引起恐惧，而非爱；但为使恐惧转变为爱，那位真正正义者被宰杀——犹太人宰杀的那只羔羊正是他的预像。他从死者中复活；从我们主的逾越节之日，如同从逾越节羔羊被宰杀之日，计算了五十天；然后圣神降临了，这时充满了爱，而非惩罚的恐惧。\BibleRef{宗 2:1}\newline{}我为什么说这些？因为我们的主复活并被光荣，正是为派遣他的圣神。我先前说过，这是如此，因为他的头在天上，他的脚在地上。若他的头在天上，他的脚在地上，那么我们主的脚在地上是什么意思？我们主在地上的圣人。谁是我们主的脚？被派遣到全世界的宗徒们。谁是我们主的脚？所有福音作者，我们的主藉着他们巡游万邦……\newline{}因此我们不必惊讶我们的主被天使的手托起升天，以免他的脚碰在石头上；惟恐那些在地上在他身体中劳苦、遍行全世界的人可能陷于法律的罪责，他取走了他们的恐惧，以爱充满他们。因恐惧，伯多禄三次否认了他，\BibleRef{玛 26:69}\newline{}因为他那时还没有领受圣神；后来，当他领受了圣神，就开始勇敢地宣讲……\newline{}我们的主这样对待他，仿佛是说：因恐惧你三次否认了我；因爱你要三次承认我。用那爱和那爱德他充满他的门徒。为什么？因为他把他的防御居所设得极高；因为他复活后被光荣时派遣了圣神，他从法律的罪责中释放了信众，使他的脚不至碰在石头上。\par

17. \BibleQuote{你要践踏毒蛇和蛇怪，你必踩踏狮子和毒龙。}\BibleRef{咏 90:13}\newline{}你们知道蛇是谁，教会如何践踏他，因为她未被征服，因为她提防他的狡猾。而他是怎样成为狮子和毒龙的，我相信你们也知道，亲爱的。狮子公开咆哮，毒龙藏匿于暗处；魔鬼具有这两种力量和权能。当殉道者被屠杀时，那是怒吼的狮子；当异端者在图谋时，那是潜行于我们脚下的毒龙。你已征服了狮子；也要征服毒龙：狮子没有压垮你，不要被毒龙欺骗……\newline{}教会中有少数妇女拥有身体上的童贞；但所有信友都拥有心灵的童贞。在信仰之事上，他害怕心灵的童贞会被魔鬼败坏；那些失去它的人，身体的童贞也是徒然。一个内心败坏的女人，她身体里保存了什么？因此，一个天主教已婚妇女优于一个异端的童贞女。前者她的身体确实不是童贞，但后者在心灵上成了已婚者；而且不是嫁给了天主作丈夫，而是嫁给了毒龙。但教会该怎么办呢？蛇怪是众蛇之王，正如魔鬼是恶灵之王。\par

18. 这些是天主对教会说的话。\BibleQuote{因为他全心爱慕我，我必拯救他。}\BibleRef{咏 90:14}\newline{}因此，不仅是那现在坐在天上的头，因为他把防御居所设得极高，没有任何邪恶能临到他，也没有灾祸靠近他的住所；而且我们这些在地上劳苦、仍活在诱惑中、唯恐脚步陷入网罗的人，也能听到我们的天主安慰我们的声音，并对我们说：\BibleQuote{因为他全心爱慕我，我必拯救他；我要将他置于高处，因为他认识我的名。}\par

19. \BibleQuote{他必呼求我，我必应允他：是的，我在困苦中与他同在}\BibleRef{咏 90:15}。当你处于困苦中时，不要害怕，好像主不在你身边。让信德与你同在，那么天主就与你同在困苦中。海上有波浪，你在小船中颠簸，因为基督睡着了。基督在船上睡着了，而人们正在丧命。\BibleRef{玛 8:24} 如果你心中的信德沉睡了，基督就好像在你的船上睡着了：因为基督藉着信德住在你内。当你开始颠簸时，唤醒沉睡的基督：唤醒你的信德，你就必确信祂不会抛弃你。但你却以为自己被抛弃了，因为祂没有在你所希望的时刻拯救你。祂从火中救出了三青年？\BibleRef{达 3:29} 那行了这事的天主，难道会抛弃玛加伯吗？\EditorialNote{玛加伯下第七章} 绝对不可！祂拯救了二者：前者是肉体上，为使不信者蒙羞；后者是精神上，为使信者效法他们。\Quote{我必搭救他，并使他得荣耀。}\par

20. \BibleQuote{我必以长寿满足他}\BibleRef{咏 90:16}。什么是长寿？永生。弟兄们，不要以为这里所说的长寿，如同日子在夏季长、冬季短那样。祂要赐给我们这样的日子吗？那长度是无穷尽的，就是永生，是那在漫长日子中所应许给我们的。而且，既然这已足够，祂有理由说：\Quote{我必满足他。}时间上长久的东西，如果有尽头，就不能使我们满足：因此它甚至不该被称为长久。如果我们贪婪，就应当贪求永生：渴慕这样一种生命，没有尽头。看，一条线，我们的贪婪可以延伸。你渴望无限的钱财吗？渴望无限的永生吧。你希望你的产业没有尽头吗？寻求永生吧。\Quote{我必将我的救恩显示给他。}我的弟兄们，这一点也不可轻描淡写地放过。\Quote{我必将我的救恩显示给他：}祂的意思是，我要将基督自己显示给他。为什么？祂不是在世上被看见了吗？祂要显示给我们的是什么大事呢？但祂显现的形像并非我们将要看见的那样。祂显现的形像使看见祂的人把祂钉在十字架上：看，那些看见祂的人把祂钉死了；我们没有看见祂，却已相信了。他们有眼睛，难道我们没有吗？是的，我们也有心灵的眼睛：但是，我们如今是藉着信德观看，并非藉着目睹。何时才能目睹？何时我们才能如宗徒所说，\BibleQuote{面对面地}\BibleRef{格前 13:12}看见祂？这是天主应许给我们的，作为我们一切劳苦的高贵赏报。无论你为什么劳苦，你劳苦都是为了这个目的：为能看见祂。我们要看见的是一件大事，因为我们的全部赏报就是看见；而我们主耶稣基督就是那极其伟大的景象。那位曾以谦卑形像显现的，将要以伟大的形像显现，并要使我们喜乐，正如祂如今被祂的天使所看见的那样。……让我们爱慕祂，效法祂：让我们跟随祂的香膏，正如在\BibleBook{}{雅歌}中所说的：\BibleQuote{因你芬芳的香膏，我们愿跟随你。}\BibleRef{歌 1:3}因为祂来了，散发出充满世界的香气。那芬芳从何而来？从天上来。那么，就向着天跟随吧，当对你们说\Quote{请举心向上}的时候，你们若不以虚假回答，就举高你们的思念、你们的爱德、你们的望德：免得它在地上腐朽。……\BibleQuote{因为你的财宝在哪里，你的心也在那里。}\BibleRef{玛 6:21}\par

\SourceRecord{Translated by J.E. Tweed.；Nicene and Post-Nicene Fathers, First Series；Vol. 8.；Edited by Philip Schaff.；Buffalo, NY: Christian Literature Publishing Co.,；1888.；<http://www.newadvent.org/fathers/1801091.htm>.}

% structure: p0001 source=h1 target=section reason=page-title

\section{圣咏第九十二篇释义}

1. …我们不是基督徒，除非是为了来生：不要有人希望现世的祝福，不要有人因为自己是基督徒而应许自己世界的幸福：而是让他按他所能的、以他所能的方式、在他所能的时候、尽其所能地使用他所拥有的幸福。当幸福存在时，让他感谢天主的安慰；当幸福缺乏时，让他感谢天主的公义。让他常怀感恩，永不不知感恩：让他感谢他的圣父，祂抚慰和爱抚他：也感谢他的圣父，当祂用鞭子管教他、教训他时：因为祂永远爱他，无论祂是爱抚还是威胁：让他说出你们在圣咏中所听到的：\BibleQuote{对主感恩是好事；歌颂你的名，至高者。} \BibleRef{咏 91:1}\par

2. 这篇圣咏的标题是：在安息日咏唱的圣咏。看，今天是安息日，犹太人现今以一种身体上的、慵懒而放纵的休息来遵守它。他们停止劳作，沉溺于琐事；尽管天主制定了安息日，他们却把它用于天主禁止的行为。我们的休息是远离恶行，他们的休息是远离善行；因为耕种比跳舞更好。他们远离善行，但不远离琐事。天主向我们宣告一个安息日。什么样的安息日？首先思考，它在哪里。它在心里，在我们内；因为许多人肢体空闲，内心却不安……那在希望之宁静中的喜悦，就是我们的安息日。这篇圣咏的主题就是赞美和歌唱一个基督徒如何在他自己内心的安息日，即在他良心的安宁、安宁和宁静中，不受干扰；因此他在这里告诉我们，人们通常在哪里受到干扰，并教导你在自己心中守安息日。\par

3. …控告你自己，你会得到宽恕。此外，许多人不是控告撒殚，而是控告他们的命运。\InnerQuote{我的命运引导了我}，有人当你问他时这样说，你为什么做了这件事？你为什么犯罪？他回答：\InnerQuote{由于我邪恶的命运。}为了避免说\InnerQuote{我做了}，他把罪恶的根源指向天主：他用舌头亵渎。他还没有公开这样说，但请听，看他是否这样说。你问他，什么是命运；他回答：\InnerQuote{邪恶的星辰。}你问，谁创造了星星？谁安排了星星？他只能回答：天主。因此，无论他是直接还是间接这样做，他都在控告天主，而当天主惩罚罪恶时，他使天主成为他自己罪恶的制造者。天主不可能惩罚祂所造成的事：祂惩罚你所做的，以便释放祂所造成的事。但有时，撇开一切，他们直接攻击天主：当他们犯罪时，他们说：\InnerQuote{天主愿意这样；如果天主不愿意，我就不会犯罪。}祂告诫你，难道不仅是为了让你不听从，从而远离罪恶，甚至是为了因为你犯罪而受到谴责吗？那么，这篇圣咏教导我们什么呢？\BibleQuote{向主认罪是好事。} 什么是向主认罪？两种情况：既在你的罪恶中，因为你行了这些事；也在你的善行中，向主认罪，因为是祂行了这些事。那么你将\BibleQuote{歌颂至高天主的名：}寻求天主的荣耀，而不是你自己的；祂的名，而不是你的。因为如果你寻求天主的名，祂也会寻求你的名；但如果你忽视了天主的名，祂也会涂抹你的名……\par

4. \BibleQuote{早晨传述你的仁慈，夜间传述你的真理} \BibleRef{咏 91:2}。这是什么意思：天主的仁慈要在早晨传述给我们，而天主的真理要在夜间传述？早晨是我们顺遂的时候；夜间是苦难的悲伤。那么他简短地说什么呢？当你兴盛时，因天主而喜乐，因为这是祂的仁慈。现在，也许你会说：如果我在兴盛时因天主而喜乐，因为这是祂的仁慈；我在悲伤、困苦时该做什么呢？我兴盛时是祂的仁慈；我逆境时是祂的残酷吗？如果我在顺境时赞美祂的仁慈，我在逆境时是否要呼喊反对祂的残酷？不。但当顺境时，赞美祂的仁慈；当逆境时，赞美祂的真理：因为祂鞭打罪恶，祂并非不义……在夜间，\Person{达尼尔}传述了天主的真理：他在祈祷中说：\BibleQuote{我们犯了罪，行了不义，作了恶。主，正义属于祢；但我们脸上惭愧。} \BibleRef{达 9:5} 他在夜间传述了天主的真理。什么是夜间传述天主的真理？不是因你遭受任何恶事而控告天主：而是将你的罪恶归于祂的惩戒；早晨传述祂的慈爱，夜间传述祂的真理。当你\Emph{这样做}时，你总是赞美天主，总是向天主认罪，并歌颂祂的名。\par

5. \BibleQuote{在十弦瑟上，以歌和竖琴} \BibleRef{咏 91:3}。你们不是第一次听到十弦瑟：它象征法律的十诫。但我们必须在那个瑟上歌唱，而不仅仅携带它。因为连犹太人也有法律：但他们携带它；他们不歌唱……\Quote{和竖琴。} 这意味着，在言语和行为上；\Quote{以歌，}在言语上；\Quote{在竖琴上，}在行为上。因此，既要善言也要善行，这样你才能有歌与竖琴一起。\par

6. \BibleQuote{因为祢，上主，藉祢的作为使我欢乐；我要因祢双手的运作而赞颂喜乐} \BibleRef{咏 91:4}。你看他说什么。祢使我生活得好，祢塑造了我：若我偶然行善，我要因祢双手的运作而喜乐；如同宗徒说：\BibleQuote{因为我们是祂的化工，因善行而受造} \BibleRef{弗 2:10}。因为除非祂塑造你行善，你便不知善行，只知恶行……因为你不能从自身得着真理，你仍需饮用从那源头流出的真理；如同你若离开光明，便陷入黑暗；如同石头并非自身发热，而是因太阳或火而热，若将它从热源移开，它就会冷却，由此显见热并非其固有，而是由太阳或火而来；同样，你若远离天主，便会变冷；你若亲近天主，便会温暖；如宗徒说\BibleQuote{心中火热} \BibleRef{罗 12:11}。论光他又说了什么？你若亲近祂，便在光中；因此圣咏说\BibleQuote{你们瞻仰祂，要喜形于色；你们的脸上不会羞愧}。因此，除非你被天主的光照亮，被天主的圣神温暖，你便不能行善；当你见自己工作有成效时，当向天主承认，并说宗徒所说的；对你自己说，免得你自高自大，\BibleQuote{你有什么不是领受的呢？} \BibleRef{格前 4:7}……\par

7. 那行善却受苦的可怜人，看见恶人亨通，便不安地说：\Quote{哦，天主，我以为恶人讨祢喜悦，祢却恨恶善人，喜爱作恶的人}……内在的人失去了安息，心灵的安宁被拒之门外，善念被驱散，他开始效法那在恶行中亨通的人，自己也转向恶行。但天主是宽容的，因为祂是永恒的，祂知道祂审判的日子，在那里祂衡量一切。\par

8. 教导我们这些，他说了什么？\BibleQuote{上主，祢的作为何其荣耀；祢的思念极其深奥} \BibleRef{咏 91:5}。的确，我的弟兄们，没有海有如天主的思念这般深，祂使恶人亨通，善人受苦：没有什么如此深奥，如此深邃；每个不信的灵魂都在那深奥、那深邃中沉没。你想渡过这深奥吗？不要离开基督十字架的木头：你不会沉没：紧紧抓住基督。我所谓紧紧抓住基督是什么意思？正因这原因，祂亲自选择在世上受苦。你曾听见，当先知被宣读时，祂\BibleQuote{没有将背转给击打祂的人，也没有将脸转给人们的唾沫，祂没有将面颊转离他们的手} \BibleRef{依 50:6}。为什么祂选择忍受这一切，岂非为了安慰受苦的人？祂本可在末日复活肉身，但那样你就没有希望的基础，因为你没有见过祂。祂没有拖延祂的复活，免得你仍存疑虑。那么，你当怀着与你在基督内所见相同的结局，忍受世上的磨难：不要让那些作恶并在今生亨通的人动摇你。\BibleQuote{祢的思念极其深奥}。天主的思念在哪里？不要像鱼因饵而欢欣：渔夫尚未收钩：鱼还含着钩。你以为长久的，其实短暂；这一切都快速过去。人的长寿与天主的永恒相比算什么？你想变得宽容忍耐吗？思想天主的永恒。因为你顾及你短暂的时日，在你的数日内你希望所有事都成就。什么事？所有恶人的定罪，所有善人的加冕：你希望这些事在你日子内成就吗？天主在祂自己的时间成就。为何你忍受疲倦？祂是永恒的：祂等候：祂是宽容的：但你却说，我不宽容，因为我是必死的。但你能够变得宽容：将你的心与天主的永恒结合，你就要与祂一起永恒……\par

9. 因此，在说了\BibleQuote{祢的思念极其深奥}之后，他随即加上：\BibleQuote{愚昧人不深思此事，蠢人也不明白} \BibleRef{咏 91:6}。愚昧人不深思、蠢人不明白的是什么事？\BibleQuote{当恶人如草一般青翠。}什么是\BibleQuote{如草一般}？他们在冬天繁茂，夏天枯萎。你看那草的花？什么更快凋谢？什么更明亮？什么更青翠？不要因其青翠而喜悦，而要畏惧其枯萎。你已听见恶人如草青翠：也要听见义人的事：\BibleQuote{因为看哪。}与此同时，思想恶人：他们如草繁茂；但谁不明白这事？愚人和蠢人。\BibleQuote{当恶人如草一般青翠，众人都观看作孽的人} \BibleRef{咏 91:7}。凡心中对天主思想不正的人，都观看恶人如草青翠之时，即他们暂时繁茂之时。为何观看他们？\BibleQuote{使他们永远灭亡。}因为他们注视其短暂的光彩，效法他们，希望与他们一同暂时繁茂，却永远丧亡：这就是\BibleQuote{使他们永远灭亡。}\par

10. \Quote{但是主，祢是永远的至高者} \BibleRef{咏 91:8}。祢在祢的永恒中在上等待，直到恶人的季节过去，义人的季节来临。\Quote{因为看啊。}听吧，弟兄们。那发言者（因为他以我们的人性说话，以基督身体的位分说话，因为基督在祂自己的身体内，即在祂的教会内发言）已经将自己与天主的永恒联合：正如我稍早前对你们所说的，天主是长久忍耐和宽仁的，容许祂所看见恶人所行的所有恶事。为什么？因为祂是永恒的，并看见祂为他们所保留的。你们也愿意长久忍耐和宽容吗？将你自己与天主的永恒联合：与祂一同等待那些低于你的事物；因为当你的心紧紧依附于至高者时，一切可朽坏的事物都将低于你；然后说出以下的话：\Quote{因为看啊，你的仇敌必灭亡。}那些现在繁荣的，日后必灭亡。谁是天主的仇敌？弟兄们，也许你们认为只有那些亵渎者才是天主的仇敌？他们确实是，而且那些在言语和思想上不断伤害天主的恶人也是。他们向永恒、至高的天主做了什么？如果你用拳头击打柱子，你会受伤：你认为你用亵渎击打天主的地方，你自己不会破碎吗？因为你不能对天主做什么。但是天主的仇敌是公开的亵渎者，而且每天都有隐藏的仇敌被发现。要当心这样的与天主为敌。因为圣经启示了一些天主的秘密仇敌：因为你心中不认识他们，你可以在天主的圣经中认识他们，并当心不要与他们一同被找到。雅各伯在他的书信中公开说：\Quote{你们不知道与世俗的友谊就是与天主为敌吗？} \BibleRef{雅 4:4}你们听见了。你们不愿意成为天主的仇敌吗？不要作这世界的朋友；因为如果你作了这世界的朋友，你将成为天主的仇敌。因为正如妻子除非与自己的丈夫为敌，否则不能成为淫妇；同样，一个因贪爱世俗事物而成为淫妇的灵魂，不能不与天主为敌。它害怕，却不爱；它畏惧惩罚，却不以正义为乐。因此，所有贪爱世界的人，所有好奇琐事的人，所有求问占卜者、占星家和邪灵的人，都是天主的仇敌。无论他们是否进入教堂，他们都是天主的仇敌。他们可能如草一时繁茂，但当祂开始眷顾他们，并对一切血肉之躯宣告判决时，他们将灭亡。将你自己与天主的圣经结合，与这篇圣咏一同说：\Quote{因为看啊，你的仇敌必灭亡} \BibleRef{咏 91:9}。不要在那里被发现，在他们灭亡的地方。\Quote{所有作恶的人必被消灭。}\par

11. ……\Quote{但是我的角必被高举，像独角兽的角} \BibleRef{咏 91:10}。祂为什么说：\Quote{像独角兽的角}？有时独角兽象征骄傲，有时意指团结的高举；因为团结被高举，一切异端将与天主的仇敌一同灭亡。\Quote{我的角必被高举，像独角兽。}什么时候如此？\Quote{我的晚年必在丰盛的慈悲中。}祂为什么说\Quote{我的晚年}？祂的意思是，我的末后日子；正如我们的晚年是我们生命的最后季节，同样，基督的身体目前在劳苦、挂虑、警醒、饥饿、干渴、绊脚石、邪恶、困苦中所遭受的一切，就是它的青年；它的晚年，即它的末后日子，将在喜乐中。弟兄们，要当心，你们不要以为谈到晚年就意味着死亡：因为人肉身衰老，正是为了死亡。教会的晚年将因善工而发白，但它不会因死亡而朽坏。正如老人头上的白发，我们的善工也将如此。你看见头发如何随着老年的临近而变白。有时你在一个按常理衰老的人头上寻找一根黑发，却找不到；同样，当我们的生命变得这样，寻求罪的黑发却找不到时，那晚年是年轻的，是青翠的，并且永远青翠。你们听说过罪人的草，现在要听义人的晚年：\Quote{我的晚年必在慈悲的父怀中。}\par

12. \Quote{我的眼已看见我的仇敌} \BibleRef{咏 91:11}。祂称谁为仇敌？所有作恶的人。不要观察你的朋友是否是恶人；等机会到来，你就能考验他。你开始反对他的邪恶，那时你就会看见，当他奉承你时，他是你的仇敌；但你还没有敲打，不是为了在他心中激起原本没有的东西，而是为了让那已有的东西爆发出来。\Quote{我的眼也观看我的仇敌：我的耳必听见恶人起来攻击我的愿望。}什么时候？在我的晚年。什么是晚年？在末后时期。我们的耳将听见什么？站在右边，我们将听见对左边的人所说的话。\par

13. 草枯干，罪人的花凋谢；义人怎样？\Quote{义人必像棕榈树一样繁茂} \BibleRef{咏 91:12}。不敬虔的人绿如草；\Quote{义人必像棕榈树一样繁茂。}祂用棕榈树表示高度。可能祂在棕榈树中有此含义，即它的末梢是美丽的：这样你就能追溯它的开始从地起，它的终点在顶端，那里是它全部的美所在。粗糙的根显现于地，美丽的枝叶朝向天。那么，你的美也将在终点。你的根扎得牢固：但我们的根是向上的。因为我们的根是基督，祂已升天。谦卑的，祂要被高举；\Quote{他必像黎巴嫩香柏一样伸展。}看看祂说的是什么树：义人必像棕榈树一样繁茂，且像黎巴嫩香柏一样伸展。太阳出来时，棕榈树会枯萎吗？香柏会死吗？但当太阳照耀数小时后，草便干枯了。因此，审判要来，使罪人枯萎，使信者繁茂。\par

14. \Quote{栽植在上主殿中的，必在我们天主的庭院里开花} \BibleRef{咏 91:12}。\Quote{他们要在丰盛的老年还要增长，并且要安静，好能显示这事} \BibleRef{咏 91:13}。这就是那安息日，不久之前我曾向你们称赞过的，诗题由此而来。\Quote{他们要安静，好能显示这事。}他们为什么安静以显示这事？恶人的青草不触动他们；香柏树和棕榈树即或是在暴风雨中也不弯曲。所以，他们安静，好能显示这事：理由充足，因为现在他们还必须向那些嘲笑此事的人显示。啊，可怜的人，世间的爱慕者！那些栽植在上主殿中的人，向你们显示：那些以歌与琴、以言与行赞美上主的人，向你们显示，并告诉你们。不要被恶人的兴盛所诱惑，不要羡慕那如草的荣华：不要羡慕那些只一时快乐、却永远悲惨的人……如果你们愿意像棕榈树一样开花，像黎巴嫩的香柏树一样伸展，而不像草在烈日下枯萎——那些人在太阳隐没时似乎开花；如果你们不愿像草，而像棕榈树和香柏树，那么你们要显示什么？\Quote{上主，我的力量，是何等真实，且祂内毫无不义。}如何毫无不义？一个人犯了如此大的罪：他平安无事，有儿子，房屋充足，充满骄傲，因尊荣而高升，向敌人报复，且行一切恶事；另一个人，无辜，专心自己的事，不抢夺他人的财物，不与人作对，却受苦于锁链、牢狱，辗转叹息于贫困。何以祂内毫无不义？安静，你将知道：因为你被搅扰，你在你的内室里使你的光昏暗。永恒的天主愿意照耀你：不要因你自己烦乱的心使天气阴霾。在你内里安静，看我对你所说的话。因为天主是永恒的，现世祂宽容恶人，引导他们悔改；祂鞭策善人，教导他们通往天国的道路：\Quote{祂内毫无不义，}不要害怕……什么？若祂现在不惩罚这个人，因为他注定要听见：\Quote{去吧，到永火里去。}但何时？当你被置于右边时，那时要对那些在左边的人说：\Quote{去吧，到那为魔鬼和他的使者预备的永火里去。}因此，不要让那些事搅扰你：安静，遵守安息日，并显示\Quote{上主，我的力量，是何等真实，且祂内毫无不义。}\par

\SourceRecord{Translated by J.E. Tweed.；Nicene and Post-Nicene Fathers, First Series；Vol. 8.；Edited by Philip Schaff.；Buffalo, NY: Christian Literature Publishing Co.,；1888.；<http://www.newadvent.org/fathers/1801092.htm>.}

% structure: p0001 source=h1 target=section reason=page-title

\section{圣咏第九十二篇释义}

1. ...它题为，\Quote{大卫自己的赞美歌，安息日前夕，大地奠定之时}。然后回想天主在这六日内所做的一切，当祂创造并安排万物，从第一日直到第六日（因为祂祝圣了第七日，在那日祂歇了祂的一切工作，祂认为都非常好），我们发现祂在第六日（这正是这里所说的\Quote{安息日前夕}）创造了地上的一切动物；最后，就在那一日，祂按自己的肖像和模样造了人。因为这些日子并非毫无理由地如此安排，而是为预示时代也要经历相似的进程，然后我们才在天主内安息。但只有当我们行善工时，我们才能安息。作为这事的预象，论到天主，经上记载说，\BibleQuote{天主在第七日休息}，那时祂完成了一切工作，都认为非常好。因为祂并不疲乏，需要休息，如今祂也没有停止工作，因为我们的主基督公开说，\BibleQuote{我父一直工作到现在}。\BibleRef{若 5:17} 祂向犹太人这样说，他们按肉体思想天主，不明白天主在安静中工作，且常工作，又常安静。同样，我们这些天主愿意在祂内预象的人，在一切善工之后也将得到安息......因为善工注定要过去，那第六日，当那些善工完成时，也有黄昏；但在安息日我们找不到黄昏，因为我们的安息将没有尽头：因为黄昏代表终结。因此，如同天主在第六日按自己的肖像造了人，同样我们发现我们的主耶稣基督进入第六时代，好使人能重新按天主的肖像被塑造。因为第一个时期，如同第一日，从亚当到诺厄；第二个时期，如同第二日，从诺厄到亚巴郎；第三个时期，如同第三日，从亚巴郎到达味；第四个时期，如同第四日，从达味到流徙巴比伦；第五个时期，如同第五日，从流徙巴比伦到若翰的宣讲；第六日从若翰的宣讲开始，直到终结；第六日过后，我们到达安息。因此，第六日如今正在过去。现在正是第六日，请看标题所说：\Quote{安息日前夕，大地奠定之时}。让我们现在来听这篇圣咏：让我们探究大地是如何造成的，或许大地是在那时造成的？但我们在\WorkTitle{创世纪}中没有读到这样的记载。那么，大地是在何时奠定的呢？岂不是当我们在宗徒书信中刚才读到的话实现之时：\BibleQuote{你们要坚定不移，毫不动摇}。\BibleRef{格前 15:58} 当天下所有相信的人都坚定在信德中时，大地就奠定了；那时人就被造在天主的肖像中。\WorkTitle{创世纪}中的第六日预表的就是这个......\par

2. \BibleQuote{主为王，祂以美为衣；主以能力为衣，且束腰带}。\BibleRef{咏 92:1} 我们看到祂穿上了两样：美丽和能力。但为什么？为奠定大地。因此接下去说，\BibleQuote{祂使大地如此稳固，永不动摇}。\BibleRef{咏 92:1} 祂为何使它如此稳固？因为祂穿上了美丽。如果祂只穿上美丽，而没有能力，就不会使它如此稳固。那么为什么美丽，为什么能力？因为祂说了二者。你们知道，弟兄们，当我们的主降生成人后，在那些祂宣讲福音的人中，祂使一些人喜悦，使另一些人不悦。因为犹太人的口舌彼此分歧：\BibleQuote{有的说：祂是好人；有的说：不，祂迷惑众人}。\BibleRef{若 7:12} 于是有的人说好话，有的人诽谤祂、撕裂祂、咬祂、侮辱祂。所以对于那些祂悦纳的人，祂\Quote{穿上美丽}；对于那些祂不悦纳的人，祂\Quote{穿上能力}。那么，效法你们的主，好使你们成为祂的衣裳：对于你们的善工所悦纳的人，要展现美丽；对于毁谤者，要展现能力......\par

3. 也许我们也应查考这句话，为什么他说：\Quote{他束上了腰。}束腰表示工作：因为每当人准备工作时，就会束腰。但他为什么用了\OriginalTerm{在前束腰}{præcinctus}这个词，而不是\OriginalTerm{束腰}{cinctus}？因为在另一篇圣咏中，他说：\BibleQuote{大能者啊，愿你腰间佩上你的剑，人民必在你以下仆倒。}用了\OriginalTerm{佩束}{accingere}，而不是\OriginalTerm{束}{cingere}，也不是\OriginalTerm{前束}{præcingere}。这个词被用来指通过束腰将某物系在腰间。主的剑，他用它杀死邪恶而征服了世界，就是天主之神在真理的天主圣言中。为什么说他将剑束在大腿上？在别处，在另一篇圣咏中，我们曾以另一种方式谈论过束腰；但既然提到，就不应忽略。剑束在大腿上是什么意思？他用大腿指肉体。因为主若没有使真理的剑进入肉体，就不能征服世界。那么为什么这里说他在前面束腰（\OriginalTerm{在前束腰}{præcinctus}）？他在前面束腰，就是把某物放在自己面前，用它来束腰；因此记载他先用毛巾束腰，然后开始洗门徒的脚。因为他用毛巾束腰时是谦卑的。他洗了自己门徒的脚。但一切力量都在谦卑中：因为一切骄傲都是脆弱的；因此当谈到力量时，他补充说：\Quote{他束上了腰。}要你记得你的天主是如何谦卑地束腰，当他洗门徒的脚时。\BibleRef{若 13:4}……他洗了他们的脚后，又坐下，对他们说：\BibleQuote{你们称我为主、为师：说得对，我实在是。若我为主、为师的，洗了你们的脚，你们也该彼此洗脚。}\par

4. ……\BibleQuote{因为他使世界坚定，不得动摇。}……那么什么是\Quote{不得动摇}的世界？他若不特别提到，就不会有一个可以动摇的世界。有一个不得动摇的世界。有一个将要被动摇的世界。因为那些在信仰中坚定不移的善人就是世界：这样没有人会说，他们只是其中的一部分；而恶人不在信仰中持守，当他们感到任何磨难时，就遍及整个世界。所以有一个动摇的世界：有一个不动摇的世界，宗徒提到它。看，动摇的世界。我问你，宗徒在这句话中说到谁？\BibleQuote{他们中有\Person{依默纳约}和\Person{斐肋托}，他们偏离了真理，说复活已过，就颠覆了一些人的信仰？}\BibleRef{弟后 2:17}这些人属于那个不得动摇的世界吗？但他们是糠秕：如他所说：\BibleQuote{他们颠覆了一些人的信仰。}……\BibleQuote{然而，天主的根基稳固，有这印记}——什么印记作为它的稳固根基？——\BibleQuote{主认识属于他的人。}这就是那不得动摇的世界；\BibleQuote{主认识属于他的人。}它有什么印记？\BibleQuote{凡称呼基督之名的人，应远离不义。}让他远离不义：因为他不能离开不义的人，因为糠秕与麦子混杂，直到被扬净……\par

5. \BibleQuote{你的宝座从那时建立，上主}\BibleRef{咏 92:2}。\Quote{从那时}是什么意思？从那时起。好像他说，什么是天主的宝座？天主坐在哪里？在他的圣人中。你愿意成为天主的宝座吗？在心里准备一个地方让他坐。什么是天主的宝座，除非是天主居住的地方？天主住在哪里，除非在他的圣殿里？是用墙围起来的吗？远非如此。也许这个世界是他的圣殿，因为它很大，而且配得上容纳天主。但它不能容纳造它的那位。他又在哪里被容纳？在安静的灵魂中，在正义的灵魂中：那就是容纳他的。……那说过\BibleQuote{在\Person{亚巴郎}之前，我是。}\BibleRef{若 8:58}的那位——不仅在\Person{亚巴郎}之前，而且在\Person{亚当}之前；不仅在\Person{亚当}之前，而且在所有天使之前，在天地之前；因为万有都是藉他而造的——他接着，免得你注意我们主诞生的日子，以为他那时才开始，\BibleQuote{你的宝座已建立，天主。}但这是什么天主？\BibleQuote{你从永远就有。}他用希腊文说。\par

\SourceRecord{Translated by J.E. Tweed.；Nicene and Post-Nicene Fathers, First Series；Vol. 8.；Edited by Philip Schaff.；Buffalo, NY: Christian Literature Publishing Co.,；1888.；<http://www.newadvent.org/fathers/1801093.htm>.}

% structure: p0001 source=h1 target=section reason=page-title

\section{\WorkTitle{圣咏第九十四篇释义}}

1. 当我们聆听圣咏诵读时，曾以极大的注意，同样，当主揭示祂屈尊在这段经文中隐藏的奥秘时，也让我们留心聆听。因为圣经中有些奥秘之所以被封闭，并非为了被否认，而是为了向敲门者开启。因此，如果你们以虔诚的情感与真诚的内心之爱敲门，那看到你们敲门的动机的主，必会为你们开门。\BibleRef{玛 7:7} 我们众人都知道（但愿我们不在其中），有人抱怨天主的忍耐，因不敬与邪恶的人活在这个世界上，或他们拥有大权而悲伤；更有甚者，坏人通常拥有对抗好人的大权，并且坏人常常压迫好人；恶人欢跃，而好人受苦；邪恶者骄傲，而好人受屈辱。观察到人类中这些事（因为它们比比皆是），不耐烦与软弱的心智遂被扭曲，仿佛为善是徒然的；因为天主转移或似乎转移祂的目光，不看虔敬与忠信者的善工，却提拔恶人于他们所爱的欢愉中。因此，软弱的人，想象自己徒然行善，便被引诱去模仿他们所见的飞黄腾达者的恶行；或者，由于身体或精神的软弱，他们因惧怕世上的刑罚律法而不敢作恶；不是因为爱正义，而是更坦白地说，由于害怕在人间的谴责，他们确实停止了恶行，却没有停止恶念。而在他们的恶念中，首当其冲的是那恶毒，使他们不虔敬地幻想天主疏忽并漠视人间事务：祂要么同等看待善恶，要么——这是一个更加有害的观念——迫害好人，眷顾恶人。如此思想的人，虽不伤害任何人，却最伤害自己，对自己不虔敬，他的恶毒不伤害天主，却杀害自己……\par

2. 这篇圣咏有这样的标题，即这个题词：\Quote{达味自己的一篇圣咏，在一周的第四日。} 这篇圣咏将教导义人受苦中的忍耐：它命令人在恶人兴旺中忍耐，并建立忍耐。这是全篇的主旨，从头到尾。那么，为什么它有这样的标题，\Quote{一周的第四日} 呢？一周的第一日是主日；第二日是第二工作日，世人称之为月亮日；第三日是第三工作日，他们称之为火星日。因此，安息日的第四日就是第四工作日，被外邦人称为水星日，也被许多基督徒如此称呼；但我不会这样称呼：我希望他们能改变得更好，停止这样做；因为他们有自己的用语，他们可以使用。因为这些名称并非通用：许多民族对此有各自不同的名称：所以教会所用的说法更适合基督徒的口。然而，如果习惯使人用舌头说出心里不认同的话，那么让他记住，所有那些名字被星星所承载的人都是人，并且星星并非在那些人开始存在时才开始在天空存在，而是早已存在；但由于对必死者的一些必死服务，那些人在他们的时代，因为他们拥有大权，在此生中卓越，因为他们被人所爱，不是因着永生，而是因着暂时的服务，而接受了神圣的尊荣。因为当时旧世界的人，在被欺骗并想要欺骗时，指着天上的星星，向那些在今生感情上对他们有过好处的人谄媚，说那是某人的星，这是另一个人的星；而那个之前没有看见他们，以至于没有看到那些人出生之前那些星星已经存在的人，就被欺骗而相信；于是这虚妄的意见被构思了。这错误的意见，魔鬼加强，基督推翻了。因此，按照我们的说法，一周的第四日就是从主日起的第四天。因此，亲爱的，请注意这个标题的含义。这里有一个伟大的奥秘，一个真正隐藏的奥秘。……因此，让我们从《圣经》创世纪中回想，第一天创造了什么；我们发现光：第二天创造了什么；我们发现了穹苍，天主称之为天：第三天创造了什么；我们发现了大地和海洋的形态，以及它们的分离，所有水的汇集称为海，所有干地称为地。第四天，主在天上造了光体：\BibleRef{创 1:3} \BibleQuote{太阳管白天，月亮和星辰管黑夜：} 这是第四天的工作。那么，为什么这篇圣咏以第四日作为标题：这篇圣咏中命令人在恶人兴旺和善人受苦中忍耐。你找到宗徒保禄讲话。\BibleRef{斐 2:14} \BibleQuote{你们做一切事，总不可抱怨、争辩，好使你们成为无可指摘和纯洁的，天主的无瑕子女，在弯曲悖谬的世代中，你们在世人中显明如光照，持守生命之言。} ……\par

3. 现在让我们来聆听这篇圣咏。\Quote{主是报复的天主；报复的天主大胆行事} \BibleRef{咏 93:1}。你以为祂不惩罚吗？\Quote{报复的天主}惩罚。什么是\Quote{报复的天主}？即惩罚的天主。你之所以抱怨，肯定是因为恶人未受惩罚：但不要抱怨，免得你成为被惩罚之列。那个人偷了东西，却活着；你抱怨天主，因为偷你东西的人没有死……因此，你若希望别人改正其手，就必须先纠正你的舌头；你希望他改正对人心的态度，你应改正对天主的心；免得当你渴望天主的报复时，若报复来临，却先找到你。因为祂必来临：祂必来临，审判那些继续作恶、对祂仁慈的延长不知感恩、对祂的忍耐不知感恩、为自己积蓄忿怒以待忿怒之日，并显明天主公义审判的人，祂将按各人的行为报应各人：\BibleRef{罗 2:4}因为，\Quote{主是报复的天主}，所以祂\Quote{大胆行事}。……我们的安全是我们的救主；祂愿将一切穷困贫乏者的希望寄托于祂。祂说什么？\Quote{我要在祂内大胆行事。}这是什么意思？祂不害怕，不姑息人的情欲和恶习。确实，如同一位忠信的医生，手中拿着宣讲的治愈之刀，祂切除了我们所有的创伤。因此，祂先前被预言和宣讲的那样，祂被发现就是那样。……那么，那被说成\Quote{祂教训他们，像有权威的人}的那一位，向他们说了何等伟大的话？\Quote{祸哉，你们经师和法利塞人，假善人！}祂当着他们的面说了何等伟大的话？祂不惧怕任何人。为什么？因为祂是报复的天主。为此，祂在言语上不宽恕他们，好使他们后来能被祂在审判中宽恕；因为他们若不愿接受祂言语的医治，后来就会招致审判者的刑罚。为什么？因为祂说了：\Quote{主是报复的天主，报复的天主大胆行事；}也就是说，祂在言语上没有宽恕任何人。那在将要受苦时在言语上不宽恕人的，在将要审判时会在审判中宽恕吗？那在谦卑时不怕任何人的，在光荣中会怕任何人吗？从祂过去这样大胆行事，想象祂在末世将如何行事。那么，不要抱怨天主，那似乎宽恕恶人的天主；但你要做好人，也许暂时祂不会宽免你的杖责，好使祂最终在审判中宽恕你……\par

4. 接下来发生了什么，因为祂大胆行事？\Quote{祢被高举，审判世界者} \BibleRef{咏 93:2}。当他们逮捕了谦卑的祂，你以为他们会逮捕被高举的祂吗？当他们审判了必死的祂，他们不会被不朽的祂审判吗？那么祂说什么？\Quote{祢被高举，}祢，那大胆行事的，祂的话的胆量恶人无法承受，反而以为他们做了一件光荣的事，当他们抓住并钉死了祢；他们本该以信德抓住祢，却以迫害抓住了祢。祢，那在恶人中大胆行事、不惧怕任何人的，因为祢受了苦，\Quote{被高举；}也就是说，复活，升天。让教会也以忍耐承受教会元首以忍耐所承受的。\Quote{祢被高举，审判世界者：按他们的应得报应骄傲的人。}祂将报应他们，弟兄们。因为这是什么，\Quote{祢被高举，审判世界者：按他们的应得报应骄傲的人}？这是预言者的预言，而非命令者的胆量。不是因为先知说\Quote{祢被高举，审判世界者}，基督才听从先知，从死者中复活，升天；而是因为基督要这样做，先知预言了它。他在精神中看见基督受贬抑，看见祂受贬抑：不惧怕任何人，在言语上不宽恕任何人，他说：\Quote{祂大胆行事。}他看见祂何等大胆行事，看见祂被捕、被钉、受辱，看见祂从死者中复活，升天，并从那里来审判那些在祂手中遭受一切邪恶的人：\Quote{祢被高举，}他说，\Quote{审判世界者，按他们的应得报应骄傲的人。}骄傲的人祂将这样报应，而非谦卑的人。谁是骄傲的人？那些作恶还算小事，甚至为自己辩护罪过的人。因为在那些钉死基督的人中，后来有奇迹发生，当犹太人中有了一些信者，基督的血也赐给了他们。他们的手是不敬的，染着基督的血。祂，他们所倾流的血，亲自洗净了他们。那些迫害了祂可见的、必死身体的人，成了祂身体的肢体，即教会。他们倾流了自己的赎价，为要饮用自己的赎价。因为后来有更多的人归化……\par

5. \Quote{主，恶人要到几时，恶人要到几时夸口？} \BibleRef{咏 93:3}。 \Quote{他们回答，并要说恶言，他们都说不义的话。} \BibleRef{咏 93:4}。他们说的话岂不是反对天主吗？当他们说：\Quote{这样活着对我们有什么益处呢？} 你会怎样回答？难道天主真的顾念我们的行为吗？因为他们活着，就以为天主不知道他们的行为。看，他们遭遇了什么恶事！如果官员知道他们在哪里，就会逮捕他们；所以他们避开官员的目光，以免立刻被捉住；但没有人能逃脱天主的眼目，因为祂不仅看见密室之内，还看见心中的隐密。甚至他们自己也相信没有什么能逃脱天主；而因为他们作恶，并且意识到自己所做的事，又看见自己活着，而天主知道，虽然如果官员发现他们，他们就不能活，他们就对自己说：\Quote{这些事是取悦天主的。} 的确，如果这些事不取悦祂，正如它们不取悦君王、不取悦审判官、不取悦总督、不取悦记录官一样，那么我们难道能逃脱天主的眼目，如同我们逃脱那些权贵的眼目吗？因此，这些事是取悦天主的……有一位义人前来，说：\Quote{不要行不义。} \Quote{为什么？} \Quote{以免你死亡。} 看，我行了不义：为什么我不死呢？那人行了正义：他却死了：为什么他死了？我行了不义：为什么天主没有把我带走？看，那人行了正义：但为什么祂这样对待他？为什么祂容许这样？他们回答；这就是\Quote{回答}一词的含义：因为他们有可答辩的理由；由于天主的宽仁使他们得以存活，他们便为自己的答辩找到了论据。祂宽恕他们是为一个理由，他们却为另一个理由而答辩，因为他们仍然活着。因为宗徒告诉我们祂为何宽恕，他解释了天主宽仁的缘由：\Quote{你这人哪，你论断行这样事的人，自己所行的却和别人一样，你以为你能逃脱天主的审判吗？还是你藐视祂丰富的恩慈、宽容和忍耐，不晓得祂的忍耐是领你悔改呢？} \Quote{但你，} 这就是说，那答辩并说：\Quote{如果我不取悦天主，祂就不会宽恕我} 的人，听一听他为自己积蓄了什么；听宗徒的话：\Quote{你竟任着你刚硬不悔改的心，为自己积蓄忿怒，以致震怒的日子，显祂公义审判的日子来到；祂要照各人的行为报应各人。} \BibleRef{罗 2:5}。祂因此增加祂的宽仁，你增加你的不义。祂的宝藏将是对那些没有藐视祂仁慈的人永远的慈悲；但你的宝藏将在忿怒中被发现，你每日一点一点积蓄的，将会在累积的全集中找到；你一粒一粒地积蓄，却会发现整个的一堆。不要忽略注视你每日最小的罪：江河是由最小的水滴汇成的。\par

6. ……\Quote{主，他们压制了你的百姓，苦害了你的产业} \BibleRef{咏 93:5}。 \Quote{他们杀害寡妇和孤儿，又杀了外方人} \BibleRef{咏 93:6}；即旅行者、朝圣者：从远方来的人，正如圣咏作者自称。这些表述的意思都太清楚，不值得详述。\par

7. \Quote{他们曾说：主必不看见} \BibleRef{咏 93:7}：祂不留意，不看顾这些事；祂关心别的事，祂不明白。这是恶人的两种断言：一种是我刚才引用的，\Quote{你行了这些事，我默不作声，你就以为不义，以为我像你一样。} \Quote{我像你一样}是什么意思？你以为我看见你的行为，并且这些事令我喜欢，因为我并没有惩罚他们。恶人还有另一种断言：因为天主既不留意这些事，也不观察，好知道我是如何生活的，天主不关心我。那么，天主难道会眷顾我吗？或者祂甚至计算我吗？或计算世人吗？不幸的人！祂曾关心你，使你存在；祂难道不关心你好好生活吗？因此，这些人的话是这样：\Quote{他们曾说：主必不看见，雅各伯的天主必不留意。}\par

8. \BibleQuote{你们愚昧的人，现在要注意：你们这些糊涂人，什么时候才能明白！}\BibleRef{咏 93:8}。祂教导祂的百姓，他们的脚可能滑跌：他们中的任何人看见恶人的发达，而自己生活在天主的圣人们中，即在教会的子女中，生活美好；他看见恶人兴旺，行不义，便嫉妒，并被引诱跟从他们的行为；因他看见，似乎他谦卑地生活、盼望在此世得到赏报，对他毫无益处。因为如果他盼望将来得到，他就不会失去；因为领受它的时间尚未来到。你在葡萄园工作：完成你的工作，你将得到你的工钱。你不会在完工之前向你的雇主索取工钱，然而你却在工作之前向天主索取工钱吗？这忍耐是你工作的一部分，你的工钱取决于你的工作：你若不选择忍耐，便选择在葡萄园做更少的工作，因为这忍耐的行为属于你的劳作本身，而劳作是要赚取工钱的。但如果你不忠，你要小心，免得你不仅得不到工钱，还要受罚，因为你选择做一个不忠的工人。当这样的工人开始作恶时，他观察雇佣他到葡萄园的雇主的目光，以便当雇主的目光移开时，他偷懒；但雇主的目光一转向他，他就勤奋工作。但是天主，雇佣你的那位，不会移开祂的目光：你不能不忠地工作：你主人的目光常在你身上：你找机会欺瞒祂，如果你能，就偷懒吧。如果你们中有人有过这样的想法，当你看见恶人兴旺，若这样的想法使你的脚在天主的道路上滑跌；这篇圣咏对你说话：但倘若你们中没有这样的人，那么它藉着你向别人说话，用这些话：\BibleQuote{现在要注意；}因为他们曾说过，\BibleQuote{上主看不见，雅各伯的天主毫不介意。}\BibleQuote{要注意，}它说，\BibleQuote{现在，你们愚昧的人，你们这些糊涂人，什么时候才能明白！}\par

9. \BibleQuote{那造耳朵的，难道不听见？那造眼睛的，难道不察看？}\BibleRef{咏 93:9} \BibleQuote{那管教万民的，难道不责备？}\BibleRef{咏 93:10}。这是天主现在所做的事：祂正在教导万民；为此祂向普世的人发送祂的圣言：祂藉着天使、圣祖、先知、仆人，通过如此多的走在审判者前面的先驱发送了它。祂也发送了祂自己的圣言本身，祂亲自发送了祂自己的圣子：祂发送了祂圣子的仆人，并在这些仆人自身中发送了祂自己的圣子。天主的圣言在普世到处被宣讲。哪里没有对人说：抛弃你们过去的邪恶，转向正路？祂宽恕，为叫你们改正自己：祂昨天没有惩罚，为使你们今天能生活得好。祂教导外邦人，难道祂因此不责备？难道祂不听祂所教导的人？难道祂不审判那些祂事先已派遣并播下警告教训的人？如果你在学校里，你会接受一个功课而不复述吗？因此当你从你的老师那里接受了它，你正在被教导：老师把功课交到你手中，难道他不会在你来复述时向你索取吗？或者当你开始复述时，你不害怕受鞭打吗？所以现在我们在接受我们的工作：以后我们被带到老师面前，以便我们把过去所有的功课交还给祂，就是，我们为现在所赐给我们的一切事交付账目。听宗徒的话：\BibleQuote{我们都必须出现在基督的审判台前，}等等。\BibleQuote{是祂教导人知识。}难道祂不知道，祂使人知道吗？\par

10. \BibleQuote{主知道人的思念，原来都是虚妄}\BibleRef{咏 93:11}。因为虽然你不知道天主的思念，它们是公义的；但\BibleQuote{祂知道人的思念，原来都是虚妄。}甚至人也曾知道天主的思念：但是那些祂与之成为朋友的人，祂向他们显示祂的计谋。兄弟们，不要轻视自己：如果你以信德接近主，你就听到天主的思念；这些你现在正在学习，这是告诉你的，并且你受教导正是为此，为什么天主在此世宽恕恶人，好使你不抱怨那教导人知识的天主。\BibleQuote{主知道人的思念，原来都是虚妄。}所以，抛弃人的思念，因为是虚妄的：好使你能持守天主的思念，因为是明智的。但谁能持守天主的思念？那被安置在天上穹苍的人。我们已经咏唱了那篇圣咏，并在其中解释了这话。\par

11. \BibleQuote{主啊，祢所管教、用祢的法律教训的人，是有福的} \BibleRef{咏 93:12}。请看，你已知道天主的旨意，为何祂容忍恶人：坑正在为罪人挖掘。你想立刻埋葬他：坑还在为他挖掘中，不要急于埋葬他。这话是什么意思：\Quote{直到为罪人挖好了坑}？祂所说的罪人是指谁？一个人吗？不是。那么是谁？所有这类罪人的种族？不是；而是骄傲的人；因为祂之前说过：\Quote{照他们所应得的报应骄傲的人}。因为那个税吏，他甚至不敢举目望天，但\BibleQuote{捶着胸膛说：天主，可怜我这个罪人吧} \BibleRef{路 18:13}，他也是一个罪人；但他并不骄傲，既然天主将报应骄傲的人，那么坑就不是为他挖的，而是为那些像他那样的人，直到祂报应骄傲的人。因此，在\Quote{直到为不虔敬的人挖好了坑}这句话中，要理解骄傲的人。谁是骄傲的人？他不通过告解自己的罪过做补赎，以便通过谦卑得到痊愈。谁是骄傲的人？他选择把那些他看来拥有的一点点好东西归功于自己，并贬低天主的仁慈。谁是骄傲的人？他虽将善行归于天主，却侮辱那些不行善的人，并抬高自己超过他们……这就是基督徒的教义：没有人能靠自己的恩宠行善。人的恶行是他自己的；善行来自天主的恩赐。当他开始行善时，不要归功于自己；当他不归功于自己时，要感谢赐予他的那一位。但当他行善时，不要侮辱那些不像他那样做的人，也不要抬高自己超过他：因为天主的恩宠并不停留在他身上，以至于不能达到另一个人。\par

12. \BibleQuote{祢使他在患难之日有忍耐，直到为不虔敬的人挖好了坑} \BibleRef{咏 93:13}。所以，每个人，如果你是基督徒，在邪恶的时期都要忍耐。邪恶的日子是那些不虔敬的人看起来兴旺，而义人受苦的日子；但义人的受苦是圣父的管教，不虔敬的人的兴旺是他们自己的陷阱。因为天主在患难中赐你忍耐，直到为不虔敬的人挖好了坑，你不要以为天使们正拿着镐站在某处，挖掘那个足以容纳所有不虔敬种族的大坑；因为你们看到恶人很多，就按肉性对自己说：真的，什么坑能容纳如此众多不虔敬的人，如此众多的罪人？这么大的坑在哪里？什么时候挖完？因此天主容忍他们。并非如此：他们的兴旺本身就是不虔敬的人的坑：因为他们将落入其中，如同落入陷阱。弟兄们，请注意，因为知道兴旺被称为陷阱是一件大事：\Quote{直到为不虔敬的人挖好了坑}。因为天主以祂隐藏的公义容忍祂所知道的不虔敬和邪恶的人；而天主的这种容忍，由于不受惩罚，使他骄傲起来……骄傲的人抬高自己反对天主；天主使他沉没：他因抬高自己反对天主而沉没。因为另一位圣咏作者这样说：\Quote{当他们被高举时，祢推倒了他们}。他不是说：因为他们被高举，祢推倒他们；或者说：他们被高举之后，祢推倒他们，以使他们的高举时期与推倒时期分开；而是在他们被高举的过程中，他们被推倒了。因为人心越骄傲，就越远离天主；若远离天主，就沉入深渊。另一方面，谦卑的心将天主从天上带下来，使祂非常接近它。的确，天主崇高，天主超越诸天，超越所有天使：这些（骄傲的人）要升得多高才能达到那位崇高者？不要因扩张自己而破裂；我给你另一个建议，免得你在扩张时因骄傲而破裂：天主崇高；你谦卑自己，祂就会降临到你身上。\par

13. ……你要在鞭打下喜乐，因为基业为你存留，\BibleQuote{因为主必不丢弃祂的百姓} \BibleRef{咏 93:14}。祂暂时管教，永不弃绝；而那些（恶人）祂暂时容忍，却要永远惩罚。你选择吧：你愿意暂时的痛苦，还是永罚？暂时的幸福，还是永生？天主威胁的是什么？永罚。祂应许的是什么？永恒的安息。祂对善人的鞭策是暂时的；祂对恶人的容忍也是暂时的。\Quote{祂也不离弃祂的产业。}\par

14. \\BibleQuote{直到义德，}他说，\\BibleQuote{转向审判，凡拥有它的人都是心地正直的}\\BibleRef{咏 93:15}。现在请听，要获得义德：因为你尚未能有审判。你应先获得义德；但你的这义德将要转向审判。宗徒们在世上有了义德，并容忍了恶人。但对他们说了什么？\\BibleQuote{你们要坐在十二个宝座上，审判以色列十二支派}\\BibleRef{玛 19:28}。因此他们的义德将转向审判。因为凡在此生是义人的，正是为此，他才能以忍耐承受邪恶：让他耐心承受受苦的时期，审判的日子就会到来。但我为何谈到天主的仆人？主自己，就是那审判所有活人和死人的，首先选择被审判，然后审判。那些目前拥有义德的人，还不是审判者。因为首先是拥有义德，然后审判：他先容忍恶人，然后审判他们。现在要有义德：以后它将转向审判。祂容忍恶人多久，如同天主所愿意的，只要天主的教会容忍他们，好使她通过他们的邪恶受到教导。然而，天主不会抛弃祂的子民，\\BibleQuote{凡拥有它的人都是心地正直的}。谁是那些心地正直的人？那些意志符合天主意志的人。祂宽恕罪人：你希望祂立即毁灭罪人。你的心是弯曲的，你的意志是乖僻的，当你的意志是一样，天主的意志是另一样时。天主愿意宽恕罪人：你不愿意宽恕罪人。天主对罪人长久忍耐：你不愿意容忍罪人……不要企图将天主的意志屈就于你的意志，而应纠正你的意志以符合祂的。天主的意志就像一把尺：看，假设你扭曲了尺：你怎能被矫正？但尺本身保持笔直：因为它是不变的。只要尺是直的，你就有地方可以转回，纠正你的乖戾；你就有方法纠正你心中的弯曲。但人们愿意什么呢？他们自己的意志弯曲还不够；他们甚至希望根据他们自己的心使天主的意志也弯曲，好让天主做他们自己愿意的事，而他们本应做天主所愿意的事……\par

15. \\BibleQuote{谁能为我起来攻击恶人？谁肯为我站立攻击作恶的人？}\\BibleRef{咏 93:16}。许多人劝我们行各种邪恶：蛇不停地向你低语，要你行不义；无论你转向何方，如果你偶尔做了好事，你寻求与某人好好生活，却几乎找不到任何人；许多恶人围绕你，因为麦粒少，糠秕多。这打谷场有其谷粒，但如今它们还在受苦。因此，整个麦堆，当与糠秕分离后，将是巨大的：谷粒虽少，但与糠秕相比，它们本身仍很多。因此当恶人从四方喊叫，说，你为什么这样生活？你是唯一的基督徒吗？你为什么不像别人那样做？你为什么不像别人那样常去剧场？为什么你不用护符和咒符？你为什么不像别人那样咨询占星家和占卜者？而你画十字，说，我是基督徒，以便你抵挡他们，无论他们是谁；但敌人紧逼，催迫攻击；更糟的是，他借着基督徒的榜样扼杀基督徒。他们劳苦工作，在炎热中：基督徒的灵魂遭受磨难：然而它有力量得胜：它自己有这样的力量吗？因此请注意他所说的话。因为他回答，我现在为自己找到护符，多得几天日子，对我有什么益处？我离开今世，往我主那里去，祂要把我投入火焰中；因为我选择了几天短暂时日胜过永生，祂要把我投入地狱。什么地狱？就是天主永恒审判的地狱。真的如此吗（敌人回答），除非你真的相信天主关心人如何生活？也许不是街上一个熟人这样对你说话，而是你家里的妻子，或可能是丈夫对忠信圣洁的妻子，她的欺骗者。如果是女人对丈夫，她之于他如同\\Person{厄娃}；如果是丈夫对妻子，他之于她如同魔鬼：要么她之于你如同\\Person{厄娃}，要么你之于她如同蛇。有时父亲会向儿子倾诉他的想法，却发现他邪恶、完全堕落：他陷入悲惨的热病中，他摇摆不定，他寻求如何制服他，他几乎被拖进去，同意了：但但愿天主与他相近……\par

16. \\BibleQuote{若非主帮助我，}他说，\\BibleQuote{我的灵魂几乎已住在地狱。}\\BibleRef{咏 93:17}。我几乎投入那为罪人所预备的坑中：也就是说，我的灵魂已住在地狱。因为他已经开始动摇，几乎同意，他回转向主。举例来说，假设他受到侮辱，被引诱行不义。因为有时恶人聚集在一起，侮辱好人；尤其当他们人数更多，并且单独抓住他时，就像常有大量糠秕围绕一粒麦子（虽然当堆被簸扬后就不会了）；他于是被许多恶人围住，受侮辱，被包围；他们想要凌驾于他之上，他们折磨他，因他的义德而侮辱他。他们说，伟大的宗徒！你已像\\Person{厄里亚}那样飞升天堂！人们做这些事，以致有时，当他听从人的言语时，他在恶人中以行善为耻。因此让他抵抗邪恶；但不要凭自己的力量，免得他变得骄傲，而当他想逃避骄傲的人时，自己反而增加了他们的数目……\par

17. \BibleQuote{我若说：我的脚滑跌了，祢的仁慈，主，扶持了我} \BibleRef{咏93:18}。请看天主多么喜爱承认。你的脚滑跌了，你却不说是我的脚滑跌了；反而当自己滑跌时，却说你是稳固的。当你开始滑跌或动摇时，就该承认那滑跌，免得你哀叹完全的跌倒；好使祂帮助你，使你的灵魂不在地狱里。天主喜爱承认，喜爱谦卑。你滑跌了，如同人一样；天主仍然帮助你：然而你要说：\BibleQuote{我的脚滑跌了。}你为何滑跌，却说你是稳固的呢？\BibleQuote{我若说：我的脚滑跌了，祢的仁慈，主，就扶持了我。}正如伯多禄自以为稳妥，却不是倚靠自己的力量。主曾显现在海面上行走，践踏此世所有骄傲人的头。祂行走在泡沫翻腾的波浪上，预示了祂自己的道路，即祂践踏骄傲人的头。教会也践踏他们：因为伯多禄就是教会本身。然而，伯多禄不敢独自在水面上行走；他说了什么？\BibleQuote{主，如果是祢，就请祢叫我从水面上走到祢那里去。}\BibleRef{玛14:28}祂凭自己的能力，伯多禄凭祂的命令；他说：\BibleQuote{请祢叫我}，\BibleQuote{到祢那里去。}祂回答说：\BibleQuote{来吧。}因为教会也践踏骄傲人的头；但由于她是教会，具有人性的软弱，为使这些话得到应验：\BibleQuote{我若说：我的脚滑跌了，}伯多禄在海面上摇晃，呼喊说：\BibleQuote{主，救我！}\BibleRef{玛14:30}所以此处所说的\BibleQuote{我若说：我的脚滑跌了，}在那里便是\BibleQuote{主，我沉没了。}而此处所说的\BibleQuote{祢的仁慈，主，扶持了我，}在那里便是\BibleQuote{耶稣立刻伸手拉住他，说：你这小信德的人，为什么怀疑？}\BibleRef{玛14:31}天主如何考验人，真是奇妙：我们的危险使那位拯救我们的那一位对我们更为甘饴。请看下文：因为他曾说：\BibleQuote{我若说：我的脚滑跌了，祢的仁慈，主，扶持了我。}主在救他脱离危险时，对他变得格外甘饴；因此，提到主的这份甘饴，他呼喊说：\BibleQuote{主，我心中有许多忧愁时，祢的安慰舒畅了我的灵魂}\BibleRef{咏93:19}。许多忧愁，却有许多安慰；苦涩的创伤，甘饴的医治。\par

18. \BibleQuote{你与邪恶的座位岂能有份？你制造痛苦作为教训。}\BibleRef{咏93:20}。他说这话的意思是：没有恶人与你同坐，你也不与邪恶的座位有份。他说明自己如何理解这话：\BibleQuote{因为你制造痛苦作为教训。}由此，因为你没有宽恕我们，我才明白你与邪恶的座位无份。这在宗徒伯多禄的书信中也有，为此他引用了圣经的见证：他说：\BibleQuote{因为时候到了，审判必须从天主的家开始；}也就是说，属于天主之家的人受审判的时候到了。如果儿子受鞭打，最邪恶的奴仆该期待什么呢？因此他接着说：\BibleQuote{如果它先从我们开始，那么那些不听从天主福音的人，结局将是怎样的呢？}他加上了这个见证：\BibleQuote{因为如果义人仅仅得救，那不虔敬和犯罪的人将在哪里出现呢？}\BibleRef{伯前4:17}。那么，邪恶之人怎能与你有份，如果你甚至不宽恕你的忠信者，为要锻炼和教导他们？\BibleRef{箴11:31}。但祂不宽恕他们，正是为了教导他们：他说：\BibleQuote{因为你制造痛苦作为教训。}\BibleQuote{制造}，即塑造：由此产生\OriginalTerm{陶工}{figulus}一词源自\OriginalTerm{塑造}{fingo}，陶匠的器皿称为\OriginalTerm{陶制的}{fictile}：不是指虚构，即谎言的意思，而是指塑造，赋予事物存在和某种形式；正如他先前所说：\BibleQuote{那\OriginalTerm{塑造了}{finxit}眼睛的，难道不能看见吗？}难道\BibleQuote{塑造了眼睛}是谎言吗？不，这话意指他造了眼睛，做了眼睛。当他造人脆弱、软弱、属地时，他难道不也是一位陶工吗？请听宗徒的话：\BibleQuote{我们在这瓦器里有这宝藏。}\BibleRef{格后4:7}……请看我们的主自己，如何显示自己是一位陶工。\BibleRef{罗4:20}因为祂曾用泥土造人，祂用泥土涂抹了那在母腹中没有眼睛的人。因此当他说：\BibleQuote{你与邪恶的座位岂能有份？}等等时，他说，你从痛苦中为我们制造教训，以致痛苦本身成为我们的教导。痛苦如何成为我们的教训？当祂鞭打那为你而死、且不在此世应许福乐、也不能欺骗的那一位，当祂不在此处赐给你所寻求的时。祂将赐给什么？祂何时赐给？那位此刻不赐予、此刻教导、制造痛苦作为教训的，将赐给多少？你的劳苦在此处，安息已应许给你。你思虑你在此处有辛劳：但你要想一想祂应许的是怎样的安息。你能想象吗？如果你能，你就会看到你在此处的劳苦与那安息相比，算不得什么……\par

19. 注意，弟兄们：这是待售品。天主对你们说，我所拥有的待售，你们买吧。祂卖什么？我有安息待售；用你们的辛劳来买它。注意，我们要因基督之名成为勇敢的基督徒：这篇圣咏所余不多，让我们不要疲倦。因为那在听上失败的人，如何能在行上坚强？主将帮助我们为你们阐释余下的内容。那么注意：天主仿佛已将天国宣告待售。你们对祂说，它的价值是什么？代价是辛劳：如果祂说，代价是金子，仅这样说还不够，你们会想知道多少金子；因为有一块金子、半盎司、一磅等。祂说了\Quote{代价}，免得你们费心去查问要多久才能找到它。商品的代价是辛劳：辛劳是多少？现在你们要寻求应为之付出多少辛劳。尚未告诉你们那辛劳注定多大，或要求你们多少辛劳：天主这样对你们说，我向你们显示那安息将有多大；你们判断应以多大尺度的辛劳去买它。\par

20. ……祂应许了安息：要忍受苦难。祂以永火警告；轻视暂时的痛苦：当基督警醒时，愿你们的心平静，使你们也能到达港湾。因为那预备了船只的，必不会不预备港湾。\Quote{祢与邪恶的宝座有何相干？祢在训诲中制造忧伤。}祂以恶人试探我们，藉他们的迫害教导我们。借着恶人的恶意，善人被鞭策，藉着奴仆，儿子被管教：这样，训诲通过忧伤得以传授。天主允许他们有权做什么，恶人就做什么，祂暂时宽容他们。\par

21. 接着是什么？\BibleQuote{他们攻击义人的灵魂} \BibleRef{咏 93:21}。他们为什么挑剔？因为他们找不到真实的控告理由。他们如何挑剔我们的主？他们捏造虚假的控告，\BibleRef{玛 26:59}，因为他们找不到真实的。\BibleQuote{并定无辜之血的罪。}为何这一切发生，他将在后续揭示。\par

22. \BibleQuote{上主作了我的避难所} \BibleRef{咏 93:22}，他说。你们若不处在危险中，就不会寻求这样的避难所；但你们之所以处于危险，正是为了寻求它：因为祂借着忧伤教导我们。祂从恶人的恶意中使我遭受磨难；被那磨难刺痛，我开始寻求一个避难所，这避难所在世俗的顺遂中我已停止寻求。因为那总是顺遂、并因现世希望而喜乐的人，有谁容易记起天主？让此生的希望退让，让天主的希望前进；使你们可以说：\BibleQuote{上主作了我的避难所}；愿我为此忧伤，好使上主成为我的避难所！\BibleQuote{我的天主是我望德的扶助}。因为上主目前是我们的望德，只要我们还在此世，我们就处于望德中，而非拥有。但免得我们在望德中失败，在我们附近有一种供应来鼓励我们，并减轻我们所受的诸恶。因为并非徒然地说：\BibleQuote{天主是信实的，必不让你们受试探超过你们所能承受的；但在受试探的同时，也要开辟一条出路，使你们能忍受得住} \BibleRef{格前 10:13}；祂将如此把我们放入磨难的熔炉中，使器皿得以坚固，而不至于破碎。\BibleQuote{上主作了我的避难所：我的天主是我望德的扶助}。那么，为什么祂在你看来好像不正义，宽容恶人呢？看，这圣咏如何被纠正，你也要与圣咏一起被纠正：因为，为此圣咏包含了你的话语。什么话语？\BibleQuote{上主，不虔敬的人要到几时，要到几时夸耀？}圣咏刚才用了你的话语：因此，你也要反过来使用圣咏的话语。\par

23. \BibleQuote{上主必按他们的行径报复他们，按他们自己的恶意；上主我们的天主必消灭他们} \BibleRef{咏 93:23}。所说的\Quote{按他们自己的恶意}并非没有意义。我因他们而受益：但这里说的是他们的恶意，而非益处。因为祂确实借着恶人试炼我们、鞭策我们。祂鞭策我们是为预备我们做什么？无疑是为了天国。\BibleQuote{因为祂鞭打所收纳的每一个儿子；哪个儿子是父亲不惩戒的呢？} \BibleRef{希 12:7}。当天主这样做时，祂是在教导我们，为得到永恒的产业；这种训练祂常常借着恶人来给予我们，借着他们，祂试炼并使我们的爱德得以成全，这爱德祂愿意扩展到甚至我们的仇敌。……同样，那些迫害殉道者的人，在地上迫害他们，却把他们送进了天国：他们明知造成了他们现世生命的损失，却无意中赐予了他们未来生命的增益；但是，尽管如此，对所有坚持对义人怀有邪恶仇恨的人，天主必将按他们自己的过犯偿还，并将在他们自己的恶意中消灭他们。因为义人的善行对恶人有害，同样，恶人的不义对义人有益。……\par

24. 因此，让义人容忍不虔敬的人；让义人暂时的苦难容忍恶人暂时的不受惩罚；因为\BibleQuote{义人因信德而生活} \BibleRef{罗 1:17}。因为此世没有人的正义，除非因信德而生活，这信德\BibleQuote{藉爱德而运作} \BibleRef{迦 5:6}。如果他因信德而生活，就让他相信：他自己在现世辛劳之后将继承安息，而他们将在现世欢腾之后遭受永罚。如果信德藉爱德运作，就让他也爱他的仇敌，并且尽己所能，有愿意造福他们的心；因为这样，当仇敌有意伤害时，他将阻止他们伤害他。每当他们可能得到伤害和欺压的能力时，让他高举他的心灵于高处，那里无人能伤害他，他要在天主的法律中受良好教导和惩戒，使他\Quote{在患难之日得到忍耐，直到为不虔敬的人挖好了坑}。……\par

25. 弟兄们，我说这话，是要你们从所听的获益，并在内心反复思想：不要让自己忘记，不仅借着反复思考这些事并谈论它们，也借着如此生活。因为按照天主的诫命度过的善生，就像一支笔，因为它听见在写在我们心中。若写在蜡上，容易擦去；写在你们心中，写在你们的品性上，就永不擦去。\par

\SourceRecord{Translated by J.E. Tweed.；Nicene and Post-Nicene Fathers, First Series；Vol. 8.；Edited by Philip Schaff.；Buffalo, NY: Christian Literature Publishing Co.,；1888.；<http://www.newadvent.org/fathers/1801094.htm>.}

% structure: p0001 source=h1 target=section reason=page-title

\section{圣咏第九十五篇释义}

1. 弟兄们，我宁愿我们是在听我们的父亲讲话：但这也是一件好事，听从我们的父亲。既然他屈尊为我们祈祷，已命令了我们，我要向你们，诸位可爱的，讲论从眼前的圣咏中，耶稣基督我们共同的主俯允赐予我们的事。现在这圣咏的标题是\Quote{达味的赞美之歌}。\Quote{赞美之歌}既表示喜乐，因为它是一首歌；又表示虔敬，因为它是赞美。因为一个人应当赞美什么，更甚于那使他如此喜悦，以致不可能不令他喜悦的事物呢？因此，在赞美天主时，我们安然赞美。在那里，赞美者是安全的，他不怕因他赞美的对象而羞愧。所以让我们既赞美又歌唱；也就是说，让我们以喜乐和欢欣赞美。但我们所要赞美的，这圣咏在接下来的经文中向我们表明。\par

2. \Quote{哦，来，让我们向主歌唱} \BibleRef{咏 94:1}。祂召唤我们参加一个伟大的喜乐盛宴，不是这世界的，而是在主内。因为如果在这生命中不存在需要与正义的喜乐相区分的邪恶之乐，那么说\Quote{来，让我们欢乐}就足够了；但祂已经简要地加以区分。什么才是正确地欢乐？就是在主内欢乐。你若愿安全地践踏世界，就当虔诚地喜乐于主。但\Quote{来}这个词是什么意思？祂从哪里召唤他们来，祂愿意与谁一同在主内欢乐？除非，当他们尚在远方时，他们可以借着来更靠近，借着更靠近而接近，借着接近而喜乐？但他们从何处是远方？一个人能从无所不在的祂那儿在空间上遥远吗？……不是因为地方，而是因为不像祂，一个人离天主遥远。不像祂是什么意思？就是品行恶劣，习惯败坏；因为如果我们因好习惯接近天主，就因坏习惯远离天主。……因此，如果我们因不像而远离天主，就要因像而接近天主。什么像？就是我们受造所依据的那肖像，因犯罪而腐化在我们内的，又通过罪的赦免而重新领受的，在我们内在的心智中得以更新，以便如同刻在钱币上第二次，即我们天主的肖像刻在我们灵魂上，使我们能回到祂的库藏。……\par

3. \Quote{让我们向天主、我们的救恩欢呼。}……亲爱的，请想想那些在普通歌曲中欢呼的人，仿佛置身于世俗欢乐的竞赛中；你看见他们在背诵写下的词句时，迸发出欢乐，以至于舌头不足以表达其程度；他们如何呐喊，用那喊声表示内心的感受，而内心所领悟的，言语无法表达。如果他们在世俗的欢乐中尚且欢呼，难道我们不该因天上的欢乐而欢呼吗？这欢乐我们实在无法用言语表达！\par

4. \Quote{让我们以宣认来到祂面前} \BibleRef{咏 94:2}。宣认在圣经中有双重含义。有赞美者的宣认，有哀叹者的宣认。赞美的宣认属于被赞美者的荣耀：哀叹的宣认属于宣认者的忏悔。因为人们赞美天主时宣认；他们指控自己时宣认；舌头没有更高贵的用途。实在，我相信这正是那誓愿，他在另一篇圣咏中说到：\Quote{我要向祢偿还我的誓愿，就是我以嘴唇所分别的。}没有比这分别更崇高的，没有比理解和实行更必要的。那么，你如何分别你向天主偿还的誓愿呢？借着赞美祂，借着指控你自己；因为祂的仁慈是宽恕我们的罪。因为如果祂选择按我们的功过对待我们，祂只会找到定罪的缘由。因此祂说\Quote{哦，来}，使我们终于能远离我们的罪，使祂不再与我们清算以往的账目；而是仿佛开始一个新账，我们一切债务的债券已被烧毁。……因此，你越因自己的过犯而绝望，就越要承认你的罪；因为赦免者的赞美越大，就如忏悔者宣认的丰盛越加充实。因此，不要以为我们在此将宣认理解为承认我们的过犯，就远离了赞美之歌：这甚至是赞美之歌的一部分；因为当我们承认自己的罪时，我们赞美天主的荣耀。\par

5. \Quote{并以圣咏向祂欢呼。} 我们已经说过什么是\Quote{欢呼}：这词被重复，好使行动得以确认；这重复本身就是劝勉。因为我们并未忘记，以致需要再次被提醒上面所说的话，即我们应当欢呼。但通常在强烈情感的段落中，一个熟悉的词会被重复，不是为了让它更熟悉，而是为了通过重复本身加强印象；重复是为了让我们理解说话者的情感……现在请听：\Quote{因为主是伟大的天主，是超越众神的伟大君王} \BibleRef{咏 94:3} \Quote{因为主必不遗弃祂的百姓。}愿赞美归于祂，愿欢乐的呼声归于祂！祂不会遗弃什么百姓？我们无权在此自行解释；因为宗徒已为我们规定，他已说明了这话指的是什么。这百姓就是犹太百姓，就是有先知们的百姓，就是有族长们的百姓，就是按血统从亚巴郎的后裔所生的百姓；就是拥有所有预示我们救主的前兆奥秘的百姓；就是在其中设立了圣殿、傅油、作为预像的司铎的百姓，好使所有这些影子过去后，光明本身来临。因此，这曾是天主的百姓；先知们被派往其中，被派遣的人在它里面诞生；天主的启示被交托并托付给它。那么怎样？整个那个百姓都被判罪了吗？决不是。宗徒称它为好橄榄树，因为它始于族长们……这就是那棵树本身；虽然有些枝条被折断了，但并非全部。因为如果所有枝条都被折断，伯多禄从哪里来？若望从哪里来？多默从哪里来？玛窦从哪里来？安德肋从哪里来？所有那些宗徒从哪里来？就是那位刚才向我们说话、并凭自己的果实为好橄榄树作证的保禄宗徒，他本人不也是来自那个百姓吗？还有那些主复活后向他显现的五百多位弟兄 \BibleRef{格前 15:6}，他们从哪里来？当宗徒们充满圣神，用各种语言说话时 \BibleRef{宗 2:4}，伯多禄的话使那么多成千的人归化，他们出于对天主荣誉的热忱和自责，以至于那些最初在愤怒中流主血的人，现在因相信而学会了喝主的血。所有这五千人都这样归化了，他们卖掉自己的财产，把价银放在宗徒们脚前。那一个富人没有做的事，当他从主口中听到并忧愁地离开主时 \BibleRef{玛 19:21}，这些亲手把基督钉在十字架上的人中，竟有这么多成千的人突然做到了。他们内心所受的创伤越深，就越急切地寻求医生。既然所有这些人都来自那里，圣咏论到他们说：\Quote{因为主必不遗弃祂的百姓。}……\par

6. 圣咏接着说什么？\Quote{大地四极在祂手中} \BibleRef{咏 94:4}：我们认出那角石：角石就是基督。没有角，除非它本身将两堵墙合而为一：它们从不同侧来到同一个角，但在角中并不彼此对立。割礼从一侧来，未受割礼从另一侧来；在基督里两个民族相遇了，因为祂成了那石头，正如经上所载：\Quote{匠人所弃的石头，已成了屋角的基石。}……\par

7. \Quote{海洋属祂，因为是祂造的} \BibleRef{咏 94:5}。因为海洋就是这个世界，但天主也造了海洋；海浪的咆哮只到岸边，那是祂划定的界限。因此，没有任何诱惑超出它的限度……\Quote{祂的手形成了干地。}你要成为干地：渴慕天主的恩宠，好使祂像甘霖一样降到你身上，在你身上找到果实。祂不让海浪淹没祂所播种的。\Quote{祂的手形成了干地。}因此，我们也应当向主欢呼。\par

8. \Quote{来，让我们俯伏朝拜祂，向创造我们的主哭泣} \BibleRef{咏 94:6}……也许你因过犯的知觉而燃烧；用眼泪抹去你罪的火燎；在主面前哭泣；无畏地在主面前哭泣，因为祂创造了你；祂不会轻视祂在你身上的亲手之作。不要以为你能靠自己修复。你靠自己可能会失落，但你不能修复自己；那创造你的会修复你。\Quote{让我们向创造我们的主哭泣：}在祂面前哭泣，向祂告明，以告明迎见祂的面。因为你在祂面前哭泣并向祂告明，你是谁？不过是祂所造的受造物。受造物对它的创造者怀有极大的信心，而且这创造并非无足轻重，而是按照祂自己的肖像和模样。\par

9. \Quote{因为祂是主，我们的天主}\BibleRef{咏 94:7}。但为使我们能无忧无惧地向祂俯伏下拜，我们是什么呢？\Quote{我们是祂牧场的百姓，是祂手下的羊。}请看，祂如何优雅地调换了词序，似乎没有给每个词配上它应有的属性；为使我们明白，这些羊也就是百姓。祂没有说：\Quote{祂手下的羊，祂牧场的百姓}——这或许更合宜，因为羊属于牧场——祂却说：\Quote{祂牧场的百姓。}所以百姓就是羊，因为祂说：\Quote{祂牧场的百姓；}百姓自己就是羊……祂在《雅歌》中也赞扬这些羊，论及一些成全者，称她们为祂净配——神圣教会——的牙齿：\Quote{你的牙齿如同剪过毛的一群母羊，洗净后上来，个个都有双胎，没有不生育的。}\BibleRef{歌 4:2}这\Quote{你的牙齿}是什么意思？是那些你借以说话的人，因为教会的牙齿就是那些她借以说话的人。你的牙齿是怎样的？\Quote{如同剪过毛的一群母羊。}为什么是\Quote{剪过毛}？因为他们卸下了世界的重担。我刚才说的那些羊，不正是剪过毛的吗？天主的命令剃去了他们的毛，祂说：\Quote{你去，变卖你所有的，分给穷人，你必有宝藏在天上；然后来跟随我。}\BibleRef{玛 19:21}他们遵行了这道命令：剪过毛而来。又因为那些信基督的人受了洗，那里说的什么？\Quote{洗净后上来；}即从洁净中上来。\Quote{个个都有双胎。}哪双胎？那两条诫命，全律法和先知都系于其上。\BibleRef{玛 22:40}\par

10. 因此，\Quote{今天若你们听从祂的声音，就不要硬着心}\BibleRef{咏 94:8}。我的人民，天主的人民！天主对祂的子民说话：不仅是祂不会弃绝的子民，也是祂所有的子民。因为祂在角石\BibleRef{弗 2:20}中对每一面墙说话：就是说，预言在基督内，对犹太人民和外邦人民说话。有一段时间你们藉着梅瑟听见祂的声音，却硬了心。那时，当你们硬着心时，祂藉着一个传报者说话；现在祂亲自说话，愿你们的心软化。祂从前派传报者在祂前面，如今屈尊亲自前来；祂用自己的口说话，祂从前藉着先知的口说话。\par

11. \Quote{如同在默黎巴，在旷野中玛撒的日子，你们的祖先在那里试探了我}\BibleRef{咏 94:9}。愿他们不再是你们的祖先：不要效法他们。他们是你们的祖先，但若你们不效法他们，他们就不是你们的祖先；然而你们由他们而生，他们是你们的祖先。如果那些从地极来的外邦人，如耶肋米亚的话：\Quote{外邦人要从地极来到祢那里，且要说：我们的祖先承受的尽是虚假、虚无，毫无益处；}\BibleRef{耶 16:19}如果外邦人离弃了他们的偶像，来归向以色列的天主；那么以色列人，他们的天主曾引领他们出埃及过红海\BibleRef{出 14:21}，在那里淹没了追赶他们的仇敌；领他们进入旷野，用玛纳养活他们\BibleRef{出 16:13}，从未收回祂管教的杖，从未停止祂仁慈的恩赐；当他们自己的天主这样对待他们时，他们岂能离弃祂，而外邦人却来归向祂？当你们的祖先试探我、考验我，并看见了我的作为……\par

12. \Quote{四十年之久，我几乎厌烦了那一代人，说：这百姓总是心迷途，他们不认识我的道路}\BibleRef{咏 94:10}。四十年与\Quote{总是}有相同的含义。因为那数字四十表示时代的圆满，仿佛时代在这数字中得以完成。因此，我们的主禁食四十天，在旷野中受试探四十天\BibleRef{玛 4:1}，复活后与门徒在一起四十天\BibleRef{宗 1:3}。前四十天祂向我们展示了试探，后四十天则展示了安慰：因为毫无疑问，当我们受试探时，我们得安慰。因为祂的身体，即教会，必在这世界受试探；但那安慰者，祂说：\Quote{看，我同你们天天在一起，直到今世的终结；}\BibleRef{玛 28:20}祂并不缺席。我之所以与他们同在四十年，为要显明这一种族的人，总是惹我发怒，直到世界的终结：因为祂用那四十年来象征这整个世界的时间。\par

13. ……我们以欢欣的喜乐开始：但这圣咏却以极大的恐惧结束：\Quote{因此我在愤怒中起誓说：他们决不得进入我的安息。}\BibleRef{咏 94:11}天主说话已是大事：祂起誓岂不更大？人起誓时你应当害怕，恐怕他因誓言而违背本心行事；何况天主起誓，祂岂能轻率起誓？祂选择起誓作为确证。天主凭谁起誓？凭自己：因为没有更大的可凭者。\BibleRef{希 6:13}祂凭自己确认祂的应许：凭自己确认祂的警告。任何人都不可心里说：祂的应许是真实的，祂的警告是虚假的；祂的应许真实，祂的警告也必定可靠。你若遵行了祂的诫命，就当同样确信安息、幸福、永恒、不朽；你若轻视了祂的诫命，就当同样确信毁灭、永火之烧、与魔鬼一同受罚。……\par

\SourceRecord{Translated by J.E. Tweed.；Nicene and Post-Nicene Fathers, First Series；Vol. 8.；Edited by Philip Schaff.；Buffalo, NY: Christian Literature Publishing Co.,；1888.；<http://www.newadvent.org/fathers/1801095.htm>.}

% structure: p0001 source=h1 target=section reason=page-title

\section{圣咏第九十六篇释义}

1. 我的大人和兄弟\Person{塞维鲁}仍推迟我们对他应许之讲论的喜悦；因为他承认自己是个债务人。因他经过的所有教会，通过他的口，主已使它们喜悦；因此，那从其中主将祂的宣讲传播到其他地方的教会，更应喜悦。但我们还能做什么，唯有服从他的旨意？然而，兄弟们，我说他推迟了，并非他欺骗了我们。所以让我们把他当作被束缚的债务人，直到他还清债务才释放他。因此，亲爱的，请留意：按主所允许的，让我们对这篇圣咏略加解说，你们其实已经知道；因为对真理的重新提及是甘甜的。当它的标题被宣读时，也许有些人惊奇地听着。因为这圣咏的题词是：\Quote{在充军之后建造房屋时}。加上这个标题后，你们或许期待在圣咏的正文中听到，从山上凿出了什么石头，什么石块被拖到工地，打下了什么根基，竖起了什么梁木，立起了什么柱子。但它的歌咏与此毫无关系……并不是这样的房屋在建造；因为请看在何处建造，不是在一处，也不是在任一特定地区。因为他这样开始：——\par

2. \BibleQuote{普世大地，请向上主歌唱新歌；请向上主歌唱，普世大地。} \BibleRef{咏 95:1} 如果普世大地歌唱新歌，它就在歌唱中建造：歌唱本身就是在建造；但条件是它不唱旧歌。肉体的欲情唱旧歌；天主的爱唱新歌……请听为什么这是新歌：主说，\BibleQuote{我给你们一条新命令：你们要彼此相爱。} \BibleRef{若 15:12} 因此普世大地唱新歌：在那里天主的殿被建造。普世大地是天主的殿。如果普世大地是天主的殿，那么不依附普世大地的人，就是废墟，不是殿；那是旧废墟，古代圣殿是其预像。因为在那里旧的被毁灭，为的是建造新的……宗徒将我们联系在一起，成为这同一个结构，并将我们在那合一中捆绑住，说：\BibleQuote{要彼此以爱德担待，竭力保持圣神的合一，藉和平的联系。} \BibleRef{弗 4:2} 哪里有这种圣神的合一，哪里就有一块石头；但那是由许多石头形成的一块石头。如何由许多形成一块？藉彼此以爱德担待。因此主我们天主的殿正在建造；这就是正在进行的工程，为此有这些话语，为此有这些读经，为此有福音在普世的宣讲；它目前仍在建造。这殿已大大增加，充满了许多邦国；然而它尚未通过所有邦国；通过它的增长，它已占据了多数，并将得胜一切：并受到那些自夸属于其家室、并说它已失势之人的反驳。它仍在增长，所有那些尚未相信的邦国注定将要相信；以致无人能说：那舌头会相信吗？蛮族会相信吗？圣神以火舌显现，\BibleRef{宗 2:3}，除了说明没有一种舌头坚硬到不能被那火软化之外，还有什么意义？因为我们知道许多蛮族已相信基督：基督已占据罗马帝国从未到达的地区；对于那些用刀剑争战的人尚且封闭的，对那用木架争战的主并不封闭。因为\Quote{上主从木架作了王。}谁是用木架争战的？基督。祂用十字架战胜了诸王，并将那十字架钉在他们的额上；他们以此夸耀，因为这是他们的救恩。这就是正在进行的工程，殿就这样增长，就这样建造：为使你们知道，请听圣咏接下来的诗句：看他们劳苦，建造这殿。\BibleQuote{普世大地，请向上主歌唱。}\par

3. \Quote{你们要向主歌唱，称颂祂的圣名；天天传扬祂的救恩} \BibleRef{咏 95:2}。建筑如何增长？他说：\Quote{天天传扬祂的救恩。}要天天宣讲；天天，他说，让它被建造；上主说，让我的房屋增长。好像工人们问：祢命令在哪里建造？祢愿意祢的房屋在哪里增长？为我们选择某个平坦、宽敞的地方，如果祢愿意为祢建造一座宽敞的房屋。祢命令我们在哪里天天传扬？祂指示地点：\Quote{在外邦人中宣扬祂的荣耀}：祂的荣耀，不是你们的。哦，你们这些建造者，\Quote{在外邦人中宣扬祂的荣耀。}如果你们选择宣扬自己的荣耀，你们将跌倒；如果宣扬祂的，你们将在建造时被建立起来。因此，那些选择宣扬自己荣耀的人，拒绝居住在房屋中；因此他们不与全地同唱新歌。因为他们不与整个圆球分享，所以他们不是在房屋中建造，而是竖起了一堵粉饰的墙。上主是如何严厉地威胁粉饰的墙？\BibleRef{则 13:10}先知有无数的见证，祂在其中咒骂粉饰的墙。粉饰的墙是什么，就是假善，即虚伪？外面明亮，里面肮脏。……一位谈论这粉饰的墙的人这样说：\Quote{正如，在一堵独立且与其他墙不相连的墙上，你若开一扇门，谁进去，谁就在门外；同样，在那拒绝与房屋同唱新歌、却选择建墙——粉饰的墙且不坚固——的部分，有一扇门又有什么用呢？}如果你进去，你发现自己是在外面。因为他们自己没有从门进去，他们的门也不让他们进入里面。因为主说：\Quote{我是门：借着我进入的。}\BibleRef{若 10:9}……\Quote{在外邦人中宣扬祂的荣耀。}什么是\Quote{在外邦人中}？也许只是少数民族；而那竖起粉饰之墙的部分仍有话可说：为什么格突利亚、努米底亚、毛利塔尼亚、比扎西乌姆不是民族呢？各省都是民族。让天主的话从假善、从粉饰的墙中除掉，在全世界建造房屋。说\Quote{在外邦人中宣扬祂的荣耀}还不够；免得你以为任何民族被排除，他加上\Quote{和祂的奇事在万民中。}\par

4. \Quote{因为上主是伟大的，应受极大的赞美}\BibleRef{咏 95:4}。哪一位主，除非是耶稣基督，\Quote{是伟大的，应受极大的赞美}？你们确实知道祂以人形出现；你们确实知道祂在妇女的子宫中受孕，你们知道祂从子宫中出生，被哺乳，被抱在怀中，受割损，为祂献祭，祂长大；最后，你们知道祂被掌掴，被吐唾沫，被荆棘冠冕，被钉十字架，死去，被长矛刺透；你们知道祂受了这一切：\Quote{祂是伟大的，应受极大的赞美。}不要轻视微小的，要理解伟大的。祂成为微小的，因为你们是这样；愿祂被承认是伟大的，并且在祂内你们将成为伟大的。……微小的舌头能对伟大的赞美说什么呢？说\Quote{无可言喻}，他已说了，并让想象力去构思：好像说，我无法说出的，你思考吧；当你思考后，仍不够。没有人思想能表达的，人的舌头能表达吗？\Quote{上主是伟大的，应受极大的赞美。}愿祂受赞美，被宣讲；祂的荣耀被宣扬，祂的房屋被建造。\par

5. ……因为他想要建造房屋的地方，本身就是多林的，正如昨天所说的：\Quote{我们在林中找到了它。}因为他正在寻找那座房屋，当他说\Quote{在林中}时。为什么那地方是多林的？人们曾崇拜偶像：不足为奇他们喂猪。因为那个离开父亲、与妓女耗尽所有的儿子，过着浪荡的生活，曾喂猪\BibleRef{路 15:12}，即崇拜魔鬼；正是通过外邦人的这种迷信，全地变成了一片树林。但建造房屋的人，要砍伐树林；因此有人说：\Quote{房屋在俘掳之后建造时。}因为人们被魔鬼俘掳，事奉魔鬼；但他们从俘掳中被赎回。他们能出卖自己，但不能赎自己。救赎主来了，付出了代价；祂倾流了祂的血，买得了全世界。你问祂买了什么？你看到祂给了什么；那么就找出祂买了什么。基督的血是代价。什么与之相等？除了全世界，还能有什么？除了所有的民族，还能有什么？他们对自己的代价非常忘恩负义，或者非常骄傲，因为他们说代价是如此之小，只能买非洲人；或者他们如此伟大，以至于只为他们的缘故而付出代价。让他们不要得意，不要骄傲：祂为全世界付出了祂所付出的。祂知道祂买了什么，因为祂知道以什么代价买的。因此因为我们被赎回，房屋在俘掳之后被建造。谁俘掳了我们？因为对他们说\Quote{宣扬祂的荣耀}的人，是树林的清除者：他们要砍伐树林，使大地从俘掳中获得自由，并建造、高举，通过宣扬上主房屋的伟大。魔鬼的树林如何被清除？除非那位超越万有的被宣讲。所有的民族当时都以魔鬼为神：他们所称为神的，就是魔鬼，正如宗徒更明确地说：\Quote{外邦人祭祀的，是祭祀魔鬼，而不是天主。}\BibleRef{格前 10:20}既然他们因向魔鬼献祭而身处俘掳，因此全地曾是树林；祂被宣扬为伟大的，超越世上的一切赞美。\par

6.……当他说：\BibleQuote{祂比众神更可敬畏}时，他又加上：\BibleQuote{外邦人的众神都是魔鬼}。……因为\BibleQuote{外邦人的众神都是魔鬼}。难道这就是对那位不能受恰当赞美的祂的赞美——祂超越外邦人的众神，即魔鬼吗？且慢，听下文：\BibleQuote{是主创造了诸天}。因此，主不仅超越众神，更超越祂所造的诸天。若他只说\BibleQuote{超越众神，因为外邦人的神是魔鬼}，而我们对主的赞美止于此，那我们对基督的思虑就比惯常的要少；但当他说\BibleQuote{但主创造了诸天}时，请看诸天与魔鬼之间的差别，以及诸天与创造诸天者之间的差别；看，主是何等崇高。他没有说：主坐在诸天之上；因为或许有人会想象是另一个人造了诸天，而主坐在其上；但他说：\BibleQuote{是主创造了诸天}。若祂造了诸天，祂也造了天使：祂亲自造了天使，亲自造了宗徒。魔鬼屈服于宗徒；但宗徒本身是诸天，他们承载着主。……哦，诸天，是祂所造的，向外邦人宣告祂的荣耀！愿祂的全地建立祂的殿，愿全地唱新歌。\par

7.\BibleQuote{颂赞与华美在祂面前}\BibleRef{咏 95:6}。你喜爱华美吗？你愿成为美丽的吗？那就颂赞！他没有说：华美与颂赞，而是颂赞与华美。你曾是污秽的；颂赞吧，使你成为洁净；你曾是罪人；颂赞吧，使你成为义人。你能使自己变形，却不能使自己变美。但我们的新郎是怎样的，祂竟爱了一个变形的人，为要使她美丽？有人说：祂怎能爱一个变形的人？祂说：\BibleQuote{我来不是召义人，而是召罪人}\BibleRef{玛 9:13}。你召谁？罪人，使他们继续作罪人吗？不，祂说。那么他们如何不再作罪人呢？\BibleQuote{颂赞与华美在祂面前}。他们藉着承认自己的罪来荣耀祂，他们吐出贪食吞下的恶物，\BibleQuote{不再像那不洁的狗回到自己的呕吐物中}\BibleRef{伯后 2:22}；于是就有了颂赞与华美：我们爱华美，让我们先选择颂赞，华美就会随之而来。又有一个人爱权力和伟大：他愿意像天使那样伟大。天使有某种伟大，有那样的权能，若天使全力施展，无人能抵挡。人人都渴望天使的权能，但并非人人都爱天使的正义。先爱正义，权能将跟随你。因为下文是什么？\BibleQuote{圣洁与威严在祂的圣所中}\BibleRef{咏 95:7}。你先前寻求伟大：先爱正义吧；当你成为正义的，你也将成为伟大的。因为若你颠倒次序，先愿成为伟大的，你将在能升起之前跌倒：因为你不能升起，你是被举起的。你更好的升起，如果那从不跌倒的祂来举起你。因为那从不跌倒的祂降到你的层面：你曾跌倒；祂降下，祂向你伸出祂的手；你无法凭自己的力量升起，握住那降下者的手，使你能被那强有力者举起。\par

8.那么怎样？若\BibleQuote{颂赞与华美在祂面前；圣洁与威严在祂的圣所中}\BibleRef{咏 95:7}。这就是我们所宣告的，当我们建造这殿时；看，已经向外邦人宣告了；外邦人该做什么？那些清除了树林的人已向外邦人宣告了主的荣耀。现在祂亲自对外邦人说：\BibleQuote{你们众民的家族，要将荣耀归于主，要将荣耀和权能归于主}。不要归于你们自己：因为那些向你们宣告的人，宣告的也不是自己的，而是祂的荣耀。那么你们\BibleQuote{要将荣耀和权能归于主}，并说：\BibleQuote{主啊，荣耀不要归于我们，不要归于我们，只归祢的名下}。不要信赖人。若你们中有人受洗，让他说：那施洗者，新郎的朋友说过：\BibleQuote{祂要以圣神施洗}。因为当你说这话时，你是将荣耀和权能归于主：\BibleQuote{要将荣耀和权能归于主}。\par

9.\BibleQuote{要将主名的荣耀归于主}\BibleRef{咏 95:8}。不要归于人的名，不要归于你们自己的名，而要归于祂的荣耀。……颂赞是献给天主的礼物。哦外邦人，若你们要进入祂的庭院，不要空手进入。\BibleQuote{带着礼物}。我们要带什么礼物呢？天主所喜悦的祭献是痛悔的精神：忧伤痛悔的心，\BibleQuote{天主啊，祢必不轻看}。怀着谦卑的心进入天主的殿，你便带着礼物进入了。但若你骄傲，你便是空手进入。因为你若不自满，何以骄傲？因为若你满了，你便不会自大。你怎能满呢？若是你带着礼物，就是你要带进主庭院的。我们不使你们久留：让我们快速浏览剩余部分。看，这殿在增长；看，这建筑遍布全地。欢喜吧，因为你们已进入庭院；欢喜吧，因为你们正被建成天主的殿。因为那些进入的人本身正被建造，他们本身就是天主的殿；祂是居住者，殿为祂建造在全地之上，而这\BibleQuote{在掳掠之后}。\BibleQuote{带着礼物，进入祂的庭院}。\par

10.\BibleQuote{要在主的圣所中朝拜主}\BibleRef{咏 95:9}：在天主教会中，这是祂的圣所。不要有人说：\BibleQuote{看，基督在这里，或在那里。因为将有假先知兴起}\BibleRef{玛 24:23}。你们要这样对他们说：\BibleQuote{这里将没有一块石头留在另一块上，而不被拆毁}。你们召我到粉饰的墙那里；我在祂的圣所中朝拜我的天主。\BibleQuote{愿全地都在祂面前战栗}。\par

11. \BibleQuote{你们要在万民中宣告，主由木架上为王；是祂使大地稳固，永不摇动} \BibleRef{咏 95:10}。这些建筑天主圣殿的见证何等明显！天上的云彩在全世界雷鸣，说天主的圣殿正在建造；而沼泽中的青蛙却叫嚷：只有我们是基督徒。我提出什么见证呢？\WorkTitle{圣咏集}的见证。我提出你们像聋子一样唱的这首圣咏：打开你们的耳朵；你们唱这首圣咏；你们和我一起唱，却不同意我；你们的舌头发出和我相同的声音，但你们的心却与我的不同。你们不唱这首圣咏吗？看，全世界的见证：\BibleQuote{愿全地在祂面前震动}；你们却说，你们没有被震动？\BibleQuote{你们要在异民中宣告，主由木架上为王}。难道人会在这里得势，说他们凭木器为王，因为他们凭自己匪帮的棍棒为王？如果你们要凭木器为王，就凭基督的十字架为王。因为你们的木器使你们成为木人；基督的木器载你们渡海。你们听见圣咏说：\BibleQuote{祂使大地稳固，永不摇动}；你们却不仅说它自从被稳固后发生了动摇，而且还说它缩减了。是你们说实话，圣咏作者说谎？是那些大喊\BibleQuote{看，基督在这里，在那里} \BibleRef{玛 24:23}的假先知说实话，而这位圣咏作者撒谎？弟兄们，面对这些最明确的言语，你们在角落里听见这样的谣言：\Quote{某某曾是背信者}，\Quote{某某曾是背信者}。你们说什么？应该听你们的话，还是天主的话？因为\Quote{是祂使大地稳固，永不摇动}。我向你们展示已建成的世界：带上你们的礼品，进入上主的庭院。你们没有礼品，因此不愿进入。这是什么意思？如果天主指定一头牛、山羊或公绵羊作为礼品，你们会找一头带来；祂指定了一颗谦卑的心，你们却不愿进入；因为你们在自己内找不到这颗心，你们因骄傲而膨胀。\Quote{祂使大地稳固，永不摇动：祂要以正义审判万民。}那时，如今拒绝喜爱正义的人将要哀号。\par

12. \BibleQuote{愿诸天欢乐，愿大地踊跃} \BibleRef{咏 95:11}。愿那传扬天主光荣的诸天欢乐；愿主所造的诸天欢乐；愿那接受诸天降雨的大地踊跃。因为诸天是宣讲者，大地是聆听者。\BibleQuote{愿海翻腾，及其中的万物。}什么海？世界。海翻腾了，及其中的万物；全世界起来反对教会，而教会正扩建并建造在全地上。关于这翻腾，你们在福音中听见：\BibleQuote{你们将被交付给议会} \BibleRef{谷 13:9}。海翻腾了；但海怎能胜过创造它的主呢？\par

13. \BibleQuote{愿平原欢乐，及其中所有的一切} \BibleRef{咏 95:12}。所有谦卑者、所有温良者、所有正义者，都是天主的\OriginalTerm{平原}{plains}。\BibleQuote{那时，林中的树木都要欢跃。}林中的树木就是外邦人。为什么他们欢跃？因为他们从野橄榄树上被砍下来，被接在好橄榄树上。\BibleRef{罗 11:17} \BibleQuote{那时，林中的树木都要欢跃}；因为巨大的香柏和柏树被砍下来，不朽的木材被买来建造圣殿。它们曾是林中的树木；但在被送到建筑之前，它们是林中的树木；在结出橄榄之前，它们是林中的树木。\par

14. \BibleQuote{在主面前，因为祂来了，祂来了要审判世界。} \BibleRef{咏 95:13} 祂起初来过，还将再来。祂先在云中来到祂的教会。承载祂的云是什么呢？就是宣讲的宗徒们，关于他们你们已听过，当书信被宣读时：\Quote{我们是基督的使者}，他说：\Quote{我们替基督请求你们：与天主和好吧。} \BibleRef{格后 5:20} 这就是祂在其中来临的云，除了祂最后的来临，那时祂将审判生者死者。祂先是在云中来临。这是祂在福音中初次发出的声音：\BibleQuote{从那时起，他们要看见人子乘云而来。} \BibleRef{谷 13:26} \Quote{从那时起}是什么意思？难道主不将在更晚的时候，当地上万族都要哀哭时来临吗？祂先在祂自己的宣讲者中来临，充满整个大地。让我们不要抗拒祂的初次来临，好使我们不至于在祂第二次来临时战栗。\BibleQuote{但在那些日子，怀孕的和哺乳的有祸了！} \BibleRef{谷 13:17} 你们刚才在福音中听到：\BibleQuote{你们要当心，因为你们不知道祂几时来到。} \BibleRef{谷 13:33} 这是比喻的说法。谁是有孕的，谁是哺乳的呢？有孕的，是那些希望寄托于世界的灵魂；但那些已得所愿的，就是所谓的\InnerQuote{哺乳者}。例如：一个人想买一处乡间别墅；他是有孕的，因为他的目的尚未达成，子宫因希望而膨胀；他买了它；他生产了，他现在就哺乳他所买的。\Quote{怀孕的和哺乳的有祸了！} 那些寄望于世的人有祸了！那些紧抓着他们因属世希望而生出之物的人有祸了！那么基督徒该做什么？他应该使用世界，而非服事世界。\BibleRef{格前 7:29} 这是什么意思？有妻子的，要像没有一样……毫无挂虑的人，毫无恐惧地等待他主的来临。因为害怕祂来临，这是对基督的哪门子爱？弟兄们，我们难道不羞愧吗？我们爱祂，却又害怕祂来临。我们确定我们爱祂吗？还是我们更爱自己的罪？因此，让我们因罪本身而恨罪，并爱那来惩罚我们罪的那一位。祂必会来临，无论我们喜欢与否：因为祂不是现在来，也不是祂根本不会来。祂必会来，且在你不知道的时候；如果祂发现你已预备好，你的无知不会伤害你。\Quote{那时林中的树木都要在主面前欢欣，因为主来了：} 这是指祂的初次来临。之后呢？\Quote{因为主来了，为审判大地。林中的树木都要欢欣。} 祂先来了，后来审判大地：祂要找到那些因信祂初次来临而欢欣的人，\Quote{因为主来了。}\par

15. \Quote{祂要以正义审判世界：} 不是一部分，因为祂赎回的不是一部分：祂要审判全部，因为祂为全部付出了代价。你们已在福音中听到，当祂来时，\BibleQuote{祂要由四方聚集祂所拣选的人。} \BibleRef{谷 13:27} 祂从四方聚集祂所有的选民：因此是从全世界。因为\Person{亚当}（我之前说过）在希腊语中意指全体世人；因为有四个字母：A、D、A、M。但按照希腊人的说法，世界的四个方位有这些首字母：他们称之为\OriginalTerm{东方}{\GreekText{Ανατολὴ}}；他们称之为\OriginalTerm{西方}{\GreekText{Δύσις}}。\par

\SourceRecord{Translated by J.E. Tweed.；Nicene and Post-Nicene Fathers, First Series；Vol. 8.；Edited by Philip Schaff.；Buffalo, NY: Christian Literature Publishing Co.,；1888.；<http://www.newadvent.org/fathers/1801096.htm>.}

% structure: p0001 source=h1 target=section reason=page-title

\section{诠释圣咏第九十七篇}

1. ……这篇圣咏的标题是：\BibleQuote{达味的一篇圣咏，当他的土地复原时。}让我们将全部内容归于基督，若我们愿持守正确领悟的道路；让我们不要偏离角石（\BibleRef{弗 2:20}），免得我们的领悟跌倒；愿在祂内，那因不稳定而摇摆的得以稳固；愿那在不确定中前后摇晃的，依靠于祂。当人聆听天主的圣经时，心中若存有疑虑，不要离开基督；当基督在话语中向他启示时，他才可确信自己已经领悟；但在达到对基督的认识之前，不可妄称已经领悟。\BibleQuote{因为基督是法律的终向，使所有信者获得正义。}（\BibleRef{罗 10:4}）这是什么意思？这些话如何在基督内理解，\BibleQuote{当他的土地复原时}？……\par

2. 土地的复原，就是肉身的复活；因为在他复活后，圣咏所歌唱的一切都实现了。让我们来听一首充满喜乐的圣咏，因土地复原的喜乐。愿主我们的天主在我们心中激发与如此大事相称的盼望与喜乐；愿祂引导我们的言谈，使之合于你们的心怀；无论我们的心在这样的景象中感受到何种喜乐，愿祂将它引到我们的舌上，再由此送入你们的耳朵，然后进入你们的心，最后进入你们的行动。\par

3. ……\BibleQuote{上主为王，愿大地踊跃；愿众多岛屿欢乐！}（\BibleRef{咏 96:1}）的确如此，因为天主的言语不仅在大陆上传扬，也在那些位于海中的岛屿上传扬；这些岛屿也充满了基督徒，充满了天主的仆人。因为大海不能阻挡那创造大海者。凡船只可到之处，天主的言语不能到达吗？岛屿已被充满。但比喻地说，岛屿可代表所有教会。为何是岛屿？因为各样诱惑的波浪在它们周围咆哮。但正如一个岛屿可能被四周冲击的波浪击打，却不会被击碎，反而它自己击碎来袭的波浪，而不是被它们击碎：同样，天主的教会，在全世界涌现，曾遭受不敬虔者的迫害，他们从四面咆哮；看，这些岛屿屹立不动，最终大海平静了。\par

4. \BibleQuote{密云和黑暗环绕祂公义和审判是祂宝座的基础。}（\BibleRef{咏 96:2}）……主亲自说：\BibleQuote{我为审判来到这世界，叫那些看不见的，可以看见；叫那些看见的，反成了瞎子。}（\BibleRef{若 9:39}）那些自以为看得见的人，那些自以为聪明、以为不需要医治的人，为了使他们成为瞎子，好叫他们无法理解。而\BibleQuote{叫那些看不见的可以看见}，叫那些承认自己盲目的人，得蒙光照。因此，\BibleQuote{密云和黑暗环绕祂}，是为那些没有理解祂的人；对于那些承认并谦卑的人，\BibleQuote{公义和审判是祂宝座的基础。}祂称那些相信祂的人为祂的宝座：因为祂从他们中为自己造了一个宝座，因为智慧坐在他们内；因为天主子是天主的智慧。但我们从圣经的另一处听到了对这一解释的有力证实。\BibleQuote{义人的灵魂是智慧的宝座。}因为那些相信祂的人，已成了义人：因信成义，他们成了祂自己的宝座；祂住在他们内，从他们中施行审判，引导他们。……\par

5. \BibleQuote{有火在祂面前先行，烧灭祂四围的敌人。}（\BibleRef{咏 96:3}）我们记得在福音中读到，祂要说：\BibleQuote{你们这些被咒骂的，离开我，到那为魔鬼和他的使者预备的永火里去吧。}（\BibleRef{玛 25:41}）我不认为这是指那火。为何不？因为祂说有一种火，在祂面前先行，在祂来审判之前。因为经上说，火在祂面前先行，烧灭祂四围的敌人，即遍及全世界。那火是在祂来临之后烧灭；而这火则是在祂之前先行。那么这火是什么呢？……看，我们已经理解那在祂面前先行的火，应理解为对不信和不敬虔之人的一种暂时的惩罚；让我们也理解这火，如果可能，也关乎被救赎者的救恩；因为我们曾这样主张。主亲自说：\BibleQuote{我来是把火投在地上，}（\BibleRef{路 12:49}）这\Quote{火}如同\Quote{刀剑}；正如在另一处祂说，祂来不是把和平，而是把刀剑带到地上。（\BibleRef{玛 10:34}）刀剑是为分开，火是为焚烧：但二者都是有益的：因为祂话语的刀剑，有益地将我们从恶习中分别出来。因为祂带来刀剑，使每一个信徒，或从他的不信基督的父亲，或从同样不信的母亲，分离出来；至少，如果我们生于基督徒父母，则从祖先分离。因为我们中间没有人没有祖辈、曾祖辈，或某些外邦祖先，和在那被天主诅咒的不信中。我们从我们从前所是的分离出来；但那分离我们的刀剑，并不杀害，而是在我们之间割开。同样，火也是：\BibleQuote{我来是把火投在地上。}信祂的人被点燃，他们领受了爱火；为此，当圣神自己被派遣到宗徒们身上时，祂这样显现：\BibleQuote{有分开的舌头，好像火焰。}（\BibleRef{宗 2:3}）被这火焚烧，他们踏上行程，走遍世界，要焚烧并点燃祂四围的敌人。谁是祂的敌人？那些离弃创造他们的天主，敬拜他们自己所造的偶像的人。……\par

6. \BibleQuote{祂的闪电照亮了世界} \BibleRef{咏 96:4}。这是极大的喜乐。我们岂没有看见？岂不清楚吗？祂的闪电照亮了全世界：祂的敌人被点燃，烧毁了。所有反对的都被烧毁，而\BibleQuote{祂的闪电照亮了世界}。它们如何照亮？为使世界最终相信。闪电从何而来？从云彩而来。天主的云彩是什么？是真理的宣讲者。但你看见一朵云，在空中朦胧而黑暗，它里面隐藏着我不知道的东西。如果云中发出闪电，就有光芒闪耀：从你所轻视的事物中，迸发出你该畏惧的事物。因此，我们的主耶稣基督派遣了祂的宗徒，作为祂的宣讲者，如同云彩：他们被视为人，受到轻视；如同云彩显现，被轻视，直到你惊异的事情从他们中闪现。因为他们首先是肉体缠累、软弱的人；然后，是地位低下、无学问、卑微的人；但内有能发光的、能闪射的。伯多禄渔夫走近，祈祷，死者就起来了。\BibleRef{宗 9:40} 他的人形是云，奇迹的光辉是闪电。同样，在他们的言语中，在他们的行为中，当他们做出令人惊奇的事，说出令人惊奇的话时，\BibleQuote{祂的闪电照亮了世界；大地看见，就战栗。}难道不是吗？整个基督教世界最终不都喊着阿们吗？因从那些云中迸发出的闪电而战栗。\par

7. \BibleQuote{山岭如蜡在主面前熔化} \BibleRef{咏 96:5}。山岭是谁？骄傲的人。所有自高反抗天主的事物，在基督和基督徒的作为面前战栗、屈服，当我已说过\BibleQuote{熔化}时，找不到更好的词了。\BibleQuote{山岭如蜡在主面前熔化。}权势的高傲在哪里？不信者的顽固在哪里？主是他们面前的火，他们如蜡在主面前熔化，一直坚硬，直到那火靠近。一切高处都被削平了；如今不敢亵渎基督：虽然外教人不信祂，却不再亵渎祂；虽然尚未成为活石，但坚硬的山已被征服。\BibleQuote{在普世的主面前：}不只是犹太人的主，也是外邦人的主，正如宗徒所说，因为祂不只是犹太人的天主，也是外邦人的天主。\BibleRef{罗 3:29} 所以祂是普世的主，主耶稣基督生于犹太，但不仅为犹太而生，因为祂在出生前就创造了众人；祂既创造了众人，也重造了众人。\par

8. \BibleQuote{诸天传扬了祂的正义，万民都看见了祂的荣耀} \BibleRef{咏 96:6}。什么诸天传扬了？\BibleQuote{诸天宣示天主的荣耀。}谁是诸天？那些成为祂宝座的人；因为天主居住在诸天中，同样祂居住在宗徒们中，同样居住在福音的宣讲者中。你，如果你愿意，也将成为天。你愿意这样吗？从你心中清除大地。如果你没有地上的贪欲，并且没有徒然地诵念回应，即你已经\Quote{举心向上}，你就将成为天。\BibleQuote{你们既然与基督一同复活了，}宗徒对信徒说，\BibleQuote{就该思念天上的事，不要思念地上的事。}\BibleRef{哥 3:1} 你已开始思念天上的事，而非地上的事；你岂不成了天？你带着肉体，但心中已是天，因为你的家乡是在天上。\BibleRef{斐 3:20} 你既是这样，你也传扬基督；因为哪位信徒不传扬基督呢？……因此整个教会传扬基督，诸天传扬祂的正义；因为所有信徒，他们关心的是为使那些尚未相信的人归向天主，并出于爱这样做，他们就是诸天。天主从他们中发出审判的恐怖雷声；不信者战栗、惊醒并相信。祂向人显示基督在全世的能力，通过劝诱他们，引导他们爱基督。因为今天有多少人把他们的朋友引向某个哑剧演员或笛手？为什么？岂不是因为喜欢他？你们也要爱基督。因为那征服了世界的，展示了这样的奇观，以致无人能说在其中找到了指责的理由。因为在剧场中每个人喜爱的角色常常在那里被击败。但在基督里没有人被击败：没有可羞耻的理由。抓住、引导、吸引你所能吸引的人；不要害怕，你正在引他们到祂那里，祂不使看见祂的人不悦；并求祂光照他们，使他们能看得好。\par

9. \Quote{愿所有崇拜雕刻偶像的人都蒙羞}\BibleRef{咏 96:7}。这不是已经应验了吗？他们不是已经蒙羞了吗？他们不是天天蒙羞吗？因为雕刻的偶像乃是人手所造的。为何所有崇拜雕刻偶像的人都蒙羞？因为万民都已看见祂的荣耀。如今所有民族都承认基督的荣耀：让那些崇拜石头的人羞愧吧。因为那些石头是死的，我们却找到了一块活石；的确，那些石头从未活过，连死都称不上；但我们的石头是活的，且永远与圣父同活，虽然祂为我们死过，却复活了，如今活着，死亡再不能统治祂。\BibleRef{罗 6:9} 祂的荣耀万民已经承认；他们离开庙宇，奔向教会。他们还在寻求崇拜雕刻的偶像吗？他们不是已经选择离弃他们的偶像了吗？他们已被他们的偶像所离弃。\Quote{凡以他们的偶像夸耀的。} 但有一位自以为有学问的辩论者说：\Quote{我不崇拜那石头，也不崇拜那没有知觉的像；……我不崇拜这像，但我敬拜我所见的，并服事我所不见的。} 那是谁？他回答说：是某个不可见的神明，主宰着那像。他们用这种解释他们的像，自以为是能干的辩士，因为他们不崇拜偶像，却仍崇拜魔鬼。\Quote{弟兄们，}宗徒说，\Quote{外邦人所献的祭，是祭给魔鬼，不是给天主；我们知道偶像算不得什么；并且外邦人所献的祭，是祭给魔鬼，不是给天主；我不愿意你们与魔鬼有份。} 因此，他们不可以此为借口，说他们不是献给无灵的偶像；他们其实是献给魔鬼，这更为危险。因为如果他们只是崇拜偶像，偶像既不能帮助他们，也不能伤害他们；但如果你崇拜和服事魔鬼，他们自己将成为你的主人……\par

10. 但请看那些如同天使一般的圣人们。当你找到某位服事天主的圣人，你若想把他当作天主来崇拜，他会禁止你：他不会将属于天主的尊荣归于自己，他不会在你面前充当天主，而是与你一起在天主之下。圣宗徒保禄和巴尔纳伯就是这样。他们在吕考尼雅宣讲天主的圣言。当他们在吕考尼雅行了奇事后，那地的人民带来祭品，想要向他们献祭，称巴尔纳伯为则乌斯，称保禄为赫尔默斯；他们并不喜悦。他们拒绝被献祭，是因为他们厌恶自己被比作魔鬼吗？不，而是因为他们战栗于将天主的尊荣归于人。他们自己的话表明了这一点：这不是我们的猜测；因为书上继续叙述他们如何被激动。\BibleRef{宗 14:14} ……正如善人们禁止那些想要把他们当作神崇拜的人，反而愿意唯独天主被崇拜，唯独天主被朝拜，唯独向天主献祭，而非向他们自己；同样，所有圣天使也寻求他们所爱之天主的荣耀；他们努力推动并激励他们所有所爱的人去默观天主：他们向他们宣告天主，而非他们自己，因为他们是天使；又因为他们是战士，他们只寻求他们统帅的荣耀；但如果他们寻求自己的荣耀，他们就被斥为篡权者。魔鬼和他的使者就是这样；他为自己要求天主的尊荣，并为所有他的魔鬼要求；他充满了异教庙宇，并说服他们向他自己献上像和祭品。崇拜圣天使不是比崇拜魔鬼更好吗？他们回答说：我们不崇拜魔鬼；我们崇拜天使，如你们所称的，伟大天主的掌权者和仆役。我倒希望你们崇拜他们：你们会轻易从他们本身学会不崇拜他们。听一位天使的教导。他在教导基督的一位门徒，并在\WorkTitle{若望默示录}中向他显示许多奇事：当某个奇异的景象向他显示后，他战栗，俯伏在那位天使脚前；但那位天使，只寻求天主的荣耀，说：\Quote{你不可这样做！因为我是你和你弟兄先知们的同仆。}\BibleRef{默 19:10} 那么，我弟兄们，该怎么办？愿无人说：\Quote{我怕天使会向我发怒，如果我不把他当作我的天主来崇拜。} 你若选择崇拜他，他倒是会向你发怒：因为他是正义的，爱天主。正如魔鬼若不受到崇拜就发怒，同样，天使若被当作天主崇拜也发怒。但恐怕那软弱战栗的心自言自语说：\Quote{如果魔鬼因不受到崇拜而被激怒，我怕得罪他们；他们的首领魔鬼能对你做什么呢？如果他对我们有能力，我们中间就没有一个人能存留。难道不是每天有许多反对他的话从基督徒口中说出，而基督徒的收成仍在增长吗？} 当你对你最邪恶的奴仆发怒时，你给他\Quote{撒殚}\Quote{魔鬼}的名字。或许你在这件事上犯了错误，因为你向一个人说这话，你的过度愤怒使你辱骂天主的肖像：然而你选择一个你深恨的词来用在他身上。如果他能，难道他不报复吗？但不被允许；他只被允许做那么多。因为当他想要试探约伯时，他必须先求得权柄才能这样做：\BibleRef{约 1:11} 若没有获得权柄，他什么也不能做。那么，你为何不无所畏惧地崇拜天主呢？没有祂的旨意，无人能伤害你；得到祂的许可，你被管教，却不被征服。因为如果主你的天主乐意允许某个人或某个邪灵伤害你：祂将管教你，使你向祂呼求：\Quote{因此，愿所有以虚无之神为乐的人都蒙羞；你们众天使，都要崇拜祂。} 让外邦人学习崇拜天主：他们想要崇拜天使：让他们效法天使，崇拜那位被天使崇拜的天主。\Quote{你们众天使，都要崇拜祂。} 让那位被派往科尔乃略的天使也崇拜祂（因为他崇拜祂，所以把科尔乃略送到伯多禄那里），他自己是伯多禄的同仆；让他崇拜基督，伯多禄的主。\Quote{众神哪，你们都要崇拜祂。}\par

11. \Quote{熙雍听见了，就欢喜} \BibleRef{咏 96:8}。熙雍听见了什么？听见祂的众天使都朝拜祂……因为那时教会尚未在外邦人中；在犹太，有些犹太人信了，而那些信了的犹太人以为只有他们属于基督：宗徒们被派往外邦人那里，科尔乃略领受了宣讲；科尔乃略信了，受了洗，与他一起的人也受了洗。\BibleRef{宗 10:47} 但你们知道所发生的事，好使他们受洗：读者确实还没有读到这一点，但有些人仍然记得；那些不记得的人，请从我这简略听一下。天使被派往科尔乃略那里；天使把科尔乃略送到伯多禄那里；伯多禄来到科尔乃略那里。因为科尔乃略和他的家人是外邦人，未受割礼：恐怕他们犹豫是否将福音传给未受割礼的人：在科尔乃略和他的家人受洗之前，圣神降临，充满了他们，他们开始说起舌音。那时，圣神还没有降在任何一个未受洗的人身上；但在这群人身上，祂在领受圣洗前就降下了。因为伯多禄可能犹豫是否给未受割礼的人施洗：圣神降临了，他们开始说起舌音；无形的恩赐被赐下，消除了对有形圣事的一切怀疑；他们全都受了洗。……熙雍听见了什么而欢喜？就是外邦人也接受了天主的话。一堵墙已经来了，但隅石尚未存在。这里的熙雍特别指在犹太的教会。\Quote{熙雍听见了，就欢喜：犹大的女子都快乐。} 经上这样记载：\Quote{在犹太的宗徒和弟兄们听见了。} 请看看犹大的女儿们是否不快乐。她们听见了什么？\Quote{外邦人也接受了天主的话。}……因此，\Quote{犹大的女子因祢的审判，上主，而快乐。} 什么是因祢的审判？因为在任何民族、任何百姓中，事奉祂的人都为祂所悦纳：因为祂不但是犹太人的天主，也是外邦人的天主。\BibleRef{罗 3:29}\par

12. 请看，这不正是犹大女子快乐的理由吗？\Quote{因为祢，上主，是至高的，超越全地} \BibleRef{咏 96:9}。不但在犹太，而且超越耶路撒冷；不但在熙雍，而且超越全地。天主的审判在这全地上得胜，以致它从各方聚集了万民：那些自绝于外的审判，他们既不听预言，也不见其应验；\Quote{因为祢，上主，是至高的，超越全地：祢被尊崇，远超众神之上。} 什么是\Quote{远超}？因为这是论及基督说的。那么\Quote{远超}是什么意思，无非是祢被承认为与圣父同等？什么是\Quote{超乎众神之上}？他们是谁？偶像没有生命，没有知觉：魔鬼有生命和知觉，但它们是邪恶的。基督被尊崇超越魔鬼，这有什么了不起？祂被尊崇超越魔鬼：但这也算不上什么大不了的；外邦人的神确实是魔鬼，但\Quote{祂远超众神之上。} 连人也被称为神：\Quote{我曾说：你们是神，都是至高者的子女：} 又经上记着：\Quote{天主站在诸神的会中，在众神中行审判。} 我们的主耶稣基督被尊崇超越一切：不仅超越偶像，不仅超越魔鬼，而且超越所有义人。这还不够；也超越所有天使：不然，\Quote{众神啊，你们都要朝拜祂！} 这句话从何而来？\Quote{祢被尊崇，远超众神之上。}\par

13. 那么我们这些在祂面前——在远超众神的那位面前聚集的人，该做什么呢？祂给了我们一条简短的诫命：\Quote{你们爱慕上主的，都当恨恶邪恶！} \BibleRef{咏 96:10}。基督不配得着你们爱祂的同时也爱贪婪。你们爱祂，就应当恨祂所恨的。有一个人是你的仇敌，他和你是同类；你们是同一造物主的作品，有相同的本性：然而，如果你的儿子与你的仇敌说话，去他家，常与他交往，你恐怕要剥夺他的继承权；因为他与你的仇敌说话。这怎样呢？因为你会似乎公正地说：你是我仇敌的朋友，还想要我的产业吗？请留心。你爱基督：贪婪是基督的仇敌；为什么你与她交谈？我不说交谈，你为什么服事她？因为基督命令你做许多事，你却不做；她命令你，你就去做。基督命令你给穷人衣服穿，你却不做；贪婪命令你欺诈，你却宁愿做这个。如果情形是这样，如果你是这样的人，就不要太自信地对自己应许基督的产业。但你说：我爱基督。由此可见，如果你被发现恨恶邪恶，你就爱慕善事。……\par

14. 因为在此之前祂曾说：\Quote{你们要憎恨邪恶的事}，免得你们因害怕邪恶，怕它会杀害你们，祂随即补充道：\Quote{主保护祂仆人的灵魂。}请听祂保护祂仆人的灵魂，并且说：\BibleQuote{你们不要害怕那些杀害肉身，却不能杀害灵魂的。}\BibleRef{玛 10:28} 那最有能力反对你们的，只杀害肉身。他向你们做了什么？他也曾对你们的主天主这样做过。你们既然害怕遭受基督所遭受的，又为何喜爱拥有基督所拥有的呢？祂来是为承受你们的生命，暂时的、软弱的、必死的。如果你们能避免死亡，你们当然害怕死亡。既然你们因本性无法避免，为何不藉信德承受呢？让那威胁的仇敌从你们那里夺去那生命吧，天主会赐给你们另一个生命：因为祂也赐给了你们这个生命，若没有祂的旨意，连这个也不会从你们那里取去；但若祂的旨意是要从你们那里取去，祂会赐给你们一个交换的生命；不要因祂的缘故害怕被剥夺。你们不愿意脱掉一件破旧的衣服吗？祂会赐给你们一件光荣的袍子。你们对我提到什么袍子？\BibleQuote{这可朽坏的必须穿上不可朽坏的，这可死的必须穿上不死。}\BibleRef{格前 15:53} 你们的这肉身本身不会朽坏。你们的仇敌可以发怒直至你们死亡：他不能超越，既不能胜过你们的灵魂，也不能胜过你们的肉身；因为即使他分散你们的肉身，也不能阻止复活。人们为自己的生命害怕；主对他们说什么？\BibleQuote{就是你们的头发，也都一一数过了。}\BibleRef{玛 10:30} 你们连一根头发也不失落，还害怕失去生命吗？一切在天主那里都是数过的。祂创造了一切，也将恢复一切。它们本来没有，而被创造；它们存在了，难道不能恢复吗？……\Quote{祂要从恶人的手中拯救他们。}\par

15. 但或许你会说：我失去这光。\BibleQuote{为义人升起了光。}\BibleRef{咏 96:11} 你害怕失去什么光？你害怕自己处于黑暗中吗？不要害怕失去光；不，要害怕的是，当你谨防失去这光时，你反而失去了那真光。因为我们看到，你所害怕失去的光给了谁，又与谁共享。当祂使太阳升起，照耀义人和不义的人，降雨给义人和不义的人时，难道只有义人看见这太阳吗？\BibleRef{玛 5:45} 恶人、强盗、不洁的人、野兽、苍蝇、虫子，都和你一起看见那光。祂将这光甚至赐给这样的人，那么祂为义人保存的是怎样的光呢？殉道者堪当以信德注视这光；因为他们藐视这太阳的光，他们的眼睛里有某种光，他们渴望那光，他们拒绝了这光。你以为他们被锁链捆绑时真的痛苦吗？对忠信者而言，监狱是宽敞的；对宣信者而言，锁链是轻的。他们在痛苦中宣讲基督，在铁椅中仍有喜乐。为义人升起了什么光？不是为不义的人升起的那光；不是祂使太阳上升照耀善人和恶人的那光。有一种不同的光为义人升起；关于这光，不义的人最终要说：\BibleQuote{因此我们偏离了真理的道路，正义的光没有照耀我们，正义的太阳也没有为我们升起。}\BibleRef{智 5:6} 看，由于热爱这太阳，他们躺在心灵的黑暗中。他们用眼睛看见这太阳，却在心灵中没有看见那光，这有什么益处呢？多俾亚是瞎眼的，但他常教导他的儿子天主的道。你们知道，多俾亚警告他的儿子，对他说：\BibleQuote{孩子，你要用你的财物施舍；因为施舍能使人不至进入黑暗。}\BibleRef{多 4:7} 连那在黑暗中的人也这样说话……你愿意认识那光吗？要心正。什么是心正？在天主面前不要心怀歪曲，抗拒祂的旨意，想要使祂屈就于你，而不去管束自己以讨祂喜欢；那么你便会感受到所有心正的人所知的那份喜乐。\par

16. \BibleQuote{义人，你们要欢喜}\BibleRef{咏 96:12}。也许信友听到\Quote{要欢喜}这话时，已经在想宴席、准备酒杯、等待玫瑰花开的季节，因为经上说：\Quote{义人，你们要欢喜！}请看下文：\Quote{你们要因主而欢喜。}你们等待春天到来，好能欢喜：你们有主作为欢喜的缘由，主永远与你们同在，祂没有特定的季节；你们在夜间有祂，在白天也有祂。要心正；那么你们便从祂那里永远得到喜乐。因为那属于世界的喜乐不是真正的喜乐。请听先知依撒意亚的话：\BibleQuote{我的天主说：恶人没有喜乐。}\BibleRef{依 57:21} 恶人所谓的喜乐不是喜乐，如同那不以他们的喜乐为意的人所知道的那样：让我们相信他吧，弟兄们。他是一个人，但他知道两种喜乐。他确实知道杯中的喜乐，因为他是一个人，他知道餐桌的喜乐，他知道婚姻的喜乐，他知道那些世俗和奢华的喜乐。知道这些的人有信心地说：\Quote{主说：恶人没有喜乐。}但这不是人在说话，而是主……但你说：我看不见依撒意亚所看见的那光。相信吧，你就会看见。因为也许你没有看见它的眼睛；因为那是一种眼睛，藉以辨别那美。正如有肉身的眼睛，藉以看见这光；同样，有心灵的眼睛，藉以感知那喜乐：也许那眼睛受伤了，因情欲、贪婪、放纵、愚昧的欲望而昏昧、烦乱；你的眼睛被扰乱了：你不能看见那光。相信吧，在你看见之前；你将会被医治，并要看见。\par

17. \Quote{并且称颂祂圣德的纪念。}现在你们既已喜乐，又在主内欢欣，就当向祂称颂；因为若不是祂的旨意，你们不会在祂内喜乐。因为主亲自说：\BibleQuote{我给你们讲了这一切，是要你们在我内有平安。但在世界上你们将受磨难。}\BibleRef{若 16:33}如果你们是基督徒，就要在世上期待磨难；不要期待更平安更好的时期。弟兄们，你们欺骗了自己；福音所没有应许你们的，你们不要向自己应许。你们知道福音所说的；我们是对基督徒说话；我们不该违背信德。福音这样说：在末期，许多恶事、许多绊脚石、许多磨难、许多不义，将大量增加；但坚持到底的，必要得救。\BibleRef{玛 24:3}\Quote{许多人的爱}，它说，\Quote{将冷淡。}所以凡是在心灵上始终热切的人，正如宗徒所说：\BibleQuote{心灵热切，}\BibleRef{罗 12:11}他的爱就不会冷淡；因为\BibleQuote{天主的爱借著所赐予我们的圣神，已倾注在我们心中了。}\BibleRef{罗 5:5}所以，任何人都不要向自己应许福音所没有应许的。看，更幸福的时期将要来到，我正在这样做、购买这个。你们听从那位不被欺骗、也从未欺骗过任何人的，祂应许你们的喜乐不在这里，而是在祂自己内，这为你们是有益的；当这里的一切都过去后，希望你们将永远与祂一同为王；免得你们渴望在这里为王，却既不能在此享受喜乐，也不能在那里找到喜乐。\par

\SourceRecord{Translated by J.E. Tweed.；Nicene and Post-Nicene Fathers, First Series；Vol. 8.；Edited by Philip Schaff.；Buffalo, NY: Christian Literature Publishing Co.,；1888.；<http://www.newadvent.org/fathers/1801097.htm>.}

% structure: p0001 source=h1 target=section reason=page-title

\section{\WorkTitle{圣咏第九十八篇注释}}

1. \BibleQuote{你们应向上主高唱新歌} \BibleRef{咏 97:1}。新人知道这歌，旧人却不知道。旧人是旧生命，新人是新生命；旧生命来自亚当，新生命在基督内形成。但在这篇圣咏中，普世都被命令高唱新歌。别处更清楚地如此说：\BibleQuote{你们应向上主高唱新歌；普世大地，请向上主歌唱}；这样，那些自绝于普世共融的人，可以明白他们不能唱新歌，因为这歌是在全体中唱，而非在部分中唱。请注意此处，也看到这话是真的。当普世大地被命令高唱新歌时，意思是和平在唱新歌。\BibleQuote{因为祂行了奇事}。什么奇事？看，刚读了福音，我们听到了主的奇事。一位寡妇的独生子死了，被抬出来；主动了怜悯，叫他们站住；他们放下他，主说：\BibleQuote{青年人，我对你说，起来。} \BibleRef{路 7:12}……\BibleQuote{主行了奇事}。什么奇事？听：\BibleQuote{祂的右手和祂的圣臂为祂施行了拯救}。主的圣臂是什么？是我们的主耶稣基督。听依撒意亚说：\BibleQuote{谁会相信我们的报道呢？上主的臂膀又向谁显示了呢？} \BibleRef{依 53:1} 因此，祂的圣臂和祂的右手就是祂自己。我们的主耶稣基督就是天主的臂膀、天主的右手；故此说\BibleQuote{为祂施行了拯救}。不是只说\BibleQuote{祂的右手拯救了世界}，而是\BibleQuote{为祂施行了拯救}。因为许多人为自己得救，而非为祂。看，多少人渴望身体健康，并从祂那里获得健康：他们被祂治愈，却不是为祂。他们如何被祂治愈却不属于祂？当他们获得健康后，变得放荡：生病时贞洁的人，痊愈后却成了奸淫者；患病时不伤害任何人的人，恢复力量后攻击和压迫无辜者：他们被治愈了，却不是为祂。谁是为祂得治愈的人？是那内心得治愈的人。谁内心得治愈？是那信赖祂的人，当他内心得治愈、更新为新人后，这终有一日会衰弱的必死肉体，最终也将恢复其最完美的健康。因此，让我们为祂得治愈。但为了我们能为祂得治愈，让我们相信祂的右手。\par

2. \BibleQuote{上主显示了自己的救恩} \BibleRef{咏 97:2}。这右手、这臂膀、这救恩，就是我们的主耶稣基督，关于祂曾说：\BibleQuote{凡有血肉的，都要看见天主的救恩；} \BibleRef{路 3:6} 那位怀抱婴儿的西默盎也说：\BibleQuote{主啊！现在可照祢的话，放祢的仆人平安去了；因为我亲眼看见了祢的救恩。} \BibleRef{路 2:28} \BibleQuote{上主显示了自己的救恩}。祂向谁显示？向一部分，还是全体？不是向任何部分特别显示。没有人可以出卖、欺骗、说：\BibleQuote{看，基督在这里，或在那里}； \BibleRef{玛 24:23} 说\BibleQuote{看，祂在这里或在那里}的人，是指某些具体地点。\BibleQuote{上主向谁显示了自己的救恩}？听下文：\BibleQuote{祂在外邦人眼前公开显示了自己的义德}。我们的主耶稣基督就是天主的右手、天主的臂膀、天主的救恩、天主的义德。\par

3. \BibleQuote{祂向雅各伯忆起了自己的仁慈，向以色列家忆起了自己的真理} \BibleRef{咏 97:3}。这\BibleQuote{祂忆起了自己的仁慈和真理}是什么意思？祂怜悯了，因而应许；因为祂应许并显示了祂的仁慈，真理随之而来：仁慈走在应许之前，应许在真理中得以实现……\par

\BibleQuote{向以色列家忆起了自己的真理}。这以色列是谁？免得你们或许以为是指一个犹太民族，听下文：\BibleQuote{普世四极看见了我们天主的救恩。} 没有说\BibleQuote{普世大地}，而是\BibleQuote{普世四极}；正如所说的，从一端到另一端。没有人可以切割，没有人可以分散；基督的合一何等坚强。祂付出了如此大的代价，买了全体：\BibleQuote{普世四极}。\par

4. 因为他们看见了，所以\BibleQuote{普世大地，请向上主欢呼} \BibleRef{咏 97:4}。你们已经知道欢呼是什么意思：欢乐并说话。如果你们无法表达自己的喜乐，就呼喊吧；若言语不能，让呼喊表明你们的喜乐；但不要让喜乐沉默；不要让你的心对它的天主沉默，不要对祂的恩赐默默无声。如果你对自己说，你是为自己得治愈；如果祂的右手为祂治愈了你，你就为那治愈你的祂说话。\BibleQuote{歌唱，欢呼，奏乐。}\par

5. \BibleQuote{用竖琴向上主歌咏：以竖琴和诗篇的声音} \BibleRef{咏 97:5}。不要只用声音赞美祂；拿起行动，使你们不仅歌唱，也工作。那歌唱并工作的，便用琴瑟和谐地奏乐。现在来看，接下来所提到的器具，是一种象征：\BibleQuote{也以铜号与角声} \BibleRef{咏 97:6}。铜号与角声是什么？铜号是黄铜制成的：它们由锤打而成；若由锤打，即被击打，你们将成为铜号，被拉伸为对天主的赞美，若你们在磨难中改进：磨难是锤打，改进则是被拉伸。\Person{约伯}是一支铜号，当他突然遭受最沉重的损失和儿子的死亡时，因如此沉重的磨难之苦而成为铜号，他发出这样的声音：\BibleQuote{上主赐予，上主收回；愿上主的名受赞美。} \BibleRef{约 1:21} 他如何发声？他的声音多么悦耳？这支铜号仍在锤打之下……我们已经听到他是如何被锤打的；让我们听他如何发声：让我们，如果你们愿意，听一听这支铜号的悦耳之声：\Quote{难道我们只从天主手里接受福，而不接受祸吗？} 啊，勇敢的，悦耳的声音！那声音岂不唤醒沉睡之人？岂不唤起对天主的信赖，使人毫无畏惧地向魔鬼进军；不是靠自己的力量，而是靠那试探他者的力量作战。因为正是祂在锤打：因为锤子本身不能这样做……看，（我敢如此说，我的弟兄们）连宗徒也曾被这同一锤子击打：他说：\BibleQuote{有一根刺加在我身上，就是撒殚的使者，来攻击我。} \BibleRef{格后 12:7} 看，他正在锤下：让我们听他是如何说的：\Quote{关于这事，我三次求主使它离开我。祂对我说：我的恩宠为你足够了，因为我的能力在软弱中才完全。} 祂，我的造物主说：我要使这支号角完美；除非我锤打它，否则我无法做到；在软弱中能力才得以成全。现在请听这支铜号正确地发声：\Quote{因为我软弱时，正是我刚强时。}……\par

6. 角管之声是什么？角高过于肉身：若高过于肉身，它就必须坚固持久，并能说话。这是从何而来？因为它超越了肉身。谁愿成为角号，就当克胜肉身。这是什么意思？克胜肉身，即超越欲望，征服肉身的私欲。请听角号……\Quote{你们要思念上面的事}是什么意思？意思是，超越肉身，不要思想属肉体的事。他们还不是角号，所以他对他们这样说：\Quote{弟兄们，我从前不能把你们当作属神的，只得把你们当作属血肉的，当作在基督里的婴孩。我喂养你们的是奶，不是干粮；因为你们那时还不能吃，就是现在还是不能。你们仍是属血肉的。} 因此，他们不是角号，因为他们没有超越肉身。角既附着于肉身，又高过于肉身；虽然从肉身生出，却超越了它。因此，如果你们是属神的，而从前是属血肉的；你们现在虽在肉身中行走于大地，但在神内却升向天堂；因为我们虽然在肉身内行走，却不按肉身作战……弟兄们，不要责备那些尚未蒙天主仁慈悔改的弟兄；要知道，你们若这样做，便是属血肉的。那号角并不是天主耳中所悦纳的：自夸的号角使战争徒劳无功。愿角号鼓起你们的勇气对抗魔鬼；不要让属血肉的号角使你们骄傲，反对你们的弟兄。\Quote{在上主君王面前欢呼吧。}\par

7. 当你们欢欣鼓舞，因着铜号与角声而喜乐时，接着是什么？\BibleQuote{愿海洋和其中的一切澎湃} \BibleRef{咏 97:7}。弟兄们，当宗徒们如同铜号与角号宣讲真理时，海洋就澎湃了，波浪翻腾，风暴加剧，教会的迫害接踵而至。海洋为何澎湃？因为在天主面前发出欢呼，感恩的圣咏正在歌唱：天主的耳朵喜悦了，海洋的波浪就被掀起。\BibleQuote{愿海洋和其中的一切澎湃；世界和其中的居民。} 愿海洋在迫害中澎湃。\BibleQuote{愿江河拍手} \BibleRef{咏 97:8}。愿海洋被激动，江河拍手；迫害兴起，圣徒们却在天主内喜乐。江河拍手是什么意思？是在行动中喜乐。拍手是喜乐；手代表行动。哪些江河？那些天主赐予那水，即圣神，使他们成为江河的人。祂说：\BibleQuote{若是有人渴了，就到我这里来喝。信我的人，从他的心中要流出活水的江河。} \BibleRef{若 7:37} 这些江河拍手，这些江河因行动而喜乐，赞美天主。\Quote{诸山要一同欢乐。}\par

8. \BibleQuote{在主面前，因为祂已来临；祂已来临，要审判大地}\BibleRef{咏 97:9}。 \Quote{山岳}象征伟大。主来审判大地，它们就欢喜。但有些山岳，当主来审判大地时，将要战栗。因此有善与恶的山岳：善的山岳是属灵的伟大；恶的山岳是骄傲的膨胀。\Quote{愿山岳在主面前一同喜乐，因为祂已来临；祂已来临，要审判大地。}祂为何来临？又将如何来临？\BibleQuote{祂要以正义审判世界，以公平治理万民}\BibleRef{咏 97:10}。因此愿山岳喜乐，因为祂不会不公正地审判。当某个人作为审判者来临，良心无法向他敞开时，即使无辜的人也可能战栗，如果期望从祂得到美德的奖赏，或惧怕刑罚的定罪；当那位不能被欺骗的祂来临之时，愿山岳喜乐，愿它们无忧无惧地喜乐；它们将被祂光照，而不是被定罪；愿它们喜乐，因为主将以公平审判世界；如果公义的山岳喜乐，愿不义的山岳战栗。但是请看，祂尚未来临：它们有什么必要战栗呢？让它们改过自新，并喜乐吧。你以何种方式等待基督的来临，这在于你。正因如此，祂延迟来临，以便祂来临时不谴责你。看，祂尚未来临：祂在天上，你在地上：祂延迟祂的来临，你不要延迟智慧。祂的来临对心硬者是严峻的，对虔诚者是温柔的。因此，甚至现在看看你是什么：如果你心硬，你可以软化；如果你温柔，甚至现在就可以因祂将来临而喜乐。因为你是一名基督徒。是的，你说，我相信你祈祷，并说：\BibleQuote{愿祢的国来临}\BibleRef{玛 6:10}。你渴望祂来临，却害怕祂的来临。改变你自己，免得你祈祷反对你自己。\par

\SourceRecord{Translated by J.E. Tweed.；Nicene and Post-Nicene Fathers, First Series；Vol. 8.；Edited by Philip Schaff.；Buffalo, NY: Christian Literature Publishing Co.,；1888.；<http://www.newadvent.org/fathers/1801098.htm>.}

% structure: p0001 source=h1 target=section reason=page-title

\section{圣咏第九十九篇讲道集}

1. 亲爱的弟兄们，身为教会的子女，并在基督的学校里通过我们古代先贤的著作——他们写下了天主的话语和天主的伟大事物——受到良好教导的你们，应当明白：他们愿意为我们这些生于此世、信仰基督的人谋取益处；祂在合宜的时候来到我们中间，第一次来是谦卑的，第二次来则是要带着荣耀而来……因为圣咏集中如此说：\BibleQuote{真理从地上发出，正义从天上俯视。}因此，我们全部的心意是，当我们听到一篇圣咏、一位先知或法律书——所有这些都是在我们的主耶稣基督降生成人之前写成的——要在其中看到基督，在其中理解基督。因此，亲爱的弟兄们，请与我一同留意这篇圣咏，让我们在这里寻找到基督；祂肯定会向那些寻求祂的人显现，祂起初曾向那些不求祂的人显现；祂不会弃绝那些渴望祂的人，祂曾救赎了那些忽略祂的人。看，圣咏开始论及祂；关于祂如此说：——\par

2. \BibleQuote{上主为王，万民当战栗} \BibleRef{咏 98:1}。因为我们的主耶稣基督开始为王，开始被传讲，是在祂从死者中复活、升天之后，在祂以圣神的勇气充满祂的门徒之后，使他们不再害怕死亡——祂已在祂自己内战胜了死亡。我们的主基督于是开始被传讲，使那些渴望救恩的人可以相信祂；而崇拜偶像的万民却愤怒了。他们因自己所造的受到敬拜而愤怒，因为那造他们的那位被宣告了。事实上，祂通过祂的门徒宣告了祂自己，祂愿意他们归向造他们的主，并远离他们自己制造的东西。他们为偶像的缘故向他们的主发怒；他们即使为偶像的缘故向自己的奴隶发怒，也当受谴责。因为他们的奴隶比他们的偶像更好：因为天主造了他们的奴隶，而木匠造了他们的偶像。他们为偶像的缘故如此愤怒，以致不怕向他们的主发怒。但\Quote{万民当战栗}这话是预言，而非命令；因为这句话是在预言中说出的：\BibleQuote{上主为王，万民当战栗。}甚至从愤怒的百姓中也产生一些好的结果：让他们愤怒吧，在他们的愤怒中，殉道者们得以加冕……你们在宗徒书信宣读之前听到耶肋米亚被宣读时（如果你们听了），你们在其中看到了我们现在所处的时代。祂说：\BibleQuote{那些没有创造天地的神祇，应从地上、从天下消灭} \BibleRef{耶 10:11}。祂没有说：那些没有创造天地的神祇，应从天上和地上消灭；因为它们从未在天上：但祂说了什么？\Quote{应从地上，从天下消灭。}好像，虽然\Quote{地}这个词重复了，但\Quote{天}这个词的重复却缺失了（因为它们从未在天上）：祂两次重复\Quote{地}，因为它是在天下。\Quote{应从地上，从天下消灭，}从它们的庙宇中。思考一下，这是否正在实现？是否绝大部分已经发生？因为还有什么，或者说多少，还存留着呢？偶像更多地存留在外邦人的心中，而非庙宇的壁龛中。\par

3. \Quote{祂坐在革鲁宾之上：}你应理解，\Quote{祂是王：愿大地震动。}……\OriginalTerm{革鲁宾}{Cherubim}是天主的宝座，正如圣经向我们显示的，是一个崇高的天上宝座，我们看不见；但天主的圣言认识它，认识它作为自己的宝座；而且天主的圣言和天主的神亲自向天主的仆人们启示了天主坐在哪里。并非天主像人一样坐着；但你，如果你愿意天主住在你内，如果你愿意成为善的，你将成为天主的宝座；因为经上如此记载：\BibleQuote{义人的灵魂是智慧的宝座} \BibleRef{箴 12:23}。因为\Quote{宝座}在我们的语言中称为座位。有些通晓希伯来语的人将\Quote{革鲁宾}翻译成拉丁语（因为它是一个希伯来词语），意为\Quote{知识的充满}。因此，因为天主超越一切知识，祂被说成坐在知识的充满之上。所以，愿知识的充满在你内，那么甚至连你自己也将成为天主的宝座……祂知道一切：因为我们的头发在天主面前都被数过 \BibleRef{玛 10:30}。但祂愿意人知道的知识的充满与此不同；祂愿意你拥有的知识，属于天主的法律。你或许会对我说：谁能完全认识法律，以至他内心拥有法律知识的充满，并能成为天主的宝座呢？不要困扰；简短地告诉你，如果你愿意拥有知识的充满并成为天主的宝座，你该做什么：因为宗徒说：\BibleQuote{爱是法律的满全} \BibleRef{罗 13:10}。那么接下来呢？你失去了所有借口。问问你的心；看看它是否有爱。若那里有爱，那里也就有法律的满全；天主已经住在你内，你已成为天主的宝座。\Quote{万民当战栗；}愤怒的百姓能对那已成为天主的宝座的人做什么呢？你留意那些向你发怒的人；但那位坐在你内的是谁，你却不留意。你已成为天，还害怕地吗？因为圣经在另一处说，上主我们的天主宣告：\BibleQuote{天是我的宝座} \BibleRef{依 66:1}。因此，如果你也通过拥有知识的充满和拥有爱，成了天主的宝座，你就成了天。因为我们用肉眼所仰望的这个天，在天主面前并不十分珍贵。圣洁的灵魂是天主的天；天使的心灵和祂所有仆人的心灵，都是天主的天。\par

4. \BibleQuote{主在熙雍为大，超乎万民之上} \BibleRef{咏 98:2}…… 我曾对你们说过，那高于革鲁宾的祂，在熙雍是伟大的。你们现在可以问：什么是熙雍？我们知道熙雍是天主的城。耶路撒冷城被称为熙雍；按照某种解释，熙雍意为守望，即凝视与默观；因为守望就是前瞻、注视或定睛观看。每一个灵魂，若它试图看见那当被看见的光，便是熙雍。因为若它注视了自己的光，它便陷入黑暗；若注视了祂的光，它便被光照。但既然清楚熙雍是天主的城，天主的城不是圣教会又是甚么？因为彼此相爱、且爱那居于他们内的天主的人们，构成了属于天主的城。因为一座城是由某种法律凝聚的；他们的法律就是爱；而那爱就是天主：因为经上清楚地写着：\BibleQuote{天主是爱。} \BibleRef{若壹 4:8} 因此，充满爱德的人充满天主；而众多充满爱德的人，构成了一座充满天主的城。那座天主的城称为熙雍；所以教会就是熙雍。天主在其中是伟大的……\par

5. 弟兄们，你们以为那些昨日乐器齐鸣的人，不会对我们的斋戒发怒吗？但我们不要向他们动怒，却要为他们守斋。因为住在我们内的主、我们的天主曾说过，祂亲自命令我们为仇人祈祷，为迫害我们的人祈祷：\BibleRef{玛 5:44} 而教会这样做了，迫害者就几乎灭绝了…… 醉酒者并不冒犯自己，却冒犯清醒的人。请向我指出一个终于在天主内得享幸福、生活庄重、叹息渴望天主所应许的永恒平安的人；你会看到，当他看见一个人随着乐器起舞时，他为那人的疯狂所感到的悲伤，远甚于为一个因热病而狂乱的人。如果我们知道他们的邪恶，想到我们自己也从那些邪恶中被释放出来，就让我们为他们悲伤；如果我们为他们悲伤，就让我们为他们祈祷；而为了我们的祈祷得到垂听，就让我们为他们守斋。因为在他们的节日里，我们并不守我们自己的斋戒。我们在临近逾越节的日子、在基督为我们确定的各个时期所守的斋戒是不同的；但在他们的节日里，我们守斋是因为当他们欢乐时，我们为他们而叹息。因为他们的欢乐激发我们的悲伤，使我们想起他们至今仍是如何可怜。但由于我们看到许多人已经从那里获得释放，我们自己也曾在其中，所以我们也不应对他们绝望。如果他们仍然愤怒，让我们祈祷；如果仍有一粒残留的尘土在我们面前扬起，让我们继续为他们哀哭，好使天主也赐给他们领悟，使他们与我们一同听见这些话语——这些我们此刻正为此欢欣的话语。\par

6. 所有这些人民，就是祢在熙雍之上为大的，\BibleQuote{愿他们称谢祢伟大的圣名。} \BibleRef{咏 98:3} 祢的名在他们发怒时是小的；如今已成为大的。现在就让他们承认吧。在何种意义上我们说基督的名曾是小的，在它被如此广大地传播之前？因为祂的名声就是指祂的名。祂的名曾是小的，如今已成了大的。有哪个民族没有听说过基督的名呢？因此，让那些从前对祢的小名发怒的人民，如今称谢祢伟大的圣名。他们为何要称谢？因为它是\Quote{奇妙而神圣的。} 祢的名本身是奇妙而神圣的。祂被宣讲为被钉十字架者，被宣讲为被贬抑者，被宣讲为受审判者，为的是祂能高升而来，为的是祂能活着而来，为的是祂能以权能审判。祂目前宽恕那亵渎祂的人民，因为\BibleQuote{天主的恩慈是引人悔改。} \BibleRef{罗 2:4} 因为那现在宽恕的，不会永远宽恕；那如今被宣讲为当受敬畏的，不会不来审判。祂必来临，我的弟兄们，祂必来临：让我们敬畏祂，并如此生活，好使我们能在祂右边被寻见。因为祂必来审判，将一些人安置在左边，一些人安置在右边。\BibleRef{玛 25:31} 祂行事并不模棱两可，以致在人群中误判，使本应站在右边的人站在左边，或使本应站在左边的人因天主的错误而站在右边：祂不可能出错，以致将恶人放在该放善人的地方，或将善人放在该放恶人的地方。若祂不会出错，那么，若我们不怕，我们便出错；但若我们在此生恐惧，那时我们便无所畏惧。\Quote{因为王上喜爱公正。}\par

7. \Quote{祢已准备好了公正；祢在雅各伯中施行了审判和正义。} 因为我们也应当有审判，应当有正义；但祂在我们内施行审判和正义，那创造了我们、并使我们能在其中行事的祂。我们应当如何也有审判和正义呢？当你区分善恶时，你就有审判；当你追随善、远离恶时，你就有正义。藉着区分，你有审判；藉着实行，你有正义。经上说：\Quote{避恶，}他说，\Quote{行善，寻求和平，追随不懈。} 你应先有审判，再有正义。什么审判？就是你先判断什么是恶，什么是善。又有什么正义？就是你要避恶行善。但这并非你能从自己获得的；请看祂所说：\Quote{祢在雅各伯中施行了审判和正义。}\par

8. \BibleQuote{你们要尊崇主我们的天主} \BibleRef{咏 98:5}。真正尊崇祂，妥善尊崇祂。让我们赞美祂，让我们尊崇那在我们内成就了我们所有义德的祂；是祂自己在我们内成就了它。因为除了那使我们成义的祂，谁能在我们内成就义德呢？因为论及基督，经上说：\BibleQuote{祂使不虔敬的人成义。} \BibleRef{罗 4:5}……\BibleQuote{在祂脚凳前俯伏朝拜，因为祂是圣的。} \BibleRef{咏 98:5} 我们要在什么面前俯伏朝拜？祂的脚凳。脚底下的东西称为脚凳，希腊文是\OriginalTerm{脚凳}{\GreekText{ὑποπόδιον}}，拉丁文是\OriginalTerm{脚凳}{Scabellum}或\OriginalTerm{脚凳}{Suppedaneum}。但是，弟兄们，考虑一下祂命令我们在什么面前俯伏朝拜。在圣经的另一处经文中说：\BibleQuote{天是我的宝座，地是我的脚凳。} \BibleRef{依 66:1} 那么，祂是否命令我们崇拜地，因为另一处经文说它是天主的脚凳呢？既然如此，我们怎能崇拜地，而圣经明明说：\BibleQuote{你要崇拜主你的天主} \BibleRef{申 6:13}？然而这里说：\BibleQuote{在祂脚凳前俯伏朝拜；} \BibleQuote{地是我的脚凳。} 我疑惑；我害怕崇拜地，唯恐创造天地的主定我的罪；然而，我又不敢不崇拜我主的脚凳，因为圣咏命令我：\BibleQuote{在祂脚凳前俯伏朝拜。} 我问，祂的脚凳是什么？圣经告诉我：\BibleQuote{地是我的脚凳。} 在犹豫中，我转向基督，因为我在其中寻找祂自己；于是我发现了如何可以虔诚地崇拜地，如何可以虔诚地崇拜祂的脚凳。因为祂从地上取了泥土；因为肉体来自泥土，而祂从玛利亚的肉体取了肉体。因为祂在此行走在真实的肉体中，并将那真实的肉体赐给我们吃，以得救恩；没有人吃那肉体，除非他先朝拜了它；我们发现了主这样的脚凳在什么意义上可以被崇拜，并且不仅朝拜它不犯罪，而且不朝拜它反而犯罪。但是肉体能赐予生命吗？主自己，当祂赞扬同一块泥土时，说：\BibleQuote{使人活的是圣神，肉体是无益的。}……但当主赞扬它时，祂是在谈论祂自己的肉体，并且祂曾说：\BibleQuote{人若不吃我的肉，就没有生命在他里面。} \BibleRef{若 6:54} 祂的一些门徒，大约七十人，感到冒犯，说：\BibleQuote{这话很难，谁能听呢？} 他们就退去了，不再与祂同行。在他们看来很难，因为祂说：\BibleQuote{你们若不吃人子的肉，就没有生命在你们里面。} 他们愚昧地接受，肉欲地思考，以为主会从祂身上割下部分，给他们；他们说：\BibleQuote{这话很难。} 是他们自己硬，而不是这话硬；因为他们若不是硬的，不温柔，他们会对自己说：\Quote{祂说这话不是没有理由的，这里面一定隐藏着奥秘。} 他们会留下与祂同在，被软化而不硬；并且会从祂那里学到那些留下的人，当其他人离开时，所学到的东西。因为当十二门徒留下与祂同在时，在他们离开后，这些留下的追随者向祂暗示，仿佛对前者的死亡悲伤，他们被祂的话冒犯并回转了。但是祂教导他们，对他们说：\BibleQuote{使人活的是神，肉体是无益的；我对你们所说的话就是神，就是生命。} \BibleRef{若 6:63} 要属灵地理解我所说的；你们不是要吃你们所看见的这身体，也不是要喝他们将要钉死我的人所流出的血。我已将一种奥秘托付给你们；属灵地理解，它将赐予生命。虽然这事需要被可见地举行，但仍必须属灵地理解。\par

9. \BibleQuote{梅瑟和亚郎在祂的司铎中，撒慕尔在呼求祂名的人中：这些人都呼号了主，祂就俯听了他们}\BibleRef{咏 98:6}。\BibleQuote{祂从云柱中对他们说话}\BibleRef{咏 98:7}……关于梅瑟，那里并没有说他是司铎。但若不是，他是什么？他能比司铎更伟大吗？这篇圣咏宣称他自己也是司铎：\BibleQuote{梅瑟和亚郎在祂的司铎中}。因此他们是主的司铎。撒慕尔后来在\BibleBook{列王}{纪}中读到：这位撒慕尔生活在达味时代，因为他傅油祝圣了圣达味。撒慕尔自幼在圣殿中长大……他提及这些人，并借此愿我们理解所有圣人。但为何他在这里特别提到这些名字？因为我们已经说过，我们应当在这里理解基督。请听，圣善的弟兄们。他上面说过：\BibleQuote{请尊崇主我们的天主，在祂的脚凳前俯伏朝拜，因为祂是圣的}，这是赞扬某一位，即我们的主耶稣基督；祂的脚凳应当受朝拜，因为祂披戴了肉躯，为出现在人类面前；并愿向我们显示，古代教父也曾宣讲过祂，因为我们的主耶稣基督就是那位真实的司铎，他提及这些人，是因为天主曾从云柱中对他们说话。\Quote{从云柱中}是什么意思？祂是借比喻说话。因为若祂在某种云中说话，那些隐晦的话语是预言某位尚不为人知、却将要显现者。这位未知者已不再未知，因为祂为我们所知，即我们的主耶稣基督……当初从云柱中说话的那位，如今亲自在祂的脚凳上对我们说话，即在地上，当祂披戴了肉躯时，因此我们朝拜祂的脚凳，因为祂是圣的。祂自己曾在云中说话，那时未被理解；祂已在祂的脚凳上说话，而祂的云中的话已被理解。\BibleQuote{他们遵守了祂的证言和祂所赐给他们的法律}……\BibleQuote{祢俯听了他们，}他说，\BibleQuote{主，我们的天主，祢宽恕了他们，天主}\BibleRef{咏 98:8}。天主被说是宽恕，只指向罪；当祂赦免罪过时，祂才宽恕。而他们身上有什么需受惩罚，以致祂在宽恕中赦免？祂在赦免他们的罪时是宽恕的，并在惩罚中也是宽恕的。因为接着是什么？\BibleQuote{祢惩罚了他们一切的私欲偏情。}即使在惩罚中，祢也对他们宽恕：因为不仅在赦免，而且也在惩罚他们的罪时，祢已是宽恕的。请思考，我的弟兄们，祂在此教导了我们什么：请听。天主向那个犯罪而不受鞭打的人发怒：因为对于祂真正宽恕的人，祂不仅赦免了罪，使它们在未来生命中不伤害他，而且亦惩戒他，使他不喜于继续犯罪。\par

10. 来，我的弟兄们；若我们问这些人如何受了惩罚，主会帮助我来告诉你们。让我们思考这三个人，\Person{梅瑟}、\Person{亚郎}和\Person{撒慕尔}：他们如何受了惩罚？因为他说：\BibleQuote{祢惩罚了他们一切的私欲偏情}，意指那些私欲偏情，是主在他们心中所知的，而人们不知的。因为他们生活在天主子民中间，无人抱怨。但我们能说什么呢？或许梅瑟的早期生活是有罪的，因为他杀了人后逃出埃及。\BibleRef{出 2:12} 亚郎的早期生活也不蒙天主喜悦，因为他允许一个疯狂而迷失的子民造一个偶像来朝拜，\BibleRef{出 32:1}竟为天主子民造了偶像来朝拜。撒慕尔呢？他自幼被献于圣殿，他一生都在天主的圣事中度过，自幼便是天主的仆人。从未有人说过撒慕尔什么，人们从未指责什么。或许天主知道那里有什么需要惩戒的，因为甚至在人眼中完美的，在那完美面前仍是不完美的。工匠将他们许多作品展示给外行人看，当外行人称其为完美时，工匠仍继续修饰，因为他们知道作品还缺少什么，以致人们惊讶于那些本以为已完美的作品又接受了如此多的修饰。这发生在建筑、绘画、刺绣以及几乎每一种艺术中。起初他们判断它已达到某种完美，以致他们的眼睛别无所求；但外行人眼睛的判断是一回事，艺术法则的判断是另一回事。同样，这些圣人在天主眼前生活，仿佛无过，仿佛完美，仿佛天使；但那位惩罚了他们一切私欲偏情的那一位，知道他们欠缺什么。但祂不是因愤怒而惩罚，而是出于仁慈：祂惩罚他们，为使他们完美祂已开始的工作，而非定他们所弃绝的罪。因此天主惩罚了他们一切的私欲偏情。祂如何惩罚了撒慕尔？惩罚在哪里？……对梅瑟所说的是一种预像，而非惩罚。对一位老人来说，死亡是什么惩罚？不能进入那地（不配的人却进去了）是什么惩罚？至于亚郎，他也年老而死，他的儿子继承了他的司铎职，他的儿子后来管理司铎职：天主又如何惩罚了亚郎？撒慕尔也圣洁地老去，留下他的儿子们作为继承人。我寻找施加给他们的惩罚，按人的看法，我找不到；但按照我所知天主的仆人每日所受的苦，他们日日受惩罚。请读，并看那些惩罚，你们这些有长进的也要承担这些惩罚。他们每日忍受固执的子民，每日忍受不虔敬的生活者；被迫生活在那些其生活每日受他们指责的人中间。这就是他们的惩罚。对某人而言这惩罚是小的，那是他尚未长进；因为别人的不虔敬折磨你，程度与你远离你自己的不虔敬成正比。……\par

11. \BibleQuote{显扬上主我们的天主！}\BibleRef{咏 98:9}。我们再次显扬祂。祂在惩罚时仍显仁慈，该如何赞美祂，该如何显扬祂？你岂能向你的儿子这样行，而天主却不能吗？因为当你抚慰你的儿子时，你并非善，当你责打他时，你并非恶。你抚慰他时，你是父亲；你责打他时，你也是他的父亲：你抚慰他，使他不要气馁；你责打他，使他不要丧亡。\Quote{显扬上主我们的天主，在祂的圣山上朝拜祂，因为上主我们的天主是圣的。}正如他上面所说：\Quote{显扬上主我们的天主，俯伏在祂的脚凳前}：现在我们明白了朝拜祂脚凳的含义；同样，现在他在显扬上主我们的天主之后，为使无人能在祂的圣山之外显扬祂，也赞美了祂的圣山。祂的圣山是什么？我们在别处读到关于这圣山的记载：一块\OriginalTerm{非人手所凿的石头}{stone cut without hands}从山上被凿出，击碎了地上所有王国，那石头本身长大了。这就是我正叙述的达尼尔异象。这块非人手所凿的石头长大了，他说：\BibleQuote{变成了一座大山，充满了全地。}\BibleRef{达 2:34}我们若希望被垂听，就当在那大山上朝拜。异端者不在那山上朝拜，因为它充满了全地；他们固守于它的一部分，却失去了整体。如果他们承认天主教会，就会与我们一同在这圣山上朝拜。因为我们已看见那块非人手所凿的石头如何长大，已覆盖多么广阔的地域，伸展到何等众多的民族。那山，即非人手所凿之石被凿出的山，是什么？首先是犹太王国，因为他们崇拜独一的天主。从那里凿出了那块石头，即我们的主耶稣基督……那块石头于是从山上非人手所凿而生；它长大了，藉着它的长大击碎了地上所有的王国。它已成为一座大山，充满了全地。这就是天主教会，你们应喜乐于与它共融。但那些不在其共融中的人，因为他们在这同一座山之外敬拜和赞美天主，他们的祈祷不获垂听而至于永生；尽管他们可能在某些暂时之事上被垂听。他们不要自夸，因为天主在某些事上垂听他们：祂也垂听外邦人在某些事上的祈求。外邦人向天主呼号，岂不降雨？为何？因为\BibleQuote{祂使太阳升起，光照善人，也照恶人；降雨给义人，也给不义的人。}\BibleRef{玛 5:45}所以，外邦人啊，不要夸耀当你向天主呼号时，天主就降雨，因为祂降雨给义人，也给不义的人。祂在暂时之事上垂听你；祂不在永恒之事上垂听你，除非你在祂的圣山上朝拜。\Quote{在祂的圣山上朝拜祂，因为上主我们的天主是圣的。}……\par

\SourceRecord{Translated by J.E. Tweed.；Nicene and Post-Nicene Fathers, First Series；Vol. 8.；Edited by Philip Schaff.；Buffalo, NY: Christian Literature Publishing Co.,；1888.；<http://www.newadvent.org/fathers/1801099.htm>.}

% structure: p0001 source=h1 target=section reason=page-title

\section{圣咏第一百篇释义}

1. 弟兄们，你们刚才听到这圣咏在咏唱；它简短而不隐晦：就如同我已向你们保证，你们不应畏惧疲惫……\par

2. 这圣咏的标题是\Quote{宣认的圣咏}。诗句虽少，但蕴含着伟大的主题；愿种子在你们心中结出果实，谷仓为主预备好收成。\par

3. 因此，\BibleQuote{向主欢呼，普世大地}（\BibleRef{咏 99:1}）。这圣咏向我们发出劝谕，要我们向主欢呼。它并非劝谕大地的某一个角落，或某一个住区或人群；而是因为它知道自己在四面八方播下了祝福，所以它在四面要求欢呼。此刻普世大地都在听我的声音吗？然而普世大地都已经听到了这声音。普世大地已经因主而欢呼；而那些尚未欢呼的，也将会欢呼。因为祝福向四方扩展，当教会从耶路撒冷开始向万民传播时（\BibleRef{路 24:47}），到处推倒不敬，到处建立虔敬：善人与恶人混杂在各地。每一块土地都充满了恶人不满的怨声，也充满了善人的欢呼。那么，什么是\Quote{欢呼}呢？这圣咏的标题尤其使我们注意这个词，因为它被称为\Quote{宣认的圣咏}。那么，带着宣认而欢呼是什么意思呢？这是另一篇圣咏所表达的情感：\BibleQuote{懂得欢呼的人民是有福的}。毫无疑问，那因理解而有福的事是伟大的。因此，愿使我们有福的主、我们的天主，赐我理解该说什么，赐你们理解所听的：\BibleQuote{懂得欢呼的人民是有福的}。所以让我们奔向这福分，让我们理解欢呼，让我们不要无理解地倾泻欢呼。当这圣咏说\BibleQuote{向主欢呼，普世大地}，而我们不理解什么是欢呼，以致我们的声音可能欢呼而内心却不如此，这有什么用呢？因为理解是心灵的言语。\par

4. 我要讲的是你们所知道的。欢呼的人并不发出言语，而是一种没有言语的喜乐之声：因为它是一种心思倾注于喜乐中的表达，尽力表达情感，但却不能完全包含那感觉。一个人因自己的欢跃而喜乐，在说出某些无法言说或理解的话语之后，爆发出没有言语的欢跃之声，以至于他似乎在以声音本身欢呼，但因充满过度的喜乐而无法用言语表达那喜乐的缘由……那些在田地里劳作的人最常欢呼；收割者、采葡萄者或采集大地果实的人，因丰收而欢喜，并因土地的富饶与丰盛而欢跃，在欢呼中歌唱；在他们以言语唱出的歌曲中，他们在欢悦的心思中插入一些没有言语的呼喊；这就是欢呼的意义……\par

5. 那么，我们何时欢呼呢？当我们赞颂那无法言说者时。因为我们观察整个受造界，大地与海洋及其中所有的一切：我们观察到每样事物都有自己的根源与原因、产生的能力、出生的次序、存续的限度、死亡的终结；世世代代毫无混乱地运行；星辰似乎从东向西滚动，完成一年的循环；我们看到月份如何计算，小时如何延伸；在这一切之中，有一种不可见的因素，我不知道是什么，但有一种统一的原则，称为精神或灵魂，存在于所有生物中，驱策它们追求快乐、避免痛苦并保全自身；人类也与天主的天使有某种共同之处；不是与牲畜共有生命、听觉、视觉等等；而是拥有能够理解天主的东西，这东西专属心灵，能够辨别正义与不义，如同眼睛辨别黑白。在我所能做的这一切受造界的思考中，让灵魂自问：是谁创造了这一切？是谁造了它们？是谁在其中造了你？……我已经尽我所能观察了整个受造界。我观察了天上地下有形的受造物，以及在我自身中无形的受造物——我，正在说话，活动我的四肢，发出声音，移动舌头，说出话语，分辨感觉。我何时能在自己内理解自己呢？那么我怎能理解在我之上的呢？然而，看见天主是应许给人心的，并且吩咐了一种净化心灵的工作；这就是圣经的训诫。在你试图看见你所爱之前，先预备看见的途径。因为除了那不敬的人，他远远地离开了祂，谁不乐意听到天主和祂的圣名呢？……\par

6. 因此，你当在虔敬上效法祂，并热心默想：因为\Quote{祂那看不见的事，藉着所造之物，就可以明明地看见，被人所领悟}\BibleRef{罗 1:20}；观看所造之物，赞美它们，寻求它们的造作者。你若不像，就会退回；你若像，就会喜乐。当你像祂，开始接近祂并感受天主时，爱在你内越发增长，既然天主是爱，你将会领悟到一些你曾试图表达却无法言说的事。在你感受天主之前，你以为你能表达\Emph{天主}；你开始感受祂，便觉得你所感受的无法表达。但当你因此发现你所感受的无法表达时，你难道会缄默不语，不赞美天主吗？你难道会在赞美天主时沉默，不为那愿意使自己为你所认识的天主献上感恩吗？你寻求时赞美祂，找到时却沉默吗？绝不；你不会忘恩负义。祂应受尊崇，应受敬畏，应受极大的赞美。请想想你自己，看看你是什么：尘土和灰烬；看看是谁配得看见，看见什么；想想你是谁，看见什么——一个人看见\Emph{天主}！我不承认这人的功绩，只承认天主的仁慈。因此，赞美那慈悲者……\par

7. \Quote{你们当欢欣地服事主。}一切奴役都充满苦楚：所有身处奴役命运的人，既是奴隶，又是不满的。不要惧怕那一位主的奴役：那里没有呻吟，没有不满，没有愤慨；无人寻求被卖到另一个主人那里，因为这侍奉是甘甜的，因为我们都蒙受了救赎。弟兄们，在那伟大的家庭中为奴，即使带着锁链，也是极大的福乐。不要害怕，被捆绑的奴隶，向主认罪：将你的锁链归于你应得的报应；在你的锁链中认罪，如果你渴望它们变成装饰……你同时是奴隶，又是自由的：是奴隶，因为你是被造的；是自由的，因为你被天主所爱，你由祂所造：是的，确实是自由的，因为你爱那造你的主。不要带着不满服事；因为你的怨言并不能使你脱离服事，反而使你成为邪恶的仆人。你是主的奴隶，你是主的自由人：不要寻求被释放而离开那释放你的家……\par

8. 因此，他说：我要与少数善人分离而居：为何要与群众共处？好吧：那极少数善人，是从什么人群中筛选出来的？若这少数人都是善良的，那么，一个人选择与那些选择了安静生活的人在一起，远离民众的喧嚣、嘈杂的人群、人生的巨浪，他们犹如在港湾中，这固然是一个美好而可称赞的意愿。那么，那里有那应许的喜乐吗？那欢声的喜乐？尚未如此；仍有呻吟，仍有诱惑的焦虑。因为港湾总有一个入口；若没有，船只无法进入；所以它必须在某侧开放：但有时风从这开放侧吹入；在没有礁石的地方，船只相互碰撞而碎裂。若连港湾都不安全，哪里还有安全？然而必须承认，确实，在港湾中的人在某种程度上比在汪洋大海中好得多。让他们彼此相爱，如同在港湾中的船只，让他们快乐地联结在一起；不要互相碰撞；让那里保持完全的平等，爱中的恒心；当偶尔风从开放侧吹入时，让那里有谨慎的驾驶。现在，一个或许管理这些地方的人——即，在所谓\Emph{隐修院}中服侍自己弟兄的人——会告诉我什么？我会谨慎：我不接纳任何恶人。你怎能不接纳恶人呢？……那些将要进入的人，连自己都不认识他们，何况是你呢？因为许多人曾向自己许诺，他们将要度那神圣的生活——一切共有，没有人称任何东西为自己，他们一心一意在主内——\BibleRef{宗 4:32}他们被投入火炉，便破裂了。那么，你怎能认识连自己都不认识的人呢？……那么，哪里是安全？这里无处可寻；在此生中，除了对天主应许的希望，无处可寻。但是，当我们到达那里时，才有完全的安全，当城门关闭，耶路撒冷城门的门闩锁紧时；那才是真正的欢欣，极大的喜乐。只是你不要感觉安全而赞美任何一种生活方式：\Quote{不要在人死之前称他有福。}\par

9. 藉此手段，人们被欺骗，以致他们要么不开始更好的生活，要么轻率地尝试；因为，当他们选择赞扬时，他们称赞时只字不提与善人混杂的恶行；而那些选择指责的人，则怀着如此嫉妒和乖戾的心态，以致对善行视而不见，只夸大实际存在或想象出来的恶行。于是便发生这样的情况：当某种职业被不恰当地（即不谨慎地）赞扬时，如果它以其名声吸引了人们，那些投奔那里的人会发现一些他们原不相信那里会有的东西；因恶人而受冒犯，便从善人那里退缩。弟兄们，请将这教训应用于你们的生活，并如此聆听，好使你们能生活。天主的教会，总的来说，被尊崇：基督徒，唯独基督徒，被称为伟大的，天主教会受尊崇；大家彼此相爱；各人尽其所能为他人服务；他们献身于祈祷、禁食、赞美诗；在整个世界，在和平的合一中，天主受到赞美。也许有人听到这些，却不知对与他们混杂的恶人只字未提；他前来，被这些赞美所吸引，却发现有恶人与之混杂，而在他来之前无人向他提起；他被假基督徒所冒犯，便逃离真基督徒。另一方面，那些仇恨并毁谤他们的人，轻率地指责他们：问道，基督徒是什么样的人？谁是基督徒？贪婪的人，放高利贷的人。那些在节日里挤满教堂的人，不也正是那些在游戏和其他表演中挤满剧院和竞技场的人吗？他们醉酒、贪食、嫉妒、彼此毁谤。确有这样的人，但并非只有这些人。而这毁谤者在其盲目中只字不提善人；那赞美者在其不谨慎中对恶人保持沉默。……同样，在隐修院存在的弟兄们的共同生活中，有伟大而圣洁的人在那里生活，每日有赞美诗、祈祷、赞美天主；他们的职业是阅读；他们亲手劳动，以此维持生计；他们不贪婪追求什么；凡虔诚弟兄们为他们带来的，他们都以知足和爱德使用；没有人声称自己没有的属于别人；大家彼此相爱，彼此容忍。你赞扬了他们；你称赞了他们；那不知内情的人，不知当风吹入时，船即使在港口也会相撞，便怀着安全的希望进去，期望找不到一个需要容忍的人；他发现有邪恶的弟兄，如果他们未被接纳，本不会存在（他们最初必须被容忍，以防他们或许会改过；除非他们首先被容忍，否则很难被排除）；而他本人变得无法忍受。谁在此要求我？我以为这里有爱德。由于少数人的乖戾而恼怒，因为他没有坚持实现自己的誓愿，他成为了如此神圣计划的逃兵，并对从未兑现的誓愿负有罪责。然后，当他本人也出去时，他也成了一个指责者、毁谤者；只记录那些（有时是真实的）他声称无法忍受的事情。但恶人真实的困扰应当为了善人的团体而忍受。圣经对他说：\BibleQuote{祸哉，那些失去忍耐的人！} \BibleRef{训 2:16} 更有甚者，他吐出愤怒的恶臭，以此阻止那些将要进入的人；因为当他进入时，他自己无法坚持。他们是什么样的人？嫉妒的、好争吵的、不容忍任何人的人、贪婪的人；说：他在那里做了这事，他在那里做了那事。邪恶者，你为什么对善人保持沉默！你对你不能容忍的人说得够多了；你对容忍你邪恶的人却只字不提。……\par

10. \BibleQuote{当喜乐地事奉主} \BibleRef{咏 99:2} ：祂对你说，无论你是谁，你在爱中忍受一切，并在希望中喜乐。\BibleQuote{事奉主}，不是在怨言的苦涩中，而是在\Quote{爱的喜乐中}。\BibleQuote{以欢欣来到祂面前。} \BibleQuote{以歌唱来到祂面前。}\par

11. \BibleQuote{要知道主就是天主} \BibleRef{咏 99:3} 。谁不知道主就是天主呢？但祂说的是那位人们认为不是天主的\Quote{主}：\BibleQuote{要知道主就是天主。} 不要让那位主在你们眼中变得卑贱：你们曾钉死祂、鞭打祂、唾弃祂、给祂戴上荆棘冠冕、给祂穿上耻辱的衣服、将祂悬在十字架上、用钉子刺穿祂、用长矛刺伤祂、在祂的坟墓前设置守卫；祂是天主。\BibleQuote{是祂创造了我们，而非我们自己。} 是祂创造了我们：\BibleQuote{凡受造的，没有一样不是藉着祂造成的。} \BibleRef{若 1:3} 你们有什么理由欢欣？有什么理由骄傲？另一位创造了你们；那造你们的，正忍受你们的所作所为。但你们自高自大，自夸，好像你们是自造的。你们最好是让那造你们的完美地塑造你们。……\BibleQuote{我们是祂的子民，是祂牧场上的羊。} 羊群和一只羊。这些羊是一只羊：我们有牧人何其慈爱！祂撇下九十九只，下来寻找那一只，祂将羊扛在自己的肩上，用自己的血赎回，\BibleRef{路 15:4} 那牧人毫无畏惧地为羊死，且在复活后重得祂的羊。\par

12. \Quote{以宣认进入祂的门} \BibleRef{咏 99:3}。门是开端：以宣认开始。因此这篇圣咏题为\Quote{宣认的圣咏：}要欢喜。要宣认你们不是由自己所造，要赞美那造了你们的那一位。愿你们的善来自祂，因为离弃祂就造成了你们的恶。\Quote{以宣认进入祂的门。}愿羊群进入门中：不要留在外面，成为豺狼的猎物。如何进入呢？\Quote{以宣认。}让门，即开端，为你们成为宣认。因此在另一篇圣咏中说：\Quote{以宣认向主开始。}那里所说的\Quote{开始，}这里称之为\Quote{门。}\Quote{以宣认进入祂的门。}什么？我们进入之后，就不再宣认了吗？要常常宣认祂：你们总有可宣认的理由。今生人很难改变到在祂身上找不出任何可责备之处的程度：你们必须责备自己，免得那将要定罪的一位责备你们。因此，即使你们进入了祂的庭院，那时也要宣认。那时不再有罪的宣认了吗？在那安息中，在那相似天使的境界中。但请思想我所说的话：那里将没有罪的宣认。我并未说，将没有宣认：因为将有赞美的宣认。你们将永远宣认：祂是天主，你是受造物；祂是你的保护者，你受保护。在祂内你仿佛被隐藏。\Quote{以赞美诗进入祂的庭院，并向祂宣认。}在门中宣认；当你们进入庭院时，以赞美诗宣认。赞美诗是赞颂。进入时责备自己；进入后赞美祂。\Quote{请给我敞开正义之门，}他在另一篇圣咏中说，\Quote{使我进去，向主宣认。}他说，进入后我就不再宣认了吗？即使进入后，他仍要宣认。因为我们的主耶稣基督在说\Quote{父啊，我宣认祢}时，宣认了什么？\BibleRef{玛 11:25}祂以赞美祂的方式宣认，而非控告自己。\Quote{称颂祂的圣名。}\par

13. \Quote{因为主是甘饴的} \BibleRef{咏 99:4}。不要以为你们在赞美祂时会衰弱。你们对祂的赞美如同食物：你们越赞美祂，就越获得力量，而你们所赞美的祂也变得越甘饴。\Quote{祂的慈悲是永远的。}因为在祂释放你们之后，祂不会停止慈悲：保护你们直到永生属于祂的慈悲。因此，\Quote{祂的慈悲直到永远：祂的真理直到万代} \BibleRef{咏 99:5}。你们要理解\Quote{直到万代，}或指每一代，或指两代，一是地上的，一是天上的。这里有一代产生可朽者；另一代造就有永生者。祂的真理既在这里，也在那里。不要以为祂的真理不在这里；如果祂的真理不在这里，祂就不会在另一篇圣咏中说：\Quote{真理从地上生起；}真理本身也不会说：\Quote{看，我同你们天天在一起，直到今世的终结。}\BibleRef{玛 28:20}\par

\SourceRecord{Translated by J.E. Tweed.；Nicene and Post-Nicene Fathers, First Series；Vol. 8.；Edited by Philip Schaff.；Buffalo, NY: Christian Literature Publishing Co.,；1888.；<http://www.newadvent.org/fathers/1801100.htm>.}

% structure: p0001 source=h1 target=section reason=page-title

\section{圣咏第一〇一篇释义}

1. 在此圣咏中，我们应在整篇中寻求我们在第一节所找到的：\Quote{上主，我要歌颂仁慈与审判}\BibleRef{咏 100:1}。愿无人因天主的仁慈而自夸永不受罚，因为也有审判；愿那改过自新的人不畏惧主的审判，因为仁慈在它之前。因为当人审判时，有时被仁慈胜过，便违背正义；他们似乎有仁慈，却没有正义。而有时，当他们想要执行严厉的审判时，便失去了仁慈。但天主既不在仁慈的慷慨中失去审判的严厉，也不在严厉审判中失去仁慈的慷慨。我们暂且按时间区分这两种，即仁慈与审判；因为可能它们并非无故按此顺序放置，以致他说的是\Quote{审判与仁慈}，而是\Quote{仁慈与审判}。因此，若我们按时间的先后区分它们，也许我们会发现现在是仁慈的季节，将来是审判的季节。何以仁慈的季节在先？首先，思量天主如何行事，好使你在祂所允许的限度内效法天父……\Quote{祂使太阳上升，光照恶人，也光照善人；降雨给义人，也给不义的人}。看，这就是仁慈。当你看见义人与不义的人享有同样的太阳、同样的光、饮同一水源、受同一雨水滋润、享同样地上果实、呼吸同样空气、平等拥有世物时，不要以为天主不公义，竟将这些同等赐予义人与不义之人。这是仁慈的季节，还不是审判的季节。因为若天主起初不藉仁慈宽恕，祂就找不到那些祂能藉审判加冕的人。因此，有一个仁慈的季节，那时天主的容忍召唤罪人悔改。\par

2. 请听宗徒区分每个季节，你也应区分……\Quote{你以为吗？}他说，\Quote{你这人哪，你判断行这样事的人，自己却照样行，你以为你能逃脱天主的审判吗？} 仿佛我们要回答说：我为何每日如此犯罪，却未遭遇灾祸？他继而向他指明仁慈的季节：\Quote{还是你藐视祂丰富的仁慈、宽容与忍耐？} 他确实藐视了；但宗徒使他忧虑。\Quote{不知道，}他说，\Quote{天主的仁慈是引领你悔改？}\BibleRef{罗 2:4} 看，这就是仁慈的季节。但为使他不以为这会永远持续，他在下一节中如何使他恐惧？现在请听审判的季节；你已听到仁慈的季节，因此\Quote{上主，我要歌颂仁慈与审判}：\Quote{可是你，}宗徒说，\Quote{凭着你硬心和不悔改的心，为自己积蓄忿怒，直到忿怒的日子，即公义审判的天主显示的日子，祂要照各人的行为报应各人}\BibleRef{罗 2:5}。看，\Quote{仁慈与审判}。但他已警告关于审判的事：那么天主的审判只应被畏惧，而不应被爱吗？恶人因惩罚而畏惧，善人因冠冕而爱慕。因为宗徒在我所引的证言中警告了恶人，请听他如何在何处给予善人关于审判的希望。他举出自己，并也表明自己是在仁慈的季节中。因为除非他找到仁慈的时期，否则审判会以何种状态找到他？一个亵渎者、迫害者、凌辱他人者。他如此说，并赞美我们现在所处的仁慈季节：\Quote{我从前是亵渎者、迫害者和凌辱者，}他说，\Quote{但我蒙了怜悯}。但或许只有他蒙了怜悯？请听他如何鼓励我们：\Quote{但基督耶稣在我这罪魁身上}他说，\Quote{首先显示祂全部的忍耐，给后来信祂得永生的人作榜样}。\Quote{显示祂全部的忍耐}是什么意思？就是使每一个罪人和恶人看见保禄得了赦免，便不绝望？看，他以自己为例，从而也鼓励了他人……但只有保禄值得这样吗？因为我曾断言，正如他以先前的证言引起我们的恐惧，他也以后面的证言鼓励我们。当他说：\Quote{主，公义的审判者，到那一日要赏给我}他加上，\Quote{不但赏给我，也赏给凡爱慕祂显现的人}\BibleRef{弟后 4:8}和祂国度的人。因此，弟兄们，我们既有仁慈的季节，就不当因此自夸或纵容自己，说：天主永远宽恕……\par

3. \BibleQuote{我要弹琴歌唱，并在无瑕之道上领悟。当祢来到我面前} \BibleRef{咏 100:2}。除非在无瑕之道上，你既不能弹琴歌唱，也不能领悟。你若愿意领悟，就在无瑕之道上歌唱，即在你的天主面前欢欣工作。什么是无瑕之道？请听以下的话：\Quote{我在我屋中，以纯洁行走。}这无瑕之道始于纯洁，也终于纯洁。何必多言？你要纯洁，就成全了正义……但谁又是纯洁的呢？那既不伤害他人，也不伤害自己的人。因为伤害自己的人，并不纯洁。有人说：\Quote{看，我没有偷窃任何人，没有压迫任何人；我要靠自己的财产，即我善良劳动所得，快乐生活；我要有精美的宴席，随意花费，与我喜欢的人尽情饮酒；我偷了谁？压迫了谁？有谁抱怨过我？}他看似纯洁。但若他败坏自己，若他推翻自身内天主的殿宇，怎能盼望他会仁慈对待他人、怜悯困苦者？那对自己残忍的人，能对他人仁慈吗？因此，全部的正义都归结为\OriginalTerm{纯洁}{innocence}一词。但喜爱不义的人，痛恨自己的灵魂。当他喜爱不义时，他以为是在伤害别人。但请想想他是否伤害了别人：\Quote{喜爱不义者，}他说，\Quote{痛恨自己的灵魂。}因此，那想要伤害别人的人，先伤害了自己；他也不能行走，因为没有空间。因为一切邪恶都受狭窄之苦，唯独纯洁有足够行走的宽度。\Quote{我在我心中的纯洁中，在我屋中行走。}他所说的\Quote{屋中}，要么指教会本身，因为\Person{基督}在教会中行走；要么指他自己的心，因为我们内在的房屋就是我们的心，正如他在前面的话中所解释的：\Quote{在我心中的纯洁}。什么是心中的纯洁？即他房屋的中间？凡在这屋中有坏房屋的人，被赶出门外。因为凡在心中被恶劣良心压迫的人，正如因水满或烟雾溢出而离开自己房屋的人一样，不容忍自己住在其内；所以那没有宁静之心的人，不能快乐地住在自己心里。这样的人在心思的倾向上走出自身，以外在事物取乐，这些事物影响肉体；他们在琐事、表演、奢侈、一切邪恶中寻求安息。他们为何希望外在安好？因为内在不安好，以致他们能在良好的良心中喜乐……\par

4. \BibleQuote{我眼中不设置邪恶之事} \BibleRef{咏 100:3}。…我没有喜爱任何邪恶之事。他解释这同一件邪恶之事：\BibleQuote{我恨恶那些行背信的人。}请听，我的弟兄们。若你们与\Person{基督}在他屋中行走，即若你们在心中享有良好安息，或在教会本身中在虔敬之道上踏上良好旅程；你们不应只恨那些外在的不忠信者，也要恨你们在内部可能发现的一切不忠信者。谁是不忠信者？他们憎恨天主的律法；他们听了，却不遵行，被称为不忠信者。要恨恶行背信的人，把他们从你那里赶走。但你们应恨不忠信者，而非人：你看，一个不忠信的人有两个名字：人和不忠信者。\Person{天主}造成他人，他使自己成不忠信者；在他身上爱\Person{天主}所造成的，迫害他自己所造成的。因为当你迫害了他的不忠信时，你就杀害了人的工作，释放了\Person{天主}的工作。\BibleQuote{我恨恶那些行背信的人。}\par

5. \Quote{邪恶之心没有紧贴于我。}…一个人的心，若不愿意任何与\Person{天主}所愿相悖之事，就被称为正直。因此，若正义之心跟随\Person{天主}，扭曲之心就反抗\Person{天主}。假设有不如意之事临到他，他呼喊道：\Quote{\Person{天主}，我对祢做了什么？我犯了什么罪？}他希望自己显得公正，\Person{天主}不公正。有什么比这更扭曲的呢？你自身扭曲还不够：你还必须认为你的准则也是扭曲的。改正你自己，你便发现祂是正直的，离开祂你使自己变得扭曲。祂行事公正，你行事不公；因此你是乖谬的，因为你称人为公，而称\Person{天主}为不公。你称哪个人为公？你自己。因为当你说\Quote{我对祢做了什么？}，你自以为公正。但让\Person{天主}回答你：\Quote{你说得对：你对我什么都没做：你一切都为自己做了；因为若你为我做了什么，你早已做了好事。因为凡是做得好的，都是为我做的；因为它是按我的诫命做的；但凡是邪恶做的，都是为你自己做的，不是为我做的；因为邪恶的人除非为自己，什么也不做，因为这不符合我的命令。}当你们看到这样的人，弟兄们，要责备他们，说服并纠正他们：如果你们不能责备或纠正他们，不要赞同他们。\par

6. \BibleQuote{恶人离开我时，我不认识他} \BibleRef{咏100:4}. 我不赞同他，不赞美他，他不使我喜悦。因为我们在圣经中有时发现\Quote{认识}一词用于\Quote{喜悦}的含义。弟兄们，有什么事向天主隐藏呢？祂认识义人，难道不认识不义的人吗？你想什么事，是祂不知道的呢？我不说你想什么；但是，你曾想过什么，是祂没有预见过的呢？所以，天主知道一切；然而，在最终的审判中，即在仁慈之后的审判中，祂对一些人说：\BibleQuote{我必向他们声明：我从来不认识你们；你们这些作恶的人，离开我吧！} \BibleRef{玛7:23} 难道有祂不认识的人吗？但\Quote{我从来不认识你们}是什么意思？我不在我的标准中认可你们。因为我认识我正义的标准：你们与它不相合，你们偏离了它，你们是弯曲的。因此，祂在这里也说：\Quote{恶人离开我时，我不认识他。}...因此，\Quote{恶人离开我时}意思是：当恶人与我不相似，不愿效法我的道路，不愿在他作恶时按照我以身作则所提出的样式生活；\Quote{我不认识他。}\Quote{我不认识他}是什么意思？不是我不知道他，而是我不认可他。\par

7. \BibleQuote{谁暗地里毁谤他的邻舍，我必要迫害他} \BibleRef{咏100:5}. 请看那义人的迫害，不是对人的迫害，而是对罪的迫害。\Quote{眼目高傲，心贪得无厌的人，我不与他同席。} \Quote{我不与他同席}是什么意思？我不与他一同吃饭；因为\Quote{同席}就是吃饭；那么，我们如何发现主自己与骄傲的人一同吃饭呢？祂不仅与税吏和罪人吃饭——因为他们是谦卑的：他们承认自己的软弱，并寻求医生——我们也发现祂与骄傲的法利塞人一同吃饭。某个骄傲的人邀请了祂：正是那个人，因为一个有罪的女人，一个城中名声不好的女人，靠近主的脚而不悦。那个法利塞人是骄傲的：主与他一同吃饭；那么，祂所说的\Quote{我不与这样的人一同吃饭}是什么意思呢？祂如何命令我们做祂自己未曾做过的事呢？祂劝勉我们效法祂：我们看见祂与骄傲的人一同吃饭；祂如何禁止我们与骄傲的人一同吃饭呢？实际上，弟兄们，我们为了规劝的缘故，避免与我们的弟兄共融，不与他们一同吃饭，好使他们改过自新？我们宁愿与陌生人、与外邦人一同吃饭，也不与那些与我们有共同信仰的人一同吃饭，如果我们看到他们生活邪恶，好让他们羞愧而悔改；正如宗徒所说：\BibleQuote{若有人不听从我们这封信上的话，要记下这人，不要与他交往，叫他自觉羞愧；可是不要以他为仇人，而要劝他如劝弟兄。} \BibleRef{得后3:14} 为了医治别人，我们通常这样做；然而，我们也常常与许多陌生人和不敬神的人一同吃饭。\par

8. 虔诚的心有它的宴席，骄傲的心也有它的宴席：正是因为骄傲之心的食物，他才说\Quote{心贪得无厌}。骄傲的心如何得喂养？如果一个人骄傲，他就嫉妒；否则不可能。骄傲是嫉妒之母：它不能不产生嫉妒，并且常常与之共存。因此，每一个骄傲的人都是嫉妒的；如果嫉妒，他就以他人的不幸为食。因此宗徒说：\BibleQuote{但如果你们彼此咬噬、吞食，你们要小心，免得被彼此消灭。} \BibleRef{迦5:15} 那么，你看见他们在吃：不要与这些人一同吃；要躲避这样的宴席：因为他们以别人的恶事为乐，无法满足，因为他们的心贪得无厌。你要小心，不要在魔鬼的圈套中被捕，就像鸟在陷阱中吃食，鱼在鱼钩上吃食一样，它们在吃的时候被捉住。因此，不敬神的人有自己的宴席，虔诚的人也有自己的宴席。听听虔诚之人的宴席：\BibleQuote{饥渴慕义的人是有福的，因为他们要得饱饫。} \BibleRef{玛5:6} 因此，如果虔诚的人以正义为食，不敬神的人以骄傲为食，那么他心中贪得无厌就不足为奇了。他吃的是不义之食；不要吃这不义之食，这样，眼目高傲、贪得无厌的人就不会与你同席了。\par

9. 你从何处得喂养？当他不与你同席时，什么使你喜悦？他说：\Quote{我的眼睛}他说，\Quote{常注视这地中信实的人，使他们与我同坐。} \BibleRef{咏100:6} 也就是说，使他们与我同坐。他们\Quote{同坐}是什么意思？\BibleQuote{你们要坐在十二宝座上，审判以色列十二支派。} \BibleRef{玛19:28} 地上的信实之人审判，因为对他们说过：\BibleQuote{你们不知道我们将审判天使吗？} \BibleRef{格前6:3} \Quote{凡行在无可指摘的道路上的人，他服事了我。} 是\Quote{服事了我}，他说，而不是服事了自己。因为许多人传福音，却是服事自己；因为他们寻求自己的事，而不是耶稣基督的事。\BibleRef{斐2:21} ...\par

10. \BibleQuote{骄傲的人没有居住在我的家中} \BibleRef{咏100:7}. 你要在心灵的意义上理解这一点。骄傲的人没有居住在我的心里：没有这样的人居住在我的心里；因为他匆忙地离开了我。唯有温良和平安的人居住在我的心里；骄傲的人不住在那里，因为不义的人不住在义人的心里。即使义人与你相距不知多少英里和驿站：你们若一心一意，就是同居共处。\Quote{行事骄傲的人没有居住在我的家中；说话不义的人没有在我眼前引导。} 这就是无可指摘的道路，我们在其中明白主何时来到我们这里。\par

11. \Quote{每早晨，我必消灭地上的一切恶徒，铲除主的城内的一切作恶的人} \BibleRef{咏 100:8}。这有些晦涩。城里有作恶的人，现在似乎被容忍了。为什么呢？因为现在是仁慈的季节；但审判的季节必将到来；因为圣咏是这样开始的：\Quote{我要歌颂仁慈与审判，主，我向你歌颂。}……\par

12. 祂现在容忍，那时将审判。但祂何时审判？当黑夜过去之后。因此祂说：\Quote{在黎明。} 当黎明最终到来，黑夜已经过去。为什么祂容忍他们直到黎明？因为是黑夜。什么是黑夜？因为那是仁慈的季节：祂仁慈，而人心是隐藏的。你看见某人生活邪恶；你容忍他：因为你不知道他将来会变成怎样；因为是黑夜；今天生活邪恶的人，明天可能生活良善；今天生活良善的人，明天可能变坏。因为是黑夜，天主容忍所有人，因祂长久忍耐；祂容忍他们，好让罪人归向祂。但那些没有在那仁慈的季节悔改的人，将被击杀。为什么？为使他们从主的城中，从耶路撒冷的共融、圣人的共融、教会的共融中被驱逐出去。但他们何时被击杀？\Quote{在黎明。} \Quote{在黎明} 是什么意思？当黑夜过去。为什么现在祂容忍？因为是仁慈的季节。为什么祂不永远容忍？因为：\Quote{我要歌颂仁慈与审判，主，我向你歌颂。} 弟兄们，不要自欺：所有作恶的人都将被击杀；基督要在黎明时击杀他们，并将他们从祂的城中消灭。但现在，在仁慈的时期，让他们听从祂。祂藉着法律、先知、圣咏、书信、福音到处呼喊：看，祂不沉默；祂容忍，祂施仁慈；但要当心，因为审判必将来临。\par

\SourceRecord{Translated by J.E. Tweed.；Nicene and Post-Nicene Fathers, First Series；Vol. 8.；Edited by Philip Schaff.；Buffalo, NY: Christian Literature Publishing Co.,；1888.；<http://www.newadvent.org/fathers/1801101.htm>.}

% structure: p0001 source=h1 target=section reason=page-title

\section{\WorkTitle{咏第一〇二篇释义}}

1. 看，一个穷人祈祷，且并非在沉默中祈祷。因此，我们可以听他，并看他究竟是谁：是否正巧是祂，宗徒所说的那位：\BibleQuote{祂本是富有的，却为你们成了贫穷的，好使你们因祂的贫穷而成为富有的。}\BibleRef{格后 8:9} 如果这是祂，那么，祂如何是贫穷的？因为祂在何种意义上是富有的，谁看不见呢？那么，有什么比祂更富有，由祂创造了财富，甚至那些并非真正财富的？因为通过祂，我们甚至拥有这些财富：能力、记忆、性格、身体的健康、感官，以及肢体的构造；因为当这些都安全时，连穷人也算是富有的。通过祂，也有那些更大的财富：信德、虔诚、正义、爱德、贞洁、善行；因为没有人拥有这些，除非通过那使不虔敬者成义的一位……看，何等富有！在如此富有的一位身上，我们如何认出这些话？\Quote{我吃了灰尘像面包一样，又以眼泪调和了我的饮料。}如此巨大的财富竟到了这般地步？前者是一种极高的状态，这却是一种极卑微的状态……然而，仍要查考这穷人是否就是祂；因为\BibleQuote{圣言成了血肉，居住在我们中间。}\BibleRef{若 1:14} 也反思这些话：\Quote{我是祢的仆人，是祢婢女的儿子。}看，这位婢女，贞洁、童贞、又是母亲；因为在那里祂接受了我们的贫穷，当祂取了奴仆的形体，空虚了自己；免得你畏惧祂的财富，在你乞丐般的状态中不敢接近祂。在那里，我说，祂取了奴仆的形体，在那里祂披上了我们的贫穷；在那里祂使自己贫穷，却使我们富有。我们现在正接近从祂理解这些事：然而我们还不可轻率地断言……\par

2. 那么，让贫穷加于贫穷：让祂将我们卑微的身体改变形状，归于祂自己；让祂作我们的元首，我们作祂的肢体，让两人成为一体。\BibleRef{斐 3:21} ……因为祂甚至屈尊把我们视为祂的肢体。忏悔的人也属于祂的肢体。因为他们没有被排除，也没有与祂的教会分离；祂也不会使教会成为祂的新娘，除非借着这样的话：\BibleQuote{你们悔改吧，因为天国临近了。}\BibleRef{玛 3:2} 那么，让我们听元首和身体的祈祷，新郎和新娘，基督和教会，两者同属一个位格；但圣言和肉体并非同一回事；圣父和圣言是同一回事；基督和教会是同一位格，一个完美的人，具有祂完满的形态……因此，让我们听基督，在我们内、与我们一同、并为我们而成为贫穷的基督。因为标题本身就指明了这穷人。最后，记住我曾猜测那穷人是谁：让我们听祂的祈祷，认出祂的位格；当你听到任何不能适用于祂元首的话时，不要弄错；正是为此，我预先作了这样的说明，使你听到这类描述时，能理解为从身体的软弱发出的声音，并在元首中认出肢体的声音。标题是：\Quote{困苦人的祈祷，当他受折磨时，在主前倾吐他的祈祷。}这同一位穷人，在别处说：\Quote{我从地极呼求祢，当我心沉重时。}祂是困苦的，因为祂也是基督；祂在先知的话中既自称为新郎，又自称为新娘：\BibleQuote{祂给我戴上冠冕，如同新郎，又装饰我，如同新娘。}\BibleRef{依 61:10} 祂自称为新郎，又自称为新娘；这是为什么？除非新郎指元首，新娘指身体。因此，他们只有一个声音，因为他们是一体。让我们听，并在这些话中认出我们自己；如果我们看见自己在外面，就让我们努力进入里面。\par

3. \BibleQuote{上主，求祢垂听我的祈祷：愿我的呼声达到祢面前}\BibleRef{咏 101:1}。\BibleQuote{上主，求祢垂听我的祈祷}等同于\Quote{愿我的呼声达到祢面前}：重复表明了恳求者的心情。\BibleQuote{求祢不要向我掩面。}天主何时向祂的儿子掩面？圣父何时向基督掩面？但为了我肢体的贫穷，\BibleQuote{求祢不要向我掩面：在我困苦的任何一日，求祢侧耳听我}\BibleRef{咏 101:2}……你今日在困苦中，我在困苦中；别人明日困苦，我困苦；在这世代以后，其他后代，接续你后代的人，在困苦中，我困苦；直到世界穷尽，凡在我身体内困苦的，我都困苦。……伯多禄祈祷了，保禄祈祷了，其余宗徒祈祷了；那些时代的信徒祈祷了，以后时代的信徒祈祷了，殉道时代的信徒祈祷了，我们时代的信徒祈祷了，我们后代时代的信徒也将祈祷。\Quote{快快应允我}：因为我现在恳求祢乐意赐予的。我不求属世的事物，像属世的人那样；而是终于从先前被掳中得赎，我渴慕天国；\Quote{快快应允我}：因为祢曾对这渴望说：\BibleQuote{你说话时，我就要说：我在这里。}\BibleRef{依 58:9} 你为什么呼求？在什么困苦中？在什么缺乏中？哦，穷人，在全富足的天主门前，你因什么渴望而乞求？你从什么匮乏中寻求解救？你因什么缺乏而敲门，好使门为你打开？\par

4. \Quote{我的岁月如烟消逝} \BibleRef{咏 101:3}。啊，岁月！若是岁月：因为哪里听到\Quote{日子}，那里就理解为光明。\Quote{我的岁月}，我的时光；那么，\Quote{如烟}，岂不是出于骄傲的膨胀？……看烟，如同骄傲，上升、膨胀、消散：因此理当衰败，而不坚固。\Quote{我的骨头如同在炉中被烤焦。}我的骨头，我的力量，基督身体的骨头，基督身体的力量，难道在任何地方比在圣宗徒那里更大吗？然而看，骨头被烤焦了。\Quote{谁软弱，我不软弱呢？} \BibleRef{格后 11:29}。他们勇敢、忠信，是圣言的有能力的解释者和宣讲者，生活如他们所言，言说如他们所听；他们显然勇敢，然而所有遭受绊脚石的人，对他们而言都是炉子。因为那里有爱，尤其在骨头中。骨头在一切肉体内，支撑一切肉体。但如果有人遭受任何绊脚石，并使自己的灵魂处于危险；骨头按其爱的程度被烤焦……\par

5. 回顾亚当，人类由此而生。因为除了从他，苦难如何传播？除了从他，这遗传的贫困从何而来？那么，那曾经在自己的身体中绝望的人，如今既已置身于基督的身体内，要怀着希望说：\Quote{我的心被击伤，如同草枯干} \BibleRef{咏 101:4}。理当如此，因为凡有血气的尽都如草。\BibleRef{依 40:6}。但这事如何临到你？\Quote{因为我忘了吃我的饼。}因为天主曾将祂的诫命赐为饼。因为灵魂的饼是什么？蛇诱惑，女人犯命，他触及了禁果（\BibleRef{创 3:6}），他忘了诫命：他的心被击伤，理所当然，如同草枯干，因为他忘了吃他的饼。忘了吃饼，他喝毒药：他的心被击伤，如同草枯干。……现在吃你曾忘记的饼。但这饼本身已经来临，在他的身体中你或许能记起你遗忘的声音，并在你的贫困中呼喊，好使你获得富足。现在吃：因为你是在祂的身体内，祂说：\Quote{我是从天上降下的活饼} \BibleRef{若 6:41}。你曾忘了吃你的饼；但祂被钉十字架后，\Quote{地极都要记念，并归向主}。遗忘之后，愿记念来临，愿从天上的饼被吃，使我们得生；不是吗哪，像他们吃了，还是死了（\BibleRef{若 6:49}）；那饼，关于它经上说：\Quote{饥渴慕义的人是有福的} \BibleRef{玛 5:6}。\par

6. \Quote{因我叹息的声音，我的骨头紧贴我的肉} \BibleRef{咏 101:5}。因为许多人叹息，我也叹息；甚至为这事我叹息，因为他们为错误的原因叹息。那人丢失了一块钱，他叹息；他丢失了信德，他不叹息；我衡量钱和信德，我发现更有理由为那人不当叹息或根本不叹息而叹息。他行欺诈，却欢喜。得了什么？失了什么？他得了钱，失了义。为此，那知道如何叹息的人就叹息；那靠近元首、正义地依附基督身体的人，为此叹息。但属血肉的人不为此叹息，而且他们使自己成为被叹息的对象，因为他们不为此叹息；我们也不能轻视他们，无论他们根本不叹息，还是为错误的原因叹息。因为我们想要纠正他们，想要改善他们，想要改造他们，当我们不能时，我们就叹息；当我们叹息时，我们并不与他们分离……\par

7. \Quote{我好像旷野中的鹈鹕，又像废墟中的猫头鹰} \BibleRef{咏 101:6}。看，三只鸟和三个地方：鹈鹕、猫头鹰和麻雀；三个地方分别是旷野、废墟和屋顶。鹈鹕在旷野，猫头鹰在废墟，麻雀在屋顶。首先必须解释鹈鹕的意义：因为它出生在一个使我们不熟悉的地区。它出生在偏僻之处，尤其是埃及尼罗河畔。无论它是哪种鸟，让我们思考圣咏想论它说什么。\Quote{它居住}，它说，\Quote{在旷野}。何必查问它的形状、肢体、声音、习性？就圣咏告诉你的，它是一种居住在孤独之处的鸟。猫头鹰是一种喜爱黑夜的鸟。\Emph{墙垣废墟}，正如我们称呼的，是没有屋顶、无人居住的墙壁，这些是猫头鹰的居所。至于屋顶和麻雀，你是熟悉的。因此，我发现在基督身体中有某个人，是圣言的宣讲者，同情软弱者，寻求基督的收益，记念将要来临的主（\BibleRef{玛 25:26}）。让我们从他的管家职分来看这三件东西。这样一个人来到那些不是基督徒的人中间吗？他是旷野中的鹈鹕。他来到那些曾是基督徒但已跌倒的人中间吗？他是废墟中的猫头鹰；因为他不抛弃甚至那些住在黑夜中的人的黑暗，他也想赢得这些人。他来到那些住在房屋中的基督徒中间，不是因为他们不信，也不是因为他们放弃了所信的，而是在所信的事上不冷不热地行走吗？麻雀向他们呼喊，不是在旷野，因为他们是基督徒；也不是在废墟，因为他们没有跌倒；而是因为他们在屋顶下；更确切地说在屋顶下，因为他们是在血肉之下。麻雀在血肉之上呼喊，不缄默天主的命令，也不成为属血肉的以致受屋顶辖制。\Quote{你们在暗中听到的，要在房顶上宣扬出来} \BibleRef{玛 10:27}。有三只鸟和三个地方；一个人可以代表三只鸟，三个人可以分别代表三只鸟；三种地方是三类人：然而旷野、废墟和屋顶不过是三类人。\par

8. ……我们不要轻易放过关于这只鸟，即鹈鹕，所说或所读到的内容；不要鲁莽地肯定什么，但也不可忽视那些已留给历代作者阅读和传述的内容。你们要这样听：如果它是真的，就能吻合；如果是假的，就不会成立。据说这种鸟用喙击打自己的雏鸟致死，并在巢中为被自己杀死的雏鸟哀悼三天；之后，据说母鸟深深创伤自己，将血倾注在雏鸟身上，雏鸟经血沐浴后复活。这可能为真，也可能为假；然而若为真，请看它多么吻合那位以自己的鲜血赐予我们生命的主。它与主吻合，在于母亲的肉身以鲜血唤回雏鸟的生命；这非常吻合。因为祂自比孵育雏鸟的母鸡。\BibleRef{玛 23:37}……如果此事确实如此，那么这种鸟就极其相似基督的肉身，我们正是因祂的血而被唤回生命。但这如何与基督吻合呢：鸟本身杀死自己的雏鸟？难道这不吻合吗？\Quote{我要杀死，我也要使人活；我击伤，我也要医治。}\BibleRef{申 32:39} 迫害者扫罗若不从天受到创伤，岂会死去？\BibleRef{宗 9:4} 除非藉着祂的鲜血所赐的生命，那位宣讲者岂能被兴起？但让那些撰写了相关主题的人去审视这些；我们不应让自己的理解立足于可疑的基础。不如让我们在旷野中认出这只鸟，正如圣咏所说：\Quote{旷野中的鹈鹕。}我想这里是指由童贞女诞生的基督。祂独自诞生，因为唯有祂是这样诞生的。诞生之后，我们来到祂的受难……诞生在旷野，因为唯独如此诞生；在犹太人的黑暗中受苦，如同在黑夜中，在他们的罪中，如同在废墟中：接下来呢？\Quote{我警醒着：}且\Quote{成了好像房顶上孤单的麻雀。}\BibleRef{咏 101:7} 你那时在废墟中安睡，曾说：\Quote{我躺下，就睡着了。}\Quote{我睡了}是什么意思？因为我愿意，我睡了：我为爱黑夜而睡；但随后是\Quote{我复活了。}因此这里说\Quote{我警醒着。}但祂警醒后，做了什么？祂升了天，藉着飞翔成了麻雀；即藉着上升；\Quote{独在房顶上，}即在天上。因此，祂藉诞生是鹈鹕，藉死亡是鸱鸺，藉再次上升是麻雀：在旷野中，祂是独自一；在毁坏之墙间，祂是被那些在建筑中站立不住的人所杀；而在这里，祂又为我们警醒飞翔，独在房顶，在那里为我们转求。\BibleRef{罗 8:34} 因为我们的头是麻雀，祂的身体是斑鸠。\Quote{麻雀为自己找到了房屋。}什么房屋？在天上，祂在那里为我们代祷。\Quote{斑鸠有巢，}天主的教会从祂十字架的木头找到了巢穴，在那里\Quote{她可以安置自己的幼雏，}即她的子女。\par

9. \Quote{我的仇敌终日辱骂我，称赞我的指着我起誓同盟。}\BibleRef{咏 101:8} 他们口里称赞，心里却为我设下网罗。听他们的称赞：\Quote{师傅，我们知道祢是真诚的，按真理教导天主的道路，也不顾忌任何人。请问，给凯撒纳税可以不可以？}\BibleRef{玛 22:16} 这恶名从何而来，岂不是因为我来到，为要使罪人成为我的肢体，使他们藉悔改能在我身体内？由此而来一切毁谤，由此而来迫害。\Quote{祢的师傅为什么同税吏和罪人一起吃饭？健康的人不需要医生，有病的人才需要。}\BibleRef{玛 9:11} 但愿你们意识到自己的病，好寻求医生；你们就不会杀害祂，并因愚蠢的骄傲在虚假的健康中灭亡。\par

10. \Quote{我吃灰尘如同吃饼，又把眼泪搀入我所喝的} \BibleRef{咏 101:9}。因为祂愿意在祂的肢体中拥有这类人，使他们得到医治与释放，因此便有恶名。如今，外教人对我们的诽谤是什么？弟兄们，你们认为他们对我们说什么？你们败坏纪律，歪曲人类道德。你们为何攻击我们？说吧，为何？我们做了什么？他们回答说：因为给人悔改的机会，应许一切罪过都得到宽恕；因此人们无所顾忌地作恶，因为当他们归正时一切都被赦免……\newline{}至于你，可怜的人，如果没有免罚的港湾，你会怎样？如果只有犯罪的自由，却没有罪过的赦免，你会在哪里？你将往何处？确实，连你也是如此，那位受苦者以灰烬为食，以眼泪为饮。这样的筵席难道现在不令你喜悦吗？然而，他回答说：人在赦免的盼望下更加犯罪。不，若他们绝望于赦免，就会更加犯罪。难道你没有看到角斗士生活在何等放纵的残忍中？这是什么原因？岂不是因为他们注定死于刀剑和祭献，便选择在流血之前满足情欲？难道你不会这样对自己说？我已经是个罪人，已经是不义之人，已经注定受罚，没有赦免的盼望：为何不做任何我喜欢的事，即使不合法？为何不尽我所能满足一切欲望，如果之后只有折磨等着我？难道你不会这样对自己说，并因这绝望而变得更糟？因此，应许赦免的那位纠正你说：\Quote{我指着我的永生起誓，主说，我决不喜悦恶人丧亡，但却喜悦恶人离开他的道路而生活。} \BibleRef{则 13:11}……\newline{}为了让人们不致因绝望而生活得更糟，祂应许了赦免的港湾；又为了不让他们因赦免的盼望而生活得更糟，祂使死亡的日子不确定：以最大的眷顾设定了两者，既作为归回者的避难所，又作为徘徊者的恐惧。以灰烬为食，以眼泪为饮；藉此筵席你将抵达天主的筵席。不要绝望，赦免已经应许给你。感谢天主，他说，因为已经应许；我紧握天主的应许。因此现在好好生活。明天，他回答，我会好好生活。天主应许了赦免；没人应许你明天……\par

11. \Quote{并且因祢的愤怒和怒气：因为祢举起我，又把我抛下} \BibleRef{咏 101:10}。这是祢的怒气，主，在亚当内：我们都在其中出生、由出生而紧附于我们的怒气；来自不义之源的怒气，来自罪过的团块的怒气；正如宗徒所说：\Quote{我们从前也都在他们中间，生来就是义怒之子。} \BibleRef{弗 2:3} 因为祂不说天主的愤怒将要临到他身上，而是说\Quote{义怒住在他身上}，因为他所出生的那怒气没有被除去……\newline{}人受尊荣，按天主的肖像受造；被高举至此尊荣，从尘土、从地上被举起，他接受了理性的灵魂；因那理性的活力，他超越一切走兽、牲畜、飞鸟和鱼类。\BibleRef{创 1:26} 因为它们中有谁有理性能理解？因为没有一个按天主的肖像受造。……\newline{}因此，\Quote{因为祢举起我，祢又抛下我}：惩罚跟随我，因为祢给了我自由的选择。因为若祢没有给我自由的选择，并因此没有使我优于牲畜，那么当我犯罪时，公正的谴责就不会跟随我。这样，祢藉着给我自由选择而举起了我，又按祢的审判抛下了我。\par

12. \Quote{我的岁月如同影儿消逝} \BibleRef{咏 101:11}……祂上面曾说：\Quote{我的岁月如烟云消散}；现在又说：\Quote{我的岁月如同影儿消逝}。在这阴影中，必须认出白日；在这阴影中，必须辨明光明；免得后来在迟延而无结果的悔改中说：\Quote{骄傲给了我们什么好处？财富与夸耀又给了我们什么？这一切都像影儿般逝去。} \BibleRef{智 5:8} 现在就说：一切都将像影儿般逝去，而你可能不像影儿般逝去。\Quote{我的岁月如同影儿消逝，我也如草枯干。} 因为上面曾说：\Quote{我的心被击伤，如草枯干。} 但被救主之血浸润的草将重新繁茂。\Quote{我如草枯干}；我，即人，在那种悖逆之后；这我因祢公义的审判而遭受：但祢是怎样的呢？\par

13. 因为并非因我跌倒，祢便衰老；因为祢强大，足以释放我，祢曾强大，足以使我谦卑。\Quote{但祢，主，却永远常存：祢的记念也存到万代} \BibleRef{咏 101:12}。\Quote{祢的记念}，因为祢不忘；\Quote{存到万代}，因我们认识生命的应许，无论今世还是来世。\BibleRef{弟前 4:8}\par

14. \Quote{求祢起来，怜悯熙雍：因为现在是怜悯她的时候} \BibleRef{咏 101:13}。什么时间？\Quote{但时期一满，天主就派遣了自己的儿子来，生于女人，生于法律之下。} 熙雍在哪里？\Quote{为赎出那些在法律之下的人。} \BibleRef{迦 4:4} 首先是犹太人：因为宗徒们来自他们，那五百多弟兄来自他们，\BibleRef{格前 15:6} 后来那众多的人来自他们，他们在天主前同心合意。\BibleRef{宗 4:32} 因此，\Quote{时候到了。} 什么时间？\Quote{看，现在正是悦纳的时候；看，现在正是救恩的日子。} \BibleRef{格后 6:2} 这是谁说的？那位天主的仆人，那位建筑师，他说：\Quote{你们是天主的建筑物。} \BibleRef{格前 3:9}\par

15. 因此，祂说什么？\Quote{因为你的仆人们喜爱她的石头} \BibleRef{咏 101:14}。谁的石头？熙雍的石头？但那里有一些不是石头。不是什么石头？接着说什么？\Quote{也怜惜其中的尘土。}我理解熙雍的石头是指所有先知：先有宣讲的声音被遣发，然后接受福音的职务，借着他们的宣讲，基督得以被人认识。因此，你的仆人们喜爱熙雍的石头。但那些不忠信、背弃天主的人，他们因恶行冒犯了造物主，便归回他们由此而来的尘土。他们成了尘土，成了不敬虔的人。但是，主啊，请等待；主啊，请容忍我们；主啊，请宽忍：不要让风冲进来，将这尘土从地面上吹走。愿你的仆人们前来，愿他们前来，愿他们在石头中认出你的声音，愿他们怜惜熙雍的尘土，愿他们按你的肖像被塑造：愿这尘土说，免得它灭亡，\Quote{求你记念我们不过是尘土。}这是关于熙雍的：那曾钉死主的，难道不是尘土吗？更糟的是，它是来自废墟的尘土；完全是尘土，但关于这尘土，并非徒然说：\Quote{父啊，赦免他们吧。}从这同样的尘土中，涌现了成千上万相信者的墙垣，他们将财产的价值放在宗徒们脚前。因此，从那里尘土中，兴起了一个被塑造而美丽的人性。外邦人中谁这样做了？与这些归化者的成千上万相比，我们称赞这样做的有几人？起初突然三千，随后五千；全都同心合意，全都变卖财产的价值，放在宗徒们脚前，以便分给每人，按各人所需要的，他们向天主有一颗心灵和一颗心。谁甚至从那同样的尘土中造成了这一切，岂非那亲自从尘土中创造了亚当的那位？因此，这是关于熙雍的，但不仅限于熙雍。\par

16. \Quote{上主，外邦人必敬畏你的圣名；地上列王必敬畏你的威严} \BibleRef{咏 101:15}。现在你已怜惜了熙雍，你的仆人们已喜爱了她的石头，因承认宗徒和先知的根基；现在他们已怜惜了她的尘土；如此，人从尘土中被塑造，或更确切地说，被重新塑造，进入生命；因此，宣讲在外邦人中增多了：愿外邦人敬畏你的圣名，愿另一堵墙也从外邦人而来，愿那角石被认出 \BibleRef{弗 2:20}，愿那来自不同地区、但在信仰上不再相异的二者，在紧密的结合中相遇。\par

17. \Quote{因为上主必建造熙雍} \BibleRef{咏 101:16}。这工程如今正在进行。哦，你们这些活石，奔向建造的工程，而非奔向废墟。熙雍正在建造，要提防倒塌的墙：塔楼正在建造，方舟正在建造；要记得洪水。这工程如今正在进行；但当熙雍建成时，会发生什么？\Quote{祂必在荣耀中显现。}为使祂能建造熙雍，为使祂成为熙雍的根基，祂曾被熙雍看见，但非在祂的荣耀中：\Quote{我们看见了祂，祂没有威仪，也没有美貌。} \BibleRef{依 53:2} 但当祂带着祂的天使来审判时 \BibleRef{玛 25:31}，那时他们岂非要仰望他们所刺透的那一位吗？\BibleRef{匝 12:10} 他们将在太迟时陷入羞愧，那些拒绝在及早和有益的忏悔中羞愧的人。\par

18. \Quote{祂垂顾了卑微穷人的祈祷，并未轻视他们的愿望} \BibleRef{咏 101:17}。这在熙雍的建造中如今正在发生：熙雍的建筑者祈祷，他们叹息：祂是一个穷人，因为穷人是许多；因为这么多民族中的成千上万的人在祂内是一，因为祂是教会和平的合一，祂是一，祂是许多：因着爱是一，因着扩展是许多。因此，我们现在祈祷，我们现在奔跑：现在，若有人向来不同，且生活乖戾，让他如吃饼般吃灰烬，以眼泪调和饮品。现在是时候，当熙雍正在建造时：现在石头正进入结构：当建筑完成，房屋奉献后，你为什么奔跑，去太迟地祈求，徒然地恳求，无用地敲门，注定与五个愚拙的童女一同留在外面？\BibleRef{玛 25:12} 因此，现在奔跑吧。\par

19. \Quote{愿这些事被记录下来，为后代的人} \BibleRef{咏 101:18}。当这些话被写下时，它们并未同样有益于那些在他们中间写下的人，因为它们是为预言新约而写下的，是在按旧约生活的人中间。但天主既赐下了那旧约，又在那应许之地安顿了自己的子民。但因\Quote{你的纪念代代相传}不属于不敬虔的人，而属于义人；\Quote{在我们的世代}属于旧约；而\Quote{在另一世代}属于新约；并且由于新约宣告了这被预言的事，\Quote{愿这些事被记录下来，为后代的人；那将要受造的百姓，要赞美上主。}不是受造的百姓，而是\Quote{那将要受造的百姓}。我的弟兄们，还有什么更清楚的呢？这里预言了那受造者，即宗徒所说的：\Quote{所以谁若在基督内，他就是一个新受造物；旧的已成过去，看，一切都成了新的。} \BibleRef{格后 5:17} \Quote{因为祂从祂崇高的圣所俯视。}祂从高处俯视，为来到卑微者中间：祂从高处成为卑微，为使卑微者升高……\par

20. \Quote{主从天上垂顾大地}\BibleRef{咏 101:19}，\Quote{为要垂听被囚者的哀叹，解救被定死罪的人}\BibleRef{咏 101:20}。我们在另一篇圣咏中读到：\Quote{愿被囚者的叹息上达于祢面前}；那里指的是殉道者的呼声。殉道者如何被囚？……但天主用这些锁链束缚他们，这些锁链固然艰难，暂时痛苦，但因祂的应许而可忍受；曾有话说：\Quote{因祢口中的言语，我谨守了艰苦的道路}。我们确实必须在这锁链中叹息，以获取天主的仁慈。不应躲避这些锁链，以免获得毁灭性的自由和今生短暂而短暂的快乐，随后是永久的痛苦。因此圣经\BibleRef{德 6:24}告诫我们不要拒绝智慧的锁链：\Quote{……那时，她的锁链要成为你的坚固保障，她的绳索要成为光荣的衣袍}。因此，被囚者应当呼喊，只要他们还在天主管教之链中，殉道者曾在此受试炼：锁链将被解开，他们将飞走，这些锁链日后将变为装饰。殉道者就是如此。因为迫害者杀害他们，除了借此松开锁链、转为荣冠之外，还成就了什么？……罪的赦免就是松绑。因为拉匝禄从坟墓出来，若不对他说\Quote{解开他，让他走}，\BibleRef{若 11:44}对他有何益处？主自己以声音唤醒他出墓，用呼喊使他复活，克服堆积在墓上的土，他便手脚缠着布带出来了；因而不是用自己的脚，而是靠那拉他出来的力量。这在忏悔者心中发生：当你听见一个人为自己的罪痛悔，他已复活；当你听见他藉告解显露良心，他已被拉出坟墓，但尚未被解开。何时被解开，由谁解开？\Quote{凡你们在地上所释放的，}祂说，\Quote{在天上也要被释放}。\BibleRef{玛 16:19}教会可合理地授予罪的赦免；但死人本身只能由内在呼喊的主唤醒，因为天主内在行事。我们向你们的耳朵说话：我们如何知道你们心中发生什么？但内在的事非我们所为，而是祂的作为。\par

21. \Quote{为使主的名在熙雍被宣扬}\BibleRef{咏 101:21}。因为起初，当被囚者被判死时，教会受压迫：经过这些苦难，主的名在熙雍大为自由地被宣扬，即在教会本身内。因为她是熙雍：不是那个起初骄傲、后被掳的地点；而是那熙雍——其预像是那个熙雍，意为瞭望塔；因为我们身处肉身时，注视眼前之事，不延伸至现在，而指向未来。因此它是瞭望塔：因为每个守望者都远眺。设防的地点称为瞭望塔：它们置于岩石上、高山上、树上，以便从更高处获得更广视野。因此，熙雍是瞭望塔，教会是瞭望塔。……如果教会是瞭望塔，主的名已在那里被宣扬。不独主的名在熙雍被宣扬，\Quote{祂的赞美，}祂说，\Quote{也在耶路撒冷}。\par

22. 如何被宣扬？\Quote{万民一同聚集，万国事奉主}\BibleRef{咏 101:22}。如何成就？除非藉被杀的鲜血？如何成就，除非藉被囚者的呻吟？因此那些在困苦和卑微中的人已被垂听，以致在我们时代，教会享有我们所见的极大光荣，连当时迫害她的列国如今也事奉主。\par

23. \Quote{她在祂力量的路上回答了祂}\BibleRef{咏 101:23}。……前面的话显示，要么\Quote{祂的赞美}，要么\Quote{耶路撒冷}回答了：因为曾说：\Quote{祂的赞美在耶路撒冷；万民一同聚集，万国事奉主。\Emph{Respondit ei.}}我们不能说\Quote{万国回答了}，因为应为\Emph{responderunt}。\Emph{Respondit ei}。不能说\Quote{万民回答了}，因为应为\Emph{responderunt}（复数）。既然它是单数\Emph{Respondit ei}，我们寻找上文单数名词，发现\Quote{祂的赞美}和\Quote{耶路撒冷}是仅有的单数词。既然有疑问是\Quote{祂的赞美}还是\Quote{耶路撒冷}，我们各按一种解释。\Quote{祂的赞美}如何回答祂？当被祂召唤的人感谢祂时。因为祂召唤，我们回答；不是用声音，而是用信德；不是用舌头，而是用生命。……从祂的选民和圣洁之人，耶路撒冷也回答祂。因为耶路撒冷也被召唤：第一个耶路撒冷拒绝听从，对她说：\Quote{看，你们的房屋将成为荒场留给你们}\BibleRef{玛 23:38}。……但那位耶路撒冷，正如经上所写：\Quote{不生育的，歌唱吧，你未曾生产}，\Quote{她回答祂了}。\Quote{她回答祂了}是什么意思？当祂召唤时，她未轻视。祂降下雨，她结出果实。\par

24. \\Quote{她回答了他：} 但哪里？ \\Quote{在他能力的路上。}…因此教会不是以软弱的方式回答了他；因为在他复活后，他从整个世界召叫了教会，不再在十字架上软弱，而是在天堂上坚强。因为基督教信仰的赞美不是他们相信基督死了，而是他们相信他从死者中复活了。就连外教人也相信他死了，并以此指责你，说你相信了一个死人。那么你的赞美是什么呢？就是你相信基督从死者中复活了，并且你希望你将通过基督从死者中复活：这就是信德的赞美。\\BibleQuote{如果你口里承认耶稣是主，心里相信天主使他从死者中复活了，你必得救。}\\BibleRef{罗 10:9}…这就是基督徒的信德。那么，在这信德中，教会得以聚集，\\Quote{她回答了他，}她按照他的诫命朝拜了他：\\Quote{在他能力的路上，}而非在他软弱的路上。\par

25. 她如何回答了他，你们已经在上文听过了。\\Quote{在万民聚集为一之中。}在此她回答了他，在合一中：谁不在合一中，就不回答他。因为他是独一，教会是合一：只有合一回答独一者。…既然有些人注定要反对她说：她曾经存在，但不再存在；\\Quote{指示我，}他说，\\Quote{我日子的短暂，}这是什么，使我不认识从我这里叛离的人低声反对我？为什么失丧的人争辩说我已经灭亡了？因为他们确实这样说：我曾经存在，但不再存在；\\Quote{指示我我日子的短暂。}我不从你那里询问那些永恒的日子：它们没有尽头，我将要在那里；我不是问那些：我问暂时性的日子；指示我我暂时性的日子；\\Quote{指示我短暂，}不是永恒，\\Quote{我日子的。}告诉我，我在这世界将有多久：因为那些说\\Quote{她曾经存在}而不再存在的人；因为那些说圣经已经应验，万民都已相信，但教会已成为背教者，并从万民中灭亡了的人…\par

26. 难道你们没有看见，仍有一些民族福音还没有被宣讲吗？既然主所说的话必须应验，向教会宣布我日子的短暂，即这福音必须在万民中被宣讲，然后结局才会来临，为什么你们说教会已经从万民中灭亡了，而福音正是为此目的被宣讲，好使它在万民中呢？因此，教会甚至存留到世界的终结，在万民中；这就是她日子的短暂，因为一切有限的都是短暂的；如此她可以从这短暂的存在进入永恒。愿异端者失丧，愿他们现在所是的失丧，并愿他们被寻获，好使他们成为他们所不是的。日子的短暂将持续到世界的终结：短暂的原因，是因为这整个时期——我所说的不是从今天到世界终结，而是从亚当直到世界终结——与永恒相比，不过是一滴水。\par

27. 因此，不要让异端者因我说了\\Quote{我日子的短暂}而自夸，好像这些日子不会持续到世界终结。因为他补充了什么？\\Quote{我的天主，不要在我日子的中途把我带走}\\BibleRef{咏 101:24}。求你不要照异端者所说的对待我。引领我直到世界的终结，不仅到我日子的中间；并完成我短暂的日子，好使你以后赐给我永恒的日子。那么你为什么询问你日子的短暂呢？为什么？你愿意听吗？\\Quote{你的年岁在万代之中。}这就是为什么我询问那些短暂的日子，因为虽然我的日子应持续到世界终结，但与你的日子相比，它们仍是短暂的。因为\\Quote{你的年岁在万代之中。}为什么他不说，你的年岁直到永世万代；因为这样才是在圣经中通常表达永恒的方式；但他却说：\\Quote{你的年岁在万代之中？}但你的年岁是什么？难道不是那些不来临就过去了吗？难道不是那些不来以至于不再停止的吗？因为在这个季节中，每一天都是这样来临又停止；每一个小时，每一个月，每一年；没有一样是固定的；在它来临之前，它将要存在；在它来临之后，它将不再存在。因此，你那些永恒的年岁，那些不改变的年岁，\\Quote{在万代之中。}有一个\\Quote{万代；}在那一代里，你的年岁将要存在。有这样一个，如果我们正确地承认它，我们就要在它里面，天主的年岁就要在我们里面。它们怎样在我们里面呢？就如天主自己要在我们里面：因此经上说：\\BibleQuote{好使天主成为万物中的万有。}\\BibleRef{格前 15:28}因为天主的年岁与天主自己并无不同：天主的年岁就是天主的永恒：永恒是天主真正的本体，没有可变之物；在那里没有过去，好似不再存在；没有未来，好似尚未存在。那里只有\\Emph{是}：那里没有\\Emph{曾是}和\\Emph{将是}，因为过去的不再存在，将来的尚未存在；但凡是那里的，单纯地是。…看这伟大的\\Emph{我是}！人的存在与这相比是什么？对这伟大的\\Emph{我是}，人无论什么，又是什么？谁能理解那\\Emph{存在}？谁能分享它？谁能渴望、向往、妄图他可以在那里？不要绝望，人性的软弱！他说：\\Quote{我是亚巴郎的天主，依撒格的天主，雅各伯的天主。}你们已经听了我自身是什么：现在听我为了你们是什么。那么，这永恒召唤了我们，圣言从永恒中迸发出来。现在它是永恒，现在它是圣言，不再是时间。\par

28. ……祢要从如此多的世代中聚集所有世代的圣洁后裔，并将形成一个世代：\Quote{在} 这 \Quote{世代之世代中有祢的岁月，}也就是说，那永恒将在那个世代中，这世代是从所有世代收集并归于一个；这将分享祢的永恒。其他世代是为了完成他们的时间而生，从这个世代中有一个永远重生；虽然被改变，它将被赋予生命，它将适合承载祢，从祢那里得到力量。\par

29. \Quote{主，祢在起初奠定了大地的根基：诸天是祢手的化工}\BibleRef{咏 101:25}……我们知道天主奠定了大地的根基；诸天是祂手的化工。因为不要想象天主一样用手做，另一样用祂的圣言做。祂用祂的圣言所做的，也是用手所做的：因为祂没有不同的身体肢体，祂说：\Quote{我是自有者。}也许祂的圣言就是祂的手，确实祂的手就是祂的能力。因为经上说：\Quote{要有穹苍}\BibleRef{创 1:6}，就有了穹苍；祂被理解为用祂的圣言创造了它；但当祂说：\Quote{让我们按我们的肖像，按我们的模样造人}\BibleRef{创 1:26}，祂似乎用手造了他。因此，请听：\Quote{诸天是祢手的化工。}看，祂用祂的圣言所造的，也用祂的手所造；因为祂通过祂的卓越、通过祂的能力造了它们。请观察祂造了什么，不要寻求知道祂如何造它们。理解祂如何造它们对你来说是太大的事，因为祂造了你，使你先是一个服从的仆人，然后或许是一个理解的朋友。\BibleRef{若 15:15}\par

30. \Quote{它们将消逝，但祢将永存}\BibleRef{咏 101:26}。宗徒伯多禄明确地说：\Quote{起初，因天主的话，就有了诸天}等等\BibleRef{伯后 3:5}。然后他说诸天已经因洪水消逝了：我们知道诸天消逝到了我们这个大气层的范围。因为水上涨，充满了鸟飞行的整个空间；这样，靠近大地的诸天消逝了；那些当我们说天上的飞鸟时所指的诸天。但还有高于这些在穹苍之上的诸天：但这些是否也要被火消逝，还是只有那些也被洪水消逝的诸天，在学者中是一个更难的问题，也不能轻易地，尤其是在有限的时间内解释。因此让我们放下或推迟它；不过，让我们知道这些事物消逝，而天主永存。……\par

31. 也许这里我们理解诸天为义人自己，天主的圣人们，天主住在他们中间，在祂的诫命中打雷，在祂的奇迹中闪亮，用真理的智慧浇灌大地，因为\Quote{诸天宣告了天主的荣耀}。但它们会消逝吗？它们会在某种意义消逝吗？在什么意义？如同衣服。什么是如同衣服？是指身体。因为身体是灵魂的衣服；我们的主称它为衣服，当祂说：\Quote{生命不是比食物更重要，身体比衣服更重要吗？}\BibleRef{玛 6:25} 那么衣服如何消逝？\Quote{虽然我们外在的人消逝，内在的人却日渐更新。}\BibleRef{格后 4:16} 因此它们将消逝：但就身体而言：\Quote{但祢将永存。}……因此这样的诸天将消逝；但不是永远消逝；它们将消逝，以便被改变。圣咏没有这么说吗？请读下文：\Quote{它们都要像衣服一样陈旧；祢将如同外衣一样改变它们，它们就被改变。}你听到衣服、外衣，难道你只理解为身体吗？我们因此也可盼望我们身体的改变，但来自那在我们之前、在我们之后永存的那位。……\Quote{但祢是同样的，祢的岁月不会穷尽}\BibleRef{咏 101:27}。但我们与这些贫困的岁月相比算什么？它们又算什么？然而我们不应绝望。祂已经在祂伟大而超越的智慧中说过：\Quote{我是自有者}，但祂仍安慰我们说：\Quote{我是亚巴郎的天主，依撒格的天主，雅各伯的天主}\BibleRef{出 3:6}；我们是亚巴郎的后裔\BibleRef{迦 3:29}；即使我们，虽然卑贱，虽然尘土和灰烬，也信赖祂。我们是仆人；但为了我们，我们的主取了仆人的形象\BibleRef{斐 2:7}；为我们这些必死的，那不死者屈尊而死，为了我们祂显示了复活的榜样。因此让我们盼望到达那些持久的岁月，在那里日子不是随着太阳的周转而消逝，而是那存留的如同它自身一样存留，因为唯独祂真正是。\par

32. \BibleQuote{你的仆人的子女要住在此中，他们的后裔要永世坚定}\BibleRef{咏 101:28} : 因为那是万世万代，永恒的世代，长存的世代。但是，\Quote{子女}他说，\Quote{你的仆人的：}难道我们怕我们是天主的仆人，而我们的子女，不是我们自己，住在那裡吗？或者，如果我们是仆人的子女，既然我们是宗徒的子女，我们该说什么？那些后来兴起的子女能有如此不幸的僭妄，竟因他们后来的继承而自夸，以致敢于说：\Quote{我们将在那裡；宗徒们将不在那裡？}愿此事远离他们作为子女的孝爱，作为小孩子的信德，作为成年人的明悟！宗徒们也将在那裡：公羊在前走，羔羊在后跟。那么，为什么说\Quote{你的仆人的子女}而不是简短地说\Quote{你的仆人}呢？他们既是你的仆人，他们的子女也是你的仆人；这些人的子女，他们的孙子，不是你的仆人又是什么呢？你若简短地说，你的仆人要住在此中，你就包括了他们全体。....\Quote{你的仆人的子女}就是指你仆人的工作；没有人能住在那裡，除非藉著他自己的工作。因此，他们的子女要居住是什么意思？谁也不可夸口说他将住在那裡，如果他自称为天主的仆人，却没有工作；因为只有子女将住在那裡。那么，\Quote{你的仆人的子女要住在那裡}是什么意思？你的仆人将藉著他们自己的工作住在那裡，你的仆人将藉著他们自己的子女住在那裡。所以，你如果渴望住在那裡，就不要不结果实；把你的子女先派遣去，好让你能跟随他们，通过先派遣他们，而不是埋葬他们。让你的子女领你到预许之地，活人之地，而不是死人之地；当你在此旅居中生活时，让他们在你前面走，让他们迎接你....\par

\SourceRecord{Translated by J.E. Tweed.；Nicene and Post-Nicene Fathers, First Series；Vol. 8.；Edited by Philip Schaff.；Buffalo, NY: Christian Literature Publishing Co.,；1888.；<http://www.newadvent.org/fathers/1801102.htm>.}

% structure: p0001 source=h1 target=section reason=page-title

\section{《圣咏第103篇释义》}

1. ……\BibleQuote{我的灵魂，请赞颂上主；我的五内，请赞颂他的圣名}\BibleRef{咏 102:1}。窃以为他并非在讲身体内部之事；我并不认为他是指我们的肺、肝等等，要发出赞颂上主的声音。确实，在我们胸中有肺，像某种风箱，不断送出气息，吸入的空气被压出成为声音和声响，当字词被清晰地发出时；任何话语从我们口中发出，无不是被压出的肺所释放；但这并非此处之意；这一切都关系到人的耳朵。天主有耳朵：心也有声音。人对自身内在说话，使它们赞颂天主，对他们说：\Quote{我的五内，请赞颂他的圣名！}你问\Quote{你的五内}是什么意思？就是你的灵魂本身。因此，当他说\Quote{我的五内，请赞颂他的圣名}时，不过是在重复前面的话：\Quote{我的灵魂，请赞颂上主}；因为\Quote{赞颂}一词是被省略的。若有听众，请用声音呼喊；若无人听你，请静默；那聆听你内在一切者从不缺乏。因此，当我们咏唱这些词语时，赞颂已从我们口中发出。我们歌唱了足够的时间，然后沉默：我们内在的心岂能对赞颂上主保持沉默？愿我们声音的响声间歇地赞颂他，而心的声音则是不断的。当你来到教堂咏唱赞美诗时，你的声音传出对天主的赞美；你已尽你所能歌唱；你离开了教堂；愿你的灵魂继续赞美天主。你从事日常事务：愿你的灵魂赞美天主。你进食时，请看宗徒所说：\BibleQuote{你们或吃或喝，无论做什么，一切都要为光荣天主而做。}\BibleRef{格前 10:31}我敢说：当你睡眠时，愿你的灵魂赞美上主。愿罪恶的思念不惊扰你，愿偷窃的图谋不惊扰你，愿诡诈交易的筹划不惊扰你。你的清白，甚至在你睡眠时，就是你灵魂的声音。\par

2. \BibleQuote{我的灵魂，请赞颂上主，不要忘记他的一切恩惠}\BibleRef{咏 102:2}。但上主的恩惠不能在你眼前，除非你的罪恶在你眼前。愿过去的罪乐不在你眼前，而愿罪的定罪在你眼前：来自你的定罪，来自天主的宽恕。因为天主如此赏报你，使你能说：\Quote{上主赐给我的一切恩惠，我该如何报答他？}这正是殉道者（我们今天纪念他们的记忆）以及所有轻视此生的圣人们所考虑的，正如你们在圣若望书信中所听到的：他们为弟兄们舍命，这是圆满的爱德，\BibleRef{若一 3:16} 一如我们的主所说：\BibleQuote{人若为自己的朋友舍命，再没有比这更大的爱了。}\BibleRef{若 15:13}因此，圣殉道者们考虑这些，就轻视了此世的生命，好能在那里找到生命，追随我们主的话：\Quote{爱惜自己生命的，必丧失生命；为我而丧失生命的，必获得永生。}……\Quote{不要忘记}，他说，\Quote{他的一切恩惠}：不是赏报，而是\Quote{恩惠}。因为所欠的是别的，而所付的并非所欠。因此也有这些话：\Quote{我该如何报答上主赐给我的一切恩惠？}你以恶报善；他以善报恶。人啊，你如何以恶报答了你的天主的善？你曾经是亵渎者、迫害者和施暴者，\BibleRef{弟前 1:13} 以亵渎报答了他。为哪些善事呢？首先，因为你存在：但石头也存在。其次，因为你活着，但牲畜也活着。你将如何报答上主，因他按自己的肖像和模样造了你，使你超越一切的牲畜和空中的飞鸟？\BibleRef{创 1:26}不要寻求如何报答他：将他的肖像归还给他；他别无要求，他只要求他的钱币。\BibleRef{玛 22:21}……\par

3. 灵魂啊，当你在思想你的一切恶行时，要思考天主的一切恩惠：因为你的罪有多少，他的恩惠也有多少。然而，你能向他献上什么礼物、什么奉献、什么祭品呢？……你将用什么报答上主？因为你曾思考，却找不到：\Quote{我要举起救恩的杯爵。}什么？难道上主自己没有赐予救恩的杯爵吗？若你能，就用你自己的来报答他？我要说：不，不要这样做；不要用你自己的来报答他；天主不愿被用你自己的来报答。若你用你自己的来报答他，你是在报答罪。因为凡你所拥有的，你都是从他所领受的；只有罪是你自己所有的。他不愿被用你所有的来报答，他只愿被用他自己的来报答。正如你若从农夫所播种的土地上带来麦穗，你是用农夫自己的出产报答了他；若带来荆棘，你是用你自己的献给了他。要报答真理，要在真理中赞美上主：若你选择用你自己的来报答他，你就在说谎。那说谎言的人，是从自己的私欲说的。\BibleRef{若 8:44}若那说谎言的人是从自己的私欲说的，那说真理的人就是从主说的。但举起救恩的杯爵是什么意思？岂不是效法我们主的苦难吗？……我要举起基督的杯爵，我要饮我主的苦难。要当心，免得你跌倒。但是，\Quote{我要呼求上主的名号。}因此，那些跌倒的人没有呼求上主；他们仗恃自己的力量。你要这样回报，如同记着你所回报的是你所领受的。因此，让你的灵魂赞颂上主，不要忘记他的一切恩惠。\par

4. 你们都要听祂的一切酬报。\BibleQuote{祂赦免你的一切罪，治愈你的一切疾苦} \BibleRef{咏 102:3}。看，祂的酬报。除了惩罚，罪人还该得到什么？除了地狱，亵渎者还该得到什么？祂没有给予这些酬报：使你们不因恐惧而战栗，并且在没有爱的情况下敬畏祂。……但你是个罪人。回转吧，领受祂的这些酬报：祂\BibleQuote{赦免你的一切罪。}……然而即使在罪赦之后，灵魂本身仍会被某些情欲所动摇；她仍在诱惑的危险之中，仍被某些建议所取悦；对有些她不喜欢，有时她又同意她所喜欢的那些：她被俘虏了。这是软弱；但祂\BibleQuote{治愈你的一切疾苦。}你的一切疾苦都将被治愈：不要害怕。你说，它们很严重；但医生更伟大。没有一种疾苦在全能的医生面前是不可治愈的：只要你让自己被治愈；不要推开祂的手；祂知道怎样对待你。不仅要在祂抚慰你时感到喜悦，也要在祂使用刀时忍耐：忍受治疗的痛苦，反思你未来的健康……你并不是在不确定中忍受：那许诺你健康的天主不会被欺骗。医生常被欺骗：并许诺人身体的健康。他为什么被欺骗？因为他所治愈的并非他自己的创造。天主造了你的身体，天主造了你的灵魂。祂知道如何恢复祂所造的，祂知道如何重塑祂已塑造的：你只需在医生的手下忍耐；因为祂厌恶拒绝祂手的人。在人医生的手下，不会发生这种事……\par

5. \BibleQuote{祂救赎你的性命脱离败坏} \BibleRef{咏 102:4}。看，\BibleQuote{那腐败的身体，重压着灵魂} \BibleRef{智 9:15}。灵魂于是在一个可朽坏的身体里有了生命。什么样的生命？它承受负荷，它肩负重量。对思想天主本身——人应当怎样思想天主——有多么大的障碍？仿佛因人类腐败的必要性而打断我们？有多少影响力召回我们，有多少事物打断我们，又有多少事物在心灵凝注高处时将其引开？何等众多幻象，何等成群的暗示！这一切在人心之中，仿佛充满了人类腐败的蛆虫。我们已经陈述了疾病的严重性，让我们也赞美医生。那么，祂岂不医治你？祂造了你，使你本来可以健康，只要你愿意遵守所接受的健康法则？……首先要想到你自身的健康。有时一个人在自己家中、在床上，被一种比平常更明显的疾病所击倒；尽管这种疾病——人们不愿正视它——也很明显；但每人都可能患上那种需要求助于人医生的疾病，并在床上因发烧而喘息；也许他想要考虑家务，对自己的财产或房子做出某种安排或处置；立刻他被朋友们明显的焦虑从这些忧虑中唤回，并被劝告放下这些事，首先关心自己的健康。那么，这是对你们和所有人说的：如果你没有生病，就考虑其他事；如果你的软弱证明你确实生病了，就首先要留意你的健康。基督就是你的健康：所以你要思想基督。领受祂救恩的健康之杯，祂\BibleQuote{治愈你的一切疾苦；}如果你选择，你将获得这健康。……因为你的生命已从腐败中被赎回：现在安息吧，一份忠信契约已经订立；没有人欺骗，没有人陷害，没有人压迫你的救赎者。祂在这里做了一次交换，祂已经付了代价，祂倾流了祂的血。天主独生子，我说，已为我们倾流了祂的血：噢，灵魂，抬起你自己，你是如此贵重……祂\BibleQuote{救赎你的性命脱离败坏。}\par

6. \BibleQuote{祂以仁爱和慈悲给你戴上冠冕。}你也许在听见\BibleQuote{祂给你戴上冠冕}这句话时，开始有点骄傲了。那么我是伟大的，我角力了。靠着谁的力量？靠你的，但由祂供应……祂给你戴上冠冕，因为祂是在加冕祂自己的恩赐，不是你的功绩。宗徒说：\BibleQuote{我比他们众人更劳碌了，}但看他接着说什么：\BibleQuote{其实不是我，而是天主的恩宠偕同我。}\BibleRef{格前 15:10}……因此，你是因祂的仁慈而受加冕；在什么事上都不要骄傲；要常常赞美主；不要忘记祂的一切酬报。当罪人、不虔敬的人被召叫，好能成义，这是一种酬报。当你被扶起并被引导，以免跌倒，这是一种酬报。当赐给你力量，使你恒心至终，这是一种酬报。甚至那曾压迫你的肉身复活，且你的一根头发也不失落，这是一种酬报。在你复活之后，你被加冕，这是一种酬报。你将永不停息地赞美天主本身，这是一种酬报……\par

7. 战斗之后，我将接受加冕；加冕之后，我将做什么？\Quote{祂以美物满足你的愿望}\BibleRef{咏 102:5}……寻求你自己的善，灵魂啊。因为一种造物有一种善，另一种造物有另一种善，所有造物都有自己的某种善，以完善和成全其本性。对于每个不完美之物，使之得以完美所不可或缺的善，各有不同；寻求你自己的善吧。\Quote{除了唯一的——天主——之外，没有善的}\BibleRef{玛 19:17}至善就是你的善。那么，享有至善的人还缺少什么呢？因为尚有次等的善，它们分别对不同造物是善的。弟兄们，对牲畜而言，什么算是善呢？不就是填饱肚子、免于匮乏、睡觉、纵欲、存在、健康、繁衍吗？这对它们就是善，并且在一定限度内，有造物主天主所赐予的定量的善。你也寻求这样的善吗？天主也赐予这些，但不要只追求它们。作为\Person{基督}的共同继承人，你能与牲畜一同在共融中喜乐吗？将你的望德提升至万善的善吧。祂将成为你的善，是祂使你在你的种类中成为善，也是祂使一切事物在它们的种类中成为善。因为天主造的一切都非常好……\par

8. 我的渴望何时才能以美物得到满足？你问：何时？\Quote{你的青春必如鹰般更新。}那么你问：你的灵魂何时能以美物得到满足？当你的青春恢复时。他又补充说，如鹰般。这里隐含着某种深意；但我们不会对有关鹰的事默然不语，因为理解它并非与我们的目的无关。只让这一点铭刻在我们心中：这并非圣神无故而言。因为它暗示给我们一种复活。事实上，鹰的青春得以恢复，却不是进入不朽，因为这里给出的只是一个比喻，尽可能从可朽之物引向不朽之物的喻意，而非证明。据说，当鹰因年老体衰时，由于它的喙过分增长，而无法进食。因为其喙的上部长期增长后，会妨碍张嘴。这种衰老对它造成的影响。于是，凭借一种天然的方法，据说鹰会把喙的上唇撞向岩石，以此磨去阻碍进食的老喙。然后它就能进食，一切恢复如初；它仿佛返老还童，肢体活力重现，羽毛光泽，翅膀有力，如从前高飞，一种复活在它身上发生。这就是比喻的目的，如同月亮的比喻，月缺之后又复圆，向我们显示复活；但月亮圆后并不长存，再次亏缺，为使这象征永不停止。同样，关于鹰所说的也是如此：鹰并非恢复到不朽，而我们却是恢复到永生；但比喻由此而来：那岩石除去了阻碍我们的东西。所以，不要依靠你的力量：岩石的坚固磨去了你的老年，因为那磐石就是\Person{基督}。\BibleRef{格前 10:4}在\Person{基督}内，我们的青春将如鹰般恢复……\par

9. \Quote{主为一切受欺压的人施行仁慈和正义}\BibleRef{咏 102:6}……有一个犯奸淫的女人被带来，按照法律要用石头砸死，但她被带到立法者本人面前。……我们的主，当她被带到祂面前时，祂低头，开始在地上写字。当祂俯身在地上时，祂在地上写字；在祂俯身之前，祂不是在地上写字，而是在石头上写字。那时，地已是肥沃的，准备从主的文字中结出果实。在石头上，祂写下了法律，暗示犹太人的心硬；祂在地上写字，象征基督徒的多产。然后，带领那妇人的那些人来了，如同怒涛冲击磐石，却被祂的回答撞得粉碎。因为祂对他们说：\Quote{你们中间谁没有罪，就先向她投石吧}\BibleRef{若 8:7}。祂又低头，继续在地上写字。这时，每个人查看自己的良心，都没有上前。并不是那个软弱的犯奸淫的妇人，而是他们自己奸淫的良心，使他们退后。他们想要惩罚，想要审判；他们来到磐石前，他们的审判却被磐石推翻……\par

10. 向恶人施行仁慈，但不是作为恶人。不要接纳恶人，只要他仍是恶人：也就是说，不要好像出于对他不义的倾向和喜爱而接纳他。因为禁止施予罪人，也禁止接纳罪人。然而，这话怎么说：\BibleQuote{凡求你的，就给他}？还有这话：\BibleQuote{如果你的仇人饿了，就给他吃}？\BibleRef{罗 12:20} 这看似矛盾：但它在那些以基督之名敲门的人面前是敞开的，在那些寻求的人面前将是清楚的。\BibleQuote{不要帮助罪人}和\BibleQuote{不要施予不虔敬的人}；\BibleRef{德 12:4} 然而又说：\BibleQuote{凡求你的，就给他}。但向我求的却是个罪人。施予，但不是作为罪人。什么时候你作为罪人施予呢？当你因那使他成为罪人的事而喜悦，以致你施予的时候……让那些施予斗兽者的人告诉我，他们为什么施予？他为什么施予这个人？他爱在他身上那构成他最大罪恶的东西；他喂养他，他给他穿衣，就是邪恶本身，被所有见证者公开宣扬。为什么施予演员、赛车手或娼妓的人施予呢？这些不也是施予人吗？但他们所注意的不是天主的化工的本性，而是人类作品的不义……因此，当你施予时，你施予的是耻辱，不是勇敢。因此，正如施予斗兽者的人，不是施予那人，而是施予一种最可耻的职业；因为如果他只是一个人，而不是斗兽者，你就不会施予；你在恶习中而非本性中尊敬他：同样，另一方面，如果你施予义人，如果你施予先知，如果你施予基督的门徒任何他所缺乏的东西，却不思想他是基督的门徒，是天主的仆人，是天主的管家；而是在那情况下思想某些暂时的利益，例如，当他或许对你的事业有必要时，他可能会因为给了你一些东西而被你收买；如果你这样施予，你并没有比施予斗兽者的人更施予那人。因此，最亲爱的，这事对我们来说是相当开放的，我认为，虽然它曾是隐晦的，但现在清楚了。主正是为此约束了你，当他说：\BibleQuote{那接待义人的}。那已经足够了。但义人可能以另一种意图被接待……他说：\BibleQuote{因义人之名接待义人的}：就是因他的义德而接待他……也就是说，因为他是基督的门徒，因为他是奥迹的管家：\BibleRef{格前 4:1} \BibleQuote{我实在告诉你们，他决不会失去他的赏报。}\BibleRef{玛 10:42} 所以你要明白，因罪人之名接待罪人的，将要失去他的赏报。\par

11. ……因此，你要无所畏惧地仁慈，甚至向你的仇敌扩展爱德：惩罚那些偶然属于你管辖的人，以情感、以爱德约束他们，为了他们的永恒救恩；免得当你顾惜肉体时，灵魂丧亡。这样做：虽然你必须忍受许多人，你对他们不能行使纪律，因为你对他们没有合法权威；忍受他们的伤害；不要忧虑。如果你曾是仁慈的，他必向你显示仁慈：你将仁慈，而你所受的伤害不会失去它们的惩罚；\BibleQuote{报复属于我，我必要报酬}，\BibleRef{申 32:35}，主说。\par

12. \BibleQuote{祂使自己的道路显示给梅瑟}\BibleRef{咏 102:7}……因为法律被赐下，正是为此目的：使病人信服自己的软弱，并呼求医生。这是天主隐藏的道路。你早已听说过：\BibleQuote{祂医治你的一切病症}。他们的病症尚隐藏在病人中；五卷书赐给了梅瑟：池子被五个门廊环绕；他带出病人，让他们躺在那里，使他们被显露，而不是被医治。五个门廊揭示了病人，却没有医治他们；池子在水被搅动时医治了下水的人：\BibleRef{若 5:2} 池子的搅动是在我们主的受难中……因此，既然那里是一个奥迹，他教导说，法律被赐下是为了使罪人被定罪，并呼求医生以领受恩宠……因此，如我刚开始所说的，因为法律中有一个伟大的奥迹：法律被赐下正是为此观点，即借由罪的增多，骄傲的人被贬抑，被贬抑的承认，承认的得医治；这些是隐藏的道路，祂使这些道路为梅瑟所知，并通过他赐下法律，藉此罪增多，恩宠更加增多……\BibleQuote{祂已将自己的美意显示给以色列子民。}对所有以色列子民吗？对真正的以色列子民；是的，对所有以色列子民。因为奸诈的、阴险的、虚伪的人，不是以色列子民。谁是真正的以色列子民？\BibleQuote{看，一个真正的以色列人，心中毫无诡诈。}\BibleRef{若 1:47}\par

13. \BibleQuote{主满怀慈悲与仁慈，含忍，且广施仁慈} \BibleRef{咏 102:8}。为何如此含忍？为何如此广施仁慈？人犯罪，却仍存活；罪上加罪，生命依旧；人每日亵渎，而\BibleQuote{祂使祂的太阳上升，光照善人，也照恶人。} \BibleRef{玛 5:45}祂从各方召叫人改过，从各方召叫人悔改：藉受造的祝福召叫，藉赐予生命的时间召叫，藉读者召叫，藉宣讲者召叫，藉内心的思想、藉管教的杖召叫，藉安慰的仁慈召叫：\BibleQuote{祂含忍，且广施仁慈。} 但你要谨慎，免得滥用天主仁慈的久长，为自己积蓄忿怒，正如宗徒所说，在忿怒之日……因为有些人准备回头，却一再拖延，在他们心中响起乌鸦的声音：\Quote{Cras! Cras!} 那从方舟放出的乌鸦，再没有回来。 \BibleRef{创 8:7}天主寻求的，不是乌鸦声音中的拖延，而是鸽子哀鸣中的告解。鸽子被放出后，回来了。要到何时，明天！明天！？你要留意你的最后明天：既然你不知道你的最后明天是哪一天，你活到今日仍为罪人，这已足够了。你曾听到，你惯常听到，你今天也听到了；你每日听到，却每日不改……\par

14. \BibleQuote{祂不会永远责斥，也不长久怀怒。} \BibleRef{咏 102:9} 既然因祂的愤怒，我们活在鞭打与可朽的腐败中：这是我们对原罪的惩罚……我的弟兄们，难道不是因祂的愤怒，才有\BibleQuote{你必汗流满面，劳苦才能吃饭，地要给你生出荆棘和蒺藜}？这是对我们祖先说的。或者，如果我们的人生与此不同，你若能转向某种享乐，在那里你不感到荆棘，就选择你所愿意的：无论是贪婪还是纵欲，仅举此二者；再加上第三个情欲，即野心；在追求荣誉中有多少荆棘？在情欲的纵乐中有多少荆棘？在贪婪的热望中有多少荆棘？在卑劣的恋爱中有多少烦恼？今世有多少可怕的焦虑？我撇开地狱不提。你要当心，免得你现在就成为自己的地狱。这一切，我的弟兄们，都是祂愤怒的结果：当你转向正义的行为时，你不得不在地上劳苦；劳苦不会在生命结束前结束。我们必须在路上劳苦，好能在我们的家乡欢乐。因此，祂用祂的应许安慰你的劳苦、你的辛劳、你的烦恼，对你说：\BibleQuote{祂不会永远责斥。}\par

15. \BibleQuote{祂没有按我们的罪恶对待我们} \BibleRef{咏 102:10}。感谢天主，因为祂恩准了这事。我们没有得到我们应得的：\BibleQuote{祂没有按我们的罪恶对待我们，也没有按我们的过犯报复我们。} \BibleQuote{就如天离地有多高，上主对敬畏祂的人的慈爱也多高} \BibleRef{咏 102:11}。你观察天：它从四面八方覆盖大地，没有一处大地不被天覆盖。人在天下犯罪，他们行一切恶事都在天下，却仍被天覆盖。从那里有眼睛的光，有空气，有呼吸，有为了果实而降在地上的雨水，有来自天上的一切仁慈。若将天对地的援助撤去，大地立刻就会荒废。因此，正如天的保护停留在大地之上，上主的保护也停留在敬畏祂的人之上。你敬畏天主，祂的保护就在你之上。但或许你受鞭打，便以为天主舍弃了你。如果天的保护舍弃了大地，天主就舍弃了你。\par

16. \BibleQuote{看，东离西有多远，祂使我们的罪离我们也有多远。} \BibleRef{咏 102:12} 认识圣事的人知道这事；不过，我只说众人能听的话。当罪被赦免时，你的罪下降，你的恩宠上升；你的罪仿佛在衰落，那释放你的恩宠在上升。\Quote{真理从地上生出。} 这是什么意思？你的恩宠诞生，你的罪下降，你在某种意义上被更新。你应看向上升的，而转离下降的。转离你的罪，转向天主的恩宠；当你的罪下降时，你上升并且获益……天的一方下降，另一方上升；但正在上升的一方在十二小时后将下降。那上升临到我们的恩宠并不如此：我们的罪永远下降，而恩宠永远常存。\par

17. \BibleQuote{就如父亲怎样怜爱自己的儿女，上主对敬畏祂的人也怎样仁慈。} \BibleRef{咏 102:13} 任凭祂随意发怒，祂是我们的父亲。祂鞭打了我们，折磨了我们，击伤了我们：祂仍是我们的父亲。儿子，你若哀哭，就在你父亲下面哀哭；不要带着愤怒，不要带着骄傲的膨胀。你所受的，你所悲伤的，是药物，而非惩罚；是管教，而非定罪。你若不愿被拒绝继承产业，就不要拒绝鞭打：不要想你在鞭打中受了什么惩罚，而要想你在遗嘱中有什么地位。\par

18. \Quote{因为祂知道我们的形成} \BibleRef{咏 102:14}：即我们的软弱。祂知道祂所创造的，如何堕落，如何得以修复，如何得以被收纳，如何得以被充实。看，我们是用泥土造的：\Quote{第一个人出于地，属于土；第二个人是主，来自天上。} \BibleRef{格前 15:47} 祂甚至派遣了祂自己的独生子，祂成了第二个人，祂在万有之先就是天主。因为祂在来临上是第二位，在回归上是第一位：祂在许多人之后死去，在所有人之前复活。\Quote{祂知道我们的形成。} 什么形成？我们自己。你为什么说祂知道？因为祂起了怜悯。\Quote{求祢记念，我们不过是尘土。} 对天主自己说话，他说：\Quote{求祢记念，} 好像天主会忘记：祂如此知道，以至不能忘记。但 \Quote{记念} 是什么意思？愿祢的仁慈继续向我们施予。祢知道我们的形成，求祢不要忘记我们的形成，免得我们忘记祢的恩宠。\par

19. \Quote{世人，他的岁月如同草} \BibleRef{咏 102:15}。愿人思想他是什么；愿人不骄傲。\Quote{他的岁月如同草。} 草为何骄傲，它现在繁茂，但很快便枯干？草为何骄傲，它只繁茂一瞬，直至日头灼热？所以仁慈临于我们是有益的，使草变成金子。\Quote{他如同田野的花一样繁盛。} 人类的一切荣华：尊位、权势、财富、骄傲、威吓，都是草的花。那座房屋繁盛，那个家族伟大，那个家族繁盛；多少人繁盛，他们活多少年！许多年于你，在天主看来不过是一瞬。天主不象你那样计算。与世代的长久寿命相比，任何家族的繁盛都如同田野的花。一年的全部华美几乎不能持续一年。无论什么在那里繁盛，那里受热温暖，那里美丽，都不能持久；不，它不能存留一整年。花朵在一段多么短暂的时光中便消逝，而这些是青草的美丽！这极为美丽的东西，很快便凋谢。\BibleRef{依 40:6} 既然祂如父亲般知道我们的形成，知道我们不过是草，只能暂时繁盛；祂就派遣了祂的圣言给我们，祂的圣言永远常存，祂使圣言成为了不常存的草的弟兄。你不要惊奇你将分享祂的永恒；祂自己首先成为了你的草的分享者。祂既然承担了来自你的卑微之事，难道会拒绝把关于你的崇高之事赐给你吗？\par

20. \Quote{风一吹过，就不复存在；它的地方也不再认识它} \BibleRef{咏 102:16}。因为祂所说的不是草，而是那甚至使圣言成为草的缘故。因为你是人，为了你，圣言成了人。\Quote{凡有血肉的，尽都如草；} \Quote{圣言成了血肉。} \BibleRef{若 1:14} 那么，草有何等大的希望，因为圣言成了血肉？那永远常存者，没有不屑于承担草，使草不致对自己绝望。\par

21. 因此，在你对自己的反思中，要思想你的卑下，要思想你的尘土：不要高傲；你若有什么更好的，那是因祂的恩宠，因祂的慈悲。因为请听下文：\Quote{但主的仁慈永远临于敬畏祂的人} \BibleRef{咏 102:17}。那不敬畏祂的人，将是草，且在草中，且在痛苦中与草一同：因为肉身将要复活，进入痛苦。愿敬畏祂的人喜乐，因为祂的仁慈临于他们。\par

22. \Quote{祂的公义归于子子孙孙} \BibleRef{咏 102:18}。祂说到报酬，\Quote{归于子子孙孙。} 有多少天主的仆人没有子女，更何况子子孙孙呢？但祂称我们的工作为我们的子女；工作的报酬为我们的 \Quote{子子孙孙。} \Quote{归于那些遵守祂盟约的人。} 愿人警惕，不可都以为这里所说的属于自己：让他们选择，趁他们还有选择的机会。\Quote{并记念祂的诫命，遵行它们。} 你本已准备好恭维自己，或许要向我背诵圣咏集，虽然我不曾记在心里，或凭记忆说出全部法律。显然，你的记忆比我更好，比任何不逐字记住法律的义人更好：但要当心，要持守诫命。但你如何持守？不是凭记忆，而是凭生活。\Quote{记念祂的诫命的人：} 不是背诵它们，而是 \Quote{遵行它们。} 现在或许每一个人的灵魂都感到不安。谁记得天主的全部诫命？谁记得天主的全部著作？看，我不仅希望将它们记在心里，也希望通过我的工作去遵行它们：但谁能全部记得它们？不要怕：祂不加重你的负担：\Quote{全部法律和先知，都系于这两条诫命。} \BibleRef{玛 22:40}\par

23. \Quote{主在天上预备了祂的宝座} \BibleRef{咏 102:19}。除了基督，谁在天上预备了祂的宝座？那下降又上升的，那死去又从死者中复活的，那将所取的人性高举到天上的，祂亲自在天上预备了祂的宝座。宝座是审判者的座位：因此你们这些听众要注意，\Quote{祂在天上预备了祂的宝座。}……王国是属于主的，祂要在万民中作元首。\Quote{祂的王国要统治万有。}\par

24. \Quote{上主的天使，你们要赞颂上主，你们是雄大能者，遵行祂的圣言} \BibleRef{咏 102:20}。所以，凭着天主的圣言，你们不是义人，也不是忠信者，除非你们遵行它。\Quote{你们是雄大能者，遵行祂的诫命，听从祂言语的声音。}\par

25. \Quote{所有祂的军旅，请讚颂上主：是祂的臣僕，奉行祂旨意者}\BibleRef{咏 102:21}。众天使，一切大能者，奉行祂圣言者，祂的所有军旅，祂的臣僕、奉行祂旨意者，请讚颂上主。因为凡生活邪恶的人，即便舌头沉默，他们的唇舌也在诅咒上主。倘若你的舌头讚颂讚歌，而你的生活却散发出亵渎，这有何益？你因生活恶劣已使许多舌头陷于亵渎。你的舌头被赋予讚歌，而目睹你之人的舌头却用于亵渎。因此，你若愿讚颂上主，就当奉行祂的圣言，承行祂的旨意……\par

26. \Quote{上主的一切化工，请在祂统治的各处，讚颂上主}\BibleRef{咏 102:22}。因此是在各处。愿祂不在祂不统治之处受讚颂：\Quote{在祂统治的各处}。或许无人会说：我无法在东方讚颂上主，因祂已去了西方；或我无法在西方讚颂祂，因祂在东方。\Quote{因为不是从东方，也不是从西方，也不是从旷野的山岭而来。为什么？因为天主是审判者。}祂处处都在，如此处处可受讚颂；祂如此周遍，使我们在各处因祂而喜乐；祂如此受讚颂于各处，使我们在各处生活美好……\Quote{在祂统治的各处：我的灵魂，请讚颂上主！}最后一节与第一节相同：讚颂在圣咏之首，讚颂在末尾；从讚颂出发，回归于讚颂，在讚颂中为王。\par

\SourceRecord{Translated by J.E. Tweed.；Nicene and Post-Nicene Fathers, First Series；Vol. 8.；Edited by Philip Schaff.；Buffalo, NY: Christian Literature Publishing Co.,；1888.；<http://www.newadvent.org/fathers/1801103.htm>.}

% structure: p0001 source=h1 target=section reason=page-title

\section{论圣咏第一百零四篇}

1. ...\Quote{我的灵魂，请讚颂主。}愿我们众人因基督合而为一的灵魂，说出这话。\Quote{主，我的天主，祢何其伟大！}\BibleRef{咏 103:1}。祢在何处受到颂扬？\Quote{祢以尊威和威严为衣。}请你们宣认，好使你们受美化，使祂穿上你们。\Quote{身披光明，好像外氅}\BibleRef{咏 103:2}。祂穿上祂的教会，因为教会因祂而成为\Quote{光明}，而她在自身原是黑暗，如宗徒所说：\Quote{你们从前原是黑暗，但现在在主内却是光明。}\BibleRef{弗 5:8}\Quote{展开苍天，如同帐幔：}要么如同你轻易展开一块皮革，若此解释为\Quote{轻易}，则你可按字面理解；要么我们以\OriginalTerm{皮帐}{skin}之名理解圣经的权威，遍布普世；因为皮革象征有死性，但全部神圣圣经的权威通过有死之人分施给我们，这些人虽已去世，声誉仍传播四方。\par

2. \Quote{以水遮盖你的宫殿}\BibleRef{咏 103:3}。什么宫殿？是天的宫殿。什么是天？我们喻指神圣圣经。神圣圣经的宫殿是什么？是爱的诫命，没有再比这更高的了。\BibleRef{谷 12:31}但为何将爱比作水？因为\Quote{天主的爱，借着所赐与我们的圣神，已倾注在我们心中了。}\BibleRef{罗 5:5}。圣神自身何时成为水？因为\Quote{耶稣站着大声喊说：谁信从我，从他的心中要流出活水的江河。}\BibleRef{若 7:37}。我们凭什么证明这是指圣神说的？让福音作者亲自说明，他接着写道：\Quote{他说这话，是指那信仰祂的人将要领受的圣神。}\Quote{乘驾风的翅膀；}即驾于众灵魂的德行之上。灵魂的德行是什么？就是爱本身。但祂如何驾于其上？因为天主对我们的爱大于我们对天主的爱。\par

3. \Quote{祢以风为使者，以火焰为仆役}\BibleRef{咏 103:4}：即那些已为神、属灵而非属血肉的人，祂使他们成为使者，派遣他们宣讲祂的福音。\Quote{以火焰为仆役。}因为传道的仆役若不火热，就不能点燃他所传道的人。\par

4. \Quote{祢奠定了大地在基础上}\BibleRef{咏 103:5}。祂建立了教会于教会的基础上。什么是教会的基础？不就是宗徒所说的：\Quote{除已奠立了的根基，即耶稣基督外，任何人不能再奠立别的根基。}\BibleRef{格前 3:11}。因此，建立于如此根基之上，她配得听见什么？\Quote{她永不动摇。}\Quote{祢奠定了大地在基础上。}即祂将教会建立在基督这个基础上。若基础动摇，教会就要动摇；但基督何时会动摇呢？基督来到我们中间，取了血肉，\Quote{万物是借着祂而造成的，凡受造的，没有一样不是由祂造成的；}\BibleRef{若 1:3}祂以大能维系万有，\BibleRef{希 1:3}以良善维系我们。既然基督不失败，\Quote{她永不动摇。}那些说教会已从世上消亡的人在哪里呢？连摇动都不可能……\par

5. \Quote{深渊如上衣是它的衣服}\BibleRef{咏 103:6}。谁的衣服？或许是天主的？但之前已说到祂的衣袍：\Quote{身披光明，好像外氅。}我听见天主披着光明，而那光明，如果我们愿意，就是我们。什么是\InnerQuote{如果我们愿意}？如果我们不再是黑暗。因此，如果天主身披光明，那么深渊又是谁的衣袍呢？因为一大片水被称为深渊。所有水、所有潮湿的本性，以及流布于海洋、江河和隐秘洞穴的物质，统称为深渊。因此我们理解大地，就是祂所说的\Quote{祂奠定了大地。}我相信祂说的正是：\Quote{深渊如上衣是它的衣服。}因为水犹如大地的衣服，环绕并覆盖它……\par

6. \Quote{大水立于山岭之上：}即大地的衣服，就是深渊，如此增加，以至水甚至立在山岭之上。我们在洪水时期读到这事的记载……先知意图预告未来之事，而非追述已往之事，因此他说这话，是因为他要让人明白教会将在迫害的洪水中。曾有一时，迫害者的洪流覆盖了天主的土地、天主的教会，甚至那些作为山岭的伟大人物也看不见了。因为他们四处逃散，岂不就不显现了吗？也许那水就是那话：\Quote{天主，求祢救我，因为水已淹至我的灵魂。}尤其是那形成海洋的水，翻腾不结果实。因为凡海洋之水覆盖的土地，它不会使之结果实，反而使之荒芜。因为水底下也有山岭，因为\Quote{大水立于山岭之上。}为何宗徒们因逃亡而隐藏？因为\Quote{大水立于山岭之上。}水的权势虽大，但能持续多久？请听下文。\par

7. \\BibleQuote{因祢的斥责，大水逃避} \\BibleRef{咏 103:7}。的确如此，诸位弟兄；因天主的斥责，大水真的飞逃了；也就是说，它们从压迫群山之处退去。如今群山本身耸立，\\Person{伯多禄}和\\Person{保禄}：他们何等巍峨！先前受迫害者压制，如今受皇帝尊崇。因为大水已从天主的斥责前飞逃，因为\\BibleQuote{君王的心在天主手中，祂随意扭转} \\BibleRef{箴 21:1}；祂命令他们赐给基督徒和平；宗徒的权威兴起并高耸……大水从天主的斥责前飞逃。\\Quote{从祢雷声的威吓，他们必恐惧。}如今有谁不畏惧呢？从天主透过宗徒发出的声音，从天主透过圣经、透过祂的云发出的声音？大海平息，众水恐惧，群山裸露，皇帝下了命令。但若非天主发雷，谁会下命令呢？因为天主愿意，他们便下令，事就成了。因此，人不可将任何事归功于自己。\par

8. \\BibleQuote{山岳上升，盆地下降，到祢为它们所定的地方} \\BibleRef{咏 103:8}。他仍在谈论水。我们不可在这里理解为陆地的山，也不是陆地的平原，而是那些大得可比作山的波浪。大海曾一度翻腾，其波浪如山，足以遮盖那些作为山的宗徒。但群山上升，平原下陷，要持续多久呢？他们曾狂怒，如今平息。他们狂怒时是山，如今平息了就成了平原：因为祂为他们奠定了一个地方。有一条水道，仿佛一个深渊，所有那些不久前狂怒的凡人之心都已退入其中……他们从前是山，如今是平原；然而，我的弟兄们，即使风平浪静也是海。为何他们现在不再猛烈？为何不再狂怒？为何他们不再试图，即使不能推翻我们的陆地，至少也要覆盖它？为何不呢？\par

9. 请听。\\BibleQuote{祢划定了界限，不得越过，不再回来掩没大地} \\BibleRef{咏 103:9}。那么，既然如今最苦涩的波浪已接受了限度，我们得以自由地宣讲这些事；因为他们已得到应有的界限，不能越过所定的界限，也不再回来覆盖大地；在大地本身发生着什么？其中有什么运作，如今大海离开它使之裸露？虽然在它的海滩上，轻微的波浪仍发出声响，虽然外教人仍在周围喃喃细语；我听见海岸的声音，却不惧怕洪水。那么，大地之中发生着什么？\\BibleQuote{谁在溪谷中发出泉水？} \\BibleRef{咏 103:10}。\\Quote{祢发出，} 他说，\\Quote{泉水在溪谷中。} 你们知道什么是小山谷，是地中的低洼处。因为与小丘和山相对，山谷和小山谷的形状相反。丘陵和山是土地的隆起；而山谷和小山谷是土地的凹陷。不要轻视低处，泉水从那里流出。\\BibleQuote{祢使泉水涌于溪谷} \\BibleRef{咏 103:10}。听一座山。宗徒说：\\Quote{我比他们更劳苦}。一种伟大呈现在我们面前；然而，为了水流，他立刻使自己成为山谷：\\BibleQuote{然而不是我，而是天主的恩宠偕同我} \\BibleRef{格前 15:10}。说他们是山同时也是山谷，并无矛盾：因为被称为山是由于他们属灵的伟大，被称为山谷是由于他们心灵的谦卑。\\Quote{不是我，} 他说，\\Quote{而是天主的恩宠偕同我。}……\par

10. 什么是\\Quote{在山中间，水流穿行}？我们已听说过谁是\\Quote{山}，是圣言伟大的宣讲者，是天主崇高的天使，虽然仍在有死的肉身中；他们并非凭自己的力量，而是凭祂的恩宠而崇高；但就他们自身而言，他们是山谷，以谦卑发出泉水。\\Quote{在中间，} 他说，\\Quote{在山中间，水流穿行。} 让我们假设这是这样说的：\\Quote{在宗徒之间，真理之言的宣讲将穿行。} 什么是在宗徒之间呢？被称为\\Quote{在中间}的是公共的。一份公共的财产，所有人都赖以生活，是在中间，不属于我，也不属于你，也不属于我……因为如果他们不在中间，就好像是私有的，不为公用而流，我有我的，他有他的，不在中间使我和他都能拥有；但和平的宣讲并非如此……因此，弟兄们，愿我们向你们的爱德所说的，因着泉水的作用达到此目的：愿它们从你们流出，你们要成为山谷，并与众人分享你们从天主所获得的。愿水在中间流，不要嫉妒任何人，饮，充满，当你们充满时流溢出去。愿天主的公共之水处处得荣耀，而非人的私谬……\par

11. 因为接着是：\\BibleQuote{林中的野兽都来饮水} \\BibleRef{咏 103:11}。我们确实也在可见的受造界中看到这一点，林中的野兽饮于泉水和山间奔流的溪水；但如今，既然天主乐意将祂自己的智慧隐藏在这类事物的形象中，并非要夺去热切寻求者，而是要向不关心者关闭、向叩门者敞开；祂也乐意亲自藉着我们劝勉你们，使我们能在所有这些似乎是关于有形可见受造界的事物中，寻求某种属灵隐藏的东西，当我们寻见时便欢喜。林中的野兽，我们理解为外邦人，圣经在许多地方为此作证……\par

12. 这些野兽，于是，饮用那些水，但却是暂时的；不是停留，而是经过；因为所有在此时分施的教导都是经过的……除非你们的爱德认为，在那一座城，就是经上所说：\Quote{耶路撒冷，你应赞美主，熙雍，你应赞美你的天主；因为祂巩固了你城门的门闩}的城里，当门闩已加固、城门已关闭，如我们不久前所说，没有朋友出去，没有敌人进来，那时我们还会有一卷书来读，或者有像现在向你们解释的道理需要解释吗？因此，现在讲解，是为了在那里可以持守；现在以音节分解，是为了在那里可以完整地默观。天主的圣言在那里不会缺少；但却不是通过文字，不是通过声音，不是通过书卷，不是通过诵读，不是通过讲解。那么如何呢？正如：\BibleQuote{在起初已有圣言}等等。\BibleRef{若 1:1} 因为祂并没有这样来到我们这里，以致从那里离开；因为祂在这世界上，世界是藉着祂造成的。我们要默观的正是这样的圣言。因为\Quote{天主的天主将在熙雍显现。}但这是何时呢？在我们朝圣之后，旅程结束时；不过，如果旅程结束后，我们没有被交给审判官，让审判官把我们送进监狱。但如果当我们旅途结束时，如同我们希望、愿意并努力的那样，我们已抵达了我们的家乡，在那里我们将默观那我们将永远赞美的那一位；那临到我们的那一位不会缺乏，我们享受的那一位也不会缺乏；吃的人不会厌腻，所吃的也不会缺少。那将是伟大而美妙的默观……\par

13. \Quote{野驴将解它们的渴。}他以野驴指某些大的野兽。因为谁不知道野驴被称为\OriginalTerm{野驴}{onagers}呢？因此，他指某些大的、未受驯服的。因为外邦人没有法律的轭：许多民族按自己的习俗生活，在骄傲的夸耀中漫游，如在旷野中。确实，所有野兽都是如此，但野驴被用来表示较大的一类。它们也要为自己的渴而饮，因为水也为它们流。兔子饮，野驴也饮：兔子小，野驴大；兔子胆怯，野驴凶猛：两种都从那里饮，但各按自己的渴……它如此忠实而温柔地流，既满足野驴，又不惊扰兔子。\Person{西塞罗}的声音响起来，\Person{西塞罗}被读，那是一卷书，是他的一个对话，或许是他自己的，或许是\Person{柏拉图}的，或任何这类作家的：一些未受过教育的人、心智较弱的人听到；谁敢渴望这样的事？那是一道水声，也许是浑浊的，但肯定流得如此猛烈，以至胆小的动物不敢靠近去饮。而圣咏响起时，谁会这么说：\Quote{这对我来说太难了？}看，现在圣咏在说什么；当然是隐藏的奥秘，但它如此响起，连孩童也喜欢听，未受教育的人也来饮，喝饱了便放声歌唱……\par

14. 然后圣咏继续其经文：\BibleQuote{天上的飞鸟要在它们之上居住}\BibleRef{咏 103:12}……因此，天空的飞鸟要在山上栖息。我们看到这些鸟住在山上，但许多鸟住在平原，许多在谷中，许多在树林，许多在花园，并非都在山上。有些飞鸟只住在山上。这个名字指某些属灵的灵魂。飞鸟是属灵的心，享受自由的空气。这些鸟喜爱天空的晴朗，但它们的食物在山上，它们要住在那里。你们知道山，已经讲过了。山是先知，山是宗徒，山是一切真理的宣讲者……\par

15. 但不要以为那些\Quote{天上的飞鸟}依从自己的权威；看圣咏说什么：\Quote{从岩石的中间发出它们的声音。}如果我对你们说：相信吧，因为\Person{西塞罗}这样说，\Person{柏拉图}这样说，\Person{毕达哥拉斯}这样说，你们中谁会不嘲笑我呢？因为我将是一只鸟，但不是从岩石发出声音。你们每个人当对我说什么？那如此受教导的人当说什么？\BibleQuote{若有人给你们宣讲福音，与你们所接受的不同，当受诅咒。}\BibleRef{迦 1:9} 你为何对我讲论\Person{柏拉图}、\Person{西塞罗}和\Person{维吉尔}？你们眼前有山岩，从岩石中间给我你们的声音。愿他们被听见，他们从岩石听见；愿他们被听见，因为也在那些许多岩石中，那同一块岩石被听见：因为\Quote{那岩石就是基督。}\BibleRef{格前 10:4} 因此，让他们被乐意地听见，从岩石中间发出他们的声音。没有什么比这样的鸟鸣更甜美。它们鸣叫，岩石回应：它们鸣叫，属灵的人讨论：岩石回应，圣经的见证给予回答。看！于是飞鸟从岩石中间发出声音，因为它们住在山上。\par

16. \BibleQuote{祂从高处灌溉群山}\BibleRef{咏 103:13}。如果有一个未受割礼的外邦人来到我们这里，将要相信基督，我们给他圣洗，并不把他召回法律的这些工作。如果一个犹太人问我们为什么这样做，我们从岩石发声，我们说：这是\Person{伯多禄}所做的，这是\Person{保禄}所做的：从岩石中间我们发出声音。但那个岩石，\Person{伯多禄}自己，那座大山，当他祈祷并看见那个异象时，是从上面被灌溉的……\par

17. \Quote{大地从祢的作为的果实中得饱足。}什么是\Quote{从祢的作为的果实}？愿没有人夸耀自己的作为，而是\Quote{夸耀的，当在主内夸耀。}\BibleRef{格前 1:31} 他得饱足时，是靠祢的恩宠得饱足：不要让他说恩宠是因自己的功德而赐予的。如果它被称为恩宠，是\Quote{白白赐予的}；如果它是因工作而偿还，那就是工资。\BibleRef{罗 4:4} 因此，白白地领受，因为你们这些不虔敬的人得以成义。\par

18. \Quote{使青草为牲畜生长，使绿菜为人的服事生长}\BibleRef{咏 103:14}。这是真实的，我明了；我认出这创造：大地确实为牲畜生出青草，为人的服事生出绿菜。但我也明了这话：\Quote{你不可笼住牛嘴，因为牛在打谷：难道天主关心的是牛吗？所以圣经说这话是为我们的缘故。}\BibleRef{格前 9:9} 那么，大地如何为牲畜生出青草？因为\Quote{主曾规定，传福音的人应靠福音生活。}祂派遣了传道者，对他们说：\Quote{你们吃他们摆上的东西，因为工人应得他的工资。}\BibleRef{路 10:7}……他们给出属灵的，收取属血肉的；给出金子，收取青草……\Quote{我们若给你们撒下了属灵之物，若收割你们的属血肉之物，就算是大事吗？}\BibleRef{格前 9:11} 这是宗徒所说的，一位如此勤劳、不知疲倦、历经考验的传道者，他把这青草给了大地。\Quote{然而，}他说，\Quote{我们没有使用这权力。}他表明这是他应得的，却未收取；他也没有谴责那些收取应得之物的人。因为应受谴责的是索取非应得之物的人，而非接受报酬的人：但他甚至放弃了自己的报酬。你不会因为有人放弃了应得之物就不再欠别人，否则你就不是那灌溉过的、为牲畜生出青草的大地……你接受属灵之物，就应回报属血肉之物：对士兵是应得的，你就归还给士兵；你是基督的军饷官。\Quote{有谁当兵而自备粮饷呢？有谁栽种葡萄园而不吃其中的果子呢？有谁牧养羊群而不吃羊奶呢？我这样说，并不是要这样对我做。}\BibleRef{格前 9:7} 曾有这样一个士兵，他甚至把食物配给让给军饷官：但军饷官仍应支付配给……\par

19. \Quote{那它可从地里生出饼来。}什么饼？基督。从什么地里？从\Person{伯多禄}、从\Person{保禄}、从其他真理的管家。请听这是从地里来的：\Quote{我们，}圣保禄说，\Quote{有这宝贝在瓦器中，为使这卓越的能力出于天主，而非出于我们。}\BibleRef{格后 4:7} 祂是从天上降下来的饼，\BibleRef{若 6:41} 为使祂从地里被生出，当祂借着祂仆人的肉身被宣讲时。大地生出青草，为的是可从地里生出饼来。什么大地生出青草？虔诚圣洁的民族。饼可以从什么地里生出？天主的话从宗徒们，从天主圣事的管家们，他们仍行走在地上，仍带着属地的身体。\par

20. \Quote{并酒悦人心}\BibleRef{咏 103:15}。愿无人预备醉酒；不，愿每人都预备醉酒。\Quote{祢的杯何其美妙，能令人沉醉！}我们不愿说：无人醉酒。要沉醉；但要谨慎，从何源头。若主的佳杯浸透了你，你的沉醉将显于你的工作，显于对正义的圣爱，最后显于你的心神转向，从属地到属天的事物。\Quote{以油润泽他的面容。}……什么是用油润泽面容？天主的恩宠；一种为彰显的照耀；如宗徒所说：\Quote{圣神赐给每个人，是为彰显。}\BibleRef{格前 12:7} 一种恩宠，人能清楚看见在人身上，以达成圣爱，被称为油，因其神圣光辉；既然它在基督身上最完美地显现，整个世界都爱祂；祂虽然在这里曾受讥讽，如今却受万邦敬拜：\Quote{因为王国属于主，祂必在万民中作治理者。}因祂的恩宠如此，许多不信祂的人赞美祂，并声明他们不愿信祂，因为无人能履行祂所命令的。那些曾经辱骂祂的人，因祂的赞美而受阻碍。然而所有人都爱祂，所有人都宣讲祂；因为祂被卓越地傅油，所以祂是\OriginalTerm{基督}{Christ}：因为祂被称为基督，源于祂所受的\OriginalTerm{圣油}{Chrism}或敷油。\OriginalTerm{默西亚}{Messiah}在希伯来语，\OriginalTerm{基督}{Christ}在希腊语，\OriginalTerm{受傅者}{Unctus}在拉丁语：但祂傅油于祂的整个身体。因此所有来到的人都领受恩宠，使他们的面容因油而喜悦。\par

21. \Quote{并饼坚固人的心。}这是什么，弟兄们？祂似乎迫使我们理解祂所说的是何种饼。因为那可见的饼坚固胃，滋养身体，却有另一种饼坚固心，因它是心的饼……因此有一种酒真正悦人心，且除了悦人心外不知做别的事。但为使你不致认为这仅应理解为属灵的酒，而非那属灵的饼；祂已表明这一点，即它也是属灵的：\Quote{并饼，}他说，\Quote{坚固人的心。}所以，要如此理解饼，正如你理解酒；内在的饥饿，内在的饥渴：\Quote{有福的，}我们的主说，\Quote{是那些饥渴慕义的人，因为他们必得饱足。}\BibleRef{玛 5:6} 那饼是义，那酒是义：它是真理，基督是真理。\BibleRef{若 14:6} \Quote{我是，}祂说，\Quote{从天降下的活饼；}\BibleRef{若 6:51} 且\Quote{我是葡萄树，你们是枝条。}\BibleRef{若 15:5}\par

22. \BibleQuote{平原的树木必得饱饫} \BibleRef{咏 103:16}，然而藉此恩宠，从地里生出。\BibleQuote{平原的树木}是指各民族中较低等级的。\BibleQuote{还有祂所种植的黎巴嫩的香柏}。黎巴嫩的香柏，即世上的有权势者，他们自己也将得饱饫。基督的饼、酒和油已到达元老、贵族、君王；平原的树木得饱饫了。首先卑微的人得饱饫；然后黎巴嫩的香柏也得饱饫，但仅是祂所种植的那些：虔诚的香柏，虔敬的信徒；因为祂正是这样种植的。因为不敬虔的人也是黎巴嫩的香柏，因为\Quote{主要折断黎巴嫩的香柏}。因为黎巴嫩是一座山：那里有这些树木，按字面说，是最长寿、最卓越的。但黎巴嫩，正如我们阅读那些论述此事之人的著作时所见，被解释为\Quote{光明}：这光明似乎属于这个世界，它现今以它的浮华发光闪耀。有黎巴嫩的香柏，是主所种植的；那些主所种植的必得饱饫……\par

23. \BibleQuote{在那里麻雀建造巢穴：它们的首领是白骨顶的巢穴} \BibleRef{咏 103:17}。麻雀在哪里建造？在黎巴嫩的香柏中……谁是麻雀？麻雀确实是鸟类，是空中的飞鸟，但小鸟常被称为麻雀。因此有一些属灵的人住在黎巴嫩的香柏中：也就是说，有一些天主的仆人，他们在福音中听到：\BibleQuote{变卖你所有的一切，施舍给穷人，你将有宝藏在天上；然后来跟随我} \BibleRef{玛 19:21}……让那舍弃许多东西的人不要骄傲。我们知道伯多禄是个渔夫：那么他能舍弃什么来跟随我们的主呢？或是他的兄弟安德肋，或是载伯德的儿子若望和雅各伯，他们自己也是渔夫；然而他们说了什么？\BibleQuote{看，我们舍弃了一切，跟随了祢} \BibleRef{玛 19:27}。我们的主没有对他说：你忘了你的贫穷；你舍弃了什么，竟要得着整个世界？我的弟兄们，那舍弃了不仅他所拥有的，而且他所渴望拥有的，舍弃了许多……\par

24. 但虽然麻雀将在黎巴嫩的香柏中建造，\Quote{它们的首领是白骨顶的巢穴}。什么是白骨顶的巢穴？白骨顶，我们都知道，是一种水鸟，栖息在沼泽或海上。它很少或从未在岸边安家；而是在水域中间的地方，因而通常在礁石小岛上，被波浪环绕。因此我们明白，磐石是白骨顶合适的家园，它从来没有比在磐石上更安全。在怎样的磐石上？一块置于海中的磐石。即使被波浪冲击，它却击碎波浪，不被它们打破：这就是海中磐石的卓越之处。多么大的波浪冲击我们的主耶稣基督？犹太人冲撞祂；他们被击碎，祂却完整无损。愿每个效法基督的人，这样居住在这世界，即这海中，他不可避免地感受到风暴和暴风雨，却能不屈服于任何风浪，反而在迎击一切时保持完整。因此，白骨顶的巢穴既是坚固的，又是卑微的。白骨顶不在高处安家；没有比那家更坚固、更卑微的了。麻雀确实因实际需要而在香柏中建造；但它们以那磐石为首领，那磐石被波浪冲击却不破碎；因为它们效法基督的苦难……\par

25. 接着是什么？\BibleQuote{最高的山岭是为公鹿的} \BibleRef{咏 103:18}。公鹿是强壮的、属灵的，在奔跑中越过所有荆棘丛生之处。\Quote{祂使我的脚像母鹿的蹄，使我稳站在高处}。让他们持守崇高的山岭，即天主的崇高诫命；让他们思想崇高的事，让他们持守圣经中最突出的内容，让他们在最高处成义：因为那些最高的山岭是为公鹿的。那么谦卑的兽类呢？野兔呢？刺猬呢？野兔是弱小动物；刺猬也是多刺的：一个是胆小的动物，另一个布满尖刺。这尖刺象征什么，岂不就是罪人？那每日犯罪的人，虽不是大罪，却覆盖着极小的尖刺。在胆怯中他是野兔；在覆盖着细微罪过中他是刺猬：他不能持守那些崇高而完美的诫命。因为\Quote{最高的山岭是为公鹿的}。那么怎样？这些会灭亡吗？不。因为\Quote{磐石是刺猬和野兔的避难所}。因为主是穷人的避难所。把那磐石放在陆地，它是刺猬和野兔的避难所；放在海上，它是白骨顶的家园。无论何处，磐石都是有用的。即使在山上也有用：因为山若没有磐石作根基，就会堕入深渊……\par

26. \BibleQuote{祂按定期定立了月亮} \BibleRef{咏 103:19}。我们以属灵的方式理解教会，从极小开始成长，并因这生命的可朽性而日渐衰老；然而，它却更接近太阳。我不是指肉眼可见的这月亮，而是指以此为名的那个。当教会尚在黑暗之中，当它尚未显现，尚未发光时，人被引入歧途，且有人说：这是教会，基督在这里；以致\BibleQuote{月亮黑暗时，他们向心里正直的人射箭}。而今月亮圆满时，谁若迷路，是何等盲目？\BibleQuote{祂按定期定立了月亮}。因为在此，教会暂时消逝：这死亡的束缚不会永久存留；盈虚终有止期；这是为某些时期所定的。\BibleQuote{太阳也知道自己的沉落}。这太阳是哪个太阳，岂不就是义德的太阳？不敬虔的人将在审判之日哀叹它从未为他们升起；他们在那天将会说：\BibleQuote{因此，我们迷误了真理之道，义德之光未曾照耀我们，太阳也没有升到我们身上。} \BibleRef{智 5:6} 那太阳为理解基督的人升起……\par

27. 弟兄们，也不要认为太阳应当被人崇拜，因为太阳有时在圣经中象征基督。人的疯狂就是这样；好像我们说受造物应受崇拜，因为说过太阳是基督的预象。那么也崇拜磐石吧，因为它也是基督的预象。\BibleRef{格前 10:4} \BibleQuote{祂被牵去如同羔羊被宰杀} \BibleRef{依 53:7}；也崇拜羔羊，因为它是基督的预象。\BibleQuote{犹大支派的狮子已获胜} \BibleRef{默 5:5}；也崇拜狮子，因为它象征基督。请看基督的预象何其多：这一切都是肖似的基督，而非本质的基督……\par

28. 那么，当太阳落下，我们的主受苦时如何？宗徒们有一种黑暗，希望破灭，在那些起初觉得祂伟大、是众人救主的人心中。何以如此？\BibleQuote{祢使黑暗，便成了夜，林中的一切野兽都在其中走动} \BibleRef{咏 103:20}……这里的林兽以不同方式使用：因为这些事总是按不同意义理解；正如我们的主自身有时被称为狮子，有时被称为羔羊。狮子和羔羊有什么不同？但那是怎样的羔羊？能胜过狼、胜过狮子的羔羊。祂是磐石，祂是牧人，祂是门。牧人从门进入；祂说：\BibleQuote{我是善牧}；又说：\BibleQuote{我是羊的门}……要这样学习理解这些比喻式的言语；免得你读到磐石象征基督\BibleRef{格前 10:4}时，以为它处处都指祂。一处有一处的意义，另一处有另一处的意义，正如我们只有看到字母的位置才能理解它的含义。\BibleQuote{狮崽咆哮觅食，向天主寻求它们的食物} \BibleRef{咏 103:21}。所以我们的主，当临近祂的沉落，这位义德的太阳意识到自己的沉落时，对门徒说，好像黑暗将至，狮子要四处巡游寻找可吞吃的人，除非得到许可，那狮子不能吞吃任何人：\BibleQuote{西满，}祂说，\BibleQuote{撒殚想要筛你们像筛麦子一样，但我已为你祈求，使你的信德不至失掉} \BibleRef{路 22:31}。当伯多禄三次否认时\BibleRef{玛 26:70}，他岂不已在狮子的牙齿之间？……\par

29. \BibleQuote{太阳升起，它们便一同聚集，躺卧在自己的洞穴里} \BibleRef{咏 103:22}。随着太阳越升越高，以致基督被普世所认识并在其中受荣耀，狮崽们便一同离去；那些魔鬼，曾在不信之子中运行，煽动人迫害天主的家，如今从对教会的迫害中退去\BibleRef{弗 2:2}。既然没有谁敢迫害教会，\BibleQuote{太阳升起，它们便一同聚集。}它们在哪里？\BibleQuote{它们躺卧在自己的洞穴里。}它们的洞穴是不信者的心。有多少人心中潜伏着狮子？它们不从那里冲出来，不攻击旅居的耶路撒冷。为何如此？因为太阳已经升起，照耀着整个世界。\par

30. 你在做什么，天主的人哪？你，天主的教会？你，基督的身体——头在天上——在做什么？人哪，祂的合一？他说：\BibleQuote{人出去工作} \BibleRef{咏 103:23}。因此，让这人在教会平安的保障中行善工，让他工作到底。因为有时会有一种普遍的黑暗，一种攻击将会发生，但在傍晚，即世界末日之时：但现在教会在平安和安宁中工作；因为\BibleQuote{人出去工作，劳碌直到傍晚。}\par

31. \Quote{主，祢的工程何其浩大！} \BibleRef{咏 103:24}。何其公义伟大，何其崇高！那些如此浩大的工程是在哪里完成的？天主当时站在何处，或坐在何处，来完成这些工程？祂在什么地方如此工作？这些如此美丽的工程最初是从哪里产生的？按字面说，每一受造物都有其秩序，按秩序运行，按秩序美丽，按秩序升起，按秩序落下，按秩序经历所有季节，这是从哪里来的？教会本身是如何兴起、成长、达至圆满的？她将如何以不朽的方式达至终局？她以怎样的宣告被宣讲？以怎样的奥迹被推荐？以怎样的预像被隐藏？以怎样的宣讲被揭示？天主在哪里做了这些事？我看见伟大的工程。\Quote{主，祢的工程何其浩大！} 我问祂在哪里完成了这些事：我找不到地方；但我看到接下来的一句：\Quote{祢以智慧创造了这一切。} 因此，祢在基督内创造了万物。……\Quote{大地充满祢的创造物。} 大地充满基督的创造物。怎么会的呢？我们明白：因为有什么不是圣父藉着圣子造的呢？一切行走和在地上爬行的，一切在水中游的，一切在空中飞的，一切在天上运行的，何况大地，整个宇宙，都是天主的工程。但在我看来，他在这里似乎是在谈论某种新的创造，关于这新的创造，宗徒说：\Quote{谁若在基督内，他就是一个新受造物：旧的已成过去，请看，一切都成了新的。一切都是出于天主。} \BibleRef{格后 5:17} 所有信基督的人，脱去旧人，穿上新人，\BibleRef{厄 4:22} 就是新受造物。\Quote{大地充满祢的工程。} 祂被钉在世上的一处地方，在那小小的地方，那颗种子落在地里，死了；却结出了丰硕的果实……\par

32. \Quote{大地充满祢的创造物。} 祢的什么创造物充满大地？一切树木和灌木，一切动物和牲畜，以及全人类；大地充满天主的创造物。我们看见、知道、阅读、认出、赞美，并在此宣讲祂；然而我们还不能完全按照我们心中在美丽默观后所洋溢的赞美来赞美这些事物。但我们更应当留意那创造物，宗徒说：\Quote{谁若在基督内，他就是一个新受造物：旧的已成过去，请看，一切都成了新的。} \BibleRef{格后 5:17} 什么是\Quote{旧的已成过去}？在外邦人中，是一切偶像崇拜；在犹太人本身，是法律的一切奴役，是那些预表现今祭献的一切祭献。那时人的旧性很多；那一位来更新祂自己的工程，熔化祂的银子，铸造祂的硬币，我们现在看见大地充满信仰天主的基督徒，他们远离过去的污秽和偶像崇拜，从过去的希望转向新时代的希望：看，它尚未实现，但已经在希望中拥有，藉着这希望我们现在歌唱并说：\Quote{大地充满祢的创造物。} 我们尚未在我们的家乡歌唱，也不在那应许的安息中，耶路撒冷城门的门闩还未坚固；但仍在我们的旅途中，注视着整个世界，注视着从四面八方奔向信仰的人们，畏惧地狱，轻视死亡，热爱永生，藐视现世，并因这样的景象满心欢喜，我们说：\Quote{大地充满祢的创造物。}\par

33. ……\Quote{海洋亦广大无边，其中有各种爬行的动物，大小野兽不计其数。} \BibleRef{咏 103:25}。他谈到可怕的海洋。陷阱在这个世界爬行，突然抓住粗心的人；谁能数算爬行的诱惑？它们爬行，但要当心，免得它们把我们抓走。让我们守住那木头；即使在水里，即使在波浪上，我们是安全的：不要让基督睡着，不要让信德睡着；如果祂睡着了，就要唤醒祂；祂将命令风，祂将平静大海；\BibleRef{玛 8:24} 航行将结束，我们将在家乡欢欣。因为我在这个可怕的海洋中仍看见不信的人；因为他们住在贫瘠苦涩的水中：但他们有大有小。我们知道这一点：这个世界的许多小人物仍不信，许多大人物也是如此：在这海洋中有生物，大大小小。他们仇恨教会：基督的名字对他们是个负担；他们不暴怒，因为不被允许；不能爆发为行动的残忍被封闭在心里。因为所有人，无论大小，\Quote{爬行的动物，大小野兽，} 他们现在因神庙被关闭、祭台被推翻、偶像被打破、禁止拜偶像的法律成为死罪而悲伤；所有为此哀悼的人，仍在海中。那么我们呢？我们又该走哪条路回我们的家乡？穿过这同一片海，但靠着木头。不要怕危险；那支撑世界的木头在承载着你。\par

34. \Quote{在那里有船只航行。} \BibleRef{咏 103:26}。看，船只漂浮在使你惊慌的水上，却不沉没。我们理解船只代表教会；它们航行在风暴中，在诱惑的狂风巨浪中，在世间的波涛中，在大小野兽之中。基督在祂十字架的木头上是舵手。\Quote{在那里有船只航行。} 让船只不要害怕，不要过于在意它们漂浮在哪里，而要看由谁在驾驶。\Quote{在那里有船只航行。} 当它们感觉基督是它们的舵手时，什么航程会使它们觉得漫长？它们将安全航行，让它们勤奋航行，它们将到达应许的港口，被引领到安息之地。\par

35. 在那海中也有一些超越一切受造物、无论大小的事物。这是什么？让我们听听圣咏：\Quote{那里有\OriginalTerm{里维雅坦}{Leviathan}，祢所造以戏弄它的。}有无数的爬行之物，有大小走兽；船只将在那里航行，而不畏惧，不仅不畏惧无数的爬行之物和大小走兽，甚至不畏惧那里的蛇；\Quote{祢所}——他向上主说——\Quote{造以戏弄它的。}这是一个伟大的奥秘；然而我即将说出你们已经知道的事。你们知道有一条蛇是教会之敌：你们没有用肉体的眼睛看见它，却用信德的眼睛看见它……\par

36. 那么这条蛇，我们古老的仇敌，怒火中烧，诡计多端，就在这大海之中。\Quote{这里就是\OriginalTerm{里维雅坦}{Leviathan}，祢所造以戏弄它的。}现今你当戏弄这条蛇：因为这条蛇被造正是为此目的。它因自己的罪恶从天上崇高领域坠落，成为魔鬼而非天使，在这广阔的大海中获得了自己的某一区域。你们所认为的它的王国，实则是它的监狱。因为许多人说：为什么魔鬼得到了如此巨大的权能，可以统治这个世界，如此猖獗，能行如此多的事？它能猖獗多少？它能做什么？除非得到许可，它什么也不能做。你当如此行事，使它不得允许攻击你；或者若它被允许试探你，愿它败退而去，不得赢取你。因为它曾被允许试探一些圣人，天主的仆人：他们战胜了它，因为他们没有离开正道，它的脚跟所窥伺的他们，并未跌倒……\par

37. 那么，我的弟兄们，谁若愿意注视蛇的头颅，并安全地渡过这海——因为这蛇必然居住在这里，而且如我开始所说的，魔鬼从天上坠落时便获得了这一区域——就让他在对世界的恐惧和对世界的贪欲方面注视它的头颅。因为正是从这里，它提出了某种恐惧或欲望的对象；它试探你的爱，或你的恐惧。你若害怕地狱，并爱慕天主的国，你必将注视它的头颅……\Quote{没有任何权柄不是出于天主。}\BibleRef{罗 13:1}那么你害怕什么？让龙在水里，让龙在海里：你将要渡过去。它被造如此，为的是被戏弄；它被注定居住在这地方，这区域是给予它的。你认为这居所对它是一件大事，因为你不知道它所坠落的天使住所：在你看来它的荣耀，实则是它的惩罚。\par

38. ……那么你怕什么？或许它要试探你的肉体：那是你主的鞭笞，而非你试探者的权能。它的意愿是要伤害那应许的救恩，但它不得许可：但为确保它不得许可，你要以基督为你的元首：击退蛇的头颅：不要同意它的建议，不要偏离你的道路。\Quote{那里有\OriginalTerm{里维雅坦}{Leviathan}，祢所造以戏弄它的。}\par

39. 你愿意看看它多么无力伤害你，除非被允许吗？\Quote{这一切，}他说，\Quote{都仰望祢，使祢按时赐给它们食物。}\BibleRef{咏 103:27}这条蛇想要吞吃，但它吞吃不了它想要吞吃的人……你已经听到了蛇的食物是什么。你不愿天主把你交给蛇吞吃，因为你不是蛇的食物：即不要离弃天主的圣言。因为当蛇被告知\Quote{你当吃土}时，那是向悖逆者说的：\Quote{你是土，还要归于土。}你不愿成为蛇的食物？不要做土。你回答说：我如何不做土？如果你不贪恋世俗之事。听听宗徒，使你不致成为土。因为你所穿戴的身体是土：但你要拒绝成为土。这是什么意思？\Quote{你们要思念天上的事，不要思念地上的事。}\BibleRef{哥 3:2}如果你不思念世俗之事，你就不是土；如果你不是土，你就不会被蛇吞吃，因为它的指定食物是土。主在祂愿意时，自愿给蛇食物：但祂判断正确，祂不能被欺骗，祂不会把金子当做土给蛇。\Quote{祢一赐给它们，它们便收集。}……\par

40. \Quote{祢一伸开手，它们都饱享美物。}\BibleRef{咏 103:28}是什么，上主，祢伸开祢的手？基督就是祢的手。\Quote{上主的膀臂向谁显示了呢？}\BibleRef{依 53:1}谁得了启示，便为它开启：因为启示就是开启。\Quote{祢一伸开手，它们都饱享美物。}当祢启示祢的基督时，\Quote{它们都饱享美物。}但它们的美物不是来自自己；这常常得到证明。\Quote{祢掩面，它们便惊慌。}\BibleRef{咏 103:29}许多饱享美物的人将所拥有的归给自己，并愿在自己本有的义德中自夸，对自己说：我是义人；我是伟大的：从而变得自满。对于这些人，宗徒说：\Quote{你有什么不是你领受的呢？}\BibleRef{格前 4:7}但是天主，愿意向人证明他所拥有的一切都是来自祂，以便他在美物中也获得谦卑，有时便使他烦恼；祂转脸离开他，他就陷入诱惑；并向他显示，他的义德和他行正道，全是在祂的管理之下……\par

41. 但祢为何如此行？为何掩面，使他们惊惶？\Quote{祢收回他们的气息，他们就气绝。}他们的气息就是他们的骄傲；他们自夸，归功于己，自充义人。因此，求祢掩面，使他们惊惶；收回他们的气息，让他们气绝；让他们向祢呼求：\Quote{上主，求祢速来俯听我，因为我的精神萎靡；求祢不要掩面不顾我。}\Quote{祢收回他们的气息，他们就气绝，归回自己的尘土。}悔罪的人发现自己本身一无能力，便向天主认罪，说自己只是尘土与灰烬。骄傲的人哪，你已归回自己的尘土，你的气息已被收回；你不再自夸，不再显耀，不再自充义人；你看出你本是尘土，而上主转面不顾你时，你便跌回自己的尘土。因此，请祈祷，承认你的尘土与软弱。\par

42. 请看下文：\BibleQuote{祢派遣祢的圣神，他们便被造成}\BibleRef{咏 103:30}。祢收回他们的气息，派遣祢自己的气息；祢收回他们的气息：他们就不再有自己的气息。难道他们被遗弃了吗？\BibleQuote{神贫的人是有福的，}\BibleRef{玛 5:3}但他们并没有被遗弃。他们拒绝拥有自己的气息，他们便领受天主的圣神。这正是吾主对将来的殉道者所说的话：\BibleRef{玛 10:20}\BibleQuote{说话的不是你们，而是你们父的圣神在你们内说话。}不要将你们的勇敢归于自己。祂说，如果这勇敢是你们的，而不是我的，那便是顽固，而非勇敢。\BibleQuote{我们原是祂的化工，}宗徒说，\BibleQuote{受造于善工之中。}\BibleRef{弗 2:10}我们从祂的圣神领受了恩宠，得以活于正义；因为\BibleQuote{祂是那使不虔敬者成义的天主}\BibleRef{罗 4:5}。\BibleQuote{祢收回他们的气息，他们就气绝；祢派遣祢的圣神，他们便被造成；祢便更新了大地的面貌：}\BibleRef{咏 103:30}即藉着新人，他们承认自己是被成义的，并非靠自己的力量为义，以使天主的恩宠存于他们内。那么，当祂收回了我们的气息，我们便又归回自己的尘土，为我们的益处看到自己的软弱，好使我们领受祂的圣神时获得更新。请看下文：\BibleQuote{愿上主的光荣永存}\BibleRef{咏 103:31}。不是你的，不是我的，不是他的，或他的；不是暂时的，而是\Quote{永存}。\BibleQuote{愿上主喜爱自己的化工。}并非喜爱你们的化工，仿佛它们是你们自己的：因为如果你们的作为是恶的，那是由于你们的过犯；如果是善的，那是由于天主的恩宠。\Quote{愿上主喜爱自己的化工。}\par

43. \BibleQuote{祂注视大地，大地便颤抖；祂触摸群山，群山便冒烟}\BibleRef{咏 103:32}。大地啊，你曾因你的美物而欢跃，将你的丰盈与富足归于自己；看，上主注视你，使你颤抖。愿祂注视你，使你颤抖：因为谦卑的颤抖胜过骄傲的自恃。……因为他说：\Quote{是天主在你们内工作。}因此，你们要战战兢兢，因为是天主在你们内工作。因为祂给了，因为你们所有的并非出于你们自己，你们便要怀着恐惧战兢工作，因为你们若不敬畏祂，祂必收回所赐的。因此，要战战兢兢地工作。听另一篇圣咏：\Quote{你们当以敬畏事奉上主，以战兢喜乐地向祂欢呼。}若我们必须以战兢喜乐，天主注视我们，便发生地震；当天主看我们时，让我们内心颤抖；然后天主会安息在那里。听祂在另一处说：\BibleQuote{我的圣神要安息在谁身上？就是那谦卑、安静、对我的话战栗的人。}\BibleRef{依 66:2}\par

\BibleQuote{祂注视大地，大地便颤抖；祂触摸群山，群山便冒烟}\BibleRef{咏 103:32}。群山骄傲，自夸自恃，天主未曾触摸它们；祂一触摸它们，它们便冒烟。群山冒烟是什么意思？就是它们向主祈祷。看，那些大山，骄傲的山，广大的山，不向天主祈祷：它们希望自己被祈求，而不祈求那在它们之上的。因为地上有哪一个有权势、傲慢、骄傲的人，肯谦卑地祈求天主呢？我是指不虔敬的人，而非\Quote{上主栽种的黎巴嫩香柏}。每一个不虔敬的人，不幸的灵魂，不知道如何祈求天主，却希望别人来祈求他。他是一座山；需要天主触摸他，使他冒烟；当他开始冒烟时，他便向天主献上祈祷，如同献上心灵的祭献。他向天主冒烟，于是捶胸顿足；因为烟会引出眼泪。\par

44. \BibleQuote{我一生要向上主歌唱}\BibleRef{咏 103:33}。什么歌唱？凡愿意的。我一生要向上主歌唱。我们现在的生命只是希望；将来的生命才是永恒：必死生命的生命，是对永生的盼望。\Quote{只要我在，我就要赞美我的天主。}既然我在祂内直到永远，只要我在，我就要赞美我的天主。不要以为当我们开始在那状态中赞美天主时，我们会有别的工作；我们的整个生命就是赞美天主。如果我们对所赞美的那位感到厌倦，我们也可能会厌倦赞美。如果祂永远是可爱的，祂就永远被我们赞美。\par

45. \Quote{愿我的言谈中悦于祂：我的喜乐将在于主} \BibleRef{咏 103:34}。人与天主的言谈，除了告明己罪，还能是什么呢？向天主承认你所是的，你便与祂交谈了。向祂言谈，行善工，并言谈。\Quote{你们要洗濯，自洁，}\Person{依撒意亚}说。\BibleRef{依 1:16}。向天主言谈是什么意思？向那位认识你的展开你自己，使祂向不认识祂的你展开祂自己。看，是你们的言谈中悦于主：你们的谦逊的奉献、你们心中的忧苦、你们生命的全燔祭，这中悦于天主。但什么中悦于你自己？\Quote{我的喜乐将在于主。}这就是我所说的在你与天主之间的言谈：向那位认识你的展示你自己，祂便向不认识祂的你展示祂自己。中悦于祂的是你们的告解，甜蜜于你们的是祂的恩宠。祂曾将祂自己向你们说了。如何说的？藉着圣言。什么圣言？\Person{基督}……\par

46. \Quote{愿罪人从地上消灭} \BibleRef{咏 103:35}。祂似乎发怒了！圣善的灵魂啊，你在此歌唱又叹息！但愿我们的灵魂与那灵魂同在！但愿它与之结合、相连、共融！当祂发怒时，它也将看见祂的慈爱。因为除了那充满爱德的，谁能明白这事？你战栗，因为他咒骂。谁在咒骂？一位圣人。无疑，他的话被垂听。但曾对圣人们说：\Quote{祝福，不要咒骂。}\BibleRef{罗 12:14}。那么，\Quote{愿罪人从地上消灭}这句话是什么意思？让他们彻底消灭；让他们的灵被带走，好使祂发出自己的神，他们便得复兴。\Quote{而不敬虔的人，以致他们不复存在。}在什么意义上他们不复存在？只是作为恶人吗？让他们因而成义，好使他们不再是不敬虔的。圣咏作者看见这事，充满喜乐，便重复圣咏的第一节：\Quote{我的灵魂，你要赞美主。}愿我们的灵魂赞美主，弟兄们，因为祂曾屈尊赐给我们理解与语言的能力，并赐给你们聆听的热忱与专注。愿每人能记起他所听的，借着彼此交谈，激发你们所领受的食粮；反刍你们所听的，不要让它沉入你们遗忘的肚腹。愿那可羡慕的宝藏\BibleRef{箴 21:20}停留在你们的唇上。这些事是经过大辛劳才寻获并发现的，也是经大辛劳才宣讲和论述的；愿我们的劳作在你们身上结果实，愿我们的灵魂赞美主。\par

\SourceRecord{Translated by J.E. Tweed.；Nicene and Post-Nicene Fathers, First Series；Vol. 8.；Edited by Philip Schaff.；Buffalo, NY: Christian Literature Publishing Co.,；1888.；<http://www.newadvent.org/fathers/1801104.htm>.}

% structure: p0001 source=h1 target=section reason=page-title

\section{\WorkTitle{圣咏第105篇释义}}

1. 这篇圣咏是首批冠以\Quote{阿肋路亚}一词的圣咏之一；该词（或更准确地说，两个词）的含义是\Quote{赞美主}。因此，他以赞美开始：\BibleQuote{你们应向主宣认，呼求祂的圣名}\BibleRef{咏 104:1}；因为这里的宣认应理解为赞美，正如吾主所说的这些话：\BibleQuote{父啊，天地的主宰，我称谢祢}\BibleRef{玛 11:25}。因为开始赞美之后，通常会呼求天主，祈祷者接着表达其渴望；因此，主祷文本身开头有几句非常简短的赞美，即：\BibleQuote{我们在天的父}\BibleRef{玛 6:9}。随后是祈求的事项。…接着又说：\Quote{告诉百姓祂所行的事；}\BibleRef{若 21:17}或者，按字面从希腊文翻译，正如其他拉丁抄本一样：\Quote{在外邦人中传布祂的福音。}这段话在预言中是对谁说的呢？难道不是对宗徒们吗？\par

2. \BibleQuote{你们应向祂歌唱，弹琴赞美祂}\BibleRef{咏 104:2}。以言行赞美祂：因为我们用声音歌唱，而弹琴是使用乐器，即用我们的手。\BibleQuote{谈论祂一切奇妙的作为。愿光荣归于祂的圣名}\BibleRef{咏 104:3}。这两节经文可以毫不牵强地视为上述两个词的释义；因此，\BibleQuote{谈论祂一切奇妙的作为}可以表达\BibleQuote{你们应向祂歌唱}；而接下来的\BibleQuote{愿光荣归于祂的圣名}则指向\BibleQuote{弹琴赞美祂}；前者关乎我们用以歌唱祂的\Quote{善言}，其中述说祂的奇妙作为；后者关乎善行，以此向祂献上美妙的音乐，使人不因自己能行善而希求受赞美。为此，在说了\BibleQuote{愿光荣归于你们}（那行善的人确实配得）之后，他加上了\BibleQuote{在祂的圣名中}，因为\BibleQuote{夸耀的，应因主而夸耀}\BibleRef{格前 1:31}。…这便是在祂圣名中的赞美。为此，我们在另一篇圣咏中也读到：\BibleQuote{我的灵魂要因主而受赞美；愿谦卑的人听见，并喜乐；}这似乎与本节相呼应，即\BibleQuote{愿寻求主的心喜乐}；因为谦卑的人喜乐，他们不嫉恨那些他们所模仿的、已经行善的人。\par

3. \BibleQuote{你们寻求主，并要坚强}\BibleRef{咏 104:4}。这是从希腊文直译而来，虽然可能不是拉丁文；因此其他抄本作\Quote{你们当被坚固}，另有些作\Quote{你们当被坚强}。…所以，\BibleQuote{你们前来靠近祂，并得光明}这些词适用于看见；而经文中的话则关乎行动：\BibleQuote{你们寻求主，并要坚强}。…但\BibleQuote{你们要时常寻求祂的面容}是什么意思？我确实知道亲近天主于我有益；但如果祂总是被寻求，何时才能找到祂呢？他所说的\Quote{时常}是否指我们在现世生活的全部？由此我们意识到应当如此寻求，因为即使找到了，祂仍需被寻求？的确，信德已经找到了祂，但望德仍寻求祂。然而爱德既通过信德找到了祂，又寻求以看见祂而拥有祂，那时祂将被寻见，足以满足我们，无需再寻求。因为若非信德在今生找到了祂，就不会说\Quote{你们寻求主}。此外，如果通过信德找到了祂，而祂仍需被殷勤寻求，那么就不会说：\BibleQuote{因为我们得救是在于希望；所希望的若已看见，就不是希望了；但若我们盼望那未看见的，就要以坚忍等待它}\BibleRef{罗 8:25}。…而这的确是\BibleQuote{你们要时常寻求祂的面容}这句话的含义，即发现不应终止那以爱德为证的寻求，相反，随着爱德的增长，对已发现者的寻求也应增长。\par

4. 他说：\BibleQuote{你们要记念祂所行的奇妙作为，祂的奇事和祂口中的判断}\BibleRef{咏 104:5}。这段经文似乎类似\BibleQuote{你要对以色列子民说：那\Emph{自有者}打发我到你们这里来}；这句话，即使极小部分，也几乎没有心智能够领会。然后，在提及祂自己的名字时，祂仁慈地将其与祂对人恩宠的眷顾相结合，说：\BibleQuote{我是亚巴郎的天主，依撒格的天主，雅各伯的天主；这是我一向的名字}\BibleRef{出 3:14}。祂借此表明，那些祂宣称自己是其天主的、以伟大爱德的力量知道要时常寻求祂面容的人，都永远与祂同活；祂这样说，即使是小孩也能明白，使他们按能力理解：\Emph{我是自有者}。\par

5. 这是对谁说的？\BibleQuote{祂的仆人亚巴郎的后裔，祂所拣选的雅各伯的子孙}\BibleRef{咏 104:6}。…接着他又说：\BibleQuote{祂是主、我们的天主；祂的判断遍及全世界}\BibleRef{咏 104:7}。难道天主只是犹太人的天主？\BibleRef{罗 3:29}断乎不是！\BibleQuote{祂是主、我们的天主}：因为教会，即祂的判断被传扬的地方，遍及全世界。…\par

6. \Quote{祂永久怀念祂的盟约}\BibleRef{咏 104:8}。其他抄本作\Quote{直到永远}；这是由希腊文的歧义所导致的。但如果我们把\Quote{永久}理解为今世而非永恒，那么，当祂解释祂怀念的是什么盟约时，为何又加上\Quote{祂向千代所说的话}呢？这可以有某种程度的限定；但随后祂又说：\Quote{就是祂与亚巴郎所立的盟约}\BibleRef{咏 104:9}：\Quote{并祂对依撒格所起的誓；又把这定为雅各伯的律例，定为以色列永远的约}\BibleRef{咏 104:10}。但若在此处，由于客纳罕地的缘故，应理解为旧约；因为圣咏的言辞这样说道：\Quote{说：我必将客纳罕地赐给你们，作为你们产业的阄分}\BibleRef{咏 104:11}：那么，它怎能被理解为永远的呢，既然那地上的产业不可能是永远的？正因如此，它被称为旧约，因为它被新约所废除。但千代似乎并不表示任何永恒的事物，因为千代有终结；然而对于这短暂的世界而言，千代又太多了。因为无论一个世代被限定为多少年，在希腊文中被称为\OriginalTerm{世代}{\GreekText{γεν}}。\par

\SourceRecord{Translated by J.E. Tweed.；Nicene and Post-Nicene Fathers, First Series；Vol. 8.；Edited by Philip Schaff.；Buffalo, NY: Christian Literature Publishing Co.,；1888.；<http://www.newadvent.org/fathers/1801105.htm>.}

% structure: p0001 source=h1 target=section reason=page-title

\section{圣咏第106篇释义}

本圣咏同样在开头冠有\Quote{阿肋路亚}的标题，而且出现了两次。但有些人说，一个\Quote{阿肋路亚}属于前一圣咏的结尾，另一个属于本圣咏的开头。他们断言，所有带有这个标题的圣咏，结尾都有\Quote{阿肋路亚}，但并非所有开头都有；因此他们不允许任何结尾没有\Quote{阿肋路亚}的圣咏，在开头有它；认为看似属于开头的，实际上属于前一圣咏的结尾。但除非他们以确凿的证据来说服我们这是真的，否则我们将遵循普遍的习俗，即每当遇到\Quote{阿肋路亚}，就将其归属于发现它在其开头的同一篇圣咏。因为极少有抄本（在我所能查阅的希腊抄本中，我没有发现任何一个）在第\Emph{150}篇的结尾有\Quote{阿肋路亚}；在那之后没有其他属于同一正典的圣咏。但即使是这一点也不能胜过习俗，即使所有抄本都如此。因为可能出于某种对天主的赞美，整部圣咏集——据说由五卷书组成（因为他们说每卷书都在写有\Quote{阿们，阿们}的地方结束）——会在所有歌唱之后以最后这个\Quote{阿肋路亚}结束；而且，鉴于第\Emph{150}篇的结尾，我不认为所有标题为\Quote{阿肋路亚}的圣咏都必须在结尾有\Quote{阿肋路亚}。但是，当一篇圣咏的开头有两个\Quote{阿肋路亚}时，为什么我们的主有时说一次\Quote{阿们}，有时说两次，同样地，\Quote{阿肋路亚}为什么不能有时使用一次，有时两次，我不知道；尤其因为，正如在第\Emph{105}篇中，两个\Quote{阿肋路亚}都放在描述圣咏编号的记号之后，而如果其中一个属于前一圣咏的结尾，它应该放在编号之前；属于这一编号的圣咏的\Quote{阿肋路亚}应该写在编号之后。但也许在这种情况下，出于无知，一种习惯占了上风，并且可能有一些我们至今尚未知晓的理由，因此真理的判断应当引导我们，而不是习俗的偏见。与此同时，在我们完全了解此事之前，每当我们在圣咏编号之后发现写有\Quote{阿肋路亚}，无论是一次还是两次，根据教会最普遍的习俗，我们将把它归属于带有同一编号的圣咏；我们承认，我们既相信圣咏所有标题的奥秘以及这些圣咏顺序的奥秘是重要的，也承认我们尚未能够如愿地深入理解它们。\par

但我发现这两篇圣咏，第\Emph{105}篇和第\Emph{106}篇如此相连，以至于在它们中的第一篇，天主的子民在蒙拣选者的身份中被赞美，这些人没有抱怨，我以为他们就是那些天主所喜悦的人；\BibleRef{格前 10:5}但在接下来的圣咏中，同一子民中那些惹怒天主的人也被提及；尽管天主的慈爱也没有缺少给这些人……因此，这篇圣咏像前一篇一样以\Quote{你们要向主称谢}开头。但在前一篇圣咏中，接下来的话是\Quote{并呼求祂的名：}而在这里，接下来的话是\Quote{因为祂是良善的，祂的慈爱永远长存}\BibleRef{咏 105:1}。因此，在这段经文中，可以理解为罪过的认罪；因为在几节之后，我们读到\Quote{我们与我们的祖先一同犯罪，我们作了恶，行了不义；}但在\Quote{因为祂是良善的，祂的慈爱永远长存}这句话中，主要是对天主的赞美，而在祂的赞美中带有认罪。尽管当任何人认罪时，他应当以赞美天主来认罪；除非没有绝望并呼求天主的慈爱，否则认罪不是虔诚的。因此它确实含有对祂的赞美，无论是在话语上，称祂为良善和仁慈，还是仅仅在情感上，当他相信这一点时……如果这里的慈爱是指那种没有天主无人能幸福的慈爱，我们可以更恰当地将它译作\Quote{永远：}但如果是指那种向不幸者显示的慈爱，使他们或在痛苦中得到安慰，甚至从中解脱，则更恰当地译作\Quote{直到世界的终结，}因为在世界上永远不会缺少不幸的人，可以向他们显示那种慈爱。除非有人胆敢说，即使是那些将要与魔鬼及其使者一同受罚的人，天主的某种慈爱也不会缺少；不是那种使他们能从中解脱的慈爱，而是可能在某种程度上为他们减轻；这样天主的慈爱就可以被称为永恒的，因为它被施行于他们永恒的痛苦上……\par

3. \BibleQuote{谁能述说主的伟大作为？}\BibleRef{咏 105:2}。他满怀着对天主作为的默想，恳求祂的仁慈时说：\BibleQuote{谁，}他说，\BibleQuote{能述说主的伟大作为，或传扬祂的一切赞美？}我们必须补充上述的话，使意思在这里完整，即：\BibleQuote{谁能使祂的一切赞美被听见？}也就是说，谁足以使祂的一切赞美被听见？\BibleQuote{使}它们\BibleQuote{被听见}，他说；也就是说，使它们被听见；表明主的伟大作为和祂的赞美应当如此被讲述，以致可以传报给那些听见的人。但谁能使\BibleQuote{一切}被听见？是否正如接下来的话：\BibleQuote{常守判断、时时行正义的人是有福的}\BibleRef{咏 105:3}；他或许指的是祂的那些赞美，即理解为祂在祂的诫命中的作为？\BibleQuote{因为是天主，}宗徒说，\BibleQuote{在你们内工作}\BibleRef{斐 2:13}……既然祂以不可言说的方式在这些事中工作。\BibleQuote{谁将实行祂的一切赞美被听见？}也就是说，谁听见了它们，就实行祂的一切赞美？这些就是祂诫命的工作。就它们被实行而言，虽然所听见的并非全部被履行，但祂应当被赞美，祂\BibleQuote{在我们内工作，使你们愿意并且为祂的美意而实行}\BibleRef{斐 2:13}。因此，他本可以说\Quote{祂的一切诫命}或\Quote{祂的一切诫命的工作}，但他更喜欢说\Quote{祂的赞美。}……\par

4. 但除非在判断与正义之间有一些区别，我们不会在另一篇圣咏中读到：\Quote{直到正义返回为判断。}圣经确实喜爱将这两个词放在一起；例如：\Quote{正义和判断是祂宝座的居所；}以及这句：\Quote{祂要使你的正义如光明，你的判断如正午；}其中显然是同一思想的重复。或许由于意义的相似，一个可能被用来代替另一个，或以判断代正义，或以正义代判断：然而，如果按其本义来说，我不怀疑有一些区别；即，那正确判断的人被称为守判断，而行为正直的人被称为行正义。而且我认为，\Quote{直到正义返回为判断}这句经文，按此理解并非荒谬：即在这里，那些在信德中守判断、在行为上行正义的人也被称为有福的。……\par

5. 其次，既然天主使人成义，即通过医治他们的不义而使人成为正义的，接着有一个祈祷：\BibleQuote{主，求祢按祢对祢子民所怀的恩惠记念我}\BibleRef{咏 105:4}：也就是说，使我们能成为祢所喜悦的人中间的一员；因为天主并不喜悦所有人。\BibleQuote{求祢以祢的救恩眷顾我。}这就是救主本身，在祂内罪过得赦免，灵魂得医治，使他们能守判断、行正义；既然在这里说话的人知道这样的人是有福的，他们为自己祈祷……\BibleQuote{求祢眷顾我们，}然后，\BibleQuote{以祢的救恩，}也就是说，以祢的基督。\BibleQuote{为见到祢选民的幸福，并在祢子民的喜乐中欢乐}\BibleRef{咏 105:5}：也就是说，因此以祢的救恩眷顾我们，使我们能见到祢选民的幸福，并在祢子民的喜乐中欢乐。因为有些抄本将\Quote{幸福}读作\Quote{甘饴；}正如在前文中，\Quote{因为祂是仁慈的；}别处读作，\Quote{因为祂是甘饴的。}在希腊文中是同一个词，正如别处所读：\Quote{主将显出甘饴；}有些人译为\Quote{幸福，}另一些人译为\Quote{慷慨。}但\Quote{求祢眷顾我们为见到祢选民的幸福}是什么意思？即祢赐给祢所选的这种福乐：只是我们可能不致仍然盲目，如那些被说的人：\BibleQuote{但如今你们说我们看见，所以你们的罪还在。}\BibleRef{若 9:41}因为主赐给盲人视力，不是因他们自己的功劳，而是因祂赐给祂选民的幸福，那正是\Quote{祢选民的幸福}的意思：如同我面容的帮助，不是出于我自己，而是我的天主。我们说到我们的日用粮，是我们的，但我们加上：\BibleQuote{求赐给我们。}\BibleRef{玛 6:11}……\BibleQuote{使祢因祢的产业而受赞美。}我感到惊奇这节经文在许多抄本中被如此解释，因为这三节经文中的希腊文短语是同一个……但既然这似乎是一个有歧义的表达，如果解释者所偏爱的\Quote{使祢受赞美}这一意义是真实的，那么前两节也必须如此理解，因为，如我所说，这三节经文中有同一个希腊文表达；因此整段应如此理解：\BibleQuote{求祢以祢的救恩眷顾我们，使祢见到祢选民的幸福；}也就是说，为这个目的眷顾我们，使祢使我们到达那里，并在那里看见我们；使\BibleQuote{祢在祢子民的喜乐中欢乐，}也就是说，祢可被称为欢乐，因为他们因祢而欢乐；使\BibleQuote{祢因祢的产业而受赞美，}也就是说，与它一同受赞美，因为它除非因祢的缘故，否则不能受赞美。……\par

6. 但让我们听听他们接着承认什么：\newline{}\BibleQuote{我们与我们的祖先一同犯了罪；我们做了错事，行了恶} \BibleRef{咏 105:6}。\newline{}\Quote{与我们的祖先一同}是什么意思？……他说：\Quote{我们的祖先}，\BibleQuote{在埃及没有留意祢的奇事} \BibleRef{咏 105:7}；他还叙述了他们许多其他的罪。或者，\Quote{我们与我们的祖先一同犯了罪}是否应理解为，我们像我们的祖先一样犯了罪，即效法他们的罪？如果是这样，那应该有这种表达方式的例证来支持；我在此刻寻找是否有人说过他与另一个人一同犯了罪或做了什么事，而那个人是在很长时间之后通过类似行为效法了前者，但我没有找到。那么，\Quote{我们的祖先没有明白祢的奇事}是什么意思？无非是他们不知道祢愿意通过这些奇迹使他们确信什么？的确，无非是永生，以及一种非暂时的、而是不变的福乐，这福乐只能通过忍耐来等待。因此，他们不耐烦地抱怨、发怒，并要求以眼前易逝的祝福来蒙福，\BibleQuote{他们也没有记念祢浩大的仁慈}。祂同时责备他们的理解和记忆。需要理解，好让他们默想天主透过这些暂时的恩召所指向的永恒祝福；也需要记忆，好让他们至少不忘记已施行的暂时奇迹，并能忠实地相信，借着他们已经体验过的同一权能，天主会把他们从敌人的迫害中拯救出来；而他们却忘记了祂在埃及借着这样的奇迹为粉碎敌人而给予他们的援助。\BibleQuote{当他们上到海边，甚至到红海时，他们触怒了主}。我们尤其应当注意圣经如何谴责那些没有理解应当理解的事、没有记住应当记住的事；人们不愿将这些归咎于自己的过错，无非是为了少祈祷、少在天主面前谦卑——他们本该在天主面前承认自己是什么样的人，并通过祈求祂的帮助而成为他们所不是的人。因为就连无知和疏忽的罪，与其辩解让它们存留，不如承认它们以便消除；与其通过激怒天主来巩固它们，不如通过呼求天主来清除它们。\par

他接着指出，天主并没有按他们的不信行事。他说：\BibleQuote{然而，祂为祂的名救了他们，为要彰显祂的权能} \BibleRef{咏 105:8}；并非因为他们自己有任何功德。\par

7. \BibleQuote{祂也斥责了红海，海就干了} \BibleRef{咏 105:9}。我们并没有读到有声音从天上发出斥责海；但他把成就此事的天主权能称为斥责：除非有人可以说，海是暗暗受了斥责，好让水能听见，而人却听不见。天主行动的权能是极为深奥奥秘的，祂使甚至无感觉的事物也立刻服从祂的旨意。\BibleQuote{祂就领他们经过深渊，如同经过旷野}。他把大量之水称为深渊。因为有些人为了表达这整节经文的含义，翻译为：\Quote{祂就领他们经过许多水之中}。那么，\Quote{经过深渊，如同经过旷野}是什么意思？无非是那地方因干燥而成了旷野，而那里从前曾是水的深渊？\par

8. \BibleQuote{祂从恨他们的人手中救了他们} \BibleRef{咏 105:10}。有些译者为了避免拉丁文中不常见的表达，用了迂回说法：\Quote{祂从恨他们的人手中救了他们，从仇敌手中救赎了他们}。在这救赎中付出了什么代价？这是否是一个预言，因为这件事是圣洗圣事的预象，在圣洗中我们以极大代价从魔鬼手中被赎，这代价就是基督的血？因此，这预象更确切地不是由任意一个海、而是由红海所表现；因为血是红色的。\par

9. \BibleQuote{至于那些扰乱他们的人，水淹没了他们：他们中没有一人留下} \BibleRef{咏 105:11}；不是所有埃及人，而是那些追赶离开的以色列人、想要抓获或杀死他们的人。\par

10. \BibleQuote{那时他们才信了祂的话，并且向祂献上赞美} \BibleRef{咏 105:12}。这个措辞似乎勉强算是拉丁文，因为他不是说\Quote{相信祂的话}或\Quote{信靠祂的话}，而是\BibleQuote{信了祂的话}；但在圣经中却很常见。\BibleQuote{并且向祂献上赞美}；这样的表达如同我们说：\Quote{他服了这劳役}，\Quote{他过了这样的一生}。这里他指的是那首著名的赞美诗，开头是：\BibleQuote{我要向主歌唱，因祂大大战胜；将马和骑者投入海中}。\par

11. \BibleQuote{他们行事急躁：忘了祂的作为} \BibleRef{咏 105:13}：其他抄本更明白地读作：\Quote{他们急忙忘了祂的作为，不愿等待祂的旨意。}因为他们本应想到，天主向他们所行的如此伟大作为并非没有目的，而是邀请他们到某种无尽的幸福中，而这幸福需要以忍耐来等待；但他们却急于用暂时的事物使自己快乐，而这些事物不能给人真正的幸福，因为它们无法平息无尽的渴望：因为主说：\BibleQuote{凡喝这水的，还要再渴} \BibleRef{若 4:4}。\par

12. 最后，\BibleQuote{他们在旷野贪求贪欲，在干旱之地试探天主} \BibleRef{咏 105:14}。\Quote{干旱之地}，即无水之地，和\Quote{旷野}是相同的；同样，\Quote{他们贪求贪欲}和\Quote{他们试探天主}也是相同的。这种说法与上面的\BibleQuote{他们献上赞美}是一样的。\par

13. \BibleQuote{祂将他们所求的赐给了他们，并使他们灵魂饱足} \BibleRef{咏 105:15}。但祂并没有因此而使他们幸福：因为那不是经上所说的饱足：\BibleQuote{饥渴慕义的人是有福的，因为他们要得饱饫。} \BibleRef{玛 5:6} 在此，他并非指理性灵魂，而是指赋予身体动物生命的灵魂；其本质属于饮食，正如福音所说：\BibleQuote{生命不是贵于食物，身体不是贵于衣服吗？} \BibleRef{玛 6:26} 仿佛灵魂属于吃，身体属于穿。\par

14. \BibleQuote{他们在帐棚中惹了\Person{梅瑟}的怒气，并激怒了主的圣者\Person{亚郎}} \BibleRef{咏 105:16}。他所说的愤怒，或如有人更直译的，是何等的挑衅，下文足以表明。\par

15. 他说：\BibleQuote{地裂开了，吞下了\Person{达堂}，掩没了\Person{阿彼兰}的党羽} \BibleRef{咏 105:17}：\Quote{吞下}对应\Quote{掩没}。\Person{达堂}和\Person{阿彼兰}都同样卷入了最亵渎的分裂。\par

16. \BibleQuote{火在他们中间燃烧起来，火焰烧灭了恶人} \BibleRef{咏 105:18}。在圣经中，这个词通常不用于那些虽然度着公义和可称赞的生活，却并非没有罪的人。相反，如同讥讽之人和讥讽者之间，抱怨之人和抱怨者之间，写作之人和写作者之间等等有区别；同样，圣经惯以\Quote{恶人}指那些非常邪恶、背负沉重罪孽的人。\par

17. \BibleQuote{他们在\Place{曷勒布}造了牛犊，并向铸像下拜} \BibleRef{咏 105:19}。\BibleQuote{于是他们将自己的荣耀，换成了吃草的牛犊的形像} \BibleRef{咏 105:20}。祂不是说\Quote{变为}形像，而是\Quote{在}形像中。这是一种表达方式，正如他所说的\Quote{他们相信了祂的话}。事实上，他极有技巧地不说他们这样做时改变了天主的荣耀；正如宗徒也说：\BibleQuote{他们将不可朽坏的天主的荣耀，变成了可朽坏的人的形像} \BibleRef{罗 1:23}；而是说\Quote{他们的荣耀}。因为天主原是他们的荣耀，只要他们听从祂的劝告，而不急于……\par

18. \BibleQuote{他们忘记了拯救他们的天主} \BibleRef{咏 105:21}。祂是如何拯救他们的？\BibleQuote{祂在\Place{埃及}行了大事，在\Place{含地}行了奇事，在\Place{红海}行了可畏的事} \BibleRef{咏 105:22}。这些奇事也是可畏的；因为没有哪件奇事不带某种恐惧。尽管这些事之所以被称为可畏，是因为它们击溃了他们的仇敌，并向他们显明当怕什么。\par

19. \BibleQuote{于是祂说，祂要消灭他们} \BibleRef{咏 105:23}。因为他们忘记了拯救他们的天主、行奇事者，并制造和敬拜了铸像，因这滔天而难以置信的邪恶，他们当得死亡。\BibleQuote{若非祂的选民\Person{梅瑟}站在祂面前，在破裂之中。}他不是说\Person{梅瑟}站在破裂中，仿佛要打破天主的愤怒，而是站在破裂的道路上，意指那将要打击他们的打击：也就是说，若他没有为他们挺身而出，说：\InnerQuote{现在，如果你赦免他们的罪——不然，求你从你所写的册上涂抹我的名。}由此证明圣人们为他人所作的转求在天主面前何等有力。因为\Person{梅瑟}，无所畏惧于天主的公义（这公义不能涂抹他），恳求仁慈，叫祂不涂抹那些祂以公义本可涂抹的人。所以他\BibleQuote{站在祂面前，在破裂之中，转回祂的烈怒，免得祂消灭他们。}\par

20. \BibleQuote{他们竟轻视了那福乐之地} \BibleRef{咏 105:24}。但他们曾看见它吗？他们怎能轻视那未见过之地呢？除非如以下经文所解释的：\BibleQuote{他们不相信祂的话。}的确，除非那被称为流奶与蜜之地 (\BibleRef{出 3:8}) 意指某种伟大事物，借着它，如同一个可见的记号，祂引导那些领悟祂奇事的人走向不可见的恩宠和天国，否则他们不会因轻视那属地之地而被责备——我们也应将它的暂时王国视为无物，好能爱那自由的耶路撒冷，她是我们的母亲 (\BibleRef{迦 4:26})，她在天上，是真正可向往的。但此处所责斥的毋宁是不信，因为他们不相信天主的话——祂正借着微小事物引领他们走向伟大事物，而他们急于通过现世事物使自己有福，他们以肉性品味这些事物——于是如以上所说，他们\Quote{没有听从祂的劝告}。\par

21. \BibleQuote{他们却在帐棚中抱怨，没有听从主的声音} \BibleRef{咏 105:25}；主曾严厉禁止他们抱怨。\par

22. \BibleQuote{于是祂向他们举起手来，要在\Place{旷野}中消灭他们} \BibleRef{咏 105:26}；\BibleQuote{使他们的后裔分散在异民中，并把他们散布在各地} \BibleRef{咏 105:27}。\par

23. \BibleQuote{他们还归附了\OriginalTerm{巴耳培敖尔}{Baalpeor}；}即是被献祭给外邦偶像；\BibleQuote{并吃了死者的祭品} \BibleRef{咏 105:28}。\BibleQuote{于是他们以自己的发明招惹祂的忿怒，毁灭便在他们中间增多} \BibleRef{咏 105:29}。仿佛祂延迟了举手要在\Place{旷野}中打倒他们、使他们的后裔分散到异民中、将他们散布在各地；正如宗徒所说：\BibleQuote{他们既然不肯在心中保留天主，天主就任凭他们陷于可耻的思念，去行那些不当行的事。} \BibleRef{罗 1:28} 因此，当他们因重罪而受到重罚时，毁灭在他们中间增多。\par

24. \BibleQuote{于是丕乃哈斯站起，平息了祂，震动就停止了}\BibleRef{咏 105:30}。他简要地叙述了全部，因为他不是在此教导无知者，而是提醒那些知道这段历史的人。这里的\InnerQuote{震动}一词与之前的\InnerQuote{破裂}相同，因为在希腊文中是同一个词。最后，他们的邪恶如此之大：奉献给偶像并吃死人的祭品（即因为外邦人向死人献祭如同向天主一样），以致天主若不通过司祭丕乃哈斯的平息，就不会被平息；他杀死了他在奸淫中所发现的一男一女。\BibleRef{户 25:8} 如果他出于对他们的仇恨，而非出于爱，在为主殿的热诚中将其吞没，那就不会算为他的义……我们的主耶稣基督在启示新约时，选择了更温和的训诫；但地狱的威胁更为严厉，而我们在天主在其暂时治理中所给予的那些威胁中，并未读到这一点。\par

25. \BibleQuote{这事就为他算为义，直到万代，永远长存}\BibleRef{咏 105:31}。天主将这算为祂的司祭的义，不仅存留到后裔存在之时，而是\InnerQuote{永远长存}；因为那洞察人心的，知道如何衡量那行为是以多大的爱德为百姓做的。\par

26. \BibleQuote{他们在米黎巴水边惹祂发怒，以致梅瑟因他们的缘故受困苦}\BibleRef{咏 105:32}；\BibleQuote{因为他们惹动了他的心神，他遂用嘴唇说了冒失的话}\BibleRef{咏 105:33}。什么是\InnerQuote{说了冒失的话}？仿佛天主，那位之前行了如此伟大奇迹的，不能使水从磐石中流出。因为他用杖触碰磐石时心存疑虑，从而使这奇迹与其他他未曾疑虑的奇迹区别开来。他如此冒犯，如此应得听到他将死而不得进入应许之地的判决。\BibleRef{申 32:49} 因为他被不信之民的怨言所扰，未能持守他应有的信心。尽管如此，天主仍在他死后，如同对祂所拣选的人一样，给了他美好的见证，使我们看到，这种信德的动摇仅受到这一惩罚：他不被允许进入他正带领百姓前往的那地……\par

27. 但是那些本圣咏所提及罪孽的人，当他们进入了那暂时的应许之地时，\BibleQuote{没有消灭上主所吩咐他们的外邦人}\BibleRef{咏 105:34}；\BibleQuote{反而与外邦人混合，学习他们的作为}\BibleRef{咏 105:35}。\BibleQuote{以致他们敬拜他们的偶像，这为他们成了绊脚石}\BibleRef{咏 105:36}。他们没有消灭他们，反而与他们混合，这成了他们的绊脚石。\par

28. \BibleQuote{他们将自己的儿女献给魔鬼}\BibleRef{咏 105:37}；\BibleQuote{流无辜人的血，就是自己儿女的血，他们献给客纳罕的偶像}\BibleRef{咏 105:38}。那段历史并未记载他们将自己的儿女献给魔鬼和偶像；但本圣咏不会撒谎，先知们也不会，他们在许多谴责的段落中肯定了这一点。此外，外邦人的文献也对他们这种习俗并未保持沉默。但接下来的句子是什么？\Quote{大地被血杀死了}。我们会以为这是抄写者的错误，他写了\OriginalTerm{被杀}{interfecta}代替\OriginalTerm{污染}{infecta}，若不是因为天主的良善，祂愿意祂的圣经以多种语言写成；若不是因为我们看见许多我们检查过的希腊抄本中都写着：\Quote{大地被血杀死了}。那么\Quote{\InnerQuote{大地被杀死}}是什么意思呢？除非这是用比喻的说法指居住在大地上的人……因为他们献上自己儿子的时候，并流那些远未同意此罪的婴儿的血时，他们是在杀害自己的灵魂；因此经上说：\Quote{他们流无辜人的血}。所以\BibleQuote{大地被血杀死，并被他们的行为玷污}\BibleRef{咏 105:39}，因为他们在灵魂上被杀，并被他们的行为玷污；\BibleQuote{他们以自己的\InnerQuote{作为}行淫}。所谓\InnerQuote{作为}，就是希腊人称为\OriginalTerm{习俗}{\GreekText{ἐπιτηδεύματα}}的东西；因为这个词语在希腊抄本中，无论是在这里还是在之前的段落中都出现了，那里说：\BibleQuote{他们以自己的作为惹祂发怒}；在两处，\InnerQuote{作为}都指他们领受自他人的事物。所以，不要让人以为\InnerQuote{作为}是指他们自己设立的、没有先例可模仿的事物。因此，其他拉丁语译者更喜欢用\InnerQuote{追求}、\InnerQuote{情感}、\InnerQuote{模仿}、\InnerQuote{快乐}来代替\InnerQuote{作为}；而且那些在此处写\InnerQuote{作为}的人，在别处写\InnerQuote{追求}。我选择提及此事，免得\InnerQuote{作为}一词被用于他们并非发明、而是模仿他人的事物，会引起困难。\par

29. \BibleQuote{因此上主的怒火向祂的子民燃起}\BibleRef{咏 105:40}。我们的译者不愿意使用\InnerQuote{愤怒}一词来对应希腊文\OriginalTerm{愤慨}{\GreekText{θυμός}}；虽然有些人用了；而另一些人则翻译为\InnerQuote{义愤}或\InnerQuote{心意}。无论采用哪种术语，情感并不影响天主；但惩罚的权能，从习惯上借用了这个名字作比喻。\par

30. \Quote{祂竟厌弃了自己的产业；将他们交在外邦人的手里：恨他们的人就管辖他们} \BibleRef{咏 105:41}：\Quote{他们的仇敌压迫他们，他们就在仇敌的手下屈服} \BibleRef{咏 105:42}。既然诗人称他们为天主的产业，显然祂厌弃了他们，并将他们交在敌人手中，并非为了毁灭他们，而是为了\OriginalTerm{纪律}{discipline}。最后，他说：\Quote{祂多次拯救他们。} \Quote{他们却以自己的谋略触犯祂} \BibleRef{咏 105:43}。这就是上文所说的：\Quote{他们没有遵守祂的谋略。}一个人的谋略对自己有害，当他只寻求自己的事，而不寻求天主的事时。\BibleRef{斐 2:21} 在那产业中，那产业本身对我们而言就是祂自己，当祂屈尊与我们同在，让我们享受祂，与圣人们同在时，我们不会因对任何事物视为己有的爱，而受团体的束缚。因为那最光荣的城，当获得所应许的产业时，其中无人会死，无人会出生，将不会有各自为自己喜乐的公民，因为\Quote{天主将成为万有中的万有。} \BibleRef{格前 15:28} 谁在此旅途中以信德和热诚渴望这个团体，他便会习惯于将公共利益置于私利之上，不寻求自己的利益，只寻求耶稣基督的；免得因在自己的事上聪明和警觉，以自己的谋略触犯天主；反而要盼望那看不见的，不要急于在可见的事物中享福；并耐心等待那看不见的永恒幸福，在祂的应许中遵循祂的谋略，在祈祷中恳求祂的助佑。这样，他在自己的宣认中也会变得谦卑，不致像那些被说成\Quote{他们因自己的邪恶而屈服}的人一样。\par

31. 然而，满有仁慈的天主并没有抛弃他们。\Quote{祂看见他们处于逆境，便听了他们的哀告} \BibleRef{咏 105:44}。\Quote{祂记念了自己的盟约，并照祂丰盛的慈爱后悔了} \BibleRef{咏 105:45}。祂说\Quote{后悔了}，因为祂改变了那似乎要毁灭他们的方式。在天主那里，万事都是安排和确定的；当祂看似出于突然动机而行动时，祂所做的无非是祂从永恒就预知要做的事；但在祂奇妙治理的受造物的暂时变化中，祂自身没有任何暂时的变化，却以突然的意志行动被称为做了祂在最隐秘谋略的不变性中所安排的在事件有定因中的事，按照这谋略，祂按时做一切事，既做现在的事，也早已做了未来的事。谁能领悟这些事呢？\BibleRef{格后 2:16} 因此，让我们听从圣经，它谦卑地讲说崇高的事，为孩童提供滋养的食物，为长者提出研究的题材：那永远的盟约\Quote{祂与亚巴郎所立的}，不是那被废弃的旧约，而是隐藏于旧约中的新约。\Quote{祂怜悯了他们}等等。祂行了所应许的事，但祂早已预知当他们在逆境中祈祷时祂会这样赐予他们；因为他们的祈祷即使尚未发出，但将要发出，无疑已为天主所知。\par

32. 所以\Quote{祂使他们在所有掳掠他们的人眼前蒙怜悯} \BibleRef{咏 105:46}。使他们不再成为\OriginalTerm{义怒之器}{vessels of wrath}，而是\OriginalTerm{怜悯之器}{vessels of mercy}。\BibleRef{罗 9:22} 祂使他们所蒙的怜悯，这里用复数，我想是因为各人从天主所得的恩赐各有不同，这人这样，那人那样。\BibleRef{格前 7:7} 所以，凡读此书者，你若藉着阅读经书和查考先知，认出了天主的恩宠，藉此我们通过我们的主耶稣基督被救赎获得永生，并看见旧约在新约中被启示，新约隐藏在旧约中；请你记住我们主耶稣基督的话，当祂把\OriginalTerm{今世之王}{prince of this world}从信友心中赶出时说：\Quote{现在这世界的元首要被赶出去}；\BibleRef{若 12:31} 以及宗徒的话：\Quote{祂拯救了我们脱离了黑暗的权势，把我们迁移到祂爱子的国度里}。\BibleRef{哥 1:13} 默想这些和类似的事，也查考旧约，看看那篇标题为\Quote{当圣殿在掳掠后重建时}的圣咏所唱的是什么：因为那里说：\Quote{你们应向主唱新歌。} 为使你不以为这仅指犹太人，他说：\Quote{普世大地，请向主歌唱：请向主歌唱，赞美祂的圣名；请宣布，} 更确切地说，\Quote{传报祂的救恩，} 或按希腊文原词说：\Quote{天天\OriginalTerm{传布福音}{evangelize}祂的救恩。} 这里提到了福音（Evangelium），在其中宣布了由日而来的日子，我们的主基督，光由光，圣子由圣父。这也是祂的救恩的意义：因为基督是天主的救恩，正如我们上面所表明的……\par

33. \BibleQuote{主、我们的天主，求祢拯救我们，从列邦中召集我们（其他抄本作\Quote{从外邦人中}），好使我们称谢祢的圣名，并以祢的赞颂为夸耀}\BibleRef{咏 105:47}。随后他简短地加上了这赞颂：\BibleQuote{愿以色列的天主从永远到永远受赞颂}\BibleRef{咏 105:48}；由此我们理解为从永远到永远，因为祂将受到那些人的无终赞颂，关于他们经上说：\Quote{住在祢殿中的，是有福的：他们将不断赞颂祢。}这就是基督的身体在第三日的圆满，那时魔鬼已被驱逐，疾病已得痊愈，直到身体本身的不朽，那些完美赞颂祂的人永远为王，因为他们完美地爱祂；他们完美地爱祂，因为他们面对面看见祂。因为那时这圣咏开端的祈祷将要完成：\Quote{主，求祢记念我们，照祢对祢子民的恩惠}等等。因为祂不仅从外邦人中召集以色列家的亡羊\BibleRef{玛 15:24}，也召集那些不属于那羊栈的，好成为一个羊群，如经上所说，并有一个牧人。但犹太人以为那预言属于他们可见的王国，因为他们不知道如何因未见之善的希望而喜乐，他们将陷入那人的网罗，主论及那人说：\BibleQuote{我因我父的名而来，你们却不接纳我；若另有人因自己的名而来，你们倒要接纳他。}\BibleRef{若 5:43}。宗徒保禄论及那人说：\BibleQuote{那罪人，那丧亡之子必要显露}等等。稍后他说：\BibleQuote{那时那恶者必要显露，主必用祂口中的神气消灭他，以祂来临的光辉摧毁他}等等\BibleRef{得后 2:3}……透过那背道者，透过那高举自己超越一切被称为神或受敬拜者的，依我看来，血肉之躯的以色列民将以为那预言应验了，即经上说的：\BibleQuote{主，求祢拯救我们，从外邦人中召集我们}；并以为在祂的引导下，在他们可见的敌人眼前，那些曾可见地俘虏他们的，他们将要获得可见的荣耀。这样，他们将相信谎言，因为他们没有接受真理的爱，好使他们爱属灵的祝福，而不是属血肉的……因为基督还有其他羊，不属于这羊栈\BibleRef{若 10:16}；但魔鬼和他的使者俘虏了所有的羊，无论在以色列人还是外邦人中。因此，魔鬼的势力从他们中被驱逐，在曾俘虏他们的恶灵眼前，他们在这预言中呼喊，要得救并永远被成全：\BibleQuote{主、我们的天主，求祢拯救我们，从外邦人中召集我们。}不是如犹太人所想像的，通过假基督应验，而是通过我们的主基督，祂以祂父的名而来，\Quote{日复一日，祂的救恩}；关于祂，这里说：\Quote{求祢以祢的救恩眷顾我们！}愿众百姓说，那受割礼和未受割礼的预定子民，圣洁的种族，蒙选的人民，\BibleQuote{阿们！阿们！}\par

\SourceRecord{Translated by J.E. Tweed.；Nicene and Post-Nicene Fathers, First Series；Vol. 8.；Edited by Philip Schaff.；Buffalo, NY: Christian Literature Publishing Co.,；1888.；<http://www.newadvent.org/fathers/1801106.htm>.}

% structure: p0001 source=h1 target=section reason=page-title

\section{\WorkTitle{圣咏第一百零七篇释义}}

1. 这篇圣咏向我们颂扬天主的仁慈，这仁慈在我们的自身经验中得到证明，因此对于有经验的人更为甘饴。除非一个人在自己身上学到了这篇圣咏中所听到的，否则很难使他感到愉悦。然而，这篇圣咏并非为某个人或某两个人所写，而是为天主的子民所写，并公之于众，好让他们如同在镜子中认识自己。它的标题现在无需详述，因为它是\Quote{阿肋路亚}，又是\Quote{阿肋路亚}。我们按教会的古老传统，在特定节日的礼仪中有咏唱它的习惯；我们在特定日子咏唱它，并非没有神圣的含义。我们确实在某些日子咏唱\Quote{阿肋路亚}，但每一天都在心中思念它。因为若此词表示对天主的赞美，纵使不在肉体的口中，也必在心灵的口中。\Quote{祂的赞美常在我口中。}但标题中\Quote{阿肋路亚}不仅出现一次，而是两次，这并非这篇圣咏独有，前一篇亦如此。从上下文看，前一篇歌唱的是以色列子民，而这一篇歌唱的是普世天主的教会，遍布整个世界。或许，它恰当地拥有两次\Quote{阿肋路亚}，因为我们呼喊\Quote{阿爸，父啊}。既然\Quote{阿爸}无非就是\Quote{父}，但宗徒说\BibleQuote{我们藉着祂呼喊：阿爸，父啊}（\BibleRef{罗 8:15}）并非毫无意义。这是因为一面墙确实来到角石前呼喊\Quote{阿爸}，而另一面墙从另一边呼喊\Quote{父}；也就是说，在那角石中，祂\BibleQuote{是我们的和平，祂使两者合而为一。}……\par

2. \BibleQuote{你们应当称谢天主，因为祂是美善的，因为祂的慈爱永远常存}（\BibleRef{咏 106:1}）。你们应当称谢祂是美善的：如果你们尝到了，就应称谢。但若未曾选择品尝，就不能称谢，因为他如何能说自己所不知道的东西是甘甜的？然而，如果你们已尝到\BibleQuote{主是如何甘饴}（\BibleRef{伯前 2:3}），\Quote{你们应当称谢天主，因为祂是美善的。}如果你们热切地尝到了，就应发出称谢之声。\Quote{因为祂的慈爱永远常存}，也就是说，直到永远。因为此处\Quote{永远}如此设定，正如在圣经其他某些地方，\Quote{永远}（希腊文称为\OriginalTerm{永远}{\GreekText{εἰς} \GreekText{αἰῶνα}}）被理解为永久。因为祂的慈爱不是暂时的，以致不能永远；因为祂现今对人所施的慈爱正是为此：使他们能与天使永远生活。\par

3. \BibleQuote{愿上主救赎的百姓这样说，祂从敌人手中救赎了他们}（\BibleRef{咏 106:2}）。的确，以色列子民似乎也从埃及地、从奴役之手、从无益的劳动、从泥泞的工作中被救赎；然而，让我们看看说这些话的人是否就是那些被主从埃及解放的人？并非如此。那么他们是谁？\Quote{就是祂所救赎的人。}但这话也可以指他们，即从敌人手中被救赎的人，例如从埃及人手中。让我们确切地指出他们是谁，因为这篇圣咏是为他们而咏唱的。\Quote{祂从各国聚集了他们}；这仍然可能是指埃及的各地，因为即使在一个省份内也有许多地方。让祂明说吧：\BibleQuote{从东方和西方，从北方和海}（\BibleRef{咏 106:3}）。现在我们就明白这些被救赎的人遍及全球。这个天主的子民，从广阔无边的埃及被释放出来，犹如经过红海（\BibleRef{出 14:22}），要在圣洗圣事中消灭他们的仇敌。因为藉着仿佛红海的圣事，即藉着基督宝血所祝圣的圣洗，追赶的埃及人——即罪恶——被洗去。……\BibleQuote{这些事都发生在他们身上，作为预像，并记载下来，为警戒我们这些末世的人}（\BibleRef{格前 10:11}）。……\par

4. \BibleQuote{他们在旷野中漂泊，在干旱之地，找不到一座可居住的城邑的路}（\BibleRef{咏 106:4}）。我们听到了悲惨的漂泊；那么缺乏什么呢？\BibleQuote{又饥又渴，他们的心灵在他们内衰竭}（\BibleRef{咏 106:5}）。但为何衰竭？为了什么益处？因为天主并不残忍，而是使自己被认识，这对我们是有益的：在我们衰竭时我们向祂呼求，祂帮助我们时我们爱祂。因此，在这漂泊、饥饿和干渴之后，\BibleQuote{他们在急难中呼求了上主，祂拯救他们脱离了困厄}（\BibleRef{咏 106:6}）。祂为他们做了什么，在他们漂泊时？\BibleQuote{祂引领他们走上正路}（\BibleRef{咏 106:7}）。他们找不到一座可居住的城邑的路，被饥渴折磨而衰竭，\BibleQuote{祂引领他们走上正路，使他们可以进入一座可居住的城邑。}祂如何解他们的饥渴，祂没有说，但你们仍可期待：\BibleQuote{愿他们称谢上主，因祂的仁慈和祂向人子所行的奇事}（\BibleRef{咏 106:8}）。你们这些有经验的人，向没有经验的人讲述吧；你们这些已在路上、已受引导去寻找城邑、终于不再饥渴的人。\BibleQuote{因祂使空虚的心灵得到满足，使饥饿的心灵充满美物}（\BibleRef{咏 106:9}）。\par

5. \Quote{他们坐在黑暗中、死影里，困苦中被铁链捆绑} \BibleRef{咏 106:10}。这是从何而来？岂不是由于你归功于自己？岂不是由于你不承认天主的恩宠？岂不是由于你拒绝了天主关于你的旨意 \BibleRef{路 7:30}？请看祂接着说什么：\Quote{因为他们因骄傲违抗了上主的话} \BibleRef{咏 106:11}，不认识天主的正义，而企图建立自己的正义，\BibleRef{罗 10:3} \Quote{他们对至高者的旨意心怀怨恨。}\Quote{他们的心因劳苦而卑微} \BibleRef{咏 106:12}。现在你与情欲争战；若天主不再援助你，你可能努力，却无法得胜。当你被你的邪恶压迫时，你的心会在劳苦中卑微，以致如今以谦卑的心学会呼喊：\Quote{我真可怜啊！谁能救我脱离这个死亡的身体呢？} \BibleRef{罗 7:24}……得释放后，你将承认上主的仁慈。\Quote{他们在困苦中呼求上主，祂就拯救他们脱离他们的患难} \BibleRef{咏 106:13}。他们从第二次诱惑中被释放。还有厌倦和厌恶的诱惑。但先看祂为获释者行了什么：\Quote{祂领他们走出黑暗和死影，折断了他们的锁链} \BibleRef{咏 106:14}。\Quote{愿他们感谢上主的仁慈，以及祂向人子所行的奇事} \BibleRef{咏 106:15}。为什么？祂克服了什么困难？\Quote{因为祂打破了铜门，砍断了铁闩} \BibleRef{咏 106:16}。\Quote{祂把他们从他们的罪恶之路上提起来，因为他们因自己的不义而卑微} \BibleRef{咏 106:17}。因为他们归荣耀给自己，而不是天主；因为他们建立自己的正义，不认识天主的正义，\BibleRef{罗 10:3}他们就卑微了。他们发现若没有祂的帮助，他们毫无能力；他们原是只依靠自己的力量的。\par

6. \Quote{他们的灵魂厌烦一切食物} \BibleRef{咏 106:18}。如今他们遭受饱足。他们因饱足而厌倦。他们因饱足而陷入危险。也许你以为他们可能因饥饿而死，却不会因饱足而死？请看接下来如何。当祂说：\Quote{他们的灵魂厌烦一切食物}，免得你认为他们因饱足而安全，却不看到他们会因饱足而死：\Quote{他们已临近}，祂说，\Quote{直到死亡之门}。那么还有什么呢？即使天主的话令你喜悦，也不可归功于自己；也不可因此骄傲自大，有了对食物的胃口，就傲慢地鄙弃那些因饱足而陷入危险的人。\Quote{他们在患难中呼求上主，祂就拯救他们脱离他们的患难} \BibleRef{咏 106:19}。因为不喜悦是一种病，\Quote{祂派遣祂的话，治愈了他们} \BibleRef{咏 106:20}。请看饱足中有何邪恶；请看祂从何处拯救，那位厌弃自己食物的人向谁呼求。\Quote{祂派遣祂的话，治愈了他们，并把他们拉上来，}从何处？不是从流浪中，不是从饥饿中，不是从克服罪恶的困难中，而是\Quote{从他们的败坏中}。厌弃甘甜之物是一种心智的败坏。因此，对于这个恩惠，如同先前的恩惠一样，\Quote{愿他们感谢上主的仁慈，以及祂向人子所行的奇事} \BibleRef{咏 106:21}。\Quote{并献上赞美之祭} \BibleRef{咏 106:22}。因为如今祂应受赞美，上主是甘甜的，\Quote{愿他们欢欣地宣扬祂的作为}。不是带着厌倦，不是带着悲伤，不是带着焦虑，不是带着厌恶，而是\Quote{欢欣地}。\par

7. ...\Quote{他们乘船下海，在大水上经营事业} \BibleRef{咏 106:23}；那就是在众多民族中间。因为\WorkTitle{若望默示录}可为证，水常代表民族，当若望问那些水是什么时，回答他说：是民族。因此，在大水上经营事业的人，\Quote{他们看见了主的作为，并祂在深渊中的奇事} \BibleRef{咏 106:24}。有什么比人心更深呢？因此常有风刮起；叛乱的暴风和纷争搅扰船只。那么在其中发生什么呢？天主愿意掌舵的和被载运的都向祂呼号，\Quote{祂一发言，暴风的气就停住了} \BibleRef{咏 106:25}。停住了是什么意思？静止了，持续了，仍然搅扰，长久颠簸；咆哮，却不离去。\Quote{因为祂一发言，暴风的气就停住了。}那么那暴风的气做了什么呢？\Quote{他们上达于天}，出于胆大；\Quote{他们下入深渊} \BibleRef{咏 106:26}，出于恐惧。\Quote{他们的灵魂因苦难而消磨。}\Quote{他们摇摇晃晃，东倒西歪，好像醉酒的人} \BibleRef{咏 106:27}。坐在舵旁的人，以及忠心爱护船只的人，能体会我所说的。当然，当他们说话、阅读、解释时，他们显得智慧。暴风真可悲！\Quote{他们的所有智慧，}他说，\Quote{都被吞没了。}有时所有人类的谋略都失败；无论人转向何方，波浪咆哮，风暴肆虐，手臂无力：船首可能撞到哪里，船舷可能暴露在哪个浪头之下，受创的船可能漂向何处，必须避开哪些礁石以免沉没，舵手都无法看见。那么除了随后的事，还有什么呢？\Quote{当他们困苦时，他们呼求了上主，祂就拯救他们脱离他们的灾难} \BibleRef{咏 106:28}。\Quote{祂命令了风暴，它便止息为晴空} \BibleRef{咏 106:29}，\Quote{它的波浪就平静了。}在这一点上，请听一个舵手的声音，他曾经遇险，被降卑，被释放。\Quote{我不愿意，}他说，你们不知道，弟兄们，我们在亚细亚所遭遇的患难：\Quote{我们被压太重，过于力量，甚至超过了限度}（我看他所有的\Quote{智慧尽失}），\Quote{以致我们}他说，\Quote{连活命的指望都绝了。}\BibleRef{格后 1:8}...\par

\Quote{他们因平静而喜乐，祂引领他们进入他们所愿的港口} \BibleRef{咏 106:30}。\Quote{愿祂的慈爱谢恩于上主，并向世人显扬祂的奇事} \BibleRef{咏 106:31}。处处毫无例外，愿我们的功德、我们的力量、我们的智慧，\Quote{谢恩于上主}，而是\Quote{祂的慈爱}。愿祂在我们的每次救援中被爱，祂曾在每次困苦中被呼求。\par

8.\Quote{愿他们在人民的集会中尊崇祂，在长老的座位上赞美祂} \BibleRef{咏 106:32}。愿他们尊崇、赞美，人民和长老、商人和舵手。因为祂在这集会中做了什么？祂建立了什么？祂从哪里拯救了它？祂赐予了它什么？正如祂抵抗骄傲的人，却赐恩宠给谦卑的人：\BibleRef{雅 4:6} 骄傲的人，就是犹太人的首批子民，傲慢无理，因自己是亚巴郎的后裔而自夸，并且因为那民族\Quote{曾受托天主的圣言} \BibleRef{罗 3:2}。这些事并没有使他们受益于健全，反而使心高气傲，宁愿膨胀而非伟大。那么天主做了什么？祂抵抗骄傲的人，却赐恩宠给谦卑的人；砍下自然的枝条因它们的骄傲；接上野橄榄枝因它的谦卑？\par

\Quote{祂使江河变为旷野} \BibleRef{咏 106:33}。那里曾有水流，预言在运行。如今在犹太人中间寻找先知，却找不到。因为\Quote{祂使水泉变为干地。}让他们说：\Quote{如今再也没有先知，祂也不再认识我们了。}\Quote{使沃土变为咸地} \BibleRef{咏 106:34}。你在那里寻找基督的信仰，找不到；寻找先知，找不到；寻找祭献，找不到；寻找圣殿，找不到。这是为什么？\Quote{因住在其中之人的邪恶。}看，祂如何抵抗骄傲的人：听，祂如何赐恩宠给谦卑的人。\Quote{祂使旷野变为水池，使旱地变为水泉} \BibleRef{咏 106:35}。\Quote{祂使饥饿的人住在那里} \BibleRef{咏 106:36}。因为曾向祂说：\Quote{你照默基瑟德的品位，永为司祭。}因为你在犹太人中间寻找祭献，你找不到按亚郎品位的祭献。你照默基瑟德的品位寻找，却不在他们中间找到，但在全世界，它在教会中被庆祝。\Quote{从日出之地到日落之处，上主的名是受赞美的。}…\Quote{他们耕种田地，栽种葡萄园，获得粮食的出产} \BibleRef{咏 106:37}：那工人因此而喜乐，他说：\Quote{我不是贪图馈赠，而是寻求果实。} \BibleRef{斐 4:17} \Quote{祂祝福了他们，他们就极其繁衍，他们的牲畜也没有减少} \BibleRef{咏 106:38}。这稳固。因为\Quote{天主的根基屹立不摇；因为主认识属于祂的人。} \BibleRef{弟后 2:19} 他们被称为\Quote{驮兽}和\Quote{牲畜}，在教会中单纯行走，却是有用的；没有大智，但充满信德。因此，无论是属神的还是属血肉的，\Quote{祂祝福了他们。}\par

9. \BibleQuote{他们变得稀少，并且受苦} \BibleRef{咏 106:39}。这是为何？由于敌对？不然，是由于内里。因为他们所以\BibleQuote{变得稀少}，是因为\BibleQuote{他们从我们中间出去，却不是属于我们的}。\BibleRef{若一 2:19}。但因此他说论到这些，就是他先前所论到的，为使他们能以理解而辨别；因为他说话如同论到同一批人，由于他们所共有的圣事。因为他们属于天主的子民，虽不是因着德能，却确因着虔诚的外表：因为关于他们，我们听到宗徒说：\BibleQuote{在末后的日子，将有艰难的时代，因为人要成为自爱的。}\BibleRef{弟后 3:2}。第一个恶就是\BibleQuote{自爱}；也就是说，以自我为乐。但愿他们不悦于自己，而悦于天主；但愿他们在困苦中呼号，并脱离他们的愁苦。但他们大大依靠自己时，\BibleQuote{他们变得稀少}。这是明显的，弟兄们：所有自合一中分离的，都变得稀少。因为他们原是众多；但在合一中，只要他们不与合一分离。因为当合一的群众开始不再属于他们时，在异端和分裂中，他们就稀少了。\BibleQuote{他们因苦难和悲伤的压迫而受苦。}\BibleQuote{轻蔑倾倒于王侯身上}\BibleRef{咏 106:40}。因为他们被天主的教会所弃绝，而更因为他们希望作王侯，故此被轻视，并成了失了味的盐，被抛在外面，任人践踏。\BibleRef{玛 5:13}。\BibleQuote{他领他们在无路之处迷途，不在道路上。}那些在上方在道路上的，那些被导向一座城，最终被领到那里，没有被领迷途；但这些，在没有路的地方，被领迷途。什么是\BibleQuote{领他们迷途}？天主\BibleQuote{将他们交付于自己内心的情欲}。\BibleRef{罗 1:24}。因为\BibleQuote{领迷途}意味着这个：将他们交付于自己。因为如果你仔细追问，正是他们自己领自己迷途……\BibleQuote{他从乞丐中帮助了穷人}\BibleRef{咏 106:41}。这意味着什么，弟兄们？王侯被轻视，穷人得帮助。骄傲者被摒弃，谦卑者得供应……\BibleQuote{使他成为如羊的家属。}你明白这一个穷人和乞丐，就是关于他所说的\BibleQuote{他帮助穷人脱离苦难}：这个穷人如今是许多家属，这个穷人是许多民族；许多教会是一个教会，一个民族，一个家属，一只羊。这些是伟大的奥秘，伟大的预像，何其深奥，何其充满隐藏的涵义；何其甜美地被发现，因久被隐藏。因此，\BibleQuote{义人将思量这事，并要喜乐；一切邪恶的口都要被堵住}\BibleRef{咏 106:42}。那妄言攻击合一、并迫使真理显明的邪恶，将被定罪，并堵住它的口。\par

10. \BibleQuote{谁有智慧？他将思量这些事，并明瞭上主的仁慈} \BibleRef{咏 106:43}……不是他自己的功绩，不是他自己的能力，不是他自己的权能；而是\BibleQuote{上主的仁慈}；他，当他迷途困乏时，领他回到道路，并喂养了他；他，当他挣扎于罪的困境、被习惯的锁链束缚时，释放并解开了他；他，当他厌恶天主的圣言、几乎因某种倦怠而垂死时，藉着发送他圣言的药物恢复了他；他，当他处于船难和风暴的危险中时，平静了大海，并带他进入港口；他，最后，将他安置在那个百姓中，在那里他赐恩宠给谦卑的人，而不是在那些他拒绝骄傲的人中；并使他成为自己的，使他留在内里得以增多，而不是出去而减少。义人看到这事，并喜乐。\BibleQuote{一切邪恶的口，}因此，\BibleQuote{都要被堵住。}\par

\SourceRecord{Translated by J.E. Tweed.；Nicene and Post-Nicene Fathers, First Series；Vol. 8.；Edited by Philip Schaff.；Buffalo, NY: Christian Literature Publishing Co.,；1888.；<http://www.newadvent.org/fathers/1801107.htm>.}

% structure: p0001 source=h1 target=section reason=page-title

\section{\WorkTitle{圣咏第一百零八篇释义}}

我不认为\Emph{第一〇八}篇圣咏需要释义，因为我已在\Emph{第五十七}篇圣咏和\Emph{第六十}篇圣咏中阐述过它，这篇圣咏是由后者的最后部分构成的。因为\Emph{第五十七}篇的最后部分是这篇的开头，直到\Quote{你的光荣高于整个大地}这句。从此处至结尾，是\Emph{第六十}篇的最后部分；正如\Emph{第一三五}篇的最后部分与\Emph{第一一五}篇相同，从\Quote{外邦人的偶像不过是金银}这句开始；又如\Emph{第十四}篇和\Emph{第五十三}篇，除中间略有改动外，从头到尾完全相同。因此，这篇第一〇八篇圣咏与构成它的那两篇相比，若有细微差异，都容易理解；正如我们在第五十七篇中看到：\Quote{我要歌唱赞美；醒来吧，我的光荣}；而在此篇中是：\Quote{我要用我的光荣歌唱赞美}。在那里说醒来，以便他用光荣歌唱赞美。同样，在那里说：\Quote{你的仁慈伟大}（或如某些译本所译，\InnerQuote{被高举}）\Quote{直达诸天}；但在此篇说：\Quote{你的仁慈伟大，高于诸天}。因为它是伟大直达诸天，为要在诸天中伟大；这正是他愿以\Quote{高于诸天}所表达的意思。同样在第六十篇中说：\Quote{我要欢欣，我将划分舍根}；在此篇说：\Quote{我将被高举，并划分舍根}。这里显示了划分舍根所预表的意义，这预言的实现应发生在主被高举之后，而这份欢欣正是指向那高举；所以他欢欣，因为他被高举。因此他在别处说：\Quote{你将我的哀叹变为舞蹈；你脱去我的麻衣，给我披上喜乐}。同样，那里说：\Quote{厄弗辣因，我头的力量}；但在此篇说：\Quote{厄弗辣因，我头的提升}。但力量来自提升，即他通过提升使人强壮，在我们内结出果实；因为厄弗辣因的意思是\Quote{结果实}。但\Quote{提升}可以理解为：当我们提升基督时，或当基督（他是教会的头）提升我们时。而前一篇圣咏中的\Quote{困扰我们的人}与这篇中的\Quote{我们的仇敌}相同。\par

这篇圣咏教导我们，那些表面似乎与历史相关的标题，按照圣咏集的创作目的，最正确地被理解为先知性的……然而这篇圣咏是由两篇标题不同的圣咏的后部分构成的。这表示每篇都指向一个共同的目标，不是历史的外表，而是预言的深处，两者的目标在这篇中合而为一，其标题是\Quote{达味的诗歌或圣咏}，除\Quote{达味}一词外，与前两个标题都不相似。因为，\Quote{多次并以多种方式}，正如\WorkTitle{希伯来书}所说，\BibleQuote{天主在古时，曾多次并以多种方式，藉着先知对祖先说过话；} \BibleRef{希 1:1}；但祂谈论的是祂后来派遣的那一位，为使先知的话得以应验：因为\BibleQuote{天主的恩许，在祂内都是是。} \BibleRef{格后 1:20}。\par

\SourceRecord{Translated by J.E. Tweed.；Nicene and Post-Nicene Fathers, First Series；Vol. 8.；Edited by Philip Schaff.；Buffalo, NY: Christian Literature Publishing Co.,；1888.；<http://www.newadvent.org/fathers/1801108.htm>.}

% structure: p0001 source=h1 target=section reason=page-title

\section{\WorkTitle{关于圣咏第109篇的讲疏}}

凡忠信诵读\WorkTitle{宗徒大事录}者，都承认这篇圣咏包含关于基督的预言；因为这里所写\Quote{愿他的岁月短少，愿别人取他的职位}明显是预言出卖基督的犹达斯……正如有些话似乎特别适用于宗徒伯多禄，但其含义若不参照他所象征预表的教会（因他在门徒中享有首席权）便不清晰；如经上所载：\BibleQuote{我要将天国的钥匙交给你}（\BibleRef{玛 16:19}）及其他类似经文；同样，犹达斯代表那些与基督为敌的犹太人——他们那时仇恨基督，如今因这种恶行代代相传仍仇恨他。关于这些人、这个民族，不仅我们在此圣咏中明白读到的内容可妥善理解，连那些更明确针对犹达斯本人的话也可如此理解。\par

圣咏随即这样开始：\BibleQuote{天主，求你不要对我的赞美缄默不语，因为邪恶与欺诈的口已经向我张开}（\BibleRef{咏 108:1}）。由此可见，不虔敬者和欺诈者不缄默的指责是虚假的，而天主不缄默的赞美才是真实的。\BibleQuote{因为天主是真实的，人都是撒谎的}（\BibleRef{罗 3:4}）；除非天主在他内说话，没有人是真实的。但至高的赞美是对天主独生子的赞美，其中甚至宣告他是什么：天主的独生子。然而，当他的软弱显露出时，这赞美并未显现，反而隐藏了，那时邪恶与欺诈的口向他张开；正因为他的美德被隐藏，他们的口才张开。他说\Quote{欺诈的口张开}，因为被诡诈掩盖的仇恨爆发成了言语。\par

\BibleQuote{他们用撒谎的舌头攻击我}（\BibleRef{咏 108:2}），最甚的是当他们用谄媚的奉承称赞他为\Quote{善师}。因此另处说：\Quote{那些赞美我的，竟联合起来攻击我。}接着，因为他们喊出\BibleQuote{钉死他，钉死他！}（\BibleRef{若 19:6}），他加上\Quote{他们也以仇恨的话围绕我。}那些用奸诈的舌头说看似爱而非恨的话\Quote{攻击我}的人，因他们诡诈地这样做，后来以并非虚假虚伪的爱，而是公开的\Quote{仇恨的话围绕我，并无故攻击我。}正如因为虔诚者无故爱基督，邪恶者也无故恨他；正如至善之人热切寻求真理本身，不以外在利益为目的，最坏之人也如此寻求邪恶。因此世俗作者论到一个极坏的人时说：\Quote{他无缘无故地邪恶残忍。}\par

他说：\BibleQuote{他们以毁谤取代我的爱}（\BibleRef{咏 108:3}）。此类行为有六种，稍加留意便易于牢记：（1）以善报恶，（2）不以恶报恶，（3）以善报善，（4）以恶报恶，（5）不以善报善，（6）以恶报善。前两种属于善人，其中第一种更优越；后两种属于恶人，其中后者更恶劣；中间两种属于某种中间类型的人，但前者接近善，后者接近恶。我们当留意圣经历中的这些事。主耶稣基督自己以善报恶，他\BibleQuote{使不虔敬者成义}（\BibleRef{罗 4:5}），并在十字架上说：\BibleQuote{父啊，宽赦他们吧，因为他们不知道他们做的是什么}（\BibleRef{路 23:34}）……\par

但在他说了\Quote{他们以毁谤取代我的爱}之后，他加上什么？\Quote{但我专务祈祷。}他并未说明祈祷什么，但我们还能更好理解为是为他们祈祷？因为那些钉死他的人大大毁谤他，嘲笑他，仿佛他是一个被他们自以为征服了的人；从这十字架上他说：\Quote{父啊，宽赦他们吧，因为他们不知道他们做的是什么}；这样，当他们在极深的恶意中以恶报善时，他在至高的良善中以善报恶……因此，天主的话藉主的榜样教导我们：当我们觉得别人对我们忘恩负义，不仅不以善报偿，甚至以恶报善时，我们应当祈祷；主确实为那些向他发怒、或处于忧伤、或在信德中面临危险的人祈祷；但我们首先当为自己祈祷，求天主仁慈和帮助，使我们克服自己的心——当有人当面或背后毁谤我们时，我们容易被报复的欲望所驱使……\par

他接着说：\BibleQuote{他们这样以恶报善}（\BibleRef{咏 108:4}）。仿佛我们问：什么恶？为什么善？他说：\Quote{以仇恨报复我的善意。}这是他们重大罪孽的总结。因为迫害者怎能伤害那自愿而非被迫受死的主？但这仇恨本身就是迫害者最大的罪行，尽管是受难者甘愿的赎罪。他充分解释了上面那句话\Quote{他们以毁谤取代我的爱}的含义，即他们不仅因一般义务而欠爱，更因回报他的爱而欠爱；在此他加上\Quote{以仇恨报复我的善意}。他在福音中提到这爱，说：\BibleQuote{我多少次愿意聚集你的子女，而你却不愿意！}（\BibleRef{玛 23:37}）\par

7. 他于是开始预言这些人因这极大不虔敬所将遭受的；他如此详尽地描述他们的命运，仿佛出于报复之愿而希望其实现。有些人不明白这种以貌似诅咒的方式预言未来的模式，便以为这是以仇恨回报仇恨，以恶念回报恶念；因为事实上，只有少数人能区分：控告者出于满足仇恨而喜悦恶人的惩罚，与法官以公义之心惩罚罪恶，两者方式截然不同。前者是以恶报恶；但法官惩罚时并非以恶报恶，因为他是在对不义者施行公义；而公义之物无疑是善的。因此，他惩罚不是出于对他人的苦难感到快乐（那是恶报恶），而是出于对公义的爱（那是以善报恶）……\par

8. \BibleQuote{愿你派一个不虔敬的人管辖他；并愿撒殚站在他右边} \BibleRef{咏 108:5}。虽然之前抱怨是关于许多人，但现在\BibleBook{}{圣咏}只论及一个人……因此，既然他在这里说的是叛徒犹达斯，他按照\BibleBook{}{宗徒大事录}的记载，将受到应得的惩罚 \BibleRef{宗 1:20}，那么，\BibleQuote{愿你派一个不虔敬的人管辖他}，除了指他在下一节提到的那个名字，即当他说的\BibleQuote{愿撒殚站在他右边}之外，还能指谁呢？因此，那拒绝服从基督的人，应当让魔鬼管辖他，也就是说，他自己应当服从魔鬼……为此，论到那些宁爱世俗享乐而不爱天主的人，他们称那拥有这样那样事物的人民为有福，故有经文说：\BibleQuote{他们的右手是欺骗的右手。} ..\par

9. \BibleQuote{任凭审判时他被定罪，愿他的祈祷变为罪恶} \BibleRef{咏 108:6}。因为祈祷若不藉着基督，就不是正义的；犹达斯犯了滔天大罪出卖了基督。不藉着基督所作的祈祷，不仅不能消除罪恶，反而自身变为罪恶。但可以追问：犹达斯是在什么情况下如此祈祷，使他的祈祷变为罪恶？我猜想，是在他出卖主之前，当他正考虑出卖祂的时候；因为他那时不能再藉着基督祈祷了。因为在他出卖祂并后悔之后，如果他藉着基督祈祷，他会恳求宽恕；如果他恳求宽恕，他就会有希望；如果有希望，他就会盼望怜悯；如果盼望怜悯，他就不会在绝望中自缢了……\par

10. \BibleQuote{愿他的日子稀少} \BibleRef{咏 108:7}。\Quote{他的日子}，指的是他作宗徒的日子，那些日子是短暂的；因为在我主受难之前，因他的罪行和死亡而结束了。仿佛有人问：那么，那最神圣的十二之数将如何？我主并非毫无意义地定意使祂最初的十二宗徒限于这个数目。祂立刻补充说：\BibleQuote{愿别人得他的职务}。这等于说：让他按自己的功过受罚，同时让他的数目得以填满。\par

11. \BibleQuote{愿他的儿女为孤儿，他的妻子为寡妇} \BibleRef{咏 108:8}。他死后，他的儿女成了孤儿，妻子成了寡妇。\BibleQuote{愿他的儿女流离失所，被掳去，乞讨食物} \BibleRef{咏 108:9}。\Quote{流离失所}意指不知何处可去，失去一切援助。\BibleQuote{愿他们从住处被赶出。}此处他解释了他上面所说的\Quote{愿他们被掳去}。\par

12. \BibleQuote{愿强夺者搜尽他的一切财物，愿外族人掠取他的劳碌所得} \BibleRef{咏 108:10}。\BibleQuote{愿无人帮助他} \BibleRef{咏 108:11}：即无人保护他的后裔；因此接着说：\BibleQuote{无人怜悯他的孤儿}。\par

13. 然而，即使孤儿在没有帮助、没有监护人的情况下，仍可在患难与贫困中增长，并通过传宗接代延续种族；他接着说：\BibleQuote{愿他的后裔断绝；在下一代，愿他的名被完全抹去} \BibleRef{咏 108:12}：也就是说，让他所生的不得再生，迅速消逝。\par

14. 但他接着说的是什么呢？\BibleQuote{愿他父亲的罪孽在主面前被记念，愿他母亲的罪恶不被涂抹} \BibleRef{咏 108:13}。这是否应理解为，甚至他父亲的罪恶也要追讨在他身上？然而，对于那些在基督内改变、不再效法恶人行为、不再作恶人之子的人，这些罪并不追讨在他们身上。\BibleRef{则 18:4} ……并且在这些话语：\BibleQuote{我必追讨父亲的罪在儿女身上} \BibleRef{出 20:5} 之后，又加上：\Quote{恨我的}；也就是说，恨我如同他们的父亲恨我一样。所以，效法善人的结果是连自己的罪也被涂去；同样，效法恶人则使人不仅承受自己应得的惩罚，也承受他们所效法之人的惩罚……\par

15. \BibleQuote{愿他们常常反对主} \BibleRef{咏 108:14}。\Quote{反对主}意思是主看在眼里：因为其他译者将这句译为\Quote{愿他们常在主的眼前}；而另一些译者译为\Quote{愿他们常在主面前}；正如别处所说：\BibleQuote{你将我们的过犯摆在你面前}。用\Quote{常常}，他意指这大罪无论在今生还是来世都不得赦免。\BibleQuote{愿他们的纪念从地上消灭}：即指他父亲和母亲的名号。这里的\Quote{纪念}意指由世代相传所保存的东西；他预言这东西将从地上消灭，因为犹达斯本人和他的儿子（即他父母的名号）没有任何后裔，如上所述，在短短一代之内就灭亡了……\par

16. \BibleQuote{他之所以如此，是因为他没有记住施行仁慈}\BibleRef{咏 108:15}；这是指犹达斯，或是指那百姓本身。但 \Quote{没有记住} 更应理解为指百姓：因为如果他们杀害了基督，他们本可怀着悔改记住这事，并仁慈对待祂的肢体，即那些他们百般坚持迫害的人。为此，他说：\BibleQuote{反而迫害了穷人和乞丐}\BibleRef{咏 108:16}。这当然也可以理解为指犹达斯；因为主并不嫌弃成为贫穷的，祂原是富有的，好使我们因祂的贫穷而致富。\BibleRef{格后 8:9} 但我又该如何理解 \Quote{乞丐} 一词呢？除非是因为祂曾对撒玛黎雅妇人说：\BibleQuote{给我水喝}\BibleRef{若 4:7}，并在十字架上说：\BibleQuote{我渴}\BibleRef{若 19:28}。至于随后的话，我看不出如何能理解为指我们的元首本身，即祂身体的救主，就是犹达斯所迫害的那位。因为在说了 \BibleQuote{他迫害了穷人和乞丐} 之后，他又加上 \BibleQuote{和杀戮}，即 \Quote{为要杀害他}，有些人如此译作 \Quote{那心中刺痛者}。这个词通常只用于指在悔改的忧伤中，对过去罪过的刺痛；正如论到那些在主升天后听了宗徒们的话而 \Quote{心中刺痛} 的人所说的，他们正是杀害主的人……\par

17. 圣咏接着说道：\BibleQuote{他喜爱咒骂，愿咒骂临到他身上}\BibleRef{咏 108:17}。虽然犹达斯喜爱咒骂，无论是在偷窃钱囊时，还是在出卖和背叛主时；然而，那百姓更加公开地喜爱咒骂，当他们说：\BibleQuote{祂的血归在我们和我们子孙身上}\BibleRef{玛 32:25}。\BibleQuote{他不喜爱祝福，愿祝福远离他。} 犹达斯确实如此，因为他不爱基督，在祂内有永久的祝福；但犹太人更坚决地拒绝祝福，那位被主开了眼的人对他们说：\BibleQuote{你们也愿意做祂的门徒吗？}\BibleRef{若 9:27}。\BibleQuote{他以咒骂为衣穿上}：指犹达斯，或那百姓。\BibleQuote{它像水一样进入他的内脏。} 因此，既在外，也在内；在外如衣服，在内如水：因为他已来到那\BibleQuote{能在地狱中毁灭身体和灵魂}者\BibleRef{玛 10:28}的审判台前；身体在外，灵魂在内。\BibleQuote{又像油入他的骨髓。} 这表明他怀着快乐行恶，为自己积蓄咒骂，即永罚；因为祝福是永生。目前，恶行是他的快乐，像水流入内脏，像油流入骨髓；但称之为咒骂，因为天主为这样的人预定了痛苦。\par

18. \BibleQuote{愿它成为他的外衣，遮盖他}\BibleRef{咏 108:18}。既然他先前已提到外衣，为何又重复？当他说：\Quote{他以咒骂为衣穿上}时，他\Quote{遮盖}的衣服与他\Quote{穿上}的衣服有所不同吗？每个人都是穿内衣，外罩大衣；除了在人前炫耀不义之外，这还能是什么呢？\Quote{又如腰带}，他说，\Quote{他时刻束着它。} 人束腰带主要是为了便于劳作，以免被衣褶妨碍。因此，凡精心策划恶事、并非一时冲动，并以如此方式学习作恶、以致随时准备下手的人，就是用咒骂束腰。\par

19. \BibleQuote{这是那些在主前诽谤我之人的工作}\BibleRef{咏 108:19}。他没有说\Quote{他们的报应}，而是说\Quote{他们的工作}：显然，他是用衣服、遮盖、水、油和腰带，来描述那些导致永远咒骂的工作本身。因此，所说的\Quote{这是那些在主前诽谤我之人的工作}，指的并非犹达斯一人，而是许多人。虽然复数可能代表单数；正如黑落德去世时，天使说：\BibleQuote{那些谋杀婴孩性命的人已经死了}\BibleRef{玛 2:20}。但谁比那些诽谤主的话本身的人更在主前诽谤基督呢？他们声称，律法和先知们预先宣告的并不是祂。\Quote{也指那些诽谤我灵魂的人：} 他们否认祂在愿意的时候能够复活；尽管祂说：\BibleQuote{我有权舍弃我的生命，也有权再取回它}\BibleRef{若 10:18}。\par

20. \BibleQuote{但是，主天主，求祢与我一同工作} \BibleRef{咏 108:20}。有些人认为应理解为\Quote{仁慈地}，有些人甚至将其添加进去；但最好的抄本是这样写的：\BibleQuote{但是，主天主，求祢为祢圣名的缘故，与我一同工作。}因此，不应忽略一个更深的意义，假设圣子这样向圣父说话：\Quote{求祢与我一同处理}，因为圣父与圣子的事工是相同的。虽然我们理解此处指的是仁慈——因为随后的经文是\BibleQuote{因为祢的仁慈是甘饴的}——因为祂没有说\Quote{在我内}或\Quote{在我之上}或类似的话，而是说\BibleQuote{求祢与我一同工作}；我们正确地理解圣父与圣子一同仁慈地对待仁慈的器皿。\BibleRef{罗 9:23}\Quote{与我一同工作}也可以理解为\Quote{帮助我}。我们在日常生活中用到这个表达，当我们谈论对自己有利的事情时，说\Quote{它对我们有利}。因为圣父帮助圣子，正如天主帮助人，因为祂取了\Quote{奴仆的形体}；对那奴仆的形体，那人也是天主，并且对那\Quote{奴仆的形体}，主也是圣父。因为在\Quote{天主的形体}中，圣子不需要帮助，因为祂与圣父同样全能，因此祂也是人的帮助者……并且因为祂说过\BibleQuote{求祢与我一同工作}后，又加上\BibleQuote{为祢圣名的缘故}，祂称赞了恩宠。因为，在没有任何预先应得的功劳之下，人性被提升到如此高度，以致整个合而为一，圣言与血肉，即天主与人，被称为天主的独生子。这样做是为了使已失落的能被创造它的那位藉着未失落的来寻找；因此接下来的话：\BibleQuote{因为祢的仁慈是甘饴的。}\par

21. \BibleQuote{求祢拯救我，因为我贫困而可怜} \BibleRef{咏 108:21}。贫困与可怜就是那软弱，藉此祂被钉在十字架上。\BibleRef{格后 13:4}\BibleQuote{我的心在我内烦乱。}这指的是当祂的苦难临近时所说的话：\BibleQuote{我的心灵忧闷得要死。} \BibleRef{玛 26:38}\par

22. \BibleQuote{我离去有如斜下的阴影} \BibleRef{咏 108:22}。藉此祂表示了死亡本身。因为如同阴影斜下意味着黑夜来临，同样死亡来自有死的肉身。\BibleQuote{我被驱赶有如蝗虫。}我认为这更适当地理解为主体的成员，即祂忠信的门徒们。为使这一点更加清楚，祂宁愿用复数形式写作\Quote{蝗虫}；尽管在单数形式下也可能理解为多数，正如那段经文：\BibleQuote{祂一发言，蝗虫就来了}；但那样会比较隐晦。门徒们被驱赶，即被迫害者赶散，祂愿意用\Quote{蝗虫}一词来指代他们的人数，或者他们从一个地方到另一个地方的迁移。\par

23. \BibleQuote{我的双膝因禁食而无力} \BibleRef{咏 108:23}。我们读到，我们的主基督曾禁食四十天：\BibleRef{玛 4:2}但禁食对祂有如此大的效力，以致双膝无力吗？或者，这更适当地理解为主体的成员，即祂的圣人们？\BibleQuote{我的肉体因油而改变}；因属神的恩宠。因此基督从希腊词\OriginalTerm{基督}{chrisma}得名，意即傅油。但肉体因油而改变，不是变坏，而是变好，即从死亡的羞辱上升到不朽的荣耀……祂的肉体尚未改变。但无论圣神是藉洗涤或灌溉的观念以水来代表，还是藉喜乐与爱德的火焰以油来代表，圣神本身并无差异，因为祂的预表不同。例如狮子和羔羊有很大差异，但基督二者都代表……\par

24. \BibleQuote{我也成了他们的辱骂} \BibleRef{咏 108:24}：藉着十字架的死。\BibleQuote{因为基督救赎了我们脱离法律的诅咒，自身成了为我们而受的诅咒。} \BibleRef{迦 3:13}\BibleQuote{他们注视着我，并摇头。}因为他们看到了祂被钉十字架，却没有看到祂的复活；他们看到祂的双膝无力，却没有看到祂的肉体改变。\par

25. \BibleQuote{上主我的天主，求祢帮助我；求祢按祢的仁慈拯救我} \BibleRef{咏 108:25}。这可指整个身体，包括头和身体：就头而言，因祂取了奴仆的形体；就身体而言，因着奴仆们自身。因为祂甚至可以在他们内对天主说：\Quote{帮助我}和\Quote{拯救我}，正如在他们内祂对保禄说：\BibleQuote{你为何迫害我？} \BibleRef{宗 9:4}接下来的话\BibleQuote{按祢的仁慈}描述了白白赐予的恩宠，并非根据功绩的回报。\par

26. \BibleQuote{愿他们知道这是祢的手，是祢，上主，作了这事} \BibleRef{咏 108:26}。祂说\Quote{愿他们知道}，是指那些在狂怒中祂甚至为他们祈祷的人；因为后来相信祂的人也在那些向祂摇头嘲笑的人群中。但那些将天主的形状归给人身的人，当懂得天主在什么意义上有一只手。因此我们要懂得，天主的手就是基督：因此别处说：\BibleQuote{上主的膀臂向谁启示了呢？} \BibleRef{依 53:1}……\par

27. \\Quote{他们咒骂，祢却赐福}\\BibleRef{咏108:27}。因此，人的咒骂是虚妄而虚假的，他们喜爱虚妄，寻求谎言；但当天主赐福时，祂必成就祂所说的。\\Quote{起来攻击我的人，愿他们受羞辱。}因为他们以为自己有些能力来攻击我，所以他们才起来攻击我；但是当我被高举于诸天之上，我的荣耀开始遍及全地时，他们必受羞辱。\\Quote{但祢的仆人却要喜乐：}或是在圣父右边，或是在祂的肢体中，他们在试炼中怀有希望而喜乐，并在试炼之后永远喜乐。\par

28. \\Quote{愿我的毁谤者以羞辱为衣}\\BibleRef{咏108:28}：就是说，让他们因毁谤我而感到羞愧。但这也可以理解为一种祝福，即他们因此被改正。\\Quote{愿他们以自己的困惑为衣，如同披上\\OriginalTerm{双倍外衣}{diplois}；}因为\\Quote{diplois}是双倍外衣；就是说，让他们内外都困惑：既在天主前，也在人前。\par

29. \\Quote{至于我，我要用口极力称谢上主}\\BibleRef{咏108:29}……祂被说成\\Quote{在众人中赞美}，是因为祂与祂的教会同在，直到今世的终结（\\BibleRef{玛28:20}），如此我们可以将\\Quote{在众人中}理解为：祂受到这众多人群的尊敬？因为居于中间的那一位，受到最高的尊崇。但如果人心可谓人的中心，那么这句话最好的解释莫过于：我要在许多人心中赞美祂。因为基督因信德住在我们心中（\\BibleRef{弗3:17}）；因此他说\\Quote{用我的口}，即用我身体的口，这身体就是教会。\par

30. \\Quote{因为祂站在穷人的右边}\\BibleRef{咏108:30}。关于犹达斯曾说过：\\Quote{愿撒殚站在他的右边；}因为他选择通过出卖基督来增加财富；但此处主站在穷人的右边，为叫主自己成为穷人的财富。\\Quote{祂站在穷人的右边，}不是为增加终有一死的年岁，也不是为增加他的积蓄，也不是为使他身体强壮，或暂时安全；\\Quote{而是，}他说，\\Quote{为从逼迫者手中拯救我的灵魂。}灵魂得以脱离逼迫者而安全，如果我们不赞同他们作恶；但当我们不赞同他们时，主站在穷人的右边，使他因自己的贫穷，即软弱，而不屈服。这帮助已赐予基督的身体，即所有神圣的殉道者。\par

\SourceRecord{Translated by J.E. Tweed.；Nicene and Post-Nicene Fathers, First Series；Vol. 8.；Edited by Philip Schaff.；Buffalo, NY: Christian Literature Publishing Co.,；1888.；<http://www.newadvent.org/fathers/1801109.htm>.}

% structure: p0001 source=h1 target=section reason=page-title

\section{对圣咏第一百一十篇的讲解}

1. ……此圣咏是那些确凿而公开预言我们的主和救主耶稣基督的许诺之一；因此，我们既是基督徒又相信福音，就全然无法怀疑基督在此圣咏中被宣告。因为当我们的主和救主耶稣基督询问犹太人，他们认为基督是谁的儿子，他们回答说\Quote{达味之子}时，他立刻答复他们说：\Quote{那么，达味在圣神里怎么还称他为主，说：\InnerQuote{主对我主说}等等？}接着他又问：\Quote{如果达味在圣神里称他为主，他怎会是达味的儿子呢？}（\BibleRef{玛 22:42}）本圣咏即以这节经文开始。\par

2. \Quote{主对我主说：你坐在我右边，等我使你的仇敌做你的脚凳}（\BibleRef{咏 109:1}）。因此，我们应当仔细思考主在圣咏开头向犹太人提出的这个问题。因为如果问我们是否承认或否认犹太人所作的回答——愿天主不许我们否认！——若有人问我们：基督是达味之子，或者不是吗？若我们回答\Quote{不是}，我们就与福音相矛盾，因为《玛窦福音》是这样开始的：\Quote{耶稣基督——达味之子的族谱}（\BibleRef{玛 1:1}）。福音圣史宣告，他是在写关于耶稣基督——达味之子的族谱书。所以，当基督问犹太人他们相信基督是谁的儿子时，他们正确地回答：\Quote{达味之子}。福音与他们的回答一致。不仅是犹太人的猜测，连基督徒的信德也如此宣告。……\Quote{如果达味在圣神里称他为主，他怎会是他的儿子呢？}犹太人对此问题默然不语；他们再找不到答复；然而他们没有寻求他为主，因为他们没有承认他自己就是那位达味之子。但是，弟兄们，让我们相信并宣告：因为\BibleQuote{心里相信成义，口里承认，以致救恩}（\BibleRef{罗 10:10}）；让我们相信，我说，并让我们同时宣告他是达味之子，也是达味的主。让我们不要因达味之子而羞愧，免得发现达味之主向我们发怒。\par

3. ……我们知道基督从死者中复活升天后，坐在圣父的右边。这已经成就了：我们没有看见它，但我们相信了它；我们在圣经中读到，听人宣讲，并以信德持守它。正因为基督是达味之子，他也成了达味的主。因为那从达味后裔所生者被如此尊荣，以至于他也是达味的主。你对此感到惊奇，仿佛人世间没有同样的事发生。如果某个普通人的儿子被立为王，他难道不成为他父亲的主吗？更奇妙的事还可能发生：不仅普通人的儿子因成为王而成为他父亲的主，而且平信徒的儿子因成为主教而成为他父亲的神父。同样，基督取了肉身，死于肉身，同这肉身复活，同这肉身升天，坐在圣父的右边——在这同一位被如此尊荣、如此荣光、改变成天上的形体的肉身中，他既是达味之子，也是达味的主。……\par

4. 因此，基督坐在天主的右边，圣子在圣父的右边，对我们隐藏。让我们相信。这里说了两件事：天主说\BibleQuote{你坐在我右边}，又加上\BibleQuote{等我使你的仇敌做你的脚凳}，即在你脚下。你看不见基督坐在圣父的右边；但你能看见这事如何成就：他的仇敌如何成了他的脚凳。当后者公开成就时，你要相信前者秘密地成就了。什么仇敌成了他的脚凳？就是那些空虚妄想的，经上对他们说：\BibleQuote{外邦人为什么激动？万民为什么图谋虚妄的事？}等等。……所以他坐在天主的右边，直到他的仇敌被置于他脚下。这正在进行，正在发生；虽然它逐渐成就，却没有止境地进行着。因为即使外邦人激动，他们联合起来反对基督，难道能阻止这些话应验吗：\BibleQuote{我必将外邦人赐你为产业，将地极赐你为田产}？……\BibleQuote{他们的纪念随着呼喊而消灭}；但\BibleQuote{上主永远常存}，正如另一篇圣咏（但不是另一位圣神）所说的。\par

5. 随后是什么？\BibleQuote{上主必从熙雍伸出你权能的权杖}（\BibleRef{咏 109:2}）。弟兄们，看来非常明显，先知并不是谈论基督在永恒中与圣父一同为王、管辖借他所造之物的那个国度；因为天主圣言（他起初就与天主同在）（\BibleRef{若 1:1}）何时不称王呢？因为经上说：\BibleQuote{愿尊贵和荣耀归于永世的君王，那不朽坏、看不见、独一的天主，直到永远}（\BibleRef{弟前 1:17}）。归给哪位永世的君王？归给那位看不见、不朽坏的。因为在此，当基督与圣父同在、是无形的和不朽坏的，因为是圣言、他的德能、他的智慧，是与天主同在的天主、万物借着他而受造，他就是\Quote{永世的君王}；然而，那藉他肉体的中介召叫我们进入永恒的现世统治，从基督徒开始；但他的统治没有终结。因此，他的仇敌成了他的脚凳，而他坐在圣父的右边，正如经上所写；这事现在正在进行，并将继续直到终结。……\par

6. 因此，当祂从熙雍发出祂权能的权杖时，会发生什么？\BibleQuote{祢甚至要在祢的仇敌中间掌权。}首先，\BibleQuote{祢要在祢的仇敌中掌权：}在狂怒的外邦人中。因为祂要在以后，当圣人领受赏报、恶人接受惩罚时，才在祂的仇敌中掌权吗？那时祂掌权有何稀奇，当义人与祂一同永远为王，恶人受永罚焚烧？那时掌权有何稀奇？如今，\BibleQuote{在祢的仇敌中，}如今在这时代的过渡中，在这人类朽坏生命的延续和继承中，如今当时间的急流滑过时，祢权能的权杖从熙雍发出，\BibleQuote{为叫祢在祢的仇敌中掌权。}祢要掌权，在外邦人、犹太人、异端者、假弟兄中掌权。祢要掌权，啊，达味之子、达味的主，在外邦人、犹太人、异端者、假弟兄中掌权。\BibleQuote{祢要在祢的仇敌中掌权。}如果我们不明白这节经文已经应验，就不能正确理解它……\par

7. \BibleQuote{在祢权能的日子，祢的元始与祢同在。}\BibleRef{咏 109:3}什么是祂权能的日子？何时有元始与祂同在？或什么是元始？在什么意义上元始与祂同在，既然祂就是元始？……\par

8. \BibleQuote{元始与祢同在}是什么意思？假设任何你喜欢的元始。对于基督自己，更应该说是\BibleQuote{祢是元始}，而不是\BibleQuote{元始与祢同在。}因为祂回答那些问祂\BibleQuote{祢是谁？}的人说：\BibleQuote{我从起初就对你们所说的，就是元始。}\BibleRef{若 8:25}因为祂的父也是元始，独生子出自于祂，在这元始内有圣言，因为圣言与天主同在。那么，如果父和子都是元始，就有两个元始吗？绝对不可！因为正如父是天主，子是天主，但父和子不是两个天主，而是一个天主；同样，父是元始，子是元始，但父和子不是两个，而是一个元始。\BibleQuote{元始与祢同在。}那么，这元始与祢同在会在什么意义上显明呢？并非现在元始不与祢同在。因为祢不也说过：\BibleQuote{看，时候要到，你们要各自散归，只留下我独自一人；但我不是独自一人，因为父与我同在}？\BibleRef{若 16:32}因此在这里，元始也与祢同在。因为祢在别处也说过：\BibleQuote{住在我内的父，祂做祂的工程。}\BibleRef{若 14:10}\BibleQuote{元始与祢同在：}父也从未与祢分离。但是当元始与祢同在显明时，那时凡变得像祢的人都要看见祢的本体；\BibleRef{若一 3:2}因为斐理伯在这里看见祢，就寻求了父。\BibleRef{若 14:8}所以，那时将看见现在所相信的：那时\BibleQuote{元始与祢同在}将在义人眼中、在圣人眼中被看见；恶人被移除，使他们不能看见主的光辉……\par

9. 解释你所说的是哪一种权能。因为在此，正如所说，祂的权能被提及，当祂权能的权杖从熙雍发出，使祂在仇敌中掌权。你说的是哪一种权能：\BibleQuote{在圣人的辉煌中}？他说：\BibleQuote{在圣人的辉煌中。}他指的是当圣人身在辉煌中时的权能；而非当他们仍带着属地的血肉，在可朽坏的身体中叹息时……\par

10. 但这被推迟，以后才会赐予：现在有什么？\BibleQuote{在晓星之前，我从胎中生了祢。}这是什么？如果天主有一个儿子，祂也怀了胎吗？不像血肉之体，祂没有；因为祂也没有怀抱；然而经上说：\BibleQuote{那在父怀里的，祂给我们详述了。}\BibleRef{若 1:18}但什么是胎，也是怀抱；怀抱和胎都用来指隐秘之处。\BibleQuote{从胎中}是什么意思？从隐秘、隐藏之处；从我自己、从我本体；这就是\BibleQuote{从胎中}的意思；因为\BibleQuote{谁能述说祂的世代？}\BibleRef{依 53:8}那么，我们理解为父对子说：\BibleQuote{在晓星之前，我从我胎中生了祢。}那么\BibleQuote{在晓星之前}是什么意思？晓星代表星辰，仿佛圣经以部分代表整体，以一颗显眼的星代表所有星辰。但这些星辰是如何被创造的？\BibleQuote{它们可以作记号、定节令、日子、年岁。}\BibleRef{创 1:14}……\BibleQuote{在晓星之前}这个说法，既是比喻性的也是字面性的，并且这样应验了。因为主在夜间从童贞玛利亚的胎中诞生；牧童的证言证实了这一点，他们当时\BibleQuote{露宿守夜看守羊群。}\BibleRef{路 2:7}\Emph{所以达味说：}啊，我主，祢坐在我主右边，祢从哪里是我子，除非因为：\BibleQuote{在晓星之前，我从胎中生了祢}？\par

11. 而祢是为着什么而诞生的呢？\Quote{主发了誓，决不反悔：祢是照默基瑟德的品位永为司铎}\BibleRef{咏 109:4}。因为祢正是为此而生于晨星之前的母腹，好使祢照默基瑟德的品位永为司铎。因为就祂由圣父所生、与天主同在、与生祂者同永的那一性体而言，祂并不是司铎；但因祂所取的血肉，因祂为我们——从我们这里领受的——所要奉献的祭品，祂才是司铎。因此，\Quote{主}发了誓。那么，主发了誓是什么意思呢？主禁止人发誓，\BibleRef{玛 5:34}，祂自己却发誓吗？或者，祂禁止人发誓，主要乃是为此：免得人陷入虚誓，而正因这个原因，主可以发誓，因为祂不可能发虚誓。因为人由于发誓的习惯，可能会滑入虚誓，所以禁止他发誓是合理的：因为他离发誓越远，就离虚誓越远。因为发誓的人可能真的发誓，也可能假的发誓；但不起誓的人根本不可能发虚誓，因为他根本不起誓。那么主为什么不可发誓呢？因为主的誓言就是应许的印证。祂发誓好了。那么，你发誓时做什么？你呼求天主作证：这就是发誓，呼求天主作证；正因如此，必须谨慎，免得你为任何虚假的事呼求天主作证。因此，如果你借着誓言呼求天主作证，那么天主为什么不可以借着誓言呼请祂自己作证呢？\Quote{我指着我的永生起誓，上主说}，这就是主的誓言……\Quote{主发了誓}，也就是说，祂确认了：\Quote{祂决不反悔}，祂不会改变。什么？\Quote{祢是永为司铎}。\Quote{永为}，因为祂决不反悔。但司铎，是什么意思？难道还要有那些祭品、圣祖们奉献的牺牲、血祭的祭台、帐幕，以及旧约的那些神圣象征吗？断乎不可！这些都已经废除了：圣殿被毁，那个司铎职已撤去，他们的牺牲和祭献也一同消失了，连犹太人也已不拥有这些。他们看见亚郎的司铎职已经消亡，却不认识默基瑟德品位的司铎职。我是对信友们说话。若慕道者有什么不明白的，让他们放下懈怠，赶紧追求知识。因此我不必在这里揭示奥秘：让圣经向你们暗示什么是照默基瑟德品位的司铎职。\par

12. \Quote{主在祢的右边}\BibleRef{咏 109:5}。主曾说：\Quote{祢坐在我的右边}；如今主却在祂的右边，仿佛他们互换了座位……那同一个基督，\Quote{主在祢的右边}，祢曾发誓，且决不反悔：祂永远作司铎做什么？祂在天主右边，为我们转求，\BibleRef{罗 8:34}，像一位司铎进入内层，进入至圣所，进入天上的奥秘，唯独祂没有罪，因此易于洁净罪恶。因此祂\Quote{在祢的右边，在祂震怒的日子，必要击伤列王}。你问，什么列王？你忘了吗？\Quote{地上的列王起来，首领们一同商议，反对主，反对祂的受傅者。}这些王被祂的荣耀击伤，祂名号的重量使王衰弱，以致他们不能实现他们的意愿。因为他们竭力要从地上消灭基督徒的名号，却不能；因为\Quote{凡跌在这石头上的，必被粉碎}\BibleRef{玛 21:44}。因此列王跌在这\Quote{绊脚石}上，因而受伤，当他们说：基督是谁？我不知道祂是个犹太人还是个加里利人，祂死了，就这样被杀！那石头就在你的脚前，可以说，卑微地躺着；因此你因藐视而绊跌，因绊跌而跌倒，因跌倒而受伤……\Quote{但那石头落在谁身上，必要把他砸得粉碎}\BibleRef{路 20:18}。因此，当有人跌在上面时，它好像躺在地上；那时它击伤；但当它要把他砸得粉碎时，它就要从上面下来。请注意在这两个词中：\Quote{击伤}和\Quote{砸碎}：他撞在上面，它却要落在他身上；这里区分了两种时期：基督的卑微与威严，隐藏的惩罚与将来的审判。祂来临的时候，不会压碎那个当祂以可鄙的外表躺着时没有被祂击伤的人……\par

13. \Quote{祂要在列邦中施行审判；祂要使那倒塌的充满}\BibleRef{咏 109:6}。不论你是谁，如果你顽固地反对基督，你已建起了一座必要倒塌的高塔。你最好自己谦卑下来，变得卑微，把自己抛在那坐在圣父右边的祂脚下，好使你在你内成一个废墟，以便被重建。因为如果你坚持在你的邪恶高位上，你将被推倒，那时你再不能重建。因为圣经在另一处论到这样的人说：\Quote{因此祂要拆毁，而不重建。}显然，祂不会说\Emph{这个}关于某些人，除非有一些人，祂拆毁他们是为再重新建造。这在目前正进行着，当基督这样在列邦中施行审判，以致充满那倒塌的。\Quote{祂要在世上击打许多头颅。}在这里，在世上，在此生中，祂要击打许多头颅。祂使他们谦卑，而非骄傲；我敢说，弟兄们，宁可谦卑地、带着受伤的头颅行走在这里，也比昂着头落入永死的审判更为有益。祂要使许多人跌倒，从而击打他们的头颅，但祂要充满他们，再建造他们。\par

14. \BibleQuote{祂要在路旁喝溪水，因此必要抬起头来} \BibleRef{咏 109:7}。让我们思考祂在途中的溪水：首先，什么是溪水？人类必死性的滚滚洪流：因为溪水由雨水汇集而成，泛滥，咆哮，奔流，而奔流就是流淌，即走完它的行程；整个人类必死性也是如此。人们出生，活着，死去，当一些人死去，另一些人出生，当他们死去，另一些人出生，他们相继而来，聚集在一起，离开而不停留。这里有什么是持守的？有什么是不流动的？有什么不是在奔向深渊，仿佛是从雨水汇集而成？因为正如一条河流突然从雨水、从阵雨的水滴汇集起来，流入大海，不再看见，而在从雨水汇集之前也未曾看见；如此这隐藏的雨水从隐藏的源头汇集，流淌；在死亡时，又行走到隐藏之处：这中间状态喧闹并消逝。祂饮了这溪水，祂没有不屑于饮这溪水；因为饮这溪水对祂而言就是出生和死亡。这溪水所拥有的是出生和死亡；基督接受了这一切，祂出生，祂死去。\BibleQuote{因此祂必要抬起头来；}也就是说，因为祂谦卑，\BibleQuote{听命至死，且死在十字架上；为此，天主极其举扬祂，并赐给祂一个名字，超越一切名字，好使因耶稣之名，一切在天上的、地上的和地底下的，都屈膝叩拜；并且一切唇舌都宣认耶稣基督是主，为光荣天主圣父。} \BibleRef{斐 2:8}。\par

\SourceRecord{Translated by J.E. Tweed.；Nicene and Post-Nicene Fathers, First Series；Vol. 8.；Edited by Philip Schaff.；Buffalo, NY: Christian Literature Publishing Co.,；1888.；<http://www.newadvent.org/fathers/1801110.htm>.}

% structure: p0001 source=h1 target=section reason=page-title

\section{圣咏第111篇释义}

1. 我们歌唱\Quote{阿肋路亚}的日子已经来到……但这些日子来了又去，去了又来，预示那不会来去的一天，因为它既不被昨日所迫而来，也不被明日所催而去……因为正如这些日子在规律性的季节中接续，带着喜乐，替代了四旬期的日子，那四旬期象征主身体复活前今生的苦难；同样，在复活之后，那将要赐给主的整个身体，即神圣的教会的一天，当今生一切烦恼和忧伤都被摒弃时，将以永福接续。但这今生要求我们克己，使我们虽然呻吟、被劳苦和挣扎所压，渴望穿上那来自天上的住所（\BibleRef{格后 5:2}），却能戒绝世俗的欢愉；这由四十这个数字象征，那是梅瑟、厄里亚和主自己守斋的时期……但在主复活后的五十天期间，我们歌唱\Quote{阿肋路亚}，所象征的不是某个时期的界限和消逝，而是那蒙福的永恒；因为五十是四十加上\OriginalTerm{十}{denary}，象征着今世劳苦的信徒所得的赏报，那是我们的圣父为最先的和最后的预备了同样的一份（\BibleRef{玛 20:10}）。因此，让我们倾听充满神圣赞美的天主之民的心声。他在这篇圣咏中代表一个在幸福喜乐中欢腾的人，他预表了那些心中充满天主之爱的子民，即基督的身体，从一切邪恶中得释放。\par

2. \Quote{我要向祢、主，宣认}他说，\Quote{用我全心}（\BibleRef{咏 110:1}）。宣认不总是认罪，而是将天主的赞美倾注于宣认的热诚中。前者哀恸，后者喜乐；前者向医师显示伤口，后者为康复而感恩。后一种宣认代表某人，不仅从一切邪恶中得释放，甚至与所有心怀恶意者分离。为此，让我们思考他在何处以全心向主宣认。他说：\Quote{在正直人的议会和集会中}：我想，是那些将要\Quote{坐在十二个宝座上，审判以色列十二支派}（\BibleRef{玛 19:28}）的人。因为在他们中不再有不义的人，犹达斯的盗窃不被允许，没有西满·马古斯受洗后想买圣神却打算出卖；没有铜匠亚历山大行许多恶事（\BibleRef{弟后 4:14}），没有人披着羊皮以伪装的手足之情潜入；教会如今必须在这类人中呻吟，并且将来必须将他们排除在外，那时所有义人都聚集在一起。\par

\Quote{上主的作为伟大，凡喜爱祂的必究察}（\BibleRef{咏 110:2}）：这些作为使仁慈不遗弃任何宣认者，没有人的邪恶不受惩罚（\BibleRef{希 12:6}）。……让随己意选择吧；上主的作为并非如此安排，以致受造物虽有自由裁量权，却能在违背造物主旨意时超越祂的旨意。天主不愿你犯罪，因为祂禁止之；但若你犯了罪，不要以为人做了他所愿的，而天主遇到了祂所不愿的。因为祂愿人不犯罪，同样祂愿宽恕罪人，使他回头并生活；祂最终愿惩罚那顽固不化的罪人，使悖逆者无法逃脱正义的权能。因此，无论你作何选择，全能者必不会不知所措地完成祂关于你的旨意。\par

3. \Quote{宣认与荣美之工是祂的作为}（\BibleRef{咏 110:3}）。还有什么比使不敬虔者成义更荣耀的事呢？但或许人的行为阻碍了天主那荣耀的作为，以致当他认了罪，才配得成义。……这是上主的荣耀之工：因为那得赦免多的，爱得也多（\BibleRef{路 7:42}）。这是上主的荣耀之工：因为\Quote{罪恶在那里增多，恩宠在那里也格外丰富}（\BibleRef{罗 5:20}）。但或许人会以为因行为而配得成义。他说：\Quote{不是出于行为，免得有人自夸。我们是祂的化工，在基督耶稣内受造，为行善事}（\BibleRef{弗 2:9}）。因为人若不先成义，就不能行义；但借着\Quote{信那称不敬虔者为义的天主}（\BibleRef{罗 4:5}），他以信德开始；好使善行不先出现表明他该得什么，而是随后表明他领受了什么……\par

4. \Quote{祂使自己的奇事成为可纪念的}（\BibleRef{咏 110:4}）：藉着贬抑这人，高举那人。祂在合宜的季节保留不寻常的奇迹，好使专注于新奇的人性软弱得以记住它们，尽管祂每日的奇迹更为伟大。祂在全地创造了如此多的树木，无人惊奇；祂用一句话使一棵枯干，凡人的心却受到雷击（\BibleRef{玛 21:19}）。因为那奇迹，由于不因频繁而变得平常，将最牢固地紧贴心灵。但奇迹有何用，除非祂被敬畏？敬畏又有何益，除非\Quote{上主是宽仁慈悲的}，赐\Quote{食物给敬畏祂的人}（\BibleRef{咏 110:5}）？那是不朽坏的食物，\Quote{从天上降下来的食粮}，祂并未因我们的任何功德而赐予。因为\Quote{基督为不敬虔者死了}（\BibleRef{罗 5:6}）。那么，除了宽仁慈悲的主，谁会赐予这样的食物呢？但若祂在此生赐予如此之多，若那将要成义的罪人领受了成为血肉的圣言；那么当他在来世得荣耀时，将领受什么呢？因为\Quote{祂永远纪念自己的盟约}。那给了抵押的，并未给全部。\par

5. \BibleQuote{祂要向祂的子民显示祂大能的作为} \BibleRef{咏 110:6}。但愿那些舍弃一切所有并跟随祂的圣洁以色列子民不要忧愁，不要悲伤说：\Quote{这样谁能得救呢？}因为\BibleQuote{骆驼穿过针眼，比富人进天主的国还容易}。因为\BibleQuote{在人这是不可能的，但在天主，一切都是可能的。} \BibleRef{玛 19:24} \BibleQuote{使祂将外邦人的产业赐给他们。}因为他们去到外邦人那里，嘱咐今世的富人\BibleQuote{不要心高气傲，也不要把希望寄托在无常的财富上，而要寄托在永生的天主身上} \BibleRef{弟前 6:17}，祂认为容易的事，在人却是困难的。因为这样有许多人被召，这样外邦人的产业被占据，这样甚至许多人在此生活中并未舍弃一切财产以便跟随祂的，却为了宣认祂的圣名而轻视了生命本身；他们如同骆驼，谦卑自己以承担患难的重负，仿佛穿过针眼一般，经由痛苦的狭窄通道进入了天国。这些效果是那位无所不能的天主成就的。\newline{}\par

6. \BibleQuote{祂手所行的是真理和公平} \BibleRef{咏 110:7}。让那些在此受审判的人持守真理。殉道者在这里被定罪，被带到审判台前，好使他们不仅审判那些审判他们的人，甚至审判天使 \BibleRef{格前 6:3}，他们在此与天使争战，即使当他们被人审判时也是如此。不要让患难、困苦、饥饿、赤身露体、刀剑使我们与基督分离。因为\BibleQuote{祂的一切诫命都是真实的} \BibleRef{罗 8:35}；祂不欺骗我们，祂赐给我们所应许的。然而，我们不应在此盼望祂所应许的；我们不应期望它；但\BibleQuote{它们永远立定，在真理和公平中施行} \BibleRef{咏 110:8}。这是公平合理的：我们在此劳苦，在那里安息；因为\BibleQuote{祂向祂的子民施行救赎} \BibleRef{咏 110:9}。但他们从什么中被救赎，岂不是从这旅居的囚禁中吗？因此，不要寻求安息，除非在天上的家乡。天主确实赐给了血肉的以色列子民一个属地的耶路撒冷，\Quote{她和她的儿女一起为奴：}但这是旧约，属于旧人。然而，那些在那时领悟了预象的人，已经是新约的继承人；因为\BibleQuote{那在上面的耶路撒冷是自由的，她是我们在天上的永生母亲。} \BibleRef{迦 4:25} 但在旧约中赐下的是暂时的应许，这由事实本身证明；然而，\Quote{祂永远确认了祂的盟约。}但哪一个盟约？岂不是新约？凡愿成为这盟约继承人的，不要自欺，不要想那流奶流蜜的地土，也不要想到愉快的田庄，或是果实丰盈、绿树成荫的园子；不要渴望获取任何这类东西，就是贪爱之眼所常贪求的。因为\BibleQuote{贪爱钱财是万恶之根} \BibleRef{弟前 6:10}，它必须被砍断，好在此被消灭；而不是被推迟，好在那里得到满足。首先要逃避惩罚，躲避地狱；在你渴求那应许的天主之前，要当心那威胁的天主。因为\Quote{祂的名是圣而可畏的。}\newline{}\par

7. ……所以，\BibleQuote{敬畏上主} \BibleQuote{是智慧的开端。} \BibleQuote{领悟是好事} \BibleRef{咏 110:10}。谁能反驳？但领悟而不实行是危险的。因此，\BibleQuote{是好事}，\BibleQuote{对遵行它的人。}并且不要让它使人心高气傲；因为\BibleQuote{祂的赞美}——敬畏祂是智慧的开端——\BibleQuote{永远长存：}这将是赏报，这将是终结，这将是永远的立足之地和居所。那里是真实的诫命，永远立定；这里是新约的产业，永远被确认。他说：\BibleQuote{有一件事我曾向上主祈求，我仍要寻求：就是一生一世居住在上主的殿宇中。}因为\BibleQuote{居住在上主的殿宇中}的人是有福的：\BibleQuote{他们要时时赞美}祂；因为\Quote{祂的赞美永远长存。}\newline{}\par

\SourceRecord{Translated by J.E. Tweed.；Nicene and Post-Nicene Fathers, First Series；Vol. 8.；Edited by Philip Schaff.；Buffalo, NY: Christian Literature Publishing Co.,；1888.；<http://www.newadvent.org/fathers/1801111.htm>.}

% structure: p0001 source=h1 target=section reason=page-title

\section{圣咏第一百一十二篇释义}

1. 弟兄们，我相信你们已经留意并记住了这篇圣咏的标题。他说：\Quote{哈盖和匝加利亚的归化。}当这些诗句被歌唱时，这两位先知尚未出现……但二人先后在一年内开始预言似乎与圣殿重建有关的事，正如许久以前所预言的……\Quote{因为天主的圣殿是圣的，这圣殿就是你们。}\BibleRef{格前 3:17}所以，凡归化自己去从事这项共同建造的工作、并怀着稳固而神圣建筑的希望，像一块活石从这悲惨世界的废墟中出来的人，就理解这篇圣咏的标题，理解\Quote{哈盖和匝加利亚的归化。}让他因此以下列诗句来歌唱，并非仅仅用舌头的发声，而是用他的生命。因为建造的完成将是那不可言喻的智慧之和平，其\Quote{开端}就是\Quote{敬畏上主：}因此，凡为此归化所建造的人，就从此开端。\par

2. \Quote{敬畏上主的人是有福的：他必在祂的诫命中极其喜乐。}\BibleRef{咏 111:1}唯独天主既诚实地又慈悲地审判，祂要看到人遵守祂诫命的程度；因为\Quote{人在世上的生命是诱惑，}\BibleRef{约 7:1}如同圣约伯所说的。但\Quote{判断我们的是主。}\BibleRef{格前 4:4}……因此祂要看到各人在祂的诫命中进步多少；然而，凡爱这共同建造之和平的人，必在其中极其喜乐；他也不该绝望，因为\Quote{在人间，善意的人享平安。}\BibleRef{路 2:14}\par

3. 接着是：\Quote{他的后裔在世上必成为大能。}\BibleRef{咏 111:2}宗徒作证，怜悯的工作是未来收成的种子，他说：\Quote{我们行善不要厌倦，因为到了时候，我们要收获；}\BibleRef{迦 6:9}又说：\Quote{我要说：那少撒种的，也必少收成。}\BibleRef{格后 9:6}但是，弟兄们，有什么比以下的事更为大能呢？不仅匝凯用他一半的财物买了天国，\BibleRef{路 19:8}甚至连那寡妇的两个小钱，\BibleRef{谷 12:42}而且每人应得同等的一份？有什么比以下的事更为大能呢？同样一个天国，对富人来说值财宝，对穷人来说值一杯凉水？……\Quote{荣耀和财富必在他的家中。}\BibleRef{咏 111:3}因为他的家就是他的心；在那里，因着赞美天主，他怀着永生的希望，生活在更大的财富中，胜过那些在奉承的人们中间、住在有大理石宫殿、天花板华丽装饰、却惧怕永死的人。\Quote{他的正义永远常存：}这就是他的荣耀，这就是他的财富。而另一个人的紫红袍、细麻衣和盛大的宴席，即使在享受时也正在消逝；当它们结束时，那燃烧的舌头要从手指尖求一滴水，却要呼喊。\BibleRef{路 16:24}\par

4. \Quote{在黑暗中，正直的人有光升起。}\BibleRef{咏 111:4}虔敬的人理应把自己的心指向他们的天主，理应同他们的天主行走，宁愿选择祂的旨意而不选择自己；他们心中没有骄傲的狂妄。因为他们记得，他们曾一度在黑暗中，但现在在主内是光。\BibleRef{弗 5:8}\Quote{上主天主是慈悲、怜悯和正义的。}祂是\Quote{慈悲和怜悯的}令我们喜悦，但上主天主是\Quote{正义的}或许令我们恐惧。不要害怕，一点也不绝望，你这敬畏上主并在祂诫命中有大喜悦的有福之人！你要温和，要慈悲，要借贷。因为主正是这样正义的：祂以不怜悯那未显示怜悯的人来审判；\BibleRef{雅 2:13}但是，\Quote{那慈悲并借贷的人是甜蜜的}\BibleRef{咏 111:5}：天主不会把他当作不甜蜜的人而从口中吐出。祂说：\Quote{宽恕，你们就必被宽恕；给予，就必给予你们。}\BibleRef{路 6:37}当你宽恕以求被宽恕时，你是慈悲的；当你给予以求被给予时，你是借贷的。因为尽管凡在苦难中协助他人的一切通常都称为怜悯，却有区别：你既不花费金钱，也不花费体力劳动，而是通过宽恕每个人对你所犯的罪，也为你自己的罪获得自由的赦免……那不愿给穷人的人寻求财富；请听所写的：\Quote{你将有宝藏在天上。}\BibleRef{玛 19:21}那么，你不会因宽恕而失去尊荣：因为战胜愤怒是一种非常值得赞美的胜利。你也不会因给予而变穷：因为天上的宝藏是更安全的拥有。前面的诗句\Quote{财富和丰裕必在他家中}已孕育着这一诗句。\par

5. 因此，实行这些事的人，\Quote{必能以明辨引导自己的言语。}他的行为本身便是言语，他将在审判中以此为自己辩护；审判不会对他缺少怜悯，因为他自己已经显示了怜悯。\BibleQuote{他必永不动摇}\BibleRef{咏 111:6}：他被召到右边，要听见这些话：\BibleQuote{你们蒙我父降福的，来吧！承受自创世以来为你们预备的国度吧！}因为那里提到的，没有别的善工，只有怜悯的工作。因此他将听见：\BibleQuote{你们蒙我父降福的，来吧！}因为\Quote{义人的族类必蒙降福。}如此，\Quote{义人必受到永久纪念。}\BibleQuote{他必不怕凶恶的传闻；因为他心意坚定，信赖上主}\BibleRef{咏 111:7}。正如他将听见对左边的人所说的话：\BibleQuote{离开我，到那为魔鬼和他的使者预备的永火里去吧！}\BibleRef{玛 25:34}。因此，谁在此世不求自己的利益，而只求耶稣基督的利益，\BibleRef{斐 2:21} 他便极其坚忍地忍受苦难，怀着信德等待恩许。他也不会被任何诱惑击倒：\BibleQuote{他的心志坚定，毫不畏缩，直到他看见仇敌遭报}\BibleRef{咏 111:8}。他的仇敌愿意在此世看到福乐，当无形的福乐被应许给他们时，他们常说：\Quote{谁能给我们显示福乐？}因此，让我们的心志坚定，毫不畏缩，直到我们看见仇敌遭报。因为他们愿在将死之地看到人的福乐；而我们却信靠要在永生之地看到上主的福乐。\par

6. 然而，使心志坚定、不被摇动，是件大事；当那些喜爱所见之物的人喜乐，并嘲笑那盼望所未见之物的人时，经上写道：\BibleQuote{天主为爱祂的人所准备的。}\BibleRef{格前 2:9} 这未见之物的价值何等大，而每个人只能按自己所能付出的代价来购买它。为此，他也\BibleQuote{慷慨施舍，赒济穷人}\BibleRef{咏 111:9}：他虽未见，却仍持续购买；但那位肯在世上穷人身上忍饥挨渴的天主，是在天上积蓄财宝。因此，若他的\Quote{仁义永存}就不足为奇了：因为创造万代的天主是他的护卫。他的\Quote{角}——其谦逊曾为骄傲者所蔑视——\Quote{必在荣耀中被高举。}\par

7. \BibleQuote{不虔敬的人看见这事，必要恼怒}\BibleRef{咏 111:10}：这就是那迟来而无果的悔改。因为当他说\BibleQuote{我们的骄傲，有什么益处？我们财富的夸耀，给了我们什么？}\BibleRef{智 5:8}时，他不是恼怒自己，而是恼怒谁呢？当他看见那位\Quote{慷慨施舍，赒济穷人}的人的角在荣耀中被高举时，\Quote{他必咬牙切齿，心力交瘁}；因为\Quote{那里要有哀号和切齿。}因为他不会再发出枝叶和开花，如同他及时悔改所会发生的那样；但那时他必要悔改，当\Quote{不虔敬者的愿望终归消灭}时，再无安慰相随。\Quote{不虔敬者的愿望终归消灭}，当\Quote{一切都如影子消逝}\BibleRef{智 5:8}，当花朵在草枯之际凋落时。\BibleQuote{惟有上主的话永远常存}\BibleRef{依 40:8}，假福乐者的虚妄曾嘲笑它，同样，当这些假福乐者真正悲惨时，它也要嘲笑他们的丧亡。\par

\SourceRecord{Translated by J.E. Tweed.；Nicene and Post-Nicene Fathers, First Series；Vol. 8.；Edited by Philip Schaff.；Buffalo, NY: Christian Literature Publishing Co.,；1888.；<http://www.newadvent.org/fathers/1801112.htm>.}

% structure: p0001 source=h1 target=section reason=page-title

\section{圣咏第113篇释义}

1. ...当你们在圣咏中听到咏唱，\BibleQuote{你们孩童，请赞美上主}\BibleRef{咏 112:1}；不要以为那劝勉与你们无关，因为你们既已度过了身体上的青春，要么正在盛年的风华之中，要么正带着年高德劭的花白头发。因为宗徒对你们众人说：\Quote{弟兄们，不要在心智上作孩童；但在邪恶上要作孩童，在心智上却要作成人。}那邪恶特别是什么？岂非骄傲？因为骄傲自恃虚假的伟大，不容人行走窄路，进入窄门；但孩童容易通过窄门进入；因此，除非成为孩童，没有人能进入天国。\Quote{请赞美上主的名。}……因此，愿祂永远被宣扬：\BibleQuote{愿上主的名受赞美，从现今直到永远}\BibleRef{咏 112:2}。\newline{}愿祂在各地被宣扬：\BibleQuote{从日出之地到日落之处，愿你们赞美上主的名}\BibleRef{咏 112:3}。\par

2. 如果在那些赞美上主之名的圣善孩童中有人问我并对我说：\Quote{直到永远}——我明白是指直到永远——但为什么\Quote{从现今，}而为什么上主的名不在现今以前、万世以前受赞美？我愿回答那个不是出于倔强而发问的幼童。对你们说的，是师傅和孩童们，对你们说的：\Quote{请赞美上主的名；愿上主的名受赞美：}愿上主的名受赞美，\Quote{从现今，}即从你们说这些话的时刻开始。因为你们开始赞美，但要无尽地赞美。……或者，既然在此段落中，他似乎更意指谦卑而非童年，其对立面是骄傲的虚妄和虚假的伟大；因此，除孩童外无人赞美上主，因为骄傲者不知如何赞美祂；愿老年如童年，你们的童年如老年；也就是说，既不可你们的智慧带有骄傲，也不可你们的谦卑没有智慧，好使你们\Quote{从现今直到永远赞美上主。}\newline{}无论基督的教会散布在哪里，在她孩童般的圣人中，都当\Quote{赞美上主的名；}也就是说，\Quote{从日出之地到日落之处。}\par

3. \BibleQuote{上主超越所有外邦人}\BibleRef{咏 112:4}。外邦人是人：若上主超越所有人，有何惊奇？他们用眼看到他们所崇拜的，在高天之上闪耀——太阳、月亮、星辰——这些受造物，他们侍奉它们却忽略了造物主。但不仅\Quote{上主超越所有外邦人}；而且\Quote{祂的荣耀}也\Quote{超越诸天。}诸天仰望那超越自身之上的祂；而谦卑的人与祂同在，他们不崇拜诸天以代替祂，尽管他们处在肉身之中、诸天之下。\par

4. \BibleQuote{谁相似上主、我们的天主？祂居住在高天，却垂顾谦卑的人？}\BibleRef{咏 112:5}。任何人都会以为祂居住在崇高的诸天，从那里祂可以垂顾地上的卑微事物；但\BibleQuote{祂垂顾诸天和地下的卑微事物}\BibleRef{咏 112:6}：那么，祂那崇高居所是什么，从那里祂垂顾诸天和地下的卑微事物？祂所垂顾的卑微事物，是否就是祂自己的崇高居所本身？因为祂如此高举谦卑的人，却又不使他们骄傲。祂既居住在高举的人中，又使他们成为祂自己的天，即祂的居所；并且看见他们不骄傲，而是恒常顺服于祂，祂在诸天本身中也看见这些非常卑微的事物，祂居住在高举的他们中。\newline{}因为圣神通过依撒意亚这样说：\Quote{至高者如此说，祂居住在高天，永远居住；上主至高，居住在圣者中。}祂用更充分的说法解释了祂借居住在高天所要表达的意思：\Quote{居住在圣者中。}……\par

5. 他也促使我们探究：上主、我们的天主在诸天和地下是否看见同样的卑微事物？还是在天上看见的卑微事物与在地上的不同？……但若上主、我们的天主在天上看见的卑微事物与在地上看见的不同；我想，祂已经在天上看见那些祂所召叫、并居住在其中的人；而在祂正在召叫、以便居住在他们中的那些人身上，祂在地上看见他们。因为祂使那些专心默想天上之事的人与祂同在，而另一些祂正在唤醒，虽然他们仍在梦着地上的事。\newline{}但是，既然连那些尚未以虔诚之心将颈项顺服于基督恩慈之轭的人也难以被称为谦卑，因为整个圣咏集中的神圣著作都警告我们要藉\Quote{谦卑}一词来理解\Quote{圣善}；还有另一种解释，亲爱的，你们可以与我一起思考。我认为，那些将要坐在十二个宝座上与主一同审判的人，现在是以\Quote{诸天}的名义被论及的；\BibleRef{玛 19:28}而\Quote{地下}的名义下，是其余蒙福的群众，他们将被置于右边，以便通过怜悯的工作，被那些他们用不义的钱财在今生为自己所结交的朋友赞美并接到永恒的居所。\BibleRef{路 16:9}……\par

6. \BibleQuote{他从尘土中提拔穷苦人，从粪土中抬举贫困人}\BibleRef{咏 112:7}；\BibleQuote{为使他与王侯同席，即与本国的王侯同席}\BibleRef{咏 112:8}。\newline{}因此，那些被高举的人的首领不可鄙弃在主的右手下成为谦卑。因为那忠信的管家虽蒙主托付钱财，被置于天主之民的王侯中，虽然他将要坐在十二个宝座上，甚至审判天使；\BibleRef{玛 19:28}然而他是从尘土中被提拔，从粪土中被抬举的。\newline{}岂不是那\Quote{曾侍奉各种私欲和享乐}的人被从粪土中抬举出来？……\par

7. 那么，弟兄们，如果我们已经听说过那些在天上的卑微者，从粪土中被高举，得以与祂子民的王子同坐；我们是否因此就未听说过主在地上所垂顾的卑微者呢？因为那些将与他们的主一同审判的朋友人数较少，而他们所接入永久居所的人则数目众多。因为虽然整个谷堆比起分离的糠秕，看起来数量很少；但就其本身而言，它却是丰盈的……教会正是在此意义上如此说话，似乎她在那群没有舍弃一切以便跟随主、并坐上十二个宝座的人群中不产生子嗣 \BibleRef{玛 19:28}。但同一群人中有多少人，他们用不义的钱财结交朋友 \BibleRef{路 16:9}，将因怜悯之工而站在右边呢？祂不仅从粪土中高举那将置于祂子民王子中的人；而且，祂还\BibleQuote{使不育的妇人安居家中，成为一位多子的欢乐母亲} \BibleRef{咏 112:9}：那居于高处、垂顾天上和地上卑微事物的那一位，亚巴郎的后裔如同天上的星辰，圣洁被高举于天上的居所；又如海边的沙粒，一个充满怜悯、不可计数的人群，从有害的波涛和不敬神的苦涩中被聚集起来。\par

\SourceRecord{Translated by J.E. Tweed.；Nicene and Post-Nicene Fathers, First Series；Vol. 8.；Edited by Philip Schaff.；Buffalo, NY: Christian Literature Publishing Co.,；1888.；<http://www.newadvent.org/fathers/1801113.htm>.}

% structure: p0001 source=h1 target=section reason=page-title

\section{\WorkTitle{圣咏第一一四篇释义}}

约旦河，当他们正要渡河进入应许之地时，被抬约柜的司铎的脚一触，河水便从上游停住，如同被缰绳束住一般，而下流的水则继续流入海中，直到全体人民在干地上走过，司铎们站在干地上。\BibleRef{苏 3:15} 我们知道这些事，然而我们不应以为，在这篇我们刚刚以歌唱阿肋路亚应答的圣咏中，圣神的目的只是要我们追忆那些过去的作为，而不再思想类似的事将来还要发生。因为正如宗徒所说：\BibleQuote{这些事} \BibleQuote{发生在他们身上，是为作鉴戒。}\BibleRef{格前 10:11}\par

\BibleQuote{以色列出了埃及，雅各伯家离开外邦人民。}\BibleRef{咏 113:1} \BibleQuote{犹大成了祂的圣所，以色列成了祂的王国。}\BibleRef{咏 113:2} \BibleQuote{海洋见了，就奔逃；约旦河也倒流。}\BibleRef{咏 113:3} 不要以为这些只是向我们叙述过去的事，不如说是预言未来。因为，虽然那些奇迹也在那民族中发生，确实发生了当时的事，但并非没有暗示未来的事……诗人以与我们所学所读不同的方式叙述了一些事，这样才不会真正被认为是重复过去的事，而非预言未来的事。首先，我们读到约旦河不是倒流，而是在靠近水源的一边停住，让百姓通过；其次，我们读到山和丘陵跳跃，这些都是诗人所加并重复的。因为在说了\BibleQuote{海洋见了，就奔逃；约旦河也倒流}之后，他又说：\BibleQuote{大山踊跃如公羊，小山舞蹈如羊羔}\BibleRef{咏 113:4}；然后问：\BibleQuote{海洋，你为什么奔逃？约旦河，你为什么倒流？}\BibleRef{咏 113:5} \BibleQuote{大山，你为什么踊跃如公羊？小山，你为什么舞蹈如羊羔？}\BibleRef{咏 113:6}\par

因此，让我们思考在这里我们所学到的东西。因为那些事迹都是我们的预像，而这些话也劝勉我们认识自己。如果我们以坚定的心持守天主所赐予我们的恩宠，我们就是以色列，亚巴郎的后裔；宗徒对我们说：\BibleQuote{因此你们是亚巴郎的后裔。}……所以，任何基督徒都不应认为自己与以色列之名无关。因为我们与那些相信的犹太人一同被联结在屋角石上，在这些犹太人中间宗徒们是首领。因此，主在另一处说：\BibleQuote{我还有别的羊，不属于这一栈，我也该把他们领来，他们要听我的声音，这样将只有一个羊群，一个牧人。}\BibleRef{若 10:16} 因此，基督徒人民才是真正的以色列，同样也是雅各伯家，因为以色列和雅各伯是同一的。但那大批的犹太人，因他们的背信而应受惩罚，为肉体的逸乐出卖了长子的名分，以致他们不属于雅各伯，而属于厄撒乌。因为你们知道，这隐义的话曾如此说：\BibleQuote{大的要服事小的。}\par

但埃及，既然据说意为苦难或使人受苦者或压迫者，常被用作今世的象征；我们必须从精神上脱离今世，以免与不信者同负一轭。\BibleRef{格后 6:14} 因为每个人正是这样成为天上耶路撒冷的适宜公民：先弃绝了今世。正如那民族必须先离开埃及，才能被领进应许之地。但他们并非凭自己离开，而是靠天主的帮助才获得自由。同样，没有人能在内心转离今世，除非蒙受天主仁慈的恩赐。因为那时所预表的，如今在每一位信者身上、在教会每日的产痛中、在这世界的终结、在这正如真福若望所写的末世中得以实现。\BibleRef{若一 2:18} 请听教导外邦人的宗徒如此教训我们：\BibleQuote{弟兄们，我不愿意你们不知道：我们的祖先都曾在云柱下，都从海中走过，都曾在云中和海中受了洗归于梅瑟，都吃过同样的神粮，都喝过同样的神饮；原来他们所饮的，是来自伴随他们的神磐石，那磐石就是基督。但在他们中，多数人未蒙天主喜悦，因而倒毙在旷野。这些事发生在他们身上，是为作我们的鉴戒。}\BibleRef{格前 10:1}\par

5. 请听更奇妙的事：古经卷中隐藏和遮蔽的奥秘在某种程度上藉古经卷得到启示。因为米该亚先知如此说：\Quote{按你出埃及的日子，我要向他显示奇妙的事，等等……}在这篇圣咏中，虽然奇妙的预言神视确实展望未来，但似乎只是在详述过去。他说：\BibleQuote{犹大成了祂的圣所；沧海看见就奔逃；}\Quote{成了}、\Quote{看见}、\Quote{奔逃}都是过去时态的词语；\Quote{约但河倒流，山岳跳跃，大地震动}同样用的是过去式表达，然而理解它们指向未来并无困难。……因为在那百姓离开埃及之后很久，并在教会这些时期之前那么久，他唱出我引用的诗句；但无论如何他见证自己无疑是在预言未来。他说：\Quote{按你出埃及地的日子，我要向他显示奇妙的事。}\Quote{万民看见，必惊慌失措。}这正是这里所说的：\BibleQuote{沧海看见就奔逃}；因为如果在这段经文中，透过过去时态的词语秘密地启示了未来，正如实情，谁敢将\Quote{看见并惊慌失措}解释为过去的事件呢？稍下文，他比光更清楚地暗示了那些追赶逃跑的我们、要杀害我们的敌人，即我们的罪恶，这些罪在圣洗中被淹没和消灭，正如埃及人淹死在海中，他说：既然\Quote{祂不永远怀怒，因为祂仁慈宽厚，祂必转回，怜悯我们，淹没我们的不义；并将他们的一切罪恶抛入深海。}\par

6. 最亲爱的，这是什么？你们知道自己是按亚巴郎的后裔为以色列人，你们是雅各伯的家，按恩许为承继人，要知道你们也已出离了埃及，因为你们弃绝了这世界；你们出离了外邦民族，因为你们藉虔诚的宣认，已与外邦人的亵渎分离。因为你们的舌头不懂赞美你们向他歌唱\OriginalTerm{阿肋路亚}{Allelujah}的天主。因为\BibleQuote{犹大}在你们中成了\BibleQuote{祂的圣所}；因为\BibleQuote{外表上的犹太人，并不是真犹太人；外表肉身的割损，也不是真割损；惟独内心的，才是真犹太人，割损也是内心的。}\BibleRef{罗 2:28} 那么，你们要省察自己的心，是否信德已割损了它们，是否宣认已洁净了它们；在你们中\BibleQuote{犹大}成了\BibleQuote{祂的圣所}，在你们中\BibleQuote{以色列}成了\BibleQuote{祂的国度}。因为\BibleQuote{祂赐给你们权能，成为天主的子女。}\BibleRef{若 1:12} ……\par

7. 但我不愿你们在自身以外寻求约但河如何倒流，不愿你们占卜任何凶兆。因为主责备那些\Quote{以背、而非以面向着}祂的人。\BibleRef{耶 2:27} 无论谁离弃他存在的根源，转身离开他的造物主，就如河流入海，滑入这世界的苦涩邪恶中。因此，他若回头，并在他回归时将他曾置于背后的天主放在面前，而将他曾置于面前、正滑向的世界之海抛在背后，那对他是有益的；好使他能忘记背后的，\Quote{竭力追求前面的，}\BibleRef{斐 3:13} 这对一旦悔改的人是有益的。……\par

8. \BibleQuote{大地啊，你应在主面前，在雅各伯的天主面前战栗。}\BibleRef{咏 113:7} 所谓\Quote{在主面前}，不就是在于那曾说\Quote{看，我天天与你们在一起，直到今世的终结}的那一位面前吗？\BibleRef{玛 28:30} 因为大地战栗了；但因为它一直怠惰，所以被震动，好能在主面前更坚固地立定。\par

9. \BibleQuote{祂使盘石变为水池，使燧石化为水泉。}\BibleRef{咏 113:8} 因为祂将自己，可说是祂的坚硬，熔化来浇灌那些信祂的人，好使祂在他们内成为\BibleQuote{涌到永生的水泉}；\BibleRef{若 4:14} 因为从前祂不为人所知时，似乎坚硬。因此，那些说\BibleQuote{这话生硬，谁能听得下去？}\BibleRef{若 6:60} 的人感到困惑，没有等待直到在圣经被揭示时祂流溢和涌流到他们身上。那盘石，那坚硬，变成了水池，那燧石变成了泉水，当祂复活后，\BibleQuote{从梅瑟及众先知开始，给他们讲解了基督必须这样受苦；}\BibleRef{路 24:26} 并派遣了圣神，关于祂曾说：\BibleQuote{谁若渴，到我这里来喝吧！}\BibleRef{若 7:37}\par

\SourceRecord{Translated by J.E. Tweed.；Nicene and Post-Nicene Fathers, First Series；Vol. 8.；Edited by Philip Schaff.；Buffalo, NY: Christian Literature Publishing Co.,；1888.；<http://www.newadvent.org/fathers/1801114.htm>.}

% structure: p0001 source=h1 target=section reason=page-title

\section{《圣咏第一一五篇释义》}

1. \Quote{主，光荣不要归于我们，不要归于我们，只归于祢的圣名}\BibleRef{咏 113:9}。因为从磐石中涌出的水（\Quote{那磐石就是基督}\BibleRef{格前 10:4}）的恩宠，不是根据先前的功行，而是出于祂的仁慈，\Quote{祂使不虔敬的人成义}\BibleRef{罗 4:5}。因为\Quote{基督为罪人死了}，好使人不再寻求自己的光荣，而只寻求主的名。\par

2. \Quote{为了祢的仁慈，为了祢的真道}\BibleRef{咏 113:10}。请注意，在圣经中，仁慈和真道这两个品质是多么频繁地结合在一起。因为祂以仁慈召叫罪人，又以真道审判那些被召却拒绝前来的人。\Quote{免得外邦人说：\InnerQuote{他们的天主在哪里？}}因为在最后，祂的仁慈和真道将彰显，那时\Quote{人子的记号要出现在天上，地上所有的支派都要哀号}\BibleRef{玛 24:30}；他们不会再说：\Quote{他们的天主在哪里？}因为祂不再被传讲为要人相信，而是显现在他们面前，令人战惧。\par

3. \Quote{我们的天主在天上}\BibleRef{咏 113:11}。不是在他们所看见的日月——即他们崇拜的天主所造的工物——的那个天上，而是\Quote{在天上}，超越所有天上和地上的物体。我们的天主也并非那样在天上，仿佛如果天从祂下面被抽走，祂就会失去宝座而坠落。\Quote{在天上和地下，祂随意创造了万有}。祂并不需要自己的工物，好像在其中有一个可以居住的地方；祂在自己永恒中持守，在那里居住，并随意创造了天上和地下的万有；因为万有并非先支撑祂，然后才被祂创造；因为除非它们被创造，它们不能支撑祂。因此，无论祂自己居住在哪里，祂可以说是包含那地方，因为那地方需要祂；祂并不被那地方所包含，好像祂需要它。或者可以这样理解：\Quote{在天上和地下，祂随意创造了万有}，无论是在祂子民中较高或较低的地位，祂都自由施予祂的恩宠，免得有人凭自己功绩的功德而自夸……\par

4. 他说：\Quote{他们的偶像乃是金银，是人手的作品}\BibleRef{咏 113:12}。也就是说，虽然我们不能向你们肉眼展示我们的天主（你们本应通过祂的工程认识祂），但你们不要因虚妄的伪装而受迷惑，因为你们可以用手指指着你们所崇拜的对象。因为你们若没有什么可指，反而更好，强如你们藉着向这些眼睛展示的东西，显露出你们内心的盲目：因为你们所展示的，除了金银，还有什么呢？他们也有铜、木、泥的偶像，以及各种不同材料的此等工物；但圣神特意提及更贵重的材料，因为当一个人为他更珍视的东西而感到羞愧时，他更容易被引导远离对更卑贱对象的崇拜。因为在圣经的另一处论到崇拜偶像的人说：\Quote{他们对木头说：\InnerQuote{你是我的父亲。}对石头说：\InnerQuote{你生了我。}}\BibleRef{耶 2:27}。但恐怕那个如此说话的对象不是石头或木头，而是金银的人，自视更为明智；让他朝这边看，让他把心灵的耳朵转向这边：\Quote{外邦人的偶像乃是金银}。这里没有提到卑贱可鄙的东西：其实对于那不属于尘世的心思，金银也是尘土，只是更美丽、更光亮、更坚实、更牢固。那么，不要用人手，从真天主所造的金属中造出一个虚假的神祇；不，是造出一个虚假的人，你们却当作真神来敬拜……\par

5. \Quote{他们有口，却不能言；有眼，却不能看}\BibleRef{咏 113:13}。\Quote{有耳，却不能听；有鼻，却不能闻}\BibleRef{咏 113:14}。\Quote{有手，却不能摸；有脚，却不能行；有喉，却不能出声}\BibleRef{咏 113:15}。因此，连制作它们的工匠也胜过它们，因为他有活动肢体和功能来塑造它们：但你们若崇拜那工匠，倒会感到羞耻。连你们也胜过它们，虽然你们没有造这些东西，因为你们能做它们所不能做的事。甚至野兽也胜过它们；因为为此又加上：\Quote{有喉，却不能出声}。因为上面已说过\Quote{有口，却不能言}，既然从头部到脚趾列举了所有肢体，何必再重复说它们不能出声呢？我想，是因为我们觉察到，提到其他肢体，是人与兽所共有的？因为兽类能看、能听、能嗅、能走，有些（如猴子）还能用手触摸。但关于口所说的话，是人所特有的：因为兽类不会说话。但恐怕有人认为这些话仅指人手的作品，且只将人置于外邦的神祇之上；于是在这一切之后，他又加上了这句话：\Quote{有喉，却不能出声}，这又是人与牲畜所共有的……那么，老鼠、蛇及其他类似的动物，对异教偶像的判断何等高明，因为它们看不到其中有人的生命，就不理会其中的人形。因此它们常在偶像中筑巢，除非被人驱赶，它们便寻找不到更安全的居所。人自己移动，从自己的神祇前赶走一只活物；却又崇拜那不能移动的神祇，仿佛他有能力，但一个比他崇拜对象更好的活物却被赶走了……连死人也胜过那既无生命又从未活过的神祇……\par

6. 然而那些自诩拥有更纯洁宗教的人说：我既不崇拜偶像，也不崇拜魔鬼；但在有形的形象中，我看到了我所应当崇拜之物的象征……他们妄图回答说，他们崇拜的不是形体本身，而是主管这些形体的神祇。因此，宗徒的一句话便证实了他们的惩罚和定罪。他说：\Quote{他们，}他说，\Quote{将天主的真理变为虚谎，且崇拜事奉受造物超过造物主，祂是永受赞美的。} \BibleRef{罗 1:25} 因为在这句话的前半部分，他谴责了偶像；在后半部分，谴责了他们对偶像的解释：因为他们用手工制作的形象来称呼天主创造之工的名称，从而将天主的真理变为虚谎；同时，他们把这些受造物本身视为神祇并加以崇拜，便事奉受造物超过了造物主，祂是永受赞美的……\par

7. 但是，有人会说，我们自己也有许多这类材质或金属制成的器具和器皿，用以举行圣事，它们因这些服务而被祝圣，被称为圣物，以尊崇那一位为我们的救恩而受崇拜的主：而这些器具或器皿本身，岂非人手所造的吗？但它们有口，却不能说话吗？有眼，却看不见吗？难道我们向它们祈祷，是因为通过它们我们向天主祈祷吗？这正是这种疯狂渎神的主要原因：那酷似活人的形象，诱使人崇拜它，在可怜之人的心中比它并非活着的明显事实更有影响力，因此它应当为活人所鄙视。\par

8. 随之而来的结果便是下一节经文所描述的：\Quote{制造它们的必与它们相似，凡信赖它们的也是如此。} \BibleRef{咏 113:16} 因此，让他们张着眼睛观看，却闭着死寂的理解去崇拜那些既看不见也无生命的偶像吧。\Quote{但以色列家族一直仰望于主。} \BibleRef{咏 113:17} \Quote{因为所见的希望不是希望；人既然看见，为什么还希望呢？但我们若希望那未看见的，就要耐心等待。} \BibleRef{罗 8:24} 但为使这忍耐能坚持到底：\Quote{祂是他们的扶助和保护者。} 也许属神的人（属血气的人藉着他们被建树于\Quote{温良之灵，} \BibleRef{迦 6:1}，因为他们作为较高者为较低者祈祷）已经看见，并且对他们而言，那对较低者是希望的东西，已经是现实吗？并非如此。因为连\Quote{亚郎家族也在主内怀有希望。} \BibleRef{咏 113:18} 因此，为使他们也能坚持不懈地奔向摆在前面的目标，并坚持不懈地奔跑，直到获得了他们所以被获得的那一位，\BibleRef{斐 3:12} 并且知道如同他们被知道一样，\BibleRef{格前 13:12} \Quote{祂是他们的扶助和保护者。} 因为二者都\Quote{敬畏主，并在主内怀有希望：祂是他们的扶助和保护者。} \BibleRef{咏 113:19}\par

9. 因为我们并非靠自己的功德抢先得到天主的仁慈；而是：\Quote{主曾眷顾我们，降福了我们。祂降福了以色列家族，降福了亚郎家族。} \BibleRef{咏 113:20} 但在降福这两者时，\Quote{祂降福了所有敬畏主的人。} \BibleRef{咏 113:21} 你会问，这两者指谁？祂回答：\Quote{无论大小。} 即以色列家族和亚郎家族，以及那民族中相信救主耶稣的人……因为在那民族中信了的人的身份上，经上说：\Quote{若非万军的上主给我们留下了苗裔，我们早已如同索多玛，相似哈摩辣了。} \BibleRef{罗 9:29} 苗裔，因为当它撒在地面上时，就增多了。\par

10. 因为亚郎家族中的伟人曾说过：\Quote{愿主使你们更加增多，使你们和你们的子女更加增多。} \BibleRef{咏 113:22} 这事就这样发生了。因为连从石头中兴起的子女也归向了亚巴郎：\BibleRef{玛 3:9} 不属于这羊栈的羊也归向了他，为能成为一牧一栈；万民的信仰加入了，数目增多了，不仅是有智慧的司祭，也有服从的人民；主不仅使他们的父辈更加增多——他们在基督内为其余应效法他们的人指引道路——也使他们的子女增多，子女追随父辈虔诚的足迹。\par

11. 因此先知对大小、高山和丘陵、公羊和羔羊如此说：\Quote{你们是主的降福者，祂创造了天与地。} \BibleRef{咏 113:23} 就如他应该说：你们是主的降福者，祂在伟大的中创造天，在微小的中创造地：不是这可见的天，缀有肉眼可见的光体。因为\Quote{诸天的天是属于主的，} \BibleRef{咏 113:24} 祂将一些圣者的心智提升到如此高度，以致他们不受任何人教导，而只受天主自己教导；与那天相比，凡肉眼所见之物皆可称为地；祂将\Quote{地赐给了世人；} 好叫人在默观它时，无论是在上照亮的那一区域（称为天），或是在下被照亮的那一区域（正确地称为地）（因为与我们称为诸天之天相比，整个，如我们所说，都是地），故此整个这地祂都赐给了世人，好叫他们通过默观它，尽可能构想出造物主，他们因心灵尚且软弱，没有那助其观念之物便不能看见祂。\par

12. 然而，因为他们获取智慧和丰富的真理，不是从人，也不是藉着人，而是藉着天主自己，他们领受了将要成为天的小子们，好使他们知道他们是诸天的天；但此时仍是地，他们为此说：\BibleQuote{我栽种了，阿颇罗浇灌了，然而使之生长的却是天主。}\BibleRef{格前 3:6}因为对于那些祂所造成天的世人，祂知道如何藉着天为地预备，祂赐给了他们可以劳作的地。愿因此，天和地，在他们的造物主天主内安居，并从祂那里生活，向祂宣认，赞美祂；因为如果他们选择从自己生活，他们就会死，正如经上所载：\BibleQuote{从死者，好似从虚无中，宣认止息。}\BibleRef{德 17:26}但\BibleQuote{死者不赞美祢，上主，也不赞美所有下到寂静中的。}\BibleRef{咏 113:25}因为圣经在另一处宣告：\Quote{罪人，当他进入邪恶的深渊时，便轻视。}\Quote{但我们这些活着的，将赞美上主，从今时直到永远。}\BibleRef{咏 113:26}\par

\SourceRecord{Translated by J.E. Tweed.；Nicene and Post-Nicene Fathers, First Series；Vol. 8.；Edited by Philip Schaff.；Buffalo, NY: Christian Literature Publishing Co.,；1888.；<http://www.newadvent.org/fathers/1801115.htm>.}

% structure: p0001 source=h1 target=section reason=page-title

\section{\WorkTitle{圣咏第一一六篇释义}}

1. \Quote{我喜爱，因为上主必将俯听我祷声。} \BibleRef{咏 114:1} 让那在远离上主的旅途中漂泊的灵魂如此歌唱，让那迷途的羊如此歌唱，让那\Quote{死了又复活、} \Quote{丢失了又被找到} \BibleRef{路 15:6}的儿子如此歌唱；弟兄们，我们最亲爱的子女们，让我们的灵魂如此歌唱。让我们受教，让我们持守，让我们与圣人们一同歌唱：\Quote{我喜爱，因为上主必将俯听我祷声。} 这难道是喜爱的理由——因为上主必将俯听我祷声？我们岂不是因祂已俯听或将俯听而喜爱？那么，\Quote{我喜爱，因为上主必将俯听} 是什么意思？是否因为他，既然希望能激发爱，便说他已喜爱，因为他希望天主会听他祷声？\par

2. 但他究竟从何怀此希望？他说：\Quote{因祂已侧耳向我；在我那一日，我就呼求了祂。} \BibleRef{咏 114:2} 我因此喜爱，因为祂将俯听；祂将俯听，\Quote{因为祂已侧耳向我。} 但你，人的灵魂，如何知道天主已侧耳向你，除非你说：\Quote{我已相信}？因此，这三者\Quote{常存：信德、望德、爱德。} \BibleRef{格前 13:13} 因为你已相信，所以你已希望；因为你已希望，所以你已喜爱。……\par

3. 而你的那一日是什么？既然你说：\Quote{在我那一日，我就呼求了祂。} 也许是这样的一日，其中\Quote{时期一满，} \Quote{天主就派遣了祂的圣子，} \BibleRef{迦 4:4} 祂早已说过：\Quote{在悦纳的时候，我俯允了你；在救恩的日子，我帮助了你。} \BibleRef{依 49:8} ……我宁可称我的日子是我苦难的日子，我必死的日子，按\Person{亚当}的日子，充满劳苦和流汗的日子，按那古老的腐朽的日子。因为\Quote{我沉沦于深泥中，} 在另一篇圣咏中我也曾呼号：\Quote{看，祢使我的日子转瞬即逝；} 在这我的日子里，我呼求了祢。因为我的日子不同于我主的日子。我称这些为我自己的日子，是我凭自己的胆量为自己造成的日子，我借此离弃了祂：而由于祂处处为王、全能全权、掌握万有，我理应受监禁；也就是说，我领受了无知的黑暗和必死的捆绑。……因为在我的这些日子里，\Quote{死亡的绳索缠绕了我，阴间的痛苦抓住了我，} \BibleRef{咏 114:3} 这些痛苦本不会抓住我，若我没有远离祢。但现在它们抓住了我；但我并未发现它们，当我在世上的福乐中喜悦时，阴间的罗网更容易欺骗人。\par

4. 但此后\Quote{我也遇到了困苦与忧愁，我呼求了上主的圣名。} \BibleRef{咏 114:4} 因为困苦与有益的忧愁，我未曾感受；这困苦中，祂赐予援助，正如有人对祂所说：\Quote{求祢做我们困苦中的援助；因为人的援助是虚妄的。} 因我以为可以在人的虚妄援助中喜乐欢腾；但当我从我的主那里听到：\Quote{哀恸的人是有福的，因为他们要受安慰。} \BibleRef{玛 5:4} 我便不再等待失去那些我曾喜乐的世俗恩赐，然后才哀恸；而是我定睛凝视我自己的这种悲惨——它使我因这些事物而喜乐，我又害怕失去它们，又不能保持它们；我坚定而勇敢地凝视它，并看到我不仅被此世的逆境折磨，甚至被其顺境束缚；于是我\Quote{遇到了困苦与忧愁，} 它们曾躲避我，\Quote{就呼求了上主的圣名：上主，我恳求祢，求祢解救我的灵魂。} 愿天主的圣民说：\Quote{我呼求了上主的圣名；} 让其余的外邦人听吧，他们尚未呼求上主的圣名；让他们倾听并寻求，好发现困苦与忧愁，并呼求上主的圣名，而得救。……\par

5. \Quote{上主是仁慈的，是公义的；是的，我们的天主是富于怜悯的。} \BibleRef{咏 114:5} 祂是仁慈的、公义的、富于怜悯的。首先祂是仁慈的，因为祂侧耳向我；而我原先并不知道天主的耳朵已贴近我的唇，直到我被那美丽的双脚唤醒，好叫我呼求主的圣名：因为除非祂先呼召，谁曾呼求祂呢？因此祂首先是\Quote{仁慈的；} 但\Quote{公义的，} 因为祂鞭打；又\Quote{富于怜悯的，} 因为祂接纳；因为\Quote{祂鞭打凡接纳的儿子；} 因此祂鞭打我，这不应使我痛苦，如同祂接纳我一样甘甜。\Quote{上主保护质朴的人，} \BibleRef{咏 114:6} 祂岂不鞭打那些当成年时祂欲使其成为继承人的？\Quote{因为哪有儿子不受父亲惩戒的呢？} \BibleRef{希 12:6} \Quote{我遭受苦厄，祂拯救了我。}\par

6. \BibleQuote{你的灵魂啊，再回到你的安息吧，因为主厚待了你}\BibleRef{咏 114:7}：不是因你的功绩，或借你的力量，而是因为主厚待了你。他说：\BibleQuote{因为他救了我的灵魂脱离死亡}\BibleRef{咏 114:8}。最亲爱的弟兄们，这是奇妙的：在他说他的灵魂要回到安息，因为主酬报了他之后，他又加上：\Quote{因为他救了我的灵魂脱离死亡}。它回到安息，是因为脱离了死亡吗？安息岂不更常用来指死亡吗？那生命是安息、死亡是烦扰的人，他的行动是怎样的呢？那么，灵魂的行动应当趋向安静的安稳，而非增加劳苦的烦扰；因为那位怜悯他的主救它脱离了死亡，说：\BibleQuote{凡劳苦和负重担的，你们都到我这里来，我要使你们安息}\BibleRef{玛 11:28}等等。因此，灵魂的行动——它趋向安息——应当是温良而谦逊的，可以说是跟随基督作为它的道路；然而，不是懒惰而懈怠的，好使它完成它的赛程，正如经上所记：\BibleQuote{在安静中完成你的工作}\BibleRef{德 3:19}。\BibleQuote{你救了我的灵魂脱离死亡，我的眼睛脱离眼泪，我的脚脱离跌倒}\BibleRef{咏 114:8}。凡感受到这肉身的锁链的人，都在希望中咏唱这些事，视作在自己身上应验。因为\BibleQuote{我备受痛苦，他便拯救了我}这话说得对；但宗徒也说，我们因希望而得救\BibleRef{罗 8:24}。我们说我们已脱离死亡，这已应验，为使我们理解不信者的死亡；论到他们，经上说：\BibleQuote{任凭死人埋葬他们的死人}\BibleRef{玛 8:22}。……那时他必擦干我们的眼睛，当他挽救我们的脚免于跌倒时。因为那时行走的脚不会再滑倒，因为软弱的肉身不再滑倒。但现在，尽管我们的道路——即基督——多么坚固，但我们把必须制服的肉身放在下面；在克制和制服它的过程中，不跌倒已是大事；但在肉身中不滑倒，谁能做到呢？\BibleQuote{我将在上主面前、在活人之地行走}\BibleRef{咏 114:9}？……我们现在确实劳苦，因为我们等待\BibleQuote{我们身体的救赎}\BibleRef{罗 8:23}；但是，\BibleQuote{当死亡被吞没在胜利中，这可朽坏的穿上不朽坏，这可死的穿上不死时}\BibleRef{格前 15:53}，那时将没有哭泣，因为没有跌倒；没有跌倒，因为没有腐朽。因此，我们那时将不再劳苦去讨他喜悦，而将完全在主面前、在活人之地讨他喜悦。\par

7. ……\BibleQuote{我信了，所以我说话。但我受了极大的屈辱}\BibleRef{咏 115:10}。因为他因他所忠信持守、忠信宣讲的言语而遭受了许多磨难；他受了极大的屈辱，正如那些宁愿爱人的赞美而不爱天主的赞美的人所害怕的。但\Quote{但我}是什么意思？他本应说：我信了，所以我说话，我受了极大的屈辱。他为什么加上\Quote{但我}？岂不是因为人可以因反对真理的人而受极大的屈辱，但他所信和所宣讲的真理本身却不能？因此，宗徒在论及他的锁链时也说：\BibleQuote{天主的言语不被束缚}\BibleRef{弟后 2:9}。所以，这位作者——由于神圣见证者（即天主的殉道者）的同一位格——也说：\Quote{我信了，所以我将说话}。\Quote{但我}；不是我所信的，不是我交付的言语；\Quote{但是我受了极大的屈辱}。\par

8. \BibleQuote{我在出神中说：众人都虚妄}\BibleRef{咏 115:11}。他所谓的出神，是指恐惧——当迫害者威胁、折磨或死亡的痛苦临近时，人性软弱所遭受的。我们这样理解，因为在这篇圣咏中听到了殉道者的声音。因为出神也有另一种含义：那时心智并非因恐惧而失常，而是被某种启示的灵感所占据。\BibleQuote{但我匆忙中说：众人都虚妄}\BibleRef{咏 115:11}。在惊惶中，他注视了自己的软弱，看出他不应信赖自己；因为就人本身而言，他是虚妄的，但借着天主的恩宠，他成了真实的；以免在敌人的压力下他可能不宣讲他所信的，反而否认它，正如伯多禄所遭遇的，因为他曾信赖自己，并要受教：人不应当信赖人。如果每个人都不应信赖人，当然也不应信赖自己，因为他是一个人。因此，他在恐惧中正确地看出每个人都是虚妄的；因为那些没有任何恐惧能夺走其镇静的人，以致他们从不因屈服于迫害者而说谎，也是天主的恩赐，而非自己的力量。……\par

9. 他问：\BibleQuote{我拿什么报答主，回报他赐给我的一切恩惠？}\BibleRef{咏 115:12}。他不是说：回报他为我所做的一切恩惠，而是\Quote{回报他回报给我的一切恩惠}。那么，在人这一方面，有什么先行的作为，使天主的这一切恩惠不是说是赐予的，而是回报的呢？在人一方面，有什么先行的作为呢？岂不就是罪过吗？因此，天主以善报恶，而人却以恶报善；因为那些说\BibleQuote{这是继承人，来，我们杀了他}\BibleRef{玛 21:38}的人，就是这样回报的。\par

10. 但此人寻求能以什么回报主，却找不到，除非出于主自己所回报的那些事物。他说：\Quote{我要举起救恩的杯，呼求主的名}\BibleRef{咏 115:13}。\Quote{我要向主还我的誓愿，在祂全体百姓面前}\BibleRef{咏 115:14}。谁给了你这救恩的杯，当你拿起它，呼求主的名，你就得以向祂回报祂所回报给你的一切酬报？除了那说\Quote{你们能饮我将要饮的杯吗？}的那位，还有谁？谁给了你效法祂的苦难，除非那先为你受苦的那位？因此，\Quote{在主眼中，祂的圣人之死极为宝贵}\BibleRef{咏 115:15}。祂以祂的血赎买了它，祂先为奴仆的救恩倾流了血，使他们不致犹豫为祂的名流自己的血；然而，这对他们自己的益处有利，而非对主有益。\par

11. 因此，让那以如此高价赎买的奴隶承认他的身分，并说：请看，主，我是祢的仆人：\Quote{我是祢的仆人，是祢婢女的儿子}\BibleRef{咏 115:16}。……因此，这是天上的耶路撒冷之子，那在上的、我们众人的自由母亲。\BibleRef{迦 4:26} 她确实脱离了罪，却是正义的婢女；对她的仍在旅居的儿子们说：\Quote{你们被召归于自由}\BibleRef{迦 5:13}；而又使他们成为仆人，当他说：\Quote{但要用爱德彼此服事。}……因此，让那仆人对天主说：许多人自称殉道者，许多人自称祢的仆人，因为他们持有祢的名号于各种异端和错误中；但既然他们站在祢的教会之外，他们就不是祢婢女的子女。但是\Quote{我是祢的仆人，是祢婢女的儿子。}\Quote{祢解开了我的锁链。}\par

12. \Quote{我要向祢献上赞颂的祭献}\BibleRef{咏 115:17}。因为我没有找到任何我自己的功绩，既然祢解开了我的锁链；所以我欠祢赞颂的祭献；因为，尽管我夸耀说我是祢的仆人、祢婢女的儿子，我却不以自己夸耀，而是以祢，我的主，祢解开了我的锁链，使我从我的背弃中回来时，能再次被缚于祢。\par

13. \Quote{我要向主还我的誓愿}\BibleRef{咏 115:18}。你要还什么誓愿？你誓许了什么祭品？什么燔祭，什么全燔祭？你是否指你先前所说的：\Quote{我要举起救恩的杯，呼求主的名}；以及\Quote{我要向祢献上感恩的祭献}？的确，凡认真思考他在向主誓许什么、他在还什么誓愿的人，让他誓许他自己，让他把自己当作誓愿还上：这正是所要求的，这正是所亏欠的。看着钱币，主说：\Quote{凯撒的归给凯撒，天主的归给天主}\BibleRef{玛 22:21}；他自己的像归给凯撒：让祂的像归给天主。\par

14. \Quote{在主的殿宇的庭院中}，他说，\Quote{在主的殿宇内}\BibleRef{咏 115:19}。主的殿宇是什么，主的婢女也是什么：天主的殿宇是什么，不就是祂的所有子民吗？因此接着是：\Quote{在祂全体百姓眼前}。如今他更公开地称呼他的母亲本身。因为祂的子民若不是下面接着的\Quote{耶路撒冷啊！在你中间}又是什么？因为那感恩的回报，若从和平中并依照和平而回报，才算是回报。但那些不是这位婢女之子的人，喜爱战争胜过和平……\par

\SourceRecord{Translated by J.E. Tweed.；Nicene and Post-Nicene Fathers, First Series；Vol. 8.；Edited by Philip Schaff.；Buffalo, NY: Christian Literature Publishing Co.,；1888.；<http://www.newadvent.org/fathers/1801116.htm>.}

% structure: p0001 source=h1 target=section reason=page-title

\section{圣咏第117篇释义}

1. \BibleQuote{列国万民，请赞美上主，一切民族，请歌颂上主} \BibleRef{咏 117:1} 。这些是主殿宇的庭院，这就是主的所有子民，这就是真正的\Place{耶路撒冷}。那些拒绝成为这城子女的人，更当听从，因为他们自绝于万民的共融。\BibleQuote{因为祂的仁慈厚加于我们，上主的忠诚必要永远常存} \BibleRef{咏 117:2} 。这就是那两件事：慈爱与真理，我在\Emph{第一一五}篇圣咏中曾劝你们要铭记。但主的仁慈厚加于我们，因为敌对民族的狂怒口舌已屈服于祂的圣名，藉此我们获得了自由；主的忠诚永远常存，无论是在祂应许给义人的事上，还是在祂警告恶人的事上。\par

\SourceRecord{Translated by J.E. Tweed.；Nicene and Post-Nicene Fathers, First Series；Vol. 8.；Edited by Philip Schaff.；Buffalo, NY: Christian Literature Publishing Co.,；1888.；<http://www.newadvent.org/fathers/1801117.htm>.}

% structure: p0001 source=h1 target=section reason=page-title

\section{圣咏第一百一十八篇释义}

1. 当我们高唱阿肋路亚（意即赞美主）时，这首圣咏教导我们，当我们听到\BibleQuote{你们应向主谢恩}\BibleRef{咏 117:1}时，就该赞美主。对天主的赞美没有比\BibleQuote{因为他是美善的}这更简短的表达了。我认为没有比这更庄严的简短话语了，因为美善如此特殊地属于天主，以致天主子本人当有人称他为\BibleQuote{善师}时——那人只看见他的肉身，不理解他天主性的圆满，仅把他视为人——他回答说：\BibleQuote{你为什么称我为善？除了天主一个外，没有谁是善的。}\BibleRef{谷 10:17}这不就等于说：你若愿称我为善，就当承认我是天主吗？但是，由于这首圣咏是对未来事物的启示，是对一个脱离了所有劳苦、漂泊和邪恶混杂的子民（这自由是藉天主的恩宠赐予的，天主不仅不以恶报恶，甚至以善报恶）而说的，因此最恰当地加上了\BibleQuote{因为他的仁慈永远常存}。\par

2. \BibleQuote{愿以色列家现在称谢他是美善的，因为他的仁慈永远常存}\BibleRef{咏 117:2}。\BibleQuote{愿亚郎家现在称谢他的仁慈永远常存}\BibleRef{咏 117:3}。\BibleQuote{愿所有敬畏主的人现在称谢他的仁慈永远常存}\BibleRef{咏 117:4}。我想，最亲爱的，你们记得什么是以色列家，什么是亚郎家，而两者都是敬畏主的人。因为他们是\Quote{小与大}，已在另一篇圣咏中被快乐地介绍到你们心中；我们所有的人都应当欢欣鼓舞，因他的恩宠（他是美善的，他的仁慈永远常存）而联合在一起；因为他们的话蒙垂听：\Quote{愿主使你们更加增多，你们和你们的子女}；这样外邦人的群众就可能加给信基督的以色列人，在他们中有我们的宗徒父辈，为了成全者的高举和小孩们的服从；以致当我们都在基督内合一，成为同一牧者下的同一羊群，并成为那头的身体，如同一人，可以说：\BibleQuote{我在患难中呼求主，主就宽阔地应允我}\BibleRef{咏 117:5}。我们患难的狭窄是有限的；但我们行走的宽阔道路没有尽头。\BibleQuote{谁能控告天主所拣选的人呢？}\BibleRef{罗 8:33}\par

3. \BibleQuote{主是我的助佑；我不害怕人能对我做什么}\BibleRef{咏 117:6}。难道人就是教会仅有的仇敌吗？一个属于血肉的人，除了血肉外，还算什么？但宗徒说：\BibleQuote{我们战斗不是对抗血肉，而是对抗}……他说：\BibleQuote{在天上的邪恶势力}\BibleRef{弗 6:12}，也就是说，魔鬼和它的使者；那魔鬼，他在别处称它为\BibleQuote{空中的掌权者}\BibleRef{弗 2:2}。因此请听下文：\BibleQuote{主是我的助佑；因此我要蔑视我的仇敌}\BibleRef{咏 117:7}。无论我的仇敌从何而来，无论是从恶人中，还是从恶天使中；在主（我们向他高唱赞美的谢歌，向他歌唱阿肋路亚）的助佑下，它们必被蔑视。\par

4. 但是，当我的仇敌被轻视后，不要让我的朋友作为善人出现在我面前，以致叫我寄望于他；因为\BibleQuote{信赖主，胜过信赖世人}\BibleRef{咏 117:8}。也不要让任何人——他在某种意义上可被称为善天使——被我视为应当信赖的对象；因为\BibleQuote{没有谁是善的，只有天主一个}\BibleRef{谷 10:18}；而当一个人或一个天使好像要帮助我们时，当他们出于真诚的爱心这么做时，是那按他们的尺度造他们为善的天主藉着他们行事。因此\BibleQuote{信赖主，胜过信赖王子}\BibleRef{咏 117:9}。因为天使也被称为王子，正如我们在\WorkTitle{达尼尔书}中读到：\BibleQuote{弥额尔，你的王子}\BibleRef{达 12:1}。\par

5. \BibleQuote{所有民族围困了我，但靠着主的名我报复了他们}\BibleRef{咏 117:10}。\BibleQuote{他们从四面围困我，我再说，他们从四面围困我；但靠着主的名我报复了他们}\BibleRef{咏 117:11}。他表明教会的劳苦和胜利；但是，仿佛有人问教会怎能克服如此巨大的邪恶，他回顾先例，宣告教会首先在她的头身上所受的苦，接着加上：\Quote{他们从四面围困我}；而\Quote{所有民族}这句话，有理由在这里不重复，因为这只是犹太人的行为。在那里，那个非常虔诚的民族（它是基督的肢体，为了它，那内在的神性通过外在的肉体以不朽的能力在必死的形态中所做的一切都已成就）曾遭受迫害者的苦，而这迫害者的种族正是那肉身所取并悬在十字架上的。\par

6. \BibleQuote{他们围困我，如同蜜蜂围困蜂巢，又像荆棘中的火焰被烧尽；我因上主的名，向他们复仇。} \BibleRef{咏 117:12} 这里语词的顺序与事件的顺序相符。因为我们正确理解，主自己，教会的元首，曾被迫害者围困，如同蜜蜂围困蜂巢。因为圣神以神秘的微妙谈论那些不知自己所作所为之人的行为。蜜蜂在蜂巢中酿蜜；而主的迫害者，虽不自觉，却因主的苦难，使我们尝到并看见主是何等甘甜，祂\BibleQuote{为我们的罪死了，为使我们成义而复活。} \BibleRef{罗 4:25} 至于接下来的\Quote{又像荆棘中的火焰被烧尽}，最好理解为祂的身体，即一个散布四方的民族，万民都环绕它，因为它从万民中聚集而来。他们以迫害的火焰烧尽了这有罪的肉体，以及这必朽生命的刺伤。\Quote{我向他们复仇：} 或者因为他们自己，那在他们身上迫害义人的邪恶被熄灭后，加入了基督的子民；或者因为其余的人，他们现今藐视那召唤者的仁慈，终将感受到那审判者的真实。\par

7. \BibleQuote{我如同沙堆被驱散，几乎跌倒，但上主扶持了我。} \BibleRef{咏 117:13} 因为虽然有一大群信徒，可比作无数沙粒，并被集合在同一共融中如同一堆；然而\Quote{人算什么，你竟眷顾他？} 他不是说外邦人的众多不能超过我队伍的丰盛，而是说\Quote{上主扶持了我。} 外邦人的迫害未能成功将居住在信仰合一中的信徒队伍推倒。\par

8. \BibleQuote{上主是我的力量，我的赞美，也成了我的救恩。} \BibleRef{咏 117:4} 那么，谁会在被推时跌倒呢？只有那些选择做自己的力量和赞美的人。因为没有人会在争战中跌倒，除非他的力量和赞美失效。因此，那以上主为力量和赞美的人，跌倒的次数不会超过上主跌倒。正因为如此，上主成了他们的救恩；并非祂变成了先前所不是的，而是因为他们信靠祂时，变成了先前所不是的，然后祂开始成为他们的救恩——当他们转向祂时，这是祂先前在他们背向时所未曾是的。\par

9. \BibleQuote{在义人的帐幕中，有欢乐和得救的声音。} \BibleRef{咏 117:15} 而那些向义人身体发怒的人，却以为那里有哀伤和毁灭的声音。因为他们不知道圣人们在对将来希望的渴望中内在的喜乐。因此宗徒也说：\BibleQuote{似乎忧愁，却常常喜乐；} \BibleRef{格后 6:10} 又说：\BibleQuote{不但如此，我们在患难中也夸耀。} \BibleRef{罗 5:3}\par

10. \BibleQuote{上主的右手行了大事。} \BibleRef{咏 117:16} 什么大事？他说：\Quote{上主的右手高举了我。} 高举谦卑者、神化必朽者、从软弱中成全、从屈辱中得荣耀、从苦难中得胜利、给予援助、从患难中提拔——使天主的真正救恩向受苦者显明，而人的救恩对迫害者毫无效用——这些是大事。但你惊讶什么？听他所重复的：\Quote{上主的右手行了大事。}\par

11. \BibleQuote{我必不死，而将活着，并宣讲上主的作为。} \BibleRef{咏 117:17} 但他们四处施加毁灭和死亡，以为基督的教会正在死去。看哪，他现在宣讲了上主的作为。基督处处是蒙福殉道者的荣耀。祂以被打胜了打祂的人；以忍受酷刑胜了施刑者；以爱胜了那些向祂发怒的人。\par

12. 然而，让我们指出，为何基督的身体、圣教会、义子之民，会遭受如此侮辱。他说：\BibleQuote{上主惩戒和管教了我，却没有将我交于死亡。} \BibleRef{咏 117:18} 因此，骄傲的恶人不可幻想有什么东西被交予他们的能力；他们不会有权柄，除非从上面赐予。家主常常命令最无用的奴仆管教他的儿子；虽然他为儿子预备产业，为奴仆预备锁链。那产业是什么？是金子、银子、珠宝、田庄、欢乐的庄园吗？请思考我们如何进入它，并了解它是什么。\par

13. 他说：\BibleQuote{请给我敞开正义的门。} \BibleRef{咏 117:19} 看，我们听到了门。门内是什么？他说：\BibleQuote{我要进入它们，并向上主谢恩。} 这是充满惊叹的赞美宣认，\Quote{直到天主的殿宇，在欢乐和赞美宣认的声音中，在守庆节的人中：} 这是义人永远的真福，他们因住在主的殿中而蒙福，永远赞美祂。\par

14. 但请思考如何进入正义的门。他说：\BibleQuote{这些是上主的门，义人要进入它们。} \BibleRef{咏 117:20} 至少不要让任何恶人进入那里，即那不接受未受割礼者的耶路撒冷，那里说：\BibleQuote{外面是狗。} \BibleRef{默 22:15} 我漫长的朝圣之旅中曾\Quote{居住在刻达尔的帐幕中}已足够；我甚至忍受了与恶人交往到底，但\Quote{这些是上主的门：义人要进入它们。}\par

15. \BibleQuote{我要称谢你，上主，因为你俯听了我，成了我的救恩}\BibleRef{咏 117:21}。这称谢往往被证实是赞美，它不向医生指出伤口，而是为所得的痊愈感恩。但医生自己就是救恩。\par

16. 但我们所说的是谁呢？\BibleQuote{匠人弃而不用的石头}\BibleRef{咏 117:22}；因为\BibleQuote{它成了屋角的基石}\BibleRef{咏 117:22}，为\BibleQuote{在祂自己身上将双方造成一个新人，成就和平；并藉着十字架使双方在一个身体内与天主和好}\BibleRef{厄 2:15}，即受割损的和未受割损的。\par

17. \BibleQuote{这是上主所作的}\BibleRef{咏 117:23}：也就是说，由主使它成为屋角的基石。因为尽管祂若不承受苦难就不会成为这样，但祂成为这样并非由于那些令祂受苦的人。因为那些建造的人弃绝了祂；但在主暗中建造的工程中，他们所弃绝的石头却成了屋角的基石。\BibleQuote{这在我们的眼中是神奇的：}在内心的眼中，在那些相信、盼望、爱慕之人的眼中；不是在那些因鄙视祂，视祂如凡人而弃绝祂的人的血肉之眼中。\par

18. \BibleQuote{这是上主所安排的一天}\BibleRef{咏 117:24}。此人想起他曾在以前的圣咏中说过：\Quote{因为祂曾侧耳听我，所以我一生要呼求祂}；提及他的往日，因此他现在说：\BibleQuote{这是上主所安排的一天}；即祂赐我救恩的日子。这是祂所说的日子：\BibleQuote{在悦纳的时候，我俯允了你；在救恩的日子，我帮助了你}\BibleRef{依 49:8}；那是一位中保成为屋角基石的日子。因此，\Quote{让我们欢欣踊跃，并在祂内喜乐。}\par

19. \BibleQuote{上主，求你拯救我；上主，求你使我道路顺利}\BibleRef{咏 117:25}。因为这是救恩的日子，\Quote{求你拯救我}；因为我们从漫长的流放归来，与那些仇恨和平的人分离，我们曾与他们和睦相处，他们却在我们与他们谈话时无故攻击我们；\Quote{求你在我们归途中使我们道路顺利}，因为你已成了我们的道路。\par

20. \BibleQuote{奉主名而来的当受赞美}\BibleRef{咏 117:26}。因此，奉自己名而来的是当受诅咒的；正如祂在福音中所说：\BibleQuote{如果别人奉自己的名而来，你们将会接受他}\BibleRef{若 5:43}。\Quote{我们从天主的殿中祝福了你。}我相信这是大者对小子说的话，即那些在精神上与天主的圣言（祂与天主同在）相通的大者，正如他们在今世所能做的；然而他们为了小子的缘故调整自己的言论，好能真诚地说出宗徒所说的话：\BibleQuote{如果我们疯狂，那是为了天主；如果我们清醒，那是为了你们。因为基督的爱催迫着我们}\BibleRef{格后 5:13}。他们从上主的殿内祝福小孩子们，在那里年复一年赞颂不止：因此请思量他们从那里宣告什么。\par

21. \BibleQuote{天主是上主，祂已赐给我们光明}\BibleRef{咏 117:27}。那位奉主名而来的主，被工匠弃绝，却成了屋角的基石\BibleRef{玛 21:9}，\Quote{天主与人类之间的中保，耶稣基督}\BibleRef{弟前 2:5}，祂是天主，与圣父同等，祂赐给我们光明，好叫我们明白我们所信的，并向你们这些尚未明白但已相信的人宣告。为使你们也能明白，\BibleQuote{在众集会中宣告节日，直到祭台的角}\BibleRef{咏 117:28}；即直到天主的殿内，我们从那里祝福了你，那里是祭台的高处。\Quote{宣告节日，}不可怠惰，而要\Quote{在众集会中。}因为这是节日喜庆的欢呼声，那些行走\Quote{在奇妙帐幕之处，直到天主的殿宇}的人所发出的。因为在那里如有属灵的祭献，永恒的赞美之祭，那么司祭是永恒的，义人平安的心是永恒的祭台……\par

22. 我们在那里将唱什么，除了赞美祂之外？我们在那里将说什么，除了\Quote{你是我的天主，我要称谢你；你是我的天主，我要赞美你。我要称谢你，因为你俯听了我，成了我的救恩。}我们不会用响亮的话语说这些；但存留在祂内的爱本身用这些话呼喊，这些话就是爱本身。因此，他以赞美开始，也以赞美结束：\BibleQuote{要称谢上主，因为祂是良善的，祂的仁慈永远长存}\BibleRef{咏 117:29}。圣咏以此开始，也以此结束；因为从那已离去的起点，也在这将要返回的终点，没有什么比赞美天主和永不停息的阿肋路亚更能令我们喜悦了。\par

\SourceRecord{Translated by J.E. Tweed.；Nicene and Post-Nicene Fathers, First Series；Vol. 8.；Edited by Philip Schaff.；Buffalo, NY: Christian Literature Publishing Co.,；1888.；<http://www.newadvent.org/fathers/1801118.htm>.}

% structure: p0001 source=h1 target=section reason=page-title

\section{《圣咏第119篇释义》}

% structure: p0002 source=h2 target=subsection reason=relative-heading-rank; base=h2

\subsection{阿肋夫}

1. 亲爱的诸位，这篇伟大的圣咏从开头就劝勉我们走向福乐，这是无人不向往的……因此，那说\Quote{品行无瑕、遵行上主法律的人，是有福的}\BibleQuote{\BibleRef{咏 119:1}}者，所教导的就是这个道理。这无异于说：我知道你们所愿望的，你们寻求福乐；那么，如果你们愿意有福，就该成为无瑕的。因为前者是众人所愿，后者却是众人所惧；然而没有后者，众人所愿的便无法获得。但何处能有人成为无瑕的呢？除非在道路上。在哪条道路上呢？除非在上主的法律中……\par

2. 现在请听他所加添的话：\Quote{那遵守祂的诫命、全心寻求祂的人，是有福的}\BibleQuote{\BibleRef{咏 119:2}}。在我看来，这话所提到的有福之人，并非与前面所说的不同。因为考察上主的诫命，全心寻求祂，就是在路上成为无瑕的，就是遵行上主的法律。他接着说：\Quote{因为行恶的人，必不行在祂的路上}\BibleQuote{\BibleRef{咏 119:3}}。然而我们知道，行恶的人之所以查考上主的诫命，是因为他们宁愿有学识而不愿成为义人；我们也知道另一些人查考上主的诫命，并非因为他们已经生活良善，而是为了知道应当如何生活。这样的人尚未无瑕地行在上主的法律中，因此他们尚不有福……\par

3. 在这篇圣咏中写着、读着、并且是真实的，即\Quote{行恶的人，不行在祂的路上}\BibleQuote{\BibleRef{咏 119:3}}。但我们必须在天主的助佑下努力——因为\Quote{我们和我们的话，都在祂手中}\BibleQuote{\BibleRef{智 7:16}}——以免那说得正确的话，因不被正确理解而混淆读者或听者。因为我们必须注意，免得所有圣人——他们说过这话：\Quote{如果我们说我们没有罪，便是欺骗自己，真理也不在我们内}\BibleQuote{\BibleRef{若一 1:8}}——要么被认为不行在上主的路上，因为罪是邪恶，而\Quote{行恶的人，不行在祂的路上}；要么，既然他们行在上主的路上是毫无疑问的，就被认为没有罪，而这显然是假的。因为这话并非仅仅为了避讳骄傲和自大而说的。否则就不会加上\Quote{真理也不在我们内}，而会说：谦卑也不在我们内；尤其是因为接下来的话更清楚地阐明了含义，消除了所有疑问的原因。因为当真福若望说了这话之后，他加上说：\Quote{如果我们承认我们的罪，祂是忠信和正义的，必赦免我们的罪，并洗净我们的一切不义}\BibleQuote{\BibleRef{若一 1:9}}……\par

4. \Quote{祢曾吩咐，叫我们极其遵守祢的诫命}\BibleQuote{\BibleRef{咏 119:4}}是什么意思？是\Quote{祢吩咐得太多}？还是\Quote{极其遵守}？无论我们理解为哪一种，其含义似乎与那句值得纪念和高尚的格言相反，就是希腊人在其智慧中所称颂的，拉丁人也一致赞扬：\Quote{勿过度}……但是拉丁语有时使用\Emph{nimis}一词，在圣经中我们见到它，也在我们的讲论中使用它，表示\Quote{非常}。在这句话\Quote{祢曾吩咐，叫我们极其遵守祢的诫命}中，如果我们正确理解，就简单地理解为\Quote{非常}；如果我们对某位至亲的朋友说\Quote{我太过爱你了}，我们并不想被理解为超过了适当的程度，而是\Quote{非常}爱。\par

5. 他说：\Quote{但愿我的道路坚定，为遵守祢的法度}\BibleQuote{\BibleRef{咏 119:5}}。祢确实吩咐过；但愿我能实现祢所吩咐的。当你听到\Quote{但愿}时，要认出这是愿望之言；认出了愿望的表达，就要放下骄傲的自恃。因为谁说他所渴望的事是在他能力之内，以至于无需任何帮助就能做到呢？因此，如果人渴望天主所吩咐的，就必须祈求天主亲自赐予祂所命令的……\par

6. \Quote{这样，我注视祢的一切诫命时，便不会羞愧}\BibleQuote{\BibleRef{咏 119:6}}。我们应当将天主的诫命——无论是读到时，还是回忆时——视为一面镜子，如同宗徒雅各伯所说的\BibleQuote{\BibleRef{雅 1:23}}。这个人希望自己成为这样的人，能如同在镜子中一般注视天主的诫命，而不致羞愧；因为他选择不仅作听者，而且作行者。为此他渴望他的道路得以坚定，以遵守天主的法度。如何得以坚定呢？除非靠着天主的恩宠？否则，他必会从天主法律中找到的，不是喜乐的缘由，而是羞愧的缘由，如果他选择了注视那并不实行的诫命。\par

7. 他说：\Quote{我要以正直的心，称谢祢；因为我学会了祢正义的典章}\BibleQuote{\BibleRef{咏 119:7}}。这不是认罪，而是赞美；正如那位没有罪的主也说的：\Quote{父啊，天地的主宰，我称谢祢}\BibleQuote{\BibleRef{玛 11:25}}；又如在\WorkTitle{德训篇}中写着：\Quote{在称谢时，你要这样说：上主的一切化工都极其美好}\BibleQuote{\BibleRef{德 39:15}}。他说：\Quote{我要以正直的心，称谢祢。}的确，如果我的道路被修直了，我要称谢祢，因为是祢成就了这事，这是祢的赞美，而不是我的……\par

8. 接着他补充说：\BibleQuote{我要遵守你的典章}（\BibleRef{咏 119:8}）……可是接下来是什么呢？\Quote{求你不要全然抛弃我！}或者，如同一些抄本所载，\Quote{不要太过分地抛弃我}，代替\Quote{全然抛弃我}。既然天主已将世界留在罪恶的报应中，他本会\Quote{全然}抛弃它，如果没有如此强大的疗效支持它，即藉着我们的主耶稣基督而来的天主的恩宠；但现在，按照基督身体的这个祈祷，他没有\Quote{全然}抛弃它；因为\BibleQuote{天主在基督内，使世界与自己和好}（\BibleRef{格后 5:19}）……\par

% structure: p0011 source=h2 target=subsection reason=relative-heading-rank; base=h2

\subsection{Beth}

9. \Quote{青年如何洁净自己的行径？就是遵守你的话}（\BibleRef{咏 119:9}）。他自问自答。\Quote{如何？}至此是提问：接着是回答，\Quote{就是遵守你的话。}但在此处，遵守天主的话必须理解为在行动上服从他的诫命：因为若不在生活中遵守，单存于记忆中是徒然的。但这里\Quote{青年}指的是什么？因为他本可以说：任何人（\OriginalTerm{任何人}{homo}）如何洁净自己的行径？或者，男子（\OriginalTerm{男子}{vir}）如何洁净自己的行径？圣经通常这样表达，以理解整个人类……但在此处，他既没有说任何人，也没有说男子，而是说\Quote{青年}。难道老年人就无望了吗？或者老年人除了以天主的话管束自己之外，还有其他方式洁净自己的行径吗？或者这也许是在提醒我们应在哪个年龄首要洁净自己的行径，正如别处所写：\BibleQuote{我儿，从你幼年时起就应接受教诲；这样你必能寻得智慧，直到你白发苍苍}（\BibleRef{德 6:18}）。另一种解释是，在这句话中认出福音中那个小儿子（\BibleRef{路 15:12}），他回到自己内说：\Quote{我要起身到我父亲那里去}（\BibleRef{路 15:18}）。他是如何洁净自己的行径呢？岂非借着按天主的话管束自己，因为他渴望父亲的饼，如同一个渴望的人……\par

10. \Quote{我全心寻求了你；求你不要使我偏离你的诫命}（\BibleRef{咏 119:10}）。看，他祈祷得蒙助佑以遵守天主的言语，他曾说青年以此洁净自己的行径。因为\Quote{求你不要使我偏离你的诫命}这句话的含义正是如此：被天主抛弃，岂非就是得不到助佑？因为人的软弱不足以服从他正义而崇高的诫命，除非他的爱预先扶助并支持。那些他不扶助的人，就正当地被称为被他抛弃……\par

11. \Quote{我将你的话藏在我心里，免得我得罪你}（\BibleRef{咏 119:11}）。他立刻寻求天主的助佑，免得天主的话徒然藏在他心里而无义德的行为相随。因为说完这话后，他接着又说：\Quote{上主，你应受赞美，求你教训我你的义德}（\BibleRef{咏 119:12}）。\Quote{求你教训我}，他说，如那些实行的人所学；而非如那些仅凭记忆以作谈论的人。那么他为何说\Quote{求你教训我你的义德}呢？岂非因为他愿通过行动，而非通过言谈或记忆来学习？既然在另一篇圣咏中读到：\Quote{赐予法律的那一位，也必赐福}；因此，他说：\Quote{上主，你应受赞美，求你教训我你的义德。}因为我已将你的话藏在我心里，免得我得罪你，你赐下了法律；也请赐下你恩宠的祝福，使我因行义而学到你所命令要教导的……\par

12. \Quote{我用我的嘴唇宣扬了你口中所出的一切判断}（\BibleRef{咏 119:13}）；这就是说，关于你的判断，你愿意藉着你的言语使我认识的一切，我无一缄默，而是毫无例外地用我的嘴唇全部宣扬。我认为他所指的是，他说并不是\Quote{你的一切判断}，而是\Quote{你口中所出的一切判断}；也就是你启示给我的那些。藉着他的口，我们可以理解他的话语，即他在众多圣人的启示以及旧约和新约中向我们揭示的；教会无时无刻不在用她的嘴唇宣扬所有这些判断。\par

13. \Quote{我喜爱你的规诫的道路，就如喜爱一切财富}（\BibleRef{咏 119:14}）。我们明白，没有什么比基督更快捷、更确实、更近、更高的天主规诫的道路，\Quote{在他内蕴藏着一切智慧和知识的宝藏}（\BibleRef{哥 2:3}）。因此他说，他在这条道路上所得的喜悦，就如获得一切财富一般。这些规诫就是他乐意向我们证明他多么爱我们的那些诫命……\par

14. \Quote{我要默想你的诫命，并注视你的道路}（\BibleRef{咏 119:15}）。因此，教会藉着博学者针对基督教及大公信仰的所有敌人所作的丰富辩论，操练自己于天主的诫命之中；这些辩论对于写作它们的人是有益的，只要在其中只注视主的道路；但\Quote{上主的一切道路}，正如经上所写，\Quote{是仁慈和真理}；这两者的丰满都在基督内寻见。通过这甜蜜的操练，也获得了随后所述：\Quote{我将在你的法令中默想，永不忘记你的话语}（\BibleRef{咏 119:16}）。\Quote{我将在其中默想}，为的是不忘记它们。正如第一篇圣咏中有福的人\Quote{昼夜默想上主的法律}……\par

% structure: p0018 source=h2 target=subsection reason=relative-heading-rank; base=h2

\subsection{Gimel}

15. 祂曾说过：\Quote{青年怎样洁净自己的行径？就是要遵守祢的言语。}看，如今祂更公开地请求帮助，为能这样做：\Quote{赏报祢的仆人，}祂说，\Quote{使我存活，并遵守祢的言语。}\BibleRef{咏 118:17}……祂所请求的正是这赏报，因为祂说：\Quote{赏报祢的仆人。}因为赏报有四种方式：或以恶报恶，如天主将以永火赏报不义的人；或以善报善，如祂将永远国度赏报义人；或以善报恶，如基督藉恩宠使不虔敬的人成义；或以恶报善，如犹达斯和犹太人因邪恶迫害基督。在这四种赏报方式中，前两种属于正义：以恶报恶，以善报善；第三种属于仁慈：以善报恶；第四种是天主所不知的，因祂绝不向任何人以恶报善。然而，我所列在第三位的，是首先必要的：因为除非天主以善报恶，就没有人可以领受祂以善报善。……\par

16. 因此，人的骄傲绝不可自高：天主将善赏报于祂自己的恩赐。……\par

17. \Quote{求你开启我的眼睛，使我默想祢法律的奇事。}\BibleRef{咏 118:18}祂接着说的\Quote{我在地上是客旅，}\BibleRef{咏 118:19}或如有些抄本所载\Quote{我在地上是寄居者，求祢不要向我隐瞒祢的诫命，}意义相同。……\par

18. 这里关于灵魂产生一个重要的问题。因为\Quote{我在地上是寄居者、客旅或异乡人}这句话，似乎不能就身体而言，因为身体源自大地。但在这个极为深奥的问题上，我不敢妄下定论。因为若这句话是针对灵魂而说（我们绝不认为灵魂来自大地），即\Quote{我是客旅}或\Quote{地上的异乡人}；或是针对整个人而言，因为他一度是乐园的居民，而说这话的人当时并不在乐园；或为更少争议，并非每个人都能说这话，而是那位被应许天上永恒家乡的人：我所知道的是，\Quote{人的生命在大地上是一种考验；}\BibleRef{约 7:1}\Quote{在亚当子孙身上有重重的轭。}\BibleRef{德 40:1}但我更愿意按照这样的解释来讨论问题：我们说我们是地上的寄居者或客旅，是因为我们已经找到了在上方的家乡，从那里我们领受了凭据，当我们抵达时，将永不再离开。……\par

19. 那些生活在天上的人，只要在此世寄居，就真是客旅。因此，愿他们祈祷，使天主的诫命不向他们隐藏，藉此他们能因爱慕天主而脱离这暂时的客居，与祂永远同在；并因爱近人而使他也能到他们将要去的那个地方。\par

20. 但若爱情本身不被爱，人又能藉爱而爱什么呢？因此，那位地上的客旅在祈求天主的诫命不向他隐藏（诫命全部或主要命令爱）之后，声明他渴望拥有对爱本身的爱，说：\Quote{我的灵魂切望时常渴慕祢的判断。}\BibleRef{咏 118:20}这种渴慕是值得赞美的，而非该受责罚。……\par

21. 但祂不仅说\Quote{切望}，还说：\Quote{我的灵魂切望渴慕祢的判断。}因为拥有天主的判断并无障碍，除非它们不被渴慕，而爱没有热情去赢得它们，尽管它们的光芒如此清晰明亮。……\par

22. \Quote{祢斥责了骄傲的人；偏离祢诫命的是该诅咒的。}\BibleRef{咏 118:21}因为骄傲的人偏离天主的诫命。因软弱或无知而不履行天主的诫命是一回事；因骄傲而偏离它们则是另一回事，正如那些在必死的状态中孕育我们进入这些邪恶的人所做的。……但请注意，在说了\Quote{祢斥责了骄傲的人}之后，祂并没有说\Quote{偏离祢诫命的是该诅咒的}，以便只想起第一个人的罪；而是说\Quote{偏离祢诫命而犯错的是该诅咒的。}因为需要让所有人都被那个例子吓住，以免偏离天主的诫命，并在今世的劳苦中藉着时时爱慕正义，重新获得我们在乐园的欢愉中所失去的。\par

23. \Quote{求祢使我免受羞辱和藐视，因为我曾寻求祢的训诲。}\BibleRef{咏 118:22}\Quote{训诲}在希腊文中称作\OriginalTerm{见证}{\GreekText{μαρτύρια}}，这个字我们现在用于拉丁文。因此，那些因对基督作证而遭受各种苦难、为真理奋斗至死的人，不称为testes，而用希腊语称作殉道者。既然你们听到这个更熟悉、更悦耳的词，就让我们将这句话视为这样说的：\Quote{求祢使我免受羞辱和藐视，因为我曾寻求祢的殉道。}当基督的身体这样说话时，它是否认为受到不虔敬和骄傲之人的斥责和羞辱是一种惩罚呢？因为它反而借此获得桂冠。那么，它为何祈求将这些作为沉重不堪的重担挪去呢？除非如我所说，它正是在为它的仇敌祈祷，因为它看到仇敌若把基督的圣名作为对基督徒的侮辱，是有害的。……因为我的仇敌，祢命令我爱他们，他们日益死去、丧失，当他们藐视祢的殉道并在谴责我身上的殉道时，若他们在我身上尊崇祢的殉道，他们确实将被唤醒生命并被寻回。事情就是这样：我们看到了。看哪，为了基督之名而殉道，在世人眼中和此世，不仅不是耻辱，而且是极大的荣耀；看哪，不仅在主眼前，也在世人眼前，\Quote{祂的圣人之死是珍贵的}；看哪，祂的殉道者非但不被鄙视，反而受到极大的尊荣。……\par

24. \BibleQuote{也曾有王公坐着议论我；但你的僕人默思你的法令。} \BibleRef{咏 118:23} 你如果想知道这是什么操练，就明白他接着所说的：\BibleQuote{你的见证是我的默想，你的法令是我的谋士。} \BibleRef{咏 118:24} 记住我先前教导你的：\OriginalTerm{见证}{testimonies}就是殉道的事迹。记住，在主的法令中，没有比爱仇人更困难、更令人钦佩的了。\BibleRef{玛 5:44} 因此，基督的身体就这样被操练，它默想那些为祂作证的殉道事迹，并爱那些因这些殉道而责备、轻视教会的人，教会因此受迫害……\par

% structure: p0029 source=h2 target=subsection reason=relative-heading-rank; base=h2

\subsection{达肋特}

25. \BibleQuote{我的灵魂紧贴地面：求你照你的话语使我活过来。} \BibleRef{咏 118:25} 他说\Quote{我的灵魂紧贴地面}是什么意思？……如果把整个世界看作一个大房子，我们看到天穹是它的拱顶，地面就是它的地板。因此，他愿意从属地的事物中被解救出来，并与宗徒一起说：\Quote{我们的家乡在天上。}所以，紧贴属地的事物是灵魂的死亡；与这种邪恶相反的是生命，他祈求生命时说：\Quote{求你使我活过来。}\par

26. ……身体本身，因为是出于土的，也合情合理地用\Quote{地面}一词来理解；因为身体仍是可朽的，且重压灵魂，\BibleRef{智 9:15} 我们义正词严地在其中叹息，对天主说：\Quote{求你使我活过来。}因为我们将来不会没有身体，当我们永远与主同在时，\BibleRef{得前 4:17} 但那时身体不再可朽，也不重压灵魂，严格说来，我们不会紧贴它们，而是它们紧贴我们，我们紧贴天主……\par

27. 因为他独自一人，在以下话中承认了他的罪：\BibleQuote{我承认了我的道路，你就俯听了我。} \BibleRef{咏 118:26} 有些抄本确实读作\Quote{你的道路}，但更多、最好的希腊文抄本读作\Quote{我的道路}，即邪恶的道路。因为他似乎对我说：我承认了我的罪，你就俯听了我；就是说，你赦免了它们。\BibleQuote{求你教导你的法令。} 我承认了我的道路；你抹去了它们；教导我你的道路。求你教导我，使我能够行动；而不仅仅是知道应该如何行动。因为就如圣经论到主，说祂不知罪，\BibleRef{格后 5:21} 这理解为祂没有犯罪；同样，真正被称为知道义德的，是那实行义德的人。这是一个正在进步的人的祈祷……\par

28. 最后他加上：\BibleQuote{求你使我熟悉你义德的途径} \BibleRef{咏 118:27} ；或者，如有些抄本所记：\Quote{求你教训我}；从希腊文更确切地表达为：\Quote{求你使我明白你义德的道理；如此我将在你的奇妙事中受操练。}他通过训诲渴望明白的这些更高诫命，称为天主的奇妙事。天主有一些如此奇妙的义德，以至于未尝试过的人可能以为人性软弱无法实行。因此，圣咏作者挣扎且疲惫于服从它们，说：\BibleQuote{我的灵魂因沉重而昏睡：求你用你的话坚固我！} \BibleRef{咏 118:28} \Quote{昏睡}是什么意思？就是说，他在曾经怀有的能够达到它们的希望中冷却了。但他接着说：\Quote{求你用你的话坚固我}，使我不会因昏睡而从那些我已感到已达成的职责中坠落；因此，求你用我已掌握的你的那些话坚固我，使我能够通过训诲达到其他职责。\par

29. \BibleQuote{求你从我这里除去不义的道路} \BibleRef{咏 118:29} 。因为既然行为之律法介入，使罪恶增多，\BibleRef{罗 5:20} 他加上：\BibleQuote{并按照你的法律怜悯我。} 按照什么法律，除非按照信德的法律？请听宗徒的话：\BibleQuote{那么，哪里还有夸耀之处？已被排除。通过什么法律？行为之律法吗？不，乃是藉着信德的法律。} \BibleRef{罗 3:27} 这就是信德的法律（\OriginalTerm{信德的法律}{law of faith}），我们藉此相信并祈祷，它因恩宠而赐予我们；使我们能成就我们凭自己无法成就的事；使我们不因不认识天主的义德，且试图建立自己的义德，而未能顺服天主的义德。\BibleRef{罗 10:3}\par

30. 但在他说了\Quote{求祢按祢的法律怜悯我}之后，他提到了他已获得的一些福分，以便祈求他尚未得到的其他福分。因为他说：\Quote{我拣选了真理的道路，祢的判断我未曾忘记}\BibleRef{咏 118:30}。\Quote{我紧紧依附祢的训诲；上主，求祢不要使我蒙羞}\BibleRef{咏 118:31}：愿我能坚持努力奔向我所奔跑的目标；愿我能抵达我所奔跑的地方！所以，\Quote{不在乎人愿意，也不在乎人奔跑，只在乎天主施恩。}\BibleRef{罗 9:16} 他接着说：\Quote{祢拓宽了我的心，我就要奔跑祢诫命的道路}\BibleRef{咏 118:32}。若祢没有拓宽我的心，我就不能奔跑。这话的意思在\Quote{我拣选了真理的道路，祢的判断我未曾忘记：我紧紧依附祢的训诲}这句中得到了清楚的解释。因为这次奔跑就是沿着天主诫命的道路。而他向上主陈述的，与其说是他自己的功德，不如说是祂的恩赐；就好像有人问他：你是如何通过拣选、不忘记天主的判断并依附祂的训诲来奔跑那条路的？你能靠自己做到这些事吗？他回答说，我不能。所以，这不是靠我自己的意志，好像它不需要祢的帮助一样；而是因为\Quote{祢拓宽了我的心。}心的拓宽就是我们喜爱正义。这是天主的恩赐，其效果是，我们不再因害怕惩罚而被祂的诫命束缚，而是因爱和喜爱正义而得以宽广。……\par

% structure: p0036 source=h2 target=subsection reason=relative-heading-rank; base=h2

\subsection{希}

31. 在这首伟大的圣咏中，接下来要处理且需要靠主的帮助来考虑和论述的是：\Quote{上主，求祢将祢律例的道路指示给我，我必始终寻求}\BibleRef{咏 118:33}。……\par

32. 为何此人仍祈求有一条律法为他规定；若没有为他规定，他就不可能怀着宽广的心奔跑天主诫命的道路？但既然说话的人是在恩宠上成长，并且知道他在恩宠上获益是天主的恩赐；那么他祈求为他规定律法，除了祈求他越来越获益之外，还祈求什么？就像你拿着一满杯水递给一个口渴的人；他通过喝来耗尽它，又通过仍然渴望它来祈求它。……\par

33. 但是\Quote{始终}是什么意思？……\Quote{始终}是否意味着只要我们活在世上，因为我们在此生中一直在恩宠中进步；但在此生之后，那在此生良好进步的人将在那里变得完美？在这里，天主的律法被不断探究，因为我们通过认知和爱慕它而进步；但在那里，它的圆满是让我们享用的，而非探究的。同样，这句话也应如此理解：\Quote{时常寻求祂的面容。}何处寻求？除了在这里？因为在将来我们\Quote{面对面地观看}\BibleRef{格前 13:12}时，我们不会再寻求天主的面容。或者，如果那以不变的情感被爱的事物被正确地说是被寻求的，我们唯一的目的是不失去它，那么我们确实将始终寻求天主的律法，即天主的真理：因为在这首圣咏中说了：\Quote{祢的律法是真理。}现在寻求它，是为了持守它；那时将持守它，是为了不失去它。……\par

34. \Quote{求祢赐我明悟，我必探究祢的律法，全心谨守}\BibleRef{咏 118:34}。因为当人探究了律法，并探究了其中包含全部意义的深奥之处时，他理应全心全意、全灵、全意爱天主，并爱邻人如己。\Quote{因为全部法律和先知都系于这两条诫命。}\BibleRef{玛 22:37} 当他说话时，似乎已承诺了此事：\Quote{是的，我必全心谨守。}\par

35. 但既然他连这件事也没有能力做到，除非得到那命令他做祂所命令的事的天主的帮助，他便补充说：\Quote{求祢引领我走祢诫命的道路，因为我的愿望就在那里}\BibleRef{咏 118:35}。我的喜爱是无力的，除非祢亲自引领我走向我所喜爱的地方。而这正是那条道路，即天主诫命的道路，他之前曾说，当主拓宽了他的心时，他已奔跑在这条路上。他称之为\Quote{道路}，因为\Quote{通往生命的门是窄的}\BibleRef{玛 7:14}；既然窄，若不怀着宽广的心，我们就不能在其中奔跑。……\par

36. 他接着说：\Quote{求祢使我的心倾向祢的训诲，而非贪财}\BibleRef{咏 118:36}。他这样祈求，是为了在意志本身上得以进步。……但宗徒说：\Quote{贪财是万恶的根源}\BibleRef{弟前 6:10}。然而在希腊文（这些词语正是从希腊文译成我们的语言）中，宗徒所用的词不是\OriginalTerm{贪婪}{\GreekText{πλεονεξία}}（这个词出现在圣咏的这段经文），而是\OriginalTerm{贪财}{\GreekText{φιλαργυρία}}，意即\Quote{爱钱}。但必须理解，宗徒用这个词是以种类指属类，即用爱钱来指普遍和一般的贪婪，这确实是万恶之根。\BibleRef{创 3:5}……因此，如果我们不使心倾向贪财，我们就只为天主而敬畏天主，这样祂就是我们事奉祂的唯一酬报。让我们爱祂本身，爱祂在我们内，爱祂在我们视为如同自己的邻人内——无论他们已拥有祂，还是为了使他们能拥有祂。……\par

37. 我们要解释的圣咏接下来的一句是：\Quote{求祢转开我的眼，免得注视虚幻；并求祢在祢的道路上赐我生命}\BibleRef{咏 118:37}。虚幻与真理彼此直接对立。这世界的欲望是虚幻；但基督，祂将我们从世界解救出来，是真理。祂也是此人希望在其中得生命的道路，因为祂也是生命：\Quote{我是道路、真理、生命，}\BibleRef{若 14:6} 这是祂亲口说的话。\par

38. …他祈求那使我们在考虑行动动机时所使用的眼睛得以转移，使它们不注视虚妄；也就是说，当他行任何善事的时候，他可以不看向虚妄作为他的动机。在这虚妄中，首要的是人对赞美的爱，由于这种爱，那些在此世被称为伟大的人成就了许多大事，他们在异教国家中备受赞誉，不是在主面前寻求光荣，而是在人中间寻求光荣，并因此表面上生活得明智、勇敢、节制和公正；当他们达到了这个，就得到了他们的酬报：虚妄的人，和虚妄的酬报。\BibleRef{玛 6:1} …此外，若为人的赞美而行善事便是虚妄，那么为得钱财、或增加钱财、或保留钱财，以及任何其他从外面来到我们这里的现世利益，又是何等虚妄呢？因为\Quote{万事都是虚妄：人在太阳之下劳碌经营，究竟有什么益处？}\BibleRef{训 1:2} 最后，我们甚至不应该为了我们现世的福利而行善事，而应该为了我们所盼望的永恒福利而行善事，在那里我们可以享受那不变的善，那是我们从主那里获得，甚至主自己就是我们的善。因为如果天主的圣人为了现世福利而行善事，那么基督的殉道者就永远不会在失去现世福利的情况下完成宣认的善工了。…\par

39. \Quote{求你使你的话在你的仆人身上坚定，使我敬畏你}\BibleRef{咏 119:38}。这岂不是：求你赐我能按你所说的而行？因为天主的话不会在那些以相反的行动将它从自身除去的人身上坚定；相反，它在那些坚定不移的人身上坚定。因此，天主将祂的话坚定在那些人身上，使他们敬畏祂，就是那些祂赐予敬畏祂之神的人；不是宗徒所说的那种恐惧：\Quote{你们没有领受使你们再陷于恐惧的奴仆之神}\BibleRef{罗 8:15}；因为\Quote{完美的爱驱除}这种\Quote{恐惧}\BibleRef{若一 4:18}，而是先知所说的\Quote{敬畏主的神}\BibleRef{依 11:2}；那种\Quote{纯洁且永远长存}的敬畏；那种害怕得罪所爱之主的敬畏。\par

40. \Quote{求你除去我所猜疑的羞辱，因为你的判断是甘甜的}\BibleRef{咏 119:39}。谁是那猜疑自己羞辱的人？谁不更清楚自己的羞辱胜过邻人的羞辱？因为一个人可能更猜疑他人而非自己；因为他不知道他所猜疑的；但在每个人自己的羞辱中，对他而言不是猜疑，而是知识，在其中良心发言。那么，\Quote{我所猜疑的责备}这句话是什么意思？它的意思必须从前一节推导出来；因为只要一个人不使他的眼睛转离，以免注视虚妄，他就在他人身上猜疑他自己里面所进行的事；从而他相信另一个人出于与自己相同的动机而崇拜天主或行善工。因为人们能看见我们所作的，但我们行动的最终目的是隐藏的。…\par

41. \Quote{看，我切慕了你的诫命：求你以你的公义使我生活}\BibleRef{咏 119:40}。看，我切慕全心、全灵、全意爱你，并爱邻人如己；但是\Quote{求你使我生活}不是在我自己的，而是\Quote{在你的公义}中，也就是说，以我所渴望的爱充满我。帮助我，使我能实行你所命令的：你亲自赐予你所吩咐的。\Quote{求你以你的公义使我生活}：因为在我自己里面，我有那导致我死亡的东西；但除了在你内，我找不到我得以生活的原因。基督就是你的公义，\Quote{祂由天主成为我们的智慧}等\BibleRef{格前 1:30}。在祂内，我找到了我所切慕的你的诫命，好使在你公义，即祂内，你使我生活。因为圣言本身就是天主；而且\Quote{圣言成了血肉}\BibleRef{若 1:14}，为叫祂自己也能成为我的邻人。\par

% structure: p0048 source=h2 target=subsection reason=relative-heading-rank; base=h2

\subsection{Vav}

42. \Quote{主，愿你的慈爱也降临到我们身上}\BibleRef{咏 119:41}。这句话似乎附属于前文：因为他不是说\Quote{让它降临到我}，而是\Quote{\Emph{也}让它降临到我。}…那么他在这里祈求什么？无非是，藉着他的慈爱（那发命令的），他可能履行他所切慕的诫命。因为他通过补充\Quote{甚至你的救恩，按你的话}（即按你的恩许），在某种程度上解释了他的意思。因此，宗徒愿意我们被理解为恩许的子女：\BibleRef{罗 9:8}使我们不致以为我们是什么是由于我们自己的行为，而是将一切都归于天主的恩宠。…基督本身就是天主的救恩，所以基督的整个身体可以说：\Quote{因天主的恩宠，我成为我今日的我}\BibleRef{格前 15:10}。\par

43. \Quote{我必要这样回答}，他说，\Quote{那些以言语斥责我的人}\BibleRef{咏 119:42}。但究竟是\Quote{以言语斥责我}还是\Quote{我将用言语回答}，是有疑问的；但两者都指向基督。那些认为被钉十字架的基督是绊脚石或愚拙的人，\BibleRef{格前 1:23}用祂来斥责我们；他们不知道\Quote{圣言成了血肉，居住在我们中间}\BibleRef{若 1:14}；那\Quote{在起初就有的}圣言，\Quote{与天主同在，且就是天主}\BibleRef{若 1:1}。但是，即使他们可能不以那为他们所不认识的圣言来斥责我们，因为祂的神性不为那些轻视祂在十字架上软弱的人所认识；我们仍要以圣言来回答，并且不要因他们的斥责而恐惧或困惑。因为\Quote{如果他们认识}圣言，\Quote{就绝不会钉死光荣之主了}\BibleRef{格前 2:8}。…因此，当圣咏作者说\Quote{我必要回答那些以言语斥责我的人}时，他立刻补充说\Quote{因为我的信赖在你的言语中}，这正是指你的恩许。\par

44. \Quote{O take not the word of Your truth away out of my mouth even exceedingly} \BibleRef{咏 118:43}。祂说，出自我口，因为说话的是身体的合一，在其肢体中也包括了那些在关键时刻因否认而跌倒，但后来藉补赎又复生的人，甚或藉更新他们的宣认，获得了已失去的殉道棕榈枝的人。因此，真理之言并未\Quote{even exceedingly}，或如某些抄本所载，even every way，即完全没有从\Person{伯多禄}口中被夺去，他是教会的预像；因为他虽在当时因恐惧而慌乱否认，却藉哭泣得以恢复（\BibleRef{玛 26:70}），并藉宣认之后获得了冠冕。因此，基督的整个身体在说话……接下来是\Quote{for I have hoped in Your judgments}。或如某些更严格地从希腊文译出的版本：\Quote{I have hoped more}；这个词虽以一种不太常见的方式复合，却满足了在翻译中传达真理的必要目的……看哪，圣人们和心谦的人，当他们信赖于祢时，在迫害中未曾跌倒；看哪，那些信赖自己而跌倒的人，却仍属于同一身体，当他们认识自己时哭泣了，并找到了祢的恩宠作为更坚固的支持，因为他们失去了自己的骄傲。\par

45. \Quote{So shall I always keep Your law} \BibleRef{咏 118:44}：就是说，若祢不将祢真理之言从我口中夺去。\Quote{Yea, unto age, and age of age:} 他表明他所谓的\Quote{alway}是何意。因为有时\Quote{alway}意即我们在此世生活之时；但这并非\Quote{unto age, and age of age}。如此翻译比某些抄本所载的\Quote{to eternity, and to age of age}更佳，因他们不能说，and to eternity of eternity。因此，该律法应被理解为宗徒所说的\Quote{爱是法律的满全}（\BibleRef{罗 13:10}）。因为这律法被圣人们所持守，真理之言未从他们口中被夺去，即基督的教会自身，不仅在此世，即直到此世终结，而且也为另一个被称为\Quote{无终世界}的世界所持守。……\par

46. \Quote{And I walked at liberty: for I sought Your precepts} \BibleRef{咏 118:45}……\Quote{And I walked at liberty}。此处并列连词\Quote{and}并非用作连接词；因为他并没有说，and I will walk，如他先前所说\Quote{and I will keep Your commandments for ever and ever}；或者若后一句是祈愿语气，and may I keep Your law，他并未加上And may I walk at liberty，好像他渴望并祈求这两件事；而是他说\Quote{And I walked at liberty}。若此连词未被使用，且句子独立于前文而引入\Quote{I walked at liberty}，读者绝不会因言谈方式之异常而想到应寻求某种隐义。无疑，他希望他所未说出的被理解，即他的祈祷已蒙垂听；然后他补充了他所成为的：好似他要说，当我为这些事祈祷时，祢垂听了，\Quote{And I walked at liberty}；同样地，他随后所添加的其余表达亦为此意。\par

47.……所以当他说了\Quote{And I walked at liberty}后，他附上了理由：\Quote{For I sought out Your commandments}。某些抄本未载\Quote{commandments}而是\Quote{testimonies}；但我们在大多数抄本中，尤其在希腊文中，发现是\Quote{commandments}；谁又愿怀疑更应相信这语言，它先于我们的语言，且这些圣咏由此译成了拉丁文？若我们欲知道他是如何寻求这些诫命的，或应如何寻求它们，让我们思想我们善师的话，祂既教导又赐予它们：\Quote{你们求，就必给你们}（\BibleRef{玛 7:7}）。稍后祂又说：\Quote{你们纵然不善，尚且知道把好东西给你们的儿女，何况你们在天堂的父，岂不更把好东西赐给求祂的人吗？}（\BibleRef{玛 7:11}）在此祂明确显示，祂所说的寻求、叩门、祈求的话，只属于祈求的恳切，即祈祷。此外，另一位福音作者并未说：祂将把好东西赐给求祂的人；这可以多种方式理解，或是属世的，或是属灵的祝福；而是排除了其他解释，并非常仔细地表达了主愿我们热切而持续祈求之事，即这些话：\Quote{何况你们的天父，岂不更将圣神赐给求祂的人吗？}……\par

48. \Quote{I spoke of Your testimonies also,} 他说，\Quote{before kings, and I was not ashamed} \BibleRef{咏 118:46}：作为一位已寻求并领受了恩宠的人，以答覆那些以言语责备他的人，并得应许真理之言不从其口被夺去。为这一真理奋斗直至死亡，甚至不在君王面前羞于谈论它。因为他所承认谈论的\Quote{testimonies}，在希腊文中称为\OriginalTerm{见证}{\GreekText{μαρτύρια}}，此词我们现今用以取代拉丁文。\Quote{殉道者}之名由此而来，耶稣曾预言他们要在君王面前承认祂（\BibleRef{玛 10:18}）。\par

49. \\BibleQuote{我默想了}，他说道，\\BibleQuote{祢的诫命，我所喜爱的} \\BibleRef{咏 118:47} 。\\BibleQuote{我也举起双手朝向祢的诫命，我所喜爱的} \\BibleRef{咏 118:48} ；或者，如一些抄本所读的，\\BibleQuote{我所极其喜爱的}，或\\BibleQuote{太多地}，或\\BibleQuote{热烈地}，正如他们选择用希腊词 \\OriginalTerm{极其}{\GreekText{σφόδρα}} 来表达。他之所以喜爱天主的诫命，是因为他行走在自由中；也就是说，藉着圣神，借着祂，\\BibleQuote{爱本身被倾注} \\BibleRef{罗 5:5} ，并扩大了信众的心。但他同时以思想和行动去爱。就思想而言，他说：\\BibleQuote{我默想了}；就行动而言，\\BibleQuote{我也举起双手}。但他对这两句话都加上了\\BibleQuote{我所喜爱的}这几个字，因为\\BibleQuote{诫命的终向是由纯洁的心发出来的爱} \\BibleRef{弟前 1:5} 。……下文\\BibleQuote{我默想了祢的规条}，与这两者相关。大多数译者更喜欢这个表达，而不是\\BibleQuote{我喜乐于}或\\BibleQuote{我谈论了}这两种翻译，有些人从希腊文 \\OriginalTerm{希腊文}{\GreekText{ἠ} \GreekText{δολ}} 译出的版本。\par

\SourceRecord{Translated by J.E. Tweed.；Nicene and Post-Nicene Fathers, First Series；Vol. 8.；Edited by Philip Schaff.；Buffalo, NY: Christian Literature Publishing Co.,；1888.；<http://www.newadvent.org/fathers/1801119.htm>.}

% structure: p0001 source=h1 target=section reason=page-title

\section{圣咏第一百二十篇释义}

1. 我们刚才听到并同声答唱的这篇圣咏很短，但极有益处。你们聆听时不会感到疲倦，劳作时也不会徒劳。因为根据其标题，它被称为\Quote{登阶之歌}。登阶可以是上升，也可以是下降。但在这篇圣咏中，所论的登阶是上升的……因此，在那梯子上既有上升的人，也有下降的人。\BibleRef{创 28:12}哪些人是上升的？那些朝向对属灵事物的理解而前进的人。哪些人是下降的？那些尽管在人力所及的范围内享受着对属灵事物的理解，却仍旧下降到婴孩中间，对他们说他们能够领受的话，以便他们在以奶喂养后，能够变得合适并强壮到足以承受属灵的食物。\par

2. 所以当一个人开始这样安排他的上升时；说得更明白些，当一位基督徒开始思想灵性的改进时，他便开始遭受反对者之舌的攻击。凡尚未遭受过这些的人，还未有进步；凡不遭受的人，甚至不愿努力改进。他愿意知道我们说的意思吗？让他同时体验一下人们对我们的传言。让他开始改进，让他开始愿意上升，愿意轻视地上的、易碎的、暂时的事物，把世俗的幸福视为虚无，只思念天主，不以获利为喜，不以损失为忧，甚至愿意变卖所有财产，分给穷人，并跟随基督；让我们看看他如何遭受诽谤者和不断反对者的舌头，以及——更大的危险——假装善意的劝告者，他们引他背离救恩……因此，那将要上升的人首先求天主抵御这些舌头：因为他说的：\Quote{我在患难中呼求了主，祂应允了我。}\BibleRef{咏 119:1}为什么祂应允了他？为了现在把他放在上升的台阶上。\par

3. \Quote{上主，求你救我的灵魂脱离不义的口唇和欺诈的舌头。}\BibleRef{咏 119:2}什么是欺诈的舌头？背信弃义的舌头，外表像劝告，实质是真正的祸害。这种人会说：你要做这件事吗？没有人做啊？你要成为唯一的基督徒吗？……有些人用劝阻来吓唬，另一些人则用赞美来阻挠。因为这种生活既然已经遍及世界一段时间，基督的权威如此之大，甚至连外邦人也不敢指责基督。无法受责备的基督被读诵。他们不能反对基督，不能反对福音，基督不能被指责；欺诈的舌头转而利用赞美作为障碍。如果你赞美，就应劝勉。为什么你以赞美来阻挠？……你转而采用另一种劝阻方式，用虚假的赞美使我远离真正的赞美；甚至用赞美基督来使我远离基督，说：这是什么呢？看，这些人做了这事；你也许做不到；你开始上升，你就坠落。它似乎在警告你；这是蛇，是欺诈的舌头，它有毒。如果你愿意上升，就向着它祈祷。\par

4. 你的主对你说：\Quote{要赐给你什么，或在你面前放置什么，以对付欺诈的舌头？}\BibleRef{咏 119:3}要赐给你什么，即作为武器来对抗欺诈的舌头，保卫自己免受欺诈之舌的伤害？\Quote{或在你面前放置什么？}祂问是为试探你，因为祂将回答自己的问题。祂顺着自己的询问回答说：\Quote{就是全能者的利箭，和使人荒凉的炭火。}\BibleRef{咏 119:4}使人荒凉的，或使人荒废的（因为不同抄本写法不同），是相同的，因为荒废，如你所见，容易导致荒凉。这些炭火是什么？首先，亲爱的弟兄们，要明白箭是什么。\Quote{全能者的利箭}是天主的话……那么，\Quote{使人荒凉的炭火}是什么？用言语来对付欺诈的舌头和不义的口唇是不够的；单单用言语辩护是不够的；我们必须也用榜样来辩护……因此，\Quote{炭火}这个词被用来表达许多罪人归向主的榜样。你听见人们惊奇地说：我认识那个人，他多么嗜酒，多么坏蛋，多么喜爱竞技场或圆形剧场，多么欺诈；现在他多么事奉天主，变得多么纯洁！不要惊奇；他是一块活炭。你为他活着而欢喜，你曾以为他死了而哀悼。但当你赞美活着的，如果你知道如何赞美，把他应用于死者，使他被点燃；任何人如果还缓慢地跟随天主，就应用那曾熄灭的炭火给他，并拥有天主之言的箭和荒废的炭火，这样你就能面对欺诈的舌头和说谎的口唇。\par

5. \BibleQuote{唉，我的寄居变得遥远！} \BibleRef{咏 119:5}。它已远离祢：我的朝圣之旅变得遥远。我尚未到达那没有恶人同住的家乡；我尚未到达那不怕冒犯的天使团体。但为何我尚未在那里？因为寄居就是朝圣之旅。寄居者乃是居住在外地、而非自己家乡的人。何时才算遥远？有时，我的弟兄们，当人远行时，他居住的人可能比他自己家乡的人更好：但当我们远离那天上的耶路撒冷时，情况并非如此。因为人变换了家乡，这异乡的寄居有时对他有益；在旅途中他找到了忠信的朋友，那是他在家乡找不到的。他因仇敌而被迫离开家乡；当他旅行时，他找到了在家乡没有的东西。但那耶路撒冷家乡并非如此，那里全是善人：凡从那里离开的，就处于恶人之中；除非他回到天使团体，回到他旅行前所在的地方，否则他无法离开恶人。在那里，所有人都是义人和圣人，他们无需阅读、无需文字就享受天主圣言：因为通过书页写给我们的，他们透过天主的圣容领悟。那是怎样的家乡！实在是一个伟大的家乡，而离开那家乡的流浪者真是可怜。\par

6. 但他所说的 \Quote{我的朝圣之旅变得遥远}，是那些人的话，即在这世上劳苦的教会本身的话。这是她的声音，在另一篇圣咏中从地极呼喊说：\BibleQuote{我从地极向祢呼号}……那么，他在何处呻吟？与谁同住？\BibleQuote{我曾在刻达尔的帐幕中居住。} 既然这是一个希伯来文词，无疑你们不懂。\Quote{我曾在刻达尔的帐幕中居住} 是什么意思？就我们对希伯来文词语的解释所记得的，\Quote{刻达尔} 意为黑暗。\Quote{刻达尔} 译成拉丁文称为 \Emph{tenebrae}。你们知道，\Person{亚巴郎}有两个儿子，宗徒曾提及他们，并宣称他们是两个盟约的预像……因此 \Person{依市玛耳} 在黑暗中，\Person{依撒格} 在光明中。凡在教会中也从天主寻求现世福乐的人，将属于 \Person{依市玛耳}。这些正是反对那些在进步的属灵人、诋毁他们、并有诡诈的舌头和不义嘴唇的人。圣咏作者在上升时曾为此祈祷，而全能者赐给他烧灭的炭火和迅疾锋利之箭作为防御。因为他在这些人中仍然生活，直到整个禾场被簸扬：因此他说：\BibleQuote{我曾居住在刻达尔的帐幕中。} \Person{依市玛耳}的帐幕被称为刻达尔的帐幕。\WorkTitle{创世纪}如此记载：刻达尔属于 \Person{依市玛耳}。\BibleRef{创 25:13} 因此 \Person{依撒格}与 \Person{依市玛耳}同在：即属于 \Person{依撒格}的人生活在属于 \Person{依市玛耳}的人中间。\BibleRef{迦 4:29} 这些人想上升，那些人想压制他们：这些人想飞向天主，那些人则尽力拔去他们的翅膀……\par

7. \BibleQuote{我的灵魂徘徊了许多} \BibleRef{咏 119:6}。免得你以为是指身体的徘徊，他说灵魂徘徊。身体在地方上徘徊，灵魂在其情感中徘徊。如果你爱大地，你就远离天主；如果你爱天主，你就上升归于天主。让我们在爱天主和爱近人中得到操练，好能回到爱德。如果我们坠向大地，我们就会枯萎腐朽。但有一位下降到那已坠落者那里，为使他能起来。他说到徘徊之时，他曾在 \OriginalTerm{刻达尔}{Kedar} 的帐幕中徘徊。为何？因为 \Quote{我的灵魂徘徊了许多}。他在他上升之处徘徊。他不是在身体中徘徊，他也不是在身体中上升。但他上升在哪里呢？他说：\BibleQuote{上升在心中}。\par

8. \BibleQuote{我与那恨和平的人和平相处} \BibleRef{咏 119:7}。但无论你们如何听，极亲爱的弟兄们，你们将不能证明你们歌唱得多么真实，除非你们已开始实行你们所歌唱的。无论我如何说，无论我用何种方式解释，无论我用何种言语转换，它都无法进入那未实行者的心中。开始行动，看看我们所说的是什么。然后每句话都会流泪，然后圣咏被咏唱，而心灵践行圣咏中所唱的……谁是恨和平的人？是那些撕裂合一的人。因为他们若不恨和平，就会留在合一之中。但他们分离了，正是为此，他们可能是义人，好使不虔诚的人不与他们混合。这些话或是我们的，或是他们的：请判断谁是谁。天主教会说：合一不可丢失，天主的教会不可被割裂。天主以后会审判恶人和善人……我们也这样说：你们要爱和平，你们要爱基督。因为若他们爱和平，他们就爱基督。因此当我们说爱和平时，我们就是说爱基督。为什么？因为宗徒论及基督说：\BibleQuote{祂是我们的和平，使双方合而为一。} \BibleRef{弗 2:14} 如果基督因此是和平，因为祂使两者合而为一：那么你们为何使一成为二？如果当基督使二成为一时，你们却使一成为二，你们如何是和平的缔造者？但既然我们这样说，我们就是与那恨和平的人和平相处；然而，那恨和平的人，当我们对他们说话时，却无故向我们开战。\par

\SourceRecord{Translated by J.E. Tweed.；Nicene and Post-Nicene Fathers, First Series；Vol. 8.；Edited by Philip Schaff.；Buffalo, NY: Christian Literature Publishing Co.,；1888.；<http://www.newadvent.org/fathers/1801120.htm>.}

% structure: p0001 source=h1 target=section reason=page-title

\section{《圣咏第一百廿一篇释义》}

1. ……愿他们\Quote{举目向山，期待从那里获得救助}\BibleRef{咏 120:1}。所谓\Quote{山已被光照}是什么意思？正义的太阳已经升起，福音已由宗徒宣讲，圣经已宣讲，一切奥秘已揭开，帐幔已撕裂，圣殿的隐秘处已显露：愿他们现在终于举目向山，从那里得到他们的救助……\Quote{从他圆满的恩宠中，我们都领受了}\BibleRef{若 1:16}，他说。因此，你的救助来自他，来自那\Quote{山}所领受的圆满，而非来自\Quote{山}；然而，除非你通过圣经举目向山，你就不会接近，以便被他光照。\par

2. 因此，请歌唱接下来的内容；如果你愿意听如何最稳妥地将脚放在台阶上，以便在攀登时不疲倦、不绊倒、不跌倒，就请这样祈祷：\Quote{求你不要让我的脚动摇！}\BibleRef{咏 120:3}。是什么使脚动摇？是什么使那曾在乐园中的人的脚动摇？但首先要思考，是什么使那曾在天使中的人的脚动摇：当他的脚动摇时，他跌倒了，从天使变成了魔鬼：因为当他的脚动摇时，他跌倒了。寻求他为何跌倒：他因骄傲而跌倒。那么，除了骄傲，没有东西能使脚动摇；除了骄傲，没有东西能使脚跌倒。爱德使脚行走、进步和上升；骄傲使脚跌倒……因此，圣咏作者听到如何上升而不跌倒，便向天主祈祷，愿他从苦难之谷得益，而不在骄傲的膨胀中失败，他如此祈祷：\Quote{求你不要让我的脚动摇！}天主回答他说：\Quote{那保护你的，不要打盹。}请留意，我所爱的。这就像是两个句子表达同一个思想：人在上升并歌唱\Quote{登阶之歌}时说：\Quote{求你不要让我的脚动摇}；天主仿佛回答说：你对我说\Quote{不要让我的脚动摇}，那么你也说\Quote{那保护你的，不要打盹}，这样你的脚就不会动摇。\par

3. 为你自己选择那位既不困倦也不睡眠的，你的脚就不会动摇。天主从不睡觉：如果你希望有一位从不睡眠的守护者，就选择天主作为你的守护者。\Quote{求你不要让我的脚动摇}你说得好，非常好；但他也对你说：\Quote{那保护你的，不要打盹。}你或许曾打算转向人作为你的守护者，并说：我能找到谁不睡觉？哪个人不打盹？我能找到谁？我往哪里去？我往哪里回？圣咏作者告诉你：\Quote{那保护以色列的，既不困倦也不睡眠}\BibleRef{咏 120:4}。你希望有一位既不打盹也不睡眠的守护者吗？看，\Quote{那保护以色列的，既不困倦也不睡眠}：因为基督保护以色列。那么，你应是以色列。\Quote{以色列}是什么意思？它被解释为\Quote{看见天主}。如何看见天主？先是借着信德，后来借着目睹。如果你还不能借着目睹看见他，就借着信德看见他……谁既不打盹也不睡眠？当你在人中间寻找时，你被欺骗了；你永远找不到。那么，不要信赖任何人：每个人都会打盹，也会睡眠。他何时打盹？当他带着软弱的肉身时。他何时睡眠？当他死去时。所以不要信赖人。凡有死的人都会打盹，他在死亡中睡眠。不要在人中寻找守护者。\par

4. 你问，谁将帮助我，除了那位不打盹也不睡眠的？请听接下来的话：\Quote{主亲自保护你}\BibleRef{咏 120:5}。所以，不是会打盹会睡眠的人，而是主保护你。他如何保护你？\Quote{主在你的右手边保护你}……在我看来，这话含有隐藏的意义；否则他本会简单地说\Quote{主将保护你}，而不加\Quote{在你的右手边}。为什么？难道主只保护我们的右手，而不保护左手？他不是创造了我们的整体吗？那创造我们右手的，不也创造了左手吗？最后，如果他认为只提右手就够了，为什么他说\Quote{在你的右手边}，而不直接说\Quote{在你的右手}？为什么这样说，除非他在这里隐藏了某些东西，需要我们敲门才能得到？因为他要么说\Quote{主将保护你}，不再加什么；要么如果他想加右手，他会说\Quote{主将保护你的右手}；或者至少，既然他加了\Quote{手}，他会说\Quote{主将保护你的手，甚至你的右手}，而不是\Quote{在你的右手边}……\par

5. 我问你们，你们如何解释福音中所说的：\Quote{不要叫左手知道右手所行的}？\BibleRef{玛 6:3} 因为如果你们明白这一点，就会发现什么是你们的右手，什么是你们的左手；同时你们也会明白，天主造了这两只手，左手和右手；然而左手不应知道右手所行的。我们的左手是指我们暂时拥有的一切；我们的右手是指我们的主所应许给我们的一切不变和永恒的事物。但是，如果那赐予永生者，也藉这些暂时的祝福安慰我们今世的生命，那么他自己就造了我们的右手和左手……\par

6. 现在让我们来看圣咏的这一节：\Quote{主在你的右手边保护你}\BibleRef{咏 120:5}。他所说的\Quote{手}是指能力。我们如何证明这一点？因为天主的能力也被称为天主的手……关于这一点，若望说：\Quote{他赐给他们权能，成为天主的子女}\BibleRef{若 1:12}。你从哪里领受了这权能？\Quote{就是那些信他名字的人}他说。因此，如果你相信，这权能就赐给了你，使你成为天主的子女。但是成为天主的子女，就是属于右手。所以，你的信德就是你右手的手：也就是说，赐给你的权能，使你成为天主的子女，就是你右手的手……\par

\BibleQuote{愿主在你的右手上庇护你}\BibleRef{咏 120:6}。我曾说过，并且我相信你们已经认出来了。因为若不是你们认出来了，并且从圣经中认出来了，你们就不会用声音表示你们明白了。既然你们明白了，诸位弟兄，请思考接下来的话：为什么主在\BibleQuote{你右手的手上}庇护你，也就是说，在你的信德中，我们藉此接受了\BibleQuote{成为天主子女的权能}，并且站在祂的右边；那么天主为什么要庇护我们？为了过犯。过犯从何而来？过犯应从两方面畏惧，因为有两条诫命，全法律和先知都系于这两条：爱天主和爱近人。\BibleRef{玛 22:37}教会因近人的缘故而受爱，但天主因祂自己的缘故而受爱。关于天主，以太阳作比喻；关于教会，以月亮作比喻。凡能犯错的人，对天主怀有不当作的想法，不相信圣父、圣子和圣神是同一本体\OriginalTerm{本体}{substantia}，就被异端者，主要是亚略派的诡计所欺骗。若他对圣子或圣神有所轻视，不如对圣父，他就因天主而受了过犯；他被太阳灼伤了。凡以为教会只存在于一个省份，而不普及全世界的人，以及相信那些说\BibleQuote{在这里}和\BibleQuote{在那里，有基督}的人，正如你们刚才听到福音诵读时所说的；因为祂付出了如此大的代价，赎买了全世界：此人可以说是在近人身上受了过犯，并被月亮灼伤。因此，凡在真理本体上犯错的人，被太阳灼伤，并在白日被灼伤；因为他在智慧本身中犯了错。天主因此造了一个太阳，它照耀善人也照耀恶人，那太阳善人和恶人都能看见；但那个太阳是另一个，非受造的，非生成的，万物藉祂而造；那里有不变真理的理智。不敬虔的人说：\BibleQuote{太阳没有照在我们身上。}凡在智慧本身中不犯错的人，不被太阳灼伤。凡在教会和主的肉体中不犯错的人，以及在我们时代为我们所做的那些事上不犯错的人，不被月亮灼伤。但每一个人，即使相信基督，也会在这方面或那方面犯错，除非在他身上实现了此处所祈求的：\BibleQuote{主是你的保护，在你右手的手上。}他接着说：\BibleQuote{这样，太阳在白日不会灼伤你，月亮在夜间也不会灼伤你。}\BibleRef{咏 120:6}因此，你的保护在你右手的手上，是为了这个缘故，使太阳在白日不灼伤你，月亮在夜间也不灼伤你。诸位弟兄，请由此理解为比喻性的说法。因为实际上，如果我们想到可见的太阳，它在白日灼伤；月亮在夜间灼伤吗？但什么是灼伤？过犯。请听宗徒的话：\BibleQuote{谁软弱，我不软弱呢？谁跌倒，我不焦急呢？}\BibleRef{格后 11:29}\par

\BibleQuote{因为主要保护你免受一切邪恶。}\BibleRef{咏 120:7}从太阳的过犯，从月亮的过犯，从一切邪恶中，祂将保护你，祂是你的保护在你右手的手上，祂不打盹也不睡觉。原因何在？因为我们处于诱惑之中：\BibleQuote{主要保护你免受一切邪恶。主保护你的灵魂：}甚至你的灵魂本身。\BibleQuote{主保护你的出入，从今时直到永远。}\BibleRef{咏 120:8}不是你的身体；因为殉道者们在身体上被消耗了；但是\BibleQuote{主保护你的灵魂}；因为殉道者们没有交出他们的灵魂。迫害者们曾对克利斯皮纳发怒，今天我们庆祝她的庆日；他们曾对一个富有的娇弱妇人发怒；但她很坚强，因为主是她右手上的保护。祂是她的保护者。我的弟兄们，在非洲有谁不认识她呢？因为她极其显赫，出身高贵，财富丰裕；但所有这些都在她的左手之下，在她的头下。敌人前来击打她的头，左手指给了敌人，那左手在她的头下。她的头在上，右手从上面拥抱了她。……\par

\SourceRecord{Translated by J.E. Tweed.；Nicene and Post-Nicene Fathers, First Series；Vol. 8.；Edited by Philip Schaff.；Buffalo, NY: Christian Literature Publishing Co.,；1888.；<http://www.newadvent.org/fathers/1801121.htm>.}

% structure: p0001 source=h1 target=section reason=page-title

\section{《圣咏第一二二篇讲道》}

污秽的爱焚烧心神，将注定丧亡的灵魂引向对尘世之物的贪欲，使牠追求朽坏之物，将牠推入至低之处，沉入深渊；同样，圣洁的爱提升我们达于天上之事，燃起我们对永恒之物的渴望，激励灵魂趋向那些永不消逝、永不死灭的事物，从地狱的深渊将我们提升至天堂。然而，一切爱都有其自身的力量，爱者的灵魂中不能有闲散的爱；爱必然牵引灵魂。但你愿意知道爱是怎样的吗？请看它引向何处……\par

这篇圣咏是一首\Quote{\OriginalTerm{登阶之歌}{Song of degrees}}，正如我们常对你们说的，因为这些阶不是下降的阶，而是上升的阶。因此他渴望上升。他渴望上升到何处？岂不是升到天堂吗？\Quote{升到天堂}是什么意思？难道他渴望升上去与太阳、月亮、星辰同在吗？断乎不是！而是在天堂里有永恒的耶路撒冷，那里有我们的同胞——天使；我们在世上相对于这些同胞是流浪者。我们在旅途中叹息，我们将在那城里欢欣。但我们在旅途中找到了一些同伴，他们已经亲眼见过这座城，他们召唤我们向它奔跑。说这话的人也以这些同伴为喜乐：\Quote{我因那些对我说\InnerQuote{我们要进入主的殿宇}的人而欢喜。}\BibleRef{咏 121:1}\par

\BibleQuote{我们的脚已站在耶路撒冷的庭院中}\BibleRef{咏 121:2}……想想你将在那里成为什么；尽管你尚在旅途，却要将此置于眼前，仿佛你已经站立，仿佛你已在天使中不停地欢欣；仿佛那经上所载的已在你身上实现：\Quote{住在祢殿中的人是有福的，他们必不断赞美祢。}\BibleQuote{我们的脚已站在耶路撒冷的庭院中。}哪一个耶路撒冷？这地上的耶路撒冷也常被称为这个名称；尽管这耶路撒冷只是那一个的影儿。站在这耶路撒冷中算得了什么大事？因为这耶路撒冷未能站立，已经沦为废墟了。那么，圣神从热爱的圣咏作者心中说出这话，难道是当作一件大事吗？岂不是那耶路撒冷，主曾对她说：\BibleQuote{耶路撒冷，耶路撒冷，你杀害先知，}\BibleRef{玛 23:37}等等。那么他渴望什么大事呢？是要站在那些杀害先知、用石头打死那些被派来的人中间吗？愿天主禁止他想到那个耶路撒冷——那位如此热爱、如此燃烧、如此渴望到达那耶路撒冷、\BibleQuote{我们的母亲}\BibleRef{迦 4:26}、宗徒说她是\BibleQuote{在天上永恒的}\BibleRef{格后 5:1}的人。\par

\BibleQuote{耶路撒冷被建筑，好似一座城}\BibleRef{咏 121:3}。弟兄们，当达味说出这些话时，那座城已经完工，不再被建筑了。所以他说的是一座正在被建筑的城，活石藉着信德奔向它，伯多禄说：\BibleQuote{你们也就像活石，被建成一座属神的殿宇。}\BibleRef{伯前 2:5}即天主的圣殿。\Quote{你们被建成活石}是什么意思？你们若相信，就活着；但你们若相信，就成为天主的圣殿；因为保禄宗徒说：\BibleQuote{天主的圣殿是圣的，这圣殿就是你们。}\BibleRef{格前 3:17}因此这座城现在正在建筑中；石头由宣讲真理者的手从山丘上砍下，被凿平，以便嵌入永恒的建筑物中。还有许多石头仍在建筑者的手中：愿它们不从祂的手中掉落，好让它们被完美地建成圣殿的结构。这就是\BibleQuote{被建筑、好似一座城的耶路撒冷}：基督是它的根基。保禄宗徒说：\BibleQuote{除已奠立的根基，即耶稣基督外，任何人不可再奠立别的根基。}\BibleRef{格前 3:11}根基奠在地上时，墙壁在上面被建造，墙壁的重量倾向最低之处，因为根基奠在底部。但如果我们的根基在天上，就让我们被建向天上吧。身体建筑了这座巴西利卡的殿宇，你们看到了它宽敞的规模；既然身体建筑了它，就把根基放在最低处；但我们既然被属神地建造，我们的根基就放在最高处。因此让我们奔向那里，好让我们被建造……但我说的是哪一个耶路撒冷？他问，难道是你所看见、耸立在墙壁结构中的那个吗？不，而是\BibleQuote{被建筑、好似一座城的耶路撒冷}。为什么不说\Quote{一座城}，而说\Quote{好似一座城}？只因那些在耶路撒冷如此建造的墙壁是一座可见的城，正如众人按字义所称的一座城；但这城被建筑\Quote{好似一座城}，因为进入其中的人如同活石；因为他们并不是字义上的石头。正如他们被称为石头，却并非如此；同样，被称为\Quote{好似一座城}的城并不是一座城；因为他说\Quote{被建筑}。他用\Quote{建筑}一词，意指身体和墙壁的结构与紧密连接。因为\Quote{城}是正确指居住在那里的人。但他说\Quote{被建筑}，向我们表明他指的是一座城镇。既然属神的建筑与身体建筑有某种相似，因此它\Quote{好似一座城被建筑}。\par

5. 但以下的话应除去一切疑惑，使我们不再按肉体理解这句话：\Quote{他们的分享是在那\Quote{自有者}内。}… \Quote{那\Quote{自有者}}是什么意思？那就是永远同一不变的；不是现在这样、以后又那样。那么，什么是\Quote{那\Quote{自有者}}，不就是那\Quote{是}者吗？什么是那\Quote{是}者？那就是永恒者……请看\Quote{那\Quote{自有者}：\InnerQuote{我是自有者}，\InnerQuote{我是}}。你无法理解；理解它已是难事，领悟它更是难事。请记住那你所不能领悟的、为你而成为的那一位。记住基督的肉身，当你因强盗的伤害而半死时，你被扶起带到客栈，在那里得到救治。因此，让我们奔向主的殿，抵达那使我们的脚可以站立的城，那城\Quote{被建造如同城邑：它的分享是在那\Quote{自有者}内。}……\par

6. 那\Quote{分享于那\Quote{自有者}}的城，分享于它的稳定；因此，奔向那里的人得以分享它的稳定，这是公义的。因为那里的一切都站立不动，没有东西经过。你也愿意站在那里而不经过吗？奔向那里吧。没有人能靠自己拥有那\Quote{自有者}。\par

7. \Quote{因为在那里，支派上去了}（\BibleRef{咏 121:4}）。我们曾问：那跌倒的人要上到哪里去？因为我们说过，这是一个正在上升的人的声音，是正在升起的教会的声音。我们能说出它升到哪里去吗？它去哪里？它被提升到哪里？他说：\Quote{那里，支派上去了。}哪里？\Quote{分享于那\Quote{自有者}}。但什么是\Quote{支派}？许多人知道，许多人不知道。因为如果我们按本义使用\Quote{库里亚}一词，我们除了每个城市中存在的\Quote{库里亚}之外，什么都不了解，\Quote{库里亚}一词衍生出\Quote{库里亚利}和\Quote{德库里奥内}，即库里亚或德库里亚的公民；你们知道每个城市都有这样的库里亚。但那些城市中曾经有或现在有人民的库里亚，一个城市有许多库里亚，如罗马有三十五个人民的库里亚。这些被称为\Quote{支派}。以色列人民有十二个这样的支派，依照雅各伯的儿子们。\par

8. 以色列人民有十二个支派：但他们中有好的，也有坏的。因为那些钉死我们主的支派是多么邪恶！那些承认主的支派是多么良善！因此，那些钉死主的支派是魔鬼的支派。所以当这里说：\Quote{因为在那里，支派上去}时，为使你不致理解为所有支派，他加上了：\Quote{甚至上主的支派。}……什么是上主的支派？\Quote{为以色列作证}。弟兄们，请听这是什么意思。\Quote{为以色列作证}：即藉此可以知道它是真正的以色列……他就是这样的人：在他内毫无诡诈。当主看见纳塔乃耳时，他说了什么？\Quote{看，这确是一个以色列人，在他内毫无诡诈。}（\BibleRef{若 1:47}）因此，如果真正的以色列人是在他内毫无诡诈的人，那么那些在他内毫无诡诈的支派就上到耶路撒冷。……他们为什么上去？\Quote{为承认你的名，上主。}没有比这更高贵的表达了。正如骄傲妄自尊大，同样谦卑承认。正如妄自尊大者希望显得与自己实际不同，同样那承认者不愿被人看见自己的本来面目，而爱那\Quote{是}者。因此，在他内毫无诡诈的以色列人上去，因为他们是真正的以色列人，因为他们内有以色列的见证。\par

9. \Quote{因为在那里设立了审判的座位}（\BibleRef{咏 121:5}）。这是一个奇妙的谜题，一个奇妙的问题，若不被理解。他称这些座位为希腊人所称的\Quote{宝座}。希腊人尊称椅子为\Quote{宝座}。因此，我的弟兄们，我们即使坐在座位上或椅子上，这并不奇怪；但是这些座位本身要坐，我们何时能理解呢？就好像有人说：让凳子或椅子坐在这里。我们坐在椅子上，坐在座位上，坐在凳子上；座位本身并不坐。那么\Quote{因为在那里设立了审判的座位}是什么意思？……因此，如果天是天主的座位，而宗徒们是天，那么他们已成了天主的座位，天主的宝座。在另一处经文中说：\Quote{义人的灵魂是智慧的宝座。}这是伟大的真理，伟大的真理被宣告：智慧的宝座是义人的灵魂；也就是说，智慧坐在义人的灵魂中，如同坐在她的椅子上，她的宝座上，然后审判她所审判的一切。因此，那里有智慧的宝座，因此主对他们说：\Quote{你们要坐在十二个宝座上，审判以色列十二个支派}（\BibleRef{玛 19:28}）。因此，他们也要坐在十二个座位上，而他们自己就是天主的座位；因为关于他们，经上说：\Quote{因为在那里设立了座位。}谁坐？\Quote{座位}。谁是那些座位？就是那些关于他们说\Quote{义人的灵魂是智慧的座位}的人。谁是那些座位？诸天。谁是诸天？天。什么是天？就是主所说的：\Quote{天是我的宝座}（\BibleRef{依 66:1}）。因此，义人自己就是座位；他们也有座位；而座位将在那耶路撒冷被安置。为什么目的？\Quote{为审判}。他说：你们要坐在十二个宝座上，你们这些宝座，审判以色列十二个支派。审判谁？那些在地上处于下方的。谁将审判？那些已成为天的人。但那些将被审判的人将分为两群：一群在右边，另一群在左边……\par

10. 他随即加上，如同对座位本身说：\BibleQuote{你们要探求耶路撒冷和平的事}\BibleRef{咏 121:6}。哦，你们这些座位，如今坐着审判，并成为审判之主的座位（因为审判者探求，被审判者被探求），他说：\BibleQuote{你们要探求，}他说，\BibleQuote{耶路撒冷和平的事。}他们询问会发现什么？发现有些人行了爱德的事，有些人没有。他们若发现有人行了爱德的事，便召他们到耶路撒冷；因为这些事是\BibleQuote{为耶路撒冷的和平}。爱德是强大的，我的弟兄们，爱德是强大的。你们愿意看看爱德有多强大吗？……如果爱德缺乏手段，以致找不到可施予穷人的，就让它爱吧：让它给出一\BibleQuote{杯凉水}\BibleRef{玛 10:42}；这将被记在他的账上，如同匝凯\BibleQuote{献给穷人一半家产}\BibleRef{路 19:8}。为何如此？一个给得这么少，另一个给得那么多，而前者也要被算得那么多吗？正是那么多。因为尽管他们的资源不等，但他们的爱德是不等的。\par

11. ……他接着说：\BibleQuote{还有丰裕，}给爱你的人。他转向\Place{耶路撒冷}本身说：爱你的人就有丰裕。匮乏之后的丰裕：这里他们贫乏，那里他们富足；这里他们软弱，那里他们刚强；这里他们缺乏，那里他们富有。他们如何变得富有？因为他们在此地献出了从天主暂时领受的，在那地领受了天主此后永久偿还的。这里，我的弟兄们，连富人也贫穷。富人若承认自己贫穷，那是好事：因为他若自以为饱足，那只是虚胀，不是丰裕。让他承认自己空虚，好被充满。他有什么？金子。他还缺什么？永生。让他想想自己有什么，看看自己没有什么。弟兄们，从他所有的，让他施舍，好领受他所没有的；让他用他所有的购买他所没有的，\BibleQuote{还有丰裕给爱你的人。}\par

12. \BibleQuote{愿和平在你的力量中}\BibleRef{咏 121:7}。哦，\Place{耶路撒冷}，哦，城市，你正被建成一座城，你的分享在于\OriginalTerm{同一者}{The Same}：\BibleQuote{愿和平在你的力量中：}和平在你的爱中；因为你的力量就是你的爱。请听\BibleBook{雅}{歌}：\BibleQuote{爱情如死亡般坚强}\BibleRef{歌 8:6}。这是个伟大的说法，弟兄们，\BibleQuote{爱情如死亡般坚强。}爱德的力量无法用比这更伟大的言辞表达，\BibleQuote{爱情如死亡般坚强。}因为谁能抵抗死亡，我的弟兄们？想一想，我的弟兄们。火、浪、剑都可抵抗：我们抵抗掌权者，抵抗君王；死亡独自来临，谁能抵抗它？没有比它更有力的。因此，爱德被比作它的力量，用这些话：\BibleQuote{爱情如死亡般坚强。}既然这爱杀死我们过去的所是，好使我们成为我们所不是的；爱在我们内创造了一种死亡。这死亡是那说\BibleQuote{世界于我已被钉在十字架上，我于世界也已钉在十字架上}\BibleRef{迦 6:14}的人所经历的；这死亡是那些他对他们说\BibleQuote{你们已经死了，你们的生命与基督一同藏在天主内}\BibleRef{哥 3:3}的人所经历的。爱情如死亡般坚强……\par

13. 因此，当他谈论爱德时，他接着说：\BibleQuote{为了我的弟兄和同伴的缘故，我论及你们说了和平}\BibleRef{咏 121:8}。哦，\Place{耶路撒冷}，你这分享在于\OriginalTerm{同一者}{The Same}的城市，我在这生命和这地上，我贫穷，他说，我作为一个陌生人而叹息，尚未完全享受你的和平，却宣讲你的和平；我不为自己的缘故宣讲，如同异端者，他们寻求自己的光荣，说：和平与你们同在；却没有他们所宣讲给人民的和平。因为他们若有和平，就不会撕裂合一。他说：\BibleQuote{我，}他说，\BibleQuote{论及你们说了和平。}但为何？\BibleQuote{为了我的弟兄和同伴的缘故：}不为我自己的荣誉，不为我自己的钱财，不为我自己的生命；因为，\Quote{为我，生活是基督，死亡是利益。}但是，\BibleQuote{我论及你们说了和平，是为了我的弟兄和同伴的缘故。}因为他渴望离去，与基督同在；但因他必须向同伴和弟兄宣讲这些事，他接着说，存留在肉身内对你们更为需要。\par

14. \BibleQuote{因上主我天主的殿宇，我为你寻求了美善}\BibleRef{咏 121:9}。不是为我自己寻求美善，因为那时我不应为你们寻求，而是为我自己；如此我便不会拥有它们，因为我不会为你们寻求；但是，\BibleQuote{因上主我天主的殿宇，}因教会，因圣人，因朝圣者；因穷人，使他们能上去；因为我们对他们说，我们要进入上主的殿宇：因上主我天主殿宇本身，我为你寻求了美善。这些冗长而必要的言语，你们要收集，弟兄们，吃它们，喝它们，变得强壮，奔跑，夺取。\par

\SourceRecord{Translated by J.E. Tweed.；Nicene and Post-Nicene Fathers, First Series；Vol. 8.；Edited by Philip Schaff.；Buffalo, NY: Christian Literature Publishing Co.,；1888.；<http://www.newadvent.org/fathers/1801122.htm>.}

% structure: p0001 source=h1 target=section reason=page-title

\section{\WorkTitle{论圣咏第一百二十三篇}}

1. ……让这位歌者上升；让这个人从你们每个人的心中歌唱，让你们每个人都成为这个人，因为当你们每个人都这样说时，既然你们都在基督内合而为一，就是一个人这样说；他说的不是：\Quote{上主，我们向祢举目}，而是：\BibleQuote{上主，我向祢举目}\BibleRef{咏 122:1}。你们确实应当想象每个人都在说话；但那一位以特殊方式说话，他散布于整个世界……\par

是什么使基督徒的心沉重？因为他是一个朝圣者，渴望着他的家乡。如果你的心为此沉重，即使你在世上一帆风顺，你仍然会叹息；如果万事都使你顺利，这个世界在各方面都向你微笑，你仍然会叹息，因为你看到你身处异乡；你感到你在愚人眼中确实有福，但尚未得到基督的恩许；你带着叹息寻求它，你带着渴望寻求它，在渴望中你上升，在上升中你唱着登阶之歌。\par

2. ……那么梯子在哪里呢？因为我们看到天地之间如此大的间隔，如此宽广的距离，如此辽阔的区域；我们想攀登那里，却看不到梯子；我们唱着登阶之歌，即上升之歌，是否在自欺呢？我们上升到天堂，如果我们思想天主，祂在心中造了上升的阶梯。什么是用心上升？就是向天主前进。正如每个跌倒的人，不是下降，而是坠落；同样每个进步的人，就是上升；但若他进步而避免骄傲；若他上升而不坠落；但若他在进步中变得骄傲，在上升中他又坠落。但为使他不要骄傲，他该做什么？他该举目仰望那住在天上的主，不要只顾自己……\par

3. 我的弟兄们，如果我们把天理解为用肉眼看到的穹苍，我们确实会错误地以为没有梯子或攀登工具就无法攀登上去；但如果我们灵性上升，我们就该灵性地理解天；如果上升在于情感，天堂在于正义。那么什么是天主的天空？所有圣洁的灵魂，所有正义的灵魂。因为宗徒们虽然肉身在地上，却是天堂；因为主临在于他们中，走遍了全世界。那么祂住在天堂里。怎样？……他们按信德成为圣殿多久？只要基督因信德住在他们心中；如宗徒所说：\BibleQuote{愿基督因信德住在你们心中。}但是那些已经能面对面看见祂的人，已经是天堂，天主已可见地住在他们内；所有圣宗徒、所有神圣的德能、掌权者、统治者，那个天上的耶路撒冷，我们这些流亡者为了它而叹息，并渴望祈祷；那里天主居住。圣咏作者已将他的信德举目于此，他在情感中上升，怀着渴望的希望：正是这渴望使灵魂涤除罪的污秽，洗净一切瑕疵，使它自己也成为天堂；因为它已举目仰望那住在天上的主。因为如果我们认定用肉眼看到的那个天是天主的居所，天主的居所将会消逝；因为\BibleQuote{天地要过去}\BibleRef{玛 24:35}。那么，在天主创造天地之前，祂住在哪里？但有人说：在天主创造圣人之前，祂住在哪里？天主住在自己内，祂与自己同住，天主与自己同在。当祂屈尊住在圣人中，圣人并不是这样的房屋，以致当天主离开时就会倒塌。因为我们住在一所房屋里是一种方式，天主住在圣人是另一种方式。你住在房子里：如果房子被移走，你就倒塌；但天主住在圣人中，如果祂自己离开，他们就会倒塌……\par

4. 既然他说了：\BibleQuote{上主，我向祢举目}\BibleRef{咏 122:2}，那么下文是什么？你是怎样举目的？\BibleQuote{看，如同仆人的眼睛望着主人的手，又如同婢女的眼睛望着主母的手：同样，我们的眼睛望着上主我们的天主，直到祂怜悯我们。}我们既是仆人，又是婢女：祂既是我们的主人，又是我们的主母。这些话是什么意思？这些比喻是什么意思？如果我们是仆人，祂是我们的主人，这并不奇怪；但奇怪的是我们是婢女，祂是我们的主母。但我们是婢女也并不奇怪；因为我们是教会；祂是我们的主母也并不奇怪；因为祂是天主的德能和智慧……因此，当你听到基督时，要举目仰望你主人的手；当你听到天主的德能和智慧时，要举目仰望你主母的手；因为你既是仆人又是婢女；仆人，因为你是子民；婢女，因为你是教会。但这婢女在天主面前获得了极大的尊荣；她成了妻子。但在她来到那灵性的拥抱中，在那里她可以无忧无虑地享受她所爱的那一位，并为那位在这个漫长而乏味的朝圣中叹息的那一位，她已订婚；并收到了一份巨大的聘礼，就是她为之叹息而无惧的新郎的血。并没有人对她说：不要爱；就像有时对任何订婚的童贞女所说的，尚未结婚；并且有正当的理由说：不要爱；当你成为妻子时再去爱；这样说是对的，因为爱一个不知道是否会嫁给他的人是一种仓促和颠倒的事情，不是贞洁的渴望。因为可能与她订婚的是一个男人，而娶她的却是另一个男人。但由于没有人能比基督更优先，让她毫无畏惧地去爱；在与他结合之前，让她去爱，从远处、从遥远的朝圣中叹息……\par

5. \Quote{我们实在饱受轻蔑} \BibleRef{咏 122:3}。凡愿按照基督虔诚生活的人，必然要受责斥（\BibleRef{弟后 3:12}），必然被那些不愿虔诚生活、只以现世为幸福的人所鄙视。他们嘲笑那些把肉眼看不见的幸福称为幸福的人，对他们说：\Quote{你相信什么，疯子？你相信你所看见的吗？有人从阴间回来，告诉你那里发生的事吗？看，我看见并享受我所爱的。}你被嘲笑，因为你盼望你所看不见的；而那个似乎拥有他所看见的人嘲笑你。仔细想想，他真的拥有它吗？……\Quote{我有我的房屋}，他夸口说。你问，他自己的房屋是哪个？\Quote{我父亲留给我的。}这房屋是他从哪里得来的？\Quote{我祖父留给他的。}再追溯到他曾祖父，再到他曾祖父的父亲，他就再也叫不出他们的名字了。你难道不因这个想法而惊恐：你看见许多人经过这房屋，却没有一个人能把它带到他永久的家？你父亲留下了它：他经过了它：这样你也将经过。因此，如果你在你的房屋中只是短暂的停留，那它便是过往客人的客栈，而非固定的居所。然而，既然我们盼望那将来的事，叹息未来的幸福，而且我们将来怎样还没有显明，尽管我们已经是\Quote{天主的子女}（\BibleRef{若一 3:2}），因为\Quote{我们的生命与基督一同藏在天主内}（\BibleRef{哥 3:3}），\Quote{我们完全被轻视}，被那些在现世寻求或享受幸福的人。\par

6. \Quote{我们的灵魂饱受困扰：成为富人的羞辱，骄傲者的轻蔑} \BibleRef{咏 122:4}。我们曾问谁是\Quote{富人}；他向你解释，因为他已说过\Quote{骄傲者}。\Quote{羞辱}和\Quote{轻蔑}是相同的；\Quote{富人}与\Quote{骄傲者}也是相同的。这是句子的重复：\Quote{成为富人的羞辱，骄傲者的轻蔑}。为什么骄傲者是富有的？因为他们希望在此世幸福。为什么？既然他们自己也是可怜的，他们是富有的吗？但也许当他们可怜时，他们并不嘲笑我们。亲爱的，请听。或许他们只在快乐时嘲笑，当他们夸耀自己财富的浮华时！当他们夸耀自己虚假荣誉的膨胀状态时：那时他们嘲笑我们，仿佛说：\Quote{看，我一切都好；我享受眼前的福乐；让那些许诺却无法展示的人离开我；我看见的，我抓住；我看见的，我享受；愿我在此生幸福。}你更应当安稳；因为基督已复活，教导你祂将在另一生命中赐予什么：你要确信祂会赐予。但那个人嘲笑你，因为他拥有他所拥有的。忍受他的嘲笑，你将嘲笑他的呻吟：因为随后会有一个时刻，这些人自己会说：\Quote{这就是我们曾一度讥笑的那个人。}（\BibleRef{智 5:3}）……\par

7. 此外，有时那些处于暂时不幸鞭打之下的人也嘲笑我们……与我们的被钉之主同钉的强盗（\BibleRef{路 23:39}）岂不也嘲笑过吗？因此，如果那些并不富有的人也嘲笑我们，为何圣咏说\Quote{成为富人的羞辱}？如果我们仔细分析，甚至这些（不幸的人）也是富有的。他们如何富有？是的；因为他们若不富有，就不会骄傲。一个人富有钱财，并因此而骄傲；另一个人富有荣誉，并因此而骄傲；另一个人自以为富有义德，因此骄傲，这更糟糕。那些似乎不富有钱财的人，自以为在天主面前富有义德；当灾难临到他们时，他们为自己辩护，指责天主，说：\Quote{我犯了什么错？我做了什么？}你回答说：\Quote{回顾一下，记起你的罪过，看看你是否没有做什么。}他的良心有所触动，回到自己，思想他的恶行；但当他思想了恶行之后，他仍然不愿意承认他应受痛苦；而是说：\Quote{看，我显然做了许多事；但我看见许多人做了更坏的事，却没有遭受任何灾祸。}他对抗天主。因此他也是富有的：他的胸膛因义德而膨胀；因为天主在他看来似乎作恶，而他似乎无辜受苦。如果你给他一艘船来驾驶，他必会连同船一起沉没：然而他却想剥夺天主管理此世的权力，自己掌握受造的舵，给所有人分配痛苦与快乐、惩罚与奖赏。可怜的灵魂！但你们为何惊奇？他是富有的，却富有于不义，富有于恶意；他越是自以为富有义德，就越富有于不义。\par

8. 然而基督徒不应富有，而应承认自己是贫穷的；如果他拥有财富，他应知道这些不是真正的财富，这样他才能渴慕别的……而我们正义的财富是什么呢？我们内无论有多少正义，相比那泉源，它不过像露珠；相比那丰盈，它只是几滴水，能软化我们的生命，放松我们顽固的罪过。只愿我们渴望被那正义的泉源充满，渴望被那丰盈的富足充满，正如圣咏所说：\BibleQuote{他们必因祢殿宇的丰盈而饱饫，祢也必叫他们喝祢喜乐的川流。}但当我们在此世时，让我们明白自己是缺乏和困乏的；不仅就那些不是真正财富的财富而言，更是就救恩本身而言。当我们健全时，让我们明白我们是软弱的。因为只要这身体饥渴，只要这身体因守夜而疲倦，因站立而疲倦，因行走而疲倦，因坐着而疲倦，因进食而疲倦；无论它转向何处以求缓解疲倦，那里就发现另一个疲乏的源头：因此，在这身体本身内也没有完全的健全。那财富就不是财富，而是贫穷；因为财富越增加，匮乏和贪欲也越增加……因此，让我们全部的饥饿、全部的干渴，都针对真正的财富、真正的健康和真正的正义。什么是真正的财富？是那天上的耶路撒冷居所。因为在此世谁被称为富有呢？当一个人被称赞为富有时，是什么意思？他很富有：什么都不缺。这称赞者称赞别人时，确实是这样说的：因为并非如此，当说\Quote{他什么都不缺}时。请思考他是否真的什么都不缺。如果他什么也不渴望，他就什么都不缺；但如果他仍渴望比他拥有的更多，那么他的财富增长的同时，他的匮乏也增长了。但在那座城里将有真正的财富，因为在那里我们什么都不缺；我们不需要任何东西，那里有真正的健康……\par

\SourceRecord{Translated by J.E. Tweed.；Nicene and Post-Nicene Fathers, First Series；Vol. 8.；Edited by Philip Schaff.；Buffalo, NY: Christian Literature Publishing Co.,；1888.；<http://www.newadvent.org/fathers/1801123.htm>.}

% structure: p0001 source=h1 target=section reason=page-title

\section{圣咏第一百二十四篇释义}

你们已经熟知，最亲爱的弟兄们，\Quote{登阶之歌，}是我们上升之歌；这上升不是靠身体的脚，而是靠心灵的意愿。我们已多次提醒你们，不必过于重复，以便为尚未说出的话留出空间。因此，你们刚才听到歌咏的这篇圣咏，题名为\Quote{登阶之歌。}这就是它的标题。他们上升时歌唱；有时仿佛是一个人歌唱，有时仿佛是许多人；因为多人为一，既然基督是一，而在基督内，基督的肢体与基督合成一个，所有这些肢体的首是在天上。虽然肢体在地上劳苦，却未与首分离；因为首从高处俯视，并眷顾肢体。\BibleRef{宗 9:4}……因此，无论是一人或多人歌唱，多人是一人，因为这是合一；而基督，如我们所说，是一，所有基督徒都是基督的肢体。……那身体中有些肢体——我们也是其中的——确实先我们而行，能够真实地歌唱。这是圣殉道者们所歌唱的：因为他们已经逃出，且与基督同享喜乐，并将在末后领受那不可能朽坏的身体，就是最初可朽坏的那同一个，其中他们曾忍受痛苦；从这同一个身体中，将为他们的义德做成装饰。因此，无论是他们实际如此，还是我们怀着希望，将我们的情感与他们的冠冕联合，并渴望如此一种在此世没有、且除非在此世渴望否则永不得着的生命，让我们一同歌唱说：\Quote{倘若上主自己没有在我们中间。}……\par

……那身体中有些肢体——我们也是其中的——确实先我们而行，能够真实地歌唱。这是圣殉道者们所歌唱的：因为他们已经逃出，且与基督同享喜乐，并将在末后领受那不可能朽坏的身体，就是最初可朽坏的那同一个，其中他们曾忍受痛苦；从这同一个身体中，将为他们的义德做成装饰。因此，无论是他们实际如此，还是我们怀着希望，将我们的情感与他们的冠冕联合，并渴望如此一种在此世没有、且除非在此世渴望否则永不得着的生命，让我们一同歌唱说：\Quote{倘若上主自己没有在我们中间。}……\par

\Quote{倘若上主自己没有在我们中间，以色列如今可以这样说} \BibleRef{咏 123:1}……何时？\Quote{当人起来攻击我们时} \BibleRef{咏 123:2}。不要惊奇：他们已被制服：因为他们是人；但主在我们内，人不在我们内：因为人起来攻击我们。然而，若非在那不能被制服的人中，不是人而是主，人就会碾压其他人。因为当你们欢喜、歌唱且稳享永福时，人能对你们做什么？当人起来反对你们时，倘若主不在你们一边，他们能做什么？\Quote{或许他们曾活活地吞噬了我们} \BibleRef{咏 123:3}。\Quote{吞噬了我们：}他们不会先杀死我们，然后吞噬我们。多么不人道、多么残忍的人！教会并不如此吞噬。对伯多禄说过：\Quote{\Emph{宰杀}了吃:} \BibleRef{宗 10:13}，而不是\Quote{活活吞噬}。因为没有人进入教会的身体，除非他先被宰杀。他所是的死去，好使他成为他所不是的。否则，那没有被宰杀、也没有被教会吞噬的人，可能存在于可见的百姓数目中；但他不能存在于天主所认识的百姓数目中，对此宗徒说：\Quote{主认识属于祂的人，} \BibleRef{弟后 2:19}，除非他先被吞噬；而被吞噬，必须先被宰杀。外邦人来了，在他内仍活着偶像崇拜；他必须被接枝到基督的肢体中：为了被接枝，他必须被吞噬；但教会不能吞噬他，除非他先被宰杀。让他弃绝世界，然后他死了；让他相信天主，然后他被吞噬……但那些主在他们内的人，被宰杀却不死。而那些同意并活着的人，当他们被吞噬时，就活着被吞，死去。但那些忍受了苦难且在考验中没有屈服的人，欢喜说：\Quote{倘若主没有在我们中间}等。\par

……\Quote{当他们的怒气对我们发怒时。}他们现在发怒，他们现在公开狂怒：\Quote{或许水已淹没我们。} \BibleRef{咏 123:4}。他用水指那不虔敬的万民；我们将在以下经文中看到是怎样的水。凡同意他们的人，水就会淹没他。因为他会死于埃及人的死，不会仿效以色列人过去。因为你们知道，弟兄们，以色列的子民经过那水，而埃及人被淹没。\BibleRef{出 14:22}但这是怎样的水？是湍流，它猛烈流淌，但终将过去……因此，我们的首最先饮了这水，正如在圣咏中论及祂所说：\Quote{祂在途中喝湍流：因此祂必将抬起头来。}因为我们的首已被高举，因为祂在途中喝了湍流；我们的主已受难。因此，若我们的首已被举起，为何肢体害怕湍流？毫无疑问，因为首已被举起，肢体也将后来说：\Quote{我们的灵魂渡过了湍流。或许我们的灵魂渡过了那虚妄之水} \BibleRef{咏 123:5}。看，他提到的是怎样的水：\Quote{水或许已淹没我们。}但\Quote{虚妄之水}是什么意思？\par

首先，\Quote{或许我们的灵魂渡过了} \BibleRef{咏 123:5}是什么意思？然而，你要明白意思是：\Quote{你以为我们的灵魂渡过了吗？}他们为什么说\Quote{你以为}？因为危险之大，使得逃脱几乎难以相信。他们忍受了巨大的死亡；他们处于巨大的危险中；他们如此被压迫，以至于几乎活着时同表同意，差点活活被吞；如今既然他们逃脱了，如今他们安全了，但仍记得那危险，巨大的危险，他们说：\Quote{或许我们的灵魂渡过了那虚妄之水吗？}\par

6. 那没有实体的水是什么，岂不是没有实体的罪的水？因为罪没有实体：它们有匮乏，没有实体；它们有欠缺，没有实体。在那没有实体的水中，小儿子丧失了他全部的财产……你愿看见这水如何没有实体吗？把你所获得的带到下界去：你将做什么？你获得了金子：你失去了你的信德：几天之后你离开此生；你不能把你用失去良好信德而获得的黄金带走；你的心，缺乏信德，出去受罚——你的心，如果充满信德，则会出去得冠冕。看，你所做的虚无：你为虚无得罪了天主。\par

7. 人们听到那常见的谚语；天主的谚语在他们心中沉睡。什么谚语？\Quote{手里有的胜过希望中的。}不幸的人，你手里有什么？你说：\Quote{手里有的胜过希望中的。}握紧它以免失去，然后再说：\Quote{手里有的胜过希望中的。}但如果你握不住它，为何不紧紧抓住你不能失去的呢？那么你手里有什么？金子。因此，把它拿在手里：如果你拿在手里，就不要让它不经你同意被拿走。但如果通过金子你也被人带到你不愿去的地方，如果更强大的强盗找寻你，因为他发现你是一个较弱的强盗；如果更强的鹰追赶你，因为你之前从他面前夺走了一只野兔：较小的猎物是你，你将成为更大的猎物。人们看不到人类事务中的这些事：他们被如此多的贪欲蒙蔽了……\par

8. 让他们逃离没有实体的水，并说：\BibleQuote{愿主受赞美，他没有把我们交给他们作牙齿的猎物}\BibleRef{咏 123:6}。因为猎人正追赶，并在他们的陷阱里放了诱饵。什么诱饵？此生的甜蜜，好使每个人为此生甜蜜的缘故，将头伸入不义，被困在陷阱中。不是他们，主在他们内，他们都说：\Quote{如果主自己不在我们内；}他们没有被困在陷阱中。让主在你内，你将不会被困在陷阱中。\par

9. \BibleQuote{我们的灵魂得以逃脱，像鸟儿脱离捕鸟人的罗网}\BibleRef{咏 123:7}。因为主在灵魂本身内，因此那灵魂得以逃脱，像鸟儿脱离捕鸟人的罗网。为什么像鸟儿？因为它曾像鸟儿一样粗心坠落；后来它能说：\Quote{天主会宽恕我。}不稳定的鸟儿，倒要将你的脚稳立在磐石上：不要陷入陷阱。你将被捉住、消灭、压碎。让主在你内，他将救你脱离更大的威胁，脱离捕鸟人的罗网。正如你看见一只鸟即将落入陷阱，你发出更大的声响使它飞离罗网；同样，当或许有些殉道者为了贪图此生的享受而伸长脖子时，在他们内的主发出了地狱的声响，那鸟便脱离了捕鸟人的罗网。那罗网就是此生的甜蜜：他们没有缠在罗网中，却被杀了；因他们的被杀，罗网被挣破；此生的甜蜜不再存留，使他们再被缠住，但罗网被压碎了。那鸟也被压碎了吗？断乎不是！因为它不在罗网中：\Quote{网罗已被挣断，我们脱离了。}\par

10. ……\BibleQuote{我们的援助是在造天地的主的名号中}\BibleRef{咏 123:8}。因为如果这不是我们的援助，罗网固然不会永存；但一旦鸟儿被捉住，它就会被压碎。因为此生将要过去；凡被其快乐所诱惑，并借这些快乐得罪了天主的人，也将与此生一同过去。因为罗网将被挣破；你们要确信：当它被指定的时限满足后，现世的一切甜蜜将不复存在；但我们不可被它束缚，以致当网被挣破时，你们能欢喜地说：\Quote{网罗已被挣断，我们脱离了。}但免得你以为能凭自己的力量做到，要思量你脱离是谁的作为（因为如果你骄傲，你就陷入罗网），并说：\BibleQuote{我们的援助是在造天地的主的名号中}\BibleRef{咏 123:8}。……\par

\SourceRecord{Translated by J.E. Tweed.；Nicene and Post-Nicene Fathers, First Series；Vol. 8.；Edited by Philip Schaff.；Buffalo, NY: Christian Literature Publishing Co.,；1888.；<http://www.newadvent.org/fathers/1801124.htm>.}

% structure: p0001 source=h1 target=section reason=page-title

\section{论圣咏第一百二十五篇}

这篇圣咏，属于登阶歌之一，教导我们，当我们以爱德与孝爱之情上升并向主我们的天主高举心神时，不要注目于在这世上亨通的人，他们的幸福是虚假、不稳、全然诱人的；在那里他们只珍惜骄傲，他们的心对天主冻结，并在祂恩宠的沛雨下变得刚硬，以致不结出果实……\par

\BibleQuote{他们信赖上主，必如熙雍山，永不动摇。}\BibleRef{咏124:1}\par

这些人是谁？\Quote{他们必永远坚定不移，那居住在耶路撒冷的。}\BibleRef{咏124:2} 如果我们理解这尘世的耶路撒冷，所有居住在其中的人都因战争和城邑的毁灭而被驱逐了；如今你在耶路撒冷城中寻找一个犹太人，却找不到。那么，为什么\Quote{那居住在耶路撒冷的永不动摇}呢？除非因为有另一个耶路撒冷，就是你们常听说的那一位。她是我们的母亲，我们在这旅途中为她叹息呻吟，为能回到她那里……因此，居住在她里面的人\Quote{永不动摇}\BibleRef{咏124:2}。但居住在那尘世耶路撒冷的人却动摇了；先在心中，后在被掳。当他们在心中动摇而跌倒时，就钉死了那天上耶路撒冷本身的君王；他们已经在精神上被摒弃在外，并将他们的君王拒之门外。因为他们把他赶出城外，在城外钉死了他。\BibleRef{若19:17} 他也将他们逐出他的城，即那永远的耶路撒冷，我们在天之母。\par

这耶路撒冷是什么？他简短地描述。\Quote{群山环绕耶路撒冷}\BibleRef{咏124:2}。这有什么了不起的，我们竟是一座群山环绕的城？难道我们的全部幸福就在于我们将拥有一座有群山环绕的城吗？我们不知道群山是什么吗？群山不就是大地的隆起吗？那么，我们所爱的那些山是与此不同的，崇高的山，真理的宣讲者，无论是天使、宗徒或先知。他们环绕耶路撒冷；他们包围她，好像给她形成一道墙。圣经常常提到这些可爱而令人喜悦的山……他们是我们在歌中所唱的山：\Quote{我举目向山，我的救助从何而来？}因为在此生我们从圣经中得到帮助。通过那领受平安的山，小山领受了正义：因为关于这些山本身，他说了什么？他没有说，他们从自己拥有平安，或他们创造平安，或产生平安；而是说，他们领受平安。主是源头，他们从那里领受平安。因此，为了平安的缘故，请你举目向山，好使你的救助从造天地的上主而来。再者，圣神提到这些山时这样说：\Quote{祢从祢的永恒之山奇妙地光照他们}。他没有说，山光照他们；而是说，祢从祢的永恒之山光照他们：通过那些祢愿意使他们成为永恒的山，宣讲福音，是祢光照他们，不是山。这样的山就是那\Quote{环绕耶路撒冷的山}。\par

为使你知道环绕耶路撒冷的是怎样的山；圣经在提到好山时，很少、几乎不曾不立刻也提到主，或同时提及他，好叫我们的希望不寄托于山……免得你再停留在山上，他立刻加上：\Quote{上主也如此环绕他的子民}：好使你的希望不在于山，而在于那光照山的主。因为当他居住在山上时，即在圣人们内时，他自己围绕他的子民；他以精神上的坚垒围护他的子民，使他永不动摇。但当圣经提到邪恶的山时，就不加上主。这样的山，我们已常告诉你们，是指某些强大而邪恶的灵魂。因为兄弟们，你们不要以为异端能通过渺小的灵魂产生。只有伟人曾是异端的创始人；但他们有多么强大，也就有多么邪恶，成为山。因为他们不是那些领受平安的山，好使小山领受正义；而是从他们的父魔鬼那里领受了纷争。因此曾有这些山：你要谨慎，不要逃向这样的山。因为人会来，并向你们说：有一个伟大的英雄，有一个伟人！那个\OriginalTerm{多纳图}{Donatus}是多么伟大！\OriginalTerm{马克西米安}{Maximian}又多么伟大！还有那个\OriginalTerm{佛提诺}{Photinus}，他是多么伟大！\OriginalTerm{亚略}{Arius}也是如此显赫！我以上提到的这些，都是山，但却是使船遇难的山……\par

6. 但要爱这样的高山，主在其中。那么，如果你不寄望于它们，那些高山也会爱你。请看，弟兄们，天主的山是什么。在另一处经文里它们也被这样称呼：\Quote{祢的正义好像天主的山。}不是他们的正义，而是\Quote{祢的正义。}请听那位伟大的山——宗徒所说：\Quote{为使我找到祂，不占有我自己的正义，即出于法律的，而是借着对基督的信德。}（\BibleRef{斐 3:9}）但那些选择凭自己的正义成为高山的人，就像某些犹太人或法利塞人的首领那样，被责备说：\Quote{他们不认识天主的正义，而企图建立自己的正义，便没有服从天主的正义。}（\BibleRef{罗 10:3}）但那些服从的人被高举却如此谦卑。因他们伟大，故是高山；因他们服从天主，故是山谷；因他们拥有虔诚的容量，故接受丰沛的和平，并将丰盈的灌溉传送到小山丘。只是要小心，现在你爱什么样的高山。如果你希望被善山所爱，就不要连善山也信赖。因为保禄是多么大的山啊！哪里能找到像他这样的人？我们说的是人的伟大。有谁能轻易地获得如此大的恩宠呢？然而，他担心那鸟会信赖他，所以说什么：\Quote{保禄为你们被钉在十字架上了吗？}（\BibleRef{格前 1:13}）但你要举目向山，帮助可以来自它们；因为，\Quote{我栽种了，阿颇罗浇灌了}，但你的帮助来自那创造天地的上主；因为，\Quote{天主使之生长。}（\BibleRef{格前 3:6}）因此，\Quote{高山环绕耶路撒冷。}但正如\Quote{高山环绕耶路撒冷，上主也环绕祂的子民，从现今直到永远。}（\BibleRef{咏 125:2}）因此，如果高山环绕耶路撒冷，而上主环绕祂的子民，上主就将祂的子民系于爱和和平的同一纽带中，使信赖上主的人像熙雍山一样，永不动摇：这就是\Quote{从现今直到永远。}\par

7. \Quote{因为上主不容恶人的棍杖停留在义人的基业上，免得义人伸手行恶。}（\BibleRef{咏 125:3}）的确，现在义人受苦有些程度，现在不义的人有时暴虐地对待义人。以什么方式？有时不义的人获得世俗的荣誉：当他们获得它们，并成为审判官或君王时；因为天主这样做是为管教祂的子民，管教祂的人民；他们权力的荣誉必须被显示出来。因为天主这样规定了祂的教会：世界上每一个被建立的权力都应当得到荣誉，有时来自那些比掌权者更好的人。为举例说明，我取一个例子；由此计算所有权力的等级。人与人之间首要且日常的权威关系就是主仆关系。几乎所有家庭都有这种权力。有主人，也有奴隶；这些是不同的名称，但人与人是平等的名称。宗徒教导奴隶要服从主人，他说什么？\Quote{奴仆们，要服从你们肉身的主人。}（\BibleRef{弗 6:5}）因为有一位按圣神是真正永久的主。但那些暂时的主人只是暂时的。当你在路上行走，当你生活在此生时，基督不愿使你骄傲。你成了基督徒，又有一个人做你的主人：你成为基督徒，并非为轻视为奴。因为当你按基督的命令服侍一个人时，你服侍的不是那人，而是那命令你的。他也说：\Quote{奴仆们，要服从你们肉身的主人。}看，他没有使人从奴仆中自由，而是使坏奴仆变好。富人欠基督多少，他为他们管理家庭！因此，如果你有一个不信主的仆人，假设基督归化了他，却不对他说：离开你的主人，你已经认识那位你真正的主；他也许是不敬虔和不义的，你现在是忠信和正义的：一个正义和忠信的人服侍一个不义和不信的主人是不相称的。他没有这样对他说，而是说：服侍他；并为了坚固那仆人，加上：像我一样服侍；我在你之前服侍了不义者……如果天地的主，万物借着祂受造，服侍了不配的人，为祂狂怒的迫害者祈求怜悯，并在祂降临时好像显示祂自己是他们的医生（因为医生，无论在技艺上还是在健康上都更好，也服侍病人）：那么，人岂不更应当全心全意、全善的意愿、全爱去服侍甚至一个坏主人吗？看，更好的服侍较差的，但只是暂时的。请理解我所说的关于主奴的话，对权力和君王，以及世上所有高位也同样适用。因为有时它们是好的权力，敬畏天主；有时它们不敬畏天主。尤利安是一个不信主的皇帝，一个背教者，一个邪恶的人，一个拜偶像者；基督徒士兵服侍一个不信主的皇帝；当他们来到基督的案件时，他们只承认那在天上的主。如果他在任何时候呼召他们去崇拜偶像、献香，他们宁可选择天主而不选择他；但每当命令他们列队行军，攻击这个或那个民族，他们立刻服从。他们分辨他们永远的主和暂时的主人；然而，为了他们永远的主的缘故，他们服从于他们暂时的主人。\par

8. 然而，难道恶人永久有权势管辖义人吗？不会如此。恶人的棍杖暂时落在义人的产业上，但不会被留在那里，不会永远如此。时候将到，基督要在祂的荣耀中显现，将万民聚集在祂面前。\BibleRef{玛 25:32} 在那里你将会看见：羊群中有许多奴隶，山羊中有许多主人；同样，羊群中也有许多主人，山羊中也有许多奴隶。因为并非所有奴隶都是好的——不要从我们给予仆人的安慰中推断这一点——也并非所有主人都是恶的，因为我们这样抑制了主人的骄傲。有好主人是信徒，也有恶主人；有好仆人是信徒，也有恶仆人。然而，只要好仆人服侍恶主人，就让他们暂时忍受。\Quote{因为天主必不让恶人的棍杖常留在义人的产业上。}为何祂不会？\Quote{免得义人伸手行恶：}为使义人暂时忍受恶人的辖制，并明白这不是永久的，好使他们预备自己，承受永远的产业....\par

9. 因此，圣咏作者接着又说：\BibleQuote{上主，求祢善待心地善良真诚的人。}\BibleRef{咏 124:4} 那些心正的人——我刚才所提到的——他们随从天主的旨意，而不随从自己的意愿——就默想此事。但那些愿意跟随天主的人，让天主在前引领，自己跟随在后；而不是自己在前，让天主在后；他们在一切事上都发现天主是良善的，无论是管教、安慰、磨练、加冕、洁净还是光照；正如宗徒所说：\BibleQuote{我们知道，万事互相效力，使爱天主的人获得益处。}\BibleRef{罗 8:28}\par

10. 因此，圣咏作者立刻加上：\BibleQuote{至于那些偏行弯曲道路的人，上主要引领他们与作恶的人一同去受绞杀。}\BibleRef{咏 124:5} 这就是说，那些效法他们行为的人，因为他们贪图眼前的快乐，不相信将来的惩罚。那么，那些心正而不后退的人将得着什么？让我们现在来谈论产业本身，弟兄们，因为我们都是儿女。我们将拥有什么？我们的产业是什么？我们的家乡是什么？它叫什么？和平。我们以此问候你们，以此向你们宣告，大山接受这和平，小山也接受这和平作为正义。和平就是基督：\BibleQuote{因为祂是我们的和平，祂使双方合而为一，拆毁了中间的隔墙。}\BibleRef{弗 2:14} 既然我们是儿女，我们将拥有产业。这产业要叫什么？不就是和平吗？要想想，那些不爱和平的人被剥夺了继承权。现今，那些分裂合一的人不爱和平。和平是虔诚者的财产，是继承人的财产。谁是继承人？儿女……。既然天主子基督就是和平，祂便来聚集祂自己的人，把他们与恶人分开。从什么样的恶人中分开？从那些仇恨耶路撒冷、仇恨和平、想要撕裂合一、不信和平、向人民宣讲虚假和平而自身却没有和平的人。当他们说：\Quote{愿和平与你们同在，}\Quote{也与你们的心灵同在，}时，回答是这样：但他们说假话，也听假话。他们对谁说\Quote{愿和平与你们同在}？是对那些他们使之与全地和平分离的人。而谁听到说\Quote{也与你们的心灵同在}？是那些拥抱纷争、仇恨和平的人。因为如果和平在他们心中，他们岂会不爱合一、放弃纷争呢？他们既然说假话，也就听假话。让我们说真话，听真话。让我们成为以色列，拥抱和平；因为耶路撒冷是和平的异象，我们是以色列，\Quote{愿和平归于以色列。}\par

\SourceRecord{Translated by J.E. Tweed.；Nicene and Post-Nicene Fathers, First Series；Vol. 8.；Edited by Philip Schaff.；Buffalo, NY: Christian Literature Publishing Co.,；1888.；<http://www.newadvent.org/fathers/1801125.htm>.}

% structure: p0001 source=h1 target=section reason=page-title

\section{对圣咏第一百二十六篇的诠释}

1. ……人如何陷入俘虏，让我们询问宗徒保禄……。因为他说道：\Quote{因为我们知道法律是属神的，但我是属血肉的，已被卖于罪恶。}\BibleRef{罗 7:14} 看，我们为何成了俘虏；因为我们被卖于罪恶。谁卖了我们？我们自己，我们同意了诱惑者。我们能出卖自己，却不能赎回自己。我们因同意罪恶而卖了自己，在信德之正义中被赎回。因为无罪的血为我们倾流，好使我们得蒙救赎。他在迫害义人时所流的血，是什么样的血？他流了义人的血，他们曾是先知、义人、我们的祖先和殉道者。他所流的血，都出自罪恶的后裔。他所流的惟独那一位并非因成义而生而为义者的血；因了这血的倾流，他失去了他所控制的人。因为那些无罪的血为他们倾流的人被赎回了，并从他们的俘虏中回头，高唱这篇圣咏。\par

2. \Quote{上主转变了锡安的俘虏时，我们就像得到了安慰的人。}\BibleRef{咏 125:1} 他以此意是说，我们变得喜乐。何时？\Quote{上主转变了锡安的俘虏时。}什么是锡安？耶路撒冷，也就是永恒的锡安。锡安如何是永恒的，又如何是俘虏的？在天使中是永恒的，在人类中是俘虏的。因为那城的所有居民并不是都成为俘虏，那些离开那里的人才成为俘虏。人是耶路撒冷的一个公民，但被卖于罪恶，成了一个旅居者。从他的后代生出了人类，锡安的俘虏充满了全地。这锡安的俘虏是如何成为那耶路撒冷的预像？那赐予犹太人的锡安的预像，在预象、在形像中，曾在巴比伦为俘虏，七十年后那民族回到了自己的城市。……但当所有时间过去后，我们返回我们的家乡，如同七十年后那民族从巴比伦俘虏中返回，因为巴比伦就是这世界；因为巴比伦解释为\Quote{混乱}。……因此，这整个人生处境就是混乱，它不属于天主。在这混乱中，在这巴比伦之地，锡安被扣押。但\Quote{上主转变了锡安的俘虏。}\Quote{我们就像，}他说，\Quote{得到了安慰的人。}这就是说，我们喜乐如同接受了安慰。安慰不给予不幸福的人，安慰只给予那些呻吟、哀悼的人。因此，\Quote{就像得到了安慰的人，}除非是因为我们仍在哀悼？我们为现今的境遇哀悼，我们在希望中得到安慰；当现今过去后，我们的哀悼将变为永久的喜乐，那时将不再需要安慰，因为我们不会为任何痛苦所伤。但他为什么说\Quote{就像}得到了安慰的人，而不说安慰？这词\Quote{就像}并不总是用于相似；当我们说\Quote{如同}时，有时指实际情形，有时指相似：这里是指实际情形。……因此，你们要在基督内行走，欢唱喜悦，如同一个得到安慰的人；因为那命令你们跟随他的那位，已走在你们前面。\par

3. \Quote{那时，我们的口充满喜乐，我们的舌充满欢欣。}\BibleRef{咏 125:2} 诸位弟兄，我们身体上的这个口，如何\Quote{充满喜乐}？它通常不\Quote{充满}，除非是食物、饮料或类似的入口之物。有时我们的口充满了；我们更要对你们的圣洁说，当我们口满时，我们不能说话。但我们有一个内在的口，就是在心里，从那里出来的，若是恶的，就玷污我们；若是善的，就洁净我们。关于这个口，你们在读福音时也听到了。因为犹太人责斥主，因为祂的门徒不用洗手就吃饭。他们责斥——那些外表洁净，内里却满污渍的人；他们责斥——那些义德只在人眼前的人。但主寻求我们内在的洁净，如果我们内在洁净，外表也必须是洁净的。\Quote{先清洁内部，}祂说，\Quote{然后外部也洁净了。}\BibleRef{玛 23:26} ……\par

4. 但让我们回到刚才从福音中读到的、与眼前经文相关的内容：\Quote{我们的口充满喜乐，我们的舌充满欢欣；}因为我们正在探究什么口和什么舌。请听，亲爱的诸位弟兄。主被讥笑，因为祂的门徒不洗手就吃饭。主恰当的回答了他们，并对祂所召来的群众说：\Quote{你们都听，且要明白：不是入于口的，污秽人；而是出于口的，那才污秽人。}\BibleRef{玛 15:10} 这是什么意思？当祂说，入于口的，祂仅指身体的口。因为食物入内，而食物不污秽人；因为\Quote{对洁净的人，一切都是洁净的；}且\Quote{天主的每一受造物都是好的，没有一样是可弃绝的，只要以感恩的心领受。}\BibleRef{弟前 4:4} ……\par

5. 要守护你内心的口免受邪恶，你便会是无辜的；你身体的口舌将是无辜的，你的手将是无辜的；甚至你的脚、你的眼、你的耳都将是无辜的；你的所有肢体将在正义之下服务，因为有一位正义的统帅——你的心。\Quote{那时，外邦人要说：上主为他们行了大事。}\par

6. \Quote{上主果然为我们行了大事，我们欢喜快乐。}\BibleRef{咏 125:3} 诸位弟兄，请思考，锡安现在不就在外邦人中说着这话吗？在整个世界中，不就看到人们正奔向教会吗？在全世界，我们的救赎被接纳了；回应了\Quote{阿们}。因此，耶路撒冷的居民，被俘虏、注定要返回、叹息渴望家乡的旅居者，在外邦人中这样说。他们说什么？\Quote{上主为我们行了大事，我们欢喜快乐。}他们为自己做了什么？他们对自己行了恶事，因为他们把自己卖于罪恶。救赎者来了，为他们行了善事。\par

7. \Quote{求主扭转我们的被掳，有如南方的洪流} \BibleRef{咏 125:4}。诸位弟兄，请思量这话的含意……就如洪流在南方被扭转，求祢扭转我们的被掳。圣经在某处劝勉我们行善时说：\Quote{你的罪过也必消融，如同晴天的寒冰。} \BibleRef{德 3:17} 因此，我们的罪曾束缚我们。如何束缚？有如寒冷冻结水流，使水不流。我们被罪的寒冰束缚，冻结了。但南风是暖风；南风吹起，寒冰融化，洪流便满溢。冬日的水流被称为洪流；因为因骤雨而满溢，奔腾有力。我们曾在那被掳中冻结；我们的罪束缚我们：南风——圣神——吹来了；我们的罪得到赦免，我们从不义的寒冰中释放；如同晴天的寒冰，我们的罪消融。让我们奔向我们的故乡，有如南方的洪流。\par

8. 因为接下来的话是：\Quote{含泪播种的，必含笑收割。} \BibleRef{咏 125:5} 在这充满泪水的生命中，让我们播种。我们播种什么？善工。怜悯之工是我们的种子；关于这些种子，宗徒说：\Quote{我们行善不要厌烦；若不松懈，到了时节自会有收获。} \BibleRef{迦 6:9} 因此，在论及施舍本身时，他说了什么？\Quote{我说：那少种的，也要少收。} \BibleRef{格后 9:6} 所以，那多种的，必多收；那少种的，必少收；那不种的，必一无所收。你们为何渴望广大的庄园，好能多种？没有比基督更宽阔的田地，祂愿意我们在祂内播种。你们的土壤是教会；尽可能多地播种。但你们没有足够的能力做这事。你们有心愿吗？正如你们所拥有的算不得什么，如果没有善意；同样，若你们有善意，也不要因缺乏而灰心。因为你们播种什么？怜悯。你们收割什么？平安。天使曾说：平安归于地上富人们吗？不，而是\Quote{平安归于地上善意的人们。} \BibleRef{路 2:14} 匝凯有强烈的意愿，匝凯有伟大的爱德。\BibleRef{路 19:8} 那么，那投入库中两个小钱的寡妇，播种得少吗？不，她与匝凯一样多。因为她的财力更有限，但意愿相同。她怀着与匝凯将一半家产捐出时同样的善意，投下了她的两个小钱。\BibleRef{路 21:1} 若你计算他们所给的，你会发现礼物不同；若你注视源头，你会发现它们相等；她给了她所有的一切，他给了他所拥有的一切……但若是那些以乞讨为业的乞丐，在他们困苦中，他们也有可彼此施与的。天主并没有如此遗弃他们，使他们没有机会因施舍而受考验。这人不能行走；能行走的，将双脚借给瘸子；能看见的，将眼睛借给瞎子；年轻力壮的，将力量借给年老或虚弱的，背负他；一个是贫穷的，另一个是富有的。\par

9. 有时富人也显得贫穷，贫穷人反而施舍给他。有人来到一条河边，他愈富有就愈娇弱；他不能渡河；若他以赤足渡河，他会着凉，会生病，会死；一个身体更强健的穷人走来；他背着富人过河；他向富人行了施舍。因此，不要以为只有没有钱财的人才是贫穷的……因此，你们要彼此相爱，彼此关怀。不要只顾及自己，也要顾及身边有需要的人。但因为这些事在此生中伴随着烦恼和忧虑，不要灰心。你们含泪播种，必将含笑收割。\par

10. 诸位弟兄，这是如何呢？当农夫带着犁和种子出去时，有时寒风刺骨，有时雨水阻挠。他仰望天空，见阴云密布，寒冷战栗，但他仍然出去播种。因为他害怕，若他留意恶劣天气，等待晴天，时光流逝，他可能找不到任何可收割的。不要拖延，诸位弟兄；在严冬中播种，甚至在哭泣中播种善工；因为\Quote{含泪播种的，必含笑收割。} 他们播种种子，即善意和善工。\Quote{他们一路前行，哭泣着撒种} \BibleRef{咏 125:6}。他们为何哭泣？因为他们处在可怜的人中，他们自己也是可怜的。诸位弟兄，与其你行施舍，不如无人可怜。然而，只要仍有行善的对象，让我们不要在这些困苦中停止撒种。虽然我们含泪播种，却将含笑收割。因为在死者复活时，各人将收到自己的禾捆，即他种子的出产，喜乐和愉悦的冠冕。那时将有欢乐的凯旋，我们将嘲笑死亡，我们之前在死亡中叹息；那时他们要对死亡说：\Quote{死亡啊，你的胜利在哪里？死亡啊，你的刺在哪里？} \BibleRef{格前 15:55} 但他们为何现在喜乐？因为\Quote{他们带着禾捆回来。}\par

11. 在这篇圣咏中，我们主要劝勉你们行施舍之事，因为这是我们上升的途径；你们看见那上升者，唱登阶之歌。要记住：不要爱下降胜于上升，而要思念你们的上升：因为那从\Place{耶路撒冷}下到\Place{耶里哥}去的，落在强盗手中。\BibleRef{路 10:30} ……那\Person{撒玛黎雅人}经过时并未轻看我们：祂医治了我们，将我们扶起放在祂的坐骑上，即祂的肉身上；祂领我们到了客店，即教会；祂将我们托付给店主，即宗徒；祂给了两个银钱，好使我们得医治，\BibleRef{路 10:35} 即爱天主和爱近人。那宗徒花费得更多；因为，虽然众宗徒作为基督的士兵，可以自基督的子民那里领取报酬，但那位宗徒还是亲手劳作，免除了子民应当供养他的义务。这一切已经实现了：如果我们曾经下降并受了伤，那么让我们上升，歌唱并前进，好能到达。\par

\SourceRecord{Translated by J.E. Tweed.；Nicene and Post-Nicene Fathers, First Series；Vol. 8.；Edited by Philip Schaff.；Buffalo, NY: Christian Literature Publishing Co.,；1888.；<http://www.newadvent.org/fathers/1801126.htm>.}

% structure: p0001 source=h1 target=section reason=page-title

\section{圣咏第一百二十七篇释义}

1. 在所有称为\Quote{登阶歌}的诗歌中，这篇圣咏在标题上另有一个附加语，就是\InnerQuote{撒罗满的}。因为它的标题是这样：\Quote{撒罗满的登阶歌}。因此，它引起了我们的注意，并促使我们探究这个附加语\InnerQuote{撒罗满的}的缘由。因为其他词语，如\Quote{登阶歌}，的解释无需重复……撒罗满在他那个时代是达味的儿子，是一位伟人，圣神通过他在圣经中成就了许多神圣的训诫、有益的劝告和神圣的奥秘。但撒罗满本人却是一个喜爱妇女的人，被天主所弃绝；这情欲成了他极大的网罗，竟被妇女引诱去祭拜偶像，\BibleRef{列上 11:7}，正如圣经为他作证的。但如果因他的堕落，经由他传下来的都被废除，那么人们就会认为这些训诫是他本人传下来的，而非经由他传下来的。因此，天主的仁慈和祂的圣神做了极好的工作，将凡通过撒罗满所宣告的美善都归于天主；而人的罪则归于人。撒罗满在天主子民中堕落有什么奇怪呢？亚当不是在乐园中堕落了吗？不是有一位天使从天上堕落，成了魔鬼吗？我们由此受教，不可寄望于任何人……撒罗满这个名字被解释为\Quote{和平者}；现在，基督是真正的和平者，宗徒论到祂说：\BibleQuote{祂是我们的和平，祂使双方合而为一。}\BibleRef{弗 2:14}……因此，既然祂是真正的撒罗满；因为那撒罗满是这位和平者的预像，当他建造圣殿时；为使你们不致认为那为天主建造殿宇的是真正的撒罗满，圣经向你们显示另一位撒罗满，这篇圣咏这样开头：\BibleQuote{若不是主建造房屋，建造的人就枉然劳力。}\BibleRef{咏 126:1}。所以，主建造房屋，主耶稣基督建造祂自己的房屋。许多人在建造中劳苦；但若不是祂建造，\Quote{建造的人就枉然劳力。}谁是那些劳苦建造的人呢？所有在教会中宣讲天主圣言的人，天主奥秘的仆役。我们都在奔跑，都在劳苦，都在现今建造；在我们之前，别人也曾奔跑、劳苦、建造；但\Quote{若不是主建造，建造的人就枉然劳力。}因此，宗徒们看见有些人跌倒，就为他们哀哭，因为他们为他们枉费力了。\BibleRef{迦 4:10}我们，因此，在外面说话，祂在里面建造。我们能观察到你们如何留心听我们；唯有祂知道你们的心思，知道你们在想什么。祂亲自建造，祂亲自劝诫，祂亲自开启悟性，祂亲自点燃你们的悟性归于信德；尽管如此，我们作为工人也劳苦；但\Quote{若不是主建造，}等等。\par

2. 但天主的家也是一座城。因为天主的家是天主的子民；因为天主的家是天主的圣殿……这就是耶路撒冷：她有守护者；正如她有建造者，劳苦建造她，她也有守护者。宗徒的这些话关系到这种守护：\BibleQuote{我恐怕，恐怕你们的心意会从基督里的单纯中受到败坏。}\BibleRef{格后 11:3}他是在守护教会。他尽力看守那些在他治下的人。主教们也这样做。因为更高的地位正是为此赐予主教的，使他们自己成为监督，如同子民的守护者。因为希腊文\OriginalTerm{主教}{Episcopus}和本地话的\Quote{监督}是同一个意思；因为主教监督，即是俯视。正如在葡萄园的照料中，园丁被分配了更高的位置，同样主教也被分配了更崇高的地位。而对这个高位，除非我们站在这里，怀着一颗使我们谦卑地俯伏在你们脚下之心，并为你们祈祷，那知道你们心意的那位亲自作你们的守护者，否则我们将面临危险的交代。因为我们能看到你们来去；但我们是如此无法看到你们内心的思想，以至于连你们在家中做什么我们都不能看到。那么，我们如何能守护你们呢？作为人：尽我们所能，尽我们所领受的能力。又因为我们像人一样守护你们，并不能完美地守护你们，难道你们就要没有守护者吗？决不可以！因为那一位在哪里呢？经上说：\BibleQuote{若不是主守护城池，守夜者就枉然警醒？}\BibleRef{咏 126:1}我们警醒守护，但我们的警醒是枉然的，除非那看到你们思想的祂守护你们。祂在你们醒着时守护，祂也在你们睡着时守护。因为祂曾一度在十字架上睡了，又复活了；祂不再睡觉。要成为以色列：因为\Quote{以色列的守护者不打盹，也不睡觉。}是的，弟兄们，如果我们愿意被保护在天主翅膀的荫下，让我们成为以色列。因为我们在管家的职分上守护你们；但我们希望与你们一同被守护。我们如同你们的牧者；但在那位牧者之下，我们与你们同是羊群。我们如同你们的教师站在这个讲台上；但在那唯一的师尊之下，我们与你们在这所学校里是同窗。\par

3. 若我们愿意受那位为我们的缘故而自谦自卑、后又为护佑我们而被高举的那一位所护卫，让我们也谦卑。任何人不可自恃己力。人没有任何善，除非他从那唯一善者领受。但谁若妄自称智，便是愚人。让他谦卑，好使智慧来临并光照他。但若智慧尚未临到他，他便自以为智，他就在光明之前升起，行走在黑暗中。他在这篇圣咏中听到了什么？\BibleQuote{你们清晨早起，实在徒劳} \BibleRef{咏 127:2}。这是什么意思？若你在光明升起之前起身，必然徒劳，因为你将在黑暗里。我们的光明、基督，已经升起；在基督之后升起是好的，不要在基督之前升起。谁在基督之前升起？是那些选择将自己置于基督之前的人。谁又愿意将自己置于基督之前？是那些愿意在此处——祂曾谦卑之处——被高举的人。因此，让他们在此处谦卑，若他们愿意在那里——基督被高举之处——被高举。……主召回\Person{载伯德}的两个儿子，使他们谦卑，并对他们说：\BibleQuote{你们能饮我将饮的杯吗？} \BibleRef{玛 20:21}。我来是为谦卑：你们却愿意在我之前被高举吗？祂说：我所走的路，你们跟随吧。因为若你们选择走我不走的路，你们在黎明前起床，便是徒劳。\Person{伯多禄}也在光明之前升起了，当时他想要给主出主意，劝阻祂为我们受苦。……但我们的主做了什么？祂使他跟随光明之后：\BibleQuote{退到我身后去，撒殚。} \BibleRef{玛 16:23}。他是撒殚，因为他想升于光明之前。\BibleQuote{退到我身后去：}好让我在前，你跟随：我去哪里，你也去那里；不要想引领我到你愿去的地方。……\par

4. 仿佛你要说：我们何时起来？我们现今被命坐下：我们的起来将在何时？当主起来之时。仰望祂，祂在你之前先行：因为若你不听从祂，\BibleQuote{你在黎明前起床，实在徒劳。} \BibleRef{咏 127:2}。祂何时被举起？当祂死了之后。因此，盼望着你死后被举起：盼望死人复活，因为祂复活了并升了天。但祂在哪里睡了？在十字架上。当祂在十字架上睡了时，祂带来了一个预象，是的，祂实现了在\Person{亚当}身上所预示的：因为当亚当睡着时，从他身上取出一根肋骨，\Person{厄娃}被创造了；\BibleRef{创 2:21}同样，当主在十字架上睡着时，祂的肋膀被长矛刺透，圣事流了出来，\BibleRef{若 19:34}教会由此诞生。因为教会，主的净配，是从祂的肋旁被创造的，正如厄娃是从亚当的肋旁被创造的。但正如厄娃是在他睡着时才从他肋旁被造，同样，教会是在祂死亡时才从祂肋旁被创造。因此，若祂不是在死后才从死人中复活，你除了在此生之后，还能盼望被高举吗？但这篇圣咏要教导你，免得你问：我何时起来？也许在我坐下之前？祂加上了：\BibleQuote{当祂赐予祂所爱的人睡眠时} \BibleRef{咏 127:2}。天主在祂所爱的人安眠之时赐下这恩赐；然后祂所爱的人，即基督的人，将要起来。因为众人确实都要起来，但并非都是祂所爱的人。有全死人的复活；但宗徒怎么说？\BibleQuote{我们众人都要复活，但并非我们众人都要改变。} \BibleRef{格前 15:51}。他们起来受罚：我们起来犹如我们的主起来，使我们能跟随我们的元首，若我们是祂的肢体。……盼望这样的复活；为此做一个基督徒，不要为此世的福乐。因为若你为此世的福乐而愿做基督徒，既然你的光明并未寻求此世的福乐；你便是在光明之前升起；你必定继续留在黑暗里。要改变，跟随你的光明；在祂复活之处起来：先坐下，然后起来，\BibleQuote{当祂赐予祂所爱的人睡眠时。}\par

5. 仿佛你又问：谁是被爱的人？\BibleQuote{看，儿女，胎妊的果实，是主的产业} \BibleRef{咏 127:3}。既然他说\BibleQuote{胎妊的果实}，这些儿女是在痛苦中诞生的。有一位妇人，在她身上，向\Person{厄娃}所说的话\BibleQuote{在痛苦中你要生育儿女}，以属灵的方式彰显出来。教会生育儿女，基督的净配；她若生育他们，必为他们经历产痛。以她为预象，厄娃也被称为\BibleQuote{众生的母亲}。那位说过\BibleQuote{我的孩子们，我为你们再受产痛，直到基督在你们内形成} \BibleRef{迦 4:19}的人，也是那经历产痛者肢体中的一员。但她并未徒然受痛，也未徒然生产：在死人复活时，将有圣洁的苗裔：现今分散在全世界的义人要增多。教会为他们呻吟，教会为他们受产痛；但在那死人复活时，教会的后裔将显现，痛苦和呻吟将要过去。……\par

6. \BibleQuote{犹如箭在能者手中，被射出之人的儿女也是如此。}（\BibleRef{咏 126:4}） 这个遗产是从哪里来的，弟兄们？如此众多的遗产是从哪里来的？有些人从主的手中像箭一样被射出，飞得很远，充满了全地，圣人就是从那里来的。因为这就是关于哪遗产所说的：\Quote{你向我请求，我必将外邦人赐你作产业，地极作你的领土。}\newline{}而这份产业如何扩展增加直到地极呢？因为\Quote{犹如箭在能者手中等等}。箭从弓射出，射出它的手臂越强壮，箭就飞得越远。但有什么比主的投掷更强壮呢？从祂的弓祂射出了祂的宗徒们：没有一处地方是这么强壮的手臂射出的箭达不到的；它已到达了地极。它不再往更远的原因是，人类种族没有更多了。因为祂有如此的力量，即使还有更远的地方，箭能够飞到，祂也会把箭射到那里。被射出之人的儿女也是如此，如同被射出之人一样。\par

7. 或许宗徒们自己也被称为被摇出之人的儿子，亦即先知的儿子。因为先知包含封闭和隐藏的奥秘：他们被摇动，以便从中显明地出来……除非涉及的预言被殷勤地筛过，隐藏的意义岂能临到我们呢？因此所有这些意义在主来临之前都是封闭的。主来了，摇动了这些隐藏的意义，它们就变得显明；先知被摇出了，宗徒就诞生了。既然他们从被摇出的先知而生，宗徒就是被摇出之人的儿子。\newline{}他们像箭放在巨人的手中，已经到达了地极……宗徒作为先知的儿子，就如同箭在能者手中。若祂是强大的，祂就用大能的手将他们摇出；若祂用大能的手将他们摇出，那么祂所摇出的人就已经到达了地极。\par

8. \BibleQuote{从他们那里满足了自己愿望的人是有福的。}（\BibleRef{咏 126:5}） 我的弟兄们，谁从他们那里满足了自己的愿望呢？那不爱世界的人。被世界的愿望充满的人，没有空间容纳他们所宣讲的进入。倒出你所携带的，以便容纳你所没有的。就是说，你渴望财富：你不能从他们那里满足你的愿望；你渴望地上的尊荣，你渴望那些天主甚至赐给牲畜的东西，即暂时的快乐、身体的健康等等；你不会从他们那里满足你的愿望……\BibleQuote{他在城门口与仇敌辩论时，必不羞愧。} 弟兄们，让我们在城门口说话，即让我们所有人都知道我们在说什么。因为选择不在城门口说话的人，他希望他所说的被隐藏，也许希望被隐藏是因为它是邪恶的。若他有信心，就让他公开在城门口说话，正如论及智慧所说：\BibleQuote{她在城门口，在城入口处呼喊。}（\BibleRef{箴 8:3}）只要他们持守无辜的正义，他们就不会羞愧：这就是在城门口宣讲。\newline{}那么谁是在城门口宣讲的人？那在基督内宣讲的人；因为基督是门，我们通过祂进入那座城。（\BibleRef{若 10:9}）……因此，那些反对基督的人在城门外，因为他们寻求自己的荣誉，而非基督的荣誉。但那在城门口宣讲的人，寻求基督的荣誉，而非自己的荣誉；因此，那在城门口宣讲的人说：不要相信我，因为你不会通过我进入，而是通过那门。而那些希望人相信自己的人，希望他们不通过那门进入；如果门向他们关闭，他们徒劳地敲门以求打开，这并不奇怪。\par

\SourceRecord{Translated by J.E. Tweed.；Nicene and Post-Nicene Fathers, First Series；Vol. 8.；Edited by Philip Schaff.；Buffalo, NY: Christian Literature Publishing Co.,；1888.；<http://www.newadvent.org/fathers/1801127.htm>.}

% structure: p0001 source=h1 target=section reason=page-title

\section{圣咏第128篇释义}

1. 费利克斯殉道者，真是费利克斯，即\Quote{幸福}，因他的名字和他的荣冠如此，他的诞辰便是今日，他曾轻视了世界。他因为敬畏上主，便因此幸福、因此有福，因为他的妻子如同地上结果的葡萄树，他的儿女围绕着他的桌子吗？这一切祝福他都完全拥有，却是在那在此被描述者的身体内；并且，因为他以此意理解它们，他轻视了现世的事物，为能领受将来的事物。弟兄们，你们知道，他受的苦并不是其他殉道者所受的死亡。因为他宣认了信仰，便被留待受酷刑；在另一天，他的尸体被发现已无生命。他们关闭了监狱，为他的身体，而不是为他的灵魂。刑吏们发现他已离去；当他们准备施刑时，他们白白地发泄了怒气。他躺卧着，死了，对他们没有知觉，以免受酷刑；对天主却有知觉，为能获得荣冠。弟兄们，他不仅以他的名字，也以永生的酬报而蒙福，难道是因为他爱了这些事吗？\par

2. \BibleQuote{凡敬畏上主、遵行祂道路的人，真是有福的} \BibleRef{咏 127:1}。他向许多人说话；但因这些许多人在基督内是一体，在接下来的话中，他用单数说话：\Quote{因为你要吃你果实中的劳苦。}……当我以复数称呼基督徒时，我在唯一的基督内理解为一。因此你们是许多，而你们是一；我们是许多，而我们是一。我们如何是许多，却又是一？因为我们紧附于祂，我们是祂的肢体；既然我们的元首在天上，为让祂的肢体跟随……让我们因此如此聆听这圣咏，把它视为论及基督的话：我们所有紧附于基督的身体、已成为基督肢体的人，都在上主的道路中行走；让我们以纯洁的敬畏、以存留到永远的敬畏来敬畏上主。……\par

3. \BibleQuote{你要吃你果实中的劳苦} \BibleRef{咏 127:2}。而你们，哦，你这许多而为一者，\Quote{你要吃你果实中的劳苦。}他向那些不理解的人说话似乎反常：因为他本应说，你要吃你劳苦的果实。因为许多人吃他们劳苦的果实。他们在葡萄园里劳作；他们吃的不是劳苦本身，而是从他们的劳苦所产生的。他们围绕结果子的树木劳作：谁会吃劳苦呢？但这些劳苦的果实，这些树木的出产；这才是令农夫喜悦的。\Quote{你要吃你果实中的劳苦} 是什么意思？目前我们有辛劳；果实将来才会来到。但因这些劳苦本身并非没有喜乐，因那望德，我们刚才稍微谈到过：\BibleQuote{在望德中要喜乐，在困苦中要忍耐} \BibleRef{罗 12:12}；目前，正是这些劳苦使我们喜悦，使我们在望德中喜乐。因此，如果我们的辛劳是可吃的，并且也能使我们喜悦；那么，我们的劳苦的果实被吃时，将会如何呢？\BibleQuote{他们边走边哭，撒播种子}，他们曾吃他们的劳苦；当他们\BibleQuote{带着禾捆，欢喜地归来}时，他们将怀着更大的喜悦来吃他们劳苦的果实吗？……\Quote{你是有福的，你必会顺利。}\Quote{你是有福的}是现在时；\Quote{你必会顺利}是将来时。当你吃你果实中的劳苦时，\Quote{你是有福的}；当你达到了你劳苦的果实，\Quote{你必会顺利。}他说了什么？因为如果一切顺利，你将是幸福的；如果你将是幸福的，你也会一切顺利。但在望德与获得之间是有分别的。如果望德如此甘甜，现实将会多么甘甜呢？\par

4. 现在我们来看这句话：\BibleQuote{你的妻子} \BibleRef{咏 127:3}：这是对基督说的。因此，祂的妻子就是教会；祂的教会，祂的妻子，就是我们自己。\BibleQuote{如同结实的葡萄树。}……\BibleQuote{在你房屋的四周。}并非所有人都被称为房屋的四周。因为我问什么是四周。我将说什么？它们是墙壁，坚固的石头，仿佛这样？如果他是说这个肉体的居所，我们也许会从四周来理解。我们所说的房屋的四周，是指那些紧附基督的人。……\par

5. \Quote{你的子女}。妻子与子女是一样的。在肉体的婚姻和婚配中，妻子是一回事，子女是另一回事；但在教会内，她是妻子，同时也是子女。因为宗徒们属于教会，是教会肢体中的成员。所以他们是在祂的妻子内，并且根据他们在祂肢体中所占的部分，就是祂的妻子。那么，为何关于他们这样说：\BibleQuote{当新郎从他们中被带走时，新郎的子女就要禁食}？\BibleRef{玛 9:15} 因此，她是妻子，也是子女。兄弟们，我讲一件奇妙的事。在主的话中，我们发现教会既是祂的兄弟姊妹，也是祂的母亲……因为玛利亚是在祂家的侧旁，而来自童贞玛利亚亲属的亲戚，那些信祂的人，也是在祂家的侧旁；不是因为他们肉体的血缘关系，而是因为他们听了天主的话并遵行了……祂接着说：\BibleQuote{凡遵行我在天之父的旨意的，就是我的兄弟、姊妹和母亲}。\BibleRef{玛 12:48} \Quote{兄弟}，或许是指教会所拥有的男性；\Quote{姊妹}，是指基督在此世祂的肢体中所拥有的女性。\Quote{母亲}，岂不是因为基督本身就在那些基督徒内，教会每日藉着圣洗将他们诞生为基督徒吗？因此，在那些你理解为妻子的人身上，你也理解了母亲，也理解了子女。\par

6. ……因此，这样的子女应当\Quote{围绕}主的\Quote{桌子，如同橄榄枝}。这是一棵完整的葡萄树，是极大的福乐：现在谁还会拒绝在那里呢？当你看见某个亵渎者有妻子、儿女、孙子，而你自己或许没有这些时，不要嫉妒他们；要明辨这应许在你身上也实现了，不过是属灵的。如果我们有，为什么有？因为我们敬畏主。\BibleQuote{看，这样敬畏主的人必蒙祝福}。\BibleRef{咏 127:4} 这个人就是那众人，而众人是一个一人；因为许多人是一体，因为基督是一。\par

7. \BibleQuote{愿主从熙雍祝福你：愿你看见耶路撒冷的美物}。\BibleRef{咏 127:5} 甚至对飞鸟也曾说：\BibleQuote{你们要滋生繁多，充满大地}。\BibleRef{创 1:22} 你愿意把赐给飞鸟的祝福当作大事吗？谁能不知道，这确实是由天主的声音所赐？但若你领受了这些恩赐，就要善用它们；更要想如何养育已出生的，而非只求更多的出生。因为幸福不在于有子女，而在于有好的子女。如果子女出生了，就要劳苦养育他们；如果没有出生，就要感谢天主……你的子女是婴孩：你抚爱婴孩，婴孩抚爱你：他们会永远这样吗？但你希望他们长大，希望他们的年岁增加。但想一想，当一个人年纪到来时，另一个年纪就逝去。当童年到来时，幼年逝去；当青年到来时，童年逝去；当成年到来时，青年逝去；当老年到来时，成年逝去；当死亡到来时，所有年纪都逝去。你希望有多少年龄的更替，就等于希望有多少年龄的死亡。因此，这些事\Quote{不是}。最后，子女为你而生，是要与你共享地上的生命，还是更要排挤你、继承你？你为那些排挤你出生的人而欢喜吗？婴儿出生时，仿佛这样对父母说：\Quote{现在，开始考虑离开这里吧，让我们也在舞台上扮演我们的角色。}因为人类整个诱惑的生命就是一场戏；因为经上说：\Quote{凡活着的人，全然是虚幻。}尽管如此，如果我们为继承我们的子女而欢喜；那么，我们岂不更应为那些我们将与他们一同留下的子女欢喜，并为那位我们为之而生的父亲欢喜——祂不会死亡，好使我们永远与祂同生？这些是耶路撒冷的美物：因为它们\Quote{存在}。我要看耶路撒冷的美物多久？\Quote{你一生所有的日子。}如果你的生命是永远的，你将永远看见耶路撒冷的美物……\par

8. 因为，\BibleQuote{如果我们只在今生寄望于基督，我们就是众人中最可怜的了}。\BibleRef{格前 15:19} 何以殉道者被判给野兽？那是什么善？能述说吗？用什么方法，或什么舌头能讲述？或什么耳朵能听见？那确实，\BibleQuote{眼睛未曾看见，耳朵未曾听见，人心也未曾想到}。\BibleRef{格前 2:9} 只让我们爱，只让我们在恩宠中成长：你们看，那么，争战并不缺少，我们与自己的私欲争战。我们向外与不信、不顺从的人争战；我们向内与肉性的建议和扰乱争战：我们处处仍在争战……那么，这是什么和平？是来自耶路撒冷的和平，因为耶路撒冷被解释为\Quote{和平的异象}。因此，\Quote{愿你看见耶路撒冷的美物}，并且\Quote{你一生所有的日子——愿你看见，}不仅是你的子女，而且是\Quote{你的子子孙孙}。\Quote{你的子女}是什么意思？你在此世所做的善工。谁是\Quote{你的子子孙孙}？你善工的果实。你施舍：这些是你的子女；因你的施舍，你得到永生，这是你的子子孙孙。\Quote{愿你看见你的子子孙孙}；并且\BibleQuote{愿平安归于以色列}，\BibleRef{咏 127:6} 这是圣咏的最后的话……\par

\SourceRecord{Translated by J.E. Tweed.；Nicene and Post-Nicene Fathers, First Series；Vol. 8.；Edited by Philip Schaff.；Buffalo, NY: Christian Literature Publishing Co.,；1888.；<http://www.newadvent.org/fathers/1801128.htm>.}

% structure: p0001 source=h1 target=section reason=page-title

\section{\WorkTitle{对圣咏第一二九篇的讲辞}}

1. 我们所咏唱的这圣咏虽然简短，但正如在福音中记载的匝凯，他是\BibleQuote{身材矮小}，\BibleRef{路 19:2}却在行为上伟大；正如那往银库里投入两文钱的寡妇，钱虽少，爱德却大；\BibleRef{谷 12:42}同样，这圣咏若论字数，是短的；若论含义，是大的。……愿天主之神发言，愿祂对我们说话，愿祂向我们歌唱；无论我们愿意与否起舞，愿祂歌唱。因为正如跳舞之人随节拍移动四肢，凡按天主诫命跳舞的人，也以行为服从声音。那拒绝这样做的，主在福音中说什么？\BibleQuote{我们给你们吹笛，你们却不跳舞；我们给你们唱哀歌，你们却不捶胸。}\BibleRef{玛 11:17}因此，愿祂歌唱；我们信赖天主的仁慈，因为将有那些祂藉以安慰我们的人。因为那些顽固不化、继续作恶的人，尽管听到天主圣言，仍以他们的过犯日日扰乱教会。这圣咏正是谈及这样的人；因为它如此开始。\par

2. \BibleQuote{从我幼年以来，他们就多次与我为敌。}\BibleRef{咏 128:1}教会谈及她所忍受的人；似乎有人问：\Quote{是现在吗？}教会是古老的：自从圣人们被召叫以来，教会就在世上。有一次，教会只在亚伯尔内，他被他邪恶的兄弟加音迫害。\BibleRef{创 4:8}有一次，教会只在哈诺客内，他被从不义者中提升。\BibleRef{创 5:24}有一次，教会只在诺厄的家中，他忍受了所有在洪水中丧亡的人，而方舟独自在波浪上漂浮，到达岸边。\EditorialNote{创世纪六至八章}有一次，教会只在亚巴郎内，我们知道他忍受了恶人什么。教会只在他兄弟的儿子罗特内，以及他在索多玛的家中，他忍受了索多玛的邪恶和乖僻，直到天主从中间释放了他。\EditorialNote{创世纪十三至二十章}教会也开始存在于以色列子民中：她忍受了法郎和埃及人。圣人的数目也开始在教会中，即在以色列子民中；梅瑟和其他圣人忍受了恶的犹太人、以色列子民。我们来到主耶稣基督：福音在圣咏中传报了……因此，免得教会现在惊奇，或免得在教会中有人惊异，谁愿作教会的好肢体，就让他听教会他母亲亲自对他说：我儿，对这些事不要惊奇：\Quote{从我幼年以来，他们就多次与我为敌。}\par

3. \BibleQuote{现在愿以色列说。}她现在似乎是在谈论自己：因为她似乎不是自己开始，而是回答。但她在回答谁？回答那些认为并说：我们忍受了多少邪恶，每天加增的绊脚石是何等大，正如恶人进入教会，我们必须忍受他们？但让教会通过某些人，即通过更强者的声音，回答弱者的抱怨，让坚固的人扶助不坚固的人，成熟的人扶助婴孩，让教会说：\BibleQuote{从我幼年以来，他们就多次与我为难。}\BibleRef{咏 128:2}让教会说这话：让她不要害怕。因为在这话\Quote{从我幼年以来}之后，加上\Quote{他们就多次与我为敌}是什么意思？目前教会的老年被攻击：但让她不要害怕。难道她没有到达老年，因为从幼年起他们就一直与她为敌？他们能把她抹去吗？让以色列安慰自己，让教会用过去的例子安慰自己。为什么他们与我为敌？\Quote{因为他们不能胜过我。}\par

4. \BibleQuote{罪人在我的背上筑巢；他们远远地施行了他们的邪恶。}\BibleRef{咏 128:3}为什么他们与我为敌？因为\Quote{他们不能胜过我。}这是什么意思？他们不能在我之上建造。我不同意他们去犯罪。因为每个恶人迫害善人，正是因为这个原因，善人不与恶人同意作恶。假设他作恶，而主教不惩戒他，主教是善人；假设主教惩戒他，主教是恶人。假设他偷窃任何东西，被盗者沉默，他是善人；只要他说话并斥责，即使他不追回财物，他都是坏透了的。那么，斥责盗贼的人是恶人，而盗贼是善人！……不要注意这样的人对你说话：是一个恶人藉着祂对你说话；但对你说话的天主之言不是恶的。控告天主：如果你能，就控告祂。\par

5. 你指控一个人贪婪，他却指控天主，因为天主造了金子。不可贪财。而你回答说，天主不应造金子。如今还剩下，因为不能克制你的恶行，你指控天主的善工：宇宙的创造者与建筑师令你不悦。祂也不应造太阳；因为许多人因窗子的采光而争讼，互相拖到法庭。哦，但愿我们能克制自身的恶习！因为一切都是好的，因为善的天主造了万物：祂的受造物赞美祂，当拥有默观精神的、虔诚与智慧精神的人思考它们的善时……\par

6. 不可放贷取利。你控告圣经，它说\BibleQuote{他没有放债取利。}这不是我写的；不是从我口中先出的：请听天主。他回答说：不要让圣职人员放贷取利。那对你说话的人，或许自己不取利；但如果他确实取利，假设他取利，那藉着他说话的天主会取利吗？如果他做了他命令你做的事，而你没有做；你将进入火焰，他进入天国。如果他不做他命令你做的事，并且和你一样行恶，宣讲他没有履行的本分；你们两人都将同样进入火焰。干草要烧毁；但\BibleQuote{主的话永远长存。}\par

7. \Quote{公义的主必砍断罪人的颈项} \BibleRef{咏 128:4}。……我们中间谁不像税吏那样双目垂地，说：\Quote{主，可怜我这个罪人吧}？ \BibleRef{路 18:13}。因此，如果众人都是罪人，没有一人是无罪的，那么所有人都必须惧怕悬在颈项上的刀剑，因为\Quote{公义的主必砍断罪人的颈项}。我的弟兄们，我并不认为这是指所有罪人；但祂所击打的肢体，正标明祂击打的是何种罪人。因为经上并未说：公义的主必砍断罪人的手或脚；而是因为经上意指的乃是骄傲的罪人，而所有骄傲的人都高昂着头颈，不仅行恶，甚至不肯承认那是恶行，且在受责备时自义……正如约伯传所载（他正论及一个不虔的罪人）：\Quote{他挺着颈项，用盾牌的厚凸面奔向天主}； \BibleRef{约 15:26}。所以此处提到颈项，正是因为你如此高抬自己，不肯双目垂地，捶胸哀痛。你本该向祂呼号，如同另一篇圣咏所呼号的：\Quote{我说：主，求你怜悯我，因为我得罪了你}。既然你不肯这样说，反因天主的言语而自义你的行为，那么经上接着的话就临到你：公义的主必砍断罪人的颈项。\par

8. \Quote{愿所有仇视熙雍的人，都蒙羞退却} \BibleRef{咏 128:5}。他们仇视熙雍，就是仇视教会：熙雍就是教会。那些伪善地进入教会的人，就是仇视教会。那些不肯遵守天主圣言的人，就是仇视教会：\Quote{他们在我背上修建}；教会除了忍耐这重担直到终结，还能做什么呢？\par

9. 然而关于他们，经上说什么呢？接下来的话是：\Quote{愿他们像屋顶上的草，尚未拔出就已枯萎} \BibleRef{咏 128:6}。屋顶上的草，就是长在屋瓦顶上的草：它看似在高处，却没有根。它若长得低些，岂不更好，开花时岂不更令人喜悦？然而它长得越高，枯萎得越快。它尚未被拔出，却已枯萎；他们尚未从天主的审判领受判决，却已失去生机的汁液。观察他们的行为，你就看见他们枯萎了……收割者将到来，但这些人却无法装满他们的禾捆。因为收割者将到来，把麦子收进谷仓，将莠子捆成捆，丢进火里。屋顶上的草也是如此被清掉，凡从其上拔起的，都丢进火里；因为它尚未拔出就已枯萎。收割者不能从那里装满自己的手。接着的话是：\Quote{收割的人不满手拿取，捆禾的人也不满怀抱取} \BibleRef{咏 128:7}。主说：\Quote{收割者是天使} \BibleRef{玛 13:39}。\par

10. \Quote{为此，过路的人连一句：愿主的祝福临于你们，我们以上主的名祝福了你们，也不说} \BibleRef{咏 128:8}。因为你们知道，弟兄们，当人路过别人工作时，习惯说：\Quote{愿主的祝福临于你们}。这在犹太民族中尤其常见。没有人路过而看见有人在田里、葡萄园、收割或任何类似工作，却不祝福便走过的……这些过路的人是谁？就是已经借这条路（即此生）而经过归于家乡的人：宗徒们是此生中的过路人，先知们也是过路人。先知和宗徒们祝福了谁？是那些在他们身上看到爱德之根的人。但那些他们发现在屋顶上高高抬起、在盾牌的厚凸面中骄傲的人，他们宣告这些人注定要成为什么，却不给他们祝福。因此，你们在圣经中读到，所有教会所背负的这些恶人，都被宣布为受诅咒的，属于假基督，属于魔鬼，属于糠秕，属于莠子……但那些说：\Quote{唯有天主圣化，人非因天主的恩赐不能成为善人}的人，他们是以上主的名祝福，不是以自己的名；因为他们是新郎的朋友 \BibleRef{若 3:29}，他们不肯成为新娘的奸夫。\par

\SourceRecord{Translated by J.E. Tweed.；Nicene and Post-Nicene Fathers, First Series；Vol. 8.；Edited by Philip Schaff.；Buffalo, NY: Christian Literature Publishing Co.,；1888.；<http://www.newadvent.org/fathers/1801129.htm>.}

% structure: p0001 source=h1 target=section reason=page-title

\section{《圣咏第一百三十篇释义》}

1. \Quote{我从深渊呼求祢，上主；主啊，求祢听我的声音}\BibleRef{咏 129:1}。约纳从深渊呼求，从鲸鱼的腹中。\BibleRef{纳 2:2}他不仅在水波之下，也在野兽的内脏中；然而，那些波涛和那身体并未阻止他的祈祷到达天主，野兽的腹中也未能容纳他祈祷的声音。它穿透了一切，冲破了所有，到达了天主的耳朵——若我们真可以说它冲破一切而到达天主的耳朵，因为祈祷者心中的耳朵就是天主的耳朵。因为，哪里没有祂临在呢？凡是以信实祈祷的人，天主就在那里。然而，我们也应当明白我们从什么深渊呼求上主。因为此尘世生命就是我们的深渊。凡认识到自己处于深渊的人，就呼号、呻吟、叹息，直到他被救出深渊，来到那坐在诸深渊之上的祂面前。……因为那些甚至不从深渊呼求的人，就在深渊的极深处。圣经说：\Quote{恶人到了罪恶的深渊，便藐视一切。}\BibleRef{箴 18:3}现在请思考，弟兄们，那藐视天主的深渊是何等样？当一个人看到自己被日复一日的罪恶所淹没，被如山般堆积的邪恶所压迫时，若有人对他说：\Quote{向天主祈祷吧}，他便会嘲笑。他会怎么说？首先，他会说：\Quote{如果罪行使天主不悦，我还能活着吗？如果天主眷顾人间事，考虑到我犯下的大罪，我岂止是活着，岂不应该兴旺发达？}因为这是那些远离深渊、在邪恶中兴旺发达之人的常见遭遇；他们越陷于深坑，就越是自以为幸福；因为虚假的幸福本身就是更大的不幸。……\par

2. \Quote{主啊，求祢听我的声音。愿祢的耳朵细听我诉怨的声音}\BibleRef{咏 129:2}。他从哪里呼求？从深渊。那么是谁在呼求？一个罪人。他带着什么希望呼求？因为那位来解除罪孽的，给了深渊中的罪人希望。因此，在这些话之后接着的是：\Quote{上主，祢若细察罪孽，主啊，谁能站立得住？}\BibleRef{咏 129:3}。这样，他揭示了他从什么深渊呼求。因为他在罪恶的重压和波涛之下呼号。……他并没有说：\Quote{我站立不住}，而是说：\Quote{谁能站立得住？}因为他看到几乎整个人类生活处处被自己的罪恶所围困，所有的良心都被自己的思想所控告，一颗依靠自己正义的清洁之心是找不到的。\par

3. 但希望在哪里？\Quote{因为祢有宽恕之能}\BibleRef{咏 129:4}。这宽恕是什么？难道不是祭献吗？祭献又是什么？岂不是那为我们所奉献的？无辜之血的倾流涂抹了所有有罪者的罪过；所偿付的如此巨大赎价，救赎了一切被仇敌掳掠的俘虏。因此，\Quote{祢有宽恕之能}。因为若祢没有怜悯，若祢只愿作审判者，拒绝仁慈，祢就会细察我们所有的罪过，追究它们。谁能承受这个？谁能站在祢面前说：\Quote{我是无辜的？}谁能站立在祢的审判前？所以只有一个希望：\Quote{为了祢的法律，我承受了祢，上主。}什么法律？那使人成为罪人的法律。因为那\Quote{圣洁、公义、良善的法律}\BibleRef{罗 7:12}曾赐给犹太人；但其效果是使他们成为罪人。法律没有被赐下而能赋予生命\BibleRef{迦 3:21}，而只是向罪人显明他的罪。因为罪人已忘记了自己，看不见自己；法律被赐给他，好让他看见自己。法律使他成为罪人，而立法者释放了他：因为立法者是至高权能。……因此有天主仁慈的法律，有天主宽恕的法律。一个是恐惧的法律，另一个是爱的法律。爱的法律赦免罪过，涂抹过去，警告将来；不遗弃路途中的同伴，成为同路人引导他走在路上。但必须与对头和解，当你与他同行在路上时。\BibleRef{玛 5:25}因为天主的圣言是你的对头，只要你不同意它。当你开始以遵行天主圣言的命令为乐时，你就同意了它。于是，曾是你的对头的，成了你的朋友；这样，当路途结束，就没有人把你交到审判官手里。因此，\Quote{为了祢的法律，我等候了祢，上主}，因为祢屈尊引入了一个仁慈的法律，赦免我所有的罪，为我将来给出警告，使我不再冒犯。……因此，\Quote{为了}这个\Quote{法律，我等候了祢，上主}。我等候直到祢来临，使我从一切匮乏中得释，因为在我匮乏中，祢也没有放弃仁慈的法律。……\Quote{我的灵魂等候了祢的圣言。}……\par

4. 因此我们毫无恐惧地信赖那不能欺骗的圣言。\Quote{我的灵魂信赖上主，从晨更直到夜晚}\BibleRef{咏 129:5}。这晨更是夜晚的结束。因此我们必须如此理解，使我们不致以为我们信赖上主只有一天。那么，弟兄们，你们认为这意思是什么？这句话的意思是：上主，通过祂我们的罪过被赦免，在晨更时从死者中复活，好使我们希望那在主内预先发生的在我们身上也要发生。因为我们的罪已经得到赦免，但我们尚未复活；若我们尚未复活，那在我们元首内预先发生的事尚未在我们身上成就。在我们元首内预先发生了什么事？因为那元首的肉体复活了；那元首的灵死了吗？祂身上死了的，复活了。祂在第三日复活了；上主仿佛这样对我们说：\Quote{你们在我身上所见到的，也要在你们自己身上盼望；也就是说，因为我从死者中复活了，你们也要复活。}\par

5. 但有人说：看，主已复活了；难道我因此也要希望我也能复活吗？当然，为此：因为主是在祂从你取来的那（身体）中复活的。倘若祂没有死，祂就不会复活；除非祂带着肉身，祂也不能死。主从你取了什么？肉身。祂自己是谁？天主的圣言，祂在万有之先，万有是藉着祂而造的。但为使祂能从你领受什么，\BibleQuote{圣言成了血肉，寄居在我们中间}。祂从你领受了，为你的缘故而奉献；正如司铎从你领受，当你想为自己的罪平息天主时，他便能为你奉献。这事已经完成了，就这样完成了。我们的司铎从我们领受了祂能为我们奉献的：因为祂从我们领受了肉身，在肉身内祂成为牺牲，成为全燔祭，成为祭献。在苦难中，祂成为祭献；在复活中，祂重造了那被宰杀的，并作为初果献于天主，并对你说：凡属于你的，如今都已祝圣了：既然这样的初果已从你献于天主，因此你要盼望那在你初果中先行的事，也将在你身上发生。\par

6. 既然祂在清晨的守望中复活，我们的灵魂便从此开始盼望：盼望到何时呢？\Quote{直到夜间}，直到我们死去；因为我们肉身的死亡就如同睡眠……\par

7. 他又回到这句话：\BibleQuote{从清晨的守望起，愿以色列寄望于上主。}不仅是\BibleQuote{愿以色列寄望}，而是\BibleQuote{从清晨的守望起，愿以色列寄望}。难道我指责世界的盼望，当它置于主内时？不；但以色列另有一种盼望。愿以色列不盼望财富作为其至善，不盼望身体的健康，不盼望地上事物的丰盈：他确实要在此处受苦，如果他的命运是为真理的缘故遭受任何患难……\par

8. \BibleQuote{因为上主有仁慈，与祂同在的还有丰富的救赎}\BibleRef{咏 129:7}。令人赞叹！这句话在其恰当的位置上，因着\Quote{从清晨的守望}这句话，再没有更好的说法了。为什么？因为主从清晨的守望中复活了；身体应当盼望那在头里先行的事。但是，恐怕有人会这样想：头之所以能复活，是因为祂没有被罪恶压垮，祂内没有罪；我们怎么办呢？我们既然被自己的罪恶压垮，还能盼望像主那样复活吗？但是请看下文：\BibleQuote{祂要救赎以色列脱离他的一切罪恶}\BibleRef{咏 129:8}。因此，他虽然被自己的罪恶压垮，天主的仁慈仍临于他。为此，祂先行于前，没有罪，为能涂抹跟随祂之人的罪。不要信赖自己，而要信赖从清晨的守望……\par

\SourceRecord{Translated by J.E. Tweed.；Nicene and Post-Nicene Fathers, First Series；Vol. 8.；Edited by Philip Schaff.；Buffalo, NY: Christian Literature Publishing Co.,；1888.；<http://www.newadvent.org/fathers/1801130.htm>.}

% structure: p0001 source=h1 target=section reason=page-title

\section{《圣咏第一百三十一篇释义》}

1. 在这篇圣咏中，天主的仆人兼忠信者的谦卑向我们推荐，藉那人的声音而歌唱；这谦卑就是基督的整个身体。因为我们常常警告你们，亲爱的，它不应被理解为一个人的歌唱的声音，而是所有在基督身体内的人的声音。既然所有人都在祂的身体内，就如一个人说话：这人是一个，但也是多人……如今，凡在天主圣殿里祈祷的，就是在教会的平安、基督身体的合一中祈祷；这基督的身体由全世界许多信者组成：因此，凡在圣殿里祈祷的，必蒙垂听。因为凡在教会的平安中祈祷的，就是在\BibleQuote{心神和真理中祈祷}\BibleRef{若 4:21}；而不是在那座作为预像的圣殿中……\par

2. \BibleQuote{主，我的心不骄矜}\BibleRef{咏 130:1}。他献上了祭献。我们如何证明他献上了祭献？因为谦卑的心是祭献……若没有祭献，就没有司铎。但如果我们有一位在天上的大司祭，祂为我们向圣父转求（因为祂进入了至圣所，在帐幔之内），……我们是安全的，因为我们有一位司铎；让我们在那里献上我们的祭献。让我们思考我们应当献上什么样的祭献；因为天主不喜悦燔祭，正如你们在圣咏中所听到的。但在那处他接着显示他所献上的：\Quote{天主的祭献就是痛悔的精神：破碎和痛悔的心，天主啊，祢必不轻看。}\par

3. \BibleQuote{主，我的心不骄矜，我的眼不高傲}\BibleRef{咏 130:1}；\BibleQuote{我未行伟大之事，也未行超过我的奇妙之事}\BibleRef{咏 130:2}。这应被更明白地讲述和倾听。我未曾骄傲：我不愿因奇能而被人知道；我也未曾寻求超过我能力的事，好使我在无知者中夸耀自己。如同那术士西满，他渴望进入超过自己的奇迹，因此宗徒的能力更使他喜悦，而非基督徒的正义……他说：超过我能力的事，我未曾寻求；我未曾在那里伸展自己，我未曾选择在那里被尊大。在丰富的恩宠中自我高举何等可怕，叫没有人能在天主的恩赐中自夸，而应保持谦卑，并遵行所记载的：\BibleQuote{你越伟大，越当谦卑自己，你便会在主前蒙恩。}\BibleRef{德 3:18} 关于天主的恩赐中的骄傲应何等惧怕，我们必须多次向你们强调……\par

4. \BibleQuote{若我没有谦卑的思想，反而高举我的灵魂，如同一个从母亲胸前被取走的婴孩，这就是对我灵魂的赏报}\BibleRef{咏 130:2}。他好像以咒诅咒束缚自己：……好像他正要说：愿这事发生在我身上。\BibleQuote{如同从母亲胸前被取走的婴孩，但愿这也是我灵魂的赏报。}你们知道，宗徒对一些软弱的弟兄说：\BibleQuote{我喂养你们的是奶，不是饭；因为直到如今，你们还不能吃，就是如今还是不能。}\BibleRef{格前 3:2} 有一些软弱的人不适合强壮的饭食；他们渴望抓取那不能接受的：如果他们果真接受了，或自以为接受了未接受的东西，他们就会因此自高自大，进而变得骄傲；他们自以为聪明。这一切都发生在所有异端者身上；他们既然属血气、属肉体，因辩护自己败坏的见解（他们看不出那是错的），就被逐出天主教会……\par

5. 确实还有人持有另一种意见，并对这些话有另一种理解……诸位弟兄，已经清楚地解释了天主愿意我们在何处谦卑，在何处崇高。谦卑，为了防备骄傲；崇高，为了接纳智慧。以奶为食，使你得滋养；得滋养，使你成长；成长，使你能吃饼。但当你开始吃饼时，你将断奶，也就是说，你将不再需要奶，而是需要固体食物。他似乎是这个意思：\BibleQuote{若我没有谦卑的思想，反而高举我的灵魂：}也就是说，如果我在心思上不是一个婴孩，我就属于邪恶。在这个意义上，他之前说过：\BibleQuote{主，我的心不骄矜，我的眼不高傲：我未行伟大之事，也未行超过我的奇妙之事。}看，在邪恶上我乃是一个婴孩。但既然我在理解上并非婴孩，\BibleQuote{若我没有谦卑的思想，反而高举我的灵魂，}愿那赐给从母亲胸前断奶的婴孩的赏报归于我，好使我能最终吃饼。\par

6. 这个解释，弟兄们，我也不反对，因为它并不与信仰冲突。然而我不能不指出，经文不仅说：\Quote{如同从乳旁被夺去的人，这可能是我灵魂的赏报}，而且还有这补充：\Quote{如同在母亲怀中从乳旁被夺去的人，这可能是我灵魂的赏报}。这里有些东西使我考虑这是一个诅咒。因为这不是一个婴孩，而是一个长大的孩子从乳旁被夺去；在他最初的婴幼期，即真正的婴儿期，他是在母亲的怀中：如果他不巧从乳旁被夺去，他就会死亡。所以加上\Quote{在母亲怀中}不是没有理由的。因为所有的人都可以因成长而断乳。一个成长并这样从乳旁被夺去的人，对他是有益的；但对一个仍在母亲怀中的人则是有害的。因此，我弟兄们，我们必须谨防并惧怕，免得有人在他时日未到之前就从乳旁被夺去……所以他不要想举起自己的灵魂，当他不适合吃干粮时，而要履行谦逊的诫命。他有可以操练自己的地方：让他相信基督，好使他能理解基督。他不能看见圣言，不能理解圣言与圣父的平等，他还不能看见圣神与圣父及圣言的平等；让他相信这事，并用这来哺育自己。他是安全的，因为当他成长后，他就能吃，这是他通过哺育成长之前所不能做的；他有一个目标可以伸展。不要探究超出你能力的事，不要寻求超越你力量的事，也就是说，那些你尚未准备好去理解的事。但我该做什么？你回答。难道我就这样待着吗？ \Quote{但主所命令你的事，你要常常思念。} \BibleRef{德 3:22} 主命令了你什么？行仁慈的工作，不破坏教会的和平，不依靠人，不因渴求奇迹而试探天主……\par

7. 因为如果你不自高，如果你不把心抬得太高，如果你不踏足于那些对你过高的大事，而是保持谦逊，天主就会向你启示你所不同意的事。\BibleRef{斐 3:15} 但如果你选择为这你所不同意的事辩护，并顽固地断言它，且反对教会的和平；那么他所描述的诅咒就会临到你身上：当你在母亲怀中并从乳旁被夺去时，你要离开母亲的怀中饿死。但如果你坚持在教会的和平中，如果你可能在某件事上与你所应当的不同，天主就会启示给你，只要你谦逊。为什么？因为\Quote{天主拒绝骄傲的人，却赐恩宠给谦逊的人。}\par

8. 因此，这篇圣咏以此作结：\Quote{以色列，你要信赖上主，从今时直到永远。} \BibleRef{咏 130:3} 但\OriginalTerm{世界}{seculum}一词并不总是指此世，有时是指永恒；因为永恒有两种理解：直到永远，即要么是永无止息，要么是直到我们抵达永恒。那么在这里应如何理解呢？直到我们抵达永恒，让我们信赖主天主；因为当我们到达永恒时，将不再有希望，而是事情本身将成为我们的。\par

\SourceRecord{Translated by J.E. Tweed.；Nicene and Post-Nicene Fathers, First Series；Vol. 8.；Edited by Philip Schaff.；Buffalo, NY: Christian Literature Publishing Co.,；1888.；<http://www.newadvent.org/fathers/1801131.htm>.}

% structure: p0001 source=h1 target=section reason=page-title

\section{圣咏第132篇释义}

极亲爱的诸位，我们的兄弟、我的同僚，当他在我们众人面前时，我们实在更应该听他讲话。但刚才他没有拒绝，只是推迟了我们；因为他向我强求，让我现在听他讲话，条件是我也可以听他讲话，因为我们在爱德本身中都在聆听那一位，他是我们在天上的唯一师尊。因此，请留意这篇圣咏，题为\WorkTitle{登阶之歌}，它比同一标题下的其他圣咏要长得多。因此，除非不得已，我们不要停留，以便倘若主允许，我们能解释全篇。因为你们也不该像未经教导的人那样听所有的话；你们应当在一定程度上从过去听过的内容帮助我们，这样就不必每一件事都像新的一样向你们阐述了。\par

2. \BibleQuote{主，请纪念达味和他的全部温良} \BibleRef{咏 131:1}。按历史的真实，达味是一个人，以色列国王，耶西的儿子。他确实是温良的，正如神圣圣经本身所标记和称赞他的，如此温良，以至于他没有向迫害者撒乌耳报恶。他对他保持了如此大的谦卑，以至他承认撒乌耳是君王，而自己是一条狗；他回答君王时并不骄傲或粗鲁，尽管他在主内更强大；他反而试图通过谦卑平息他，而不是通过骄傲激怒他。撒乌耳甚至被交在他手中，这是主天主所为，好让他能任意对待他；但由于没有被命令杀死他，而只是拥有这个权力（人可以使用他的权力），他反而将主给他的权力转向了慈悲。……因此，达味的谦卑被称赞，达味的温良被称赞；并且向天主说：\BibleQuote{主，请纪念达味和他的全部温良}。为了什么目的？\BibleQuote{他怎样向主发誓，向雅各伯的全能天主许愿} \BibleRef{咏 131:2}。因此纪念这一点，好使他能完成他所许诺的。达味自己发誓好像他有权能，他祈求天主实现他的誓愿：在誓愿中有热忱，但在祈祷中有谦卑。没有人应该妄想靠自己的力量完成他所许的愿。那劝你许愿的，他自己帮助你实现。因此，让我们看看他许了什么愿，由此我们明白达味如何应在预像中被理解。\Person{达味}被解释为\OriginalTerm{强壮的手}{Strong of hand}，因为他是一位伟大的战士。的确，他信赖主他的天主，发动了所有战争，推倒了所有敌人，天主帮助他，按照那王国的安排；然而他预先描绘了一位强壮的手，要摧毁他的敌人，魔鬼和他的天使。这些敌人是教会与之争战并战胜的。……那么，\Quote{他怎样发誓}等等是什么意思？让我们看看这是什么愿。我们没有什么比发誓更能取悦天主的。如今发誓就是坚定地应许。考虑这个愿，即他以怎样的热情发誓，以怎样的爱，以怎样的渴望；然而，他用这些话祈求主来实现它：\BibleQuote{主啊，请纪念达味和他的全部温良。} 在这种心情下，他发了他的愿，并且应有一个天主的殿：\BibleQuote{我绝不进入我房屋的帐幕，也不上床休息} \BibleRef{咏 131:3}。\BibleQuote{我不让我的眼睛睡觉，也不让我的眼睑打盹} \BibleRef{咏 131:4}。这似乎不够；他加上：\BibleQuote{也不让我的太阳穴歇息，直到我为主寻得一个地方，为雅各伯的天主寻得一个居所} \BibleRef{咏 131:5}。他在哪里寻求天主的处所？如果他是温良的，他在自己内寻求。因为一个人如何成为天主的处所？听先知的话：\BibleQuote{我的圣神要安息在谁身上？在他那贫穷和忏悔的精神，并在我的话语前战栗的人身上} \BibleRef{依 66:2}。你希望成为天主的处所吗？你就要成为贫穷、忏悔、在天主的话语前战栗的人，你自己将成为你所寻求的。因为如果你所寻求的没有在你身上实现，在别人身上又对你有何益处呢？……\par

3. 我的诸位弟兄，有多少千人相信了，当他们把财产的价值放在宗徒脚前时！但圣经对他们说了什么？他们确实成为了天主的圣殿；不仅每个人各自成为天主的圣殿，而且所有人一起成为天主的圣殿。因此，他们成了给主的处所。而且，为使你们知道一处处所在众人中为天主所成，圣经说：他们一心一意朝向天主。但许多人，为了不为主预备处所，寻求自己的事，爱自己的事，以自己的权力为乐，贪图自己的私利。然而，那愿为主预备处所的人，不应以私利为乐，而应以公益为乐。……\par

4. 因此，诸位弟兄，让我们戒绝私有财产的拥有；或者，如果我们不能戒绝拥有，至少戒绝对其的爱恋；这样我们就为天主预备处所。有人说：这对我来说太难了。但请你想想你是谁，你即将为天主预备处所。如果某个议员想要在你家被款待，我不是说议员，而是这世上某个大人物的代表，并且说：你家中有东西冒犯了我；尽管你喜爱它，你仍会将其移除，以免冒犯他，因为你正追求他的友谊。而人的友谊对你有什么益处呢？……毫无畏惧地渴望基督的友谊：祂愿意在你家被款待；为祂腾出空间。什么是为祂腾出空间？不要爱你自己，要爱祂。如果你爱你自己，你就向祂关闭了门；如果你爱祂，你就向祂敞开：并且如果你敞开，祂进入，你就不会因爱自己而失落，反而会在爱你的祂那里找到你自己。……\par

5. \BibleQuote{看，我们听说约柜在\OriginalTerm{厄弗辣大}{Ephrata}} \BibleRef{咏 131:6}。什么？为主预备一个地方。\BibleQuote{我们听说在厄弗辣大，在森林平原找到了它。} 他是听见了又在那里找到了吗？还是在别处听见，在另一处找到？因此让我们探究厄弗辣大是什么，他在那里听见；也探究森林平原是什么意思，他在那里找到。厄弗辣大是一个希伯来词，在拉丁文中译为\OriginalTerm{镜子}{Speculum}，正如圣经中希伯来词的译者所传给我们，让我们明白。他们从希伯来语译成希腊语，从希腊语我们有了拉丁文译本。因为有人在圣经中守望。因此如果厄弗辣大意即镜子，那么在森林平原找到的那座房屋，是在镜子里听说的。镜子有影像：一切预言都是未来事物的影像。因此，未来的天主的房屋在预言的影像中被宣告。\BibleQuote{我们在森林平原找到了它。} 什么是\BibleQuote{森林平原}？\OriginalTerm{saltus}{Saltus}这里不是通常的意思，指几百亩的土地；\OriginalTerm{saltus}{Saltus}本义指尚未耕种、多树木的地方。因为有些读本作\BibleQuote{森林中的平原}。那么什么是森林平原，不就是尚未开化的万民吗？它们是什么，不就是尚被偶像崇拜的荆棘覆盖的区域吗？因此，虽然那里有偶像崇拜的荆棘，我们仍在那里为主找到一个地方，为雅各伯的天主安置帐幕。向犹太人宣告的影像，在外邦人的信德中显现出来。\par

6. \BibleQuote{我们要进入祂的帐幕} \BibleRef{咏 131:7}。谁的？雅各伯的天主。进去居住的人，就是那些进去以便被居住的人。你进入你的房屋，以便居住其中；进入天主的房屋，以便被祂居住。因为主更好，当祂开始住在你内时，祂会使你幸福。因为如果你不被祂居住，你将是可怜的。那个儿子说：\BibleQuote{父亲，请把我应得的家业分给我}，等等，\BibleRef{路 15:12}，他希望自己做主。在父亲手中保管得好，免得被浪费在娼妓身上。他得到了，给了他自己的权力；去远方，全都挥霍在娼妓身上。最后他挨饿，想起父亲；回来，为得饱食。因此进入，以便被居住；可以说不属于自己，而属于祂：\Quote{我们要进入祂的帐幕。我们要在祂脚站立的地方朝拜。} 谁的脚？主的，还是那房屋本身的？因为那是主的房屋，祂说应该在祂里面朝拜。在祂的房屋之外，主不听以至于永生；因为属于天主房屋的人，是在爱德中被活石建筑在里面的。而谁没有爱德，就会跌倒；当他跌倒时，房屋站立……\par

7. 但如果你倾向于理解为房屋本身，即那房屋的脚所站立之处；让你的脚站立在基督内。它们就会站立，如果你在基督内持守。因为关于魔鬼说了什么？\BibleQuote{他从起初就是凶手，不站在真理中。} \BibleRef{若 8:44} 魔鬼的脚因此没有站立。也关于骄傲之人说了什么？\Quote{求祢不要让骄傲的脚来攻击我，不要让恶人的手使我颠仆。那里他们跌倒了，所有作恶的人；他们被推倒，不能站立。} 那么这就是天主的房屋，它的脚站立。因此若望喜乐地说：什么？\Quote{有新娘的是新郎；但新郎的朋友站着听祂。} 如果他不站着，他就听不到祂。他恰当地站着，因为\Quote{他为新郎的声音而喜乐。} 因此现在你们看到为什么他们跌倒，他们因自己的声音而喜乐。那位新郎的朋友说：\BibleQuote{那一位就是施洗者。} \BibleRef{若 1:33} 有些人说：我们施洗；因自己的声音而喜乐，他们不能站立；也不属于那座房屋，关于这房屋说：\Quote{祂的脚所站立之处。}\par

8. \BibleQuote{上主，请起来，进入祢的安息之所} \BibleRef{咏 131:8}。祂对沉睡的主说：\Quote{起来。} 你已经知道谁睡了，谁又复活了。……\BibleQuote{祢和祢圣洁的约柜：} 即，请起来，使祢所圣化的圣洁约柜也起来。祂是我们的元首；祂的约柜是祂的教会：祂先起来，教会也将起来。身体不敢应许自己复活，除非元首先起来。基督的身体，由玛利亚所生，被一些人理解为圣洁的约柜；所以这话的意思是，祢与祢的身体一同起来，使那些不信的人可以触摸。\par

9. \BibleQuote{愿祢的司祭披上正义，愿祢的圣民欢欣歌唱。} \BibleRef{咏 131:9} 当祢从死者中复活，到祢父那里去时，愿那王者的司祭职披上信德，因为\BibleQuote{义人因信德而生活} \BibleRef{罗 1:17}；并且，领受了圣神的抵押，愿肢体在元首先行复活中欢欣喜乐：因为宗徒对他们说：\BibleQuote{在希望中喜乐。} \BibleRef{罗 12:12}\par

10. \Quote{为祢的仆人达味之故，请不要转面不顾祢的受傅者。}\BibleRef{咏 131:10} 这些话是对天主圣父说的。\Quote{为祢的仆人达味之故，请不要转面不顾祢的受傅者。}主在犹太地被钉十字架；祂被犹太人钉死；受他们折磨后，祂安眠了。祂复活起来审判那些在祂安眠时以残暴之手对待祂的人；并且祂在别处说：\Quote{求祢再将我举起，我必酬报他们。}祂已经酬报了他们，也必将继续酬报。犹太人自己很清楚，主死后他们遭受了多大的苦难。他们全都被驱逐出那座他们杀死祂的城。那么怎样？所有从达味根系和犹大支派出来的人都灭亡了吗？不：因为那支派中有些人相信了，事实上那支派中有成千上万的人相信了，并且是在主复活之后。他们曾狂怒地钉死了祂；后来开始看到因被钉者之名所行的奇迹；他们更加战栗，因为祂的名字竟有如此大能，而当初在祂手中时似乎无法行任何奇迹；他们心中刺痛，终于相信在他们曾视同常人者内隐藏着某种神性，于是向宗徒们求教，他们得到回答：\BibleQuote{你们悔改，并要以我们的主耶稣基督之名受洗。}\BibleRef{宗 2:38} 既然基督复活是为审判那些曾钉死祂的人，并转面不顾犹太人，而将祂的容貌转向外邦人；如此，天主似乎是代替以色列的遗民受恳求；而有人向祂说：\Quote{为祢的仆人达味之故，请不要转面不顾祢的受傅者。}如果糠秕被定罪，愿麦子被聚集。愿遗民得救，正如依撒意亚所说：\Quote{遗民}显然\Quote{已得救：}\BibleRef{依 10:21} 因为他们中有十二位宗徒，有五百多位弟兄，主复活后曾显现给他们：\BibleRef{格前 15:6} 他们中有成千上万的人受洗，\BibleRef{宗 2:41} 并把财产的价格放在宗徒脚前。如此，向天主所作的祈祷便应验了：\Quote{为祢的仆人达味之故，请不要转面不顾祢的受傅者。}\par

11. \Quote{主向达味发了忠诚的誓，并且决不后悔，}\BibleRef{咏 131:11} 所谓\Quote{发了誓}是什么意思？是指祂藉着祂自己确认了一个许诺。所谓\Quote{决不后悔}是什么意思？是指祂不会改变。因为天主不会遭受悔恨的痛苦，也不会在任何事情上受骗，以致祂会愿意纠正自己错误之处。但是当一个人后悔某事时，他愿意改变他所做的；因此，当你听到天主后悔时，应当期待一种实际的变化。天主做事的方式与你不同，尽管祂称之为后悔；因为你这样做是因为你犯了错误；而祂这样做是因为祂施行报复或解救。当祂后悔时，祂改变了撒乌耳的王位，正如经上所记：而在圣经说\Quote{祂后悔了}的那段话中，稍后又说\Quote{因为祂不是人，以致后悔}。因此，当祂通过自己不改变的计划改变祂的作为时，祂被称为后悔，乃是因为这个变化，不是出于计划的改变，而是出于作为的改变。但祂许诺了这一点，以致不会改变它。正如这段经文也说：\BibleQuote{主发了誓，决不后悔：你永远按照默基瑟德的品位做司铎；}\BibleRef{咏 110:4} 同样，既然这许诺如此确定而不会改变，因为它必须发生且永久不变；祂说：\Quote{主向达味发了忠诚的誓，并且决不后悔：我要从你腹中的果实安置在你的宝座上。}祂本可以说\Quote{从你腰间的果实}，但祂为什么选择说\Quote{从你腹中的果实}？如果祂那样说，也是真的；但祂选择用更深的意义说\OriginalTerm{出自你腹中的果实}{Ex fructu ventris}，因为基督是由一个女人而生，没有男人。\par

那么如何？\BibleQuote{主向\Person{达味}立了忠信之誓，决不反悔：我必从你身所生的，安置于你的宝座上。如果你的子孙肯遵守我的盟约，遵行我教训他们的法度，他们的子孙也必永远坐在你的宝座上。}\BibleRef{咏 131:12}如果你的子孙遵守我的盟约，他们的子孙也必永远坐在你的宝座上。父母为自己的子孙立下了功绩。假如他的子孙遵守了盟约，而他们的子孙却不遵守，那又如何呢？为何子孙的幸福应许与他们父母的功德相关？因为祂说什么？\Quote{如果你的子孙遵守我的盟约，他们的子孙也必永远坐在你的宝座上}——祂不是这样说：如果你的子孙遵守我的盟约，他们就要坐在你的宝座上；如果他们的子孙遵守我的盟约，他们也要坐在你的宝座上。而是说：\Quote{如果你的子孙遵守我的盟约，他们的子孙也必永远坐在你的宝座上}——除非祂在此愿意以\Quote{子孙}来理解他们的果实？\Quote{如果你的子孙，}祂说，\Quote{肯遵守我的盟约，并遵守我教训他们的法度，他们的子孙也必坐在你的宝座上}：这就是说，他们的果实就是坐在你的宝座上。因为在此生，弟兄们，我们所有在基督内劳苦的人，所有因祂的言语而战栗的人，所有以各种方式努力遵行祂的旨意，并在祈祷祂的帮助时叹息，为能完成祂所命令的人，我们是否已经坐在那些应许给我们的福乐宝座上了呢？不，但我们遵守祂的诫命，希望这会实现。这希望以\Quote{儿子}的形象来描述；因为儿子是今生之人的希望，儿子是他的果实。为此，人们在为自己的贪婪辩解时，声称他们积蓄是为了留给子孙；由于不愿施舍给穷人，便以孝爱为名辩解，因为他们的子孙是他们的希望。因为所有按此世生活的人，都宣称他们的希望是成为能留下后代之父。因此，祂普遍以\Quote{子孙}之名来描述希望，并说：\Quote{如果你的子孙肯遵守我的盟约，并遵守我教训他们的法度，他们的子孙也必永远坐在你的宝座上}：这就是说，他们将获得这样的果实，使他们的希望不致落空，使他们能到达他们所希望到达之处。因此，目前他们如父亲，是寄望于未来的人；但一旦得到了他们所希望的，他们就成了儿子，因为他们已在自己行为中产生并结出了所获的果实。这为他们保留到将来，因为将来本身常以\Quote{子孙}来象征。\par

或者，如果你理解\Quote{子孙}是指实际的人，那么\Quote{如果你的子孙肯遵守我的盟约和我要教训他们的法度}这句话可能意指\Quote{如果你的子孙肯遵守我的盟约和我要教训他们的法度，并且他们的子孙也}，也就是说，如果他们也都遵守我的盟约；所以在此你必须稍作停顿，然后推断出\Quote{他们必永远坐在你的宝座上}，这就是说，你的子孙和他们的子孙，但全都必须遵守我的盟约。那么，如果他们不遵守呢？天主的应许落空了吗？不，但这样说和应许，是因为天主预见了什么？岂不是他们将要相信吗？但为使没有人好像要威胁天主的应许，并宁愿将天主所应许的实现置于自己的权下，所以祂说：\Quote{祂立下了誓}，以此表明这事必定成就。那么，祂在此如何说：\Quote{如果他们肯遵守我的盟约}？不要因应许而夸耀，而忽略你不遵守盟约的事。那么，如果你遵守盟约，你便是\Person{达味}的儿子；但如果你不遵守，你便不是\Person{达味}的儿子。天主曾应许给\Person{达味}的子孙。不要说，我是\Person{达味}的儿子，如果你堕落的话。如果那些从这血统出生的人都不这样说（不，他们这样说，但他们受了迷惑。因为主公开说：\BibleQuote{如果你们是\Person{亚巴郎}的子女，你们就会做\Person{亚巴郎}的事。}\BibleRef{若 8:39}祂因此否认他们是子女，因为他们不做那些事），我们这些人，按肉体并非\Person{达味}的后裔，又怎能自称为\Person{达味}的子女呢？那么，除非我们效法他的信德，除非我们如同他一样敬拜天主，否则我们便不是子女。因此，如果你不因血统而希望得到，却不愿以行为去获得，那么坐\Person{达味}宝座的事怎能成就在你身上？如果它不能成就在你身上，你以为它根本不会成就吗？祂如何在那林地的小径上找到了它？祂的双脚如何站立？无论你可能是谁，那房屋将屹立不倒。\par

\BibleQuote{因为主拣选了\Place{熙雍}作为自己的居所}\BibleRef{咏 131:13}。\Place{熙雍}就是教会本身；她也是那\Place{耶路撒冷}，我们正向她的和平奔跑，她不是在天使中，而是在我们中间作客旅，在她更好的部分等待着将要归回的部分；由此有书信来到我们这里，每日诵读。这城就是\Place{熙雍}，主所拣选的。\par

\BibleQuote{这是我永远的安息之所}\BibleRef{咏 131:14}。这些话是天主的话。\Quote{我的安息}：我在那里安息。天主多么爱我们啊，弟兄们，因为当我们安息时，祂也说祂安息！因为祂自己并非时而烦恼，祂也不像我们那样安息；但祂说祂在那里安息，是因为我们将在祂内获得安息。\Quote{我要住在这里，因为我喜爱这里。}\par

16. \BibleQuote{我必以祝福祝福她的寡妇，用食粮满足她的穷人}\BibleRef{咏 131:15}。每一个意识到自己除了天主以外，毫无所有援助的灵魂，就是寡妇。因为宗徒如何描述寡妇？\BibleQuote{那真正孤寡的寡妇，独自信赖天主。}\BibleRef{弟前 5:5} 他是在谈论我们在教会中称为寡妇的那些人。他说：\BibleQuote{那好宴乐的寡妇，虽生犹死；}他不将她们列入寡妇之中。但在描述真正的寡妇时，他说什么？\BibleQuote{那真正孤寡的寡妇，独自信赖天主，并日夜在恳求和祈祷中坚持。}这里他加上：\BibleQuote{但那好宴乐的，虽生犹死。}那么，是什么使一个妇女成为寡妇？就是她从任何其他来源都没有援助，只从天主而来。有丈夫的，以丈夫的保护为荣；寡妇看似孤苦无依，但她们的援助反倒更为坚强。因此，整个教会就是一个寡妇，无论是男人还是女人，无论是已婚男人还是已婚女人，无论是青年人还是老年人，还是贞女：整个教会就是一个寡妇，在这世界上孤苦无依，如果她感受到这一点，如果她意识到自己的寡妇身份，那么援助就在她身边。我的弟兄们，你们在福音中没有认出这个寡妇吗？当主说\BibleQuote{人应当时常祈祷，不要灰心}时？\BibleQuote{某城中有一个判官，他说：不敬畏天主，也不尊重人。那城中有一个寡妇，她天天来到他那里，说：\InnerQuote{请制裁我的对头，为我伸冤。}}这个寡妇因每日的恳求而说服了他；因为判官心里说：\BibleQuote{我虽不敬畏天主，也不尊重人，但因这寡妇烦扰我，我要给她伸冤。}\BibleRef{路 18:1} 如果邪恶的判官听了寡妇的话，以免被搅扰；天主岂不听祂所劝勉祈祷的教会吗？\par

17. 此外，\BibleQuote{我要用食粮满足她的穷人；}弟兄们，这意味着什么？让我们成为穷人，那么我们就将被满足。许多信赖世界、骄傲自大的人，是基督徒；他们敬拜基督，却没有被满足；因为他们已被满足，且以自己的骄傲为满足。关于他们，经上说：\BibleQuote{我们的灵魂饱受富足人的讥讽和骄傲人的轻慢。}这些人富足有余，因此他们吃，却不满足。关于他们，圣咏中怎么说？\BibleQuote{地上所有肥胖的人，都吃了，且朝拜了。}他们敬拜基督，尊敬基督，向基督祈祷；但他们并没有因祂的智慧和正义而满足。为什么？因为他们不是穷人。因为穷人，即心里谦卑的人，越是饥饿，吃得越多；他们越是空虚于世，就越饥饿。一个饱足的人，你给他任何东西，他都会拒绝，因为他已经饱了。给我一个饥饿的人；给我一个经上所说的人：\BibleQuote{饥渴慕义的人是有福的，因为他们要得饱饫；}\BibleRef{玛 5:6}这些人就是穷人，他刚才说\BibleQuote{我要用食粮满足她的穷人。}因为正是在那篇圣咏中，经上说\BibleQuote{地上所有肥胖的人，都吃了，且朝拜了}；对穷人也有同样的说法，而且与这篇圣咏完全相同：\BibleQuote{穷人吃了，且得饱饫；寻求上主的人将赞美祂。}那里说\BibleQuote{地上所有肥胖的人都吃了，且朝拜了}；这里说\BibleQuote{穷人吃了，且得饱饫}。为什么当说到富人朝拜时，不说他们得饱饫；而当提到穷人时，却说他们得饱饫？他们从哪里得饱饫？弟兄们，这种饱饫的性质是什么？天主自己就是他们的食粮。这食粮降在地上，成为我们的乳，并对自己的人说：\BibleQuote{我是从天上降下的生活的食粮。}\BibleRef{若 6:51}因此圣咏中有这些话：\BibleQuote{穷人吃了，且得饱饫。}他们从哪里得到饱饫？听下面的话：\BibleQuote{寻求上主的人将赞美祂。}\par

18. 因此，要成为穷人，要成为那寡妇的肢体，让你的帮助唯独在于天主。金钱毫无价值；你从它得不到援助。许多人因金钱而跌倒，许多人因金钱而丧亡；许多人因他们的财富而被强盗盯上；他们原本是安全的，如果没有那些使人追踪他们的东西的话。许多人依赖他们更有权势的朋友；他们所依赖的人跌倒了，并使那些信赖他们的人也一同丧亡。回顾人类历史中的实例。我告诉你们的，难道是罕见的事吗？我们不仅从这些圣经中讲这些事；要在全世界中阅读它们。注意，不要信赖金钱、朋友、世界的荣誉和夸耀。远离这一切；但如果你拥有它们，要感谢天主，如果你轻视它们。但如果你因它们而骄傲，不要想你什么时候会成为人的猎物；你早已是魔鬼的猎物。但如果你没有信赖这些事物，你将成为那寡妇——即教会——的肢体，关于她经上说\BibleQuote{我必以祝福祝福她的寡妇}；你也是穷人，是那些经上所说的人之一：\BibleQuote{我要用食粮满足她的穷人。}\par

19. 但有时候——此事不可不提——你会见到一个穷人骄傲，一个富人谦卑：我们每日都忍受这样的人。你听见穷人在富人面前呻吟，当那更有权势的富人压迫他时，你便见他谦卑；有时甚至那时也不谦卑，反而仍然骄傲；由此你可知他若拥有财产，会是什么样子。天主的穷人因此是神贫的，不是囊中羞涩。有时人出门，拥有满屋、富饶的土地、许多田产、大量金银；他知道不可信赖这些，在天主前谦卑自己，用它们行善；这样他的心被举向天主，以致他意识到不仅财富本身对他无益，甚至还会阻碍他的脚步，除非天主治理并援助它们；于是他就算在那些以食粮得饱足的穷人之中。你找到另一个骄傲的乞丐，或者他只是因为一无所有才不骄傲，但仍寻求可以使他骄傲的凭借。天主不看人所拥有的手段，而是看他所有的意愿，并依照他对暂时福乐的意愿审判他，而非依照他无缘拥有的手段。因此宗徒论及富人说：\BibleQuote{你要嘱咐那些今世富裕的人，不要心高气傲，也不要寄望于无常的财富，而要寄望于永生的天主，祂多多赐给我们一切享用。} 那么他们该用财富做什么呢？他接着说：\BibleQuote{要他们行善，在善事上致富，甘心施舍，乐意与人分享。} 并看他们在今世是贫穷的：\BibleQuote{为自己积蓄}，他补充道，\BibleQuote{美好的根基，以备将来，好使他们把握永生。} \BibleRef{弟前 6:17} 当他们把握了永生时，那时他们才会富有；但既然尚未拥有，就当知道自己贫穷。如此，天主将一切内心谦卑、双爱坚固的人——\BibleRef{玛 22:37}——无论他们在今世拥有什么——都算作祂的穷人，祂以食粮使他们饱足。\par

20. \BibleQuote{我要使她的司祭披上救恩，她的圣民要欢欣歌唱。} \BibleRef{咏 132:16} 我们如今已到圣咏末尾；亲爱的诸位，请再听片刻。\BibleQuote{我要使她的司祭披上救恩，她的圣民要欢欣歌唱。} 我们的救恩是谁？岂不是我们的基督？那么\BibleQuote{我要使她的司祭披上救恩}是什么意思？\BibleQuote{你们凡受洗归于基督的，就是穿上了基督。} \BibleRef{迦 3:27} \BibleQuote{她的圣民要欢欣歌唱。} 他们从哪里欢欣歌唱？因为他们披上了救恩：不是凭自己。因为他们成了光，但是在主内；他们从前原是黑暗。\BibleRef{弗 5:8} 因此他加上说：\BibleQuote{在那里我要使达味的角茁壮。} \BibleRef{咏 132:17} 这就是达味的高举，即信赖基督。因为角象征高举：但这是什么高举？不是属血肉的。因此，当所有骨头都裹在肉中时，角却伸出肉外。属神的高度就是角。但什么是属神的崇高呢？岂非信赖基督？不是说：这是我的工作，我施洗；而是\BibleQuote{那一位才是施洗者。} \BibleRef{若 1:33} 这就是达味的角；为使你知道那是达味的角，请注意下文：\BibleQuote{我为我的受傅者预备了一盏灯。} 灯是什么？你们已经知道主论及若翰的话：\BibleQuote{他是点着的明灯。} \BibleRef{若 5:35} 若翰说什么？\BibleQuote{那一位才是施洗者。} 因此圣民要在此欢欣，司祭要在此欢欣：因为他们自己的一切善，都不是出于自己，而是出于那拥有施洗权柄者。因此，凡领受洗礼的人，都无所畏惧地来到祂的圣殿；因为那不是人的，而是那使达味的角昌盛者的。\par

21. \BibleQuote{我的圣化要在祂身上昌盛。} \BibleRef{咏 132:18} 在谁身上？在我的受傅者身上。因为当祂说\BibleQuote{我的受傅者}时，乃是父的声音，祂说：\BibleQuote{我要大大祝福她的寡妇，以食粮使她的穷人饱足。我要使她的司祭披上救恩，她的圣民要欢欣歌唱。} 说\BibleQuote{在那里我要使达味的角茁壮}的那一位是天主。祂亲自说：\BibleQuote{我为我的受傅者预备了一盏灯}，因为基督既是我们的，也是父的：祂是我们的基督，当祂拯救并治理我们时，祂也是我们的主：祂是父的子，但既是我们也是父的基督。因为倘若祂不是父的基督，就不会在前面说：\BibleQuote{为了祢仆人达味的缘故，求祢不要转脸不顾祢的受傅者。} \BibleQuote{我的圣化要在祂身上昌盛。} 圣化在基督身上昌盛。任何人不可将此归于自己，认为自己能圣化别人；否则\BibleQuote{我的圣化要在祂身上昌盛}就不真实了。圣化的荣耀要昌盛。因此基督的圣化在基督自己身上，就是天主的圣化能力在基督内。祂说\BibleQuote{要昌盛}，是指祂的荣耀：因为树木昌盛时，便美丽。因此圣化在洗礼中：从那里昌盛并发光。世界为何顺服于这美丽？因为它昌盛在基督内；若将之置于人的权下，它又如何昌盛呢？因为\BibleQuote{凡有血肉的都如草，一切荣美都如草上的花。}\par

\SourceRecord{Translated by J.E. Tweed.；Nicene and Post-Nicene Fathers, First Series；Vol. 8.；Edited by Philip Schaff.；Buffalo, NY: Christian Literature Publishing Co.,；1888.；<http://www.newadvent.org/fathers/1801132.htm>.}

% structure: p0001 source=h1 target=section reason=page-title

\section{圣咏第一三三篇释义}

这是一篇短小的圣咏，却是一篇众所周知且常被引用的。\BibleQuote{看，兄弟们同居共处，多么美好，多么快乐！}（见：\BibleRef{咏 132:1}）。如此甜蜜的声响，连那些不认识圣咏的人，也吟唱这句诗。……\par

对于这些圣咏词句，这甜蜜的声响，这美妙的旋律，既是心灵的又是赞美诗的，甚至产生了隐修院。这声音唤醒了渴望同居共处的兄弟们。这句诗是他们的号角。它响彻大地，使那些分散的人聚集起来。天主的召唤、圣神的召唤、先知的召唤，在\Place{犹大}未被听闻，却在全世界被听见。他们在那声音被唱响之处是聋的，却有人耳朵张开，正如经上所说：\Quote{\BibleQuote{那些没有被告知他的人，将要看见他；那些没有听说过的人，将要明白。}}（见：\BibleRef{依 65:1}）然而，最亲爱的，如果我们思考，这祝福本身是从那割损的墙垣发出的。因为所有的犹太人都灭亡了吗？宗徒们、先知之子们、被掳者的儿子们是从哪里来的呢？他是对那些知道的人说话。那五百位在主复活后看见他的人，是宗徒\Person{保禄}所纪念的，是从哪里来的呢？（见：\BibleRef{格前 15:6}）那一百二十位（见：\BibleRef{宗 1:15}）在主复活并升天后聚集在一处，圣神在五旬节那天降临在他们身上，从天降下，正如所应许的，是从哪里来的呢？（见：\BibleRef{宗 2:1}）所有的人都是从那里来的，他们首先同居共处；他们变卖了一切所有，把价银放在宗徒脚前，正如在宗徒大事录中所读的（见：\BibleRef{宗 4:34}）。按照每个人的需要分配（见：\BibleRef{宗 2:45}），没有人称什么是他自己的，而是万有共有。那么，\Quote{同居共处}是什么意思呢？经上说，他们一心一意，专注于天主。（见：\BibleRef{宗 4:32}）所以他们是首先听到\Quote{看，兄弟们同居共处，多么美好，多么快乐！}的人。他们是首先听到的，但并非唯一听到的。……\par

圣咏的话中取用了\Quote{隐修士}的名字，这样就没有人能因这名字而责备你们这些天主教徒了。当你们正当地以\OriginalTerm{瑟西利约派}{Circelliones}来责备异端者，好使他们因羞耻而得救时，他们却以隐修士来责备你们。……\par

此外，亲爱的，有些人也是假隐修士，我们知道这样的人；但虔敬的兄弟之情并不因那些自称是而实非的人而废除。有假隐修士，正如在圣职人员中有假人，在信徒中也有假人。……\par

圣咏既说：\Quote{\BibleQuote{看，兄弟们同居共处，多么美好，多么快乐！}}那么我们为何不称隐修士为如此呢？因为\OriginalTerm{独一}{Monos}是\Quote{一}。不是任何一种一，因为人群中的一个人也是一，但他虽然能与他人一起被称作一，却不能称为\Emph{Monos}，即\Quote{独一}，因为\Emph{Monos}的意思是\Quote{单独的一}。那么，那些如此生活以至于成为一体的人，以至于他们真正拥有经上所写的\Quote{一心一意}（见：\BibleRef{宗 4:32}），许多身体，但不是许多心思；许多身体，但不是许多心灵；便可以正当地被称为\Emph{Monos}，即\Quote{独一}。……\par

让圣咏告诉我们他们像什么。\Quote{\BibleQuote{如同头上的香液，流到胡须上，流到\Person{亚郎}的胡须上，流到他衣领的边上。}}（见：\BibleRef{咏 132:2}）\Person{亚郎}是谁？一位司祭。谁是司祭，除非那位进入至圣所的独一司祭？那位司祭是谁，除了同时是祭品和司祭的那位？除了那位在世上找不到洁净之物可献，而奉献了自己的那位？香液在他的头上，因为基督与教会是一整体，但香液是从头而来的。我们的头是被钉死并被埋葬的基督；他复活了，升了天；圣神从头上降下。降到哪里？到胡须上。胡须象征勇敢者；胡须区分出成年男子，即认真、活跃、强壮的人。因此当我们描述这样的人时，我们说，他是个有胡须的人。所以那香液首先降临到宗徒们身上，降临到那些承担世界最初冲击的人身上，因此圣神降临在他们身上。因为那些首先开始同居共处的人，遭受了迫害，但因为香液流到胡须上，他们虽然受苦，却没有被打败。……\par

\Quote{\BibleQuote{如同\Place{赫尔孟}的甘露，降落在\Place{熙雍}的山岭上。}}（见：\BibleRef{咏 132:3}）他愿意让我们明白，我的弟兄们，兄弟们同居共处是因天主的恩宠。……\par

8. 但你们应当知道\OriginalTerm{赫尔孟}{Hermon}是什么。它是一座远离\Place{耶路撒冷}，即远离\Place{熙雍}的山。因此，他说：像赫尔孟的甘露，降落在熙雍山上，就显得奇怪，因为\Place{赫尔孟山}远离耶路撒冷，据说它在\Place{约旦}对岸。那么，让我们为赫尔孟寻求一种解释。这个词是希伯来语，我们从通晓这种语言的人那里得知其意义。据说赫尔孟一词意为\Quote{置于高处的光}。因为甘露来自\Person{基督}。没有光被置于高处，除非是基督。祂是怎样被置于高处的？首先在\Place{十字架}上，之后在\Place{天上}。当祂被贬抑时，被置于十字架的高处；被贬抑，但祂的贬抑不能不崇高。人的职务逐渐减少，正如在\Person{若翰}身上所预表的；\Person{主耶稣基督}的天主职务则增长，正如在他们诞生时所显示的。前者出生，按教会传统，是在六月二十四日，此时白昼开始变短。主诞生于十二月二十五日，此时白昼开始变长。在这里，若翰亲自承认：\BibleQuote{祂必须增长，我却必须减少。}\BibleRef{若 3:30} 两人的苦难也表明了这一点。主被高举在十字架上；若翰被斩首而减少。因此，置于高处的光就是基督，是赫尔孟甘露的来源……但如果他有了赫尔孟的甘露，即降落在熙雍山上的，他就是安静、和平、谦卑、顺从的，倾吐祈祷而不是抱怨。因为抱怨者在圣经某处被恰当地描述为：\BibleQuote{愚人的心如同车轮。}\BibleRef{德 33:5} \BibleQuote{愚人的心如同车轮}是什么意思？它载着干草，并且吱嘎作响。车轮无法停止吱嘎。同样，有许多弟兄，他们只是在身体上同居。但谁才是真正的同居者？就是那些经上所说的：\BibleQuote{他们一心一意朝向天主。}\BibleRef{宗 4:32}\par

9. \Quote{因为在那里主命令了祝福。} 祂在哪里命令的？在那些同居的弟兄中间。在那里祂命令了祝福，在那里他们同心祝福天主。因为你不会在分裂的心中祝福天主……你在世上受困吗？离开，去天上居住。你会说，我一个身穿血肉、被血肉奴役的人，怎么能在天上居住？首先在内心前往，然后身体才能跟随。不要用聋耳聆听\Quote{高举你们的心。} 保持你心高举，就没有人在天上使你受困。\par

\SourceRecord{Translated by J.E. Tweed.；Nicene and Post-Nicene Fathers, First Series；Vol. 8.；Edited by Philip Schaff.；Buffalo, NY: Christian Literature Publishing Co.,；1888.；<http://www.newadvent.org/fathers/1801133.htm>.}

% structure: p0001 source=h1 target=section reason=page-title

\section{圣咏第一三四篇释义}

\Quote{看，如今你们当赞美主，主的一切仆人}\BibleRef{咏 133:1}，\Quote{你们站在上主的殿宇里，在我们天主的殿宇的庭院中}\BibleRef{咏 133:2}。他为何加上\Quote{在庭院中}？庭院指房屋中较宽阔的空间。站在庭院中的人不受狭隘，不被限制，在某种程度上是宽阔的。留在这宽阔中，你就能爱你的仇敌，因为你所爱的不是仇敌能限制你的东西。你如何被理解为站在庭院中？站在爱德中，你就站在庭院中。宽阔在于爱德，狭窄在于仇恨。\par

2. \Quote{在圣所里举起你们的手，在夜间赞美主}\BibleRef{咏 133:2}。在白天赞美是容易的。什么是\Quote{在白天}？在顺境中。因为夜间是悲伤之事，白天是快乐之事。当你顺利时，你赞美主。你的儿子生病，他痊愈了，你赞美主。你的儿子生病，或许你寻求了占星者、占卜者，或许对主的咒骂已来，不是出于你的口舌，而是出于你的行为，出于你的行为和生活。不要自夸，因为你用舌头赞美，却用生活诅咒。因此你们当赞美主。何时？在夜间。\Person{约伯}何时赞美？当那是一个悲伤的夜间时。他所拥有的一切都被夺走；他为存留财物而生的子女也被夺走。他的夜间何等悲惨！然而让我们看看他是否在夜间不赞美。\BibleQuote{主赐予，主收回；主所愿，愿主的名受赞美}\BibleRef{约 1:21}。黑夜是何等黑暗……。\par

3. \Quote{那造天造地的主由熙雍祝福你}\BibleRef{咏 133:3}。他劝勉众人赞美，而他自己却祝福那一个人，因为他从众人中造就了一个人，正如\Quote{弟兄们共聚一处是美好而愉快的}。\Quote{弟兄}是复数，然而\Quote{共聚一处}却是单数。你们中任何人不要说：这与我无关。你们知道他所说的是谁吗？\Quote{主由熙雍祝福你}。他祝福了一个人。成为一个人，祝福就临于你。\par

\SourceRecord{Translated by J.E. Tweed.；Nicene and Post-Nicene Fathers, First Series；Vol. 8.；Edited by Philip Schaff.；Buffalo, NY: Christian Literature Publishing Co.,；1888.；<http://www.newadvent.org/fathers/1801134.htm>.}

% structure: p0001 source=h1 target=section reason=page-title

\section{圣咏第一三五篇释义}

1. 这篇圣咏所劝勉我们的，对我们来说应当是非常甘饴的，我们应当因为它的甘饴而喜乐。因为它说：\BibleQuote{赞美主的圣名}\BibleRef{咏 134:1}。并且它立即附加了理由，为什么赞美主的圣名是正义的。\Quote{赞美主，你们这些仆人。}还有什么比这更正义？有什么比这更相称？有什么比这更知恩？……因为如果祂教导那些事奉祂且蒙祂恩待的、祂圣言的宣讲者、祂教会的管理者、祂圣名的朝拜者、祂诫命的服从者，叫他们在自己的良心中拥有生命的甘甜，以免因人的赞美而败坏或受人的责备而气馁；那么，那位至高无上、不变的天主，祂教导这些事，既不因你赞美而更大，也不因你责备而更小，岂不更当如此？因为你赞美你的主，如同仆人，这绝非不当。纵然你们永远只是仆人，也应当赞美主；你们这些仆人岂不更应当赞美主，好使你们将来能获得儿子的恩宠？\par

2. ……因此，\BibleQuote{你们站在主的殿中，在我们天主的殿的庭院中，要赞美主}\BibleRef{咏 134:2}。要知恩；你们从前在外面，现在却站在里面。既然你们站着，那么你们认为应当在何处赞美那位扶起跌倒的你们、使你们站在祂殿中、认识祂、赞美祂的主，这难道是小事吗？我们站在主的殿中，这难道是小恩惠吗？……若有人想到这事，而不忘恩，他就会因主的爱而完全轻视自己，因为主为他行了如此大事。既然他没有什么可以回报天主如此浩大的恩惠，除了感谢祂而非回报祂，还能做什么呢？感恩的行动本身，就在于\Quote{接受主的杯，呼求祂的名}。因为仆人能用什么回报主所赐给他的一切呢？\par

3. 我为何要给你们理由赞美祂？\BibleQuote{因为主是良善的}\BibleRef{咏 134:3}。简言之，一句话就说明了我们天主、主的赞美。\Quote{主是良善的}；祂的良善不同于祂所造之物的良善。因为天主造了一切，样样都好；\BibleRef{创 1:31}不仅是好，而且是极好。祂造了天和地，以及其中所有的一切，都是好的，并且造得极好。若祂所造的这一切都是好的，那么造它们的那位该是怎样的一位呢？……\par

4. 我们对祂的良善能讲到多远呢？谁能心中领会或理解主是何等良善呢？然而，让我们回到我们自身，在我们内认出祂，并在祂的受造物中赞美造物主，因为我们还不配直接默观祂。怀着希望，当我们因信德净化心灵后，将来能在真理中喜乐；如今既然我们看不见祂，就让我们观看祂的化工，好叫我们活着不能没有赞美祂。因此我说：\Quote{赞美主，因为祂是良善的；歌颂祂的圣名，因为祂是甘饴的}。……祂是中保，因此是甘饴的。有什么比天使之粮更甘饴呢？既然人曾吃过天使之粮，天主怎能不成为甘饴的呢？因为人与天使并非靠不同的食物生活。那就是真理，那就是智慧，那就是天主的良善，但你不能像天使那样享受它。……为使人能吃天使之粮，天使的造物主成了人。你们若尝过，就当歌颂；你们若尝过主是何等甘饴，就当歌颂；若你们所尝的滋味美好，就当赞美它；谁对厨师或供应者如此忘恩，以致在品尝时若不因食物合口而赞美感谢呢？若我们在那种场合不沉默，岂能在赐给我们一切的主面前沉默呢？……\par

5. \BibleQuote{因为主拣选了雅各伯归自己，以色列归为自己的产业}\BibleRef{咏 134:4}。……因此，让雅各伯不要自夸，不要自傲，或归功于自己的功德。他早已被认识，早已被预定，早已被拣选，不是因自己的功德而被拣选，而是被寻见，且因天主的恩宠被赐予生命。外邦人也是如此；因为野橄榄有何功德，能接在真橄榄上，脱离其果实的苦味和野生的贫瘠呢？它本是荒野的树木，而非主的田园里的，然而主因自己的仁慈，将野橄榄接在橄榄上。但直到如今，野橄榄尚未被接上。\par

6. ……\BibleQuote{因为，}他说，\BibleQuote{我知道主是伟大的，我们的天主超越一切神祇}\BibleRef{咏 134:5}。如果我们对他说：我们求问你，请向我们解释祂的伟大；难道他不会回答我们说：我所看见的祂，如果我能解释，祂就不是那么伟大？那么，让他回到祂的化工，告诉我们吧。让他心中持守他所看见的天主的伟大，并将此交托给我们的信德，因为他不能引导我们的眼目到那里，却能列举一些主在此所行的事；好使那些不能像他那样看见天主伟大的我们，通过我们能理解的化工，而品尝主的甘甜。……\par

7. \BibleQuote{主所愿意的一切，祂都在天上、地上、海中及一切深渊中造成了}\BibleRef{咏 134:6}。谁能领会这些事？谁能数算主在天上、地上、海中及一切深渊中的化工呢？纵然我们不能完全领会，也应当毫无疑惑地相信和持守，因为天上、地上、海中及一切深渊中的任何受造物，都是主所造的。……\par

8. \BibleQuote{从地极升起云彩}\BibleRef{咏 134:7} 我们在祂的受造界中看到天主的这些工程。因为云彩从地极升起，来到地中央，降雨；你无法察验它们从何而来。因此先知以此表明，从\BibleQuote{地极}，无论是从底层还是从地极的周边，凡祂愿意，祂就从地上升起云彩。\BibleQuote{祂使闪电转为雨}。因为没有雨的闪电只会恐吓你，却不给你任何东西。\BibleQuote{祂使闪电转为雨}。闪电时你战栗；下雨时你欢欣。\BibleQuote{祂使闪电转为雨}。那恐吓你的，却使你喜悦。\BibleQuote{祂从祂的宝库中吹出风}，它们的原因隐藏着，你不知从何而来。当风吹起，你感觉到它；它为何吹起，或从祂智慧的哪座宝库中吹出，你却不知道；\BibleRef{若 3:8} 但你还欠天主信德的敬拜，因为若不是那创造风的祂命令，若不是那创造风的祂叫它出来，它不会吹起。\par

9. \BibleQuote{击杀埃及的长子}\BibleRef{咏 134:8} 但那些神圣作为被述说，是你可能爱慕的；那些不被述说的，是你可能畏惧的。留心，并看到当祂发怒时，祂仍做祂愿意的事。\BibleQuote{从人到兽，祂在你们中间，\Place{埃及}啊，施行了神迹和奇事}\BibleRef{咏 134:9} 你知道，你读过了主的手藉\Person{梅瑟}在\Place{埃及}所行的，为压碎并打倒骄傲的\Place{埃及人}：\BibleQuote{对\Person{法郎}和他的众仆役}。祂在\Place{埃及}还做得少吗？祂带出祂的子民之后，又做了什么？\BibleQuote{击杀了许多民族}\BibleRef{咏 134:10}，这些民族曾占有那块土地，就是天主愿意赐给祂子民的土地。\BibleQuote{杀死了强壮的君王，\Person{息红}阿摩黎人的王，和\Person{敖格}巴商的王，以及\Place{客纳罕}所有的王国}\BibleRef{咏 134:11} 圣咏所简单记述的这一切，我们在主其他的经卷中也同样读到，在那里主的手是伟大的。当你看到恶人所受的，要当心，免得你也被同样对待....但是当善人看到恶人所受的苦，让他洁净自己脱离一切不义，免得陷入类似的刑罚、类似的惩戒。那么你就彻底明白了这些事。天主那时做了什么？祂驱逐了恶人，\BibleQuote{将他们的土地赐为产业，就是赐给祂的仆人\Place{以色列}为产业}\BibleRef{咏 134:12}。\par

10. 接着是大声赞美祂的欢呼。\BibleQuote{主，祢的名号永远常存}\BibleRef{咏 134:13}，在祢所做的一切之后。因为我看到祢做了什么？我注视祢在天上所创造的受造物，我注视这较低的部分，我们居住的地方，在这里我看到祢恩赐的云彩、风和雨。我注视祢的子民；祢带领他们离开了为奴之家，并对他们的仇敌行了神迹和奇事。祢惩罚了那些使他们受苦的人，祢驱逐了恶人离开他们的土地，祢杀死了他们的王，祢将他们的土地赐给了祢的子民：我看到了这一切，充满喜乐地说：\BibleQuote{主，祢的名号永远常存}...\par

11. 这一切，天主那时在身体上推翻了，当时我们的祖先被领出\Place{埃及}地；现在在精神上也是如此。祂的手直到终结也不停息。所以不要认为这些天主的伟业那时就结束并停止了。\BibleQuote{主，祢的名号}他说，\BibleQuote{永远长存}。也就是说，祢的慈爱不停息，祢的手永不停止做这些事，就是祢先前在预像中已宣布的事。\BibleQuote{但这些事记录下来，为劝诫我们这些末世的人}\BibleRef{格前 10:11}。一代和另一代；那使我们成为信众、藉圣洗而重生的世代；那使我们从死者中复活、与天使永远生活的世代。主，祢的纪念高于这一代，也高于那一代；因为祂现在不会忘记召唤我们，将来也不会忘记给我们加冕。\par

12. \BibleQuote{主审判了祂的子民，并在祂的仆人中被呼求} \BibleRef{咏 134:14}。祂已经审判了那人民。除了最后的审判，犹太人民已被审判。什么是\Quote{审判}？义人被带走，不义的人被留下。但如果我撒谎，或被认为撒谎，因为我说了它已被审判，请听主说：\BibleQuote{我为审判来到这世界上，叫那看不见的能看见，那能看见的反而成为瞎子。} \BibleRef{若 9:39}骄傲的人成为瞎子，谦卑的人被光照。因此，\Quote{祂审判了祂的子民。}依撒意亚预言了这审判。\BibleQuote{现在，雅各伯家，来罢！让我们在上主的光明中行走。} \BibleRef{依 2:5}这是小事；但接下来说什么？\Quote{因为祂抛弃了祂的子民，以色列家。}雅各伯家就是以色列家；因为那雅各伯也就是以色列……因此，天主审判了祂的子民，通过分开恶人和善人；也就是说，\Quote{祂将在祂的仆人中被呼求。}被谁？被外邦人。因为有多少民族因信德而来！现在有多少田地和旷野来到我们这里？他们从那里来，无人能说有多少；他们愿意相信。我们对他们说：你们要什么？他们回答：要知道天主的荣耀。相信吧，我的弟兄们，我们对这些乡野之人的如此要求感到惊奇和喜悦。他们不知从何处来，被不知谁所唤醒。我怎么说，我不知道被谁？我确实知道被谁，因为祂说：\BibleQuote{除非父吸引他，没有人能到我这里来。} \BibleRef{若 6:44}他们突然从树林、旷野、最遥远的高山来到教会；他们中许多人，甚至几乎所有人，都这样说，以致我们真看见天主在他们内教导他们。圣经的预言应验了，当它说：\BibleQuote{他们都要蒙受天主的训诲。} \BibleRef{依 54:13}我们对他们说：你们渴望什么？他们回答：要看天主的荣耀。\BibleRef{若 6:45}\BibleQuote{因为众人都犯了罪，亏缺了天主的荣耀。} \BibleRef{罗 3:23}他们相信，他们被圣化，他们愿意有神职人员为他们祝圣。这岂不是应验了：\Quote{祂将在祂的仆人中被呼求}？\par

13. 最后，经过这一切安排和措施，天主之神转而责备和嘲笑那些偶像，现在连它们的崇拜者也嘲笑它们。\BibleQuote{外邦人的偶像不过是金银} \BibleRef{咏 134:15}。既然天主造了这一切，祂按自己的旨意创造了天上和地上的一切，那么人所造的，除了成为嘲笑的对象，还能是崇拜的对象吗？或许祂正要谈论\Quote{外邦人的偶像}，使我们鄙视它们全部？祂要谈论异教徒的偶像，石头和木头，灰泥和未烧过的砖？我不说这些，它们是卑贱的材料。我说的是他们特别珍爱、特别尊崇的东西。\BibleQuote{外邦人的偶像是金银，是人手的制作。}诚然是金子，诚然是银子：因为银子闪耀，金子闪耀，难道它们因此就有眼睛，或能看见吗？……但既然这些东西没有感觉，为什么你们把金银制成的人当作神？你们难道没看见你们所造的神看不见？\BibleQuote{他们有口，却不能言；有眼，却不能看} \BibleRef{咏 134:16}；\BibleQuote{有耳，却不能听；口中也没有气息} \BibleRef{咏 134:17}；\BibleQuote{有鼻，却不能闻；有手，却不能操作；有脚，却不能行走。}（缺引用）这一切，木匠、银匠、金匠都能制作，包括眼、耳、鼻、口、手、脚，但他不能使眼睛看见，耳朵听见，口说话，鼻子闻嗅，手活动，脚行走。\par

14. 人啊，如果你知道你被谁所造，你无疑会嘲笑你所造的东西。但对于那些不知道的人，经上怎么说？\BibleQuote{制造它们的人，以及所有信赖它们的人，都要变得像它们一样} \BibleRef{咏 134:18}。你们相信，弟兄们，这种与偶像的相似不是表现在他们的肉体上，而是在他们的内心。因为\Quote{他们有耳，却听不见。}\Emph{天主}确实呼召他们，\Quote{有耳可听的，就听罢！} \BibleRef{玛 11:15}他们有眼，却看不见，因为他们有肉体的眼，却没有信德的眼。最后，这预言在万民中应验了……难道没有应验吗？难道没有像所写的那样被看见吗？而那些留下的人有眼，却看不见；有鼻，却闻不到。他们嗅不到那香气。\Quote{我们是基督的馨香，} \BibleRef{格后 2:15}正如宗徒到处所说的。他们有鼻，却闻不到基督那如此甘甜的馨香，这有什么益处呢？这实在是发生在了他们身上，并且经上论到他们说：\Quote{制造它们的人，以及所有信赖它们的人，都要变得像它们一样}等等。\par

15. 如今人们因我主基督的奇迹而相信；如今盲人的眼睛、聋子的耳朵被开启，无知者的鼻孔被吹入气息，哑子的舌头被松开，瘫痪者的手被加强，瘸子的脚被引导；从这些石头中给亚巴郎兴起子孙来，\BibleRef{玛 3:9}对他们所有人都要说：\BibleQuote{你们以色列家，当赞美上主} \BibleRef{咏 134:19}。所有人都是亚巴郎的子孙；如果从这些石头中给亚巴郎兴起子孙来，显然他们是以色列家，属于以色列家，不是按血肉，而是按信德作亚巴郎的后裔。但即使这话是对那家说的，指的是以色列人民，难道从那里不是有宗徒和成千上万的受了割损的人相信吗？\BibleQuote{你们亚郎家，当赞美上主。你们肋未家，当赞美上主} \BibleRef{咏 134:20}。万民啊，赞美上主，这就是一般的\Quote{以色列家}；领袖们啊，赞美祂，这就是\Quote{亚郎家}；仆人们啊，赞美祂，这就是\Quote{肋未家}。至于其他的民族？\Quote{你们敬畏上主的，当赞美上主。}（缺引用）\par

16. 让我们也异口同声地说出接下来的一句：\Quote{愿主由熙雍受赞颂，祂居住在耶路撒冷。}\BibleRef{咏 134:21} 熙雍也就是耶路撒冷。熙雍意为\Quote{守望}，耶路撒冷意为\Quote{和平的异象}。如今主要居住在怎样的耶路撒冷？不是那已倾覆的，而是我们的母亲，那在天上的，正如经上所载：\Quote{那荒凉的，比有丈夫的儿女更多。} 如今主来自熙雍，因为我们守望祂的来临；只要我们生活在望德中，我们就在熙雍。当我们的道路终结时，我们将居住在那永不倾塌的城中，因为主住在其中并守护她，那是和平的异象，永恒的耶路撒冷。我的弟兄们，言辞不足以赞美她；在那里我们找不到仇敌，无论是在教会内还是教会外，无论是在我们的肉身中还是我们的思想里。因为\Quote{死亡被胜利吞灭了，}\BibleRef{格前 15:54} 我们将获得自由，在永恒的和平中得见天主，成为耶路撒冷——天主之城——的公民。\par

\SourceRecord{Translated by J.E. Tweed.；Nicene and Post-Nicene Fathers, First Series；Vol. 8.；Edited by Philip Schaff.；Buffalo, NY: Christian Literature Publishing Co.,；1888.；<http://www.newadvent.org/fathers/1801135.htm>.}

% structure: p0001 source=h1 target=section reason=page-title

\section{\WorkTitle{圣咏第一三六篇释义}}

1. \BibleQuote{你们应当称谢上主，因为祂是良善的，因为祂的仁慈永远常存}\BibleRef{咏 135:1}。这篇圣咏赞颂天主，每一节都以同样的方式结束。因此，虽然这里叙述了许多颂扬天主的事，但祂的仁慈最受称扬；因为若没有这明确的称扬，那位圣神藉以说出此圣咏的人就不会让任何一节结束。虽然在世界的终结时，审判区分义人和不义者，义人进入永生，不义者进入永火\BibleRef{玛 25:46}，之后便没有那些天主将怜悯的对象了，然而，祂对祂的圣徒和信友所施予的未来仁慈，仍可正确地被理解为永远常存，这是因为他们不会永远悲惨而因此需要祂永远的仁慈，而是因为那极度的福乐——祂仁慈地赐给悲惨者，使他们不再悲惨，开始幸福——将没有终结，因此\Quote{祂的仁慈永远常存}。因为我们从为义人成为义人，从病弱者成为健全，从死亡者成为活人，从会死者成为不死者，从不幸者成为幸福者，都是出于祂的仁慈。但我们将成为的这状态将永远常存，因此\Quote{祂的仁慈永远常存}。所以，\Quote{你们应当称谢上主}，也就是说，以感谢赞颂上主，\Quote{因为祂是良善的}：你们从这宣认中不会得到任何暂时的好处，因为\Quote{祂的仁慈永远常存}，即祂仁慈赐予你们的恩惠是永恒的。\par

2. 接着是：\Quote{你们应当称谢\OriginalTerm{万神之神}{God of gods}，因为祂的仁慈永远常存}\BibleRef{咏 135:2}。\Quote{你们应当称谢\OriginalTerm{万主之主}{Lord of lords}，因为祂的仁慈永远常存}\BibleRef{咏 135:3}。我们很可以探究：那些以真天主为其天主及主宰的\Quote{诸神}和\Quote{诸主}是谁？我们在另一篇圣咏中读到，甚至人也被称为神。主也在福音中引用这一见证，说：\Quote{你们的法律上岂不是写着：\InnerQuote{我说过：你们是神}吗？}\BibleRef{若 10:34}……因此，他们被称为神，并非因为他们都是善的，而是因为\Quote{天主的话临到了他们}。因为若是因他们都是善的，祂就不会这样区分他们。祂说：\Quote{祂在诸神之中施行审判}。接着是：\Quote{你们要审判不义要到几时呢？}等等，祂所说的显然不是对所有人，而是对某些人，因为祂在区分时说这话，然而祂是在诸神之间进行区分。\par

3. 但若有人问：如果那些领受天主之言的人被称为神，那么天使是否也应被称为神？因为应许给义人和圣人的最大赏报是与天使同等？在圣经中，我不知是否常见——至少不易找到——天使被明确称为神；但当论及上主天主时，说：\Quote{祂是可敬畏的，超越诸神}，接着他仿佛解释为什么说这话时，补充道：\Quote{因为外邦人的神是魔鬼}，以便我们理解希伯来文所表达的\Quote{外邦人的神是偶像}，而其实指的是居于偶像中的魔鬼。至于那些在希腊文称为偶像（idola）、我们拉丁文沿用此名的雕像，它们有眼却不能看，以及一切关于它们的描述，因为它们是完全没有知觉的；因此它们不能被惊吓，因为没有任何无知觉的东西能被惊吓。那么，怎么能说上主\Quote{超越诸神是可敬畏的，因为外邦人的神是偶像}，除非在这些雕像中理解到可被惊吓的魔鬼？因此，宗徒也说：\Quote{我们知道偶像算不了什么}\BibleRef{格前 8:4}，这是指其土质无感觉的材料。但为了不使任何人以为没有一种有生命、有感觉的本性喜悦外邦人的祭献，他补充说：\Quote{但是外邦人所祭献的，是祭献给魔鬼，而不是给天主；我不愿意你们与魔鬼有份}\BibleRef{格前 10:20}。因此，如果我们从未在天主的言语中发现圣天使被称为神，我认为最好的理由是：免得人们因这个名称而把那宗教的职事和服务——希腊文称为\OriginalTerm{礼仪}{\GreekText{λειτουργία}}或\OriginalTerm{崇拜}{\GreekText{λατρία}}——献给圣天使，而圣天使自己也决不会让人这样事奉，除非献给那同为祂们和人的天主的那一位。因此，祂们更恰当地被称为\Quote{天使}，拉丁文是\OriginalTerm{使者}{Nuntii}，以便通过祂们职务的名称——而非祂们的本质——让我们清楚地明白，祂们要我们崇拜那祂们所宣告的天主。宗徒已经简要地阐明了这整个问题，他说：\Quote{因为虽然有被称为神的，不论在天上或在地上，就如有许多神和许多主，可是我们只有一个天主圣父，万物都出于祂，我们也归于祂；也只有一个主耶稣基督，万物藉祂而有，我们也藉祂而有}\BibleRef{格前 8:5}。\par

4. 因此，让我们\Quote{感谢众神之神，万主之主，因为祂的仁慈}等等。\Quote{唯有祂独自行了奇事}\BibleRef{咏 135:4}。正如每一节的末尾所写的：\Quote{因为祂的仁慈永远常存}，我们也必须在每节的开头理解，尽管没有写出来，\Quote{感谢}。这在希腊文里非常清楚。如果我们的译者能够使用那种表达，拉丁文也会如此。其实他们本可以在这一节这样做，如果他们说了：\Quote{对那位行奇事者}。因为我们有的地方是\Quote{祂行了奇事}，希腊文是\OriginalTerm{行奇事的那一位}{\GreekText{τῷ} \GreekText{ποιήσαντι}}，我们必须理解为\Quote{感谢}。我希望他们加上了代词，说\Quote{对那位行了}，或\Quote{对那位正行的}，或\Quote{对那位使之坚固的}；因为那样就能容易理解\Quote{让我们感谢}。因为现在译得很模糊，以至于不知道或不愿查考希腊文手稿的人可能会以为：\Quote{祂造了天，祂铺设了地在水上，祂造了大的光体，因为祂的仁慈永远常存}之所以这样说，是因为祂做这些事的原因是\Quote{因为祂的仁慈永远常存}：而实际上，那些祂从苦难中解放出来的人属于祂的仁慈：但我们不应相信祂出于仁慈造了天、地和光体；因为这些都是祂善心的标记，祂造了万有，样样都甚好。\BibleRef{创 1:31} 因为祂创造了万有，使它们得以存在；\BibleRef{智 1:14} 但清洁我们的罪并拯救我们脱离永苦，则是祂仁慈的工程。因此，圣咏这样对我们说：\Quote{感谢众神之神，感谢万主之主。}感谢祂，\Quote{唯有祂独自行了奇事}；感谢祂，\Quote{祂以智慧造了诸天}；感谢祂，\Quote{祂将地铺在水上}；感谢祂，\Quote{唯有祂造了大的光体}。但我们为何要赞美，他在所有节的末尾都写明了：\Quote{因为祂的仁慈永远常存}。\par

5. 但\Quote{唯有祂独自行了奇事}是什么意思？是因为祂藉着天使和人行了许多奇事？有些奇事是天主独自行的，祂列举了这些，说：\Quote{祂以智慧造了诸天}\BibleRef{咏 135:5}，\Quote{祂将地铺在水上}\BibleRef{咏 135:6}，\Quote{唯有祂造了大的光体}\BibleRef{咏 135:7}。因此他在这一节也加了\Quote{独自}，因为祂将要述说的其他奇事，是天主藉着人行的。因为说了\Quote{唯有祂造了大的光体}之后，祂接着解释这些是什么：\Quote{太阳管白天}\BibleRef{咏 135:8}，\Quote{月亮和星辰管黑夜}\BibleRef{咏 135:9}；然后祂开始述说祂藉着天使和人行的奇事：\Quote{祂击杀了埃及人及其长子}\BibleRef{咏 135:10}，等等。显然，整个受造界是天主独自造的，不是藉着任何受造物；关于这受造界，祂提到了一些更突出的部分，让我们想到整体；我们能理解诸天，我们看见大地。既然也有可见的诸天，祂藉着提及其中的光体，使我们看见祂造了整个天穹。\par

6. 然而，祂所说的\Quote{祂以理解造了诸天}，或如其他人所译的\Quote{以理智}，祂是否意指我们所能理解的诸天，或者祂在自己的理解或理智中，即祂的智慧中造了诸天（就如别处所写的：\Quote{祢以智慧造了这一切}），暗示那独生的圣言，这是一个问题。但若是如此，即我们应理解\Quote{天主以祂的智慧造了诸天}，为何祂只提及诸天？而天主以同一位智慧造了万有？这是因为只需要在那里表达出来，以便在其他部分中虽未写明，也能理解。那么，如果\Quote{以理解}或\Quote{以理智}是指\Quote{以祂的智慧}，即藉着独生的圣言，又怎能说是\Quote{独自}呢？是不是因为天主圣三不是三个天主，而是一个天主，所以祂说天主独自行了这些事，因为祂造了受造界，并非藉着任何受造物？\par

7. 但\Quote{祂将地铺在水上}是什么意思？这是一个难题，因为地似乎是更重的，所以不应相信它浮在水上，而应相信它承载水。为了不让我们似乎固执地维护我们的圣经，对抗那些自以为已经根据可靠原理发现了这些事的人，我们可以给出第二种解释：人类所居住、包含地上生物的地，是\Quote{铺在水上}，因为它耸立于环绕它的水之上。因为当我们说一座城建于海上\Quote{在水的上面}时，并不是说海在它下面，如同水在洞穴居室下面或船只航行其上那样；而是说它\Quote{在}海\Quote{上}，因为它耸立于下面的海之上。\par

8. 但如果这些话还意味着别的与我们有更密切关系的事，那么，天主\Quote{以祂的智慧造了诸天}，即祂的圣人、属神的人——祂不仅赏赐他们相信，而且理解神性之事；那些尚不能达到这一点、只坚定持守信德、如同在诸天之下的，以\Quote{地}之名来象征。因此说：\Quote{祂将大地铺展在水上。}再者，因为关于我们的主耶稣基督，经上记载：\Quote{在祂内蕴藏着智慧和知识的一切宝藏}【\BibleRef{哥 2:3}】，并且智慧和知识这两者彼此有所区别，圣经其他章节也证明这一点，尤其是在圣约伯的话中，两者都以某种方式被定义；那么，我们领悟智慧在于认识和爱慕那永恒不变者——即天主，也并非不合宜。因为那处说：\Quote{虔诚是智慧}，在希腊文是\OriginalTerm{神}{\GreekText{θεοσ}}。\par

\SourceRecord{Translated by J.E. Tweed.；Nicene and Post-Nicene Fathers, First Series；Vol. 8.；Edited by Philip Schaff.；Buffalo, NY: Christian Literature Publishing Co.,；1888.；<http://www.newadvent.org/fathers/1801136.htm>.}

% structure: p0001 source=h1 target=section reason=page-title

\section{《圣咏第一百三十六篇释义》}

1. ……但我们今日咏唱说：\BibleQuote{我们坐在巴比伦河畔，一想起熙雍就泪流满面}\BibleRef{咏 136:1}。……\par

2. 请注意\Quote{巴比伦的河水}。\Quote{巴比伦的河水}是指在此世受人喜爱而转瞬即逝的一切事物。例如，有人喜爱务农，藉以致富，倾心于此，并从中取乐；且看他结局，便知他所爱的并非耶路撒冷的根基，而是巴比伦的河流。另一人说：做军人是大事：所有的农夫都惧怕军人……\par

3. 但是，圣城耶路撒冷的其他居民，明白自己的被掳，觉察到人本性的欲望和各式各样的贪欲怎样催促、拉扯他们东奔西跑，驱入大海；他们看见这事，便不投身于巴比伦的河水，反而\BibleQuote{坐下哭泣}，或为那些被河水冲走的人，或为他们自己，因为他们的行为应得身在巴比伦，但他们是坐着的，即谦卑自下。啊，圣洁的熙雍！那里一切稳固，没有流动的东西！是谁使我们头朝下坠入此地？我们为何离开了你的创立者和你的团体？看啊，置身于一切流动消逝之处，几乎无人能抓住树木，从洪流中被救出而逃脱。因此，我们在被掳中自卑，让我们\BibleQuote{坐在巴比伦的河畔}，不要胆敢投身于那些河流，也不要在我们被掳的邪恶与悲伤中骄傲自大，而要坐下，并这样哭泣。让我们坐在巴比伦的河水\Quote{旁}，而不是在河水之下；愿我们的谦卑如此，不至淹没我们。是坐在河水\Quote{旁}，不是\Quote{在水中}，也不是\Quote{在水下}；但仍要坐着，以谦卑的方式，不要像在耶路撒冷那样说话……\par

4. 因为许多人以巴比伦的哭泣而哭泣，因为他们也以巴比伦的喜乐而喜乐。当人们因获益而喜乐，因损失而哭泣时，两者都属于巴比伦。你应当哭泣，但要在纪念熙雍时哭泣。如果你在纪念熙雍时哭泣，即使在巴比伦境遇顺遂时，你也应当哭泣……\par

5. \BibleQuote{我们把竖琴挂在岸边的柳树上}\BibleRef{咏 136:2}。耶路撒冷的子民有他们的\Quote{竖琴}，即天主的圣经、天主的诫命、天主的恩许、对来生的默想；但当他们住在\Quote{巴比伦}时，他们\Quote{把竖琴挂起来}。柳树是不结果实的树，在此处如此安排，以致不能从它们领悟任何善事；他处或许可以。此处指不结果实的树，生长在巴比伦的河畔。这些树受巴比伦河水的浇灌，却不结果实；正如有些人贪婪、吝啬、在善工上不结果实，他们是巴比伦的公民，甚至成为该地的树；他们在此世短暂之乐的享受中得滋养，好像被\Quote{巴比伦的河水}浇灌。你在他们身上寻找果子，却无处可寻。……因此，由于暂不把圣经应用于他们，\Quote{我们把竖琴挂在柳树上}。因为我们不认为他们配得携带我们的竖琴。所以我们不把竖琴插入他们中间并绑在他们身上，而是推迟使用，这样便挂起来。因为柳树是巴比伦不结果实的树，以暂时的快乐为滋养，如同以\Quote{巴比伦的河水}为滋养。\par

6. \BibleQuote{俘虏我们的人要我们唱歌，迫害我们的人要我们欢乐}\BibleRef{咏 136:3}。俘虏我们的人向我们索取歌词和一首歌。……我们受世上快乐的诱惑，每日与非法快感的暗示搏斗；甚至在祈祷中也难得自由呼吸；我们明白自己是俘虏。但谁俘虏了我们？什么人？什么民族？什么君王？如果我们被救赎，我们曾为俘虏。谁救赎了我们？基督。祂从谁手中救赎了我们？从魔鬼。因此魔鬼和他的使者俘虏了我们；他们不会俘虏我们，除非我们同意。……\par

7. 那么，\Quote{那些}\Quote{俘虏我们的人}，即魔鬼和他的使者，何时曾对我们说：\Quote{给我们唱一支熙雍的歌吧}？我们如何回答？巴比伦怀胎你们，巴比伦容纳你们，巴比伦滋养你们，巴比伦藉你们的口说话；你们只知道接受眼前闪光的东西，不知道默想永恒之事，你们不理解你们所求的事。\BibleQuote{我们怎能在外乡唱主的歌呢？}\BibleRef{咏 136:4}。的确，弟兄们，事实如此。当你开始愿意按你所知道的限度宣讲真理时，你就会看到你必须忍受多少这样的嘲弄者、真理的迫害者、充满谎言的人。当他们向你索取他们不能理解的东西时，你要回答他们，并满怀信德唱你的圣歌：\BibleQuote{我们怎能在外乡唱主的歌呢！}\par

8. 但是，天主的子民、基督的身体、高贵的流亡团体（因为你们的家乡不在这里，而在别处）要小心你们如何住在他们中间，免得因爱他们、追求他们的友谊、害怕得罪这样的人，巴比伦开始令你们欢喜，而你们忘记了耶路撒冷。因此，出于对这事恐惧，请看圣咏作者接着说什么，请看下文。\BibleQuote{耶路撒冷！我若忘记你}\BibleRef{咏 136:5}，在那些俘虏我的人的话语中，在诡诈之人的话语中，在那些存心不良、询问却又不愿意学习的人的话语中……那么，\BibleQuote{耶路撒冷！我若忘记你，情愿右手忘记技巧。}\par

9. \BibleQuote{愿我的舌头贴住上颚，如果我不怀念你} \BibleRef{咏 136:6}。这就是说，愿我成为哑巴，如果我不怀念你。因为那不说熙雍之歌的人，他说了什么话、发出了什么声音呢？那是我们的舌头，耶路撒冷之歌。对这世界的爱之歌是一种陌生的语言，一种野蛮的语言，我们在被掳中学会了它。那么，忘记耶路撒冷的人在天主面前将是哑巴。并且光怀念还不够：因为她的敌人也怀念她，想要推翻她。\Quote{那是什么城？}他们说；\Quote{基督徒是什么样的人？基督徒是怎样的？但愿他们不是基督徒。}如今被掳的团体征服了掳掠者；他们仍然在喃喃自语、发怒，想要杀死那在他们中间寄居的圣城。所以，光怀念还不够：注意你是如何怀念的。因为有些事情我们出于憎恨而怀念，有些事情出于爱。因此，当他说：\Quote{如果我忘记了你，耶路撒冷}等等时，他立刻添加：\Quote{如果我不在喜乐的顶点中高举耶路撒冷。}因为喜乐的顶点是我们在那里享受天主，我们在那里安全地享有共融的兄弟情谊和公民的合一。那里没有试探者攻击我们，没有人能引诱我们到任何诱惑：那里只有善使我们喜悦：那里一切匮乏都将消失，那里完全的幸福将照耀我们。\par

10. 然后他转向天主，向那座城市的敌人祈祷。\BibleQuote{上主，求你记念厄东的子孙} \BibleRef{咏 136:7}。厄东就是也称作厄撒乌的那位：因为你们刚刚听到宗徒的话被宣读：\BibleQuote{我爱了雅各伯，但我恨了厄撒乌。} \BibleRef{罗 9:13}。……厄撒乌因此代表所有属血肉的人，雅各伯代表所有属神的人。……所有属血肉的人都是属神之人的敌人，因为所有这些贪图现世事物的人，迫害那些他们视为渴求永恒事物的人。针对这些人，圣咏作者望着耶路撒冷，恳求天主让他脱离被掳，说——什么？\Quote{上主，求你记念厄东的子孙。}求你从属血肉的人中解救我们，从那些效法厄撒乌的人中，他们是兄长，却是敌人。他们是长子，但次子赢得了优先地位，因为肉欲的贪念打倒了前者，对贪念的蔑视抬举了后者。后者活着，嫉妒并且迫害。\Quote{在耶路撒冷的日子。}耶路撒冷的日子，就是它受考验、被掳的日子，还是耶路撒冷幸福的日子，即它被释放、达到目标、分享永恒的日子？\Quote{求你记念，}他说，\Quote{上主，}不要忘记那些\Quote{说：\InnerQuote{拆毁，拆毁，直到根基。}的人。}所以这意味着记念他们企图推翻耶路撒冷的那一天。教会受了多么大的迫害啊！厄东的子孙，即属血肉的人、魔鬼及其天使的奴仆，他们崇拜木石，追随肉欲的贪念，他们怎么说？\Quote{铲除基督徒，消灭基督徒，一个也不留，把他们连根拔起！}这些事难道没有说过吗？当这些话被说出时，迫害者被拒绝，殉道者被加冕……。\par

11. 然后他转向她：\Quote{不幸的巴比伦女儿，}在你的狂喜、你的傲慢、你的敌意中不幸；\BibleQuote{不幸的巴比伦女儿！} \BibleRef{咏 136:8}。这城既被称为巴比伦，又称为巴比伦的女儿：正如他们谈论\Quote{耶路撒冷}和\Quote{耶路撒冷女儿}，\Quote{熙雍}和\Quote{熙雍女儿}，\Quote{教会}和\Quote{教会女儿}。因为它接续前者，所以称为\Quote{女儿}；因为它优于前者，所以称为\Quote{母亲}。曾有一个先前的巴比伦；人民还留在其中吗？因为它接续巴比伦，所以称为巴比伦的女儿。哦，巴比伦的女儿，\Quote{不幸}的你！……\par

12. \Quote{那报复你、照你对待我们的方式对待你的，是有福的。}他指的是什么报复？圣咏以此结束：\BibleQuote{那抓住你的婴孩摔在磐石上的，是有福的} \BibleRef{咏 136:9}。他称她为不幸，却称那照她对待我们的方式报复她的人为有福。我们要问，什么报应？这就是报复。因为那个巴比伦对我们做了什么？我们已经在另一篇圣咏中唱过：\Quote{恶人的话压制了我们。}因为当我们出生时，这个世界的混乱就找到了我们，并在我们还是婴儿时就用各种错误的空洞观念窒息了我们。那生来注定成为耶路撒冷公民、在天主的预定中已是公民、但暂时被俘的婴儿，除了父母在他耳边低语的，又能学到爱什么？他们教导并训练他贪婪、盗窃、天天说谎、崇拜各种偶像和魔鬼、使用邪术和护身符等非法疗法。一个尚在襁褓中的婴儿，一个柔弱的灵魂，看着长辈们所做的一切，除了效法他所见的，还能做什么？巴比伦于是当我们年幼时迫害了我们，但天主在我们长大时赐予我们认识自己，使我们不再追随父母的错误。……他们如何报复她？照她对待我们的方式。让她的小孩子也被窒息：是的，让她的小孩子也被摔死。什么是巴比伦的小孩子？就是刚出生的邪恶欲望。因为有些人必须与根深蒂固的情欲作战。当情欲诞生时，在邪恶习惯获得力量攻击你之前，当情欲还小时，绝不要让它获得邪恶习惯的力量；当它还小的时候，就摔碎它。但你害怕，唯恐虽然摔了它却不死；\BibleQuote{将它摔在磐石上；而那磐石就是基督。} \BibleRef{格前 10:4}\par

13. 弟兄们，不要让你们的乐器在工作时休息：要彼此唱起\Place{熙雍}的歌谣。你们已欣然听了；更乐意地实行你们所听到的，如果你们不愿作\Place{巴比伦}的柳树，被它的河水滋养却不结果实。但要为那永恒的\Place{耶路撒冷}叹息：你们的望德向那里先行，愿你们的生命跟随；在那里我们将与\Person{基督}同在。\Person{基督}如今是我们的元首；祂从高处统治我们；在那城中，祂将把我们拥入祂自己；我们将与天主的众天使同等。我们不敢凭自己想象这事，若非真理亲自应许了。因此，弟兄们，要渴慕这事，日夜思念。无论世界怎样快慰地照耀你们，不可妄自尊大，不可甘愿与你们的情欲谈和。它是一个长大的仇敌吗？让它被击杀在磐石上。它是一个小仇敌吗？让它被摔碎在磐石上。将长大的击杀在磐石上，将小小的摔碎在磐石上。让磐石得胜。要建造在磐石上，如果你们不愿被溪水、或风、或雨冲去。如果你们愿在这世上武装起来对抗诱惑，就让那永恒的\Place{耶路撒冷}的渴望在你们心中增长并坚强。你们的被掳必将过去，你们的幸福必将来临；最后的仇敌将被消灭，我们将与我们的君王一起凯旋，不再有死亡。\par

\SourceRecord{Translated by J.E. Tweed.；Nicene and Post-Nicene Fathers, First Series；Vol. 8.；Edited by Philip Schaff.；Buffalo, NY: Christian Literature Publishing Co.,；1888.；<http://www.newadvent.org/fathers/1801137.htm>.}

% structure: p0001 source=h1 target=section reason=page-title

\section{圣咏第一三八篇释义}

1. 这篇圣咏的标题简短而朴实，无需我们久留；因为我们知道达味所预表的是谁，并且在他身上我们也认出自己，因为我们也是那身体的肢体。整个标题是：\Quote{致达味本人}。让我们来看看，什么是\Quote{致达味本人}。圣咏的标题通常告诉我们其中所论的内容；但在此，既然标题并未告知这一点，只告诉我们它是向谁咏唱的，第一节便告诉我们整篇圣咏所论的内容：\Quote{我要称谢祢。}那么，让我们来听这称谢。但我首先提醒你们，圣经中\Quote{称谢}一词，当我们论及对天主的称谢时，有两种用法：罪过的告明与赞美的颂谢。罪过的告明人人都知道，赞美的颂谢却少有人留意。罪过的告明如此为人熟知，以致无论在圣经何处听到\Quote{主，我要称谢祢}或\Quote{我们要称谢祢}这些话，我们立即出于习惯这样理解，手就急忙捶打胸膛：人们完全习惯于不将\Quote{称谢}理解为别的，只理解为罪过的告明。但难道我们的主耶稣基督自己也是个罪人吗？祂在福音中说：\Quote{父啊，天地的主宰，我称谢祢}？祂接着说明祂称谢什么，好让我们明白祂的称谢是赞美的颂谢，而非罪过的告明：\Quote{父啊，天地的主宰，我称谢祢，因为祢将这些事瞒住了智慧及明达的人，而启示给了小孩子。}祂赞美了父，祂赞美了天主，因为祂不轻视谦卑的人，却轻视骄傲的人。我们现在要听的正是这样的称谢：对天主的赞美，感恩的称谢。\Quote{以我全心}。我将我的全心放在祢赞美的祭台上，向祢献上全燔的赞美祭……\Quote{主，我要以全心称谢祢，因为祢俯听了我的口中的言语}\BibleRef{咏 137:1}。什么口？岂不是我心的口？因为我们有那声音，是天主所听的，人的耳朵完全听不到。我们有一个内在的口，在那里我们祈求，从那里我们祈求，若我们为天主准备了居所或房屋，我们在那里说话，在那里被听。\Quote{因为祂离我们每一个人都不远，因为我们生活、行动、存在都在祂内}\BibleRef{宗 17:27}。除了罪以外，没有什么能使你远离天主。推倒罪的隔墙，你便与你所祈求的祂同在。\par

2. \Quote{且在天使面前我要歌颂祢}。不在人前我要歌颂，却在天使面前。我的歌是我的喜乐；但我对下界之物的喜乐是在人前，我对上界之物的喜乐是在天使面前。因为恶人不认识义人的喜乐：\Quote{我的天主说：恶人决无喜乐}。恶人在他的酒店里喜乐，殉道者在锁链中喜乐。那位今日庆祝其节庆的圣女克利斯比纳，她在什么中喜乐呢？当她被逮捕时，当她被带往审判官前时，当她被投入监狱时，当她被捆绑押出时，当她被抬上刑台时，当她受审时，当她被定罪时——她在这一切事中都喜乐；而那些可怜的人以为她可怜，其实她正在天使面前喜乐。\par

3. \Quote{我要向祢的圣殿下拜}\BibleRef{咏 137:2}。什么圣殿？我们将要居住的，我们将要朝拜的。因为我们急忙赶去朝拜。我们的心怀了孕，临产，寻找在哪里生产。哪里是朝拜天主的地方？……\Quote{天主的殿是圣的，}宗徒说，\Quote{这殿就是你们}\BibleRef{格前 3:17}。但很显然，天主居住在天使中。因此，当我们属灵的、而非属世的喜乐向主高唱诗歌，要在天使面前歌唱时，那天使的集会就是天主的圣殿，我们向着天主的圣殿下拜。下面有一个教会，上面也有一个教会；下面的教会是在所有信徒中，上面的教会是在所有天使中。但天使的天主降到了下面的教会，天使在地上服侍了祂\BibleRef{玛 4:11}，而祂服侍了我们；因为\Quote{我来不是为受服侍，而是为服侍}\BibleRef{玛 20:28}……天使的主为人类而死。因此，\Quote{我要向祢的圣殿下拜}；我的意思不是人手所造的殿，而是祢为自己所造的殿。\par

4. \Quote{我要因祢的仁慈和祢的真理称谢祢的圣名}……这些祢所赐给我的，我也按我的能力回报祢：仁慈，在帮助他人上；真理，在判断上。天主藉这些扶助我们，我们藉这些赢得天主的恩宠。因此，\Quote{主的一切道路都是仁慈和真理}。没有别的道路祂能到我们这里来，也没有别的道路我们能到祂那里去。\Quote{因为祢使祢的圣名超越一切}。弟兄们，这是怎样的感恩呢？祂使祂的圣名超越亚巴郎。由亚巴郎生依撒格；天主在那家中被尊崇；然后是雅各伯；天主被尊崇，祂说：\Quote{我是亚巴郎的天主，依撒格的天主，雅各伯的天主。}然后来了他的十二个儿子。主的名在以色列中被尊崇。然后来了童贞玛利亚。然后我们的主基督，\Quote{为我们的罪死了，为使我们成义而复活了}\BibleRef{罗 4:25}，以祂的圣神充满信徒，派遣人向外邦人宣讲：\Quote{你们悔改吧}等等\BibleRef{玛 3:2}。看，\Quote{祂使祂的圣名超越一切}。\par

5. \Quote{在我呼求祢的每一日，求祢快快俯听我}\BibleRef{咏 137:3}。为什么\Quote{快快}？因为祢曾说：\Quote{你还在说话时，我就要说：我在这里}\BibleRef{依 58:9}。为什么\Quote{快快}？因为如今我不寻求现世的幸福，我已从新约中学得了神圣的渴望。我不寻求土地，也不寻求地上的丰裕，也不寻求暂时的健康，也不寻求仇敌的倾覆，也不寻求财富，也不寻求高位：这些我一样也不寻求；因此\Quote{快快俯听我}。既然祢已教导我该寻求什么，就赐给我我所寻求的……\par

6. 让我们看看他所求的是什么，他有何权利说：\Quote{求你快听我。}因为你求什么，就应当速速蒙听？\Quote{祢必加增我。}加增可从多方面理解……因为人灵魂中因忧虑而加增：一个灵魂似乎加增的人，其实连恶习也加增了。那是匮乏的加增，而非丰盈的加增。那么，你这说过\Quote{求你快听我}的人，且已从身体、从一切尘世事物、从一切尘世欲望中完全引退，以致对天主说\Quote{祢必加增我的灵魂}，你所渴望的是什么？请进一步解释你所求的。他说：祢必以德行加增我的灵魂。\par

7. \BibleQuote{愿地上所有的君王都称谢祢，上主}\BibleRef{咏 137:4}。必将如此，也正如此，且每日如此；并且显示这话并非徒然，只是因为它指向未来。但他们在称谢祢、赞美祢时，也不可向祢求尘世之事。地上的君王还能求什么呢？他们不是已经拥有主权了吗？在地上，人还能渴望什么？主权是其渴望的顶点。他还能奢求什么？必然是某种更高的尊位。但或许越高越危险。因此，君王在地上地位越高，就越应在天主面前谦卑。他们做什么？\Quote{因为他们听见了祢口中的一切言语。}在一个民族中隐藏着法律和先知，\Quote{祢口中的一切言语}；唯独在犹太民族中才有\Quote{祢口中的一切言语}，宗徒称赞这民族说：\BibleQuote{犹太人有何优势？多方面，首先是因为天主的圣言交托了他们。}\BibleRef{罗 3:1}……这些是天主的话。基德红的羊毛意味着什么？它好比世界中的犹太民族，拥有圣事的恩宠，但不是公开彰显，而是藏在云雾或帐幔中，如同羊毛中的甘露。\BibleRef{民 6:37}时候到了，甘露要在禾场上显明；已显明了，不再隐藏。唯独基督是甘露的甘饴：你们在圣经中认不出祂，而圣经正是为祂而写。然而，\Quote{他们听见了祢口中的一切言语。}\par

8. \BibleQuote{愿他们在上主的道路上歌唱，因为上主的荣耀伟大}\BibleRef{咏 137:5}。愿地上所有的君王都在上主的道路上歌唱。在什么道路上？就是上面所说的\Quote{在祢的仁慈和祢的真理中}。因此地上的君王不可骄傲，要谦卑。他们若谦卑，就在上主的道路上歌唱：他们要爱，就必歌唱。我们知道旅人歌唱；他们歌唱，并急忙赶往旅程的终点。有邪恶的歌曲，属于旧人；新人有一首新歌。故此，愿地上的君王也在祢的道路上行走，在祢的道路上行走并歌唱。歌唱什么？歌唱\Quote{上主的荣耀伟大}，而非君王的荣耀。\par

9. 请看他是如何希望君王在路上歌唱，谦卑地背负主，不抬举自己反对主。因为若他们抬举自己，结果如何？\BibleQuote{因为主至高，却垂顾卑微的人}\BibleRef{咏 137:6}。那么君王希望祂垂顾他们吗？让他们谦卑。然后呢？若他们骄傲自高，能逃脱祂的眼目吗？恐怕因为你听见\Quote{祂垂顾卑微的人}，你便选择骄傲，心里说：天主垂顾卑微的人，不垂顾我，我想做什么就做什么。愚昧的人啊！你若知道应当爱什么，还会说这话吗？看，即使天主不愿看你，你岂不正是害怕这事，即祂不愿看你？……看来，祂不垂顾高傲的人，因祂垂顾谦卑的人。\Quote{高傲的人}——什么？\Quote{祂从远处察知。}那么骄傲者得到什么？从远处被看见，而非逃避被看见。不要以为因此你就安全，因为从远处看你的人看得不那么清楚。其实，你从远处看确实看不清；而天主虽然从远处看你，却看得完全，只是祂不与你同在。你所得的，不是被看得不完全，而是你不在看你的那位面前。但谦卑者得什么？\Quote{主亲近痛悔的人。}所以，让骄傲者任意抬高自己吧，天主确实居于高处，天主在天上：你希望祂亲近你吗？自卑吧。因你越抬举自己，祂就越在你之上。\par

10. \BibleQuote{我若在患难中行走，祢必使我复苏}\BibleRef{咏 137:7}。确实如此：无论你处于何种患难，承认吧，呼求祂；祂释放你，祂使你复苏……爱另一个生命，你就会看见此生命是患难，无论它闪耀着何等繁荣，洋溢着何等欢愉；因为我们尚未拥有那最终最安全、无任何诱惑的喜乐，那是天主为我们保留的，无疑这是患难。让我们也理解这里说的患难是什么意思，弟兄们。并非如同他说：\Quote{若有任何患难临到我，祢必从中释放我。}而是他怎么说？\Quote{我若行走}等等；这就是说，除非我在患难中行走，否则祢不会使我复苏。\par

11. \Quote{祢在仇敌的怒气上伸出祢的手，祢的右手必使我安全。}让我的仇敌发怒：他们能做什么？他们可以夺我钱财，剥我衣物，流放我，放逐我；以忧伤和折磨加害我；最后，若获准，甚至杀我：他们还能做什么？但在仇敌所能做的之上，祢已伸出祢的手。因为仇敌不能将我与祢分开：祢越迟延，越加倍报仇……然而不是使我绝望；因为接着有：\Quote{祢的右手必使我安全。}\par

12. \Quote{主，祢要为我报复}\BibleRef{咏 137:8}。我不报复；祢要报复。让我的敌人任意发怒吧：祢要报复我所不能报复的……\Quote{亲爱的诸位，你们不要报复自己，}宗徒说，\Quote{但要给忿怒留有余地，因为经上记载：报复属于我，我必要偿还——主说。}\BibleRef{罗 12:19}这里还有另一个不容忽视、或许更值得采用的解释。\Quote{主}基督，\Quote{祢要为我偿还。}因为如果我偿还，我便侵夺了；祢却偿还了祢未曾侵夺的。主，祢要\Quote{为我偿还}。看，祂为我们偿还了。那些前来征收丁税的人来到祂跟前：\BibleRef{玛 17:24}他们向来要求每人缴纳一个\OriginalTerm{双德拉克马}{didrachma}作为丁税，即每人两个德拉克马；他们来到主那里要求交税；或者更确切地说，不是到祂那里，而是到祂的门徒那里，对他们说：\Quote{你们的老师不交税吗？}他们来告诉了祂。祂对伯多禄说：\Quote{为了不得罪他们，你到海里去，投下钓钩，拿那首先上来的鱼，开了它的口，你就会找到一块\OriginalTerm{斯塔特}{stater}。你拿去，为我和你缴纳。}那首先从海中上来的，正是从死者中首生者。在祂的口中，我们找到了两个双德拉克马，即四个德拉克马；在祂的口中，我们找到了四福音。藉着这四个德拉克马，我们脱离了这世界的索债；藉着四位福音作者，我们不再欠债；因为那里偿还了我们所有罪债。那么，祂已经为我们偿还了，感谢祂的仁慈。祂一无所欠；祂不是为自己偿还；祂是为我们偿还……\par

13. \Quote{主，祢的仁慈永远长存。}……我不愿只是暂时得到释放。\Quote{祢的仁慈永远长存，}祢以此仁慈释放了殉道者，并迅速将他们从这生命带走。\Quote{不要轻视祢亲手所造的工程。}主，我并非说\Quote{不要轻视我双手的工程：}我不夸耀我自己的工程。诚然，\Quote{我曾在夜间在祂面前以双手寻求主，并没有受骗；}然而我仍不赞美自己双手的工程；我怕祢审视它们时，在它们中发现的罪过多于功绩。请看在我身上的，是祢的工程，不是我的；因为祢若见到我的，祢就定我的罪；祢若见到祢的，祢就加冠。因为凡我所有的善工，都是从祢而来，所以它们更属于祢而非我。\BibleRef{弗 2:8}因此，无论就我们是人而言，还是就我们已从邪恶中改变而成义而言，主，\Quote{不要轻视祢亲手所造的工程。}\par

\SourceRecord{Translated by J.E. Tweed.；Nicene and Post-Nicene Fathers, First Series；Vol. 8.；Edited by Philip Schaff.；Buffalo, NY: Christian Literature Publishing Co.,；1888.；<http://www.newadvent.org/fathers/1801138.htm>.}

% structure: p0001 source=h1 target=section reason=page-title

\section{圣咏第一百三十九篇释义}

1. ……我们的主耶稣基督在先知中发言，有时以祂自己的名义，有时以我们的名义，因为祂使自己与我们合一；正如经上所说：\Quote{二人成为一体。}因此，主在福音中论及婚姻时也说：\Quote{从此他们不是两个，而是一体了。}成为一体，是因为祂为了我们的有死之性而取了肉身；并非成为同一神性，因为祂是造物主，我们是被造物。所以，我们的主以祂所取肉身的身份所说的一切，既属于那已经升天的元首，也属于那些仍在尘世漂泊中劳苦的肢体。那么，让我们听我们的主耶稣基督在预言中说话吧。因为圣咏早在主由玛利亚诞生之前就已咏唱，但并非在主成为主之前：因为祂从永远就是万物的创造者，但在时期满足时由祂的受造者所生。让我们相信那神性，并尽可能理解祂与圣父同等。但那与圣父同等的天主性，分享了我们的有死之性，不是出于祂自己的富足，而是出于我们的；好使我们也能分享祂的天主性，不是出于我们的富足，而是出于祂的。\par

2. \BibleQuote{主，祢已经试炼了我，并认识了我}\BibleRef{咏 138:1}。让主耶稣基督自己来说这话吧；也让祂向圣父说：\Quote{主}。因为圣父不是祂的主，除非因为祂屈尊按肉身出生。祂是天主的父，也是人的主。你愿意知道祂是谁的父吗？是那同等的圣子的父。宗徒说：\BibleQuote{祂虽具有天主的形体，却没有以自己与天主同等为强夺的。}\BibleRef{斐 2:6}对此\Quote{形体}，天主是父，那\Quote{形体}与祂同等，是独生子，由祂的本体所生。但为了我们，好使我们得以再造，成为祂天主性的分享者，藉着更新而得永生，祂成为了我们有死之性的分享者；那么宗徒论到祂说了什么呢？他说：\Quote{祂却空虚了自己，取了仆人的形体，成了人的模样，且以人的形象出现。}祂具有天主的形体，与圣父同等；祂取了仆人的形体，为要在其中比圣父小……\par

3. \BibleQuote{祢已经知道我的坐下和起来}\BibleRef{咏 138:2}。这里什么是\Quote{坐下}，什么是\Quote{起来}？那坐下的，就是谦卑自己。主在祂的苦难中\Quote{坐下}，在祂的复活中\Quote{起来}。经上说：\Quote{祢}已经知道这事；也就是说，祢愿意，祢认可；按照祢的旨意成就了。但如果你愿意将元首的话理解为在身体的身份上：人坐下时，是在忏悔中谦卑自己；他起来时，是罪过得赦免，并被提升至永生的希望中。不要高抬自己，除非你先被贬抑。因为许多人希望在坐下之前就起来，希望在承认自己是罪人之前就显为义……\par

4. \BibleQuote{祢从远处已经知道了我的思念；祢已经细查了我的道路和我的界限}\BibleRef{咏 138:3}；\BibleQuote{并且祢已经预见了我的所有途径}\BibleRef{咏 138:4}。什么叫做\Quote{从远处}？当我还在旅途中，尚未到达我的真正家乡之前，祢已经知道了我的思念……那小儿子去了远方。在他劳苦、受苦、受难、匮乏之后，他想起了父亲，渴望回去，便说：\Quote{我要起来，到我父亲那里去。}他说：\Quote{我要起来}，因为他之前是坐着的。因此，你可以认出他在说：\Quote{祢已经知道我的坐下和起来。}我坐着，在匮乏中；我起来，在渴望祢的食粮中。\Quote{祢从远处已经知道了我的思念。}因为我的确走到了远处；但我所离开的祂，何处不在呢？所以主在福音中说，他的父亲在他回来时迎接了他。确实如此；因为\Quote{祂从远处已经知道了他的思念。}\Quote{我的道路}，他说；那是什么，岂不是一条坏路，就是他离开父亲所走的路？……什么是\Quote{我的道路}？我所行经的路。什么是\Quote{我的界限}？我所达到的境地。\Quote{祢已经细查了我的道路和我的界限。}那我的界限，虽在远方，却不远离祢的眼目。我已经走得远，但祢仍在那里。\Quote{并且祢已经预见了我的所有途径。}祂没有说\Quote{已经看见}，而是说\Quote{已经预见}。在我行经它们之前，在我行走在它们之前，祢已经预见了它们；祢允许我在劳苦中走自己的道路，好使我不愿劳苦，就回到祢的道路中。\Quote{因为我的舌头上没有诡诈。}祂这话是什么意思？看，我向祢承认，我走了自己的道路，我远离了祢，我离开了祢——与祢同在时，我原是好的，而没有祢，对我来说是坏的……\par

5. \BibleQuote{看，主，祢已经知道了我所有的末后作为和往日的作为}\BibleRef{咏 138:5}。祢知道我末后的作为，就是我牧猪的时候；祢知道我往日的作为，就是我向祢要求我的家产分的时候。往日的作为对我而言是末后灾祸的开始：往日的罪，当我们堕落时；末后的惩罚，当我们进入这劳苦而危险的有死性时。但愿这对我们来说是\Quote{末后}；如果我们现在愿意回头，就会如此。因为对某些恶人还有另一个\Quote{末后}，他们将要听到：\BibleQuote{你们这些受咒骂的，到永火里去吧！}\BibleRef{玛 25:41}……\Quote{祢造就了我，并将祢的手加于我。}\Quote{造就了我}，在哪里？在这有死性中；就在我们众人所要生的劳苦中。因为没有人出生，天主不曾在母胎中造就他；也没有任何受造物是天主不是其造就者的。但祢\Quote{在这劳苦中造就了我}，\Quote{并将祢的手加于我}，祢惩罚的手，压制骄傲的人。因为祂健康地推倒了骄傲的人，为要高举谦卑的人。\par

6. \Quote{祢的技能在我身上奇妙地显为伟大：它变得强盛；我将不能达到它} \BibleRef{咏 138:6}。现在请听，且要了解一些隐晦的事，但它能带来理解上不小的乐趣。天主的圣仆梅瑟，天主曾藉云彩与他说话，因为，用人情的方式来说，祂必须藉祂所取的某种受造物与祂的仆人说话……他切愿看见天主的真容，并对那与他交谈的天主说：\Quote{如果我在祢眼前蒙恩，求祢将祢自己显示给我。} \BibleRef{出 33:13} 当他热切渴望此事，并几乎要藉那种友善的熟悉（如果可以这样说）从天主那里强求出来，就是天主屈尊与他交往时，好使他能看见祂的荣耀和祂的面容——按我们所能谈论的天主的面容——天主就对他说：\Quote{你不能看见我的面容；因为没有人看见了我而能存活。} \BibleRef{出 33:20} 但我将把你放在磐石穴中，我要经过，并用手遮盖你；当我经过以后，你要看见我的后背。从这些话中又产生另一个谜，即真理的一个隐晦形象。\Quote{当我经过以后，} 天主说，\Quote{你要看见我的后背；} 好像祂有一面是面容，另一面是后背。愿我们远离对那威严的任何这种想法！因为凡对天主有这种想法的人，即使圣殿关闭，对他又有何益？他是在自己心中建造偶像。这些话中蕴藏着伟大的奥秘。……那些曾向他们所见的基督发怒的人，如今寻求得救的办法；有人对他们说：\Quote{你们悔改，并要因耶稣基督的名受洗，你们的罪就得赦免。} \BibleRef{宗 2:38} 看，他们看见了祂的后背，而他们本来不能看见祂的面容。因为祂的手曾放在他们的眼睛上，并非永远，而是当祂经过时。祂经过后，就移开祂的手。当手从他们眼睛上移开时，他们对门徒们说：\Quote{我们当做什么？} 起初他们凶狠，后来变得仁爱；起初愤怒，后来恐惧；起初刚硬，后来温和；起初盲目，后来明悟。……\par

7. 你看，那个在远方逃走的逃亡者，也不能逃脱他所逃离的那位的眼睛。如今他还能往哪里去，他的\Quote{界限已被追踪}？看，他说什么？\Quote{我从祢的圣神那里往哪里去？} \BibleRef{咏 138:7} 谁能从充满世界的那圣神那里逃离呢？ \BibleRef{智 1:7} \Quote{我从祢的面容那里往哪里逃？} 他寻找一个地方，好从天主的义怒中逃离。什么能庇护天主的逃亡者？那些庇护逃亡者的人，问他们从谁那里逃来；当他们发现某人是某个不如自己强大的主人的奴隶时，他们似乎毫无恐惧地庇护他，心里说：\Quote{他的主人不能追踪他。} 但当他们听说一个强大的主人时，他们要么不庇护，要么极其恐惧地庇护，因为即使一个强大的人也可能被欺骗。天主在哪里不在？谁能欺骗天主？天主不看谁？天主不向谁追讨祂的逃亡者？那么，那个逃亡者从天主的面容那里往哪里去呢？他再三转向这里那里，好像寻找一个可逃的地方。\par

8. \Quote{如果我升到天上，祢在那里；如果我下到阴府，祢也在那里} \BibleRef{咏 138:8}。可怜的逃亡者，你终于明白了，你绝不能使你远离那位你曾想远远离开的祂。看，祂无所不在；你，你要往哪里去？他找到了主意，而且是由那位如今屈尊召回他的祂所启示的……如果我因犯罪而下到邪恶的深处，并拒绝承认，说：\Quote{谁看见我？}（因为\Quote{在阴府中谁能承认祢？}）那里也有祢在惩罚。那么，我往哪里去才能逃避祢的面容，即不遇见祢的义怒呢？他找到了这个主意：我要这样逃避祢的面容，这样逃避祢的圣神；逃避祢惩罚的圣神，祢惩罚的面容，这样逃避。如何？\Quote{若我再取双翼向前，并住在海的极处} \BibleRef{咏 138:9}。这样我能逃避祢的面容。若他逃到海的极处以逃避天主的面容，他所逃离的那位岂不也在那里？……因为\Quote{海的极处}是什么，不就是世界的终结吗？让我们在希望和渴慕中现在就飞到那里，藉着双重的爱之翼；愿我们只在\Quote{海的极处}得安息。因为若我们在别处寻求安息，我们将被投入海中。让我们甚至飞到海的极处，让我们藉着双重的爱之翼高飞；同时，让我们在希望中飞向天主，并在忠信的希望中默想那\Quote{海的极处}。\par

9. 现在请听谁可能带我们到那里。正是那位我们想要逃避其义怒的面容的同一位置。因为下文说：\Quote{就是在那里，祢的手要引导我，祢的右手要带领我} \BibleRef{咏 138:10}。亲爱的弟兄们，让我们默想这话，愿这是我们的希望，我们的安慰。让我们藉着爱再次取得我们因贪欲而失去的翅膀。因为贪欲是我们翅膀的粘胶，它把我们从天上的自由，即天主圣神的自由微风中击落。被击落后我们失去了翅膀，并可以说被囚禁在捕鸟者的权下；我们从那要捉拿我们而逃的那位那里，被祂以自己的血救赎。祂用祂的诫命为我们制造翅膀；我们现在高举它们，脱离粘胶……因此我们必须有翅膀，也必须由祂引导，因为祂是我们的助佑。我们有自由意志；但即使有那自由意志，若那位命令我们的不帮助我们，我们能做什么呢？\par

10. 他考虑路途的遥远，对自己说了什么？\Quote{我说，恐怕黑暗会淹没我}11. 看，如今我已信了基督，如今我藉着双重的爱被高举……关于路途的遥远，我对自己说：\Quote{黑夜在我欢愉中成了光明}。黑夜对我来说成了光明，因为在黑夜中我曾绝望，无法穿越如此广阔的大海，无法克服如此漫长的旅途，无法恒心坚持到底而达到终极。感谢那寻找我这逃犯、用鞭打击打我的背、以召叫使我脱离毁灭、使我的黑夜成为光明的那一位。因为当我们穿越此生时，一直是黑夜。黑夜如何成了光明？因为基督降临到了黑夜之中……\par

11. \BibleQuote{因为黑暗在祢面前不会变暗} \BibleRef{咏 138:12}。所以你不要使你的黑暗变暗；天主不使它变暗，反而更加光照它；因为另一篇圣咏向祂说：\Quote{主，祢要照亮我的灯：我的天主要光照我的黑暗}。但那些使他们的黑暗变暗的人，天主不使之变暗的，是谁呢？是恶人，是悖逆的人；当他们犯罪时，他们确实是黑暗；当他们不承认自己所犯的罪，反而继续辩护时，他们就\Quote{使他们的黑暗变暗}。因此，若你现在犯了罪，你就在黑暗中；但通过承认你的黑暗，你将获得你的黑暗被光照；若你辩护你的黑暗，你将\Quote{使你的黑暗变暗}。而你将逃到哪里去躲避双重的黑暗呢？你已经在单一的黑暗中艰难……让我们不要通过辩护我们的罪来\Quote{使我们的黑暗变暗}，并\Quote{黑夜在我们欢愉中成了光明}。\par

12. \Quote{黑夜要光明如白昼}。\Quote{黑夜如白昼}。\Quote{白昼}对我们而言是世俗的顺境，黑夜是此世的逆境；但是，若我们领悟到我们因自己的罪过而遭受逆境，并且我们父亲的鞭挞对我们来说是甘甜的，以免审判者的判决对我们苦涩，那么我们就将发现这黑夜的黑暗，仿佛就是这黑夜的光明……但当我们的主基督来临，并藉信德居住于灵魂内，应许了另一种光明，激励并赐予忍耐，警告人不要因顺境而得意，也不要因逆境而沮丧，那人因有信德，便开始漠视此世；不因顺境而得意，也不因逆境而沮丧，反而在一切事上赞美天主，不仅在他富足时，也在他失落时；不仅在他健康时，也在他生病时……\Quote{祂的黑暗如何，祂的光明也如何}。祂的黑暗不会压倒我，因为祂的光明不会高举我。\par

13. \BibleQuote{因为主，祢已占有了我的脏腑} \BibleRef{咏 138:13}。占有者在内；祂不仅占据内心，也占据脏腑；不仅占据思想，也占据喜乐：祂占有那使我为世上任何光明而感到喜乐的根源；祂占据了我的脏腑：除了祂智慧的内在光明，我不知道别的喜乐。那么如何？你不因你的事务非常顺利、时机有利而喜乐吗？你不因荣誉、财富、家庭而喜乐吗？\Quote{我不，}他说。为什么？因为\Quote{主，祢已占有了我的脏腑；祢从母腹中提携了我}。当我在母腹中时，我并不漠视那夜的黑夜和那夜的光明……如今，从我们那位母亲的腹中被提携出来后，我们漠视它们，说：\Quote{祂的黑暗如何，祂的光明也如何}。地上的顺境不能使我们幸福，地上的逆境也不能使我们不幸。我们必须持守正义，珍爱信德，寄望于天主，爱天主，也爱我们的邻人。在这些辛劳之后，我们将拥有不灭的光明，没有黄昏的白昼。这黑夜的一切光明与黑暗都是转瞬即逝的。\par

14. \BibleQuote{我要向祢称颂，主，因为祢被奇妙地造成了可畏：祢的作为奇妙，我的灵魂深知此点} \BibleRef{咏 138:14}。先前，\Quote{祢的知识从我身上成了奇妙，它已变得伟大，我无法企及}。因此，\Quote{它已变得伟大}。现在，\Quote{我的灵魂}如何\Quote{深知此点}，若不是因为\Quote{黑夜在我欢愉中成了光明？}若不是因为祢的恩宠临到我，光照了我的黑暗？若不是因为祢占有了我的脏腑？若不是因为祢从母腹中提携了我？\par

15. \Quote{我的骨头，你为我藏在隐秘处所造的，不向你隐藏}\BibleRef{咏 138:15}。他说：\Quote{他的骨头。}百姓称之为\Emph{ossum}，在拉丁文中称为\Emph{os}。这正是希腊文中的用词。因为我们可能以为这里的\Emph{os}是复数\Emph{ora}的那个词，而不是短音\Emph{os}（它构成\Emph{ossa}）。于是他说：我有一块骨头（\Emph{ossum}）在隐秘处。我们宁愿使用这个词；学者们责怪我们，总比百姓不懂我们好。\Quote{因此，}他说，\Quote{我有一块骨头，在里面，是隐藏的；你在我里面隐秘地造了骨头，但在你面前它并不隐藏。你隐秘地造了它，难道你就因此将它向你自己隐藏了吗？}你造了我的这块骨头，人看不见，人不知道：你这位造物主知道。弟兄们，那么这\Quote{骨头}指什么？让我们寻找它，它是\Quote{在隐秘处}。但因为我们作为基督徒，奉主的名向基督徒说话，现在我们发现这种骨头是什么。它是一种内在的力量；因为力量与刚毅被认为是存在于骨头中的。那么，灵魂有一种内在的力量，它不会折断。无论什么折磨、什么磨难、什么逆境在周围肆虐，天主在我们内隐秘地造的那坚固的，不能被折断，也不屈服。因为天主造了某种忍耐的力量，在另一篇圣咏中提到：\Quote{我的灵魂，请只安息于天主，因为我的忍耐是由他而来。}……你因什么而夸耀？\Quote{在磨难中，知道磨难生忍耐。}\BibleRef{罗 4:5}请看这力量是如何在他心中形成的：\Quote{因为天主的爱，借着所赐与我们的圣神，已倾注在我们心中了。}就这样，那隐藏的骨头被形成并变得坚强，以至于使我们甚至在磨难中也夸耀。但在人看来我们似乎不幸，因为我们内在的东西对他们来说是隐藏的。\Quote{我的形体在地的下层。}看，我的形体在肉身内，但我体内有一块骨头，是你所造的，它使我绝不屈服于这低处的任何迫害，我的形体仍在这里。因为如果一位天使勇敢，那算什么大事？如果肉身勇敢，这才是大事。肉身从哪里勇敢，瓦器从哪里勇敢，除非因为其中有一块骨头在隐秘处被造？\par

16. ……\Quote{你的眼睛看见了我不完全的形体，并且在你的书上，一切都要被写下}\BibleRef{咏 138:16}，不仅是完全的，也包括不完全的。让不完全的人不要害怕，只让他们前进。然而，因为我曾说：\Quote{让他们不要害怕}，不要让他们喜爱自己的不完全，并停留在他们被发现的境地。让他们尽力前进。每天让他们增添，每天让他们接近；但不要让他们从主的身体上脱落，好使他们融合在一个身体里，并在这些肢体中，被堪称为配得上那句话。\Quote{他们必在白日流离，其中没有一人。}\Quote{白日}仍在地上，就是我们的主耶稣基督。因此他说：\Quote{你们趁着白日行走。}\BibleRef{若 12:35}但\Quote{在白日}他的不完全者\Quote{必流离}。他们也曾以为我们的主耶稣基督只是一个人，他内中没有隐藏的天主性，他并不是隐秘的天主，而只是那所见的；他们也这样想……但\Quote{他们在白日流离}是什么意思？他们会灭亡吗？那么\Quote{在你的书上，一切都要被写下}在哪里呢？那么他们何时\Quote{在白日流离}？当他们不明白那安置在地上的主的时候。接着是什么？\Quote{但对我来说，天主，你的朋友极其尊贵}\BibleRef{咏 138:17}；那些曾在\Quote{白日流离，其中无人}的，成了你的朋友，对我变得极其尊贵。在主复活之后，那骨头在他们内隐秘地被造，他们为他的名而受苦，他们曾对他的死亡感到惊愕。\Quote{他们的首领被大力坚固。}他们成了宗徒，成了教会的领袖，成了领导羊群的首羊，\Quote{被大力坚固。}\par

17. \Quote{我要数算他们，他们要比沙更多}\BibleRef{咏 138:18}。借着那些在\Quote{白日流离}的人，看啊！生出了所有这伟大的群众，现在如同沙子般不可胜数，除非由天主。因为他说：\Quote{他们要比沙更多}，然而他曾说：\Quote{我要数算他们。}正是那些被数算的，\Quote{要比沙更多}。因为由他，沙子被数算，由他，\Quote{我们头上的头发都被数过}。\BibleRef{玛 10:30}\Quote{我复活了，我仍与你同在。}我已经受苦了，他说，我已经被埋葬了；看啊！我复活了，但他们还不明白我同他们在一起。\Quote{我仍与你同在，}就是说，还不与他们同在，因为还不认识我。因为我们在福音中读到，我们的主耶稣基督复活后，当他显现给他们时，他们并没有立刻认出他。还有另一个意义：\Quote{我复活了，我仍与你同在，}仿佛他要指明这如今的时候，他仍隐藏在父的右边，直到他在光荣中显现，那时他要来审判生者死者。\par

18. 然后祂讲述在此期间，即当祂已经复活，仍然与圣父同在的整个时期，祂在祂的身体——教会内，因罪人的混杂和异端者的分离而受的苦楚。\Quote{天主，你若杀戮恶人（因为你将在你的思想中说：离开我，你们这些嗜血的人），他们必虚妄地接收他们的城市}（第19—20节）。这些话似乎按此顺序相连：\Quote{天主，你若杀戮恶人，他们必虚妄地接收他们的城市。}恶人就是这样被杀戮，因为\Quote{他们的理智受了蒙蔽，与天主的生命隔绝了。}\BibleRef{弗 4:18}因为由于高傲，他们丧失了忏悔，因而被杀戮，在他们身上应验了圣经所说的：\Quote{忏悔从死者身上消失，如同从不存在的人一样。}\BibleRef{德 17:28}因此\Quote{他们虚妄地接收他们的城市}，即他们虚妄的民众，这些民众追随他们的虚妄；当他们被正义之名所膨胀时，便说服人们打破合一的纽带，盲目无知地跟随他们，以为他们更正义……但现在基督的身体——教会说：骄傲的人为什么诬告我，好像我被别人的罪玷污了，以致他们藉着分离自己，\Quote{虚妄地接收他们的城市}？\Quote{主，我岂不恨那恨你的人？}\BibleRef{咏 138:21}那些自身更坏的人为什么要求我也在身体上与恶人分离，如同在精神上一样，以致在收割时节之前，连麦子一同拔除莠子，在簸扬之前失去忍受糠秕的能力，在各类鱼被带到世界末日如同岸边一样分离之前，撕裂和平与合一的网？我所领受的圣事，难道是恶人的吗？我是否藉着同意而参与他们的生活和行为？……但是，\Quote{爱你们的仇敌}在哪里？是因为祂说\Quote{你们的}，而不是\Quote{天主的}吗？\Quote{善待那恨你们的人。}\BibleRef{玛 5:44}祂不是说\Quote{那恨天主的人}。因此他效法此榜样，说：\Quote{主，我岂不恨那恨你的人？}他不是说\Quote{那恨我的人}。\Quote{我因你的敌人而消瘦。}他说\Quote{你的}，而不是\Quote{我的}。但那恨我们、与我们为敌的人，仅仅因我们事奉祂，他们所做的岂不正是恨祂，成为祂的敌人？我们难道应当爱这样的敌人吗？或者，那些为天主而受迫害的人，不是听到他们所说\Quote{为那迫害你们的人祈祷}吗？那么请看下文。\Quote{我以完全的仇恨恨他们}\BibleRef{咏 138:22}。什么是\Quote{以完全的仇恨}？我在他们身上恨他们的不义，我爱你的受造。这就是以完全的仇恨去恨，既不因恶习而恨其人，也不因其人而爱恶习。请看他还加上：\Quote{他们成了我的敌人。}他现在不仅把他们描述为天主的敌人，也描述为他自己的敌人。那么，他如何在他们身上既实现自己所说的\Quote{主，我岂不恨那恨你的人}，又实现主的命令\Quote{爱你们的仇敌}呢？除非藉着那\Quote{完全的仇恨}，即恨他们的邪恶，爱他们的人性，他如何成就这事？即使在旧约时代，当血肉之民受可见惩罚的约束时，天主的仆人梅瑟——他因理解而属于新约——是如何恨罪人的？他既为他们祈祷，又杀戮他们，岂不正是\Quote{以完全的仇恨恨他们}？因为他如此完全地恨他所惩罚的不义，以至于爱他所祈祷的人性。\par

19. 既然基督的身体最终也要在身体上与不圣洁和邪恶的人分离，但此刻仍在他们中间呻吟，那\Quote{在女子中的爱，如同百合花在荆棘中}\BibleRef{歌 2:2}又是什么？她的话是什么？她的良心是什么？那\Quote{王女在里面的华美}是什么？请听她所说的话。\Quote{天主，求你察验我，知道我的心}\BibleRef{咏 138:23}。天主，求你察验，求你知晓；不是人，不是异端者，他既不懂得如何察验，也不能知晓我的心；惟独你察验并知道我没有同意恶人的行为，尽管他们以为我会被别人的罪玷污；因此，当我在漫长的漂泊中，做我在另一篇圣咏中所哀叹的事，即\Quote{我在那恨和平的人中为和平劳苦}，直到我来到那被称为耶路撒冷的和平神视，\Quote{它是我们众人的母亲}，那\Quote{在天上永恒的}城邑；他们争论、诬告并分离自己，\Quote{虚妄地接收}——显然不是在永恒中，而是——\Quote{他们的城市}。为何如此？请看下文。\par

20. \Quote{且看，}他说，\Quote{在我内是否有邪恶的道路，并引导我走永生的道路}\BibleRef{咏 138:24}。\Quote{察验}，他说，\Quote{我的路径}，即我的计谋和思念。这岂非说\Quote{引导我走向基督}？因为谁是\Quote{永生的道路}，除非那永生者？永生的祂曾说：\Quote{我是道路、真理、生命。}\BibleRef{若 14:6}因此，如果你在我的路上发现什么令你眼目不悦的事，因我的道路是必朽的，求你\Quote{引导我走永生的道路}，其中没有不义；因为\Quote{若有人犯罪，我们在父那里有一位辩护者——正义的耶稣基督；祂是为我们的罪作赎罪祭的；}\BibleRef{若一 2:1}祂是\Quote{永生的道路}，没有罪；祂是永生的生命，没有惩罚。\par

弟兄们，这些是伟大的奥秘。天主圣神如何与我们说话？祂如何使我们在黑夜中获得喜悦？弟兄们，我们问你们，这是什么？为何越晦暗，就越甘甜？祂按照祂的爱，以某些奇妙的方式为我们调配饮品。祂使祂自己的话变得奇妙，以致当我们讲论你们已经知道的事时，然而由于这是从看似晦涩的段落中发掘出来的，知识本身似乎变得焕然一新。弟兄们，难道你们不知道，在教会内应当容忍恶人，且不应制造分裂？难道你们不知道，在那网住好鱼和坏鱼的网中，我们必须坚持到岸边，不可撕破网，因为在岸边，好的将被收进器皿，坏的将被丢弃？你们已经知道这些；但你们不明白这篇圣咏的这几节；你们所不明白的被解释了；你们所知道的已被更新。\par

\SourceRecord{Translated by J.E. Tweed.；Nicene and Post-Nicene Fathers, First Series；Vol. 8.；Edited by Philip Schaff.；Buffalo, NY: Christian Literature Publishing Co.,；1888.；<http://www.newadvent.org/fathers/1801139.htm>.}

% structure: p0001 source=h1 target=section reason=page-title

\section{圣咏第一百四十篇释义}

1. 我们的诸位主上吩咐了我，弟兄们，而在他们内，万有的主，使我能按天主赐予我的程度，将这篇圣咏带至你们的理解中。愿祂藉你们的祈祷帮助我，使我能说出我该说的话，你们能聆听，好让我们所有人都能受益于天主的圣言。因为并非所有人都受益，因为\BibleQuote{并非所有人都具有信德}。\BibleRef{得前 3:2}...\par

2. 这篇圣咏所包含的内容，我相信当它被咏唱时，你们已经察觉；因为基督的教会，处于恶人之中，在其中抱怨、叹息，并向天主倾泻祈祷。因为在每如此这般预言中，她的声音是一个需要和匮乏之人的声音，尚未满足，\BibleQuote{饥渴慕义}，\BibleRef{玛 5:6}，对于他们，在终结时已应许并保留了某种圆满……\par

3. \Quote{至于终结，达味自己的圣咏}。你不可寻求其他终结，除了宗徒亲自为你们指定的。因为\BibleQuote{基督就是终结}。\BibleRef{罗 10:4}……祂是达味的后裔，不是按祂的天主性，藉此祂是达味的创造者，而是按血肉；因此祂屈尊在预言中被称为达味：注视这个\Quote{终结}，因为这首圣咏是\Quote{向达味本身}咏唱的；听祂奥体之声；成为祂奥体的一部分。让您所听见的声音成为您的，并祈祷，说出接下来的话。\par

4. \BibleQuote{上主，求祢从恶人手中拯救我} \BibleRef{咏 139:1}。不是仅仅从一个恶人，而是从整个类别；不是仅仅从那些器皿，而是从他们的首领本身，即魔鬼。为什么说\Quote{从人}，如果他的意思是脱离魔鬼？因为魔鬼在比喻中也被称为人。\BibleRef{玛 13:24}……如今既然成了光，不是凭自己，而是在主内，\BibleRef{弗 5:8}，让我们不仅对抗黑暗，即对抗仍被魔鬼占据的罪人，也要对抗他们的首领，魔鬼本身，他运行在悖逆之子身上。\Quote{求祢从恶人手中拯救我}。这与\Quote{从恶人}相同。因为他称他为恶人是因为他不义，免得你或许认为任何不义的人可能是好人。因为许多不义之人看似无害；他们不凶猛、不野蛮、不迫害、不欺压；然而他们是不义的，因为追随某种其他习性，他们奢侈、酗酒、耽于享乐……因此，每一个不义之人都是恶人，他必须是有害的，无论他是温和还是凶猛。无论谁落入他的道路，无论谁被他陷阱所擒，都会发现他以为无害的东西是多么有害。因为，弟兄们，即使是荆棘，其根部也不会刺人。将荆棘从地里拔起，触摸它们的根，看看你是否感到疼痛。然而，在生长过程中让你疼痛的东西，却是从那根生出的。因此，不要让那些看似温和善良却贪爱肉欲、追随污秽情欲的人取悦你；不要让他们取悦你。虽然他们现在看似温和，却是荆棘的根……因此，我的弟兄们，基督的奥体，基督的肢体，在如此恶人中叹息，无论你们发现谁一头栽进邪恶的情欲和致命的享乐，要立刻斥责，立刻惩罚，立刻焚烧。让根被烧尽，就没有再生出荆棘的场所。如果你们不能做到，必定会拥有他们作为敌人。他们可能沉默，可能隐藏敌意，但无法爱你们。但既然他们不能爱你们，既然恨你们的人必定会寻求你们的伤害，就不要让你们的舌头和心灵迟于向天主说：\Quote{上主，求祢从恶人手中拯救我}。\par

5. \BibleQuote{他们心中图谋不义} \BibleRef{咏 139:2}……从他们中拯救我，从他们中，愿祢的手最有力拯救我。因为躲避公开的敌意是容易的，避开一个公开明显的敌人是容易的，当不义在他嘴唇和心里时；他是一个麻烦的敌人，他是隐秘的，难以躲避，他在口中含着好东西，却在心里隐藏恶事。\Quote{他们终日制造战争}。什么是\Quote{战争}？他们为我制造了我终日要与之争战的东西。因为从那里，从这样的心中，产生了基督徒所对抗的一切。无论是叛乱、分裂、异端、还是骚乱的反对，无不源于这些隐藏的念头，当他们口说好话时，\Quote{他们终日制造战争}。你听见和平的话语，然而制造战争从未离开他们的思想。因为\Quote{终日}一词表示不间断地，在整个时间中。\BibleQuote{他们磨快舌头如同蛇} \BibleRef{咏 139:3}。如果你还想辨认那恶人，看一个比喻：在一切野兽中，蛇最狡猾诡诈去伤害；因此它爬行。它甚至没有脚，所以它来时的脚步声能被听见。在它的行进中，它仿佛轻轻拖着身体，却不笔直。因此他们这样爬行蠕动去伤害，甚至在温柔的触碰下隐藏着毒液。所以接着是：\Quote{虺蛇的毒液在他们的嘴唇之下}。看，是\Quote{在}他们的嘴唇之下，好让我们能察觉，他们嘴唇之下是一种东西，嘴唇之中是另一种……\par

6. \BibleQuote{上主，求你救我脱离恶人的手，保护我脱离强暴的人}\BibleRef{咏 139:4}。在此他们露出了真面目，他们被认出来了；这里我们无需理解，只需行动：我们需要祈祷，而不是问他们是谁。但你当如何向这样的人祈祷，他在下文解释。因为许多人不知如何向恶人祈祷。\BibleQuote{他们图谋，}他说，\BibleQuote{使我的脚步滑跌。}至此可以从肉体上理解。每个人都有敌人，他们在交易中设法欺骗他，夺走他的钱财，在共同从事的业务中；每个人都有某个邻居是他的敌人，他谋划如何给其家庭带来祸害，以某种方式损害其财产，他无疑是通过诡计、欺诈、恶魔般的计谋来达成此事：没有人能怀疑这一点。然而，防备他们并非为了这些原因，而是恐怕他们埋伏你，把你拉向自己，即，使你与基督的身体分离，并使他们成为你的身体。因为基督是善人的头，魔鬼是他们的头。\Quote{使我的脚步滑跌}是什么意思？并非指你在与他交易时被欺骗，或他在法庭上与你打官司时骗了你。他\Quote{使你的脚步滑跌}，如果他在天主的道路上阻碍了你；致使你所正确引导的脚步绊倒，或离开道路，或在路上跌倒，或从道路后退，或在路上止步，或回到原来的地方。凡是对你做了这些的，就是使你滑跌，欺骗了你。你要向这样的陷阱祈祷，免得你失去天上的产业，免得你失去基督，你的共同继承人，因为你注定要与祂永远生活，祂使你成为继承人。因为你成为继承人，并非要继承一个死后的人，而是要与一位永远生活的主一同生活。\par

7. \BibleQuote{骄傲的人为我暗设网罗}\BibleRef{咏 139:5}。当他提到\Quote{骄傲的人}时，他简要地描述了整个魔鬼的身体。因此，他们多半自称义人，其实是不义的。因此，没有什么比承认自己的罪更令他们痛苦。他们是那些假义人，必然嫉妒真正的义人。因为没有人会嫉妒另一个人在他不愿成为或看似成为的方面……由此产生所有引诱和使人跌倒的行为。这是魔鬼最初的愿望，当他堕落时，嫉妒那站立的凡人……\par

8. 但那些\Quote{骄傲的人为我暗设网罗；}他们试图使我的脚步滑跌。他们做了什么？\Quote{并张设绳索为网罗。}什么绳索？这个词在圣经中广为人知，我们在别处也能找到\Quote{绳索}的含义。因为经上说：\BibleQuote{各人被自己的罪恶束缚}\BibleRef{箴 5:22}。依撒意亚明确说：\BibleQuote{祸哉，那些用绳索牵引罪孽的人}\BibleRef{依 5:18}。为什么称为\Quote{绳索}？因为每个坚持犯罪的罪人，罪上加罪；当他本应通过控告自己的罪来改正时，却通过辩护使本可通过忏悔消除的罪加倍，并且常常试图用其他罪来巩固自己已有的罪……但他们将这些罪\Quote{铺开}，当劝说义人做他们自己所行的恶事时，就铺在义人面前。因此他说：\Quote{他们张设绳索和网罗；}即，他们用他们的罪想要推翻我。他们在哪里这样做？\Quote{在我的路径旁，他们为我设下绊脚石：}不是在路径上，而是\Quote{在路径旁}。你的\Quote{路径}是天主的诫命。他们\Quote{在路径旁设下绊脚石}；你不要离开路径，就不会撞上绊脚石。然而我并不愿你说：\Quote{天主应当阻止他们在我的路径旁设绊脚石，这样他们就不会设了。}恰恰相反，天主允许他们\Quote{在你的路径旁设绊脚石}，为叫你不可离开路径。\par

9. 还剩下什么？在这样的祸患、这样的诱惑、这样的危险中，有什么补救？\BibleQuote{我曾对上主说：祢是我的天主}\BibleRef{咏 139:6}。祈祷的声音是响亮的，它激发信心。祂不是别人的天主吗？谁不是那真天主的天主？然而祂尤其属于那些享受祂、事奉祂、甘心服从祂的人。因为恶人，即使不情愿，也服从祂……\Quote{求祢侧耳听我祈祷的声音。}他没有说：\Quote{求祢侧耳听我的祈祷；}而是，仿佛更清楚地表达他内心的情感：\Quote{我祈祷\Emph{的声音}}，我祈祷的生命，我祈祷的灵魂，不是我言语中发出声响的，而是赋与我言语生命的。因为所有其他无生命的声音都可以称为声响，而不是言词。言词属于有灵魂的、有生命的东西。但有多少人向天主祈祷，却既没有对天主的感知，也没有对天主的正确想法！这些人可能有祈祷的声响，却不能有声音，因为他们没有生命。这是一个活人的祈祷的声音，因为他明白天主是他的天主，看见他被谁释放，意识到他从谁那里被释放。\par

10. 将这交于天主的耳中，让他说：\Quote{主，主。}祢是主——主，即最真实的主，不像那些以钱袋买人的主，而是以祂的血买人的主。\Quote{主，主，我救恩的力量}（\BibleRef{咏 139:7}），即赐给我救恩力量的那位。\Quote{我救恩的力量}是什么意思？他抱怨罪人的绊脚石和陷阱，恶人、魔鬼的工具，他们在他周围狂吠，在他周围设下陷阱，那些嫉妒义人的骄傲者。但他立刻加上了安慰：\Quote{那坚持到底的，必要得救。}他注意到了这一点并感到害怕，因罪恶众多而苦恼，转而投向希望。我必将得救，如果坚持到底；但坚持到底以得救，是属于力量；祢是\Quote{我救恩的力量}，祢使我坚持，使我达到救恩。……于是，在这场战斗中劳苦，他回头仰望天主的恩宠；因为他已经开始感到灼热和干渴，便找到了一个荫庇之所，得以存活。\Quote{在战争之日，祢荫庇了我的头}：即，在炎热中，免得我发热，免得我被烤焦。\par

11. \Quote{主，求祢不要因我的私欲而把我交给罪人}（\BibleRef{咏 139:8}）。看，祢的荫庇对我有何益处：使我不受自身的炎热。那\Quote{罪人}即使狂怒，又能对我做什么？因为恶人曾向殉道者发怒，拖走他们，用锁链捆绑，囚禁在监狱，用刀剑杀死，抛给野兽，用火焚烧：他们做了这一切；但天主并未把他们交给罪人，因为他们没有被自己的私欲所交。因此，你要全力祈求，求天主\Quote{不要因你的私欲而把你交给罪人}。因为你自己的私欲给魔鬼留了地步。看，魔鬼在你面前摆上利益，引诱你行不义；你不能得到这利益，除非行不义：利益是诱饵，不义是陷阱。你要这样看诱饵，也要看到陷阱；因为你若不行不义，就不能得利益；你若行不义，就会被捕。……因此，在战争之日，你的头被荫庇。因为私欲引起炎热，但主的荫庇调和私欲，使我们能控制那使我们被冲走的东西，免得我们被热得被拖入陷阱。\Quote{他们图谋攻击我；求祢不要丢弃我，免得他们自高自大。}你在另一处看到：\Quote{压迫我的人若见我动摇，就会欢欣。}他们就是这样，因为魔鬼自己也是如此。……\par

12. \Quote{他们环绕的头，他们嘴唇的劳苦必遮蔽他们}（\BibleRef{咏 139:9}）。他说，我，祢翅膀的荫影必遮蔽我；因为\Quote{在战争之日，祢遮蔽了我}。遮蔽他们的将是什么？\Quote{他们环绕的头}，即骄傲。\Quote{他们环绕}是什么意思？他们如何环绕而不站立，如何在错误的圈中行走，那里是没有终点的旅程。走直线的人，从某点开始，在某点结束；走圆圈的人，永无终点。这就是恶人的劳苦，在另一篇圣咏中更清楚地表明：\Quote{恶人行走在圆圈中。}但\Quote{他们环绕的头}是骄傲，因为骄傲是一切罪的开始。但骄傲如何是\Quote{他们嘴唇的劳苦}？每个骄傲的人都是虚假的，每个虚假的人都是说谎者。人在说谎言时劳苦；而说真理，他们可以极其轻松地说出。因为劳苦的是那制造自己言语的人；愿意说真理的人不劳苦，因为真理自己说话毫无劳苦。……\par

13. \Quote{炭火必落在他们身上，祢要将他们抛在地上}（\BibleRef{咏 139:10}）。\Quote{在地上}是什么意思？在这里，甚至在此生，这里\Quote{炭火要落在他们身上。}\Quote{炭火}是什么？我们知道这些炭。它们与我们将要说的那些不同吗？因为我看这些用于惩罚，那些我将要说的用于救恩。因为我们曾谈到某些炭，当人寻求对抗诡诈的舌头时。……\Quote{炭}的榜样被加到箭的伤口上（因为我不怕说\Quote{伤口}，既然新娘自己说\Quote{我被爱所伤}），然后干草被烧尽，所以它们被称为\Quote{吞噬的炭}。干草被吞噬，但金子被炼净，人用死亡换取生命，并开始自己也成为燃烧的炭；如同宗徒那样的炭，\Quote{他先前是亵渎者、迫害者和施暴者}，是又黑又灭的炭；但当他获得了怜悯，就从天上被点燃，基督的声音点燃了他，他里面的一切黑色都消失了，他开始在神火中炽热，用他被点燃的火点燃别人。……\par

14. \BibleQuote{一个多言的人，在地上必不得引导} \BibleRef{咏 139:11} 。\Quote{一个多言的人}喜爱谎言。因为他除了说话还有什么快乐呢？他不关心自己说什么，只要他说就行。这样的人不可能被引导。那么，被这些\Quote{炭火}点燃、且自己也成为救恩之炭的天主之仆该做什么呢？他应该宁愿听而不愿说；正如经上所载：\BibleQuote{各人要快快地听，慢慢地说} \BibleRef{雅 1:19} 。如果可以，他该渴望这样：不是永远不必说话，而是只在必须履行职责时说话；让他乐于沉默，在意志上，当他必须说话时，为教导而说。因为何时你必须说话和教导？当你遇到一个无知的人，一个无学问的人时。如果教导常令你喜悦，你就希望总有可教导的无知者……\BibleQuote{邪恶必追捕不义的人，直到毁灭} \BibleRef{咏 139:11} 。邪恶来临，他站立不住；因此说\Quote{他们必追捕他直到毁灭}。因为许多好人、义人也遭遇了邪恶，邪恶似乎找到了他们。因此，当邪恶追捕善人时，即我们的殉道者，当他们被捉住时，他们被\Quote{追捕}，但不是\Quote{直到毁灭}。因为肉身被压制，精神却得了冠冕；精神被逐出身体，但对肉身所行的一切，都不能阻碍它将来的复活。让肉身被焚烧、鞭打、撕裂吧；难道它因为被交在迫害者手中，就离开了它的创造主吗？那从无中创造了它的主，岂不会把它再造得比以前更好吗？\newline{}\par

15. \BibleQuote{我知道上主要为穷苦人伸冤} \BibleRef{咏 139:12} 。这个\Quote{穷苦人}不是\Quote{多言的人}；因为多言的人想要富足，不知道饥饿。\Quote{穷苦人}是指经上所说的：\BibleQuote{饥渴慕义的人是有福的，因为他们要得饱饫} \BibleRef{玛 5:6} 。他们在恶人的绊脚石中叹息，向他们的元首祈祷，求从恶人手中被救出。\Quote{为贫苦人的案件}。这些人的案件，主必不忽略；虽然他们现在受苦，但当他们的元首显现时，他们的荣耀也必显现。因为经上对仍在现世的人说：\BibleQuote{你们已经死了，你们的生命与基督一同藏在天主内} \BibleRef{哥 3:3} 。这样，我们是穷苦的，我们的生命是隐藏的；让我们向那作为我们食粮的主呼喊 \BibleRef{若 6:51} ……\newline{}\par

16. \BibleQuote{但义人必称谢祢的圣名} \BibleRef{咏 139:13} 。无论祢何时为他们辩护，何时为他们伸冤，他们\Quote{必称谢祢的圣名}；决不将任何事归功于自己的功劳，一切都只归功于祢的仁慈……因此请看下文，看作者以什么作为结束。\Quote{正直的人必享见祢的圣容}。因为他们在自己的面容上本是恶劣的；在祢的圣容中，他们将得福乐。因为他们爱慕自己的面容时，\BibleQuote{必须汗流满面，才有饭吃} \BibleRef{创 3:19} 。祢的圣容将带着丰盛而来，使他们满足。他们不再寻求别的，因为没有更好的；他们不再离弃祢，也不会被祢离弃。因为论及主的复活，经上怎么说？\Quote{祢必将喜乐充满我的面容}。没有祢的圣容，祂不会赐给我们喜乐。我们正是为此洁净自己的面容，好能在祢的圣容中喜乐 \BibleRef{若一 3:2} ……因为，\BibleQuote{心里洁净的人是有福的，因为他们要看见天主} \BibleRef{玛 5:8}；祂将人形赐给了善人与恶人，却将天主的形象保留给了纯洁和善良的人，好使我们能在祂内喜乐，并因祢的圣容而永远幸福。\newline{}\par

\SourceRecord{Translated by J.E. Tweed.；Nicene and Post-Nicene Fathers, First Series；Vol. 8.；Edited by Philip Schaff.；Buffalo, NY: Christian Literature Publishing Co.,；1888.；<http://www.newadvent.org/fathers/1801140.htm>.}

% structure: p0001 source=h1 target=section reason=page-title

\section{圣咏第一百四十一篇释义}

1. ……我们刚才咏唱的这篇圣咏，在许多部分有些隐晦。当藉着主的帮助，开始解释和说明所说的内容时，你们将看到你们在听你们已经知道的事。但正因如此，这些事以多种方式被说出，以便表达的多样性可以消除对真理的一切厌倦……\par

2.  \BibleQuote{上主，我向祢呼求，求祢听我} \BibleRef{咏 140:1}。这一点我们都能说。这不单是我说：整个基督都说。但这话更以身体的名义说出：因为祂，当祂在此世并带着我们的肉身时，也曾祈祷；当祂祈祷时，血滴从祂整个身体流下。福音中正是如此记载： \BibleQuote{耶稣恳切祈祷，祂的汗珠如同大血滴。} \BibleRef{路 22:44} 这从祂全身流出的汗珠是什么，岂非全体教会中殉道者的苦难？ \Quote{求祢倾听我祷告的呼声，当我向祢呼号时。} 你以为当你说了 \Quote{我向祢呼求} 时，呼号的工作已经完成了吗？你已经呼号了，但不要自以为安全。如果苦难结束了，呼号就结束了；但如果苦难为了教会、为了基督的身体，直到世界终结仍然存在，那么它不只应说 \Quote{我向祢呼求}，也应说 \Quote{求祢倾听我祷告的呼声}。\par

3.  \BibleQuote{愿我的祈祷如馨香上达于祢面前，愿我的双手高举，如献晚祭} \BibleRef{咏 140:2}。每个基督徒都承认，这通常被理解为指头本身。因为当日已向晚低垂时，主在十字架上 \BibleQuote{舍了祂的生命，为再取回} \BibleRef{若 10:17}，并非违背祂的意愿而失去。但我们也在那里被预表。因为挂在树上的祂，除了取自我们者外，还有什么？圣父怎么可能离开并舍弃祂的独生子，尤其是当祂与祂同是一位天主时？然而，将我们的软弱钉在十字架上，在那里，如宗徒所说， \BibleQuote{我们的旧人已与祂同钉十字架} \BibleRef{罗 6:6}，祂以我们那 \Quote{旧人} 的声音呼喊说： \Quote{祢为什么舍弃了我？} 那就是 \Quote{晚祭}，主的苦难，主的十字架，救恩祭献的奉献，蒙天主悦纳的全燔祭。那 \Quote{晚祭} 在祂的复活中产生了一个清晨的奉献。所以，从忠信心中纯正发出的祈祷，如同馨香从神圣的祭台上升起。没有什么比主的馨香更令人愉悦：愿所有相信的人都拥有这样的馨香。\par

4. …… \BibleQuote{上主，求祢在我的口前设置警卫，在我唇边围上一道门} \BibleRef{咏 140:3}。祂没说 \Quote{一道栅栏}，而是说 \Quote{一道门}。门既可开也可关。如果它是一道 \Quote{门}，就让它既能开也能关：开，为认罪；关，为辩解罪。这样它就成为 \Quote{一道门}，而不是毁灭之门。因为这道 \Quote{门} 对我们有什么益处？基督以祂身体的名义祈祷什么？ \BibleQuote{求祢不要让我的心偏向恶言} \BibleRef{咏 140:4}。什么是 \Quote{我的心}？我教会的心，即我身体的心……\par

5. 但是，当你，基督的肢体，你的心没有偏向 \BibleQuote{恶言，以罪过为借口，与作恶的人同流合污} 时，你也不会与他们的选民联合。因为接下来是： \Quote{我不与他们的选民联合。} 谁是 \Quote{他们的选民}？那些自以为义的人。谁是他们的选民？那些 \Quote{自以为义，而轻视别人} 的人，如同法利塞人在圣殿中所说： \BibleQuote{天主，我感谢祢，因为我不像其他人一样。} \BibleRef{路 18:11} 谁是他们的选民？ \BibleQuote{这人若是先知，必知道摸祂的是怎样的女人。} \BibleRef{路 7:39} 你在此认出那另一个法利塞人的话，他曾邀请主到他家；当那城中的女人，一个罪妇，来靠近祂的脚时……\par

因为就连这女人自己，如果她的心 \Quote{偏向恶言}，她也不会缺少为她自己的罪辩护的言辞。难道不是每天都有与她同样污秽、但不像她那样告明的妇女——娼妓、奸妇、行可耻之事者——为她们的罪辩护吗？如果她们没有被看见，她们就否认；如果被抓住并定罪，或公开行了恶事，她们就辩护。她们的辩护何等容易，何等现成，却又何等轻率；何等常见，却又何等亵渎！ \Quote{如果不是天主愿意，我不会做：是天主愿意的；是命运愿意的；是命运注定的。} ……这些是这世界 \Quote{选民} 的辩护。但让基督的肢体、基督的身体说，让基督以祂身体的名义说： \BibleQuote{求祢不要使我的心偏向恶言} 等等， \Quote{我不与他们的选民联合。} ……\par

6. \Quote{与行恶的人一同。} 是什么恶？让我提一些他们犯罪的恶行。让我告诉你一件他们公开承认的罪恶之恶。他们说：人最好是做盘剥者，而不是农夫。你问理由，他们便给出一个理由....他折磨基督的肢体，就是那用犁铧洁净大地的人；他折磨基督的肢体，就是那从地上拔草的人；他折磨基督的肢体，就是那从树上摘苹果的人。为了避免在农场中犯下他们想象中的谋杀，他们却在盘剥中犯下真实的谋杀。他不给穷人食物。请看，还有什么不义比这种义更大呢？他不给饥饿的人食物。你问，为什么？免得乞丐领受那在面包中的生命（他们称之为天主的肢体、天主的实体），并将其束缚在肉里。那么你在做什么？你为什么吃？你没有肉吗？有；但是他们，他们说，既然我们因信仰\Person{摩尼}而得光照，因我们的祈祷和我们的圣咏，既然我们是选民，我们便借此洁净那面包，并将其送入天上的宝库。这样的选民，他们不是要被天主拯救，而是天主的拯救者。他们说，这就是基督，被钉在普世之中。我在福音中领受了拯救者基督，但你们在你们的书卷中却是基督的拯救者。显然，你们是基督的亵渎者，因此不能由基督拯救。因此，为了避免给饥饿的人一点碎屑，并且在碎屑中基督的肢体受折磨，饥饿的人就要饿死吗？对碎屑的虚假怜悯导致了对人的真正谋杀。但谁是他们的选民？\Quote{不要使我的心偏向恶言，我不愿与他们的选民联合。}\par

7. \BibleQuote{正义者必在仁慈中修正我，并责斥我} \BibleRef{咏 140:5}。看，罪人在忏悔。他希望在仁慈中被修正，而不是在欺骗中被赞美....\BibleQuote{必责斥我}，但是\BibleQuote{在仁慈中}：责斥，却并不憎恨：是的，他更会责斥，因为他不憎恨。那么他为何因此感恩？因为\BibleQuote{责斥明智人，他必爱你。} \BibleRef{箴 9:8} \BibleQuote{正义者必修正我。} 因为他逼迫你吗？千万不要。反而是在憎恨中修正的人，更需要修正自己。那么他为何修正？\BibleQuote{在仁慈中。并要责斥我。} 在何事上？\BibleQuote{在仁慈中。因为罪人的油必不润泽我的头。} 我的头必不因谄媚而增长。过度的赞美是谄媚：谄媚者的过度赞美就是\BibleQuote{罪人的油}。因此，当人们以虚假的赞美嘲弄某人时，他们就说：\Quote{我膏了他的头。} 那么，你们要喜爱\Quote{被正义者在仁慈中责斥}；不要喜爱被罪人在嘲弄中赞美。你们自己要有油，你们就不会寻求\BibleQuote{罪人的油。}..\par

8. 你对我说：我在做什么？我被谄媚者包围；他们不停地围攻我；他们赞扬我所不愿的，他们赞扬我所不重视的；我珍视的，他们却加以责备；这些谄媚者，背信弃义，虚伪欺诈。例如：\Quote{\Person{Gaiuseius}是个大人物，伟大、博学、智慧；但他为什么是基督徒？因为他的学问伟大，他的阅读广泛，他的智慧高超。}如果他的智慧高超，那就应该赞同他是基督徒；如果他的学问伟大，那他是有学问地选择了。总之，你所责备的，正是你所赞扬的那个人所喜欢的。但那又如何？那赞美并不甘美：它是\Quote{罪人的油}。然而他仍不停地这样说。不可借此\Quote{养肥你的头}；也就是说，不要因这些事而喜乐；不要同意这些事；不要认同这些事；不要因这些事而喜乐；这样，即使他向你涂抹了谄媚的油，你的头仍然保持原样，没有膨胀，没有肿胀……\Quote{因为我的话仍将令他们悦纳。}再等片刻：现在他们亵渎我，基督说。在基督徒的早期，基督徒处处受到责备。还要等待；\Quote{我的话将令他们悦纳。}时候将到，他们将征服成千上万的人，这些人会捶胸说：\Quote{宽恕我们的罪债，如同我们宽恕得罪我们的人。}即使现在，还有多少人以捶胸为耻？那就让他们责备我们吧，让我们忍受。让他们责备；让他们仇恨、控告、诽谤；\Quote{我的话仍将令他们悦纳；}时候将到，我的话将令他们喜悦……哦，不义的雄辩辩护！确实，现在整个民族都在说这话，无数人捶胸的雷声不断。云中如今住着天主，因此云正确地发出雷声。现在那饶舌、那吹嘘\Quote{我是正义的，我什么恶事也没有做}在哪里？确实，当你在圣经中默观正义的律法时，无论你进步多远，你都会发现自己是个罪人……我现在说的是什么样的人呢，弟兄们？我说的是那唯独敬拜天主、宣认基督、认识圣父、圣子和圣神是独一天主的人；他不向他犯邪淫；他不敬拜魔鬼；他不向魔鬼寻求帮助；他持守天主教会；没有人控告他欺诈；没有软弱的邻舍因他的压迫而呻吟；他不觊觎他人的妻子；他满足于自己的妻子，或甚至没有妻子，只要合法，并符合宗徒的纪律，经双方同意，\BibleRef{格前 7:5}或女方尚未结婚。即使是这样的人，也会被我提到的这些事所绊倒。对于所有这些日常的罪，我们的希望除了在天主经中以谦卑的心说——我们不辩护自己的罪，而是承认它们——\Quote{宽恕我们的罪债，如同我们宽恕得罪我们的人}，\BibleRef{玛 6:12}以及\Quote{在父那里有一位护慰者，正义的耶稣基督}，使祂成为\Quote{我们罪的赎罪祭}，\BibleRef{若一 2:1}还能是什么呢？看接下来的经文：\Quote{他们的审判者被吞没在磐石旁。}\BibleRef{咏 140:6} 什么是\Quote{被吞没在磐石旁}？\Quote{那磐石就是基督。}\BibleRef{格前 10:4} \Quote{他们被吞没在磐石旁。}\Quote{旁}，即比较，作为审判者，作为有权力、有学问、有智慧的人；他们被称为\Quote{他们的审判者}，因为他们判断道德并发表自己的意见。这是\Person{亚里士多德}说的。将他置于磐石旁，他就被吞没了。\Person{亚里士多德}是谁？让他听，\Quote{基督说过，}他就在死者中颤抖。这是\Person{毕达哥拉斯}说的，那是\Person{柏拉图}说的。将他们置于磐石旁，将他们的权威与福音的权威比较，将骄傲的人与被钉十字架者比较。让我们对他们说：\Quote{你们在骄傲的人心中写下你们的话；祂在君王心中竖起祂的十字架；最后，祂死了，又复活了；你们是死的，我不问你们如何复活。}因此\Quote{他们的审判者被吞没在}那\Quote{磐石}旁。因此，他们的话似乎有些价值，直到它们与磐石相比较。所以，如果其中有人被发现说了基督也说过的话，我们祝贺他，但我们不跟随他。但他先于基督而来。如果有人说了真理，难道他就先于真理本身吗？看基督，人啊，不是当他来到你这里的时候，而是当他创造你的时候。病人也可能说：\Quote{但是我在医生来之前就卧病在床了。}为什么？正是因为他先病了，所以祂最后才来。\newline{}\par

9. \BibleQuote{他们要听我的话，因为他们已经得胜。}我的话胜过了他们的话。他们说得巧妙，我说真理。称赞能言善辩是一回事，称赞讲真理的是另一回事。\BibleQuote{他们要听我的话，因为他们已经得胜。}他们如何得胜？他们中有谁在献祭时被拿住——法律禁止这样的事——而没有否认？他们中有谁在拜偶像时被拿住，而没有喊叫\Quote{我没有做}，并且害怕被定罪？魔鬼有这类的仆人。但主的话如何得胜？\BibleQuote{看，我派遣你们好像羊进入狼群。不要害怕那些杀害身体的，}等等。祂给了他们畏惧，祂给了他们希望，祂点燃了爱。\BibleQuote{你们不要害怕死亡，}祂说。你们害怕死亡吗？我先死。你们害怕连一根头发也不受损吗？我先在肉身中毫无损伤地复活。你们确实听了祂的话，因为它们得胜了。他们说话，就被杀害；他们倒下，却又站立。这么多的殉道者之死带来了什么结果？岂非那些话得胜了，大地——可说是被基督见证人的血所浇灌——教会的十字架到处生长？他们如何\Quote{得胜}？我们已经说过，当这些话由不怕的人宣讲时。不怕什么？不怕放逐、损失、死亡、甚至钉十字架：他们不仅不怕死亡，甚至不怕那种被认为是最可咒骂的死——钉十字架。主承受了它，为叫祂的门徒不但不怕死，而且连那种死也不怕。那么，当这些事由不怕的人说出时，它们就得胜了。\par

10. 那么，那些殉道者的所有死亡成就了什么？请听：\BibleQuote{如同肥油铺满地面，我们的骨头被撒在阴间旁边} \BibleRef{咏140:7} 。\Quote{骨头}指殉道者，即基督见证人的身体。殉道者被杀害，杀害他们的人似乎得胜了。他们藉着迫害得胜，好让基督的话藉着宣讲得胜。圣人们死亡的结果是什么？\Quote{肥油铺满地面}是什么意思？我们知道，一切废物都是地上的肥油。那些在人眼中可鄙的东西，却使地肥沃。……\BibleQuote{主视祂圣者的死为珍贵。}如同它被世界轻视，所以它对农夫是珍贵的。因为他知道它的用途和丰富的汁液；他知道他渴望什么，寻求什么，肥沃的庄稼从何而来；但这世界轻视它。你们岂不知道\BibleQuote{天主拣选了世上可鄙的，和那不存在的，如同那存在的，为要使那存在的归于虚无}？\BibleRef{格前1:27} 伯多禄和保禄是从粪堆中被提拔起来的；当他们被处死时，他们被轻视：现在，大地因他们而肥沃，教会的十字架生长起来，看，世界上所有尊贵和首要的人，甚至皇帝本人，来到罗马，他往哪里赶？去皇帝的庙宇，还是渔夫的纪念堂？\par

11. \BibleQuote{因为我的眼目仰望祢，主；我寄望于祢，求祢不要夺去我的生命} \BibleRef{咏140:8} 。因为他们在迫害中受折磨，许多人跌倒了。他想起了许多人跌倒，许多人处于危险中，在迫害的苦难中，仿佛发出了祈祷的声音：\BibleQuote{因为我的眼目仰望祢，主：}我不在乎周围人威胁什么，\BibleQuote{我的眼目仰望祢，主。}我更注视祢的应许，胜过他们的威胁。我知道祢为我受了什么，祢给我许诺了什么。\par

12. \BibleQuote{求祢救我脱离他们为我设下的罗网} \BibleRef{咏140:9} 。罗网是什么？\Quote{你若同意，我就饶你一命。}在罗网中放置了现世生活的诱饵；如果飞鸟喜爱这饵，就落入网罗；但如果飞鸟能说：\BibleQuote{我不贪求人的日子，祢知道：} \BibleRef{耶17:16} \BibleQuote{他必从网中拔出脚来，}等等。他提到了两件事，要彼此区别：他说罗网是由迫害者设下的；绊脚石来自那些已经同意并背教的人；他愿从两者中被保护。一方面他们威胁发怒，另一方面他们同意跌倒：我害怕前者使我恐惧，后者使我效仿。\Quote{你若不同意，我就这样待你。}\BibleQuote{求祢救我脱离罗网，}等等。\Quote{看，你的兄弟已经同意了。}\Quote{也脱离绊脚石，}等等。\par

13. \Quote{罪人必落入他的网罗} \BibleRef{咏 140:10}。不是所有罪人，而是某些罪人，他们如此之大，竟至爱此生命到宁愿它超过永生的地步，\Quote{必落入他的陷阱}。但你说什么？你以为这样的人会落入他的网罗吗？你的门徒呢，哦基督？看，当迫害猖獗时，当他们全都\Quote{离开你独自一人，各归自己的地方}时：\BibleRef{若 16:32}。看！那些最亲近你的人，在你的考验和迫害中，当你的敌人要求把你钉在十字架上时，他们遗弃了你。还有那个勇敢的人，他许诺要与你同死，却从医生那里听到了病人身上正在发生的事。因为他发烧时，他说自己健康；但主触摸了他心灵的脉搏。然后考验来了；然后审判来了；然后控告来了；而现在，不是被什么大人物质问，而是被一个卑微的奴隶，且是一个女人，被一个婢女质问，他屈服了；他三次否认……经上说：\Quote{他痛哭流涕}。他尚未适宜受苦。他曾被告知：\Quote{你日后要跟随我} \BibleRef{若 13:36}。此后他将坚固，因主的复活而坚强。那时还不是那些\Quote{骨头}被\Quote{散在坑旁}的时候。因为看哪，有多少人失败了，甚至那些最初挂在他嘴边的人；他们也失败了。为何？\Quote{我独自一人，直到我逾越}：因为这是诗篇接着说的……\par

14. \OriginalTerm{逾越节}{Pascha}，正如那些知道并解释给我们听的人所说的，意思是\Quote{逾越}。当主的苦难即将来临时，福音作者仿佛要用这个词，说道：\Quote{耶稣知道时候到了，他要\Emph{逾越}到父那里去} \BibleRef{若 13:1}。因此我们在这节诗中听到逾越节，\Quote{我独自一人，直到我逾越}。逾越节后我不再独自一人，过去之后我不再独自一人。许多人将效法我，许多人将跟随我。如果后来他们要跟随，现在情况如何？\Quote{我独自一人，直到我逾越}。主在这篇诗篇中说的\Quote{我独自一人，直到我逾越}是什么意思？我们解释了什么？如果我们明白了，就听他在福音中的话：他说：\BibleQuote{一粒麦子如果不落在地里死了，它仍旧是一粒；如果死了，就结出许多子粒来} \BibleRef{若 12:24}。……因此，他在受死之前是独自一人。……在那一粒麦子落在地里之前，竟无人为这名而死，即为承认基督之名而死，甚至就在他之前被杀的若翰，被一个邪恶的君王交给一个跳舞的女子，也不是因为承认基督而被杀。当然他可能因此而被杀，且被许多人杀。如果因其他原因被一个人杀，更何况被那些杀害基督的人杀呢？因为若翰为基督作证。那些听见基督的人想要杀他；为他作证的人他们却没有杀。……若翰自由地为基督作证，而犹太人杀了基督，却没有杀若翰；他被黑落德所杀，因为他说：\Quote{你占有你兄弟的妻子是不合法的} \BibleRef{玛 14:4}。因为他的兄弟并非无嗣而死。为了真理的法律，为了公平，为了正义，他死了；因此他是圣人，因此是殉道者；但他并非为我们成为基督徒的那名而死，这是为何？除非是为了应验这句话：\Quote{我独自一人，直到我逾越}。\par

\SourceRecord{Translated by J.E. Tweed.；Nicene and Post-Nicene Fathers, First Series；Vol. 8.；Edited by Philip Schaff.；Buffalo, NY: Christian Literature Publishing Co.,；1888.；<http://www.newadvent.org/fathers/1801141.htm>.}

% structure: p0001 source=h1 target=section reason=page-title

\section{《圣咏第一四二篇释义》}

1. ……\Quote{我以我的声音呼号了上主}\newline{}\BibleRef{咏 141:1}。说\Quote{以声音}已足够；也许加上\Quote{我的}并非无因。因为许多人呼号上主，并非用他们自己的声音，而是用他们身体的声音。那么，让那\Quote{内在的人}——基督已开始\Quote{因信德住在其中}\newline{}\BibleRef{弗 3:17}——呼号上主，不是用嘴唇的嘈杂，而是用内心的情感。天主听不见人听见之处：除非你用肺腑和舌头的呼声呼喊，人听不见你；你的思念，就是你对上主的呼号。\Quote{我以我的声音祈求了上主。}他用\Quote{我呼号}所表达的，用\Quote{我祈求}加以解释。因为连那些亵渎的人也呼号上主。在前段他写下了他的呼号，在后段他解释了那是什么。仿佛被问：你以怎样的呼号呼号了上主？他说：我祈求了上主。我的呼号是我的祈祷，不是辱骂，不是怨言，不是亵渎。\par

2. \Quote{我将在祂面前倾吐我的祈祷}\newline{}\BibleRef{咏 141:2}。什么是\Quote{在祂面前}？在祂的视线内。什么是\Quote{在祂的视线内}？在祂看见的地方。但祂哪里看不见呢？因为我们说\Quote{在祂看见的地方}，似乎祂在某处看不见。但在这一堆物质实体中，人也看见，动物也看见：祂看见人看不见之处。因为无人看见你的思想，但天主看见。那么，在那里倾吐你的祈祷，在祂独自看见、祂独自赏报之处。因为主耶稣基督命令你暗中祈祷：但如果你知道什么是\Quote{你的内室}，并洁净它，你就在那里向天主祈祷。\Quote{至于你，}祂说，\Quote{当你祈祷时，要进入你的内室，关上门，向你在暗中的父祈祷；那在暗中看见的父必将报答你。}\newline{}\BibleRef{玛 6:6}。如果人将报答你，就在人前倾吐你的祈祷；如果天主将报答你，就在祂面前倾吐你的祈祷；并关上门，免得诱惑者进入。所以宗徒——因为关上门，即我们心灵的门，而非墙壁的门，在于我们，我们的\Quote{内室}就在其中——因为关这门在于我们，他说：\Quote{不要给魔鬼留有余地。}\newline{}\BibleRef{弗 4:27}。但什么是\Quote{关上门}？这门像有两扇：欲望和恐惧。要么你渴求某件地上的事物，他从这一扇进入；要么你恐惧某件地上的事物，他从那一扇进入。那么，对魔鬼关上恐惧和欲望的门，对基督敞开它。你如何对基督敞开这扇对开门？通过渴求天国，通过恐惧地狱之火。通过此世的欲望魔鬼进入，通过永生的欲望基督进入；通过暂时的惩罚的恐惧魔鬼进入，通过永恒之火的恐惧基督进入……\par

3. \Quote{我要在祂面前述说我的困苦。}这里有一个重复，无论是在前面的两个句子中，还是在后面的这些句子中：两个思想，但两者都表达了两次……因为\Quote{在祂面前}等同于\Quote{在祂眼前}；\Quote{我要述说我的困苦}等同于\Quote{我要倾吐我的祈祷}。你何时这样做？他说：在迫害之中时，\Quote{当我的精神在我内衰弱时}\newline{}\BibleRef{咏 141:3}。为何你的精神衰弱了，殉道者，困苦中的你？为使我不会把我的力量据为己有，为使我认识另一位在我内成就我所拥有的善。也许人们听到了我的精神在我内衰弱，就对我绝望，并说：\Quote{我们俘虏了他，我们制胜了他}；\Quote{而祢却认识了我的路径。}他们以为我被推倒，祢见我站立正直。那些迫害我并已捉住我的人以为我的脚缠住了，\Quote{但他们的脚却缠住了，他们跌倒了，我们却站起来，直立着。}因为我的眼睛常仰望上主，因为祂必将我的脚从网罗中拔出。我坚持行走，因为\Quote{那坚持到底的，必然得救。}\newline{}\BibleRef{玛 10:22}。他们以为我被制伏，但我继续行走。我行了在哪里？在他们看不见的路径上，那些以为我被俘的人，在祢正义的路径上，在祢诫命的路径上……因为每条路径都是一条路，但并非每条路都是路径。为何那些路被称为路径，除非因为它们狭窄？恶人的路是宽敞的，义人的路是狭窄的。那\Quote{路}也是\Quote{诸路}，正如\Quote{教会}也是\Quote{诸教会}，\Quote{天}也是\Quote{诸天}：它们被称为复数，也被称为单数。由于教会的合一，它是一个教会：\Quote{我的鸽子是唯一的，是她母亲的独一者。}\newline{}\BibleRef{歌 6:8}。由于弟兄们在各地的集会，有许多教会。保禄说：\Quote{犹太在基督内的各教会都欢庆，并在我内光荣天主。}\newline{}\BibleRef{迦 1:22}。他这样说到诸教会；关于一个教会，他这样说：\Quote{你们不可成为犹太人，或外邦人，或天主的教会跌倒的原因。}……\par

4. \Quote{在我所走的这条路上，他们为我设了陷阱。}这条\Quote{我所走的路，}就是基督；他们为基督之名迫害我，就在基督里为我设了陷阱。他们\Quote{为我藏了陷阱。}他们恨我什么？他们迫害我什么？就是我是基督徒……因为异端者也想在基督的名里为我们隐藏绊脚石，而他们自己却被欺骗了。他们以为放在路上的东西，却放在了路外，因为他们自己就在路外。他们不能在自身不在的地方设陷阱……外教人想在路上设绊脚石，当他对我说：\Quote{你崇拜一位被钉在十字架上的天主。}他指责基督的十字架，但他并不理解。他以为他在基督里放置了什么，其实他放在路旁。我不离开基督，这样我就不会从路上掉进陷阱。让他嘲笑被钉的基督吧，让我看到基督的十字架在君王们的额头上。他所嘲笑的，正是我从中得救的。没有什么比一个病人嘲笑自己的药更骄傲的了。如果他不嘲笑它，他就会服用它而得痊愈。十字架是谦卑的标志，但他因极度的骄傲而不承认那能治愈他灵魂肿胀的十字架。然而，如果我承认，我就在路上行走。我远不以十字架为耻，以至于我不在暗处保存基督的十字架，而是将它戴在额头上。我们领受许多圣事，有的以一种方式，有的以另一种方式：有些你们知道，我们用口领受，有些我们接受在全身上。但因为额头是羞耻的座位，那说过\Quote{凡在人前以我为耻的，我在我在天之父面前也必以他为耻}（\BibleRef{路 9:26}）的主，可以说，把那外邦人所嘲笑的羞辱本身，置于我们羞耻的部位。你听到有人攻击一个无耻的人说：\Quote{他没有额头。}什么是\Quote{他没有额头}？他没有羞耻。让我不要有光秃的额头，让我主的十字架遮盖它……\par

5. \Quote{我向右看，并看见了}（\BibleRef{咏 141:4}）。他向右看，并看见了：谁向左看，就被蒙蔽。什么是向右看？就是那些将要听到\Quote{你们这些蒙我父降福的，来吧！}（\BibleRef{玛 25:34}）的人……他继续说：\Quote{却没有人认识我。}因为当你惧怕万事时，谁知道你注视什么，是向右看还是向左看？如果你在忍受中寻求人的赞美，你就向左看了；如果你在忍受中寻求天主的应许，你就向右看了。你若向右看，就会看见；你若向左看，就会被蒙蔽。但即使你向右看见，也没有人认识你。因为除了主，谁能安慰你呢？\Quote{逃避已远离了我。}他说话时好像被围困。让迫害者因他而喜乐吧；他被压倒，被捉拿，被围困，被征服。\Quote{逃避已远离了}那不逃跑的人。但那不逃跑的人，他为基督承受一切所能承受的：也就是说，他灵魂不逃跑。因为在身体上逃跑是合法的，是被允许的，是许可的；因为主说：\Quote{有人在这城迫害你们，你们就逃到另一城去}（\BibleRef{玛 10:23}）。所以那灵魂不逃跑的人，从他那里\Quote{逃避已远离了。}但重要的是他为什么不逃跑：是因为被围困、被捉拿，还是因为他勇敢。因为对于被捉拿的人来说，逃避已远离了；对于勇敢的人来说，逃避也已远离了。那么，哪种逃避应当避免？哪种逃避我们要让它远离我们？就是主在福音中所说的：\Quote{善牧为羊舍掉自己的性命。但雇工，不是牧人，看见狼来，就逃跑了。}……人的生命有两种方式被寻找：要么被他的迫害者，要么被他的爱慕者。所以他说\Quote{没有人寻求我的生命，}他说的是他们；他们确实迫害我的生命，却不寻求我的生命。但如果他们寻求我的生命，他们会发现它依附于你；如果他们知道如何寻求它，他们也知道如何效法它。\par

6. \Quote{上主，我向祢呼号：我说：祢是我的希望}（\BibleRef{咏 141:5}）。当我忍受时，当我受困苦时，\Quote{我说：祢是我的希望。}我的希望在此世，所以我忍受。但\Quote{我的福分}不在此世，而是在\Quote{生活之地}。天主在生活之地赐予福分；但不是从祂自身之外赐予什么。祂将赐给那爱祂的人什么呢？除了祂自己。\par

7. \BibleQuote{求祢垂听我的祈祷，因为我已受尽卑微} \BibleRef{咏 141:6}。迫害者使他卑微，告解使他卑微。他在人眼前自卑，在仇敌眼前却被卑微。因此，那一位在可见与不可见中提升他。不可见中，殉道者已被提升；可见中，他们将在死人复活时被提升，\Quote{当这可朽坏的穿上不朽时}；那时，迫害者所能攻击的这部分将被更新。\BibleQuote{不要害怕那些杀害身体，却不能杀害灵魂的。} \BibleRef{玛 10:28} 什么会丧亡？他们杀的是什么？……那么，你为什么还担心你其余的肢体，连一根头发也不会失落呢？\Quote{求祢从迫害我的人中拯救我。} 你认为他求从谁那里获得拯救？是从迫害他的人吗？是这样吗？难道只有人是我们的仇敌吗？我们还有其他无形的仇敌，以另一种方式迫害我们。人迫害，是为了杀害身体；另一种迫害，是为了诱惑灵魂。\BibleRef{弗 2:2} ……我们也有其他仇敌，我们应当祈求天主从他们中拯救我们，免得他们藉着世俗的苦难压垮我们，或藉着其诱惑引诱我们。这些仇敌是谁？让我们看看，是否有主的仆人、现已成全的战士曾与它们交战，清楚描述过它们。请听宗徒说：\BibleQuote{我们战斗不是对抗血肉}，\BibleRef{弗 6:12} 好像要说：不要将仇恨转向人，不要认为他们是你的仇敌，不要以为你是因他们的敌意而受折磨；你所惧怕的人只是血肉……\Quote{因为他们在我身上得了势。}谁说了\Quote{他们在我身上得了势}？基督的身体在呼号，那是教会的声音；基督的肢体在呼号：\Quote{罪人的数目大增。}\BibleQuote{因为罪恶增加，许多人的爱德冷淡。}\BibleRef{玛 24:12}\par

8. \BibleQuote{求祢从狱中引领我的灵魂出来，好使它向祢的圣名宣认} \BibleRef{咏 141:7}。这个\Quote{狱}被以前的作者有不同的理解。也许它就是标题中称为\Quote{山洞}的那个。因为这篇圣咏的标题是这样写的：\Quote{达味的训诲诗，当他在山洞时的祈祷。}山洞也就是狱。我们面前有两件事要了解，但当我们了解了一件，两件就都了解了。人的功过造就了狱。因为在同一住处，一个人找到的是房屋，另一个人找到的是狱……因此，有些人认为\Quote{山洞}和\Quote{狱}就是这个世界；而教会祈祷，求被带出狱，即脱离这个世界，脱离太阳之下，那里尽是虚幻。超越这个世界，天主应许我们将在某种安息中；因此，也许我们为此呼喊：\Quote{求祢从狱中引领我的灵魂出来。}我们的灵魂因信德和望德而在基督内；\BibleQuote{你们的生命与基督一同藏在天主内}。但我们的身体却在这个狱中，在这个世界里……但也有人说，这个狱和山洞就是身体，因此\Quote{求祢从狱中引领我的灵魂出来}的意思就是求灵魂离开身体。但这种解释也有问题。因为说\Quote{求祢从狱中引领我的灵魂出来}，从身体里出来，这算是什么大事呢？强盗和恶人的灵魂不也是离开身体，进入比在这里所受的更重的惩罚吗？那么，\Quote{求祢从狱中引领我的灵魂出来}又是什么了不起的祈求呢？因为迟早总要出来的。也许义人会说：\Quote{现在就让我死吧，从这身体的狱中引领我的灵魂出来。}如果他过于急切，就没有爱德。他确实应该渴慕与向往，正如宗徒所说：\BibleQuote{渴望离世与基督同在，这实在是好得无比的}。但爱德在哪里呢？因此接着又说：\BibleQuote{但留在肉身中，为你们是必要的。}那么，让天主按祂的旨意带领我们离开身体吧。我们的身体也可称为狱，不是因为天主造的身体是狱，而是因为它处于惩罚之下并可能死亡。因为在我们身体中要考虑两件事：天主的工程，和它应得的惩罚……也许他所说的\Quote{求祢从狱中引领我的灵魂出来}，意思是从朽坏中引领我的灵魂出来。如果我们这样理解，那并不是亵渎，而且意义一致。最后，弟兄们，我认为他的意思是：\Quote{求祢从狱中引领我的灵魂出来}，即从狭窄中领出来。因为对欢乐的人来说，即使是狱也是宽广的；对忧伤的人来说，田野也是狭窄的。因此他祈求被带出狭窄。因为虽然在希望中他有广阔，但在现实此刻他受拘束……不是身体重压灵魂，而是有朽坏的身体。因此，不是身体构成狱，而是朽坏。\Quote{求祢从狱中引领我的灵魂出来，好向祢的圣名称谢。}接下来的话似乎来自元首，我们的主耶稣基督。这些与昨天最后的话相同。如果你记得，昨天最后的话是：\Quote{我独自一人，直到我过去。}这里最后的话是什么？\Quote{义人必扶持我，直到祢赏报我。}\par

\SourceRecord{Translated by J.E. Tweed.；Nicene and Post-Nicene Fathers, First Series；Vol. 8.；Edited by Philip Schaff.；Buffalo, NY: Christian Literature Publishing Co.,；1888.；<http://www.newadvent.org/fathers/1801142.htm>.}

% structure: p0001 source=h1 target=section reason=page-title

\section{\WorkTitle{圣咏第一四三篇释义}}

1. \BibleQuote{达味所作，当他的儿子追赶他时。}我们从列王纪\EditorialNote{撒慕尔纪下第十五章}中知道此事发生了：……但我们必须在此认出另一个达味，祂的确是\Quote{手中有力者}，这正是达味之名的解释——即我们的主耶稣基督。因为那一切过去的事件，都是未来之事的预象。因此，让我们在这篇圣咏中寻求我们的主、救主耶稣基督，祂在祂的预言中预先宣告自己，并用很久以前所行之事预告了此时将要发生的事。因为祂亲自在先知们中预言了自己：祂是天主的圣言。先知们若不是充满圣言，就不可能说出这类事。他们充满基督，便宣告基督；他们在祂将要来临之前走在祂前面，而祂也没有离弃那些走在前面的人……\par

2. 那么，让我们的主说话；让基督与我们、整个基督说话。\BibleQuote{上主，求祢俯听我的祈祷，求祢侧耳听我的恳求}\BibleRef{咏 142:1} 。\Quote{听}和\Quote{侧耳倾听}是同一回事。这是重复，这是确认。\Quote{在祢的真理中俯听我，在祢的正义中。}当祂说\Quote{在祢的正义中}时，不可轻描淡写地对待。因为这是恩宠的赞美，免得我们任何人以为自己的正义是属于自己的。因为这是天主的正义，是天主赐给你们拥有的。关于那些夸耀自己的正义的人，宗徒说了什么？论及犹太人，他说：\BibleQuote{他们对天主有热心，但不是依照真知}\BibleRef{罗 10:2}。……你是乖戾的，因为你把做错的事归咎于天主，把做对的事归于自己；你将是对的，当你把做错的事归于自己，把做对的事归于天主……看，\Quote{在祢的正义中俯听我。}因为当我审视自己时，除了罪以外，我找不到任何属于自己的东西。\BibleQuote{我是在罪恶中怀孕的。}\par

3. \BibleQuote{求祢不要与祢的仆人进入审判}\BibleRef{咏 142:2}。谁愿意与祂进入审判呢？除非是那些\BibleQuote{不认识天主的正义，而想建立自己的正义的人？} \BibleQuote{为什么我们禁食，祢没有看见？为什么我们克苦己心，祢不理会？}\BibleRef{依 58:3} 好像他们会说：\Quote{我们做了祢所吩咐的事，为什么祢不把祢所应许的给我们呢？} 天主回答你们：我给你们去接受我所应许的；我已经给了你们，使你们能做那使你们能接受的事。最后，对这些骄傲的人，先知说：\BibleQuote{你们为什么与我争辩？你们都违背了我——上主说。}\BibleRef{耶 2:29} 为什么你们要与我进入审判，并数算你们自己的义行？……\BibleQuote{因为在祢面前，凡活着的人都不能称为义。} \Quote{凡活着的人；} 活着，即在这里，在肉身中活着，在死亡的期待中活着；生为人，从人领受生命；从亚当所出，一个活着的亚当；每一个这样活着的人，或许在他自己面前被称为义，但在祢面前却不能。如何在自己面前？因取悦自己，而不取悦祢。那么，我的上主天主，求祢不要与祢的仆进入审判。不管我在自己看来多么正直，祢从祢的宝库中拿出标准，祢用标准衡量我，我就发现自己是弯的。说得很好：\Quote{与祢的仆人。} 祢与祢的仆人、甚至与祢的朋友进入审判，是不相称的。\BibleRef{玛 5:40} ……至于宗徒们自己呢？……为使你们立刻明白，他们学会了我们祈祷的经文：天上的导师给了他们祈祷的模式。祂说：\BibleQuote{你们要这样祈祷}\BibleRef{玛 6:9}。在首先写下某些事之后，祂也把这句放在牧羊人的首领、同一羊群的牧者和召集者的主要成员的口中；甚至他们也学会了说：\BibleQuote{宽免我们的罪债。}\BibleRef{玛 6:12} 他们没有说：\Quote{感谢祢，祢已经宽免了我们的罪债，如同我们也宽免了我们的欠债人，} 而是：\Quote{宽免我们，如同我们宽免他人。} 但确实，当时信友们祈祷，宗徒们祈祷，因为这篇主祷文更是赐给信友们的。如果那些罪债仅指由圣洗所赦免的，那就更适合由慕道者来说：\Quote{宽免我们的罪债。} 那么就让宗徒们说，是的，让他们说：\Quote{宽免我们的罪债。} 当有人对他们说：\Quote{你们为什么这么说？你们的罪债是什么？} 让他们回答：\BibleQuote{因为在祢面前，凡活着的人都不能称为义。}\par

4. \BibleQuote{因为敌人迫害我的性命，将我的生命压制在地上}\BibleRef{咏 142:3}。这里我们发言，这里我们的元首为我们发言。显然，魔鬼迫害了基督的灵魂，犹大迫害了他导师的灵魂；现在，同样的魔鬼仍然继续迫害基督的身体，一个犹大接替另一个犹大。因此，身体也不缺少谁能说：\BibleQuote{因为敌人迫害了我的灵魂。} 每一个迫害我们的人，他所努力的，岂不是要使我们放弃天上的希望，而贪恋地上的事物，屈服于迫害者，并爱慕世俗之物吗？\BibleQuote{他们把我放在黑暗之处，如同世上的死人。} 你们更容易从元首那里听到这个；你们更容易在元首身上觉察这个。因为祂确实为我们死了，但祂并不是\BibleQuote{世上的死人}之一。谁是\BibleQuote{世上的死人}？祂怎么不是\BibleQuote{世上的死人}之一？\BibleQuote{世上的死人}是那些因自己的罪过而死，承受不义的报应，从遗传的罪中接受死亡的人，正如经上所说：\BibleQuote{我是在罪恶中怀孕的。} ……祂说：在死去时，我成就我父的旨意，但我并不该死。我没有做过任何该死的事，然而我死是我自己的作为，好使一个无罪者的死亡，使那些本来该死的人获得自由。\BibleQuote{他们把我安置在处所}，仿佛在阴府，仿佛在坟墓中，仿佛在祂的苦难本身，\BibleQuote{如同世上的死人。}\par

5. \BibleQuote{我的灵在我里面}，祂说，\BibleQuote{困倦不堪}\BibleRef{咏 142:4}。请记住：\BibleQuote{我的心灵忧闷得要死。}\BibleRef{玛 26:38}。在此我们看到同一个声音。我们岂不显然看到从元首到肢体、从肢体到元首的转变吗？……\par

6. 但我们也在那里。祂走向肢体。\BibleQuote{我追想古昔之日}\BibleRef{咏 142:5}。那创造每一天的祂，岂会\BibleQuote{追想古昔之日}？不，是那身体在说话，每一个因祂恩宠而成义、在爱中及虔诚的谦卑中住在祂内的人，在说话并说：\BibleQuote{我默思祢的一切作为：}显然因为祢造了万物皆善，凡非祢所立的，无一能存立。祢的受造物在我面前成了景观：我在作品中追寻工匠，在受造物中追寻造物主。为何如此？有何目的？无非是为使人明白，凡他内的善都是祂所造的。……那么，请回望你生命的塑造者、你本质的作者、你的正义与救恩的作者：\BibleQuote{默思祂手的作为}，因为那在你内的正义，你也会发现属于祂的手。请听教导你的宗徒说：\BibleQuote{不是出于功行，免得有人自夸。}我们岂没有善功？显然我们有：但请看下文：\BibleQuote{因为我们是祂的化工}\BibleRef{弗 2:9}，他说。\BibleQuote{我们是祂的化工：}或许他如此说\Quote{化工}，是指我们作为人的本性？显然不是：他是在说功行。但我们不要猜测；让经文继续：\BibleQuote{因为我们是祂的化工，在基督耶稣内受造，为行善功。}那么，不要以为你自己能做什么，除非在你邪恶的方面。……\BibleQuote{你们要怀着恐惧战栗，努力成就你们的救恩，}...\BibleQuote{因为是天主在你们内工作，使你们愿意，并使你们行事，为成就祂的美意。}\BibleRef{斐 2:12-13}。因此要\BibleQuote{怀着恐惧战栗}，好让我们的造物主乐于在谦卑的幽谷中工作。……\par

7. \BibleQuote{我向祢伸开双手；我的灵魂向祢有如干旱之地}\BibleRef{咏 142:6}。请降雨在我身上，好从我结出善果。\BibleQuote{上主必赐下甘霖，使我们的土地出产果实。}\BibleQuote{我向祢伸开双手；我的灵魂向祢有如干旱之地}，不是向我，而是\BibleQuote{向祢}。我能渴慕祢，却不能滋润自己。\par

8. \BibleQuote{上主，求祢速速俯听我}\BibleRef{咏 142:7}。既然我已如此急切地干渴，又何须延迟以激发我的渴呢？祢延迟雨水，是为使我饮、吸收而非拒绝祢的灌注。为此，祢若延迟，现在就赐下吧；因为\BibleQuote{我的精神已经衰颓。}愿祢的圣神充满我。这就是祢应速速俯听我的理由。我现在已成为\BibleQuote{神贫的人}，求祢使我\BibleQuote{在天国里有福}\BibleRef{玛 5:3}。因为那在自己精神内活着的人，是骄傲的，以自己的精神向天主自高。……\par

9. \BibleQuote{求祢不要转面不顾我。}我骄傲时，祢曾转面不顾我。因为我曾饱满，在我的饱满中我自高自大。我曾\BibleQuote{在我的饱满中说：我永不动摇。}\BibleQuote{我在饱满中说：我必不动摇，}不明祢的正义，而想建立自己的正义；但\BibleQuote{上主，祢以祢的旨意使我的华美有力。}\BibleQuote{我在饱满中说：我必不动摇，}但我的饱满全是来自祢。为向我证明它来自祢，\BibleQuote{祢转面不顾我，我便惊慌失措。}经此困扰——我被投入其中，因为祢转面不顾我——经过我的精神困倦，经过我的心在我内烦乱，因为祢转面不顾我，然后我成了\BibleQuote{向祢有如干旱之地：求祢不要转面不顾我。}祢在我骄傲时转面不顾我；如今我谦卑了，求祢将面容还给我。因为若祢转面不顾我，\BibleQuote{我就要与下到深坑的人同类。}\BibleQuote{下到深坑}是什么意思？当罪人陷入罪孽的深处，他必藐视。\BibleQuote{下到深坑}的那些人，甚至失去了告解；反此，有话说：\BibleQuote{不要让深坑在我上合口。}圣经常将这种深处称为\Quote{深坑}，罪人一旦陷入这深处，\Quote{便藐视。}\Quote{便藐视}是什么意思？他不再相信上智安排，或即便相信，也认为自己与它再无关系。……\par

10. \BibleQuote{求祢在早晨使我聆听祢的仁慈，因为我寄望于祢。}\BibleRef{咏 142:8} 看，我处于黑夜，却仍\Quote{寄望于祢，}直到黑夜的罪孽过去。\Quote{因为我们有，}正如伯多禄所说，\BibleQuote{更确切的预言，你们对此留意是好的，如同留意照在暗处的灯，直到破晓，晨星在你们心中升起。} 然后他称\Quote{早晨}为世界终末后的时刻，那时我们将看见今世所相信的一切。但在此，直到早晨来到之前呢？因为单单盼望早晨是不够的；我们必须做些什么。为何要做些什么？在黑夜中，人应以双手寻求天主。什么是\Quote{以双手}？就是善工。既然我们必须这样盼望早晨，忍受此黑夜，在忍耐中坚持直到破晓，那么在此期间我们在此该做什么呢？以免你自以为能靠自己做些什么，从而赚得被带往早晨。\Quote{求祢使我认识当行的路，主。} 因此祂点燃了预言之灯，因此祂派遣主降在肉身器皿中，祂甚至要说：\Quote{我的力量如瓦片般枯干。} 要照着预言行走，照着未来之事预言之灯行走，照着天主圣言行走……\par

11. \BibleQuote{求祢从我的仇敌中拯救我，主，因为我逃避到祢那里。}\BibleRef{咏 142:9} 我曾一度从祢面前逃走，如今却逃向祢。因为亚当从天主的容颜前逃走，藏在乐园的树木中，以致在\BibleBook{约伯}{传}中论到他说：\BibleQuote{如同仆人逃避他的主人，找到荫影。} 他从他主人的容颜前逃走，找到一片荫影。他有祸了，如果他继续留在荫下，免得后来有人说：\BibleQuote{一切都像影子般消逝。}\BibleRef{智 5:9} 这世界的统治者，这黑暗的统治者，恶者的统治者；你们要与这些搏斗。你们的战斗是巨大的：不见敌人而仍得胜。对抗这世界的统治者、这黑暗的统治者，就是魔鬼及其使者；不是那经上说\BibleQuote{世界是藉着祂造的}的世界，而是那经上说\BibleQuote{世界不认识祂}的世界。\BibleRef{若 1:10} \BibleQuote{因为我逃避到祢那里。}…… 我该逃到哪里？\Quote{我何处可躲避祢的神？}\par

12. \BibleQuote{求祢教导我承行祢的旨意，因为祢是我的天主。}\BibleRef{咏 142:10} 光荣的宣认！光荣的法则！\Quote{因为祢，}他说，\Quote{是我的天主。} 若我曾由另一位造成，我必急忙奔向另一位被重塑。祢是我的一切，\Quote{因为祢是我的天主。} 我要寻求一位父亲以获得产业吗？\Quote{祢是我的天主，}不单是产业之赐予者，更是产业本身。\Quote{上主是我产业的份。} 我要寻求一位保护者以获得救赎吗？\Quote{祢是我的天主。} 最后，既受造，我渴望被重塑吗？\Quote{祢是我的天主，}我的造物主，祢曾藉着祢的圣言造了我，也藉着祢的圣言重塑了我。\Quote{求祢教导我：}因为祢是我的天主，而我不可能作自己的主人。看，恩宠如何被推荐给我们。要持守，要吸收，莫让任何人从你们心中驱走它，免得你们有\BibleQuote{对天主的热诚，但不按照真知。}\BibleRef{罗 10:2} 那么，这样说：\Quote{祢的善神，}不是我的恶神，\Quote{祢的善神要领我进入正直之地。} 因为我的恶神已领我进入弯曲之地。我有什么功德？没有祢的帮助，我能算出什么善工，藉以获得、配得被祢的圣神领入正直之地呢？\par

13. 那么，用你们全部力量聆听恩宠的推荐，你们白白得救，无需代价。\BibleQuote{上主，因祢圣名的缘故，求祢因祢的正义使我活。}\BibleRef{咏 142:11} 不是因我的正义：不是因为我配得，而是因为祢仁慈。因为若我展示自己的功德，从祢那里我什么也不配得，只配惩罚。祢已从我身上剪除我自己的功绩；祢已接上祢自己的恩赐。\BibleQuote{求祢从患难中领出我的灵魂。} \BibleQuote{并因祢的仁慈消灭我的仇敌：祢要毁灭一切磨难我灵魂的人，因为我是祢的仆人。}\BibleRef{咏 142:12}\par

\SourceRecord{Translated by J.E. Tweed.；Nicene and Post-Nicene Fathers, First Series；Vol. 8.；Edited by Philip Schaff.；Buffalo, NY: Christian Literature Publishing Co.,；1888.；<http://www.newadvent.org/fathers/1801143.htm>.}

% structure: p0001 source=h1 target=section reason=page-title

\section{\WorkTitle{《聖詠第一四四篇》釋義}}

1. 這篇聖詠的題名，字數雖簡，但其奧義的分量卻很沉重。\Quote{致達味本人，反對哥肋雅。}這場戰事發生在我們祖先的時代，親愛的各位，請你們同我一起從《聖經》中回憶此事……達味將五塊石頭放入他的行囊，卻只擲出了一塊。五卷書被選出，但合一得勝了。然後，他擊倒並打敗了敵人，拿起敵人的劍，砍下了他的頭。我們的這位達味也同樣做了：祂用敵人自己的武器推翻了魔鬼；當那些在魔鬼權下的大人物——魔鬼藉著他們殺害其他靈魂——信了，他們便用自己的舌頭轉而反對魔鬼，這樣，哥肋雅的頭就被他自己的劍砍下了。\par

2. \BibleQuote{願主，我的天主，受讚頌，祂教導我的手作戰，我的手指打仗}（\BibleRef{詠 143:1}）。這是我們的言語，如果我們是基督的身體。這似乎是意思的重複：\Quote{我的手作戰}和\Quote{我的手指打仗}是一樣的。或者\Quote{手}和\Quote{手指}之間有某種區別？當然，手和手指都工作。那麼，我們把\Quote{手指}當作\Quote{手}來理解，並非沒有理由。但在\Quote{手指}中，我們仍然認識到行動的分工，然而仍有一種合一。看哪，那恩寵！宗徒說：這人得這個，那人得那個；\Quote{恩賜各有區別，卻只有同一個聖神在運作一切}；那就是合一的根。因此，基督的身體就用這些\Quote{手指}作戰，出去\Quote{打仗}，出去\Quote{交戰}……藉著慈悲的工作，我們的仇敵被征服；如果我們沒有愛德，我們就不可能行慈悲的工作；除非我們藉著聖神領受了愛德，否則我們不可能有愛德；因此，祂\Quote{教導我們的手作戰，我們的手指打仗}：我們正當地對祂說：\Quote{我的慈悲}，因為我們從祂領受了我們也成為慈悲的：\BibleQuote{因為那不憐憫人的，將沒有憐憫的審判}（\BibleRef{雅 2:13}）。\par

3. \BibleQuote{我的慈悲，我的避難所，我的支持者，我的拯救者}（\BibleRef{詠 143:2}）。這位戰士勞苦甚多，他的肉情慾反抗著他的心神。保持你所擁有的。那麼，當\BibleQuote{死亡被吞滅於勝利中}（\BibleRef{格前 15:54}），你將完全擁有你所願望的；那時這可朽壞的身體被復活，改變成天使的狀況，升騰到天上的性質……那裏有生命，那裏有美好的日子，在那裏沒有什麼反抗心神，那裏不是說\Quote{戰鬥}，而是說\Quote{喜樂}。但誰渴望這些日子呢？當然每個人都說：\Quote{我渴望。}請聽接下來的話。我看到你在勞苦，我看到你正在爭戰，處在危險中；請聽接下來的話……\Quote{離棄邪惡，行善吧}：不要讓窮人先在你之下哭泣，好讓窮人能因你而喜樂。因為有什麼報酬呢？既然你現在在戰鬥。\Quote{尋求和平，並追隨它。}學習並說：\Quote{我的慈悲，我的避難所，我的支持者，我的拯救者，我的保護者：\InnerQuote{我的支持者}，免得我跌倒；\InnerQuote{我的拯救者}，免得我陷入；\InnerQuote{我的保護者}，免得我受打擊。}在这一切事中，在我一切的勞苦中，在我一切的戰役中，在我一切的困難中，我寄望於祂，\Quote{祂使我的百姓臣服於我}。看，我們的頭同我們一起說話。\par

4. \BibleQuote{主，人是什麼，以致你讓他認識你？}（\BibleRef{詠 143:3}）\Quote{或人子，以致你重視他？}你重視他，就是說，你使他如此重要，你把他看作如此的寶貴，你知道你把他放在什麼之下，你把他置於什麼之上。因為重視就是估量一件東西的價值。祂為人何等重視，竟為他傾流了祂獨生子的血！因為天主重視人不像一個人重視另一個人一樣：人若買一個奴隸，他為一匹馬付出的價錢比為一個人更高。請思考祂何等重視你，好使你能說：\Quote{若是天主偕同我們，誰能反對我們呢？}祂何等重視你，\Quote{祂沒有憐惜自己的兒子} \BibleQuote{祂豈不也把一切同祂一起賜給我們嗎？}（\BibleRef{羅 8:31}）祂將這食物賜給作戰的人，那麼祂為得勝者保留的是什麼呢？……\par

5. \BibleQuote{人相似虛幻：他的歲月如同影兒飛逝}（\BibleRef{詠 143:4}）。什麼虛幻？時間，它逝去，流走。因為這個\Quote{虛幻}是與永恆不變、永不衰竭的真理相比較而言的；因為它（時間）也是祂手中的一個作品，按其程度而言。正如經上所載：\BibleQuote{天主以祂的美好之物充滿了大地}（\BibleRef{德 16:29}）。\Quote{祂的}是什麼？就是與祂相合的事物。但這一切屬地、短暫、易逝的事物，若與那真理相比，就是那裏所說的\BibleQuote{我是自有者}（\BibleRef{出 3:14}），那麼這一切流逝之物就被稱為\Quote{虛幻}。因為藉著時間，它如空氣中的波動一樣消失。我何必多說那宗徒雅各伯所說的呢？他願意將驕傲的人引入謙卑：\BibleQuote{你們的生命是什麼？不過是一片雲霧，出現片刻，以後就消失了}（\BibleRef{雅 4:14}）……那麼，即使是在夜間，也要用你的手工作，就是藉著善工尋求天主，免得那應使你歡樂的日子來臨，卻成了使你憂傷的日子。因為你看，你若尋求祂而不被祂離棄，你的工作是多麼安全，\BibleQuote{好使你那在暗中看見的父，必要報答你}（\BibleRef{瑪 6:4}）……\par

6. \Quote{主啊，求祢屈尊诸天，降临；触摸群山，群山冒烟}（\BibleRef{咏 143:5}）。\Quote{求祢发出闪电，使仇敌四散；射出祢的箭，使他们混乱}（\BibleRef{咏 143:6}）。\Quote{求祢从高处伸出祢的手，拯救我，把我从多水中拉出来}（\BibleRef{咏 143:7}）。基督的身体，谦卑的达味，充满恩宠，依靠天主，在世上争战，呼求天主的助佑。什么是\Quote{屈尊的诸天}？是谦卑的宗徒。因为那些\Quote{诸天宣扬天主的荣耀}；关于这些宣扬天主荣耀的诸天，现在说道：\Quote{没有言语，没有语言，但他们的声音在其中被听到}等等。当这些诸天将他们的声音传遍大地，并行了奇事，而主从他们中藉着奇迹和诫命发出闪电和雷声时，人们以为诸神从天降到了人间。因为某些外邦人，这么想，甚至想要向他们献祭……但他们向这些人推荐了主耶稣基督，谦卑自己，为叫天主受赞美；因为\Quote{诸天}被\Quote{屈尊}，使\Quote{天主}可以\Quote{降临}。……\Quote{触摸群山，群山冒烟}。只要它们不被触摸，它们在自视甚高；它们现在快要说：\Quote{主啊，祢何其伟大}；群山也快要说：\Quote{唯有祢是至高者，超越全地}。\par

7. 但有些人共谋，\Quote{聚集起来反对主，反对祂的受傅者}。他们聚集，他们共谋。\Quote{发出祢的闪电，祢就驱散他们}。使祢的奇迹丰盛，他们的阴谋必被粉碎。……\Quote{射出祢的箭，祢就使他们混乱}。让不健康的人受伤，好使他们因受重伤而变得健康；让他们，现在已在教会中、在基督的身体内，与教会一同说：\Quote{我为爱受伤}。\Quote{从高处伸出祢的手}。此后如何？结局如何？基督的身体如何得胜？藉着天上的援助。\Quote{因为主必亲自随同总领天使的呼声和天主的号角降下}（\BibleRef{得前 4:16}），祂自己是身体的救主，天主的手。\Quote{从多水中}是什么？从许多民族中。什么民族？外邦人、不信者，无论是从外攻击我们，还是在内设伏。把我从多水中拉出来，祢在其中管教了我，在其中滚动了我，为使我脱离污秽。这就是\Quote{争闹之水}（\BibleRef{户 20:13}）。……\Quote{从外邦之子手中}。请听，弟兄们，我们身处他们中间，与他们一同生活，渴望从中被解救出来。\Quote{他们的口说虚妄的话}（\BibleRef{咏 143:8}）。你们所有人，今天若没有聚集在这里观看天主圣言的这些神圣演出，此刻没有参与其中，你们会听到多么大的虚妄！\Quote{他们的口说虚妄的话}：简而言之，他们何时说虚妄的话，会听到你们说虚妄的话？\Quote{他们的右手是邪恶的右手}。你带着装有五块石头的牧羊袋在他们中间做什么？请你以另一种形式对我说：你所用以五块石头表示的同一条法律，也以某种其他方式表示它。\Quote{天主，我要向祢唱新歌}（\BibleRef{咏 143:9}）。\Quote{新歌}关乎恩宠；\Quote{新歌}关乎新人；\Quote{新歌}关乎新约。但恐怕你以为恩宠离开了法律，其实恩宠反而成全了法律，\Quote{我要在十弦琴上向祢歌唱}。在十诫的法律上：在那里我向祢歌唱；在那里我向祢欢欣；在那里\Quote{我向祢唱新歌}；因为\Quote{爱是法律的满全}（\BibleRef{罗 13:10}）。但没有爱的人可以拿着十弦琴，却不能歌唱。矛盾不能使我的十弦琴沉默。\par

8. \Quote{祢把救恩赐给君王，救赎了祢的仆人达味}（\BibleRef{咏 143:10}）。你们知道达味是谁；你们自己成为达味。从哪里\Quote{祂救赎了祂的仆人达味}？从哪里祂救赎了基督？从哪里祂救赎了基督的身体？\Quote{救我脱离邪恶意图的剑。}\Quote{脱离剑}是不够的；他加上\Quote{邪恶意图的}。无疑地，有一种善意之剑。什么是善意之剑？就是主所说的：\Quote{我来不是把平安带到地上，而是带来剑}（\BibleRef{玛 10:34}）。因为祂要分开信徒与不信者、儿子与父母，并斩断所有其他联系，剑砍去有病的部分，却治愈了基督的肢体。所以善意之剑是双刃的，两刃都有力，即旧约和新约，包含过去的历史和未来的应许。那就是善意之剑；但另一种是邪恶意图的，人用它说虚妄的话；因为那剑是善意的，天主用它说出真理。\BibleQuote{诚然人们的牙齿是枪和箭，他们的舌头是锋利的剑}。从这\Quote{剑}救我（\BibleRef{咏 143:11}）。\Quote{把我从外邦之子手中取出，他们的口说虚妄的话}：和先前一样。至于接下来的\Quote{他们的右手是邪恶的右手}，他先前也提到过，当他称他们为\Quote{多水}时。因为恐怕你以为\Quote{多水}是好的水，他就用\Quote{邪恶意图的剑}来解释它们。\par

9. \BibleQuote{他们的儿子好比幼小的葡萄树，在青年期被稳固地栽种} \BibleRef{咏 143:12}。他愿意讲述他们的幸福。请观察，你们光明之子，和平之子：请观察，你们教会的子女，基督的肢体；请观察他称谁为\Quote{外人}，称谁为\Quote{外邦之子}，称谁为\Quote{悖逆之水}，称谁为\Quote{恶意的剑}。请观察，我恳求你们，因为你们在他们中间处于危险之中，在你们的舌头之间，你们与肉身的欲望作战，在他们的舌头之间，那被置于魔鬼手中，魔鬼用以作战。\BibleRef{弗 6:12} ……他们的口说了何等虚妄的事，他们的右手是怎样的不义之手？\Quote{他们的女儿被装饰打扮，如同圣殿的模样。}\BibleQuote{他们的仓库充满，从一仓溢出到另一仓；他们的羊群繁殖，在街上增多} \BibleRef{咏 143:13}：\BibleQuote{他们的牛犊肥壮；他们的篱笆没有倒塌，他们的道路没有毁坏，他们的街上也没有哭号} \BibleRef{咏 143:14}。这难道不是幸福吗？我问你们，天国之子，永远复活的后裔，基督的身体，基督的肢体，天主的圣殿。这难道不是幸福吗——拥有安全的儿子，美丽的女儿，满仓的谷物，众多的牲畜，没有倒塌——我不说城墙，甚至篱笆也没有——街上没有骚动和喧哗，只有宁静、和平、富足，家中和城里一切丰盛？这难道不是幸福吗？或者，义人应当逃避它吗？难道你们没有发现义人的家也充满了这一切，充满了这种幸福吗？亚巴郎的家不是富有金银、儿女、仆役、牲畜吗？我们怎么说？这不是幸福吗？就算是，它仍然在左手边。什么在左手边？暂时的、可朽的、肉身的。我并不愿你逃避它，而是不要以为它在右手边。……因为他们应当把什么放在右手边呢？天主、永恒、永不衰竭的天主之年，论及此，经上说：\Quote{你的年岁没有穷尽。}右手边应当在那里，我们的渴望应当在那里。让我们暂时使用左手，让我们切望永恒的战斗。\Quote{若财富增加，不要一心依恋它们。}。..\par

10. \BibleQuote{他们称拥有这些的人民为有福的} \BibleRef{咏 143:15}。啊，说话虚妄的人！他们失去了真正的右手，邪恶而悖逆，他们反转了天主的恩惠。啊，邪恶的人，说话虚妄的人，外邦之子！那在左手边的，他们放在了右边。你们呢，\Person{达味}？你们呢，基督的身体？你们呢，基督的肢体？你们呢，不是外邦之子，而是天主的子女？……你们说什么？你们要同我们一起说：\BibleQuote{那以主为其天主的人民，是有福的。}\par

\SourceRecord{Translated by J.E. Tweed.；Nicene and Post-Nicene Fathers, First Series；Vol. 8.；Edited by Philip Schaff.；Buffalo, NY: Christian Literature Publishing Co.,；1888.；<http://www.newadvent.org/fathers/1801144.htm>.}

% structure: p0001 source=h1 target=section reason=page-title

\section{圣咏第一百四十五篇释义}

1. ...标题是：\Quote{赞美，归于达味自己。}赞美归于基督自己。既然那被称为达味的、出自达味后裔来到我们中间者，是我们的君王，治理我们，引领我们进入他的王国，因此\Quote{赞美归于达味自己}应理解为赞美归于基督自己。基督按肉躯是达味，因为他是达味之子；但按他的天主性，他是达味的造物主和达味的主。\BibleQuote{我要高举你，我的天主，我的君王；我要赞美你的圣名，直到永世，代代无穷}\BibleRef{咏 144:1}。你看，天主的赞美从这里开始，这赞美一直延续到圣咏的结束……那么现在就开始赞美吧，如果你想要永远赞美。那不愿在这短暂的\Quote{世代}中赞美的，当\Quote{代代无穷}来临时，他将缄默。但恐怕有人以别样方式理解他所说的\Quote{我要赞美你的圣名，直到永世}，并寻求另一个世代去赞美，他便说：\BibleQuote{我要日日称颂你}\BibleRef{咏 144:2}。所以，要日日赞美祝福主你的天主，好使当一个个日子过去，那无尽的一日来临时，你能从赞美到赞美，如同\Quote{从力量到力量}。没有一日会消逝，而我不称颂你。如果你在喜乐之日赞美主，这不足为奇。但若偶然某日忧伤降临于你，按我们必死本性的状况，罪过众多，诱惑倍增，那么，若你遭遇什么悲伤之事，作为人，你会停止赞美天主吗？你会停止祝福你的造物主吗？若你停止，你便在说\Quote{每日}的事上说了谎。但若你不停止，虽然在忧愁之日你似乎处境不佳，但在你的天主内，你必会安好……\par

2.\BibleQuote{上主伟大，极堪赞美}\BibleRef{咏 144:3}。他想要说多少呢？他要寻求怎样的言辞呢？他在\Quote{极堪}一词中包含了何等广阔的概念？随你想象，因为那无法被包含的，怎能被想象呢？\Quote{他极堪赞美，他的伟大没有穷尽}；因此他说\Quote{极堪}：免得你或许开始想要赞美，并以为你能达到他赞美的尽头，而他的伟大是没有尽头的。那么不要以为他，这位伟大没有尽头者，能被你足够赞美。既然他没有尽头，你的赞美也不应有尽头，这样岂不更好？他的伟大没有尽头；愿你的赞美也没有尽头……\par

3. 此外，他那无穷的良善和无限的伟大还创造了多少我们不知道的宏伟事物！当我们举目仰望天空，然后从太阳、月亮、星辰收回目光至大地，视线所及尽是空间；谁能将心灵的目光，更不必说肉躯的目光，延伸到诸天之外呢？那么就按我们所知的他的化工，让我们因他的化工而赞美他。\BibleRef{罗 1:20}\BibleQuote{世世代代宣扬你的化工}\BibleRef{咏 144:4}。每一代都要宣扬你的化工。因为\Quote{世代和世代}或许意指每一代？……他是否特意用重复来暗示两个世代？因为在这一世代，我们是天主的儿子；在另一世代，我们将是复活之子。圣经称我们为\Quote{复活之子}；它又称复活本身为\Quote{重生}。\BibleQuote{在重生时}，它说，\BibleQuote{当人子坐在他的威严中}\BibleRef{玛 19:28}。同样在另一处；\BibleQuote{因为他们不娶不嫁，因为他们是复活之子}\BibleRef{路 20:35}。因此\Quote{世世代代宣扬你的化工……他们要传述你的尊威。}因为他们宣扬你的化工，正是为了\Quote{传述你的尊威}啊。学校的孩童被教导去赞美，所有这类事物被摆在他们面前供赞美，正如天主所行的一样：一个必死之人被教导去赞美太阳、天空、大地；甚至更小的事物，去赞美一朵玫瑰或一棵月桂；这些都是天主的化工：它们被提出、被承担、被赞美；作品受赞扬，而对造物主却默不作声。我渴望在化工中赞美造物主：我不喜欢忘恩的赞美者。你赞美他所造的，却对造物主缄默不语？在你所见的事物中，你赞美什么？是形状、功用、某种美德、还是事物中的某种力量？若美令你喜悦，有什么比造物主更美？若功用被赞美，有什么比造物主更有用？若尊威被赞美，有什么比造物主更尊威？……\par

4. \BibleQuote{他們要述說你聖德的尊威光榮，傳揚你的奇妙作為} \BibleRef{詠 144:5}。\BibleQuote{他們要講述你可畏的作為的優越；他們要宣揚你的偉大} \BibleRef{詠 144:6}。\BibleQuote{他們要湧出你豐厚甘飴的紀念} \BibleRef{詠 144:7}：惟獨你的甘飴。請看，這個人在默想你的作為時，是否從工人轉向了工作？是否從造物主沉淪到他所造之物？對於他所造之物，他把它們當作階梯攀登到他那裡，而不是從他那裡墮落到它們之中。因為你若愛這些超過愛他，你就不會擁有他。若工人離你而去，擁有滿滿的工作又有何益？你確實應當愛它們，但更要愛他，為他的緣故而愛它們。因為他既給予應許，也給予警告：若他不給予應許，就沒有鼓勵；若他不給予警告，就沒有糾正。因此，讚美你的人也要\BibleQuote{講述你可畏的作為的優越}；他們要講述你那懲罰和管教的手所作的卓越工作，他們不會緘默：因為他們不會在宣揚你的永恆國度時，對你的永恆烈火保持沉默。因為對天主的讚美，既引導你走上正道，就該向你展示你應愛什麼，應怕什麼；你應追求什麼，應躲避什麼；你應選擇什麼，應避免什麼。選擇的時候是現在，領受的時候將在將來。因此，讓你可畏之事的優越為人所傳揚。雖然它是無限的，但\Quote{你的偉大無窮無盡}，他們不會對它沉默。他們如何敘述它，如果它是無窮的？他們在讚美時敘述它；因為它是無窮的，所以對他的讚美也將永無窮盡。\par

5. \BibleQuote{他們要湧出你豐厚甘飴的紀念。} 何等幸福的盛宴！那些這樣\Quote{湧出}的人要吃什麼！……你要這樣吃，好叫你也能湧出；這樣領受，好叫你能給予。你吃，當你學習時；你湧出，當你教導時；你吃，當你聆聽時；你湧出，當你宣講時，但你湧出的是你首先吃下的。最後，那最熱切的赴宴者\Person{若望}，連主的餐桌都不足以滿足他，除非他靠在主的懷裡，從他內心深處汲取神聖的奧秘；他湧出了什麼？\BibleQuote{在起初已有聖言，聖言與天主同在。} \BibleRef{若 1:1} 為什麼僅僅說\Quote{你的紀念}或\Quote{你豐厚的紀念}還不夠呢？因為，如果豐厚卻不甜蜜，又有何用？同樣，如果甜蜜卻太少，也是惱人的。\par

6. ……藉著\Quote{湧出}這一點，他的宣講者\Quote{要因他的正義而歡躍}，而不是因他們自己的正義。上主，我們所讚美的主，你為我們做了什麼，使我們得以存在，使我們得以讚美，使我們\Quote{因你的正義而歡躍}，使我們\Quote{說出你豐厚甘飴的紀念}？讓我們述說，並在述說時讚美。\par

7. \BibleQuote{上主慈悲為懷，寬仁大方，緩於發怒，仁愛無量} \BibleRef{詠 144:8}。\BibleQuote{上主對一切人慈愛，他的慈愛散佈在他的一切造化上} \BibleRef{詠 144:9}。若非如此，就沒有人尋求挽回我們。請看你自己：罪人，你配得什麼？輕視天主的人，你配得什麼？看看有什麼臨到你，除了刑罰，還有什麼？你看到你應得的，以及他所賜予的——他白白賜予。罪人得赦免，得賜成義之神，得賜愛德，使你行一切善工；此外，他還將賜給你永生和與天使的共融：這都是他的憐憫。……請聽聖經：\BibleQuote{我不願惡人死亡，而願他歸依，得以生存。} \BibleRef{則 33:11}。天主的這些話使人重獲希望；但還有一個陷阱需要警惕，免得這希望反使人更放肆犯罪。那麼，你這因希望而更放肆犯罪的人，你會說什麼？\Quote{我隨時歸依，天主就會寬恕我一切；我要隨意行事。} 不要說：\Quote{明天我會歸依，明天我會悅樂天主；今天和昨天的一切行為都將被赦免。} 你說得對：天主應許了你的歸依得赦免；他沒有應許你的拖延有明天。\BibleRef{德 5:7}\par

8. \BibleQuote{上主對一切人慈愛，他的慈愛散佈在他的一切造化上。} 那麼，他為何定罪？為何鞭打？他所定罪、鞭打的，難道不是他的造化嗎？的確是。你願知道\Quote{他的慈愛散佈在他的一切造化上}嗎？正因他的寬容，\BibleQuote{他使太陽上升，照惡人，也照善人。} \BibleRef{瑪 5:45} \BibleQuote{他的慈愛豈不是散佈在他的一切造化上，降雨給義人，也給不義的人}？他因寬容等待罪人，說：\BibleQuote{你們歸向我，我就轉向你們。} \BibleRef{匝 1:3} \Quote{他的慈愛豈不是散佈在他的一切造化上}？當他說：\BibleQuote{你們到那為魔鬼和他的使者預備的永火裡去} \BibleRef{瑪 25:41}，這不是他的慈愛，而是他的嚴厲。他的慈愛賜給他的造化；他的嚴厲不是針對他的造化，而是針對你的惡行。最後，你若除去你自己的惡行，在你內只剩下他的造化，他的慈愛必不離開你；但你若不除去你的惡行，嚴厲將臨到你的惡行，而非他的造化。\par

9. \Quote{主，愿祢的一切受造物都称谢祢，愿祢的圣人祝福祢}\BibleRef{咏 144:10}。这是怎么回事？大地不是祂的受造物吗？树木不是祂的受造物吗？牲畜、走兽、鱼、飞禽，难道不都是祂的受造物吗？显然它们也是。那么，它们又如何称谢祂呢？我确实看到，在天使中祂的受造物称谢祂，因为天使是祂的受造物；人类也是祂的受造物；当人向祂称谢时，祂的受造物便在称谢祂；但树木和石头有称谢的声音吗？是的，实在如此：\Quote{愿一切}受造物\Quote{都称谢}祂。你说什么？连大地和树木也……但在赞美的事上，也出现与称谢相同的问题。因为如果大地和一切无感觉之物不能称谢，因它们没有声音来称谢；那么它们也不能赞美，因它们没有声音来宣告。但那三个孩子在无害的火焰中行走时，不也列举了一切受造物吗？那时他们有时间不仅不害怕，甚至赞美天主。他们对天上和地上的一切说：\Quote{你们要赞美上主，颂扬祂，永远尊崇祂。}看，他们就这样赞美。不要以为无言的石头或无言的动物有理智可以理解天主。这样想的人，已远离真理。天主安排了一切，造了一切：有些祂赋予了知觉、理智和永生，如天使；有些祂赋予了知觉和理智与可朽性，如人类；有些祂赋予了身体的感觉，却没有给予它们理智或永生，如牲畜；有些祂既没有赋予感觉，也没有赋予理智或永生，如草木、石头；然而这些也不能在它们的种类中缺少，祂按一定的等级安排了祂的受造物，从地上到天上，从可见到不可见，从可朽到不朽。这受造物的框架，这最完美有序的美，从最低升至最高，从最高降至最低，从未中断，却由不同事物调和而成，全都赞美天主。那么，为什么一切都赞美天主呢？因为当你观看它，见到它的美，你便在它内赞美天主。大地的美，是无声大地的一种声音……而你在它内所发现的，正是它称谢的声音，使你赞美造物主。当你思想这世界的普遍之美时，难道它的美本身不就像用同一个声音回答你：\Quote{不是我造了自己，是天主造了我}？\par

10. 因为当祢的圣人祝福祢时，他们说什么？\Quote{他们要讲论祢王国的荣耀，谈论祢的大能}\BibleRef{咏 144:11}。天主多么有能力，祂造了大地！天主多么有能力，祂以美物充满了大地！天主多么有能力，祂赐给牲畜各自的生命！天主多么有能力，祂将不同的种子放入大地的腹中，使它们生出如此多样的嫩芽，如此美丽的树木！天主的权能何其伟大！你询问，受造物回答，借着它的回答，借着受造物的称谢，你，天主的圣人，祝福天主，\Quote{谈论祂的大能}。\par

11. \Quote{为使祢的能力被世人知晓，祢王国美之伟大的荣耀}\BibleRef{咏 144:12}。因此，祢的圣人传扬\Quote{祢王国美之伟大的荣耀}，其美之伟大的荣耀。有一种\Quote{祢王国之美的伟大}：也就是说，祢的王国具有美，且是伟大的美。因为凡有美的事物，其美都来自于祢，祢的整个王国该有多么伟大的美！愿那王国不要吓倒我们：它也有美，足以令我们喜悦。因为那将来圣人所要享有的美是什么？他们将听到：\Quote{来，我父所祝福的，承受王国吧}\BibleRef{玛 25:34}。他们从哪里来？要往哪里去？弟兄们，请看，如果你们能够，尽你们所能，思想那将要来临的王国之美；我们的祈祷说：\Quote{愿祢的国来临}。我们渴望那王国来临，圣人宣告那王国将要来临。看这世界：它是美丽的。大地、海洋、空气、天空、星辰多么美丽。难道这一切不会使思考它们的人感到畏惧吗？它们的美如此显著，似乎不可能找到更美的东西。在这美丽、这几乎不可言说的优美中，虫、老鼠和地上一切爬行之物与你一同生活，它们与你一同生活在这全部的美中。那唯有天使与祢一同生活的王国，其美何其伟大！那里有一种伟大之美；愿它在被看见之前就被爱慕，好使当它被看见时，得以持守。\par

12. \Quote{祢的王国。}我说的什么王国？\Quote{万世的王国。}因为今世的王国也有它自身的美，但其中没有那样伟大的美，如同在\Quote{万世的王国}中。\Quote{祢的统治存留于世世代代}\BibleRef{咏144:13}。这就是我们所注意到的重复，表示每一世代，或指此世之后的世代。\Quote{主在祂的话上是信实的，在祂的一切工作中是圣洁的。}\Quote{主在祂的话上是信实的：}因为祂所应许的，哪一样没有赐下？\Quote{主在祂的话上是信实的。}至此，有些事祂曾应许，但尚未赐下；但要从祂所赐下的事上来相信祂。我们理应相信祂，即使祂只是说话；祂不愿意我们只相信祂说话，而要我们手中有祂的圣经：……如同一份天主的契约，所有路过的人都能阅读，并遵守其应许的道路。按照那份契约，祂已经偿还了多么大的事！人们还犹豫是否相信祂关于死者复活和来世生命的事吗？那是唯一尚待偿还的，如果祂来与不信者算账，不信者岂不脸红？如果天主对你说：你有我的契约；我应许了审判、善恶的分离、给信徒的永生，你还不信吗？在我的契约中，读一读我所应许的一切，与我算账：确实，数算我已偿还的，你就能相信我会偿还仍然所欠的。在那契约中，你得到了我独生子的应许，\Quote{祂没有怜惜，反而为你们众人把祂交出}\BibleRef{罗8:32}；那么，把这算作已偿还的。读那契约：我应许了藉我子赐下圣神的凭据：算作已偿还。我应许了光荣殉道者的血和冠冕；愿白弥撒提醒你，我欠的债已经偿还。……祂在众人眼前陈列祂债务的偿还：有些在我们祖先的时代偿还了，我们没有看见；有些在我们的时代偿还了，他们没有看见；在历世历代中，祂偿还了所写的一切。还剩下什么？当祂偿还了这一切之后，人还不信祂吗？还剩下什么？看，你已经算过了：这一切祂都偿还了；难道祂会因为剩下的几件事而变得不忠信吗？决不会！为什么？因为\Quote{主在祂的话上是信实的，在祂的一切工作中是圣洁的。}\par

13. \Quote{主扶持一切跌倒的人}\BibleRef{咏144:14}。但谁才是\Quote{一切跌倒的人}？所有的人一般意义上都跌倒，但祂指的是那些以特定方式跌倒的人。因为许多人离开祂跌倒，许多人从自己的幻想中跌倒。如果他们怀有恶念，就从这些恶念中跌倒，而\Quote{天主扶持一切跌倒的人。}有人在世上失去什么，却仍是圣洁的，他们在世上仿佛受辱，从富有变为贫穷，从尊贵变为卑贱，然而他们是天主的圣人；他们仿佛在跌倒。但\Quote{天主扶持。}因为\Quote{义人七次跌倒，仍必起来；但恶人必在祸患中软弱}\BibleRef{箴24:16}。当灾祸临到恶人，他们因此软弱；当灾祸临到义人，\Quote{主扶持一切跌倒的人。}……\Quote{并扶起所有被推倒的：}即所有属于祂的人；因为\Quote{天主抵制骄傲的人}\BibleRef{雅4:6}。\par

14. \Quote{众目都仰望祢，祢按时赐给他们食物}\BibleRef{咏144:15}。正如你按时慰问病人，当病人应当接受时，祢就赐予，并且赐予他所应当接受的。有时人们长久渴望，祂却不赐予；那位照料者知道赐予的时间。我为何说这话，弟兄们？免得有人心里沮丧，如果偶然他祈求天主一些正义的事，却未被垂听。因为当他祈求不义的事时，他会被垂听而受罚；但当他祈求天主一些正义的事时，如果偶然未被垂听，不要灰心，不要沮丧，让他的眼睛等待那食物，祂按时赐予。当祂不赐予时，正是因而不赐予，免得所赐的造成伤害。\BibleRef{格后12:7}……\Quote{祢按时赐给他们食粮。}\par

15. \Quote{祢张开祢的手，使一切生物都充满福泽}\BibleRef{咏144:16}。虽然有时祢不赐予，但\Quote{按时}祢仍赐予；祢延迟，并非拒绝，而且是按时。\Quote{主在祂的一切道路上正义，在祂的一切工作中圣洁}\BibleRef{咏144:17}。无论祂责打还是医治，祂都是正义的，在祂内没有不义。最后，祂所有的圣人在患难中时，都先赞美祂的正义，然后寻求祂的福泽。他们先说了：\Quote{祢所做的是正义的。}达尼尔和其他圣人也如此祈求：\Quote{祢的审判是正义的：我们受得对，我们遭得应当。}他们不把不义归给天主，也不把不公和愚昧归给祂。他们先赞美祂的责罚，然后才感到祂的喂养。\par

16. \Quote{主亲近所有呼求祂的人} \BibleRef{咏 144:18} 。那么，那句话在哪里呢：\Quote{那时他们呼求我，我必不俯听}？\BibleRef{箴 1:28} 再看下文：\Quote{一切以真理呼求祂的人。}因为许多人呼求祂，但不是以真理。他们向祂寻求别的东西，却不寻求祂本身。你为什么爱天主？\Quote{因为祂使我痊愈。}这很清楚：是祂使你如此。因为健康无不由祂而来。\Quote{因为祂赐给我，}另一个人说，\Quote{一位富有的妻子，而我先前一无所有，且一位顺服我的。}这也是祂所赐：你说得对。\Quote{祂赐给我，}另一个人说，\Quote{许多好儿子，祂赐给我一户人家，祂赐给我一切好东西。}你为此而爱祂吗？……因此，若天主是良善的，祂将你所拥有的赐给了你，当祂将祂自己赐给你时，你将更为有福！你已向祂祈求了这一切：我恳求你也祈求祂本身。因为这些并不比祂更甜蜜，也绝不可与祂相比。那么，凡将天主本身置于他所领受的一切之上，并因所领受的而欢欣的人，他就\Quote{以真理呼求天主。}……\par

17. \Quote{祂必成就敬畏祂者的意愿} \BibleRef{咏 144:19} 。祂必成就，祂必成就：虽然祂不立刻成就，但祂必将成就。的确，若你敬畏天主，好承行祂的旨意，看，祂也以一种方式服务于你；祂成就你的意愿。\Quote{祂必俯听他们的祈祷，并拯救他们。}你看见，医生俯听是为拯救。何时？请听宗徒告诉你们。因为我们因希望而得救：但看得见的希望不是希望；但我们若盼望那看不见的，就应坚忍等待：\BibleRef{罗 8:24} \Quote{那救恩，}即伯多禄所称的\Quote{在末次时期要显现的救恩。}\BibleRef{伯前 1:5}\par

18. \Quote{主保守所有爱祂的人，而祂要消灭一切罪人} \BibleRef{咏 144:20} 。你看，在那如此甘甜的主那里也有严厉。祂要拯救所有寄望于祂的人，所有信众，所有敬畏祂的人，所有以真理呼求祂的人：\Quote{并要消灭一切罪人。}什么是\Quote{一切罪人}？就是那些坚持犯罪的人；敢于责难天主而不责难自己；天天与天主争辩；绝望于罪过的赦免，并因这绝望而堆积罪过；或邪僻地自许赦免，并因这自许而不离弃他们的罪恶和不敬。时候将到，这一切都要被分开，分为两群，一在右边，一在左边；义人承受永生王国，恶人进入永火。既然如此，我们既听到了主的祝福、主的作为、主的奇事、主的仁慈、主的严厉、祂对万有的眷顾、万有的承认，请看祂如何以赞美结束：\Quote{我的口要述说主的赞美，愿一切血肉之躯赞美祂的圣名，直到永远。}\BibleRef{咏 144:21}\par

\SourceRecord{Translated by J.E. Tweed.；Nicene and Post-Nicene Fathers, First Series；Vol. 8.；Edited by Philip Schaff.；Buffalo, NY: Christian Literature Publishing Co.,；1888.；<http://www.newadvent.org/fathers/1801145.htm>.}

% structure: p0001 source=h1 target=section reason=page-title

\section{圣咏第一四六篇诠解}

1. 看，圣咏在响；这是某个人的声音（而如果你愿意，这个人就是你），他鼓励自己的灵魂赞颂天主，对自己说：\BibleQuote{我的灵魂，请赞颂上主}\BibleRef{咏 145:1}。因为有时在此生当前的磨难与诱惑中，无论我们愿意与否，我们的灵魂受到困扰；他在另一篇圣咏中论及了这种困扰。但为了消除这困扰，他提出了喜乐；这喜乐尚未成为现实，而是出于望德；当灵魂受到困扰、焦虑、悲伤、哀痛时，他对它说：\Quote{请寄望于天主，因为我还要赞美祂。}……\par

2. 但谁在说这话，又对谁说呢？弟兄们，我们该说什么？是肉躯在说：\BibleQuote{我的灵魂，你当赞颂上主}吗？肉躯岂能给灵魂出好主意？无论肉躯多么被克制，借着主所赐的力量被制服如奴仆，使它完全如奴隶般服事我们，只要它不妨碍我们，就足够我们用了……因为身体就其作为身体而言，甚至低于灵魂；而任何灵魂，无论多么卑贱，都被发现比最卓越的身体更优越。你们不要觉得这很奇妙，就连任何卑贱有罪的灵魂也比任何伟大超绝的身体更好。它更好，不是在于功德，而是在于本性。灵魂确实有罪，被某些欲望的污点所玷污；然而黄金，即便生锈，也比最光的铅更好。那么让你们的思绪漫游过整个受造界，你们会看到我们所说的并非不可信：一个灵魂，无论多么应受责备，却比一个值得称赞的身体更值得称赞。有两样东西：灵魂和身体。我责骂灵魂，我称赞身体；我责骂灵魂，因为它有罪；我称赞身体，因为它健康。然而我是按它的种类称赞灵魂，按它的种类责备灵魂；同样，我也是按它的种类称赞身体或责备它。若你问我哪一个更好，是我所责备的，还是我所称赞的，你将得到奇妙的答复……于是你说论及最好的马和最坏的人：然而你宁可将你所指责的人置于你所称赞的马之上。灵魂的本性比身体的本性更卓越；它远超于身体，它是属灵的、非物质的，与天主的本质相似。它是无形可见的，它管理身体，推动肢体，引导感官，预备思想，实行行动，收纳无数事物的图像；总之，亲爱的弟兄们，谁足以赞扬灵魂呢？然而给予它的恩宠如此之大，以致这人说：\BibleQuote{我的灵魂，请赞颂上主。}……不是肉躯在说这话。愿身体如天使一样，它仍低于灵魂，不能给它的上级出主意。当肉躯适当服从时，它是灵魂的婢女：灵魂统治，身体服从；灵魂命令，身体执行；那么肉躯怎能给灵魂出这个主意呢？那么也许正是灵魂自身对自己说话，并好像在命令自己、劝勉自己、询问自己？因为她本性的某一部分因某些激情而摇摆；但在另一部分，他们称之为理性心智，即她用以思考的智慧，她依附于天主，如今向祂叹息，她察觉到自己本性中某些较低的部分被世俗情绪所困扰，并被某种尘世欲望的激动带向外在事物，离开了在内的天主；于是她将自己从外在的事物召回到内在，从较低的事物上到较高的事物，说：\BibleQuote{我的灵魂，请赞颂上主。}……灵魂借着理性心智从天主的光照中获得自我劝勉，由此她理解植根于造物主永恒本性中的智慧。她在那里读到某些该敬畏、该赞颂、该爱慕、该渴望、该追求的事物；但她尚未抓住它，尚未理解它；她仿佛被光辉眩目，没有力量停留在那里。因此她好像收聚自身进入一种健康状态，说：\BibleQuote{我的灵魂，请赞颂上主。}……然后灵魂，好像被压重，无法适当地站起，回答心智说：\BibleQuote{我一生要赞颂上主}\BibleRef{咏 145:2}。\BibleQuote{在我一生中}是什么意思？因为现在我在死亡中。因此首先鼓励你自己，说：\BibleQuote{我的灵魂，请赞颂上主。}你的灵魂回答你：我尽我所能地赞颂，微弱地、贫乏地、软弱地。为什么？因为\BibleQuote{我们住在此身内，与主远离}\BibleRef{格后 5:6}。……\par

3. \BibleQuote{在我一生中。}现在它有什么？它可能回答你：\Quote{我的死亡。}从何而来的\Quote{我的死亡}？因为我远离了主。因为若依附于祂是生命，那么离开祂就是死亡。但什么安慰你？望德。现在你在望德中生活；在望德中赞颂，在望德中歌唱。你的死亡来自现世的悲伤，你在对未来生命的望德中生活。你将如何赞颂你的主？\BibleQuote{只要我存在，我要向我的天主歌唱}。这是怎样的赞颂？\BibleQuote{只要我存在，我要向我的天主歌唱}？看哪，弟兄们，这将是怎样的存在；那里将有永恒的赞颂，也将有永恒的存在。看，现在你存在：你只要存在，你就向天主歌唱吗？看，你正在歌唱，然后转向一些事务，你不再歌唱，然而你存在：你存在，却不再歌唱。或许你的欲望也将你转向某事；你不仅不歌唱，甚至冒犯祂的耳朵，然而你存在。那将是怎样的赞颂，当你赞颂直到你存在？但\BibleQuote{只要我存在}是什么意思？会有他不再存在的时候吗？不，相反的，那个\Quote{长}将是永恒的，因此它将是真正的\Quote{长}。因为任何在时间中有终结的东西，无论它如何延长，却仍不是\Quote{长}。……\par

4. \Quote{不要信赖王侯}\BibleRef{咏 146:3}。弟兄们，我们在此领受了一项重大的任务；这声音来自天上，从高处向我们响起。因为如今，由于某种软弱，人的灵魂每当在此世遭受困苦时，便对天主失望，而选择依赖人。若有人身处某种患难中，有人对他说：\Quote{有一位伟人，他能使你得自由；}他便微笑，欢喜，得意洋洋。但若有人对他说：\Quote{天主释放你，}他却因绝望而几乎冷淡。应许了必死之人的援助，你便欢喜；应许了不朽者的援助，你反而忧伤？应许你将要被一个也需要与你一同被释放的人释放，你便雀跃，仿佛得到了重大帮助；应许你那位不需要任何人释放他的释放者，你反倒绝望，仿佛这只是无稽之谈。这样的思想有祸了！它们远离正道；其中确实有悲惨而重大的死亡。靠近吧，开始渴望，开始寻求并认识那造你的那一位。因为祂不会离弃祂的工程，只要祂的工程不离弃祂。\par

5. ...\Quote{他的气息一断，他归于尘土；在那日，他的思念也归于灭亡}\BibleRef{咏 146:4}。傲慢在哪里？骄傲在哪里？夸耀在哪里？但也许他去了一个好地方，如果他真的过去了的话。因为我不知那如此说话的人往哪里去了。因为他傲慢地说话；我不知道这样的人往哪里去，只有查看另一篇圣咏，看到他们的逝去是不善的。\Quote{我见过恶人高耸如黎巴嫩香柏，我经过时，看，他已不在了；我寻找他，却寻不着他的处所。}那善人经过，找不到恶人，他到达了一个恶人不在的地方。因此，弟兄们，让我们都听从：弟兄们，天主的爱子们，让我们都听从；无论在任何患难中，无论对天上恩赐有何等渴望，\Quote{让我们不要信赖王侯，也不信赖世人，在他们内没有救恩。}这一切都是可朽的、短暂的、会消逝的。\par

那么，我们该做什么，如果我们不寄望于人子，也不寄望于王侯？我们该做什么？\BibleQuote{以雅各伯的天主为助佑的人是有福的}（\BibleRef{咏145:5}）：不是这个人或那个人；不是这个天使或那个天使；而是，\BibleQuote{以雅各伯的天主为助佑的人是有福的}：因为对雅各伯也是如此伟大的助佑，竟使他成为以色列。哦！强大的助佑！现在他是以色列，\Quote{看见天主}。当你们还在此地，身为流浪者，尚未看见天主时，若你们有雅各伯的天主为助佑，你们将从雅各伯成为以色列，将要\Quote{看见天主}，一切劳苦和呻吟都将终结，啃噬的忧虑将止息，幸福的赞美将接续。\BibleQuote{以雅各伯的天主为助佑的人是有福的}；这雅各伯。为何有福？同时，当还在今生呻吟时，\Quote{他的希望是在主、他的天主内}。……这位是谁，\Quote{主，他的天主}？……\BibleQuote{对我们来说，只有一个天主——圣父，万物出于祂；一个主耶稣基督，万物藉着祂。}（\BibleRef{格前8:6}）因此，愿祂成为你的希望，就是主、你的天主；愿你的希望在祂内。崇拜撒顿的人，他的希望也在主，他的神内；崇拜乃普敦或默丘利的人，他的希望也在主，他的神内；不仅如此，我还要说，崇拜肚腹的人，如经上所说：\BibleQuote{他们的神是肚腹。}（\BibleRef{斐3:19}）这个人的神是那个，那个人的是另一个。谁是这位\Quote{有福的}？因为\Quote{他的希望是在主，他的天主内。}但祂是谁？\BibleQuote{创造了上天、大地、海洋和其中的一切的那位}（\BibleRef{咏145:6}）。我的弟兄们，我们有一位伟大的天主；让我们赞美祂的圣名，因为祂屈尊使我们成为祂的产业。你们尚未看见天主；你们无法完全爱那尚未看见的。你们所看见的一切，都是祂造的。你赞叹世界，为何不赞叹世界的创造者？你仰望上天，惊奇不已：你思量大地，颤抖不已；你何时能容纳海洋的浩瀚？看看数不胜数的星辰，看看各种种子，各类动物，水中游的，地上爬的，天上飞的，空中翱翔的；这一切何其伟大，何其美丽，何其壮观，何其奇妙！看，创造这一切的，是你的天主。把你的希望寄托在祂内，好让你有福。\Quote{他的希望是在主，他的天主内。}我的弟兄们，请看这位伟大的天主，良善的天主，创造这一切的。……如果祂只提到这些东西，或许你会回答我：\Quote{天主创造了上天、大地和海洋，是一位伟大的天主：但祂会想到我吗？}会有人对你说：\Quote{祂创造了你。}怎么？我是上天，或是大地，或是海洋？显然，我不是上天，也不是大地，也不是海洋，但我在地上。至少你承认这一点，你在地上。那么请听：天主不仅创造了上天、大地和海洋，因为祂\Quote{创造了上天、大地、海洋和其中的一切}。如果祂创造了其中的一切，祂也创造了你。说\Quote{你}还不够；麻雀、蝗虫、蠕虫，祂没有不创造的，祂关心一切。祂的关怀不涉及祂的诫命，因为这条诫命祂只赐给了人。……就诫命的实质而言，\BibleQuote{天主不关心牛。}（\BibleRef{格前9:9}）就祂对宇宙的眷顾，祂创造万物并治理世界而言，\BibleQuote{上主，祢拯救人和兽。}这里或许有人会对我说：\Quote{天主不关心牛，}来自新约；\Quote{上主，祢拯救人和兽，}来自旧约。有些人挑剔说，这两部经彼此不合。……让我们听主自己，宗徒的元首和导师：\Quote{你们看，}祂说，\Quote{天空的飞鸟，它们不播种，也不收割，也不收藏在仓里，你们的天父养活它们。}（\BibleRef{玛6:26}）因此，即使在人类之外，这些动物也是天主所关心的对象，得以喂养，而不是接受法律。就赐予法律而言，\BibleQuote{天主不关心牛；}就创造、喂养、治理、主宰而言，万物都与天主有关。\BibleQuote{两只麻雀不是卖一个铜钱吗？}主耶稣基督说，\BibleQuote{但一只也不会不经过你们父的旨意掉在地上；你们比它们贵重多了。}（\BibleRef{玛10:29}）或许你说，在这庞大的数目中，天主不算我。接着福音里有奇妙的一节：\BibleQuote{你们连头发都被数过了。}（\BibleRef{玛10:30}）\par

6. \BibleQuote{谁永久保存真理。} 什么\Quote{永远的真理}？哪个\Quote{真理}是祂所\Quote{保存，}并且祂在哪里\Quote{永久保存}它？\BibleQuote{为受虐待的人伸冤的主}\BibleRef{咏 145:7}。祂为受冤屈的人报仇。宗徒的声音立刻传到你们耳中：\BibleQuote{为什么你们不宁愿受冤屈呢？}\BibleRef{格前 6:7}。祂劝你们不要受委屈，而要受冤屈：因为并非每一次委屈都是冤屈。因为凡合法所受的都不是冤屈；免得你或许会说：我也是那些受冤屈的人之一，因为我在某处因某事而受了某种苦。你要考虑你是否受了冤屈。强盗受很多苦，但他们不受冤屈。恶人、作恶者、入室行窃者、奸夫、淫棍，所有这些人都受很多苦，但都不是冤屈。受冤屈是一回事，受磨难、刑罚、委屈或惩罚是另一回事。要看你在哪里；看你做了什么；看你为什么受苦；然后你才能看清你受的是什么苦。义与不义是相反的。义就是正义。因为并非所有被称为义的都是义。如果有人为你制定了不义的法律呢？不义的事确实不能称为义。真正的义也就是正义。要看你做了什么，而不是你受了什么苦。如果你行了义，你受的就是冤屈；如果你行了恶，你受的就不是冤屈....\par

7. \BibleQuote{谁给饥饿的人食物。} 看，我不指望你们什么：\BibleQuote{天主给饥饿的人食物。} 谁是\Quote{饥饿的人}？所有人。\Quote{所有人}是什么意思？祂给所有有生命的东西、所有人食物：祂不为祂所爱的人预留一些食物吗？如果他们另有饥饿，他们也就另有食物。让我们先探询他们的饥饿是什么，然后我们将找到他们的食物。\BibleQuote{饥渴慕义的人是有福的，因为他们要得饱饫。}\BibleRef{玛 5:6}。我们应当是天主的饥饿者....\BibleQuote{主释放被囚的人；主扶起被摔倒的人；主使瞎子有智慧。}\BibleRef{咏 145:8}。祂藉着最后这句话完美地向我们解释了前面所有的：免得当祂说了\BibleQuote{主释放被囚的人，}我们可能将它理解为那些因犯罪而被主人用铁链捆绑的囚犯：并且当祂说了\BibleQuote{祂扶起被摔倒的人，}我们心目中会出现某个绊倒、跌落或从马上摔下来的人。有另一种跌倒，有另一种锁链，正如有另一种黑暗和另一种光明。然而祂说\BibleQuote{祂使瞎子有智慧；}祂不会说祂光照瞎子，免得你们也按肉身来理解，就如主用唾沫和泥抹在瞎子的眼睛上治好他那样：免得你们在祂谈论属灵事物时寻求此类东西，祂指向一种智慧之光，瞎子借此被光照。因此，正如瞎子被智慧之光照亮，同样，被囚者得释放，被摔倒的人被扶起。那么我们是如何被囚禁的？如何被摔倒的？我们的身体曾经是我们的装饰：现在，我们犯了罪，因此被加上了锁链。我们的锁链是什么？我们的必死性....\BibleQuote{主喜爱义人。}谁是\Quote{义人}？他们现在有多义？就和你一样；\BibleQuote{主保护外方人}\BibleRef{咏 145:9}。\Quote{外方人}就是陌生人。每一个外邦人的教会都是陌生人。因为它归入先祖们，不是出于他们的血肉，而是因效法他们成为他们的女儿。然而保护他们的是主，不是任何人。\BibleQuote{孤儿和寡妇祂必扶持。}不要有人认为祂扶持孤儿是为了他的遗产，或扶持寡妇是为了她的什么事务。的确，天主确实帮助他们；在人类的一切责任中，照顾孤儿、不离弃寡妇的人行的是善工：但在某种意义上，我们都是孤儿，不是因为我们的父死了，而是因为祂不在.. ..\par

8. \Quote{祂要铲除罪人的道路。} 什么是\Quote{罪人的道路}？就是嘲笑我们所说的这些事。谁是孤儿，谁是寡妇？有什么天国，有什么地狱的惩罚？这些都是基督徒的寓言。我只看眼前，就为眼前活着：\BibleQuote{我们吃喝吧，因为明天我们就要死了。} \BibleRef{格前 15:32} 当心不要让这类人说服你什么：不要让他们从你的耳朵进入你的心；让他们在你耳中遇到荆棘；让试图这样进入的人被刺穿而去，因为\BibleQuote{恶劣的交往败坏良好的品行。} \BibleRef{格前 15:33} 但你或许会说：\Quote{那么他们为何顺利呢？看，他们不敬拜天主，每日作种种恶事，却丰裕地拥有那些因缺乏而劳苦的东西。} 不要嫉妒罪人。他们得到了什么，你看见了；他们将来会怎样，你难道看不见吗？……你难道不信你的主、你的天主，祂说：\BibleQuote{通往毁灭的路是宽而大的，有许多人从那里走？} \BibleRef{玛 7:13} 这条\Quote{路主要铲除。} 当\Quote{罪人的道路}被\Quote{铲除}之后，给我们留下什么？\BibleQuote{来吧，我父所祝福的，来享受王国。} \BibleRef{玛 25:34} \BibleQuote{主要永远为王} \BibleRef{咏 145:10}。 \Quote{哦，熙雍，你的天主}要永远为王；当然，你的天主不会没有你而统治。\Quote{世世代代。} 他说了两次，因为他不能说\Quote{永远}。不要以为永恒能被有限的言语所界定。\Quote{永恒}这个词由四个音节组成；它本身是没有止境的。只能这样向你推荐：\Quote{世世代代。} 他说得太少了：如果他整天说，也嫌太窄：如果他一生说，难道他最终不能静默吗？热爱永恒：你将没有止境地统治，如果基督是你的终向，你与祂一起统治，永永远远。阿们。\par

\SourceRecord{Translated by J.E. Tweed.；Nicene and Post-Nicene Fathers, First Series；Vol. 8.；Edited by Philip Schaff.；Buffalo, NY: Christian Literature Publishing Co.,；1888.；<http://www.newadvent.org/fathers/1801146.htm>.}

% structure: p0001 source=h1 target=section reason=page-title

\section{对圣咏第一四七篇的诠释}

1. 曾对我们说：\Quote{你们要讚美上主}\BibleRef{咏 146:1}。这是对万民说的，不仅是对我们。这些话由读经者在各处发出声音，每个教会分别听到；但天主同一个声音向众人宣告，要我们讚美祂。好像我们问为什么应该讚美上主，看，祂提出了什么理由：\Quote{你们要讚美上主，}祂说，\Quote{因为圣咏是好的。}这就是所有讚美者的赏报吗？……\Quote{圣咏}就是讚美天主。所以祂说：\Quote{你们要讚美上主，因为讚美上主是好的。}我们不要就这样忽略对上主的讚美。话说过了，就过去了；事情做了，我们就沉默了；我们讚美了，然后就休息了；我们歌唱了，然后休息了。我们出去办理等待我们的某些事务，当其他工作找到我们时，对上主的讚美会在我们心中停止吗？不：你的舌头只讚美一时，让你的生命永远讚美。因此\Quote{圣咏是好的。}\par

2. 因为\Quote{圣咏}是一首歌，不是任何一种歌，而是用琴瑟唱的歌。琴瑟是一种乐器，如七弦琴和竖琴，以及这类为音乐发明的乐器。因此，那唱圣咏的，不仅用声音唱，而且另外用一种称为琴瑟的乐器，用手配合声音。那么，你要唱圣咏吗？不要只让你的声音发出讚美天主的声音；也要让你的行为与你的声音和谐……为了悦耳，你用声音唱；但用心不要缄默，用生命不要静止。你在心中不图谋欺诈：你就向天主唱圣咏。当你吃喝时，唱圣咏：不是混合着悦耳的甜美声音，而是节制、简朴、有度地吃喝；因为正如宗徒所说：\Quote{你们或吃或喝，或无论做什么，一切都要为光荣天主而做。}\BibleRef{格前 10:31}……如果你因无度的贪食，超出了自然的应有界限，在过量的酒中放纵自己，那么无论你的舌头发出多么伟大的对天主的讚美，你的生命却在亵渎祂。饭后饮后，你躺下睡觉：在你的床上不要行任何污秽之事，也不要超出天主法律所许可的范围；与你的妻子保持婚床的贞洁；如果你渴望生养子女，也不要有放纵的情欲：在你的床上尊敬你的妻子，\BibleRef{伯前 3:7}因为你们都是基督的肢体，都是祂所造的，都是祂的血所更新的；这样做，你就讚美天主，你的讚美也不会全然无声。当睡眠临到你时又如何？不要让邪恶的良心将你从安眠中惊醒：这样，你睡眠的无辜就讚美天主……\par

3. \Quote{让讚美成为我们天主的甘饴。}如何做到？如果祂被我们的善生所讚美。听，那时讚美才会令祂甘饴。在另一处说：\Quote{讚美在罪人口中是不合宜的。}\BibleRef{德 15:9}如果讚美在罪人口中不合宜，那么它也不甘饴，因为只有合宜的才是甘饴的……因为讚美可能使人甘饴，当他听到某人用巧妙优雅的辞藻和甜美的声音讚美时；但\Quote{让讚美成为我们天主的甘饴}，祂的耳朵不是向着口，而是向着心；不是向着舌头，而是向着那讚美者的生命。\par

4. 谁是\Quote{我们的天主}，使讚美成为祂的甘饴？祂使自己对我们甘饴，祂把自己推荐给我们；感谢祂的屈尊就卑……\Quote{但是天主把自己的爱显示给我们}……\Quote{当我们还是罪人的时候，基督就为我们死了。}\BibleRef{罗 5:8}……让我们看看，是否如宗徒所说，基督为罪人和不敬虔的人死了：\Quote{那建造耶路撒冷、聚集以色列被赶散的人的上主}\BibleRef{咏 146:2}。因为耶路撒冷的百姓就是以色列百姓。它是\Quote{天上永恒的耶路撒冷}，天使也是那里的公民……因此，那座城的所有公民，通过\Quote{看见天主}，在那伟大、广阔、天上的城中喜乐；他们凝视着天主本身。但我们却是从那城中流亡出来的，因罪恶被驱逐，不能再留在那里；被必朽性所压，不能再回去。天主眷顾了我们的流亡，祂\Quote{建造耶路撒冷}，恢复了那已倒塌的部分。祂怎样恢复那倒塌的部分呢？……祂于是派遣祂的儿子作为救赎者来到我们被掳的境地。祂说：你拿一个袋子，里面装着赎俘虏的代价。因为祂穿上了我们必朽的血肉，其中就有宝血，藉着流这血，我们得以被赎。用那血，祂\Quote{聚集了以色列被赶散的人}。如果祂聚集了从前被赶散的人，那么当现在被赶散的人，我们该如何努力使他们被聚集呢？如果被赶散的人已被聚集，好在建造者手中被塑造成建筑物，那么那些因不安而从建造者手中掉落的，又该如何被聚集呢？看，我们讚美的是谁；看，我们一生所欠的讚美是归于谁的。\par

5. 祂如何聚集？祂聚集时做什么？\BibleQuote{祂治愈心灵破碎的人}\BibleRef{咏 146:3}。请看以色列流散者如何被聚集：借着治愈心灵破碎的人。那些心不破碎的人，不被治愈。什么叫作破碎心灵？弟兄们，要明白，要去做，这样你们才能被治愈。因为在圣经许多地方都说过：\Quote{天主的祭献是忧苦的精神，破碎和痛悔的心，天主必不轻视。}祂就这样治愈心灵破碎的人，因为祂亲近他们以治愈他们；正如另一处所说：\Quote{主亲近那些破碎自己心灵的人。}谁有\Quote{破碎的心}？谦卑的人。谁没有\Quote{破碎的心}？骄傲的人。破碎的心将被治愈，膨胀的心将被击碎。或许正是为此被击碎，好使它因破碎而得以痊愈。所以，弟兄们，不要让我们的心在尚未正直之前就渴望被扶正。未经矫正就升高，是不好的……\par

6. 祂用什么方法来\Quote{包扎他们的伤口}？正如医生包扎骨折。因为有时（请留意这一点，亲爱的；那些观察过的人，或从医生那里听说过的人都知道），有时当肢体完好，但弯曲畸形时，医生会折断它们以便矫正，造成新的伤口，因为原本完好的形状是歪的……\par

7. 祂用什么方法包扎？就是今生的圣事，在此期间我们获得慰藉；而所有我们对你们说的话，那些有声音且会过去的话语，以及现今在教会中所做的一切，都是\Quote{祂包扎我们伤口}的方法。正如当肢体完全痊愈时，医生取下绷带；同样，在我们自己的城市耶路撒冷，当我们与天使平等时，你们认为我们会在那里接受我们在这里所接受的吗？那时还需要为我们朗读福音，使我们的信德得以持守吗？或者需要任何主教为我们覆手吗？这一切都是包扎骨折的方法；当我们获得完全的痊愈时，它们将被取下；但如果我们没有经过包扎，我们永远无法达到痊愈。\par

8. \BibleQuote{祂数算星辰的数目，一一称呼它们的名字}\BibleRef{咏 146:4}。对天主来说，\Quote{数算星辰的数目}算得了什么呢？连人都曾试图这样做；他们能否成功，是他们的事；他们若不是认为自己能做到，就不会尝试。我们且不论他们能做什么，达到了什么程度；我认为对天主来说，数算所有星辰并非什么大事。难道祂需要一一数过，以免忘记吗？由祂来数算星辰，其中\BibleQuote{连你们的头发也都数过了}\BibleRef{玛 10:30}，这算得了什么呢？这些星辰是教会中的光体，安慰我们的黑夜；宗徒论及这一切说：\BibleQuote{在弯曲悖谬的世代中，你们显在这世代中，好像明光照耀，将生命的话表明出来}\BibleRef{斐 2:15}。这些星辰，天主都在数算；所有将与祂一同为王的人，所有将聚集在祂独生子身体里的人，祂已数算，仍在数算。不配的人，甚至不被数算。也有许多人相信，或者不如说，以一种信仰的虚假外表，依附于祂的子民；但祂知道祂数算什么，祂簸扬掉什么。因为福音是如此崇高，以至于正如经上所说：\Quote{我宣扬过，我说过，他们多过数不清。}所以人民中可以说有一些超额者。我所说的超额者是什么意思？就是比将在那里的人数更多。在这围墙之内，比将进入天主的国、天上的耶路撒冷的人更多；这些人超出了数目。你们每个人都要省察自己是否在黑暗中发光，是否拒绝被世界的黑暗不义引入歧途；若他不被迷惑，也不被战胜，他就好比一颗星辰，天主已经将他数算。\Quote{一一称呼它们的名字}，祂说。这便是我们的全部赏报。我们可以在天主那里有某些名字，使天主知道我们的名字，这是我们应当渴望的，为此而行动，为此而尽力；不要因其他事物而欢喜，甚至不要因某些神恩而欢喜……当门徒们完成任务回来，欢欣地说：\BibleQuote{主！因祢的名号，连魔鬼都服从我们}\BibleRef{路 10:17}——那时，祂（知道许多人会说\Quote{我们不是因祢的名号驱过魔鬼吗？}，祂将对他们说\Quote{我不认识你们}）说：\BibleQuote{不要因为魔鬼服从你们而欢喜，但应当欢喜，因为你们的名字记录在天上}\BibleRef{路 10:20}。\par

9. \BibleQuote{我们的主伟大} \BibleRef{咏 146:5}。圣咏作者充满喜乐，他奇妙地倾吐了言语：然而他有些无法言说，又怎能思量呢？\BibleQuote{他的能力伟大，他的智慧不可测量。}他\BibleQuote{数点星辰}，自己却不可数。谁能解释这一点？谁能恰当地想象\BibleQuote{他的智慧没有数字}的含义？……因此，凡这世界所含的无限事物，虽对人无限，对天主却不然：说对天主太小了：甚至天使也能数算。他的智慧超越一切计算者；我们无法数算。数字本身谁能数算？那么在天主那里有什么？他造了一切，用何物，在何处？对祂曾说：\BibleQuote{你以度量、数目和重量安排了万物} \BibleRef{智 11:20}。或者谁能数算、度量、称量度量、数目和重量本身——天主以此安排了万物？因此，\BibleQuote{他的智慧没有数字。}让人的声音静默，人的思想安静：不要伸展到不可理解的事物，好像能理解，而要仿佛要分享它们，因为我们将是分享者……那么我们将是分享者：无人怀疑：圣经如此说。我们将分享什么，好像天主内有这些部分，好像天主被分成部分？谁能解释多少人分享一个单一实体？不要追问我认为你们看到无法恰当言说的事：而要回到救主的医治，破碎你们的心。祂将引导它，包扎它破碎之处，使它完全健康；那时那些对我们不可能的事将不再不可能。因为谁渴望达到天主性，承认软弱是好的。\par

10. \BibleQuote{主提拔温良的人} \BibleRef{咏 146:6}。例如：你不明白，你不理解，不能达到：要尊敬天主的圣经，尊敬天主的圣言，即使不清晰：以敬畏等待理解。不要放肆地指责圣经的隐晦或表面矛盾。其中毫无矛盾：有些隐晦，并非为了拒绝你，而是为锻炼那将要领受的人。当它隐晦时，这是医生的作为，要你叩门。祂愿你通过叩门得到锻炼；祂愿祂在你叩门时为你开门。通过叩门你将受锻炼；受锻炼，你将扩大；扩大，你就能容纳所赐的。不要因关闭而愤怒；要温和，要柔顺。不要踢踹黑暗，也不要说：若这样说，更好。因为你怎么能这样说，或判断怎样说才适宜？它正如所适宜的那样说的。病人不要试图修正他的药方：医生知道如何调配：相信那照顾你的祂。因此接下来是什么？……\BibleQuote{主提拔温良的人，却将罪人贬抑至地。}他从前面提到的温良，意指某类罪人。因此此处罪人，我们理解为凶暴者，以及不温良者。祂为何\BibleQuote{将他们贬抑至地}？他们挑剔可理解的对象，他们只感知属世的事物。\par

11. \BibleQuote{以宣认向主开始} \BibleRef{咏 146:7}。以此开始，若你愿达到清晰理解真理。若你愿从信德之路被带到真实之宣认，\BibleQuote{以宣认开始。}先控告自己：控告自己，赞美天主。宣认之后呢？让善行跟随。\BibleQuote{弹着竖琴歌颂我们的天主。}什么是\BibleQuote{弹着竖琴}？正如我解释过的，如同咏于瑟琴，竖琴也是如此：不仅以声音，也以行动。\par

12. ...\Quote{谁以云彩遮盖高天，为大地预备雨水} \BibleRef{咏 146:8}。现在你惊慌，因为你不能看见天：下雨后你将收获果实，并将看见晴朗的天空。或许我们的天主正是如此做了。因为若不是圣经的晦暗给了我们机会，我们本不会向你们说那些使你们喜乐的事。那么这或许就是你们喜乐的雨水。若不是天主用预像的云彩遮盖了圣经的天，我们的口舌就不可能向你们表达它。为此，祂愿意先知的话晦暗不明，为使天主的仆人日后能有藉以阐释的，从而流向人的耳朵和心灵，使他们从天主的云彩中领受神性喜乐的丰腴。\Quote{祂使草生长在山上，又为人的服务生出菜蔬。}看哪，这就是雨水的果实。经上说：\Quote{祂使草生长在山上。}难道草就不长在低洼之处吗？是的，但草长在\Quote{山上}是一件大事。……因为没有什么比坚硬的山更贫瘠了。\Quote{又为人的服务生出菜蔬。}什么\Quote{服务}？请听保禄自己说的。他说：\Quote{我们自己，为耶稣基督的缘故，是你们的仆人。} \BibleRef{格后 4:5} 那位说：\Quote{我们若给你们播下了神性的事物，若收获你们物质的事物还算大事吗？}却说自己是一个\Quote{仆人}。因为我们是你们的仆人，弟兄们。我们中谁也不可自夸，好像胜过你们。我们若更谦卑，就会更伟大。\Quote{但谁若愿意在你们中间成为伟大的}（这是主的话），\Quote{就应当作你们的仆人。} \BibleRef{玛 20:26} 宗徒保禄，的确，靠自己的双手生活，甚至拒绝接受\Quote{山上的草}；他选择了匮乏；然而，山还是给出了\Quote{草}。因为他选择不接受，难道山就不该给出，从而保持荒芜吗？果实归雨水，食物归仆人，正如主说：\Quote{你们吃他们所供给的：}并且他们不该以为自己给了什么属于自己的东西，祂又说：\Quote{因为工人配得他的工资。} \BibleRef{路 10:7}\par

13. ...刚才读到了：\Quote{凡求你的，就给他} \BibleRef{路 6:30}；而在另一处，圣经说：\Quote{让你的施舍在你手中出汗，直到你找到一个义人可以给他。}有一个人来找你，另一个人你应当去寻找。确实，不要使来找你的人空手而回，因为\Quote{凡求你的，就给他}；然而，还有另一个人你应当去寻找：\Quote{找一个义人，把施舍给他。}除非你从你的财产中，按照各人家庭的需要，如同给国库缴纳的固定款项，划出一部分，你就永远不会这样做。如果基督没有自己的国库，祂也没有库房。……那么，从你的年收入或日常收益中，砍下并修剪出一个固定的数额，否则你似乎是在从资本中给予，你的手必踌躇，当你伸向你没有许诺的东西时。从你的收入中划出一部分；若你愿意，十分之一，尽管那也很少。因为经上说，法利塞人捐了十分之一：\Quote{我每周禁食两次，凡我所得的都捐献十分之一。} \BibleRef{路 18:12} 主说了什么？\Quote{除非你们的义德超过经师和法利塞人的义德，你们决不能进入天国。} \BibleRef{玛 5:20} 那位你应当超过其义德的人，捐了十分之一；你连千分之一也不捐。既然你比不上他，怎能超过他呢？\Quote{谁为大地预备雨水。}\par

14. \Quote{祂赐给牲畜他们的食物} \BibleRef{咏 146:9}。祂所说的牲畜，就是天主的羊群。天主不会借着人，亏欠祂羊群的食物，因为\Quote{祂使草生长，是为了人的服务。} \Quote{也赐给啼叫的幼鸦。}难道我们该如此想，以为乌鸦啼叫向天主求食？不要以为无理性的受造物会呼求天主：没有受造物知道如何呼求天主，惟独有理性的受造物才知道。你要把这当作比喻来听，免得你像某些恶人所说，人的灵魂会转入牲畜、狗、猪、乌鸦之中。不要让这在你心中或信仰中占一席之地。人的灵魂是按照天主的肖像造的：祂不会把自己的肖像给狗或猪。谁又是\Quote{幼鸦}？以色列人惯常说，只有他们是义人，因为法律赐给了他们；其他所有民族的每一个人，他们常称为罪人。事实上，所有民族都沉沦于罪恶、偶像崇拜、崇拜石头和木头：但他们是持续如此吗？虽然乌鸦本身，我们的祖先，并没有如此，但我们，\Quote{幼鸦}，却呼求天主。\BibleRef{伯前 1:18}……因为\Quote{幼鸦}，他们似乎敬拜祖先的偶像，已经前进，转向了天主。现在你听见\Quote{幼鸦}呼求唯一的天主。那么，你向\Quote{幼鸦}说：\Quote{你离开了你的父亲吗？}的确，我离开了，他说；因为他是乌鸦，不呼求天主。我，\Quote{幼鸦}，却呼求天主。\par

15. \BibleQuote{祂不喜爱马的力大} \BibleRef{咏 146:10}。\Quote{马的力大}就是骄傲。因为马似乎天生适合把人类高高驮起，使人在行进中越发高升。事实上，它的颈项正象征着一种骄傲之情。人不可因自己的价值而自高，不可因自己的尊荣而自视崇高；他们要小心，免得被一匹未驯服的马抛落…… \Quote{也不喜爱人的帐幕。} 主的帐幕也就是圣教会，遍及普世。异端者脱离教会的帐幕，为自己搭起了帐幕。因为如果有哪位善良虔诚、承认自己的软弱、是\Quote{呼求天主的乌鸦的雏}，却未能享有世俗的荣耀，他不会离开教会，不会在教会之外为自己搭起一座天主不会喜爱的帐幕。但他说什么？\Quote{我宁愿被弃于天主的殿中，也不愿住在恶人的帐幕里。}\par

16. 但他说了什么？\BibleQuote{上主喜爱敬畏祂的人和仰望祂慈爱的人。} \BibleRef{咏 146:11}。强盗令人恐惧，野兽令人恐惧，不义而有权势的人更令人恐惧。\Quote{上主喜爱仰望祂慈爱的人。} 看，出卖我主的\Person{犹达斯}也恐惧了，但他没有仰望祂的慈爱。……你恐惧确实是好的，但只要你信靠你所恐惧的那一位的慈爱。他绝望了，\Quote{出去吊死了。} 因此，你要这样敬畏上主，使你在祂的慈爱中抱有希望。\par

17. \BibleQuote{耶路撒冷，你要齐声赞美你的天主。} \BibleRef{咏 147:12}。他们仍在为奴之地，却看见那些羊群，更确切地说，是那唯一的一群，所有城民从四方汇集到那座城里；他们看见群众的喜乐，在打谷和簸扬之后，已收藏在仓里，不再惧怕，不再经受劳苦和磨难；而他们自己仍在这里，在打谷之中，将希望的喜乐寄予未来，并渴慕它，仿佛将自己的心与天主的众天使，以及与那群将永远与他们同享喜乐的子民联合在一起。因为那时你要做什么呢，\Place{耶路撒冷}？诚然，劳苦和叹息都将过去。你要做什么？你会犁地、播种、栽种葡萄、航行、经商吗？你要做什么？你仍然有义务从事你现在所做的善工，这些善工虽好，且出于怜悯吗？数一数你的同伴，环顾你四周的团体：看看有没有人饥饿，需要你给面包吃；有没有人口渴，需要你给一杯凉水喝；有没有人是旅客，需要你接待；有没有人患病，需要你探望；有没有人争吵，需要你劝和；有没有人临终，需要你埋葬。那么，你要做什么？\Quote{耶路撒冷，你要齐声赞美你的天主。} 看，这就是你的事业。正如铭文上常说的：\Quote{用吧，并快乐。}\par

18. 你要成为\Place{耶路撒冷}；记住这话是对谁说的：\Quote{主啊，在祢的城市里，祢将他们的形象化为虚无。} 这些人就是现今为这样的虚荣而快乐的人；他们中有人今天没有来，因为今天有表演。这是给谁的礼物？这是谁的损失？或者，为什么是礼物？为什么是损失？因为不仅举办这些表演的人遭受损失，而且以观看为乐的人遭受更大的损失。前者钱袋中的金子被掏空，后者胸中的义德财富被掠夺。多数举办表演的人哀叹变卖产业；罪人该当怎样哀叹失去灵魂呢？难道主在主日曾呼喊\Quote{你们要惊醒}，就是为了让众人今天如此惊醒吗？我恳求你们，\Place{耶路撒冷}的公民们，我因\Place{耶路撒冷}的和平恳求你们，因\Place{耶路撒冷}的救赎主、建设者、统治者恳求你们，请你们为这些人向天主祈祷。愿他们看见，愿他们感到自己是在荒废时光；他们虽然专注于取悦他们的景象，但最终要看看自己，并对自己不满。因为我们为许多人已经这样做了而欢喜；而且我们自己也曾经坐在那里，并且疯狂。现在有谁正坐在那里，我们想，将来他们不仅会成为基督徒，而且会成为主教！从过去我们推测未来；从已经发生的，我们预先宣告天主将做的事。愿你们的祈祷保持惊醒，你们的叹息不会落空。那些已经逃脱的人，为仍在危险中的人祈祷，因为他们也曾置身于危险中，他们的祈祷被垂听；天主必将祂的子民从巴比伦的掳掠中领出来；祂必定会救赎和拯救他们，那批背负天主肖像的圣徒数目将达至圆满。……\BibleQuote{我们虽多}，宗徒说，\BibleQuote{在基督内却是一个身体。} \BibleRef{格前 10:17} 因此，我们既是许多人，\Quote{就齐声赞美；}既然我们是一个，\Quote{就赞美。}众是一又是多，因为他们在其中合而为一的那一位永是独一。\par

19. 因此，耶路撒冷说，我为何同声赞颂主，如同熙雍，赞颂我的天主？耶路撒冷与熙雍是同一处。此二名因不同理由而有。耶路撒冷意为\Quote{和平的异象}；熙雍意为\Quote{守望}。看看这些词语是否不像异象？免得外邦人以为他们有异象而我们没有。有时，在剧院或露天剧场散场后，当失丧的人群开始从那巢穴中被呕吐出来时，有时，他们心中保留着那虚无娱乐的形象，以不仅无用而且有害之物喂养记忆，以之为甘甜而欢喜，实则致命；他们常常，也许，看到天主的仆人经过，凭其外衣或头饰认出他们，或凭面貌认识他们，便彼此，或心中说：\InnerQuote{可怜的人，他们损失了多少啊！}弟兄们，让我们以向主的祈祷回报他们的善意（因他们确是好意）。他们愿我们好；但\Quote{喜爱不义的人，憎恨自己的灵魂。}若他憎恨自己的灵魂，怎能爱我的灵魂呢？然而，以那乖谬、空洞、虚妄的善意——若真可称为善意——他们为我们失去他们所爱的而悲伤：让我们祈求他们不失去我们所爱的。看哪，那座耶路撒冷将是何种品性，它受劝勉赞美，或更说预见它将赞美。因为那城市的赞美，当我们看见、爱慕并赞美时，将不需被先知之声催促和激发；但现在先知如此说，是为了在他们尚存肉身时尽可能汲取未来真福的喜乐，然后传入我们耳中，以激发我们对那城市的爱慕。让我们热切渴望，不要心灵懈怠。\Quote{赞美你的天主，哦，熙雍。}\par

20. 他说：\Quote{祂坚固了你城门的门闩}\BibleRef{咏 147:13}。坚固门闩并非为敞开的门，而是关闭的门；因此大多数抄本写道：\Quote{祂坚固了你城门的门闩。}请留心，亲爱的诸位。祂命令关闭的耶路撒冷赞颂主。我们现今同声赞颂，现今赞颂；但这是在许多冒犯之中。许多我们所不愿的人进入；许多我们所不愿的人出去：因此冒犯频繁。\Quote{由于罪恶增多，}真理说，\Quote{许多人的爱情冷淡了。}\BibleRef{玛 24:12}因我们无法分辨进入的人，因我们无法留住出去的人。为何如此？因为尚未有圆满，尚未有那未来的福乐。为何如此？因为此时仍是打谷场，并非谷仓。那么那时将有何事，除了不怕发生此类之事？祂不仅说\Quote{祂设立了}，而且说\Quote{祂坚固了你城门的门闩。}不要让任何人出去，不要让任何人进来。不让任何人出去，我们欢喜；不让任何人进来，我们害怕。不，不要怕此事：当你进入时，将被告知：你只要在那些携带油灯的童女之列。\par

21. \Quote{祂在你内祝福你的儿女。}谁？祂\Quote{使你的边界成为和平。}你们何等欢欣！请爱和平，我的弟兄们。当和平之爱从你们心中喊出时，我们大大喜悦。它使你们何等喜悦！我什么也没说：什么也没解释：我只读了这节经文，而你们呼喊。在你们内呼喊的是什么？是和平之爱……天国的儿女，耶路撒冷的居民啊，在耶路撒冷有和平的异象：所有爱和平的人在她内蒙祝福，当门关闭、门闩坚固时，他们进入。这，你们只听到名字就如此爱慕、珍视的，你们要追随，要渴望：要在你们的家中、事务中、妻子中、儿子中、奴隶中、朋友中、仇敌中爱此和平。\par

22. 你们方才一提到和平就呼喊，你们因渴望而呼喊：你们的呼喊出于干渴，而非饱足；因为完全的正义将在完全的和平之处。现在我们饥渴慕义。\Quote{他们将得饱饫。}\BibleRef{玛 5:6}他们如何得饱饫？当我们到达和平之时。因此，当祂说：\Quote{祂使你的边界成为和平，}因那里有饱足、没有匮乏，祂立刻加上\Quote{并用麦子的脂油充满你}\BibleRef{咏 147:14}。\par

23. \Quote{祂向大地发出祂的圣言}\BibleRef{咏 147:15}。看哪，我们在地上辛劳、疲倦、软弱、迟缓、寒冷；我们何时能被提升到那满足的麦子脂油，若非祂向大地发出祂的圣言——我们被它压向大地，我们被它阻挠而无法返回？祂发出祂的圣言，祂即使在旷野中也没有抛弃我们，祂从天上降下玛纳。\Quote{祂向大地发出祂的圣言}；祂的圣言来到大地。如何？祂的圣言是什么？\Quote{祂的圣言甚至跑到急速。}祂不是说：\Quote{祂的圣言是快速的，}而是：\Quote{祂的圣言甚至跑到急速。}让我们明白，我的弟兄们：祂不能选一个更好的词语。热者因热而热，冷者因冷而冷，快者因快而快……那么它跑到何种程度？\Quote{甚至到急速。}你尽可增加圣言的速度，说它像此或彼那样快，如飞鸟、如风、如天使；其中有任何一个与急速本身一样大吗？\Quote{甚至到急速}？什么是急速本身，弟兄们？它是无处不在的；它不是局部性的。这属于天主的圣言：不在局部，而是祂自己作为圣言遍及各处，祂就是\Quote{天主的德能、天主的智慧}\BibleRef{格前 1:24}，在祂取得肉身之前。若我们思想在天主形式内的天主，那与圣父平等的圣言，这就是天主的圣言，关于它说：\Quote{它从一端到另一端，强有力地运行。}\BibleRef{智 8:1}何等强大的速度！\Quote{它从一端到另一端，强有力地运行。}\par

24. 我們既被這寒冷身體的遲鈍，以及這屬地而可朽壞的生命之束縛所重壓，難道就沒有望德領受那\Quote{聖言}，它\Quote{迅速地奔馳}嗎？或者，雖然我們因身體而被壓到最深處，祂就拋棄了我們嗎？在我們出生於這可朽而遲鈍的身體之前，難道祂不是已經預定了我們嗎？那麼，這位預定我們的天主，就將雪賜給了大地，也就是賜給了我們自己。因為現在讓我們來看看聖詠中那些較為隱晦的經句，讓那些糾結開始被解開。看，我們在這地上是遲鈍的，彷彿在這裡凍結了。正如雪花所發生的：它們在高處凝結，然後落下；同樣，當愛德變冷時，人的本性就墮落到這地上，並陷入遲鈍的身體中，變得像雪一樣。但在那雪中，有預定的天主子民。因為，\BibleQuote{祂賜雪如羊毛}（\BibleRef{咏147:16}）。\Quote{如羊毛}是什麼意思？意思是，對於祂所賜的雪，這些在靈性上仍然緩慢而寒冷、祂所預定的人，祂將要有所作為。因為羊毛是衣物的材料：當我們看到羊毛時，我們視之為某種衣物的預備。因此，既然祂預定了這些目前寒冷、在地上爬行、尚未因愛德之火而發熱的人（因為祂仍在談論預定），天主就將他們賜下作為某種羊毛：祂將用他們做一件衣物。基督的\BibleQuote{衣服在山上發光如雪}（\BibleRef{玛17:2}）。基督的衣服如雪般發光，好像用那雪已經做了一件衣服：而那羊毛，也就是祂如羊毛般賜下的雪，他們既被預定，本是遲鈍的；但請等待，看接下來的事。既然祂將他們賜下作為羊毛，一件衣服就從他們製成。因為正如教會被稱為基督的身體，同樣，教會也被稱為基督的衣服：由此而來宗徒所說的：\BibleQuote{為將她呈獻給自己，做一個光榮的教會，沒有污點或皺紋}（\BibleRef{弗5:27}）。那麼，願祂將她呈獻給自己，做一個光榮的教會，沒有污點或皺紋；願祂用那在雪中所預定的羊毛為自己做一件衣服。當人們仍然不信、寒冷、遲鈍時，願祂用這羊毛做一件衣服。願它被洗去污點，藉著信德潔淨；願它沒有皺紋，願它被伸展在十字架上……\par

25. \BibleQuote{祂散佈雲霧如灰燼}。\Quote{\InnerQuote{祂散佈，}聖詠作者說，\InnerQuote{雲霧如灰燼}}。誰？這位\BibleQuote{祂賜雪如羊毛}者。因為祂所預定的人，祂就召叫他們悔改；因為\BibleQuote{祂所預定的人，祂也召叫了}。但\Quote{灰燼}與悔改相連。請聽祂召叫人悔改，當祂責備某些城市時說：\BibleQuote{禍哉，你，苛辣匝因！禍哉，你，貝特賽達！因為在你們那裏所行的奇能，如果行在提洛和漆冬，他們早已披麻蒙灰悔改了}（\BibleRef{玛11:21}）。因此，\Quote{祂散佈雲霧如灰燼}。\Quote{祂散佈雲霧如灰燼}是什麼意思？當一個人被召叫來認識天主，並對他說：\Quote{接受真理}；他開始願意接受真理，卻不能；他發現自己處於某種黑暗之中，這是他先前所看不見的……不要在雲霧中徘徊，要跟隨信德。但既然你竭力去看卻不能，就為你的罪悔改吧，因為雲霧如灰燼般被散佈。你現在要後悔曾頑固地反對天主，後悔跟隨了你自己的邪道。你已進入這種境界，其中你難以看見有福的景象，而這雲霧對你將是有益的，因為天主如灰燼般散佈它。你自己仍是一團雲霧，但如灰燼一般。因為那些悔改的人，我的弟兄們，仍在灰燼中翻滾，如同作證他們與灰燼相似，向天主說：\Quote{我是灰燼。}因為某段經文說：\Quote{我輕視了自己，並消耗殆盡，我把自己當作塵土和灰燼。}這就是悔改者的謙卑。當亞巴郎與他的天主說話，希望索多瑪的焚燒被啟示給他時，他說：\BibleQuote{我不過是塵土和灰燼}（\BibleRef{创18:27}）。這種謙卑在偉大而神聖的人身上是如何被發現的啊！\par

26. \BibleQuote{他抛下他的\OriginalTerm{水晶}{crystal}像面包屑}\BibleRef{咏147:17}。我们无需再费力解释什么是水晶。我们已经说过，相信你们，亲爱的，不会忘记。那么，\Quote{他抛下他的水晶像面包屑}是什么意思？什么是\Quote{水晶}？它很坚硬，凝结得非常紧密；不能像雪那样容易融化。雪经过多年累积和时代传承，变得坚硬，就称为\Quote{水晶}，并且\Quote{他抛下像面包屑}。这是什么意思？他们太硬了，不再适合比作雪，而是比作水晶；但他们也被预定和蒙召，甚至其中一些能养活别人，对他人也有用。还需要列举许多我们碰巧认识的人吗？这个和那个？每个人思考时都能回忆起他所认识的一些人曾经是多么顽固和固执，如何反抗真理；但现在他们宣讲真理，他们成了面包屑。那唯一的面包是谁？\Quote{我们虽多，}宗徒说，\Quote{在基督内是一个身体}\BibleRef{罗12:5}；他也说，\Quote{我们虽多，只是一个饼，一个身体}\BibleRef{格前10:17}。因此，如果基督的整个身体是一个饼，基督的肢体就是面包屑。他从一些坚硬的人中制作自己的肢体，使他们对喂养他人有用……看，宗徒保禄曾是水晶，坚硬，抗拒真理，对福音高声喊叫，好像把自己硬化以对抗太阳……既然他是水晶，他看起来清澈洁白，但坚硬且非常冰冷。他如何明亮洁白？\Quote{希伯来人中的希伯来人；就法律说，是法利塞人。}看水晶的光亮。现在听水晶的坚硬。\Quote{就热诚说，迫害教会}\BibleRef{斐3:5}。在圣斯德望殉道者被石头砸死的人中，他是坚硬的，也许比所有都硬。\Quote{因为他看守所有投石者的衣服}\BibleRef{宗22:20}，所以他借众人的手投石。\par

27. 于是我们看到\Quote{雪、雾、\OriginalTerm{水晶}{crystal}}：幸好他吹气融化它们。因为他若不吹气，若不亲自融化这冰的坚硬，\Quote{在他面前，谁能抵挡他的寒冷？}他舍弃一个罪人，看，他不召叫他；看，他不开启他的悟性；看，他不倾注恩宠；让人自己从愚昧的冰中融化吧。他不能。为什么不能？\Quote{在他面前，谁能抵挡他的寒冷？}看他越发坚硬，说：\Quote{我这个人真不幸！谁能救我脱离这该死的肉身？}看，我变冷了，看，我变硬了，什么热力能融化我使我奔跑？\Quote{谁能救我脱离这该死的肉身？……在他面前，谁能抵挡他的寒冷？}如果天主舍弃他，谁能自救？谁来解救？\Quote{天主的恩宠，藉着我们的主耶稣基督}\BibleRef{罗7:24}。那么我们要绝望吗？决不。因为接着：\BibleQuote{他发出他的言语，融化它们}\BibleRef{咏147:18}。那么，雪不要绝望，雾不要绝望，水晶也不要绝望。因为雪如同羊毛，正被制成衣服。那雾在悔改中找到安全：因为\Quote{他所预定的人，也召叫了他们。}但即使他们是被预定中最硬的，长期硬化，成了水晶，在天主的仁慈面前也不会坚硬。\BibleQuote{他发出他的言语，融化它们。}\Quote{融化}是什么意思？不要误解\Quote{融化}为贬义：意思是，他将使之液化，融化它们。因为他们因骄傲而坚硬。骄傲被称为迟钝是恰当的：因为凡迟钝的，也是冷的……对水晶也不要绝望。听水晶的话：\Quote{我从前是亵渎者、迫害者和施暴者}\BibleRef{弟前1:13}。但天主为何融化水晶？为叫雪不自暴自弃。因为他说：\Quote{我蒙受了怜悯，为叫耶稣基督在我身上显示他的全部忍耐，为给将来信靠他得永生的人作榜样}\BibleRef{弟前1:16}。于是天主召唤外邦人：\Quote{融化吧，水晶；来吧，雪。}\Quote{他的神吹动，水就流动。}看，\Quote{水晶}和\Quote{雪}融化了，变成了水，\Quote{凡口渴的，请来喝吧}\BibleRef{若7:37}。扫禄，坚硬如水晶，迫害斯德望至死；保禄，如今在活水中\BibleRef{若4:14}，召唤外邦人到泉源……\par

\BibleQuote{向雅各伯宣布祂的圣言，向以色列宣布祂的义和祂的典章} \BibleRef{咏 147:19}。何为\Quote{义}，何为\Quote{典章}？因为在此之前，当人类经历\Quote{雪}、\Quote{雾}和\Quote{冰}时，凡人类所遭受的，都是因其骄傲和对抗天主所应得的报应。让我们回到我们堕落的根源，看看圣咏中唱得何等真切：\Quote{我受困苦以前，曾经失足。}但说\Quote{我受困苦以前，曾经失足}的那一位也说：\Quote{祢压迫了我，为叫我能学习祢的诫命，这为我有益。}雅各伯从天主学习了这些义行，天主使他与一位天使搏斗；在那位天使的形像下，天主亲自与他搏斗。他抓住了祂，用力抓住祂，他得胜而抓住了祂：祂出于慈悲，而非软弱，让自己被抓住。因此雅各伯搏斗并得胜：他抓住了祂，当他似乎已征服了祂时，却求祂祝福。他如何明白与谁搏斗、抓住了谁？为什么他用力搏斗并抓住祂？因为\BibleQuote{天国是以猛力夺取的，以猛力夺取的人就攫取了它。} \BibleRef{玛 11:12} 那么他为何搏斗？因为要付出劳苦。我们为何难以持守那轻易失去的？免得我们轻易取回所失的，就学会失去我们所持守的。让人劳苦去持守：他只有劳苦之后才持守的，他会牢固地持守。因此，天主向雅各伯和以色列启示了这些祂的典章……\par

\BibleQuote{祂没有这样对待所有民族} \BibleRef{咏 147:20}。不要让人欺骗你们：这审判并非向任何民族宣布的，即义人与不义的人如何受苦，所有人如何因自己的功过而受苦，义人本身如何因天主的恩宠而非自己的功德而得释放。这并不是向所有民族宣布的，而只向雅各伯，只向以色列。那么，如果祂没有向所有民族宣布，而只向雅各伯和以色列，我们又当如何？我们将在哪里？在雅各伯内。\BibleQuote{祂没有向他们启示祂的典章。}向谁？向所有外邦人。那么当冰融化时，那些\Quote{雪}如何被召叫？外邦人如何被召叫，既然保禄成了义人？除了在雅各伯内，如何得救？野橄榄枝从原树干砍下，被接在橄榄树上：现在他们属于橄榄树，不再应被称为外邦人，而是在基督内的一个民族，雅各伯的民族，以色列的民族……什么是以色列？\Quote{看见天主。}他在哪里看见天主？在和平中。什么和平？耶路撒冷的和平；因为祂说：\BibleQuote{祂将和平置于你的边界上。}在那里我们要赞美：在那里我们众人将成为一体，在独一者内，归于独一者：因为那时，虽然众多，却不再分散。\par

\SourceRecord{Translated by J.E. Tweed.；Nicene and Post-Nicene Fathers, First Series；Vol. 8.；Edited by Philip Schaff.；Buffalo, NY: Christian Literature Publishing Co.,；1888.；<http://www.newadvent.org/fathers/1801147.htm>.}

% structure: p0001 source=h1 target=section reason=page-title

\section{\WorkTitle{圣咏第148篇释义}}

1. 我们在此生默想的主题，应为对天主的讚颂；因为来世永恒崇高之生命，便是对天主的讚颂，而无人能适於那来世，除非他现今就操练自己。所以，现今我们讚颂天主，但也向祂祈祷。我们的讚颂充满喜乐，我们的祈祷则伴着呻吟……论及这两个时期——一个即是现今，在这生命的试探与磨难中；另一个则是将来，在永恒的安息与欢跃中——我们也确立了两种庆典：复活节前与复活节后。复活节前象征我们现今所处的磨难；而我们现今正在庆祝的复活节后，则象征我们将来的福乐。因此，复活节前所守的庆典，正是我们现今所行的；而复活节后所守的，则象征我们尚未拥有的。於是，那段时期我们用于斋戒与祈祷；而现今这段时期，我们则用于讚颂，并放松斋戒。这就是我们所唱的\Quote{阿肋路亚}，正如你们所知，它在拉丁文中的意思是\Quote{讚颂上主}。因此，那段时期是在主复活之前，这段是在祂复活之后；由此时期象征那尚未拥有的未来希望：因为我们在主复活后所代表的，将在我们自己的复活后拥有。因为在我们元首内，两者都已预示，两者都已展现。主的圣洗向我们展示了这充满考验的现世生命，因为在此我们必须劳苦、受折磨，最终死去；但主的复活与受荣则向我们展示了将来要有的生命，那时祂将前来赏赐应得的报酬，恶人得恶，善人得善。而现今，所有恶人也与我们同唱\Quote{阿肋路亚}；但若他们执迷於恶，或许能以嘴唇唱出来世生命的歌曲，但生命本身——那将在未来成为现实、现今只是预象的——他们却得不到，因为他们不愿在它来临之前操练自己，也不愿抓住那将要来的。\par

2. \Quote{阿肋路亚。}你向你的邻人说：\Quote{讚颂上主，}他也向你说：当众人互相劝勉时，众人都在行他们所劝勉之事。但要全心讚颂：也就是说，不只你的口舌与声音讚颂天主，你的良心、你的生活、你的行为也要讚颂。因为现今，当我们聚集在教会内时，我们讚颂；但当我们各自外出处理事务时，我们似乎停止了讚颂。但愿人不停止善生，这样他就常讚颂天主……一个人的行为不可能邪恶，若他的思想是善的。因为行为出自思想：除非思想先发出命令，人不能做任何事，或动其肢体去做任何事；就如你在各省所见的一切事，无论皇帝下令什么，都从他的宫殿内部传达至整个罗马帝国。皇帝端坐宫中，一声令下便引起多大的骚动！他说话时只动动嘴唇；当他的话被执行时，整个省份都震动了。同样，在每个人之内，皇帝也在内里，他的宝座是在心中。若他是善的，发出善令，则善事得以实行；若他是恶的，发出恶令，则恶事得以实行。当基督坐在那里，祂能命令什么，除非是善？当魔鬼占据时，他能命令什么，除非是恶？但天主愿意让这成为你的选择：你愿为谁准备居所，为天主或为魔鬼？当你准备好，那占据者也将统治。因此，弟兄们，不仅要注意声音；当你讚颂天主时，要全心讚颂：愿你的声音、你的生活、你的行为，都一同歌唱。\par

3. \BibleQuote{请你们从天上讚颂上主} \BibleRef{咏 148:1}。好像他发现天上万物在讚颂上主方面缄默不语，於是劝勉它们起来讚颂。天上万物从未停止讚颂它们的造物主，地上万物也从未停止讚颂天主。但显然，有些受造物拥有气息，以天主所喜悦的那种心境讚颂天主。因为无人讚颂任何事物，除非它令牠喜悦。另有些生物没有生命的气息和理解力来讚颂天主，但因为它们也是善的，且按其适当秩序妥善安排，並构成天主所创造的宇宙之美，所以虽然它们自己不以声音和心灵讚颂天主，但当它们被有智慧的人审视时，天主便在它们内受到讚颂；而当天主在它们内受到讚颂时，它们本身也以某种方式讚颂天主……\par

4. \BibleQuote{请你们从天上讚颂上主：请你们在高处讚颂祂。}他先说\Quote{从天上，}然后说从地上；因为所讚颂的是天主，祂造了天和地。天上一切平静安宁：常有喜乐，没有死亡，没有疾病，没有烦恼；那里有福者常讚颂天主；但我们仍在下界：然而，当我们思想天主在那里如何受讚颂时，让我们把心放在那里，且不要自听\Quote{请举心向上。}让我们把心高举向上，免得它在世上腐化：因为我们以天使在那里所行的事为乐。我们现在以希望行之；将来当我们到达那里时，则以实际行之。因此，\Quote{请你们在高处讚颂祂。}\par

5. \BibleQuote{所有祂的天使，请赞美祂，所有祂的万军，请赞美祂} \BibleRef{咏 148:2} 。\BibleQuote{太阳和月亮，请赞美祂，所有星辰和光明，请赞美祂} \BibleRef{咏 148:3} 。\BibleQuote{诸天之上的天，以及天上的水，请赞美祂} \BibleRef{咏 148:4} 。\BibleQuote{愿他们赞美上主的圣名} \BibleRef{咏 148:5} 。他何时能在枚举中展开一切呢？然而他仿佛概括地触及了所有，将天上的一切都包含在赞美造物主之中。如同有人对他说：\Quote{他们为什么赞美祂？祂赐给了他们什么，使他们应该赞美祂？} 他接着说：\BibleQuote{因为祂一说，万物即成；祂一命，万物受造。} 受造物赞美造物主并不奇怪，被造之物赞美制造者并不奇怪，受造界赞美它的创造主并不奇怪。在此，基督也被提及，尽管我们似乎没有听到祂的名字……万物是藉什么造成的？藉着圣言？\BibleRef{若 1:1} 在这篇圣咏中，他如何显示万物是藉圣言造成的？\BibleQuote{祂一说，万物即成；祂一命，万物受造。} 没有人说话，没有人下令，除非藉着言语。\par

6. \BibleQuote{祂将它们安置在永远，世世代代} \BibleRef{咏 148:6} 。天上的一切，高天的一切，万军和天使，一座至高的城，美善、圣洁、蒙福；从那里我们因流徙而悲惨，因我们将返回而怀有望德蒙福；当我们返回后，将确实蒙福；\BibleQuote{祂赐给他们一条永不废止的律法。} 你们想，天上的事物和圣天使接受了怎样的命令？天主赐给了他们什么命令？岂不是要他们赞美祂？他们是有福的，他们的职务就是赞美天主！他们不耕耘，不播种，不磨面，不烹饪；因为这些都是必需的工作，而那里没有必需。他们不偷窃，不抢劫，不行奸淫；因为这些都是邪恶的行为，而那里没有邪恶。他们不为饥饿者擘饼，不给赤身者衣服，不接待旅客，不探望病人，不为争讼者和解，不埋葬死者；因为这些都是慈悲的工作，而那里没有苦难需要慈悲。哦，有福的他们！我们想我们也会像这样吗？啊！让我们叹息，让我们在叹息中呻吟。我们算什么，竟能到那里？必死的、被遗弃的、卑贱的、尘土和灰烬！但那位应许者是全能的……\par

7. 那么，既然他已经述说了天上之物的赞美，现在让他转向地上的事物。\BibleQuote{请你们从地上赞美上主} \BibleRef{咏 148:7} 。因为他先前是如何开始的？\BibleQuote{请你们从天上赞美上主：} 他列举了天上的事物；现在请听地上的事物。\BibleQuote{巨龙和一切深渊。} \Quote{深渊} 是水的深处：所有的海洋和这云雾缭绕的大气层都属于 \Quote{深渊。} 那里有云，有风暴，有雨，闪电，雷，雹，雪，以及天主愿意在地之上藉此潮湿多雾的大气成就的一切，他都以\Quote{地}之名提及，因为它非常易变而可朽；除非你们认为雨是从星辰之上降下的。这一切都发生在这里，靠近地面。有时人们甚至在山顶上，看见云在他们下面，常常下雨：而所有因大气扰动产生的骚动，那些细心观察的人看到它们发生在这里，在宇宙的这个较低部分……那么你们看到这一切是怎样的——易变的、扰乱的、可怕的、可朽的：然而它们有自己的位置，有自己的等级，它们也按自己的程度充实宇宙之美，因此它们赞美上主。于是他转向它们，仿佛也要劝勉它们，或劝勉我们，藉着默观它们而赞美上主。\Quote{巨龙} 生活在水边，从洞穴中出来，飞过空中；空气因它们而扰动：\Quote{巨龙} 是一种巨大的生物，在地上没有比它们更大的。因此他以它们开始：\BibleQuote{巨龙和一切深渊。} 有隐藏的水穴，泉源和溪流从中涌出：有些涌出流在地上，有些秘密地流在地下；所有这一切，连同海洋和这低层的空气，被称为深渊，或 \Quote{深渊,} 那里有巨龙生活并赞美天主。什么？我们认为巨龙组成唱诗班赞美天主吗？绝非如此。但你们，当你们想到巨龙时，注视巨龙的设计者、巨龙的创造者：那么，当你们赞叹巨龙，说：\Quote{造这一切的上主是伟大的，} 那时，巨龙藉你们的声音赞美天主。\par

8. \Quote{火、冰雹、雪、冰、执行祂话的暴风} \BibleRef{咏 148:8}。为何他在这里加上\Quote{执行祂的话}呢？许多愚蠢的人，不能默观和辨识受造物——它们在不同的地方和等级中，依天主的示意和命令而运行——他们认为天主确实统治着天上的一切，却鄙视、撇弃、遗弃了下界的事物，以致祂既不关心、也不引导、更不管理它们；相反，它们被偶然性所统治，随其所能、尽其所能地运行；他们有时被相互之间的话语所影响，例如他们说：\Quote{如果是天主赐雨，祂会降雨到海里吗？这是怎样的眷顾，}他们说，\Quote{这是怎么回事？\Place{革突里亚}干旱，雨却降在海里。}他们以为这事处理得巧妙。有人应对他们说：\Quote{\Place{革突里亚}确实口渴，你却不口渴。}因为对你们来说，更好是对天主说：\Quote{我的灵魂渴慕祢。}因为这样争辩的人已经自满；他自以为有学问，不愿学习，所以他不口渴。因为如果他口渴，他就愿学习，并会发现地上的一切都出于天主的眷顾，他还会惊叹于就连跳蚤的肢体安排。请注意，亲爱的。谁安排了跳蚤和蚊子的肢体，使它们有各自的秩序、生命、运动？试想一个渺小的生物，哪怕是最小的，随你选。如果你考虑它肢体的秩序，以及使之运动的生命力：它如何逃避死亡，热爱生命，寻求快乐，避免痛苦，运用各种感官，有力地使用适合自己的运动！谁给了蚊子螫针，用以吸血？那吸管多么细窄！谁安排了这一切？谁造了这一切？你们为最小的事物惊讶；要赞美那伟大的祂。因此，请持守这一点，我的弟兄们，不要让任何人动摇你们的信德或纯正的道理。那造了天上天使的，也同样造了地上的虫子：天使在天上居住，虫子在地上存留。祂没有造天使在泥中爬行，也没有造虫子在天空移动。祂已将居住者分派到不同的住所；不朽的住所祂分派给了不朽的，可朽的住所给了可朽的。观察整体，赞美整体。那么，那安排了虫子肢体的，难道不管理云彩吗？祂为何降雨入海？难道海里没有靠雨水滋养的东西吗？难道祂没有在其中造鱼，没有在其中造生物吗？看鱼如何游向淡水。那人说，祂为何降雨给鱼，有时却不降雨给我？这是为使你意识到你身在荒野之地、在生命的旅途中；好使现世生活对你变得苦涩，使你渴慕来世；或者使你受鞭打、受惩罚、得改正。祂对各地的特性分配得多么妥当。看，既然我们提到了\Place{革突里亚}，那里几乎每年降雨，每年出产谷物；这里的谷物不能储存，很快就腐烂，因为每年都出产；那边，因为出产稀少，不仅出产多，而且能久存。但或许你认为天主在那里遗弃了人，或者那里的人不以自己的欢乐方式赞美和光荣天主？把一个\Place{革突里亚}人从本国带来，置于我们宜人的树林中；他会想逃走，回到他那贫瘠的\Place{革突里亚}。因此，对一切地方、区域、季节，天主已分派和安排了适合它们的事物。谁能阐明呢？然而，有眼的人在其中看到许多事物：看见时，它们令人喜欢；令人喜欢时，它们被赞美；其实不是它们，而是造它们的那位；这样，万物都要赞美天主。\par

先知的神念及此，便加上一句：\Quote{遵行祂的话。}因此，不要以为这些事物是偶然运动的，它们的一切运行都听从天主。天主意愿火往哪里去，火就往那里去；云往哪里急驰，就往那里去，无论其中含着雨、雪还是冰雹。为什么闪电有时击中山岭，却不击中强盗呢？……也许祂仍在寻求强盗的悔改，因此击打不畏惧的山，为使畏惧的人得以改变。你有时在维护纪律时，也会击打地面以惊吓孩子。有时祂也击打祂所愿的人。但你会对我说：看，祂击打更无辜的人，却放过了更有罪的人。不要奇怪；死亡，无论来自何处，对善人都是好的。你怎么知道，若那更有罪的人不愿悔改，为他暗中保留了怎样的惩罚？难道他们不是更适合被闪电灼烧，因为他们最终将被宣告：\Quote{往永火里去}？\BibleRef{玛 25:41}要紧的是，你要纯良。为何如此？死于船难是坏事，死于发热是好事吗？无论他这样死还是那样死，要问他是怎样的人；要问他死后往哪里去，而不是他如何离开生命……因此，凡发生在此世违背我们意愿的事，你要知道，除非是天主的旨意、藉祂的天主眷顾、藉祂的安排、藉祂的示意、藉祂的法律，否则不会发生；如果我们不明白某事为何发生，让我们承认，它并非无因之果：这样我们才不会成为亵渎者。因为当我们开始争论天主的工作时，\Quote{为何如此？} \Quote{为何那样？} \Quote{祂本不该做这事，} \Quote{祂做错了；}这时天主的赞美在哪里？你失去了你的阿肋路亚。要以悦乐天主、赞美造物主的方式看待万事。因为如果你偶然进入一个铁匠的作坊，你不会敢于挑剔他的风箱、他的铁砧、他的锤子。但若有一个无知的人，不知道每样东西的用途，他就会挑剔一切。但如果他没有工匠的技能，却有人类的推理能力，他会对自己说什么？风箱放在这里不是无缘无故的：工匠知道为什么，尽管我不知道。在作坊里他不敢挑剔铁匠，却在宇宙中敢于挑剔天主。因此，正如\Quote{火、冰雹、雪、冰、风暴的旋风，遵行祂的话，}一样，自然界中一切在愚人看来似乎是随意产生的事物，都单纯地\Quote{遵行祂的话}，因为除了奉祂的命令，它们不会产生。\par

10. 然后他提到，为赞美主，\Quote{山岭和丘陵，果树和所有香柏木} \BibleRef{咏 148:9}：\Quote{走兽和一切牲畜，爬虫和有翅的飞禽} \BibleRef{咏 148:10}。然后他论到人：\Quote{地上的君王和万民，首领和地上的所有判官} \BibleRef{咏 148:11}：\Quote{青年和少女，老人和幼童，愿他们赞美主的名字} \BibleRef{咏 148:12}。天上的赞美结束了，地上的赞美结束了。\Quote{唯有祂的名被尊崇} \BibleRef{咏 148:13}。愿没有人寻求抬高自己的名。你要被高举吗？你要服从那不能贬抑者。\Quote{对祂的赞美在天地之中} \BibleRef{咏 148:14}。什么是\Quote{对祂的赞美}？难道是祂用以赞美自己的赞美吗？不，而是万物借此赞美祂，万物高呼：万物的美丽仿佛是它们的声音，借此它们赞美天主。天空向天主呼喊：\Quote{祢造了我，不是我造了自己。}大地呼喊：\Quote{祢创造了我，不是我造了自己。}它们如何呼喊？当你注视它们，并发现这一点，它们就藉你的声音呼喊，藉你的注视呼喊。注视天空，它是美丽的；观察大地，它是美丽的；二者一起非常美丽。祂造了它们，祂管理它们，凭祂的点头它们被支配，祂安排它们的季节，祂更新它们的运动，祂亲自更新它们。所有这些事都赞美祂，无论静止或运动，无论来自地下的大地或来自天上的穹苍，无论处于旧状态或更新之中。当你看见这一切，欢喜并提升到造物主面前，凝视\Quote{祂那看不见的事，可由所造之物明白} \BibleRef{罗 1:20}，\Quote{对祂的赞美在天地之中}：也就是说，你从地上的事物向祂歌颂，你从天上的事物向祂歌颂。既然祂造了万物，没有比祂更好的，凡祂所造的都比祂小，凡在这些事物中取悦你的，都比祂小。因此，不要让祂所造之物如此取悦你，以致把你从造物主那里引开：如果你爱祂所造之物，更要爱那造物者。如果祂造之物是美丽的，那么造它们的那位是多么更美丽。\Quote{祂必高举祂子民的角}。请看哈盖和匝加利亚预言了什么。如今\Quote{祂子民的角}在困苦中、在磨难中、在试探中、在捶胸中是卑微的；祂何时才\Quote{高举祂子民的角}？当主来到，我们的太阳升起，不是那肉眼所见、\Quote{上升照善人也照恶人} \BibleRef{玛 5:45}的太阳，而是那关于祂所说的话，对你们这些听从天主的人：\Quote{正义的太阳将要升起，在祂的翼下有医治} \BibleRef{拉 4:2}；后来骄傲和邪恶的人将要说：\Quote{正义之光没有照耀我们，正义的太阳没有在我们身上升起} \BibleRef{智 5:6}。这将是我们的夏季。如今在冬季天气，果实不在枝条上显现；你观察，可以说，在冬季像枯死的树木。那不能正确观看的人，以为葡萄树死了；也许旁边有一株确实死了的；两者在冬季相似；一株是活的，另一株是死的，但生命和死亡都隐藏着：夏天来临；于是那株的生命光辉灿烂，那株的死彰显出来：树叶的华丽，果实的丰盛，显现出来，葡萄树从它枝条内的东西披上外表。因此，弟兄们，现在我们和其他人一样：就如他们出生、吃喝、穿衣、度日，圣人们也是如此。有时真理本身欺骗人，他们说：\Quote{看，他成了基督徒：他不再头痛了？}或者，\Quote{因为他成了基督徒，我从他得到什么利益？}啊，枯死的葡萄树，你旁边有一株在冬季确实光秃、却未死的葡萄树。夏天将要来到，主将要来到，我们的光辉，那隐藏在枝条内的，然后\Quote{祂必高举祂子民的角}，在我们生于此世的囚禁之后……\par

11. \Quote{一首献给所有祂的圣人的圣诗。}你们知道什么是圣诗吗？圣诗是一首赞美天主的歌曲。如果你赞美天主而不歌唱，你就不是咏唱圣诗；如果你歌唱而不赞美天主，你也不是咏唱圣诗；如果你赞美别的什么，与赞美天主无关，尽管你歌唱并赞美，你也不是咏唱圣诗。因此，圣诗包含这三样：歌曲、赞美，而且赞美对象是天主。所以，赞美天主的歌曲才称为圣诗。那么\Quote{一首献给所有祂的圣人的圣诗}是什么意思？愿祂的圣人们接受一首圣诗：愿祂的圣人们咏唱一首圣诗：因为这就是他们最终要接受的，一首永恒的圣诗……\par

\SourceRecord{Translated by J.E. Tweed.；Nicene and Post-Nicene Fathers, First Series；Vol. 8.；Edited by Philip Schaff.；Buffalo, NY: Christian Literature Publishing Co.,；1888.；<http://www.newadvent.org/fathers/1801148.htm>.}

% structure: p0001 source=h1 target=section reason=page-title

\section{\WorkTitle{圣咏第一四九篇释义}}

1. 我们要用声音、理解和善行来赞美主；并如这篇圣咏所劝勉的，向祂唱一首新歌。它开头说：\BibleQuote{你们要向主唱新歌。祂的赞美在圣人的教会中} \BibleRef{咏 149:1}。旧人有旧歌，新人有新歌。旧约是旧歌，新约是新歌。旧约中有现世和地上的应许。凡爱慕地上事物的人就唱旧歌；但愿那渴望唱新歌的人，爱慕永恒的事物。爱德本身是崭新且永恒的；因此它常新，因为它永不陈旧……这首是和平之歌，是爱德之歌。谁若自绝于圣人的团体之外，便不唱新歌；因为他追随的是旧的纷争，而非新的爱德。在新的爱德中有什么？平安、圣洁团体的纽带、属灵的合一、活石的建筑。这在哪里？不在一个地方，而是在普世。另一篇圣咏说：\BibleQuote{你们要向主歌唱，全地都要唱。}由此可知，谁不与全地同唱，他唱的就是旧歌，无论从他口中说出什么话……弟兄们，我们已经说过，全地都唱新歌。谁不与全地同唱新歌，就任他唱他所愿的，任他的舌头高呼\Quote{阿肋路亚}，任他昼夜不停地呼喊，我的耳朵并不那么专注于歌者的声音，而是寻求行善者的行为。因为我询问并说：\Quote{你唱的是什么？}他回答：\Quote{阿肋路亚。}\Quote{阿肋路亚}是什么？\Quote{赞美主。}来，让我们一起赞美主。若你赞美主，我也赞美主，为何我们不合一？爱德赞美主，纷争亵渎主……\par

2. 主的田地是世界，不是阿非利加洲。这与我们这些田地不同，在\Place{格突利亚}有的出产六十倍或百倍，\Place{努米底亚}只有十倍：但各处都为祂结果实，有百倍的、六十倍的、三十倍的；只要你选择你将成为什么，如果你自认属于主的十字架。因此，\Quote{圣人的教会}就是天主教会。圣人的教会不是异端者的教会。圣人的教会是天主在它被看见之前预先表明，然后彰显出来让人看见的教会。圣人的教会从前只在经卷中，如今在列国中；圣人的教会从前只被读到，如今既被读到也被看见。当它只被读到的时候，人们相信它；如今它被看见，却受人反对。祂的赞美在\Quote{天国之子}中，即\Quote{圣人的教会}。\par

3. \BibleQuote{愿以色列因造他的主而喜乐} \BibleRef{咏 149:2}。\Quote{以色列}是什么？\Quote{看见天主。}看见天主的人，因造他的主而喜乐。那么，弟兄们，这意味着什么？我们说过我们属于圣人的教会：我们已看见天主了吗？我们又如何是以色列，若我们看不见？有一种属于现时的看见，另有属于将来的看见；现在的看见是藉着信德，将来的看见将是实体。如果我们相信，我们就看见；如果我们爱，我们就看见：看见什么？天主。请教若望：\BibleQuote{天主是爱} \BibleRef{若一 4:16}；让我们赞美祂的圣名，并因喜乐于爱中而在天主内喜乐。谁若有爱，为何我们送他到远方去见天主？让他省察自己的良心，在那里他就看见天主……\BibleQuote{愿熙雍的子女因他们的君王而欢欣。}教会的子女就是以色列。因为熙雍原是一座城，已经倒塌；在其废墟中，某些圣人按肉身居住过；但真正的熙雍，真正的耶路撒冷（因为熙雍和耶路撒冷是一体），是\BibleQuote{在天上永存的房屋} \BibleRef{格后 5:1}，是\BibleQuote{我们的母亲} \BibleRef{迦 4:26}。她生了我们，她是圣人的教会，抚养了我们——她一部分在旅途中，一部分安居在天上。安居在天上的那部分有天使的福乐，在世上漂泊的那部分有义人的希望。关于前者说：\BibleQuote{光荣归于至高天主；在地上平安归于善良的人} \BibleRef{路 2:14}。那么，那些在此生中叹息、渴望故国的人，要藉着爱奔跑，而非藉着身体的脚；他们不要寻找船只，而要寻找翅膀，要抓住爱的双翼。爱的双翼是什么？爱天主和爱近人。因为如今我们是旅客，我们叹息，我们哀伤。有一封信从我们的故国寄来：我们读给你们听。\BibleQuote{愿熙雍的子女因他们的君王而欢欣。}天主子造了我们，又成为我们中的一员：祂统治我们如君王，因为祂是我们的创造者，造了我们。但造我们的那一位，也就是统治我们的那一位，我们是基督徒，因为祂是基督。祂被称为基督，源自\OriginalTerm{傅油}{Chrism}，即受傅……要给司祭献上一些祭品。人能找到什么洁净的祭品献上？什么祭品？罪人能献上什么洁净之物？啊，不义的人，罪人，你所献的一切都是不洁的，必须为你献上某种洁净之物……那么，让洁净的司祭献上祂自己，并洁净你。这就是基督所做的。祂在人身上找不到任何洁净之物可以为人类献上：祂将自己作为洁净的祭品献上。有福的祭品，真正的祭品，无玷的奉献。祂所献的，不是我们给祂的；更确切地说，祂献上了祂从我们身上取来的，并洁净地献上。因为祂从我们取了肉身，并献上这肉身。但祂在哪里取的？在童贞玛利亚的胎中，为使我们这些不洁的人，祂能洁净地献上。祂是我们的君王，祂是我们的司祭，我们要在祂内喜乐。\par

4. \BibleQuote{愿他们在合唱中赞美祂的名}\BibleRef{咏 149:3}。\Quote{合唱}是什么意思？许多人知道什么是\Quote{合唱}：不，既然我们在城镇中讲话，几乎人人皆知。\Quote{合唱}是歌唱者的联合。如果我们\Quote{合唱}歌唱，就让我们和谐地歌唱。若任何人的声音在合唱中不协调，它便会刺耳，扰乱合唱。若一个人的不谐之音扰乱了歌唱者的和谐，那么异端的不谐之音又如何扰乱赞美者的和谐呢？现今整个世界都是基督的合唱。基督的合唱从东到西和谐地响起。\BibleQuote{愿他们击鼓弹琴歌颂祂}\BibleRef{咏 149:3}。他为何拿起\Quote{鼓和瑟}？为叫不仅声音赞美，而且行为也赞美。当拿起鼓和瑟时，手与声音和谐一致。所以，当你歌唱\Quote{阿肋路亚}时，也要把你的饼分给饥饿的人，给赤身者衣服穿，收留旅客：那时不仅你的声音响起，而且你的手也与它和谐地响起，因为你的行为与你的言语一致。你拿起了一件乐器，你的手指与你的舌头一致。我们也不可忽略\Quote{鼓和瑟}的奥义。鼓上绷着皮，瑟上绷着肠线；两种乐器上都钉着肉。那说\BibleQuote{世界已被钉死在我身上，我也被钉死在世界上}\BibleRef{迦 6:14}的人，岂不是很好地\Quote{在鼓和瑟上咏诗}吗？这位喜爱新歌、教导你们、对你们说\BibleQuote{谁若愿意跟随我，该弃绝自己，背着自己的十字架来跟随我}\BibleRef{玛 16:24}的那一位，愿你们拿起这瑟或鼓。愿他不放下他的瑟，不放下他的鼓，愿他把自己伸展在木头上，并从肉欲中干燥。弦绷得越紧，声音就越清脆。那么，宗徒\Person{保禄}为使他的瑟发出清脆响声，说了什么？\BibleQuote{向着前面伸展}等等。\BibleRef{斐 3:13}他伸展了自己：基督触摸了他；真理的甘甜便回响了。\par

5. \BibleQuote{因为上主善待了祂的百姓}\BibleRef{咏 149:4}。有哪一个仁慈的作为，比为不虔敬者而死更伟大？有哪一个仁慈的作为，比用正义的宝血涂去罪人身上的罪状更伟大？有哪一个仁慈的作为，比说\Quote{我不看重你们过去是什么，现在成为你们过去所不是的}更伟大？祂仁慈地使背离者归正，援助争战者，为胜利者加冕。\BibleQuote{温良的人祂必举扬于救恩中}\BibleRef{咏 149:4}。因为骄傲者也被举起，但不是升入救恩：温良者被举起升入救恩，骄傲者被举起进入死亡：也就是说，骄傲者高举自己，天主却贬抑他们；温良者谦卑自己，天主却高举他们。\par

6. \BibleQuote{圣徒们要在光荣中欢跃}\BibleRef{咏 149:5}。关于圣人的光荣，我愿意说些重要的事。因为没有人不爱光荣。但愚人的光荣，即所谓的世俗光荣，有迷惑人的陷阱，使人受虚浮之人的赞美所影响，愿意以这样的方式生活，以便被人谈论，无论他们是谁，无论以什么方式。因此，人们变得疯狂，被骄傲冲昏头脑，内心空虚，外表肿胀，宁愿将自己的财产浪费在舞台演员、戏子、斗兽者、战车御者身上。他们给出多少，花费多少！他们不仅耗尽祖产，也耗尽心智。他们轻视穷人，因为民众并不呼喊应该施舍给穷人，但民众确实呼喊应该施舍给斗兽者。因此，当没有人向他们呼喊时，他们拒绝花费；当疯狂者向他们呼喊时，他们也疯狂：不，所有人都疯狂，表演者、观看者和施予者都一样。这种疯狂的光荣受到主的责备，在全能者的眼中是可憎的……你选择给斗兽者衣服穿，他可能被打败，使你脸红：基督从未被征服；祂已征服了魔鬼，祂为你、向你、在你内征服了；像这样一位征服者，你却不选择给他衣服穿。为什么？因为关于祂的呼喊较少，疯狂也较少。因此，那些喜爱这种光荣的人，良心是空虚的。正如他们掏空钱箱，送衣服作礼物，同样他们也掏空良心，以致里面没有任何宝贵之物。\par

7. 但那些\Quote{在光荣中欢跃}的圣人，我们无需说明他们如何欢跃：只需听诗篇接下来的诗句：\BibleQuote{圣徒们要在光荣中欢跃，他们在床榻上欢欣}\BibleRef{咏 149:5}。不是在剧院、或露天剧场、或竞技场、或愚行、或市场，而是\Quote{在他们的内室}。什么是\Quote{在他们的内室}？在他们的内心。听宗徒\Person{保禄}在他的内室中欢跃：\BibleQuote{我们的光荣就是良心的见证}\BibleRef{格后 1:12}。另一方面，有理由担心有人会自满，因而显得骄傲，夸耀自己的良心。因为每个人都应当怀着敬畏欢跃，因为他所欢跃的是天主的恩赐，而非自己的功劳。因为有许多人自娱自乐，自以为义；另有一段经文反对他们，说：\BibleQuote{谁能说自己的心洁净，自己是纯洁无罪的？}\BibleRef{箴 20:9}。因此，可以说，在良心中夸耀有一个限度，即知道你的信德真诚，望德坚定，爱德无伪。\BibleQuote{天主的颂赞在他们口中}\BibleRef{咏 149:6}。他们应当这样\Quote{在他们的内室中欢欣}，不把他们的善归给自己，而是赞美那使他们成为他们所是的那一位，祂召叫他们去获得他们所不是的，并盼望从祂那里得着圆满，向祂感恩，因为祂已开始了。\par

8. \Quote{他们手中拿着两面锋利的剑。}这种兵器含有极大的奥秘意义，因为它两面锋利。所谓\Quote{两面锋利的剑}，我们理解为圣言：\BibleRef{希 4:12}它是一把剑，但之所以被称为许多把，是因为有许多圣人的口和许多圣人的舌头。它怎么是两面锋利的呢？它论及暂时之事，也论及永恒之事。在两方面的情形下，它都证明自己所说的，凡被它击中的，它就使之与世界分离。这不正是主所说的那把剑吗？\BibleQuote{我来不是为把平安带到地上，而是带一把剑。}\BibleRef{玛 10:34}请看祂来如何分划，如何分离。祂分划圣人，分划不敬虔的人，祂从你那里分离那阻碍你的事物。儿子愿意事奉天主，父亲却不愿意：剑来了，圣言来了，把儿子从父亲那里分离出来……那么，为什么剑在他们手中，而不是在他们舌头上呢？经上说：\Quote{他们手中拿着两面锋利的剑。}所谓\Quote{在他们手中}，意思是在权力中。他们领受了圣言，并拥有权力，要在哪里说就在哪里说，要对谁说就对谁说，既不惧怕权势，也不轻视贫穷。因为他们手中有一把剑：他们愿意在哪里挥舞，就在哪里挥舞，用它击打；这一切都在宣讲者的权力之中。因为如果圣言不在他们手中，为什么经上写着\BibleQuote{圣言交给了\Person{哈盖}先知的手中}？的确，弟兄们，天主并没有把祂的圣言放在祂的手指里。所谓\Quote{放在他手中}是什么意思？就是交给他宣讲圣言的权力。最后，我们也可以用另一种方式理解这些\Quote{手}。因为说话的人将圣言放在舌头上，而写字的人则放在手中。\par

9. 如今，弟兄们，你们看见圣人武装起来了：看那屠杀，看他们光荣的战役。因为有统帅，就必须有士兵；有士兵，就有敌人；有战斗，就有胜利。这些手中拿着两面锋利的剑的人做了什么？\Quote{为向列邦报仇。}看看是否没有向列邦报仇。天天都在做：我们自己说话时就在做。看巴比伦的列邦如何被击杀。她受到了双倍的报应：因为经上论到她这样说：\Quote{她怎样行事，就怎样加倍回报她。}\BibleRef{默 18:6}她怎么是加倍回报呢？圣人作战，拔出\Quote{两面锋利的剑}；随之而来的是失败、屠杀、分离：她怎么是加倍回报呢？当她有权力迫害基督徒时，她固然杀害了肉体，但没有压碎天主：现在她得到了双倍的报应，因为外教人被消灭，偶像被打破……并且，免得你以为人真的被剑所击，血真的流出，肉体受伤，他继续解释说：\Quote{谴责在万民中。}\Quote{谴责}是什么？是谴责。让\Quote{两面锋利的剑}从你口中发出，不要迟延。对你的朋友说——如果你还有一个朋友可以对他说——\Quote{你是什么样的人，竟抛弃了造你的天主，反而崇拜祂所造的东西？工匠比祂的作品更好。}当他开始脸红、开始痛悔时，你就用你的剑刺伤了他，剑已触及心脏，他即将死去，以便生存。\par

10. \BibleQuote{他们要把他们的君王用锁链捆绑，把他们的显贵用铁镣束缚}\BibleRef{咏 149:8}。 \BibleQuote{要执行在他们身上所写的审判}\BibleRef{咏 149:9}。 外邦人的君王将被锁链捆绑，\Quote{他们的显贵也被锁链捆绑}，而且是\Quote{铁的}。……我们现在开始解释的这些经文是隐晦的。因为天主有意这样隐晦地写下祂的一些经文，并不是要从中挖出什么新的东西，而是要使那已经众所周知的事，因被隐晦地表述而成为新的。我们知道君王已成为基督徒；我们知道外邦人的显贵已成为基督徒。他们如今正在成为基督徒；他们过去是，将来也是；那\Quote{双刃的剑}在圣人们手中并未闲置。那么，我们如何理解他们被锁链和铁链捆绑呢？你们知道，亲爱的、有学识的弟兄们（我称你们为有学识的，因为你们在教会中受滋养，习惯于听天主的圣言诵读），\Quote{天主选择了世上懦弱的人，来羞辱那强壮的；天主又选择了世上愚拙的人，来羞辱那有智慧的；并且选择了那不存在的，如同那存在的，为了使那存在的归于无有。}……主说：\BibleQuote{你若愿意是成全的，去，变卖你所有的，施舍给穷人，然后来跟随我，你将有宝藏在天上}\BibleRef{玛 19:21}。许多显贵这样做过，但他们不再作外邦人的显贵，宁愿在世上作穷人，在基督内作尊贵的人。但许多人保留了他们从前的尊贵，保留了他们君王的权力，却仍然是基督徒。这些人，可以说，\Quote{被锁链和铁链捆绑}。何以如此？他们接受了锁链，以阻止他们走向非法之事，那\BibleQuote{智慧的锁链}\BibleRef{德 6:25}，天主圣言的锁链。那么，为什么是铁链而不是金链呢？只要他们畏惧，就是铁的；让他们爱，就会变成金的。亲爱的，请注意我所说的。你们刚才听到宗徒\Person{若望}说：\BibleQuote{爱内没有恐惧，完美的爱驱除恐惧，因为恐惧含有痛苦}\BibleRef{若壹 4:18}。这就是铁链。然而，一个人若不藉着敬畏天主开始，就不会达到爱。\Quote{敬畏上主是智慧的开端。}因此，开端是铁链，结局是金链。因为论及智慧，曾说：\BibleQuote{你颈上的金链}\BibleRef{德 6:24}。……有一个在世上有权势的人来找我们，他的妻子得罪了他，或许他贪恋了别人更美丽的妻子，或另一个更富有的女人，他想休掉现在的妻子，却没有这样做。他听了天主仆人的话，听了先知的话，听了宗徒的话，就没有做；那手持\Quote{双刃剑}的人告诉他：你不可这样做；这对你是不合法的；天主不允许你休妻，\BibleQuote{除非为奸淫的缘故}\BibleRef{玛 5:32}。他听了这话，害怕了，就没有做……年轻人，听着；这些链子是铁的，不要把你们的脚伸进去；如果你们伸进去，你们将被更紧地捆绑。主教的双手为你们加固了这样的锁链。难道那些这样被捆绑的人，不是逃到教会，在这里被解开吗？人们确实逃到这里，想要摆脱他们的妻子：在这里他们被绑得更紧：没有人解开这些锁链。\BibleQuote{天主所结合的，人不可拆散}\BibleRef{玛 19:6}。但这些链子是硬的。谁不知道呢？宗徒们对这艰难感到忧心，说：\BibleQuote{如果妻子是这样的情况，那倒不如不娶了}\BibleRef{玛 19:10}。如果链子是铁的，就不该把脚伸进去。主说：\BibleQuote{这话不是人人都能领悟的，但谁能领悟，就领悟吧}\BibleRef{玛 19:11}。\Quote{你被妻子束缚了吗？不要寻求解脱}，因为你被铁链捆绑。\Quote{你没有妻子吗？不要寻找妻子}，不要用铁链捆绑自己。\par

11. \Quote{在他们身上执行所写的审判。}这就是圣人们在万民中所执行的审判。为什么说\Quote{所写的}？因为这些事从前被写下来，如今得以应验。看哪，如今它们正在实行；从前它们被诵读，却没有实行。他这样总结说：\Quote{这荣耀是属于祂所有的圣人的。}在全世界的万民中，圣人们这样做，他们就这样被荣耀，就这样\Quote{用他们的口尊崇天主}，就这样\Quote{在他们的床上欢悦}，就这样\Quote{在他们的荣耀中欢欣}，就这样\Quote{在救恩中被高举}，就这样\Quote{唱新歌}，就这样在内心、言语和生活中说：阿肋路亚。阿们。\par

\SourceRecord{Translated by J.E. Tweed.；Nicene and Post-Nicene Fathers, First Series；Vol. 8.；Edited by Philip Schaff.；Buffalo, NY: Christian Literature Publishing Co.,；1888.；<http://www.newadvent.org/fathers/1801149.htm>.}

% structure: p0001 source=h1 target=section reason=page-title

\section{对圣咏第150篇的讲论}

1. 虽然圣咏的编排（在我看来它隐藏着一个伟大奥秘的秘密）至今尚未向我启示，但事实上它们总共有一百五十篇，这也已向我们这尚未用心智之眼洞察其全部编排之深度的人提示了一些东西，使我们在不自恃过甚、按天主所赐的情况下能够谈论。首先，十五这个数字（一百五十是它的倍数），我说十五这个数字，象征两约的契合。因为在旧约中遵守安息日，它象征安息；而在新约中遵守主日，它象征复活。安息日是第七天，但主日，在第七天之后，必然是第八天，同时也算作第一天。因为它被称为周的第一天，从它起算第二天、第三天、第四天，直到周的第七天，即安息日。但从主日到主日有八天，这宣告了新约的启示，这启示在旧约中仿佛被尘世的许诺所遮蔽。此外，七加八等于十五。同样数目的圣咏被称为\Quote{登阶的}，因为那是圣殿台阶的数目。此外，五十这个数字本身也含有一个伟大的奥秘。因为它由七个星期组成，再加上一个第八天凑成五十之数。因为七乘七得四十九，加上一成为五十。这个五十的数字意义如此重大，以至于从主复活算起，满了那些日子后，正好在第五十天，圣神降临到那些在基督内聚集的人身上。而在圣经中，圣神尤其以数字七来指代，无论是在\WorkTitle{依撒意亚}还是\WorkTitle{默示录}中，那里极其直接地提到天主的七神，这是由于同一圣神七重的作用。\BibleRef{默 1:20} 这七重的作用在\WorkTitle{依撒意亚}中也提到。\BibleRef{依 11:2}……因此圣神也被以数字七来称呼。但主将这个五十天分为四十和十：因为祂复活后第四十天升天，然后满了十天后祂派遣了圣神：以四十这个数字向我们预示在世上寄居的时期。因为四十中占主导的是数字四，而世界和年各有四个部分；加上数字十，作为因善工完成法律而附加的某种赏报，则预表永恒本身。这个五十，一百五十包含三次，仿佛它被乘以圣三。因此，我们也看出这篇圣咏的这个数字并非不合适。\par

2. 现在，有些人相信圣咏被分为五卷，他们是被下面的情况所引导，即圣咏结尾处多次出现\Quote{阿们，阿们}这样的话。但当我试图弄明白这种划分的原则时，我未能做到；因为五个部分彼此不等，无论在内容上还是在圣咏数目上，都不至让每卷包含三十篇。而且如果每卷都以\Quote{阿们，阿们}结束，我们有理由问，为什么第五卷也就是最后一卷没有同样的结尾。然而，我们遵循正典圣经的权威，那里说：\BibleQuote{因为在\WorkTitle{圣咏集}上写着：}\BibleRef{宗 1:20} 我们知道只有一卷\WorkTitle{圣咏集}。我确实看到这如何可以是真的，而另一种说法也可以是真的，而不与之相悖。因为可能有些希伯来文学的习俗，据此将一部书称为一卷，而它却由多卷组成，正如许多教会组成一个教会，许多天组成一个天，……许多地组成一个地。因为我们日常习惯说\Quote{地球的球体}和\Quote{大地的球体}。而当说\Quote{写在\WorkTitle{圣咏集}中}时，尽管习惯说法似乎暗示只有一卷，但对此可以回答说，那句话的意思是\Quote{在某一卷\WorkTitle{圣咏集}中}，也就是说，在五卷中的任何一卷。这在日常语言中如此前所未闻，或至少如此罕见，以至于我们之所以确信\WorkTitle{十二先知书}成为一卷，是因为我们同样读到：\BibleQuote{正如在\WorkTitle{先知书}上写着：}\BibleRef{宗 7:42} 也有些人把所有正典圣经一起称为一卷，因为它们以一种非常奇妙神圣的合一彼此一致……\par

3. 无论哪一种理解正确，这本由每五十篇组成的书卷都给出了一个重要且非常值得思考的答复。因为我觉得并非无意义：第五十篇是关于补赎的，第一百篇是关于仁慈和审判的，第一百五十篇是关于在天主圣人中赞美天主的。因为这样我们才得以进入永福的生命：首先谴责自己的罪，然后生活正义，这样，谴责了不良的生活并活出了善生之后，我们才能获得永生。我们的预定并非在我们自身内成就，而是在祂内隐秘地、在祂的预知中成就。但我们借着宣讲补赎而被召叫。我们借着仁慈的召叫和对审判的恐惧而成义。那已预先获得救恩的人不畏审判。被召叫后，我们借补赎弃绝魔鬼，以免继续在它的轭下；成义后，我们借仁慈得医治，以免畏惧审判；得光荣后，我们进入永生，在那里我们无休止地赞美天主……这篇圣咏结尾的经文就是永生的声音。\par

4. \BibleQuote{在祂的圣人中赞美主}，即，在祂所光荣的人中：\BibleQuote{在祂能力的穹苍中赞美祂} \BibleRef{咏 150:1} 。\BibleQuote{在祂强有力的作为中赞美祂；} 或如别人所解释的，\BibleQuote{在祂大能的作为中赞美祂：照祂伟大的众多赞美祂} \BibleRef{咏 150:2} 。所有这些就是祂的圣人；如宗徒所说：\BibleQuote{使我们在祂内成为天主的正义。} \BibleRef{格后 5:21} 如果他们是天主的正义，是祂在他们内所行的，为何他们不是基督的力量，是祂在他们内所行的，使他们从死者中复活？因为在基督的复活中，\BibleQuote{力量} 特别显示给我们，因为在祂的受难中是软弱，如宗徒所说。而它说得很好：\BibleQuote{祂能力的穹苍。} 因为这就是\BibleQuote{祂能力的穹苍}，即祂\BibleQuote{不再死，死亡再没有统治祂的权柄} \BibleRef{罗 6:9} 。为何他们不被称为天主\BibleQuote{力量的工作}，就是祂在他们内所行的？不仅如此，他们本身就是祂力量的工作；正如所说：\BibleQuote{我们在祂内成为天主的正义。} 有什么比祂永远为王，将一切仇敌置于祂脚下更有力量的？为何他们不也是\BibleQuote{祂伟大的众多}？不是指祂本身伟大，而是指祂使他们伟大，他们人数众多，即千千万万。正如正义也有两种方式理解：一种是指祂本身的正义，另一种是指祂在我们内所行的，使我们成为祂的正义。这些圣人都由随后的一切乐器所象征，为赞美天主。因为圣咏作者以\BibleQuote{在祂的圣人中赞美主} 开始，然后以各种方式象征这些祂的圣人来继续。\par

5. \BibleQuote{在号角声中赞美祂} \BibleRef{咏 150:3} ：由于他们赞美之声超乎寻常的嘹亮。\BibleQuote{在瑟和琴中赞美祂。} \BibleRef{咏 150:3} \OriginalTerm{瑟}{psaltery} 从高处赞美天主，\OriginalTerm{琴}{harp} 从低处赞美天主；我是指，从天上和地上的事物，如同祂创造了天地。我们在另一篇圣咏中已经解释过，\OriginalTerm{瑟}{psaltery} 有那块板，琴弦列于其上，以便发出更好的声音，它在上面，而\OriginalTerm{琴}{harp} 的板在下面。\BibleQuote{在鼓和歌舞中赞美祂} \BibleRef{咏 150:4} 。\OriginalTerm{鼓}{timbrel} 赞美天主，当肉体已经改变，其中没有属土腐败的软弱。因为鼓是由晾干并加强的皮革制成的。\OriginalTerm{歌舞}{choir} 赞美天主，当和平的社会赞美祂。\BibleQuote{在弦和箫中赞美祂。} \BibleRef{咏 150:4} \OriginalTerm{弦}{strings} 和\OriginalTerm{箫}{organ}。前面提到的\OriginalTerm{瑟}{psaltery} 和\OriginalTerm{琴}{harp} 都有弦。但\OriginalTerm{箫}{organ} 是所有乐器的总称，虽然现今的习惯已使那些用风箱吹奏的乐器特别被称为\OriginalTerm{管风琴}{organ}。但我认为这里不是指这种。因为\OriginalTerm{箫}{organ} 是一个希腊词，一般用于所有乐器，正如我所说；而这种用风箱的乐器，希腊人另有名称。但称其为\OriginalTerm{管风琴}{organ} 更多是拉丁语和口语用法。因此，当他说\BibleQuote{在弦和箫中}，在我看来，他有意指某种有弦的乐器。因为不仅有瑟和琴有弦；并且因为在\OriginalTerm{瑟}{psaltery} 和\OriginalTerm{琴}{harp} 中，由于来自低处和高处的声音，我们已经发现一些可按此区分来理解的东西，所以他提示我们在弦本身中寻求别的意义：因为它们也是肉体，但现在是已从朽坏中解放的肉体。也许他为此加上了\OriginalTerm{箫}{organ}，以表明它们不是各自单独发声，而是在最和谐的多样性中一起发声，就像在乐器中排列的那样。因为即使那时，天主的圣人也会有他们的差异，和谐而非不和谐，即一致而非分歧，正如最甜美的和声来自不同但并不对立的音调。\par

6. \BibleQuote{在响声和谐的钹上赞美祂，在欢颂的钹上赞美祂} \BibleRef{咏 150:5} 。\OriginalTerm{钹}{cymbals} 相互碰触以发声，因此有人将它们比作我们的嘴唇。但我认为更好的理解是，当每人被邻人尊崇，而非自夸，然后互相尊崇，他们赞美天主时，天主在某种程度上被\OriginalTerm{钹}{cymbals} 赞美。但恐怕有人会理解这样的钹是没有生命的发声，因此我认为他加上了\BibleQuote{在欢颂的钹上}。因为\OriginalTerm{欢颂}{jubilation}，即不可言喻的赞美，只出自生命。我也不认为我该忽略音乐家所说的：有三种声音：由声音、由气息、由弹拨。由声音：由喉咙和气管发出，当人不用任何乐器歌唱时；由气息：如笛子或类似之物；由弹拨：如竖琴或类似之物。这里没有遗漏任何一种：因为\OriginalTerm{歌舞}{choir} 中有声音，\OriginalTerm{号角}{trumpet} 中有气息，\OriginalTerm{琴}{harp} 中有弹拨，代表心智、灵性、身体，但只是类推，而非字面意义。因此，当他提出\BibleQuote{在祂的圣人中赞美天主}，他向谁说的，岂非他们自己？他们应在哪里赞美天主，岂非在他们自己内？因为你，他说，乃是\BibleQuote{祂的圣人}；你是\BibleQuote{祂的力量}，但那是祂在你内所行的；你是\BibleQuote{祂大能的作为，和祂伟大的众多}，祂在你内所行并彰显的。你们是\BibleQuote{号角、瑟、琴、鼓、歌舞、弦和箫、欢颂的钹，响声和谐}，因为和谐发声。所有这些都是你们；不要想任何卑贱、短暂、可笑的事物。既然随从肉体的是死亡，\BibleQuote{让每个有气息的赞美主} \BibleRef{咏 150:6}。\par

% structure: p0008 source=h2 target=subsection reason=relative-heading-rank; base=h2

\subsection{\Person{圣奥斯定}的祈祷}

\Emph{这是他讲道和讲座后惯常使用的。}\par

我们转向主天主，全能圣父，并以纯洁的心，按我们卑微之所能，向祂献上伟大而真诚的感谢，全心祈求祂的无上仁慈：愿祂按祂的善意，垂听我们的祈祷；愿祂以大能，从我们的言行中驱逐仇敌；愿祂增加我们的信德，引导我们的理智，赐予我们属神的思想，并引领我们进入祂的真福。藉着耶稣基督，祂的圣子，我们的主，祂与圣父及圣神同享永生，永王永统，唯一的天主，无穷世之世。阿们。\par

\SourceRecord{Translated by J.E. Tweed.；Nicene and Post-Nicene Fathers, First Series；Vol. 8.；Edited by Philip Schaff.；Buffalo, NY: Christian Literature Publishing Co.,；1888.；<http://www.newadvent.org/fathers/1801150.htm>.}
